% Generated by matins/tools/build_latex.py from data/corpus.json — do not edit.
\documentclass[10pt,twoside,openright]{book}
\usepackage{matins}
\begin{document}
\BodySize
\frontmatter
\thispagestyle{empty}
\vspace*{0.25\textheight}
{\centering
{\Huge\rubr{Lectiones ad Matutinum}}\par\vspace{0.4em}
{\large\itshape ex Breviario Romano anni MCMLIV}\par\vspace{2em}
{\LARGE The Lessons of Matins}\par\vspace{0.4em}
{\large\itshape from the Roman Breviary of 1954, in Latin and English}\par\vspace{3em}
{\Large\rubr{Pars Autumnalis}}\par\vspace{0.3em}{\large\itshape Autumn}\par\vspace{2em}
{\small }\par}
\clearpage
\thispagestyle{empty}
{\footnotesize\setlength{\parindent}{0pt}\setlength{\parskip}{0.35em}\raggedright
\par\addvspace{0.6em}\textsc{Divinum Officium}\par
The texts, and the program that arranges them by the rubrics, are the work of the Divinum Officium Project (divinumofficium.com), created by the late Laszlo Kiss and continued since by the Project's team of volunteers. Its code and data are freely available at github.com/DivinumOfficium/divinum-officium under the MIT License.\par
\par\addvspace{0.6em}\textsc{Latin}\par
The Latin text derives from the Ratisbon 1888 edition of the Breviarium Romanum, scanned by the University of St Michael's College, Toronto, amended from the 1906 edition, and corrected by Laszlo Kiss and the Project team against later printings: Pustet (Regensburg, 1943), Benziger Brothers (New York, 1945), Burns Oates \& Washbourne (London, 1946), La Presse Catholique Panaméricaine (Montréal, 1943), and Desclée and Mame (1962).\par
The accented text of the Sixto-Clementine Vulgate was supplied by Frère Romain-Marie of the Abbaye Saint-Joseph de Clairval, Flavigny-sur-Ozerain.\par
\par\addvspace{0.6em}\textsc{English}\par
Holy Scripture is given in the Douay-Rheims translation.\par
The other lessons and the responsories are from the translation of the Roman Breviary by John, Marquess of Bute (edition of 1908), made available by the University of St Michael's College, Toronto.\par
\par\addvspace{0.6em}\textsc{This edition}\par
The lessons were gathered by running the Divinum Officium program, unchanged, for every day from 1950 to 2100 under the rubrics of 1954, and tracing each lesson it read to the office it belongs to. Errors of arrangement in this edition are not the Project's.\par
\par\addvspace{0.6em}\textsc{Typeface}\par
Set in EB Garamond, designed by Georg Duffner and Octavio Pardo, under the SIL Open Font License 1.1.\par
}
\mainmatter
\PartTitle{Proprium de Tempore}{Proper of Time}

\addcontentsline{toc}{chapter}{Proprium de Tempore \textit{— Proper of Time}}

\SeasonTitle{Tempus post Pentecosten}{Time after Pentecost}

\addcontentsline{toc}{section}{Tempus post Pentecosten \textit{— Time after Pentecost}}

\Office{Dominica XI Post Pentecosten}{Eleventh Sunday after Pentecost}{Semiduplex}
\addcontentsline{toc}{subsection}{Dominica XI Post Pentecosten \textit{— Eleventh Sunday after Pentecost}}
\begin{paracol}{2}
\LA
\RefLine{Lectiones I–VI: de Scriptura occurrente.}
\switchcolumn
\EN
\RefLine{Lessons I–VI: from the Scripture of the day.}
\switchcolumn*

\LA
\LessonHead{Lectio I}
\switchcolumn
\EN
\LessonHead{Lesson I}
\switchcolumn*

\LA
\LessonTitle{De libro quarto Regum}
\switchcolumn
\EN
\LessonTitle{Lesson from the second book of Kings}
\switchcolumn*

\LA
\Cite{4 Reg 20:1-3}
\switchcolumn
\EN
\Cite{2 Kgs 20:1-3}
\switchcolumn*

\LA
\noindent In diébus illis ægrotáret Ezechías usque ad mortem, et venit ad eum Isaías fílius Amos prophéta dixítque ei: Hæc dicit Dóminus Deus: Prǽcipe dómui tuæ, moriéris enim tu et non vives. Qui convértit fáciem suam ad paríetem et orávit Dóminum, dicens: Obsecro, Dómine: meménto, quæso, quómodo ambuláverim coram te in veritáte et in corde perfécto et quod plácitum est coram te fécerim. Flevit ítaque Ezechías fletu magno.\par
\switchcolumn
\EN
\noindent In those days Ezechias was sick unto death: and Isaias the son of Amos the prophet came and said to him: Thus saith the Lord God: Give charge concerning thy house, for thou shalt die, and not live. And he turned his face to the wall, and prayed to the Lord, saying: I beseech thee, O Lord, remember how I have walked before thee in truth, and with a perfect heart, and have done that which is pleasing before thee. And Ezechias wept with much weeping.\par
\switchcolumn*

\LA
\begin{Resp}\Rsign Præparáte corda vestra Dómino, et servíte illi soli: \Star{} Et liberábit vos de mánibus inimicórum vestrórum.\end{Resp}
\begin{Resp}\Vsign Convertímini ad eum in toto corde vestro, et auférte deos aliénos de médio vestri.\end{Resp}
\begin{Resp}\Rsign Et liberábit vos de mánibus inimicórum vestrórum.\end{Resp}
\switchcolumn
\EN
\begin{Resp}\Rsign Prepare your hearts unto the Lord, and serve Him Only \Star{} And He will deliver you out of the hand of your enemies.\end{Resp}
\begin{Resp}\Vsign Return unto Him with all your hearts, and put away the strange gods from among you.\end{Resp}
\begin{Resp}\Rsign And He will deliver you out of the hand of your enemies.\end{Resp}
\switchcolumn*

\LA
\LessonHead{Lectio II}
\switchcolumn
\EN
\LessonHead{Lesson II}
\switchcolumn*

\LA
\Cite{4 Reg 20:4-7}
\switchcolumn
\EN
\Cite{2 Kgs 20:4-7}
\switchcolumn*

\LA
\noindent Et ántequam egrederétur Isaías médiam partem átrii, factus est sermo Dómini ad eum, dicens: Revértere et dic Ezechíæ duci pópuli mei: Hæc dicit Dóminus Deus David patris tui: Audívi oratiónem tuam et vidi lácrimas tuas et ecce sanávi te: die tértio ascéndes templum Dómini, et addam diébus tuis quíndecim annos; sed et de manu regis Assyriórum liberábo te et civitátem hanc et prótegam urbem istam propter me et propter David servum meum. Dixítque Isaías: Afférte massam ficórum. Quam cum attulíssent et posuíssent super ulcus ejus, curátus est.\par
\switchcolumn
\EN
\noindent And before Isaias was gone out of the middle of the court, the word of the Lord came to him, saying: Go back, and tell Ezechias the captain of my people: Thus saith the Lord the God of David thy father: I have heard thy prayer, and I have seen thy tears: and behold I have healed thee; on the third day thou shalt go up to the temple of the Lord. And I will add to thy days fifteen years: and I will deliver thee and this city out of the hand of the king of the Assyrians, and I will protect this city for my own sake, and for David my servant's sake. And Isaias said: Bring me a lump of figs. And when they had brought it, and laid it upon his boil. he was healed.\par
\switchcolumn*

\LA
\begin{Resp}\Rsign Deus ómnium exaudítor est: ipse misit Angelum suum, et tulit me de óvibus patris mei: \Star{} Et unxit me unctióne misericórdiæ suæ.\end{Resp}
\begin{Resp}\Vsign Dóminus, qui erípuit me de ore leónis, et de manu béstiæ liberávit me.\end{Resp}
\begin{Resp}\Rsign Et unxit me unctióne misericórdiæ suæ.\end{Resp}
\switchcolumn
\EN
\begin{Resp}\Rsign God, Which heareth all, even He sent His Angel, and took me from keeping my father's sheep, and \Star{} Anointed me with the oil of His mercy.\end{Resp}
\begin{Resp}\Vsign The Lord That delivered me out of the mouth of the lion, and out of the paw of the bear\end{Resp}
\begin{Resp}\Rsign And anointed me with the oil of His mercy.\end{Resp}
\switchcolumn*

\LA
\LessonHead{Lectio III}
\switchcolumn
\EN
\LessonHead{Lesson III}
\switchcolumn*

\LA
\Cite{4 Reg 20:8-11}
\switchcolumn
\EN
\Cite{2 Kgs 20:8-11}
\switchcolumn*

\LA
\noindent Díxerat autem Ezechías ad Isaíam: Quod erit signum quia Dóminus me sanábit, et quia ascensúrus sum die tértia templum Dómini? Cui ait Isaías: Hoc erit signum a Dómino, quod factúrus sit Dóminus sermónem quem locútus est: vis ut ascéndat umbra decem líneis, an ut revertátur tótidem grádibus? Et ait Ezechías: Fácile est umbram créscere decem líneis, nec hoc volo ut fiat; sed ut revertátur retrórsum decem grádibus. Invocávit ítaque Isaías prophéta Dóminum, et redúxit umbram per líneas, quibus jam descénderat in horológio Achaz, retrórsum decem grádibus.\par
\switchcolumn
\EN
\noindent And Ezechias had said to Isaias: What shall be the sign that the Lord will heal me, and that I shall go up to the temple of the Lord the third day? And Isaias said to him: This shall be the sign from the Lord, that the Lord will do the word which he hath spoken: Wilt thou that the shadow go forward ten lines, or that it go back so many degrees? And Ezechias said: It is an easy matter for the shadow to go forward ten lines: and I do not desire that this be done, but let it return back ten degrees. And Isaias the prophet called upon the Lord, and he brought the shadow ten degrees backwards by the lines, by which it had already gone down in the dial of Achaz.\par
\switchcolumn*

\LA
\begin{Resp}\Rsign Dóminus, qui erípuit me de ore leónis, et de manu béstiæ liberávit me, \Star{} Ipse me erípiet de mánibus inimicórum meórum.\end{Resp}
\begin{Resp}\Vsign Misit Deus misericórdiam suam et veritátem suam: ánimam meam erípuit de médio catulórum leónum.\end{Resp}
\begin{Resp}\Rsign Ipse me erípiet de mánibus inimicórum meórum.\end{Resp}
\begin{Resp}\Vsign Glória Patri, et Fílio, \Star{} et Spirítui Sancto.\end{Resp}
\begin{Resp}\Rsign Ipse me erípiet de mánibus inimicórum meórum.\end{Resp}
\switchcolumn
\EN
\begin{Resp}\Rsign The Lord That delivered me out of the mouth of the lion, and out of the paw of the bear \Star{} He will deliver me out of the hand of mine enemies.\end{Resp}
\begin{Resp}\Vsign God hath sent forth His mercy and His truth, and delivered my soul from among the lion's whelps.\end{Resp}
\begin{Resp}\Rsign He will deliver me out of the hand of mine enemies.\end{Resp}
\begin{Resp}\Vsign Glory be to the Father, and to the Son, \Star{} and to the Holy Ghost.\end{Resp}
\begin{Resp}\Rsign He will deliver me out of the hand of mine enemies.\end{Resp}
\switchcolumn*

\LA
\LessonHead{Lectio IV}
\switchcolumn
\EN
\LessonHead{Lesson IV}
\switchcolumn*

\LA
\LessonTitle{Ex Expositióne sancti Hierónymi Presbýteri in Isaíam Prophétam}
\switchcolumn
\EN
\LessonTitle{From the Exposition of the Prophet Isaiah written by St. Jerome, Priest at Bethlehem.}
\switchcolumn*

\LA
\Source{Lib. 11 in Isaíæ cap. 38}
\switchcolumn
\EN
\Source{Bk. ii. on Is. xxxviii.}
\switchcolumn*

\LA
\noindent Ne elevarétur cor Ezechíæ post incredíbiles triúmphos et de média captivitáte victóriam, infirmitáte córporis sui visitátur, et audit se esse moritúrum; ut convérsus ad Dóminum flectat senténtiam ejus. Quod quidem et in Jona prophéta légimus, et in comminatiónibus contra David: quæ dicúntur futúra, nec facta sunt, non Deo mutánte senténtiam, sed provocánte humánum genus ad notítiam sui. Dóminus enim pœ́nitens est super malítiis. Convertítque Ezechías fáciem suam ad paríetem, quia ad templum ire non póterat. Ad paríetem autem templi, juxta quod Sálomon palátium exstrúxerat; vel absolúte ad paríetem, ne lácrimas suas assidéntibus ostentáre viderétur.\par
\switchcolumn
\EN
\noindent Lest the heart of Hezekiah should be puffed up by his strange and unlooked for triumphs, and by his victory when he was but a prisoner, he was visited by bodily weakness, and told that he was to die that he might betake himself to the Lord, and turn Him from carrying out the sentence. We read of a like case in the history of the Prophet Jonah (ii. 4, 10.) And in regard to the threatening made against David (2 Kgs (Sam.) xxiv. 16) - when punishments were foretold which were not brought to pass. This is not because that God is a Being capable of changing His mind, but because He willeth to mankind to know Him, how that “He [is gracious and merciful, slow to anger, and of great kindness, and] repenteth Him of the evil.” (Joel ii. 13.) Hezekiah turned his face unto the wall, not being able to go up to the Temple. This may either mean that he turned towards the wall of the Temple, hard by which Solomon had built a palace, or simply, that he turned his face to the wall, so as not to parade his tears before his attendants.\par
\switchcolumn*

\LA
\begin{Resp}\Rsign Percússit Saul mille, et David decem míllia: \Star{} Quia manus Dómini erat cum illo: percússit Philisthǽum, et ábstulit oppróbrium ex Israël.\end{Resp}
\begin{Resp}\Vsign Nonne iste est David, de quo canébant in choro, dicéntes: Saul percússit mille, et David decem míllia?\end{Resp}
\begin{Resp}\Rsign Quia manus Dómini erat cum illo: percússit Philisthǽum, et ábstulit oppróbrium ex Israël.\end{Resp}
\switchcolumn
\EN
\begin{Resp}\Rsign Saul hath slain his thousands, and David his ten thousands. \Star{} Because the hand of the Lord was with him, he smote the Philistine, and took away the reproach from Israel.\end{Resp}
\begin{Resp}\Vsign Is not this David? Did they not sing one to another of him in dances, saying: Saul hath slain his thousands, and David his ten thousands?\end{Resp}
\begin{Resp}\Rsign Because the hand of the Lord was with him, he smote the Philistine, and took away the reproach from Israel.\end{Resp}
\switchcolumn*

\LA
\LessonHead{Lectio V}
\switchcolumn
\EN
\LessonHead{Lesson V}
\switchcolumn*

\LA
\noindent Audiénsque se esse moritúrum, non precátur vitam et annos plúrimos, sed in Dei judício, quid velit præstáre, dimíttit. Nóverat enim idcírco Deo placuísse Salomónem, quod annos vitæ non petíerit amplióres; sed itúrus ad Dóminum narrat ópera sua, quómodo ambuláverit coram eo in veritáte et in corde perfécto. Felix consciéntia, quæ afflictiónis témpore bonórum óperum recordátur. Beáti enim mundo corde: quóniam ipsi Deum vidébunt. Et quómodo álibi scríbitur: Quis gloriábitur purum habére se cor? Quod ita sólvitur: perfectiónem cordis in eo nunc dici, quod idóla destrúxerit, templi valvas aperúerit, serpéntem ǽneum comminúerit, et cétera fécerit quæ Scriptúra commémorat.\par
\switchcolumn
\EN
\noindent Having been told that he was about to die, he prayed not for life and many years, but left it to God to do as in His good judgment He was pleased to will. He knew how this had pleased God on the part of Solomon. (3 (1) Kings iii. 11.) So, when he betook him to the Lord, he only made mention of his works, how he had walked before Him in truth, and with a perfect heart. Happy is he whose conscience in the hour of affliction can assure him of good works. Yea, “blessed are the pure in heart, for they shall see God” (Matth. v. 8.) It is indeed written in another place “Who can say, I have made my heart clean, I am pure from my sin?” (Prov. xx. 9.) How then could Hezekiah say that he had walked with a perfect heart But the answer is, that by this is meant that he had destroyed the idols, opened the doors of the Temple, broken in pieces the brazen serpent, and done the rest of the things whereof the Scripture maketh mention.\par
\switchcolumn*

\LA
\begin{Resp}\Rsign Montes Gélboë, nec ros nec plúvia véniant super vos, \Star{} Ubi cecidérunt fortes Israël.\end{Resp}
\begin{Resp}\Vsign Omnes montes, qui estis in circúitu ejus, vísitet Dóminus: a Gélboë autem tránseat.\end{Resp}
\begin{Resp}\Rsign Ubi cecidérunt fortes Israël.\end{Resp}
\switchcolumn
\EN
\begin{Resp}\Rsign Ye mountains of Gilboa, let there be no dew, neither let there be rain upon you \Star{} For there are the mighty of Israel fallen.\end{Resp}
\begin{Resp}\Vsign All ye mountains that stand round about, the Lord look upon you but let Him pass by Gilboa\end{Resp}
\begin{Resp}\Rsign For there are the mighty of Israel fallen.\end{Resp}
\switchcolumn*

\LA
\LessonHead{Lectio VI}
\switchcolumn
\EN
\LessonHead{Lesson VI}
\switchcolumn*

\LA

\switchcolumn
\EN
\LessonTitle{“And Hezekiah wept sore.” He had then no children, and it seemed as though the promise which God had made unto David, [that “his seed should endure for ever, and his throne as the sun before” Him, (Ps. lxxxviii. 29,)] was about to fail in his own death. It is written that “Manasseh was twelve years old when he began to reign,” (xxi. 1) - whence it is evident that Hezekiah begat him not till after three years of his new lease of life. Sore therefore wept he, when all hope was torn from him that the Messiah should spring from his seed. Others again remark that he wept sore, since death terrifieth sometimes even the saints, since they know not what sentence is about to be pronounced upon them, and what place shall be allotted them in the inscrutable judgment.}
\switchcolumn*

\LA
\noindent Flevit autem fletu magno, propter promissiónem Dómini ad David, quam vidébat in sua morte peritúram. Eo enim témpore Ezechías fílios non habébat; nam post mortem ejus Manásses cum duódecim esset annórum, regnáre cœpit in Judǽa: ex quo perspícuum est, post tértium annum concéssæ vitæ Manássen esse generátum. Ergo iste omnis est fletus, quod desperábat Christum de suo sémine nascitúrum. Alii ásserunt, quamvis sanctos viros morte terréri, propter incértum judícii et ignoratiónem senténtiæ Dei, quam sedem habitúri sint.\par
\switchcolumn
\EN

\switchcolumn*

\LA
\begin{Resp}\Rsign Ego te tuli de domo patris tui, dicit Dóminus, et pósui te páscere gregem pópuli mei: \Star{} Et fui tecum in ómnibus ubicúmque ambulásti, firmans regnum tuum in ætérnum.\end{Resp}
\begin{Resp}\Vsign Fecíque tibi nomen grande, juxta nomen magnórum, qui sunt in terra: et réquiem dedi tibi ab ómnibus inimícis tuis.\end{Resp}
\begin{Resp}\Rsign Et fui tecum in ómnibus ubicúmque ambulásti, firmans regnum tuum in ætérnum.\end{Resp}
\begin{Resp}\Vsign Glória Patri, et Fílio, \Star{} et Spirítui Sancto.\end{Resp}
\begin{Resp}\Rsign Et fui tecum in ómnibus ubicúmque ambulásti, firmans regnum tuum in ætérnum.\end{Resp}
\switchcolumn
\EN
\begin{Resp}\Rsign Thus saith the Lord: I took thee out of thy father's house, and appointed thee to be ruler over My people, over Israel. \Star{} And I was with thee whithersoever thou wentest, to establish thy kingdom for ever.\end{Resp}
\begin{Resp}\Vsign And I have made thee a great name, like unto the name of the great men that are in the earth, and have caused thee to rest from all thine enemies.\end{Resp}
\begin{Resp}\Rsign And I was with thee whithersoever thou wentest, to establish thy kingdom for ever.\end{Resp}
\begin{Resp}\Vsign Glory be to the Father, and to the Son, \Star{} and to the Holy Ghost.\end{Resp}
\begin{Resp}\Rsign And I was with thee whithersoever thou wentest, to establish thy kingdom for ever.\end{Resp}
\switchcolumn*

\LA
\LessonHead{Lectio VII}
\switchcolumn
\EN
\LessonHead{Lesson VII}
\switchcolumn*

\LA
\LessonTitle{Léctio sancti Evangélii secúndum Marcum}
\switchcolumn
\EN
\LessonTitle{From the Holy Gospel according to Mark}
\switchcolumn*

\LA
\Cite{Marc 7:31-37}
\switchcolumn
\EN
\Cite{Mark 7:31-37}
\switchcolumn*

\LA
\noindent In illo témpore: Exiens Jesus de fínibus Tyri venit per Sidónem ad mare Galilǽæ inter médios fines Decapóleos. Et réliqua.\par
\switchcolumn
\EN
\noindent At that time: Jesus, departing from the coasts of Tyre and Sidon, came unto the sea of Galilee, through the midst of the coasts of Decapolis. And so on.\par
\switchcolumn*

\LA
\LessonTitle{Homilía sancti Gregórii Papæ}
\switchcolumn
\EN
\LessonTitle{Homily by Pope St. Gregory the Great}
\switchcolumn*

\LA
\Source{Homilia 10, liber 1 in Ezech., ante medium}
\switchcolumn
\EN
\Source{Hom. x. Bk. i. on Ezekiel}
\switchcolumn*

\LA
\noindent Quid est quod creátor ómnium Deus, cum surdum et mutum sanáre voluísset, in aures illíus suos dígitos misit, et éxspuens linguam ejus tétigit? Quid per dígitos Redemptóris, nisi dona Sancti Spíritus designántur? Unde cum in álio loco ejecísset dæmónium, dixit: Si in dígito Dei eício dæmónia, profécto pervénit in vos regnum Dei. Qua de re per Evangelístam álium dixísse descríbitur: Si ego in Spíritu Dei eício dǽmones, ígitur pervénit in vos regnum Dei. Ex quo utróque loco collígitur, quia dígitus Spíritus vocátur. Dígitos ergo in aurículas míttere, est per dona Spíritus Sancti mentem surdi ad obediéndum aperíre.\par
\switchcolumn
\EN
\noindent What signifieth it that when God, the Maker of all, would heal a deaf and dumb man, “He put His Fingers into his ears, and He spit, and touched his tongue.” What is figured by the Fingers of the Redeemer but the gifts of the Holy Ghost? Hence it is written in another place (Luke xi. 20) that after He had cast out an evil spirit, He said: “If I with the finger of God cast out devils, no doubt the kingdom of God is come upon you.” Which words are thus given by another Evangelist (Matth. xii. 28): “If I cast out devils by the Spirit of God, then the kingdom of God is come unto you.” By setting these two passages together we see that the Spirit is called the Finger. For our Lord, then, to put His Fingers into the deaf man's ears was by the gift of the Holy Spirit to enlighten his dark mind unto obedience.\par
\switchcolumn*

\LA
\begin{Resp}\Rsign Dómine, Pater et Deus vitæ meæ, ne derelínquas me in cogitátu malígno: extolléntiam oculórum meórum ne déderis mihi, et desidérium malígnum avérte a me, Dómine; aufer a me concupiscéntiam, \Star{} Et ánimo irreverénti et infruníto ne tradas me, Dómine.\end{Resp}
\begin{Resp}\Vsign Ne derelínquas me, Dómine, ne accréscant ignorántiæ meæ, nec multiplicéntur delícta mea.\end{Resp}
\begin{Resp}\Rsign Et ánimo irreverénti et infruníto ne tradas me, Dómine.\end{Resp}
\switchcolumn
\EN
\begin{Resp}\Rsign O Lord, Father and God of my life, leave me not to evil counsels; give me not a proud look, but turn away from me an haughty mind, O Lord turn away from me concupiscence, \Star{} And give me not over unto an impudent and froward mind, O Lord!\end{Resp}
\begin{Resp}\Vsign Leave me not, O Lord, lest mine ignorance increase, and my sins abound.\end{Resp}
\begin{Resp}\Rsign And give me not over unto an impudent and froward mind, O Lord.\end{Resp}
\switchcolumn*

\LA
\LessonHead{Lectio VIII}
\switchcolumn
\EN
\LessonHead{Lesson VIII}
\switchcolumn*

\LA
\noindent Quid est vero, quod éxspuens linguam ejus tétigit? Salíva nobis est ex ore Redemptóris, accépta sapiéntia in elóquio divíno. Salíva quippe ex cápite défluit in ore. Ea ergo sapiéntia, quæ ipse est, dum lingua nostra tángitur, mox ad prædicatiónis verba formátur. Qui suspíciens in cælum, ingémuit: non quod ipse necessárium gémitum habéret, qui dabat quod postulábat; sed nos ad eum gémere, qui cælo prǽsidet, dócuit: ut et aures nostræ per donum Spíritus Sancti aperíri, et lingua per salívam oris, id est, per sciéntiam divínæ locutiónis, solvi débeat ad verba prædicatiónis.\par
\switchcolumn
\EN
\noindent That signifieth it also that “He spit and touched his tongue.” We receive spittle out of the Redeemer's mouth upon our tongues when we receive wisdom to speak God's truth. Spittle is a secretion of the head which floweth into the mouth. And so, that wisdom, which is Himself, the great Head of His Church, as soon as it hath touched our tongue, doth straightway take the form of preaching. “And looking up to heaven, He sighed,” not that He had any need to sigh, Who gave whatsoever He asked, but that He was fain to teach us to look up and sigh toward Him Whose throne is in heaven, confessing our need, that our ears should be opened by the gift of the Holy Spirit, and our tongue loosed by the spittle of our Saviour's Mouth, that is, by knowledge of His Divine Word, before we can use it to preach to others.\par
\switchcolumn*

\LA
\begin{Resp}\Rsign Duo Séraphim clamábant alter ad álterum: \Star{} Sanctus, sanctus, sanctus Dóminus Deus Sábaoth: * Plena est omnis terra glória ejus.\end{Resp}
\begin{Resp}\Vsign Tres sunt qui testimónium dant in cælo: Pater, Verbum, et Spíritus Sanctus: et hi tres unum sunt.\end{Resp}
\begin{Resp}\Rsign Sanctus, sanctus, sanctus Dóminus Deus Sábaoth.\end{Resp}
\begin{Resp}\Vsign Glória Patri, et Fílio, \Star{} et Spirítui Sancto.\end{Resp}
\begin{Resp}\Rsign Plena est omnis terra glória ejus.\end{Resp}
\switchcolumn
\EN
\begin{Resp}\Rsign One Seraph cried unto another: \Star{} Holy, Holy, Holy is the Lord God of hosts; the whole earth is full of His glory.\end{Resp}
\begin{Resp}\Vsign There are Three That bear record in heaven: the Father, the Word, and the Holy Ghost, and these Three are One.\end{Resp}
\begin{Resp}\Rsign Holy, Holy, Holy is the Lord God of hosts.\end{Resp}
\begin{Resp}\Vsign Glory be to the Father, and to the Son, \Star{} and to the Holy Ghost.\end{Resp}
\begin{Resp}\Rsign The whole earth is full of His glory.\end{Resp}
\switchcolumn*

\LA
\LessonHead{Lectio IX}
\switchcolumn
\EN
\LessonHead{Lesson IX}
\switchcolumn*

\LA

\switchcolumn
\EN
\LessonTitle{“And He said unto him: Ephphatha, that is, be opened. And straightway his ears were opened, and the string of his tongue was loosed.” Herein we must remark the command, “Be opened” was addressed to the deaf ears, but the tongue also was immediately loosed. Just so, when the ears of a man's heart have been opened to learn the obedience of faith, the string of his tongue also is thereupon loosed, that he may exhort others to do the good things which himself doth. It is well added “And he spake plain.” He only doth well preach obedience to others who hath first learnt himself to obey.}
\switchcolumn*

\LA
\noindent Cui mox, Ephphétha, id est, Adaperíre dícitur: et statim apértæ sunt aures ejus, et solútum est vínculum linguæ ejus. Qua in re notándum est, quia propter clausas aures dictum est, Adaperíre. Sed cui aures cordis ad obediéndum apértæ fúerint, ex subsequénti procul dúbio étiam linguæ ejus vínculum sólvitur; ut bona quæ ipse fécerit, étiam faciénda áliis loquátur. Ubi bene ádditur: Et loquebátur recte. Ille enim recte lóquitur, qui prius obediéndo fécerit quæ loquéndo ádmonet esse faciénda.\par
\switchcolumn
\EN

\switchcolumn*

\LA
\TeDeum{Te Deum}
\switchcolumn
\EN
\TeDeum{Te Deum}
\switchcolumn*

\end{paracol}

\Office{Dominica XII Post Pentecosten}{Twelfth Sunday after Pentecost}{Semiduplex}
\addcontentsline{toc}{subsection}{Dominica XII Post Pentecosten \textit{— Twelfth Sunday after Pentecost}}
\begin{paracol}{2}
\LA
\RefLine{Lectiones I–VI: de Scriptura occurrente.}
\switchcolumn
\EN
\RefLine{Lessons I–VI: from the Scripture of the day.}
\switchcolumn*

\LA
\LessonHead{Lectio VII}
\switchcolumn
\EN
\LessonHead{Lesson VII}
\switchcolumn*

\LA
\LessonTitle{Léctio sancti Evangélii secúndum Lucam}
\switchcolumn
\EN
\LessonTitle{From the Holy Gospel according to Luke}
\switchcolumn*

\LA
\Cite{Luc 10:23-37}
\switchcolumn
\EN
\Cite{Luke 10:23-37}
\switchcolumn*

\LA
\noindent In illo témpore: Dixit Jesus discípulis suis: Beáti óculi qui vident quæ vos vidétis; dico enim vobis quod multi prophétæ et reges voluérunt vidére quæ vos vidétis, et non vidérunt. Et réliqua.\par
\switchcolumn
\EN
\noindent At that time, Jesus said unto His disciples: Blessed are the eyes which see the things that ye see. For I tell you that many prophets and kings have desired to see those things which ye see, and have not seen them. And so on.\par
\switchcolumn*

\LA
\LessonTitle{Homilía sancti Bedæ Venerábilis Presbýteri}
\switchcolumn
\EN
\LessonTitle{Homily by the Venerable Bede, Priest (at Jarrow.)}
\switchcolumn*

\LA
\Source{Liber 3, cap. 43 in Lucæ 10}
\switchcolumn
\EN
\Source{Bk. iii. ch. 43 on Luke x.}
\switchcolumn*

\LA
\noindent Non óculi scribárum et pharisæórum, qui corpus tantum Dómini vidére; sed illi beáti óculi, qui ejus possunt cognóscere sacraménta, de quibus dícitur: Et revelásti ea párvulis. Beáti óculi parvulórum, quibus et se et Patrem Fílius reveláre dignátur. Abraham exsultávit ut vidéret diem Christi; et vidit, et gavísus est. Isaías quoque, et Michǽas, et multi álii prophétæ vidérunt glóriam Dómini, qui et proptérea Vidéntes sunt appelláti; sed hi omnes, a longe aspiciéntes et salutántes, per spéculum et in ænígmate vidérunt.\par
\switchcolumn
\EN
\noindent Blessed were the eyes not of Scribes and Pharisees, which saw but the Body of the Lord, but those eyes, eyes blessed indeed, which were able to see those things whereof it is written “Thou hast hid these things from the wise and prudent, and hast revealed them unto babes.” Blessed are the eyes of those little ones unto whom it seemeth good in the eyes of the Son to reveal Himself and the Father also. Abraham rejoiced to see the day of Christ and he saw it, and was glad. (John viii. 56.) Isaiah, and Micah, and many among the Prophets, saw the glory of the Lord, wherefore also they be called Seers, but all they beheld it and hailed it afar off, seeing but as through a glass, darkly. (1 Cor. xiii. 12.)\par
\switchcolumn*

\LA
\begin{Resp}\Rsign Dómine, Pater et Deus vitæ meæ, ne derelínquas me in cogitátu malígno: extolléntiam oculórum meórum ne déderis mihi, et desidérium malígnum avérte a me, Dómine; aufer a me concupiscéntiam, \Star{} Et ánimo irreverénti et infruníto ne tradas me, Dómine.\end{Resp}
\begin{Resp}\Vsign Ne derelínquas me, Dómine, ne accréscant ignorántiæ meæ, nec multiplicéntur delícta mea.\end{Resp}
\begin{Resp}\Rsign Et ánimo irreverénti et infruníto ne tradas me, Dómine.\end{Resp}
\switchcolumn
\EN
\begin{Resp}\Rsign O Lord, Father and God of my life, leave me not to evil counsels; give me not a proud look, but turn away from me an haughty mind, O Lord turn away from me concupiscence, \Star{} And give me not over unto an impudent and froward mind, O Lord!\end{Resp}
\begin{Resp}\Vsign Leave me not, O Lord, lest mine ignorance increase, and my sins abound.\end{Resp}
\begin{Resp}\Rsign And give me not over unto an impudent and froward mind, O Lord.\end{Resp}
\switchcolumn*

\LA
\LessonHead{Lectio VIII}
\switchcolumn
\EN
\LessonHead{Lesson VIII}
\switchcolumn*

\LA
\noindent Apóstoli autem, in præsentiárum habéntes Dóminum, convescentésque ei, et quæcúmque voluíssent interrogándo discéntes, nequáquam per Angelos aut várias visiónum spécies opus habébant docéri. Quos vero Lucas multos prophétas et reges dicit, Matthǽus apértius prophétas et justos appéllat. Ipse sunt enim reges magni; quia tentatiónum suárum mótibus non consentiéndo succúmbere, sed regéndo præésse novérunt.\par
\switchcolumn
\EN
\noindent Otherwise were the Apostles, who saw the Lord face to Face, eating with Him, and learning from Him by asking whatsoever they listed. For them there was no need to be taught by Angels, or the shifting fabric of visions. They whom Luke doth call Prophets and kings, Matthew nameth as “Prophets and righteous men” (xiii. 17.) Righteous men are indeed mighty kings, who know how to lord it over their own rebellious temptations, instead of falling under them to become their slaves.\par
\switchcolumn*

\LA
\begin{Resp}\Rsign Duo Séraphim clamábant alter ad álterum: \Star{} Sanctus, sanctus, sanctus Dóminus Deus Sábaoth: * Plena est omnis terra glória ejus.\end{Resp}
\begin{Resp}\Vsign Tres sunt qui testimónium dant in cælo: Pater, Verbum, et Spíritus Sanctus: et hi tres unum sunt.\end{Resp}
\begin{Resp}\Rsign Sanctus, sanctus, sanctus Dóminus Deus Sábaoth.\end{Resp}
\begin{Resp}\Vsign Glória Patri, et Fílio, \Star{} et Spirítui Sancto.\end{Resp}
\begin{Resp}\Rsign Plena est omnis terra glória ejus.\end{Resp}
\switchcolumn
\EN
\begin{Resp}\Rsign One Seraph cried unto another: \Star{} Holy, Holy, Holy is the Lord God of hosts; the whole earth is full of His glory.\end{Resp}
\begin{Resp}\Vsign There are Three That bear record in heaven: the Father, the Word, and the Holy Ghost, and these Three are One.\end{Resp}
\begin{Resp}\Rsign Holy, Holy, Holy is the Lord God of hosts.\end{Resp}
\begin{Resp}\Vsign Glory be to the Father, and to the Son, \Star{} and to the Holy Ghost.\end{Resp}
\begin{Resp}\Rsign The whole earth is full of His glory.\end{Resp}
\switchcolumn*

\LA
\LessonHead{Lectio IX}
\switchcolumn
\EN
\LessonHead{Lesson IX}
\switchcolumn*

\LA

\switchcolumn
\EN
\LessonTitle{“And, behold, a certain lawyer stood up, and tempted Him, saying Master, what shall I do to inherit eternal life?” This lawyer, who stood up to ask the Lord a tempting question touching eternal life, took the subject of his asking, as I think, from the words which the Lord had just uttered, when He said “Rejoice, because your names are written in heaven”. But his attempt was a proof of the truth of that which the Lord immediately added “I thank thee, O Father, Lord of heaven and earth, that Thou hast hid these things from the wise and prudent, and hast revealed them unto babes!”}
\switchcolumn*

\LA
\noindent Et ecce quidam legisperítus surréxit, tentans eum et dicens: Magíster, quid faciéndo vitam ætérnam possidébo? Legisperítus, qui de vita ætérna Dóminum tentans intérrogat, occasiónem, ut reor, tentándi de ipsis Dómini sermónibus sumpsit, ubi ait: Gaudéte autem quod nómina vestra scripta sunt in cælis. Sed ipsa sua tentatióne declárat quam vera sit illa Dómini conféssio, qua Patri lóquitur: Quod abscondísti hæc a sapiéntibus, et prudéntibus, et revelásti ea párvulis.\par
\switchcolumn
\EN

\switchcolumn*

\LA
\TeDeum{Te Deum}
\switchcolumn
\EN
\TeDeum{Te Deum}
\switchcolumn*

\end{paracol}

\Office{Dominica XIII Post Pentecosten}{Thirteenth Sunday after Pentecost}{Semiduplex}
\addcontentsline{toc}{subsection}{Dominica XIII Post Pentecosten \textit{— Thirteenth Sunday after Pentecost}}
\begin{paracol}{2}
\LA
\RefLine{Lectiones I–VI: de Scriptura occurrente.}
\switchcolumn
\EN
\RefLine{Lessons I–VI: from the Scripture of the day.}
\switchcolumn*

\LA
\LessonHead{Lectio VII}
\switchcolumn
\EN
\LessonHead{Lesson VII}
\switchcolumn*

\LA
\LessonTitle{Léctio sancti Evangélii secúndum Lucam}
\switchcolumn
\EN
\LessonTitle{From the Holy Gospel according to Luke}
\switchcolumn*

\LA
\Cite{Luc 17:11-19}
\switchcolumn
\EN
\Cite{Luke 17:11-19}
\switchcolumn*

\LA
\noindent In illo témpore: Dum iret Jesus in Jerúsalem, transíbat per médiam Samaríam et Galilǽam. Et cum ingrederétur quoddam castéllum, occurrérunt ei decem viri leprósi. Et réliqua.\par
\switchcolumn
\EN
\noindent It came to pass, as Jesus went to Jerusalem, that He passed through the midst of Samaria and Galilee. And, as He entered into a certain village, there met Him ten men that were lepers. And so on.\par
\switchcolumn*

\LA
\LessonTitle{Homilía sancti Augustíni Epíscopi}
\switchcolumn
\EN
\LessonTitle{Homily by St. Augustine, Bishop of Hippo}
\switchcolumn*

\LA
\Source{Lib. 2 quæst. Evang. cap. 40}
\switchcolumn
\EN
\Source{Bk. ii, Gospel Questions, ch. 40}
\switchcolumn*

\LA
\noindent De decem leprósis, quos Dóminus ita mundávit, cum ait: Ite, osténdite vos sacerdótibus; quæri potest, cur eos ad sacerdótes míserit, ut cum irent, mundaréntur. Nullum enim eórum, quibus hæc corporália benefícia prǽstitit, invenítur misísse ad sacerdótes, nisi leprósos. Nam et illum a lepra mundáverat, cui dixit: Vade, osténde te sacerdótibus, et offer pro te sacrifícium, quod præcépit Móyses, in testimónium illis. Quæréndum ígitur est, quid ipsa lepra signíficet: non enim sanáti, sed mundáti dicúntur, qui ea caruérunt. Colóris quippe vítium est, non valetúdinis, aut integritátis sénsuum atque membrórum.\par
\switchcolumn
\EN
\noindent The ten lepers “lifted up their voices and said: Jesus, Master, have mercy on us. And when He saw them, He said unto them: Go, show yourselves unto the Priests. And it came to pass that, as they went, they were cleansed.” Question: why did the Lord send them unto the Priests, that, as they went, they might be cleansed? Lepers were the only class among those upon whose bodies He worked mercy, whom we find that He sent unto the Priests. It is written in another place that He said to a leper whom He had cleansed: “Go, and show thyself to the Priest, and offer for thy cleansing according as Moses commanded, for a testimony unto them” (Luke v. 14, Lev. xiv. seq.) We ask then, of what leprosy was a type, whereof they that were ridded were called, not “healed,” but “cleansed.” It is a disease which doth first appear in the skin, but destroyeth not immediately the strength, nor the use of feeling and the limbs.\par
\switchcolumn*

\LA
\begin{Resp}\Rsign Quis mihi tríbuat, ut in inférno prótegas me et abscóndas me, donec pertránseat furor tuus, Dómine, nisi tu, qui solus es Deus? \Star{} Et constítuas mihi tempus, in quo recordéris mei?\end{Resp}
\begin{Resp}\Vsign Numquid sicut dies hóminis dies tui, ut quæras iniquitátem meam; cum sit nemo, qui de manu tua possit erúere.\end{Resp}
\begin{Resp}\Rsign Et constítuas mihi tempus, in quo recordéris mei?\end{Resp}
\switchcolumn
\EN
\begin{Resp}\Rsign O that Thou wouldest hide me in the grave; that Thou wouldest keep me secret, until thy wrath be past, even thine, O Lord, Thou That alone art God; \Star{} That Thou wouldest appoint me a set time, and remember me!\end{Resp}
\begin{Resp}\Vsign Are thy days as the days of man, that Thou inquirest after mine iniquity, and there is none that can deliver out of thine hand.\end{Resp}
\begin{Resp}\Rsign That Thou wouldest appoint me a set time, and remember me!\end{Resp}
\switchcolumn*

\LA
\LessonHead{Lectio VIII}
\switchcolumn
\EN
\LessonHead{Lesson VIII}
\switchcolumn*

\LA
\noindent Leprósi ergo non absúrde intélligi possunt, qui sciéntiam veræ fídei non habéntes, várias doctrínas profiténtur erróris. Non enim abscóndunt imperítiam suam; sed pro summa perítia próferunt in lucem, et jactántia sermónis osténtant. Nulla porro falsa doctrína est, quæ non áliqua vera intermísceat. Vera ergo falsis inordináte permíxta, in una disputatióne vel narratióne hóminis, tamquam in uníus córporis colóre apparéntia, signíficant lepram, tamquam veris falsísque colórum variántem atque maculántem.\par
\switchcolumn
\EN
\noindent The lepers, therefore, we may not absurdly suppose such to be figured as have not the knowledge of the true faith, but do show forth diverse-coloured teachings of error. They hide not their witlessness, but do use all such wit as they have to make it manifest, and proclaim it in high-sounding phrases. There is no false doctrine but hath some truth mixed up with it. A man's discourse then, with some truths in it unequally mingled with falsehoods, and all confounded in one mass, is like to the body of one that is stricken with leprosy, whereon all manner of foul colours do appear in this and that place along with the true colour of skin.\par
\switchcolumn*

\LA
\begin{Resp}\Rsign Duo Séraphim clamábant alter ad álterum: \Star{} Sanctus, sanctus, sanctus Dóminus Deus Sábaoth: * Plena est omnis terra glória ejus.\end{Resp}
\begin{Resp}\Vsign Tres sunt qui testimónium dant in cælo: Pater, Verbum, et Spíritus Sanctus: et hi tres unum sunt.\end{Resp}
\begin{Resp}\Rsign Sanctus, sanctus, sanctus Dóminus Deus Sábaoth.\end{Resp}
\begin{Resp}\Vsign Glória Patri, et Fílio, \Star{} et Spirítui Sancto.\end{Resp}
\begin{Resp}\Rsign Plena est omnis terra glória ejus.\end{Resp}
\switchcolumn
\EN
\begin{Resp}\Rsign One Seraph cried unto another: \Star{} Holy, Holy, Holy is the Lord God of hosts; the whole earth is full of His glory.\end{Resp}
\begin{Resp}\Vsign There are Three That bear record in heaven: the Father, the Word, and the Holy Ghost, and these Three are One.\end{Resp}
\begin{Resp}\Rsign Holy, Holy, Holy is the Lord God of hosts.\end{Resp}
\begin{Resp}\Vsign Glory be to the Father, and to the Son, \Star{} and to the Holy Ghost.\end{Resp}
\begin{Resp}\Rsign The whole earth is full of His glory.\end{Resp}
\switchcolumn*

\LA
\LessonHead{Lectio IX}
\switchcolumn
\EN
\LessonHead{Lesson IX}
\switchcolumn*

\LA
\noindent Hi autem tam vitándi sunt Ecclésiæ, ut, si fíeri potest, lóngius remóti, magno clamóre Christum interpéllent; sicut isti decem stetérunt a longe, et levavérunt vocem, dicéntes: Jesu præcéptor, miserére nostri. Nam et quod præceptórem vocant, quo nómine néscio utrum quisquam Dóminum interpelláverit pro medicína corporáli; satis puto significáre, lepram falsam esse doctrínam, quam bonus præcéptor abstérgit.\par
\switchcolumn
\EN
\noindent Such men as these are banished out of the walls of the Church, to the end that haply when they stand afar off they may lift up their voices and cry to Christ for pardon, just as those ten men that were lepers, which stood afar off, outside the village, lifted up their voices and said “Jesus, Master, have mercy on us.” That they styled Him Master, by which title I know not if any besought the Lord for bodily healing, I think doth sufficiently show that leprosy signifieth false doctrine, whereof the Good Master doth cleanse us.\par
\switchcolumn*

\LA
\TeDeum{Te Deum}
\switchcolumn
\EN
\TeDeum{Te Deum}
\switchcolumn*

\end{paracol}

\Office{Dominica XIV Post Pentecosten}{Fourteenth Sunday after Pentecost}{Semiduplex}
\addcontentsline{toc}{subsection}{Dominica XIV Post Pentecosten \textit{— Fourteenth Sunday after Pentecost}}
\begin{paracol}{2}
\LA
\RefLine{Lectiones I–VI: de Scriptura occurrente.}
\switchcolumn
\EN
\RefLine{Lessons I–VI: from the Scripture of the day.}
\switchcolumn*

\LA
\LessonHead{Lectio VII}
\switchcolumn
\EN
\LessonHead{Lesson VII}
\switchcolumn*

\LA
\LessonTitle{Léctio sancti Evangélii secúndum Matthǽum}
\switchcolumn
\EN
\LessonTitle{From the Holy Gospel according to Matthew}
\switchcolumn*

\LA
\Cite{Matt 6:24-33}
\switchcolumn
\EN
\Cite{Matt 6:24-33}
\switchcolumn*

\LA
\noindent In illo témpore: Dixit Jesus discípulis suis: Nemo potest duóbus dóminis servíre. Et réliqua.\par
\switchcolumn
\EN
\noindent At that time, Jesus said unto His disciples: No man can serve two masters. And so on.\par
\switchcolumn*

\LA
\LessonTitle{Homilía sancti Augustíni Epíscopi}
\switchcolumn
\EN
\LessonTitle{Homily by St. Augustine, Bishop of Hippo.}
\switchcolumn*

\LA
\Source{Lib. 2 de Sermone Dómini in monte, cap. 14}
\switchcolumn
\EN
\Source{Bk. ii. on the Lord's Sermon on the Mount, ch. xiv.}
\switchcolumn*

\LA
\noindent Nemo potest duóbus dóminis servíre. Ad hanc ipsam intentiónem referéndum est quod consequénter expónit, dicens: Aut enim unum ódio habébit, et álterum díliget; aut álterum patiétur, et álterum contémnet. Quæ verba diligénter consideránda sunt: nam, qui sint duo dómini, deínceps osténdit, cum dicit: Non potéstis Deo servíre, et mammónæ. Mammóna apud Hebrǽos divítiæ appellári dicúntur. Cóngruit et Púnicum nomen; nam lucrum Púnice mammon dícitur.\par
\switchcolumn
\EN
\noindent “No man can serve two masters,” and this is further explained “for either he will hate the one, and love the other; or else he will hold to the one, and despise the other.” These words we ought carefully to weigh, for the Lord showeth straightway who be the two masters whom we have choice of: “Ye cannot serve God and Mammon.” Mammon is a term which the Hebrews are said to use for riches. It is also a Carthaginian word for the Punic for “gain” is “mammon.”\par
\switchcolumn*

\LA
\begin{Resp}\Rsign Quis mihi tríbuat, ut in inférno prótegas me et abscóndas me, donec pertránseat furor tuus, Dómine, nisi tu, qui solus es Deus? \Star{} Et constítuas mihi tempus, in quo recordéris mei?\end{Resp}
\begin{Resp}\Vsign Numquid sicut dies hóminis dies tui, ut quæras iniquitátem meam; cum sit nemo, qui de manu tua possit erúere.\end{Resp}
\begin{Resp}\Rsign Et constítuas mihi tempus, in quo recordéris mei?\end{Resp}
\switchcolumn
\EN
\begin{Resp}\Rsign O that Thou wouldest hide me in the grave; that Thou wouldest keep me secret, until thy wrath be past, even thine, O Lord, Thou That alone art God; \Star{} That Thou wouldest appoint me a set time, and remember me!\end{Resp}
\begin{Resp}\Vsign Are thy days as the days of man, that Thou inquirest after mine iniquity, and there is none that can deliver out of thine hand.\end{Resp}
\begin{Resp}\Rsign That Thou wouldest appoint me a set time, and remember me!\end{Resp}
\switchcolumn*

\LA
\LessonHead{Lectio VIII}
\switchcolumn
\EN
\LessonHead{Lesson VIII}
\switchcolumn*

\LA
\noindent Sed qui servit mammónæ, illi útique servit, qui rebus istis terrénis mérito suæ perversitátis præpósitus, magistrátus hujus sǽculi a Dómino dícitur. Aut enim unum ódio habébit homo, et álterum díliget, id est, Deum; aut álterum patiétur, et álterum contémnet. Patiétur enim durum et perniciósum dóminum, quisquis servit mammónæ; sua enim cupiditáte implicátus, súbditur diábolo, et non eum díligit. Quis enim est qui díligat diábolum? sed tamen pátitur.\par
\switchcolumn
\EN
\noindent He which serveth mammon, serveth that evil one who hath perversely chosen to be lord of these earthly things, and is called by the Lord “the prince of this world.” (John xiv. 30.) Of these two masters, either a man will hate the one and love the other, that is God or he will hold to the one and despise the other. He which serveth mammon holdeth to an hard and destroying master, for he is led captive by his lust, and sold a slave to the devil, and him loveth no man is there any man that loveth the devil And yet there be that hold to him.\par
\switchcolumn*

\LA
\begin{Resp}\Rsign Duo Séraphim clamábant alter ad álterum: \Star{} Sanctus, sanctus, sanctus Dóminus Deus Sábaoth: * Plena est omnis terra glória ejus.\end{Resp}
\begin{Resp}\Vsign Tres sunt qui testimónium dant in cælo: Pater, Verbum, et Spíritus Sanctus: et hi tres unum sunt.\end{Resp}
\begin{Resp}\Rsign Sanctus, sanctus, sanctus Dóminus Deus Sábaoth.\end{Resp}
\begin{Resp}\Vsign Glória Patri, et Fílio, \Star{} et Spirítui Sancto.\end{Resp}
\begin{Resp}\Rsign Plena est omnis terra glória ejus.\end{Resp}
\switchcolumn
\EN
\begin{Resp}\Rsign One Seraph cried unto another: \Star{} Holy, Holy, Holy is the Lord God of hosts; the whole earth is full of His glory.\end{Resp}
\begin{Resp}\Vsign There are Three That bear record in heaven: the Father, the Word, and the Holy Ghost, and these Three are One.\end{Resp}
\begin{Resp}\Rsign Holy, Holy, Holy is the Lord God of hosts.\end{Resp}
\begin{Resp}\Vsign Glory be to the Father, and to the Son, \Star{} and to the Holy Ghost.\end{Resp}
\begin{Resp}\Rsign The whole earth is full of His glory.\end{Resp}
\switchcolumn*

\LA
\LessonHead{Lectio IX}
\switchcolumn
\EN
\LessonHead{Lesson IX}
\switchcolumn*

\LA

\switchcolumn
\EN
\LessonTitle{“Therefore, I say unto you, Take no thought for your life, what ye shall eat, or what ye shall drink; nor yet for your body, what ye shall put on”, lest, albeit such things are not idle, but needful to be sought after, yet the seeking for things even needful should divide the heart and our intention should be corrupted when we do something as it were mercifully that is, lest, when we would seem to be seeking another's good, it should be profit to ourselves, rather than benefit to him, that we seek and therefore we seem not to ourselves to sin, because we would seek things not idle, but needful.}
\switchcolumn*

\LA
\noindent Ideo, inquit, dico vobis, non habére sollicitúdinem ánimæ vestræ quid edátis, neque córpori vestro quid induátis; ne forte, quamvis jam supérflua non quærántur, propter ipsa necessária cor duplicétur, et ad ipsa conquirénda, nostra detorqueátur inténtio, cum áliquid quasi misericórditer operámur: id est, ut cum consúlere alícui vidéri vólumus, nostrum emoluméntum ibi pótius, quam illíus utilitátem attendámus; et ídeo nobis non videámur peccáre, quia non supérflua, sed necessária sunt, quæ cónsequi vólumus.\par
\switchcolumn
\EN

\switchcolumn*

\LA
\TeDeum{Te Deum}
\switchcolumn
\EN
\TeDeum{Te Deum}
\switchcolumn*

\end{paracol}

\Office{Dominica XV Post Pentecosten}{Fifteenth Sunday after Pentecost}{Semiduplex}
\addcontentsline{toc}{subsection}{Dominica XV Post Pentecosten \textit{— Fifteenth Sunday after Pentecost}}
\begin{paracol}{2}
\LA
\RefLine{Lectiones I–VI: de Scriptura occurrente.}
\switchcolumn
\EN
\RefLine{Lessons I–VI: from the Scripture of the day.}
\switchcolumn*

\LA
\LessonHead{Lectio VII}
\switchcolumn
\EN
\LessonHead{Lesson VII}
\switchcolumn*

\LA
\LessonTitle{Léctio sancti Evangélii secúndum Lucam}
\switchcolumn
\EN
\LessonTitle{From the Holy Gospel according to Luke}
\switchcolumn*

\LA
\Cite{Luc 7:11-16}
\switchcolumn
\EN
\Cite{Luke 7:11-16}
\switchcolumn*

\LA
\noindent In illo témpore: Ibat Jesus in civitátem, quæ vocátur Naim; et ibant cum eo discípuli ejus, et turba copiósa. Et réliqua.\par
\switchcolumn
\EN
\noindent At that time: Jesus went into a city called Nain and His disciples went with Him, and much people. And so on.\par
\switchcolumn*

\LA
\LessonTitle{Homilía sancti Augustíni Epíscopi}
\switchcolumn
\EN
\LessonTitle{Homily by St. Augustine, Bishop of Hippo}
\switchcolumn*

\LA
\Source{Sermo 44 de verbis Domini, circa initium}
\switchcolumn
\EN
\Source{44th Discourse on the Words of the Lord}
\switchcolumn*

\LA
\noindent De júvene illo resuscitáto gavísa est mater vídua; de homínibus in spíritu cotídie suscitátis gaudet mater Ecclésia. Ille quidem mórtuus erat córpore; illi autem mente. Illíus mors visíbilis visibíliter plangebátur; illórum mors invisíbilis nec quærebátur, nec videbátur. Quæsívit ille, qui nóverat mórtuos. Ille solus nóverat mórtuos, qui póterat fácere vivos. Nisi enim ad mórtuos suscitándos venísset, Apóstolus non díceret: Surge, qui dormis, et exsúrge a mórtuis, et illuminábit te Christus.\par
\switchcolumn
\EN
\noindent That her son was called again to life was the joy of that widowed mother; that souls of men are every day called to life is the joy of our Mother the Church. He was dead in body they have been dead in mind. His death was outward, and was outwardly bewailed; their inward. Death hath been neither mourned for nor seen. But He hath sought for them, Who hath seen that they are dead, and He only hath seen that they are dead, Who hath been able to make them alive. If He had not come to raise the dead, the Apostle had not said: “Awake, thou that sleepest, and arise from the dead, and Christ shall give thee light.” (Eph. v. 14.)\par
\switchcolumn*

\LA
\begin{Resp}\Rsign Quis mihi tríbuat, ut in inférno prótegas me et abscóndas me, donec pertránseat furor tuus, Dómine, nisi tu, qui solus es Deus? \Star{} Et constítuas mihi tempus, in quo recordéris mei?\end{Resp}
\begin{Resp}\Vsign Numquid sicut dies hóminis dies tui, ut quæras iniquitátem meam; cum sit nemo, qui de manu tua possit erúere.\end{Resp}
\begin{Resp}\Rsign Et constítuas mihi tempus, in quo recordéris mei?\end{Resp}
\switchcolumn
\EN
\begin{Resp}\Rsign O that Thou wouldest hide me in the grave; that Thou wouldest keep me secret, until thy wrath be past, even thine, O Lord, Thou That alone art God; \Star{} That Thou wouldest appoint me a set time, and remember me!\end{Resp}
\begin{Resp}\Vsign Are thy days as the days of man, that Thou inquirest after mine iniquity, and there is none that can deliver out of thine hand.\end{Resp}
\begin{Resp}\Rsign That Thou wouldest appoint me a set time, and remember me!\end{Resp}
\switchcolumn*

\LA
\LessonHead{Lectio VIII}
\switchcolumn
\EN
\LessonHead{Lesson VIII}
\switchcolumn*

\LA
\noindent Tres autem mórtuos invenímus a Dómino resuscitátos visibíliter, míllia invisibíliter. Quot autem mórtuos visibíliter suscitáverit, quis novit? Non enim ómnia, quæ fecit, scripta sunt. Joánnes hoc dixit: Multa ália fecit Jesus, quæ si scripta essent, árbitror totum mundum non posse libros cápere. Multi ergo sunt álii sine dúbio suscitáti, sed non tres frustra commemoráti. Dóminus enim noster Jesus Christus ea quæ faciébat corporáliter, étiam spiritáliter volébat intélligi. Neque enim tantum mirácula propter mirácula faciébat; sed ut illa, quæ faciébat, mira essent vidéntibus, vera essent intelligéntibus.\par
\switchcolumn
\EN
\noindent We find written how the Lord raised from the dead three persons visibly, but thousands invisibly. But how many they may have been whom He raised visibly, who knoweth For all the things which He did are not written. John saith thus: “There are also many other things which Jesus did, the which, if they should be written every one, I suppose that even the world itself could not contain the books that should be written.” (xxi. 25.) There were then, doubtless, many more raised to life, but it is not meaningless that three are recorded. For our Lord Jesus Christ hath willed that those things which He did carnally, we should understand also spiritually. He worked not miracles only for the sake of working wonders, but that His works might be at once wonderful to them that beheld, and true to them that understand them.\par
\switchcolumn*

\LA
\begin{Resp}\Rsign Duo Séraphim clamábant alter ad álterum: \Star{} Sanctus, sanctus, sanctus Dóminus Deus Sábaoth: * Plena est omnis terra glória ejus.\end{Resp}
\begin{Resp}\Vsign Tres sunt qui testimónium dant in cælo: Pater, Verbum, et Spíritus Sanctus: et hi tres unum sunt.\end{Resp}
\begin{Resp}\Rsign Sanctus, sanctus, sanctus Dóminus Deus Sábaoth.\end{Resp}
\begin{Resp}\Vsign Glória Patri, et Fílio, \Star{} et Spirítui Sancto.\end{Resp}
\begin{Resp}\Rsign Plena est omnis terra glória ejus.\end{Resp}
\switchcolumn
\EN
\begin{Resp}\Rsign One Seraph cried unto another: \Star{} Holy, Holy, Holy is the Lord God of hosts; the whole earth is full of His glory.\end{Resp}
\begin{Resp}\Vsign There are Three That bear record in heaven: the Father, the Word, and the Holy Ghost, and these Three are One.\end{Resp}
\begin{Resp}\Rsign Holy, Holy, Holy is the Lord God of hosts.\end{Resp}
\begin{Resp}\Vsign Glory be to the Father, and to the Son, \Star{} and to the Holy Ghost.\end{Resp}
\begin{Resp}\Rsign The whole earth is full of His glory.\end{Resp}
\switchcolumn*

\LA
\LessonHead{Lectio IX}
\switchcolumn
\EN
\LessonHead{Lesson IX}
\switchcolumn*

\LA
\noindent Quemádmodum qui videt lítteras in códice óptime scripto, et non novit légere, laudat quidem antiquárii manum, admírans ápicum pulchritúdinem; sed quid sibi velint, quid índicent illi ápices, nescit, et est óculis laudátor, mente non cógnitor. Alius autem et laudat artifícium, et capit intelléctum: ille útique, qui non solum vidére quod commúne est ómnibus, potest, sed étiam légere; quod qui non dídicit, non potest. Ita qui vidérunt Christi mirácula, et non intellexérunt quid sibi vellent, et quid intelligéntibus quodámmodo innúerent, miráti sunt tantum quia facta sunt; álii vero et facta miráti, et intellécta assecúti. Tales nos in schola Christi esse debémus.\par
\switchcolumn
\EN
\noindent Even as one that looketh upon a scroll right fairly written, and knoweth not how to read therein, praiseth the hand of the old scribe when he seeth the beauty of the points, but what it saith, what those points mean, he knoweth not, and praiseth by the eye, without understanding by the mind, and as, on the other hand, he that can not only gaze on it, as can all men, but also can read it, praiseth the penmanship, and catcheth the sense likewise, which the unlearned cannot do even so, there were some that saw the miracles which Christ did, and understood not what they meant, nor what they, as it were, hinted to such as did understand them, and these only marvelled to see them wrought. And other some there were which saw the works, and marvelled, and understood them, and profited by them. And it is as these last that we ought to be in the school of Christ.\par
\switchcolumn*

\LA
\TeDeum{Te Deum}
\switchcolumn
\EN
\TeDeum{Te Deum}
\switchcolumn*

\end{paracol}

\Office{Dominica XVI Post Pentecosten}{Sixteenth Sunday after Pentecost}{Semiduplex}
\addcontentsline{toc}{subsection}{Dominica XVI Post Pentecosten \textit{— Sixteenth Sunday after Pentecost}}
\begin{paracol}{2}
\LA
\RefLine{Lectiones I–VI: de Scriptura occurrente.}
\switchcolumn
\EN
\RefLine{Lessons I–VI: from the Scripture of the day.}
\switchcolumn*

\LA
\LessonHead{Lectio VII}
\switchcolumn
\EN
\LessonHead{Lesson VII}
\switchcolumn*

\LA
\LessonTitle{Léctio sancti Evangélii secúndum Lucam}
\switchcolumn
\EN
\LessonTitle{From the Holy Gospel according to Luke}
\switchcolumn*

\LA
\Cite{Luc 14:1-11}
\switchcolumn
\EN
\Cite{Luke 14:1-11}
\switchcolumn*

\LA
\noindent In illo témpore: Cum intráret Jesus in domum cujúsdam príncipis pharisæórum sábbato manducáre panem, et ipsi observábant eum. Et ecce homo quidam hydrópicus erat ante illum. Et réliqua.\par
\switchcolumn
\EN
\noindent At that time: As Jesus went into the house of one of the chief Pharisees, to eat bread on the Sabbath-day, they watched Him. And, behold, there was a certain man before Him, which had the dropsy. And so on.\par
\switchcolumn*

\LA
\LessonTitle{Homilía sancti Ambrósii Epíscopi}
\switchcolumn
\EN
\LessonTitle{Homily by St. Ambrose, Bishop (of Milan.)}
\switchcolumn*

\LA
\Source{Liber 7 in Lucæ cap. 14}
\switchcolumn
\EN
\Source{Bk. 7 on Luke XIV}
\switchcolumn*

\LA
\noindent Curátur hydrópicus, in quo fluxus carnis exúberans ánimæ gravábat offícia, spíritus exstinguébat ardórem. Deínde docétur humílitas, dum in illo convívio nuptiáli appeténtia loci superióris arcétur; cleménter tamen, ut persuásio humanitátis aspertitátem coërcitiónis exclúderet, rátio profíceret ad persuasiónis efféctum, et corréctio emendáret afféctum. Huic quasi próximo límine humánitas copulátur: quæ ita Domínicæ senténtiæ definitióne distínguitur, si in páuperes et débiles conferátur: nam hospitálem esse remuneratúris, afféctus avarítiæ est.\par
\switchcolumn
\EN
\noindent Now is healed this man sick of the dropsy, in whom too much watery matter had well-nigh drowned the functions of life, and quenched the fire of understanding. Anon, a lesson is given in lowly-mindedness, when it is forbidden to the guests at a marriage feast to go and sit down unasked in the highest room, albeit the Lord spake gently, that the teaching of courtesy might forestall a harsh rebuke, reason prevail by dint of persuasion, and the desires be bent to follow the instruction. And upon this, as next-door neighbour, cometh courtesy, which is so called by the Lord, when it is shown to the poor and weak, since to show it to them from whom we are to receive aught, is but a movement of self-interest.\par
\switchcolumn*

\LA
\begin{Resp}\Rsign Quis mihi tríbuat, ut in inférno prótegas me et abscóndas me, donec pertránseat furor tuus, Dómine, nisi tu, qui solus es Deus? \Star{} Et constítuas mihi tempus, in quo recordéris mei?\end{Resp}
\begin{Resp}\Vsign Numquid sicut dies hóminis dies tui, ut quæras iniquitátem meam; cum sit nemo, qui de manu tua possit erúere.\end{Resp}
\begin{Resp}\Rsign Et constítuas mihi tempus, in quo recordéris mei?\end{Resp}
\switchcolumn
\EN
\begin{Resp}\Rsign O that Thou wouldest hide me in the grave; that Thou wouldest keep me secret, until thy wrath be past, even thine, O Lord, Thou That alone art God; \Star{} That Thou wouldest appoint me a set time, and remember me!\end{Resp}
\begin{Resp}\Vsign Are thy days as the days of man, that Thou inquirest after mine iniquity, and there is none that can deliver out of thine hand.\end{Resp}
\begin{Resp}\Rsign That Thou wouldest appoint me a set time, and remember me!\end{Resp}
\switchcolumn*

\LA
\LessonHead{Lectio VIII}
\switchcolumn
\EN
\LessonHead{Lesson VIII}
\switchcolumn*

\LA
\noindent Postrémo jam quasi eméritæ milítiæ viro contemnendárum stipéndium præscríbitur facultátum: quod neque ille, qui stúdiis inténtus inferióribus possessiónes sibi terrénas coémit, regnum cæli possit adipísci, cum Dóminus dicat: Vende ómnia tua, et séquere me: nec ille, qui emit boves, cum Eliséus occíderit, et pópulo divíserit quos habébat: et ille qui duxit uxórem, cógitet quæ mundi sunt, non quæ Dei. Non quo conjúgium reprehendátur, sed quia ad majórem honórem vocétur intégritas: quóniam múlier innúpta et vídua cógitat quæ sunt Dómini, ut sit sancta córpore et spíritu.\par
\switchcolumn
\EN
\noindent Lastly, as to a soldier that hath served his full time, is apportioned a reward for esteeming lightly of riches so he only can inherit the kingdom of God, whose soul is not given to seek after lower ends, and who purchaseth not to himself earthly possessions whereas the Lord saith: “Sell that thou hast, and follow Me.” (Matth. xix. 21.) Neither can he gain it that buyeth oxen, which beasts Elisha slew and gave unto the people. (3 Kings xix. 21.) Neither can he win it which hath married a wife and therefore cannot come, for “he that is unmarried careth for the things that belong to the Lord, how he may please the Lord but he that is married careth for the things that are of the world, how he may please his wife.” (1 Cor. vii. 32, 33.) Not that this is to be taken for blame of marriage, but only that virginity is the more honourable way, since “the unmarried woman” and the widow “careth for the things of the Lord, that she may be holy both in body and in spirit.” (34-)\par
\switchcolumn*

\LA
\begin{Resp}\Rsign Duo Séraphim clamábant alter ad álterum: \Star{} Sanctus, sanctus, sanctus Dóminus Deus Sábaoth: * Plena est omnis terra glória ejus.\end{Resp}
\begin{Resp}\Vsign Tres sunt qui testimónium dant in cælo: Pater, Verbum, et Spíritus Sanctus: et hi tres unum sunt.\end{Resp}
\begin{Resp}\Rsign Sanctus, sanctus, sanctus Dóminus Deus Sábaoth.\end{Resp}
\begin{Resp}\Vsign Glória Patri, et Fílio, \Star{} et Spirítui Sancto.\end{Resp}
\begin{Resp}\Rsign Plena est omnis terra glória ejus.\end{Resp}
\switchcolumn
\EN
\begin{Resp}\Rsign One Seraph cried unto another: \Star{} Holy, Holy, Holy is the Lord God of hosts; the whole earth is full of His glory.\end{Resp}
\begin{Resp}\Vsign There are Three That bear record in heaven: the Father, the Word, and the Holy Ghost, and these Three are One.\end{Resp}
\begin{Resp}\Rsign Holy, Holy, Holy is the Lord God of hosts.\end{Resp}
\begin{Resp}\Vsign Glory be to the Father, and to the Son, \Star{} and to the Holy Ghost.\end{Resp}
\begin{Resp}\Rsign The whole earth is full of His glory.\end{Resp}
\switchcolumn*

\LA
\LessonHead{Lectio IX}
\switchcolumn
\EN
\LessonHead{Lesson IX}
\switchcolumn*

\LA
\noindent Sed ut in grátiam, ut supra cum víduis, ita nunc étiam cum conjúgibus revertámur: non refúgimus opiniónem, quam sequúntur pleríque, ut tria génera hóminum a consórtio magnæ illíus cœnæ æstimémus exclúdi: Gentílium, Judæórum, Hæreticórum. Et ídeo Apóstolus avarítiam dicit esse fugiéndam; ne impedíti more Gentíli, iniquitáte, malítia, impudicítia, avarítia, ad regnum Christi perveníre nequeámus. Omnis enim immúndus aut avárus, quod est idolórum sérvitus, non habet hereditátem in regno Christi et Dei.\par
\switchcolumn
\EN
\noindent But in all fairness, having thus spoken concerning widows, let us betake ourselves again among the married, and join with them in entertaining the opinion which is held by so many, that there are only three classes of men who are shut out from the great supper named in the gospel, which three classes are Heathens, Jews, and Heretics. And therefore it is that the Apostle warneth us that we “walk not as other Gentiles walk,” in malice and bitterness, and uncleanness, and covetousness, and so have no entry into the kingdom of Christ, since “no unclean person, nor covetous man, who is an idolater, hath any inheritance in the kingdom of Christ and of God.” (Eph. iv. 17, v. 5.)\par
\switchcolumn*

\LA
\TeDeum{Te Deum}
\switchcolumn
\EN
\TeDeum{Te Deum}
\switchcolumn*

\end{paracol}

\Office{Dominica XVII Post Pentecosten}{Seventeenth Sunday after Pentecost}{Semiduplex}
\addcontentsline{toc}{subsection}{Dominica XVII Post Pentecosten \textit{— Seventeenth Sunday after Pentecost}}
\begin{paracol}{2}
\LA
\RefLine{Lectiones I–VI: de Scriptura occurrente.}
\switchcolumn
\EN
\RefLine{Lessons I–VI: from the Scripture of the day.}
\switchcolumn*

\LA
\LessonHead{Lectio VII}
\switchcolumn
\EN
\LessonHead{Lesson VII}
\switchcolumn*

\LA
\LessonTitle{Léctio sancti Evangélii secúndum Matthǽum}
\switchcolumn
\EN
\LessonTitle{From the Holy Gospel according to Matthew}
\switchcolumn*

\LA
\Cite{Matt 22:34-46}
\switchcolumn
\EN
\Cite{Matt 22:34-46}
\switchcolumn*

\LA
\noindent In illo témpore: Accessérunt ad Jesum pharisǽi, et interrogávit eum unus ex eis legis doctor tentans eum: Magíster, quod est mandátum magnum in lege? Et réliqua.\par
\switchcolumn
\EN
\noindent At that time, the Pharisees came unto Jesus, and one of them, which was a lawyer, asked Him a question, tempting Him, saying: Master, which is the great commandment in the Law? And so on.\par
\switchcolumn*

\LA
\LessonTitle{Homilía sancti Joánnis Chrysóstomi}
\switchcolumn
\EN
\LessonTitle{Homily by St. John Chrysostom, Patriarch (of Constantinople.)}
\switchcolumn*

\LA
\Source{Homilia 72 in Matthæum}
\switchcolumn
\EN
\Source{72nd on Matthew.}
\switchcolumn*

\LA
\noindent Sadducǽis confúsis, pharisǽi rursus aggrediúntur; cumque quiéscere oportéret, decertáre voluérunt: et legis perítiam profiténtem præmíttunt, non díscere, sed tentáre cupiéntes; ac ita intérrogant: Quodnam primum mandátum in lege sit. Nam cum primum illud sit, Díliges Dóminum Deum tuum: putántes causas sibi allatúrum ad mandátum hoc corrigéndum, áliquid addéndo, quóniam Deum se faciébat, hoc modo intérrogant. Quid ígitur Christus? Ut osténdat idcírco ad hæc eos devenísse, quia nulla in eis esset cáritas, sed invídiæ livóre tabéscerent: Díliges, inquit, Dóminum Deum tuum: hoc primum et magnum mandátum est. Secúndum autem símile huic: Díliges próximum tuum sicut teípsum.\par
\switchcolumn
\EN
\noindent When the Pharisees had heard that Christ had put the Sadducees to silence, they gathered themselves together for a fresh attack just when it behoved them to be quiet, they willed to contend and so they put forward one of themselves who professed skill in the law, not wishing to learn, but to lay a snare. This person therefore proposed the question: “Which is the great commandment in the law?” The first and great commandment is: “Thou shalt love the Lord thy God,” but they expected that He would make some exception or addition to this in His Own case, since He made Himself God. (John x. 33.) With this expectation they asked Him the question, but what said Christ? To show that they had adopted this course, because they were loveless, and sick with envy, He answered: “Thou shalt love the Lord thy God with all thy heart, and with all thy soul, and with all thy mind. This is the first and great commandment. And the second is like unto it: “Thou shalt love thy neighbour as thyself.”\par
\switchcolumn*

\LA
\begin{Resp}\Rsign Refúlsit sol in clípeos áureos, et resplenduérunt montes ab eis: \Star{} Et fortitúdo géntium dissipáta est.\end{Resp}
\begin{Resp}\Vsign Erat enim exércitus magnus valde et fortis: et appropiávit Judas et exércitus ejus in prǽlio.\end{Resp}
\begin{Resp}\Rsign Et fortitúdo géntium dissipáta est.\end{Resp}
\switchcolumn
\EN
\begin{Resp}\Rsign The sun shone upon the shields of gold, and the mountains glistered therewith; \Star{} And the army of the heathens was spread abroad.\end{Resp}
\begin{Resp}\Vsign For the army was very great and mighty then Judas and his host drew near and entered into battle.\end{Resp}
\begin{Resp}\Rsign And the army of the heathens was spread abroad.\end{Resp}
\switchcolumn*

\LA
\LessonHead{Lectio VIII}
\switchcolumn
\EN
\LessonHead{Lesson VIII}
\switchcolumn*

\LA
\noindent Quam ob rem símile est huic? Quóniam hoc illud indúcit, et ab illo rursus munítur. Quicúmque enim male agit, ódio habet lucem, et non venit ad lucem. Et rursus: Dixit insípiens in corde suo, Non est Deus. Deínde séquitur: Corrúpti sunt, et abominábiles facti sunt in stúdiis suis. Et íterum: Radix ómnium malórum avarítia est; quam quidam appeténtes, erravérunt a fide. Et, Qui díligit me, mandáta mea servábit: quorum caput et radix est: Díliges Dóminum Deum tuum, et próximum tuum sicut teípsum.\par
\switchcolumn
\EN
\noindent Why is this second commandment like unto the first Because the first is the second's source and sanction. “For every one that doeth evil hateth the light, neither cometh to the light.” (John iii. 20.) And again: “The fool hath said in his heart There is no God” and there followeth: “They are corrupt, and become abominable in their works.” (Ps. xiii. 1.) And yet again: “The love of money is the root of all evil which while some coveted after, they have erred from the faith.” (1 Tim. vi. 10.) And yet once more: “If ye love Me, keep My commandments,” (John xiv. 15,) of which commandments the head and root is “Thou shalt love the Lord thy God and thy neighbour as thyself.”\par
\switchcolumn*

\LA
\begin{Resp}\Rsign Duo Séraphim clamábant alter ad álterum: \Star{} Sanctus, sanctus, sanctus Dóminus Deus Sábaoth: * Plena est omnis terra glória ejus.\end{Resp}
\begin{Resp}\Vsign Tres sunt qui testimónium dant in cælo: Pater, Verbum, et Spíritus Sanctus: et hi tres unum sunt.\end{Resp}
\begin{Resp}\Rsign Sanctus, sanctus, sanctus Dóminus Deus Sábaoth.\end{Resp}
\begin{Resp}\Vsign Glória Patri, et Fílio, \Star{} et Spirítui Sancto.\end{Resp}
\begin{Resp}\Rsign Plena est omnis terra glória ejus.\end{Resp}
\switchcolumn
\EN
\begin{Resp}\Rsign One Seraph cried unto another: \Star{} Holy, Holy, Holy is the Lord God of hosts; the whole earth is full of His glory.\end{Resp}
\begin{Resp}\Vsign There are Three That bear record in heaven: the Father, the Word, and the Holy Ghost, and these Three are One.\end{Resp}
\begin{Resp}\Rsign Holy, Holy, Holy is the Lord God of hosts.\end{Resp}
\begin{Resp}\Vsign Glory be to the Father, and to the Son, \Star{} and to the Holy Ghost.\end{Resp}
\begin{Resp}\Rsign The whole earth is full of His glory.\end{Resp}
\switchcolumn*

\LA
\LessonHead{Lectio IX}
\switchcolumn
\EN
\LessonHead{Lesson IX}
\switchcolumn*

\LA
\noindent Si ergo dilígere Deum, dilígere próximum est: (nam si díligis me, o Petre, inquit, pasce oves meas) si étiam diléctio próximi facit ut mandáta custódias: mérito ait in his totam legem et prophétas pendére. Et quemádmodum in superióribus, cum de resurrectióne interrogarétur, plus dócuit quam tentántes petébant; sic in hoc loco de primo interrogátus mandáto, secúndum étiam non valde quam primum inférius, sponte áttulit; secúndum enim est primo símile. Ita occúlte insinuávit, ódio illos ad quæréndum incitári. Cáritas enim, inquit, non æmulátur.\par
\switchcolumn
\EN
\noindent If therefore, to love God is to love our neighbour also, (as it appeared, where it is written: “Simon, son of Jonas, lovest thou Me? And he said unto Him: 'Lord, Thou knowest all things, Thou knowest that I love thee.' Jesus saith unto him: 'Feed My sheep',” John xxi. 17,) and if “love is the fulfilling of the law,” (Rom. xiii. 10,) justly doth the Lord say that “on these two commandments hang all the law and the Prophets.” And even as when, before this, (23-32,) being interrogated about the Resurrection, He answered them more than they asked, so, now, being interrogated concerning the first and great commandment, He answereth them, of His own accord, touching that second one also, which is little lower than the first, for “the second is like unto it.” Herein He would have them understand that it was hatred stirred them up to question Him. For “Charity,” saith the Apostle, “envieth not.” (1 Cor. xiii. 4.)\par
\switchcolumn*

\LA
\TeDeum{Te Deum}
\switchcolumn
\EN
\TeDeum{Te Deum}
\switchcolumn*

\end{paracol}

\Office{Dominica XVIII Post Pentecosten}{Eighteenth Sunday after Pentecost}{Semiduplex}
\addcontentsline{toc}{subsection}{Dominica XVIII Post Pentecosten \textit{— Eighteenth Sunday after Pentecost}}
\begin{paracol}{2}
\LA
\RefLine{Lectiones I–VI: de Scriptura occurrente.}
\switchcolumn
\EN
\RefLine{Lessons I–VI: from the Scripture of the day.}
\switchcolumn*

\LA
\LessonHead{Lectio VII}
\switchcolumn
\EN
\LessonHead{Lesson VII}
\switchcolumn*

\LA
\LessonTitle{Léctio sancti Evangélii secúndum Matthǽum}
\switchcolumn
\EN
\LessonTitle{From the Holy Gospel according to Matthew}
\switchcolumn*

\LA
\Cite{Matt 9:1-8}
\switchcolumn
\EN
\Cite{Matt 9:1-8}
\switchcolumn*

\LA
\noindent In illo témpore: Ascéndens Jesus in navículam transfretávit et venit in civitátem suam. Et réliqua.\par
\switchcolumn
\EN
\noindent At that time: Jesus entered into a ship, and passed over, and came into His own city. And so on.\par
\switchcolumn*

\LA
\LessonTitle{Homilía sancti Petri Chrysólogi}
\switchcolumn
\EN
\LessonTitle{Homily by St. Peter Chrysologus, Archbishop (of Ravenna.)}
\switchcolumn*

\LA
\Source{Sermo 50}
\switchcolumn
\EN
\Source{Sermon 50}
\switchcolumn*

\LA
\noindent Christum in humánis áctibus divína gessísse mystéria, et in rebus visibílibus invisibília exercuísse negótia, léctio hodiérna monstrávit. Ascéndit, inquit, in navículam, et transfretávit, et venit in civitátem suam. Nonne ipse est, qui, fugátis flúctibus, maris profúnda nudávit, ut Israëlíticus pópulus inter stupéntes undas sicco vestígio velut móntium cóncava pertransíret? Nonne hic est, qui Petri pédibus marínos vórtices inclinávit, ut iter líquidum humánis gréssibus sólidum præbéret obséquium.\par
\switchcolumn
\EN
\noindent This day's reading hath shown us an instance of how Christ, in those things which He did as Man, worked deep works of God, and by things which were seen wrought things which were not seen. The Evangelist saith Jesus “entered into a ship, and passed over, and came into His Own city.” Was not This He Who had once parted the waves hither and thither, and made the dry ground appear at the bottom of the sea, so that His people Israel passed dry-shod between masses of water standing still, as through an hollow glen in a mountain? Was not This He Who made the depths of the sea solid under the feet of Peter, so that the watery path offered a firm way for human footsteps?\par
\switchcolumn*

\LA
\begin{Resp}\Rsign Refúlsit sol in clípeos áureos, et resplenduérunt montes ab eis: \Star{} Et fortitúdo géntium dissipáta est.\end{Resp}
\begin{Resp}\Vsign Erat enim exércitus magnus valde et fortis: et appropiávit Judas et exércitus ejus in prǽlio.\end{Resp}
\begin{Resp}\Rsign Et fortitúdo géntium dissipáta est.\end{Resp}
\switchcolumn
\EN
\begin{Resp}\Rsign The sun shone upon the shields of gold, and the mountains glistered therewith; \Star{} And the army of the heathens was spread abroad.\end{Resp}
\begin{Resp}\Vsign For the army was very great and mighty then Judas and his host drew near and entered into battle.\end{Resp}
\begin{Resp}\Rsign And the army of the heathens was spread abroad.\end{Resp}
\switchcolumn*

\LA
\LessonHead{Lectio VIII}
\switchcolumn
\EN
\LessonHead{Lesson VIII}
\switchcolumn*

\LA
\noindent Et quid est, quod ipse sibi sic maris dénegat servitútem, ut brevíssimi lacus tránsitum sub mercéde náutica transfretáret? Ascéndit, inquit, in navículam, et transfretávit. Et quid mirum, fratres? Christus venit suscípere infirmitátes nostras, et suas nobis conférre virtútes; humána quǽrere, præstáre divína; accípere injúrias, réddere dignitátes; ferre tǽdia, reférre sanitátes: quia médicus, qui non fert infirmitátes, curáre nescit; et qui non fúerit cum infírmo infirmátus, infírmo non potest conférre sanitátem.\par
\switchcolumn
\EN
\noindent Therefore then denied He unto Himself a like service from the sea, but crossed over that narrow lake at the cost of a voyage on shipboard “He entered into a ship, and passed over.” What wonder, my brethren: Christ came to take our weakness upon Him, that He might n make us partakers of His strength to seek the things of men, that He might give to men the things of God to receive insults, that He might bestow honours to bear weariness, that He might grant rest for the physician that is himself beset by no frailties, knoweth not how to treat the frailties of others, nor he that is not weak with the weak, how to make the weak strong.\par
\switchcolumn*

\LA
\begin{Resp}\Rsign Duo Séraphim clamábant alter ad álterum: \Star{} Sanctus, sanctus, sanctus Dóminus Deus Sábaoth: * Plena est omnis terra glória ejus.\end{Resp}
\begin{Resp}\Vsign Tres sunt qui testimónium dant in cælo: Pater, Verbum, et Spíritus Sanctus: et hi tres unum sunt.\end{Resp}
\begin{Resp}\Rsign Sanctus, sanctus, sanctus Dóminus Deus Sábaoth.\end{Resp}
\begin{Resp}\Vsign Glória Patri, et Fílio, \Star{} et Spirítui Sancto.\end{Resp}
\begin{Resp}\Rsign Plena est omnis terra glória ejus.\end{Resp}
\switchcolumn
\EN
\begin{Resp}\Rsign One Seraph cried unto another: \Star{} Holy, Holy, Holy is the Lord God of hosts; the whole earth is full of His glory.\end{Resp}
\begin{Resp}\Vsign There are Three That bear record in heaven: the Father, the Word, and the Holy Ghost, and these Three are One.\end{Resp}
\begin{Resp}\Rsign Holy, Holy, Holy is the Lord God of hosts.\end{Resp}
\begin{Resp}\Vsign Glory be to the Father, and to the Son, \Star{} and to the Holy Ghost.\end{Resp}
\begin{Resp}\Rsign The whole earth is full of His glory.\end{Resp}
\switchcolumn*

\LA
\LessonHead{Lectio IX}
\switchcolumn
\EN
\LessonHead{Lesson IX}
\switchcolumn*

\LA
\noindent Christus ergo, si in suis mansísset virtútibus, commúne cum homínibus nil habéret; et nisi implésset carnis órdinem, carnis in illo esset otiósa suscéptio. Ascéndit, inquit, in navículam, et transfretávit, et venit in civitátem suam. Creátor rerum orbis Dóminus, posteáquam se propter nos nostra angustávit in carne, cœpit habére humánam pátriam, cœpit civitátis Judáicæ esse civis, paréntes habére cœpit paréntum ómnium ipse parens; ut invitáret amor, attráheret cáritas, vincíret afféctio, suadéret humánitas, quos fúgerat dominátio, metus dispérserat, fécerat vis potestátis extórres.\par
\switchcolumn
\EN
\noindent Therefore, if Christ had abode still in His strength, He had in no wise been a fellow of men if in Him Flesh had not run the way of flesh, then had it been idle for Him to have taken Flesh at all. “He entered into a ship, and passed over, and came into His Own city.” The Lord, the Maker of the world, and of all things that are therein, having been pleased for our sakes to prison Himself in our flesh, began to have an human home, and to be a citizen of a Jewish city Himself the Father of all, to have parents and all, that His love might invite, His charity draw, His tenderness bind, His gentleness persuade them whom His Kingship had scared, His awfulness scattered, and His power terrified out of His dominion.\par
\switchcolumn*

\LA
\TeDeum{Te Deum}
\switchcolumn
\EN
\TeDeum{Te Deum}
\switchcolumn*

\end{paracol}

\Office{Dominica XIX Post Pentecosten}{Nineteenth Sunday after Pentecost}{Semiduplex}
\addcontentsline{toc}{subsection}{Dominica XIX Post Pentecosten \textit{— Nineteenth Sunday after Pentecost}}
\begin{paracol}{2}
\LA
\RefLine{Lectiones I–VI: de Scriptura occurrente.}
\switchcolumn
\EN
\RefLine{Lessons I–VI: from the Scripture of the day.}
\switchcolumn*

\LA
\LessonHead{Lectio VII}
\switchcolumn
\EN
\LessonHead{Lesson VII}
\switchcolumn*

\LA
\LessonTitle{Léctio sancti Evangélii secúndum Matthǽum}
\switchcolumn
\EN
\LessonTitle{From the Holy Gospel according to Matthew}
\switchcolumn*

\LA
\Cite{Matt 22:1-14}
\switchcolumn
\EN
\Cite{Matt 22:1-14}
\switchcolumn*

\LA
\noindent In illo témpore: Loquebátur Jesus princípibus sacerdótum et pharisǽis in parábolis, dicens: Símile factum est regnum cælórum hómini regi, qui fecit núptias fílio suo. Et réliqua.\par
\switchcolumn
\EN
\noindent At that time, Jesus spoke by parables unto the chief priests and Pharisees, and said: The kingdom of heaven is like unto a certain king, which made a marriage for his son. And so on.\par
\switchcolumn*

\LA
\LessonTitle{Homilía sancti Gregórii Papæ}
\switchcolumn
\EN
\LessonTitle{Homily by Pope St. Gregory (the Great.)}
\switchcolumn*

\LA
\Source{Homilia 38 in Evangelia, post initium}
\switchcolumn
\EN
\Source{38th on the Gospels}
\switchcolumn*

\LA
\noindent Sæpe jam me dixísse mémini, quod plerúmque in sancto Evangélio regnum cælórum præsens Ecclésia nominátur: congregátio quippe justórum, regnum cælórum dícitur. Quia enim per prophétam Dóminus dicit: Cælum mihi sedes est; et Sálomon ait: Anima justi sedes sapiéntiæ; Paulus étiam dicit Christum Dei virtútem, et Dei sapiéntiam: líquido collígere debémus, quia si Deus sapiéntia, ánima autem justi, sedes sapiéntiæ, dum cælum dícitur sedes Dei, cælum ergo est ánima justi. Hinc per Psalmístam de sanctis prædicatóribus dícitur: Cæli enárrant glóriam Dei.\par
\switchcolumn
\EN
\noindent I remember that I have often said that, in the Holy Gospel, the Church as she now is, is called the kingdom of heaven, for the kingdom of heaven is indeed the assembly of the righteous. The Lord hath said by the mouth of His Prophet: The heaven is My throne. (Isa. lxvi. 1.) Solomon saith: The throne of wisdom is the soul of the righteous. And Paul saith that Christ is the power of God and the wisdom of God. (1 Cor. i. 24.) From these passages we may clearly gather that if wisdom be God, and wisdom's throne be the soul of the righteous, and God's throne be the heaven, then the soul of the righteous is heaven. Hence also the Psalmist saith, speaking of holy preachers: The heavens declare the glory of God. (xviii. 2.)\par
\switchcolumn*

\LA
\begin{Resp}\Rsign Refúlsit sol in clípeos áureos, et resplenduérunt montes ab eis: \Star{} Et fortitúdo géntium dissipáta est.\end{Resp}
\begin{Resp}\Vsign Erat enim exércitus magnus valde et fortis: et appropiávit Judas et exércitus ejus in prǽlio.\end{Resp}
\begin{Resp}\Rsign Et fortitúdo géntium dissipáta est.\end{Resp}
\switchcolumn
\EN
\begin{Resp}\Rsign The sun shone upon the shields of gold, and the mountains glistered therewith; \Star{} And the army of the heathens was spread abroad.\end{Resp}
\begin{Resp}\Vsign For the army was very great and mighty then Judas and his host drew near and entered into battle.\end{Resp}
\begin{Resp}\Rsign And the army of the heathens was spread abroad.\end{Resp}
\switchcolumn*

\LA
\LessonHead{Lectio VIII}
\switchcolumn
\EN
\LessonHead{Lesson VIII}
\switchcolumn*

\LA
\noindent Regnum ergo cælórum est Ecclésia justórum: quia dum eórum corda in terra nil ámbiunt, per hoc quod ad supérna suspírant, jam in eis Dóminus quasi in cæléstibus regnat. Dicátur ergo: Símile est regnum cælórum hómini regi, qui fecit núptias fílio suo. Jam intélligit cáritas vestra, quis est iste Rex, Regis fílii pater: ille nimírum, cui Psalmísta ait: Deus judícium tuum Regi da, et justítiam tuam fílio Regis. Qui fecit núptias fílio suo. Tunc enim Deus Pater Deo Fílio suo núptias fecit, quando hunc in útero Vírginis humánæ natúræ conjúnxit, quando Deum ante sǽcula fíeri vóluit hóminem in fine sæculórum.\par
\switchcolumn
\EN
\noindent The kingdom of heaven, therefore, is the Church of the righteous, even of them whose hearts seek not for anything upon earth, but who sigh so continually after the things which are above, that God doth already reign in them as He doth in heaven. Let it then be said, The kingdom of heaven is like unto a certain king, which made a marriage for his son. Ye already understand, my loving friends, who is that Royal Father of a Royal Son. It is indeed no other than He to Whom the Psalmist saith: Give the King thy judgments, O God, and thy righteousness unto the King's son. (lxxi. 2.) Which made a marriage for his son. God the Father made a marriage for God the Son, when He wedded Him to the manhood in the womb of the Virgin, when He willed that He Who is God before all ages, should in the end of the ages become Man.\par
\switchcolumn*

\LA
\begin{Resp}\Rsign Duo Séraphim clamábant alter ad álterum: \Star{} Sanctus, sanctus, sanctus Dóminus Deus Sábaoth: * Plena est omnis terra glória ejus.\end{Resp}
\begin{Resp}\Vsign Tres sunt qui testimónium dant in cælo: Pater, Verbum, et Spíritus Sanctus: et hi tres unum sunt.\end{Resp}
\begin{Resp}\Rsign Sanctus, sanctus, sanctus Dóminus Deus Sábaoth.\end{Resp}
\begin{Resp}\Vsign Glória Patri, et Fílio, \Star{} et Spirítui Sancto.\end{Resp}
\begin{Resp}\Rsign Plena est omnis terra glória ejus.\end{Resp}
\switchcolumn
\EN
\begin{Resp}\Rsign One Seraph cried unto another: \Star{} Holy, Holy, Holy is the Lord God of hosts; the whole earth is full of His glory.\end{Resp}
\begin{Resp}\Vsign There are Three That bear record in heaven: the Father, the Word, and the Holy Ghost, and these Three are One.\end{Resp}
\begin{Resp}\Rsign Holy, Holy, Holy is the Lord God of hosts.\end{Resp}
\begin{Resp}\Vsign Glory be to the Father, and to the Son, \Star{} and to the Holy Ghost.\end{Resp}
\begin{Resp}\Rsign The whole earth is full of His glory.\end{Resp}
\switchcolumn*

\LA
\LessonHead{Lectio IX}
\switchcolumn
\EN
\LessonHead{Lesson IX}
\switchcolumn*

\LA
\noindent Sed quia ex duábus persónis fíeri solet ista nuptiális conjúnctio; absit hoc ab intelléctibus nostris, ut persónam Dei et hóminis Redemptóris nostri Jesu Christi, ex duábus persónis credámus unítam. Ex duábus quippe atque in duábus hunc natúris exsístere dícimus; sed ex duábus persónis compósitum credi, ut nefas vitámus. Apértius ergo atque secúrius dici potest, quia in hoc Pater Regi Fílio núptias fecit, quo ei per incarnatiónis mystérium sanctam Ecclésiam sociávit. Uterus autem Genitrícis Vírginis, hujus sponsi thálamus fuit. Unde et Psalmísta dicit: In sole pósuit tabernáculum suum, et ipse tamquam sponsus procédens de thálamo suo.\par
\switchcolumn
\EN
\noindent The marriage union is the union of two persons, but God forbid that we should imagine that the One Person of our Redeemer Jesus Christ, Who is both God and Man, is formed by a union of an human person with a Divine Person. We profess concerning Him that He is of, and in two natures, but we shrink from the blasphemy of saying that He is compounded of two persons. It will therefore be clearer and safer to say that the marriage which the Father made for His Royal Son was the wedding Him, through the mystery of the Incarnation, to His mystic Bride the Holy Church. The womb of the Maiden Mother was the marriage chamber in which this union took place. Hence it is that the Psalmist saith: In the sun hath He set His tabernacle, Who is as a bridegroom coming out of his chamber. (xviii. 6.2)\par
\switchcolumn*

\LA
\TeDeum{Te Deum}
\switchcolumn
\EN
\TeDeum{Te Deum}
\switchcolumn*

\end{paracol}

\Office{Dominica XX Post Pentecosten}{Twentieth Sunday after Pentecost}{Semiduplex}
\addcontentsline{toc}{subsection}{Dominica XX Post Pentecosten \textit{— Twentieth Sunday after Pentecost}}
\begin{paracol}{2}
\LA
\RefLine{Lectiones I–VI: de Scriptura occurrente.}
\switchcolumn
\EN
\RefLine{Lessons I–VI: from the Scripture of the day.}
\switchcolumn*

\LA
\LessonHead{Lectio VII}
\switchcolumn
\EN
\LessonHead{Lesson VII}
\switchcolumn*

\LA
\LessonTitle{Léctio sancti Evangélii secúndum Joánnem}
\switchcolumn
\EN
\LessonTitle{From the Holy Gospel according to John}
\switchcolumn*

\LA
\Cite{Joannes 4:46-53}
\switchcolumn
\EN
\Cite{John 4:46-53}
\switchcolumn*

\LA
\noindent In illo témpore: Erat quidam régulus, cujus fílius infirmabátur Caphárnaum. Et réliqua.\par
\switchcolumn
\EN
\noindent At that time: There was a certain nobleman, whose son was sick at Capernaum. And so on.\par
\switchcolumn*

\LA
\LessonTitle{Homilía sancti Gregórii Papæ}
\switchcolumn
\EN
\LessonTitle{Homily by Pope St. Gregory (the Great.)}
\switchcolumn*

\LA
\Source{Homilia 28 in Evangelia}
\switchcolumn
\EN
\Source{28th Homily on the Gospels}
\switchcolumn*

\LA
\noindent Léctio sancti Evangélii, quam modo, fratres, audístis, expositióne non índiget: sed ne hanc táciti præteriísse videámur, exhortándo pótius quam exponéndo in ea áliquid loquámur. Hoc autem nobis solúmmodo de expositióne vídeo esse requiréndum, cur is, qui ad salútem fílio peténdam vénerat, audívit: Nisi signa et prodígia vidéritis, non créditis. Qui enim salútem fílio quærébat, proculdúbio credébat; neque enim ab eo quǽreret salútem, quem non créderet Salvatórem. Quare ergo dícitur: Nisi signa et prodígia vidéritis, non créditis: qui ante crédidit, quam signa vidéret?\par
\switchcolumn
\EN
\noindent My brethren, the passage from the Holy Gospel, which ye have just now heard, standeth in need of no explanation. But lest I should seem to pass the same by in idle silence, I will say somewhat thereon, but that rather by way of exhortation than of explanation. Indeed, there seemeth to me to be but one point which calleth for explanation, and that point is this Wherefore was it that when the nobleman went unto the Lord, and besought Him that He would come down and heal his son, Jesus said unto him: “Except ye see signs and wonders, ye will not believe.” The very fact that he had come to beseech Christ to heal his son, putteth it beyond all doubt that this nobleman believed if he had not believed Him to be a Saviour, he would not have asked Him to save his son. Wherefore then said Jesus unto him: “Except ye see signs and wonders, ye will not believe,” since he was one who had not seen, and yet had believed?\par
\switchcolumn*

\LA
\begin{Resp}\Rsign Refúlsit sol in clípeos áureos, et resplenduérunt montes ab eis: \Star{} Et fortitúdo géntium dissipáta est.\end{Resp}
\begin{Resp}\Vsign Erat enim exércitus magnus valde et fortis: et appropiávit Judas et exércitus ejus in prǽlio.\end{Resp}
\begin{Resp}\Rsign Et fortitúdo géntium dissipáta est.\end{Resp}
\switchcolumn
\EN
\begin{Resp}\Rsign The sun shone upon the shields of gold, and the mountains glistered therewith; \Star{} And the army of the heathens was spread abroad.\end{Resp}
\begin{Resp}\Vsign For the army was very great and mighty then Judas and his host drew near and entered into battle.\end{Resp}
\begin{Resp}\Rsign And the army of the heathens was spread abroad.\end{Resp}
\switchcolumn*

\LA
\LessonHead{Lectio VIII}
\switchcolumn
\EN
\LessonHead{Lesson VIII}
\switchcolumn*

\LA
\noindent Sed mementóte quid pétiit; et apérte cognoscétis, quia in fide dubitávit. Popóscit namque, ut descénderet et sanáret fílium ejus. Corporálem ergo præséntiam Dómini quærébat, qui per spíritum nusquam déerat. Minus ítaque in illum crédidit, quem non putávit posse salútem dare, nisi præsens esset et córpore. Si enim perfécte credidísset, proculdúbio sciret, quia non esset locus ubi non esset Deus.\par
\switchcolumn
\EN
\noindent But bethink you what was his prayer, and then shall ye understand clearly wherein his faith was shaky. He “besought Him that He would come down and heal his son.” He asked for the bodily presence of Him Who is spiritually always present everywhere. He believed not therefore enough in Christ, for he thought that He could not heal unless He were bodily present. Had his faith been perfect, he would doubtless have known that God is everywhere.\par
\switchcolumn*

\LA
\begin{Resp}\Rsign Duo Séraphim clamábant alter ad álterum: \Star{} Sanctus, sanctus, sanctus Dóminus Deus Sábaoth: * Plena est omnis terra glória ejus.\end{Resp}
\begin{Resp}\Vsign Tres sunt qui testimónium dant in cælo: Pater, Verbum, et Spíritus Sanctus: et hi tres unum sunt.\end{Resp}
\begin{Resp}\Rsign Sanctus, sanctus, sanctus Dóminus Deus Sábaoth.\end{Resp}
\begin{Resp}\Vsign Glória Patri, et Fílio, \Star{} et Spirítui Sancto.\end{Resp}
\begin{Resp}\Rsign Plena est omnis terra glória ejus.\end{Resp}
\switchcolumn
\EN
\begin{Resp}\Rsign One Seraph cried unto another: \Star{} Holy, Holy, Holy is the Lord God of hosts; the whole earth is full of His glory.\end{Resp}
\begin{Resp}\Vsign There are Three That bear record in heaven: the Father, the Word, and the Holy Ghost, and these Three are One.\end{Resp}
\begin{Resp}\Rsign Holy, Holy, Holy is the Lord God of hosts.\end{Resp}
\begin{Resp}\Vsign Glory be to the Father, and to the Son, \Star{} and to the Holy Ghost.\end{Resp}
\begin{Resp}\Rsign The whole earth is full of His glory.\end{Resp}
\switchcolumn*

\LA
\LessonHead{Lectio IX}
\switchcolumn
\EN
\LessonHead{Lesson IX}
\switchcolumn*

\LA
\noindent Ex magna ergo parte diffísus est, qui virtútem non dedit majestáti, sed præséntiæ corporáli. Salútem ítaque fílio pétiit, et tamen in fide dubitávit; quia eum ad quem vénerat, et poténtem ad curándum crédidit, et tamen moriénti fílio esse abséntem putávit. Sed Dóminus, qui rogátur ut vadat, quia non desit ubi invitátur, índicat: solo jussu salútem réddidit, qui voluntáte ómnia creávit.\par
\switchcolumn
\EN
\noindent This was therefore a grievously imperfect faith, in attributing the virtue not to Christ's Majesty, but to His bodily presence. Thus it was that his faith was still unsound, even while he was asking for his son's health. For, though he believed concerning Him unto Whom he came that He was mighty to save, yet he thought also that at that moment He was absent from his dying child. But the Lord, being asked to go, showed that, wherever He is called on, He is there, and being He Who, by a simple act of will, brought all things into being, gave health by a simple command.\par
\switchcolumn*

\LA
\TeDeum{Te Deum}
\switchcolumn
\EN
\TeDeum{Te Deum}
\switchcolumn*

\end{paracol}

\Office{Dominica XXI Post Pentecosten}{Twenty-first Sunday after Pentecost}{Semiduplex}
\addcontentsline{toc}{subsection}{Dominica XXI Post Pentecosten \textit{— Twenty-first Sunday after Pentecost}}
\begin{paracol}{2}
\LA
\RefLine{Lectiones I–VI: de Scriptura occurrente.}
\switchcolumn
\EN
\RefLine{Lessons I–VI: from the Scripture of the day.}
\switchcolumn*

\LA
\LessonHead{Lectio VII}
\switchcolumn
\EN
\LessonHead{Lesson VII}
\switchcolumn*

\LA
\LessonTitle{Léctio sancti Evangélii secúndum Matthǽum}
\switchcolumn
\EN
\LessonTitle{From the Holy Gospel according to Matthew}
\switchcolumn*

\LA
\Cite{Matt 18:23-35}
\switchcolumn
\EN
\Cite{Matt 18:23-35}
\switchcolumn*

\LA
\noindent In illo témpore: Dixit Jesus discípulis suis parábolam hanc: Assimilátum est regnum cælórum hómini regi, qui vóluit ratiónem pónere cum servis suis. Et réliqua.\par
\switchcolumn
\EN
\noindent At that time, Jesus spoke unto His disciples this parable: The kingdom of heaven is likened unto a certain king, which would take account of his servants. And so on.\par
\switchcolumn*

\LA
\LessonTitle{Homilía sancti Hierónymi Presbýteri}
\switchcolumn
\EN
\LessonTitle{Homily by St. Jerome, Priest (at Bethlehem.)}
\switchcolumn*

\LA
\Source{Lib. 3 Comm. in cap. 18 Matth.}
\switchcolumn
\EN
\Source{Bk. III Comment. on Matth. XVIII}
\switchcolumn*

\LA
\noindent Familiáre est Syris, et máxime Palæstínis, ad omnem sermónem suum parábolas júngere: ut quod per simplex præcéptum tenéri ab auditóribus non potest, per similitúdinem exempláque teneátur. Præcépit ítaque Dóminus Petro sub comparatióne regis et dómini, et servi, qui débitor decem míllium talentórum a dómino rogans véniam impetráverat, ut ipse quoque dimíttat consérvis suis minóra peccántibus. Si enim ille rex et dóminus servo debitóri decem míllia talentórum tam fácile dimísit: quanto magis servi consérvis suis debent minóra dimíttere?\par
\switchcolumn
\EN
\noindent It is a way much in use with the Syrians, and especially with the inhabitants of Palestine, to illustrate their discourse with parables, that what their hearers may not be able to catch so easily when spoken plainly, they may lay hold on by dint of comparisons and examples. Thus it was that the Lord, by an allegory about a Royal master and a servant who owed him ten thousand talents, and who obtained by entreaty forgiveness of the debt, taught Peter how it was his duty to forgive his fellow-servants their comparatively trifling offences. For if that Royal master so readily forgave his servant his debt of ten thousand talents, should not his servants much more forgive lesser debts unto their fellows?\par
\switchcolumn*

\LA
\begin{Resp}\Rsign Refúlsit sol in clípeos áureos, et resplenduérunt montes ab eis: \Star{} Et fortitúdo géntium dissipáta est.\end{Resp}
\begin{Resp}\Vsign Erat enim exércitus magnus valde et fortis: et appropiávit Judas et exércitus ejus in prǽlio.\end{Resp}
\begin{Resp}\Rsign Et fortitúdo géntium dissipáta est.\end{Resp}
\switchcolumn
\EN
\begin{Resp}\Rsign The sun shone upon the shields of gold, and the mountains glistered therewith; \Star{} And the army of the heathens was spread abroad.\end{Resp}
\begin{Resp}\Vsign For the army was very great and mighty then Judas and his host drew near and entered into battle.\end{Resp}
\begin{Resp}\Rsign And the army of the heathens was spread abroad.\end{Resp}
\switchcolumn*

\LA
\LessonHead{Lectio VIII}
\switchcolumn
\EN
\LessonHead{Lesson VIII}
\switchcolumn*

\LA
\noindent Quod ut maniféstius fiat, dicámus sub exémplo. Si quis nostrum commíserit adultérium, homicídium, sacrilégium; majóra crímina decem míllium talentórum rogántibus dimittúntur, si et ipsi dimíttant minóra peccántibus. Sin autem ob factam contuméliam simus implacábiles, et propter amárum verbum pérpetes habeámus discórdias; nonne nobis vidémur recte redigéndi in cárcerem, et sub exémplo óperis nostri hoc ágere, ut majórum nobis delictórum vénia non relaxétur?\par
\switchcolumn
\EN
\noindent Let put this more clearly, let us take a case. If one of us were to commit adultery, or murder, or sacrilege, our sin, great like a debt of ten thousand talents, would be forgiven us in answer to prayer, if we also from our heart forgive our brethren their trespasses against us. But if we refuse to forgive a slight, and keep up unceasing enmity because of an unkind word, how just doth it appear that we should be cast into prison, and entail on ourselves, by the example of our own deeds, that our great debt should not be forgiven unto us.\par
\switchcolumn*

\LA
\begin{Resp}\Rsign Duo Séraphim clamábant alter ad álterum: \Star{} Sanctus, sanctus, sanctus Dóminus Deus Sábaoth: * Plena est omnis terra glória ejus.\end{Resp}
\begin{Resp}\Vsign Tres sunt qui testimónium dant in cælo: Pater, Verbum, et Spíritus Sanctus: et hi tres unum sunt.\end{Resp}
\begin{Resp}\Rsign Sanctus, sanctus, sanctus Dóminus Deus Sábaoth.\end{Resp}
\begin{Resp}\Vsign Glória Patri, et Fílio, \Star{} et Spirítui Sancto.\end{Resp}
\begin{Resp}\Rsign Plena est omnis terra glória ejus.\end{Resp}
\switchcolumn
\EN
\begin{Resp}\Rsign One Seraph cried unto another: \Star{} Holy, Holy, Holy is the Lord God of hosts; the whole earth is full of His glory.\end{Resp}
\begin{Resp}\Vsign There are Three That bear record in heaven: the Father, the Word, and the Holy Ghost, and these Three are One.\end{Resp}
\begin{Resp}\Rsign Holy, Holy, Holy is the Lord God of hosts.\end{Resp}
\begin{Resp}\Vsign Glory be to the Father, and to the Son, \Star{} and to the Holy Ghost.\end{Resp}
\begin{Resp}\Rsign The whole earth is full of His glory.\end{Resp}
\switchcolumn*

\LA
\LessonHead{Lectio IX}
\switchcolumn
\EN
\LessonHead{Lesson IX}
\switchcolumn*

\LA

\switchcolumn
\EN
\LessonTitle{“So likewise shall My heavenly Father do also unto you, if ye from your hearts forgive not every one his brother their trespasses.” God's awful purpose can be turned and changed but if we will not forgive unto our brethren small things, God will not forgive us great things. And if we forgive them, it must be from our hearts. Any one can say: I have nothing against such-an-one he knoweth what he hath done, and God will judge him for it I do not care what he doeth I have forgiven him. But the Lord maketh His sentence clear, and destroyeth such a mockery of peace as this, where He saith: “So likewise shall My heavenly Father do also unto you, if ye from your hearts forgive not every one his brother their trespasses.”}
\switchcolumn*

\LA
\noindent Sic et Pater meus cæléstis fáciet vobis, si non remiséritis unusquísque fratri suo de córdibus vestris. Formidolósa senténtia, si juxta nostram mentem senténtia Dei fléctitur, atque mutátur: si parva frátribus non dimíttimus, magna nobis a Deo non dimitténtur. Et quia potest unusquísque dícere: Nihil hábeo contra eum, ipse novit, habet Deum júdicem; non mihi curæ est quid velit ágere, ego ignóvi ei: confírmat senténtiam suam, et omnem simulatiónem fictæ pacis evértit, dicens: Si non remiséritis unusquísque fratri suo de córdibus vestris.\par
\switchcolumn
\EN

\switchcolumn*

\LA
\TeDeum{Te Deum}
\switchcolumn
\EN
\TeDeum{Te Deum}
\switchcolumn*

\end{paracol}

\Office{Dominica XXII Post Pentecosten}{Twenty-second Sunday after Pentecost}{Semiduplex}
\addcontentsline{toc}{subsection}{Dominica XXII Post Pentecosten \textit{— Twenty-second Sunday after Pentecost}}
\begin{paracol}{2}
\LA
\RefLine{Lectiones I–VI: de Scriptura occurrente.}
\switchcolumn
\EN
\RefLine{Lessons I–VI: from the Scripture of the day.}
\switchcolumn*

\LA
\LessonHead{Lectio VII}
\switchcolumn
\EN
\LessonHead{Lesson VII}
\switchcolumn*

\LA
\LessonTitle{Léctio sancti Evangélii secúndum Matthǽum}
\switchcolumn
\EN
\LessonTitle{From the Holy Gospel according to Matthew}
\switchcolumn*

\LA
\Cite{Matt 22:15-21}
\switchcolumn
\EN
\Cite{Matt 22:15-21}
\switchcolumn*

\LA
\noindent In illo témpore: Abeúntes pharisǽi consílium iniérunt ut cáperent Jesum in sermóne. Et réliqua.\par
\switchcolumn
\EN
\noindent At that time: The Pharisees went and took counsel how they might entangle Jesus in His talk. And so on.\par
\switchcolumn*

\LA
\LessonTitle{Homilía sancti Hilárii Epíscopi}
\switchcolumn
\EN
\LessonTitle{Homily by St. Hilary, Bishop (of Poitiers.)}
\switchcolumn*

\LA
\Source{Comment. in Matth. can. 23}
\switchcolumn
\EN
\Source{Comm. on Matth. Can. 23}
\switchcolumn*

\LA
\noindent Frequénter pharisǽi commovéntur, et occasiónem insimulándi eum habére ex prætéritis non possunt. Cádere enim vítium in gesta ejus et dicta non póterat; sed de malítiæ afféctu, in omnem se inquisitiónem reperiúndæ accusatiónis exténdunt. Namque a sǽculi vítiis, atque a superstitiónibus humanárum religiónum, univérsos ad spem regni cæléstis vocábat. Igitur an violáret sǽculi potestátem, de propósitæ interrogatiónis condicióne perténtant; an vidélicet reddi tribútum Cǽsari oportéret.\par
\switchcolumn
\EN
\noindent The Pharisees had oftentimes been put to confusion, and were not able to find any ground to accuse Him out of anything that He had hitherto said or done. His words and works are, of necessity, faultless, but still, from spite, they set themselves to seek in every direction for some cause to accuse Him. He was calling all to turn away from the corruptions of the world, and the superstitious practices of devotion invented by men, and to fix their hopes upon the kingdom of heaven. They therefore arranged a question calculated to entrap Him into an offense against civil government, namely: “Is it lawful to give tribute unto Caesar or not?”\par
\switchcolumn*

\LA
\begin{Resp}\Rsign Laudábilis pópulus, \Star{} Quem Dóminus exercítuum benedíxit dicens: Opus mánuum meárum tu es, heréditas mea Israël.\end{Resp}
\begin{Resp}\Vsign Beáta gens, cujus est Dóminus Deus, pópulus eléctus in hereditátem.\end{Resp}
\begin{Resp}\Rsign Quem Dóminus exercítuum benedíxit dicens: Opus mánuum meárum tu es, heréditas mea Israël.\end{Resp}
\switchcolumn
\EN
\begin{Resp}\Rsign Blessed is the people \Star{} Whom the Lord of hosts hath blessed, saying O Israel thou art the work of Mine own hands, thou art Mine own inheritance.\end{Resp}
\begin{Resp}\Vsign Blessed is the nation whose God is the Lord, and the people whom He hath chosen for His own inheritance.\end{Resp}
\begin{Resp}\Rsign Whom the Lord of hosts hath blessed, saying O Israel thou art the work of Mine own hands, thou art Mine own inheritance.\end{Resp}
\switchcolumn*

\LA
\LessonHead{Lectio VIII}
\switchcolumn
\EN
\LessonHead{Lesson VIII}
\switchcolumn*

\LA
\noindent Qui intérna cognitiónum secréta cognóscens (Deus enim nihil eórum quæ intra hóminem sunt abscónsa, non speculátur) afférri sibi denárium jussit, et quæsívit cujus et inscríptio esset et forma. Pharisǽi respondérunt: Cǽsaris eam esse. Quibus ait: Cǽsari redhibénda esse quæ Cǽsaris sunt; Deo autem reddénda esse, quæ Dei sunt. O plenam miráculi responsiónem, et perféctam dicti cæléstis absolutiónem! Ita ómnia inter contémptum sǽculi, et contuméliam lædéndi Cǽsaris temperávit, ut curis ómnibus et offíciis humánis devótas Deo mentes absólveret, cum Cǽsari quæ ejus essent, reddénda decérnit.\par
\switchcolumn
\EN
\noindent But Jesus perceived their wickedness, for in sooth there is nothing hidden in the heart of man, but what God seeth it, “and said Why tempt ye Me, ye hypocrites? Show Me the tribute-money. And they brought unto Him a penny. And He saith unto them Whose is this image and superscription They say unto Him Caesar's. Then saith He unto them: Render therefore unto Caesar the things which are Caesar's, and unto God the things that are God's.” How wonderful is this answer? How perfect the fulfillment of the Divine Law herein prescribed So beautifully doth He here strike the balance between caring not for the things of the world, on the one hand, and the offense of injuring Caesar, on the other, that He proveth the perfect freedom of minds, however devoted to God, to discharge all human cares and duties, by commanding them to render unto Caesar the things which are Caesar's.\par
\switchcolumn*

\LA
\begin{Resp}\Rsign Duo Séraphim clamábant alter ad álterum: \Star{} Sanctus, sanctus, sanctus Dóminus Deus Sábaoth: * Plena est omnis terra glória ejus.\end{Resp}
\begin{Resp}\Vsign Tres sunt qui testimónium dant in cælo: Pater, Verbum, et Spíritus Sanctus: et hi tres unum sunt.\end{Resp}
\begin{Resp}\Rsign Sanctus, sanctus, sanctus Dóminus Deus Sábaoth.\end{Resp}
\begin{Resp}\Vsign Glória Patri, et Fílio, \Star{} et Spirítui Sancto.\end{Resp}
\begin{Resp}\Rsign Plena est omnis terra glória ejus.\end{Resp}
\switchcolumn
\EN
\begin{Resp}\Rsign One Seraph cried unto another: \Star{} Holy, Holy, Holy is the Lord God of hosts; the whole earth is full of His glory.\end{Resp}
\begin{Resp}\Vsign There are Three That bear record in heaven: the Father, the Word, and the Holy Ghost, and these Three are One.\end{Resp}
\begin{Resp}\Rsign Holy, Holy, Holy is the Lord God of hosts.\end{Resp}
\begin{Resp}\Vsign Glory be to the Father, and to the Son, \Star{} and to the Holy Ghost.\end{Resp}
\begin{Resp}\Rsign The whole earth is full of His glory.\end{Resp}
\switchcolumn*

\LA
\LessonHead{Lectio IX}
\switchcolumn
\EN
\LessonHead{Lesson IX}
\switchcolumn*

\LA
\noindent Si enim nihil ejus penes nos reséderit, conditióne reddéndi ei quæ sua sunt, non tenébimur. Porro autem si rebus illíus incubámus, si jure potestátis ejus útimur, et nos tamquam mercenários aliéni patrimónii procuratióni subjícimus; extra querélam injúriæ est, Cǽsari redhibéri quod Cǽsaris est, Deo autem quæ ejus sunt própria, réddere nos oportére, corpus, ánimam, voluntátem. Ab eo enim hæc profécta atque aucta retinémus: proínde condígnum est, ut ei se totum reddant, cui debére se récolunt et oríginem et proféctum.\par
\switchcolumn
\EN
\noindent If we have nothing which is Caesar's, then we have nothing which we are bound to render unto him. But if we are concerned with the things which are his, if we are entrusted by him with the use of delegated power, if we are subject to him as paid servants to take care of property which is not our own, there can be no dispute but that it is our duty to render unto Caesar the things which are Caesar's. But unto God all of us are bound always to render the things that are God's, that is to say, our body, soul, and will. These are things which we hold from Him, and whereof He is the Author and Maker. It is therefore simply just that they, who acknowledge that they owe to Him their being and creation, should render to Him all that they are.\par
\switchcolumn*

\LA
\TeDeum{Te Deum}
\switchcolumn
\EN
\TeDeum{Te Deum}
\switchcolumn*

\end{paracol}

\Office{Dominica XXIII Post Pentecosten}{Twenty-third Sunday after Pentecost}{Semiduplex}
\addcontentsline{toc}{subsection}{Dominica XXIII Post Pentecosten \textit{— Twenty-third Sunday after Pentecost}}
\begin{paracol}{2}
\LA
\RefLine{Lectiones I–VI: de Scriptura occurrente.}
\switchcolumn
\EN
\RefLine{Lessons I–VI: from the Scripture of the day.}
\switchcolumn*

\LA
\LessonHead{Lectio VII}
\switchcolumn
\EN
\LessonHead{Lesson VII}
\switchcolumn*

\LA
\LessonTitle{Léctio sancti Evangélii secúndum Matthǽum}
\switchcolumn
\EN
\LessonTitle{From the Holy Gospel according to Matthew}
\switchcolumn*

\LA
\Cite{Matt 9:18-26}
\switchcolumn
\EN
\Cite{Matt 9:18-26}
\switchcolumn*

\LA
\noindent In illo témpore: Loquénte Jesu ad turbas, ecce princeps unus accéssit et adorábat eum dicens: Dómine, fília mea modo defúncta est. Et réliqua.\par
\switchcolumn
\EN
\noindent At that time: While Jesus spake these things unto the multitudes, behold, there came a certain ruler, and worshipped Him, saying Lord, my daughter is even now dead. And so on.\par
\switchcolumn*

\LA
\LessonTitle{Homilía sancti Hierónymi Presbýteri}
\switchcolumn
\EN
\LessonTitle{Homily by St. Jerome, Priest (at Bethlehem.)}
\switchcolumn*

\LA
\Source{Liber 1 Comment. in cap. 9 Matthæi}
\switchcolumn
\EN
\Source{Bk. I Comment. on Matth. IX}
\switchcolumn*

\LA
\noindent Octávum signum est, in quo princeps suscitári póstulat fíliam suam, nolens de mystério veræ circumcisiónis exclúdi: sed subíntrat múlier sánguine fluens, et octávo sanátur loco; ut príncipis fília de hoc exclúsa número, véniat ad nonum, juxta illud quod in Psalmis dícitur: Æthiópia prævéniet manus ejus Deo; et, Cum intráverit plenitúdo géntium, tunc omnis Israël salvus fiet.\par
\switchcolumn
\EN
\noindent The eighth miracle took place upon the occasion when a certain ruler, desiring not to be kept out of the mystery of the true circumcision, besought Christ to recall his daughter to life. But a woman, which was diseased with an issue of blood, thrust herself in, and her cure occupieth the eighth place, so that the resurrection of the ruler's daughter is postponed and made the ninth in enumeration even as it is written in the Psalms: “Ethiopia shall soon stretch forth her hands unto God.” (lxvii. 32.) And again: “Blindness in part is happened to Israel, until the fullness of the Gentiles is come in and so all Israel shall be saved.” (Rom. xi. 25, 26.)\par
\switchcolumn*

\LA
\begin{Resp}\Rsign Laudábilis pópulus, \Star{} Quem Dóminus exercítuum benedíxit dicens: Opus mánuum meárum tu es, heréditas mea Israël.\end{Resp}
\begin{Resp}\Vsign Beáta gens, cujus est Dóminus Deus, pópulus eléctus in hereditátem.\end{Resp}
\begin{Resp}\Rsign Quem Dóminus exercítuum benedíxit dicens: Opus mánuum meárum tu es, heréditas mea Israël.\end{Resp}
\switchcolumn
\EN
\begin{Resp}\Rsign Blessed is the people \Star{} Whom the Lord of hosts hath blessed, saying O Israel thou art the work of Mine own hands, thou art Mine own inheritance.\end{Resp}
\begin{Resp}\Vsign Blessed is the nation whose God is the Lord, and the people whom He hath chosen for His own inheritance.\end{Resp}
\begin{Resp}\Rsign Whom the Lord of hosts hath blessed, saying O Israel thou art the work of Mine own hands, thou art Mine own inheritance.\end{Resp}
\switchcolumn*

\LA
\LessonHead{Lectio VIII}
\switchcolumn
\EN
\LessonHead{Lesson VIII}
\switchcolumn*

\LA
\noindent Et ecce múlier, quæ sánguinis fluxum patiebátur duódecim annis, accéssit retro, et tétigit fímbriam vestiménti ejus. In Evangélio secúndum Lucam scríbitur, quod príncipis fília duódecim annos habéret ætátis. Nota ergo, quod eo témpore hæc múlier, id est, Géntium pópulus cœ́perit ægrotáre, quo genus Judæórum credíderat. Nisi enim ex comparatióne virtútum, vítium non osténditur.\par
\switchcolumn
\EN
\noindent And, behold, a woman, which was diseased with an issue of blood, twelve years, came behind Him, and touched the hem of His garment.” In the Gospel according to Luke (viii. 42) it is written that the ruler's daughter was about twelve years of age. Note therefore that this woman, who typifies the Gentiles, had been diseased for the same time that the Jewish nation, typified by the ruler's daughter, had been living in faith. We see not clearly the hideousness of evil, until we compare it with good.\par
\switchcolumn*

\LA
\begin{Resp}\Rsign Duo Séraphim clamábant alter ad álterum: \Star{} Sanctus, sanctus, sanctus Dóminus Deus Sábaoth: * Plena est omnis terra glória ejus.\end{Resp}
\begin{Resp}\Vsign Tres sunt qui testimónium dant in cælo: Pater, Verbum, et Spíritus Sanctus: et hi tres unum sunt.\end{Resp}
\begin{Resp}\Rsign Sanctus, sanctus, sanctus Dóminus Deus Sábaoth.\end{Resp}
\begin{Resp}\Vsign Glória Patri, et Fílio, \Star{} et Spirítui Sancto.\end{Resp}
\begin{Resp}\Rsign Plena est omnis terra glória ejus.\end{Resp}
\switchcolumn
\EN
\begin{Resp}\Rsign One Seraph cried unto another: \Star{} Holy, Holy, Holy is the Lord God of hosts; the whole earth is full of His glory.\end{Resp}
\begin{Resp}\Vsign There are Three That bear record in heaven: the Father, the Word, and the Holy Ghost, and these Three are One.\end{Resp}
\begin{Resp}\Rsign Holy, Holy, Holy is the Lord God of hosts.\end{Resp}
\begin{Resp}\Vsign Glory be to the Father, and to the Son, \Star{} and to the Holy Ghost.\end{Resp}
\begin{Resp}\Rsign The whole earth is full of His glory.\end{Resp}
\switchcolumn*

\LA
\LessonHead{Lectio IX}
\switchcolumn
\EN
\LessonHead{Lesson IX}
\switchcolumn*

\LA
\noindent Hæc autem múlier sánguine fluens, non in domo, non in urbe accédit ad Dóminum, quia juxta legem úrbibus excludebátur; sed in itínere, ambulánte Dómino; ut, dum pergit ad áliam, ália curarétur. Unde dicunt et Apóstoli: Vobis quidem oportébat prædicári verbum Dei: sed quóniam vos judicástis indígnos salúte, transgrédimur ad gentes.\par
\switchcolumn
\EN
\noindent This woman with the issue of blood came not to the Lord in an house or in a city, for such as she were by the Law banished out of cities, (Lev. xv. 25,) but in the way, as He walked so that the Lord healed one, even while He was on the road to heal another. Whence also the Apostles said: “It was necessary that the word of God should first have been spoken to you but, seeing ye put it from you, and judge yourselves unworthy of everlasting life, lo, we turn to the Gentiles.” (Acts xiii. 46.)\par
\switchcolumn*

\LA
\TeDeum{Te Deum}
\switchcolumn
\EN
\TeDeum{Te Deum}
\switchcolumn*

\end{paracol}

\Office{Dominica XXIV et Ultima Post Pentecosten}{Dominica XXIV et ultima Post Pentecosten}{Semiduplex}
\addcontentsline{toc}{subsection}{Dominica XXIV et Ultima Post Pentecosten \textit{— Dominica XXIV et ultima Post Pentecosten}}
\begin{paracol}{2}
\LA
\RefLine{Lectiones I–VI: de Scriptura occurrente.}
\switchcolumn
\EN
\RefLine{Lessons I–VI: from the Scripture of the day.}
\switchcolumn*

\LA
\LessonHead{Lectio VII}
\switchcolumn
\EN
\LessonHead{Lesson VII}
\switchcolumn*

\LA
\LessonTitle{Léctio sancti Evangélii secúndum Matthǽum}
\switchcolumn
\EN
\LessonTitle{From the Holy Gospel according to Matthew}
\switchcolumn*

\LA
\Cite{Matt 24:15-35}
\switchcolumn
\EN
\Cite{Matt 24:15-35}
\switchcolumn*

\LA
\noindent In illo témpore: Dixit Jesus discípulis suis: Cum vidéritis abominatiónem desolatiónis, quæ dicta est a Daniéle prophéta, stantem in loco sancto: qui legit, intéllegat. Et réliqua.\par
\switchcolumn
\EN
\noindent At that time, Jesus said unto his disciples: When ye shall see the abomination of desolation, spoken of by Daniel the Prophet, stand in the Holy Place, (whoso readeth, let him understand.) And so on.\par
\switchcolumn*

\LA
\LessonTitle{Homilía sancti Hierónymi Presbýteri}
\switchcolumn
\EN
\LessonTitle{Homily by St. Jerome, Priest (at Bethlehem.)}
\switchcolumn*

\LA
\Source{Liber 4 Comment. in cap. 24 Matthæi}
\switchcolumn
\EN
\Source{Bk. IV Comm. on Matth. XXIV}
\switchcolumn*

\LA
\noindent Quando ad intellegéntiam provocámur, mýsticum monstrátur esse quod dictum est. Légimus autem in Daniéle hoc modo: Et in dimídio hebdómadis auferétur sacrifícium et libámina; et in templo abominátio desolatiónum erit, usque ad consummatiónem témporis, et consummátio dábitur super solitúdinem. De hoc et Apóstolus lóquitur, quod homo iniquitátis, et adversárius elevándus sit contra omne quod dícitur Deus et cólitur; ita ut áudeat stare in templo Dei, et osténdere quod ipse sit Deus: cujus advéntus secúndum operatiónem sátanæ déstruat eos, et ad Dei solitúdinem rédigat, qui se suscéperint.\par
\switchcolumn
\EN
\noindent This injunction to whoso readeth, to understand, showeth that there is here something mysterious. In Daniel we read as followeth: “And in the midst of the week the sacrifice and the oblations shall be taken away and in the temple there shall be the abomination of desolation, even until the consummation of the time and a consummation shall be given to the desolation.” (ix. 27.) It is of this same thing that the Apostle speaketh, when he saith that a man of iniquity, even an adversary, shall be exalted against whatsoever is called God, or is worshipped so that he shall even dare to stand in the temple of God, and to show himself as God whose coming shall, according to the working of Satan, destroy and banish away from God all who shall receive him.\par
\switchcolumn*

\LA
\begin{Resp}\Rsign Laudábilis pópulus, \Star{} Quem Dóminus exercítuum benedíxit dicens: Opus mánuum meárum tu es, heréditas mea Israël.\end{Resp}
\begin{Resp}\Vsign Beáta gens, cujus est Dóminus Deus, pópulus eléctus in hereditátem.\end{Resp}
\begin{Resp}\Rsign Quem Dóminus exercítuum benedíxit dicens: Opus mánuum meárum tu es, heréditas mea Israël.\end{Resp}
\switchcolumn
\EN
\begin{Resp}\Rsign Blessed is the people \Star{} Whom the Lord of hosts hath blessed, saying O Israel thou art the work of Mine own hands, thou art Mine own inheritance.\end{Resp}
\begin{Resp}\Vsign Blessed is the nation whose God is the Lord, and the people whom He hath chosen for His own inheritance.\end{Resp}
\begin{Resp}\Rsign Whom the Lord of hosts hath blessed, saying O Israel thou art the work of Mine own hands, thou art Mine own inheritance.\end{Resp}
\switchcolumn*

\LA
\LessonHead{Lectio VIII}
\switchcolumn
\EN
\LessonHead{Lesson VIII}
\switchcolumn*

\LA
\noindent Potest autem simpliciter aut de Antichristo accipi, aut de imagine Cæsaris, quam Pilátus pósuit in templo, aut de Hadriáni equestri státua, quæ in ipso sancto sanctórum loco usque in præséntem diem stetit. Abominátio quoque secúndum veterem Scripturam idolum nuncupátur; et idcirco additur, desolatiónis, quod in desolato templo atque destructo idolum positum sit.\par
\switchcolumn
\EN
\noindent This prophecy may be understood either (first) simply of Antichrist, (secondly) of the statue of Caesar, which Pilate set up in the Temple, or (thirdly) of the statue of Hadrian on horse-back, which hath stood, even until our own day, upon the site of the Holy of Holies. In the Scriptures of the Old Testament “abomination” is a word very often used for an idol, and the farther title “of desolation” is added to identify an idol erected upon the site of the desolate and ruined temple.\par
\switchcolumn*

\LA
\begin{Resp}\Rsign Duo Séraphim clamábant alter ad álterum: \Star{} Sanctus, sanctus, sanctus Dóminus Deus Sábaoth: * Plena est omnis terra glória ejus.\end{Resp}
\begin{Resp}\Vsign Tres sunt qui testimónium dant in cælo: Pater, Verbum, et Spíritus Sanctus: et hi tres unum sunt.\end{Resp}
\begin{Resp}\Rsign Sanctus, sanctus, sanctus Dóminus Deus Sábaoth.\end{Resp}
\begin{Resp}\Vsign Glória Patri, et Fílio, \Star{} et Spirítui Sancto.\end{Resp}
\begin{Resp}\Rsign Plena est omnis terra glória ejus.\end{Resp}
\switchcolumn
\EN
\begin{Resp}\Rsign One Seraph cried unto another: \Star{} Holy, Holy, Holy is the Lord God of hosts; the whole earth is full of His glory.\end{Resp}
\begin{Resp}\Vsign There are Three That bear record in heaven: the Father, the Word, and the Holy Ghost, and these Three are One.\end{Resp}
\begin{Resp}\Rsign Holy, Holy, Holy is the Lord God of hosts.\end{Resp}
\begin{Resp}\Vsign Glory be to the Father, and to the Son, \Star{} and to the Holy Ghost.\end{Resp}
\begin{Resp}\Rsign The whole earth is full of His glory.\end{Resp}
\switchcolumn*

\LA
\LessonHead{Lectio IX}
\switchcolumn
\EN
\LessonHead{Lesson IX}
\switchcolumn*

\LA
\noindent Abominátio desolatiónis intelligi potest et omne dogma perversum: quod cum vidérimus stare in loco sancto, hoc est in Ecclésia, et se osténdere Deum, debémus fugere de Judæa ad montes, hoc est, dimissa occidénte littera, et Judáica pravitáte, appropinquare móntibus ætérnis, de quibus illúminat mirabíliter Deus; et esse in tecto et in dómate, quo non possint igníta diaboli jácula perveníre, nec descéndere et tóllere aliquid de domo conversatiónis prístinæ, nec quærere quæ retrórsum sunt, sed magis sérere in agro spiritualium Scripturárum, ut fructus capiámus ex eo; nec tóllere álteram túnicam, quam Apóstoli habére prohibentur.\par
\switchcolumn
\EN
\noindent But we may also understand by the abomination of desolation, any bad doctrine and when we see such a thing get a standing in the Holy Place, that is, in the Church, and showing itself that it is God, that is, pretending that it is His revealed truth, then will be the time when it will be our duty to flee from Judea into the mountains, that is to say, to leave the letter, which passeth away, and all guise of Jewish superstition, and to hie us unto the everlasting hills, from whence God doth right wondrously cause His light to shine forth. (Ps. lxxv. 5.) Then will it be our duty to find ourselves under a roof and in an house, wherethrough the fiery darts of the wicked one can never pierce to smite us, and not to come down to take anything out of the house of our old conversation, or to have regard unto those things which are behind but rather to sow in the field of the spiritual Scriptures, that we may reap thereof a bountiful harvest neither to have two coats, that thing forbidden to Apostles. (Matth. x. 10.)\par
\switchcolumn*

\LA
\TeDeum{Te Deum}
\switchcolumn
\EN
\TeDeum{Te Deum}
\switchcolumn*

\end{paracol}

\SeasonTitle{Dominicæ et Feriæ mensis Septembris}{Sundays and Ferias of September}

\addcontentsline{toc}{section}{Dominicæ et Feriæ mensis Septembris \textit{— Sundays and Ferias of September}}

\Office{Dominica I Septembris}{First Sunday of September}{}
\addcontentsline{toc}{subsection}{Dominica I Septembris \textit{— First Sunday of September}}
\begin{paracol}{2}
\LA
\LessonHead{Lectio I}
\switchcolumn
\EN
\LessonHead{Lesson I}
\switchcolumn*

\LA
\LessonTitle{Incipit liber Job}
\switchcolumn
\EN
\LessonTitle{Lesson from the book of Job}
\switchcolumn*

\LA
\Cite{Job 1:1-3}
\switchcolumn
\EN
\Cite{Job 1:1-3}
\switchcolumn*

\LA
\noindent Vir erat in terra Hus nómine Job, et erat vir ille simplex et rectus ac timens Deum et recédens a malo. Natíque sunt ei septem fílii et tres fíliæ. Et fuit posséssio ejus septem míllia óvium et tria míllia camelórum, quingénta quoque juga boum et quingéntæ ásinæ ac família multa nimis: erátque vir ille magnus inter omnes Orientáles.\par
\switchcolumn
\EN
\noindent There was a man in the land of Hus, whose name was Job, and that man was simple and upright, and fearing God, and avoiding evil. And there were born to him seven sons and three daughters. And his possession was seven thousand sheep, and three thousand camels, and five hundred yoke of oxen, and five hundred she asses, and a family exceeding great: and this man was great among all the people of the east.\par
\switchcolumn*

\LA
\begin{Resp}\Rsign Si bona suscépimus de manu Dei, mala autem quare non sustineámus? \Star{} Dóminus dedit, Dóminus ábstulit; sicut Dómino plácuit, ita factum est: sit nomen Dómini benedíctum.\end{Resp}
\begin{Resp}\Vsign Nudus egréssus sum de útero matris meæ et nudus revértar illuc.\end{Resp}
\begin{Resp}\Rsign Dóminus dedit, Dóminus ábstulit; sicut Dómino plácuit, ita factum est: sit nomen Dómini benedíctum.\end{Resp}
\switchcolumn
\EN
\begin{Resp}\Rsign What shall we receive good at the hand of God, and shall we not receive evil? \Star{} The Lord gave and the Lord hath taken away. As the Lord hath pleased, so hath it befallen. Blessed be the Name of the Lord.\end{Resp}
\begin{Resp}\Vsign Naked came I out of my mother's womb, and naked shall I return thither.\end{Resp}
\begin{Resp}\Rsign The Lord gave and the Lord hath taken away. As the Lord hath pleased, so hath it befallen. Blessed be the Name of the Lord.\end{Resp}
\switchcolumn*

\LA
\LessonHead{Lectio II}
\switchcolumn
\EN
\LessonHead{Lesson II}
\switchcolumn*

\LA
\Cite{Job 1:4-5}
\switchcolumn
\EN
\Cite{Job 1:4-5}
\switchcolumn*

\LA
\noindent Et ibant fílii ejus et faciébant convívium per domos, unusquísque in die suo. Et mitténtes vocábant tres soróres suas, ut coméderent et bíberent cum eis. Cumque in orbem transíssent dies convívii, mittébat ad eos Job et sanctificábat illos, consurgénsque dilúculo offerébat holocáusta pro síngulis; dicébat enim: Ne forte peccáverint fílii mei et benedíxerint Deo in córdibus suis. Sic faciébat Job cunctis diébus.\par
\switchcolumn
\EN
\noindent And his sons went, and made a feast by houses every one in his day. And sending they called their three sisters to eat and drink with them. And when the days of their feasting were gone about, Job sent to them, and sanctified them: and rising up early offered holocausts for every one of them. For he said: Lest perhaps my sons have sinned, and have blessed God in their hearts. So did Job all days.\par
\switchcolumn*

\LA
\begin{Resp}\Rsign Antequam cómedam, suspíro, et tamquam inundántes aquæ sic rugítus meus; quia timor, quem timébam, evénit mihi, et quod verébar áccidit. Nonne dissimulávi? nonne sílui? nonne quiévi? \Star{} Et venit super me indignátio.\end{Resp}
\begin{Resp}\Vsign Ecce non est auxílium mihi in me, et necessárii quoque mei recessérunt a me.\end{Resp}
\begin{Resp}\Rsign Et venit super me indignátio.\end{Resp}
\switchcolumn
\EN
\begin{Resp}\Rsign My sighing cometh before I eat, and my roarings are poured out like the waters, for the thing which I greatly feared is come upon me, and that which I was afraid of is come unto me. Was not I silent? Held not I my peace? Was not I at rest? \Star{} And trouble came.\end{Resp}
\begin{Resp}\Vsign Behold, I cannot help myself, and they that were needful unto me have forsaken me.\end{Resp}
\begin{Resp}\Rsign And trouble came.\end{Resp}
\switchcolumn*

\LA
\LessonHead{Lectio III}
\switchcolumn
\EN
\LessonHead{Lesson III}
\switchcolumn*

\LA
\Cite{Job 1:6-11}
\switchcolumn
\EN
\Cite{Job 1:6-11}
\switchcolumn*

\LA
\noindent Quadam autem die, cum veníssent fílii Dei ut assísterent coram Dómino, áffuit inter eos étiam Satan. Cui dixit Dóminus: Unde venis? Qui respóndens ait: Circuívi terram et perambulávi eam. Dixítque Dóminus ad eum: Numquid considerásti servum meum Job, quod non sit ei símilis in terra, homo simplex et rectus ac timens Deum et recédens a malo? Cui respóndens Satan ait: Numquid Job frustra timet Deum? nonne tu vallásti eum ac domum ejus, universámque substántiam per circúitum opéribus mánuum ejus benedixísti, et posséssio ejus crevit in terra? Sed exténde páululum manum tuam et tange cuncta quæ póssidet, nisi in fáciem benedíxerit tibi.\par
\switchcolumn
\EN
\noindent Now on a certain day when the sons of God came to stand before the Lord, Satan also was present among them. And the Lord said to him: Whence comest thou And he answered and said: I have gone round about the earth, and walked through it. And the Lord said to him: Hast thou considered my servant Job, that there is none like him in the earth, a simple and upright man, and fearing God, and avoiding evil? And Satan answering, said: Doth Job fear God in vain Hast not thou made a fence for him, and his house, and all his substance round about, blessed the works of his hands, and his possession hath increased on the earth But stretch forth thy hand a little, and touch all that he hath, and see if he blesseth thee not to thy face.\par
\switchcolumn*

\LA
\begin{Resp}\Rsign Quare detraxístis sermónibus veritátis? ad increpándum verba compónitis et subvértere nitímini amícum vestrum: \Star{} Verúmtamen quæ cogitástis, expléte.\end{Resp}
\begin{Resp}\Vsign Quod justum est, judicáte; et non inveniétis in lingua mea iniquitátem.\end{Resp}
\begin{Resp}\Rsign Verúmtamen quæ cogitástis, expléte.\end{Resp}
\begin{Resp}\Vsign Glória Patri, et Fílio, \Star{} et Spirítui Sancto.\end{Resp}
\begin{Resp}\Rsign Verúmtamen quæ cogitástis, expléte.\end{Resp}
\switchcolumn
\EN
\begin{Resp}\Rsign Why do ye argue against the words of truth? Do ye imagine words to reprove me and strive to confound one that is your friend? \Star{} Nevertheless, finish that ye have in mind.\end{Resp}
\begin{Resp}\Vsign Judge that which is just, and ye shall find no iniquity in my tongue.\end{Resp}
\begin{Resp}\Rsign Nevertheless, finish that ye have in mind.\end{Resp}
\begin{Resp}\Vsign Glory be to the Father, and to the Son, \Star{} and to the Holy Ghost.\end{Resp}
\begin{Resp}\Rsign Nevertheless, finish that ye have in mind.\end{Resp}
\switchcolumn*

\LA
\LessonHead{Lectio IV}
\switchcolumn
\EN
\LessonHead{Lesson IV}
\switchcolumn*

\LA
\LessonTitle{Ex libro Morálium sancti Gregórii Papæ}
\switchcolumn
\EN
\LessonTitle{From the Book of Moral (Reflections upon Job) written by Pope St. Gregory (the Great.)}
\switchcolumn*

\LA
\Source{Lectio iv, Lib. 2, Cap. 1}
\switchcolumn
\EN
\Source{Bk. ii. ch. i.}
\switchcolumn*

\LA
\noindent Scriptúra sacra mentis óculis quasi quoddam spéculum oppónitur, ut intérna nostra fácies in ipsa videátur. Ibi étenim fœda, ibi pulchra nostra cognóscimus: ibi sentímus quantum profícimus, ibi a proféctu quam longe distámus. Narrat autem gesta Sanctórum, et ad imitatiónem corda próvocat infirmórum; dumque illórum victrícia facta commémorat, contra vitiórum prǽlia, debília nostra confírmat: fitque verbis illíus, ut eo mens minus inter certámina trépidet, quo ante se pósitos tot virórum fórtium triúmphos videt.\par
\switchcolumn
\EN
\noindent The Holy Scripture is put before the eyes of our mind somewhat after the fashion of a looking-glass, that we may see therein the aspect of our inward man. Therein we see what are our unsightly, and what our comely traits thereby we judge how we are growing, and how far yet we are from fullness of stature. The Holy Scripture telleth of the doings of the Saints, and stirreth up the heart of us weaklings to follow them. While it maketh memorial of their victorious deeds, it strengtheneth our frailty to strive against sin. And so by the words of the Scripture it cometh to pass that the soul trembleth less at the battle, for that she seeth how many times the enemies before her have been beaten by brave men.\par
\switchcolumn*

\LA
\begin{Resp}\Rsign Indúta est caro mea putrédine, et sórdibus púlveris cutis mea áruit et contrácta est: \Star{} Meménto mei, Dómine, quóniam ventus est vita mea.\end{Resp}
\begin{Resp}\Vsign Dies mei velócius transiérunt quam a texénte tela succíditur, et consúmpti sunt absque ulla spe.\end{Resp}
\begin{Resp}\Rsign Meménto mei, Dómine, quóniam ventus est vita mea.\end{Resp}
\switchcolumn
\EN
\begin{Resp}\Rsign My flesh is clothed with worms and clods of dust. My skin is dry and drawn together. \Star{} Remember me, O Lord, for my life is wind.\end{Resp}
\begin{Resp}\Vsign My days are swifter than a weaver's shuttle, and are spent without hope.\end{Resp}
\begin{Resp}\Rsign Remember me, Lord, for my life is wind.\end{Resp}
\switchcolumn*

\LA
\LessonHead{Lectio V}
\switchcolumn
\EN
\LessonHead{Lesson V}
\switchcolumn*

\LA
\noindent Nonnúmquam vero non solum nobis eórum virtútes ásserit, sed étiam casus innotéscit; ut et in victória fórtium, quod imitándo arrípere, et rursum videámus in lápsibus quid debeámus timére. Ecce enim Job descríbitur tentatióne auctus, sed David tentatióne prostrátus; ut et majórum virtus spem nostram fóveat, et majórum casus ad cautélam nos humilitátis accíngat: quátenus dum illa gaudéntes súblevant, ista metuéntes premant; et audiéntis ánimus illinc spei fidúcia, hinc humilitáte timóris erudítus, nec temeritáte supérbiat, quia formídine prémitur; nec pressus timóre despéret, quia ad spei fidúciam virtútis exémplo roborátur.\par
\switchcolumn
\EN
\noindent And some while the Scripture showeth unto us, not only how the Saints fought bravely, but also how they fell, that we may see by the example of the mighty, not only what weapons we must take, if we would conquer, but also what snares we must keep clear of, if we would avoid falling. For example, here is Job on the one hand, waxing nobler under trial, and on the other hand, David, tried, and failing utterly. And so the glory of the great strengtheneth our hope, and the backsliding of the same doth stir us up to be watchful and lowly the one cheering us with gladness, and the other putting us on our guard through fear, so that the soul of him which heareth of these things may by the one gain sure and certain hope, and by the other fearfulness and watchfulness, and so neither be rashly puffed up, nor hopelessly cast down, nor may faint under the weight of dread, forasmuch as she is stirred up to trustfulness by the example of him who triumphed.\par
\switchcolumn*

\LA
\begin{Resp}\Rsign Páucitas diérum meórum finiétur brevi; dimítte me, Dómine, ut plangam páululum dolórem meum, \Star{} Antequam vadam ad terram tenebrósam et opértam mortis calígine.\end{Resp}
\begin{Resp}\Vsign Manus tuæ, Dómine, fecérunt me, et plasmavérunt me totum in circúitu; et sic repénte præcípitas me?\end{Resp}
\begin{Resp}\Rsign Antequam vadam ad terram tenebrósam et opértam mortis calígine.\end{Resp}
\switchcolumn
\EN
\begin{Resp}\Rsign My days are few, and in a short while they will be ended. Let me alone, then, O Lord, that I may bewail my sorrow a little; \Star{} Before I go to the land of darkness and of the shadow of death.\end{Resp}
\begin{Resp}\Vsign Thine hands, O Lord, have made me, and fashioned me together round about, and yet dost Thou forthwith destroy me.\end{Resp}
\begin{Resp}\Rsign Before I go to the land of darkness and of the shadow of death.\end{Resp}
\switchcolumn*

\LA
\LessonHead{Lectio VI}
\switchcolumn
\EN
\LessonHead{Lesson VI}
\switchcolumn*

\LA
\Source{Lib. 1 Moral., cap. 1}
\switchcolumn
\EN
\Source{Bk. i. ch. 1}
\switchcolumn*

\LA
\noindent Vir erat in terra Hus, nómine Job. Idcírco sanctus vir ubi habitáverit dícitur, ut ejus méritum virtútis exprimátur. Hus namque quis nésciat quod sit in terra Gentílium? Gentílitas autem eo obligáta vítiis éxstitit, quo cognitiónem sui Conditóris ignorávit. Dicátur ítaque ubi habitáverit; ut hoc ejus láudibus profíciat, quod bonus inter malos fuit. Neque enim valde laudábile est, bonum esse cum bonis, sed bonum esse cum malis. Sicut enim gravióris culpæ est, inter bonus bonum non esse; ita imménsi est præcónii, bonum étiam inter malos exstitísse.\par
\switchcolumn
\EN
\noindent There was a man in the land of Uz, whose name was Job. We are told where this holy man lived, that thereby we may gauge the worth of his bravery. Who knoweth not that Uz is a place in the countries of the Gentiles The Gentile world had been so degraded and corrupted by sin, that they had ceased to know that they had a Maker. Therefore is it told us where Job dwelt, that it may redound to his praise that he was good in the midst of the wicked. It is not very praiseworthy to be good among the good, but to be good among the bad. For even as it is more grievous to be bad among the good, so is it right praiseworthy to have remained good among the bad.\par
\switchcolumn*

\LA
\begin{Resp}\Rsign Non abscóndas me, Dómine, a fácie tua: manum tuam longe fac a me, \Star{} Et formído tua non me térreat.\end{Resp}
\begin{Resp}\Vsign Córripe me, Dómine, in misericórdia, non in furóre tuo, ne forte ad níhilum rédigas me.\end{Resp}
\begin{Resp}\Rsign Et formído tua non me térreat.\end{Resp}
\begin{Resp}\Vsign Glória Patri, et Fílio, \Star{} et Spirítui Sancto.\end{Resp}
\begin{Resp}\Rsign Et formído tua non me térreat.\end{Resp}
\switchcolumn
\EN
\begin{Resp}\Rsign Hide not thy face from me, O Lord. Withdraw not thine hand far from me, \Star{} And let not thy dread make me afraid.\end{Resp}
\begin{Resp}\Vsign O Lord, correct me but in mercy not in thine anger, lest Thou bring me to nothing.\end{Resp}
\begin{Resp}\Rsign And let not thy dread make me afraid.\end{Resp}
\begin{Resp}\Vsign Glory be to the Father, and to the Son, \Star{} and to the Holy Ghost.\end{Resp}
\begin{Resp}\Rsign And let not thy dread make me afraid.\end{Resp}
\switchcolumn*

\end{paracol}

\Office{Feria II infra Hebdomadam I Septembris}{Monday of the first week of September}{}
\addcontentsline{toc}{subsection}{Feria II infra Hebdomadam I Septembris \textit{— Monday of the first week of September}}
\begin{paracol}{2}
\LA
\LessonHead{Lectio I}
\switchcolumn
\EN
\LessonHead{Lesson I}
\switchcolumn*

\LA
\LessonTitle{De libro Job}
\switchcolumn
\EN
\LessonTitle{Lesson from the book of Job}
\switchcolumn*

\LA
\Cite{Job 1:13-16}
\switchcolumn
\EN
\Cite{Job 1:13-16}
\switchcolumn*

\LA
\noindent Cum autem quadam die fílii et fíliæ ejus coméderent et bíberent vinum in domo fratris sui primogéniti, núntius venit ad Job, qui díceret: Boves arábant et ásinæ pascebántur juxta eos, et irruérunt Sabǽi tulerúntque ómnia et púeros percussérunt gládio, et evási ego solus ut nuntiárem tibi. Cumque adhuc ille loquerétur, venit alter et dixit: Ignis Dei cécidit e cælo, et tactas oves puerósque consúmpsit, et effúgi ego solus ut nuntiárem tibi.\par
\switchcolumn
\EN
\noindent Now upon a certain day when his sons and daughters were eating and drinking wine in the house of their eldest brother, There came a messenger to Job, and said: The oxen were ploughing, and the asses feeding beside them, And the Sabeans rushed in, and took all away, and slew the servants with the sword, and I alone have escaped to tell thee. And while he was yet speaking, another came, and said: The fire of God fell from heaven, and striking the sheep and the servants, hath consumed them, and I alone have escaped to tell thee.\par
\switchcolumn*

\LA
\begin{Resp}\Rsign Versa est in luctum cíthara mea, et órganum meum in vocem fléntium: \Star{} Parce mihi, Dómine, nihil enim sunt dies mei.\end{Resp}
\begin{Resp}\Vsign Cutis mea denigráta est super me, et ossa mea aruérunt.\end{Resp}
\begin{Resp}\Rsign Parce mihi, Dómine, nihil enim sunt dies mei.\end{Resp}
\switchcolumn
\EN
\begin{Resp}\Rsign My harp is turned to mourning, and my organ into the voice of them that weep. \Star{} Let me alone, O Lord, for my days are vanity.\end{Resp}
\begin{Resp}\Vsign My skin is black upon me, and my bones are burned with heat.\end{Resp}
\begin{Resp}\Rsign Let me alone, O Lord, for my days are vanity.\end{Resp}
\switchcolumn*

\LA
\LessonHead{Lectio II}
\switchcolumn
\EN
\LessonHead{Lesson II}
\switchcolumn*

\LA
\Cite{Job 1:17-19}
\switchcolumn
\EN
\Cite{Job 1:17-19}
\switchcolumn*

\LA
\noindent Sed, et illo adhuc loquénte, venit álius et dixit: Chaldǽi fecérunt tres turmas et invasérunt camélos et tulérunt eos, necnon et púeros percussérunt gládio, et ego fugi solus, ut nuntiárem tibi. Adhuc loquebátur ille, et ecce álius intrávit, et dixit: Fíliis tuis et filiábus vescéntibus et bibéntibus vinum in domo fratris sui primogéniti, repénte ventus véhemens írruit a regióne desérti, et concússit quátuor ángulos domus, quæ córruens oppréssit líberos tuos, et mórtui sunt, et effúgi ego solus ut nuntiárem tibi.\par
\switchcolumn
\EN
\noindent And while he also was yet speaking, there came another, and said: The Chaldeans made three troops, and have fallen upon the camels, and taken them, moreover they have slain the servants with the sword, and I alone have escaped to tell thee. He was yet speaking, and behold another came in, and said: thy sons and daughters were eating and drinking wine in the house of their elder brother: A violent wind came on a sudden from the side of the desert, and shook the four corners of the house, and it fell upon thy children and they are dead, and I alone have escaped to fell thee.\par
\switchcolumn*

\LA
\begin{Resp}\Rsign Utinam appenderéntur peccáta mea, quibus iram mérui, \Star{} Et calámitas, quam pátior, in statéra.\end{Resp}
\begin{Resp}\Vsign Quasi aréna maris hæc grávior apparéret, unde et verba mea dolóre sunt plena.\end{Resp}
\begin{Resp}\Rsign Et calámitas, quam pátior, in statéra.\end{Resp}
\switchcolumn
\EN
\begin{Resp}\Rsign O that my sins, whereby I have deserved wrath; \Star{} And the calamity whereunder I suffer, were laid in the balances together.\end{Resp}
\begin{Resp}\Vsign For now it would appear heavier than the sand of the sea, therefore also my words are full of sorrow.\end{Resp}
\begin{Resp}\Rsign And the calamity, whereunder I suffer, were laid in the balances together.\end{Resp}
\switchcolumn*

\LA
\LessonHead{Lectio III}
\switchcolumn
\EN
\LessonHead{Lesson III}
\switchcolumn*

\LA
\Cite{Job 1:20-22}
\switchcolumn
\EN
\Cite{Job 1:20-22}
\switchcolumn*

\LA
\noindent Tunc surréxit Job et scidit vestiménta sua, et, tonso cápite, córruens in terram adorávit, et dixit: Nudus egréssus sum de útero matris meæ et nudus revértar illuc. Dóminus dedit, Dóminus ábstulit; sicut Dómino plácuit, ita factum est: sit nomen Dómini benedíctum. In ómnibus his non peccávit Job lábiis suis, neque stultum quid contra Deum locútus est.\par
\switchcolumn
\EN
\noindent Then Job rose up, and rent his garments, and having shaven his head fell down upon the ground and worshipped, And said: Naked came I out of my mother's womb, and naked shall I return thither: the Lord gave, and the Lord hath taken away: as it hath pleased the Lord so is it done: blessed be the name of the Lord. In all these things Job sinned not by his lips, nor spoke he any foolish thing against God.\par
\switchcolumn*

\LA
\begin{Resp}\Rsign Quare detraxístis sermónibus veritátis? ad increpándum verba compónitis et subvértere nitímini amícum vestrum: \Star{} Verúmtamen quæ cogitástis, expléte.\end{Resp}
\begin{Resp}\Vsign Quod justum est, judicáte; et non inveniétis in lingua mea iniquitátem.\end{Resp}
\begin{Resp}\Rsign Verúmtamen quæ cogitástis, expléte.\end{Resp}
\begin{Resp}\Vsign Glória Patri, et Fílio, \Star{} et Spirítui Sancto.\end{Resp}
\begin{Resp}\Rsign Verúmtamen quæ cogitástis, expléte.\end{Resp}
\switchcolumn
\EN
\begin{Resp}\Rsign Why do ye argue against the words of truth? Do ye imagine words to reprove me and strive to confound one that is your friend? \Star{} Nevertheless, finish that ye have in mind.\end{Resp}
\begin{Resp}\Vsign Judge that which is just, and ye shall find no iniquity in my tongue.\end{Resp}
\begin{Resp}\Rsign Nevertheless, finish that ye have in mind.\end{Resp}
\begin{Resp}\Vsign Glory be to the Father, and to the Son, \Star{} and to the Holy Ghost.\end{Resp}
\begin{Resp}\Rsign Nevertheless, finish that ye have in mind.\end{Resp}
\switchcolumn*

\end{paracol}

\Office{Feria III infra Hebdomadam I Septembris}{Tuesday of the first week of September}{}
\addcontentsline{toc}{subsection}{Feria III infra Hebdomadam I Septembris \textit{— Tuesday of the first week of September}}
\begin{paracol}{2}
\LA
\LessonHead{Lectio I}
\switchcolumn
\EN
\LessonHead{Lesson I}
\switchcolumn*

\LA
\LessonTitle{De libro Job}
\switchcolumn
\EN
\LessonTitle{Lesson from the book of Job}
\switchcolumn*

\LA
\Cite{Job 2:1-5}
\switchcolumn
\EN
\Cite{Job 2:1-5}
\switchcolumn*

\LA
\noindent Factum est autem, cum quadam die veníssent fílii Dei et starent coram Dómino, venísset quoque Satan inter eos et staret in conspéctu ejus, ut díceret Dóminus ad Satan: Unde venis? Qui respóndens ait: Circuívi terram et perambulávi eam. Et dixit Dóminus ad Satan: Numquid considerásti servum meum Job, quod non sit ei símilis in terra, vir simplex et rectus ac timens Deum et recédens a malo et adhuc rétinens innocéntiam? Tu autem commovísti me advérsus eum, ut afflígerem eum frustra. Cui respóndens Satan ait: Pellem pro pelle, et cuncta, quæ habet homo, dabit pro ánima sua. Alióqui mitte manum tuam, et tange os ejus et carnem; et tunc vidébis quod in fáciem benedícat tibi.\par
\switchcolumn
\EN
\noindent And it came to pass, when on a certain day the sons of God came, and stood before the Lord, and Satan came among them, and stood in his sight, That the Lord said to Satan: Whence comest thou? And he answered and said: I have gone round about the earth, and walked through it. And the Lord said to Satan: Hast thou considered my servant Job, that there is none like him in the earth, a man simple, and upright, and fearing God, and avoiding evil, and still keeping his innocence. But thou hast moved me against him, that I should afflict him without cause. And Satan answered, and said: Skin for skin, and all that a man hath he will give for his life: But put forth thy hand, and touch his bone and his flesh, and then thou shalt see that he will bless thee to thy face.\par
\switchcolumn*

\LA
\begin{Resp}\Rsign Indúta est caro mea putrédine, et sórdibus púlveris cutis mea áruit et contrácta est: \Star{} Meménto mei, Dómine, quóniam ventus est vita mea.\end{Resp}
\begin{Resp}\Vsign Dies mei velócius transiérunt quam a texénte tela succíditur, et consúmpti sunt absque ulla spe.\end{Resp}
\begin{Resp}\Rsign Meménto mei, Dómine, quóniam ventus est vita mea.\end{Resp}
\switchcolumn
\EN
\begin{Resp}\Rsign My flesh is clothed with worms and clods of dust. My skin is dry and drawn together. \Star{} Remember me, O Lord, for my life is wind.\end{Resp}
\begin{Resp}\Vsign My days are swifter than a weaver's shuttle, and are spent without hope.\end{Resp}
\begin{Resp}\Rsign Remember me, Lord, for my life is wind.\end{Resp}
\switchcolumn*

\LA
\LessonHead{Lectio II}
\switchcolumn
\EN
\LessonHead{Lesson II}
\switchcolumn*

\LA
\Cite{Job 2:6-10}
\switchcolumn
\EN
\Cite{Job 2:6-10}
\switchcolumn*

\LA
\noindent Dixit ergo Dóminus ad Satan: Ecce in manu tua est, verúmtamen ánimam illíus serva. Egréssus ígitur Satan a fácie Dómini percússit Job úlcere péssimo a planta pedis usque ad vérticem ejus; qui testa sániem radébat sedens in sterquilínio. Dixit autem illi uxor sua: Adhuc tu pérmanes in simplicitáte tua? Bénedic Deo et mórere. Qui ait ad illam: Quasi una de stultis muliéribus locúta es: si bona suscépimus de manu Dei, mala quare non suscipiámus? In ómnibus his non peccávit Job lábiis suis.\par
\switchcolumn
\EN
\noindent And the Lord said to Satan: Behold he is in thy hand, but yet save his life. So Satan went forth from the presence of the Lord, and struck Job with a very grievous ulcer, from the sole of the foot even to the top of his head: And he took a potsherd and scraped the corrupt matter, sitting on a dunghill. And his wife said to him: Dost thou still continue in thy simplicity? Bless God and die. And he said to her: Thou hast; spoken like one of the foolish women: if we have received good things at the hand of God, why should we not receive evil? In all these things Job did not sin with his lips.\par
\switchcolumn*

\LA
\begin{Resp}\Rsign Páucitas diérum meórum finiétur brevi; dimítte me, Dómine, ut plangam páululum dolórem meum, \Star{} Antequam vadam ad terram tenebrósam et opértam mortis calígine.\end{Resp}
\begin{Resp}\Vsign Manus tuæ, Dómine, fecérunt me, et plasmavérunt me totum in circúitu; et sic repénte præcípitas me?\end{Resp}
\begin{Resp}\Rsign Antequam vadam ad terram tenebrósam et opértam mortis calígine.\end{Resp}
\switchcolumn
\EN
\begin{Resp}\Rsign My days are few, and in a short while they will be ended. Let me alone, then, O Lord, that I may bewail my sorrow a little; \Star{} Before I go to the land of darkness and of the shadow of death.\end{Resp}
\begin{Resp}\Vsign Thine hands, O Lord, have made me, and fashioned me together round about, and yet dost Thou forthwith destroy me.\end{Resp}
\begin{Resp}\Rsign Before I go to the land of darkness and of the shadow of death.\end{Resp}
\switchcolumn*

\LA
\LessonHead{Lectio III}
\switchcolumn
\EN
\LessonHead{Lesson III}
\switchcolumn*

\LA
\Cite{Job 2:11-13}
\switchcolumn
\EN
\Cite{Job 2:11-13}
\switchcolumn*

\LA
\noindent Igitur audiéntes tres amíci Job omne malum quod accidísset ei, venérunt sínguli de loco suo, Eliphaz Themanítes et Baldad Suhítes et Sophar Naamathítes. Condíxerant enim ut páriter veniéntes visitárent eum et consolaréntur. Cumque elevássent procul óculos suos, non cognovérunt eum et exclamántes ploravérunt, scissísque véstibus, sparsérunt púlverem super caput suum in cælum. Et sedérunt cum eo in terra septem diébus et septem nóctibus, et nemo loquebátur ei verbum; vidébant enim dolórem esse veheméntem.\par
\switchcolumn
\EN
\noindent Now when Job's three friends heard all the evil that had befallen him, they came every one from his own place, Alphas the Themanite, and Baldad the Suhite, and Sophar the Naamathite. For they had made an appointment to come together and visit him, and comfort him. And when they had lifted up their eyes afar off, they knew him not, and crying out they wept, and rending their garments they sprinkled dust upon their heads towards heaven. And they sat with him on the ground seven days and seven nights, and no man spoke to him a word: for they saw that his grief was very great.\par
\switchcolumn*

\LA
\begin{Resp}\Rsign Non abscóndas me, Dómine, a fácie tua: manum tuam longe fac a me, \Star{} Et formído tua non me térreat.\end{Resp}
\begin{Resp}\Vsign Córripe me, Dómine, in misericórdia, non in furóre tuo, ne forte ad níhilum rédigas me.\end{Resp}
\begin{Resp}\Rsign Et formído tua non me térreat.\end{Resp}
\begin{Resp}\Vsign Glória Patri, et Fílio, \Star{} et Spirítui Sancto.\end{Resp}
\begin{Resp}\Rsign Et formído tua non me térreat.\end{Resp}
\switchcolumn
\EN
\begin{Resp}\Rsign Hide not thy face from me, O Lord. Withdraw not thine hand far from me, \Star{} And let not thy dread make me afraid.\end{Resp}
\begin{Resp}\Vsign O Lord, correct me but in mercy not in thine anger, lest Thou bring me to nothing.\end{Resp}
\begin{Resp}\Rsign And let not thy dread make me afraid.\end{Resp}
\begin{Resp}\Vsign Glory be to the Father, and to the Son, \Star{} and to the Holy Ghost.\end{Resp}
\begin{Resp}\Rsign And let not thy dread make me afraid.\end{Resp}
\switchcolumn*

\end{paracol}

\Office{Feria IV infra Hebdomadam I Septembris}{Wednesday of the first week of September}{}
\addcontentsline{toc}{subsection}{Feria IV infra Hebdomadam I Septembris \textit{— Wednesday of the first week of September}}
\begin{paracol}{2}
\LA
\LessonHead{Lectio I}
\switchcolumn
\EN
\LessonHead{Lesson I}
\switchcolumn*

\LA
\LessonTitle{De libro Job}
\switchcolumn
\EN
\LessonTitle{Lesson from the book of Job}
\switchcolumn*

\LA
\Cite{Job 3:1-5}
\switchcolumn
\EN
\Cite{Job 3:1-5}
\switchcolumn*

\LA
\noindent Post hæc apéruit Job os suum et maledíxit diéi suo, et locútus est: Péreat dies in qua natus sum, et nox in qua dictum est: Concéptus est homo. Dies ille vertátur in ténebras: non requírat eum Deus désuper, et non illustrétur lúmine; obscúrent eum ténebræ et umbra mortis; óccupet eum calígo, et involvátur amaritúdine.\par
\switchcolumn
\EN
\noindent After this Job opened his mouth, and cursed his day, and he said: Let the day perish wherein I was born, and the night in which it was said: A man child is conceived. Let that day be turned into darkness, let not God regard it from above, and let not the light shine upon it. Let darkness, and the shadow of death cover it, let a mist overspread it, and let it be wrapped up in bitterness.\par
\switchcolumn*

\LA
\begin{Resp}\Rsign Quis mihi tríbuat, ut in inférno prótegas me et abscóndas me, donec pertránseat furor tuus, Dómine, nisi tu, qui solus es Deus? \Star{} Et constítuas mihi tempus, in quo recordéris mei?\end{Resp}
\begin{Resp}\Vsign Numquid sicut dies hóminis dies tui, ut quæras iniquitátem meam; cum sit nemo, qui de manu tua possit erúere.\end{Resp}
\begin{Resp}\Rsign Et constítuas mihi tempus, in quo recordéris mei?\end{Resp}
\switchcolumn
\EN
\begin{Resp}\Rsign O that Thou wouldest hide me in the grave; that Thou wouldest keep me secret, until thy wrath be past, even thine, O Lord, Thou That alone art God; \Star{} That Thou wouldest appoint me a set time, and remember me!\end{Resp}
\begin{Resp}\Vsign Are thy days as the days of man, that Thou inquirest after mine iniquity, and there is none that can deliver out of thine hand.\end{Resp}
\begin{Resp}\Rsign That Thou wouldest appoint me a set time, and remember me!\end{Resp}
\switchcolumn*

\LA
\LessonHead{Lectio II}
\switchcolumn
\EN
\LessonHead{Lesson II}
\switchcolumn*

\LA
\Cite{Job 3:6-10}
\switchcolumn
\EN
\Cite{Job 3:6-10}
\switchcolumn*

\LA
\noindent Noctem illam tenebrósus turbo possídeat; non computétur in diébus anni, nec numerétur in ménsibus. Sit nox illa solitária, nec laude digna; maledícant ei qui maledícunt diéi, qui paráti sunt suscitáre levíathan. Obtenebréntur stellæ calígine ejus, exspéctet lucem et non vídeat, nec ortum surgéntis auróræ; quia non conclúsit óstia ventris, qui portávit me, nec ábstulit mala ab óculis meis.\par
\switchcolumn
\EN
\noindent Let a darksome whirlwind seize upon that night, let it not be counted in the days of the year, nor numbered in the months. Let that night be solitary, and not worthy of praise. Let them curse it who curse the day. who are ready to raise up a leviathan: Let the stars be darkened with the mist thereof: let it expect light and not see it, nor the rising of the dawning of the day: Because it shut not up the doors of the womb that bore me, nor took away evils from my eyes.\par
\switchcolumn*

\LA
\begin{Resp}\Rsign Utinam appenderéntur peccáta mea, quibus iram mérui, \Star{} Et calámitas, quam pátior, in statéra.\end{Resp}
\begin{Resp}\Vsign Quasi aréna maris hæc grávior apparéret, unde et verba mea dolóre sunt plena.\end{Resp}
\begin{Resp}\Rsign Et calámitas, quam pátior, in statéra.\end{Resp}
\switchcolumn
\EN
\begin{Resp}\Rsign O that my sins, whereby I have deserved wrath; \Star{} And the calamity whereunder I suffer, were laid in the balances together.\end{Resp}
\begin{Resp}\Vsign For now it would appear heavier than the sand of the sea, therefore also my words are full of sorrow.\end{Resp}
\begin{Resp}\Rsign And the calamity, whereunder I suffer, were laid in the balances together.\end{Resp}
\switchcolumn*

\LA
\LessonHead{Lectio III}
\switchcolumn
\EN
\LessonHead{Lesson III}
\switchcolumn*

\LA
\Cite{Job 3:11-16}
\switchcolumn
\EN
\Cite{Job 3:11-16}
\switchcolumn*

\LA
\noindent Quare non in vulva mórtuus sum? egréssus ex útero non statim périi? quare excéptus génibus? cur lactátus ubéribus? Nunc enim dórmiens silérem et somno meo requiéscerem cum régibus et consúlibus terræ, qui ædíficant sibi solitúdines, aut cum princípibus, qui póssident aurum et replent domos suas argénto; aut sicut abortívum abscónditum non subsísterem, vel qui concépti non vidérunt lucem.\par
\switchcolumn
\EN
\noindent Why did I not die in the womb, why did I not perish when I came out of the belly? Why received upon the knees? why suckled at the breasts? For now I should have been asleep and still, and should have rest in my sleep. With kings and consuls of the earth, who build themselves solitudes: Or with princes, that possess gold, and All their houses with silver: Or as a hidden untimely birth I should not be, or as they that being conceived have not seen the light.\par
\switchcolumn*

\LA
\begin{Resp}\Rsign Quare detraxístis sermónibus veritátis? ad increpándum verba compónitis et subvértere nitímini amícum vestrum: \Star{} Verúmtamen quæ cogitástis, expléte.\end{Resp}
\begin{Resp}\Vsign Quod justum est, judicáte; et non inveniétis in lingua mea iniquitátem.\end{Resp}
\begin{Resp}\Rsign Verúmtamen quæ cogitástis, expléte.\end{Resp}
\begin{Resp}\Vsign Glória Patri, et Fílio, \Star{} et Spirítui Sancto.\end{Resp}
\begin{Resp}\Rsign Verúmtamen quæ cogitástis, expléte.\end{Resp}
\switchcolumn
\EN
\begin{Resp}\Rsign Why do ye argue against the words of truth? Do ye imagine words to reprove me and strive to confound one that is your friend? \Star{} Nevertheless, finish that ye have in mind.\end{Resp}
\begin{Resp}\Vsign Judge that which is just, and ye shall find no iniquity in my tongue.\end{Resp}
\begin{Resp}\Rsign Nevertheless, finish that ye have in mind.\end{Resp}
\begin{Resp}\Vsign Glory be to the Father, and to the Son, \Star{} and to the Holy Ghost.\end{Resp}
\begin{Resp}\Rsign Nevertheless, finish that ye have in mind.\end{Resp}
\switchcolumn*

\end{paracol}

\Office{Feria V infra Hebdomadam I Septembris}{Thursday of the first week of September}{}
\addcontentsline{toc}{subsection}{Feria V infra Hebdomadam I Septembris \textit{— Thursday of the first week of September}}
\begin{paracol}{2}
\LA
\LessonHead{Lectio I}
\switchcolumn
\EN
\LessonHead{Lesson I}
\switchcolumn*

\LA
\LessonTitle{De libro Job}
\switchcolumn
\EN
\LessonTitle{Lesson from the book of Job}
\switchcolumn*

\LA
\Cite{Job 4:1-6}
\switchcolumn
\EN
\Cite{Job 4:1-6}
\switchcolumn*

\LA
\noindent Respóndens autem Eliphaz Themanítes, dixit: Si cœpérimus loqui tibi, fórsitan moléste accípies; sed concéptum sermónem tenére quis póterit? Ecce docuísti multos et manus lassas roborásti, vacillántes confirmavérunt sermónes tui et génua treméntia confortásti. Nunc autem venit super te plaga, et defecísti; tétigit te, et conturbátus es. Ubi est timor tuus, fortitúdo tua, patiéntia tua et perféctio viárum tuárum?\par
\switchcolumn
\EN
\noindent Then Eliphaz the Themanite answered, and said: If we begin to speak to thee, perhaps thou wilt take it ill, but who can withhold the words he hath conceived? Behold thou hast taught many, and thou hast strengthened the weary hands: thy words have confirmed them that were staggering, and thou hast strengthened the trembling knees: But now the scourge is come upon thee, and thou faintest: it hath touched thee, and thou art troubled. Where is thy fear, thy fortitude, thy patience, and the perfection of thy ways?\par
\switchcolumn*

\LA
\begin{Resp}\Rsign Si bona suscépimus de manu Dei, mala autem quare non sustineámus? \Star{} Dóminus dedit, Dóminus ábstulit; sicut Dómino plácuit, ita factum est: sit nomen Dómini benedíctum.\end{Resp}
\begin{Resp}\Vsign Nudus egréssus sum de útero matris meæ et nudus revértar illuc.\end{Resp}
\begin{Resp}\Rsign Dóminus dedit, Dóminus ábstulit; sicut Dómino plácuit, ita factum est: sit nomen Dómini benedíctum.\end{Resp}
\switchcolumn
\EN
\begin{Resp}\Rsign What shall we receive good at the hand of God, and shall we not receive evil? \Star{} The Lord gave and the Lord hath taken away. As the Lord hath pleased, so hath it befallen. Blessed be the Name of the Lord.\end{Resp}
\begin{Resp}\Vsign Naked came I out of my mother's womb, and naked shall I return thither.\end{Resp}
\begin{Resp}\Rsign The Lord gave and the Lord hath taken away. As the Lord hath pleased, so hath it befallen. Blessed be the Name of the Lord.\end{Resp}
\switchcolumn*

\LA
\LessonHead{Lectio II}
\switchcolumn
\EN
\LessonHead{Lesson II}
\switchcolumn*

\LA
\Cite{Job 4:7-11}
\switchcolumn
\EN
\Cite{Job 4:7-11}
\switchcolumn*

\LA
\noindent Recordáre, óbsecro te, quis umquam ínnocens périit? aut quando recti deléti sunt? Quin pótius vidi eos, qui operántur iniquitátem et séminant dolóres et metunt eos, flante Deo perísse et spíritu iræ ejus esse consúmptos. Rugítus leónis et vox leǽnæ et dentes catulórum leónum contríti sunt; Tigris périit, eo quod non habéret prædam, et cátuli leónis dissipáti sunt.\par
\switchcolumn
\EN
\noindent Remember, I pray thee, who ever perished being innocent? or when were the just destroyed? On the contrary I have seen those who work iniquity, and sow sorrows, and reap them, Perishing by the blast of God, and consumed by the spirit of his wrath. The roaring of the lion, and the voice of the lioness, and the teeth of the whelps of lions are broken: The tiger hath perished for want of prey, and the young lions are scattered abroad.\par
\switchcolumn*

\LA
\begin{Resp}\Rsign Antequam cómedam, suspíro, et tamquam inundántes aquæ sic rugítus meus; quia timor, quem timébam, evénit mihi, et quod verébar áccidit. Nonne dissimulávi? nonne sílui? nonne quiévi? \Star{} Et venit super me indignátio.\end{Resp}
\begin{Resp}\Vsign Ecce non est auxílium mihi in me, et necessárii quoque mei recessérunt a me.\end{Resp}
\begin{Resp}\Rsign Et venit super me indignátio.\end{Resp}
\switchcolumn
\EN
\begin{Resp}\Rsign My sighing cometh before I eat, and my roarings are poured out like the waters, for the thing which I greatly feared is come upon me, and that which I was afraid of is come unto me. Was not I silent? Held not I my peace? Was not I at rest? \Star{} And trouble came.\end{Resp}
\begin{Resp}\Vsign Behold, I cannot help myself, and they that were needful unto me have forsaken me.\end{Resp}
\begin{Resp}\Rsign And trouble came.\end{Resp}
\switchcolumn*

\LA
\LessonHead{Lectio III}
\switchcolumn
\EN
\LessonHead{Lesson III}
\switchcolumn*

\LA
\Cite{Job 4:12-18}
\switchcolumn
\EN
\Cite{Job 4:12-18}
\switchcolumn*

\LA
\noindent Porro ad me dictum est verbum abscónditum, et quasi furtíve suscépit auris mea venas susúrri ejus. In horróre visiónis noctúrnæ, quando solet sopor occupáre hómines, Pavor ténuit me et tremor, et ómnia ossa mea pertérrita sunt, Et cum spíritus, me præsénte, transíret, inhorruérunt pili carnis meæ. Stetit quidam, cujus non agnoscébam vultum, imágo coram óculis meis, et vocem quasi auræ lenis audívi: Numquid homo Dei comparatióne justificábitur? aut Factóre suo púrior erit vir? Ecce qui sérviunt ei, non sunt stábiles, et in Angelis suis réperit pravitátem.\par
\switchcolumn
\EN
\noindent Now there was a word spoken to me in private, and my ears by stealth as it were received the veins of its whisper. In the horror of a vision by night, when deep sleep is wont to hold men, Fear seized upon me, and trembling, and all my bones were affrighted: And when a spirit passed before me, the hair of my flesh stood up. There stood one whose countenance I knew not, an image before my eyes, and I heard the voice as it were of a gentle wind: Shall man be justified in comparison of God, or shall a man be more pure than his maker? Behold they that serve him are not steadfast, and in his angels he found wickedness:\par
\switchcolumn*

\LA
\begin{Resp}\Rsign Quare detraxístis sermónibus veritátis? ad increpándum verba compónitis et subvértere nitímini amícum vestrum: \Star{} Verúmtamen quæ cogitástis, expléte.\end{Resp}
\begin{Resp}\Vsign Quod justum est, judicáte; et non inveniétis in lingua mea iniquitátem.\end{Resp}
\begin{Resp}\Rsign Verúmtamen quæ cogitástis, expléte.\end{Resp}
\begin{Resp}\Vsign Glória Patri, et Fílio, \Star{} et Spirítui Sancto.\end{Resp}
\begin{Resp}\Rsign Verúmtamen quæ cogitástis, expléte.\end{Resp}
\switchcolumn
\EN
\begin{Resp}\Rsign Why do ye argue against the words of truth? Do ye imagine words to reprove me and strive to confound one that is your friend? \Star{} Nevertheless, finish that ye have in mind.\end{Resp}
\begin{Resp}\Vsign Judge that which is just, and ye shall find no iniquity in my tongue.\end{Resp}
\begin{Resp}\Rsign Nevertheless, finish that ye have in mind.\end{Resp}
\begin{Resp}\Vsign Glory be to the Father, and to the Son, \Star{} and to the Holy Ghost.\end{Resp}
\begin{Resp}\Rsign Nevertheless, finish that ye have in mind.\end{Resp}
\switchcolumn*

\end{paracol}

\Office{Feria VI infra Hebdomadam I Septembris}{Friday of the first week of September}{}
\addcontentsline{toc}{subsection}{Feria VI infra Hebdomadam I Septembris \textit{— Friday of the first week of September}}
\begin{paracol}{2}
\LA
\LessonHead{Lectio I}
\switchcolumn
\EN
\LessonHead{Lesson I}
\switchcolumn*

\LA
\LessonTitle{De libro Job}
\switchcolumn
\EN
\LessonTitle{Lesson from the book of Job}
\switchcolumn*

\LA
\Cite{Job 6:1-4}
\switchcolumn
\EN
\Cite{Job 6:1-4}
\switchcolumn*

\LA
\noindent Respóndens autem Job, dixit: Utinam appenderéntur peccáta mea, quibus iram mérui, et calámitas, quam pátior, in statéra. Quasi aréna maris hæc grávior apparéret, unde et verba mea dolóre sunt plena; quia sagíttæ Dómini in me sunt, quarum indignátio ébibit spíritum meum, et terróres Dómini mílitant contra me.\par
\switchcolumn
\EN
\noindent But Job answered, and said: O that my sins, whereby I have deserved wrath, and the calamity that I suffer, were weighed in a balance. As the sand of the sea this would appear heavier: therefore my words are full of sorrow: For the arrows of the Lord are in me, the rage whereof drinketh up my spirit, and the terrors of the Lord war against me.\par
\switchcolumn*

\LA
\begin{Resp}\Rsign Indúta est caro mea putrédine, et sórdibus púlveris cutis mea áruit et contrácta est: \Star{} Meménto mei, Dómine, quóniam ventus est vita mea.\end{Resp}
\begin{Resp}\Vsign Dies mei velócius transiérunt quam a texénte tela succíditur, et consúmpti sunt absque ulla spe.\end{Resp}
\begin{Resp}\Rsign Meménto mei, Dómine, quóniam ventus est vita mea.\end{Resp}
\switchcolumn
\EN
\begin{Resp}\Rsign My flesh is clothed with worms and clods of dust. My skin is dry and drawn together. \Star{} Remember me, O Lord, for my life is wind.\end{Resp}
\begin{Resp}\Vsign My days are swifter than a weaver's shuttle, and are spent without hope.\end{Resp}
\begin{Resp}\Rsign Remember me, Lord, for my life is wind.\end{Resp}
\switchcolumn*

\LA
\LessonHead{Lectio II}
\switchcolumn
\EN
\LessonHead{Lesson II}
\switchcolumn*

\LA
\Cite{Job 6:5-7}
\switchcolumn
\EN
\Cite{Job 6:5-7}
\switchcolumn*

\LA
\noindent Numquid rúgiet ónager, cum habúerit herbam? aut múgiet bos, cum ante præsépe plenum stéterit? Aut póterit cómedi insúlsum, quod non est sale condítum? aut potest áliquis gustáre, quod gustátum affert mortem? Quæ prius nolébat tángere ánima mea, nunc præ angústia, cibi mei sunt.\par
\switchcolumn
\EN
\noindent Will the wild ass bray when he hath grass? or will the ox low when he standeth before a full manger? Or can an unsavoury thing be eaten, that is not seasoned with salt? or can a man taste that which when tasted bringeth death? The things which before my soul would not touch, now, through anguish are my meats.\par
\switchcolumn*

\LA
\begin{Resp}\Rsign Páucitas diérum meórum finiétur brevi; dimítte me, Dómine, ut plangam páululum dolórem meum, \Star{} Antequam vadam ad terram tenebrósam et opértam mortis calígine.\end{Resp}
\begin{Resp}\Vsign Manus tuæ, Dómine, fecérunt me, et plasmavérunt me totum in circúitu; et sic repénte præcípitas me?\end{Resp}
\begin{Resp}\Rsign Antequam vadam ad terram tenebrósam et opértam mortis calígine.\end{Resp}
\switchcolumn
\EN
\begin{Resp}\Rsign My days are few, and in a short while they will be ended. Let me alone, then, O Lord, that I may bewail my sorrow a little; \Star{} Before I go to the land of darkness and of the shadow of death.\end{Resp}
\begin{Resp}\Vsign Thine hands, O Lord, have made me, and fashioned me together round about, and yet dost Thou forthwith destroy me.\end{Resp}
\begin{Resp}\Rsign Before I go to the land of darkness and of the shadow of death.\end{Resp}
\switchcolumn*

\LA
\LessonHead{Lectio III}
\switchcolumn
\EN
\LessonHead{Lesson III}
\switchcolumn*

\LA
\Cite{Job 6:8-13}
\switchcolumn
\EN
\Cite{Job 6:8-13}
\switchcolumn*

\LA
\noindent Quis det ut véniat petítio mea, et quod exspécto tríbuat mihi Deus? Et qui cœpit ipse me cónterat, solvat manum suam et succídat me? Et hæc mihi sit consolátio, ut afflígens me dolóre non parcat, nec contradícam sermónibus Sancti. Quæ est enim fortitúdo mea, ut sustíneam? aut quis finis meus, ut patiénter agam? Nec fortitúdo lápidum fortitúdo mea, nec caro mea ǽnea est. Ecce non est auxílium mihi in me, et necessárii quoque mei recessérunt a me.\par
\switchcolumn
\EN
\noindent Who will grant that my request may come: and that God may give me what I look for? And that he that hath begun may destroy me, that he may let loose his hand, and cut me off? And that this may be my comfort, that afflicting me with sorrow, he spare not, nor I contradict the words of the Holy One. For what is my strength, that I can hold out? or what is my end that I should keep patience? My strength is not the strength of stones, nor is my flesh of brass. Behold there is no help for me in myself, and my familiar friends also are departed from me.\par
\switchcolumn*

\LA
\begin{Resp}\Rsign Non abscóndas me, Dómine, a fácie tua: manum tuam longe fac a me, \Star{} Et formído tua non me térreat.\end{Resp}
\begin{Resp}\Vsign Córripe me, Dómine, in misericórdia, non in furóre tuo, ne forte ad níhilum rédigas me.\end{Resp}
\begin{Resp}\Rsign Et formído tua non me térreat.\end{Resp}
\begin{Resp}\Vsign Glória Patri, et Fílio, \Star{} et Spirítui Sancto.\end{Resp}
\begin{Resp}\Rsign Et formído tua non me térreat.\end{Resp}
\switchcolumn
\EN
\begin{Resp}\Rsign Hide not thy face from me, O Lord. Withdraw not thine hand far from me, \Star{} And let not thy dread make me afraid.\end{Resp}
\begin{Resp}\Vsign O Lord, correct me but in mercy not in thine anger, lest Thou bring me to nothing.\end{Resp}
\begin{Resp}\Rsign And let not thy dread make me afraid.\end{Resp}
\begin{Resp}\Vsign Glory be to the Father, and to the Son, \Star{} and to the Holy Ghost.\end{Resp}
\begin{Resp}\Rsign And let not thy dread make me afraid.\end{Resp}
\switchcolumn*

\end{paracol}

\Office{Sabbato infra Hebdomadam I Septembris}{Saturday of the first week of September}{}
\addcontentsline{toc}{subsection}{Sabbato infra Hebdomadam I Septembris \textit{— Saturday of the first week of September}}
\begin{paracol}{2}
\LA
\LessonHead{Lectio I}
\switchcolumn
\EN
\LessonHead{Lesson I}
\switchcolumn*

\LA
\LessonTitle{De libro Job}
\switchcolumn
\EN
\LessonTitle{Lesson from the book of Job}
\switchcolumn*

\LA
\Cite{Job 7:1-4}
\switchcolumn
\EN
\Cite{Job 7:1-4}
\switchcolumn*

\LA
\noindent Milítia est vita hóminis super terram, et sicut dies mercenárii, dies ejus. Sicut servus desíderat umbram, et sicut mercenárius præstolátur finem óperis sui, sic et ego hábui menses vácuos et noctes laboriósas enumerávi mihi. Si dormíero, dicam: Quando consúrgam? Et rursum exspectábo vésperam, et replébor dolóribus usque ad ténebras.\par
\switchcolumn
\EN
\noindent The life of man upon earth is a warfare, and his days are like the days of a hireling. As a servant longeth for the shade, as the hireling looketh for the end of his work; So I also have had empty months, and have numbered to myself wearisome nights. If I lie down to sleep, I shall say: When shall I arise? and again I shall look for the evening, and shall be filled with sorrows even till darkness.\par
\switchcolumn*

\LA
\begin{Resp}\Rsign Quis mihi tríbuat, ut in inférno prótegas me et abscóndas me, donec pertránseat furor tuus, Dómine, nisi tu, qui solus es Deus? \Star{} Et constítuas mihi tempus, in quo recordéris mei?\end{Resp}
\begin{Resp}\Vsign Numquid sicut dies hóminis dies tui, ut quæras iniquitátem meam; cum sit nemo, qui de manu tua possit erúere.\end{Resp}
\begin{Resp}\Rsign Et constítuas mihi tempus, in quo recordéris mei?\end{Resp}
\switchcolumn
\EN
\begin{Resp}\Rsign O that Thou wouldest hide me in the grave; that Thou wouldest keep me secret, until thy wrath be past, even thine, O Lord, Thou That alone art God; \Star{} That Thou wouldest appoint me a set time, and remember me!\end{Resp}
\begin{Resp}\Vsign Are thy days as the days of man, that Thou inquirest after mine iniquity, and there is none that can deliver out of thine hand.\end{Resp}
\begin{Resp}\Rsign That Thou wouldest appoint me a set time, and remember me!\end{Resp}
\switchcolumn*

\LA
\LessonHead{Lectio II}
\switchcolumn
\EN
\LessonHead{Lesson II}
\switchcolumn*

\LA
\Cite{Job 7:5-8}
\switchcolumn
\EN
\Cite{Job 7:5-8}
\switchcolumn*

\LA
\noindent Indúta est caro mea putrédine, et sórdibus púlveris cutis mea áruit et contrácta est. Dies mei velócius transiérunt quam a texénte tela succíditur, et consúmpti sunt absque ulla spe. Meménto quia ventus est vita mea, et non revertétur óculus meus ut vídeat bona. Nec aspíciet me visus hóminis; óculi tui in me, et non subsístam.\par
\switchcolumn
\EN
\noindent My flesh is clothed with rottenness and the filth of dust, my skin is withered and drawn together. My days have passed more swiftly than the web is cut by the weaver, and are consumed without any hope. Remember that my life is but wind, and my eyes shall not return to see good things. Nor shall the sight of man behold me: thy eyes are upon me, and I shall be no more.\par
\switchcolumn*

\LA
\begin{Resp}\Rsign Utinam appenderéntur peccáta mea, quibus iram mérui, \Star{} Et calámitas, quam pátior, in statéra.\end{Resp}
\begin{Resp}\Vsign Quasi aréna maris hæc grávior apparéret, unde et verba mea dolóre sunt plena.\end{Resp}
\begin{Resp}\Rsign Et calámitas, quam pátior, in statéra.\end{Resp}
\begin{Resp}\Vsign Glória Patri, et Fílio, \Star{} et Spirítui Sancto.\end{Resp}
\begin{Resp}\Rsign Et calámitas, quam pátior, in statéra.\end{Resp}
\switchcolumn
\EN
\begin{Resp}\Rsign O that my sins, whereby I have deserved wrath; \Star{} And the calamity whereunder I suffer, were laid in the balances together.\end{Resp}
\begin{Resp}\Vsign For now it would appear heavier than the sand of the sea, therefore also my words are full of sorrow.\end{Resp}
\begin{Resp}\Rsign And the calamity, whereunder I suffer, were laid in the balances together.\end{Resp}
\begin{Resp}\Vsign Glory be to the Father, and to the Son, \Star{} and to the Holy Ghost.\end{Resp}
\begin{Resp}\Rsign And the calamity, whereunder I suffer, were laid in the balances together.\end{Resp}
\switchcolumn*

\LA
\LessonHead{Lectio III}
\switchcolumn
\EN
\LessonHead{Lesson III}
\switchcolumn*

\LA
\Cite{Job 7:9-12}
\switchcolumn
\EN
\Cite{Job 7:9-12}
\switchcolumn*

\LA
\noindent Sicut consúmitur nubes, et pertránsit, sic qui descénderit ad ínferos non ascéndet; nec revertétur ultra in domum suam, neque cognóscet eum ámplius locus ejus. Quaprópter et ego non parcam ori meo; loquar in tribulatióne spíritus mei, confabulábor cum amaritúdine ánimæ meæ. Numquid mare ego sum aut cetus, quia circumdedísti me cárcere?\par
\switchcolumn
\EN
\noindent As a cloud is consumed, and passeth away: so he that shall go down to hell shall not come up. Nor shall he return any more into his house, neither shall his place know him any more. Wherefore I will not spare my month, I will speak in the affliction of my spirit: I will talk with the bitterness of my soul. Am I a sea, or a whale, that thou hast enclosed me in a prison?\par
\switchcolumn*

\LA
\begin{Resp}\Rsign Quare detraxístis sermónibus veritátis? ad increpándum verba compónitis et subvértere nitímini amícum vestrum: \Star{} Verúmtamen quæ cogitástis, expléte.\end{Resp}
\begin{Resp}\Vsign Quod justum est, judicáte; et non inveniétis in lingua mea iniquitátem.\end{Resp}
\begin{Resp}\Rsign Verúmtamen quæ cogitástis, expléte.\end{Resp}
\begin{Resp}\Vsign Glória Patri, et Fílio, \Star{} et Spirítui Sancto.\end{Resp}
\begin{Resp}\Rsign Verúmtamen quæ cogitástis, expléte.\end{Resp}
\switchcolumn
\EN
\begin{Resp}\Rsign Why do ye argue against the words of truth? Do ye imagine words to reprove me and strive to confound one that is your friend? \Star{} Nevertheless, finish that ye have in mind.\end{Resp}
\begin{Resp}\Vsign Judge that which is just, and ye shall find no iniquity in my tongue.\end{Resp}
\begin{Resp}\Rsign Nevertheless, finish that ye have in mind.\end{Resp}
\begin{Resp}\Vsign Glory be to the Father, and to the Son, \Star{} and to the Holy Ghost.\end{Resp}
\begin{Resp}\Rsign Nevertheless, finish that ye have in mind.\end{Resp}
\switchcolumn*

\end{paracol}

\Office{Dominica II Septembris}{Second Sunday of September}{}
\addcontentsline{toc}{subsection}{Dominica II Septembris \textit{— Second Sunday of September}}
\begin{paracol}{2}
\LA
\LessonHead{Lectio I}
\switchcolumn
\EN
\LessonHead{Lesson I}
\switchcolumn*

\LA
\LessonTitle{De libro Job}
\switchcolumn
\EN
\LessonTitle{Lesson from the book of Job}
\switchcolumn*

\LA
\Cite{Job 9:1-5}
\switchcolumn
\EN
\Cite{Job 9:1-5}
\switchcolumn*

\LA
\noindent Et respóndens Job, ait: Vere scio quod ita sit, et quod non justificétur homo compósitus Deo. Si volúerit conténdere cum eo, non póterit ei respondére unum pro mille. Sápiens corde est et fortis róbore. Quis réstitit ei, et pacem hábuit? Qui tránstulit montes, et nesciérunt hi quos subvértit in furóre suo.\par
\switchcolumn
\EN
\noindent And Job answered, and said: Indeed I know it is so, and that man cannot be justified compared with God, If he will contend with him, he cannot answer him one for a thousand. He is wise in heart, and mighty in strength: who hath resisted him, and hath had peace Who hath removed mountains, and they whom he overthrew in his wrath, knew it not.\par
\switchcolumn*

\LA
\begin{Resp}\Rsign Si bona suscépimus de manu Dei, mala autem quare non sustineámus? \Star{} Dóminus dedit, Dóminus ábstulit; sicut Dómino plácuit, ita factum est: sit nomen Dómini benedíctum.\end{Resp}
\begin{Resp}\Vsign Nudus egréssus sum de útero matris meæ et nudus revértar illuc.\end{Resp}
\begin{Resp}\Rsign Dóminus dedit, Dóminus ábstulit; sicut Dómino plácuit, ita factum est: sit nomen Dómini benedíctum.\end{Resp}
\switchcolumn
\EN
\begin{Resp}\Rsign What shall we receive good at the hand of God, and shall we not receive evil? \Star{} The Lord gave and the Lord hath taken away. As the Lord hath pleased, so hath it befallen. Blessed be the Name of the Lord.\end{Resp}
\begin{Resp}\Vsign Naked came I out of my mother's womb, and naked shall I return thither.\end{Resp}
\begin{Resp}\Rsign The Lord gave and the Lord hath taken away. As the Lord hath pleased, so hath it befallen. Blessed be the Name of the Lord.\end{Resp}
\switchcolumn*

\LA
\LessonHead{Lectio II}
\switchcolumn
\EN
\LessonHead{Lesson II}
\switchcolumn*

\LA
\Cite{Job 9:6-10}
\switchcolumn
\EN
\Cite{Job 9:6-10}
\switchcolumn*

\LA
\noindent Qui cómmovet terram de loco suo, et colúmnæ ejus concutiúntur; qui prǽcipit soli, et non óritur, et stellas claudit quasi sub signáculo; qui exténdit cælos suos, et gráditur super fluctus maris; qui facit Arctúrum et Oríona et Hýadas et interióra Austri; qui facit magna et incomprehensibília et mirabília, quorum non est númerus.\par
\switchcolumn
\EN
\noindent Who shaketh the earth out of her place, and the pillars thereof tremble. Who commandeth the sun and it riseth not: and shutteth up the stars as it were under a seal: Who alone spreadeth out the heavens, and walketh upon the waves of the sea. Who maketh Arcturus, and Orion, and Hyades, and the inner parts of the south. Who doth things great and incomprehensible, and wonderful, of which there is no number.\par
\switchcolumn*

\LA
\begin{Resp}\Rsign Antequam cómedam, suspíro, et tamquam inundántes aquæ sic rugítus meus; quia timor, quem timébam, evénit mihi, et quod verébar áccidit. Nonne dissimulávi? nonne sílui? nonne quiévi? \Star{} Et venit super me indignátio.\end{Resp}
\begin{Resp}\Vsign Ecce non est auxílium mihi in me, et necessárii quoque mei recessérunt a me.\end{Resp}
\begin{Resp}\Rsign Et venit super me indignátio.\end{Resp}
\switchcolumn
\EN
\begin{Resp}\Rsign My sighing cometh before I eat, and my roarings are poured out like the waters, for the thing which I greatly feared is come upon me, and that which I was afraid of is come unto me. Was not I silent? Held not I my peace? Was not I at rest? \Star{} And trouble came.\end{Resp}
\begin{Resp}\Vsign Behold, I cannot help myself, and they that were needful unto me have forsaken me.\end{Resp}
\begin{Resp}\Rsign And trouble came.\end{Resp}
\switchcolumn*

\LA
\LessonHead{Lectio III}
\switchcolumn
\EN
\LessonHead{Lesson III}
\switchcolumn*

\LA
\Cite{Job 9:11-17}
\switchcolumn
\EN
\Cite{Job 9:11-17}
\switchcolumn*

\LA
\noindent Si vénerit ad me, non vidébo eum; si abíerit, non intéllegam; si repénte intérroget, quis respondébit ei vel quis dícere potest: Cur ita facis? Deus, cujus iræ nemo resístere potest, et sub quo curvántur qui portant orbem; quantus ergo sum ego, ut respóndeam ei, et loquar verbis meis cum eo? Qui, étiam si habúero quíppiam justum, non respondébo, sed meum júdicem deprecábor. Et, cum invocántem exaudíerit me, non credo quod audíerit vocem meam. In túrbine enim cónterit me, et multiplicábit vúlnera mea, étiam sine causa.\par
\switchcolumn
\EN
\noindent If he come to me, I shall not see him: if he depart I shall not understand. If he examine on a sudden, who shall answer him? or who can say: Why dost thou so? God, whose wrath no man can resist, and under whom they stoop that bear up the world. What am I then, that I should answer him, and have words with him? I, who although I should have any just thing, would not answer, but would make supplication to my judge. And if he should hear me when I call, I should not believe that he had heard my voice. For he shall crush me in a whirlwind, and multiply my wounds even without cause.\par
\switchcolumn*

\LA
\begin{Resp}\Rsign Quare detraxístis sermónibus veritátis? ad increpándum verba compónitis et subvértere nitímini amícum vestrum: \Star{} Verúmtamen quæ cogitástis, expléte.\end{Resp}
\begin{Resp}\Vsign Quod justum est, judicáte; et non inveniétis in lingua mea iniquitátem.\end{Resp}
\begin{Resp}\Rsign Verúmtamen quæ cogitástis, expléte.\end{Resp}
\begin{Resp}\Vsign Glória Patri, et Fílio, \Star{} et Spirítui Sancto.\end{Resp}
\begin{Resp}\Rsign Verúmtamen quæ cogitástis, expléte.\end{Resp}
\switchcolumn
\EN
\begin{Resp}\Rsign Why do ye argue against the words of truth? Do ye imagine words to reprove me and strive to confound one that is your friend? \Star{} Nevertheless, finish that ye have in mind.\end{Resp}
\begin{Resp}\Vsign Judge that which is just, and ye shall find no iniquity in my tongue.\end{Resp}
\begin{Resp}\Rsign Nevertheless, finish that ye have in mind.\end{Resp}
\begin{Resp}\Vsign Glory be to the Father, and to the Son, \Star{} and to the Holy Ghost.\end{Resp}
\begin{Resp}\Rsign Nevertheless, finish that ye have in mind.\end{Resp}
\switchcolumn*

\LA
\LessonHead{Lectio IV}
\switchcolumn
\EN
\LessonHead{Lesson IV}
\switchcolumn*

\LA
\LessonTitle{Ex libro Morálium sancti Gregórii Papæ}
\switchcolumn
\EN
\LessonTitle{From the Book of Moral (Reflections upon Job) written by Pope St. Gregory (the Great.)}
\switchcolumn*

\LA
\Source{Liber 9, Cap. 2.}
\switchcolumn
\EN
\Source{Bk. ix. ch. 2}
\switchcolumn*

\LA
\noindent Vero scio quod ita sit, et quod non justificábitur homo compósitus Deo. Homo quippe Deo non compósitus justítiam pércipit, compósitus amíttit; quia quisquis se auctóri bonórum cómparat, bono se, quod accéperat, privat. Qui enim accépta bona sibi árrogat, suis contra Deum donis pugnat. Unde ergo despéctus erígitur, dignum est ut eréctus inde destruátur. Sanctus autem vir, quia omne virtútis nostræ méritum esse vítium cónspicit, si ab intérno árbitro distrícte judicétur, recte subjúngit: Si volúerit conténdere cum eo, non póterit respondére ei unum pro mille.\par
\switchcolumn
\EN
\noindent We know that it is so of a truth, and that a man cannot be justified as against God. When God is put out of the consideration, a man may be considered to be just, but considered as against God, his righteousness vanisheth away. When a man measureth himself by his relation to Him, Who is the Author of all good, he doth thereby acknowledge that of himself he hath no good in him, but hath received from God whatsoever he hath. He that glorifieth himself because of good which hath been given him, fighteth against God with God's own gifts. It is just therefore that the grounds upon which he ought to have been humbled, but upon which he hath puffed himself up, should be used to humble his vain-glory. But an holy man, because he perceiveth that the worth of our own good deeds falleth short, when he considereth his own spiritual man, justly saith If He will contend with him, he cannot answer Him one of a thousand.\par
\switchcolumn*

\LA
\begin{Resp}\Rsign Indúta est caro mea putrédine, et sórdibus púlveris cutis mea áruit et contrácta est: \Star{} Meménto mei, Dómine, quóniam ventus est vita mea.\end{Resp}
\begin{Resp}\Vsign Dies mei velócius transiérunt quam a texénte tela succíditur, et consúmpti sunt absque ulla spe.\end{Resp}
\begin{Resp}\Rsign Meménto mei, Dómine, quóniam ventus est vita mea.\end{Resp}
\switchcolumn
\EN
\begin{Resp}\Rsign My flesh is clothed with worms and clods of dust. My skin is dry and drawn together. \Star{} Remember me, O Lord, for my life is wind.\end{Resp}
\begin{Resp}\Vsign My days are swifter than a weaver's shuttle, and are spent without hope.\end{Resp}
\begin{Resp}\Rsign Remember me, Lord, for my life is wind.\end{Resp}
\switchcolumn*

\LA
\LessonHead{Lectio V}
\switchcolumn
\EN
\LessonHead{Lesson V}
\switchcolumn*

\LA
\noindent In Scriptúra sancta millenárius númerus pro universitáte solet intélligi. Hinc étenim Psalmísta ait: Verbi, quod mandávit in mille generatiónes: cum profécto constet, quod ab ipso mundi exórdio usque ad Redemptiónis advéntum per Evangelístam non ámplius quam septuagínta et septem propágines numeréntur. Quid ergo in millenário número nisi ad proferéndam novam sóbolem perfécta univérsitas præscítæ generatiónis exprímitur? Hinc et per Joánnem dícitur: Et regnábunt cum eo mille annis: quia vidélicet regnum sanctæ Ecclésiæ universitátis perfectióne solidátur.\par
\switchcolumn
\EN
\noindent In the Holy Scriptures the numeral a thousand is used to be taken as signifying a generalization. Thus, the Psalmist saith The word which He commanded to a thousand generations (Ps. civ. 8), whereas it is notorious that the Evangelist doth not reckon more then seventy-and-seven generations between the very beginning of the world and the coming of our Redeemer. What therefore is to be understood here by a thousand The general ripeness of the old generation to bring forth a new offspring. Hence also it is said by John And shall reign with Him a thousand years (Apoc. xx. 6,) because the reign of the Holy Church will be over all mankind made perfect.\par
\switchcolumn*

\LA
\begin{Resp}\Rsign Páucitas diérum meórum finiétur brevi; dimítte me, Dómine, ut plangam páululum dolórem meum, \Star{} Antequam vadam ad terram tenebrósam et opértam mortis calígine.\end{Resp}
\begin{Resp}\Vsign Manus tuæ, Dómine, fecérunt me, et plasmavérunt me totum in circúitu; et sic repénte præcípitas me?\end{Resp}
\begin{Resp}\Rsign Antequam vadam ad terram tenebrósam et opértam mortis calígine.\end{Resp}
\switchcolumn
\EN
\begin{Resp}\Rsign My days are few, and in a short while they will be ended. Let me alone, then, O Lord, that I may bewail my sorrow a little; \Star{} Before I go to the land of darkness and of the shadow of death.\end{Resp}
\begin{Resp}\Vsign Thine hands, O Lord, have made me, and fashioned me together round about, and yet dost Thou forthwith destroy me.\end{Resp}
\begin{Resp}\Rsign Before I go to the land of darkness and of the shadow of death.\end{Resp}
\switchcolumn*

\LA
\LessonHead{Lectio VI}
\switchcolumn
\EN
\LessonHead{Lesson VI}
\switchcolumn*

\LA
\noindent Quia vero monas décies multiplicáta in denárium dúcitur, denárius per semetípsum ductus in centenárium dilatátur, qui rursus per denárium ductus in millenárium ténditur; cum ab uno incípimus, ut ad millenárium veniámus, quid hoc loco uníus appellatióne, nisi bene vivéndi inítium? quid millenárii númeri amplitúdine, nisi ejúsdem bonæ vitæ perféctio designátur? Cum Deo autem conténdere, est, non ei tribúere, sed sibi glóriam suæ virtútis arrogáre. Sed sanctus vir conspíciat, quia et qui summa jam dona percépit, si de accéptis extóllitur, cuncta quæ accéperat, amíttit.\par
\switchcolumn
\EN
\noindent When times one is ten, and ten times ten is an hundred, and ten times an hundred is a thousand. Observing therefore this connection between one and a thousand, what are we to understand by the one (in the text, connected as it is with the thousand whereby we understand perfection)? Is it not the beginning of a good life, even as the thousand representeth perfection? The contending with God (which is spoken of in the text) is the non-acknowledgment of that which is owed to Him, and the vain-glorying instead in our own strength. But an holy man should see, that even if one had received the gifts of perfection, and were to make them the grounds of self-glorifying, such an one would thereby lose all that he had received.\par
\switchcolumn*

\LA
\begin{Resp}\Rsign Non abscóndas me, Dómine, a fácie tua: manum tuam longe fac a me, \Star{} Et formído tua non me térreat.\end{Resp}
\begin{Resp}\Vsign Córripe me, Dómine, in misericórdia, non in furóre tuo, ne forte ad níhilum rédigas me.\end{Resp}
\begin{Resp}\Rsign Et formído tua non me térreat.\end{Resp}
\begin{Resp}\Vsign Glória Patri, et Fílio, \Star{} et Spirítui Sancto.\end{Resp}
\begin{Resp}\Rsign Et formído tua non me térreat.\end{Resp}
\switchcolumn
\EN
\begin{Resp}\Rsign Hide not thy face from me, O Lord. Withdraw not thine hand far from me, \Star{} And let not thy dread make me afraid.\end{Resp}
\begin{Resp}\Vsign O Lord, correct me but in mercy not in thine anger, lest Thou bring me to nothing.\end{Resp}
\begin{Resp}\Rsign And let not thy dread make me afraid.\end{Resp}
\begin{Resp}\Vsign Glory be to the Father, and to the Son, \Star{} and to the Holy Ghost.\end{Resp}
\begin{Resp}\Rsign And let not thy dread make me afraid.\end{Resp}
\switchcolumn*

\end{paracol}

\Office{Feria II infra Hebdomadam II Septembris}{Monday of the second week of September}{}
\addcontentsline{toc}{subsection}{Feria II infra Hebdomadam II Septembris \textit{— Monday of the second week of September}}
\begin{paracol}{2}
\LA
\LessonHead{Lectio I}
\switchcolumn
\EN
\LessonHead{Lesson I}
\switchcolumn*

\LA
\LessonTitle{De libro Job}
\switchcolumn
\EN
\LessonTitle{Lesson from the book of Job}
\switchcolumn*

\LA
\Cite{Job 27:1-5}
\switchcolumn
\EN
\Cite{Job 27:1-5}
\switchcolumn*

\LA
\noindent Addidit quoque Job assúmens parábolam suam, et dixit: Vivit Deus, qui ábstulit judícium meum, et Omnípotens, qui ad amaritúdinem addúxit ánimam meam. Quia, donec súperest hálitus in me, et spíritus Dei in náribus meis, non loquéntur lábia mea iniquitátem, nec lingua mea meditábitur mendácium. Absit a me, ut justos vos esse júdicem; donec defíciam, non recédam ab innocéntia mea.\par
\switchcolumn
\EN
\noindent Job also added, taking up his parable, and said: As God liveth, who hath taken away my judgment, and the Almighty, who hath brought my soul to bitterness, As long as breath remaineth in me, and the spirit of God in my nostrils, My lips shall not speak iniquity, neither shall my tongue contrive lying. God forbid that I should judge you to be just: till I die I will not depart from my innocence.\par
\switchcolumn*

\LA
\begin{Resp}\Rsign Versa est in luctum cíthara mea, et órganum meum in vocem fléntium: \Star{} Parce mihi, Dómine, nihil enim sunt dies mei.\end{Resp}
\begin{Resp}\Vsign Cutis mea denigráta est super me, et ossa mea aruérunt.\end{Resp}
\begin{Resp}\Rsign Parce mihi, Dómine, nihil enim sunt dies mei.\end{Resp}
\switchcolumn
\EN
\begin{Resp}\Rsign My harp is turned to mourning, and my organ into the voice of them that weep. \Star{} Let me alone, O Lord, for my days are vanity.\end{Resp}
\begin{Resp}\Vsign My skin is black upon me, and my bones are burned with heat.\end{Resp}
\begin{Resp}\Rsign Let me alone, O Lord, for my days are vanity.\end{Resp}
\switchcolumn*

\LA
\LessonHead{Lectio II}
\switchcolumn
\EN
\LessonHead{Lesson II}
\switchcolumn*

\LA
\Cite{Job 27:6-10}
\switchcolumn
\EN
\Cite{Job 27:6-10}
\switchcolumn*

\LA
\noindent Justificatiónem meam, quam cœpi tenére, non déseram; neque enim reprehéndit me cor meum in omni vita mea. Sit ut ímpius inimícus meus, et adversárius meus quasi iníquus. Quæ est enim spes hypócritæ, si aváre rápiat, et non líberet Deus ánimam ejus? Numquid Deus áudiet clamórem ejus, cum vénerit super eum angústia? Aut póterit in Omnipoténte delectári, et invocáre Deum omni témpore?\par
\switchcolumn
\EN
\noindent My justification, which I have begun to hold, I will not forsake: for my heart doth not reprehend me in all my life. Let my enemy be as the ungodly, and my adversary as the wicked one. For what is the hope of the hypocrite if through covetousness he take by violence, and God deliver not his soul? Will God hear his cry, when distress shall come upon him? Or can he delight himself in the Almighty, and call upon God at all times?\par
\switchcolumn*

\LA
\begin{Resp}\Rsign Utinam appenderéntur peccáta mea, quibus iram mérui, \Star{} Et calámitas, quam pátior, in statéra.\end{Resp}
\begin{Resp}\Vsign Quasi aréna maris hæc grávior apparéret, unde et verba mea dolóre sunt plena.\end{Resp}
\begin{Resp}\Rsign Et calámitas, quam pátior, in statéra.\end{Resp}
\switchcolumn
\EN
\begin{Resp}\Rsign O that my sins, whereby I have deserved wrath; \Star{} And the calamity whereunder I suffer, were laid in the balances together.\end{Resp}
\begin{Resp}\Vsign For now it would appear heavier than the sand of the sea, therefore also my words are full of sorrow.\end{Resp}
\begin{Resp}\Rsign And the calamity, whereunder I suffer, were laid in the balances together.\end{Resp}
\switchcolumn*

\LA
\LessonHead{Lectio III}
\switchcolumn
\EN
\LessonHead{Lesson III}
\switchcolumn*

\LA
\Cite{Job 27:11-15}
\switchcolumn
\EN
\Cite{Job 27:11-15}
\switchcolumn*

\LA
\noindent Docébo vos per manum Dei quæ Omnípotens hábeat, nec abscóndam. Ecce vos omnes nostis, et quid sine causa vana loquímini? Hæc est pars hóminis ímpii apud Deum, et heréditas violentórum, quam ab Omnipoténte suscípient: si multiplicáti fúerint fílii ejus, in gládio erunt, et nepótes ejus non saturabúntur pane; qui réliqui fúerint ex eo sepeliéntur in intéritu, et víduæ illíus non plorábunt.\par
\switchcolumn
\EN
\noindent I will teach you by the hand of God, what the Almighty hath, and I will not conceal it. Behold you all know it, and why do you speak vain things without cause? This is the portion of a wicked man with God, and the inheritance of the violent, which they shall receive of the Almighty. If his sons be multiplied, they shall be for the sword, and his grandsons shall not be filled with bread. They that shall remain of him, shall be buried in death, and his widows shall not weep.\par
\switchcolumn*

\LA
\begin{Resp}\Rsign Quare detraxístis sermónibus veritátis? ad increpándum verba compónitis et subvértere nitímini amícum vestrum: \Star{} Verúmtamen quæ cogitástis, expléte.\end{Resp}
\begin{Resp}\Vsign Quod justum est, judicáte; et non inveniétis in lingua mea iniquitátem.\end{Resp}
\begin{Resp}\Rsign Verúmtamen quæ cogitástis, expléte.\end{Resp}
\begin{Resp}\Vsign Glória Patri, et Fílio, \Star{} et Spirítui Sancto.\end{Resp}
\begin{Resp}\Rsign Verúmtamen quæ cogitástis, expléte.\end{Resp}
\switchcolumn
\EN
\begin{Resp}\Rsign Why do ye argue against the words of truth? Do ye imagine words to reprove me and strive to confound one that is your friend? \Star{} Nevertheless, finish that ye have in mind.\end{Resp}
\begin{Resp}\Vsign Judge that which is just, and ye shall find no iniquity in my tongue.\end{Resp}
\begin{Resp}\Rsign Nevertheless, finish that ye have in mind.\end{Resp}
\begin{Resp}\Vsign Glory be to the Father, and to the Son, \Star{} and to the Holy Ghost.\end{Resp}
\begin{Resp}\Rsign Nevertheless, finish that ye have in mind.\end{Resp}
\switchcolumn*

\end{paracol}

\Office{Feria III infra Hebdomadam II Septembris}{Tuesday of the second week of September}{}
\addcontentsline{toc}{subsection}{Feria III infra Hebdomadam II Septembris \textit{— Tuesday of the second week of September}}
\begin{paracol}{2}
\LA
\LessonHead{Lectio I}
\switchcolumn
\EN
\LessonHead{Lesson I}
\switchcolumn*

\LA
\LessonTitle{De libro Job}
\switchcolumn
\EN
\LessonTitle{Lesson from the book of Job}
\switchcolumn*

\LA
\Cite{Job 28:12-16}
\switchcolumn
\EN
\Cite{Job 28:12-16}
\switchcolumn*

\LA
\noindent Sapiéntia ubi invenítur? et quis est locus intellegéntiæ? Nescit homo prétium ejus, nec invenítur in terra suáviter vivéntium. Abýssus dicit: Non est in me, et mare lóquitur: Non est mecum. Non dábitur aurum obrýzum pro ea, nec appendétur argéntum in commutatióne ejus; non conferétur tinctis Indiæ colóribus, nec lápidi sardónycho pretiosíssimo vel sapphíro.\par
\switchcolumn
\EN
\noindent But where is wisdom to be found, and where is the place of understanding? Man knoweth not the price thereof, neither is it found in the land of them that live in delights. The depth saith: It is not in me: and the sea saith: It is not with me. The finest gold shall not purchase it, neither shall silver be weighed in exchange for it. It shall not be compared with the dyed colours of India, or with the most precious stone sardonyx, or the sapphire.\par
\switchcolumn*

\LA
\begin{Resp}\Rsign Indúta est caro mea putrédine, et sórdibus púlveris cutis mea áruit et contrácta est: \Star{} Meménto mei, Dómine, quóniam ventus est vita mea.\end{Resp}
\begin{Resp}\Vsign Dies mei velócius transiérunt quam a texénte tela succíditur, et consúmpti sunt absque ulla spe.\end{Resp}
\begin{Resp}\Rsign Meménto mei, Dómine, quóniam ventus est vita mea.\end{Resp}
\switchcolumn
\EN
\begin{Resp}\Rsign My flesh is clothed with worms and clods of dust. My skin is dry and drawn together. \Star{} Remember me, O Lord, for my life is wind.\end{Resp}
\begin{Resp}\Vsign My days are swifter than a weaver's shuttle, and are spent without hope.\end{Resp}
\begin{Resp}\Rsign Remember me, Lord, for my life is wind.\end{Resp}
\switchcolumn*

\LA
\LessonHead{Lectio II}
\switchcolumn
\EN
\LessonHead{Lesson II}
\switchcolumn*

\LA
\Cite{Job 28:17-22}
\switchcolumn
\EN
\Cite{Job 28:17-22}
\switchcolumn*

\LA
\noindent Non adæquábitur ei aurum vel vitrum, nec commutabúntur pro ea vasa auri. Excélsa et eminéntia non memorabúntur comparatióne ejus; tráhitur autem sapiéntia de occúltis. Non adæquábitur ei topázius de Æthiópia, nec tinctúræ mundíssimæ componétur. Unde ergo sapiéntia venit? et quis est locus intellegéntiæ? Abscóndita est ab óculis ómnium vivéntium; vólucres quoque cæli latet. Perdítio et mors dixérunt: Auribus nostris audívimus famam ejus.\par
\switchcolumn
\EN
\noindent Gold or crystal cannot equal it, neither shall any vessels of gold be changed for it. High and eminent things shall not be mentioned in comparison of it: but wisdom is drawn out of secret places. The topaz of Ethiopia shall not be equal to it, neither shall it be compared to the cleanest dyeing. Whence then cometh wisdom? and where is the place of understanding? It is hid from the eyes of all living. and the fowls of the air know it not. Destruction and death have said: With our ears we have heard the fame thereof.\par
\switchcolumn*

\LA
\begin{Resp}\Rsign Páucitas diérum meórum finiétur brevi; dimítte me, Dómine, ut plangam páululum dolórem meum, \Star{} Antequam vadam ad terram tenebrósam et opértam mortis calígine.\end{Resp}
\begin{Resp}\Vsign Manus tuæ, Dómine, fecérunt me, et plasmavérunt me totum in circúitu; et sic repénte præcípitas me?\end{Resp}
\begin{Resp}\Rsign Antequam vadam ad terram tenebrósam et opértam mortis calígine.\end{Resp}
\switchcolumn
\EN
\begin{Resp}\Rsign My days are few, and in a short while they will be ended. Let me alone, then, O Lord, that I may bewail my sorrow a little; \Star{} Before I go to the land of darkness and of the shadow of death.\end{Resp}
\begin{Resp}\Vsign Thine hands, O Lord, have made me, and fashioned me together round about, and yet dost Thou forthwith destroy me.\end{Resp}
\begin{Resp}\Rsign Before I go to the land of darkness and of the shadow of death.\end{Resp}
\switchcolumn*

\LA
\LessonHead{Lectio III}
\switchcolumn
\EN
\LessonHead{Lesson III}
\switchcolumn*

\LA
\Cite{Job 28:23-28}
\switchcolumn
\EN
\Cite{Job 28:23-28}
\switchcolumn*

\LA
\noindent Deus intéllegit viam ejus, et ipse novit locum illíus; ipse enim fines mundi intuétur, et ómnia, quæ sub cælo sunt, réspicit. Qui fecit ventis pondus, et aquas appéndit in mensúra. Quando ponébat plúviis legem et viam procéllis sonántibus, tunc vidit illam et enarrávit et præparávit et investigávit. Et dixit hómini: Ecce timor Dómini, ipsa est sapiéntia, et recédere a malo, intellegéntia.\par
\switchcolumn
\EN
\noindent God understandeth the way of it, and he knoweth the place thereof. For he beholdeth the ends of the world: and looketh on all things that are under heaven. Who made a weight for the winds and weighed the waters by measure. When he gave a law for the rain, and a way for the sounding storms. Then he saw it, and declared, and prepared, and searched it. And he said to man: Behold the fear of the Lord, that is wisdom: and to depart from evil, is understanding.\par
\switchcolumn*

\LA
\begin{Resp}\Rsign Non abscóndas me, Dómine, a fácie tua: manum tuam longe fac a me, \Star{} Et formído tua non me térreat.\end{Resp}
\begin{Resp}\Vsign Córripe me, Dómine, in misericórdia, non in furóre tuo, ne forte ad níhilum rédigas me.\end{Resp}
\begin{Resp}\Rsign Et formído tua non me térreat.\end{Resp}
\begin{Resp}\Vsign Glória Patri, et Fílio, \Star{} et Spirítui Sancto.\end{Resp}
\begin{Resp}\Rsign Et formído tua non me térreat.\end{Resp}
\switchcolumn
\EN
\begin{Resp}\Rsign Hide not thy face from me, O Lord. Withdraw not thine hand far from me, \Star{} And let not thy dread make me afraid.\end{Resp}
\begin{Resp}\Vsign O Lord, correct me but in mercy not in thine anger, lest Thou bring me to nothing.\end{Resp}
\begin{Resp}\Rsign And let not thy dread make me afraid.\end{Resp}
\begin{Resp}\Vsign Glory be to the Father, and to the Son, \Star{} and to the Holy Ghost.\end{Resp}
\begin{Resp}\Rsign And let not thy dread make me afraid.\end{Resp}
\switchcolumn*

\end{paracol}

\Office{Feria IV infra Hebdomadam II Septembris}{Wednesday of the second week of September}{}
\addcontentsline{toc}{subsection}{Feria IV infra Hebdomadam II Septembris \textit{— Wednesday of the second week of September}}
\begin{paracol}{2}
\LA
\LessonHead{Lectio I}
\switchcolumn
\EN
\LessonHead{Lesson I}
\switchcolumn*

\LA
\LessonTitle{De libro Job}
\switchcolumn
\EN
\LessonTitle{Lesson from the book of Job}
\switchcolumn*

\LA
\Cite{Job 31:1-6}
\switchcolumn
\EN
\Cite{Job 31:1-6}
\switchcolumn*

\LA
\noindent Pépigi fœdus cum óculis meis, ut ne cogitárem quidem de vírgine. Quam enim partem habéret in me Deus désuper, et hereditátem Omnípotens de excélsis? Numquid non perdítio est iníquo, et alienátio operántibus injustítiam? Nonne ipse consíderat vias meas, et cunctos gressus meos dinúmerat? Si ambulávi in vanitáte, et festinávit in dolo pes meus, appéndat me in statéra justa, et sciat Deus simplicitátem meam.\par
\switchcolumn
\EN
\noindent I made a covenant with my eyes, that I would not so much as think upon a virgin. For what part should God from above have in me, and what inheritance the Almighty from on high? Is not destruction to the wicked, and aversion to them that work iniquity? Doth not he consider my ways, and number all my steps? If I have walked in vanity, and my foot hath made haste to deceit: Let him weigh me in a just balance, and let God know my simplicity.\par
\switchcolumn*

\LA
\begin{Resp}\Rsign Quis mihi tríbuat, ut in inférno prótegas me et abscóndas me, donec pertránseat furor tuus, Dómine, nisi tu, qui solus es Deus? \Star{} Et constítuas mihi tempus, in quo recordéris mei?\end{Resp}
\begin{Resp}\Vsign Numquid sicut dies hóminis dies tui, ut quæras iniquitátem meam; cum sit nemo, qui de manu tua possit erúere.\end{Resp}
\begin{Resp}\Rsign Et constítuas mihi tempus, in quo recordéris mei?\end{Resp}
\switchcolumn
\EN
\begin{Resp}\Rsign O that Thou wouldest hide me in the grave; that Thou wouldest keep me secret, until thy wrath be past, even thine, O Lord, Thou That alone art God; \Star{} That Thou wouldest appoint me a set time, and remember me!\end{Resp}
\begin{Resp}\Vsign Are thy days as the days of man, that Thou inquirest after mine iniquity, and there is none that can deliver out of thine hand.\end{Resp}
\begin{Resp}\Rsign That Thou wouldest appoint me a set time, and remember me!\end{Resp}
\switchcolumn*

\LA
\LessonHead{Lectio II}
\switchcolumn
\EN
\LessonHead{Lesson II}
\switchcolumn*

\LA
\Cite{Job 31:7-12}
\switchcolumn
\EN
\Cite{Job 31:7-12}
\switchcolumn*

\LA
\noindent Si declinávit gressus meus de via, et si secútum est óculos meos cor meum, et si mánibus meis adhǽsit mácula, seram, et álius cómedat, et progénies mea eradicétur. Si decéptum est cor meum super mulíere, et si ad óstium amíci mei insidiátus sum, scortum altérius sit uxor mea, et super illam incurvéntur álii. Hoc enim nefas est et iníquitas máxima; ignis est usque ad perditiónem dévorans et ómnia eradícans genímina.\par
\switchcolumn
\EN
\noindent If my step hath turned out of the way, and if my heart hath followed my eyes, and if a spot hath cleaved to my hands: Then let me sow and let another eat: and let my offspring be rooted out. If my heart hath been deceived upon a woman, and if I have laid wait at my friend's door: Let my wife be the harlot of another, and let other men lie with her. For this is a heinous crime, and a most grievous iniquity. It is a fire that devoureth even to destruction, and rooteth up all things that spring.\par
\switchcolumn*

\LA
\begin{Resp}\Rsign Utinam appenderéntur peccáta mea, quibus iram mérui, \Star{} Et calámitas, quam pátior, in statéra.\end{Resp}
\begin{Resp}\Vsign Quasi aréna maris hæc grávior apparéret, unde et verba mea dolóre sunt plena.\end{Resp}
\begin{Resp}\Rsign Et calámitas, quam pátior, in statéra.\end{Resp}
\switchcolumn
\EN
\begin{Resp}\Rsign O that my sins, whereby I have deserved wrath; \Star{} And the calamity whereunder I suffer, were laid in the balances together.\end{Resp}
\begin{Resp}\Vsign For now it would appear heavier than the sand of the sea, therefore also my words are full of sorrow.\end{Resp}
\begin{Resp}\Rsign And the calamity, whereunder I suffer, were laid in the balances together.\end{Resp}
\switchcolumn*

\LA
\LessonHead{Lectio III}
\switchcolumn
\EN
\LessonHead{Lesson III}
\switchcolumn*

\LA
\Cite{Job 31:13-18}
\switchcolumn
\EN
\Cite{Job 31:13-18}
\switchcolumn*

\LA
\noindent Si contémpsi subíre judícium cum servo meo et ancílla mea, cum disceptárent advérsum me: quid enim fáciam, cum surréxerit ad judicándum Deus? et, cum quæsíerit, quid respondébo illi? Numquid non in útero fecit me, qui et illum operátus est? et formávit me in vulva unus? Si negávi, quod volébant, paupéribus et óculos víduæ exspectáre feci; si comédi buccéllam meam solus, et non comédit pupíllus ex ea (quia ab infántia mea crevit mecum miserátio, et de útero matris meæ egréssa est mecum).\par
\switchcolumn
\EN
\noindent If I have despised to abide judgment with my manservant, or my maidservant, when they had any controversy against me: For what shall I do when God shall rise to judge? and when he shall examine, what shall I answer him? Did not he that made me in the womb make him also: and did not one and the same form me in the womb? If I have denied to the poor what they desired, and have made the eyes of the widow wait: If I have eaten my morsel alone, and the fatherless hath not eaten thereof: (For from my infancy mercy grew up with me: and it came out with me from my mother's womb)\par
\switchcolumn*

\LA
\begin{Resp}\Rsign Quare detraxístis sermónibus veritátis? ad increpándum verba compónitis et subvértere nitímini amícum vestrum: \Star{} Verúmtamen quæ cogitástis, expléte.\end{Resp}
\begin{Resp}\Vsign Quod justum est, judicáte; et non inveniétis in lingua mea iniquitátem.\end{Resp}
\begin{Resp}\Rsign Verúmtamen quæ cogitástis, expléte.\end{Resp}
\begin{Resp}\Vsign Glória Patri, et Fílio, \Star{} et Spirítui Sancto.\end{Resp}
\begin{Resp}\Rsign Verúmtamen quæ cogitástis, expléte.\end{Resp}
\switchcolumn
\EN
\begin{Resp}\Rsign Why do ye argue against the words of truth? Do ye imagine words to reprove me and strive to confound one that is your friend? \Star{} Nevertheless, finish that ye have in mind.\end{Resp}
\begin{Resp}\Vsign Judge that which is just, and ye shall find no iniquity in my tongue.\end{Resp}
\begin{Resp}\Rsign Nevertheless, finish that ye have in mind.\end{Resp}
\begin{Resp}\Vsign Glory be to the Father, and to the Son, \Star{} and to the Holy Ghost.\end{Resp}
\begin{Resp}\Rsign Nevertheless, finish that ye have in mind.\end{Resp}
\switchcolumn*

\end{paracol}

\Office{Feria V infra Hebdomadam II Septembris}{Thursday of the second week of September}{}
\addcontentsline{toc}{subsection}{Feria V infra Hebdomadam II Septembris \textit{— Thursday of the second week of September}}
\begin{paracol}{2}
\LA
\LessonHead{Lectio I}
\switchcolumn
\EN
\LessonHead{Lesson I}
\switchcolumn*

\LA
\LessonTitle{De libro Job}
\switchcolumn
\EN
\LessonTitle{Lesson from the book of Job}
\switchcolumn*

\LA
\Cite{Job 38:1-7}
\switchcolumn
\EN
\Cite{Job 38:1-7}
\switchcolumn*

\LA
\noindent Respóndens autem Dóminus Job de túrbine, dixit: Quis est iste invólvens senténtias sermónibus imperítis? Accínge, sicut vir, lumbos tuos; interrogábo te et respónde mihi. Ubi eras quando ponébam fundaménta terræ? Indica mihi, si habes intellegéntiam. Quis pósuit mensúras ejus, si nosti? vel quis teténdit super eam líneam? Super quo bases illíus solidátæ sunt? aut quis demísit lápidem angulárem ejus, cum me laudárent simul astra matutína, et jubilárent omnes fílii Dei?\par
\switchcolumn
\EN
\noindent Then the Lord answered Job out of a whirlwind, and said: Who is this that wrappeth up sentences in unskillful words? Gird up thy loins like a man I will ask thee, and answer thou me. Where wast thou when I laid up the foundations of the earth tell me if thou hast understanding. Who hath laid the measures thereof, if thou knowest? or who hath stretched the line upon it? Upon what are its bases grounded? or who laid the corner stone thereof, When the morning stars praised me together, and all the sons of God made a joyful melody?\par
\switchcolumn*

\LA
\begin{Resp}\Rsign Si bona suscépimus de manu Dei, mala autem quare non sustineámus? \Star{} Dóminus dedit, Dóminus ábstulit; sicut Dómino plácuit, ita factum est: sit nomen Dómini benedíctum.\end{Resp}
\begin{Resp}\Vsign Nudus egréssus sum de útero matris meæ et nudus revértar illuc.\end{Resp}
\begin{Resp}\Rsign Dóminus dedit, Dóminus ábstulit; sicut Dómino plácuit, ita factum est: sit nomen Dómini benedíctum.\end{Resp}
\switchcolumn
\EN
\begin{Resp}\Rsign What shall we receive good at the hand of God, and shall we not receive evil? \Star{} The Lord gave and the Lord hath taken away. As the Lord hath pleased, so hath it befallen. Blessed be the Name of the Lord.\end{Resp}
\begin{Resp}\Vsign Naked came I out of my mother's womb, and naked shall I return thither.\end{Resp}
\begin{Resp}\Rsign The Lord gave and the Lord hath taken away. As the Lord hath pleased, so hath it befallen. Blessed be the Name of the Lord.\end{Resp}
\switchcolumn*

\LA
\LessonHead{Lectio II}
\switchcolumn
\EN
\LessonHead{Lesson II}
\switchcolumn*

\LA
\Cite{Job 38:8-13}
\switchcolumn
\EN
\Cite{Job 38:8-13}
\switchcolumn*

\LA
\noindent Quis conclúsit óstiis mare, quando erumpébat quasi de vulva procédens, cum pónerem nubem vestiméntum ejus, et calígine illud, quasi pannis infántiæ, obvólverem? Circúmdedi illud términis meis, et pósui vectem et óstia, et dixi: Usque huc vénies, et non procédes ámplius, et hic confrínges tuméntes fluctus tuos. Numquid post ortum tuum præcepísti dilúculo, et ostendísti auróræ locum suum? Et tenuísti concútiens extréma terræ, et excussísti ímpios ex ea?\par
\switchcolumn
\EN
\noindent Who shut up the sea with doors, when it broke forth as issuing out of the womb? when I made a cloud the garment thereof, and wrapped it in a mist as in swaddling bands? I set my bounds around it, and made it bars and doors. And I said: Hitherto thou shalt come, and shalt go no further, and here thou shalt break thy swelling waves. Didst thou since thy birth command the morning, and shew the dawning of the day its place? And didst thou hold the extremities of the earth shaking them, and hast thou shaken the ungodly out of it?\par
\switchcolumn*

\LA
\begin{Resp}\Rsign Antequam cómedam, suspíro, et tamquam inundántes aquæ sic rugítus meus; quia timor, quem timébam, evénit mihi, et quod verébar áccidit. Nonne dissimulávi? nonne sílui? nonne quiévi? \Star{} Et venit super me indignátio.\end{Resp}
\begin{Resp}\Vsign Ecce non est auxílium mihi in me, et necessárii quoque mei recessérunt a me.\end{Resp}
\begin{Resp}\Rsign Et venit super me indignátio.\end{Resp}
\switchcolumn
\EN
\begin{Resp}\Rsign My sighing cometh before I eat, and my roarings are poured out like the waters, for the thing which I greatly feared is come upon me, and that which I was afraid of is come unto me. Was not I silent? Held not I my peace? Was not I at rest? \Star{} And trouble came.\end{Resp}
\begin{Resp}\Vsign Behold, I cannot help myself, and they that were needful unto me have forsaken me.\end{Resp}
\begin{Resp}\Rsign And trouble came.\end{Resp}
\switchcolumn*

\LA
\LessonHead{Lectio III}
\switchcolumn
\EN
\LessonHead{Lesson III}
\switchcolumn*

\LA
\Cite{Job 38:14-20}
\switchcolumn
\EN
\Cite{Job 38:14-20}
\switchcolumn*

\LA
\noindent Restituétur ut lutum signáculum, et stabit sicut vestiméntum: auferétur ab ímpiis lux sua, et brácchium excélsum confringétur? Numquid ingréssus es profúnda maris, et in novíssimis abýssi deambulásti? Numquid apértæ sunt tibi portæ mortis, et óstia tenebrósa vidísti? Numquid considerásti latitúdinem terræ? Indica mihi, si nosti, ómnia: in qua via lux hábitet, et tenebrárum quis locus sit; ut ducas unumquódque ad términos suos, et intéllegas sémitas domus ejus.\par
\switchcolumn
\EN
\noindent The seal shall be restored as clay, and shall stand as a garment. From the wicked their light shall be taken away, and the high arm shall be broken. Hast thou entered into the depths of the sea, and walked in the lowest parts of the deep? Have the gates of death been opened to thee, and hast thou seen the darksome doors? Hast thou considered the breadth of the earth? tell me, if thou knowest all things? Where is the way where light dwelleth, and where is the place of darkness? That thou mayst bring every thing to its own bounds, and understand the paths of the house thereof.\par
\switchcolumn*

\LA
\begin{Resp}\Rsign Quare detraxístis sermónibus veritátis? ad increpándum verba compónitis et subvértere nitímini amícum vestrum: \Star{} Verúmtamen quæ cogitástis, expléte.\end{Resp}
\begin{Resp}\Vsign Quod justum est, judicáte; et non inveniétis in lingua mea iniquitátem.\end{Resp}
\begin{Resp}\Rsign Verúmtamen quæ cogitástis, expléte.\end{Resp}
\begin{Resp}\Vsign Glória Patri, et Fílio, \Star{} et Spirítui Sancto.\end{Resp}
\begin{Resp}\Rsign Verúmtamen quæ cogitástis, expléte.\end{Resp}
\switchcolumn
\EN
\begin{Resp}\Rsign Why do ye argue against the words of truth? Do ye imagine words to reprove me and strive to confound one that is your friend? \Star{} Nevertheless, finish that ye have in mind.\end{Resp}
\begin{Resp}\Vsign Judge that which is just, and ye shall find no iniquity in my tongue.\end{Resp}
\begin{Resp}\Rsign Nevertheless, finish that ye have in mind.\end{Resp}
\begin{Resp}\Vsign Glory be to the Father, and to the Son, \Star{} and to the Holy Ghost.\end{Resp}
\begin{Resp}\Rsign Nevertheless, finish that ye have in mind.\end{Resp}
\switchcolumn*

\end{paracol}

\Office{Feria VI infra Hebdomadam II Septembris}{Friday of the second week of September}{}
\addcontentsline{toc}{subsection}{Feria VI infra Hebdomadam II Septembris \textit{— Friday of the second week of September}}
\begin{paracol}{2}
\LA
\LessonHead{Lectio I}
\switchcolumn
\EN
\LessonHead{Lesson I}
\switchcolumn*

\LA
\LessonTitle{De libro Job}
\switchcolumn
\EN
\LessonTitle{Lesson from the book of Job}
\switchcolumn*

\LA
\Cite{Job 40:1-5}
\switchcolumn
\EN
\Cite{Job 40:1-5}
\switchcolumn*

\LA
\noindent Respóndens autem Dóminus Job de túrbine, dixit: Accínge sicut vir lumbos tuos; interrogábo te, et índica mihi. Numquid írritum fácies judícium meum, et condemnábis me, ut tu justificéris? Et, si habes brácchium sicut Deus, et si voce símili tonas, circúmda tibi decórem, et in sublíme erígere, et esto gloriósus et speciósis indúere véstibus.\par
\switchcolumn
\EN
\noindent And the Lord answering Job out of the whirlwind, said: Gird up thy loins like a man I will ask thee, and do thou tell me. Wilt thou make void my judgment and condemn me, that thou mayst be justified? And hast thou an arm like God, and canst thou thunder with a voice like him? Clothe thyself with beauty, and set thyself up on high and be glorious, and put on goodly garments.\par
\switchcolumn*

\LA
\begin{Resp}\Rsign Indúta est caro mea putrédine, et sórdibus púlveris cutis mea áruit et contrácta est: \Star{} Meménto mei, Dómine, quóniam ventus est vita mea.\end{Resp}
\begin{Resp}\Vsign Dies mei velócius transiérunt quam a texénte tela succíditur, et consúmpti sunt absque ulla spe.\end{Resp}
\begin{Resp}\Rsign Meménto mei, Dómine, quóniam ventus est vita mea.\end{Resp}
\switchcolumn
\EN
\begin{Resp}\Rsign My flesh is clothed with worms and clods of dust. My skin is dry and drawn together. \Star{} Remember me, O Lord, for my life is wind.\end{Resp}
\begin{Resp}\Vsign My days are swifter than a weaver's shuttle, and are spent without hope.\end{Resp}
\begin{Resp}\Rsign Remember me, Lord, for my life is wind.\end{Resp}
\switchcolumn*

\LA
\LessonHead{Lectio II}
\switchcolumn
\EN
\LessonHead{Lesson II}
\switchcolumn*

\LA
\Cite{Job 40:6-11}
\switchcolumn
\EN
\Cite{Job 40:6-11}
\switchcolumn*

\LA
\noindent Dispérge supérbos in furóre tuo et respíciens omnem arrogántem humília; réspice cunctos supérbos et confúnde eos, et cóntere ímpios in loco suo; abscónde eos in púlvere simul, et fácies eórum demérge in fóveam; et ego confitébor quod salváre te possit déxtera tua. Ecce béhemoth, quem feci tecum, fænum, quasi bos, cómedet. Fortitúdo ejus in lumbis ejus, et virtus illíus in umbilíco ventris ejus.\par
\switchcolumn
\EN
\noindent Scatter the proud in thy indignation, and behold every arrogant man, and humble him. Look on all that are proud, and confound them, and crush the wicked in their place. Hide them in the dust together, and plunge their faces into the pit. Then I will confess that thy right hand is able to save thee. Behold behemoth whom I made with thee, he eateth grass like an ox. His strength is in his loins, and his force in the navel of his belly.\par
\switchcolumn*

\LA
\begin{Resp}\Rsign Páucitas diérum meórum finiétur brevi; dimítte me, Dómine, ut plangam páululum dolórem meum, \Star{} Antequam vadam ad terram tenebrósam et opértam mortis calígine.\end{Resp}
\begin{Resp}\Vsign Manus tuæ, Dómine, fecérunt me, et plasmavérunt me totum in circúitu; et sic repénte præcípitas me?\end{Resp}
\begin{Resp}\Rsign Antequam vadam ad terram tenebrósam et opértam mortis calígine.\end{Resp}
\switchcolumn
\EN
\begin{Resp}\Rsign My days are few, and in a short while they will be ended. Let me alone, then, O Lord, that I may bewail my sorrow a little; \Star{} Before I go to the land of darkness and of the shadow of death.\end{Resp}
\begin{Resp}\Vsign Thine hands, O Lord, have made me, and fashioned me together round about, and yet dost Thou forthwith destroy me.\end{Resp}
\begin{Resp}\Rsign Before I go to the land of darkness and of the shadow of death.\end{Resp}
\switchcolumn*

\LA
\LessonHead{Lectio III}
\switchcolumn
\EN
\LessonHead{Lesson III}
\switchcolumn*

\LA
\Cite{Job 42:1-6}
\switchcolumn
\EN
\Cite{Job 42:1-6}
\switchcolumn*

\LA
\noindent Respóndens autem Job Dómino, dixit: Scio quia ómnia potes, et nulla te latet cogitátio. Quis est iste qui celat consílium absque sciéntia? Ideo insipiénter locútus sum, et quæ ultra modum excéderent sciéntiam meam. Audi, et ego loquar; interrogábo te, et respónde mihi. Audítu auris audívi te; nunc autem óculus meus videt te. Idcírco ipse me reprehéndo, et ago pœniténtiam in favílla et cínere.\par
\switchcolumn
\EN
\noindent Then Job answered the Lord, and said: I know that thou canst do all things, and no thought is hid from thee. Who is this that hideth counsel without knowledge? Therefore I have spoken unwisely, and things that above measure exceeded my knowledge. Hear, and I will speak I will ask thee, and do thou tell me. With the hearing of the ear, I have heard thee, but now my eye seeth thee. Therefore I reprehend myself, and do penance in dust and ashes.\par
\switchcolumn*

\LA
\begin{Resp}\Rsign Non abscóndas me, Dómine, a fácie tua: manum tuam longe fac a me, \Star{} Et formído tua non me térreat.\end{Resp}
\begin{Resp}\Vsign Córripe me, Dómine, in misericórdia, non in furóre tuo, ne forte ad níhilum rédigas me.\end{Resp}
\begin{Resp}\Rsign Et formído tua non me térreat.\end{Resp}
\begin{Resp}\Vsign Glória Patri, et Fílio, \Star{} et Spirítui Sancto.\end{Resp}
\begin{Resp}\Rsign Et formído tua non me térreat.\end{Resp}
\switchcolumn
\EN
\begin{Resp}\Rsign Hide not thy face from me, O Lord. Withdraw not thine hand far from me, \Star{} And let not thy dread make me afraid.\end{Resp}
\begin{Resp}\Vsign O Lord, correct me but in mercy not in thine anger, lest Thou bring me to nothing.\end{Resp}
\begin{Resp}\Rsign And let not thy dread make me afraid.\end{Resp}
\begin{Resp}\Vsign Glory be to the Father, and to the Son, \Star{} and to the Holy Ghost.\end{Resp}
\begin{Resp}\Rsign And let not thy dread make me afraid.\end{Resp}
\switchcolumn*

\end{paracol}

\Office{Sabbato infra Hebdomadam II Septembris}{Saturday of the second week of September}{}
\addcontentsline{toc}{subsection}{Sabbato infra Hebdomadam II Septembris \textit{— Saturday of the second week of September}}
\begin{paracol}{2}
\LA
\LessonHead{Lectio I}
\switchcolumn
\EN
\LessonHead{Lesson I}
\switchcolumn*

\LA
\LessonTitle{De libro Job}
\switchcolumn
\EN
\LessonTitle{Lesson from the book of Job}
\switchcolumn*

\LA
\Cite{Job 42:7-8}
\switchcolumn
\EN
\Cite{Job 42:7-8}
\switchcolumn*

\LA
\noindent Postquam autem locútus est Dóminus verba hæc ad Job, dixit ad Eliphaz Themanítem: Irátus est furor meus in te et in duos amícos tuos, quóniam non estis locúti coram me rectum, sicut servus meus Job. Súmite ergo vobis septem tauros et septem aríetes et ite ad servum meum Job et offérte holocáustum pro vobis; Job autem servus meus orábit pro vobis. Fáciem ejus suscípiam, ut non vobis imputétur stultítia; neque enim locúti estis ad me recta, sicut servus meus Job.\par
\switchcolumn
\EN
\noindent And after the Lord had spoken these words to Job, he said to Eliphaz the Themanite: My wrath is kindled against thee, and against thy two friends, because you have not spoken the thing that is right before me, as my servant Job hath. Take unto you therefore seven oxen, and seven rams, and go to my servant Job, and offer for yourselves a holocaust and my servant Job shall pray for you; his face I will accept, that folly be not imputed to you for you have not spoken right things before me, as my servant Job hath.\par
\switchcolumn*

\LA
\begin{Resp}\Rsign Hódie nata est beáta Virgo María ex progénie David; \Star{} Per quam salus mundi credéntibus appáruit, cujus vita gloriósa lucem dedit sǽculo.\end{Resp}
\begin{Resp}\Vsign Nativitátem beátæ Maríæ Vírginis cum gáudio celebrémus.\end{Resp}
\begin{Resp}\Rsign Per quam salus mundi credéntibus appáruit, cujus vita gloriósa lucem dedit sǽculo.\end{Resp}
\switchcolumn
\EN
\begin{Resp}\Rsign This day was the Blessed Virgin Mary born of the lineage of David. \Star{} The same is she through whom the salvation of the world hath been manifested before the eyes of all believers. This is she whose glorious life hath given light to the world.\end{Resp}
\begin{Resp}\Vsign Let us keep with rejoicing the Birthday of the Blessed Virgin Mary.\end{Resp}
\begin{Resp}\Rsign The same is she through whom the salvation of the world hath been manifested before the eyes of all believers. This is she whose glorious life hath given light to the world.\end{Resp}
\switchcolumn*

\LA
\LessonHead{Lectio II}
\switchcolumn
\EN
\LessonHead{Lesson II}
\switchcolumn*

\LA
\Cite{Job 42:9-11}
\switchcolumn
\EN
\Cite{Job 42:9-11}
\switchcolumn*

\LA
\noindent Abiérunt ergo Eliphaz Themanítes et Baldad Suhítes et Sophar Naamathítes, et fecérunt sicut locútus fúerat Dóminus ad eos, et suscépit Dóminus fáciem Job. Dóminus quoque convérsus est ad pœniténtiam Job, cum oráret ille pro amícis suis; et áddidit Dóminus ómnia, quæcúmque fúerant Job, duplícia. Venérunt autem ad eum omnes fratres sui et univérsæ soróres suæ et cuncti, qui nóverant eum prius, et comedérunt cum eo panem in domo ejus; et movérunt super eum caput, et consoláti sunt eum super omni malo, quod intúlerat Dóminus super eum, et dedérunt ei unusquísque ovem unam et ináurem áuream unam.\par
\switchcolumn
\EN
\noindent So Eliphaz the Themanite, and Baldad the Suhite, and Sophar the Naamathite went, and did as the Lord had spoken to them, and the Lord accepted the face of Job. The Lord also was turned at the penance of Job, when he prayed for his friends. And the Lord gave Job twice as much as he had before. And all his brethren came to him, and all his sisters, and all that knew him before, and they ate bread with him in his house and bemoaned him, and comforted him upon all the evil that God had brought upon him. And every man gave him one ewe, and one earring of gold.\par
\switchcolumn*

\LA
\begin{Resp}\Rsign Beatíssimæ Vírginis Maríæ Nativitátem devotíssime celebrémus, \Star{} Ut ipsa pro nobis intercédat ad Dóminum Jesum Christum.\end{Resp}
\begin{Resp}\Vsign Cum jucunditáte Nativitátem beátæ Maríæ Vírginis devotíssime celebrémus.\end{Resp}
\begin{Resp}\Rsign Ut ipsa pro nobis intercédat ad Dóminum Jesum Christum.\end{Resp}
\begin{Resp}\Vsign Glória Patri, et Fílio, \Star{} et Spirítui Sancto.\end{Resp}
\begin{Resp}\Rsign Ut ipsa pro nobis intercédat ad Dóminum Jesum Christum.\end{Resp}
\switchcolumn
\EN
\begin{Resp}\Rsign Let us keep right heartily the Birthday of the Most Blessed Virgin Mary, \Star{} That she may pray for us to our Lord Jesus Christ.\end{Resp}
\begin{Resp}\Vsign Let us keep with right hearty rejoicing the Birthday of the Blessed Virgin Mary.\end{Resp}
\begin{Resp}\Rsign That she may pray for us to our Lord Jesus Christ.\end{Resp}
\begin{Resp}\Vsign Glory be to the Father, and to the Son, \Star{} and to the Holy Ghost.\end{Resp}
\begin{Resp}\Rsign That she may pray for us to our Lord Jesus Christ.\end{Resp}
\switchcolumn*

\LA
\LessonHead{Lectio III}
\switchcolumn
\EN
\LessonHead{Lesson III}
\switchcolumn*

\LA
\Cite{Job 42:12-16}
\switchcolumn
\EN
\Cite{Job 42:12-16}
\switchcolumn*

\LA
\noindent Dóminus autem benedíxit novíssimis Job magis quam princípio ejus; et facta sunt ei quatuórdecim míllia óvium et sex míllia camelórum et mille juga boum et mille ásinæ, et fuérunt ei septem fílii et tres fíliæ, et vocávit nomen uníus Diem et nomen secúndæ Cássiam et nomen tértiæ Cornustíbii. Non sunt autem invéntæ mulíeres speciósæ sicut fíliæ Job in univérsa terra; dedítque eis pater suus hereditátem inter fratres eárum. Vixit autem Job post hæc centum quadragínta annis, et vidit fílios suos et fílios filiórum suórum usque ad quartam generatiónem; et mórtuus est senex et plenus diérum.\par
\switchcolumn
\EN
\noindent And the Lord blessed the latter end of Job more than his beginning. And he had fourteen thousand sheep, and six thousand camels, and a thousand yoke of oxen, and a thousand she asses. And he had seven sons, and three daughters. And he called the names of one Dies, and the name of the second Cassia, and the name of the third Cornustibil. And there were not found in all the earth women so beautiful as the daughters of Job and their father gave them inheritance among their brethren. And Job lived after these things, a hundred and forty years, and he saw his children, and his children's children, unto the fourth generation, and he died an old man, and full of days.\par
\switchcolumn*

\LA
\begin{Resp}\Rsign Quare detraxístis sermónibus veritátis? ad increpándum verba compónitis et subvértere nitímini amícum vestrum: \Star{} Verúmtamen quæ cogitástis, expléte.\end{Resp}
\begin{Resp}\Vsign Quod justum est, judicáte; et non inveniétis in lingua mea iniquitátem.\end{Resp}
\begin{Resp}\Rsign Verúmtamen quæ cogitástis, expléte.\end{Resp}
\begin{Resp}\Vsign Glória Patri, et Fílio, \Star{} et Spirítui Sancto.\end{Resp}
\begin{Resp}\Rsign Verúmtamen quæ cogitástis, expléte.\end{Resp}
\switchcolumn
\EN
\begin{Resp}\Rsign Why do ye argue against the words of truth? Do ye imagine words to reprove me and strive to confound one that is your friend? \Star{} Nevertheless, finish that ye have in mind.\end{Resp}
\begin{Resp}\Vsign Judge that which is just, and ye shall find no iniquity in my tongue.\end{Resp}
\begin{Resp}\Rsign Nevertheless, finish that ye have in mind.\end{Resp}
\begin{Resp}\Vsign Glory be to the Father, and to the Son, \Star{} and to the Holy Ghost.\end{Resp}
\begin{Resp}\Rsign Nevertheless, finish that ye have in mind.\end{Resp}
\switchcolumn*

\end{paracol}

\Office{Dominica III Septembris}{Third Sunday of September}{}
\addcontentsline{toc}{subsection}{Dominica III Septembris \textit{— Third Sunday of September}}
\begin{paracol}{2}
\LA
\LessonHead{Lectio I}
\switchcolumn
\EN
\LessonHead{Lesson I}
\switchcolumn*

\LA
\LessonTitle{Incipit liber Tobíæ}
\switchcolumn
\EN
\LessonTitle{Lesson from the book of Tobias}
\switchcolumn*

\LA
\Cite{Tob 1:1-4}
\switchcolumn
\EN
\Cite{Tob 1:1-4}
\switchcolumn*

\LA
\noindent Tobías ex tribu et civitáte Néphthali, quæ est in superióribus Galilǽæ supra Naásson, post viam quæ ducit ad Occidéntem in sinístro habens civitátem Sephet, cum captus esset in diébus Salmánasar regis Assyriórum, in captivitáte tamen pósitus viam veritátis non deséruit; ita ut ómnia, quæ habére póterat, cotídie concaptívis frátribus, qui erant ex ejus génere, impertíret. Cumque esset júnior ómnibus in tribu Néphtali, nihil tamen pueríle gessit in ópere.\par
\switchcolumn
\EN
\noindent Tobias of the tribe and city of Nephtali, (which is in the upper parts of Galilee above Naasson, beyond the way that leadeth to the west, having on the right hand the city of Sephet,) When he was made captive in the days of Salmanasar king of the Assyrians, even in his captivity, forsook not the way of truth, But every day gave all he could get to his brethren his fellow captives, that were of his kindred. And when he was younger than any of the tribe of Nephtali, yet did he no childish thing in his work.\par
\switchcolumn*

\LA
\begin{Resp}\Rsign Peto, Dómine, ut de vínculo impropérii hujus absólvas me, aut certe désuper terram erípias me: \Star{} Ne reminiscáris delícta mea vel paréntum meórum, neque vindíctam sumas de peccátis meis: quia éruis sustinéntes te, Dómine.\end{Resp}
\begin{Resp}\Vsign Omnia enim judícia tua justa sunt, et omnes viæ tuæ misericórdia et véritas: et nunc, Dómine, meménto mei.\end{Resp}
\begin{Resp}\Rsign Ne reminiscáris delícta mea vel paréntum meórum, neque vindíctam sumas de peccátis meis: quia éruis sustinéntes te, Dómine.\end{Resp}
\switchcolumn
\EN
\begin{Resp}\Rsign I entreat thee, O Lord, that thou wouldest loose me from this reproach, or that thou wouldest take me away from the earth: \Star{} Remember not mine offences nor the offences of my forefathers, neither take thou vengeance of my sins; for thou, O Lord, art a Redeemer unto all that hope in thee.\end{Resp}
\begin{Resp}\Vsign For all thy judgements are just, and all thy ways are mercy and truth; and now, O Lord, remember me.\end{Resp}
\begin{Resp}\Rsign Remember not mine offences nor the offences of my forefathers, neither take thou vengeance of my sins; for thou, O Lord, art a Redeemer unto all that hope in thee.\end{Resp}
\switchcolumn*

\LA
\LessonHead{Lectio II}
\switchcolumn
\EN
\LessonHead{Lesson II}
\switchcolumn*

\LA
\Cite{Tob 1:5-10}
\switchcolumn
\EN
\Cite{Tob 1:5-10}
\switchcolumn*

\LA
\noindent Dénique, cum irent omnes ad vítulos áureos, quos Jeróboam fécerat rex Israël, hic solus fugiébat consórtia ómnium; sed pergébat in Jerúsalem ad templum Dómini et ibi adorábat Dóminum Deum Israël, ómnia primitíva sua et décimas suas fidéliter ófferens; ita ut in tértio anno prosélytis et ádvenis ministráret omnem decimatiónem. Hæc et his simília secúndum legem Dei puérulus observábat. Cum vero factus esset vir, accépit uxórem Annam de tribu sua genuítque ex ea fílium, nomen suum impónens ei. Quem ab infántia timére Deum dócuit, et abstinére ab omni peccáto.\par
\switchcolumn
\EN
\noindent Moreover when all went to the golden calves which Jeroboam king of Israel had made, he alone fled the company of all, And went to Jerusalem to the temple of the Lord, and there adored the Lord God of Israel, offering faithfully all his firstfruits, and his tithes, So that in the third year he gave all his tithes to the proselytes, and strangers. These and such like things did he observe when but a boy according to the law of God. But when he was a man, he took to wife Anna of his own tribe, and had a son by her, whom he called after his own name, And from his infancy he taught him to fear God, and to abstain from all sin.\par
\switchcolumn*

\LA
\begin{Resp}\Rsign Omni témpore bénedic Deum, et pete ab eo ut vias tuas dírigat, \Star{} Et in omni témpore consília tua in ipso permáneant.\end{Resp}
\begin{Resp}\Vsign Inquíre ut fácias quæ plácita sunt illi in veritáte, et in tota virtúte tua.\end{Resp}
\begin{Resp}\Rsign Et in omni témpore consília tua in ipso permáneant.\end{Resp}
\switchcolumn
\EN
\begin{Resp}\Rsign In all seasons bless God, and ask of Him to order thy goings; \Star{} And in all seasons let thy counsels be steadfastly in Him.\end{Resp}
\begin{Resp}\Vsign Seek faithfully and with all thy strength to do such things as please Him.\end{Resp}
\begin{Resp}\Rsign And in all seasons let thy counsels be steadfastly in Him.\end{Resp}
\switchcolumn*

\LA
\LessonHead{Lectio III}
\switchcolumn
\EN
\LessonHead{Lesson III}
\switchcolumn*

\LA
\Cite{Tob 1:11-15}
\switchcolumn
\EN
\Cite{Tob 1:11-15}
\switchcolumn*

\LA
\noindent Igitur, cum per captivitátem devenísset cum uxóre sua et fílio in civitátem Níniven cum omni tribu sua, cum omnes éderent ex cibis gentílium, iste custodívit ánimam suam et numquam contaminátus est in escis eórum. Et quóniam memor fuit Dómini in toto corde suo, dedit illi Deus grátiam in conspéctu Salmánasar regis, et dedit illi potestátem quocúmque vellet ire, habens libertátem quæcúmque fácere voluísset. Pergébat ergo ad omnes, qui erant in captivitáte, et mónita salútis dabat eis.\par
\switchcolumn
\EN
\noindent And when by the captivity he with his wife and his son and all his tribe was come to the city of Ninive, (When all ate of the meats of the Gentiles) he kept his soul and never was defiled with their meats. And because he was mindful of the Lord with all his heart, God gave him favour in the sight of Salmanasar the king. And he gave him leave to go whithersoever he would, with liberty to do whatever he had a mind. He therefore went to all that were in captivity, and gave them wholesome admonitions.\par
\switchcolumn*

\LA
\begin{Resp}\Rsign Memor esto, fili, quóniam páuperem vitam gérimus: \Star{} Habébis multa bona, si timúeris Deum.\end{Resp}
\begin{Resp}\Vsign In mente habéto eum, et cave nequándo prætermíttas præcépta ejus.\end{Resp}
\begin{Resp}\Rsign Habébis multa bona, si timúeris Deum.\end{Resp}
\begin{Resp}\Vsign Glória Patri, et Fílio, \Star{} et Spirítui Sancto.\end{Resp}
\begin{Resp}\Rsign Habébis multa bona, si timúeris Deum.\end{Resp}
\switchcolumn
\EN
\begin{Resp}\Rsign My son, remember that we have but a frail life. \Star{} If thou fear God thou shalt have great goods.\end{Resp}
\begin{Resp}\Vsign Be mindful of Him, and beware lest ever thou transgress His commandments.\end{Resp}
\begin{Resp}\Rsign If thou fear God thou shalt have great goods.\end{Resp}
\begin{Resp}\Vsign Glory be to the Father, and to the Son, \Star{} and to the Holy Ghost.\end{Resp}
\begin{Resp}\Rsign If thou fear God thou shalt have great goods.\end{Resp}
\switchcolumn*

\LA
\LessonHead{Lectio IV}
\switchcolumn
\EN
\LessonHead{Lesson IV}
\switchcolumn*

\LA
\LessonTitle{Sermo sancti Leónis Papæ}
\switchcolumn
\EN
\LessonTitle{Sermon of St. Leo the Great, Pope.}
\switchcolumn*

\LA
\Source{Sermo 9 de jejúnio 7 mensis}
\switchcolumn
\EN
\Source{Serm. 9. de jejunio 7. mensis.}
\switchcolumn*

\LA
\noindent Scio quidem, dilectíssimi, plúrimos vestrum ita in iis, quæ ad observántiam christiánæ fídei pértinent, esse devótos, ut nostris cohortatiónibus non indígeant admonéri. Quod enim dudum et tradítio decrévit, et consuetúdo firmávit; nec erudítio ignórat, nec píetas prætermíttit. Sed quia sacerdotális offícii est, erga omnes Ecclésiæ fílios curam habére commúnem, in id quod et rúdibus prosit et doctis, quos simul dilígimus, páriter incitámus; ut jejúnium, quod nobis séptimi mensis recúrsus indícit, fide álacri per castigatiónem ánimi et córporis celebrémus.\par
\switchcolumn
\EN
\noindent Well do I know, dearly beloved, that many of you are fervent in your observance of all those practices which belong to the Christian Faith, so that ye have no need to be admonished by our exhortations. For what tradition hath laid down, and custom well established, is neither unknown to the learned nor neglected by the devout. But because it appertaineth to the priestly office to exercise the same general care over all the Church's children in all such matters as be profitable alike to the learned and to the simple (both of whom are equally dear to us), we do now exhort the both of you to celebrate, with lively faith, and all due discipline of soul and body, the Quarterly Fast, which the seventh month [that is, September] doth once again bring to us in its yearly round.\par
\switchcolumn*

\LA
\begin{Resp}\Rsign Sufficiébat nobis paupértas nostra, ut divítiæ computaréntur: numquam fuísset pecúnia ipsa, pro qua misísti fílium nostrum, \Star{} Báculum senectútis nostræ!\end{Resp}
\begin{Resp}\Vsign Heu me, fili mi, ut quid te mísimus peregrinári, lumen oculórum nostrórum?\end{Resp}
\begin{Resp}\Rsign Báculum senectútis nostræ!\end{Resp}
\switchcolumn
\EN
\begin{Resp}\Rsign Our poverty was enough for us, that it might have been accounted riches. O that the money had never been, for which thou hast sent away our son, \Star{} The staff of our old age.\end{Resp}
\begin{Resp}\Vsign Alas, my son, wherefore have we sent thee wandering, even thee, the light of our eyes?\end{Resp}
\begin{Resp}\Rsign The staff of our old age.\end{Resp}
\switchcolumn*

\LA
\LessonHead{Lectio V}
\switchcolumn
\EN
\LessonHead{Lesson V}
\switchcolumn*

\LA
\noindent Ideo enim ipsa continéntiæ observántia quátuor est assignáta tempóribus, ut in idípsum totíus anni redeúnte decúrsu, cognoscerémus nos indesinénter purificatiónibus indigére; sempérque esse niténdum, dum hujus vitæ varietáte jactámur, ut peccátum, quod fragilitáte carnis et cupiditátum pollutióne contráhitur, jejúniis atque eleemósynis deleátur. Esuriámus páululum, dilectíssimi, et aliquántulum, quod juvándis possit prodésse paupéribus, nostræ consuetúdini subtrahámus.\par
\switchcolumn
\EN
\noindent The Ember Days of fasting are appointed to the four seasons, in order that their quarterly recurrence in the course of the year may teach us how unceasingly we need to be purified, and how, as long as we are tossed about by the changes and chances of this life, we need through fasting and alms-deeds to be cleansed from the stain of that sin which we have contracted by the frailty of our flesh and our concupiscence. Let us diminish a little, beloved, what we are accustomed to use for ourselves, in order that we have somewhat more to use for the relief of the poor and needy.\par
\switchcolumn*

\LA
\begin{Resp}\Rsign Benedícite Deum cæli et coram ómnibus vivéntibus confitémini ei, \Star{} Quia fecit vobíscum misericórdiam suam.\end{Resp}
\begin{Resp}\Vsign Ipsum benedícite et cantáte illi: et enarráte ómnia mirabília ejus.\end{Resp}
\begin{Resp}\Rsign Quia fecit vobíscum misericórdiam suam.\end{Resp}
\switchcolumn
\EN
\begin{Resp}\Rsign Bless the God of heaven, and confess Him before all living: \Star{} For He hath had mercy upon you.\end{Resp}
\begin{Resp}\Vsign Bless Him, and sing praises unto Him, and tell of all His marvelous works.\end{Resp}
\begin{Resp}\Rsign For He hath had mercy upon you.\end{Resp}
\switchcolumn*

\LA
\LessonHead{Lectio VI}
\switchcolumn
\EN
\LessonHead{Lesson VI}
\switchcolumn*

\LA
\noindent Delectétur consciéntia benignórum frúctibus largitátis: et gáudia tríbuens, quo es lætificándus, accípies. Diléctio próximi, diléctio Dei est, qui plenitúdinem legis et prophetárum in hac géminæ caritátis unitáte constítuit; ut nemo ambígeret, Deo se offérre, quod hómini contulísset, dicénte Dómino Salvatóre, cum de aléndis juvandísque paupéribus loquerétur: Quod uni eórum fecístis, mihi fecístis. Quarta ígitur et sexta féria jejunémus; sábbato vero apud beátum Petrum Apóstolum vigílias celebrémus: cujus nos méritis et oratiónibus crédimus adjuvándos, ut misericórdi Deo jejúnio nostro et devotióne placeámus.\par
\switchcolumn
\EN
\noindent The conscience of the generous can thus be made glad by the fruits of their own liberality. Whilst thou art giving happiness thou shalt receive joy. Thy love for thy neighbor is a unity with thy love for God; and he hath taught us that in the unity of this twofold charity is to be found the fulfillment of all the Law and the Prophets. Further, if anyone doubt that what is given to man is offered to God, we have the saying of our Lord and Saviour, when he spake of feeding and helping the poor: Inasmuch as ye have done it unto one of the least of these, my brethren, ye have done it unto me. Wherefore, let us fast on Ember Wednesday and Friday; and on Ember Saturday let us also keep vigil at the shrine of blessed Peter the Apostle; by whose merits and prayers we believe that we shall be aided, so that we may please our merciful God in our fasting and prayer.\par
\switchcolumn*

\LA
\begin{Resp}\Rsign Tempus est ut revértar ad eum qui misit me; \Star{} Vos autem benedícite Deum, et enarráte ómnia mirabília ejus.\end{Resp}
\begin{Resp}\Vsign Confitémini ei coram ómnibus vivéntibus, quia fecit vobíscum misericórdiam suam.\end{Resp}
\begin{Resp}\Rsign Vos autem benedícite Deum, et enarráte ómnia mirabília ejus.\end{Resp}
\begin{Resp}\Vsign Glória Patri, et Fílio, \Star{} et Spirítui Sancto.\end{Resp}
\begin{Resp}\Rsign Vos autem benedícite Deum, et enarráte ómnia mirabília ejus.\end{Resp}
\switchcolumn
\EN
\begin{Resp}\Rsign It is time for me to return unto Him That sent me; \Star{} But bless ye God, and tell of all His marvelous works.\end{Resp}
\begin{Resp}\Vsign Confess Him before all living, for He hath had mercy upon you.\end{Resp}
\begin{Resp}\Rsign But bless ye God, and tell of all His marvelous works.\end{Resp}
\begin{Resp}\Vsign Glory be to the Father, and to the Son, \Star{} and to the Holy Ghost.\end{Resp}
\begin{Resp}\Rsign But bless ye God, and tell of all His marvelous works.\end{Resp}
\switchcolumn*

\end{paracol}

\Office{Feria II infra Hebdomadam III Septembris}{Monday of the third week of September}{}
\addcontentsline{toc}{subsection}{Feria II infra Hebdomadam III Septembris \textit{— Monday of the third week of September}}
\begin{paracol}{2}
\LA
\LessonHead{Lectio I}
\switchcolumn
\EN
\LessonHead{Lesson I}
\switchcolumn*

\LA
\LessonTitle{De libro Tobíæ}
\switchcolumn
\EN
\LessonTitle{Lesson from the book of Tobias}
\switchcolumn*

\LA
\Cite{Tob 2:1-4}
\switchcolumn
\EN
\Cite{Tob 2:1-4}
\switchcolumn*

\LA
\noindent Post hæc vero, cum esset dies festus Dómini, et factum esset prándium bonum in domo Tobíæ, dixit fílio suo: Vade et adduc áliquos de tribu nostra timéntes Deum, ut epuléntur nobíscum. Cumque abiísset, revérsus nuntiávit ei unum ex fíliis Israël jugulátum jacére in platéa. Statímque exsíliens de accúbitu suo relínquens prándium, jejúnus pervénit ad corpus, tollénsque illud portávit ad domum suam occúlte, ut, dum sol occubuísset, caute sepelíret eum.\par
\switchcolumn
\EN
\noindent But after this, when there was a festival of the Lord, and a good dinner was prepared in Tobias's house, He said to his son: Go, and bring some of our tribe that fear God, to feast with us. And when he had gone, returning he told him, that one of the children of Israel lay slain in the street. And he forthwith leaped up from his place at the table, and left his dinner, and came fasting to the body: And taking it up carried it privately to his house, that after the sun was down, he might bury him cautiously.\par
\switchcolumn*

\LA
\begin{Resp}\Rsign Peto, Dómine, ut de vínculo impropérii hujus absólvas me, aut certe désuper terram erípias me: \Star{} Ne reminiscáris delícta mea vel paréntum meórum, neque vindíctam sumas de peccátis meis: quia éruis sustinéntes te, Dómine.\end{Resp}
\begin{Resp}\Vsign Omnia enim judícia tua justa sunt, et omnes viæ tuæ misericórdia et véritas: et nunc, Dómine, meménto mei.\end{Resp}
\begin{Resp}\Rsign Ne reminiscáris delícta mea vel paréntum meórum, neque vindíctam sumas de peccátis meis: quia éruis sustinéntes te, Dómine.\end{Resp}
\switchcolumn
\EN
\begin{Resp}\Rsign I entreat thee, O Lord, that thou wouldest loose me from this reproach, or that thou wouldest take me away from the earth: \Star{} Remember not mine offences nor the offences of my forefathers, neither take thou vengeance of my sins; for thou, O Lord, art a Redeemer unto all that hope in thee.\end{Resp}
\begin{Resp}\Vsign For all thy judgements are just, and all thy ways are mercy and truth; and now, O Lord, remember me.\end{Resp}
\begin{Resp}\Rsign Remember not mine offences nor the offences of my forefathers, neither take thou vengeance of my sins; for thou, O Lord, art a Redeemer unto all that hope in thee.\end{Resp}
\switchcolumn*

\LA
\LessonHead{Lectio II}
\switchcolumn
\EN
\LessonHead{Lesson II}
\switchcolumn*

\LA
\Cite{Tob 2:8-12}
\switchcolumn
\EN
\Cite{Tob 2:8-12}
\switchcolumn*

\LA
\noindent Arguébant autem eum omnes próximi ejus dicéntes: Jam hujus rei causa intérfici jussus es et vix effugísti mortis impérium; et íterum sépelis mórtuos? Sed Tobías plus timens Deum quam regem, rapiébat córpora occisórum et occultábat in domo sua et médiis nóctibus sepeliébat ea. Cóntigit autem ut quadam die fatigátus a sepultúra, véniens in domum suam jactásset se juxta paríetem et obdormísset, et ex nido hirúndinum dormiénti illi cálida stércora incíderent super óculos ejus fierétque cæcus. Hanc autem tentatiónem ídeo permísit Dóminus eveníre illi, ut pósteris darétur exémplum patiéntiæ ejus, sicut et sancti Job.\par
\switchcolumn
\EN
\noindent Now all his neighbours blamed him, saying: Once already commandment was given for thee to be slain because of this matter, and thou didst scarce escape the sentence of death, and dost thou again bury the dead? But Tobias fearing God more than the king, carried off the bodies of them that were slain, and hid them in his house, and at midnight buried them. Now it happened one day, that being wearied with burying, he came to his house, and cast himself down by the wall and slept, And as he was sleeping, hot dung out of a swallow's nest fell upon his eyes, and he was made blind. Now this trial the Lord therefore permitted to happen to him, that an example might be given to posterity of his patience, as also of holy Job.\par
\switchcolumn*

\LA
\begin{Resp}\Rsign Omni témpore bénedic Deum, et pete ab eo ut vias tuas dírigat, \Star{} Et in omni témpore consília tua in ipso permáneant.\end{Resp}
\begin{Resp}\Vsign Inquíre ut fácias quæ plácita sunt illi in veritáte, et in tota virtúte tua.\end{Resp}
\begin{Resp}\Rsign Et in omni témpore consília tua in ipso permáneant.\end{Resp}
\switchcolumn
\EN
\begin{Resp}\Rsign In all seasons bless God, and ask of Him to order thy goings; \Star{} And in all seasons let thy counsels be steadfastly in Him.\end{Resp}
\begin{Resp}\Vsign Seek faithfully and with all thy strength to do such things as please Him.\end{Resp}
\begin{Resp}\Rsign And in all seasons let thy counsels be steadfastly in Him.\end{Resp}
\switchcolumn*

\LA
\LessonHead{Lectio III}
\switchcolumn
\EN
\LessonHead{Lesson III}
\switchcolumn*

\LA
\Cite{Tob 2:13-18}
\switchcolumn
\EN
\Cite{Tob 2:13-18}
\switchcolumn*

\LA
\noindent Nam, cum ab infántia sua semper Deum timúerit et mandáta ejus custodíerit, non est contristátus contra Deum quod plaga cæcitátis evénerit ei; sed immóbilis in Dei timóre permánsit agens grátias Deo ómnibus diébus vitæ suæ. Nam sicut beáto Job insultábant reges, ita isti paréntes et cognáti ejus irridébant vitam ejus, dicéntes: Ubi est spes tua, pro qua eleemósynas et sepultúras faciébas? Tobías vero increpábat eos, dicens: Nolíte ita loqui; quóniam fílii sanctórum sumus, et vitam illam exspectámus, quam Deus datúrus est his qui fidem suam nunquam mutant ab eo.\par
\switchcolumn
\EN
\noindent For whereas he had always feared God from his infancy, and kept his commandments, he repined not against God because the evil of blindness had befallen him, But continued immoveable in the fear of God, giving thanks to God all the days of his life. For as the kings insulted over holy Job: so his relations and kinsmen mocked at his life, saying: Where is thy hope, for which thou gavest alms, and buriedst the dead? But Tobias rebuked them, saying: Speak not so: For we are the children of the saints, and look for that life which God will give to those that never change their faith from him.\par
\switchcolumn*

\LA
\begin{Resp}\Rsign Memor esto, fili, quóniam páuperem vitam gérimus: \Star{} Habébis multa bona, si timúeris Deum.\end{Resp}
\begin{Resp}\Vsign In mente habéto eum, et cave nequándo prætermíttas præcépta ejus.\end{Resp}
\begin{Resp}\Rsign Habébis multa bona, si timúeris Deum.\end{Resp}
\begin{Resp}\Vsign Glória Patri, et Fílio, \Star{} et Spirítui Sancto.\end{Resp}
\begin{Resp}\Rsign Habébis multa bona, si timúeris Deum.\end{Resp}
\switchcolumn
\EN
\begin{Resp}\Rsign My son, remember that we have but a frail life. \Star{} If thou fear God thou shalt have great goods.\end{Resp}
\begin{Resp}\Vsign Be mindful of Him, and beware lest ever thou transgress His commandments.\end{Resp}
\begin{Resp}\Rsign If thou fear God thou shalt have great goods.\end{Resp}
\begin{Resp}\Vsign Glory be to the Father, and to the Son, \Star{} and to the Holy Ghost.\end{Resp}
\begin{Resp}\Rsign If thou fear God thou shalt have great goods.\end{Resp}
\switchcolumn*

\end{paracol}

\Office{Feria III infra Hebdomadam III Septembris}{Tuesday of the third week of September}{}
\addcontentsline{toc}{subsection}{Feria III infra Hebdomadam III Septembris \textit{— Tuesday of the third week of September}}
\begin{paracol}{2}
\LA
\LessonHead{Lectio I}
\switchcolumn
\EN
\LessonHead{Lesson I}
\switchcolumn*

\LA
\LessonTitle{De libro Tobíæ}
\switchcolumn
\EN
\LessonTitle{Lesson from the book of Tobias}
\switchcolumn*

\LA
\Cite{Tob 2:19-21}
\switchcolumn
\EN
\Cite{Tob 2:19-21}
\switchcolumn*

\LA
\noindent Anna vero uxor ejus ibat ad opus textrínum cotídie, et de labóre mánuum suárum victum, quem cónsequi póterat, deferébat. Unde factum est ut hædum caprárum accípiens detulísset domi; cujus cum vocem balántis vir ejus audísset, dixit: Vidéte, ne forte furtívus sit; réddite eum dóminis suis, quia non licet nobis aut édere ex furto áliquid aut contíngere.\par
\switchcolumn
\EN
\noindent Now Anna his wife went daily to weaving work, and she brought home what she could get for their living by the labour of her hands. Whereby it came to pass, that she received a young kid, and brought it home: And when her husband heard it bleating, he said: Take heed, lest perhaps it be stolen: restore ye it to its owners, for it is not lawful for us either to eat or to touch any thing that cometh by theft.\par
\switchcolumn*

\LA
\begin{Resp}\Rsign Sufficiébat nobis paupértas nostra, ut divítiæ computaréntur: numquam fuísset pecúnia ipsa, pro qua misísti fílium nostrum, \Star{} Báculum senectútis nostræ!\end{Resp}
\begin{Resp}\Vsign Heu me, fili mi, ut quid te mísimus peregrinári, lumen oculórum nostrórum?\end{Resp}
\begin{Resp}\Rsign Báculum senectútis nostræ!\end{Resp}
\switchcolumn
\EN
\begin{Resp}\Rsign Our poverty was enough for us, that it might have been accounted riches. O that the money had never been, for which thou hast sent away our son, \Star{} The staff of our old age.\end{Resp}
\begin{Resp}\Vsign Alas, my son, wherefore have we sent thee wandering, even thee, the light of our eyes?\end{Resp}
\begin{Resp}\Rsign The staff of our old age.\end{Resp}
\switchcolumn*

\LA
\LessonHead{Lectio II}
\switchcolumn
\EN
\LessonHead{Lesson II}
\switchcolumn*

\LA
\Cite{Tob 2:22-23; 3:1-3}
\switchcolumn
\EN
\Cite{Tob 2:22-23; 3:1-3}
\switchcolumn*

\LA
\noindent Ad hæc uxor ejus iráta respóndit: Maniféste vana facta est spes tua et eleemósynæ tuæ modo apparuérunt. Atque his et áliis hujuscémodi verbis exprobrábat ei. Tunc Tobías ingémuit et cœpit oráre cum lácrimis, dicens: Justus es, Dómine, et ómnia judícia tua justa sunt, et omnes viæ tuæ misericórdia, et véritas, et judícium. Et nunc, Dómine, memor esto mei, et ne vindíctam sumas de peccátis meis, neque reminiscáris delícta mea, vel paréntum meórum.\par
\switchcolumn
\EN
\noindent At these words his wife being angry answered: It is evident thy hope is come to nothing, and thy alms now appear. And with these, and other such like words she upbraided him. Then Tobias sighed, and began to pray with tears, Saying: Thou art just, O Lord, and all thy judgments are just, and all thy ways mercy, and truth, and judgment: And now, O Lord, think of me, and take not revenge of my sins, neither remember my offenses, nor those of my parents.\par
\switchcolumn*

\LA
\begin{Resp}\Rsign Benedícite Deum cæli et coram ómnibus vivéntibus confitémini ei, \Star{} Quia fecit vobíscum misericórdiam suam.\end{Resp}
\begin{Resp}\Vsign Ipsum benedícite et cantáte illi: et enarráte ómnia mirabília ejus.\end{Resp}
\begin{Resp}\Rsign Quia fecit vobíscum misericórdiam suam.\end{Resp}
\switchcolumn
\EN
\begin{Resp}\Rsign Bless the God of heaven, and confess Him before all living: \Star{} For He hath had mercy upon you.\end{Resp}
\begin{Resp}\Vsign Bless Him, and sing praises unto Him, and tell of all His marvelous works.\end{Resp}
\begin{Resp}\Rsign For He hath had mercy upon you.\end{Resp}
\switchcolumn*

\LA
\LessonHead{Lectio III}
\switchcolumn
\EN
\LessonHead{Lesson III}
\switchcolumn*

\LA
\Cite{Tob 3:4-6}
\switchcolumn
\EN
\Cite{Tob 3:4-6}
\switchcolumn*

\LA
\noindent Quóniam non obedívimus præcéptis tuis, ídeo tráditi sumus in direptiónem, et captivitátem, et mortem, et in fábulam, et in impropérium ómnibus natiónibus, in quibus dispersísti nos. Et nunc, Dómine, magna judícia tua, quia non égimus secúndum præcépta tua et non ambulávimus sincériter coram te. Et nunc, Dómine, secúndum voluntátem tuam fac mecum, et prǽcipe in pace récipi spíritum meum; éxpedit enim mihi mori magis quam vívere.\par
\switchcolumn
\EN
\noindent For we have not obeyed thy commandments, therefore are we delivered to spoil and to captivity, and death, and are made a fable, and a reproach to all nations, amongst which thou hast scattered us. And now, O Lord, great are thy judgments, because we have not done according to thy precepts, and have not walked sincerely before thee: And now, O Lord, do with me according to thy will, and command my spirit to be received in peace: for it is better for me to die, than to live.\par
\switchcolumn*

\LA
\begin{Resp}\Rsign Tempus est ut revértar ad eum qui misit me; \Star{} Vos autem benedícite Deum, et enarráte ómnia mirabília ejus.\end{Resp}
\begin{Resp}\Vsign Confitémini ei coram ómnibus vivéntibus, quia fecit vobíscum misericórdiam suam.\end{Resp}
\begin{Resp}\Rsign Vos autem benedícite Deum, et enarráte ómnia mirabília ejus.\end{Resp}
\begin{Resp}\Vsign Glória Patri, et Fílio, \Star{} et Spirítui Sancto.\end{Resp}
\begin{Resp}\Rsign Vos autem benedícite Deum, et enarráte ómnia mirabília ejus.\end{Resp}
\switchcolumn
\EN
\begin{Resp}\Rsign It is time for me to return unto Him That sent me; \Star{} But bless ye God, and tell of all His marvelous works.\end{Resp}
\begin{Resp}\Vsign Confess Him before all living, for He hath had mercy upon you.\end{Resp}
\begin{Resp}\Rsign But bless ye God, and tell of all His marvelous works.\end{Resp}
\begin{Resp}\Vsign Glory be to the Father, and to the Son, \Star{} and to the Holy Ghost.\end{Resp}
\begin{Resp}\Rsign But bless ye God, and tell of all His marvelous works.\end{Resp}
\switchcolumn*

\end{paracol}

\Office{Feria IV infra Hebdomadam III Septembris}{Wednesday of the third week of September}{Feria major}
\addcontentsline{toc}{subsection}{Feria IV infra Hebdomadam III Septembris \textit{— Wednesday of the third week of September}}
\begin{paracol}{2}
\LA
\LessonHead{Lectio I}
\switchcolumn
\EN
\LessonHead{Lesson I}
\switchcolumn*

\LA
\Rubric{Commemoratio Feriae}
\switchcolumn
\EN
\Rubric{Commemoration of the day}
\switchcolumn*

\LA
\LessonTitle{Léctio sancti Evangélii secúndum Marcum}
\switchcolumn
\EN
\LessonTitle{Continuation of the Holy Gospel according to Mark}
\switchcolumn*

\LA
\Cite{Marc 9:16-28}
\switchcolumn
\EN
\Cite{Mark 9:16-28}
\switchcolumn*

\LA
\noindent In illo témpore: Respóndens unus de turba dixit ad Jesum: Magíster, áttuli fílium meum ad te habéntem spíritum mutum. Et réliqua.\par
\switchcolumn
\EN
\noindent At that time, one of the multitude answered and said unto Jesus: Master, I have brought unto thee my son, which hath a dumb spirit. And so on.\par
\switchcolumn*

\LA
\LessonTitle{Homilía sancti Bedæ Venerábilis Presbýteri}
\switchcolumn
\EN
\LessonTitle{Homily by the Venerable Bede, Priest (at Jarrow.)}
\switchcolumn*

\LA
\Source{Liber 3, cap. 38 in Marcum 9}
\switchcolumn
\EN
\Source{Bk. III, ch. 38, on Mark IX}
\switchcolumn*

\LA
\noindent Dæmoníacum hunc, quem descéndens de monte Dóminus sanávit, Marcus quidem surdum mutúmque, Matthǽus vero lunáticum fuísse commémorat. Signíficat autem eos, de quibus scriptum est: Stultus ut luna mutátur. Qui nunquam in eódem statu permanéntes, nunc ad illa vítia mutáti, crescunt atque decréscunt. Qui muti sunt, non confiténdo fidem; surdi, nec ipsum aliquátenus veritátis audiéndo sermónem. Spumant autem, cum stultítia tabéscunt; stultórum namque et languéntium atque hébetum est, spumas salivárum ex ore dimíttere. Strident déntibus, cum iracúndiæ furóre flamméscunt: aréscunt, cum ótio torpénte languéscunt, et nulla virtútis indústria confortáti enérviter vivunt.\par
\switchcolumn
\EN
\noindent Concerning this possessed person whom the Lord healed, after that He was come down from the mount, Mark saith that he was deaf and dumb, and Matthew (xvii. 15) that he was lunatic. He was a figure of them of whom it is said: A fool changeth as the moon, (Ecclus. xxvii. 12.) These are they who continue never in one stay, but change now to one sin, and now to another, waxing and waning dumb, in that they confess not the faith deaf, in that they have no ears for the word of truth. They foam at the mouth also, and pine away with folly. For it is the way with idiots, and swooners, and stupified, to foam their spittle out at their mouths. They gnash their teeth when they are inflamed with the heat of passion. They wither up in the paralysis of sloth and live nerveless lives unbraced by any strong exercise.\par
\switchcolumn*

\LA
\TeDeum{Te Deum}
\switchcolumn
\EN
\TeDeum{Te Deum}
\switchcolumn*

\end{paracol}

\Office{Feria V infra Hebdomadam III Septembris}{Thursday of the third week of September}{}
\addcontentsline{toc}{subsection}{Feria V infra Hebdomadam III Septembris \textit{— Thursday of the third week of September}}
\begin{paracol}{2}
\LA
\LessonHead{Lectio II}
\switchcolumn
\EN
\LessonHead{Lesson II}
\switchcolumn*

\LA
\Cite{Tob 12:5-10}
\switchcolumn
\EN
\Cite{Tob 12:5-10}
\switchcolumn*

\LA
\noindent Et vocántes eum, pater scílicet et fílius, tulérunt eum in partem; et rogáre cœpérunt ut dignarétur dimídiam partem ómnium quæ attúlerant, accéptam habére. Tunc dixit eis occúlte: Benedícite Deum cæli et coram ómnibus vivéntibus confitémini ei, quia fecit vobíscum misericórdiam suam. Etenim sacraméntum regis abscóndere bonum est, ópera autem Dei reveláre et confitéri honoríficum est. Bona est orátio cum jejúnio, et eleemósyna magis quam thesáuros auri recóndere; quóniam eleemósyna a morte líberat, et ipsa est quæ purgat peccáta, et facit inveníre misericórdiam et vitam ætérnam. Qui autem fáciunt peccátum et iniquitátem hostes sunt ánimæ suæ.\par
\switchcolumn
\EN
\noindent So the father and the son, calling him, took him aside: and began to desire him that he would vouchsafe to accept of half of all things that they had brought. Then he said to them secretly: Bless ye the God of heaven, give glory to him in the sight of all that live, because he hath shewn his mercy to you. For it is good to hide the secret of a king: but honourable to reveal and confess the works of God. Prayer is good with fasting and alms more than to lay up treasures of gold: For alms delivereth from death, and the same is that which purgeth away sins, and maketh to find mercy and life everlasting. But they that commit sin and iniquity, are enemies to their own soul.\par
\switchcolumn*

\LA
\begin{Resp}\Rsign Omni témpore bénedic Deum, et pete ab eo ut vias tuas dírigat, \Star{} Et in omni témpore consília tua in ipso permáneant.\end{Resp}
\begin{Resp}\Vsign Inquíre ut fácias quæ plácita sunt illi in veritáte, et in tota virtúte tua.\end{Resp}
\begin{Resp}\Rsign Et in omni témpore consília tua in ipso permáneant.\end{Resp}
\switchcolumn
\EN
\begin{Resp}\Rsign In all seasons bless God, and ask of Him to order thy goings; \Star{} And in all seasons let thy counsels be steadfastly in Him.\end{Resp}
\begin{Resp}\Vsign Seek faithfully and with all thy strength to do such things as please Him.\end{Resp}
\begin{Resp}\Rsign And in all seasons let thy counsels be steadfastly in Him.\end{Resp}
\switchcolumn*

\LA
\LessonHead{Lectio III}
\switchcolumn
\EN
\LessonHead{Lesson III}
\switchcolumn*

\LA
\Cite{Tob 12:11-17}
\switchcolumn
\EN
\Cite{Tob 12:11-17}
\switchcolumn*

\LA
\noindent Manifésto ergo vobis veritátem et non abscóndam a vobis occúltum sermónem. Quando orábas cum lácrimis, et sepeliébas mórtuos, et derelinquébas prándium tuum, et mórtuos abscondébas per diem in domo tua, et nocte sepeliébas eos, ego óbtuli oratiónem tuam Dómino. Et quia accéptus eras Deo, necésse fuit ut tentátio probáret te. Et nunc misit me Dóminus ut curárem te et Saram uxórem fílii tui a dæmónio liberárem; ego enim sum Ráphaël Angelus, unus ex septem qui astámus ante Dóminum. Cumque hæc audíssent, turbáti sunt et treméntes cecidérunt super terram in fáciem suam. Dixítque eis Angelus: Pax vobis, nolíte timére.\par
\switchcolumn
\EN
\noindent I discover then the truth unto you, and I will not hide the secret from you. When thou didst pray with tears, and didst bury the dead, and didst leave thy dinner, and hide the dead by day in thy house, and bury them by night, I offered thy prayer to the Lord. And because thou wast acceptable to God, it was necessary that temptation should prove thee. And now the Lord hath sent me to heal thee, and to deliver Sara thy son's wife from the devil. For I am the angel Raphael, one of the seven, who stand before the Lord. And when they had heard these things, they were troubled, and being seized with fear they fell upon the ground on their face. And the angel said to them: Peace be to you, fear not.\par
\switchcolumn*

\LA
\begin{Resp}\Rsign Memor esto, fili, quóniam páuperem vitam gérimus: \Star{} Habébis multa bona, si timúeris Deum.\end{Resp}
\begin{Resp}\Vsign In mente habéto eum, et cave nequándo prætermíttas præcépta ejus.\end{Resp}
\begin{Resp}\Rsign Habébis multa bona, si timúeris Deum.\end{Resp}
\begin{Resp}\Vsign Glória Patri, et Fílio, \Star{} et Spirítui Sancto.\end{Resp}
\begin{Resp}\Rsign Habébis multa bona, si timúeris Deum.\end{Resp}
\switchcolumn
\EN
\begin{Resp}\Rsign My son, remember that we have but a frail life. \Star{} If thou fear God thou shalt have great goods.\end{Resp}
\begin{Resp}\Vsign Be mindful of Him, and beware lest ever thou transgress His commandments.\end{Resp}
\begin{Resp}\Rsign If thou fear God thou shalt have great goods.\end{Resp}
\begin{Resp}\Vsign Glory be to the Father, and to the Son, \Star{} and to the Holy Ghost.\end{Resp}
\begin{Resp}\Rsign If thou fear God thou shalt have great goods.\end{Resp}
\switchcolumn*

\end{paracol}

\Office{Feria VI infra Hebdomadam III Septembris}{Friday of the third week of September}{Feria major}
\addcontentsline{toc}{subsection}{Feria VI infra Hebdomadam III Septembris \textit{— Friday of the third week of September}}
\begin{paracol}{2}
\LA
\LessonHead{Lectio I}
\switchcolumn
\EN
\LessonHead{Lesson I}
\switchcolumn*

\LA
\Rubric{Commemoratio Feriae}
\switchcolumn
\EN
\Rubric{Commemoration of the day}
\switchcolumn*

\LA
\LessonTitle{Léctio sancti Evangélii secúndum Lucam}
\switchcolumn
\EN
\LessonTitle{Continuation of the Holy Gospel according to Luke}
\switchcolumn*

\LA
\Cite{Luc 7:36-50}
\switchcolumn
\EN
\Cite{Luke 7:36-50}
\switchcolumn*

\LA
\noindent In illo témpore: Rogábat Jesum quidam de pharisǽis, ut manducáret cum illo. Et ingréssus domum pharisǽi discúbuit. Et réliqua.\par
\switchcolumn
\EN
\noindent At that time, one of the Pharisees desired that Jesus would eat with him. And He went into the Pharisee's house, and sat down to meat. And so on.\par
\switchcolumn*

\LA
\LessonTitle{Homilía sancti Gregórii Papæ}
\switchcolumn
\EN
\LessonTitle{Homily by Pope St. Gregory (the Great.)}
\switchcolumn*

\LA
\Source{Homilia 33 in Evangelia, post initium}
\switchcolumn
\EN
\Source{33rd on the Gospels.}
\switchcolumn*

\LA
\noindent Quem pharisǽus de falsa justítia præsúmens, nisi Judáicum pópulum; quem peccátrix múlier ad vestígia Dómini véniens et plorans, nisi convérsam gentilitátem desígnat? Quæ cum alabástro venit, unguéntum fudit, retro secus pedes Dómini stetit, lácrimis pedes rigávit, capíllis tersit, eosdémque quos infundébat et tergébat pedes osculári non désiit. Nos ergo, nos illa múlier expréssit, si toto corde ad Dóminum post peccáta redeámus, si ejus pœniténtiæ luctus imitémur. Quid namque unguénto, nisi bonæ odor opiniónis exprímitur? Unde et Paulus dicit: Christi bonus odor sumus Deo in omni loco.\par
\switchcolumn
\EN
\noindent Of what is the Pharisee that was exalted by self-righteousness a type, but of the Jewish people? And of what the woman which was a sinner and came and wept at the Lord's feet, but of the conversion of the Gentiles? She brought an alabaster box of ointment, and stood at His feet behind Him weeping, and began to wash His Feet with tears, and did wipe them with the hairs of her head, and kissed His Feet, and anointed them with the ointment. Of us, therefore, even of us, was that woman a type, if after our sins we turn unto the Lord with all our heart, and imitate the example of her repentant grief. And of what is the ointment a type, but of the sweet savour of a good reputation? Whence also Paul saith: In every place we are unto God a sweet savour of Christ. (2 Cor. ii. 15.)\par
\switchcolumn*

\LA
\TeDeum{Te Deum}
\switchcolumn
\EN
\TeDeum{Te Deum}
\switchcolumn*

\end{paracol}

\Office{Sabbato infra Hebdomadam III Septembris}{Saturday of the third week of September}{Feria major}
\addcontentsline{toc}{subsection}{Sabbato infra Hebdomadam III Septembris \textit{— Saturday of the third week of September}}
\begin{paracol}{2}
\LA
\LessonHead{Lectio I}
\switchcolumn
\EN
\LessonHead{Lesson I}
\switchcolumn*

\LA
\Rubric{Commemoratio Feriae}
\switchcolumn
\EN
\Rubric{Commemoration of the day}
\switchcolumn*

\LA
\LessonTitle{Léctio sancti Evangélii secúndum Lucam}
\switchcolumn
\EN
\LessonTitle{Continuation of the Holy Gospel according to Luke}
\switchcolumn*

\LA
\Cite{Luc 13:6-17}
\switchcolumn
\EN
\Cite{Luke 13:6-17}
\switchcolumn*

\LA
\noindent In illo témpore: Dicébat Jesus turbis hanc similitúdinem: Arborem fici habébat quidam plantátam in vínea sua et venit quærens fructum in illa et non invénit. Et réliqua.\par
\switchcolumn
\EN
\noindent At that time, Jesus spoke unto the multitudes this parable: A certain man had a fig-tree planted in his vineyard, and he came and sought fruit thereon, and found none. And so on.\par
\switchcolumn*

\LA
\LessonTitle{Homilía sancti Gregórii Papæ}
\switchcolumn
\EN
\LessonTitle{Homily by Pope St. Gregory (the Great.)}
\switchcolumn*

\LA
\Source{Homilia 31 in Evangelia}
\switchcolumn
\EN
\Source{31st on the Gospels.}
\switchcolumn*

\LA
\noindent Dóminus ac Redémptor noster per Evangélium suum aliquándo verbis, aliquándo rebus lóquitur: aliquándo áliud verbis, atque áliud rebus; aliquándo autem hoc verbis, quod rebus. Duas étenim res ex Evangélio, fratres, audístis: ficúlneam infructuósam, et mulíerem curvam; et utríque rei est píetas impénsa. Illud autem dixit per similitúdinem: istud egit per exhibitiónem. Sed hoc signíficat ficúlnea infructuósa, quod múlier inclináta; et hoc ficúlnea reserváta, quod múlier erécta.\par
\switchcolumn
\EN
\noindent Our Lord and Redeemer speaketh unto us sometimes by words, and sometimes by deeds, sometimes one thing by words, and another by deeds, and sometimes the same thing both by word and deed. In the portion of the Gospel which hath this day been read, ye have heard, my brethren, two things, the parable of the fig-tree and the history of the woman which was bowed together. In both is a manifestation of the Lord's mercy, but in the one by a parable, in the other by an example. But the barren fig-tree signifieth the same thing as doth the woman bowed together, and the patience shown to the fig-tree the same thing as doth the healing of the woman bowed together.\par
\switchcolumn*

\LA
\TeDeum{Te Deum}
\switchcolumn
\EN
\TeDeum{Te Deum}
\switchcolumn*

\end{paracol}

\Office{Dominica IV Septembris}{Fourth Sunday of September}{}
\addcontentsline{toc}{subsection}{Dominica IV Septembris \textit{— Fourth Sunday of September}}
\begin{paracol}{2}
\LA
\LessonHead{Lectio I}
\switchcolumn
\EN
\LessonHead{Lesson I}
\switchcolumn*

\LA
\LessonTitle{Incipit liber Judith}
\switchcolumn
\EN
\LessonTitle{Lesson from the book of Judith}
\switchcolumn*

\LA
\Cite{Jdt 1:1-4}
\switchcolumn
\EN
\Cite{Jdt 1:1-4}
\switchcolumn*

\LA
\noindent Arpháxad ítaque rex Medórum subjugáverat multas gentes império suo, et ipse ædificávit civitátem potentíssimam, quam appellávit Ecbátanis, ex lapídibus quadrátis et sectis. Fecit muros ejus in altitúdinem cubitórum septuagínta et in latitúdinem cubitórum trigínta; turres vero ejus pósuit in altitúdinem cubitórum centum. Per quadrum vero eárum latus utrúmque vicenórum pedum spátio tendebátur; posuítque portas ejus in altitúdinem túrrium; et gloriabátur quasi potens in poténtia exércitus sui et in glória quadrigárum suárum.\par
\switchcolumn
\EN
\noindent Now Arphaxad king of the Medes had brought many nations under his dominions, and he built a very strong city, which he called Ecbatana, Of stones squared and hewed: he made the walls thereof seventy cubits broad, and thirty cubits high, and the towers thereof he made a hundred cubits high. But on the square of them, each side was extended the space of twenty feet. And he made the gates thereof according to the height of the towers: And he gloried as a mighty one in the force of his army and in the glory of his chariots.\par
\switchcolumn*

\LA
\begin{Resp}\Rsign Adonái, Dómine, Deus magne et mirábilis, qui dedísti salútem in manu féminæ, \Star{} Exáudi preces servórum tuórum.\end{Resp}
\begin{Resp}\Vsign Benedíctus es, Dómine, qui non derelínquis præsuméntes de te, et de sua virtúte gloriántes humílias.\end{Resp}
\begin{Resp}\Rsign Exáudi preces servórum tuórum.\end{Resp}
\switchcolumn
\EN
\begin{Resp}\Rsign O Adonai, O Lord God, Thou art great and glorious. Who hast given salvation into the hand of a woman: \Star{} Graciously hear the prayers of thy servants.\end{Resp}
\begin{Resp}\Vsign Blessed art Thou, O Lord, Who failest none that put their trust in thee, and humblest such as boast themselves in their own strength.\end{Resp}
\begin{Resp}\Rsign Graciously hear the prayers of thy servants.\end{Resp}
\switchcolumn*

\LA
\LessonHead{Lectio II}
\switchcolumn
\EN
\LessonHead{Lesson II}
\switchcolumn*

\LA
\Cite{Jdt 1:5-9}
\switchcolumn
\EN
\Cite{Jdt 1:5-9}
\switchcolumn*

\LA
\noindent Anno ígitur duodécimo regni sui Nabuchodónosor rex Assyriórum, qui regnábat in Nínive civitáte magna, pugnávit contra Arpháxad et obtínuit eum in campo magno qui appellátur Rágau, circa Euphráten et Tigrin et Jádason, in campo Erioch regis Elicórum. Tunc exaltátum est regnum Nabuchodónosor, et cor ejus elevátum est; et misit ad omnes, qui habitábant in Cilícia et Damásco et Líbano, et ad gentes quæ sunt in Carmélo et Cedar et inhabitántes Galilǽam in campo magno Esdrelon, et ad omnes qui erant in Samaría et trans flumen Jordánem usque ad Jerúsalem, et omnem terram Jesse, quoúsque perveniátur ad términos Æthiópiæ.\par
\switchcolumn
\EN
\noindent Now in the twelfth year of his reign, Nabuchodonosor king of the Assyrians, who reigned in Ninive the great city, fought against Arphaxad and overcame him, In the great plain which is called Ragua, about the Euphrates, and the Tigris, and the Jadason, in the plain of Erioch the king of the Elicians. Then was the kingdom of Nabuchodonosor exalted, and his heart was elevated: and he sent to all that dwelt in Cilicia and Damascus, and Libanus, And to the nations that are in Carmelus, and Cedar, and to the inhabitants of Galilee in the great plain of Asdrelon, And to all that were in Samaria, and beyond the river Jordan even to Jerusalem, and all the land of Jesse till you come to the borders of Ethiopia.\par
\switchcolumn*

\LA
\begin{Resp}\Rsign Tribulatiónes civitátum audívimus, quas passæ sunt, et defécimus: timor et hebetúdo mentis cécidit super nos et super líberos nostros: ipsi montes nolunt recípere fugam nostram: \Star{} Dómine, miserére.\end{Resp}
\begin{Resp}\Vsign Peccávimus cum pátribus nostris, injúste égimus, iniquitátem fécimus.\end{Resp}
\begin{Resp}\Rsign Dómine, miserére.\end{Resp}
\switchcolumn
\EN
\begin{Resp}\Rsign We have heard of the tribulation of those cities, which they have suffered, and we have fainted. Fear and confusion of mind are fallen upon us. Even the mountains will not give us a refuge. \Star{} Lord, have mercy.\end{Resp}
\begin{Resp}\Vsign We have sinned like our forefathers, we have done unjustly, and wrought iniquity.\end{Resp}
\begin{Resp}\Rsign Lord, have mercy.\end{Resp}
\switchcolumn*

\LA
\LessonHead{Lectio III}
\switchcolumn
\EN
\LessonHead{Lesson III}
\switchcolumn*

\LA
\Cite{Jdt 1:10-12; 2:1-3}
\switchcolumn
\EN
\Cite{Jdt 1:10-12; 2:1-3}
\switchcolumn*

\LA
\noindent Ad hos omnes misit núntios Nabuchodónosor rex Assyriórum; qui omnes uno ánimo contradixérunt et remisérunt eos vácuos et sine honóre abjecérunt. Tunc indignátus Nabuchodónosor rex advérsus omnem terram illam, jurávit per thronum et regnum suum, quod defénderet se de ómnibus regiónibus his. Anno tertiodécimo Nabuchodónosor regis, vigésima et secúnda die mensis primi factum est verbum in domo Nabuchodónosor regis Assyriórum, ut defénderet se. Vocavítque omnes majóres natu, omnésque duces et bellatóres suos, et hábuit cum eis mystérium consílii sui: dixítque cogitatiónem suam in eo esse, ut omnem terram suo subjugáret império.\par
\switchcolumn
\EN
\noindent To all these Nabuchodonosor king of the Assyrians, sent messengers: But they all with one mind refused, and sent them back empty, and rejected them without honour. Then king Nabuchodonosor being angry against all that land, swore by his throne and kingdom that he would revenge himself of all those countries. In the thirteenth year of the reign of Nabuchodonosor, the two and twentieth day of the first month, the word was given out in the house of Nabuchodonosor king of the Assyrians, that he would revenge himself. And he called all the ancients, and all the governors, and his officers of war, and communicated to them the secret of his counsel: And he said that his thoughts were to bring all the earth under his empire.\par
\switchcolumn*

\LA
\begin{Resp}\Rsign Benedícat te Dóminus in virtúte sua, qui per te ad níhilum redégit inimícos nostros: \Star{} Ut non defíciat laus tua de ore hóminum.\end{Resp}
\begin{Resp}\Vsign Benedíctus Dóminus qui creávit cælum et terram; quia hódie nomen tuum ita magnificávit.\end{Resp}
\begin{Resp}\Rsign Ut non defíciat laus tua de ore hóminum.\end{Resp}
\begin{Resp}\Vsign Glória Patri, et Fílio, \Star{} et Spirítui Sancto.\end{Resp}
\begin{Resp}\Rsign Ut non defíciat laus tua de ore hóminum.\end{Resp}
\switchcolumn
\EN
\begin{Resp}\Rsign The Lord bless thee by His power, Who hath brought our enemies to nought through thee. \Star{} And may the praise of thee never fail from the mouth of men.\end{Resp}
\begin{Resp}\Vsign Blessed be the Lord, Who hath created the heaven and the earth, because that He hath so glorified thy name this day.\end{Resp}
\begin{Resp}\Rsign And may the praise of thee never fail from the mouth of men.\end{Resp}
\begin{Resp}\Vsign Glory be to the Father, and to the Son, \Star{} and to the Holy Ghost.\end{Resp}
\begin{Resp}\Rsign And may the praise of thee never fail from the mouth of men.\end{Resp}
\switchcolumn*

\LA
\LessonHead{Lectio IV}
\switchcolumn
\EN
\LessonHead{Lesson IV}
\switchcolumn*

\LA
\LessonTitle{Ex libro sancti Ambrósii Epíscopi de Elía et jejúnio}
\switchcolumn
\EN
\LessonTitle{From the Book upon Elijah and Fasting, written by St. Ambrose, Bishop (of Milan.)}
\switchcolumn*

\LA
\Source{Lectio iv, Cap. 9}
\switchcolumn
\EN
\Source{ch. 9}
\switchcolumn*

\LA
\noindent Poténtes vinum prohibéntur bíbere, ne, cum bíberint, obliviscántur sapiéntiam. Dénique bibébant vinum in ebrietáte poténtes, qui Holoférni príncipi milítiæ regis Assyriórum se trádere gestiébant; sed non bibébant fémina Judith, jejúnans ómnibus diébus viduitátis suæ, præter festórum diérum solemnitátes. His armis muníta procéssit et omnem Assyriórum circumvénit exércitum. Sóbrii vigóre consílii ábstulit Holoférnis caput, servávit pudicítiam, victóriam reportávit.\par
\switchcolumn
\EN
\noindent It is not for kings to drink wine, nor for princes strong drink, lest they drink and forget the law. (Prov. xxxi. 5.) The rulers drank wine even unto drunkenness, who planned to deliver themselves into the hand of Holofernes, captain of the host of the King of the Assyrians but the woman Judith drank not, who fasted all the days of her widowhood, saving the solemn Feast-days. She went forth in the harness of this abstinence, and over - reached the whole army of the Assyrians. By the clear thought of her soberness she took away the head of Holofernes, kept her chastity, and carried off the victory.\par
\switchcolumn*

\LA
\begin{Resp}\Rsign Nos álium Deum nescímus præter Dóminum, in quo sperámus: \Star{} Qui non déspicit nos, nec ámovet salútem suam a génere nostro.\end{Resp}
\begin{Resp}\Vsign Indulgéntiam ipsíus fusis lácrimis postulémus, et humiliémus illi ánimas nostras.\end{Resp}
\begin{Resp}\Rsign Qui non déspicit nos, nec ámovet salútem suam a génere nostro.\end{Resp}
\switchcolumn
\EN
\begin{Resp}\Rsign We know no strange God before the Lord. In Him we trust. \Star{} He despiseth us not, neither putteth He away His salvation from our nation.\end{Resp}
\begin{Resp}\Vsign His mercy let us seek with tears, and humble our souls before Him.\end{Resp}
\begin{Resp}\Rsign He despiseth us not, neither putteth He away His salvation from our nation.\end{Resp}
\switchcolumn*

\LA
\LessonHead{Lectio V}
\switchcolumn
\EN
\LessonHead{Lesson V}
\switchcolumn*

\LA
\noindent Hæc enim succíncta jejúnio in castris prætendébat aliénis; ille vino sepúltus jacébat, ut ictum vúlneris sentíre non posset. Ítaque uníus mulíeris jejúnium innúmeros stravit exércitus Assyriórum. Esther quoque púlchrior facta est jejúnio; Dóminus enim grátiam sóbriæ mentis augébat. Omne genus suum, id est, totum pópulum Judæórum a persecutiónis acerbitáte liberávit, ita ut regem sibi fáceret esse subjéctum.\par
\switchcolumn
\EN
\noindent Armed with fasting, she entered the camp of the strangers he lay soaked in wine, so that he could not feel the blow that slew him. And thus the fast of one woman overthrew the countless armies of the Assyrians. Esther also became fairer by fasting for the Lord gave favour unto her for her soberness. She delivered all her nation, that is, the whole people of the Jews, from the fierceness of persecution, so that she brought down the King himself under her will.\par
\switchcolumn*

\LA
\begin{Resp}\Rsign Dominátor, Dómine, cælórum et terræ, Creátor aquárum, Rex univérsæ creatúræ: \Star{} Exáudi oratiónem servórum tuórum.\end{Resp}
\begin{Resp}\Vsign Tu, Dómine, cui humílium semper et mansuetórum plácuit deprecátio.\end{Resp}
\begin{Resp}\Rsign Exáudi oratiónem servórum tuórum.\end{Resp}
\switchcolumn
\EN
\begin{Resp}\Rsign O Lord, Ruler of the heavens and of the earth, Maker of the waters, King of every creature, \Star{} Graciously hear the prayer of thy servants.\end{Resp}
\begin{Resp}\Vsign Thou, O Lord, unto Whom the supplications of the humble and meek are alway well-pleasing\end{Resp}
\begin{Resp}\Rsign Graciously hear the prayer of thy servants.\end{Resp}
\switchcolumn*

\LA
\LessonHead{Lectio VI}
\switchcolumn
\EN
\LessonHead{Lesson VI}
\switchcolumn*

\LA
\noindent Itaque illa, quæ tríduo jejunávit contínuo, et corpus suum aqua lavit, plus plácuit, et vindíctam rétulit. Aman autem, dum se regáli jactat convívio, inter ipsa vina pœnam suæ ebrietátis exsólvit. Est ergo jejúnium reconciliatiónis sacrifícium, virtútis increméntum, quod fecit étiam féminas fortióres augménto grátiæ. Jejúnium nescit fæneratórem, non sortem fǽnoris novit: non rédolet usúras mensa jejunántium. Etiam ipsis jejúnium convíviis dat grátiam: dulcióres post famem épulæ fiunt, quæ assiduitáte fastídio sunt, et diutúrna continuatióne viléscunt. Condiméntum cibi jejúnium est: quanto avídior appeténtia, tanto esca jucúndior.\par
\switchcolumn
\EN
\noindent Thus also (Esther) who fasted three days, and washed her body with water, found greater favour, and obtained vengeance, whereas Haman, who boasted himself at the King's table, paid the penalty of his drunkenness, even while yet he was in his cups. Fasting, therefore, is a sacrifice of reconciliation, a means of strength, whereby in the might of grace, women wax manful. Fasting knoweth not usury, nor the gain of the usurer the faster's table smelleth not of usury, but the fast itself giveth favour to them that sit at meat. A banquet is all the pleasanter after hunger, whereas by constant use it becometh unattractive, and when it is long carried on cometh to be lightly esteemed. Fasting is a good sauce for meat. The keener the appetite, the more toothsome the food.\par
\switchcolumn*

\LA
\begin{Resp}\Rsign Dómine Deus, qui cónteris bella ab inítio, éleva brácchium tuum super gentes, quæ cógitant servis tuis mala: \Star{} Et déxtera tua glorificétur in nobis.\end{Resp}
\begin{Resp}\Vsign Allíde virtútem eórum in virtúte tua; cadat robur eórum in iracúndia tua.\end{Resp}
\begin{Resp}\Rsign Et déxtera tua glorificétur in nobis.\end{Resp}
\begin{Resp}\Vsign Glória Patri, et Fílio, \Star{} et Spirítui Sancto.\end{Resp}
\begin{Resp}\Rsign Et déxtera tua glorificétur in nobis.\end{Resp}
\switchcolumn
\EN
\begin{Resp}\Rsign O Lord God, That breakest the battles from of old, lift up thine arm against the Gentiles, that devise evil against thy servants. \Star{} And let thy right hand be glorified in us.\end{Resp}
\begin{Resp}\Vsign Throw down their strength in thy power, and bring down their force in thy wrath.\end{Resp}
\begin{Resp}\Rsign And let thy right hand be glorified in us.\end{Resp}
\begin{Resp}\Vsign Glory be to the Father, and to the Son, \Star{} and to the Holy Ghost.\end{Resp}
\begin{Resp}\Rsign And let thy right hand be glorified in us.\end{Resp}
\switchcolumn*

\end{paracol}

\Office{Feria II infra Hebdomadam IV Septembris}{Monday of the fourth week of September}{}
\addcontentsline{toc}{subsection}{Feria II infra Hebdomadam IV Septembris \textit{— Monday of the fourth week of September}}
\begin{paracol}{2}
\LA
\LessonHead{Lectio I}
\switchcolumn
\EN
\LessonHead{Lesson I}
\switchcolumn*

\LA
\LessonTitle{De libro Judith}
\switchcolumn
\EN
\LessonTitle{Lesson from the book of Judith}
\switchcolumn*

\LA
\Cite{Jdt 4:1-4}
\switchcolumn
\EN
\Cite{Jdt 4:1-4}
\switchcolumn*

\LA
\noindent Fílii Israël, qui habitábant in terra Juda, timuérunt valde a fácie Holoférnis; tremor et horror invásit sensus eórum, ne hoc fáceret Jerúsalem et templo Dómini, quod fécerat céteris civitátibus et templis eárum. Et misérunt in omnem Samaríam per circúitum usque Jéricho et præoccupavérunt omnes vértices móntium et muris circumdedérunt vicos suos et congregavérunt fruménta in præparatiónem pugnæ.\par
\switchcolumn
\EN
\noindent Then the children of Israel, who dwelt in the land of Juda, hearing these things, were exceedingly afraid of him. Dread and horror seized upon their minds, lest he should do the same to Jerusalem and to the temple of the Lord, that he had done to other cities and their temples. And they sent into all Samaria round about, as far as Jericho, and seized upon all the tops of the mountains: And they compassed their towns with walls, and gathered together corn for provision for war.\par
\switchcolumn*

\LA
\begin{Resp}\Rsign Adonái, Dómine, Deus magne et mirábilis, qui dedísti salútem in manu féminæ, \Star{} Exáudi preces servórum tuórum.\end{Resp}
\begin{Resp}\Vsign Benedíctus es, Dómine, qui non derelínquis præsuméntes de te, et de sua virtúte gloriántes humílias.\end{Resp}
\begin{Resp}\Rsign Exáudi preces servórum tuórum.\end{Resp}
\switchcolumn
\EN
\begin{Resp}\Rsign O Adonai, O Lord God, Thou art great and glorious. Who hast given salvation into the hand of a woman: \Star{} Graciously hear the prayers of thy servants.\end{Resp}
\begin{Resp}\Vsign Blessed art Thou, O Lord, Who failest none that put their trust in thee, and humblest such as boast themselves in their own strength.\end{Resp}
\begin{Resp}\Rsign Graciously hear the prayers of thy servants.\end{Resp}
\switchcolumn*

\LA
\LessonHead{Lectio II}
\switchcolumn
\EN
\LessonHead{Lesson II}
\switchcolumn*

\LA
\Cite{Jdt 4:5-8}
\switchcolumn
\EN
\Cite{Jdt 4:5-8}
\switchcolumn*

\LA
\noindent Sacérdos étiam Elíachim scripsit ad univérsos qui erant contra Esdrelon, quæ est contra fáciem campi magni juxta Dóthain, et univérsos per quos viæ tránsitus esse póterat, ut obtinérent ascénsus móntium, per quos via esse póterat ad Jerúsalem, et illic custodírent ubi angústum iter esse póterat inter montes. Et fecérunt fílii Israël secúndum quod constitúerat eis sacérdos Dómini Elíachim. Et clamávit omnis pópulus ad Dóminum instántia magna, et humiliavérunt ánimas suas in jejúniis et oratiónibus, ipsi et mulíeres eórum.\par
\switchcolumn
\EN
\noindent And Eliachim the priest wrote to all that were over against Esdrelon, which faceth the great plain near Dothain, and to all by whom there might be a passage of way, that they should take possession of the ascents of the mountains, by which there might be any way to Jerusalem, and should keep watch where the way was narrow between the mountains. And the children of Israel did as the priest of the Lord Eliachim had appointed them, And all the people cried to the Lord with great earnestness, and they humbled their souls in fastings, and prayers, both they and their wives. And the priests put on haircloths, and they caused the little children to lie prostrate before the temple of the Lord, and the altar of the Lord they covered with haircloth.\par
\switchcolumn*

\LA
\begin{Resp}\Rsign Tribulatiónes civitátum audívimus, quas passæ sunt, et defécimus: timor et hebetúdo mentis cécidit super nos et super líberos nostros: ipsi montes nolunt recípere fugam nostram: \Star{} Dómine, miserére.\end{Resp}
\begin{Resp}\Vsign Peccávimus cum pátribus nostris, injúste égimus, iniquitátem fécimus.\end{Resp}
\begin{Resp}\Rsign Dómine, miserére.\end{Resp}
\switchcolumn
\EN
\begin{Resp}\Rsign We have heard of the tribulation of those cities, which they have suffered, and we have fainted. Fear and confusion of mind are fallen upon us. Even the mountains will not give us a refuge. \Star{} Lord, have mercy.\end{Resp}
\begin{Resp}\Vsign We have sinned like our forefathers, we have done unjustly, and wrought iniquity.\end{Resp}
\begin{Resp}\Rsign Lord, have mercy.\end{Resp}
\switchcolumn*

\LA
\LessonHead{Lectio III}
\switchcolumn
\EN
\LessonHead{Lesson III}
\switchcolumn*

\LA
\Cite{Jdt 4:9-12}
\switchcolumn
\EN
\Cite{Jdt 4:9-12}
\switchcolumn*

\LA
\noindent Et induérunt se sacerdótes cilíciis, et infántes prostravérunt contra fáciem templi Dómini, et altáre Dómini operuérunt cilício; et clamavérunt ad Dóminum Deum Israël unanímiter, ne daréntur in prædam infántes eórum, et uxóres eórum in divisiónem, et civitátes eórum in extermínium, et sancta eórum in pollutiónem, et fíerent oppróbrium géntibus. Tunc Elíachim sacérdos Dómini magnus circuívit omnem Israël allocutúsque est eos, dicens: Scitóte quóniam exáudiet Dóminus preces vestras, si manéntes permanséritis in jejúniis et oratiónibus in conspéctu Dómini.\par
\switchcolumn
\EN
\noindent And they cried to the Lord the God of Israel with one accord, that their children might not be made a prey, and their wives carried off, and their cities destroyed, and their holy things profaned, and that they might not be made a reproach to the Gentiles. Then Eliachim the high priest of the Lord went about all Israel and spoke to them, Saying: Know ye that the Lord will hear your prayers, if you continue with perseverance in fastings and prayers in the sight of the Lord. Remember Moses the servant of the Lord, who overcame Amalec that trusted in his own strength, and in his power, and in his army, and in his shields, and in his chariots, and in his horsemen, not by fighting with the sword, but by holy prayers:\par
\switchcolumn*

\LA
\begin{Resp}\Rsign Benedícat te Dóminus in virtúte sua, qui per te ad níhilum redégit inimícos nostros: \Star{} Ut non defíciat laus tua de ore hóminum.\end{Resp}
\begin{Resp}\Vsign Benedíctus Dóminus qui creávit cælum et terram; quia hódie nomen tuum ita magnificávit.\end{Resp}
\begin{Resp}\Rsign Ut non defíciat laus tua de ore hóminum.\end{Resp}
\begin{Resp}\Vsign Glória Patri, et Fílio, \Star{} et Spirítui Sancto.\end{Resp}
\begin{Resp}\Rsign Ut non defíciat laus tua de ore hóminum.\end{Resp}
\switchcolumn
\EN
\begin{Resp}\Rsign The Lord bless thee by His power, Who hath brought our enemies to nought through thee. \Star{} And may the praise of thee never fail from the mouth of men.\end{Resp}
\begin{Resp}\Vsign Blessed be the Lord, Who hath created the heaven and the earth, because that He hath so glorified thy name this day.\end{Resp}
\begin{Resp}\Rsign And may the praise of thee never fail from the mouth of men.\end{Resp}
\begin{Resp}\Vsign Glory be to the Father, and to the Son, \Star{} and to the Holy Ghost.\end{Resp}
\begin{Resp}\Rsign And may the praise of thee never fail from the mouth of men.\end{Resp}
\switchcolumn*

\end{paracol}

\Office{Feria III infra Hebdomadam IV Septembris}{Tuesday of the fourth week of September}{}
\addcontentsline{toc}{subsection}{Feria III infra Hebdomadam IV Septembris \textit{— Tuesday of the fourth week of September}}
\begin{paracol}{2}
\LA
\LessonHead{Lectio I}
\switchcolumn
\EN
\LessonHead{Lesson I}
\switchcolumn*

\LA
\LessonTitle{De libro Judith}
\switchcolumn
\EN
\LessonTitle{Lesson from the book of Judith}
\switchcolumn*

\LA
\Cite{Jdt 8:1-4}
\switchcolumn
\EN
\Cite{Jdt 8:1-4}
\switchcolumn*

\LA
\noindent Et factum est, cum audísset hæc verba Judith vidua, quæ erat fília Merári, fílii Idox, fílii Joseph, fílii Ozíæ, fílii Elai, fílii Jammor, fílii Gédeon, fílii Ráphaim, fílii Achitob, fílii Melchíæ, fílii Enan, fílii Nathaníæ, fílii Saláthiel, fílii Símeon, fílii Ruben, et vir ejus fuit Manásses, qui mórtuus est in diébus messis hordeáceæ; instábat enim super alligántes manípulos in campo, et venit æstus super caput ejus, et mórtuus est in Bethúlia civitáte sua et sepúltus est illic cum pátribus suis. Erat autem Judith relícta ejus vídua jam annis tribus et ménsibus sex.\par
\switchcolumn
\EN
\noindent Now it came to pass, when Judith a widow had heard these words, who was the daughter of Merari, the son of Idox, the son of Joseph, the son of Ozias, the son of Elai, the son of Jamnor, the son of Gedeon, the son of Raphaim, the son of Achitob, the son of Melehias, the son of Enan, the son of Nathanias, the son of Salathiel, the son of Simeon, the son of Ruben: And her husband was Manasses, who died in the time of the barley harvest: For he was standing over them that bound sheaves in the field and the heat came upon his head, and he died in Bethulia his own city, and was buried there with his fathers. And Judith his relict was a widow now three years and six months.\par
\switchcolumn*

\LA
\begin{Resp}\Rsign Nos álium Deum nescímus præter Dóminum, in quo sperámus: \Star{} Qui non déspicit nos, nec ámovet salútem suam a génere nostro.\end{Resp}
\begin{Resp}\Vsign Indulgéntiam ipsíus fusis lácrimis postulémus, et humiliémus illi ánimas nostras.\end{Resp}
\begin{Resp}\Rsign Qui non déspicit nos, nec ámovet salútem suam a génere nostro.\end{Resp}
\switchcolumn
\EN
\begin{Resp}\Rsign We know no strange God before the Lord. In Him we trust. \Star{} He despiseth us not, neither putteth He away His salvation from our nation.\end{Resp}
\begin{Resp}\Vsign His mercy let us seek with tears, and humble our souls before Him.\end{Resp}
\begin{Resp}\Rsign He despiseth us not, neither putteth He away His salvation from our nation.\end{Resp}
\switchcolumn*

\LA
\LessonHead{Lectio II}
\switchcolumn
\EN
\LessonHead{Lesson II}
\switchcolumn*

\LA
\Cite{Jdt 8:5-8}
\switchcolumn
\EN
\Cite{Jdt 8:5-8}
\switchcolumn*

\LA
\noindent Et in superióribus domus suæ fecit sibi secrétum cubículum, in quo cum puéllis suis clausa morabátur. Et habens super lumbos suos cilícium, jejunábat ómnibus diébus vitæ suæ præter sábbata et neoménias et festa domus Israël. Erat autem elegánti aspéctu nimis; cui vir suus relíquerat divítias multas et famíliam copiósam ac possessiónes arméntis boum et grégibus óvium plenas. Et erat hæc in ómnibus famosíssima, quóniam timébat Dóminum valde, nec erat qui loquerétur de illa verbum malum.\par
\switchcolumn
\EN
\noindent And she made herself a private chamber in the upper part of her house, in which she abode shut up with her maids. And she wore haircloth upon her loins, and fasted all the days of her life, except the sabbaths, and new moons, and the feasts of the house of Israel. And she was exceedingly beautiful, and her husband left her great riches, and very many servants, and large possessions of herds of oxen, and flocks of sheep. And she was greatly renowned among all, because she feared the Lord very much, neither was there any one that spoke an ill word of her.\par
\switchcolumn*

\LA
\begin{Resp}\Rsign Dominátor, Dómine, cælórum et terræ, Creátor aquárum, Rex univérsæ creatúræ: \Star{} Exáudi oratiónem servórum tuórum.\end{Resp}
\begin{Resp}\Vsign Tu, Dómine, cui humílium semper et mansuetórum plácuit deprecátio.\end{Resp}
\begin{Resp}\Rsign Exáudi oratiónem servórum tuórum.\end{Resp}
\switchcolumn
\EN
\begin{Resp}\Rsign O Lord, Ruler of the heavens and of the earth, Maker of the waters, King of every creature, \Star{} Graciously hear the prayer of thy servants.\end{Resp}
\begin{Resp}\Vsign Thou, O Lord, unto Whom the supplications of the humble and meek are alway well-pleasing\end{Resp}
\begin{Resp}\Rsign Graciously hear the prayer of thy servants.\end{Resp}
\switchcolumn*

\LA
\LessonHead{Lectio III}
\switchcolumn
\EN
\LessonHead{Lesson III}
\switchcolumn*

\LA
\Cite{Jdt 8:9-11}
\switchcolumn
\EN
\Cite{Jdt 8:9-11}
\switchcolumn*

\LA
\noindent Hæc ítaque, cum audísset quóniam Ozías promisísset quod, transácto quinto die, tráderet civitátem, misit ad presbýteros Chabri et Charmi. Et venérunt ad illam, et dixit illis: Quod est hoc verbum, in quo consénsit Ozías, ut tradat civitátem Assýriis, si intra quinque dies non vénerit vobis adjutórium? Et qui estis vos, qui tentátis Dóminum?\par
\switchcolumn
\EN
\noindent When therefore she had heard that Ozias had promised that he would deliver up the city after the fifth day, she sent to the ancients Chabri and Charmi. And they came to her, and she said to them: What is this word, by which Ozias hath consented to give up the city to the Assyrians, if within five days there come no aid to us? And who are you that tempt the Lord?\par
\switchcolumn*

\LA
\begin{Resp}\Rsign Dómine Deus, qui cónteris bella ab inítio, éleva brácchium tuum super gentes, quæ cógitant servis tuis mala: \Star{} Et déxtera tua glorificétur in nobis.\end{Resp}
\begin{Resp}\Vsign Allíde virtútem eórum in virtúte tua; cadat robur eórum in iracúndia tua.\end{Resp}
\begin{Resp}\Rsign Et déxtera tua glorificétur in nobis.\end{Resp}
\begin{Resp}\Vsign Glória Patri, et Fílio, \Star{} et Spirítui Sancto.\end{Resp}
\begin{Resp}\Rsign Et déxtera tua glorificétur in nobis.\end{Resp}
\switchcolumn
\EN
\begin{Resp}\Rsign O Lord God, That breakest the battles from of old, lift up thine arm against the Gentiles, that devise evil against thy servants. \Star{} And let thy right hand be glorified in us.\end{Resp}
\begin{Resp}\Vsign Throw down their strength in thy power, and bring down their force in thy wrath.\end{Resp}
\begin{Resp}\Rsign And let thy right hand be glorified in us.\end{Resp}
\begin{Resp}\Vsign Glory be to the Father, and to the Son, \Star{} and to the Holy Ghost.\end{Resp}
\begin{Resp}\Rsign And let thy right hand be glorified in us.\end{Resp}
\switchcolumn*

\end{paracol}

\Office{Feria IV infra Hebdomadam IV Septembris}{Wednesday of the fourth week of September}{}
\addcontentsline{toc}{subsection}{Feria IV infra Hebdomadam IV Septembris \textit{— Wednesday of the fourth week of September}}
\begin{paracol}{2}
\LA
\LessonHead{Lectio I}
\switchcolumn
\EN
\LessonHead{Lesson I}
\switchcolumn*

\LA
\LessonTitle{De libro Judith}
\switchcolumn
\EN
\LessonTitle{Lesson from the book of Judith}
\switchcolumn*

\LA
\Cite{Jdt 10:1-4}
\switchcolumn
\EN
\Cite{Jdt 10:1-4}
\switchcolumn*

\LA
\noindent Factum est autem, cum cessásset clamáre ad Dóminum, surréxit de loco in quo jacúerat prostráta ad Dóminum; vocavítque abram suam, et descéndens in domum suam ábstulit a se cilícium et éxuit se vestiméntis viduitátis suæ, et lavit corpus suum et unxit se myro óptimo et discriminávit crinem cápitis sui et impósuit mitram super caput suum et índuit se vestiméntis jucunditátis suæ induítque sandália pédibus suis assumpsítque dextralíola et lília et ináures et ánulos et ómnibus ornaméntis suis ornávit se. Cui étiam Dóminus cóntulit splendórem.\par
\switchcolumn
\EN
\noindent And it came to pass, when she had ceased to cry to the Lord, that she rose from the place wherein she lay prostrate before the Lord. And she called her maid, and going down into her house she took off her haircloth, and put away the garments of her widowhood, And she washed her body, and anointed herself with the best ointment, and plaited the hair of her head, and put a bonnet upon her head, and clothed herself with the garments of her gladness, and put sandals on her feet, and took her bracelets, and lilies, and earlets, and rings, and adorned herself with all her ornaments. And the Lord also gave her more beauty.\par
\switchcolumn*

\LA
\begin{Resp}\Rsign Confórta me, Rex, Sanctórum principátum tenens: \Star{} Et da sermónem rectum et bene sonántem in os meum.\end{Resp}
\begin{Resp}\Vsign Dómine, Rex univérsæ potestátis, convérte consílium eórum super eos.\end{Resp}
\begin{Resp}\Rsign Et da sermónem rectum et bene sonántem in os meum.\end{Resp}
\switchcolumn
\EN
\begin{Resp}\Rsign Strengthen me, O King, Who reignest over the holy ones. \Star{} Put Thou in my mouth clear and well-sounding words.\end{Resp}
\begin{Resp}\Vsign O Lord, King of all forces, turn back their device upon themselves.\end{Resp}
\begin{Resp}\Rsign Put Thou in my mouth clear and well-sounding words.\end{Resp}
\switchcolumn*

\LA
\LessonHead{Lectio II}
\switchcolumn
\EN
\LessonHead{Lesson II}
\switchcolumn*

\LA
\Cite{Jdt 10:11-12}
\switchcolumn
\EN
\Cite{Jdt 10:11-12}
\switchcolumn*

\LA
\noindent Factum est autem, cum descénderet montem circa ortum diéi, occurrérunt ei exploratóres Assyriórum et tenuérunt eam dicéntes: Unde venis aut quo vadis? Quæ respóndit: Fília sum Hebræórum; ídeo ego fugi a fácie eórum, quóniam futúrum agnóvi quod dentur vobis in deprædatiónem, pro eo quod contemnéntes vos noluérunt ultro trádere seípsos, ut invenírent misericórdiam in conspéctu vestro.\par
\switchcolumn
\EN
\noindent And it came to pass, when she went down the hill, about break of day, that the watchmen of the Assyrians met her and stopped her, saying: Whence comest thou? or whither goest thou? And she answered: I am a daughter of the Hebrews, and I am fled from them, because I knew they would be made a prey to you, because they despised you, and would not of their own accord yield themselves, that they might find mercy in your sight.\par
\switchcolumn*

\LA
\begin{Resp}\Rsign Tribulatiónes civitátum audívimus, quas passæ sunt, et defécimus: timor et hebetúdo mentis cécidit super nos et super líberos nostros: ipsi montes nolunt recípere fugam nostram: \Star{} Dómine, miserére.\end{Resp}
\begin{Resp}\Vsign Peccávimus cum pátribus nostris, injúste égimus, iniquitátem fécimus.\end{Resp}
\begin{Resp}\Rsign Dómine, miserére.\end{Resp}
\switchcolumn
\EN
\begin{Resp}\Rsign We have heard of the tribulation of those cities, which they have suffered, and we have fainted. Fear and confusion of mind are fallen upon us. Even the mountains will not give us a refuge. \Star{} Lord, have mercy.\end{Resp}
\begin{Resp}\Vsign We have sinned like our forefathers, we have done unjustly, and wrought iniquity.\end{Resp}
\begin{Resp}\Rsign Lord, have mercy.\end{Resp}
\switchcolumn*

\LA
\LessonHead{Lectio III}
\switchcolumn
\EN
\LessonHead{Lesson III}
\switchcolumn*

\LA
\Cite{Jdt 10:16-20}
\switchcolumn
\EN
\Cite{Jdt 10:16-20}
\switchcolumn*

\LA
\noindent Duxerúntque illam ad tabernáculum Holoférnis, annuntiántes eam. Cumque intrásset ante fáciem ejus, statim captus est in suis óculis Holoférnes. Dixerúntque ad eum satéllites ejus: Quis contémnat pópulum Hebræórum, qui tam decóras mulíeres habent, ut non pro his mérito pugnáre contra eos debeámus? Videns ítaque Judith Holoférnem sedéntem in conopéo, quod erat ex púrpura et auro et smarágdo et lapídibus pretiósis intéxtum, et, cum in fáciem ejus intendísset, adorávit eum prostérnens se super terram.\par
\switchcolumn
\EN
\noindent And they brought her to the tent of Holofernes, telling him of her. And when she was come into his presence, forthwith Holofernes was caught by his eyes. And his officers said to him: Who can despise the people of the Hebrews who have such beautiful women, that we should not think it worth our while for their sakes to fight against them? And Judith seeing Holofernes sitting under a canopy, which was woven of purple and gold, with emeralds and precious stones: After she had looked on his face bowed down to him, prostrating herself to the ground.\par
\switchcolumn*

\LA
\begin{Resp}\Rsign Benedícat te Dóminus in virtúte sua, qui per te ad níhilum redégit inimícos nostros: \Star{} Ut non defíciat laus tua de ore hóminum.\end{Resp}
\begin{Resp}\Vsign Benedíctus Dóminus qui creávit cælum et terram; quia hódie nomen tuum ita magnificávit.\end{Resp}
\begin{Resp}\Rsign Ut non defíciat laus tua de ore hóminum.\end{Resp}
\begin{Resp}\Vsign Glória Patri, et Fílio, \Star{} et Spirítui Sancto.\end{Resp}
\begin{Resp}\Rsign Ut non defíciat laus tua de ore hóminum.\end{Resp}
\switchcolumn
\EN
\begin{Resp}\Rsign The Lord bless thee by His power, Who hath brought our enemies to nought through thee. \Star{} And may the praise of thee never fail from the mouth of men.\end{Resp}
\begin{Resp}\Vsign Blessed be the Lord, Who hath created the heaven and the earth, because that He hath so glorified thy name this day.\end{Resp}
\begin{Resp}\Rsign And may the praise of thee never fail from the mouth of men.\end{Resp}
\begin{Resp}\Vsign Glory be to the Father, and to the Son, \Star{} and to the Holy Ghost.\end{Resp}
\begin{Resp}\Rsign And may the praise of thee never fail from the mouth of men.\end{Resp}
\switchcolumn*

\end{paracol}

\Office{Feria V infra Hebdomadam IV Septembris}{Thursday of the fourth week of September}{}
\addcontentsline{toc}{subsection}{Feria V infra Hebdomadam IV Septembris \textit{— Thursday of the fourth week of September}}
\begin{paracol}{2}
\LA
\LessonHead{Lectio I}
\switchcolumn
\EN
\LessonHead{Lesson I}
\switchcolumn*

\LA
\LessonTitle{De libro Judith}
\switchcolumn
\EN
\LessonTitle{Lesson from the book of Judith}
\switchcolumn*

\LA
\Cite{Jdt 12:10-13}
\switchcolumn
\EN
\Cite{Jdt 12:10-13}
\switchcolumn*

\LA
\noindent Et factum est, in quarto die Holoférnes fecit cenam servis suis et dixit ad Vágao eunúchum suum: Vade, et suáde Hebrǽam illam ut sponte conséntiat habitáre mecum. Fœdum est enim apud Assýrios, si fémina irrídeat virum agéndo ut immúnis ab eo tránseat. Tunc introívit Vágao ad Judith et dixit: Non vereátur bona puélla introíre ad dóminum meum, ut honorificétur ante fáciem ejus, ut mandúcet cum eo et bibat vinum in jucunditáte. Cui Judith respóndit: Quæ ego sum, ut contradícam dómino meo?\par
\switchcolumn
\EN
\noindent And it came to pass on the fourth day, that Holofernes made a supper for his servants, and said to Vagao his eunuch: go, and persuade that Hebrew woman, to consent of her own accord to dwell with me. For it is looked upon as shameful among the Assyrians, if a woman mock a man, by doing so as to pass free from him. Then Vagao went in to Judith, and said: Let not my good maid be afraid to go in to my lord, that she may be honoured before his face, that she may eat with him and drink wine and be merry. And Judith answered him: Who am I, that I should gainsay my lord?\par
\switchcolumn*

\LA
\begin{Resp}\Rsign Adonái, Dómine, Deus magne et mirábilis, qui dedísti salútem in manu féminæ, \Star{} Exáudi preces servórum tuórum.\end{Resp}
\begin{Resp}\Vsign Benedíctus es, Dómine, qui non derelínquis præsuméntes de te, et de sua virtúte gloriántes humílias.\end{Resp}
\begin{Resp}\Rsign Exáudi preces servórum tuórum.\end{Resp}
\switchcolumn
\EN
\begin{Resp}\Rsign O Adonai, O Lord God, Thou art great and glorious. Who hast given salvation into the hand of a woman: \Star{} Graciously hear the prayers of thy servants.\end{Resp}
\begin{Resp}\Vsign Blessed art Thou, O Lord, Who failest none that put their trust in thee, and humblest such as boast themselves in their own strength.\end{Resp}
\begin{Resp}\Rsign Graciously hear the prayers of thy servants.\end{Resp}
\switchcolumn*

\LA
\LessonHead{Lectio II}
\switchcolumn
\EN
\LessonHead{Lesson II}
\switchcolumn*

\LA
\Cite{Jdt 13:1-7}
\switchcolumn
\EN
\Cite{Jdt 13:1-7}
\switchcolumn*

\LA
\noindent Ut autem sero factum est, festinavérunt servi illíus ad hospítia sua; et conclúsit Vágao óstia cubículi et ábiit. Erant autem omnes fatigáti a vino; erátque Judith sola in cubículo. Porro Holoférnes jacébat in lecto nímia ebrietáte sopítus. Dixítque Judith puéllæ suæ ut staret foris ante cubículum, et observáret. Stetítque Judith ante lectum orans cum lácrimis et labiórum motu in siléntio, dicens: Confírma me, Dómine Deus Israël, et réspice in hac hora ad ópera mánuum meárum, ut, sicut promisísti, Jerúsalem civitátem tuam érigas, et hoc quod credens per te posse fíeri cogitávi, perfíciam.\par
\switchcolumn
\EN
\noindent And when it was grown late, his servants made haste to their lodgings, and Vagao shut the chamber doors, and went his way. And they were all overcharged with wine. And Judith was alone in the chamber. But Holofernes lay on his bed, fast asleep, being exceedingly drunk. And Judith spoke to her maid to stand without before the chamber, and to watch: And Judith stood before the bed praying with tears, and the motion of her lips in silence, Saying: Strengthen me, O Lord God of Israel, and in this hour look on the works of my hands, that as thou hast promised, thou mayst raise up Jerusalem thy city: and that I may bring to pass that which I have purposed, having a belief that it might be done by thee.\par
\switchcolumn*

\LA
\begin{Resp}\Rsign Tribulatiónes civitátum audívimus, quas passæ sunt, et defécimus: timor et hebetúdo mentis cécidit super nos et super líberos nostros: ipsi montes nolunt recípere fugam nostram: \Star{} Dómine, miserére.\end{Resp}
\begin{Resp}\Vsign Peccávimus cum pátribus nostris, injúste égimus, iniquitátem fécimus.\end{Resp}
\begin{Resp}\Rsign Dómine, miserére.\end{Resp}
\switchcolumn
\EN
\begin{Resp}\Rsign We have heard of the tribulation of those cities, which they have suffered, and we have fainted. Fear and confusion of mind are fallen upon us. Even the mountains will not give us a refuge. \Star{} Lord, have mercy.\end{Resp}
\begin{Resp}\Vsign We have sinned like our forefathers, we have done unjustly, and wrought iniquity.\end{Resp}
\begin{Resp}\Rsign Lord, have mercy.\end{Resp}
\switchcolumn*

\LA
\LessonHead{Lectio III}
\switchcolumn
\EN
\LessonHead{Lesson III}
\switchcolumn*

\LA
\Cite{Jdt 13:8-11}
\switchcolumn
\EN
\Cite{Jdt 13:8-11}
\switchcolumn*

\LA
\noindent Et, cum hæc dixísset, accéssit ad colúmnam, quæ erat ad caput léctuli ejus, et pugiónem ejus, qui in ea ligátus pendébat, exsólvit. Cumque evaginásset illum, apprehéndit comam cápitis ejus et ait: Confírma me, Dómine Deus, in hac hora. Et percússit bis in cervícem ejus et abscídit caput ejus et ábstulit conopéum ejus a colúmnis et evólvit corpus ejus truncum. Et post pusíllum exívit, et trádidit caput Holoférnis ancíllæ suæ et jussit ut mítteret illud in peram suam.\par
\switchcolumn
\EN
\noindent And when she had said this, she went to the pillar that was at his bed's head, and loosed his sword that hung tied upon it. And when she had drawn it out, she took him by the hair of his head, and said: Strengthen me, O Lord God, at this hour. And she struck twice upon his neck, and out off his head, and took off his canopy from the pillars, and rolled away his headless body. And after a while she went out, and delivered the head of Holofernes to her maid, and bade her put it into her wallet.\par
\switchcolumn*

\LA
\begin{Resp}\Rsign Benedícat te Dóminus in virtúte sua, qui per te ad níhilum redégit inimícos nostros: \Star{} Ut non defíciat laus tua de ore hóminum.\end{Resp}
\begin{Resp}\Vsign Benedíctus Dóminus qui creávit cælum et terram; quia hódie nomen tuum ita magnificávit.\end{Resp}
\begin{Resp}\Rsign Ut non defíciat laus tua de ore hóminum.\end{Resp}
\begin{Resp}\Vsign Glória Patri, et Fílio, \Star{} et Spirítui Sancto.\end{Resp}
\begin{Resp}\Rsign Ut non defíciat laus tua de ore hóminum.\end{Resp}
\switchcolumn
\EN
\begin{Resp}\Rsign The Lord bless thee by His power, Who hath brought our enemies to nought through thee. \Star{} And may the praise of thee never fail from the mouth of men.\end{Resp}
\begin{Resp}\Vsign Blessed be the Lord, Who hath created the heaven and the earth, because that He hath so glorified thy name this day.\end{Resp}
\begin{Resp}\Rsign And may the praise of thee never fail from the mouth of men.\end{Resp}
\begin{Resp}\Vsign Glory be to the Father, and to the Son, \Star{} and to the Holy Ghost.\end{Resp}
\begin{Resp}\Rsign And may the praise of thee never fail from the mouth of men.\end{Resp}
\switchcolumn*

\end{paracol}

\Office{Feria VI infra Hebdomadam IV Septembris}{Friday of the fourth week of September}{}
\addcontentsline{toc}{subsection}{Feria VI infra Hebdomadam IV Septembris \textit{— Friday of the fourth week of September}}
\begin{paracol}{2}
\LA
\LessonHead{Lectio I}
\switchcolumn
\EN
\LessonHead{Lesson I}
\switchcolumn*

\LA
\LessonTitle{De libro Judith}
\switchcolumn
\EN
\LessonTitle{Lesson from the book of Judith}
\switchcolumn*

\LA
\Cite{Jdt 15:1-3}
\switchcolumn
\EN
\Cite{Jdt 15:1-3}
\switchcolumn*

\LA
\noindent Cumque omnis exércitus decollátum Holoférnem audísset, fugit mens et consílium ab eis, et solo tremóre et metu agitáti fugæ præsídium sumunt, ita ut nullus loquerétur cum próximo suo, sed inclináto cápite, relíctis ómnibus, evádere festinábant Hebrǽos, quos armátos super se veníre audiébant, fugiéntes per vias campórum et sémitas cóllium. Vidéntes ítaque fílii Israël fugiéntes, secúti sunt illos, descenderúntque clangéntes tubis et ululántes post ipsos.\par
\switchcolumn
\EN
\noindent And when all the army heard that Holofernes was beheaded, courage and counsel fled from them, and being seized with trembling and fear they thought only to save themselves by flight: So that no one spoke to his neighbor, but hanging down the head, leaving all things behind, they made haste to escape from the Hebrews, who, as they heard, were coming armed upon them, and fled by the ways of the fields, and the paths of the hills. So the children of Israel seeing them fleeing, followed after them. And they went down sounding with trumpets and shouting after them.\par
\switchcolumn*

\LA
\begin{Resp}\Rsign Nos álium Deum nescímus præter Dóminum, in quo sperámus: \Star{} Qui non déspicit nos, nec ámovet salútem suam a génere nostro.\end{Resp}
\begin{Resp}\Vsign Indulgéntiam ipsíus fusis lácrimis postulémus, et humiliémus illi ánimas nostras.\end{Resp}
\begin{Resp}\Rsign Qui non déspicit nos, nec ámovet salútem suam a génere nostro.\end{Resp}
\switchcolumn
\EN
\begin{Resp}\Rsign We know no strange God before the Lord. In Him we trust. \Star{} He despiseth us not, neither putteth He away His salvation from our nation.\end{Resp}
\begin{Resp}\Vsign His mercy let us seek with tears, and humble our souls before Him.\end{Resp}
\begin{Resp}\Rsign He despiseth us not, neither putteth He away His salvation from our nation.\end{Resp}
\switchcolumn*

\LA
\LessonHead{Lectio II}
\switchcolumn
\EN
\LessonHead{Lesson II}
\switchcolumn*

\LA
\Cite{Jdt 15:5-7}
\switchcolumn
\EN
\Cite{Jdt 15:5-7}
\switchcolumn*

\LA
\noindent Misit ítaque Ozías núntios per omnes civitátes et regiónes Israël. Omnis ítaque regio omnísque urbs eléctam juventútem armátam misit post eos; et persecúti sunt eos in ore gládii, quoúsque pervenírent ad extremitátem fínium suórum. Réliqui autem qui erant in Bethúlia, ingréssi sunt castra Assyriórum, et prædam, quam fugiéntes Assýrii relíquerant, abstulérunt, et onustáti sunt valde.\par
\switchcolumn
\EN
\noindent And Ozias sent messengers through all the cities and countries of Israel. And every country, and every city, sent their chosen young men armed after them, and they pursued them with the edge of the sword until they came to the extremities of their confines. And the rest that were in Bethulia went into the camp of the Assyrians, and took away the spoils, which the Assyrians in their flight had left behind them, and they were laden exceedingly.\par
\switchcolumn*

\LA
\begin{Resp}\Rsign Dominátor, Dómine, cælórum et terræ, Creátor aquárum, Rex univérsæ creatúræ: \Star{} Exáudi oratiónem servórum tuórum.\end{Resp}
\begin{Resp}\Vsign Tu, Dómine, cui humílium semper et mansuetórum plácuit deprecátio.\end{Resp}
\begin{Resp}\Rsign Exáudi oratiónem servórum tuórum.\end{Resp}
\switchcolumn
\EN
\begin{Resp}\Rsign O Lord, Ruler of the heavens and of the earth, Maker of the waters, King of every creature, \Star{} Graciously hear the prayer of thy servants.\end{Resp}
\begin{Resp}\Vsign Thou, O Lord, unto Whom the supplications of the humble and meek are alway well-pleasing\end{Resp}
\begin{Resp}\Rsign Graciously hear the prayer of thy servants.\end{Resp}
\switchcolumn*

\LA
\LessonHead{Lectio III}
\switchcolumn
\EN
\LessonHead{Lesson III}
\switchcolumn*

\LA
\Cite{Jdt 15:9-12}
\switchcolumn
\EN
\Cite{Jdt 15:9-12}
\switchcolumn*

\LA
\noindent Jóachim autem summus póntifex de Jerúsalem venit in Bethúliam cum univérsis presbýteris suis, ut vidéret Judith. Quæ cum exísset ad illum, benedixérunt eam omnes una voce, dicéntes: Tu glória Jerúsalem, tu lætítia Israël, tu honorificéntia pópuli nostri. Quia fecísti viríliter, et confortátum est cor tuum, eo quod castitátem amáveris et, post virum tuum, álterum nescíeris, ídeo et manus Dómini confortávit te, et ídeo eris benedícta in ætérnum. Et dixit omnis pópulus: Fiat, fiat.\par
\switchcolumn
\EN
\noindent And Joachim the high priest came from Jerusalem to Bethulia with all his ancients to see Judith. And when she was come out to him, they all blessed her with one voice, saying: Thou art the glory of Jerusalem, thou art the joy of Israel, thou art the honour of our people: For thou hast done manfully, and thy heart has been strengthened, because thou hast loved chastity, and after thy husband hast not known any other: therefore also the hand of the Lord hath strengthened thee, and therefore thou shalt be blessed for ever. And all the people said: So be it, so be it.\par
\switchcolumn*

\LA
\begin{Resp}\Rsign Dómine Deus, qui cónteris bella ab inítio, éleva brácchium tuum super gentes, quæ cógitant servis tuis mala: \Star{} Et déxtera tua glorificétur in nobis.\end{Resp}
\begin{Resp}\Vsign Allíde virtútem eórum in virtúte tua; cadat robur eórum in iracúndia tua.\end{Resp}
\begin{Resp}\Rsign Et déxtera tua glorificétur in nobis.\end{Resp}
\begin{Resp}\Vsign Glória Patri, et Fílio, \Star{} et Spirítui Sancto.\end{Resp}
\begin{Resp}\Rsign Et déxtera tua glorificétur in nobis.\end{Resp}
\switchcolumn
\EN
\begin{Resp}\Rsign O Lord God, That breakest the battles from of old, lift up thine arm against the Gentiles, that devise evil against thy servants. \Star{} And let thy right hand be glorified in us.\end{Resp}
\begin{Resp}\Vsign Throw down their strength in thy power, and bring down their force in thy wrath.\end{Resp}
\begin{Resp}\Rsign And let thy right hand be glorified in us.\end{Resp}
\begin{Resp}\Vsign Glory be to the Father, and to the Son, \Star{} and to the Holy Ghost.\end{Resp}
\begin{Resp}\Rsign And let thy right hand be glorified in us.\end{Resp}
\switchcolumn*

\end{paracol}

\Office{Sabbato infra Hebdomadam IV Septembris}{Saturday of the fourth week of September}{}
\addcontentsline{toc}{subsection}{Sabbato infra Hebdomadam IV Septembris \textit{— Saturday of the fourth week of September}}
\begin{paracol}{2}
\LA
\LessonHead{Lectio I}
\switchcolumn
\EN
\LessonHead{Lesson I}
\switchcolumn*

\LA
\LessonTitle{De libro Judith}
\switchcolumn
\EN
\LessonTitle{Lesson from the book of Judith}
\switchcolumn*

\LA
\Cite{Jdt 16:22-23}
\switchcolumn
\EN
\Cite{Jdt 16:22-23}
\switchcolumn*

\LA
\noindent Et factum est post hæc, omnis pópulus post victóriam venit in Jerúsalem adoráre Dóminum; et mox ut purificáti sunt, obtulérunt omnes holocáusta, et vota, et repromissiónes suas. Porro Judith univérsa vasa béllica Holoférnis, quæ dedit illi pópulus, et conopéum, quod ipsa sustúlerat de cubíli ipsíus, óbtulit in anáthema obliviónis.\par
\switchcolumn
\EN
\noindent And it came to pass after these things, that all the people, after the victory, came to Jerusalem to adore the Lord: and as soon as they were purified, they all offered holocausts, and vows, and their promises. And Judith offered for an anathema of oblivion all the arms of Holofernes, which the people gave her, and the canopy that she had taken away out of his chamber.\par
\switchcolumn*

\LA
\begin{Resp}\Rsign Confórta me, Rex, Sanctórum principátum tenens: \Star{} Et da sermónem rectum et bene sonántem in os meum.\end{Resp}
\begin{Resp}\Vsign Dómine, Rex univérsæ potestátis, convérte consílium eórum super eos.\end{Resp}
\begin{Resp}\Rsign Et da sermónem rectum et bene sonántem in os meum.\end{Resp}
\switchcolumn
\EN
\begin{Resp}\Rsign Strengthen me, O King, Who reignest over the holy ones. \Star{} Put Thou in my mouth clear and well-sounding words.\end{Resp}
\begin{Resp}\Vsign O Lord, King of all forces, turn back their device upon themselves.\end{Resp}
\begin{Resp}\Rsign Put Thou in my mouth clear and well-sounding words.\end{Resp}
\switchcolumn*

\LA
\LessonHead{Lectio II}
\switchcolumn
\EN
\LessonHead{Lesson II}
\switchcolumn*

\LA
\Cite{Jdt 16:24-27}
\switchcolumn
\EN
\Cite{Jdt 16:24-27}
\switchcolumn*

\LA
\noindent Erat autem pópulus jucúndus secúndum fáciem sanctórum; et per tres menses gáudium hujus victóriæ celebrátum est cum Judith. Post dies autem illos unusquísque rédiit in domum suam, et Judith magna facta est in Bethúlia et præclárior erat univérsæ terræ Israël. Erat étiam virtúti cástitas adjúncta, ita ut non cognósceret virum ómnibus diébus vitæ suæ, ex quo defúnctus est Manásses vir ejus. Erat autem, diébus festis, procédens cum magna glória.\par
\switchcolumn
\EN
\noindent And the people were joyful in the sight of the sanctuary, and for three months the joy of this victory was celebrated with Judith. And after those days every man returned to his house, and Judith was made great in Bethulia, and she was most renowned in all the land of Israel. And chastity was joined to her virtue, so that she knew no man all the days of her life, after the death of Manasses her husband. And on festival days she came forth with great glory.\par
\switchcolumn*

\LA
\begin{Resp}\Rsign Tribulatiónes civitátum audívimus, quas passæ sunt, et defécimus: timor et hebetúdo mentis cécidit super nos et super líberos nostros: ipsi montes nolunt recípere fugam nostram: \Star{} Dómine, miserére.\end{Resp}
\begin{Resp}\Vsign Peccávimus cum pátribus nostris, injúste égimus, iniquitátem fécimus.\end{Resp}
\begin{Resp}\Rsign Dómine, miserére.\end{Resp}
\begin{Resp}\Vsign Glória Patri, et Fílio, \Star{} et Spirítui Sancto.\end{Resp}
\begin{Resp}\Rsign Dómine, miserére.\end{Resp}
\switchcolumn
\EN
\begin{Resp}\Rsign We have heard of the tribulation of those cities, which they have suffered, and we have fainted. Fear and confusion of mind are fallen upon us. Even the mountains will not give us a refuge. \Star{} Lord, have mercy.\end{Resp}
\begin{Resp}\Vsign We have sinned like our forefathers, we have done unjustly, and wrought iniquity.\end{Resp}
\begin{Resp}\Rsign Lord, have mercy.\end{Resp}
\begin{Resp}\Vsign Glory be to the Father, and to the Son, \Star{} and to the Holy Ghost.\end{Resp}
\begin{Resp}\Rsign Lord, have mercy.\end{Resp}
\switchcolumn*

\end{paracol}

\Office{Dominica V Septembris}{Fifth Sunday of September}{}
\addcontentsline{toc}{subsection}{Dominica V Septembris \textit{— Fifth Sunday of September}}
\begin{paracol}{2}
\LA
\LessonHead{Lectio I}
\switchcolumn
\EN
\LessonHead{Lesson I}
\switchcolumn*

\LA
\LessonTitle{Incipit liber Esther}
\switchcolumn
\EN
\LessonTitle{Lesson from the book of Esther}
\switchcolumn*

\LA
\Cite{Esth 1:1-4}
\switchcolumn
\EN
\Cite{Esth 1:1-4}
\switchcolumn*

\LA
\noindent In diébus Assuéri, qui regnávit ab India usque Æthiópiam super centum vigínti septem províncias, quando sedit in sólio regni sui, Susan cívitas regni ejus exórdium fuit. Tértio ígitur anno impérii sui, fecit grande convívium cunctis princípibus et púeris suis fortíssimis Persárum et Medórum ínclitis et præféctis provinciárum coram se, ut osténderet divítias glóriæ regni sui, ac magnitúdinem atque jactántiam poténtiæ suæ, multo témpore, centum vidélicet et octogínta diébus.\par
\switchcolumn
\EN
\noindent In the days of Assuerus, who reigned from India to Ethiopia over a hundred and twenty-seven provinces: When he sat on the throne of his kingdom, the city Susan was the capital of his kingdom. Now in the third year of his reign he made a great feast for all the princes, and for his servants, for the most mighty of the Persians, and the nobles of the Medes, and the governors of the provinces in his sight, That he might shew the riches of the glory of his kingdom, and the greatness, and boasting of his power, for a long time, to wit, for a hundred and fourscore days.\par
\switchcolumn*

\LA
\begin{Resp}\Rsign Dómine, mi Rex omnípotens, in dicióne tua cuncta sunt pósita, et non est qui possit resístere voluntáti tuæ: \Star{} Líbera nos propter nomen tuum.\end{Resp}
\begin{Resp}\Vsign Exáudi oratiónem nostram, et convérte luctum nostrum in gáudium.\end{Resp}
\begin{Resp}\Rsign Líbera nos propter nomen tuum.\end{Resp}
\switchcolumn
\EN
\begin{Resp}\Rsign O my Lord, the King Almighty All things are in thy power, and there is no man that can gainsay thy will; \Star{} Deliver us for thy Name's sake.\end{Resp}
\begin{Resp}\Vsign Hear our prayer, and turn our sorrow into joy.\end{Resp}
\begin{Resp}\Rsign Deliver us for thy Name's sake.\end{Resp}
\switchcolumn*

\LA
\LessonHead{Lectio II}
\switchcolumn
\EN
\LessonHead{Lesson II}
\switchcolumn*

\LA
\Cite{Esth 1:5-6}
\switchcolumn
\EN
\Cite{Esth 1:5-6}
\switchcolumn*

\LA
\noindent Cumque impleréntur dies convívii, invitávit omnem pópulum, qui invéntus est in Susan, a máximo usque ad mínimum; et jussit septem diébus convívium præparári in vestíbulo horti et némoris, quod régio cultu et manu cónsitum erat. Et pendébant ex omni parte tentória aérii colóris et carbásini ac hyacínthini sustentáta fúnibus býssinis atque purpúreis, qui ebúrneis círculis insérti erant et colúmnis marmóreis fulciebántur. Léctuli quoque áurei et argéntei super paviméntum smarágdino et pário stratum lápide dispósiti erant, quod mira varietáte pictúra decorábat.\par
\switchcolumn
\EN
\noindent And when the days of the feast were expired, he invited all the people that were found in Susan, from the greatest to the least: and commanded a feast to be made seven days in the court of the garden, and of the wood, which was planted by the care and the hand of the king. And there were hung up on every side sky coloured, and green, and violet hangings, fastened with cords of silk, and of purple, which were put into rings of ivory, and were held up with marble pillars. The beds also were of gold and silver, placed in order upon a floor paved with porphyry and white marble: which was embellished with painting of wonderful variety.\par
\switchcolumn*

\LA
\begin{Resp}\Rsign Confórta me, Rex, Sanctórum principátum tenens: \Star{} Et da sermónem rectum et bene sonántem in os meum.\end{Resp}
\begin{Resp}\Vsign Dómine, Rex univérsæ potestátis, convérte consílium eórum super eos.\end{Resp}
\begin{Resp}\Rsign Et da sermónem rectum et bene sonántem in os meum.\end{Resp}
\switchcolumn
\EN
\begin{Resp}\Rsign Strengthen me, O King, Who reignest over the holy ones. \Star{} Put Thou in my mouth clear and well-sounding words.\end{Resp}
\begin{Resp}\Vsign O Lord, King of all forces, turn back their device upon themselves.\end{Resp}
\begin{Resp}\Rsign Put Thou in my mouth clear and well-sounding words.\end{Resp}
\switchcolumn*

\LA
\LessonHead{Lectio III}
\switchcolumn
\EN
\LessonHead{Lesson III}
\switchcolumn*

\LA
\Cite{Esth 1:7-9}
\switchcolumn
\EN
\Cite{Esth 1:7-9}
\switchcolumn*

\LA
\noindent Bibébant autem, qui invitáti erant, áureis póculis et áliis atque áliis vasis cibi inferebántur. Vinum quoque, ut magnificéntia régia dignum erat, abúndans et præcípuum ponebátur. Nec erat qui noléntes cógeret ad bibéndum, sed, sicut rex statúerat præpónens mensis síngulos de princípibus suis, ut súmeret unusquísque quod vellet. Vasthi quoque regína fecit convívium feminárum in palátio, ubi rex Assuérus manére consuéverat.\par
\switchcolumn
\EN
\noindent And they that were invited, drank in golden cups, and the meats were brought in diverse vessels one after another. Wine also in abundance and of the best was presented, as was worthy of a king's magnificence. Neither was there any one to compel them to drink that were not willing, but as the king had appointed, who set over every table one of his nobles, that every man might take what he would. Also Vasthi the queen made a feast for the women in the palace, where king Assuerus was used to dwell.\par
\switchcolumn*

\LA
\begin{Resp}\Rsign Spem in álium nunquam hábui, prætérquam in te, Deus Israël: \Star{} Qui irásceris, et propítius eris, et ómnia peccáta hóminum in tribulatióne dimíttis.\end{Resp}
\begin{Resp}\Vsign Dómine Deus, Creátor cæli et terræ, réspice ad humilitátem nostram.\end{Resp}
\begin{Resp}\Rsign Qui irásceris, et propítius eris, et ómnia peccáta hóminum in tribulatióne dimíttis.\end{Resp}
\begin{Resp}\Vsign Glória Patri, et Fílio, \Star{} et Spirítui Sancto.\end{Resp}
\begin{Resp}\Rsign Qui irásceris, et propítius eris, et ómnia peccáta hóminum in tribulatióne dimíttis.\end{Resp}
\switchcolumn
\EN
\begin{Resp}\Rsign I have had no hope in any other but in thee, O God of Israel, Who art angry, and wilt again show mercy, \Star{} Thou that forgivest all the sins of men when they are in affliction.\end{Resp}
\begin{Resp}\Vsign O Lord God, Maker of heaven and earth, look upon our low estate;\end{Resp}
\begin{Resp}\Rsign Thou that forgivest all the sins of men when they are in affliction.\end{Resp}
\begin{Resp}\Vsign Glory be to the Father, and to the Son, \Star{} and to the Holy Ghost.\end{Resp}
\begin{Resp}\Rsign Thou that forgivest all the sins of men when they are in affliction.\end{Resp}
\switchcolumn*

\LA
\LessonHead{Lectio IV}
\switchcolumn
\EN
\LessonHead{Lesson IV}
\switchcolumn*

\LA
\LessonTitle{Ex libro Officiórum sancti Ambrósii Epíscopi}
\switchcolumn
\EN
\LessonTitle{From the Book On Duties written by St. Ambrose, Bishop (of Milan.)}
\switchcolumn*

\LA
\Source{Lectio iv, Liber 3, Cap. 15}
\switchcolumn
\EN
\Source{Bk. iii. ch. 15}
\switchcolumn*

\LA
\noindent Quid Esther regína, nonne ut pópulum suum perículo erúeret (quod erat decórum atque honéstum) morti se óbtulit, nec immítis regis trepidávit furórem? Ipse quoque rex Persárum, ferox atque túmido corde, tamen decórum judicávit índici insidiárum, quæ sibi parátæ forent, grátiam repræsentáre, populúmque líberum a servitúte erípere, erúere neci, nec párcere neci ejus, qui tam indecóra suasísset. Dénique quem secúndum a se, ac præcípuum inter omnes amícos habéret, cruci trádidit, quod dehonestátum se ejus frauduléntis consíliis animadvertísset.\par
\switchcolumn
\EN
\noindent What did Queen Esther? Did she not, to save her people from danger (a beautiful and noble object) put herself in jeopardy of death, and face the anger of the cruel King The King of the Persians, cruel and violent as he was, nevertheless, thought it seemly to show grace unto him that told him of the plot that was made against him, to free the people from bondage, and to deliver them from death, but not to spare him that had persuaded such iniquity. In the end he went up to the gallows, whom he had held second only to himself, and chiefest among all his friends, because he found himself dishonoured through his false counsels.\par
\switchcolumn*

\LA
\begin{Resp}\Rsign Meménto mei, Dómine Deus, in bonum: \Star{} Et ne déleas miseratiónes meas, quas feci in domo Dei mei et in cæremóniis ejus.\end{Resp}
\begin{Resp}\Vsign Recordáre mei, Dómine, Deus meus.\end{Resp}
\begin{Resp}\Rsign Et ne déleas miseratiónes meas, quas feci in domo Dei mei et in cæremóniis ejus.\end{Resp}
\switchcolumn
\EN
\begin{Resp}\Rsign Remember me for good, O Lord God, \Star{} And put not away my works of mercy, which I have wrought in the house of my God, and in the times of His solemn rites.\end{Resp}
\begin{Resp}\Vsign Remember me, O Lord my God\end{Resp}
\begin{Resp}\Rsign And put not away my works of mercy, which I have wrought in the house of my God, and in the times of His solemn rites.\end{Resp}
\switchcolumn*

\LA
\LessonHead{Lectio V}
\switchcolumn
\EN
\LessonHead{Lesson V}
\switchcolumn*

\LA
\Source{Cap. 16}
\switchcolumn
\EN
\Source{Cap. 16}
\switchcolumn*

\LA
\noindent Ea enim amicítia probábilis, quæ honestátem tuétur, præferénda sane ópibus, honóribus, potestátibus; honestáti vero præférri non solet, sed honestátem sequi. Qualis fuit Jónathæ, qui pro pietáte nec offénsam patris, nec salútis perículum refugiébat. Qualis fuit Achímelech, qui pro hospitális grátiæ offíciis necem pótius sibi, quam proditiónem fugiéntis amíci, subeúndam arbitrabátur. Nihil ígitur præferéndum honestáti; quæ tamen ne amicítiæ stúdio prætereátur, étiam hoc Scriptúra ádmonet.\par
\switchcolumn
\EN
\noindent That true friendship, which careth for honour, careth less for riches, or dignities, or power than for itself, but for honour before itself. Such was the friendship of Jonathan, which caused him to risk the anger of his father, and danger to himself. Such was the friendship of Achimelech, who chose to earn death for himself by giving relief to David, rather than to betray the outlaw. But before honour nothing is to be put, and friendship must not be allowed to outrun it, even as we are warned by the Scriptures.\par
\switchcolumn*

\LA
\begin{Resp}\Rsign Tribulatiónes civitátum audívimus, quas passæ sunt, et defécimus: timor et hebetúdo mentis cécidit super nos et super líberos nostros: ipsi montes nolunt recípere fugam nostram: \Star{} Dómine, miserére.\end{Resp}
\begin{Resp}\Vsign Peccávimus cum pátribus nostris, injúste égimus, iniquitátem fécimus.\end{Resp}
\begin{Resp}\Rsign Dómine, miserére.\end{Resp}
\switchcolumn
\EN
\begin{Resp}\Rsign We have heard of the tribulation of those cities, which they have suffered, and we have fainted. Fear and confusion of mind are fallen upon us. Even the mountains will not give us a refuge. \Star{} Lord, have mercy.\end{Resp}
\begin{Resp}\Vsign We have sinned like our forefathers, we have done unjustly, and wrought iniquity.\end{Resp}
\begin{Resp}\Rsign Lord, have mercy.\end{Resp}
\switchcolumn*

\LA
\LessonHead{Lectio VI}
\switchcolumn
\EN
\LessonHead{Lesson VI}
\switchcolumn*

\LA
\noindent Sunt enim plerǽque philosophórum quæstiónes: Utrum amíci causa quisquam contra pátriam sentíre necne débeat, ut amíco obédiat: utrum opórteat ut fidem déserat, dum indúlget atque inténdit amíci commoditátibus. Et Scriptúra quidem ait: Clava, et gládius, et sagítta ferráta, sic homo est testimónium dans falsum advérsus amícum suum. Sed consídera quid ástruat. Non testimónium reprehéndit dictum in amícum, sed falsum testimónium. Quid enim si Dei causa, quid si pátriæ, cogátur áliquis dícere testimónium? Numquid præponderáre debet amicítia religióni, præponderáre caritáti vítium?\par
\switchcolumn
\EN
\noindent The Philosophers have started diverse questions whether friendship can, or cannot justify disloyalty to a man's own country whether friendship can, or cannot justify serving a friend at the cost of breach of faith. Scripture indeed saith A man that beareth false witness against his neighbour, is a maul, and a sword, and a sharp arrow. (Prov. xxv. 18.) But mark that what is here condemned is not witness by itself, but false witness. How if a man be compelled to give such witness, for the sake of God or for the sake of his country? Ought friendship to outweigh religion Is not to say this, as much as to say that a sinful weakness is to outweigh a duty?\par
\switchcolumn*

\LA
\begin{Resp}\Rsign Benedícat te Dóminus in virtúte sua, qui per te ad níhilum redégit inimícos nostros: \Star{} Ut non defíciat laus tua de ore hóminum.\end{Resp}
\begin{Resp}\Vsign Benedíctus Dóminus qui creávit cælum et terram; quia hódie nomen tuum ita magnificávit.\end{Resp}
\begin{Resp}\Rsign Ut non defíciat laus tua de ore hóminum.\end{Resp}
\begin{Resp}\Vsign Glória Patri, et Fílio, \Star{} et Spirítui Sancto.\end{Resp}
\begin{Resp}\Rsign Ut non defíciat laus tua de ore hóminum.\end{Resp}
\switchcolumn
\EN
\begin{Resp}\Rsign The Lord bless thee by His power, Who hath brought our enemies to nought through thee. \Star{} And may the praise of thee never fail from the mouth of men.\end{Resp}
\begin{Resp}\Vsign Blessed be the Lord, Who hath created the heaven and the earth, because that He hath so glorified thy name this day.\end{Resp}
\begin{Resp}\Rsign And may the praise of thee never fail from the mouth of men.\end{Resp}
\begin{Resp}\Vsign Glory be to the Father, and to the Son, \Star{} and to the Holy Ghost.\end{Resp}
\begin{Resp}\Rsign And may the praise of thee never fail from the mouth of men.\end{Resp}
\switchcolumn*

\end{paracol}

\Office{Feria II infra Hebdomadam V Septembris}{Monday of the fifth week of September}{}
\addcontentsline{toc}{subsection}{Feria II infra Hebdomadam V Septembris \textit{— Monday of the fifth week of September}}
\begin{paracol}{2}
\LA
\LessonHead{Lectio I}
\switchcolumn
\EN
\LessonHead{Lesson I}
\switchcolumn*

\LA
\LessonTitle{De libro Esther}
\switchcolumn
\EN
\LessonTitle{Lesson from the book of Esther}
\switchcolumn*

\LA
\Cite{Esth 2:5-7}
\switchcolumn
\EN
\Cite{Esth 2:5-7}
\switchcolumn*

\LA
\noindent Erat vir Judǽus in Susan civitáte, vocábulo Mardochǽus, fílius Jair, fílii Sémei, fílii Cis, de stirpe Jémini, qui translátus fúerat de Jerúsalem eo témpore quo Jechoníam regem Juda Nabuchodónosor rex Babylónis transtúlerat. Qui fuit nutrítius fíliæ fratris sui Edíssæ, quæ áltero nómine vocabátur Esther, et utrúmque paréntem amíserat: pulchra nimis et decóra fácie. Mortuísque patre ejus ac matre, Mardochǽus sibi eam adoptávit in fíliam.\par
\switchcolumn
\EN
\noindent There was a man in the city of Susan, a Jew, named Mardochai, the son of Jair, the son of Semei, the son of Cis, of the race of Jemini, Who had been carried away from Jerusalem at the time that Nabuchodonosor king of Babylon carried away Jechonias king of Juda, And he had brought up his brother's daughter Edissa, who by another name was called Esther: now she had lost both her parents: and was exceeding fair and beautiful. And her father and mother being dead, Mardochai adopted her for his daughter.\par
\switchcolumn*

\LA
\begin{Resp}\Rsign Dómine, mi Rex omnípotens, in dicióne tua cuncta sunt pósita, et non est qui possit resístere voluntáti tuæ: \Star{} Líbera nos propter nomen tuum.\end{Resp}
\begin{Resp}\Vsign Exáudi oratiónem nostram, et convérte luctum nostrum in gáudium.\end{Resp}
\begin{Resp}\Rsign Líbera nos propter nomen tuum.\end{Resp}
\switchcolumn
\EN
\begin{Resp}\Rsign O my Lord, the King Almighty All things are in thy power, and there is no man that can gainsay thy will; \Star{} Deliver us for thy Name's sake.\end{Resp}
\begin{Resp}\Vsign Hear our prayer, and turn our sorrow into joy.\end{Resp}
\begin{Resp}\Rsign Deliver us for thy Name's sake.\end{Resp}
\switchcolumn*

\LA
\LessonHead{Lectio II}
\switchcolumn
\EN
\LessonHead{Lesson II}
\switchcolumn*

\LA
\Cite{Esth 2:8-11}
\switchcolumn
\EN
\Cite{Esth 2:8-11}
\switchcolumn*

\LA
\noindent Cumque percrebuísset regis impérium, et juxta mandátum illíus multæ pulchræ vírgines adduceréntur Susan et Egéo traderéntur eunúcho; Esther quoque inter céteras puéllas ei trádita est, ut servarétur in número feminárum. Quæ plácuit ei et invénit grátiam in conspéctu illíus. Et præcépit eunúcho ut acceleráret mundum mulíebrem et tráderet ei partes suas et septem puéllas speciosíssimas de domo regis, et tam ipsam quam pedíssequas ejus ornáret atque excóleret. Quæ nóluit indicáre ei pópulum et pátriam suam; Mardochǽus enim præcéperat ei, ut de hac re omníno reticéret. Qui deambulábat cotídie ante vestíbulum domus, in qua eléctæ vírgines servabántur, curam agens salútis Esther.\par
\switchcolumn
\EN
\noindent And when the king's ordinance was noised abroad, and according to his commandment many beautiful virgins were brought to Susan, and were delivered to Egeus the eunuch: Esther also among the rest of the maidens was delivered to him to be kept in the number of the women. And she pleased him, and found favour in his sight. And he commanded the eunuch to hasten the women's ornaments, and to deliver to her her part, and seven of the most beautiful maidens of the king's house, and to adorn and deck out both her and her waiting maids. And she would not tell him her people nor her country. For Mardochai had charged her to say nothing at all of that: And he walked every day before the court of the house, in which the chosen virgins werre kept, having a care for Esther's welfare, and desiring to know what would befall her.\par
\switchcolumn*

\LA
\begin{Resp}\Rsign Confórta me, Rex, Sanctórum principátum tenens: \Star{} Et da sermónem rectum et bene sonántem in os meum.\end{Resp}
\begin{Resp}\Vsign Dómine, Rex univérsæ potestátis, convérte consílium eórum super eos.\end{Resp}
\begin{Resp}\Rsign Et da sermónem rectum et bene sonántem in os meum.\end{Resp}
\switchcolumn
\EN
\begin{Resp}\Rsign Strengthen me, O King, Who reignest over the holy ones. \Star{} Put Thou in my mouth clear and well-sounding words.\end{Resp}
\begin{Resp}\Vsign O Lord, King of all forces, turn back their device upon themselves.\end{Resp}
\begin{Resp}\Rsign Put Thou in my mouth clear and well-sounding words.\end{Resp}
\switchcolumn*

\LA
\LessonHead{Lectio III}
\switchcolumn
\EN
\LessonHead{Lesson III}
\switchcolumn*

\LA
\Cite{Esth 2:15-17}
\switchcolumn
\EN
\Cite{Esth 2:15-17}
\switchcolumn*

\LA
\noindent Evolúto autem témpore per órdinem, instábat dies quo Esther fília Abíhail fratris Mardochǽi, quam sibi adoptáverat in fíliam, debéret intráre ad regem. Quæ non quæsívit mulíebrem cultum, sed quæcúmque vóluit Egéus eunúchus custos vírginum, hæc ei ad ornátum dedit; erat enim formósa valde et incredíbili pulchritúdine: ómnium óculis gratiósa et amábilis videbátur. Ducta est ítaque ad cubículum regis Assuéri mense décimo, qui vocátur Tebeth, séptimo anno regni ejus. Et adamávit eam rex plus quam omnes mulíeres, habuítque grátiam et misericórdiam coram eo super omnes mulíeres: et pósuit diadéma regni in cápite ejus fecítque eam regnáre in loco Vasthi.\par
\switchcolumn
\EN
\noindent And as the time came orderly about, the day was at hand, when Esther, the daughter of Abihail the brother of Mardochai, whom he had adopted for his daughter, was to go in to the king. But she sought not women's ornaments, but whatsoever Egeus the eunuch the keeper of the virgins had a mind, he gave her to adorn her. For she was exceeding fair, and her incredible beauty made her appear agreeable and amiable in the eyes of all. So she was brought to the chamber of king Assuerus the tenth month, which is called Tebeth, in the seventh year of his reign. And the king loved her more than all the women, and she had favour and kindness before him above all the women, and he set the royal crown on her head, and made her queen instead of Vasthi.\par
\switchcolumn*

\LA
\begin{Resp}\Rsign Spem in álium nunquam hábui, prætérquam in te, Deus Israël: \Star{} Qui irásceris, et propítius eris, et ómnia peccáta hóminum in tribulatióne dimíttis.\end{Resp}
\begin{Resp}\Vsign Dómine Deus, Creátor cæli et terræ, réspice ad humilitátem nostram.\end{Resp}
\begin{Resp}\Rsign Qui irásceris, et propítius eris, et ómnia peccáta hóminum in tribulatióne dimíttis.\end{Resp}
\begin{Resp}\Vsign Glória Patri, et Fílio, \Star{} et Spirítui Sancto.\end{Resp}
\begin{Resp}\Rsign Qui irásceris, et propítius eris, et ómnia peccáta hóminum in tribulatióne dimíttis.\end{Resp}
\switchcolumn
\EN
\begin{Resp}\Rsign I have had no hope in any other but in thee, O God of Israel, Who art angry, and wilt again show mercy, \Star{} Thou that forgivest all the sins of men when they are in affliction.\end{Resp}
\begin{Resp}\Vsign O Lord God, Maker of heaven and earth, look upon our low estate;\end{Resp}
\begin{Resp}\Rsign Thou that forgivest all the sins of men when they are in affliction.\end{Resp}
\begin{Resp}\Vsign Glory be to the Father, and to the Son, \Star{} and to the Holy Ghost.\end{Resp}
\begin{Resp}\Rsign Thou that forgivest all the sins of men when they are in affliction.\end{Resp}
\switchcolumn*

\end{paracol}

\Office{Feria III infra Hebdomadam V Septembris}{Tuesday of the fifth week of September}{}
\addcontentsline{toc}{subsection}{Feria III infra Hebdomadam V Septembris \textit{— Tuesday of the fifth week of September}}
\begin{paracol}{2}
\LA
\LessonHead{Lectio I}
\switchcolumn
\EN
\LessonHead{Lesson I}
\switchcolumn*

\LA
\LessonTitle{De libro Esther}
\switchcolumn
\EN
\LessonTitle{Lesson from the book of Esther}
\switchcolumn*

\LA
\Cite{Esth 3:1-3}
\switchcolumn
\EN
\Cite{Esth 3:1-3}
\switchcolumn*

\LA
\noindent Post hæc rex Assuérus exaltávit Aman fílium Amadáthi, qui erat de stirpe Agag, et pósuit sólium ejus super omnes príncipes quos habébat. Cunctíque servi regis, qui in fóribus palátii versabántur, flectébant génua et adorábant Aman; sic enim præcéperat eis imperátor. Solus Mardochǽus non flectébat genu neque adorábat eum. Cui dixérunt púeri regis, qui ad fores palátii præsidébant: Cur præter céteros non obsérvas mandátum regis?\par
\switchcolumn
\EN
\noindent After these things, king Assuerus advanced Aman, the son of Amadathi, who was of the race of Agag: and he set his throne above all the princes that were with him. And all the king's servants, that were at the doors of the palace, bent their knees, and worshipped Aman: for so the emperor had commanded them, only Mardochai did not bend his knee, nor worship him. And the king's servants that were chief at the doors of the palace, said to him: Why dost thou alone not observe the king's commandment?\par
\switchcolumn*

\LA
\begin{Resp}\Rsign Meménto mei, Dómine Deus, in bonum: \Star{} Et ne déleas miseratiónes meas, quas feci in domo Dei mei et in cæremóniis ejus.\end{Resp}
\begin{Resp}\Vsign Recordáre mei, Dómine, Deus meus.\end{Resp}
\begin{Resp}\Rsign Et ne déleas miseratiónes meas, quas feci in domo Dei mei et in cæremóniis ejus.\end{Resp}
\switchcolumn
\EN
\begin{Resp}\Rsign Remember me for good, O Lord God, \Star{} And put not away my works of mercy, which I have wrought in the house of my God, and in the times of His solemn rites.\end{Resp}
\begin{Resp}\Vsign Remember me, O Lord my God\end{Resp}
\begin{Resp}\Rsign And put not away my works of mercy, which I have wrought in the house of my God, and in the times of His solemn rites.\end{Resp}
\switchcolumn*

\LA
\LessonHead{Lectio II}
\switchcolumn
\EN
\LessonHead{Lesson II}
\switchcolumn*

\LA
\Cite{Esth 3:4-6}
\switchcolumn
\EN
\Cite{Esth 3:4-6}
\switchcolumn*

\LA
\noindent Cumque hoc crébrius dícerent, et ille nollet audíre, nuntiavérunt Aman scire cupiéntes utrum perseveráret in senténtia; díxerat enim eis se esse Judǽum. Quod cum audísset Aman et experiménto probásset quod Mardochǽus non flécteret sibi genu nec se adoráret, irátus est valde et pro níhilo duxit in unum Mardochǽum míttere manus suas.\par
\switchcolumn
\EN
\noindent And when they were saying this often, and he would not hearken to them; they told Aman, desirous to know whether he would continue in his resolution: for he had told them that he was a Jew. Now when Aman had heard this, and had proved by experience that Mardochai did not bend his knee to him, nor worship him, he was exceeding angry. And he counted it nothing to lay his hands upon Mardochai alone: for he had heard that he was of the nation of the Jews, and he chose rather to destroy all the nation of the Jews that were in the kingdom of Assuerus.\par
\switchcolumn*

\LA
\begin{Resp}\Rsign Tribulatiónes civitátum audívimus, quas passæ sunt, et defécimus: timor et hebetúdo mentis cécidit super nos et super líberos nostros: ipsi montes nolunt recípere fugam nostram: \Star{} Dómine, miserére.\end{Resp}
\begin{Resp}\Vsign Peccávimus cum pátribus nostris, injúste égimus, iniquitátem fécimus.\end{Resp}
\begin{Resp}\Rsign Dómine, miserére.\end{Resp}
\switchcolumn
\EN
\begin{Resp}\Rsign We have heard of the tribulation of those cities, which they have suffered, and we have fainted. Fear and confusion of mind are fallen upon us. Even the mountains will not give us a refuge. \Star{} Lord, have mercy.\end{Resp}
\begin{Resp}\Vsign We have sinned like our forefathers, we have done unjustly, and wrought iniquity.\end{Resp}
\begin{Resp}\Rsign Lord, have mercy.\end{Resp}
\switchcolumn*

\LA
\LessonHead{Lectio III}
\switchcolumn
\EN
\LessonHead{Lesson III}
\switchcolumn*

\LA
\Cite{Esth 3:6-7}
\switchcolumn
\EN
\Cite{Esth 3:6-7}
\switchcolumn*

\LA
\noindent Audíerat enim quod esset gentis Judǽæ, magísque vóluit omnem Judæórum, qui erant in regno Assuéri, pérdere natiónem. Mense primo, cujus vocábulum est Nisan, anno duodécimo regni Assuéri, missa est sors in urnam, quæ Hebráice dícitur phur, coram Aman, quo die et quo mense gens Judæórum debéret intérfici; et exívit mensis duodécimus qui vocátur Adar.\par
\switchcolumn
\EN
\noindent And he counted it nothing to lay his hands upon Mardochai alone: for he had heard that he was of the nation of the Jews, and he chose rather to destroy all the nation of the Jews that were in the kingdom of Assuerus. In the first month (which is called Nisan) in the twelfth year of the reign of Assuerus, the lot was cast into an urn, which in Hebrew is called Phur, before Aman, on what day and what month the nation of the Jews should be destroyed: and there came out the twelfth month, which is called Adar.\par
\switchcolumn*

\LA
\begin{Resp}\Rsign Benedícat te Dóminus in virtúte sua, qui per te ad níhilum redégit inimícos nostros: \Star{} Ut non defíciat laus tua de ore hóminum.\end{Resp}
\begin{Resp}\Vsign Benedíctus Dóminus qui creávit cælum et terram; quia hódie nomen tuum ita magnificávit.\end{Resp}
\begin{Resp}\Rsign Ut non defíciat laus tua de ore hóminum.\end{Resp}
\begin{Resp}\Vsign Glória Patri, et Fílio, \Star{} et Spirítui Sancto.\end{Resp}
\begin{Resp}\Rsign Ut non defíciat laus tua de ore hóminum.\end{Resp}
\switchcolumn
\EN
\begin{Resp}\Rsign The Lord bless thee by His power, Who hath brought our enemies to nought through thee. \Star{} And may the praise of thee never fail from the mouth of men.\end{Resp}
\begin{Resp}\Vsign Blessed be the Lord, Who hath created the heaven and the earth, because that He hath so glorified thy name this day.\end{Resp}
\begin{Resp}\Rsign And may the praise of thee never fail from the mouth of men.\end{Resp}
\begin{Resp}\Vsign Glory be to the Father, and to the Son, \Star{} and to the Holy Ghost.\end{Resp}
\begin{Resp}\Rsign And may the praise of thee never fail from the mouth of men.\end{Resp}
\switchcolumn*

\end{paracol}

\Office{Feria IV infra Hebdomadam V Septembris}{Wednesday of the fifth week of September}{}
\addcontentsline{toc}{subsection}{Feria IV infra Hebdomadam V Septembris \textit{— Wednesday of the fifth week of September}}
\begin{paracol}{2}
\LA
\LessonHead{Lectio I}
\switchcolumn
\EN
\LessonHead{Lesson I}
\switchcolumn*

\LA
\LessonTitle{De libro Esther}
\switchcolumn
\EN
\LessonTitle{Lesson from the book of Esther}
\switchcolumn*

\LA
\Cite{Esth 4:1-5}
\switchcolumn
\EN
\Cite{Esth 4:1-5}
\switchcolumn*

\LA
\noindent Quæ cum audísset Mardochǽus, scidit vestiménta sua et indútus est sacco spargens cínerem cápiti; et in platéa médiæ civitátis voce magna clamábat osténdens amaritúdinem ánimi sui, et hoc ejulátu usque ad fores palátii grádiens; non enim erat lícitum indútum sacco aulam regis intráre. In ómnibus quoque provínciis, óppidis ac locis, ad quæ crudéle regis dogma pervénerat, planctus ingens erat apud Judǽos, jejúnium, ululátus et fletus, sacco et cínere multis pro strato uténtibus. Ingréssæ autem sunt puéllæ Esther et eunúchi nuntiaverúntque ei. Quod áudiens consternáta est et vestem misit, ut, abláto sacco, indúerent eum: quam accípere nóluit. Accitóque Athach eunúcho, quem rex minístrum ei déderat, præcépit ei ut iret ad Mardochǽum et dísceret ab eo cur hoc fáceret.\par
\switchcolumn
\EN
\noindent Now when Mardochai had heard these things, he rent his garments, and put on sackcloth, strewing ashes on his head: and he cried with a loud voice in the street in the midst of the city, shewing the anguish of his mind. And he came lamenting in this manner even to the gate of the palace: for no one clothed with sackcloth might enter the king's court. And in all provinces, towns, and places, to which the king's cruel edict was come, there was great mourning among the Jews, with fasting, wailing, and weeping, many using sackcloth and ashes for their bed. Then Esther's maids and her eunuchs went in, and told her. And when she heard it she was in a consternation: and she sent a garment, to clothe him, and to take away the sackcloth: but he would not receive it. And she called for Athach the eunuch, whom the king had appointed to attend upon her, and she commanded him to go to Mardochai, and learn of him why he did this.\par
\switchcolumn*

\LA
\begin{Resp}\Rsign Nos álium Deum nescímus præter Dóminum, in quo sperámus: \Star{} Qui non déspicit nos, nec ámovet salútem suam a génere nostro.\end{Resp}
\begin{Resp}\Vsign Indulgéntiam ipsíus fusis lácrimis postulémus, et humiliémus illi ánimas nostras.\end{Resp}
\begin{Resp}\Rsign Qui non déspicit nos, nec ámovet salútem suam a génere nostro.\end{Resp}
\switchcolumn
\EN
\begin{Resp}\Rsign We know no strange God before the Lord. In Him we trust. \Star{} He despiseth us not, neither putteth He away His salvation from our nation.\end{Resp}
\begin{Resp}\Vsign His mercy let us seek with tears, and humble our souls before Him.\end{Resp}
\begin{Resp}\Rsign He despiseth us not, neither putteth He away His salvation from our nation.\end{Resp}
\switchcolumn*

\LA
\LessonHead{Lectio II}
\switchcolumn
\EN
\LessonHead{Lesson II}
\switchcolumn*

\LA
\Cite{Esth 4:6-11}
\switchcolumn
\EN
\Cite{Esth 4:6-11}
\switchcolumn*

\LA
\noindent Egressúsque Athach, ivit ad Mardochǽum stantem in platéa civitátis, ante óstium palátii. Qui indicávit ei ómnia quæ accíderant, quómodo Aman promisísset, ut in thesáuros regis pro Judæórum nece inférret argéntum; exémplar quoque edícti, quod pendébat in Susan, dedit ei ut regínæ osténderet et monéret eam ut intráret ad regem et deprecarétur eum pro pópulo suo. Regréssus Athach nuntiávit Esther ómnia quæ Mardochǽus díxerat. Quæ respóndit ei, et jussit ut díceret Mardochǽo: Omnes servi regis et cunctæ quæ sub dicióne ejus sunt norunt provínciæ, quod sive vir sive múlier non vocátus intérius átrium regis intráverit, absque ulla cunctatióne statim interficiátur, nisi forte rex áuream virgam ad eum teténderit pro signo cleméntiæ, atque ita possit vívere. Ego ígitur quómodo ad regem intráre pótero, quæ trigínta jam diébus non sum vocáta ad eum?\par
\switchcolumn
\EN
\noindent And Athach going out went to Mardochai, who was standing in the street of the city, before the palace gate: And Mardochai told him all that had happened, how Aman had promised to pay money into the king's treasures, to have the Jews destroyed. He gave him also a copy of the edict which was hanging up in Susan, that he should shew it to the queen, and admonish her to go in to the king, and to en- treat him for her people. And Athach went back and told Esther all that Mardochai had said. She answered him, and bade him say to Mardochai: All the king's servants, and all the provinces that are under his dominion, know, that whosoever, whether man or woman, cometh into the king's inner court, who is not called for, is immediately to be put to death without any delay: except the king shall hold out the golden sceptre to him, in token of clemency, that so he may live. How then can I go in to the king, who for these thirty days now have not been called unto him?\par
\switchcolumn*

\LA
\begin{Resp}\Rsign Confórta me, Rex, Sanctórum principátum tenens: \Star{} Et da sermónem rectum et bene sonántem in os meum.\end{Resp}
\begin{Resp}\Vsign Dómine, Rex univérsæ potestátis, convérte consílium eórum super eos.\end{Resp}
\begin{Resp}\Rsign Et da sermónem rectum et bene sonántem in os meum.\end{Resp}
\switchcolumn
\EN
\begin{Resp}\Rsign Strengthen me, O King, Who reignest over the holy ones. \Star{} Put Thou in my mouth clear and well-sounding words.\end{Resp}
\begin{Resp}\Vsign O Lord, King of all forces, turn back their device upon themselves.\end{Resp}
\begin{Resp}\Rsign Put Thou in my mouth clear and well-sounding words.\end{Resp}
\switchcolumn*

\LA
\LessonHead{Lectio III}
\switchcolumn
\EN
\LessonHead{Lesson III}
\switchcolumn*

\LA
\Cite{Esth 4:12-17}
\switchcolumn
\EN
\Cite{Esth 4:12-17}
\switchcolumn*

\LA
\noindent Quod cum audísset Mardochǽus, rursum mandávit Esther dicens: Ne putes quod ánimam tuam tantum líberes, quia in domo regis es, præ cunctis Judǽis; si enim nunc silúeris, per áliam occasiónem liberabúntur Judǽi, et tu et domus patris tui períbitis. Et quis novit utrum idcírco ad regnum véneris, ut in tali témpore pararéris? Rursúmque Esther hæc Mardochǽo verba mandávit: Vade et cóngrega omnes Judǽos, quos in Susan repéreris, et oráte pro me. Non comedátis et non bibátis tribus diébus et tribus nóctibus, et ego cum ancíllis meis simíliter jejunábo; et tunc ingrédiar ad regem contra legem fáciens non vocáta tradénsque me morti et perículo. Ivit ítaque Mardochǽus, et fecit ómnia quæ ei Esther præcéperat.\par
\switchcolumn
\EN
\noindent And when Mardochai had heard this, He sent word to Esther again, saying: Think not that thou mayst save thy life only, because thou art in the king a house, more than all the Jews: For if thou wilt now hold thy peace, the Jews shall be delivered by some other occasion: and thou, and thy father's house shall perish. And who knoweth whether thou art not therefore come to the kingdom, that thou mightest be ready in such a time as this? And again Esther sent to Mardochai in these words: Go, and gather together all the Jews whom thou shalt find in Susan, and pray ye for me. Neither eat nor drink for three days and three nights: and I with my handmaids will fast in like manner, and then I will go in to the king, against the law, not being called, and expose myself to death and to danger. So Mardochai went, and did all that Esther had commanded him.\par
\switchcolumn*

\LA
\begin{Resp}\Rsign Spem in álium nunquam hábui, prætérquam in te, Deus Israël: \Star{} Qui irásceris, et propítius eris, et ómnia peccáta hóminum in tribulatióne dimíttis.\end{Resp}
\begin{Resp}\Vsign Dómine Deus, Creátor cæli et terræ, réspice ad humilitátem nostram.\end{Resp}
\begin{Resp}\Rsign Qui irásceris, et propítius eris, et ómnia peccáta hóminum in tribulatióne dimíttis.\end{Resp}
\begin{Resp}\Vsign Glória Patri, et Fílio, \Star{} et Spirítui Sancto.\end{Resp}
\begin{Resp}\Rsign Qui irásceris, et propítius eris, et ómnia peccáta hóminum in tribulatióne dimíttis.\end{Resp}
\switchcolumn
\EN
\begin{Resp}\Rsign I have had no hope in any other but in thee, O God of Israel, Who art angry, and wilt again show mercy, \Star{} Thou that forgivest all the sins of men when they are in affliction.\end{Resp}
\begin{Resp}\Vsign O Lord God, Maker of heaven and earth, look upon our low estate;\end{Resp}
\begin{Resp}\Rsign Thou that forgivest all the sins of men when they are in affliction.\end{Resp}
\begin{Resp}\Vsign Glory be to the Father, and to the Son, \Star{} and to the Holy Ghost.\end{Resp}
\begin{Resp}\Rsign Thou that forgivest all the sins of men when they are in affliction.\end{Resp}
\switchcolumn*

\end{paracol}

\Office{Feria V infra Hebdomadam V Septembris}{Thursday of the fifth week of September}{}
\addcontentsline{toc}{subsection}{Feria V infra Hebdomadam V Septembris \textit{— Thursday of the fifth week of September}}
\begin{paracol}{2}
\LA
\LessonHead{Lectio I}
\switchcolumn
\EN
\LessonHead{Lesson I}
\switchcolumn*

\LA
\LessonTitle{De libro Esther}
\switchcolumn
\EN
\LessonTitle{Lesson from the book of Esther}
\switchcolumn*

\LA
\Cite{Esth 5:1-5}
\switchcolumn
\EN
\Cite{Esth 5:1-5}
\switchcolumn*

\LA
\noindent Die autem tértio indúta est Esther regálibus vestiméntis et stetit in átrio domus régiæ, quod erat intérius contra basílicam regis. At ille sedébat super sólium suum in consistório palátii contra óstium domus. Cumque vidísset Esther regínam stantem, plácuit óculis ejus et exténdit contra eam virgam áuream quam tenébat manu. Quæ accédens osculáta est summitátem virgæ ejus, dixítque ad eam rex: Quid vis, Esther regína? quæ est petítio tua? Etiam si dimídiam partem regni petíeris, dábitur tibi. At illa respóndit: Si regi placet, óbsecro ut vénias ad me hódie et Aman tecum ad convívium quod parávi. Statímque rex, Vocáte, inquit, cito Aman ut Esther obédiat voluntáti.\par
\switchcolumn
\EN
\noindent And on the third day Esther put on her royal apparel, and stood in the inner court of the king's house, over against the king's hall: now he sat upon his throne in the hall of the palace, over against the door of the house. And when he saw Esther the queen standing, she pleased his eyes, and he held out toward her the golden sceptre, which he held in his hand: and she drew near, and kissed the top of his sceptre. And the king said to her: What wilt then, queen Esther? what is thy request? if thou shouldst even ask one half of the kingdom, it shall be given to thee. But she answered: If it please the king. I beseech thee to come to me this day, and Aman with thee to the banquet which I have prepared. And the king said forthwith: Call ye Aman quickly, that he may obey Esther's will. So the king and Aman came to the banquet which the queen had prepared for them.\par
\switchcolumn*

\LA
\begin{Resp}\Rsign Dómine, mi Rex omnípotens, in dicióne tua cuncta sunt pósita, et non est qui possit resístere voluntáti tuæ: \Star{} Líbera nos propter nomen tuum.\end{Resp}
\begin{Resp}\Vsign Exáudi oratiónem nostram, et convérte luctum nostrum in gáudium.\end{Resp}
\begin{Resp}\Rsign Líbera nos propter nomen tuum.\end{Resp}
\switchcolumn
\EN
\begin{Resp}\Rsign O my Lord, the King Almighty All things are in thy power, and there is no man that can gainsay thy will; \Star{} Deliver us for thy Name's sake.\end{Resp}
\begin{Resp}\Vsign Hear our prayer, and turn our sorrow into joy.\end{Resp}
\begin{Resp}\Rsign Deliver us for thy Name's sake.\end{Resp}
\switchcolumn*

\LA
\LessonHead{Lectio II}
\switchcolumn
\EN
\LessonHead{Lesson II}
\switchcolumn*

\LA
\Cite{Esth 5:9-13}
\switchcolumn
\EN
\Cite{Esth 5:9-13}
\switchcolumn*

\LA
\noindent Egréssus est ítaque illo die Aman lætus et álacer. Cumque vidísset Mardochǽum sedéntem ante fores palátii, et non solum non assurrexísse sibi, sed nec motum quidem de loco sessiónis suæ, indignátus est valde; et, dissimuláta ira, revérsus in domum suam convocávit ad se amícos suos et Zares uxórem suam, et expósuit illis magnitúdinem divitiárum suárum filiorúmque turbam et quanta eum glória super omnes príncipes et servos suos rex elevásset. Et post hæc ait: Regína quoque Esther nullum álium vocávit ad convívium cum rege, præter me, apud quam étiam cras cum rege pransúrus sum; et cum hæc ómnia hábeam, nihil me habére puto, quámdiu vídero Mardochǽum Judǽum sedéntem ante fores régias.\par
\switchcolumn
\EN
\noindent So Aman went out that day joyful and merry. And when he saw Mardochai sitting before the gate of the palace, and that he not only did not rise up to honour him, but did not so much as move from the place where he sat, he was exceedingly angry: But dissembling his anger, and returning into his house, he called together to him his friends, and Zares his wife: And he declared to them the greatness of his riches, and the multitude of his children, and with how great glory the king had advanced him above all his princes and servants. And after this he said: Queen Esther also hath invited no other to the banquet with the king, but me: and with her I am also to dine to morrow with the king: And whereas I have all these things, I think I have nothing, so long as I see Mardochai the Jew sitting before the king's gate.\par
\switchcolumn*

\LA
\begin{Resp}\Rsign Confórta me, Rex, Sanctórum principátum tenens: \Star{} Et da sermónem rectum et bene sonántem in os meum.\end{Resp}
\begin{Resp}\Vsign Dómine, Rex univérsæ potestátis, convérte consílium eórum super eos.\end{Resp}
\begin{Resp}\Rsign Et da sermónem rectum et bene sonántem in os meum.\end{Resp}
\begin{Resp}\Vsign Glória Patri, et Fílio, \Star{} et Spirítui Sancto.\end{Resp}
\begin{Resp}\Rsign Et da sermónem rectum et bene sonántem in os meum.\end{Resp}
\switchcolumn
\EN
\begin{Resp}\Rsign Strengthen me, O King, Who reignest over the holy ones. \Star{} Put Thou in my mouth clear and well-sounding words.\end{Resp}
\begin{Resp}\Vsign O Lord, King of all forces, turn back their device upon themselves.\end{Resp}
\begin{Resp}\Rsign Put Thou in my mouth clear and well-sounding words.\end{Resp}
\begin{Resp}\Vsign Glory be to the Father, and to the Son, \Star{} and to the Holy Ghost.\end{Resp}
\begin{Resp}\Rsign Put Thou in my mouth clear and well-sounding words.\end{Resp}
\switchcolumn*

\LA
\LessonHead{Lectio III}
\switchcolumn
\EN
\LessonHead{Lesson III}
\switchcolumn*

\LA
\Cite{Esth 5:14}
\switchcolumn
\EN
\Cite{Esth 5:14}
\switchcolumn*

\LA
\noindent Responderúntque ei Zares uxor ejus et céteri amíci: Jube parári excélsam trabem habéntem altitúdinis quinquagínta cúbitos, et dic mane regi ut appendátur super eam Mardochǽus; et sic ibis cum rege lætus ad convívium. Plácuit ei consílium et jussit excélsam parári crucem.\par
\switchcolumn
\EN
\noindent Then Zares his wife, and the rest of his friends answered him: Order a great beam to be prepared, fifty cubits high, and in the morning speak to the king, that Mardochai may be hanged upon it, and so thou shalt go full of joy with the king to the banquet. The counsel pleased him, and he commanded a high gibbet to be prepared.\par
\switchcolumn*

\LA
\begin{Resp}\Rsign Spem in álium nunquam hábui, prætérquam in te, Deus Israël: \Star{} Qui irásceris, et propítius eris, et ómnia peccáta hóminum in tribulatióne dimíttis.\end{Resp}
\begin{Resp}\Vsign Dómine Deus, Creátor cæli et terræ, réspice ad humilitátem nostram.\end{Resp}
\begin{Resp}\Rsign Qui irásceris, et propítius eris, et ómnia peccáta hóminum in tribulatióne dimíttis.\end{Resp}
\begin{Resp}\Vsign Glória Patri, et Fílio, \Star{} et Spirítui Sancto.\end{Resp}
\begin{Resp}\Rsign Qui irásceris, et propítius eris, et ómnia peccáta hóminum in tribulatióne dimíttis.\end{Resp}
\switchcolumn
\EN
\begin{Resp}\Rsign I have had no hope in any other but in thee, O God of Israel, Who art angry, and wilt again show mercy, \Star{} Thou that forgivest all the sins of men when they are in affliction.\end{Resp}
\begin{Resp}\Vsign O Lord God, Maker of heaven and earth, look upon our low estate;\end{Resp}
\begin{Resp}\Rsign Thou that forgivest all the sins of men when they are in affliction.\end{Resp}
\begin{Resp}\Vsign Glory be to the Father, and to the Son, \Star{} and to the Holy Ghost.\end{Resp}
\begin{Resp}\Rsign Thou that forgivest all the sins of men when they are in affliction.\end{Resp}
\switchcolumn*

\end{paracol}

\Office{Feria VI infra Hebdomadam V Septembris}{Friday of the fifth week of September}{}
\addcontentsline{toc}{subsection}{Feria VI infra Hebdomadam V Septembris \textit{— Friday of the fifth week of September}}
\begin{paracol}{2}
\LA
\LessonHead{Lectio I}
\switchcolumn
\EN
\LessonHead{Lesson I}
\switchcolumn*

\LA
\LessonTitle{De libro Esther}
\switchcolumn
\EN
\LessonTitle{Lesson from the book of Esther}
\switchcolumn*

\LA
\Cite{Esth 6:1-5}
\switchcolumn
\EN
\Cite{Esth 6:1-5}
\switchcolumn*

\LA
\noindent Noctem illam duxit rex insómnem jussítque sibi afférri histórias et annáles priórum témporum. Quæ cum illo præsénte legeréntur, ventum est ad illum locum, ubi scriptum erat quómodo nuntiásset Mardochǽus insídias Bágathan et Thares eunuchórum regem Assuérum juguláre cupiéntium. Quod cum audísset rex, ait: Quid pro hac fide honóris ac prǽmii Mardochǽus consecútus est? Dixérunt ei servi illíus ac minístri: Nihil omníno mercédis accépit. Statímque rex, Quis est, inquit, in átrio? Aman quippe intérius átrium domus régiæ intráverat, ut suggéreret regi et jubéret Mardochǽum affígi patíbulo, quod ei fúerat præparátum. Respondérunt púeri: Aman stat in átrio. Dixítque rex: Ingrediátur.\par
\switchcolumn
\EN
\noindent That night the king passed without sleep, and he commanded the histories and chronicles of former times to be brought him. And when they were reading them before him, They came to that place where it was written, how Mardochai had discovered the treason of Bagathan and Thares the eunuchs, who sought to kill king Assuerus. And when the king heard this, he said: What honour and reward hath Mardochai received for this fidelity? His servants and ministers said to him: He hath received no reward at all. And the king said immediately: Who is in the court? for Aman was coming in to the inner court of the king's house, to speak to the king, that he might order Mardochai to be hanged upon the gibbet which was prepared for him. The servants answered: Aman standeth in the court, and the king said: Let him come in.\par
\switchcolumn*

\LA
\begin{Resp}\Rsign Meménto mei, Dómine Deus, in bonum: \Star{} Et ne déleas miseratiónes meas, quas feci in domo Dei mei et in cæremóniis ejus.\end{Resp}
\begin{Resp}\Vsign Recordáre mei, Dómine, Deus meus.\end{Resp}
\begin{Resp}\Rsign Et ne déleas miseratiónes meas, quas feci in domo Dei mei et in cæremóniis ejus.\end{Resp}
\switchcolumn
\EN
\begin{Resp}\Rsign Remember me for good, O Lord God, \Star{} And put not away my works of mercy, which I have wrought in the house of my God, and in the times of His solemn rites.\end{Resp}
\begin{Resp}\Vsign Remember me, O Lord my God\end{Resp}
\begin{Resp}\Rsign And put not away my works of mercy, which I have wrought in the house of my God, and in the times of His solemn rites.\end{Resp}
\switchcolumn*

\LA
\LessonHead{Lectio II}
\switchcolumn
\EN
\LessonHead{Lesson II}
\switchcolumn*

\LA
\Cite{Esth 6:6-9}
\switchcolumn
\EN
\Cite{Esth 6:6-9}
\switchcolumn*

\LA
\noindent Cumque esset ingréssus, ait illi: Quid debet fíeri viro, quem rex honoráre desíderat? Cógitans autem in corde suo Aman et réputans quod nullum álium rex nisi se vellet honoráre, respóndit: Homo quem rex honoráre cupit, debet índui véstibus régiis et impóni super equum qui de sella regis est, et accípere régium diadéma super caput suum, et primus de régiis princípibus ac tyránnis téneat equum ejus et per platéam civitátis incédens clamet, et dicat: Sic honorábitur quemcúmque volúerit rex honoráre.\par
\switchcolumn
\EN
\noindent And when he was come in, he said to him: What ought to be done to the man whom the king is desirous to honour? But Aman thinking in his heart, and supposing that the king would honour no other but himself, Answered: The man whom the king desireth to honour, Ought to be clothed with the king's apparel, and to be set upon the horse that the king rideth upon, and to have the royal crown upon his head, And let the first of the king's princes and nobles hold his horse, and going through the street of the city, proclaim before him and say: Thus shall he be honoured, whom the king hath a mind to honour.\par
\switchcolumn*

\LA
\begin{Resp}\Rsign Tribulatiónes civitátum audívimus, quas passæ sunt, et defécimus: timor et hebetúdo mentis cécidit super nos et super líberos nostros: ipsi montes nolunt recípere fugam nostram: \Star{} Dómine, miserére.\end{Resp}
\begin{Resp}\Vsign Peccávimus cum pátribus nostris, injúste égimus, iniquitátem fécimus.\end{Resp}
\begin{Resp}\Rsign Dómine, miserére.\end{Resp}
\begin{Resp}\Vsign Glória Patri, et Fílio, \Star{} et Spirítui Sancto.\end{Resp}
\begin{Resp}\Rsign Dómine, miserére.\end{Resp}
\switchcolumn
\EN
\begin{Resp}\Rsign We have heard of the tribulation of those cities, which they have suffered, and we have fainted. Fear and confusion of mind are fallen upon us. Even the mountains will not give us a refuge. \Star{} Lord, have mercy.\end{Resp}
\begin{Resp}\Vsign We have sinned like our forefathers, we have done unjustly, and wrought iniquity.\end{Resp}
\begin{Resp}\Rsign Lord, have mercy.\end{Resp}
\begin{Resp}\Vsign Glory be to the Father, and to the Son, \Star{} and to the Holy Ghost.\end{Resp}
\begin{Resp}\Rsign Lord, have mercy.\end{Resp}
\switchcolumn*

\LA
\LessonHead{Lectio III}
\switchcolumn
\EN
\LessonHead{Lesson III}
\switchcolumn*

\LA
\Cite{Esth 6:10-13}
\switchcolumn
\EN
\Cite{Esth 6:10-13}
\switchcolumn*

\LA
\noindent Dixítque ei rex: Festína et, sumpta stola et equo, fac, ut locútus es, Mardochǽo Judǽo, qui sedet ante fores palátii. Cave ne quidquam de his, quæ locútus es, prætermíttas. Tulit ítaque Aman stolam et equum, indutúmque Mardochǽum in platéa civitátis et impósitum equo præcedébat atque clamábat: Hoc honóre condígnus est quemcúmque rex volúerit honoráre. Reversúsque est Mardochǽus ad jánuam palátii, et Aman festinávit ire in domum suam lugens et opérto cápite. Narravítque Zares uxóri suæ et amícis ómnia quæ eveníssent sibi.\par
\switchcolumn
\EN
\noindent And the king said to him: Make haste and take the robe and the horse, and do as thou hast spoken to Mardochai the Jew, who sitteth before the gates of the palace. Beware thou pass over any of those things which thou hast spoken. So Aman took the robe and the horse, and arraying Mardochai in the street of the city, and setting him on the horse, went before him, and proclaimed: This honour is he worthy of, whom the king hath a mind to honour. But Mardochai returned to the palace gate: and Aman made haste to go to his house, mourning and having his head covered: And he told Zares his wife, and his friends, all that had befallen him. And the wise men whom he had in counsel, and his wife answered him: If Mardochai be of the seed of the Jews, before whom thou hast begun to fall, thou canst not resist him, but thou shalt fall in his sight.\par
\switchcolumn*

\LA
\begin{Resp}\Rsign Benedícat te Dóminus in virtúte sua, qui per te ad níhilum redégit inimícos nostros: \Star{} Ut non defíciat laus tua de ore hóminum.\end{Resp}
\begin{Resp}\Vsign Benedíctus Dóminus qui creávit cælum et terram; quia hódie nomen tuum ita magnificávit.\end{Resp}
\begin{Resp}\Rsign Ut non defíciat laus tua de ore hóminum.\end{Resp}
\begin{Resp}\Vsign Glória Patri, et Fílio, \Star{} et Spirítui Sancto.\end{Resp}
\begin{Resp}\Rsign Ut non defíciat laus tua de ore hóminum.\end{Resp}
\switchcolumn
\EN
\begin{Resp}\Rsign The Lord bless thee by His power, Who hath brought our enemies to nought through thee. \Star{} And may the praise of thee never fail from the mouth of men.\end{Resp}
\begin{Resp}\Vsign Blessed be the Lord, Who hath created the heaven and the earth, because that He hath so glorified thy name this day.\end{Resp}
\begin{Resp}\Rsign And may the praise of thee never fail from the mouth of men.\end{Resp}
\begin{Resp}\Vsign Glory be to the Father, and to the Son, \Star{} and to the Holy Ghost.\end{Resp}
\begin{Resp}\Rsign And may the praise of thee never fail from the mouth of men.\end{Resp}
\switchcolumn*

\end{paracol}

\Office{Sabbato infra Hebdomadam V Septembris}{Saturday of the fifth week of September}{}
\addcontentsline{toc}{subsection}{Sabbato infra Hebdomadam V Septembris \textit{— Saturday of the fifth week of September}}
\begin{paracol}{2}
\LA
\LessonHead{Lectio I}
\switchcolumn
\EN
\LessonHead{Lesson I}
\switchcolumn*

\LA
\LessonTitle{De libro Esther}
\switchcolumn
\EN
\LessonTitle{Lesson from the book of Esther}
\switchcolumn*

\LA
\Cite{Esth 7:1-4}
\switchcolumn
\EN
\Cite{Esth 7:1-4}
\switchcolumn*

\LA
\noindent Intrávit ítaque rex et Aman, ut bíberent cum regína. Dixítque ei rex étiam secúnda die, postquam vino incalúerat: Quæ est petítio tua, Esther, ut detur tibi, et quid vis fíeri? Etiam si dimídiam partem regni mei petíeris, impetrábis. Ad quem illa respóndit: Si invéni grátiam in óculis tuis, o rex, et si tibi placet, dona mihi ánimam meam, pro qua rogo, et pópulum meum, pro quo óbsecro. Tráditi enim sumus ego et pópulus meus, ut conterámur, jugulémur et pereámus. Atque, útinam in servos et fámulas venderémur: esset tolerábile malum et gemens tacérem; nunc autem hostis noster est cujus crudélitas redúndat in regem.\par
\switchcolumn
\EN
\noindent So the king and Aman went in, to drink with the queen. And the king said to her again the second day, after he was warm with wine: What is thy petition, Esther, that it may be granted thee? and what wilt thou have done: although thou ask the half of my kingdom, thou shalt have it. Then she answered: If I have found Favour in thy sight, O king, and if it please thee, give me my life for which I ask, and my people for which I request. For we are given up, I and my people, to be destroyed, to be slain, and to perish. And would God we were sold for bondmen and bondwomen: the evil might be borne with, and I would have mourned in silence: but now we have an enemy, whose cruelty redoundeth upon the king.\par
\switchcolumn*

\LA
\begin{Resp}\Rsign Nos álium Deum nescímus præter Dóminum, in quo sperámus: \Star{} Qui non déspicit nos, nec ámovet salútem suam a génere nostro.\end{Resp}
\begin{Resp}\Vsign Indulgéntiam ipsíus fusis lácrimis postulémus, et humiliémus illi ánimas nostras.\end{Resp}
\begin{Resp}\Rsign Qui non déspicit nos, nec ámovet salútem suam a génere nostro.\end{Resp}
\switchcolumn
\EN
\begin{Resp}\Rsign We know no strange God before the Lord. In Him we trust. \Star{} He despiseth us not, neither putteth He away His salvation from our nation.\end{Resp}
\begin{Resp}\Vsign His mercy let us seek with tears, and humble our souls before Him.\end{Resp}
\begin{Resp}\Rsign He despiseth us not, neither putteth He away His salvation from our nation.\end{Resp}
\switchcolumn*

\LA
\LessonHead{Lectio II}
\switchcolumn
\EN
\LessonHead{Lesson II}
\switchcolumn*

\LA
\Cite{Esth 7:5-7}
\switchcolumn
\EN
\Cite{Esth 7:5-7}
\switchcolumn*

\LA
\noindent Respondénsque rex Assuérus, ait: Quis est iste et cujus poténtiæ, ut hæc áudeat fácere? Dixítque Esther: Hostis et inimícus noster péssimus iste est Aman. Quod ille áudiens íllico obstúpuit vultum regis ac regínæ ferre non sústinens. Rex autem irátus surréxit et de loco convívii intrávit in hortum arbóribus cónsitum. Aman quoque surréxit ut rogáret Esther regínam pro ánima sua; intelléxit enim a rege sibi parátum malum.\par
\switchcolumn
\EN
\noindent And king Assuerus answered and said: Who is this, and of what power, that he should do these things? And Esther said: It is this Aman that is our adversary and most wicked enemy. Aman hearing this was forthwith astonished, not being able to bear the countenance of the king and of the queen. But the king being angry rose up, and went from the place of the banquet into the garden set with trees. Aman also rose up to entreat Esther the queen for his life, for he understood that evil was prepared for him by the king.\par
\switchcolumn*

\LA
\begin{Resp}\Rsign Confórta me, Rex, Sanctórum principátum tenens: \Star{} Et da sermónem rectum et bene sonántem in os meum.\end{Resp}
\begin{Resp}\Vsign Dómine, Rex univérsæ potestátis, convérte consílium eórum super eos.\end{Resp}
\begin{Resp}\Rsign Et da sermónem rectum et bene sonántem in os meum.\end{Resp}
\switchcolumn
\EN
\begin{Resp}\Rsign Strengthen me, O King, Who reignest over the holy ones. \Star{} Put Thou in my mouth clear and well-sounding words.\end{Resp}
\begin{Resp}\Vsign O Lord, King of all forces, turn back their device upon themselves.\end{Resp}
\begin{Resp}\Rsign Put Thou in my mouth clear and well-sounding words.\end{Resp}
\switchcolumn*

\LA
\LessonHead{Lectio III}
\switchcolumn
\EN
\LessonHead{Lesson III}
\switchcolumn*

\LA
\Cite{Esth 7:8-10}
\switchcolumn
\EN
\Cite{Esth 7:8-10}
\switchcolumn*

\LA
\noindent Qui cum revérsus esset de horto nemóribus cónsito et intrásset convívii locum, réperit Aman super léctulum corruísse in quo jacébat Esther et ait: Etiam regínam vult opprímere, me præsénte, in domo mea. Necdum verbum de ore regis exíerat, et statim operuérunt fáciem ejus. Dixítque Harbóna unus de eunúchis qui stabant in ministério regis: En lignum, quod paráverat Mardochǽo qui locútus est pro rege, stat in domo Aman habens altitúdinis quinquagínta cúbitos. Cui dixit rex: Appéndite eum in eo. Suspénsus est ítaque Aman in patíbulo, quod paráverat Mardochǽo, et regis ira quiévit.\par
\switchcolumn
\EN
\noindent And when the king came back out of the garden set with trees, and entered into the place of the banquet, he found Aman was fallen upon the bed on which Esther lay, and he said: He will force the queen also in my presence, in my own house. The word was not yet gone out of the king's mouth, and immediately they covered his face. And Harbona, one of the eunuchs that stood waiting on the king, said: Behold the gibbet which he hath prepared for Mardochai, who spoke for the king, standeth in Aman's house, being fifty cubits high. And the king said to him: Hang him upon it. So Aman was hanged on the gibbet, which he had prepared for Mardochai: and the king's wrath ceased.\par
\switchcolumn*

\LA
\begin{Resp}\Rsign Spem in álium nunquam hábui, prætérquam in te, Deus Israël: \Star{} Qui irásceris, et propítius eris, et ómnia peccáta hóminum in tribulatióne dimíttis.\end{Resp}
\begin{Resp}\Vsign Dómine Deus, Creátor cæli et terræ, réspice ad humilitátem nostram.\end{Resp}
\begin{Resp}\Rsign Qui irásceris, et propítius eris, et ómnia peccáta hóminum in tribulatióne dimíttis.\end{Resp}
\begin{Resp}\Vsign Glória Patri, et Fílio, \Star{} et Spirítui Sancto.\end{Resp}
\begin{Resp}\Rsign Qui irásceris, et propítius eris, et ómnia peccáta hóminum in tribulatióne dimíttis.\end{Resp}
\switchcolumn
\EN
\begin{Resp}\Rsign I have had no hope in any other but in thee, O God of Israel, Who art angry, and wilt again show mercy, \Star{} Thou that forgivest all the sins of men when they are in affliction.\end{Resp}
\begin{Resp}\Vsign O Lord God, Maker of heaven and earth, look upon our low estate;\end{Resp}
\begin{Resp}\Rsign Thou that forgivest all the sins of men when they are in affliction.\end{Resp}
\begin{Resp}\Vsign Glory be to the Father, and to the Son, \Star{} and to the Holy Ghost.\end{Resp}
\begin{Resp}\Rsign Thou that forgivest all the sins of men when they are in affliction.\end{Resp}
\switchcolumn*

\end{paracol}

\SeasonTitle{Dominicæ et Feriæ mensis Octobris}{Sundays and Ferias of October}

\addcontentsline{toc}{section}{Dominicæ et Feriæ mensis Octobris \textit{— Sundays and Ferias of October}}

\Office{Dominica I Octobris}{First Sunday of October}{}
\addcontentsline{toc}{subsection}{Dominica I Octobris \textit{— First Sunday of October}}
\begin{paracol}{2}
\LA
\LessonHead{Lectio I}
\switchcolumn
\EN
\LessonHead{Lesson I}
\switchcolumn*

\LA
\LessonTitle{Incipit liber primus Machabæórum}
\switchcolumn
\EN
\LessonTitle{Lesson from the first book of Machabees}
\switchcolumn*

\LA
\Cite{1 Mac 1:1-7}
\switchcolumn
\EN
\Cite{1 Mac 1:1-7}
\switchcolumn*

\LA
\noindent Et factum est, postquam percússit Alexánder Philíppi Mácedo, qui primus regnávit in Grǽcia, egréssus de terra Cethim, Daríum regem Persárum et Medórum, constítuit prǽlia multa et obtínuit ómnium munitiónes et interfécit reges terræ, et pertránsiit usque ad fines terræ et accépit spólia multitúdinis géntium, et síluit terra in conspéctu ejus. Et congregávit virtútem et exércitum fortem nimis, et exaltátum est et elevátum cor ejus, et obtínuit regiónes géntium et tyránnos, et facti sunt illi in tribútum. Et post hæc décidit in lectum, et cognóvit quia morerétur, et vocávit púeros suos nóbiles, qui secum erant nutríti a juventúte, et divísit illis regnum suum, cum adhuc víveret.\par
\switchcolumn
\EN
\noindent Now it came to pass, after that Alexander the son of Philip the Macedonian, who first reigned in Greece, coming out of the land of Cethim, had overthrown Darius king of the Persians and Medes: He fought many battles, and took the strongholds of all, and slew the kings of the earth: And he went through even to the ends of the earth, and took the spoils of many nations: and the earth was quiet before him. And he gathered a power, and a very strong army: and his heart was exalted and lifted up. And he subdued countries of nations, and princes: and they became tributaries to him. And after these things, he fell down upon his bed, and knew that he should die. And he called his servants the nobles that were brought up with him from his youth: and he divided his kingdom among them, while he was yet alive.\par
\switchcolumn*

\LA
\begin{Resp}\Rsign Adapériat Dóminus cor vestrum in lege sua et in præcéptis suis et fáciat pacem in diébus vestris: \Star{} Concédat vobis salútem, et rédimat vos a malis.\end{Resp}
\begin{Resp}\Vsign Exáudiat Dóminus oratiónes vestras, et reconciliétur vobis, nec vos déserat in témpore malo.\end{Resp}
\begin{Resp}\Rsign Concédat vobis salútem, et rédimat vos a malis.\end{Resp}
\switchcolumn
\EN
\begin{Resp}\Rsign The Lord open your hearts in His law and commandments, and send peace in your days. \Star{} May He grant you salvation and redeem you out of all evil.\end{Resp}
\begin{Resp}\Vsign The Lord hear your prayers, and be at one with you, and never forsake you in the time of trouble.\end{Resp}
\begin{Resp}\Rsign May He grant you salvation and redeem you out of all evil.\end{Resp}
\switchcolumn*

\LA
\LessonHead{Lectio II}
\switchcolumn
\EN
\LessonHead{Lesson II}
\switchcolumn*

\LA
\Cite{1 Mac 1:8-11}
\switchcolumn
\EN
\Cite{1 Mac 1:8-11}
\switchcolumn*

\LA
\noindent Et regnávit Alexánder annis duódecim et mórtuus est. Et obtinuérunt púeri ejus regnum, unusquísque in loco suo, et imposuérunt omnes sibi diadémata post mortem ejus, et fílii eórum post eos annis multis. Et multiplicáta sunt mala in terra. Et éxiit ex eis radix peccátrix, Antíochus illústris, fílius Antíochi regis qui fúerat Romæ obses, et regnávit in anno centésimo trigésimo séptimo regni Græcórum.\par
\switchcolumn
\EN
\noindent And Alexander reigned twelve years, and he died. And his servants made themselves kings every one in his place: And they all put crowns upon themselves after his death, and their sons after them many years, and evils were multiplied in the earth. And there came out of them a wicked root, Antiochus the Illustrious, the son of king Antiochus, who had been a hostage at Rome: and he reigned in the hundred and thirty-seventh year of the kingdom of the Greeks.\par
\switchcolumn*

\LA
\begin{Resp}\Rsign Exáudiat Dóminus oratiónes vestras, et reconciliétur vobis, nec vos déserat in témpore malo, \Star{} Dóminus, Deus noster.\end{Resp}
\begin{Resp}\Vsign Det vobis cor ómnibus, ut colátis eum et faciátis ejus voluntátem.\end{Resp}
\begin{Resp}\Rsign Dóminus, Deus noster.\end{Resp}
\switchcolumn
\EN
\begin{Resp}\Rsign The Lord hear your prayers, and be at one with you, and never forsake you in the time of trouble \Star{} Even He, the Lord our God.\end{Resp}
\begin{Resp}\Vsign Give you all an heart to serve Him, and to do His will.\end{Resp}
\begin{Resp}\Rsign Even He, the Lord our God.\end{Resp}
\switchcolumn*

\LA
\LessonHead{Lectio III}
\switchcolumn
\EN
\LessonHead{Lesson III}
\switchcolumn*

\LA
\Cite{1 Mac 1:12-16}
\switchcolumn
\EN
\Cite{1 Mac 1:12-16}
\switchcolumn*

\LA
\noindent In diébus illis exiérunt ex Israël fílii iníqui et suasérunt multis dicéntes: Eámus et disponámus testaméntum cum géntibus, quæ circa nos sunt, quia, ex quo recéssimus ab eis, invenérunt nos multa mala. Et bonus visus est sermo in óculis eórum. Et destinavérunt áliqui de pópulo, et abiérunt ad regem, et dedit illis potestátem ut fácerent justítiam géntium; et ædificavérunt gymnásium in Jerosólymis secúndum leges natiónum, et fecérunt sibi præpútia et recessérunt a testaménto sancto, et juncti sunt natiónibus, et venúmdati sunt ut fácerent malum.\par
\switchcolumn
\EN
\noindent In those days there went out of Israel wicked men, and they persuaded many, saying: Let us go, and make a covenant with the heathens that are round about us: for since we departed from them, many evils have befallen us. And the word seemed good in their eyes. And some of the people determined to do this, and went to the king: and he gave them license to do after the ordinances of the heathens. And they built a place of exercise in Jerusalem, according to the laws of the nations: And they made themselves prepuces, and departed from the holy covenant, and joined themselves to the heathens, and were sold to do evil.\par
\switchcolumn*

\LA
\begin{Resp}\Rsign Congregáti sunt inimíci nostri, et gloriántur in virtúte sua: cóntere fortitúdinem illórum, Dómine, et dispérge illos: \Star{} Ut cognóscant quia non est álius qui pugnet pro nobis, nisi tu, Deus noster.\end{Resp}
\begin{Resp}\Vsign Dispérge illos in virtúte tua, et déstrue eos, protéctor noster, Dómine.\end{Resp}
\begin{Resp}\Rsign Ut cognóscant quia non est álius qui pugnet pro nobis, nisi tu, Deus noster.\end{Resp}
\begin{Resp}\Vsign Glória Patri, et Fílio, \Star{} et Spirítui Sancto.\end{Resp}
\begin{Resp}\Rsign Ut cognóscant quia non est álius qui pugnet pro nobis, nisi tu, Deus noster.\end{Resp}
\switchcolumn
\EN
\begin{Resp}\Rsign Our enemies are gathered together, and make their boast of their own strength. O Lord, break their power, and scatter them \Star{} That they may know that there is none other that fighteth for us, but only Thou, O our God!\end{Resp}
\begin{Resp}\Vsign Scatter them in thy strength, and destroy them, O Lord our Shield.\end{Resp}
\begin{Resp}\Rsign That they may know that there is none other that fighteth for us, but only Thou, O our God!\end{Resp}
\begin{Resp}\Vsign Glory be to the Father, and to the Son, \Star{} and to the Holy Ghost.\end{Resp}
\begin{Resp}\Rsign That they may know that there is none other that fighteth for us, but only Thou, O our God!\end{Resp}
\switchcolumn*

\LA
\LessonHead{Lectio IV}
\switchcolumn
\EN
\LessonHead{Lesson IV}
\switchcolumn*

\LA
\LessonTitle{Ex libro Officiórum sancti Ambrósii Epíscopi}
\switchcolumn
\EN
\LessonTitle{From the Book upon Duties written by St. Ambrose, Bishop (of Milan).}
\switchcolumn*

\LA
\Source{Lib. 1, Cap. 40}
\switchcolumn
\EN
\Source{Bk. i. ch. 40}
\switchcolumn*

\LA
\noindent Fortásse áliquos béllica defíxos glória tenet, ut putent solam esse præliárem fortitúdinem; et ídeo me ad ista deflexísse, quia illa nostris déforet. Quam fortis Jesus Nave, ut uno prǽlio quinque reges captos stérneret cum pópulis suis! Deínde cum advérsum Gabaonítas urgéret prǽlium, et vererétur ne nox impedíret victóriam, magnitúdine mentis et fídei clamávit: Stet sol, et stetit, donec victória consummarétur. Gédeon in trecéntis viris de ingénti pópulo et acérbo hoste revéxit triúmphum. Jónathas adoléscens virtútem magnam fecit in prǽlio.\par
\switchcolumn
\EN
\noindent There may perchance be some who are so blinded by the glory of war as to think there is no valour but warlike valour, and that the reason why I have taken up other subjects is that among us there is no warlike valour whereof to speak. But what was the valour of Josue the son of Nun, when in one battle he laid low five nations, and took prisoners their kings when he was fighting against the Gibeonites, and feared lest the closing in of night should cut short his victory, he cried aloud in the greatness of his mind and of his faith? And he said, in the sight of Israel: Sun stand thou still over against Gibeon, and thou Moon over against the valley of Ajalon and the sun stood still, and the moon stayed, until the people had avenged themselves upon their enemies (Josue x. 12, 13). Gideon, with three hundred men, won the victory over the vast people, and the savage enemy. The lad Jonathan waxed valiant in fight.\par
\switchcolumn*

\LA
\begin{Resp}\Rsign Impetum inimicórum ne timuéritis: mémores estóte, quómodo salvi facti sunt patres nostri: \Star{} Et nunc clamémus in cælum et miserébitur nostri Deus noster.\end{Resp}
\begin{Resp}\Vsign Mementóte mirabílium ejus, quæ fecit pharaóni et exercítui ejus in Mari Rubro.\end{Resp}
\begin{Resp}\Rsign Et nunc clamémus in cælum et miserébitur nostri Deus noster.\end{Resp}
\switchcolumn
\EN
\begin{Resp}\Rsign Be ye not afraid of the assault of the enemy; remember how our fathers were delivered. \Star{} Now, therefore, let us cry unto heaven, and our God will have mercy upon us.\end{Resp}
\begin{Resp}\Vsign Remember His marvellous works that He hath done unto Pharaoh and his host in the Red Sea.\end{Resp}
\begin{Resp}\Rsign Now, therefore, let us cry unto heaven, and our God will have mercy upon us.\end{Resp}
\switchcolumn*

\LA
\LessonHead{Lectio V}
\switchcolumn
\EN
\LessonHead{Lesson V}
\switchcolumn*

\LA
\noindent Quid de Machabǽis loquar? Sed prius de pópulo dicam patrum; qui cum essent paráti ad repugnándum pro templo Dei et pro legítimis suis, dolo hóstium die lacessíti sábbati, maluérunt vulnéribus offérre nuda córpora, quam repugnáre, ne violárent sábbatum. Itaque omnes læti se obtulérunt morti. Sed Machabǽi considerántes quod hoc exémplo gens omnis posset períre, sábbato étiam, cum ipsi in bellum provocaréntur, ulti sunt innocéntium necem fratrum suórum. Unde póstea stimulátus rex Antíochus, cum bellum accénderet per duces suos Lýsiam, Nicánorem, Górgiam, ita cum Orientálibus suis et Assýriis attrítus est cópiis, ut quadragínta et octo míllia in médio campi a tribus míllibus prosterneréntur.\par
\switchcolumn
\EN
\noindent Shall I speak of the Maccabees But before I speak of them, I will speak of their fathers, even of them who, when they were ready to fight for the Temple of God and for their own rights, were assailed by a trick of their enemies upon the Sabbath day, and were willing rather to offer their bodies naked to the sword than to strike back again and break the Sabbath, and so they gave themselves up gladly to death but when the Maccabees bethought them that the whole nation might thus perish, avenged the innocent blood of their brethren even upon the Sabbath day when they were provoked to battle, and afterward, when King Antiochus had been stirred up to make war on them by his generals, even Lysias and Nicanor and Gorgias, he and his Eastern and Assyrian forces were so crushed that forty and eight thousand were laid low on the field by three thousand.\par
\switchcolumn*

\LA
\begin{Resp}\Rsign Congregátæ sunt gentes in multitúdine, ut dímicent contra nos, et ignorámus quid ágere debeámus: \Star{} Dómine Deus, ad te sunt óculi nostri, ne pereámus.\end{Resp}
\begin{Resp}\Vsign Tu scis quæ cógitant in nos: quómodo potérimus subsístere ante fáciem illórum, nisi tu ádjuves nos?\end{Resp}
\begin{Resp}\Rsign Dómine Deus, ad te sunt óculi nostri, ne pereámus.\end{Resp}
\switchcolumn
\EN
\begin{Resp}\Rsign The heathen are assembled together to fight against us, and we know not what we should do. \Star{} Our eyes look unto thee, O Lord our God, that we should not perish.\end{Resp}
\begin{Resp}\Vsign What things they imagine against us, Thou knowest. How shall we be able to stand against them, except Thou be our help.\end{Resp}
\begin{Resp}\Rsign Our eyes look unto thee, O Lord our God, that we should not perish.\end{Resp}
\switchcolumn*

\LA
\LessonHead{Lectio VI}
\switchcolumn
\EN
\LessonHead{Lesson VI}
\switchcolumn*

\LA
\noindent Virtútem ducis Judæ Machabǽi de uno ejus mílite consideráte. Namque Eleázarus, cum supereminéntem céteris elephántem loríca vestítum régia advérteret, arbitrátus quod in eo esset rex, cursu cóncito in médium legiónis se prorúpit: et, abjécto clýpeo, utráque manu interficiébat, donec perveníret ad béstiam, atque intrávit sub eam, et subjécto gládio interémit eam. Itaque cadens béstia oppréssit Eleázarum, atque ita mórtuus est. Quanta ígitur virtus ánimi! primo, ut mortem non timéret; deínde, ut circumfúsus legiónibus inimicórum, in confértos raperétur hostes, médium penetráret agmen, et contémpta morte ferócior, abjécto clýpeo, utráque manu vulnerátæ molem béstiæ subíret ac sustinéret: post infra ipsam succéderet, quo plenióri feríret ictu; cujus ruína inclúsus magis quam oppréssus, suo est sepúltus triúmpho.\par
\switchcolumn
\EN
\noindent What was the valour of Judas the Maccabean leader we may judge by the type of one of his men. When Eleazar saw an elephant bigger than the rest, and adorned with the King's harness, he thought that the King was riding thereon, and he threw himself into the midst of the enemy, and cast away his shield and slew on either hand until he was come to the beast, and ran underneath it, and killed it with his sword, and so the beast fell upon Eleazar and crushed him, and he died. What valour was here To begin with, he feared not to die, and when the enemy surrounded him he cast himself into the midst of their ranks, pierced their column, and becoming all the fiercer through his mockery of death, he threw away his shield and upheld with both hands the huge bulk of the wounded monster beneath which he had gone the better to spite it, so that when he died with it he might well have been said not so much to be crushed as to be swallowed up in victory.\par
\switchcolumn*

\LA
\begin{Resp}\Rsign Tua est poténtia, tuum regnum, Dómine: tu es super omnes gentes: \Star{} Da pacem, Dómine, in diébus nostris.\end{Resp}
\begin{Resp}\Vsign Creátor ómnium, Deus, terríbilis et fortis, justus et miséricors.\end{Resp}
\begin{Resp}\Rsign Da pacem, Dómine, in diébus nostris.\end{Resp}
\begin{Resp}\Vsign Glória Patri, et Fílio, \Star{} et Spirítui Sancto.\end{Resp}
\begin{Resp}\Rsign Da pacem, Dómine, in diébus nostris.\end{Resp}
\switchcolumn
\EN
\begin{Resp}\Rsign Thine, O Lord, is the power, thine is the kingdom, O Lord, and Thou art exalted above all the heathen. \Star{} Give peace in our time, O Lord.\end{Resp}
\begin{Resp}\Vsign O Lord God, Creator of all things, Who art fearful and strong, righteous and merciful.\end{Resp}
\begin{Resp}\Rsign Give peace in our time, O Lord.\end{Resp}
\begin{Resp}\Vsign Glory be to the Father, and to the Son, \Star{} and to the Holy Ghost.\end{Resp}
\begin{Resp}\Rsign Give peace in our time, O Lord.\end{Resp}
\switchcolumn*

\end{paracol}

\Office{Feria II infra Hebdomadam I Octobris}{Monday of the first week of October}{}
\addcontentsline{toc}{subsection}{Feria II infra Hebdomadam I Octobris \textit{— Monday of the first week of October}}
\begin{paracol}{2}
\LA
\LessonHead{Lectio I}
\switchcolumn
\EN
\LessonHead{Lesson I}
\switchcolumn*

\LA
\LessonTitle{De libro primo Machabæórum}
\switchcolumn
\EN
\LessonTitle{Lesson from the first book of Machabees}
\switchcolumn*

\LA
\Cite{1 Mac 1:17-20}
\switchcolumn
\EN
\Cite{1 Mac 1:17-20}
\switchcolumn*

\LA
\noindent Et parátum est regnum in conspéctu Antíochi, et cœpit regnáre in terra Ægýpti, ut regnáret super duo regna, et intrávit in Ægýptum in multitúdine gravi, in cúrribus et elephántis et equítibus et copiósa návium multitúdine, et constítuit bellum advérsus Ptolemǽum regem Ægýpti, et véritus est Ptolemǽus a fácie ejus et fugit, et cecidérunt vulneráti multi, et comprehéndit civitátes munítas in terra Ægýpti et accépit spólia terræ Ægýpti.\par
\switchcolumn
\EN
\noindent And the kingdom was established before Antiochus, and he had a mind to reign over the land of Egypt, that he might reign over two kingdoms. And he entered into Egypt with a great multitude, with chariots and elephants, and horsemen, and a great number of ships: And he made war against Ptolemee king of Egypt, but Ptolemee was afraid at his presence, and fled, and many were wounded unto death. And he took the strong cities in the land of Egypt: and he took the spoils of the land of Egypt.\par
\switchcolumn*

\LA
\begin{Resp}\Rsign Dixit Judas Simóni fratri suo: Elige tibi viros et vade, líbera fratres tuos in Galilǽam; ego autem et Jónathas frater tuus íbimus in Galaadítim: \Star{} Sicut fúerit volúntas in cælo, sic fiat.\end{Resp}
\begin{Resp}\Vsign Accingímini, fílii poténtes, et estóte paráti: quóniam mélius est nobis, mori in bello, quam vidére mala gentis nostræ et sanctórum.\end{Resp}
\begin{Resp}\Rsign Sicut fúerit volúntas in cælo, sic fiat.\end{Resp}
\switchcolumn
\EN
\begin{Resp}\Rsign Judas said unto Simon his brother: Choose thee out men, and go, and deliver thy brethren that are in Galilee and I, and Jonathan thy brother, will go into the country of Galaad. \Star{} As the Will is in heaven, so let it be.\end{Resp}
\begin{Resp}\Vsign Arm yourselves, ye valiant men, and be in readiness for it is better for us to die in battle, than to behold the calamities of our people, and our sanctuary.\end{Resp}
\begin{Resp}\Rsign As the Will is in heaven, so let it be.\end{Resp}
\switchcolumn*

\LA
\LessonHead{Lectio II}
\switchcolumn
\EN
\LessonHead{Lesson II}
\switchcolumn*

\LA
\Cite{1 Mac 1:21-23}
\switchcolumn
\EN
\Cite{1 Mac 1:21-23}
\switchcolumn*

\LA
\noindent Et convértit Antíochus, postquam percússit Ægýptum in centésimo et quadragésimo tértio anno, et ascéndit ad Israël. Et ascéndit Jerosólymam in multitúdine gravi, et intrávit in sanctificatiónem cum supérbia, et accépit altáre áureum, et candelábrum lúminis, et univérsa vasa ejus, et mensam propositiónis, et libatória, et phíalas, et mortaríola áurea, et velum, et corónas, et ornaméntum áureum, quod in fácie templi erat: et commínuit ómnia.\par
\switchcolumn
\EN
\noindent And after Antiochus had ravaged Egypt in the hundred and forty-third year, he returned and went up against Israel. And he went up to Jerusalem with a great multitude. And he proudly entered into the sanctuary, and took away the golden altar, and the candlestick of light, and all the vessels thereof, and the table of proposition, and the pouring vessels, and the vials, and the little mortars of gold, and the veil, and the crowns, and the golden ornament that was before the temple: and he broke them all in pieces.\par
\switchcolumn*

\LA
\begin{Resp}\Rsign Ornavérunt fáciem templi corónis áureis, et dedicavérunt altáre Dómino: \Star{} Et facta est lætítia magna in pópulo.\end{Resp}
\begin{Resp}\Vsign In hymnis et confessiónibus benedicébant Dóminum.\end{Resp}
\begin{Resp}\Rsign Et facta est lætítia magna in pópulo.\end{Resp}
\begin{Resp}\Vsign Glória Patri, et Fílio, \Star{} et Spirítui Sancto.\end{Resp}
\begin{Resp}\Rsign Et facta est lætítia magna in pópulo.\end{Resp}
\switchcolumn
\EN
\begin{Resp}\Rsign They decked the fore-front of the Temple with crowns of gold, and dedicated the Altar unto the Lord. \Star{} And there was very great gladness among the people.\end{Resp}
\begin{Resp}\Vsign They praised the Lord with Psalms and thanksgiving.\end{Resp}
\begin{Resp}\Rsign And there was very great gladness among the people.\end{Resp}
\begin{Resp}\Vsign Glory be to the Father, and to the Son, \Star{} and to the Holy Ghost.\end{Resp}
\begin{Resp}\Rsign And there was very great gladness among the people.\end{Resp}
\switchcolumn*

\LA
\LessonHead{Lectio III}
\switchcolumn
\EN
\LessonHead{Lesson III}
\switchcolumn*

\LA
\Cite{1 Mac 1:24-29}
\switchcolumn
\EN
\Cite{1 Mac 1:24-29}
\switchcolumn*

\LA
\noindent Et accépit argéntum, et aurum, et vasa concupiscibília: et accépit thesáuros occúltos, quos invénit: et, sublátis ómnibus, ábiit in terram suam; et fecit cædem hóminum, et locútus est in supérbia magna, et factus est planctus magnus in Israël, et in omni loco eórum: et ingemuérunt príncipes et senióres, vírgines et júvenes infirmáti sunt, et speciósitas mulíerum immutáta est. Omnis marítus sumpsit laméntum, et quæ sedébant in thoro maritáli, lugébant: et commóta est terra super habitántes in ea, et univérsa domus Jacob índuit confusiónem.\par
\switchcolumn
\EN
\noindent And he took the silver and gold, and the precious vessels: and he took the hidden treasures which he found: and when he had taken all away he departed into his own country. And he made a great slaughter of men, and spoke very proudly. And there was great mourning in Israel, and in every place where they were. And the princes, and the ancients mourned, and the virgins and the young men were made feeble, and the beauty of the women was changed. Every bridegroom took up lamentation: and the bride that set in the marriage bed, mourned: And the land was moved for the inhabitants thereof, and all the house of Jacob was covered with confusion.\par
\switchcolumn*

\LA
\begin{Resp}\Rsign In hymnis et confessiónibus benedicébant Dóminum, \Star{} Qui magna fecit in Israël, et victóriam dedit illis Dóminus omnípotens.\end{Resp}
\begin{Resp}\Vsign Ornavérunt fáciem templi corónis áureis, et dedicavérunt altáre Dómino.\end{Resp}
\begin{Resp}\Rsign Qui magna fecit in Israël, et victóriam dedit illis Dóminus omnípotens.\end{Resp}
\begin{Resp}\Vsign Glória Patri, et Fílio, \Star{} et Spirítui Sancto.\end{Resp}
\begin{Resp}\Rsign Qui magna fecit in Israël, et victóriam dedit illis Dóminus omnípotens.\end{Resp}
\switchcolumn
\EN
\begin{Resp}\Rsign They praised the Lord with psalms and thanksgiving \Star{} Who had done so great things in Israel, and given them the victory the Lord Almighty.\end{Resp}
\begin{Resp}\Vsign They decked the fore-front of the Temple with crowns of gold, and dedicated the Altar unto the Lord.\end{Resp}
\begin{Resp}\Rsign Who had done so great things for Israel, and given them the victory the Lord Almighty.\end{Resp}
\begin{Resp}\Vsign Glory be to the Father, and to the Son, \Star{} and to the Holy Ghost.\end{Resp}
\begin{Resp}\Rsign Who had done so great things for Israel, and given them the victory the Lord Almighty.\end{Resp}
\switchcolumn*

\end{paracol}

\Office{Feria III infra Hebdomadam I Octobris}{Tuesday of the first week of October}{}
\addcontentsline{toc}{subsection}{Feria III infra Hebdomadam I Octobris \textit{— Tuesday of the first week of October}}
\begin{paracol}{2}
\LA
\LessonHead{Lectio I}
\switchcolumn
\EN
\LessonHead{Lesson I}
\switchcolumn*

\LA
\LessonTitle{De libro primo Machabæórum}
\switchcolumn
\EN
\LessonTitle{Lesson from the first book of Machabees}
\switchcolumn*

\LA
\Cite{1 Mac 2:1-6}
\switchcolumn
\EN
\Cite{1 Mac 2:1-6}
\switchcolumn*

\LA
\noindent In diébus illis surréxit Mathathías fílius Joánnis, fílii Simeónis, sacérdos ex fíliis Jóarib, ab Jerúsalem, et consédit in monte Modin. Et habébat fílios quinque, Joánnem, qui cognominabátur Gaddis, et Simónem, qui cognominabátur Thasi, et Judam, qui vocabátur Machabǽus, et Eleázarum, qui cognominabátur Abaron, et Jónathan, qui cognominabátur Apphus. Hi vidérunt mala quæ fiébant in pópulo Juda, et in Jerúsalem.\par
\switchcolumn
\EN
\noindent In those days arose Mathathias the son of John, the son of Simeon, a priest of the sons of Joarib, from Jerusalem, and he abode in the mountain of Modin. And he had five sons: John who was surnamed Gaddis: And Simon, who was surnamed Thasi: And Judas, who was called Machabeus: And Eleazar, who was surnamed Abaron: and Jonathan, who was surnamed Apphus. These saw the evils that were done in the people of Juda, and in Jerusalem.\par
\switchcolumn*

\LA
\begin{Resp}\Rsign Hic est fratrum amátor et pópuli Israël: \Star{} Hic est, qui multum orat pro pópulo et univérsa sancta civitáte Jerúsalem.\end{Resp}
\begin{Resp}\Vsign Vir iste in pópulo suo mitíssimus appáruit.\end{Resp}
\begin{Resp}\Rsign Hic est, qui multum orat pro pópulo et univérsa sancta civitáte Jerúsalem.\end{Resp}
\switchcolumn
\EN
\begin{Resp}\Rsign This is a lover of the brethren, and of the people of Israel \Star{} This is one who prayeth much for the people, and for all the Holy City, Jerusalem.\end{Resp}
\begin{Resp}\Vsign There appeared a man most gentle toward all his people.\end{Resp}
\begin{Resp}\Rsign This is one who prayeth much for the people, and for all the Holy City, Jerusalem.\end{Resp}
\switchcolumn*

\LA
\LessonHead{Lectio II}
\switchcolumn
\EN
\LessonHead{Lesson II}
\switchcolumn*

\LA
\Cite{1 Mac 2:7-10}
\switchcolumn
\EN
\Cite{1 Mac 2:7-10}
\switchcolumn*

\LA
\noindent Et dixit Mathathías: Væ mihi! ut quid natus sum vidére contritiónem pópuli mei et contritiónem civitátis sanctæ et sedére illic, cum datur in mánibus inimicórum? Sancta in manu extraneórum facta sunt, templum ejus sicut homo ignóbilis, vasa glóriæ ejus captíva abdúcta sunt, trucidáti sunt senes ejus in platéis, et júvenes ejus cecidérunt in gládio inimicórum. Quæ gens non hereditávit regnum ejus et non obtínuit spólia ejus?\par
\switchcolumn
\EN
\noindent And Mathathias said: Woe is me, wherefore was I born to see the ruin of my people, and the ruin of the holy city, and to dwell there, when it is given into the hands of the enemies? The holy places are come into the hands of strangers: her temple is become as a man without honour. The vessels of her glory are carried away captive: her old men are murdered in the streets, and her young men are fallen by the sword of the enemies. What nation hath not inherited her kingdom, and gotten of her spoils?\par
\switchcolumn*

\LA
\begin{Resp}\Rsign Tu, Dómine universórum, qui nullam habes indigéntiam, voluísti templum tuum fíeri in nobis; \Star{} Consérva domum istam immaculátam in ætérnum, Dómine.\end{Resp}
\begin{Resp}\Vsign Tu elegísti, Dómine, domum istam ad invocándum nomen tuum in ea, ut esset domus oratiónis et obsecratiónis pópulo tuo.\end{Resp}
\begin{Resp}\Rsign Consérva domum istam immaculátam in ætérnum, Dómine.\end{Resp}
\switchcolumn
\EN
\begin{Resp}\Rsign Thou, O Lord of all things, Who hast need of nothing, wast pleased that the Temple of thine habitation should be among us. \Star{} Therefore now, O Lord, keep this house ever undefiled\end{Resp}
\begin{Resp}\Vsign Thou, O Lord, didst choose this house, that thy Name should be called on therein, and to be an house of prayer and petition for thy people.\end{Resp}
\begin{Resp}\Rsign Therefore now, O Lord, keep this house ever undefiled.\end{Resp}
\switchcolumn*

\LA
\LessonHead{Lectio III}
\switchcolumn
\EN
\LessonHead{Lesson III}
\switchcolumn*

\LA
\Cite{1 Mac 2:14-16}
\switchcolumn
\EN
\Cite{1 Mac 2:14-16}
\switchcolumn*

\LA
\noindent Et scidit vestiménta sua Mathathías, et fílii ejus: et operuérunt se cilíciis, et planxérunt valde. Et venérunt illuc qui missi erant a rege Antíocho, ut cógerent eos, qui confúgerant in civitátem Modin, immoláre, et accéndere thura, et a lege Dei discédere, et multi de pópulo Israël consentiéntes accessérunt ad eos: sed Mathathías et fílii ejus constánter stetérunt.\par
\switchcolumn
\EN
\noindent And Mathathias and his sons rent their garments, and they covered themselves with haircloth, and made great lamentation. And they that were sent from king Antiochus came thither, to compel them that were fled into the city of Modin, to sacrifice, and to burn incense, and to depart from the law of God. And many of the people of Israel consented, and came to them: but Mathathias and his sons stood firm.\par
\switchcolumn*

\LA
\begin{Resp}\Rsign Aperi óculos tuos, Dómine, et vide afflictiónem nostram: circumdedérunt nos gentes ad puniéndum nos: \Star{} Sed tu, Dómine, exténde brácchium tuum, et líbera ánimas nostras.\end{Resp}
\begin{Resp}\Vsign Afflíge oppriméntes nos et contuméliam faciéntes in supérbiam; et custódi partem tuam.\end{Resp}
\begin{Resp}\Rsign Sed tu, Dómine, exténde brácchium tuum, et líbera ánimas nostras.\end{Resp}
\begin{Resp}\Vsign Glória Patri, et Fílio, \Star{} et Spirítui Sancto.\end{Resp}
\begin{Resp}\Rsign Sed tu, Dómine, exténde brácchium tuum, et líbera ánimas nostras.\end{Resp}
\switchcolumn
\EN
\begin{Resp}\Rsign Open thine eyes, O Lord, and behold our affliction for the heathen are come round about us to punish us. \Star{} But Thou, O Lord, stretch forth thine arm, and deliver our souls.\end{Resp}
\begin{Resp}\Vsign Punish them that oppress us and with pride do us wrong, and keep thine own portion.\end{Resp}
\begin{Resp}\Rsign But Thou, O Lord, stretch forth thine arm, and deliver our souls.\end{Resp}
\begin{Resp}\Vsign Glory be to the Father, and to the Son, \Star{} and to the Holy Ghost.\end{Resp}
\begin{Resp}\Rsign But Thou, O Lord, stretch forth thine arm, and deliver our souls.\end{Resp}
\switchcolumn*

\end{paracol}

\Office{Feria IV infra Hebdomadam I Octobris}{Wednesday of the first week of October}{}
\addcontentsline{toc}{subsection}{Feria IV infra Hebdomadam I Octobris \textit{— Wednesday of the first week of October}}
\begin{paracol}{2}
\LA
\LessonHead{Lectio I}
\switchcolumn
\EN
\LessonHead{Lesson I}
\switchcolumn*

\LA
\LessonTitle{De libro primo Machabæórum}
\switchcolumn
\EN
\LessonTitle{Lesson from the first book of Machabees}
\switchcolumn*

\LA
\Cite{1 Mac 2:19-22}
\switchcolumn
\EN
\Cite{1 Mac 2:19-22}
\switchcolumn*

\LA
\noindent Et respóndit Mathathías et dixit magna voce: Etsi omnes gentes regi Antíocho obédiunt, ut discédat unusquísque a servitúte legis patrum suórum, et conséntiat mandátis ejus: ego et fílii mei, et fratres mei, obediémus legi patrum nostrórum. Propítius sit nobis Deus, non est nobis útile relínquere legem et justítias Dei: non audiémus verba regis Antíochi, nec sacrificábimus transgrediéntes legis nostræ mandáta, ut eámus áltera via.\par
\switchcolumn
\EN
\noindent Then Mathathias answered, and said with a loud voice: Although all nations obey king Antiochus, so as to depart every man from the service of the law of his fathers, and consent to his commandments: I and my sons, and my brethren will obey the law of our fathers. God be merciful unto us: it is not profitable for us to forsake the law, and the justices of God: We will not hearken to the words of king Antiochus, neither will we sacrifice, and transgress the commandments of our law, to go another way.\par
\switchcolumn*

\LA
\begin{Resp}\Rsign Refúlsit sol in clípeos áureos, et resplenduérunt montes ab eis: \Star{} Et fortitúdo géntium dissipáta est.\end{Resp}
\begin{Resp}\Vsign Erat enim exércitus magnus valde et fortis: et appropiávit Judas et exércitus ejus in prǽlio.\end{Resp}
\begin{Resp}\Rsign Et fortitúdo géntium dissipáta est.\end{Resp}
\switchcolumn
\EN
\begin{Resp}\Rsign The sun shone upon the shields of gold, and the mountains glistered therewith; \Star{} And the army of the heathens was spread abroad.\end{Resp}
\begin{Resp}\Vsign For the army was very great and mighty then Judas and his host drew near and entered into battle.\end{Resp}
\begin{Resp}\Rsign And the army of the heathens was spread abroad.\end{Resp}
\switchcolumn*

\LA
\LessonHead{Lectio II}
\switchcolumn
\EN
\LessonHead{Lesson II}
\switchcolumn*

\LA
\Cite{1 Mac 2:23-26}
\switchcolumn
\EN
\Cite{1 Mac 2:23-26}
\switchcolumn*

\LA
\noindent Et, ut cessávit loqui verba hæc, accéssit quidam Judǽus in ómnium óculis sacrificáre idólis super aram in civitáte Modin, secúndum jussum regis. Et vidit Mathathías et dóluit, et contremuérunt renes ejus, et accénsus est furor ejus secúndum judícium legis, et insíliens trucidávit eum super aram. Sed et virum, quem rex Antíochus míserat, qui cogébat immoláre, occídit in ipso témpore, et aram destrúxit: et zelátus est legem, sicut fecit Phínees Zamri fílio Salómi.\par
\switchcolumn
\EN
\noindent Now as he left off speaking these words, there came a certain Jew in the sight of all to sacrifice to the idols upon the altar in the city of Modin, according to the king's commandment. And Mathathias saw and was grieved, and his reins trembled, and his wrath was kindled according to the judgment of the law, and running upon him he slew him upon the altar: Moreover the man whom king Antiochus had sent, who compelled them to sacrifice, he slew at the same time, and pulled down the altar. And shewed zeal for the law, as Phinees did by Zamri the son of Salomi.\par
\switchcolumn*

\LA
\begin{Resp}\Rsign Ornavérunt fáciem templi corónis áureis, et dedicavérunt altáre Dómino: \Star{} Et facta est lætítia magna in pópulo.\end{Resp}
\begin{Resp}\Vsign In hymnis et confessiónibus benedicébant Dóminum.\end{Resp}
\begin{Resp}\Rsign Et facta est lætítia magna in pópulo.\end{Resp}
\switchcolumn
\EN
\begin{Resp}\Rsign They decked the fore-front of the Temple with crowns of gold, and dedicated the Altar unto the Lord. \Star{} And there was very great gladness among the people.\end{Resp}
\begin{Resp}\Vsign They praised the Lord with Psalms and thanksgiving.\end{Resp}
\begin{Resp}\Rsign And there was very great gladness among the people.\end{Resp}
\switchcolumn*

\LA
\LessonHead{Lectio III}
\switchcolumn
\EN
\LessonHead{Lesson III}
\switchcolumn*

\LA
\Cite{1 Mac 2:27-30}
\switchcolumn
\EN
\Cite{1 Mac 2:27-30}
\switchcolumn*

\LA
\noindent Et exclamávit Mathathías voce magna in civitáte, dicens: Omnis qui zelum habet legis státuens testaméntum éxeat post me. Et fugit ipse, et fílii ejus in montes, et reliquérunt quæcúmque habébant in civitáte. Tunc descendérunt multi quæréntes judícium et justítiam in desértum. Et sedérunt ibi ipsi, et fílii eórum, et mulíeres eórum, et pécora eórum, quóniam inundavérunt super eos mala.\par
\switchcolumn
\EN
\noindent And Mathathias cried out in the city with a loud voice, saying: Every one that hath zeal for the law, and maintaineth the testament, let him follow me. So he, and his sons fled into the mountains, and left all that they had in the city. Then many that sought after judgment, and justice, went down into the desert: And they abode there, they and their children, and their wives, and their cattle: because afflictions increased upon them.\par
\switchcolumn*

\LA
\begin{Resp}\Rsign In hymnis et confessiónibus benedicébant Dóminum, \Star{} Qui magna fecit in Israël, et victóriam dedit illis Dóminus omnípotens.\end{Resp}
\begin{Resp}\Vsign Ornavérunt fáciem templi corónis áureis, et dedicavérunt altáre Dómino.\end{Resp}
\begin{Resp}\Rsign Qui magna fecit in Israël, et victóriam dedit illis Dóminus omnípotens.\end{Resp}
\begin{Resp}\Vsign Glória Patri, et Fílio, \Star{} et Spirítui Sancto.\end{Resp}
\begin{Resp}\Rsign Qui magna fecit in Israël, et victóriam dedit illis Dóminus omnípotens.\end{Resp}
\switchcolumn
\EN
\begin{Resp}\Rsign They praised the Lord with psalms and thanksgiving \Star{} Who had done so great things in Israel, and given them the victory the Lord Almighty.\end{Resp}
\begin{Resp}\Vsign They decked the fore-front of the Temple with crowns of gold, and dedicated the Altar unto the Lord.\end{Resp}
\begin{Resp}\Rsign Who had done so great things for Israel, and given them the victory the Lord Almighty.\end{Resp}
\begin{Resp}\Vsign Glory be to the Father, and to the Son, \Star{} and to the Holy Ghost.\end{Resp}
\begin{Resp}\Rsign Who had done so great things for Israel, and given them the victory the Lord Almighty.\end{Resp}
\switchcolumn*

\end{paracol}

\Office{Feria V infra Hebdomadam I Octobris}{Thursday of the first week of October}{}
\addcontentsline{toc}{subsection}{Feria V infra Hebdomadam I Octobris \textit{— Thursday of the first week of October}}
\begin{paracol}{2}
\LA
\LessonHead{Lectio I}
\switchcolumn
\EN
\LessonHead{Lesson I}
\switchcolumn*

\LA
\LessonTitle{De libro primo Machabæórum}
\switchcolumn
\EN
\LessonTitle{Lesson from the first book of Machabees}
\switchcolumn*

\LA
\Cite{1 Mac 2:49-54}
\switchcolumn
\EN
\Cite{1 Mac 2:49-54}
\switchcolumn*

\LA
\noindent Et appropinquavérunt dies Mathathíæ moriéndi, et dixit fíliis suis: Nunc confortáta est supérbia, et castigátio, et tempus eversiónis, et ira indignatiónis. Nunc ergo, o fílii, æmulatóres estóte legis, et date ánimas vestras pro testaménto patrum vestrórum, et mementóte óperum patrum, quæ fecérunt in generatiónibus suis, et accipiétis glóriam magnam, et nomen ætérnum. Abraham, nonne in tentatióne invéntus est fidélis, et reputátum est ei ad justítiam? Joseph in témpore angústiæ suæ custodívit mandátum, et factus est dóminus Ægýpti; Phínees pater noster, zelándo zelum Dei, accépit testaméntum sacerdótii ætérni.\par
\switchcolumn
\EN
\noindent Now the days drew near that Mathathias should die, and he said to his sons: Now hath pride and chastisement gotten strength, and the time of destruction, and the wrath of indignation: Now therefore, O my sons, be ye zealous for the law, and give your lives for the covenant of your fathers. And call to remembrance the works of the fathers, which they have done in their generations: and you shall receive great glory, and an everlasting name. Was not Abraham found faithful in temptation, and it was reputed to him unto justice? Joseph in the time of his distress kept the commandment, and he was made lord of Egypt. Phinees our father, by being fervent in the zeal of God, received the covenant of an everlasting priesthood.\par
\switchcolumn*

\LA
\begin{Resp}\Rsign Adapériat Dóminus cor vestrum in lege sua et in præcéptis suis et fáciat pacem in diébus vestris: \Star{} Concédat vobis salútem, et rédimat vos a malis.\end{Resp}
\begin{Resp}\Vsign Exáudiat Dóminus oratiónes vestras, et reconciliétur vobis, nec vos déserat in témpore malo.\end{Resp}
\begin{Resp}\Rsign Concédat vobis salútem, et rédimat vos a malis.\end{Resp}
\switchcolumn
\EN
\begin{Resp}\Rsign The Lord open your hearts in His law and commandments, and send peace in your days. \Star{} May He grant you salvation and redeem you out of all evil.\end{Resp}
\begin{Resp}\Vsign The Lord hear your prayers, and be at one with you, and never forsake you in the time of trouble.\end{Resp}
\begin{Resp}\Rsign May He grant you salvation and redeem you out of all evil.\end{Resp}
\switchcolumn*

\LA
\LessonHead{Lectio II}
\switchcolumn
\EN
\LessonHead{Lesson II}
\switchcolumn*

\LA
\Cite{1 Mac 2:55-63}
\switchcolumn
\EN
\Cite{1 Mac 2:55-63}
\switchcolumn*

\LA
\noindent Jesus, dum implévit verbum, factus est dux in Israël; Caleb, dum testificátur in ecclésia, accépit hereditátem; David in sua misericórdia consecútus est sedem regni in sǽcula; Elías, dum zelat zelum legis, recéptus est in cælum; Ananías et Azarías et Misaël credéntes, liberáti sunt de flamma; Dániel in sua simplicitáte liberátus est de ore leónum; et ita cogitáte per generatiónem et generatiónem, quia omnes qui sperant in eum, non infirmántur. Et a verbis viri peccatóris ne timuéritis, quia glória ejus stercus et vermis est: hódie extóllitur et cras non inveniéntur, quia convérsus est in terram suam, et cogitátio ejus périit.\par
\switchcolumn
\EN
\noindent Jesus, whilst he fulfilled the word, was made ruler in Israel. Caleb, for bearing witness before the congregation, received an inheritance. David by his mercy obtained the throne of an everlasting kingdom. Elias, while he was full of zeal for the law, was taken up into heaven. Ananias and Azarias and Misael by believing, were delivered out of the flame. Daniel in his innocency was delivered out of the mouth of the lions. And thus consider through all generations: that none that trust in him fail in strength. And fear not the words of a sinful man, for his glory is dung, and worms: To day he is lifted up, and to morrow he shall not be found, because he is returned into his earth; and his thought is come to nothing.\par
\switchcolumn*

\LA
\begin{Resp}\Rsign Exáudiat Dóminus oratiónes vestras, et reconciliétur vobis, nec vos déserat in témpore malo, \Star{} Dóminus, Deus noster.\end{Resp}
\begin{Resp}\Vsign Det vobis cor ómnibus, ut colátis eum et faciátis ejus voluntátem.\end{Resp}
\begin{Resp}\Rsign Dóminus, Deus noster.\end{Resp}
\switchcolumn
\EN
\begin{Resp}\Rsign The Lord hear your prayers, and be at one with you, and never forsake you in the time of trouble \Star{} Even He, the Lord our God.\end{Resp}
\begin{Resp}\Vsign Give you all an heart to serve Him, and to do His will.\end{Resp}
\begin{Resp}\Rsign Even He, the Lord our God.\end{Resp}
\switchcolumn*

\LA
\LessonHead{Lectio III}
\switchcolumn
\EN
\LessonHead{Lesson III}
\switchcolumn*

\LA
\Cite{1 Mac 2:64-69}
\switchcolumn
\EN
\Cite{1 Mac 2:64-69}
\switchcolumn*

\LA
\noindent Vos ergo, fílii, confortámini et viríliter ágite in lege, quia in ipsa gloriósi éritis. Et ecce Simon frater vester, scio quod vir consílii est: ipsum audíte semper, et ipse erit vobis pater; et Judas Maccabǽus fortis víribus a juventúte sua sit vobis princeps milítiæ, et ipse aget bellum pópuli, et adducétis ad vos omnes factóres legis: et vindicáte vindíctam pópuli vestri; retribúite retributiónem géntibus, et inténdite in præcéptum legis. Et benedíxit eos, et appósitus est ad patres suos.\par
\switchcolumn
\EN
\noindent You therefore, my sons, take courage, and behave manfully in the law: for by it you shall be glorious. And behold, I know that your brother Simon is a man of counsel: give ear to him always, and he shall be a father to you. And Judas Machabeus who is valiant and strong from his youth up, let him be the leader of your army, and he shall manage the war of the people. And you shall take to you all that observe the law: and revenge ye the wrong of your people. Render to the Gentiles their reward, and take heed to the precepts of the law. And he blessed them, and was joined to his fathers.\par
\switchcolumn*

\LA
\begin{Resp}\Rsign Congregáti sunt inimíci nostri, et gloriántur in virtúte sua: cóntere fortitúdinem illórum, Dómine, et dispérge illos: \Star{} Ut cognóscant quia non est álius qui pugnet pro nobis, nisi tu, Deus noster.\end{Resp}
\begin{Resp}\Vsign Dispérge illos in virtúte tua, et déstrue eos, protéctor noster, Dómine.\end{Resp}
\begin{Resp}\Rsign Ut cognóscant quia non est álius qui pugnet pro nobis, nisi tu, Deus noster.\end{Resp}
\begin{Resp}\Vsign Glória Patri, et Fílio, \Star{} et Spirítui Sancto.\end{Resp}
\begin{Resp}\Rsign Ut cognóscant quia non est álius qui pugnet pro nobis, nisi tu, Deus noster.\end{Resp}
\switchcolumn
\EN
\begin{Resp}\Rsign Our enemies are gathered together, and make their boast of their own strength. O Lord, break their power, and scatter them \Star{} That they may know that there is none other that fighteth for us, but only Thou, O our God!\end{Resp}
\begin{Resp}\Vsign Scatter them in thy strength, and destroy them, O Lord our Shield.\end{Resp}
\begin{Resp}\Rsign That they may know that there is none other that fighteth for us, but only Thou, O our God!\end{Resp}
\begin{Resp}\Vsign Glory be to the Father, and to the Son, \Star{} and to the Holy Ghost.\end{Resp}
\begin{Resp}\Rsign That they may know that there is none other that fighteth for us, but only Thou, O our God!\end{Resp}
\switchcolumn*

\end{paracol}

\Office{Feria VI infra Hebdomadam I Octobris}{Friday of the first week of October}{}
\addcontentsline{toc}{subsection}{Feria VI infra Hebdomadam I Octobris \textit{— Friday of the first week of October}}
\begin{paracol}{2}
\LA
\LessonHead{Lectio I}
\switchcolumn
\EN
\LessonHead{Lesson I}
\switchcolumn*

\LA
\LessonTitle{De libro primo Machabæórum}
\switchcolumn
\EN
\LessonTitle{Lesson from the first book of Machabees}
\switchcolumn*

\LA
\Cite{1 Mac 2:70; 3:1-3; 3:5-6}
\switchcolumn
\EN
\Cite{1 Mac 2:70; 3:1-3; 3:5-6}
\switchcolumn*

\LA
\noindent Defúnctus est Mathathías anno centésimo quadragésimo sexto: et sepúltus est a fíliis suis in sepúlcris patrum suórum in Modin, et planxérunt eum omnis Israël planctu magno. Et surréxit Judas, qui vocabátur Machabǽus, fílius ejus, pro eo: et adjuvábant eum omnes fratres ejus, et univérsi qui se conjúnxerant patri ejus, et præliabántur prǽlium Israël cum lætítia. Et dilatávit glóriam pópulo suo, et persecútus est iníquos perscrútans eos et, qui conturbábant pópulum suum, eos succéndit flammis, et repúlsi sunt inimíci ejus præ timóre ejus, et omnes operárii iniquitátis conturbáti sunt: et dirécta est salus in manu ejus.\par
\switchcolumn
\EN
\noindent And he died in the hundred and forty-sixth year: and he was buried by his sons in the sepulchres of his fathers in Modin, and all Israel mourned for him with great mourning. Then his son Judas, called Machabeus, rose up in his stead. And all his brethren helped him, and all they that had joined themselves to his father, and they fought with cheerfulness the battle of Israel. And he got his people great honour, and put on a breastplate as a giant, and girt his warlike armour about him in battles, and protected the camp with his sword. And he pursued the wicked and sought them out, and them that troubled his people he burnt with fire: And his enemies were driven away for fear of him, and all the workers of iniquity were troubled: and salvation prospered in his hand.\par
\switchcolumn*

\LA
\begin{Resp}\Rsign Impetum inimicórum ne timuéritis: mémores estóte, quómodo salvi facti sunt patres nostri: \Star{} Et nunc clamémus in cælum et miserébitur nostri Deus noster.\end{Resp}
\begin{Resp}\Vsign Mementóte mirabílium ejus, quæ fecit pharaóni et exercítui ejus in Mari Rubro.\end{Resp}
\begin{Resp}\Rsign Et nunc clamémus in cælum et miserébitur nostri Deus noster.\end{Resp}
\switchcolumn
\EN
\begin{Resp}\Rsign Be ye not afraid of the assault of the enemy; remember how our fathers were delivered. \Star{} Now, therefore, let us cry unto heaven, and our God will have mercy upon us.\end{Resp}
\begin{Resp}\Vsign Remember His marvellous works that He hath done unto Pharaoh and his host in the Red Sea.\end{Resp}
\begin{Resp}\Rsign Now, therefore, let us cry unto heaven, and our God will have mercy upon us.\end{Resp}
\switchcolumn*

\LA
\LessonHead{Lectio II}
\switchcolumn
\EN
\LessonHead{Lesson II}
\switchcolumn*

\LA
\Cite{1 Mac 3:7-12}
\switchcolumn
\EN
\Cite{1 Mac 3:7-12}
\switchcolumn*

\LA
\noindent Et exacerbábat reges multos, et lætificábat Jacob in opéribus suis, et in sǽculum memória ejus in benedictióne, et perambulávit civitátes Juda, et pérdidit ímpios ex eis, et avértit iram ab Israël; et nominátus est usque ad novíssimum terræ, et congregávit pereúntes. Et congregávit Apollónius gentes, et a Samaría virtútem multam et magnam ad bellándum contra Israël. Et cognóvit Judas, et éxiit óbviam illi: et percússit, et occídit illum: et cecidérunt vulneráti multi, et réliqui fugérunt. Et accépit spólia eórum: et gládium Apollónii ábstulit Judas, et erat pugnans in eo ómnibus diébus.\par
\switchcolumn
\EN
\noindent And he grieved many kings, and made Jacob glad with his works, and his memory is blessed for ever. And he went through the cities of Juda, and destroyed the wicked out of them, and turned away wrath from Israel. And he was renowned even to the utmost part of the earth, and he gathered them that were perishing. And Apollonius gathered together the Gentiles, and a numerous and great army from Samaria, to make war against Israel. And Judas understood it, and went forth to meet him: and he overthrew him, and killed him: and many fell down slain, the rest fled away. And he took their spoils, and Judas took the sword of Apollonius, and fought with it all his lifetime.\par
\switchcolumn*

\LA
\begin{Resp}\Rsign Congregátæ sunt gentes in multitúdine, ut dímicent contra nos, et ignorámus quid ágere debeámus: \Star{} Dómine Deus, ad te sunt óculi nostri, ne pereámus.\end{Resp}
\begin{Resp}\Vsign Tu scis quæ cógitant in nos: quómodo potérimus subsístere ante fáciem illórum, nisi tu ádjuves nos?\end{Resp}
\begin{Resp}\Rsign Dómine Deus, ad te sunt óculi nostri, ne pereámus.\end{Resp}
\switchcolumn
\EN
\begin{Resp}\Rsign The heathen are assembled together to fight against us, and we know not what we should do. \Star{} Our eyes look unto thee, O Lord our God, that we should not perish.\end{Resp}
\begin{Resp}\Vsign What things they imagine against us, Thou knowest. How shall we be able to stand against them, except Thou be our help.\end{Resp}
\begin{Resp}\Rsign Our eyes look unto thee, O Lord our God, that we should not perish.\end{Resp}
\switchcolumn*

\LA
\LessonHead{Lectio III}
\switchcolumn
\EN
\LessonHead{Lesson III}
\switchcolumn*

\LA
\Cite{1 Mac 3:25-28}
\switchcolumn
\EN
\Cite{1 Mac 3:25-28}
\switchcolumn*

\LA
\noindent Et cécidit timor Judæ ac fratrum ejus, et formído super omnes gentes in circúitu eórum; et pervénit ad regem nomen ejus, et de prǽliis Judæ narrábant omnes gentes. Ut audívit autem rex Antíochus sermónes istos, irátus est ánimo et misit et congregávit exércitum univérsi regni sui, castra fórtia valde, et apéruit ærárium suum, et dedit stipéndia exercítui in annum: et mandávit illis, ut essent paráti ad ómnia.\par
\switchcolumn
\EN
\noindent And the fear of Judas and of his brethren, and the dread of them fell upon all the nations round about them. And his fame came to the king, and all nations told of the battles of Judas. Now when king Antiochus heard these words, he was angry in his mind: and he sent and gathered the forces of all his kingdom, an exceeding strong army. And he opened his treasury, and gave out pay to the army for a year: and he commanded them, that they should be ready for all things.\par
\switchcolumn*

\LA
\begin{Resp}\Rsign Tua est poténtia, tuum regnum, Dómine: tu es super omnes gentes: \Star{} Da pacem, Dómine, in diébus nostris.\end{Resp}
\begin{Resp}\Vsign Creátor ómnium, Deus, terríbilis et fortis, justus et miséricors.\end{Resp}
\begin{Resp}\Rsign Da pacem, Dómine, in diébus nostris.\end{Resp}
\begin{Resp}\Vsign Glória Patri, et Fílio, \Star{} et Spirítui Sancto.\end{Resp}
\begin{Resp}\Rsign Da pacem, Dómine, in diébus nostris.\end{Resp}
\switchcolumn
\EN
\begin{Resp}\Rsign Thine, O Lord, is the power, thine is the kingdom, O Lord, and Thou art exalted above all the heathen. \Star{} Give peace in our time, O Lord.\end{Resp}
\begin{Resp}\Vsign O Lord God, Creator of all things, Who art fearful and strong, righteous and merciful.\end{Resp}
\begin{Resp}\Rsign Give peace in our time, O Lord.\end{Resp}
\begin{Resp}\Vsign Glory be to the Father, and to the Son, \Star{} and to the Holy Ghost.\end{Resp}
\begin{Resp}\Rsign Give peace in our time, O Lord.\end{Resp}
\switchcolumn*

\end{paracol}

\Office{Sabbato infra Hebdomadam I Octobris}{Saturday of the first week of October}{}
\addcontentsline{toc}{subsection}{Sabbato infra Hebdomadam I Octobris \textit{— Saturday of the first week of October}}
\begin{paracol}{2}
\LA
\LessonHead{Lectio I}
\switchcolumn
\EN
\LessonHead{Lesson I}
\switchcolumn*

\LA
\LessonTitle{De libro primo Machabæórum}
\switchcolumn
\EN
\LessonTitle{Lesson from the first book of Machabees}
\switchcolumn*

\LA
\Cite{1 Mac 3:42-45}
\switchcolumn
\EN
\Cite{1 Mac 3:42-45}
\switchcolumn*

\LA
\noindent Et vidit Judas et fratres ejus, quia multiplicáta sunt mala, et exércitus applicábant ad fines eórum: et cognovérunt verba regis, quæ mandávit pópulo fácere in intéritum et consummatiónem; et dixérunt unusquísque ad próximum suum: Erigámus dejectiónem pópuli nostri, et pugnémus pro pópulo nostro, et sanctis nostris. Et congregátus est convéntus, ut essent paráti in prǽlium, et ut orárent et péterent misericórdiam et miseratiónes. Et Jerúsalem non habitabátur, sed erat sicut desértum: non erat qui ingrederétur et egrederétur de natis ejus, et sanctum conculcabátur: et fílii alienigenárum erant in arce: ibi erat habitátio géntium, et abláta est volúptas a Jacob, et defécit ibi tíbia et cíthara.\par
\switchcolumn
\EN
\noindent And Judas and his brethren saw that evils were multiplied, and that the armies approached to their borders: and they knew the orders the king had given to destroy the people and utterly abolish them. And they said every man to his neighbour: Let us raise up the low condition of our people, and let us fight for our people, and our sanctuary. And the assembly was gathered that they might be ready for battle: and that they might pray, and ask mercy and compassion. Now Jerusalem was not inhabited, but was like a desert: there was none of her children that went in or out: and the sanctuary was trodden down: and the children of strangers were in the castle, there was the habitation of the Gentiles: and joy was taken away from Jacob, and the pipe and harp ceased there.\par
\switchcolumn*

\LA
\begin{Resp}\Rsign Refúlsit sol in clípeos áureos, et resplenduérunt montes ab eis: \Star{} Et fortitúdo géntium dissipáta est.\end{Resp}
\begin{Resp}\Vsign Erat enim exércitus magnus valde et fortis: et appropiávit Judas et exércitus ejus in prǽlio.\end{Resp}
\begin{Resp}\Rsign Et fortitúdo géntium dissipáta est.\end{Resp}
\switchcolumn
\EN
\begin{Resp}\Rsign The sun shone upon the shields of gold, and the mountains glistered therewith; \Star{} And the army of the heathens was spread abroad.\end{Resp}
\begin{Resp}\Vsign For the army was very great and mighty then Judas and his host drew near and entered into battle.\end{Resp}
\begin{Resp}\Rsign And the army of the heathens was spread abroad.\end{Resp}
\switchcolumn*

\LA
\LessonHead{Lectio II}
\switchcolumn
\EN
\LessonHead{Lesson II}
\switchcolumn*

\LA
\Cite{1 Mac 3:46-53}
\switchcolumn
\EN
\Cite{1 Mac 3:46-53}
\switchcolumn*

\LA
\noindent Et congregáti sunt, et venérunt in Maspha contra Jerúsalem, quia locus oratiónis erat in Maspha ante in Israël, et jejunavérunt illa die, et induérunt se cilíciis, et cínerem imposuérunt cápiti suo, et discidérunt vestiménta sua: et expandérunt libros legis, de quibus scrutabántur gentes similitúdinem simulacrórum suórum: et attulérunt ornaménta sacerdotália, et primítias, et décimas: et suscitavérunt Nazarǽos qui impléverant dies, et clamavérunt voce magna in cælum, dicéntes: Quid faciémus istis et quo eos ducémus? Et sancta tua conculcáta sunt, et contamináta sunt, et sacerdótes tui facti sunt in luctum et in humilitátem: et ecce natiónes convenérunt advérsum nos, ut nos dispérdant: tu scis quæ cógitant in nos. Quómodo potérimus subsístere ante fáciem eórum, nisi tu, Deus, ádjuves nos?\par
\switchcolumn
\EN
\noindent And they assembled together, and came to Maspha over against Jerusalem: for in Maspha was a place of prayer heretofore in Israel. And they fasted that day, and put on haircloth, and put ashes upon their heads: and they rent their garments: And they laid open the books of the law, in which the Gentiles searched for the likeness of their idols: And they brought the priestly ornaments, and the firstfruits and tithes, and stirred up the Nazarites that had fulfilled their days: And they cried with a loud voice toward heaven, saying: What shall we do with these, and whither shall we carry them? For thy holies are trodden down, and are profaned, and thy priests are in mourning, and are brought low. And behold the nations are come together against us to destroy us: thou knowest what they intend against us. How shall we be able to stand before their face, unless thou, O God, help us?\par
\switchcolumn*

\LA
\begin{Resp}\Rsign Ornavérunt fáciem templi corónis áureis, et dedicavérunt altáre Dómino: \Star{} Et facta est lætítia magna in pópulo.\end{Resp}
\begin{Resp}\Vsign In hymnis et confessiónibus benedicébant Dóminum.\end{Resp}
\begin{Resp}\Rsign Et facta est lætítia magna in pópulo.\end{Resp}
\switchcolumn
\EN
\begin{Resp}\Rsign They decked the fore-front of the Temple with crowns of gold, and dedicated the Altar unto the Lord. \Star{} And there was very great gladness among the people.\end{Resp}
\begin{Resp}\Vsign They praised the Lord with Psalms and thanksgiving.\end{Resp}
\begin{Resp}\Rsign And there was very great gladness among the people.\end{Resp}
\switchcolumn*

\LA
\LessonHead{Lectio III}
\switchcolumn
\EN
\LessonHead{Lesson III}
\switchcolumn*

\LA
\Cite{1 Mac 3:54-60}
\switchcolumn
\EN
\Cite{1 Mac 3:54-60}
\switchcolumn*

\LA
\noindent Et tubis exclamavérunt voce magna. Et post hæc constítuit Judas duces pópuli, tribúnos, et centuriónes et pentacontárchos, et decuriónes, et dixit his, qui ædificábant domos, et sponsábant uxóres, et plantábant víneas, et formidolósis, ut redírent unusquísque in domum suam secúndum legem. Et movérunt castra et collocavérunt ad Austrum Emmaum. Et ait Judas: Accingímini, et estóte fílii poténtes, et estóte paráti in mane, ut pugnétis advérsus natiónes has, quæ convenérunt advérsus nos dispérdere nos, et sancta nostra; quóniam mélius est nos mori in bello, quam vidére mala gentis nostræ, et sanctórum. Sicut autem fúerit volúntas in cælo, sic fiat.\par
\switchcolumn
\EN
\noindent Then they sounded with trumpets, and cried out with a loud voice. And after this Judas appointed captains over the people, over thousands, and over hundreds, and over fifties, and over tens. And he said to them that were building houses, or had betrothed wives, or were planting vineyards, or were fearful, that they should return every man to his house, according to the law. So they removed the camp, and pitched on the south side of Emmaus. And Judas said: Gird yourselves, and be valiant men, and be ready against the morning, that you may fight with these nations that are assembled against us to destroy us and our sanctuary. For it is better for us to die in battle, than to see the evils of our nation, and of the holies: Nevertheless as it shall be the will of God in heaven so be it done.\par
\switchcolumn*

\LA
\begin{Resp}\Rsign In hymnis et confessiónibus benedicébant Dóminum, \Star{} Qui magna fecit in Israël, et victóriam dedit illis Dóminus omnípotens.\end{Resp}
\begin{Resp}\Vsign Ornavérunt fáciem templi corónis áureis, et dedicavérunt altáre Dómino.\end{Resp}
\begin{Resp}\Rsign Qui magna fecit in Israël, et victóriam dedit illis Dóminus omnípotens.\end{Resp}
\begin{Resp}\Vsign Glória Patri, et Fílio, \Star{} et Spirítui Sancto.\end{Resp}
\begin{Resp}\Rsign Qui magna fecit in Israël, et victóriam dedit illis Dóminus omnípotens.\end{Resp}
\switchcolumn
\EN
\begin{Resp}\Rsign They praised the Lord with psalms and thanksgiving \Star{} Who had done so great things in Israel, and given them the victory the Lord Almighty.\end{Resp}
\begin{Resp}\Vsign They decked the fore-front of the Temple with crowns of gold, and dedicated the Altar unto the Lord.\end{Resp}
\begin{Resp}\Rsign Who had done so great things for Israel, and given them the victory the Lord Almighty.\end{Resp}
\begin{Resp}\Vsign Glory be to the Father, and to the Son, \Star{} and to the Holy Ghost.\end{Resp}
\begin{Resp}\Rsign Who had done so great things for Israel, and given them the victory the Lord Almighty.\end{Resp}
\switchcolumn*

\end{paracol}

\Office{Dominica II Octobris}{Second Sunday of October}{}
\addcontentsline{toc}{subsection}{Dominica II Octobris \textit{— Second Sunday of October}}
\begin{paracol}{2}
\LA
\LessonHead{Lectio I}
\switchcolumn
\EN
\LessonHead{Lesson I}
\switchcolumn*

\LA
\LessonTitle{De libro primo Machabæórum}
\switchcolumn
\EN
\LessonTitle{Lesson from the first book of Machabees}
\switchcolumn*

\LA
\Cite{1 Mac 4:36-40}
\switchcolumn
\EN
\Cite{1 Mac 4:36-40}
\switchcolumn*

\LA
\noindent Dixit autem Judas et fratres ejus: Ecce contríti sunt inimíci nostri; ascendámus nunc mundáre sancta, et renováre. Et congregátus est omnis exércitus, et ascendérunt in montem Sion, et vidérunt sanctificatiónem desértam et altáre profanátum et portas exústas et in átriis virgúlta nata, sicut in saltu vel in móntibus, et pastophória díruta. Et scidérunt vestiménta sua et planxérunt planctu magno et imposuérunt cínerem super caput suum, et cecidérunt in fáciem super terram, et exclamavérunt tubis signórum, et clamavérunt in cælum.\par
\switchcolumn
\EN
\noindent Then Judas, and his brethren said: Behold our enemies are discomfited: let us go up now to cleanse the holy places and to repair them. And all the army assembled together, and they went up into mount Sion. And they saw the sanctuary desolate, and the altar profaned, and the gates burnt, and shrubs growing up in the courts as in a forest, or on the mountains, and the chambers joining to the temple thrown down. And they rent their garments, and made great lamentation, and put ashes on their heads: And they fell face down to the ground on their faces, and they sounded with the trumpets of alarm, and they cried towards heaven.\par
\switchcolumn*

\LA
\begin{Resp}\Rsign Adapériat Dóminus cor vestrum in lege sua et in præcéptis suis et fáciat pacem in diébus vestris: \Star{} Concédat vobis salútem, et rédimat vos a malis.\end{Resp}
\begin{Resp}\Vsign Exáudiat Dóminus oratiónes vestras, et reconciliétur vobis, nec vos déserat in témpore malo.\end{Resp}
\begin{Resp}\Rsign Concédat vobis salútem, et rédimat vos a malis.\end{Resp}
\switchcolumn
\EN
\begin{Resp}\Rsign The Lord open your hearts in His law and commandments, and send peace in your days. \Star{} May He grant you salvation and redeem you out of all evil.\end{Resp}
\begin{Resp}\Vsign The Lord hear your prayers, and be at one with you, and never forsake you in the time of trouble.\end{Resp}
\begin{Resp}\Rsign May He grant you salvation and redeem you out of all evil.\end{Resp}
\switchcolumn*

\LA
\LessonHead{Lectio II}
\switchcolumn
\EN
\LessonHead{Lesson II}
\switchcolumn*

\LA
\Cite{1 Mac 4:41-46}
\switchcolumn
\EN
\Cite{1 Mac 4:41-46}
\switchcolumn*

\LA
\noindent Tunc ordinávit Judas viros ut pugnárent advérsus eos, qui erant in arce, donec emundárent sancta. Et elégit sacerdótes sine mácula voluntátem habéntes in lege Dei, et mundavérunt sancta et tulérunt lápides contaminatiónis in locum immúndum. Et cogitávit de altári holocaustórum, quod profanátum erat, quid de eo fáceret. Et íncidit illis consílium bonum ut destrúerent illud: ne forte illis esset in oppróbrium, quia contaminavérunt illud gentes, et demolíti sunt illud, et reposuérunt lápides in monte domus in loco apto, quoadúsque veníret prophéta et respondéret de eis.\par
\switchcolumn
\EN
\noindent Then Judas appointed men to fight against them that were in the castle, till they had cleansed the holy places. And he chose priests without blemish, whose will was set upon the law of God: And they cleansed the holy places, and took away the stones that had been defiled into an unclean place. And he considered about the altar of holocausts that had been profaned, what he should do with it. And a good counsel came into their minds, to pull it down: lest it should be a reproach to them, because the Gentiles had defiled it; so they threw it down. And they laid up the stones in the mountain of the temple in a convenient place, till there should come a prophet, and give answer concerning them.\par
\switchcolumn*

\LA
\begin{Resp}\Rsign Exáudiat Dóminus oratiónes vestras, et reconciliétur vobis, nec vos déserat in témpore malo, \Star{} Dóminus, Deus noster.\end{Resp}
\begin{Resp}\Vsign Det vobis cor ómnibus, ut colátis eum et faciátis ejus voluntátem.\end{Resp}
\begin{Resp}\Rsign Dóminus, Deus noster.\end{Resp}
\switchcolumn
\EN
\begin{Resp}\Rsign The Lord hear your prayers, and be at one with you, and never forsake you in the time of trouble \Star{} Even He, the Lord our God.\end{Resp}
\begin{Resp}\Vsign Give you all an heart to serve Him, and to do His will.\end{Resp}
\begin{Resp}\Rsign Even He, the Lord our God.\end{Resp}
\switchcolumn*

\LA
\LessonHead{Lectio III}
\switchcolumn
\EN
\LessonHead{Lesson III}
\switchcolumn*

\LA
\Cite{1 Mac 4:47-51}
\switchcolumn
\EN
\Cite{1 Mac 4:47-51}
\switchcolumn*

\LA
\noindent Et accepérunt lápides íntegros secúndum legem et ædificavérunt altáre novum secúndum illud quod fuit prius: et ædificavérunt sancta, et quæ intra domum erant intrínsecus: et ædem, et átria sanctificavérunt, et fecérunt vasa sancta nova, et intulérunt candelábrum, et altáre incensórum, et mensam in templum, et incénsum posuérunt super altáre, et accendérunt lucérnas, quæ super candelábrum erant, et lucébant in templo. Et posuérunt super mensam panes, et appendérunt vela, et consummavérunt ómnia ópera quæ fécerant.\par
\switchcolumn
\EN
\noindent Then they took whole stones according to the law, and built a new altar according to the former: And they built up the holy places, and the things that were within the temple: and they sanctified the temple, and the courts. And they made new holy vessels, and brought in the candlestick, and the altar of incense, and the table into the temple. And they put incense upon the altar, and lighted up the lamps that were upon the candlestick, and they gave light in the temple. And they set the loaves upon the table, and hung up the veils, and finished all the works that they had begun to make.\par
\switchcolumn*

\LA
\begin{Resp}\Rsign Congregáti sunt inimíci nostri, et gloriántur in virtúte sua: cóntere fortitúdinem illórum, Dómine, et dispérge illos: \Star{} Ut cognóscant quia non est álius qui pugnet pro nobis, nisi tu, Deus noster.\end{Resp}
\begin{Resp}\Vsign Dispérge illos in virtúte tua, et déstrue eos, protéctor noster, Dómine.\end{Resp}
\begin{Resp}\Rsign Ut cognóscant quia non est álius qui pugnet pro nobis, nisi tu, Deus noster.\end{Resp}
\begin{Resp}\Vsign Glória Patri, et Fílio, \Star{} et Spirítui Sancto.\end{Resp}
\begin{Resp}\Rsign Ut cognóscant quia non est álius qui pugnet pro nobis, nisi tu, Deus noster.\end{Resp}
\switchcolumn
\EN
\begin{Resp}\Rsign Our enemies are gathered together, and make their boast of their own strength. O Lord, break their power, and scatter them \Star{} That they may know that there is none other that fighteth for us, but only Thou, O our God!\end{Resp}
\begin{Resp}\Vsign Scatter them in thy strength, and destroy them, O Lord our Shield.\end{Resp}
\begin{Resp}\Rsign That they may know that there is none other that fighteth for us, but only Thou, O our God!\end{Resp}
\begin{Resp}\Vsign Glory be to the Father, and to the Son, \Star{} and to the Holy Ghost.\end{Resp}
\begin{Resp}\Rsign That they may know that there is none other that fighteth for us, but only Thou, O our God!\end{Resp}
\switchcolumn*

\LA
\LessonHead{Lectio IV}
\switchcolumn
\EN
\LessonHead{Lesson IV}
\switchcolumn*

\LA
\LessonTitle{Ex libro sancti Augustíni Epíscopi de Civitáte Dei}
\switchcolumn
\EN
\LessonTitle{From the City of God of St. Augustine, Bishop of Hippo.}
\switchcolumn*

\LA
\Source{Liber 18, cap. 45}
\switchcolumn
\EN
\Source{Book 18. ch. 45}
\switchcolumn*

\LA
\noindent Posteáquam gens Judǽa cœpit non habére prophétas, proculdúbio detérior facta est, eo scílicet témpore, quo se sperábat instauráto templo post captivitátem, quæ fuit in Babylónia, futúram esse meliórem. Sic quippe intelligébat pópulus ille carnális, quod prænuntiátum est per Aggæum prophétam dicéntem: Magna erit glória domus istíus novíssimæ, plus quam primæ. Quod de novo testaménto dictum esse, paulo supérius demonstrávit, ubi ait apérte Christum promíttens: Et movébo omnes gentes, et véniet Desiderátus cunctis Géntibus.\par
\switchcolumn
\EN
\noindent The Jewish nation no doubt became worse after it ceased to have prophets, just at the very time when, on the rebuilding of the temple after the captivity in Babylon, it hoped to become better. For so, indeed, did that carnal people understand what was foretold by Haggai the prophet, saying, The glory of this latter house shall be greater than that of the former. Now, that this is said of the new testament, he showed a little above, where he says, evidently promising Christ, And I will move all nations, and the desired One shall come to all nations.\par
\switchcolumn*

\LA
\begin{Resp}\Rsign Impetum inimicórum ne timuéritis: mémores estóte, quómodo salvi facti sunt patres nostri: \Star{} Et nunc clamémus in cælum et miserébitur nostri Deus noster.\end{Resp}
\begin{Resp}\Vsign Mementóte mirabílium ejus, quæ fecit pharaóni et exercítui ejus in Mari Rubro.\end{Resp}
\begin{Resp}\Rsign Et nunc clamémus in cælum et miserébitur nostri Deus noster.\end{Resp}
\switchcolumn
\EN
\begin{Resp}\Rsign Be ye not afraid of the assault of the enemy; remember how our fathers were delivered. \Star{} Now, therefore, let us cry unto heaven, and our God will have mercy upon us.\end{Resp}
\begin{Resp}\Vsign Remember His marvellous works that He hath done unto Pharaoh and his host in the Red Sea.\end{Resp}
\begin{Resp}\Rsign Now, therefore, let us cry unto heaven, and our God will have mercy upon us.\end{Resp}
\switchcolumn*

\LA
\LessonHead{Lectio V}
\switchcolumn
\EN
\LessonHead{Lesson V}
\switchcolumn*

\LA
\noindent Tálibus enim eléctis Géntium, domus Dei ædificátur per Testaméntum novum lapídibus vivis, longe gloriósior, quam templum illud fuit, quod a rege Salomóne constrúctum est, et post captivitátem instaurátum. Propter hoc ergo nec prophétas ex illo témpore hábuit illa gens; sed multis cládibus afflícta est ab alienígenis régibus, ipsísque Románis, ne hanc Aggæi prophetíam in illa instauratióne templi opinarétur implétam. Non multo enim post, adveniénte Alexándro, subjugáta est; quando, etsi nulla facta est vastátio, quóniam non sunt ausi ei resístere, et ídeo placátum facíllime súbditi recepérunt; non erat tamen glória tanta domus illíus, quanta fuit in suórum regum líbera potestáte.\par
\switchcolumn
\EN
\noindent For by such chosen ones of the nations there is built, through the new testament, with living stones, a house of God far more glorious than that temple was which was constructed by king Solomon, and rebuilt after the captivity. For this reason, then, that nation had no prophets from that time, but was afflicted with many plagues by kings of alien race, and by the Romans themselves, lest they should fancy that this prophecy of Haggai was fulfilled by that rebuilding of the temple.For not long after, on the arrival of Alexander, it was subdued, when, although there was no pillaging, because they dared not resist him, and thus, being very easily subdued, received him peaceably, yet the glory of that house was not so great as it was when under the free power of their own kings.\par
\switchcolumn*

\LA
\begin{Resp}\Rsign Congregátæ sunt gentes in multitúdine, ut dímicent contra nos, et ignorámus quid ágere debeámus: \Star{} Dómine Deus, ad te sunt óculi nostri, ne pereámus.\end{Resp}
\begin{Resp}\Vsign Tu scis quæ cógitant in nos: quómodo potérimus subsístere ante fáciem illórum, nisi tu ádjuves nos?\end{Resp}
\begin{Resp}\Rsign Dómine Deus, ad te sunt óculi nostri, ne pereámus.\end{Resp}
\switchcolumn
\EN
\begin{Resp}\Rsign The heathen are assembled together to fight against us, and we know not what we should do. \Star{} Our eyes look unto thee, O Lord our God, that we should not perish.\end{Resp}
\begin{Resp}\Vsign What things they imagine against us, Thou knowest. How shall we be able to stand against them, except Thou be our help.\end{Resp}
\begin{Resp}\Rsign Our eyes look unto thee, O Lord our God, that we should not perish.\end{Resp}
\switchcolumn*

\LA
\LessonHead{Lectio VI}
\switchcolumn
\EN
\LessonHead{Lesson VI}
\switchcolumn*

\LA
\noindent Deínde Ptolemǽus Lagi fílius, post Alexándri mortem captívos inde in Ægýptum tránstulit, quos ejus succéssor Ptolemǽus Philadélphus benevolentíssime inde dimísit: per quem factum est, ut Septuagínta intérpretum Scriptúras haberémus. Deínde contríti sunt bellis, quæ in Machabæórum libris explicántur. Post hæc capti a rege Alexandríæ Ptolemǽo qui est appellátus Epíphanes, inde ab Antíocho rege Sýriæ multis et gravíssimis malis ad idóla colénda compúlsi: templúmque ipsum replétum sacrílegis superstitiónibus Géntium, quod tamen dux eórum strenuíssimus Judas, qui étiam Machabǽus dictus est, Antíochi dúcibus pulsis, ab omni illa idololatríæ contaminatióne mundávit.\par
\switchcolumn
\EN
\noindent Then Ptolemy son of Lagus, after Alexander's death carried them captive into Egypt. His successor, Ptolemy Philadelphus, most benevolently dismissed them; and by him it was brought about, as I have narrated a little before, that we should have the Septuagint version of the Scriptures. Then they were crushed by the wars which are explained in the books of the Maccabees. Afterward they were taken captive by Ptolemy king of Alexandria, who was called Epiphanes. Then Antiochus king of Syria compelled them by many and most grievous evils to worship idols, and filled the temple itself with the sacrilegious superstitions of the Gentiles. Yet their most vigorous leader Judas, who is also called Maccabæus, after beating the generals of Antiochus, cleansed it from all that defilement of idolatry.\par
\switchcolumn*

\LA
\begin{Resp}\Rsign Tua est poténtia, tuum regnum, Dómine: tu es super omnes gentes: \Star{} Da pacem, Dómine, in diébus nostris.\end{Resp}
\begin{Resp}\Vsign Creátor ómnium, Deus, terríbilis et fortis, justus et miséricors.\end{Resp}
\begin{Resp}\Rsign Da pacem, Dómine, in diébus nostris.\end{Resp}
\begin{Resp}\Vsign Glória Patri, et Fílio, \Star{} et Spirítui Sancto.\end{Resp}
\begin{Resp}\Rsign Da pacem, Dómine, in diébus nostris.\end{Resp}
\switchcolumn
\EN
\begin{Resp}\Rsign Thine, O Lord, is the power, thine is the kingdom, O Lord, and Thou art exalted above all the heathen. \Star{} Give peace in our time, O Lord.\end{Resp}
\begin{Resp}\Vsign O Lord God, Creator of all things, Who art fearful and strong, righteous and merciful.\end{Resp}
\begin{Resp}\Rsign Give peace in our time, O Lord.\end{Resp}
\begin{Resp}\Vsign Glory be to the Father, and to the Son, \Star{} and to the Holy Ghost.\end{Resp}
\begin{Resp}\Rsign Give peace in our time, O Lord.\end{Resp}
\switchcolumn*

\end{paracol}

\Office{Feria II infra Hebdomadam II Octobris}{Monday of the second week of October}{}
\addcontentsline{toc}{subsection}{Feria II infra Hebdomadam II Octobris \textit{— Monday of the second week of October}}
\begin{paracol}{2}
\LA
\LessonHead{Lectio I}
\switchcolumn
\EN
\LessonHead{Lesson I}
\switchcolumn*

\LA
\LessonTitle{De libro primo Machabæórum}
\switchcolumn
\EN
\LessonTitle{Lesson from the first book of Machabees}
\switchcolumn*

\LA
\Cite{1 Mac 4:52-55}
\switchcolumn
\EN
\Cite{1 Mac 4:52-55}
\switchcolumn*

\LA
\noindent Ante matutínum surrexérunt quinta et vigésima die mensis noni (hic est mensis Cásleu) centésimi quadragésimi octávi anni: et obtulérunt sacrifícium secúndum legem super altáre holocaustórum novum, quod fecérunt. Secúndum tempus et secúndum diem, in qua contaminavérunt illud gentes, in ipsa renovátum est in cánticis, et cítharis, et cínyris, et cýmbalis. Et cécidit omnis pópulus in fáciem, et adoravérunt, et benedixérunt in cælum eum, qui prosperávit eis.\par
\switchcolumn
\EN
\noindent And they arose before the morning on the five and twentieth day of the ninth month (which is the month of Casleu) in the hundred and forty-eighth year. And they offered sacrifice according to the law upon the new altar of holocausts which they had made. According to the time, and according to the day wherein the heathens had defiled it, in the same was it dedicated anew with canticles, and harps, and lutes, and cymbals. And all the people fell upon their faces, and adored, and blessed up to heaven, him that had prospered them.\par
\switchcolumn*

\LA
\begin{Resp}\Rsign Dixit Judas Simóni fratri suo: Elige tibi viros et vade, líbera fratres tuos in Galilǽam; ego autem et Jónathas frater tuus íbimus in Galaadítim: \Star{} Sicut fúerit volúntas in cælo, sic fiat.\end{Resp}
\begin{Resp}\Vsign Accingímini, fílii poténtes, et estóte paráti: quóniam mélius est nobis, mori in bello, quam vidére mala gentis nostræ et sanctórum.\end{Resp}
\begin{Resp}\Rsign Sicut fúerit volúntas in cælo, sic fiat.\end{Resp}
\switchcolumn
\EN
\begin{Resp}\Rsign Judas said unto Simon his brother: Choose thee out men, and go, and deliver thy brethren that are in Galilee and I, and Jonathan thy brother, will go into the country of Galaad. \Star{} As the Will is in heaven, so let it be.\end{Resp}
\begin{Resp}\Vsign Arm yourselves, ye valiant men, and be in readiness for it is better for us to die in battle, than to behold the calamities of our people, and our sanctuary.\end{Resp}
\begin{Resp}\Rsign As the Will is in heaven, so let it be.\end{Resp}
\switchcolumn*

\LA
\LessonHead{Lectio II}
\switchcolumn
\EN
\LessonHead{Lesson II}
\switchcolumn*

\LA
\Cite{1 Mac 4:56-59}
\switchcolumn
\EN
\Cite{1 Mac 4:56-59}
\switchcolumn*

\LA
\noindent Et fecérunt dedicatiónem altáris diébus octo, et obtulérunt holocáusta cum lætítia, et sacrifícium salutáris et laudis, et ornavérunt fáciem templi corónis áureis et scútulis, et dedicavérunt portas et pastophória, et imposuérunt eis jánuas. Et facta est lætítia in pópulo magna valde, et avérsum est oppróbrium géntium. Et státuit Judas et fratres ejus et univérsa ecclésia Israël, ut agátur dies dedicatiónis altáris in tempóribus suis ab anno in annum per dies octo, a quinta et vigésima die mensis Cásleu, cum lætítia et gáudio.\par
\switchcolumn
\EN
\noindent And they kept the dedication of the altar eight days, and they offered holocausts with joy, and sacrifices of salvation, and of praise. And they adorned the front of the temple with crowns of gold, and escutcheons, and they renewed the gates, and the chambers, and hanged doors upon them. And there was exceeding great joy among the people, and the reproach of the Gentiles was turned away. And Judas, and his brethren, and all the church of Israel decreed, that the day of the dedication of the altar should be kept in its season from year to year for eight days, from the five and twentieth day of the month of Casleu, with joy and gladness.\par
\switchcolumn*

\LA
\begin{Resp}\Rsign Ornavérunt fáciem templi corónis áureis, et dedicavérunt altáre Dómino: \Star{} Et facta est lætítia magna in pópulo.\end{Resp}
\begin{Resp}\Vsign In hymnis et confessiónibus benedicébant Dóminum.\end{Resp}
\begin{Resp}\Rsign Et facta est lætítia magna in pópulo.\end{Resp}
\switchcolumn
\EN
\begin{Resp}\Rsign They decked the fore-front of the Temple with crowns of gold, and dedicated the Altar unto the Lord. \Star{} And there was very great gladness among the people.\end{Resp}
\begin{Resp}\Vsign They praised the Lord with Psalms and thanksgiving.\end{Resp}
\begin{Resp}\Rsign And there was very great gladness among the people.\end{Resp}
\switchcolumn*

\LA
\LessonHead{Lectio III}
\switchcolumn
\EN
\LessonHead{Lesson III}
\switchcolumn*

\LA
\Cite{1 Mac 4:60-61}
\switchcolumn
\EN
\Cite{1 Mac 4:60-61}
\switchcolumn*

\LA
\noindent Et ædificavérunt in témpore illo montem Sion, et per circúitum muros altos et turres firmas, nequándo venírent gentes et conculcárent eum, sicut ántea fecérunt. Et collocávit illic exércitum, ut servárent eum, et munívit eum ad custodiéndam Bethsúram, ut habéret pópulus munitiónem contra fáciem Idumǽæ.\par
\switchcolumn
\EN
\noindent They built up also at that time mount Sion, with high walls, and strong towers round about, lest the Gentiles should at any time come, and tread it down as they did before. And he placed a garrison there to keep it, and he fortified it to secure Bethsura, that the people might have a defence against Idumea.\par
\switchcolumn*

\LA
\begin{Resp}\Rsign In hymnis et confessiónibus benedicébant Dóminum, \Star{} Qui magna fecit in Israël, et victóriam dedit illis Dóminus omnípotens.\end{Resp}
\begin{Resp}\Vsign Ornavérunt fáciem templi corónis áureis, et dedicavérunt altáre Dómino.\end{Resp}
\begin{Resp}\Rsign Qui magna fecit in Israël, et victóriam dedit illis Dóminus omnípotens.\end{Resp}
\begin{Resp}\Vsign Glória Patri, et Fílio, \Star{} et Spirítui Sancto.\end{Resp}
\begin{Resp}\Rsign Qui magna fecit in Israël, et victóriam dedit illis Dóminus omnípotens.\end{Resp}
\switchcolumn
\EN
\begin{Resp}\Rsign They praised the Lord with psalms and thanksgiving \Star{} Who had done so great things in Israel, and given them the victory the Lord Almighty.\end{Resp}
\begin{Resp}\Vsign They decked the fore-front of the Temple with crowns of gold, and dedicated the Altar unto the Lord.\end{Resp}
\begin{Resp}\Rsign Who had done so great things for Israel, and given them the victory the Lord Almighty.\end{Resp}
\begin{Resp}\Vsign Glory be to the Father, and to the Son, \Star{} and to the Holy Ghost.\end{Resp}
\begin{Resp}\Rsign Who had done so great things for Israel, and given them the victory the Lord Almighty.\end{Resp}
\switchcolumn*

\end{paracol}

\Office{Feria III infra Hebdomadam II Octobris}{Tuesday of the second week of October}{}
\addcontentsline{toc}{subsection}{Feria III infra Hebdomadam II Octobris \textit{— Tuesday of the second week of October}}
\begin{paracol}{2}
\LA
\LessonHead{Lectio I}
\switchcolumn
\EN
\LessonHead{Lesson I}
\switchcolumn*

\LA
\LessonTitle{De libro primo Machabæórum}
\switchcolumn
\EN
\LessonTitle{Lesson from the first book of Machabees}
\switchcolumn*

\LA
\Cite{1 Mac 5:1-5}
\switchcolumn
\EN
\Cite{1 Mac 5:1-5}
\switchcolumn*

\LA
\noindent Et factum est, ut audiérunt gentes in circúitu quia ædificátum est altáre et sanctuárium sicut prius, irátæ sunt valde: et cogitábant tóllere genus Jacob, qui erant inter eos, et cœpérunt occídere de pópulo, et pérsequi. Et debellábat Judas fílios Esau in Idumǽa, et eos qui erant in Acrabatháne, quia circumsedébant Israëlítas, et percússit eos plaga magna, et recordátus est malítiam filiórum Bean, qui erant pópulo in láqueum et in scándalum, insidiántes ei in via. Et conclúsi sunt ab eo in túrribus, et applícuit ad eos, et anathematizávit eos et incéndit turres eórum igni cum ómnibus qui in eis erant.\par
\switchcolumn
\EN
\noindent Now it came to pass, when the nations round about heard that the altar and the sanctuary were built up as before, that they were exceeding angry. And they thought to destroy the generation of Jacob that were among them, and they began to kill some of the people, and to persecute them. Then Judas fought against the children of Esau in Idumea, and them that were in Acrabathane: because they beset the Israelites around about, and he made a great slaughter of them. And he remembered the malice of the children of Bean: who were a snare and a stumblingblock to the people, by lying in wait for them in the way. And they were shut up by him in towers, and he set upon them, and devoted them to utter destruction, and burnt their towers with fire, and all that were in them.\par
\switchcolumn*

\LA
\begin{Resp}\Rsign Hic est fratrum amátor et pópuli Israël: \Star{} Hic est, qui multum orat pro pópulo et univérsa sancta civitáte Jerúsalem.\end{Resp}
\begin{Resp}\Vsign Vir iste in pópulo suo mitíssimus appáruit.\end{Resp}
\begin{Resp}\Rsign Hic est, qui multum orat pro pópulo et univérsa sancta civitáte Jerúsalem.\end{Resp}
\switchcolumn
\EN
\begin{Resp}\Rsign This is a lover of the brethren, and of the people of Israel \Star{} This is one who prayeth much for the people, and for all the Holy City, Jerusalem.\end{Resp}
\begin{Resp}\Vsign There appeared a man most gentle toward all his people.\end{Resp}
\begin{Resp}\Rsign This is one who prayeth much for the people, and for all the Holy City, Jerusalem.\end{Resp}
\switchcolumn*

\LA
\LessonHead{Lectio II}
\switchcolumn
\EN
\LessonHead{Lesson II}
\switchcolumn*

\LA
\Cite{1 Mac 5:6-9}
\switchcolumn
\EN
\Cite{1 Mac 5:6-9}
\switchcolumn*

\LA
\noindent Et transívit ad fílios Ammon, et invénit manum fortem, et pópulum copiósum, et Timótheum ducem ipsórum: et commísit cum eis prǽlia multa, et contríti sunt in conspéctu eórum, et percússit eos: et cepit Gazer civitátem et fílias ejus, et revérsus est in Judǽam. Et congregátæ sunt gentes quæ sunt in Gálaad, advérsus Israëlítas, qui erant in fínibus eórum, ut tóllerent eos. Et fugérunt in Dátheman munitiónem.\par
\switchcolumn
\EN
\noindent Then he passed over to the children of Ammon, where he found a mighty power, and much people, and Timotheus was their captain: And he fought many battles with them, and they were discomfited in their sight, and he smote them: And he took the city of Gazer and her towns, and returned into Judea. And the Gentiles that were in Galaad, assembled themselves together against the Israelites that were in their quarters to destroy them: and they fled into the fortress of Datheman.\par
\switchcolumn*

\LA
\begin{Resp}\Rsign Tu, Dómine universórum, qui nullam habes indigéntiam, voluísti templum tuum fíeri in nobis; \Star{} Consérva domum istam immaculátam in ætérnum, Dómine.\end{Resp}
\begin{Resp}\Vsign Tu elegísti, Dómine, domum istam ad invocándum nomen tuum in ea, ut esset domus oratiónis et obsecratiónis pópulo tuo.\end{Resp}
\begin{Resp}\Rsign Consérva domum istam immaculátam in ætérnum, Dómine.\end{Resp}
\switchcolumn
\EN
\begin{Resp}\Rsign Thou, O Lord of all things, Who hast need of nothing, wast pleased that the Temple of thine habitation should be among us. \Star{} Therefore now, O Lord, keep this house ever undefiled\end{Resp}
\begin{Resp}\Vsign Thou, O Lord, didst choose this house, that thy Name should be called on therein, and to be an house of prayer and petition for thy people.\end{Resp}
\begin{Resp}\Rsign Therefore now, O Lord, keep this house ever undefiled.\end{Resp}
\switchcolumn*

\LA
\LessonHead{Lectio III}
\switchcolumn
\EN
\LessonHead{Lesson III}
\switchcolumn*

\LA
\Cite{1 Mac 5:10-13}
\switchcolumn
\EN
\Cite{1 Mac 5:10-13}
\switchcolumn*

\LA
\noindent Et misérunt lítteras ad Judam et fratres ejus, dicéntes: Congregátæ sunt advérsum nos gentes per circúitum, ut nos áuferant, et parant veníre et occupáre munitiónem, in quam confúgimus: et Timótheus est dux exércitus eórum; nunc ergo veni, et éripe nos de mánibus eórum, quia cécidit multitúdo de nobis. Et omnes fratres nostri, qui erant in locis Tubin, interfécti sunt, et captívas duxérunt uxóres eórum, et natos, et spólia, et peremérunt illic fere mille viros.\par
\switchcolumn
\EN
\noindent And they sent letters to Judas and his brethren, saying, The heathens that are round about are gathered together against us, to destroy us: And they are preparing to come, and to take the fortress into which we are fled: and Timotheus is the captain of their host. Now therefore come, and deliver us out of their hands, for many of us are slain. And all our brethren that were in the places of Tubin, are killed: and they have carried away their wives, and their children, captives, and taken their spoils, and they have slain there almost a thousand men.\par
\switchcolumn*

\LA
\begin{Resp}\Rsign Aperi óculos tuos, Dómine, et vide afflictiónem nostram: circumdedérunt nos gentes ad puniéndum nos: \Star{} Sed tu, Dómine, exténde brácchium tuum, et líbera ánimas nostras.\end{Resp}
\begin{Resp}\Vsign Afflíge oppriméntes nos et contuméliam faciéntes in supérbiam; et custódi partem tuam.\end{Resp}
\begin{Resp}\Rsign Sed tu, Dómine, exténde brácchium tuum, et líbera ánimas nostras.\end{Resp}
\begin{Resp}\Vsign Glória Patri, et Fílio, \Star{} et Spirítui Sancto.\end{Resp}
\begin{Resp}\Rsign Sed tu, Dómine, exténde brácchium tuum, et líbera ánimas nostras.\end{Resp}
\switchcolumn
\EN
\begin{Resp}\Rsign Open thine eyes, O Lord, and behold our affliction for the heathen are come round about us to punish us. \Star{} But Thou, O Lord, stretch forth thine arm, and deliver our souls.\end{Resp}
\begin{Resp}\Vsign Punish them that oppress us and with pride do us wrong, and keep thine own portion.\end{Resp}
\begin{Resp}\Rsign But Thou, O Lord, stretch forth thine arm, and deliver our souls.\end{Resp}
\begin{Resp}\Vsign Glory be to the Father, and to the Son, \Star{} and to the Holy Ghost.\end{Resp}
\begin{Resp}\Rsign But Thou, O Lord, stretch forth thine arm, and deliver our souls.\end{Resp}
\switchcolumn*

\end{paracol}

\Office{Feria IV infra Hebdomadam II Octobris}{Wednesday of the second week of October}{}
\addcontentsline{toc}{subsection}{Feria IV infra Hebdomadam II Octobris \textit{— Wednesday of the second week of October}}
\begin{paracol}{2}
\LA
\LessonHead{Lectio I}
\switchcolumn
\EN
\LessonHead{Lesson I}
\switchcolumn*

\LA
\LessonTitle{De libro primo Machabæórum}
\switchcolumn
\EN
\LessonTitle{Lesson from the first book of Machabees}
\switchcolumn*

\LA
\Cite{1 Mac 5:55-58}
\switchcolumn
\EN
\Cite{1 Mac 5:55-58}
\switchcolumn*

\LA
\noindent In diébus, quibus erat Judas et Jónathas in terra Gálaad, et Simon frater ejus in Galilǽa contra fáciem Ptolemáidis, audívit Joséphus Zacharíæ fílius et Azarías princeps virtútis res bene gestas et prǽlia, quæ facta sunt, et dixit: Faciámus et ipsi nobis nomen, et eámus pugnáre advérsus gentes, quæ in circúitu nostro sunt. Et præcépit his, qui erant in exércitu suo, et abiérunt Jámniam.\par
\switchcolumn
\EN
\noindent Now in the days that Judas and Jonathan were in the land of Galaad, and Simon his brother in Galilee before Ptolemais, Joseph the son of Zacharias, and Azarias captain of the soldiers, heard of the good success, and the battles that were fought. And he said: Let us also get us a name, and let us go fight against the Gentiles that are round about us. And he gave charge to them that were in his army, and they went towards Jamnia.\par
\switchcolumn*

\LA
\begin{Resp}\Rsign Refúlsit sol in clípeos áureos, et resplenduérunt montes ab eis: \Star{} Et fortitúdo géntium dissipáta est.\end{Resp}
\begin{Resp}\Vsign Erat enim exércitus magnus valde et fortis: et appropiávit Judas et exércitus ejus in prǽlio.\end{Resp}
\begin{Resp}\Rsign Et fortitúdo géntium dissipáta est.\end{Resp}
\switchcolumn
\EN
\begin{Resp}\Rsign The sun shone upon the shields of gold, and the mountains glistered therewith; \Star{} And the army of the heathens was spread abroad.\end{Resp}
\begin{Resp}\Vsign For the army was very great and mighty then Judas and his host drew near and entered into battle.\end{Resp}
\begin{Resp}\Rsign And the army of the heathens was spread abroad.\end{Resp}
\switchcolumn*

\LA
\LessonHead{Lectio II}
\switchcolumn
\EN
\LessonHead{Lesson II}
\switchcolumn*

\LA
\Cite{1 Mac 5:59-62}
\switchcolumn
\EN
\Cite{1 Mac 5:59-62}
\switchcolumn*

\LA
\noindent Et exívit Górgias de civitáte et viri ejus óbviam illis in pugnam, et fugáti sunt Joséphus et Azarías usque in fines Judǽæ, et cecidérunt illo die de pópulo Israël ad duo míllia viri, et facta est fuga magna in pópulo: quia non audiérunt Judam et fratres ejus existimántes fórtiter se factúros. Ipsi autem non erant de sémine virórum illórum, per quos salus facta est in Israël.\par
\switchcolumn
\EN
\noindent And Gorgias and his men went out of the city, to give them battle. And Joseph and Azarias were put to flight, and were pursued unto the borders of Judea: and there fell, on that day, of the people of Israel about two thousand men, and there was a great overthrow of the people: Because they did not hearken to Judas, and his brethren, thinking that they should do manfully. But they were not of the seed of those men by whom salvation was brought to Israel.\par
\switchcolumn*

\LA
\begin{Resp}\Rsign Ornavérunt fáciem templi corónis áureis, et dedicavérunt altáre Dómino: \Star{} Et facta est lætítia magna in pópulo.\end{Resp}
\begin{Resp}\Vsign In hymnis et confessiónibus benedicébant Dóminum.\end{Resp}
\begin{Resp}\Rsign Et facta est lætítia magna in pópulo.\end{Resp}
\switchcolumn
\EN
\begin{Resp}\Rsign They decked the fore-front of the Temple with crowns of gold, and dedicated the Altar unto the Lord. \Star{} And there was very great gladness among the people.\end{Resp}
\begin{Resp}\Vsign They praised the Lord with Psalms and thanksgiving.\end{Resp}
\begin{Resp}\Rsign And there was very great gladness among the people.\end{Resp}
\switchcolumn*

\LA
\LessonHead{Lectio III}
\switchcolumn
\EN
\LessonHead{Lesson III}
\switchcolumn*

\LA
\Cite{1 Mac 5:63-67}
\switchcolumn
\EN
\Cite{1 Mac 5:63-67}
\switchcolumn*

\LA
\noindent Et viri Juda magnificáti sunt valde in conspéctu omnis Israël, et géntium ómnium, ubi audiebátur nomen eórum. Et convenérunt ad eos fausta acclamántes. Et exívit Judas et fratres ejus, et expugnábant fílios Esau in terra quæ ad Austrum est, et percússit Chebron et fílias ejus: et muros ejus, et turres succéndit igni in circúitu, et movit castra, ut iret in terram alienigenárum, et perambulábat Samaríam. In die illa cecidérunt sacerdótes in bello, dum volunt fórtiter fácere, dum sine consílio éxeunt in prǽlium.\par
\switchcolumn
\EN
\noindent And the men of Juda were magnified exceedingly in the sight of all Israel, and of all the nations where their name was heard. And people assembled to them with joyful acclamations. Then Judas and his brethren went forth and attacked the children of Esau, in the land toward the south, and he took Chebron, and her towns: and he burnt the walls thereof and the towers all round it. And he removed his camp to go into the land of the aliens, and he went through Samaria. In that day some priests fell in battle, while desiring to do manfully they went out unadvisedly to fight.\par
\switchcolumn*

\LA
\begin{Resp}\Rsign In hymnis et confessiónibus benedicébant Dóminum, \Star{} Qui magna fecit in Israël, et victóriam dedit illis Dóminus omnípotens.\end{Resp}
\begin{Resp}\Vsign Ornavérunt fáciem templi corónis áureis, et dedicavérunt altáre Dómino.\end{Resp}
\begin{Resp}\Rsign Qui magna fecit in Israël, et victóriam dedit illis Dóminus omnípotens.\end{Resp}
\begin{Resp}\Vsign Glória Patri, et Fílio, \Star{} et Spirítui Sancto.\end{Resp}
\begin{Resp}\Rsign Qui magna fecit in Israël, et victóriam dedit illis Dóminus omnípotens.\end{Resp}
\switchcolumn
\EN
\begin{Resp}\Rsign They praised the Lord with psalms and thanksgiving \Star{} Who had done so great things in Israel, and given them the victory the Lord Almighty.\end{Resp}
\begin{Resp}\Vsign They decked the fore-front of the Temple with crowns of gold, and dedicated the Altar unto the Lord.\end{Resp}
\begin{Resp}\Rsign Who had done so great things for Israel, and given them the victory the Lord Almighty.\end{Resp}
\begin{Resp}\Vsign Glory be to the Father, and to the Son, \Star{} and to the Holy Ghost.\end{Resp}
\begin{Resp}\Rsign Who had done so great things for Israel, and given them the victory the Lord Almighty.\end{Resp}
\switchcolumn*

\end{paracol}

\Office{Feria V infra Hebdomadam II Octobris}{Thursday of the second week of October}{}
\addcontentsline{toc}{subsection}{Feria V infra Hebdomadam II Octobris \textit{— Thursday of the second week of October}}
\begin{paracol}{2}
\LA
\LessonHead{Lectio I}
\switchcolumn
\EN
\LessonHead{Lesson I}
\switchcolumn*

\LA
\LessonTitle{De libro primo Machabæórum}
\switchcolumn
\EN
\LessonTitle{Lesson from the first book of Machabees}
\switchcolumn*

\LA
\Cite{1 Mac 6:1-6}
\switchcolumn
\EN
\Cite{1 Mac 6:1-6}
\switchcolumn*

\LA
\noindent Et rex Antíochus perambulábat superióres regiónes, et audívit esse civitátem Elymáidem in Pérside nobilíssimam, et copiósam in argénto et auro, templúmque in ea lócuples valde et illic velámina áurea, et lorícæ, et scuta, quæ relíquit Alexánder Philíppi rex Mácedo, qui regnávit primus in Grǽcia. Et venit, et quærébat cápere civitátem, et deprædári eam: et non pótuit, quóniam innótuit sermo his qui erant in civitáte: et insurrexérunt in prǽlium, et fugit inde, et ábiit cum tristítia magna, et revérsus est in Babylóniam. Et venit qui nuntiáret ei in Pérside, quia fugáta sunt castra quæ erant in terra Juda: et quia ábiit Lýsias cum virtúte forti in primis, et fugátus est a fácie Judæórum.\par
\switchcolumn
\EN
\noindent Now king Antiochus was going through the higher countries, and he heard that the city of Elymais in Persia was greatly renowned, and abounding in silver and gold. And that there was in it a temple, exceeding rich: and coverings of gold, and breastplates, and shields which king Alexander, son of Philip the Macedonian that reigned first in Greece, had left there. Lo, he came, and sought to take the city and to pillage it: But he was not able, because the design was known to them that were in the city. And they rose up against him in battle, and he fled away from thence, and departed with great sadness, and returned towards Babylonia. And whilst he was in Persia, there came one that told him, how the armies that were in the land of Juda were put to flight: And that Lysias went with a very great power, and was put to flight before the face of the Jews,\par
\switchcolumn*

\LA
\begin{Resp}\Rsign Adapériat Dóminus cor vestrum in lege sua et in præcéptis suis et fáciat pacem in diébus vestris: \Star{} Concédat vobis salútem, et rédimat vos a malis.\end{Resp}
\begin{Resp}\Vsign Exáudiat Dóminus oratiónes vestras, et reconciliétur vobis, nec vos déserat in témpore malo.\end{Resp}
\begin{Resp}\Rsign Concédat vobis salútem, et rédimat vos a malis.\end{Resp}
\switchcolumn
\EN
\begin{Resp}\Rsign The Lord open your hearts in His law and commandments, and send peace in your days. \Star{} May He grant you salvation and redeem you out of all evil.\end{Resp}
\begin{Resp}\Vsign The Lord hear your prayers, and be at one with you, and never forsake you in the time of trouble.\end{Resp}
\begin{Resp}\Rsign May He grant you salvation and redeem you out of all evil.\end{Resp}
\switchcolumn*

\LA
\LessonHead{Lectio II}
\switchcolumn
\EN
\LessonHead{Lesson II}
\switchcolumn*

\LA
\Cite{1 Mac 6:6-9}
\switchcolumn
\EN
\Cite{1 Mac 6:6-9}
\switchcolumn*

\LA
\noindent Et invaluérunt armis, et víribus, et spóliis multis, quæ cepérunt de castris, quæ excidérunt: et quia diruérunt abominatiónem, quam ædificáverat super altáre, quod erat in Jerúsalem, et sanctificatiónem, sicut prius, circumdedérunt muris excélsis, sed et Bethsúram civitátem suam. Et factum est, ut audívit rex sermónes istos, expávit et commótus est valde: et décidit in lectum, et íncidit in languórem præ tristítia, quia non factum est ei sicut cogitábat. Et erat illic per dies multos, quia renováta est in eo tristítia magna, et arbitrátus est se mori.\par
\switchcolumn
\EN
\noindent And that thy were grown strong by the armour, and power, and store of spoils, which they had gotten out of the camps which they had destroyed: And that they had thrown down the abomination which he had set up upon the altar in Jerusalem, and that they had compassed about the sanctuary with high walls as before, and Bethsura also his city. And it came to pass when the king heard these words, that he was struck with fear, and exceedingly moved: and he laid himself down upon his bed, and fell sick for grief, because it had not fallen out to him as he imagined. And he remained there many days: for great grief came more and more and more upon him, and he made account that he should die.\par
\switchcolumn*

\LA
\begin{Resp}\Rsign Exáudiat Dóminus oratiónes vestras, et reconciliétur vobis, nec vos déserat in témpore malo, \Star{} Dóminus, Deus noster.\end{Resp}
\begin{Resp}\Vsign Det vobis cor ómnibus, ut colátis eum et faciátis ejus voluntátem.\end{Resp}
\begin{Resp}\Rsign Dóminus, Deus noster.\end{Resp}
\switchcolumn
\EN
\begin{Resp}\Rsign The Lord hear your prayers, and be at one with you, and never forsake you in the time of trouble \Star{} Even He, the Lord our God.\end{Resp}
\begin{Resp}\Vsign Give you all an heart to serve Him, and to do His will.\end{Resp}
\begin{Resp}\Rsign Even He, the Lord our God.\end{Resp}
\switchcolumn*

\LA
\LessonHead{Lectio III}
\switchcolumn
\EN
\LessonHead{Lesson III}
\switchcolumn*

\LA
\Cite{1 Mac 6:10-13}
\switchcolumn
\EN
\Cite{1 Mac 6:10-13}
\switchcolumn*

\LA
\noindent Et vocávit omnes amícos suos, et dixit illis: Recéssit somnus ab óculis meis, et cóncidi, et córrui corde præ sollicitúdine: et dixi in corde meo: In quantam tribulatiónem devéni, et in quos fluctus tristítiæ, in qua nunc sum, qui jucúndus eram, et diléctus in potestáte mea! Nunc vero reminíscor malórum, quæ feci in Jerúsalem, unde et ábstuli ómnia spólia áurea et argéntea quæ erant in ea, et misi auférre habitántes Judǽam sine causa. Cognóvi ergo quia proptérea invenérunt me mala ista: et ecce péreo tristítia magna in terra aliéna.\par
\switchcolumn
\EN
\noindent And he called for all his friends, and said to them: Sleep is gone from my eyes, and I am fallen away, and my heart is cast down for anxiety. And I said in my heart: Into how much tribulation am I come, and into what floods of sorrow, wherein now I am: I that was pleasant and beloved in my power! But now I remember the evils that I have done in Jerusalem, from whence also I took away all the spoils of gold, and of silver that were in it, and I sent to destroy the inhabitants of Juda without cause. I know therefore that for this cause these evils have found me: and behold I perish with great grief in a strange land.\par
\switchcolumn*

\LA
\begin{Resp}\Rsign Congregáti sunt inimíci nostri, et gloriántur in virtúte sua: cóntere fortitúdinem illórum, Dómine, et dispérge illos: \Star{} Ut cognóscant quia non est álius qui pugnet pro nobis, nisi tu, Deus noster.\end{Resp}
\begin{Resp}\Vsign Dispérge illos in virtúte tua, et déstrue eos, protéctor noster, Dómine.\end{Resp}
\begin{Resp}\Rsign Ut cognóscant quia non est álius qui pugnet pro nobis, nisi tu, Deus noster.\end{Resp}
\begin{Resp}\Vsign Glória Patri, et Fílio, \Star{} et Spirítui Sancto.\end{Resp}
\begin{Resp}\Rsign Ut cognóscant quia non est álius qui pugnet pro nobis, nisi tu, Deus noster.\end{Resp}
\switchcolumn
\EN
\begin{Resp}\Rsign Our enemies are gathered together, and make their boast of their own strength. O Lord, break their power, and scatter them \Star{} That they may know that there is none other that fighteth for us, but only Thou, O our God!\end{Resp}
\begin{Resp}\Vsign Scatter them in thy strength, and destroy them, O Lord our Shield.\end{Resp}
\begin{Resp}\Rsign That they may know that there is none other that fighteth for us, but only Thou, O our God!\end{Resp}
\begin{Resp}\Vsign Glory be to the Father, and to the Son, \Star{} and to the Holy Ghost.\end{Resp}
\begin{Resp}\Rsign That they may know that there is none other that fighteth for us, but only Thou, O our God!\end{Resp}
\switchcolumn*

\end{paracol}

\Office{Feria VI infra Hebdomadam II Octobris}{Friday of the second week of October}{}
\addcontentsline{toc}{subsection}{Feria VI infra Hebdomadam II Octobris \textit{— Friday of the second week of October}}
\begin{paracol}{2}
\LA
\LessonHead{Lectio I}
\switchcolumn
\EN
\LessonHead{Lesson I}
\switchcolumn*

\LA
\LessonTitle{De libro primo Machabæórum}
\switchcolumn
\EN
\LessonTitle{Lesson from the first book of Machabees}
\switchcolumn*

\LA
\Cite{1 Mac 7:1; 7:4-7}
\switchcolumn
\EN
\Cite{1 Mac 7:1; 7:4-7}
\switchcolumn*

\LA
\noindent Anno centésimo quinquagésimo primo, éxiit Demétrius Seléuci fílius ab urbe Roma, et ascéndit cum paucis viris in civitátem marítimam, et regnávit illic. Et sedit Demétrius super sedem regni sui. Et venérunt ad eum viri iníqui et ímpii ex Israël, et Alcimus dux eórum, qui volébat fíeri sacérdos, et accusavérunt pópulum apud regem, dicéntes: Pérdidit Judas et fratres ejus omnes amícos tuos, et nos dispérsit de terra nostra. Nunc ergo mitte virum, cui credis, ut eat, et vídeat extermínium omne quod fecit nobis, et regiónibus regis, et púniat omnes amícos ejus, et adjutóres eórum.\par
\switchcolumn
\EN
\noindent In the hundred and fifty-first year Demetrius the son of Seleucus departed from the city of Rome, and came up with a few men into a city of the sea coast, and reigned there. So the army slew them. And Demetrius sat upon the throne of his kingdom: And there came to him the wicked and ungodly men of Israel: And Alcimus was at the head of them, who desired to be made high priest. And they accused the people to the king, saying: Judas and his brethren have destroyed all thy friends, and he hath driven us out of our land. Now therefore send some man whom thou trustest, and let him go, and see all the havock he hath made amongst us, and in the king's lands: and let him punish all his friends and their helpers.\par
\switchcolumn*

\LA
\begin{Resp}\Rsign Impetum inimicórum ne timuéritis: mémores estóte, quómodo salvi facti sunt patres nostri: \Star{} Et nunc clamémus in cælum et miserébitur nostri Deus noster.\end{Resp}
\begin{Resp}\Vsign Mementóte mirabílium ejus, quæ fecit pharaóni et exercítui ejus in Mari Rubro.\end{Resp}
\begin{Resp}\Rsign Et nunc clamémus in cælum et miserébitur nostri Deus noster.\end{Resp}
\switchcolumn
\EN
\begin{Resp}\Rsign Be ye not afraid of the assault of the enemy; remember how our fathers were delivered. \Star{} Now, therefore, let us cry unto heaven, and our God will have mercy upon us.\end{Resp}
\begin{Resp}\Vsign Remember His marvellous works that He hath done unto Pharaoh and his host in the Red Sea.\end{Resp}
\begin{Resp}\Rsign Now, therefore, let us cry unto heaven, and our God will have mercy upon us.\end{Resp}
\switchcolumn*

\LA
\LessonHead{Lectio II}
\switchcolumn
\EN
\LessonHead{Lesson II}
\switchcolumn*

\LA
\Cite{1 Mac 7:8-11}
\switchcolumn
\EN
\Cite{1 Mac 7:8-11}
\switchcolumn*

\LA
\noindent Et elégit rex ex amícis suis Bácchidem, qui dominabátur trans flumen magnum in regno, et fidélem regi: Et misit eum, ut vidéret extermínium, quod fecit Judas; sed et Alcimum ímpium constítuit in sacerdótium, et mandávit ei fácere ultiónem in fílios Israël. Et surrexérunt, et venérunt cum exércitu magno in terram Juda: et misérunt núntios, et locúti sunt ad Judam et ad fratres ejus verbis pacíficis in dolo. Et non intendérunt sermónibus eórum: vidérunt enim quia venérunt cum exércitu magno.\par
\switchcolumn
\EN
\noindent Then the king chose Bacchides, one of his friends that ruled beyond the great river in the kingdom, and was faithful to the king: and he sent him, To see the havock that Judas had made: and the wicked Alcimus he made high priest, and commanded him to take revenge upon the children of Israel. And they arose, and came with a great army into the land of Juda: and they sent messengers, and spoke to Judas and his brethren with peaceable words deceitfully. But they gave no heed to their words: for they saw that they were come with a great army.\par
\switchcolumn*

\LA
\begin{Resp}\Rsign Congregátæ sunt gentes in multitúdine, ut dímicent contra nos, et ignorámus quid ágere debeámus: \Star{} Dómine Deus, ad te sunt óculi nostri, ne pereámus.\end{Resp}
\begin{Resp}\Vsign Tu scis quæ cógitant in nos: quómodo potérimus subsístere ante fáciem illórum, nisi tu ádjuves nos?\end{Resp}
\begin{Resp}\Rsign Dómine Deus, ad te sunt óculi nostri, ne pereámus.\end{Resp}
\switchcolumn
\EN
\begin{Resp}\Rsign The heathen are assembled together to fight against us, and we know not what we should do. \Star{} Our eyes look unto thee, O Lord our God, that we should not perish.\end{Resp}
\begin{Resp}\Vsign What things they imagine against us, Thou knowest. How shall we be able to stand against them, except Thou be our help.\end{Resp}
\begin{Resp}\Rsign Our eyes look unto thee, O Lord our God, that we should not perish.\end{Resp}
\switchcolumn*

\LA
\LessonHead{Lectio III}
\switchcolumn
\EN
\LessonHead{Lesson III}
\switchcolumn*

\LA
\Cite{1 Mac 7:12-17}
\switchcolumn
\EN
\Cite{1 Mac 7:12-17}
\switchcolumn*

\LA
\noindent Et convenérunt ad Alcimum et Bácchidem congregátio scribárum requírere quæ justa sunt: et primi, Assidǽi, qui erant in fíliis Israël, et exquirébant ab eis pacem. Dixérunt enim: Homo sacérdos de sémine Aaron venit: non decípiet nos. Et locútus est cum eis verba pacífica, et jurávit illis, dicens: Non inferémus vobis malum, neque amícis vestris. Et credidérunt ei: et comprehéndit ex eis sexagínta viros, et occídit eos in una die, secúndum verbum quod scriptum est: Carnes sanctórum tuórum, et sánguinem ipsórum effudérunt in circúitu Jerúsalem, et non erat qui sepelíret.\par
\switchcolumn
\EN
\noindent Then there assembled to Alcimus and Bacchides a company of the scribes to require things that are just: And first the Assideans that were among the children of Israel, and they sought peace of them. For they said: One that is a priest of the seed of Aaron is come, he will not deceive us. And he spoke to them peaceably: and he swore to them, saying: We will do you no harm nor your friends. And they believed him. And he took threescore of them, and slew them in one day, according to the word that is written: The flesh of thy saints, and the blood of them they have shed round about Jerusalem, and there was none to bury them.\par
\switchcolumn*

\LA
\begin{Resp}\Rsign Tua est poténtia, tuum regnum, Dómine: tu es super omnes gentes: \Star{} Da pacem, Dómine, in diébus nostris.\end{Resp}
\begin{Resp}\Vsign Creátor ómnium, Deus, terríbilis et fortis, justus et miséricors.\end{Resp}
\begin{Resp}\Rsign Da pacem, Dómine, in diébus nostris.\end{Resp}
\begin{Resp}\Vsign Glória Patri, et Fílio, \Star{} et Spirítui Sancto.\end{Resp}
\begin{Resp}\Rsign Da pacem, Dómine, in diébus nostris.\end{Resp}
\switchcolumn
\EN
\begin{Resp}\Rsign Thine, O Lord, is the power, thine is the kingdom, O Lord, and Thou art exalted above all the heathen. \Star{} Give peace in our time, O Lord.\end{Resp}
\begin{Resp}\Vsign O Lord God, Creator of all things, Who art fearful and strong, righteous and merciful.\end{Resp}
\begin{Resp}\Rsign Give peace in our time, O Lord.\end{Resp}
\begin{Resp}\Vsign Glory be to the Father, and to the Son, \Star{} and to the Holy Ghost.\end{Resp}
\begin{Resp}\Rsign Give peace in our time, O Lord.\end{Resp}
\switchcolumn*

\end{paracol}

\Office{Sabbato infra Hebdomadam II Octobris}{Saturday of the second week of October}{}
\addcontentsline{toc}{subsection}{Sabbato infra Hebdomadam II Octobris \textit{— Saturday of the second week of October}}
\begin{paracol}{2}
\LA
\LessonHead{Lectio I}
\switchcolumn
\EN
\LessonHead{Lesson I}
\switchcolumn*

\LA
\LessonTitle{De libro primo Machabæórum}
\switchcolumn
\EN
\LessonTitle{Lesson from the first book of Machabees}
\switchcolumn*

\LA
\Cite{1 Mac 8:1-4}
\switchcolumn
\EN
\Cite{1 Mac 8:1-4}
\switchcolumn*

\LA
\noindent Et audívit Judas nomen Romanórum, quia sunt poténtes víribus, et acquiéscunt ad ómnia quæ postulántur ab eis, et quicúmque accessérunt ad eos, statuérunt cum eis amicítias, et quia sunt poténtes víribus. Et audiérunt prǽlia eórum, et virtútes bonas, quas fecérunt in Galátia, quia obtinuérunt eos, et duxérunt sub tribútum: et quanta fecérunt in regióne Hispániæ, et quod in potestátem redegérunt metálla argénti et auri, quæ illic sunt, et possedérunt omnem locum consílio suo, et patiéntia: lócaque quæ longe erant valde ab eis, et reges, qui supervenérunt eis ab extrémis terræ, contrivérunt, et percussérunt eos plaga magna; céteri autem dant eis tribútum ómnibus annis.\par
\switchcolumn
\EN
\noindent Now Judas heard of the fame of the Romans, that they are powerful and strong, and willingly agree to all things that are requested of them: and that whosoever have come to them, they have made amity with them, and that they are mighty in power. And they heard of their battles, and their noble acts, which they had done in Galatia, how they conquered them, and brought them under tribute: And how great things they had done in the land of Spain, and that they had brought under their power the mines of silver and of gold that are there, and had gotten possession of all the place by their counsel and patience: And had conquered places that were very far off from them, and kings that came against them from the ends of the earth, and had overthrown them with great slaughter: and the rest pay them tribute every year.\par
\switchcolumn*

\LA
\begin{Resp}\Rsign Refúlsit sol in clípeos áureos, et resplenduérunt montes ab eis: \Star{} Et fortitúdo géntium dissipáta est.\end{Resp}
\begin{Resp}\Vsign Erat enim exércitus magnus valde et fortis: et appropiávit Judas et exércitus ejus in prǽlio.\end{Resp}
\begin{Resp}\Rsign Et fortitúdo géntium dissipáta est.\end{Resp}
\switchcolumn
\EN
\begin{Resp}\Rsign The sun shone upon the shields of gold, and the mountains glistered therewith; \Star{} And the army of the heathens was spread abroad.\end{Resp}
\begin{Resp}\Vsign For the army was very great and mighty then Judas and his host drew near and entered into battle.\end{Resp}
\begin{Resp}\Rsign And the army of the heathens was spread abroad.\end{Resp}
\switchcolumn*

\LA
\LessonHead{Lectio II}
\switchcolumn
\EN
\LessonHead{Lesson II}
\switchcolumn*

\LA
\Cite{1 Mac 8:17-22}
\switchcolumn
\EN
\Cite{1 Mac 8:17-22}
\switchcolumn*

\LA
\noindent Et elégit Judas Eupólemum fílium Joánnis, fílii Jacob, et Jásonem fílium Eleázari, et misit eos Romam constitúere cum illis amicítiam et societátem: et, ut auférrent ab eis jugum Græcórum, quia vidérunt quod in servitútem prémerent regnum Israël. Et abiérunt Romam viam multam valde, et introiérunt cúriam, et dixérunt: Judas Machabǽus, et fratres ejus, et pópulus Judæórum misérunt nos ad vos statúere vobíscum societátem et pacem, et conscríbere nos sócios et amícos vestros. Et plácuit sermo in conspéctu eórum. Et hoc rescríptum est, quod rescripsérunt in tábulis ǽreis, et misérunt in Jerúsalem, ut esset apud eos ibi memoriále pacis et societátis:\par
\switchcolumn
\EN
\noindent So Judas chose Eupolemus the son of John, the son of Jacob, and Jason the son of Eleazar, and he sent them to Rome to make a league of amity and confederacy with them. And that they might take off from them the yoke of the Grecians, for they saw that they oppressed the kingdom of Israel with servitude. And they went to Rome, a very long journey, and they entered into the senate house, and said: Judas Machabeus, and his brethren, and the people of the Jews have sent us to you, to make alliance and peace with you, and that we may be registered your confederates and friends. And the proposal was pleasing in their sight. And this is the copy of the writing that they wrote back again, graven in tables of brass, and sent to Jerusalem, that it might be with them there for a memorial of the peace and alliance.\par
\switchcolumn*

\LA
\begin{Resp}\Rsign Ornavérunt fáciem templi corónis áureis, et dedicavérunt altáre Dómino: \Star{} Et facta est lætítia magna in pópulo.\end{Resp}
\begin{Resp}\Vsign In hymnis et confessiónibus benedicébant Dóminum.\end{Resp}
\begin{Resp}\Rsign Et facta est lætítia magna in pópulo.\end{Resp}
\switchcolumn
\EN
\begin{Resp}\Rsign They decked the fore-front of the Temple with crowns of gold, and dedicated the Altar unto the Lord. \Star{} And there was very great gladness among the people.\end{Resp}
\begin{Resp}\Vsign They praised the Lord with Psalms and thanksgiving.\end{Resp}
\begin{Resp}\Rsign And there was very great gladness among the people.\end{Resp}
\switchcolumn*

\LA
\LessonHead{Lectio III}
\switchcolumn
\EN
\LessonHead{Lesson III}
\switchcolumn*

\LA
\Cite{1 Mac 8:23-27}
\switchcolumn
\EN
\Cite{1 Mac 8:23-27}
\switchcolumn*

\LA
\noindent Bene sit Románis, et genti Judæórum, in mari et in terra in ætérnum, gladiúsque et hostis procul sit ab eis. Quod si instíterit bellum Románis prius, aut ómnibus sóciis eórum in omni dominatióne eórum, auxílium feret gens Judæórum, prout tempus dictáverit, corde pleno; et præliántibus non dabunt, neque subministrábunt tríticum, arma, pecúniam, naves, sicut plácuit Románis; et custódient mandáta eórum, nihil ab eis accipiéntes. Simíliter autem et, si genti Judæórum prius accíderit bellum, adjuvábunt Románi ex ánimo, prout eis tempus permíserit.\par
\switchcolumn
\EN
\noindent Good success be to the Romans, and to the people of the Jews, by sea and by land for ever: and far be the sword and enemy from them. But if there come first any war upon the Romans, or any of their confederates, in all their dominions: The nation of the Jews shall help them according as the time shall direct, with all their heart: Neither shall they give them, whilst they are fighting, or furnish them with wheat, or arms, or money, or ships, as it hath seemed good to the Romans: and they shall obey their orders, without taking any thing of them. In like manner also if war shall come first upon the nation of the Jews, the Romans shall help them with all their heart, according as the time shall permit them.\par
\switchcolumn*

\LA
\begin{Resp}\Rsign In hymnis et confessiónibus benedicébant Dóminum, \Star{} Qui magna fecit in Israël, et victóriam dedit illis Dóminus omnípotens.\end{Resp}
\begin{Resp}\Vsign Ornavérunt fáciem templi corónis áureis, et dedicavérunt altáre Dómino.\end{Resp}
\begin{Resp}\Rsign Qui magna fecit in Israël, et victóriam dedit illis Dóminus omnípotens.\end{Resp}
\begin{Resp}\Vsign Glória Patri, et Fílio, \Star{} et Spirítui Sancto.\end{Resp}
\begin{Resp}\Rsign Qui magna fecit in Israël, et victóriam dedit illis Dóminus omnípotens.\end{Resp}
\switchcolumn
\EN
\begin{Resp}\Rsign They praised the Lord with psalms and thanksgiving \Star{} Who had done so great things in Israel, and given them the victory the Lord Almighty.\end{Resp}
\begin{Resp}\Vsign They decked the fore-front of the Temple with crowns of gold, and dedicated the Altar unto the Lord.\end{Resp}
\begin{Resp}\Rsign Who had done so great things for Israel, and given them the victory the Lord Almighty.\end{Resp}
\begin{Resp}\Vsign Glory be to the Father, and to the Son, \Star{} and to the Holy Ghost.\end{Resp}
\begin{Resp}\Rsign Who had done so great things for Israel, and given them the victory the Lord Almighty.\end{Resp}
\switchcolumn*

\end{paracol}

\Office{Dominica III Octobris}{Third Sunday of October}{}
\addcontentsline{toc}{subsection}{Dominica III Octobris \textit{— Third Sunday of October}}
\begin{paracol}{2}
\LA
\LessonHead{Lectio I}
\switchcolumn
\EN
\LessonHead{Lesson I}
\switchcolumn*

\LA
\LessonTitle{De libro primo Machabæórum}
\switchcolumn
\EN
\LessonTitle{Lesson from the first book of Machabees}
\switchcolumn*

\LA
\Cite{1 Mac 9:1-6}
\switchcolumn
\EN
\Cite{1 Mac 9:1-6}
\switchcolumn*

\LA
\noindent Intérea, ut audívit Demétrius quia cécidit Nicánor et exércitus ejus in prǽlio, appósuit Bácchidem et Alcimum rursum míttere in Judǽam, et dextrum cornu cum illis. Et abiérunt viam quæ ducit ad Gálgala et castra posuérunt in Másaloth, quæ est in Arbéllis: et occupavérunt eam, et peremérunt ánimas hóminum multas. In mense primo anni centésimi et quinquagésimi secúndi applicuérunt exércitum ad Jerúsalem: et surrexérunt et abiérunt in Béream vigínti míllia virórum, et duo míllia équitum. Et Judas posúerat castra in Láisa, et tria míllia viri elécti cum eo. Et vidérunt multitúdinem exércitus, quia multi sunt, et timuérunt valde, et multi subtraxérunt se de castris, et non remansérunt ex eis nisi octingénti viri.\par
\switchcolumn
\EN
\noindent In the mean time when Demetrius heard that Nicanor and his army were fallen in battle, he sent again Bacchides and Alcimus into Judea; and the right wing of his army with them. And they took the road that leadeth to Galgal, and they camped in Masaloth, which is in Arabella: and they made themselves masters of it, and slew many people. In the first month of the hundred and fifty-second year they brought the army to Jerusalem: And they arose, and went to Berea with twenty thousand men, and two thousand horsemen. Now Judas had pitched his tents in Laisa, and three thousand chosen men with him: And they saw the multitude of the army that they were many, and they were seized with great fear: and many withdrew themselves out of the camp, and there remained of them no more than eight hundred men.\par
\switchcolumn*

\LA
\begin{Resp}\Rsign Adapériat Dóminus cor vestrum in lege sua et in præcéptis suis et fáciat pacem in diébus vestris: \Star{} Concédat vobis salútem, et rédimat vos a malis.\end{Resp}
\begin{Resp}\Vsign Exáudiat Dóminus oratiónes vestras, et reconciliétur vobis, nec vos déserat in témpore malo.\end{Resp}
\begin{Resp}\Rsign Concédat vobis salútem, et rédimat vos a malis.\end{Resp}
\switchcolumn
\EN
\begin{Resp}\Rsign The Lord open your hearts in His law and commandments, and send peace in your days. \Star{} May He grant you salvation and redeem you out of all evil.\end{Resp}
\begin{Resp}\Vsign The Lord hear your prayers, and be at one with you, and never forsake you in the time of trouble.\end{Resp}
\begin{Resp}\Rsign May He grant you salvation and redeem you out of all evil.\end{Resp}
\switchcolumn*

\LA
\LessonHead{Lectio II}
\switchcolumn
\EN
\LessonHead{Lesson II}
\switchcolumn*

\LA
\Cite{1 Mac 9:7-11}
\switchcolumn
\EN
\Cite{1 Mac 9:7-11}
\switchcolumn*

\LA
\noindent Et vidit Judas quod deflúxit exércitus suus, et bellum perurgébat eum, et confráctus est corde, quia non habébat tempus congregándi eos, et dissolútus est. Et dixit his qui resídui erant: Surgámus et eámus ad adversários nostros, si potérimus pugnáre advérsus eos. Et avertébant eum, dicéntes: Non potérimus, sed liberémus ánimas nostras modo, et revertámur ad fratres nostros, et tunc pugnábimus advérsus eos; nos autem pauci sumus. Et ait Judas: Absit istam rem fácere, ut fugiámus ab eis: et, si appropiávit tempus nostrum, moriámur in virtúte propter fratres nostros, et non inferámus crimen glóriæ nostræ. Et movit exércitus de castris, et stetérunt illis óbviam, et divísi sunt équites in duas partes, et fundibulárii et sagittárii præíbant exércitum, et primi certáminis omnes poténtes.\par
\switchcolumn
\EN
\noindent And Judas saw that his army slipped away, and the battle pressed upon him, and his heart was cast down: because he had not time to gather them together, and he was discouraged. Then he said to them that remained: Let us arise, and go against our enemies, if we may be able to fight against them. But they dissuaded him, saying: We shall not be able, but let us save our lives now, and return to our brethren, and then we will fight against them: for we are but few. Then Judas said: God forbid we should do this thing, and flee away from them: but if our time be come, let us die manfully for our brethren, and let us not stain our glory. And the army removed out of the camp, and they stood over against them: and the horsemen were divided into two troops, and the slingers, and the archers went before the army, and they that were in the front were all men of valour.\par
\switchcolumn*

\LA
\begin{Resp}\Rsign Exáudiat Dóminus oratiónes vestras, et reconciliétur vobis, nec vos déserat in témpore malo, \Star{} Dóminus, Deus noster.\end{Resp}
\begin{Resp}\Vsign Det vobis cor ómnibus, ut colátis eum et faciátis ejus voluntátem.\end{Resp}
\begin{Resp}\Rsign Dóminus, Deus noster.\end{Resp}
\switchcolumn
\EN
\begin{Resp}\Rsign The Lord hear your prayers, and be at one with you, and never forsake you in the time of trouble \Star{} Even He, the Lord our God.\end{Resp}
\begin{Resp}\Vsign Give you all an heart to serve Him, and to do His will.\end{Resp}
\begin{Resp}\Rsign Even He, the Lord our God.\end{Resp}
\switchcolumn*

\LA
\LessonHead{Lectio III}
\switchcolumn
\EN
\LessonHead{Lesson III}
\switchcolumn*

\LA
\Cite{1 Mac 9:12-20}
\switchcolumn
\EN
\Cite{1 Mac 9:12-20}
\switchcolumn*

\LA
\noindent Bácchides autem erat in dextro cornu, et proximávit légio ex duábus pártibus, et clamábant tubis. Exclamavérunt autem et hi qui erant ex parte Judæ, étiam ipsi, et commóta est terra a voce exercítuum: et commíssum est prǽlium a mane usque ad vésperam. Et vidit Judas quod fírmior est pars exércitus Bácchidis in dextris, et convenérunt cum ipso omnes constántes corde; et contríta est déxtera pars ab eis, et persecútus est eos usque ad montem Azóti; et qui in sinístro cornu erant, vidérunt quod contrítum est dextrum cornu, et secúti sunt post Judam, et eos, qui cum ipso erant, a tergo: et ingravátum est prǽlium, et cecidérunt vulneráti multi ex his et ex illis, et Judas cécidit, et céteri fugérunt. Et Jónathas et Simon tulérunt Judam fratrem suum, et sepeliérunt eum in sepúlcro patrum suórum in civitáte Modin, et flevérunt eum omnis pópulus Israël planctu magno.\par
\switchcolumn
\EN
\noindent And Bacchides was in the right wing, and the legion drew near on two sides, and they sounded the trumpets: And they also were on Judas' side, even they also cried out, and the earth shook at the noise of the armies: and the battle was fought from morning even unto the evening. And Judas perceived that the stronger part of the army of Bacchides was on the right side, and all the stout of heart came together with him: And the right wing was discomfited by them, and he pursued them even to the mount Azotus. And they that were in the left wing saw that the right wing was discomfited, and they followed after Judas, and them that were with him, at their back: And the battle was hard fought, and there fell many wounded of the one side and of the other. And Judas was slain, and the rest fled away. And Jonathan and Simon took Judas their brother, and buried him in the sepulchre of their fathers in the city of Modin. And all the people of Israel bewailed him with great lamentation, and they mourned for him many days.\par
\switchcolumn*

\LA
\begin{Resp}\Rsign Congregáti sunt inimíci nostri, et gloriántur in virtúte sua: cóntere fortitúdinem illórum, Dómine, et dispérge illos: \Star{} Ut cognóscant quia non est álius qui pugnet pro nobis, nisi tu, Deus noster.\end{Resp}
\begin{Resp}\Vsign Dispérge illos in virtúte tua, et déstrue eos, protéctor noster, Dómine.\end{Resp}
\begin{Resp}\Rsign Ut cognóscant quia non est álius qui pugnet pro nobis, nisi tu, Deus noster.\end{Resp}
\begin{Resp}\Vsign Glória Patri, et Fílio, \Star{} et Spirítui Sancto.\end{Resp}
\begin{Resp}\Rsign Ut cognóscant quia non est álius qui pugnet pro nobis, nisi tu, Deus noster.\end{Resp}
\switchcolumn
\EN
\begin{Resp}\Rsign Our enemies are gathered together, and make their boast of their own strength. O Lord, break their power, and scatter them \Star{} That they may know that there is none other that fighteth for us, but only Thou, O our God!\end{Resp}
\begin{Resp}\Vsign Scatter them in thy strength, and destroy them, O Lord our Shield.\end{Resp}
\begin{Resp}\Rsign That they may know that there is none other that fighteth for us, but only Thou, O our God!\end{Resp}
\begin{Resp}\Vsign Glory be to the Father, and to the Son, \Star{} and to the Holy Ghost.\end{Resp}
\begin{Resp}\Rsign That they may know that there is none other that fighteth for us, but only Thou, O our God!\end{Resp}
\switchcolumn*

\LA
\LessonHead{Lectio IV}
\switchcolumn
\EN
\LessonHead{Lesson IV}
\switchcolumn*

\LA
\LessonTitle{Ex libro Officiórum sancti Ambrósii Epíscopi}
\switchcolumn
\EN
\LessonTitle{On the duties of the Clergy Ambrose bishop of Milan}
\switchcolumn*

\LA
\Source{Lib. 1. cap. 41}
\switchcolumn
\EN
\Source{Book 1. ch. 41}
\switchcolumn*

\LA
\noindent Quia fortitúdo non solum secúndis rebus, sed étiam advérsis probátur, spectémus Judæ Machabǽi éxitum. Is enim post victum Nicánorem, regis Demétrii ducem, secúrior advérsus vigínti míllia exércitus regis, cum octingéntis viris bellum adórsus, voléntibus his cédere, ne multitúdine opprimeréntur, gloriósam magis mortem quam turpem fugam suásit: Ne crimen, inquit, nostræ relinquámus glóriæ. Ita commísso prǽlio, cum a primo ortu diéi in vésperam dimicarétur, dextrum cornu, in quo validíssimam manum advértit hóstium, aggréssus fácile avértit. Sed dum fugiéntes séquitur, a tergo vúlneri locum prǽbuit; ítaque gloriosiórem triúmphis mortem invénit.\par
\switchcolumn
\EN
\noindent 209. But as fortitude is proved not only by prosperity but also in adversity, let us now consider the death of Judas Maccabæus. For he, after Nicanor, the general of King Demetrius, was defeated, boldly engaged 20,000 of the king's army with 800 men who were anxious to retire for fear of being overcome by so great a multitude, but whom he persuaded to endure a glorious death rather than to retire in disgraceful flight. Let us not leave, he says, any stain upon our glory. Thus, then, engaging in battle after having fought from sunrise till evening, he attacks and quickly drives back the right wing, where he sees the strongest troop of the enemy to be. But while pursuing the fugitives from the rear he gave a chance for a wound to be inflicted. (1 Macc 9:8) Thus he found the spot of death more full of glory for himself than any triumph.\par
\switchcolumn*

\LA
\begin{Resp}\Rsign Impetum inimicórum ne timuéritis: mémores estóte, quómodo salvi facti sunt patres nostri: \Star{} Et nunc clamémus in cælum et miserébitur nostri Deus noster.\end{Resp}
\begin{Resp}\Vsign Mementóte mirabílium ejus, quæ fecit pharaóni et exercítui ejus in Mari Rubro.\end{Resp}
\begin{Resp}\Rsign Et nunc clamémus in cælum et miserébitur nostri Deus noster.\end{Resp}
\switchcolumn
\EN
\begin{Resp}\Rsign Be ye not afraid of the assault of the enemy; remember how our fathers were delivered. \Star{} Now, therefore, let us cry unto heaven, and our God will have mercy upon us.\end{Resp}
\begin{Resp}\Vsign Remember His marvellous works that He hath done unto Pharaoh and his host in the Red Sea.\end{Resp}
\begin{Resp}\Rsign Now, therefore, let us cry unto heaven, and our God will have mercy upon us.\end{Resp}
\switchcolumn*

\LA
\LessonHead{Lectio V}
\switchcolumn
\EN
\LessonHead{Lesson V}
\switchcolumn*

\LA
\noindent Quid Jónatham fratrem ejus attéxam, qui cum parva manu advérsus exércitus régios pugnans, desértus a suis et cum duóbus tantum relíctus, reparávit bellum, avértit hostem, fugitántes suos ad societátem revocávit triúmphi? Habes fortitúdinem béllicam, in qua non medíocris honésti ac decóri forma est, quod mortem servitúti prǽferat ac turpitúdini. Quid autem de Mártyrum dicam passiónibus? Et ne lóngius evagémur, non minórem de supérbo rege Antíocho Machabǽi púeri revexérunt triúmphum, quam paréntes próprii; síquidem illi armáti, isti sine armis vicérunt.\par
\switchcolumn
\EN
\noindent Why need I further mention his brother Jonathan, who fought against the king's force, with but a small troop. (1 Macc 11:68) Though forsaken by his men, and left with only two, he retrieved the battle, drove back the enemy, and recalled his own men, who were flying in every direction, to share in his triumph. Here, then, is fortitude in war, which bears no light impress of what is virtuous and seemly upon it, for it prefers death to slavery and disgrace. But what am I to say of the sufferings of the martyrs? Not to go too far abroad, did not the children of Maccabæus gain triumphs over the proud King Antiochus, as great as those of their fathers? The latter in truth were armed, but they conquered without arms.\par
\switchcolumn*

\LA
\begin{Resp}\Rsign Congregátæ sunt gentes in multitúdine, ut dímicent contra nos, et ignorámus quid ágere debeámus: \Star{} Dómine Deus, ad te sunt óculi nostri, ne pereámus.\end{Resp}
\begin{Resp}\Vsign Tu scis quæ cógitant in nos: quómodo potérimus subsístere ante fáciem illórum, nisi tu ádjuves nos?\end{Resp}
\begin{Resp}\Rsign Dómine Deus, ad te sunt óculi nostri, ne pereámus.\end{Resp}
\switchcolumn
\EN
\begin{Resp}\Rsign The heathen are assembled together to fight against us, and we know not what we should do. \Star{} Our eyes look unto thee, O Lord our God, that we should not perish.\end{Resp}
\begin{Resp}\Vsign What things they imagine against us, Thou knowest. How shall we be able to stand against them, except Thou be our help.\end{Resp}
\begin{Resp}\Rsign Our eyes look unto thee, O Lord our God, that we should not perish.\end{Resp}
\switchcolumn*

\LA
\LessonHead{Lectio VI}
\switchcolumn
\EN
\LessonHead{Lesson VI}
\switchcolumn*

\LA
\noindent Stetit invícta septem puerórum cohors, régiis cincta legiónibus: defecérunt supplícia, cessérunt tortóres, non defecérunt Mártyres. Alius córium cápitis exútus, spéciem mutáverat, virtútem áuxerat. Alius linguam jussus amputándam prómere, respóndit: Non solos Dóminus audit loquéntes, qui audiébat Móysen tacéntem; plus audit tácitas cogitatiónes suórum, quam voces ómnium. Linguæ flagéllum times, flagéllum sánguinis non times? Habet et sanguis vocem suam, qua clamat ad Deum, sicut clamávit in Abel.\par
\switchcolumn
\EN
\noindent The company of the seven brothers stood unconquered, though surrounded by the legions of the king — tortures failed, tormentors ceased; but the martyrs failed not. One, having had the skin of his head pulled off, though changed in appearance, grew in courage. Another, bidden to put forth his tongue, so that it might be cut off, answered: The Lord hears not only those who speak, for He heard Moses when silent. He hears better the silent thoughts of His own than the voice of all others. Do you fear the scourge of my tongue— and do you not fear the scourge of blood spilt upon the ground? Blood, too, has a voice whereby it cries aloud to God— as it did in the case of Abel.\par
\switchcolumn*

\LA
\begin{Resp}\Rsign Tua est poténtia, tuum regnum, Dómine: tu es super omnes gentes: \Star{} Da pacem, Dómine, in diébus nostris.\end{Resp}
\begin{Resp}\Vsign Creátor ómnium, Deus, terríbilis et fortis, justus et miséricors.\end{Resp}
\begin{Resp}\Rsign Da pacem, Dómine, in diébus nostris.\end{Resp}
\begin{Resp}\Vsign Glória Patri, et Fílio, \Star{} et Spirítui Sancto.\end{Resp}
\begin{Resp}\Rsign Da pacem, Dómine, in diébus nostris.\end{Resp}
\switchcolumn
\EN
\begin{Resp}\Rsign Thine, O Lord, is the power, thine is the kingdom, O Lord, and Thou art exalted above all the heathen. \Star{} Give peace in our time, O Lord.\end{Resp}
\begin{Resp}\Vsign O Lord God, Creator of all things, Who art fearful and strong, righteous and merciful.\end{Resp}
\begin{Resp}\Rsign Give peace in our time, O Lord.\end{Resp}
\begin{Resp}\Vsign Glory be to the Father, and to the Son, \Star{} and to the Holy Ghost.\end{Resp}
\begin{Resp}\Rsign Give peace in our time, O Lord.\end{Resp}
\switchcolumn*

\end{paracol}

\Office{Feria II infra Hebdomadam III Octobris}{Monday of the third week of October}{}
\addcontentsline{toc}{subsection}{Feria II infra Hebdomadam III Octobris \textit{— Monday of the third week of October}}
\begin{paracol}{2}
\LA
\LessonHead{Lectio I}
\switchcolumn
\EN
\LessonHead{Lesson I}
\switchcolumn*

\LA
\LessonTitle{De libro primo Machabæórum}
\switchcolumn
\EN
\LessonTitle{Lesson from the first book of Machabees}
\switchcolumn*

\LA
\Cite{1 Mac 9:28-32}
\switchcolumn
\EN
\Cite{1 Mac 9:28-32}
\switchcolumn*

\LA
\noindent Et congregáti sunt omnes amíci Judæ et dixérunt Jónathæ: Ex quo frater tuus Judas defúnctus est, vir símilis ei non est, qui éxeat contra inimícos nostros, Bácchidem et eos qui inimíci sunt gentis nostræ. Nunc ítaque, te hódie elégimus esse pro eo nobis in príncipem, et ducem ad bellándum bellum nostrum. Et suscépit Jónathas témpore illo principátum et surréxit loco Judæ fratris sui. Et cognóvit Bácchides et quærébat eum occídere.\par
\switchcolumn
\EN
\noindent And all the friends of Judas came together, and said to Jonathan: Since thy brother Judas died, there is not a man like him to go forth against our enemies, Bacchides, and them that are the enemies of our nation. Now therefore we have chosen thee this day to be our prince, and captain in his stead to fight our battles. So Jonathan took upon him the government at that time, and rose up in the place of Judas his brother. And Bacchides had knowledge of it, and sought to kill him.\par
\switchcolumn*

\LA
\begin{Resp}\Rsign Dixit Judas Simóni fratri suo: Elige tibi viros et vade, líbera fratres tuos in Galilǽam; ego autem et Jónathas frater tuus íbimus in Galaadítim: \Star{} Sicut fúerit volúntas in cælo, sic fiat.\end{Resp}
\begin{Resp}\Vsign Accingímini, fílii poténtes, et estóte paráti: quóniam mélius est nobis, mori in bello, quam vidére mala gentis nostræ et sanctórum.\end{Resp}
\begin{Resp}\Rsign Sicut fúerit volúntas in cælo, sic fiat.\end{Resp}
\switchcolumn
\EN
\begin{Resp}\Rsign Judas said unto Simon his brother: Choose thee out men, and go, and deliver thy brethren that are in Galilee and I, and Jonathan thy brother, will go into the country of Galaad. \Star{} As the Will is in heaven, so let it be.\end{Resp}
\begin{Resp}\Vsign Arm yourselves, ye valiant men, and be in readiness for it is better for us to die in battle, than to behold the calamities of our people, and our sanctuary.\end{Resp}
\begin{Resp}\Rsign As the Will is in heaven, so let it be.\end{Resp}
\switchcolumn*

\LA
\LessonHead{Lectio II}
\switchcolumn
\EN
\LessonHead{Lesson II}
\switchcolumn*

\LA
\Cite{1 Mac 9:33-36}
\switchcolumn
\EN
\Cite{1 Mac 9:33-36}
\switchcolumn*

\LA
\noindent Et cognóvit Jónathas, et Simon frater ejus, et omnes qui cum eo erant, et fugérunt in desértum Thécuæ et consedérunt ad aquam lacus Asphar; et cognóvit Bácchides et die sabbatórum venit ipse et omnis exércitus ejus trans Jordánem. Et Jónathas misit fratrem suum ducem pópuli, et rogávit Nabuthǽos amícos suos, ut commodárent illis apparátum suum, qui erat copiósus. Et exiérunt fílii Jambri ex Mádaba, et comprehendérunt Joánnem et ómnia quæ habébat, et abiérunt habéntes ea.\par
\switchcolumn
\EN
\noindent And Jonathan and Simon his brother, knew it, and all that were with them: and they fled into the desert of Thecua, and they pitched by the water of the lake of Asphar, And Bacchides understood it, and he came himself with all his army over the Jordan on the sabbath day. And Jonathan sent his brother a captain of the people, to desire the Nabutheans his friends, that they would lend them their equipage, which was copious. And the children of Jambri came forth out of Madaba, and took John, and all that he had, and went away with them.\par
\switchcolumn*

\LA
\begin{Resp}\Rsign Ornavérunt fáciem templi corónis áureis, et dedicavérunt altáre Dómino: \Star{} Et facta est lætítia magna in pópulo.\end{Resp}
\begin{Resp}\Vsign In hymnis et confessiónibus benedicébant Dóminum.\end{Resp}
\begin{Resp}\Rsign Et facta est lætítia magna in pópulo.\end{Resp}
\switchcolumn
\EN
\begin{Resp}\Rsign They decked the fore-front of the Temple with crowns of gold, and dedicated the Altar unto the Lord. \Star{} And there was very great gladness among the people.\end{Resp}
\begin{Resp}\Vsign They praised the Lord with Psalms and thanksgiving.\end{Resp}
\begin{Resp}\Rsign And there was very great gladness among the people.\end{Resp}
\switchcolumn*

\LA
\LessonHead{Lectio III}
\switchcolumn
\EN
\LessonHead{Lesson III}
\switchcolumn*

\LA
\Cite{1 Mac 9:37-40}
\switchcolumn
\EN
\Cite{1 Mac 9:37-40}
\switchcolumn*

\LA
\noindent Post hæc verba renuntiátum est Jónathæ et Simóni fratri ejus, quia fílii Jambri fáciunt núptias magnas, et ducunt sponsam ex Mádaba fíliam uníus de magnis princípibus Chánaan cum ambitióne magna. Et recordáti sunt sánguinis Joánnis fratris sui: et ascendérunt, et abscondérunt se sub teguménto montis, et elevavérunt óculos suos, et vidérunt: et ecce tumúltus, et apparátus multus: et sponsus procéssit, et amíci ejus, et fratres ejus óbviam illis cum týmpanis, et músicis, et armis multis; et surrexérunt ad eos ex insídiis, et occidérunt eos, et cecidérunt vulneráti multi, et resídui fugérunt in montes, et accepérunt ómnia spólia eórum.\par
\switchcolumn
\EN
\noindent After this it was told Jonathan, and Simon his brother, that the children of Jambri made a great marriage, and were bringing the bride out of Madaba, the daughter of one of the great princes of Chanaan, with great pomp. And the remembered the blood of John their brother: and they went up, and hid themselves under the covert of the mountain. And they lifted up their eyes, and saw: and behold a tumult, and great preparation: and the bridegroom came forth, and his friends, and his brethren to meet them with timbrels, and musical instruments, and many weapons. And they rose up against them from the place where they lay in ambush, and slew them, and there fell many wounded, and the rest fled into the mountains, and they took all their spoils:\par
\switchcolumn*

\LA
\begin{Resp}\Rsign In hymnis et confessiónibus benedicébant Dóminum, \Star{} Qui magna fecit in Israël, et victóriam dedit illis Dóminus omnípotens.\end{Resp}
\begin{Resp}\Vsign Ornavérunt fáciem templi corónis áureis, et dedicavérunt altáre Dómino.\end{Resp}
\begin{Resp}\Rsign Qui magna fecit in Israël, et victóriam dedit illis Dóminus omnípotens.\end{Resp}
\begin{Resp}\Vsign Glória Patri, et Fílio, \Star{} et Spirítui Sancto.\end{Resp}
\begin{Resp}\Rsign Qui magna fecit in Israël, et victóriam dedit illis Dóminus omnípotens.\end{Resp}
\switchcolumn
\EN
\begin{Resp}\Rsign They praised the Lord with psalms and thanksgiving \Star{} Who had done so great things in Israel, and given them the victory the Lord Almighty.\end{Resp}
\begin{Resp}\Vsign They decked the fore-front of the Temple with crowns of gold, and dedicated the Altar unto the Lord.\end{Resp}
\begin{Resp}\Rsign Who had done so great things for Israel, and given them the victory the Lord Almighty.\end{Resp}
\begin{Resp}\Vsign Glory be to the Father, and to the Son, \Star{} and to the Holy Ghost.\end{Resp}
\begin{Resp}\Rsign Who had done so great things for Israel, and given them the victory the Lord Almighty.\end{Resp}
\switchcolumn*

\end{paracol}

\Office{Feria III infra Hebdomadam III Octobris}{Tuesday of the third week of October}{}
\addcontentsline{toc}{subsection}{Feria III infra Hebdomadam III Octobris \textit{— Tuesday of the third week of October}}
\begin{paracol}{2}
\LA
\LessonHead{Lectio I}
\switchcolumn
\EN
\LessonHead{Lesson I}
\switchcolumn*

\LA
\LessonTitle{De libro primo Machabæórum}
\switchcolumn
\EN
\LessonTitle{Lesson from the first book of Machabees}
\switchcolumn*

\LA
\Cite{1 Mac 12:1-4}
\switchcolumn
\EN
\Cite{1 Mac 12:1-4}
\switchcolumn*

\LA
\noindent Et vidit Jónathas quia tempus eum juvat, et elégit viros, et misit eos Romam statúere et renováre cum eis amicítiam; et ad Spartiátas, et ad ália loca misit epístolas secúndum eámdem formam. Et abiérunt Romam, et intravérunt cúriam, et dixérunt: Jónathas summus sacérdos, et gens Judæórum misérunt nos, ut renovarémus amicítiam et societátem secúndum prístinum. Et dedérunt illis epístolas ad ipsos per loca, ut dedúcerent eos in terram Juda cum pace.\par
\switchcolumn
\EN
\noindent And Jonathan saw that the time served him, and he chose certain men and sent them to Rome, to confirm and to renew the amity with them: And he sent letters to the Spartans, and to other places according to the same form. And they went to Rome, and entered into the senate house, and said: Jonathan the high priest, and the nation of the Jews have sent us to renew the amity, and alliance as it was before. And they gave them letters to their governors in every place, to conduct them into the land of Juda with peace.\par
\switchcolumn*

\LA
\begin{Resp}\Rsign Hic est fratrum amátor et pópuli Israël: \Star{} Hic est, qui multum orat pro pópulo et univérsa sancta civitáte Jerúsalem.\end{Resp}
\begin{Resp}\Vsign Vir iste in pópulo suo mitíssimus appáruit.\end{Resp}
\begin{Resp}\Rsign Hic est, qui multum orat pro pópulo et univérsa sancta civitáte Jerúsalem.\end{Resp}
\switchcolumn
\EN
\begin{Resp}\Rsign This is a lover of the brethren, and of the people of Israel \Star{} This is one who prayeth much for the people, and for all the Holy City, Jerusalem.\end{Resp}
\begin{Resp}\Vsign There appeared a man most gentle toward all his people.\end{Resp}
\begin{Resp}\Rsign This is one who prayeth much for the people, and for all the Holy City, Jerusalem.\end{Resp}
\switchcolumn*

\LA
\LessonHead{Lectio II}
\switchcolumn
\EN
\LessonHead{Lesson II}
\switchcolumn*

\LA
\Cite{1 Mac 12:5-8}
\switchcolumn
\EN
\Cite{1 Mac 12:5-8}
\switchcolumn*

\LA
\noindent Et hoc est exémplum epistolárum, quas scripsit Jónathas Spartiátis: Jónathas summus sacérdos, et senióres gentis, et sacerdótes, et réliquus pópulus Judæórum, Spartiátis frátribus salútem. Jamprídem missæ erant epístolæ ad Oníam summum sacerdótem ab Arío, qui regnábat apud vos, quóniam estis fratres nostri, sicut rescríptum cóntinet, quod subjéctum est. Et suscépit Onías virum, qui missus fúerat, cum honóre, et accépit epístolas, in quibus significabátur de societáte et amicítia.\par
\switchcolumn
\EN
\noindent And this is a copy of the letters which Jonathan wrote to the Spartans: Jonathan the high priest, and the ancients of the nation, and the priests, and the rest of the people of the Jews, to the Spartans, their brethren, greeting. There were letters sent long ago to Onias the high priest from Arius who reigned then among you, to signify that you are our brethren, as the copy here underwritten doth specify. And Onias received the ambassador with honour: and received the letters wherein there was mention made of the alliance, and amity.\par
\switchcolumn*

\LA
\begin{Resp}\Rsign Tu, Dómine universórum, qui nullam habes indigéntiam, voluísti templum tuum fíeri in nobis; \Star{} Consérva domum istam immaculátam in ætérnum, Dómine.\end{Resp}
\begin{Resp}\Vsign Tu elegísti, Dómine, domum istam ad invocándum nomen tuum in ea, ut esset domus oratiónis et obsecratiónis pópulo tuo.\end{Resp}
\begin{Resp}\Rsign Consérva domum istam immaculátam in ætérnum, Dómine.\end{Resp}
\switchcolumn
\EN
\begin{Resp}\Rsign Thou, O Lord of all things, Who hast need of nothing, wast pleased that the Temple of thine habitation should be among us. \Star{} Therefore now, O Lord, keep this house ever undefiled\end{Resp}
\begin{Resp}\Vsign Thou, O Lord, didst choose this house, that thy Name should be called on therein, and to be an house of prayer and petition for thy people.\end{Resp}
\begin{Resp}\Rsign Therefore now, O Lord, keep this house ever undefiled.\end{Resp}
\switchcolumn*

\LA
\LessonHead{Lectio III}
\switchcolumn
\EN
\LessonHead{Lesson III}
\switchcolumn*

\LA
\Cite{1 Mac 12:9-11}
\switchcolumn
\EN
\Cite{1 Mac 12:9-11}
\switchcolumn*

\LA
\noindent Nos, cum nullo horum indigerémus, habéntes solátio sanctos libros, qui sunt in mánibus nostris, malúimus míttere ad vos renováre fraternitátem et amicítiam, ne forte aliéni efficiámur a vobis; multa enim témpora transiérunt, ex quo misístis ad nos. Nos ergo in omni témpore sine intermissióne in diébus solémnibus, et céteris quibus opórtet, mémores sumus vestri in sacrifíciis quæ offérimus, et in observatiónibus, sicut fas est, et decet meminísse fratrum.\par
\switchcolumn
\EN
\noindent We, though we needed none of these things, having for our comfort the holy books that are in our hands, Chose rather to send to you to renew the brotherhood and friendship, lest we should become strangers to you altogether: for there is a long time passed since you sent to us. We therefore at all times without ceasing, both in our festivals, and other days, wherein it is convenient, remember you in the sacrifices that we offer, and in our observances, as it is meet, and becoming to remember brethren.\par
\switchcolumn*

\LA
\begin{Resp}\Rsign Aperi óculos tuos, Dómine, et vide afflictiónem nostram: circumdedérunt nos gentes ad puniéndum nos: \Star{} Sed tu, Dómine, exténde brácchium tuum, et líbera ánimas nostras.\end{Resp}
\begin{Resp}\Vsign Afflíge oppriméntes nos et contuméliam faciéntes in supérbiam; et custódi partem tuam.\end{Resp}
\begin{Resp}\Rsign Sed tu, Dómine, exténde brácchium tuum, et líbera ánimas nostras.\end{Resp}
\begin{Resp}\Vsign Glória Patri, et Fílio, \Star{} et Spirítui Sancto.\end{Resp}
\begin{Resp}\Rsign Sed tu, Dómine, exténde brácchium tuum, et líbera ánimas nostras.\end{Resp}
\switchcolumn
\EN
\begin{Resp}\Rsign Open thine eyes, O Lord, and behold our affliction for the heathen are come round about us to punish us. \Star{} But Thou, O Lord, stretch forth thine arm, and deliver our souls.\end{Resp}
\begin{Resp}\Vsign Punish them that oppress us and with pride do us wrong, and keep thine own portion.\end{Resp}
\begin{Resp}\Rsign But Thou, O Lord, stretch forth thine arm, and deliver our souls.\end{Resp}
\begin{Resp}\Vsign Glory be to the Father, and to the Son, \Star{} and to the Holy Ghost.\end{Resp}
\begin{Resp}\Rsign But Thou, O Lord, stretch forth thine arm, and deliver our souls.\end{Resp}
\switchcolumn*

\end{paracol}

\Office{Feria IV infra Hebdomadam III Octobris}{Wednesday of the third week of October}{}
\addcontentsline{toc}{subsection}{Feria IV infra Hebdomadam III Octobris \textit{— Wednesday of the third week of October}}
\begin{paracol}{2}
\LA
\LessonHead{Lectio I}
\switchcolumn
\EN
\LessonHead{Lesson I}
\switchcolumn*

\LA
\LessonTitle{De libro primo Machabæórum}
\switchcolumn
\EN
\LessonTitle{Lesson from the first book of Machabees}
\switchcolumn*

\LA
\Cite{1 Mac 12:39-43}
\switchcolumn
\EN
\Cite{1 Mac 12:39-43}
\switchcolumn*

\LA
\noindent Et, cum cogitásset Tryphon regnáre Asiæ, et assúmere diadéma, et exténdere manum in Antíochum regem: timens ne forte non permítteret eum Jónathas, sed pugnáret advérsus eum, quærébat comprehéndere eum, et occídere, et exsúrgens ábiit in Bethsan. Et exívit Jónathas óbviam illi cum quadragínta míllibus virórum electórum in prǽlium, et venit Bethsan. Et vidit Tryphon quia venit Jónathas cum exércitu multo, ut exténderet in eum manus: tímuit; et excépit eum cum honóre, et commendávit eum ómnibus amícis suis, et dedit ei múnera: et præcépit exercítibus suis ut obedírent ei, sicut sibi.\par
\switchcolumn
\EN
\noindent Now when Tryphon had conceived a design to make himself king of Asia, and to take the crown, and to stretch out his hand against king Antiochus: Fearing lest Jonathan would not suffer him, but would fight against him: he sought to seize upon him, and to kill him. So he rose up and came to Bethsan. And Jonathan went out to meet him with forty thousand men chosen for battle, and came to Bethsan. Now when Tryphon saw that Jonathan came with a great army, he durst not stretch forth his hand against him, But received him with honour, and commended him to all his friends, and gave him presents: and he commanded his troops to obey him, as himself.\par
\switchcolumn*

\LA
\begin{Resp}\Rsign Refúlsit sol in clípeos áureos, et resplenduérunt montes ab eis: \Star{} Et fortitúdo géntium dissipáta est.\end{Resp}
\begin{Resp}\Vsign Erat enim exércitus magnus valde et fortis: et appropiávit Judas et exércitus ejus in prǽlio.\end{Resp}
\begin{Resp}\Rsign Et fortitúdo géntium dissipáta est.\end{Resp}
\switchcolumn
\EN
\begin{Resp}\Rsign The sun shone upon the shields of gold, and the mountains glistered therewith; \Star{} And the army of the heathens was spread abroad.\end{Resp}
\begin{Resp}\Vsign For the army was very great and mighty then Judas and his host drew near and entered into battle.\end{Resp}
\begin{Resp}\Rsign And the army of the heathens was spread abroad.\end{Resp}
\switchcolumn*

\LA
\LessonHead{Lectio II}
\switchcolumn
\EN
\LessonHead{Lesson II}
\switchcolumn*

\LA
\Cite{1 Mac 12:44-47}
\switchcolumn
\EN
\Cite{1 Mac 12:44-47}
\switchcolumn*

\LA
\noindent Et dixit Jónathæ: Ut quid vexásti univérsum pópulum, cum bellum nobis non sit? Et nunc remítte eos in domos suas; élige autem tibi viros paucos, qui tecum sint, et veni mecum Ptolemáidam, et tradam eam tibi, et réliqua præsídia, et exércitum, et univérsos præpósitos negótii: et convérsus abíbo; proptérea enim veni. Et crédidit ei, et fecit sicut dixit: et dimísit exércitum, et abiérunt in terram Juda. Retínuit autem secum tria míllia virórum: ex quibus remísit in Galilǽam duo míllia; mille autem venérunt cum eo.\par
\switchcolumn
\EN
\noindent And he said to Jonathan: Why hast thou troubled all the people, whereas we have no war? Now therefore send them back to their own houses: and choose thee a few men that may be with thee, and come with me to Ptolemais, and I will deliver it to thee, and the rest of the strongholds, and the army, and all that have any charge, and I will return and go away: for this is the cause of my coming. And Jonathan believed him, and did as he said: and sent away his army, and they departed into the land of Juda: But he kept with him three thousand men: of whom he sent two thousand into Galilee, and one thousand went with him.\par
\switchcolumn*

\LA
\begin{Resp}\Rsign Ornavérunt fáciem templi corónis áureis, et dedicavérunt altáre Dómino: \Star{} Et facta est lætítia magna in pópulo.\end{Resp}
\begin{Resp}\Vsign In hymnis et confessiónibus benedicébant Dóminum.\end{Resp}
\begin{Resp}\Rsign Et facta est lætítia magna in pópulo.\end{Resp}
\switchcolumn
\EN
\begin{Resp}\Rsign They decked the fore-front of the Temple with crowns of gold, and dedicated the Altar unto the Lord. \Star{} And there was very great gladness among the people.\end{Resp}
\begin{Resp}\Vsign They praised the Lord with Psalms and thanksgiving.\end{Resp}
\begin{Resp}\Rsign And there was very great gladness among the people.\end{Resp}
\switchcolumn*

\LA
\LessonHead{Lectio III}
\switchcolumn
\EN
\LessonHead{Lesson III}
\switchcolumn*

\LA
\Cite{1 Mac 12:48-52}
\switchcolumn
\EN
\Cite{1 Mac 12:48-52}
\switchcolumn*

\LA
\noindent Ut autem intrávit Ptolemáidam Jónathas, clausérunt portas civitátis Ptoleménses, et comprehendérunt eum: et omnes qui cum eo intráverant, gládio interfecérunt. Et misit Tryphon exércitum et équites in Galilǽam et in campum magnum, ut pérderent omnes sócios Jónathæ. At illi, cum cognovíssent quia comprehénsus est Jónathas, et périit, et omnes qui cum eo erant, hortáti sunt semetípsos, et exiérunt paráti in prǽlium. Et vidéntes hi qui insecúti fúerant, quia pro ánima res est illis, revérsi sunt; illi autem venérunt omnes cum pace in terram Juda. Et planxérunt Jónathan et eos qui cum ipso fúerant valde.\par
\switchcolumn
\EN
\noindent Now as soon as Jonathan entered into Ptolemais, they of Ptolemais shut the gates of the city, and took him: and all them that came in with him they slew with the sword. Then Tryphon sent an army and horsemen into Galilee, and into the great plain to destroy all Jonathan's company. But they, when they understood that Jonathan and all that were with him were taken and slain, encouraged one another, and went out ready for battle. Then they that had come after them, seeing that they stood for their lives, returned back. Whereupon they all came peaceably into the land of Juda. And they bewailed Jonathan, and them that had been with him, exceedingly: and Israel mourned with great lamentation.\par
\switchcolumn*

\LA
\begin{Resp}\Rsign In hymnis et confessiónibus benedicébant Dóminum, \Star{} Qui magna fecit in Israël, et victóriam dedit illis Dóminus omnípotens.\end{Resp}
\begin{Resp}\Vsign Ornavérunt fáciem templi corónis áureis, et dedicavérunt altáre Dómino.\end{Resp}
\begin{Resp}\Rsign Qui magna fecit in Israël, et victóriam dedit illis Dóminus omnípotens.\end{Resp}
\begin{Resp}\Vsign Glória Patri, et Fílio, \Star{} et Spirítui Sancto.\end{Resp}
\begin{Resp}\Rsign Qui magna fecit in Israël, et victóriam dedit illis Dóminus omnípotens.\end{Resp}
\switchcolumn
\EN
\begin{Resp}\Rsign They praised the Lord with psalms and thanksgiving \Star{} Who had done so great things in Israel, and given them the victory the Lord Almighty.\end{Resp}
\begin{Resp}\Vsign They decked the fore-front of the Temple with crowns of gold, and dedicated the Altar unto the Lord.\end{Resp}
\begin{Resp}\Rsign Who had done so great things for Israel, and given them the victory the Lord Almighty.\end{Resp}
\begin{Resp}\Vsign Glory be to the Father, and to the Son, \Star{} and to the Holy Ghost.\end{Resp}
\begin{Resp}\Rsign Who had done so great things for Israel, and given them the victory the Lord Almighty.\end{Resp}
\switchcolumn*

\end{paracol}

\Office{Feria V infra Hebdomadam III Octobris}{Thursday of the third week of October}{}
\addcontentsline{toc}{subsection}{Feria V infra Hebdomadam III Octobris \textit{— Thursday of the third week of October}}
\begin{paracol}{2}
\LA
\LessonHead{Lectio I}
\switchcolumn
\EN
\LessonHead{Lesson I}
\switchcolumn*

\LA
\LessonTitle{De libro primo Machabæórum}
\switchcolumn
\EN
\LessonTitle{Lesson from the first book of Machabees}
\switchcolumn*

\LA
\Cite{1 Mac 13:1-6}
\switchcolumn
\EN
\Cite{1 Mac 13:1-6}
\switchcolumn*

\LA
\noindent Et audívit Simon quod congregávit Tryphon exércitum copiósum, ut veníret in terram Juda et attéreret eam. Videns quia in tremóre pópulus est, et in timóre, ascéndit Jerúsalem, et congregávit pópulum: et adhórtans, dixit: Vos scitis quanta ego, et fratres mei, et domus patris mei, fécimus pro légibus et pro sanctis, prǽlia, et angústias quales vídimus; horum grátia periérunt fratres mei omnes propter Israël, et relíctus sum ego solus. Et nunc non mihi contíngat párcere ánimæ meæ in omni témpore tribulatiónis; non enim mélior sum frátribus meis; vindicábo ítaque gentem meam, et sancta, natos quoque nostros, et uxóres, quia congregátæ sunt univérsæ gentes contérere nos inimicítiæ grátia.\par
\switchcolumn
\EN
\noindent Now Simon heard that Tryphon was gathering together a very great army, to invade the land of Juda, and to destroy it. And seeing that the people was in dread, and in fear, he went up to Jerusalem, and assembled the people: And exhorted them, saying: You know what great battles I and my brethren, and the house of my father, have fought for the laws, and the sanctuary, and the distresses that we have seen: By reason whereof all my brethren have lost their lives for Israel's sake, and I am left alone. And now far be it from me to spare my life in any time of trouble: for I am not better than my brethren. I will avenge then my nation and the sanctuary, and our children, and wives: for all the heathens are gathered together to destroy us out of mere malice.\par
\switchcolumn*

\LA
\begin{Resp}\Rsign Adapériat Dóminus cor vestrum in lege sua et in præcéptis suis et fáciat pacem in diébus vestris: \Star{} Concédat vobis salútem, et rédimat vos a malis.\end{Resp}
\begin{Resp}\Vsign Exáudiat Dóminus oratiónes vestras, et reconciliétur vobis, nec vos déserat in témpore malo.\end{Resp}
\begin{Resp}\Rsign Concédat vobis salútem, et rédimat vos a malis.\end{Resp}
\switchcolumn
\EN
\begin{Resp}\Rsign The Lord open your hearts in His law and commandments, and send peace in your days. \Star{} May He grant you salvation and redeem you out of all evil.\end{Resp}
\begin{Resp}\Vsign The Lord hear your prayers, and be at one with you, and never forsake you in the time of trouble.\end{Resp}
\begin{Resp}\Rsign May He grant you salvation and redeem you out of all evil.\end{Resp}
\switchcolumn*

\LA
\LessonHead{Lectio II}
\switchcolumn
\EN
\LessonHead{Lesson II}
\switchcolumn*

\LA
\Cite{1 Mac 13:7-13}
\switchcolumn
\EN
\Cite{1 Mac 13:7-13}
\switchcolumn*

\LA
\noindent Et accénsus est spíritus pópuli, simul ut audívit sermónes istos, et respondérunt voce magna, dicéntes: Tu es dux noster loco Judæ, et Jónathæ fratris tui; pugna prǽlium nostrum; et ómnia, quæcúmque díxeris nobis, faciémus. Et cóngregans omnes viros bellatóres, accelerávit consummáre univérsos muros Jerúsalem, et munívit eam in gyro. Et misit Jónathan fílium Absalómi, et cum eo exércitum novum in Joppen, et, ejéctis his qui erant in ea, remánsit illic ipse. Et movit Tryphon a Ptolemáida cum exércitu multo, ut veníret in terram Juda, et Jónathas cum eo in custódia; Simon autem applícuit in Addus contra fáciem campi.\par
\switchcolumn
\EN
\noindent And the spirit of the people was enkindled as soon as they heard these words. And they answered with a loud voice, saying: Thou art our leader in the place of Judas, and Jonathan thy brother. Fight thou our battles, and we will do whatsoever thou shalt say to us. So gathering together all the men of war, he made haste to finish all the walls of Jerusalem, and he fortified it round about. And he sent Jonathan the son of Absalom, and with him a new army into Joppe, and he cast out them that were in it, and himself remained there. And Tryphon removed from Ptolemais with a great army, to invade the land of Juda, and Jonathan was with him in custody. But Simon pitched in Addus, over against the plain.\par
\switchcolumn*

\LA
\begin{Resp}\Rsign Exáudiat Dóminus oratiónes vestras, et reconciliétur vobis, nec vos déserat in témpore malo, \Star{} Dóminus, Deus noster.\end{Resp}
\begin{Resp}\Vsign Det vobis cor ómnibus, ut colátis eum et faciátis ejus voluntátem.\end{Resp}
\begin{Resp}\Rsign Dóminus, Deus noster.\end{Resp}
\switchcolumn
\EN
\begin{Resp}\Rsign The Lord hear your prayers, and be at one with you, and never forsake you in the time of trouble \Star{} Even He, the Lord our God.\end{Resp}
\begin{Resp}\Vsign Give you all an heart to serve Him, and to do His will.\end{Resp}
\begin{Resp}\Rsign Even He, the Lord our God.\end{Resp}
\switchcolumn*

\LA
\LessonHead{Lectio III}
\switchcolumn
\EN
\LessonHead{Lesson III}
\switchcolumn*

\LA
\Cite{1 Mac 13:14-19}
\switchcolumn
\EN
\Cite{1 Mac 13:14-19}
\switchcolumn*

\LA
\noindent Et ut cognóvit Tryphon, quia surréxit Simon loco fratris sui Jónathæ, et quia commissúrus esset cum eo prǽlium, misit ad eum legátos, dicens: Pro argénto, quod debébat frater tuus Jónathas in ratióne regis propter negótia quæ hábuit, detinúimus eum. Et nunc mitte argénti talénta centum, et duos fílios ejus óbsides, ut non dimíssus fúgiat a nobis, et remittémus eum. Et cognóvit Simon quia cum dolo loquerétur secum: jussit tamen dari argéntum et púeros, ne inimicítiam magnam súmeret ad pópulum Israël, dicéntem: Quia non misit ei argéntum, et púeros, proptérea périit. Et misit púeros, et centum talénta: et mentítus est, et non dimísit Jónathas.\par
\switchcolumn
\EN
\noindent And when Tryphon understood that Simon was risen up in the place of his brother Jonathan, and that he meant to join battle with him, he sent messengers to him, Saying: We have detained thy brother Jonathan for the money that he owed in the king's account, by reason of the affairs which he had the management of. But now send a hundred talents of silver, and his two sons for hostages, that when he is set at liberty he may not revolt from us, and we will release him. Now Simon knew that he spoke deceitfully to him, nevertheless he ordered the money, and the children to be sent: lest he should bring upon himself a great hatred of the people of Israel, who might have said: Because he sent not the money, and the children, therefore is he lost. So he sent the children, and the hundred talents: and he lied, and did not let Jonathan go.\par
\switchcolumn*

\LA
\begin{Resp}\Rsign Congregáti sunt inimíci nostri, et gloriántur in virtúte sua: cóntere fortitúdinem illórum, Dómine, et dispérge illos: \Star{} Ut cognóscant quia non est álius qui pugnet pro nobis, nisi tu, Deus noster.\end{Resp}
\begin{Resp}\Vsign Dispérge illos in virtúte tua, et déstrue eos, protéctor noster, Dómine.\end{Resp}
\begin{Resp}\Rsign Ut cognóscant quia non est álius qui pugnet pro nobis, nisi tu, Deus noster.\end{Resp}
\begin{Resp}\Vsign Glória Patri, et Fílio, \Star{} et Spirítui Sancto.\end{Resp}
\begin{Resp}\Rsign Ut cognóscant quia non est álius qui pugnet pro nobis, nisi tu, Deus noster.\end{Resp}
\switchcolumn
\EN
\begin{Resp}\Rsign Our enemies are gathered together, and make their boast of their own strength. O Lord, break their power, and scatter them \Star{} That they may know that there is none other that fighteth for us, but only Thou, O our God!\end{Resp}
\begin{Resp}\Vsign Scatter them in thy strength, and destroy them, O Lord our Shield.\end{Resp}
\begin{Resp}\Rsign That they may know that there is none other that fighteth for us, but only Thou, O our God!\end{Resp}
\begin{Resp}\Vsign Glory be to the Father, and to the Son, \Star{} and to the Holy Ghost.\end{Resp}
\begin{Resp}\Rsign That they may know that there is none other that fighteth for us, but only Thou, O our God!\end{Resp}
\switchcolumn*

\end{paracol}

\Office{Feria VI infra Hebdomadam III Octobris}{Friday of the third week of October}{}
\addcontentsline{toc}{subsection}{Feria VI infra Hebdomadam III Octobris \textit{— Friday of the third week of October}}
\begin{paracol}{2}
\LA
\LessonHead{Lectio I}
\switchcolumn
\EN
\LessonHead{Lesson I}
\switchcolumn*

\LA
\LessonTitle{De libro primo Machabæórum}
\switchcolumn
\EN
\LessonTitle{Lesson from the first book of Machabees}
\switchcolumn*

\LA
\Cite{1 Mac 14:16-19}
\switchcolumn
\EN
\Cite{1 Mac 14:16-19}
\switchcolumn*

\LA
\noindent Et audítum est Romæ quia defúnctus esset Jónathas, et usque in Spartiátas: et contristáti sunt valde; ut audiérunt autem quod Simon frater ejus factus esset summus sacérdos loco ejus, et ipse obtinéret omnem regiónem, et civitátes in ea, scripsérunt ad eum in tábulis ǽreis, ut renovárent amicítias et societátem, quam fécerant cum Juda et cum Jónatha frátribus ejus, et lectæ sunt in conspéctu ecclésiæ in Jerúsalem. Et hoc exémplum epistolárum, quas Spartiátæ misérunt:\par
\switchcolumn
\EN
\noindent And it was heard at Rome, and as far as Sparta, that Jonathan was dead: and they were very sorry. But when they heard that Simon his brother was made high priest in his place, and was possessed of all the country, and the cities therein: They wrote to him in tables of brass, to renew the friendship and alliance which they had made with Judas, and with Jonathan his brethren. And they were read before the assembly in Jerusalem. And this is the copy of the letters that the Spartans sent.\par
\switchcolumn*

\LA
\begin{Resp}\Rsign Impetum inimicórum ne timuéritis: mémores estóte, quómodo salvi facti sunt patres nostri: \Star{} Et nunc clamémus in cælum et miserébitur nostri Deus noster.\end{Resp}
\begin{Resp}\Vsign Mementóte mirabílium ejus, quæ fecit pharaóni et exercítui ejus in Mari Rubro.\end{Resp}
\begin{Resp}\Rsign Et nunc clamémus in cælum et miserébitur nostri Deus noster.\end{Resp}
\switchcolumn
\EN
\begin{Resp}\Rsign Be ye not afraid of the assault of the enemy; remember how our fathers were delivered. \Star{} Now, therefore, let us cry unto heaven, and our God will have mercy upon us.\end{Resp}
\begin{Resp}\Vsign Remember His marvellous works that He hath done unto Pharaoh and his host in the Red Sea.\end{Resp}
\begin{Resp}\Rsign Now, therefore, let us cry unto heaven, and our God will have mercy upon us.\end{Resp}
\switchcolumn*

\LA
\LessonHead{Lectio II}
\switchcolumn
\EN
\LessonHead{Lesson II}
\switchcolumn*

\LA
\Cite{1 Mac 14:20-23}
\switchcolumn
\EN
\Cite{1 Mac 14:20-23}
\switchcolumn*

\LA
\noindent Spartianórum príncipes et civitátes, Simóni sacerdóti magno, et senióribus, et sacerdótibus, et réliquo pópulo Judæórum frátribus salútem. Legáti, qui missi sunt ad pópulum nostrum, nuntiavérunt nobis de vestra glória, et honóre, ac lætítia, et gavísi sumus in intróitu eórum, et scrípsimus quæ ab eis erant dicta in concíliis pópuli, sic: Numénius Antíochi, et Antípater Jásonis fílius, legáti Judæórum, venérunt ad nos, renovántes nobíscum amicítiam prístinam. Et plácuit pópulo excípere viros glorióse, et pónere exémplum sermónum eórum in segregátis pópuli libris, ut sit ad memóriam pópulo Spartiatárum. Exémplum autem horum scrípsimus Simóni magno sacerdóti.\par
\switchcolumn
\EN
\noindent The princes and the cities of the Spartans to Simon the high priest, and to the ancients, and the priests, and the rest of the people of the Jews their brethren, greeting. The ambassadors that were sent to our people, have told us of your glory, and honour, and joy: and we rejoice at their coming. And we registered what was said by them in the councils of the people in this manner: Numenius the son of Antiochus, and Antipater the son of Jason, ambassadors of the Jews, came to us to renew the former friendship with us. And it pleased the people to receive the men honourably, and to put a copy of their words in the public records, to be a memorial to the people of the Spartans. And we have written a copy of them to Simon the high priest.\par
\switchcolumn*

\LA
\begin{Resp}\Rsign Congregátæ sunt gentes in multitúdine, ut dímicent contra nos, et ignorámus quid ágere debeámus: \Star{} Dómine Deus, ad te sunt óculi nostri, ne pereámus.\end{Resp}
\begin{Resp}\Vsign Tu scis quæ cógitant in nos: quómodo potérimus subsístere ante fáciem illórum, nisi tu ádjuves nos?\end{Resp}
\begin{Resp}\Rsign Dómine Deus, ad te sunt óculi nostri, ne pereámus.\end{Resp}
\switchcolumn
\EN
\begin{Resp}\Rsign The heathen are assembled together to fight against us, and we know not what we should do. \Star{} Our eyes look unto thee, O Lord our God, that we should not perish.\end{Resp}
\begin{Resp}\Vsign What things they imagine against us, Thou knowest. How shall we be able to stand against them, except Thou be our help.\end{Resp}
\begin{Resp}\Rsign Our eyes look unto thee, O Lord our God, that we should not perish.\end{Resp}
\switchcolumn*

\LA
\LessonHead{Lectio III}
\switchcolumn
\EN
\LessonHead{Lesson III}
\switchcolumn*

\LA
\Cite{1 Mac 14:24-26}
\switchcolumn
\EN
\Cite{1 Mac 14:24-26}
\switchcolumn*

\LA
\noindent Post hæc autem misit Simon Numénium Romam, habéntem clípeum áureum magnum, pondo mnarum mille, ad statuéndam cum eis societátem. Cum autem audísset pópulus Románus sermónes istos, dixérunt: Quam gratiárum actiónem reddémus Simóni, et fíliis ejus? Restítuit enim ipse fratres suos, et expugnávit inimícos Israël ab eis. Et statuérunt ei libertátem, et descripsérunt in tábulis ǽreis, et posuérunt in títulis in monte Sion.\par
\switchcolumn
\EN
\noindent And after this Simon sent Numenius to Rome, with a great shield of gold the weight of a thousand pounds, to confirm the league with them. And when the people of Rome had heard These words, they said: What thanks shall we give to Simon, and his sons? For he hath restored his brethren, and hath driven away in fight the enemies of Israel from them: and they decreed him liberty, and registered it in tables of brass, and set it upon pillars in mount Sion.\par
\switchcolumn*

\LA
\begin{Resp}\Rsign Tua est poténtia, tuum regnum, Dómine: tu es super omnes gentes: \Star{} Da pacem, Dómine, in diébus nostris.\end{Resp}
\begin{Resp}\Vsign Creátor ómnium, Deus, terríbilis et fortis, justus et miséricors.\end{Resp}
\begin{Resp}\Rsign Da pacem, Dómine, in diébus nostris.\end{Resp}
\begin{Resp}\Vsign Glória Patri, et Fílio, \Star{} et Spirítui Sancto.\end{Resp}
\begin{Resp}\Rsign Da pacem, Dómine, in diébus nostris.\end{Resp}
\switchcolumn
\EN
\begin{Resp}\Rsign Thine, O Lord, is the power, thine is the kingdom, O Lord, and Thou art exalted above all the heathen. \Star{} Give peace in our time, O Lord.\end{Resp}
\begin{Resp}\Vsign O Lord God, Creator of all things, Who art fearful and strong, righteous and merciful.\end{Resp}
\begin{Resp}\Rsign Give peace in our time, O Lord.\end{Resp}
\begin{Resp}\Vsign Glory be to the Father, and to the Son, \Star{} and to the Holy Ghost.\end{Resp}
\begin{Resp}\Rsign Give peace in our time, O Lord.\end{Resp}
\switchcolumn*

\end{paracol}

\Office{Sabbato infra Hebdomadam III Octobris}{Saturday of the third week of October}{}
\addcontentsline{toc}{subsection}{Sabbato infra Hebdomadam III Octobris \textit{— Saturday of the third week of October}}
\begin{paracol}{2}
\LA
\LessonHead{Lectio I}
\switchcolumn
\EN
\LessonHead{Lesson I}
\switchcolumn*

\LA
\LessonTitle{De libro primo Machabæórum}
\switchcolumn
\EN
\LessonTitle{Lesson from the first book of Machabees}
\switchcolumn*

\LA
\Cite{1 Mac 16:14-17}
\switchcolumn
\EN
\Cite{1 Mac 16:14-17}
\switchcolumn*

\LA
\noindent Simon autem, perámbulans civitátes quæ erant in regióne Judǽæ, et sollicitúdinem gerens eárum, descéndit in Jéricho ipse, et Mathathías fílius ejus, et Judas, anno centésimo septuagésimo séptimo, mense undécimo: hic est mensis Sabath. Et suscépit eos fílius Abóbi in munitiúnculam, quæ vocátur Doch, cum dolo, quam ædificávit; et fecit eis convívium magnum, et abscóndit illic viros et, cum inebriátus esset Simon et fílii ejus, surréxit Ptolemǽus cum suis, et sumpsérunt arma sua, et intravérunt in convívium: et occidérunt eum, et duos fílios ejus, et quosdam púeros ejus: et fecit deceptiónem magnam in Israël, et réddidit mala pro bonis.\par
\switchcolumn
\EN
\noindent Now Simon, as he was going through the cities that were in the country of Judea, and taking care for the good ordering of them, went down to Jericho, he and Mathathias and Judas his sons, in the year one hundred and seventy-seven, the eleventh month: the same is the month Sabath. And the son of Abobus received them deceitfully into a little fortress, that is called Doch which he had built: and he made them a great feast, and hid men there. And when Simon and his sons had drunk plentifully, Ptolemee and his men rose up and took their weapons, and entered into the banqueting place, and slew him, and his two sons, and some of his servants. And he committed a great treachery in Israel, and rendered evil for good.\par
\switchcolumn*

\LA
\begin{Resp}\Rsign Refúlsit sol in clípeos áureos, et resplenduérunt montes ab eis: \Star{} Et fortitúdo géntium dissipáta est.\end{Resp}
\begin{Resp}\Vsign Erat enim exércitus magnus valde et fortis: et appropiávit Judas et exércitus ejus in prǽlio.\end{Resp}
\begin{Resp}\Rsign Et fortitúdo géntium dissipáta est.\end{Resp}
\switchcolumn
\EN
\begin{Resp}\Rsign The sun shone upon the shields of gold, and the mountains glistered therewith; \Star{} And the army of the heathens was spread abroad.\end{Resp}
\begin{Resp}\Vsign For the army was very great and mighty then Judas and his host drew near and entered into battle.\end{Resp}
\begin{Resp}\Rsign And the army of the heathens was spread abroad.\end{Resp}
\switchcolumn*

\LA
\LessonHead{Lectio II}
\switchcolumn
\EN
\LessonHead{Lesson II}
\switchcolumn*

\LA
\Cite{1 Mac 16:18-21}
\switchcolumn
\EN
\Cite{1 Mac 16:18-21}
\switchcolumn*

\LA
\noindent Et scripsit hæc Ptolemǽus, et misit regi, ut mítteret ei exércitum in auxílium, et tráderet ei regiónem, et civitátes eórum, et tribúta. Et misit álios in Gázaram tóllere Joánnem: et tribúnis misit epístolas, ut venírent ad se, et daret eis argéntum, et aurum, et dona. Et álios misit occupáre Jerúsalem et montem templi. Et præcúrrens quidam, nuntiávit Joánni in Gázara, quia périit pater ejus et fratres ejus, et quia: Misit te quoque intérfici.\par
\switchcolumn
\EN
\noindent And Ptolemee wrote these things and sent to the king that he should send him an army to aid him, and he would deliver him the country, and their cities, and tributes. And he sent others to Gazara to kill John: and to the tribunes he sent letters to come to him, and that he would give them silver, and gold, and gifts. And he sent others to take Jerusalem, and the mountain of the temple. Now one running before, told John in Gazara, that his father and his brethren were slain, and that he hath sent men to kill thee also.\par
\switchcolumn*

\LA
\begin{Resp}\Rsign Ornavérunt fáciem templi corónis áureis, et dedicavérunt altáre Dómino: \Star{} Et facta est lætítia magna in pópulo.\end{Resp}
\begin{Resp}\Vsign In hymnis et confessiónibus benedicébant Dóminum.\end{Resp}
\begin{Resp}\Rsign Et facta est lætítia magna in pópulo.\end{Resp}
\begin{Resp}\Vsign Glória Patri, et Fílio, \Star{} et Spirítui Sancto.\end{Resp}
\begin{Resp}\Rsign Et facta est lætítia magna in pópulo.\end{Resp}
\switchcolumn
\EN
\begin{Resp}\Rsign They decked the fore-front of the Temple with crowns of gold, and dedicated the Altar unto the Lord. \Star{} And there was very great gladness among the people.\end{Resp}
\begin{Resp}\Vsign They praised the Lord with Psalms and thanksgiving.\end{Resp}
\begin{Resp}\Rsign And there was very great gladness among the people.\end{Resp}
\begin{Resp}\Vsign Glory be to the Father, and to the Son, \Star{} and to the Holy Ghost.\end{Resp}
\begin{Resp}\Rsign And there was very great gladness among the people.\end{Resp}
\switchcolumn*

\LA
\LessonHead{Lectio III}
\switchcolumn
\EN
\LessonHead{Lesson III}
\switchcolumn*

\LA
\Cite{1 Mac 16:22-24}
\switchcolumn
\EN
\Cite{1 Mac 16:22-24}
\switchcolumn*

\LA
\noindent Ut audívit autem, veheménter expávit: et comprehéndit viros, qui vénerant pérdere eum, et occídit eos; cognóvit enim quia quærébant eum pérdere. Et cétera sermónum Joánnis, et bellórum ejus, et bonárum virtútum, quibus fórtiter gessit, et ædifícii murórum, quos exstrúxit, et rerum gestárum ejus, ecce hæc scripta sunt in libro diérum sacerdótii ejus, ex quo factus est princeps sacerdótum post patrem suum.\par
\switchcolumn
\EN
\noindent But when he heard it he was exceedingly afraid: and he apprehended the men that came to kill him, and he put them to death: for he knew that they sought to make him away. And as concerning the rest of the acts of John, and his wars, and the worthy deeds, which he bravely achieved, and the building of the walls, which he made, and the things that he did: Behold these are written in the book of the days of his priesthood, from the time he was made high priest after his father.\par
\switchcolumn*

\LA
\begin{Resp}\Rsign In hymnis et confessiónibus benedicébant Dóminum, \Star{} Qui magna fecit in Israël, et victóriam dedit illis Dóminus omnípotens.\end{Resp}
\begin{Resp}\Vsign Ornavérunt fáciem templi corónis áureis, et dedicavérunt altáre Dómino.\end{Resp}
\begin{Resp}\Rsign Qui magna fecit in Israël, et victóriam dedit illis Dóminus omnípotens.\end{Resp}
\begin{Resp}\Vsign Glória Patri, et Fílio, \Star{} et Spirítui Sancto.\end{Resp}
\begin{Resp}\Rsign Qui magna fecit in Israël, et victóriam dedit illis Dóminus omnípotens.\end{Resp}
\switchcolumn
\EN
\begin{Resp}\Rsign They praised the Lord with psalms and thanksgiving \Star{} Who had done so great things in Israel, and given them the victory the Lord Almighty.\end{Resp}
\begin{Resp}\Vsign They decked the fore-front of the Temple with crowns of gold, and dedicated the Altar unto the Lord.\end{Resp}
\begin{Resp}\Rsign Who had done so great things for Israel, and given them the victory the Lord Almighty.\end{Resp}
\begin{Resp}\Vsign Glory be to the Father, and to the Son, \Star{} and to the Holy Ghost.\end{Resp}
\begin{Resp}\Rsign Who had done so great things for Israel, and given them the victory the Lord Almighty.\end{Resp}
\switchcolumn*

\end{paracol}

\Office{Dominica IV Octobris}{Fourth Sunday of October}{}
\addcontentsline{toc}{subsection}{Dominica IV Octobris \textit{— Fourth Sunday of October}}
\begin{paracol}{2}
\LA
\LessonHead{Lectio I}
\switchcolumn
\EN
\LessonHead{Lesson I}
\switchcolumn*

\LA
\LessonTitle{Incipit liber secúndus Machabæórum}
\switchcolumn
\EN
\LessonTitle{Lesson from the second book of Machabees}
\switchcolumn*

\LA
\Cite{2 Mac 1:1-6}
\switchcolumn
\EN
\Cite{2 Mac 1:1-6}
\switchcolumn*

\LA
\noindent Frátribus qui sunt per Ægýptum Judǽis, salútem dicunt fratres qui sunt in Jerosólymis Judǽi, et qui in regióne Judǽæ, et pacem bonam. Benefáciat vobis Deus, et memínerit testaménti sui, quod locútus est ad Abraham, et Isaac, et Jacob servórum suórum fidélium: et det vobis cor ómnibus ut colátis eum, et faciátis ejus voluntátem, corde magno et ánimo volénti. Adapériat cor vestrum in lege sua, et in præcéptis suis, et fáciat pacem; exáudiat oratiónes vestras, et reconciliétur vobis, nec vos déserat in témpore malo. Et nunc hic sumus orántes pro vobis.\par
\switchcolumn
\EN
\noindent To the brethren the Jews that are throughout Egypt, the brethren, the Jews that are in Jerusalem, and in the land of Judea, send health, and good peace. May God be gracious to you, and remember his covenant that he made with Abraham, and Isaac, and Jacob, his faithful servants: And give you all a heart to worship him, and to do his will with a great heart, and a willing mind. May he open your heart in his law, and in his commandments, and send you peace. May he hear your prayers, and be reconciled unto you, and never forsake you in the evil time. And now here we are praying for you.\par
\switchcolumn*

\LA
\begin{Resp}\Rsign Adapériat Dóminus cor vestrum in lege sua et in præcéptis suis et fáciat pacem in diébus vestris: \Star{} Concédat vobis salútem, et rédimat vos a malis.\end{Resp}
\begin{Resp}\Vsign Exáudiat Dóminus oratiónes vestras, et reconciliétur vobis, nec vos déserat in témpore malo.\end{Resp}
\begin{Resp}\Rsign Concédat vobis salútem, et rédimat vos a malis.\end{Resp}
\switchcolumn
\EN
\begin{Resp}\Rsign The Lord open your hearts in His law and commandments, and send peace in your days. \Star{} May He grant you salvation and redeem you out of all evil.\end{Resp}
\begin{Resp}\Vsign The Lord hear your prayers, and be at one with you, and never forsake you in the time of trouble.\end{Resp}
\begin{Resp}\Rsign May He grant you salvation and redeem you out of all evil.\end{Resp}
\switchcolumn*

\LA
\LessonHead{Lectio II}
\switchcolumn
\EN
\LessonHead{Lesson II}
\switchcolumn*

\LA
\Cite{2 Mac 1:18-19}
\switchcolumn
\EN
\Cite{2 Mac 1:18-19}
\switchcolumn*

\LA
\noindent Factúri ígitur quinta et vigésima die mensis Cásleu, purificatiónem templi, necessárium dúximus significáre vobis: ut et vos quoque agátis diem scenopégiæ, et diem ignis, qui datus est quando Nehemías, ædificáto templo et altári, óbtulit sacrifícia. Nam, cum in Pérsidem duceréntur patres nostri, sacerdótes qui tunc cultóres Dei erant, accéptum ignem de altári occúlte abscondérunt in valle, ubi erat púteus altus et siccus, et in eo contutáti sunt eum, ita ut ómnibus ignótus esset locus.\par
\switchcolumn
\EN
\noindent Therefore whereas we purpose to keep the purification of the temple on the five and twentieth day of the month of Casleu, we thought it necessary to signify it to you: that you also may keep the day of Scenopegia, and the day of the fire, that was given when Nehemias offered sacrifice, after the temple and the altar was built. For when our fathers were led in Persia, the priests that then were worshippers of God took privately the fire from the altar, and hid it in a valley where there was a deep pit without water, and there they kept it safe, so that the place was unknown to all men.\par
\switchcolumn*

\LA
\begin{Resp}\Rsign Exáudiat Dóminus oratiónes vestras, et reconciliétur vobis, nec vos déserat in témpore malo, \Star{} Dóminus, Deus noster.\end{Resp}
\begin{Resp}\Vsign Det vobis cor ómnibus, ut colátis eum et faciátis ejus voluntátem.\end{Resp}
\begin{Resp}\Rsign Dóminus, Deus noster.\end{Resp}
\switchcolumn
\EN
\begin{Resp}\Rsign The Lord hear your prayers, and be at one with you, and never forsake you in the time of trouble \Star{} Even He, the Lord our God.\end{Resp}
\begin{Resp}\Vsign Give you all an heart to serve Him, and to do His will.\end{Resp}
\begin{Resp}\Rsign Even He, the Lord our God.\end{Resp}
\switchcolumn*

\LA
\LessonHead{Lectio III}
\switchcolumn
\EN
\LessonHead{Lesson III}
\switchcolumn*

\LA
\Cite{2 Mac 1:20-22}
\switchcolumn
\EN
\Cite{2 Mac 1:20-22}
\switchcolumn*

\LA
\noindent Cum autem præteríssent anni multi, et plácuit Deo ut mitterétur Nehemías a rege Pérsidis, nepótes sacerdótum illórum, qui abscónderant, misit ad requiréndum ignem: et, sicut narravérunt nobis, non invenérunt ignem, sed aquam crassam. Et jussit eos hauríre, et afférre sibi: et sacrifícia, quæ impósita erant, jussit sacérdos Nehemías aspérgi ipsa aqua: et ligna et quæ erant superpósita. Utque hoc factum est, et tempus áffuit quo sol refúlsit, quo prius erat in núbilo, accénsus est ignis magnus, ita ut omnes miraréntur.\par
\switchcolumn
\EN
\noindent But when many years had passed, and it pleased God that Nehemias should be sent by the king of Persia, he sent some of the posterity of those priests that had hid it, to seek for the fire: and as they told us, they found no fire, but thick water. Then he bade them draw it up, and bring it to him: and the priest Nehemias commanded the sacrifices that were laid on, to be sprinkled with the same water, both the wood, and the things that were laid upon it. And when this was done, and the time came that the sun shone out, which before was in a cloud, there was a great fire kindled, so that all wondered.\par
\switchcolumn*

\LA
\begin{Resp}\Rsign Congregáti sunt inimíci nostri, et gloriántur in virtúte sua: cóntere fortitúdinem illórum, Dómine, et dispérge illos: \Star{} Ut cognóscant quia non est álius qui pugnet pro nobis, nisi tu, Deus noster.\end{Resp}
\begin{Resp}\Vsign Dispérge illos in virtúte tua, et déstrue eos, protéctor noster, Dómine.\end{Resp}
\begin{Resp}\Rsign Ut cognóscant quia non est álius qui pugnet pro nobis, nisi tu, Deus noster.\end{Resp}
\begin{Resp}\Vsign Glória Patri, et Fílio, \Star{} et Spirítui Sancto.\end{Resp}
\begin{Resp}\Rsign Ut cognóscant quia non est álius qui pugnet pro nobis, nisi tu, Deus noster.\end{Resp}
\switchcolumn
\EN
\begin{Resp}\Rsign Our enemies are gathered together, and make their boast of their own strength. O Lord, break their power, and scatter them \Star{} That they may know that there is none other that fighteth for us, but only Thou, O our God!\end{Resp}
\begin{Resp}\Vsign Scatter them in thy strength, and destroy them, O Lord our Shield.\end{Resp}
\begin{Resp}\Rsign That they may know that there is none other that fighteth for us, but only Thou, O our God!\end{Resp}
\begin{Resp}\Vsign Glory be to the Father, and to the Son, \Star{} and to the Holy Ghost.\end{Resp}
\begin{Resp}\Rsign That they may know that there is none other that fighteth for us, but only Thou, O our God!\end{Resp}
\switchcolumn*

\LA
\LessonHead{Lectio IV}
\switchcolumn
\EN
\LessonHead{Lesson IV}
\switchcolumn*

\LA
\noindent Ex Tractátu sancti Joánnis Chrysóstomi super Psalmum quadragésimum tértium\par
\switchcolumn
\EN
\noindent We have heard with our ears, O God, our fathers have told us, what thou hast done in their time of old. The Prophet speaketh thus in the Psalm, yet not in his own person, but in the person of the Maccabees, relating and foretelling what events were to happen at the time. For such are the Prophets : they outrun all times, past, present, and future. But in order that our discussion of the subject may be more intelligible, we must first state who were these Maccabees, and what they suffered, and what they did. For when Antiochus, surnamed Epiphanes, had invaded Judaea, and laid everthing waste, and had forced many who then dwelt there to fall away from the laws of their fathers, the Maccabees remained unsullied by these temptations.\par
\switchcolumn*

\LA
\begin{Resp}\Rsign Impetum inimicórum ne timuéritis: mémores estóte, quómodo salvi facti sunt patres nostri: \Star{} Et nunc clamémus in cælum et miserébitur nostri Deus noster.\end{Resp}
\begin{Resp}\Vsign Mementóte mirabílium ejus, quæ fecit pharaóni et exercítui ejus in Mari Rubro.\end{Resp}
\begin{Resp}\Rsign Et nunc clamémus in cælum et miserébitur nostri Deus noster.\end{Resp}
\switchcolumn
\EN
\begin{Resp}\Rsign Be ye not afraid of the assault of the enemy; remember how our fathers were delivered. \Star{} Now, therefore, let us cry unto heaven, and our God will have mercy upon us.\end{Resp}
\begin{Resp}\Vsign Remember His marvellous works that He hath done unto Pharaoh and his host in the Red Sea.\end{Resp}
\begin{Resp}\Rsign Now, therefore, let us cry unto heaven, and our God will have mercy upon us.\end{Resp}
\switchcolumn*

\LA
\LessonHead{Lectio V}
\switchcolumn
\EN
\LessonHead{Lesson V}
\switchcolumn*

\LA
\noindent Et quando grave quidem bellum ingruébat, nec quidquam possent fácere quod prodésset, se abscondébant; nam hoc quoque fecérunt Apóstoli. Non enim semper apparéntes in média irruébant perícula, sed nonnúmquam et fugiéntes, et laténtes secedébant. Postquam autem parum respirárunt, tamquam generósi quidam cátuli ex antris exsiliéntes et e látebris emergéntes, statuérunt non se ámplius solos serváre, sed étiam álios quoscúmque possent: et civitátem et omnem regiónem obeúntes, collegérunt quotquot invenérunt adhuc sanos et íntegros; et multos étiam qui laborábant et corrúpti erant, in statum prístinum redegérunt, eis persuadéntes redíre ad legem pátriam.\par
\switchcolumn
\EN
\noindent And when a serious war broke out, and they could do nothing to help themselves, they hid themselves, as also in aftertimes the Apostles did. For they did not always rush openly into the midst of dangers, but sometimes fled, withdrawing thus to hide. However, after one such short respite, they were like eager animals leaping out of their caves and coming forth from their lairs, and they thereupon resolved for the future, not to win safety for themselves only, but for others, whomsoever they could. And going through all that city and country, they gathered together as many as they found who were still healthy and steadfast; and even many who were weak, and had been corrupted, they persuaded to return to the Law of their fathers.\par
\switchcolumn*

\LA
\begin{Resp}\Rsign Congregátæ sunt gentes in multitúdine, ut dímicent contra nos, et ignorámus quid ágere debeámus: \Star{} Dómine Deus, ad te sunt óculi nostri, ne pereámus.\end{Resp}
\begin{Resp}\Vsign Tu scis quæ cógitant in nos: quómodo potérimus subsístere ante fáciem illórum, nisi tu ádjuves nos?\end{Resp}
\begin{Resp}\Rsign Dómine Deus, ad te sunt óculi nostri, ne pereámus.\end{Resp}
\switchcolumn
\EN
\begin{Resp}\Rsign The heathen are assembled together to fight against us, and we know not what we should do. \Star{} Our eyes look unto thee, O Lord our God, that we should not perish.\end{Resp}
\begin{Resp}\Vsign What things they imagine against us, Thou knowest. How shall we be able to stand against them, except Thou be our help.\end{Resp}
\begin{Resp}\Rsign Our eyes look unto thee, O Lord our God, that we should not perish.\end{Resp}
\switchcolumn*

\LA
\LessonHead{Lectio VI}
\switchcolumn
\EN
\LessonHead{Lesson VI}
\switchcolumn*

\LA
\noindent Deum enim dicébant esse benígnum et cleméntem, nec umquam adímere salútem, quæ proficíscitur ex pœniténtia. Hæc autem dicéntes, habuérunt deléctum fortissimórum virórum. Non enim pro uxóribus, líberis, et ancíllis, patriǽque eversióne et captivitáte, sed pro lege et pátria república pugnábant. Eórum autem dux erat Deus. Cum ergo áciem dirígerent, et suas ánimas prodígerent, fundébant adversários, non armis fidéntes, sed loco omnis armatúræ, pugnæ causam suffícere ducéntes. Ad bellum autem eúntes non tragœ́dias excitábant, non pæána canébant, sicut nonnúlli fáciunt: non ascivérunt tibícines, ut fit in áliis castris: sed Dei supérne auxílium invocábant, ut adésset, opem ferret et manum præbéret, propter quem bellum gerébant, pro cujus glória decertábant.\par
\switchcolumn
\EN
\noindent For they told them that God is merciful and gracious, and that he hath never deprived men of that salvation which is obtained by penítence. And, so saying, they raised a levy of the most valiant men. For they fought not for their wives, their children and servants, or because of the ruin and captivity of their fatherland, but for the Law, and the religion of their fathers. Now their leader was God. Therefore, when they arrayed their battle line, and put their lives in jeopardy, they overthrew their adversaries because they trusted not in arms, but considered that the just cause of their war was in itself a good armour. Moreover, when they went forth to the conflict, they uttered no bombast, nor sang battle songs, as some do ; nor did they call together musicians, as is done in other armies; but they invoked the help of the Most High God, that he might be with them, and aid them, and strengthen their hand, because that war which they waged was for his glory.\par
\switchcolumn*

\LA
\begin{Resp}\Rsign Tua est poténtia, tuum regnum, Dómine: tu es super omnes gentes: \Star{} Da pacem, Dómine, in diébus nostris.\end{Resp}
\begin{Resp}\Vsign Creátor ómnium, Deus, terríbilis et fortis, justus et miséricors.\end{Resp}
\begin{Resp}\Rsign Da pacem, Dómine, in diébus nostris.\end{Resp}
\begin{Resp}\Vsign Glória Patri, et Fílio, \Star{} et Spirítui Sancto.\end{Resp}
\begin{Resp}\Rsign Da pacem, Dómine, in diébus nostris.\end{Resp}
\switchcolumn
\EN
\begin{Resp}\Rsign Thine, O Lord, is the power, thine is the kingdom, O Lord, and Thou art exalted above all the heathen. \Star{} Give peace in our time, O Lord.\end{Resp}
\begin{Resp}\Vsign O Lord God, Creator of all things, Who art fearful and strong, righteous and merciful.\end{Resp}
\begin{Resp}\Rsign Give peace in our time, O Lord.\end{Resp}
\begin{Resp}\Vsign Glory be to the Father, and to the Son, \Star{} and to the Holy Ghost.\end{Resp}
\begin{Resp}\Rsign Give peace in our time, O Lord.\end{Resp}
\switchcolumn*

\end{paracol}

\Office{Feria II infra Hebdomadam IV Octobris}{Monday of the fourth week of October}{}
\addcontentsline{toc}{subsection}{Feria II infra Hebdomadam IV Octobris \textit{— Monday of the fourth week of October}}
\begin{paracol}{2}
\LA
\LessonHead{Lectio I}
\switchcolumn
\EN
\LessonHead{Lesson I}
\switchcolumn*

\LA
\LessonTitle{De libro secúndo Machabæórum}
\switchcolumn
\EN
\LessonTitle{Lesson from the second book of Machabees}
\switchcolumn*

\LA
\Cite{2 Mac 2:1-3}
\switchcolumn
\EN
\Cite{2 Mac 2:1-3}
\switchcolumn*

\LA
\noindent Invenítur autem in descriptiónibus Jeremíæ prophétæ, quod jussit eos ignem accípere qui transmigrábant, ut significátum est, et ut mandávit transmigrátis, et dedit illis legem, ne oblivisceréntur præcépta Dómini, et ut non exerrárent méntibus, vidéntes simulácra áurea et argéntea, et ornaménta eórum; et ália hujúsmodi dicens, hortabátur ne legem amovérent a corde suo.\par
\switchcolumn
\EN
\noindent Now it is found in the descriptions of Jeremias the prophet, that he commanded them that went into captivity, to take the fire, as it hath been signified, and how he gave charge to them that were carried away into captivity. And how he gave them the law that they should not forget the commandments of the Lord, and that they should not err in their minds, seeing the idols of gold, and silver, and the ornaments of them. And with other such like speeches, he exhorted them that they would not remove the law from their heart.\par
\switchcolumn*

\LA
\begin{Resp}\Rsign Dixit Judas Simóni fratri suo: Elige tibi viros et vade, líbera fratres tuos in Galilǽam; ego autem et Jónathas frater tuus íbimus in Galaadítim: \Star{} Sicut fúerit volúntas in cælo, sic fiat.\end{Resp}
\begin{Resp}\Vsign Accingímini, fílii poténtes, et estóte paráti: quóniam mélius est nobis, mori in bello, quam vidére mala gentis nostræ et sanctórum.\end{Resp}
\begin{Resp}\Rsign Sicut fúerit volúntas in cælo, sic fiat.\end{Resp}
\switchcolumn
\EN
\begin{Resp}\Rsign Judas said unto Simon his brother: Choose thee out men, and go, and deliver thy brethren that are in Galilee and I, and Jonathan thy brother, will go into the country of Galaad. \Star{} As the Will is in heaven, so let it be.\end{Resp}
\begin{Resp}\Vsign Arm yourselves, ye valiant men, and be in readiness for it is better for us to die in battle, than to behold the calamities of our people, and our sanctuary.\end{Resp}
\begin{Resp}\Rsign As the Will is in heaven, so let it be.\end{Resp}
\switchcolumn*

\LA
\LessonHead{Lectio II}
\switchcolumn
\EN
\LessonHead{Lesson II}
\switchcolumn*

\LA
\Cite{2 Mac 2:4-6}
\switchcolumn
\EN
\Cite{2 Mac 2:4-6}
\switchcolumn*

\LA
\noindent Erat autem in ipsa scriptúra, quómodo tabernáculum et arcam jussit prophéta, divíno respónso ad se facto, comitári secum, úsquequo éxiit in montem, in quo Móyses ascéndit, et vidit Dei hereditátem. Et véniens ibi Jeremías, invénit locum spelúncæ: et tabernáculum, et arcam, et altáre incénsi íntulit illuc, et óstium obstrúxit. Et accessérunt quidam simul, qui sequebántur, ut notárent sibi locum: et non potuérunt inveníre.\par
\switchcolumn
\EN
\noindent It was also contained in the same writing, how the prophet, being warned by God, commanded that the tabernacle and the ark should accompany him, till he came forth to the mountain where Moses went up, and saw the inheritance of God. And when Jeremias came thither he found a hollow cave: and he carried in thither the tabernacle, and the ark, and the altar of incense, and so stopped the door. Then some of them that followed him, came up to mark the place: but they could not find it.\par
\switchcolumn*

\LA
\begin{Resp}\Rsign Ornavérunt fáciem templi corónis áureis, et dedicavérunt altáre Dómino: \Star{} Et facta est lætítia magna in pópulo.\end{Resp}
\begin{Resp}\Vsign In hymnis et confessiónibus benedicébant Dóminum.\end{Resp}
\begin{Resp}\Rsign Et facta est lætítia magna in pópulo.\end{Resp}
\switchcolumn
\EN
\begin{Resp}\Rsign They decked the fore-front of the Temple with crowns of gold, and dedicated the Altar unto the Lord. \Star{} And there was very great gladness among the people.\end{Resp}
\begin{Resp}\Vsign They praised the Lord with Psalms and thanksgiving.\end{Resp}
\begin{Resp}\Rsign And there was very great gladness among the people.\end{Resp}
\switchcolumn*

\LA
\LessonHead{Lectio III}
\switchcolumn
\EN
\LessonHead{Lesson III}
\switchcolumn*

\LA
\Cite{2 Mac 2:7-9}
\switchcolumn
\EN
\Cite{2 Mac 2:7-9}
\switchcolumn*

\LA
\noindent Ut autem cognóvit Jeremías, culpans illos dixit: quod ignótus erit locus, donec cóngreget Deus congregatiónem pópuli, et propítius fiat; et tunc Dóminus osténdet hæc, et apparébit majéstas Dómini, et nubes erit, sicut et Móysi manifestabátur, et sicut, cum Sálomon pétiit ut locus sanctificarétur magno Deo, manifestábat hæc; magnífice étenim sapiéntiam tractábat: et ut sapiéntiam habens, óbtulit sacrifícium dedicatiónis et consummatiónis templi.\par
\switchcolumn
\EN
\noindent And when Jeremias perceived it, he blamed them, saying: The place shall be unknown, till God gather together the congregation of the people, and receive them to mercy. And then the Lord will shew these things, and the majesty of the Lord shall appear, and there shall be a cloud as it was also shewed to Moses, and he shewed it when Solomon prayed that the place might be sanctified to the great God. For he treated wisdom in a magnificent manner: and like a wise man, he offered the sacrifice of the dedication, and of the finishing of the temple.\par
\switchcolumn*

\LA
\begin{Resp}\Rsign In hymnis et confessiónibus benedicébant Dóminum, \Star{} Qui magna fecit in Israël, et victóriam dedit illis Dóminus omnípotens.\end{Resp}
\begin{Resp}\Vsign Ornavérunt fáciem templi corónis áureis, et dedicavérunt altáre Dómino.\end{Resp}
\begin{Resp}\Rsign Qui magna fecit in Israël, et victóriam dedit illis Dóminus omnípotens.\end{Resp}
\begin{Resp}\Vsign Glória Patri, et Fílio, \Star{} et Spirítui Sancto.\end{Resp}
\begin{Resp}\Rsign Qui magna fecit in Israël, et victóriam dedit illis Dóminus omnípotens.\end{Resp}
\switchcolumn
\EN
\begin{Resp}\Rsign They praised the Lord with psalms and thanksgiving \Star{} Who had done so great things in Israel, and given them the victory the Lord Almighty.\end{Resp}
\begin{Resp}\Vsign They decked the fore-front of the Temple with crowns of gold, and dedicated the Altar unto the Lord.\end{Resp}
\begin{Resp}\Rsign Who had done so great things for Israel, and given them the victory the Lord Almighty.\end{Resp}
\begin{Resp}\Vsign Glory be to the Father, and to the Son, \Star{} and to the Holy Ghost.\end{Resp}
\begin{Resp}\Rsign Who had done so great things for Israel, and given them the victory the Lord Almighty.\end{Resp}
\switchcolumn*

\end{paracol}

\Office{Feria III infra Hebdomadam IV Octobris}{Tuesday of the fourth week of October}{}
\addcontentsline{toc}{subsection}{Feria III infra Hebdomadam IV Octobris \textit{— Tuesday of the fourth week of October}}
\begin{paracol}{2}
\LA
\LessonHead{Lectio I}
\switchcolumn
\EN
\LessonHead{Lesson I}
\switchcolumn*

\LA
\LessonTitle{De libro secúndo Machabæórum}
\switchcolumn
\EN
\LessonTitle{Lesson from the second book of Machabees}
\switchcolumn*

\LA
\Cite{2 Mac 3:1-4}
\switchcolumn
\EN
\Cite{2 Mac 3:1-4}
\switchcolumn*

\LA
\noindent Igitur, cum sancta cívitas habitarétur in omni pace, leges étiam adhuc óptime custodiréntur, propter Oníæ pontíficis pietátem, et ánimos ódio habéntes mala, fiébat ut et ipsi reges et príncipes locum summo honóre dignum dúcerent, et templum máximis munéribus illustrárent: ita ut Seléucus Asiæ rex de reddítibus suis præstáret omnes sumptus ad ministérium sacrificiórum pertinéntes. Simon autem de tribu Bénjamin, præpósitus templi constitútus, contendébat, obsisténte sibi príncipe sacerdótum, iníquum áliquid in civitáte molíri.\par
\switchcolumn
\EN
\noindent Therefore when the holy city was inhabited with all peace, and the laws as yet were very well kept, because of the godliness of Onias the high priest, and the hatred his soul had of evil, It came to pass that even the kings themselves, and the princes esteemed the place worthy of the highest honour, and glorified the temple with very great gifts: So that Seleucus king of Asia allowed out of his revenues all the charges belonging to the ministry of the sacrifices. But one Simon of the tribe of Benjamin, who was appointed overseer of the temple, strove in opposition to the high priest, to bring about some unjust thing in the city.\par
\switchcolumn*

\LA
\begin{Resp}\Rsign Hic est fratrum amátor et pópuli Israël: \Star{} Hic est, qui multum orat pro pópulo et univérsa sancta civitáte Jerúsalem.\end{Resp}
\begin{Resp}\Vsign Vir iste in pópulo suo mitíssimus appáruit.\end{Resp}
\begin{Resp}\Rsign Hic est, qui multum orat pro pópulo et univérsa sancta civitáte Jerúsalem.\end{Resp}
\switchcolumn
\EN
\begin{Resp}\Rsign This is a lover of the brethren, and of the people of Israel \Star{} This is one who prayeth much for the people, and for all the Holy City, Jerusalem.\end{Resp}
\begin{Resp}\Vsign There appeared a man most gentle toward all his people.\end{Resp}
\begin{Resp}\Rsign This is one who prayeth much for the people, and for all the Holy City, Jerusalem.\end{Resp}
\switchcolumn*

\LA
\LessonHead{Lectio II}
\switchcolumn
\EN
\LessonHead{Lesson II}
\switchcolumn*

\LA
\Cite{2 Mac 3:5-8}
\switchcolumn
\EN
\Cite{2 Mac 3:5-8}
\switchcolumn*

\LA
\noindent Sed, cum víncere Oníam non posset, venit ad Apollónium Tharsǽæ fílium, qui eo témpore erat dux Cælesýriæ et Phœnícis: et nuntiávit ei pecúniis innumerabílibus plenum esse ærárium Jerosólymis, et commúnes cópias imménsas esse, quæ non pértinent ad ratiónem sacrificiórum; esse autem possíbile sub potestáte regis cádere univérsa. Cumque retulísset ad regem Apollónius de pecúniis quæ delátæ erant, ille accítum Heliodórum, qui erat super negótia ejus, misit cum mandátis, ut prædíctam pecúniam transportáret. Statímque Heliodórus iter est aggréssus, spécie quidem quasi per Cœlesýriam et Phœnícen civitátes esset peragratúrus, re vera autem regis propósitum perfectúrus.\par
\switchcolumn
\EN
\noindent And when he could not overcome Onias he went to Apollonius the son of Tharseas, who at that time was governor of Celesyria and Phenicia: And told him, that the treasury in Jerusalem was full of immense sums of money, and the common store was infinite, which did not belong to the account of the sacrifices: and that it was possible to bring all into the king's hands. Now when Apollonius had given the king notice concerning the money that he was told of, he called for Heliodorus, who had the charge over his affairs, and sent him with commission to bring him the foresaid money. So Heliodorus forthwith began his journey, under a colour of visiting the cities of Celesyria and Phenicia, but indeed to fulfil the king's purpose.\par
\switchcolumn*

\LA
\begin{Resp}\Rsign Tu, Dómine universórum, qui nullam habes indigéntiam, voluísti templum tuum fíeri in nobis; \Star{} Consérva domum istam immaculátam in ætérnum, Dómine.\end{Resp}
\begin{Resp}\Vsign Tu elegísti, Dómine, domum istam ad invocándum nomen tuum in ea, ut esset domus oratiónis et obsecratiónis pópulo tuo.\end{Resp}
\begin{Resp}\Rsign Consérva domum istam immaculátam in ætérnum, Dómine.\end{Resp}
\begin{Resp}\Vsign Glória Patri, et Fílio, \Star{} et Spirítui Sancto.\end{Resp}
\begin{Resp}\Rsign Consérva domum istam immaculátam in ætérnum, Dómine.\end{Resp}
\switchcolumn
\EN
\begin{Resp}\Rsign Thou, O Lord of all things, Who hast need of nothing, wast pleased that the Temple of thine habitation should be among us. \Star{} Therefore now, O Lord, keep this house ever undefiled\end{Resp}
\begin{Resp}\Vsign Thou, O Lord, didst choose this house, that thy Name should be called on therein, and to be an house of prayer and petition for thy people.\end{Resp}
\begin{Resp}\Rsign Therefore now, O Lord, keep this house ever undefiled.\end{Resp}
\begin{Resp}\Vsign Glory be to the Father, and to the Son, \Star{} and to the Holy Ghost.\end{Resp}
\begin{Resp}\Rsign Therefore now, O Lord, keep this house ever undefiled.\end{Resp}
\switchcolumn*

\LA
\LessonHead{Lectio III}
\switchcolumn
\EN
\LessonHead{Lesson III}
\switchcolumn*

\LA
\Cite{2 Mac 3:9-12}
\switchcolumn
\EN
\Cite{2 Mac 3:9-12}
\switchcolumn*

\LA
\noindent Sed, cum venísset Jerosólymam, et benígne a summo sacerdóte in civitáte esset excéptus, narrávit de dato indício pecuniárum, et cujus rei grátia adésset, apéruit; interrogábat autem si vere hæc ita essent. Tunc summus sacérdos osténdit depósita esse hæc, et victuália viduárum et pupillórum; quædam vero esse Hircáni Tobíæ viri valde eminéntis, in his quæ detúlerat ímpius Simon: univérsa autem argénti talénta esse quadragínta, et auri ducénta: décipi vero eos qui credidíssent loco et templo, quod per univérsum mundum honorátur, pro sui veneratióne et sanctitáte, omníno impossíbile esse.\par
\switchcolumn
\EN
\noindent And when he was come to Jerusalem, and had been courteously received in the city by the high priest, he told him what information had been given concerning the money: and declared the cause for which he was come: and asked if these things were so indeed. Then the high priest told him that these were sums deposited, and provisions for the subsistence of the widows and the fatherless. And that some part of that which wicked Simon had given intelligence of, belonged to Hircanus son of Tobias, a man of great dignity: and that the whole was four hundred talents of silver, and two hundred of gold: But that to deceive them who had trusted to the place and temple which is honoured throughout the whole world, for the reverence and holiness of it, was a thing which could not by any means be done.\par
\switchcolumn*

\LA
\begin{Resp}\Rsign Aperi óculos tuos, Dómine, et vide afflictiónem nostram: circumdedérunt nos gentes ad puniéndum nos: \Star{} Sed tu, Dómine, exténde brácchium tuum, et líbera ánimas nostras.\end{Resp}
\begin{Resp}\Vsign Afflíge oppriméntes nos et contuméliam faciéntes in supérbiam; et custódi partem tuam.\end{Resp}
\begin{Resp}\Rsign Sed tu, Dómine, exténde brácchium tuum, et líbera ánimas nostras.\end{Resp}
\begin{Resp}\Vsign Glória Patri, et Fílio, \Star{} et Spirítui Sancto.\end{Resp}
\begin{Resp}\Rsign Sed tu, Dómine, exténde brácchium tuum, et líbera ánimas nostras.\end{Resp}
\switchcolumn
\EN
\begin{Resp}\Rsign Open thine eyes, O Lord, and behold our affliction for the heathen are come round about us to punish us. \Star{} But Thou, O Lord, stretch forth thine arm, and deliver our souls.\end{Resp}
\begin{Resp}\Vsign Punish them that oppress us and with pride do us wrong, and keep thine own portion.\end{Resp}
\begin{Resp}\Rsign But Thou, O Lord, stretch forth thine arm, and deliver our souls.\end{Resp}
\begin{Resp}\Vsign Glory be to the Father, and to the Son, \Star{} and to the Holy Ghost.\end{Resp}
\begin{Resp}\Rsign But Thou, O Lord, stretch forth thine arm, and deliver our souls.\end{Resp}
\switchcolumn*

\end{paracol}

\Office{Feria IV infra Hebdomadam IV Octobris}{Wednesday of the fourth week of October}{}
\addcontentsline{toc}{subsection}{Feria IV infra Hebdomadam IV Octobris \textit{— Wednesday of the fourth week of October}}
\begin{paracol}{2}
\LA
\LessonHead{Lectio I}
\switchcolumn
\EN
\LessonHead{Lesson I}
\switchcolumn*

\LA
\LessonTitle{De libro secúndo Machabæórum}
\switchcolumn
\EN
\LessonTitle{Lesson from the second book of Machabees}
\switchcolumn*

\LA
\Cite{2 Mac 3:23-25}
\switchcolumn
\EN
\Cite{2 Mac 3:23-25}
\switchcolumn*

\LA
\noindent Heliodórus autem, quod decréverat, perficiébat eódem loco ipse cum satellítibus circa ærárium præsens. Sed spíritus omnipoténtis Dei magnam fecit suæ ostensiónis evidéntiam, ita ut omnes qui ausi fúerant parére ei, ruéntes Dei virtúte, in dissolutiónem et formídinem converteréntur. Appáruit enim illis quidam equus terríbilem habens sessórem, óptimis operiméntis adornátus. Isque cum ímpetu Heliodóro prióres calces elísit; qui autem ei sedébat videbátur arma habére áurea.\par
\switchcolumn
\EN
\noindent But Heliodorus executed that which he had resolved on, himself being present in the same place with his guard about the treasury. But the spirit of the almighty God gave a great evidence of his presence, so that all that had presumed to obey him, falling down by the power of God, were struck with fainting and dread. For there appeared to them a horse with a terrible rider upon him, adorned with a very rich covering: and he ran fiercely and struck Heliodorus with his fore feet, and he that sat upon him seemed to have armour of gold.\par
\switchcolumn*

\LA
\begin{Resp}\Rsign Refúlsit sol in clípeos áureos, et resplenduérunt montes ab eis: \Star{} Et fortitúdo géntium dissipáta est.\end{Resp}
\begin{Resp}\Vsign Erat enim exércitus magnus valde et fortis: et appropiávit Judas et exércitus ejus in prǽlio.\end{Resp}
\begin{Resp}\Rsign Et fortitúdo géntium dissipáta est.\end{Resp}
\switchcolumn
\EN
\begin{Resp}\Rsign The sun shone upon the shields of gold, and the mountains glistered therewith; \Star{} And the army of the heathens was spread abroad.\end{Resp}
\begin{Resp}\Vsign For the army was very great and mighty then Judas and his host drew near and entered into battle.\end{Resp}
\begin{Resp}\Rsign And the army of the heathens was spread abroad.\end{Resp}
\switchcolumn*

\LA
\LessonHead{Lectio II}
\switchcolumn
\EN
\LessonHead{Lesson II}
\switchcolumn*

\LA
\Cite{2 Mac 3:26-29}
\switchcolumn
\EN
\Cite{2 Mac 3:26-29}
\switchcolumn*

\LA
\noindent Alii étiam apparuérunt duo júvenes virtúte decóri, óptimi glória, speciosíque amíctu: qui circumstetérunt eum, et ex utráque parte flagellábant, sine intermissióne multis plagis verberántes. Súbito autem Heliodórus cóncidit in terram, eúmque multa calígine circumfúsum rapuérunt, atque in sella gestatória pósitum ejecérunt. Et is, qui cum multis cursóribus et satellítibus prædíctum ingréssus est ærárium, portabátur nullo sibi auxílium ferénte, manifésta Dei cógnita virtúte. Et ille quidem per divínam virtútem jacébat mutus, atque omni spe et salúte privátus.\par
\switchcolumn
\EN
\noindent Moreover there appeared two other young men beautiful and strong, bright and glorious, and in comely apparel: who stood by him, on either side, and scourged him without ceasing with many stripes. And Heliodorus suddenly fell to the ground, and they took him up covered with great darkness, and having put him into a litter they carried him out. So he that came with many servants, and all his guard into the aforesaid treasury, was carried out, no one being able to help him, the manifest power of God being known. And he indeed by the power of God lay speechless, and without all hope of recovery.\par
\switchcolumn*

\LA
\begin{Resp}\Rsign Ornavérunt fáciem templi corónis áureis, et dedicavérunt altáre Dómino: \Star{} Et facta est lætítia magna in pópulo.\end{Resp}
\begin{Resp}\Vsign In hymnis et confessiónibus benedicébant Dóminum.\end{Resp}
\begin{Resp}\Rsign Et facta est lætítia magna in pópulo.\end{Resp}
\switchcolumn
\EN
\begin{Resp}\Rsign They decked the fore-front of the Temple with crowns of gold, and dedicated the Altar unto the Lord. \Star{} And there was very great gladness among the people.\end{Resp}
\begin{Resp}\Vsign They praised the Lord with Psalms and thanksgiving.\end{Resp}
\begin{Resp}\Rsign And there was very great gladness among the people.\end{Resp}
\switchcolumn*

\LA
\LessonHead{Lectio III}
\switchcolumn
\EN
\LessonHead{Lesson III}
\switchcolumn*

\LA
\Cite{2 Mac 3:32-34}
\switchcolumn
\EN
\Cite{2 Mac 3:32-24}
\switchcolumn*

\LA
\noindent Consíderans autem summus sacérdos ne forte rex suspicarétur malítiam áliquam ex Judǽis circa Heliodórum consummátam, óbtulit pro salúte viri hóstiam salutárem. Cumque summus sacérdos exoráret, iídem júvenes eísdem véstibus amícti astántes Heliodóro, dixérunt: Oníæ sacerdóti grátias age; nam propter eum Dóminus tibi vitam donávit. Tu autem a Deo flagellátus, núntia ómnibus magnália Dei, et potestátem. Et, his dictis, non comparuérunt.\par
\switchcolumn
\EN
\noindent So the high priest considering that the king might perhaps suspect that some mischief had been done to Heliodorus by the Jews, offered a sacrifice of health for the recovery of the man. And when the high priest was praying, the same young men in the same clothing stood by Heliodorus, and said to him: Give thanks to Onias the priest: because for his sake the Lord hath granted thee life. And thou having been scourged by God, declare unto all men the great works and the power of God. And having spoken thus, they appeared no more.\par
\switchcolumn*

\LA
\begin{Resp}\Rsign In hymnis et confessiónibus benedicébant Dóminum, \Star{} Qui magna fecit in Israël, et victóriam dedit illis Dóminus omnípotens.\end{Resp}
\begin{Resp}\Vsign Ornavérunt fáciem templi corónis áureis, et dedicavérunt altáre Dómino.\end{Resp}
\begin{Resp}\Rsign Qui magna fecit in Israël, et victóriam dedit illis Dóminus omnípotens.\end{Resp}
\begin{Resp}\Vsign Glória Patri, et Fílio, \Star{} et Spirítui Sancto.\end{Resp}
\begin{Resp}\Rsign Qui magna fecit in Israël, et victóriam dedit illis Dóminus omnípotens.\end{Resp}
\switchcolumn
\EN
\begin{Resp}\Rsign They praised the Lord with psalms and thanksgiving \Star{} Who had done so great things in Israel, and given them the victory the Lord Almighty.\end{Resp}
\begin{Resp}\Vsign They decked the fore-front of the Temple with crowns of gold, and dedicated the Altar unto the Lord.\end{Resp}
\begin{Resp}\Rsign Who had done so great things for Israel, and given them the victory the Lord Almighty.\end{Resp}
\begin{Resp}\Vsign Glory be to the Father, and to the Son, \Star{} and to the Holy Ghost.\end{Resp}
\begin{Resp}\Rsign Who had done so great things for Israel, and given them the victory the Lord Almighty.\end{Resp}
\switchcolumn*

\end{paracol}

\Office{Feria V infra Hebdomadam IV Octobris}{Thursday of the fourth week of October}{}
\addcontentsline{toc}{subsection}{Feria V infra Hebdomadam IV Octobris \textit{— Thursday of the fourth week of October}}
\begin{paracol}{2}
\LA
\LessonHead{Lectio I}
\switchcolumn
\EN
\LessonHead{Lesson I}
\switchcolumn*

\LA
\LessonTitle{De libro secúndo Machabæórum}
\switchcolumn
\EN
\LessonTitle{Lesson from the second book of Machabees}
\switchcolumn*

\LA
\Cite{2 Mac 4:1-5}
\switchcolumn
\EN
\Cite{2 Mac 4:1-5}
\switchcolumn*

\LA
\noindent Simon autem prædíctus, pecuniárum et pátriæ delátor, male loquebátur de Onía, tamquam ipse Heliodórum instigásset ad hæc, et ipse fuísset incéntor malórum: provisorémque civitátis, ac defensórem gentis suæ, et æmulatórem legis Dei, audébat insidiatórem regni dícere. Sed, cum inimicítiæ in tantum procéderent, ut étiam per quosdam Simónis necessários homicídia fíerent, consíderans Onías perículum contentiónis, et Apollónium insaníre, útpote ducem Cœlesýriæ et Phœnícis, ad augéndam malítiam Simónis ad regem se cóntulit, non ut cívium accusátor, sed commúnem utilitátem apud semetípsum univérsæ multitúdinis consíderans.\par
\switchcolumn
\EN
\noindent But Simon, of whom we spoke before, and of his country, spoke ill of Onias, as though he had incited Heliodorus to do these things, and had been the promoter of evils: And he presumed to call him a traitor to the kingdom, who provided for the city, and defended his nation, and was zealous for the law of God. But when the enmities proceeded so far, that murders also were committed by some of Simon's friends: Onias considering the danger of this contention, and that Apollonius, who was the governor of Celesyria and Phenicia, was outrageous, which increased the malice of Simon, went to the king, not to be an accuser of his countrymen, but with a view to the common good of all the people.\par
\switchcolumn*

\LA
\begin{Resp}\Rsign Adapériat Dóminus cor vestrum in lege sua et in præcéptis suis et fáciat pacem in diébus vestris: \Star{} Concédat vobis salútem, et rédimat vos a malis.\end{Resp}
\begin{Resp}\Vsign Exáudiat Dóminus oratiónes vestras, et reconciliétur vobis, nec vos déserat in témpore malo.\end{Resp}
\begin{Resp}\Rsign Concédat vobis salútem, et rédimat vos a malis.\end{Resp}
\switchcolumn
\EN
\begin{Resp}\Rsign The Lord open your hearts in His law and commandments, and send peace in your days. \Star{} May He grant you salvation and redeem you out of all evil.\end{Resp}
\begin{Resp}\Vsign The Lord hear your prayers, and be at one with you, and never forsake you in the time of trouble.\end{Resp}
\begin{Resp}\Rsign May He grant you salvation and redeem you out of all evil.\end{Resp}
\switchcolumn*

\LA
\LessonHead{Lectio II}
\switchcolumn
\EN
\LessonHead{Lesson II}
\switchcolumn*

\LA
\Cite{2 Mac 4:6-9}
\switchcolumn
\EN
\Cite{2 Mac 4:6-9}
\switchcolumn*

\LA
\noindent Vidébat enim sine regáli providéntia impossíbile esse pacem rebus dari, nec Simónem posse cessáre a stultítia sua. Sed post Seléuci vitæ excéssum, cum suscepísset regnum Antíochus, qui Nóbilis appellabátur, ambiébat Jason frater Oníæ summum sacerdótium. Adito rege, promíttens ei argénti talénta trecénta sexagínta, et ex reddítibus áliis talénta octogínta, super hæc promittébat et ália centum quinquagínta, si potestáti ejus concederétur, gymnásium et ephebíam sibi constitúere, et eos qui in Jerosólymis erant, Antiochénos scríbere.\par
\switchcolumn
\EN
\noindent For he saw that, except the king took care, it was impossible that matters should be settled in peace, or that Simon would cease from his folly. But after the death of Seleucus, when Antiochus, who was called the Illustrious, had taken possession of the kingdom, Jason the brother of Onias ambitiously sought the high priesthood: And went to the king, promising him three hundred and sixty talents of silver, and out of other revenues four- score talents. Besides this he promised also a hundred and fifty more, if he might have license to set him up a place for exercise, and a place for youth, and to entitle them, that were at Jerusalem, Antiochians.\par
\switchcolumn*

\LA
\begin{Resp}\Rsign Exáudiat Dóminus oratiónes vestras, et reconciliétur vobis, nec vos déserat in témpore malo, \Star{} Dóminus, Deus noster.\end{Resp}
\begin{Resp}\Vsign Det vobis cor ómnibus, ut colátis eum et faciátis ejus voluntátem.\end{Resp}
\begin{Resp}\Rsign Dóminus, Deus noster.\end{Resp}
\begin{Resp}\Vsign Glória Patri, et Fílio, \Star{} et Spirítui Sancto.\end{Resp}
\begin{Resp}\Rsign Dóminus, Deus noster.\end{Resp}
\switchcolumn
\EN
\begin{Resp}\Rsign The Lord hear your prayers, and be at one with you, and never forsake you in the time of trouble \Star{} Even He, the Lord our God.\end{Resp}
\begin{Resp}\Vsign Give you all an heart to serve Him, and to do His will.\end{Resp}
\begin{Resp}\Rsign Even He, the Lord our God.\end{Resp}
\begin{Resp}\Vsign Glory be to the Father, and to the Son, \Star{} and to the Holy Ghost.\end{Resp}
\begin{Resp}\Rsign Even He, the Lord our God.\end{Resp}
\switchcolumn*

\LA
\LessonHead{Lectio III}
\switchcolumn
\EN
\LessonHead{Lesson III}
\switchcolumn*

\LA
\Cite{2 Mac 4:10-11}
\switchcolumn
\EN
\Cite{2 Mac 4:10-11}
\switchcolumn*

\LA
\noindent Quod cum rex annuísset, et obtinuísset principátum, statim ad gentílem ritum contribúles suos transférre cœpit. Et, amótis his quæ humanitátis causa Judǽis a régibus fúerant constitúta, per Joánnem patrem Eupólemi, qui apud Romános de amicítia et societáte functus est legatióne legítima, cívium jura destítuens, prava institúta sanciébat.\par
\switchcolumn
\EN
\noindent Which when the king had granted, and he had gotten the rule into his hands, forthwith he began to bring over his countrymen to the fashion of the heathens. And abolishing those things, which had been decreed of special favour by the kings in behalf of the Jews, by the means of John the father of that Eupolemus, who went ambassador to Rome to make amity and alliance, he disannulled the lawful ordinances of the citizens, and brought in fashions that were perverse.\par
\switchcolumn*

\LA
\begin{Resp}\Rsign Congregáti sunt inimíci nostri, et gloriántur in virtúte sua: cóntere fortitúdinem illórum, Dómine, et dispérge illos: \Star{} Ut cognóscant quia non est álius qui pugnet pro nobis, nisi tu, Deus noster.\end{Resp}
\begin{Resp}\Vsign Dispérge illos in virtúte tua, et déstrue eos, protéctor noster, Dómine.\end{Resp}
\begin{Resp}\Rsign Ut cognóscant quia non est álius qui pugnet pro nobis, nisi tu, Deus noster.\end{Resp}
\begin{Resp}\Vsign Glória Patri, et Fílio, \Star{} et Spirítui Sancto.\end{Resp}
\begin{Resp}\Rsign Ut cognóscant quia non est álius qui pugnet pro nobis, nisi tu, Deus noster.\end{Resp}
\switchcolumn
\EN
\begin{Resp}\Rsign Our enemies are gathered together, and make their boast of their own strength. O Lord, break their power, and scatter them \Star{} That they may know that there is none other that fighteth for us, but only Thou, O our God!\end{Resp}
\begin{Resp}\Vsign Scatter them in thy strength, and destroy them, O Lord our Shield.\end{Resp}
\begin{Resp}\Rsign That they may know that there is none other that fighteth for us, but only Thou, O our God!\end{Resp}
\begin{Resp}\Vsign Glory be to the Father, and to the Son, \Star{} and to the Holy Ghost.\end{Resp}
\begin{Resp}\Rsign That they may know that there is none other that fighteth for us, but only Thou, O our God!\end{Resp}
\switchcolumn*

\end{paracol}

\Office{Feria VI infra Hebdomadam IV Octobris}{Friday of the fourth week of October}{}
\addcontentsline{toc}{subsection}{Feria VI infra Hebdomadam IV Octobris \textit{— Friday of the fourth week of October}}
\begin{paracol}{2}
\LA
\LessonHead{Lectio I}
\switchcolumn
\EN
\LessonHead{Lesson I}
\switchcolumn*

\LA
\LessonTitle{De libro secúndo Machabæórum}
\switchcolumn
\EN
\LessonTitle{Lesson from the second book of Machabees}
\switchcolumn*

\LA
\Cite{2 Mac 5:1-4}
\switchcolumn
\EN
\Cite{2 Mac 5:1-4}
\switchcolumn*

\LA
\noindent Eódem témpore Antíochus secúndam profectiónem parávit in Ægýptum. Cóntigit autem per univérsam Jerosolymórum civitátem vidéri diébus quadragínta per áëra équites discurréntes, aurátas stolas habéntes et hastis, quasi cohórtes, armátos: et cursus equórum per órdines digéstos, et congressiónes fíeri cóminus, et scutórum motus, et galeatórum multitúdinem gládiis distríctis, et telórum jactus, et aureórum armórum splendórem, omnísque géneris loricárum. Quaprópter omnes rogábant in bonum monstra convérti.\par
\switchcolumn
\EN
\noindent At the same time Antiochus prepared for a second journey into Egypt. And it came to pass that through the whole city of Jerusalem for the space of forty days there were seen horsemen running in the air, in gilded raiment, and armed with spears, like bands of soldiers. And horses set in order by ranks, running one against another, with the shakings of shields, and a multitude of men in helmets, with drawn swords, and casting of darts, and glittering of golden armour, and of harnesses of all sorts. Wherefore all men prayed that these prodigies might turn to good.\par
\switchcolumn*

\LA
\begin{Resp}\Rsign Impetum inimicórum ne timuéritis: mémores estóte, quómodo salvi facti sunt patres nostri: \Star{} Et nunc clamémus in cælum et miserébitur nostri Deus noster.\end{Resp}
\begin{Resp}\Vsign Mementóte mirabílium ejus, quæ fecit pharaóni et exercítui ejus in Mari Rubro.\end{Resp}
\begin{Resp}\Rsign Et nunc clamémus in cælum et miserébitur nostri Deus noster.\end{Resp}
\switchcolumn
\EN
\begin{Resp}\Rsign Be ye not afraid of the assault of the enemy; remember how our fathers were delivered. \Star{} Now, therefore, let us cry unto heaven, and our God will have mercy upon us.\end{Resp}
\begin{Resp}\Vsign Remember His marvellous works that He hath done unto Pharaoh and his host in the Red Sea.\end{Resp}
\begin{Resp}\Rsign Now, therefore, let us cry unto heaven, and our God will have mercy upon us.\end{Resp}
\switchcolumn*

\LA
\LessonHead{Lectio II}
\switchcolumn
\EN
\LessonHead{Lesson II}
\switchcolumn*

\LA
\Cite{2 Mac 5:5-7}
\switchcolumn
\EN
\Cite{2 Mac 5:5-7}
\switchcolumn*

\LA
\noindent Sed, cum falsus rumor exísset, tamquam vita excessísset Antíochus, assúmptis Jason non minus mille viris, repénte aggréssus est civitátem: et, cívibus ad murum convolántibus, ad últimum apprehénsa civitáte, Meneláus fugit in arcem. Jason vero non parcébat in cæde cívibus suis, nec cogitábat prosperitátem advérsum cognátos malum esse máximum, árbitrans hóstium et non cívium se trophǽa captúrum. Et principátum quidem non obtínuit, finem vero insidiárum suárum confusiónem accépit, et prófugus íterum ábiit in Ammaníten.\par
\switchcolumn
\EN
\noindent Now when there was gone forth a false rumour, as though Antiochus had been dead, Jason taking with him no fewer than a thousand men, suddenly assaulted the city: and though the citizens ran together to the wall, the city at length was taken, and Menelaus fled into the castle. But Jason slew his countrymen without mercy, not considering that prosperity against one's own kindred is a very great evil, thinking they had been enemies, and not citizens, whom he conquered. Yet he did not get the principality, but received confusion at the end, for the reward of his treachery, and fled again into the country of the Ammonites.\par
\switchcolumn*

\LA
\begin{Resp}\Rsign Congregátæ sunt gentes in multitúdine, ut dímicent contra nos, et ignorámus quid ágere debeámus: \Star{} Dómine Deus, ad te sunt óculi nostri, ne pereámus.\end{Resp}
\begin{Resp}\Vsign Tu scis quæ cógitant in nos: quómodo potérimus subsístere ante fáciem illórum, nisi tu ádjuves nos?\end{Resp}
\begin{Resp}\Rsign Dómine Deus, ad te sunt óculi nostri, ne pereámus.\end{Resp}
\begin{Resp}\Vsign Glória Patri, et Fílio, \Star{} et Spirítui Sancto.\end{Resp}
\begin{Resp}\Rsign Dómine Deus, ad te sunt óculi nostri, ne pereámus.\end{Resp}
\switchcolumn
\EN
\begin{Resp}\Rsign The heathen are assembled together to fight against us, and we know not what we should do. \Star{} Our eyes look unto thee, O Lord our God, that we should not perish.\end{Resp}
\begin{Resp}\Vsign What things they imagine against us, Thou knowest. How shall we be able to stand against them, except Thou be our help.\end{Resp}
\begin{Resp}\Rsign Our eyes look unto thee, O Lord our God, that we should not perish.\end{Resp}
\begin{Resp}\Vsign Glory be to the Father, and to the Son, \Star{} and to the Holy Ghost.\end{Resp}
\begin{Resp}\Rsign Our eyes look unto thee, O Lord our God, that we should not perish.\end{Resp}
\switchcolumn*

\end{paracol}

\Office{Sabbato infra Hebdomadam IV Octobris}{Saturday of the fourth week of October}{}
\addcontentsline{toc}{subsection}{Sabbato infra Hebdomadam IV Octobris \textit{— Saturday of the fourth week of October}}
\begin{paracol}{2}
\LA
\LessonHead{Lectio I}
\switchcolumn
\EN
\LessonHead{Lesson I}
\switchcolumn*

\LA
\LessonTitle{De libro secúndo Machabæórum}
\switchcolumn
\EN
\LessonTitle{Lesson from the second book of Machabees}
\switchcolumn*

\LA
\Cite{2 Mac 6:1-4}
\switchcolumn
\EN
\Cite{2 Mac 6:1-4}
\switchcolumn*

\LA
\noindent Sed non post multum témporis, misit rex Antíochus senem quemdam Antiochénum, qui compélleret Judǽos, ut se transférrent a pátriis et Dei légibus: contamináre étiam, quod in Jerosólymis erat, templum; et cognomináre Jovis Olýmpii: et in Garízim, prout erant hi qui locum inhabitábant, Jovis hospitális. Péssima autem et univérsis gravis erat malórum incúrsio; nam templum luxúria et comessatiónibus géntium erat plenum, et scortántium cum meretrícibus: sacratísque ǽdibus mulíeres se ultro ingerébant, intro feréntes ea quæ non licébat.\par
\switchcolumn
\EN
\noindent But not long after the king sent a certain old man of Antioch, to compel the Jews to depart from the laws of their fathers and of God: And to defile the temple that was in Jerusalem, and to call it the temple of Jupiter Olympius: and that in Gazarim of Jupiter Hospitalis, according as they were that inhabited the place. And very bad was this invasion of evils and grievous to all. For the temple was full of the riot and revellings of the Gentiles: and of men lying with lewd women. And women thrust themselves of their accord into the holy places, and brought in things that were not lawful.\par
\switchcolumn*

\LA
\begin{Resp}\Rsign Refúlsit sol in clípeos áureos, et resplenduérunt montes ab eis: \Star{} Et fortitúdo géntium dissipáta est.\end{Resp}
\begin{Resp}\Vsign Erat enim exércitus magnus valde et fortis: et appropiávit Judas et exércitus ejus in prǽlio.\end{Resp}
\begin{Resp}\Rsign Et fortitúdo géntium dissipáta est.\end{Resp}
\switchcolumn
\EN
\begin{Resp}\Rsign The sun shone upon the shields of gold, and the mountains glistered therewith; \Star{} And the army of the heathens was spread abroad.\end{Resp}
\begin{Resp}\Vsign For the army was very great and mighty then Judas and his host drew near and entered into battle.\end{Resp}
\begin{Resp}\Rsign And the army of the heathens was spread abroad.\end{Resp}
\switchcolumn*

\LA
\LessonHead{Lectio II}
\switchcolumn
\EN
\LessonHead{Lesson II}
\switchcolumn*

\LA
\Cite{2 Mac 6:5-9}
\switchcolumn
\EN
\Cite{2 Mac 6:5-9}
\switchcolumn*

\LA
\noindent Altáre étiam plenum erat illícitis, quæ légibus prohibebántur; neque autem sábbata custodiebántur, neque dies solémnes pátrii servabántur, nec simplíciter Judǽum se esse quisquam confitebátur. Ducebántur autem cum amára necessitáte in die natális regis ad sacrifícia: et, cum Líberi sacra celebraréntur, cogebántur hédera coronáti Líbero circuíre. Decrétum autem éxiit in próximas gentílium civitátes, suggeréntibus Ptolemǽis, ut pari modo et ipsi advérsus Judǽos ágerent, ut sacrificárent; eos autem, qui nollent transíre ad institúta géntium, interfícerent: erat ergo vidére misériam.\par
\switchcolumn
\EN
\noindent The altar also was filled with unlawful things, which were forbidden by the laws. And neither were the sabbaths kept, nor the solemn days of the fathers observed, neither did any man plainly profess himself to be a Jew. But they were led by bitter constraint on the king's birthday to the sacrifices: and when the feast of Bacchus was kept, they were compelled to go about crowned with ivy in honour of Bacchus. And there went out a decree into the neighbouring cities of the Gentiles, by the suggestion of the Ptolemeans, that they also should act in like manner against the Jews, to oblige them to sacrifice: And whosoever would not conform themselves to the ways of the Gentiles, should be put to death: then was misery to be seen.\par
\switchcolumn*

\LA
\begin{Resp}\Rsign Ornavérunt fáciem templi corónis áureis, et dedicavérunt altáre Dómino: \Star{} Et facta est lætítia magna in pópulo.\end{Resp}
\begin{Resp}\Vsign In hymnis et confessiónibus benedicébant Dóminum.\end{Resp}
\begin{Resp}\Rsign Et facta est lætítia magna in pópulo.\end{Resp}
\begin{Resp}\Vsign Glória Patri, et Fílio, \Star{} et Spirítui Sancto.\end{Resp}
\begin{Resp}\Rsign Et facta est lætítia magna in pópulo.\end{Resp}
\switchcolumn
\EN
\begin{Resp}\Rsign They decked the fore-front of the Temple with crowns of gold, and dedicated the Altar unto the Lord. \Star{} And there was very great gladness among the people.\end{Resp}
\begin{Resp}\Vsign They praised the Lord with Psalms and thanksgiving.\end{Resp}
\begin{Resp}\Rsign And there was very great gladness among the people.\end{Resp}
\begin{Resp}\Vsign Glory be to the Father, and to the Son, \Star{} and to the Holy Ghost.\end{Resp}
\begin{Resp}\Rsign And there was very great gladness among the people.\end{Resp}
\switchcolumn*

\end{paracol}

\Office{Dominica V Octobris}{Fifth Sunday of October}{}
\addcontentsline{toc}{subsection}{Dominica V Octobris \textit{— Fifth Sunday of October}}
\begin{paracol}{2}
\LA
\LessonHead{Lectio I}
\switchcolumn
\EN
\LessonHead{Lesson I}
\switchcolumn*

\LA
\LessonTitle{De libro secúndo Machabæórum}
\switchcolumn
\EN
\LessonTitle{Lesson from the second book of Machabees}
\switchcolumn*

\LA
\Cite{2 Mac 6:18-22}
\switchcolumn
\EN
\Cite{2 Mac 6:18-22}
\switchcolumn*

\LA
\noindent Igitur Eleázarus, unus de primóribus scribárum, vir ætáte provéctus, et vultu decórus, apérto ore hians compellebátur carnem porcínam manducáre; at ille gloriosíssimam mortem magis quam odíbilem vitam compléctens, voluntárie præíbat ad supplícium. Intuens autem quemádmodum oportéret accédere, patiénter sústinens, destinávit non admíttere illícita propter vitæ amórem. Hi autem qui astábant, iníqua miseratióne commóti propter antíquam viri amicítiam, tolléntes eum secréto rogábant afférri carnes quibus vesci ei licébat, ut simularétur manducásse, sicut rex imperáverat, de sacrifícii cárnibus, ut hoc facto, a morte liberarétur: et propter véterem viri amicítiam, hanc in eo faciébant humanitátem.\par
\switchcolumn
\EN
\noindent Eleazar one of the chief of the scribes, a man advanced in years, and of a comely countenance, was pressed to open his mouth to eat swine's flesh. But he, choosing rather a most glorious death than a hateful life, went forward voluntarily to the torment. And considering in what manner he was come to it, patiently bearing, he determined not to do any unlawful things for the love of life. But they that stood by, being moved with wicked pity, for the old friendship they had with the man, taking him aside, desired that flesh might be brought, which it was lawful for him to eat, that he might make as if he had eaten, as the king had commanded of the flesh of the sacrifice: That by so doing he might be delivered from death: and for the sake of their old friendship with the man they did him this courtesy.\par
\switchcolumn*

\LA
\begin{Resp}\Rsign Adapériat Dóminus cor vestrum in lege sua et in præcéptis suis et fáciat pacem in diébus vestris: \Star{} Concédat vobis salútem, et rédimat vos a malis.\end{Resp}
\begin{Resp}\Vsign Exáudiat Dóminus oratiónes vestras, et reconciliétur vobis, nec vos déserat in témpore malo.\end{Resp}
\begin{Resp}\Rsign Concédat vobis salútem, et rédimat vos a malis.\end{Resp}
\switchcolumn
\EN
\begin{Resp}\Rsign The Lord open your hearts in His law and commandments, and send peace in your days. \Star{} May He grant you salvation and redeem you out of all evil.\end{Resp}
\begin{Resp}\Vsign The Lord hear your prayers, and be at one with you, and never forsake you in the time of trouble.\end{Resp}
\begin{Resp}\Rsign May He grant you salvation and redeem you out of all evil.\end{Resp}
\switchcolumn*

\LA
\LessonHead{Lectio II}
\switchcolumn
\EN
\LessonHead{Lesson II}
\switchcolumn*

\LA
\Cite{2 Mac 6:23-28}
\switchcolumn
\EN
\Cite{2 Mac 6:23-28}
\switchcolumn*

\LA
\noindent At ille cogitáre cœpit ætátis ac senectútis suæ eminéntiam dignam, et ingénitæ nobilitátis canítiem, atque a púero óptimæ conversatiónis actus: et secúndum sanctæ et a Deo cónditæ legis constitúta, respóndit cito, dicens præmítti se velle in inférnum. Non enim ætáti nostræ dignum est, inquit, fíngere: ut multi adolescéntium, arbitrántes Eleázarum nonagínta annórum transísse ad vitam alienigenárum, et ipsi propter meam simulatiónem, et propter módicum corruptíbilis vitæ tempus decipiántur, et per hoc máculam atque exsecratiónem meæ senectúti conquíram. Nam, etsi in præsénti témpore supplíciis hóminum erípiar, sed manum Omnipoténtis nec vivus, nec defúnctus, effúgiam. Quam ob rem fórtiter vita excedéndo, senectúte quidem dignus apparébo: adolescéntibus autem exémplum forte relínquam, si prompto ánimo ac fortiter pro gravíssimis ac sanctíssimis légibus honésta morte perfúngar. His dictis, conféstim ad supplícium trahebátur.\par
\switchcolumn
\EN
\noindent But he began to consider the dignity of his age, and his ancient years, and the inbred honour of his grey head, and his good life and conversation from a child: and he answered without delay, according to the ordinances of the holy law made by God, saying, that he would rather be sent into the other world. For it doth not become our age, said he, to dissemble: whereby many young persons might think that Eleazar, at the age of fourscore and ten years, was gone over to the life of the heathens: And so they, through my dissimulation, and for a little time of a corruptible life, should be deceived, end hereby I should bring a stain and a curse upon my old age. For though, for the present time, I should be delivered from the punishments of men, yet should I not escape the hand of the Almighty neither alive nor dead. Wherefore by departing manfully out of this life, I shall shew myself worthy of my old age: And I shall leave an example of fortitude to young men, if with a ready mind and constancy I suffer an honourable death, for the most venerable and most holy laws. And having spoken thus, he was forthwith carried to execution.\par
\switchcolumn*

\LA
\begin{Resp}\Rsign Tu, Dómine universórum, qui nullam habes indigéntiam, voluísti templum tuum fíeri in nobis; \Star{} Consérva domum istam immaculátam in ætérnum, Dómine.\end{Resp}
\begin{Resp}\Vsign Tu elegísti, Dómine, domum istam ad invocándum nomen tuum in ea, ut esset domus oratiónis et obsecratiónis pópulo tuo.\end{Resp}
\begin{Resp}\Rsign Consérva domum istam immaculátam in ætérnum, Dómine.\end{Resp}
\switchcolumn
\EN
\begin{Resp}\Rsign Thou, O Lord of all things, Who hast need of nothing, wast pleased that the Temple of thine habitation should be among us. \Star{} Therefore now, O Lord, keep this house ever undefiled\end{Resp}
\begin{Resp}\Vsign Thou, O Lord, didst choose this house, that thy Name should be called on therein, and to be an house of prayer and petition for thy people.\end{Resp}
\begin{Resp}\Rsign Therefore now, O Lord, keep this house ever undefiled.\end{Resp}
\switchcolumn*

\LA
\LessonHead{Lectio III}
\switchcolumn
\EN
\LessonHead{Lesson III}
\switchcolumn*

\LA
\Cite{2 Mac 7:1-5}
\switchcolumn
\EN
\Cite{2 Mac 7:1-5}
\switchcolumn*

\LA
\noindent Cóntigit autem et septem fratres una cum matre sua apprehénsos compélli a rege édere contra fas carnes porcínas, flagris et táureis cruciátos. Unus autem ex illis, qui erat primus, sic ait: Quid quæris et quid vis díscere a nobis? Paráti sumus mori, magis quam pátrias Dei leges prævaricári. Irátus ítaque rex jussit sartágines et ollas ǽneas succéndi. Quibus statim succénsis, jussit ei qui prior fúerat locútus, amputári linguam, et, cute cápitis abstrácta, summas quoque manus et pedes ei præscíndi, céteris ejus frátribus et matre inspiciéntibus. Et, cum jam per ómnia inútilis factus esset, jussit ignem admovéri, et adhuc spirántem torréri in sartágine: in qua, cum diu cruciarétur, céteri una cum matre ínvicem se hortabántur mori fórtiter.\par
\switchcolumn
\EN
\noindent It came to pass also, that seven brethren, together with their mother, were apprehended, and compelled by the king to eat swine's flesh against the law, for which end they were tormented with whips and scourges. But one of them, who was the eldest, said thus: What wouldst thou ask, or learn of us? we are ready to die rather than to transgress the laws of God, received from our fathers. Then the king being angry commanded fryingpans, and brazen caldrons to be made hot: which forthwith being heated, He commanded to cut out the tongue of him that had spoken first: and the skin of his head being drawn off, to chop off also the extremities of his hands and feet, the rest of his brethren, and his mother, looking on. And when he was now maimed in all parts, he commanded him, being yet alive, to be brought to the fire, and to be fried in the fryingpan: and while he was suffering therein long torments, the rest, together with the mother, exhorted one another to die manfully.\par
\switchcolumn*

\LA
\begin{Resp}\Rsign Aperi óculos tuos, Dómine, et vide afflictiónem nostram: circumdedérunt nos gentes ad puniéndum nos: \Star{} Sed tu, Dómine, exténde brácchium tuum, et líbera ánimas nostras.\end{Resp}
\begin{Resp}\Vsign Afflíge oppriméntes nos et contuméliam faciéntes in supérbiam; et custódi partem tuam.\end{Resp}
\begin{Resp}\Rsign Sed tu, Dómine, exténde brácchium tuum, et líbera ánimas nostras.\end{Resp}
\begin{Resp}\Vsign Glória Patri, et Fílio, \Star{} et Spirítui Sancto.\end{Resp}
\begin{Resp}\Rsign Sed tu, Dómine, exténde brácchium tuum, et líbera ánimas nostras.\end{Resp}
\switchcolumn
\EN
\begin{Resp}\Rsign Open thine eyes, O Lord, and behold our affliction for the heathen are come round about us to punish us. \Star{} But Thou, O Lord, stretch forth thine arm, and deliver our souls.\end{Resp}
\begin{Resp}\Vsign Punish them that oppress us and with pride do us wrong, and keep thine own portion.\end{Resp}
\begin{Resp}\Rsign But Thou, O Lord, stretch forth thine arm, and deliver our souls.\end{Resp}
\begin{Resp}\Vsign Glory be to the Father, and to the Son, \Star{} and to the Holy Ghost.\end{Resp}
\begin{Resp}\Rsign But Thou, O Lord, stretch forth thine arm, and deliver our souls.\end{Resp}
\switchcolumn*

\end{paracol}

\Office{Feria III infra Hebdomadam V Octobris}{Tuesday of the fifth week of October}{}
\addcontentsline{toc}{subsection}{Feria III infra Hebdomadam V Octobris \textit{— Tuesday of the fifth week of October}}
\begin{paracol}{2}
\LA
\LessonHead{Lectio I}
\switchcolumn
\EN
\LessonHead{Lesson I}
\switchcolumn*

\LA
\LessonTitle{De libro secúndo Machabæórum}
\switchcolumn
\EN
\LessonTitle{Lesson from the second book of Machabees}
\switchcolumn*

\LA
\Cite{2 Mac 7:7-12}
\switchcolumn
\EN
\Cite{2 Mac 7:7-12}
\switchcolumn*

\LA
\noindent Mórtuo ítaque illo primo hoc modo, sequéntem deducébant ad illudéndum: et, cute cápitis ejus cum capíllis abstrácta, interrogábant, si manducáret prius quam toto córpore per membra síngula punirétur. At ille respóndens pátria voce, dixit: Non fáciam. Propter quod et iste, sequénti loco, primi torménta suscépit: et in último spíritu constitútus, sic ait: Tu quidem, scelestíssime, in præsénti vita nos perdis: sed Rex mundi defúnctos nos pro suis légibus in ætérnæ vitæ resurrectióne suscitábit. Post hunc tértius illúditur, et linguam postulátus cito prótulit, et manus constánter exténdit: et cum fidúcia ait: E cælo ista possídeo, sed propter Dei leges nunc hæc ipsa despício, quóniam ab ipso me ea receptúrum spero; ita ut rex, et qui cum ipso erant, miraréntur adolescéntis ánimum, quod tamquam níhilum dúceret cruciátus.\par
\switchcolumn
\EN
\noindent So when the first was dead after this manner, they brought the next to make him a mocking stock: and when they had pulled off the skin of his head with the hair, they asked him if he would eat, before he were punished throughout the whole body in every limb. But he answered in his own language, and said: I will not do it. Wherefore he also in the next place, received the torments of the first: And when he was at the last gasp, he said thus: Thou indeed, O most wicked man, destroyest us out of this present life: but the King of the world will raise us up, who die for his laws, in the resurrection of eternal life. After him the third was made a mocking stock, and when he was required, he quickly put forth his tongue, and courageously stretched out his hands: And said with confidence: These I have from heaven, but for the laws of God I now despise them, because I hope to receive them again from him. So that the king, and they that were with him, wondered at the young man's courage, because he esteemed the torments as nothing.\par
\switchcolumn*

\LA
\begin{Resp}\Rsign Hic est fratrum amátor et pópuli Israël: \Star{} Hic est, qui multum orat pro pópulo et univérsa sancta civitáte Jerúsalem.\end{Resp}
\begin{Resp}\Vsign Vir iste in pópulo suo mitíssimus appáruit.\end{Resp}
\begin{Resp}\Rsign Hic est, qui multum orat pro pópulo et univérsa sancta civitáte Jerúsalem.\end{Resp}
\switchcolumn
\EN
\begin{Resp}\Rsign This is a lover of the brethren, and of the people of Israel \Star{} This is one who prayeth much for the people, and for all the Holy City, Jerusalem.\end{Resp}
\begin{Resp}\Vsign There appeared a man most gentle toward all his people.\end{Resp}
\begin{Resp}\Rsign This is one who prayeth much for the people, and for all the Holy City, Jerusalem.\end{Resp}
\switchcolumn*

\LA
\LessonHead{Lectio II}
\switchcolumn
\EN
\LessonHead{Lesson II}
\switchcolumn*

\LA
\Cite{2 Mac 7:13-19}
\switchcolumn
\EN
\Cite{2 Mac 7:13-19}
\switchcolumn*

\LA
\noindent Et hoc ita defúncto, quartum vexábant simíliter torquéntes; et, cum jam esset ad mortem, sic ait: Pótius est ab homínibus morti datos spem exspectáre a Deo, íterum ab ipso resuscitándos: tibi enim resurréctio ad vitam non erit. Et, cum admovíssent quintum, vexábant eum. At ille respíciens in eum, dixit: Potestátem inter hómines habens, cum sis corruptíbilis, facis quod vis: noli autem putáre genus nostrum a Deo esse derelíctum. Tu autem patiénter sústine, et vidébis magnam potestátem ipsíus, quáliter te et semen tuum torquébit. Post hunc ducébant sextum, et is, mori incípiens, sic ait: Noli frustra erráre; nos enim propter nosmetípsos hæc pátimur, peccántes in Deum nostrum, et digna admiratióne facta sunt in nobis. Tu autem ne exístimes tibi impúne futúrum, quod contra Deum pugnáre tentáveris.\par
\switchcolumn
\EN
\noindent And after he was thus dead, they tormented the fourth in the like manner And when he was now ready to die, he spoke thus: It is better, being put to death by men, to look for hope from God, to be raised up again by him: for, as to thee thou shalt have no resurrection unto life. And when they had brought the fifth, they tormented him. But he looking upon the king, Said: Whereas thou hast power among men, though thou art corruptible, thou dost what thou wilt: but think not that our nation is forsaken by God. But stay patiently a while, and thou shalt see his great power, in what manner he will torment thee and thy seed. After him they brought the sixth, and he being ready to die, spoke thus: Be not deceived without cause: for we suffer these things for ourselves, having sinned against our God, and things worthy of admiration are done to us: But do not think that thou shalt escape unpunished, for that thou attempted to fight against God.\par
\switchcolumn*

\LA
\begin{Resp}\Rsign Tu, Dómine universórum, qui nullam habes indigéntiam, voluísti templum tuum fíeri in nobis; \Star{} Consérva domum istam immaculátam in ætérnum, Dómine.\end{Resp}
\begin{Resp}\Vsign Tu elegísti, Dómine, domum istam ad invocándum nomen tuum in ea, ut esset domus oratiónis et obsecratiónis pópulo tuo.\end{Resp}
\begin{Resp}\Rsign Consérva domum istam immaculátam in ætérnum, Dómine.\end{Resp}
\switchcolumn
\EN
\begin{Resp}\Rsign Thou, O Lord of all things, Who hast need of nothing, wast pleased that the Temple of thine habitation should be among us. \Star{} Therefore now, O Lord, keep this house ever undefiled\end{Resp}
\begin{Resp}\Vsign Thou, O Lord, didst choose this house, that thy Name should be called on therein, and to be an house of prayer and petition for thy people.\end{Resp}
\begin{Resp}\Rsign Therefore now, O Lord, keep this house ever undefiled.\end{Resp}
\switchcolumn*

\LA
\LessonHead{Lectio III}
\switchcolumn
\EN
\LessonHead{Lesson III}
\switchcolumn*

\LA
\Cite{2 Mac 7:20-23}
\switchcolumn
\EN
\Cite{2 Mac 7:20-23}
\switchcolumn*

\LA
\noindent Supra modum autem mater mirábilis, et bonórum memória digna, quæ pereúntes septem fílios sub uníus diéi témpore conspíciens, bono ánimo ferébat propter spem quam in Deum habébat; síngulos illórum hortabátur voce pátria fórtiter, repléta sapiéntia: et femíneæ cogitatióni masculínum ánimum ínserens, dixit ad eos: Néscio quáliter in útero meo apparuístis, neque enim ego spíritum et ánimam donávi vobis et vitam et singulórum membra non ego ipsa compégi; sed enim mundi Creátor, qui formávit hóminis nativitátem, quique ómnium invénit oríginem, et spíritum vobis íterum cum misericórdia reddet et vitam, sicut nunc vosmetípsos despícitis propter leges ejus.\par
\switchcolumn
\EN
\noindent Now the mother was to be admired above measure, and worthy to be remembered by good men, who beheld seven sons slain in the space of one day, and bore it with a good courage, for the hope that she had in God: And she bravely exhorted every one of them in her own language, being filled with wisdom: and joining a man's heart to a woman's thought, She said to them: I know not how you were formed in my womb: for I neither gave you breath, nor soul, nor life, neither did I frame the limbs of every one of you. But the Creator of the world, that formed the nativity of man, and that found out the origin of all, he will restore to you again in his mercy, both breath and life, as now you despise yourselves for the sake of his laws.\par
\switchcolumn*

\LA
\begin{Resp}\Rsign Aperi óculos tuos, Dómine, et vide afflictiónem nostram: circumdedérunt nos gentes ad puniéndum nos: \Star{} Sed tu, Dómine, exténde brácchium tuum, et líbera ánimas nostras.\end{Resp}
\begin{Resp}\Vsign Afflíge oppriméntes nos et contuméliam faciéntes in supérbiam; et custódi partem tuam.\end{Resp}
\begin{Resp}\Rsign Sed tu, Dómine, exténde brácchium tuum, et líbera ánimas nostras.\end{Resp}
\begin{Resp}\Vsign Glória Patri, et Fílio, \Star{} et Spirítui Sancto.\end{Resp}
\begin{Resp}\Rsign Sed tu, Dómine, exténde brácchium tuum, et líbera ánimas nostras.\end{Resp}
\switchcolumn
\EN
\begin{Resp}\Rsign Open thine eyes, O Lord, and behold our affliction for the heathen are come round about us to punish us. \Star{} But Thou, O Lord, stretch forth thine arm, and deliver our souls.\end{Resp}
\begin{Resp}\Vsign Punish them that oppress us and with pride do us wrong, and keep thine own portion.\end{Resp}
\begin{Resp}\Rsign But Thou, O Lord, stretch forth thine arm, and deliver our souls.\end{Resp}
\begin{Resp}\Vsign Glory be to the Father, and to the Son, \Star{} and to the Holy Ghost.\end{Resp}
\begin{Resp}\Rsign But Thou, O Lord, stretch forth thine arm, and deliver our souls.\end{Resp}
\switchcolumn*

\end{paracol}

\SeasonTitle{Dominicæ et Feriæ mensis Novembris}{Sundays and Ferias of November}

\addcontentsline{toc}{section}{Dominicæ et Feriæ mensis Novembris \textit{— Sundays and Ferias of November}}

\Office{Dominica I Novembris}{First Sunday of November}{}
\addcontentsline{toc}{subsection}{Dominica I Novembris \textit{— First Sunday of November}}
\begin{paracol}{2}
\LA
\LessonHead{Lectio IV}
\switchcolumn
\EN
\LessonHead{Lesson IV}
\switchcolumn*

\LA
\LessonTitle{De Expositióne sancti Gregórii Papæ in Ezechiélem Prophétam}
\switchcolumn
\EN
\LessonTitle{From the Exposition of the Prophet Ezekiel written by Pope St. Gregory (the Great.)}
\switchcolumn*

\LA
\Source{Liber 1, Homilía 2.}
\switchcolumn
\EN
\Source{Bk, i. Horn. 2}
\switchcolumn*

\LA
\noindent Usus prophéticæ locutiónis est, ut prius persónam, tempus, locúmque descríbat, et póstmodum dícere mystéria prophetíæ incípiat; quátenus ad veritátem solídius ostendéndam, ante, históriæ radícem figat, et post, fructus spíritus per signa et allegorías próferat. Ezéchiel ítaque ætátis suæ tempus índicat, dicens: Et factum est in trigésimo anno, in quarto mense, in quinta mensis. Locum quoque denúntians, adjúngit: Cum essem in médio captivórum juxta flumen Chobar, apérti sunt cæli, et vidi visiónes Dei. Tempus étiam insínuat, subdens, In quinta mensis, ipse est annus quintus transmigratiónis regis Jóachin. Qui ut bene persónam índicet, étiam genus narrat, cum súbditur: Et factum est verbum Dómini ad Ezechiélem, fílium Buzi, sacerdótem.\par
\switchcolumn
\EN
\noindent It is the use of the Prophetic writers first to give name, date, and place, and then to begin to unfold the mysteries of the prophecy thus, to give certainty of trustworthiness, a foundation is laid before, and afterward the fruits of the Spirit are set forth by signs and in figures. Thus Ezekiel saith concerning the date: And it came to pass in the thirtieth year, in the fourth month, in the fifth day of the month. And to show the place, he addeth further: “As I was among the captives by the river of Chebar, the heavens were opened, and I saw visions of God.” Then he defineth the time even more exactly, saying In the fifth day of the month, which was the fifth year of King Jehoiachim's captivity. And he who had thus clearly indicated his individuality, goeth on farther to state his kin, saying The word of the Lord came unto Ezekiel, the son of Buzi, the Priest.\par
\switchcolumn*

\LA
\begin{Resp}\Rsign Super muros tuos, Jerúsalem, constítui custódes; \Star{} Tota die et nocte non tacébunt laudáre nomen Dómini.\end{Resp}
\begin{Resp}\Vsign Prædicábunt pópulis fortitúdinem meam, et annuntiábunt géntibus glóriam meam.\end{Resp}
\begin{Resp}\Rsign Tota die et nocte non tacébunt laudáre nomen Dómini.\end{Resp}
\switchcolumn
\EN
\begin{Resp}\Rsign I have set watchmen upon thy walls, O Jerusalem; \Star{} Which shall never hold their peace day nor night, to praise the name of the Lord.\end{Resp}
\begin{Resp}\Vsign They shall proclaim My might unto the nations, and declare My glory unto the Gentiles.\end{Resp}
\begin{Resp}\Rsign Which shall never hold their peace day nor night, to praise the name of the Lord.\end{Resp}
\switchcolumn*

\LA
\LessonHead{Lectio V}
\switchcolumn
\EN
\LessonHead{Lesson V}
\switchcolumn*

\LA
\noindent Sed prima quǽstio nobis óritur, cur is qui nihil adhuc díxerat, ita exórsus est dicens: Et factum est in trigésimo anno. Et, namque sermo conjunctiónis est: et scimus quia non conjúngitur sermo súbsequens, nisi sermóni præcedénti. Qui ígitur nihil díxerat, cur dicit, Et factum est: cum non sit sermo, cui hoc quod íncipit, subjúngat? Qua in re intuéndum est, quia sicut nos corporália, sic prophétæ sensu spiritália aspíciunt; eísque et illa sunt præséntia, quæ nostræ ignorántiæ abséntia vidéntur. Unde fit, ut in mente prophetárum ita conjúncta sint exterióribus interióra, quátenus simul útraque vídeant, simúlque in eis fiat et intus verbum quod áudiunt, et foris quod dicunt.\par
\switchcolumn
\EN
\noindent But the first question which meeteth us is: Wherefore doth the Prophet, having hitherto said nothing, begin with the words: “And it came to pass in the thirtieth year.” Now, this word “And” is a conjunction, and we know that it is so called because it conjoineth that which cometh after it with that which goeth before it. Wherefore, then, doth he who hath hitherto been silent, commence by And, when there is nothing going before for the conjunction to join to that which cometh after. To explain this, we must consider that our senses perceive only things bodily, while those of Prophets perceive also things ghostly, and to them things exist which to our ignorance seem not to do so. Hence it cometh that in the mind of a Prophet, things outer and things inner are so joined that he seeth both together, and the word which he heareth within him and that which he uttereth come together.\par
\switchcolumn*

\LA
\begin{Resp}\Rsign Muro tuo inexpugnábili circumcínge nos, Dómine, et armis tuæ poténtiæ prótege nos semper: \Star{} Líbera, Dómine, Deus Israël, clamántes ad te.\end{Resp}
\begin{Resp}\Vsign Erue nos in mirabílibus tuis, et da glóriam nómini tuo.\end{Resp}
\begin{Resp}\Rsign Líbera, Dómine, Deus Israël, clamántes ad te.\end{Resp}
\switchcolumn
\EN
\begin{Resp}\Rsign Hedge us about with thy wall that cannot be broken down, O Lord, and shield us continually with the arms of thy might. \Star{} O Lord God of Israel, deliver them that cry unto thee.\end{Resp}
\begin{Resp}\Vsign Deliver us also according to thy marvellous works, and give glory to thy Name.\end{Resp}
\begin{Resp}\Rsign O Lord God of Israel, deliver them that cry unto thee.\end{Resp}
\switchcolumn*

\LA
\LessonHead{Lectio VI}
\switchcolumn
\EN
\LessonHead{Lesson VI}
\switchcolumn*

\LA
\noindent Patet ígitur causa, cur qui nihil díxerat, inchoávit dicens: Et factum est in trigésimo anno: quia hoc verbum quod foris prótulit, illi verbo quod intus audíerat, conjúnxit. Continuávit ergo verba quæ prótulit visióni íntimæ, et idcírco íncipit dicens, Et factum est. Subjúnxit enim hoc quod extérius loqui ínchoat; ac si et illud foris sit, quod intus vidit. Hoc autem quod dícitur, quia in trigésimo anno spíritum prophetíæ accéperit, índicat áliquid nobis considerándum; vidélicet quia juxta ratiónis usum, doctrínæ sermo non súppetit, nisi in ætáte perfécta. Unde et ipse Dóminus anno duodécimo ætátis suæ in médio doctórum in templo sedens, non docens sed intérrogans vóluit inveníri.\par
\switchcolumn
\EN
\noindent It appeareth plainly, therefore, that he which had hitherto been silent, beginneth by the words: “And it came to pass in the thirtieth year”, because his first utterance was but the continuation of something to which he had already been listening in his own mind. The words which he spoke were merely a continuation of the vision already going on within, and therefore the first are, And it came to pass. His language is framed as though his inner revelation had been an open one. That it was in the thirtieth year that the word of the Lord came unto Ezekiel, causeth us to remark that in the ordinary use of human understanding, men receive not a call to teach until they be of full age. Hence also even the Lord Himself, when He sat in the Temple in the midst of the doctors, in the twelfth year of His age, was pleased to be found, not teaching, but hearing them and asking them questions. (Luke ii. 46.)\par
\switchcolumn*

\LA
\begin{Resp}\Rsign Sustinúimus pacem, et non venit: quæsívimus bona, et ecce turbátio: cognóvimus, Dómine, peccáta nostra; \Star{} Non in perpétuum obliviscáris nos.\end{Resp}
\begin{Resp}\Vsign Peccávimus, ímpie géssimus, iniquitátem fécimus, Dómine, in omnem justítiam tuam.\end{Resp}
\begin{Resp}\Rsign Non in perpétuum obliviscáris nos.\end{Resp}
\begin{Resp}\Vsign Glória Patri, et Fílio, \Star{} et Spirítui Sancto.\end{Resp}
\begin{Resp}\Rsign Non in perpétuum obliviscáris nos.\end{Resp}
\switchcolumn
\EN
\begin{Resp}\Rsign We looked for peace, and it came not we asked for good, and behold trouble. We acknowledge, O Lord, our wickedness. \Star{} Forget us not for ever.\end{Resp}
\begin{Resp}\Vsign O Lord (our God) we have sinned, we have done ungodly, we have dealt unrighteously in all thine ordinances.\end{Resp}
\begin{Resp}\Rsign Forget us not for ever.\end{Resp}
\begin{Resp}\Vsign Glory be to the Father, and to the Son, \Star{} and to the Holy Ghost.\end{Resp}
\begin{Resp}\Rsign Forget us not for ever.\end{Resp}
\switchcolumn*

\end{paracol}

\Office{Feria II infra Hebdomadam I Novembris}{Monday of the first week of November}{}
\addcontentsline{toc}{subsection}{Feria II infra Hebdomadam I Novembris \textit{— Monday of the first week of November}}
\begin{paracol}{2}
\LA
\LessonHead{Lectio I}
\switchcolumn
\EN
\LessonHead{Lesson I}
\switchcolumn*

\LA
\LessonTitle{De Ezechiéle Prophéta}
\switchcolumn
\EN
\LessonTitle{Lesson from the book of Ezekiel}
\switchcolumn*

\LA
\Cite{Ezek 2:2-5}
\switchcolumn
\EN
\Cite{Ezek 2:2-5}
\switchcolumn*

\LA
\noindent Et audívi loquéntem ad me, et dicéntem: Fili hóminis, mitto ego te ad fílios Israël, ad gentes apostatríces quæ recessérunt a me; ipse et patres eórum prævaricáti sunt pactum meum usque ad diem hanc; et fílii dura fácie et indomábili corde sunt, ad quos ego mitto te, et dices ad eos: Hæc dicit Dóminus Deus: Si forte vel ipsi áudiant, et si forte quiéscant, quóniam domus exásperans est: et scient quia prophéta fúerit in médio eórum.\par
\switchcolumn
\EN
\noindent And the spirit entered into me after that he spoke to me, and he set me upon my feet: and I heard him speaking to me, And saying: Son of man, I send thee to the children of Israel, to a rebellious people, that hath revolted from me, they, and their fathers, have transgressed my covenant even unto this day. And they to whom I send thee are children of a hard face, and of an obstinate heart: and thou shalt say to them: Thus saith the Lord God: If so be they at least will hear, and if so be they will forbear, for they are a provoking house: and they shall know that there hath been a prophet in the midst of them.\par
\switchcolumn*

\LA
\begin{Resp}\Rsign Redémit pópulum suum et liberávit eum, et vénient et exsultábunt in monte Sion et gaudébunt de bonis Dómini super fruménto, vino et óleo, \Star{} Et ultra non esúrient.\end{Resp}
\begin{Resp}\Vsign Erítque ánima eórum quasi hortus irríguus.\end{Resp}
\begin{Resp}\Rsign Et ultra non esúrient.\end{Resp}
\switchcolumn
\EN
\begin{Resp}\Rsign He hath redeemed His people, and ransomed them therefore they shall come and sing in the height of Zion, and shall rejoice in the goodness of the Lord, for wheat, and for wine, and for oil \Star{} And they shall hunger no more.\end{Resp}
\begin{Resp}\Vsign And their soul shall be as a watered garden.\end{Resp}
\begin{Resp}\Rsign And they shall hunger no more.\end{Resp}
\switchcolumn*

\LA
\LessonHead{Lectio II}
\switchcolumn
\EN
\LessonHead{Lesson II}
\switchcolumn*

\LA
\Cite{Ezek 2:6-7}
\switchcolumn
\EN
\Cite{Ezek 2:6-7}
\switchcolumn*

\LA
\noindent Tu ergo, fili hóminis, ne tímeas eos, neque sermónes eórum métuas, quóniam incréduli et subversóres sunt tecum, et cum scorpiónibus hábitas. Verba eórum ne tímeas, et vultus eórum ne formídes, quia domus exásperans est. Loquéris ergo verba mea ad eos, si forte áudiant et quiéscant, quóniam irritatóres sunt.\par
\switchcolumn
\EN
\noindent And thou, O son of man, fear not, neither be thou afraid of their words: for thou art among unbelievers and destroyers, and thou dwellest with scorpions. Fear not their words, neither be thou dismayed at their looks: for they are a provoking house. And thou shalt speak my words to them, if perhaps they will hear, and forbear: for they provoke me to anger.\par
\switchcolumn*

\LA
\begin{Resp}\Rsign Angústiæ mihi sunt úndique, et quid éligam ignóro; \Star{} Mélius est mihi incídere in manus hóminum, quam derelínquere legem Dei mei.\end{Resp}
\begin{Resp}\Vsign Si enim hoc égero, mors mihi est; si autem non égero, non effúgiam manus vestras.\end{Resp}
\begin{Resp}\Rsign Mélius est mihi incídere in manus hóminum, quam derelínquere legem Dei mei.\end{Resp}
\switchcolumn
\EN
\begin{Resp}\Rsign I am straitened on every side, and know not what to choose. \Star{} It is better for me to fall into the hands of men, than to sin against the law of my God.\end{Resp}
\begin{Resp}\Vsign For if I do this thing, it is death unto me and if I do it not, I cannot escape your hands.\end{Resp}
\begin{Resp}\Rsign It is better for me to fall into the hands of men, than to sin against the law of my God.\end{Resp}
\switchcolumn*

\LA
\LessonHead{Lectio III}
\switchcolumn
\EN
\LessonHead{Lesson III}
\switchcolumn*

\LA
\Cite{Ezek 2:8-9}
\switchcolumn
\EN
\Cite{Ezek 2:8-9}
\switchcolumn*

\LA
\noindent Tu autem, fili hóminis, audi quæcúmque loquor ad te, et noli esse exásperans, sicut domus exasperátrix est: áperi os tuum, et cómede quæcúmque ego do tibi. Et vidi, et ecce manus missa ad me, in qua erat involútus liber: et expándit illum coram me, qui erat scriptus intus et foris, et scriptæ erant in eo lamentatiónes, et carmen, et væ.\par
\switchcolumn
\EN
\noindent But thou, O son of man, hear all that I say to thee: and do not thou provoke me, as that house provoketh me: open thy mouth, and eat what I give thee. And I looked, and behold, a hand was sent to me, wherein was a book rolled up: and he spread it before me, and it was written within and without: and there were written in it lamentations, and canticles, and woe.\par
\switchcolumn*

\LA
\begin{Resp}\Rsign Misit Dóminus Angelum suum et conclúsit ora leónum, \Star{} Et non contaminavérunt, quia coram eo injustítia invénta non est in me.\end{Resp}
\begin{Resp}\Vsign Misit Deus misericórdiam suam et veritátem suam: ánimam meam erípuit de médio catulórum leónum.\end{Resp}
\begin{Resp}\Rsign Et non contaminavérunt, quia coram eo injustítia invénta non est in me.\end{Resp}
\begin{Resp}\Vsign Glória Patri, et Fílio, \Star{} et Spirítui Sancto.\end{Resp}
\begin{Resp}\Rsign Et non contaminavérunt: quia coram eo injustítia invénta non est in me.\end{Resp}
\switchcolumn
\EN
\begin{Resp}\Rsign The Lord hath sent His angel, and hath shut the lions' mouths, that \Star{} They have not hurt me forasmuch as before Him innocency was found in me.\end{Resp}
\begin{Resp}\Vsign God hath sent forth His mercy and His truth, (and delivered) my soul from among the lions' whelps.\end{Resp}
\begin{Resp}\Rsign They have not hurt me forasmuch as before Him innocency was found in me.\end{Resp}
\begin{Resp}\Vsign Glory be to the Father, and to the Son, \Star{} and to the Holy Ghost.\end{Resp}
\begin{Resp}\Rsign They have not hurt me forasmuch as before Him innocency was found in me.\end{Resp}
\switchcolumn*

\end{paracol}

\Office{Feria III infra Hebdomadam I Novembris}{Tuesday of the first week of November}{}
\addcontentsline{toc}{subsection}{Feria III infra Hebdomadam I Novembris \textit{— Tuesday of the first week of November}}
\begin{paracol}{2}
\LA
\LessonHead{Lectio I}
\switchcolumn
\EN
\LessonHead{Lesson I}
\switchcolumn*

\LA
\LessonTitle{De Ezechiéle Prophéta}
\switchcolumn
\EN
\LessonTitle{Lesson from the book of Ezekiel}
\switchcolumn*

\LA
\Cite{Ezek 3:1-4}
\switchcolumn
\EN
\Cite{Ezek 3:1-4}
\switchcolumn*

\LA
\noindent Et dixit ad me: Fili hóminis, quodcúmque invéneris, cómede: cómede volúmen istud, et vadens lóquere ad fílios Israël. Et apérui os meum, et cibávit me volúmine illo: et dixit ad me: Fili hóminis, venter tuus cómedet, et víscera tua complebúntur volúmine isto quod ego do tibi. Et comédi illud, et factum est in ore meo sicut mel dulce. Et dixit ad me: Fili hóminis, vade ad domum Israël, et loquéris verba mea ad eos.\par
\switchcolumn
\EN
\noindent And he said to me: Son of man, eat all that thou shalt find: eat this book, and go speak to the children of Israel. And I opened my mouth, and he caused me to eat that book. And he said to me: Son of man, thy belly shall eat, and thy bowels shall be filled with this book, which I give thee. And I did eat it: and it was sweet as honey in my mouth. And he said to me: Son of man, go to the house of Israel, and thou shalt speak my words to them.\par
\switchcolumn*

\LA
\begin{Resp}\Rsign A fácie furóris tui, Deus, conturbáta est omnis terra: \Star{} Sed miserére, Dómine, et ne fácias consummatiónem.\end{Resp}
\begin{Resp}\Vsign Dómine, Dóminus noster, quam admirábile est nomen tuum!\end{Resp}
\begin{Resp}\Rsign Sed miserére, Dómine, et ne fácias consummatiónem.\end{Resp}
\switchcolumn
\EN
\begin{Resp}\Rsign Before the face of thine anger, O God, the whole earth is troubled; \Star{} But Thou, O Lord, have mercy, and make not an end utterly.\end{Resp}
\begin{Resp}\Vsign Lord our Ruler, how excellent is thy Name in all the earth!\end{Resp}
\begin{Resp}\Rsign But Thou, O Lord, have mercy, and make not an end utterly.\end{Resp}
\switchcolumn*

\LA
\LessonHead{Lectio II}
\switchcolumn
\EN
\LessonHead{Lesson II}
\switchcolumn*

\LA
\Cite{Ezek 3:5-9}
\switchcolumn
\EN
\Cite{Ezek 3:5-9}
\switchcolumn*

\LA
\noindent Non enim ad pópulum profúndi sermónis et ignótæ linguæ, tu mítteris ad domum Israël; neque ad pópulos multos profúndi sermónis et ignótæ linguæ, quorum non possis audíre sermónes; et, si ad illos mitteréris, ipsi audírent te; domus autem Israël nolunt audíre te, quia nolunt audíre me; omnis quippe domus Israël attríta fronte est, et duro corde. Ecce dedi fáciem tuam valentiórem faciébus eórum, et frontem tuam duriórem fróntibus eórum: ut adamántem et ut sílicem dedi fáciem tuam; ne tímeas eos neque métuas a fácie eórum, quia domus exásperans est.\par
\switchcolumn
\EN
\noindent For thou art not sent to a people of a profound speech, and of an unknown tongue, but to the house of Israel: Nor to many nations of a strange speech, and of an unknown tongue, whose words thou canst not understand: and if thou wert sent to them, they would hearken to thee. But the house of Israel will not hearken to thee: because they will not hearken to me: for all the house of Israel are of a hard forehead and an obstinate heart. Behold I have made thy face stronger than their faces: and thy forehead harder than their foreheads. I have made thy face like an adamant and like flint: fear them not, neither be thou dismayed at their presence: for they are a provoking house.\par
\switchcolumn*

\LA
\begin{Resp}\Rsign Civitátem istam tu circúmda, Dómine: et Angeli tui custódiant muros ejus. \Star{} Exáudi, Dómine, pópulum tuum cum misericórdia.\end{Resp}
\begin{Resp}\Vsign Avertátur furor tuus, Dómine, a pópulo tuo et a civitáte sancta tua.\end{Resp}
\begin{Resp}\Rsign Exáudi, Dómine, pópulum tuum cum misericórdia.\end{Resp}
\switchcolumn
\EN
\begin{Resp}\Rsign Fence Thou this city; O Lord, and let thine angels keep the walls thereof. \Star{} O Lord, hearken unto thy people with mercy.\end{Resp}
\begin{Resp}\Vsign O Lord, let thine anger be turned away from thy people, and from thine holy city.\end{Resp}
\begin{Resp}\Rsign O Lord, hearken unto thy people with mercy.\end{Resp}
\switchcolumn*

\LA
\LessonHead{Lectio III}
\switchcolumn
\EN
\LessonHead{Lesson III}
\switchcolumn*

\LA
\Cite{Ezek 3:10-13}
\switchcolumn
\EN
\Cite{Ezek 3:10-13}
\switchcolumn*

\LA
\noindent Et dixit ad me: fili hóminis, omnes sermónes meos, quos ego loquor ad te, assúme in corde tuo et áuribus tuis audi; et vade, ingrédere ad transmigratiónem, ad fílios pópuli tui, et loquéris ad eos; et dices eis: Hæc dicit Dóminus Deus: Si forte áudiant et quiéscant. Et assúmpsit me spíritus, et audívi post me vocem commotiónis magnæ: Benedícta glória Dómini de loco suo; et vocem alárum animálium percutiéntium álteram ad álteram, et vocem rotárum sequéntium animália, et vocem commotiónis magnæ.\par
\switchcolumn
\EN
\noindent And he said to me: Son of man, receive in thy heart, and hear with thy ears, all the words that I speak to thee: And go get thee in to them of the captivity, to the children of thy people, and thou shalt speak to them, and shalt say to them: Thus saith the Lord: If so be they will hear and will forbear. And the spirit took me up, and I heard behind me the voice of a great commotion, saying: Blessed be the glory of the Lord, from his place. And the noise of the wings of the living creatures striking one against another, and the noise of the wheels following the living creatures, and the noise of a great commotion.\par
\switchcolumn*

\LA
\begin{Resp}\Rsign Genti peccatríci, pópulo pleno peccáto miserére, \Star{} Dómine Deus.\end{Resp}
\begin{Resp}\Vsign Esto placábilis super nequítiam pópuli tui.\end{Resp}
\begin{Resp}\Rsign Dómine Deus.\end{Resp}
\begin{Resp}\Vsign Glória Patri, et Fílio, \Star{} et Spirítui Sancto.\end{Resp}
\begin{Resp}\Rsign Dómine Deus.\end{Resp}
\switchcolumn
\EN
\begin{Resp}\Rsign O Lord God! have mercy upon \Star{} The sinful nation, upon the people laden with iniquity.\end{Resp}
\begin{Resp}\Vsign Let it repent thee concerning the transgression of thy people.\end{Resp}
\begin{Resp}\Rsign The people laden with iniquity.\end{Resp}
\begin{Resp}\Vsign Glory be to the Father, and to the Son, \Star{} and to the Holy Ghost.\end{Resp}
\begin{Resp}\Rsign The people laden with iniquity.\end{Resp}
\switchcolumn*

\end{paracol}

\Office{Feria IV infra Hebdomadam I Novembris}{Wednesday of the first week of November}{}
\addcontentsline{toc}{subsection}{Feria IV infra Hebdomadam I Novembris \textit{— Wednesday of the first week of November}}
\begin{paracol}{2}
\LA
\LessonHead{Lectio I}
\switchcolumn
\EN
\LessonHead{Lesson I}
\switchcolumn*

\LA
\LessonTitle{De Ezechiéle Prophéta}
\switchcolumn
\EN
\LessonTitle{Lesson from the book of Ezekiel}
\switchcolumn*

\LA
\Cite{Ezek 7:1-4}
\switchcolumn
\EN
\Cite{Ezek 7:1-4}
\switchcolumn*

\LA
\noindent Et factus est sermo Dómini ad me, dicens: Et tu, fili hóminis, hæc dicit Dóminus Deus terræ Israël: Finis venit, venit finis super quátuor plagas terræ. Nunc finis super te, et immíttam furórem meum in te: et judicábo te juxta vias tuas, et ponam contra te omnes abominatiónes tuas. Et non parcet óculus meus super te, et non miserébor; sed vias tuas ponam super te, et abominatiónes tuæ in médio tui erunt, et sciétis quia ego Dóminus.\par
\switchcolumn
\EN
\noindent And the word of the Lord came to me, saying: And thou son of man, thus saith the Lord God to the land of Israel: The end is come, the end is come upon the four quarters of the land. Now is an end come upon thee, and I will send my wrath upon thee, and I will judge thee according to thy ways: and I will set all thy abominations against thee. And my eye shall not spare thee, and I will shew thee no pity: but I will lay thy ways upon thee, and thy abominations shall be in the midst of thee: and you shall know that I am the Lord.\par
\switchcolumn*

\LA
\begin{Resp}\Rsign Indicábo tibi, homo, quid sit bonum aut quid Dóminus requírat a te: \Star{} Fácere judícium et justítiam et sollícitum ambuláre cum Deo tuo.\end{Resp}
\begin{Resp}\Vsign Spera in Dómino, et fac bonitátem, et inhábita terram.\end{Resp}
\begin{Resp}\Rsign Fácere judícium et justítiam et sollícitum ambuláre cum Deo tuo.\end{Resp}
\switchcolumn
\EN
\begin{Resp}\Rsign I will show thee, O man, what is good, and what doth the Lord require of thee, but to \Star{} Do justice and judgment, and to walk humbly with thy God.\end{Resp}
\begin{Resp}\Vsign Trust in the Lord, and do good, and dwell in the land.\end{Resp}
\begin{Resp}\Rsign Do justice and judgment, and walk humbly with thy God.\end{Resp}
\switchcolumn*

\LA
\LessonHead{Lectio II}
\switchcolumn
\EN
\LessonHead{Lesson II}
\switchcolumn*

\LA
\Cite{Ezek 7:5-9}
\switchcolumn
\EN
\Cite{Ezek 7:5-9}
\switchcolumn*

\LA
\noindent Hæc dicit Dóminus Deus: Afflíctio una, afflíctio ecce venit. Finis venit, venit finis, evigilávit advérsum te, ecce venit: venit contrítio super te, qui hábitas in terra; venit tempus, prope est dies occisiónis, et non glóriæ móntium. Nunc de propínquo effúndam iram meam super te, et complébo furórem meum in te: et judicábo te juxta vias tuas, et impónam tibi ómnia scélera tua, et non parcet óculus meus, nec miserébor; sed vias tuas impónam tibi, et abominatiónes tuæ in médio tui erunt, et sciétis quia ego sum Dóminus percútiens.\par
\switchcolumn
\EN
\noindent Thus saith the Lord God: One affliction, behold an affliction is come. An end is come, the end is come, it hath awaked against thee: behold it is come. Destruction is come upon thee that dwellest in the land: the time is come, the day of slaughter is near, and not of the joy of mountains. Now very shortly I will pour out my wrath upon thee, and I will accomplish my anger in thee: and I will judge thee according to thy ways, and I will lay upon thee all thy crimes. And my eye shall not spare, neither will I shew mercy: but I will lay thy ways upon thee, and thy abominations shall be in the midst of thee: and you shall know that I am the Lord that strike.\par
\switchcolumn*

\LA
\begin{Resp}\Rsign Angústiæ mihi sunt úndique, et quid éligam ignóro; \Star{} Mélius est mihi incídere in manus hóminum, quam derelínquere legem Dei mei.\end{Resp}
\begin{Resp}\Vsign Si enim hoc égero, mors mihi est; si autem non égero, non effúgiam manus vestras.\end{Resp}
\begin{Resp}\Rsign Mélius est mihi incídere in manus hóminum, quam derelínquere legem Dei mei.\end{Resp}
\switchcolumn
\EN
\begin{Resp}\Rsign I am straitened on every side, and know not what to choose. \Star{} It is better for me to fall into the hands of men, than to sin against the law of my God.\end{Resp}
\begin{Resp}\Vsign For if I do this thing, it is death unto me and if I do it not, I cannot escape your hands.\end{Resp}
\begin{Resp}\Rsign It is better for me to fall into the hands of men, than to sin against the law of my God.\end{Resp}
\switchcolumn*

\LA
\LessonHead{Lectio III}
\switchcolumn
\EN
\LessonHead{Lesson III}
\switchcolumn*

\LA
\Cite{Ezek 7:10-13}
\switchcolumn
\EN
\Cite{Ezek 7:10-13}
\switchcolumn*

\LA
\noindent Ecce dies, ecce venit; egréssa est contrítio, flóruit virga, germinávit supérbia, iníquitas surréxit in virga impietátis; non ex eis et non ex pópulo neque ex sónitu eórum, et non erit réquies in eis. Venit tempus, appropinquávit dies: qui emit non lætétur, et qui vendit, non lúgeat: quia ira super omnem pópulum ejus; quia qui vendit, ad id quod véndidit non revertétur: et adhuc in vivéntibus vita eórum.\par
\switchcolumn
\EN
\noindent Behold the day, behold it is come: destruction is gone forth, the rod hath blossomed, pride hath budded. Iniquity is risen up into a rod of impiety: nothing of them shall remain, nor of their people, nor of the noise of them: and there shall be no rest among them. The time is come, the day is at hand: let not the buyer rejoice: nor the seller mourn: for wrath is upon all the people thereof. For the seller shall not return to that which he hath sold, although their life be yet among the living. For the vision which regardeth all the multitude thereof, shall not go back: neither shall man be strengthened in the iniquity of his life.\par
\switchcolumn*

\LA
\begin{Resp}\Rsign Misit Dóminus Angelum suum et conclúsit ora leónum, \Star{} Et non contaminavérunt, quia coram eo injustítia invénta non est in me.\end{Resp}
\begin{Resp}\Vsign Misit Deus misericórdiam suam et veritátem suam: ánimam meam erípuit de médio catulórum leónum.\end{Resp}
\begin{Resp}\Rsign Et non contaminavérunt, quia coram eo injustítia invénta non est in me.\end{Resp}
\begin{Resp}\Vsign Glória Patri, et Fílio, \Star{} et Spirítui Sancto.\end{Resp}
\begin{Resp}\Rsign Et non contaminavérunt: quia coram eo injustítia invénta non est in me.\end{Resp}
\switchcolumn
\EN
\begin{Resp}\Rsign The Lord hath sent His angel, and hath shut the lions' mouths, that \Star{} They have not hurt me forasmuch as before Him innocency was found in me.\end{Resp}
\begin{Resp}\Vsign God hath sent forth His mercy and His truth, (and delivered) my soul from among the lions' whelps.\end{Resp}
\begin{Resp}\Rsign They have not hurt me forasmuch as before Him innocency was found in me.\end{Resp}
\begin{Resp}\Vsign Glory be to the Father, and to the Son, \Star{} and to the Holy Ghost.\end{Resp}
\begin{Resp}\Rsign They have not hurt me forasmuch as before Him innocency was found in me.\end{Resp}
\switchcolumn*

\end{paracol}

\Office{Feria V infra Hebdomadam I Novembris}{Thursday of the first week of November}{}
\addcontentsline{toc}{subsection}{Feria V infra Hebdomadam I Novembris \textit{— Thursday of the first week of November}}
\begin{paracol}{2}
\LA
\LessonHead{Lectio I}
\switchcolumn
\EN
\LessonHead{Lesson I}
\switchcolumn*

\LA
\LessonTitle{De Ezechiéle Prophéta}
\switchcolumn
\EN
\LessonTitle{Lesson from the book of Ezekiel}
\switchcolumn*

\LA
\Cite{Ezek 13:1-6}
\switchcolumn
\EN
\Cite{Ezek 13:1-6}
\switchcolumn*

\LA
\noindent Et factus est sermo Dómini ad me, dicens: Fili hóminis, vaticináre ad prophétas Israël, qui prophétant, et dices prophetántibus de corde suo: Audíte verbum Dómini. Hæc dicit Dóminus Deus: Væ prophétis insipiéntibus, qui sequúntur spíritum suum, et nihil vident! Quasi vulpes in desértis prophétæ tui, Israël, erant. Non ascendístis ex advérso, neque opposuístis murum pro domo Israël, ut starétis in prǽlio in die Dómini. Vident vana, et divínant mendácium, dicéntes: Ait Dóminus, cum Dóminus non míserit eos, et perseveravérunt confirmáre sermónem.\par
\switchcolumn
\EN
\noindent And the word of the Lord came to me, saying: Son of man, prophesy thou against the prophets of Israel that prophesy: and thou shalt say to them that prophesy out of their own heart: Hear ye the word of the Lord: Thus saith the Lord God: Woe to the foolish prophets that follow their own spirit, and see nothing. Thy prophets, O Israel, were like foxes in the deserts. You have not gone up to face the enemy, nor have you set up a wall for the house of Israel, to stand in battle in the day of the Lord. They see vain things, and they foretell lies, saying: The Lord saith: whereas the Lord hath not sent them: and they have persisted to confirm what they have said.\par
\switchcolumn*

\LA
\begin{Resp}\Rsign Vidi Dóminum sedéntem super sólium excélsum et elevátum, et plena erat omnis terra majestáte ejus: \Star{} Et ea, quæ sub ipso erant, replébant templum.\end{Resp}
\begin{Resp}\Vsign Séraphim stabant super illud: sex alæ uni, et sex alæ álteri.\end{Resp}
\begin{Resp}\Rsign Et ea, quæ sub ipso erant, replébant templum.\end{Resp}
\switchcolumn
\EN
\begin{Resp}\Rsign I saw the Lord sitting upon a throne, high and lifted up, and the whole earth was full of His glory; \Star{} And His train filled the temple.\end{Resp}
\begin{Resp}\Vsign Above it stood the Seraphim each one had six wings.\end{Resp}
\begin{Resp}\Rsign And His train filled the temple.\end{Resp}
\switchcolumn*

\LA
\LessonHead{Lectio II}
\switchcolumn
\EN
\LessonHead{Lesson II}
\switchcolumn*

\LA
\Cite{Ezek 13:7-10}
\switchcolumn
\EN
\Cite{Ezek 13:7-10}
\switchcolumn*

\LA
\noindent Numquid non visiónem cassam vidístis, et divinatiónem mendácem locúti estis, et dícitis: Ait Dóminus, cum ego non sim locútus? Proptérea hæc dicit Dóminus Deus: Quia locúti estis vana, et vidístis mendácium, ídeo ecce ego ad vos, dicit Dóminus Deus. Et erit manus mea super prophétas, qui vident vana et divínant mendácium: in consílio pópuli mei non erunt, et in scriptúra domus Israël non scribéntur, nec in terram Israël ingrediéntur, et sciétis quia ego Dóminus Deus. Eo quod decéperint pópulum meum, dicéntes: Pax, et non est pax, et ipse ædificábat paríetem, illi autem liniébant eum luto absque páleis.\par
\switchcolumn
\EN
\noindent Have you not seen a vain vision and spoken a lying divination: and you say: The Lord saith: whereas I have not spoken. Therefore thus saith the Lord God: Because you have spoken vain things, and have seen lies: therefore behold I come against you, saith the Lord God. And my hand shall be upon the prophets that see vain things, and that divine lies: they shall not be in the council of my people, nor shall they be written in the writing of the house of Israel, neither shall they enter into the land of Israel, and you shall know that I am the Lord God. Because they have deceived my people, saying: Peace, and there is no peace: and the people built up a wall, and they daubed it with dirt without straw.\par
\switchcolumn*

\LA
\begin{Resp}\Rsign Aspice, Dómine, de sede sancta tua, et cógita de nobis: inclína, Deus meus, aurem tuam et audi: \Star{} Aperi óculos tuos et vide tribulatiónem nostram.\end{Resp}
\begin{Resp}\Vsign Qui regis Israël, inténde, qui dedúcis velut ovem Joseph.\end{Resp}
\begin{Resp}\Rsign Aperi óculos tuos et vide tribulatiónem nostram.\end{Resp}
\switchcolumn
\EN
\begin{Resp}\Rsign Look down, O Lord, from the dwelling-place of thine holiness, and take thought for us. O my God, incline thine ear, and hear. \Star{} Open thine eyes, and behold our desolation.\end{Resp}
\begin{Resp}\Vsign Give ear, O Shepherd of Israel, Thou That leadest Joseph like a flock.\end{Resp}
\begin{Resp}\Rsign Open thine eyes, and behold our desolation.\end{Resp}
\switchcolumn*

\LA
\LessonHead{Lectio III}
\switchcolumn
\EN
\LessonHead{Lesson III}
\switchcolumn*

\LA
\Cite{Ezek 13:11-14}
\switchcolumn
\EN
\Cite{Ezek 13:11-14}
\switchcolumn*

\LA
\noindent Dic ad eos qui líniunt absque temperatúra, quod casúrus sit; erit enim imber inúndans et dabo lápides prægrándes désuper irruéntes et ventum procéllæ dissipántem. Síquidem ecce cécidit páries; numquid non dicétur vobis: Ubi est litúra quam linístis? Proptérea hæc dicit Dóminus Deus: Et erúmpere fáciam spíritum tempestátum in indignatióne mea, et imber inúndans in furóre meo erit, et lápides grandes in ira in consumptiónem, et déstruam paríetem, quem linístis absque temperaménto.\par
\switchcolumn
\EN
\noindent Say to them that daub without tempering, that it shall fall: for there shall be an overflowing shower, and I will cause great hailstones to fall violently from above, and a stormy wind to throw it down. Behold, when the wall is fallen: shall it not be said to you: Where is the daubing wherewith you have daubed it? Therefore thus saith the Lord God: Lo, I will cause a stormy wind to break forth in my indignation, and there shall be an overflowing shower in my anger: and great hailstones in my wrath to consume. And I will break down the wall that you have daubed with untempered mortar.\par
\switchcolumn*

\LA
\begin{Resp}\Rsign Aspice, Dómine, quia facta est desoláta cívitas plena divítiis, sedet in tristítia dómina géntium: \Star{} Non est qui consolétur eam, nisi tu, Deus noster.\end{Resp}
\begin{Resp}\Vsign Plorans plorávit in nocte, et lácrimæ ejus in maxíllis ejus.\end{Resp}
\begin{Resp}\Rsign Non est qui consolétur eam, nisi tu, Deus noster.\end{Resp}
\begin{Resp}\Vsign Glória Patri, et Fílio, \Star{} et Spirítui Sancto.\end{Resp}
\begin{Resp}\Rsign Non est qui consolétur eam, nisi tu, Deus noster.\end{Resp}
\switchcolumn
\EN
\begin{Resp}\Rsign Consider, O Lord, how that the city sitteth solitary that was full of riches; how is she become as a widow, she that was great among the nations; \Star{} She hath none to comfort her, save thee, O our God.\end{Resp}
\begin{Resp}\Vsign She weepeth sore in the night, and her tears are on her cheeks.\end{Resp}
\begin{Resp}\Rsign She hath none to comfort her, save thee, O our God.\end{Resp}
\begin{Resp}\Vsign Glory be to the Father, and to the Son, \Star{} and to the Holy Ghost.\end{Resp}
\begin{Resp}\Rsign She hath none to comfort her, save thee, O our God.\end{Resp}
\switchcolumn*

\end{paracol}

\Office{Feria VI infra Hebdomadam I Novembris}{Friday of the first week of November}{}
\addcontentsline{toc}{subsection}{Feria VI infra Hebdomadam I Novembris \textit{— Friday of the first week of November}}
\begin{paracol}{2}
\LA
\LessonHead{Lectio I}
\switchcolumn
\EN
\LessonHead{Lesson I}
\switchcolumn*

\LA
\LessonTitle{De Ezechiéle Prophéta}
\switchcolumn
\EN
\LessonTitle{Lesson from the book of Ezekiel}
\switchcolumn*

\LA
\Cite{Ezek 15:1-5}
\switchcolumn
\EN
\Cite{Ezek 15:1-5}
\switchcolumn*

\LA
\noindent Et factus est sermo Dómini ad me, dicens: Fili hóminis, quid fiet de ligno vitis, ex ómnibus lignis némorum, quæ sunt inter ligna silvárum? Numquid tollétur de ea lignum, ut fiat opus, aut fabricábitur de ea paxíllus, ut depéndeat in eo quodcúmque vas? Ecce igni datum est in escam: utrámque partem ejus consúmpsit ignis, et medíetas ejus redácta est in favíllam. Numquid útile erit ad opus? Etiam cum esset íntegrum, non erat aptum ad opus: quanto magis cum illud ignis devoráverit et combússerit: nihil ex eo fiet óperis?\par
\switchcolumn
\EN
\noindent And the word of the Lord came to me, saying: Son of man, what shall be made of the wood of the vine, out of all the trees of the woods that are among the trees of the forests? Shall wood be taken of if, to do any work, or shall a pin be made of it for any vessel to hang thereon? Behold it is cast into the fire for fuel: the fire hath consumed both ends thereof, and the midst thereof is reduced to ashes: shall it be useful for any work? Even when it was whole it was not fit for work: how much less, when the fire hath devoured and consumed it, shall any work be made of it?\par
\switchcolumn*

\LA
\begin{Resp}\Rsign Super muros tuos, Jerúsalem, constítui custódes; \Star{} Tota die et nocte non tacébunt laudáre nomen Dómini.\end{Resp}
\begin{Resp}\Vsign Prædicábunt pópulis fortitúdinem meam, et annuntiábunt géntibus glóriam meam.\end{Resp}
\begin{Resp}\Rsign Tota die et nocte non tacébunt laudáre nomen Dómini.\end{Resp}
\switchcolumn
\EN
\begin{Resp}\Rsign I have set watchmen upon thy walls, O Jerusalem; \Star{} Which shall never hold their peace day nor night, to praise the name of the Lord.\end{Resp}
\begin{Resp}\Vsign They shall proclaim My might unto the nations, and declare My glory unto the Gentiles.\end{Resp}
\begin{Resp}\Rsign Which shall never hold their peace day nor night, to praise the name of the Lord.\end{Resp}
\switchcolumn*

\LA
\LessonHead{Lectio II}
\switchcolumn
\EN
\LessonHead{Lesson II}
\switchcolumn*

\LA
\Cite{Ezek 15:6-8}
\switchcolumn
\EN
\Cite{Ezek 15:6-8}
\switchcolumn*

\LA
\noindent Proptérea hæc dicit Dóminus Deus: Quómodo lignum vitis inter ligna silvárum, quod dedi igni ad devorándum, sic tradam habitatóres Jerúsalem. Et ponam fáciem meam in eos: de igne egrediéntur, et ignis consúmet eos: et sciétis quia ego Dóminus, cum posúero fáciem meam in eos, et dédero terram ínviam et desolátam, eo quod prævaricatóres exstíterint, dicit Dóminus Deus.\par
\switchcolumn
\EN
\noindent Therefore thus saith the Lord God: As the vine tree among the trees of the forests which I have given to the fire to be consumed, so will I deliver up the inhabitants of Jerusalem. And I will set my face against them: they shall go out from fire, and fire shall consume them and you shall know that I am the Lord, when I shall have set my face against them. And I shall have made their land a wilderness, and desolate, because they have been transgressors, saith the Lord God.\par
\switchcolumn*

\LA
\begin{Resp}\Rsign Muro tuo inexpugnábili circumcínge nos, Dómine, et armis tuæ poténtiæ prótege nos semper: \Star{} Líbera, Dómine, Deus Israël, clamántes ad te.\end{Resp}
\begin{Resp}\Vsign Erue nos in mirabílibus tuis, et da glóriam nómini tuo.\end{Resp}
\begin{Resp}\Rsign Líbera, Dómine, Deus Israël, clamántes ad te.\end{Resp}
\switchcolumn
\EN
\begin{Resp}\Rsign Hedge us about with thy wall that cannot be broken down, O Lord, and shield us continually with the arms of thy might. \Star{} O Lord God of Israel, deliver them that cry unto thee.\end{Resp}
\begin{Resp}\Vsign Deliver us also according to thy marvellous works, and give glory to thy Name.\end{Resp}
\begin{Resp}\Rsign O Lord God of Israel, deliver them that cry unto thee.\end{Resp}
\switchcolumn*

\LA
\LessonHead{Lectio III}
\switchcolumn
\EN
\LessonHead{Lesson III}
\switchcolumn*

\LA
\Cite{Ezek 16:1-5}
\switchcolumn
\EN
\Cite{Ezek 16:1-5}
\switchcolumn*

\LA
\noindent Et factus est sermo Dómini ad me, dicens: Fili hóminis, notas fac Jerúsalem abominatiónes suas, et dices: Hæc dicit Dóminus Deus Jerúsalem: Radix tua et generátio tua de terra Chánaan, pater tuus Amorrhǽus, et mater tua Cethǽa. Et quando nata es, in die ortus tui, non est præcísus umbilícus tuus, et aqua non es lota in salútem, nec sale salíta, nec involúta pannis. Non pepércit super te óculus, ut fáceret tibi unum de his, misértus tui, sed projécta es super fáciem terræ in abjectióne ánimæ tuæ, in die qua nata es.\par
\switchcolumn
\EN
\noindent And the word of the Lord came to me, saying: Son of man, make known to Jerusalem her abominations. And thou shalt say: Thus saith the Lord God to Jerusalem: thy root, and thy nativity is of the land of Chanaan, thy father was an Amorrhite, and thy mother a Cethite. And when thou wast born, in the day of thy nativity thy navel was not cut, neither wast thou washed with water for thy health, nor salted with salt, nor swaddled with clouts. No eye had pity on thee to do any of these things for thee, out of compassion to thee: but thou wast cast out upon the face of the earth in the abjection of thy soul, in the day that thou wast born.\par
\switchcolumn*

\LA
\begin{Resp}\Rsign Sustinúimus pacem, et non venit: quæsívimus bona, et ecce turbátio: cognóvimus, Dómine, peccáta nostra; \Star{} Non in perpétuum obliviscáris nos.\end{Resp}
\begin{Resp}\Vsign Peccávimus, ímpie géssimus, iniquitátem fécimus, Dómine, in omnem justítiam tuam.\end{Resp}
\begin{Resp}\Rsign Non in perpétuum obliviscáris nos.\end{Resp}
\begin{Resp}\Vsign Glória Patri, et Fílio, \Star{} et Spirítui Sancto.\end{Resp}
\begin{Resp}\Rsign Non in perpétuum obliviscáris nos.\end{Resp}
\switchcolumn
\EN
\begin{Resp}\Rsign We looked for peace, and it came not we asked for good, and behold trouble. We acknowledge, O Lord, our wickedness. \Star{} Forget us not for ever.\end{Resp}
\begin{Resp}\Vsign O Lord (our God) we have sinned, we have done ungodly, we have dealt unrighteously in all thine ordinances.\end{Resp}
\begin{Resp}\Rsign Forget us not for ever.\end{Resp}
\begin{Resp}\Vsign Glory be to the Father, and to the Son, \Star{} and to the Holy Ghost.\end{Resp}
\begin{Resp}\Rsign Forget us not for ever.\end{Resp}
\switchcolumn*

\end{paracol}

\Office{Sabbato infra Hebdomadam I Novembris}{Saturday of the first week of November}{}
\addcontentsline{toc}{subsection}{Sabbato infra Hebdomadam I Novembris \textit{— Saturday of the first week of November}}
\begin{paracol}{2}
\LA
\LessonHead{Lectio I}
\switchcolumn
\EN
\LessonHead{Lesson I}
\switchcolumn*

\LA
\LessonTitle{De Ezechiéle Prophéta}
\switchcolumn
\EN
\LessonTitle{Lesson from the book of Ezekiel}
\switchcolumn*

\LA
\Cite{Ezek 19:1-7}
\switchcolumn
\EN
\Cite{Ezek 19:1-7}
\switchcolumn*

\LA
\noindent Et tu assúme planctum super príncipes Israël, et dices: Quare mater tua leǽna inter leónes cubávit? in médio leunculórum enutrívit cátulos suos? Et edúxit unum de leúnculis suis, et leo factus est: et dídicit cápere prædam, hominémque comédere. Et audiérunt de eo gentes: et non absque vulnéribus suis cepérunt eum, et adduxérunt eum in caténis in terram Ægýpti. Quæ cum vidísset quóniam infirmáta est, et périit exspectátio ejus, tulit unum de leúnculis suis, leónem constítuit eum, qui incedébat inter leónes, et factus est leo: et dídicit prædam cápere et hómines devoráre, dídicit víduas fácere, et civitátes eórum in desértum addúcere: et desoláta est terra et plenitúdo ejus a voce rugítus illíus.\par
\switchcolumn
\EN
\noindent Moreover take thou up a lamentation for the princes of Israel, And say: Why did thy mother the lioness lie down among the lions, and bring up her whelps in the midst of young lions? And she brought out one of her whelps, and he became a lion: and he learned to catch the prey, and to devour men. And the nations heard of him, and took him, but not without receiving wounds: and they brought him in chains into the land of Egypt. But she seeing herself weakened, and that her hope was lost, took one of her young lions, and set him up for a lion. And he went up and down among the lions, and became a lion: and he learned to catch the prey, and to devour men. He learned to make widows, and to lay waste their cities: and the land became desolate, and the fulness thereof by the noise of his roaring.\par
\switchcolumn*

\LA
\begin{Resp}\Rsign Laudábilis pópulus, \Star{} Quem Dóminus exercítuum benedíxit dicens: Opus mánuum meárum tu es, heréditas mea Israël.\end{Resp}
\begin{Resp}\Vsign Beáta gens, cujus est Dóminus Deus, pópulus eléctus in hereditátem.\end{Resp}
\begin{Resp}\Rsign Quem Dóminus exercítuum benedíxit dicens: Opus mánuum meárum tu es, heréditas mea Israël.\end{Resp}
\switchcolumn
\EN
\begin{Resp}\Rsign Blessed is the people \Star{} Whom the Lord of hosts hath blessed, saying O Israel thou art the work of Mine own hands, thou art Mine own inheritance.\end{Resp}
\begin{Resp}\Vsign Blessed is the nation whose God is the Lord, and the people whom He hath chosen for His own inheritance.\end{Resp}
\begin{Resp}\Rsign Whom the Lord of hosts hath blessed, saying O Israel thou art the work of Mine own hands, thou art Mine own inheritance.\end{Resp}
\switchcolumn*

\LA
\LessonHead{Lectio II}
\switchcolumn
\EN
\LessonHead{Lesson II}
\switchcolumn*

\LA
\Cite{Ezek 19:8-11}
\switchcolumn
\EN
\Cite{Ezek 19:8-11}
\switchcolumn*

\LA
\noindent Et convenérunt advérsus eum gentes úndique de provínciis, et expandérunt super eum rete suum, in vulnéribus eárum captus est, et misérunt eum in cáveam, in caténis adduxérunt eum ad regem Babylónis, miserúntque eum in cárcerem, ne audirétur vox ejus ultra super montes Israël. Mater tua quasi vínea in sánguine tuo super aquam plantáta est; fructus ejus et frondes ejus crevérunt ex aquis multis, et factæ sunt ei virgæ sólidæ in sceptra dominántium, et exaltáta est statúra ejus inter frondes, et vidit altitúdinem suam in multitúdine pálmitum suórum.\par
\switchcolumn
\EN
\noindent And the nations came together against him on every side out of the provinces, and they spread their net over him, in their wounds he was taken. And they put him into a cage, they brought him in chains to the king of Babylon: and they cast him into prison, that his voice should no more be heard upon the mountains of Israel. thy mother is like a vine in thy blood planted by the water: her fruit and her branches have grown out of many waters. And she hath strong rods to make sceptres for them that bear rule, and her stature was exalted among the branches: and she saw her height in the multitude of her branches.\par
\switchcolumn*

\LA
\begin{Resp}\Rsign Angústiæ mihi sunt úndique, et quid éligam ignóro; \Star{} Mélius est mihi incídere in manus hóminum, quam derelínquere legem Dei mei.\end{Resp}
\begin{Resp}\Vsign Si enim hoc égero, mors mihi est; si autem non égero, non effúgiam manus vestras.\end{Resp}
\begin{Resp}\Rsign Mélius est mihi incídere in manus hóminum, quam derelínquere legem Dei mei.\end{Resp}
\switchcolumn
\EN
\begin{Resp}\Rsign I am straitened on every side, and know not what to choose. \Star{} It is better for me to fall into the hands of men, than to sin against the law of my God.\end{Resp}
\begin{Resp}\Vsign For if I do this thing, it is death unto me and if I do it not, I cannot escape your hands.\end{Resp}
\begin{Resp}\Rsign It is better for me to fall into the hands of men, than to sin against the law of my God.\end{Resp}
\switchcolumn*

\LA
\LessonHead{Lectio III}
\switchcolumn
\EN
\LessonHead{Lesson III}
\switchcolumn*

\LA
\Cite{Ezek 19:12-14}
\switchcolumn
\EN
\Cite{Ezek 19:12-14}
\switchcolumn*

\LA
\noindent Et evúlsa est in ira, in terrámque projécta, et ventus urens siccávit fructum ejus; marcuérunt et arefáctæ sunt virgæ róboris ejus, ignis comédit eam; et nunc transplantáta est in desértum, in terra ínvia et sitiénti, et egréssus est ignis de virga ramórum ejus, qui fructum ejus comédit, et non fuit in ea virga fortis, sceptrum dominántium. Planctus est, et erit in planctum.\par
\switchcolumn
\EN
\noindent But she was plucked up in wrath, and cast on the ground, and the burning wind dried up her fruit: her strong rods are withered, and dried up: the fire hath devoured her. And now she is transplanted into the desert, in a land not passable, and dry. And a fire is gone out from a rod of her branches, which hath devoured her fruit: so that she now hath no strong rod, to be a sceptre of rulers. This is a lamentation, and it shall be for a lamentation.\par
\switchcolumn*

\LA
\begin{Resp}\Rsign Misit Dóminus Angelum suum et conclúsit ora leónum, \Star{} Et non contaminavérunt, quia coram eo injustítia invénta non est in me.\end{Resp}
\begin{Resp}\Vsign Misit Deus misericórdiam suam et veritátem suam: ánimam meam erípuit de médio catulórum leónum.\end{Resp}
\begin{Resp}\Rsign Et non contaminavérunt, quia coram eo injustítia invénta non est in me.\end{Resp}
\begin{Resp}\Vsign Glória Patri, et Fílio, \Star{} et Spirítui Sancto.\end{Resp}
\begin{Resp}\Rsign Et non contaminavérunt: quia coram eo injustítia invénta non est in me.\end{Resp}
\switchcolumn
\EN
\begin{Resp}\Rsign The Lord hath sent His angel, and hath shut the lions' mouths, that \Star{} They have not hurt me forasmuch as before Him innocency was found in me.\end{Resp}
\begin{Resp}\Vsign God hath sent forth His mercy and His truth, (and delivered) my soul from among the lions' whelps.\end{Resp}
\begin{Resp}\Rsign They have not hurt me forasmuch as before Him innocency was found in me.\end{Resp}
\begin{Resp}\Vsign Glory be to the Father, and to the Son, \Star{} and to the Holy Ghost.\end{Resp}
\begin{Resp}\Rsign They have not hurt me forasmuch as before Him innocency was found in me.\end{Resp}
\switchcolumn*

\end{paracol}

\Office{Dominica II Novembris}{Second Sunday of November}{}
\addcontentsline{toc}{subsection}{Dominica II Novembris \textit{— Second Sunday of November}}
\begin{paracol}{2}
\LA
\LessonHead{Lectio I}
\switchcolumn
\EN
\LessonHead{Lesson I}
\switchcolumn*

\LA
\LessonTitle{De Ezechiéle Prophéta}
\switchcolumn
\EN
\LessonTitle{Lesson from the book of Ezekiel}
\switchcolumn*

\LA
\Cite{Ezek 21:1-5}
\switchcolumn
\EN
\Cite{Ezek 21:1-5}
\switchcolumn*

\LA
\noindent Et factus est sermo Dómini ad me, dicens: Fili hóminis, pone fáciem tuam ad Jerúsalem, et stilla ad sanctuária, et prophéta contra humum Israël, et dices terræ Israël: Hæc dicit Dóminus Deus: Ecce ego ad te, et eíciam gládium meum de vagína sua, et occidam in te justum et ímpium. Pro eo autem quod occidi in te justum et ímpium, idcírco egrediétur gládius meus de vagína sua ad omnem carnem, ab Austro usque ad Aquilónem, ut sciat omnis caro quia ego Dóminus, edúxi gládium meum de vagína sua irrevocábilem.\par
\switchcolumn
\EN
\noindent And the word of the Lord came to me, saying: Son of man, set thy face toward Jerusalem, and let thy speech flow towards the holy places, and prophesy against the land of Israel: And say to the land of Israel: Thus saith the Lord God: Behold I come against thee, and I will draw forth my sword out of its sheath, and will cut off in thee the just, and the wicked. And forasmuch as I have cut off in thee the just, and the wicked, therefore shall my sword go forth out of its sheath against all flesh, from the south even to the north. That all flesh may know that I the Lord have drawn my sword out of its sheath not to be turned back.\par
\switchcolumn*

\LA
\begin{Resp}\Rsign Vidi Dóminum sedéntem super sólium excélsum et elevátum, et plena erat omnis terra majestáte ejus: \Star{} Et ea, quæ sub ipso erant, replébant templum.\end{Resp}
\begin{Resp}\Vsign Séraphim stabant super illud: sex alæ uni, et sex alæ álteri.\end{Resp}
\begin{Resp}\Rsign Et ea, quæ sub ipso erant, replébant templum.\end{Resp}
\switchcolumn
\EN
\begin{Resp}\Rsign I saw the Lord sitting upon a throne, high and lifted up, and the whole earth was full of His glory; \Star{} And His train filled the temple.\end{Resp}
\begin{Resp}\Vsign Above it stood the Seraphim each one had six wings.\end{Resp}
\begin{Resp}\Rsign And His train filled the temple.\end{Resp}
\switchcolumn*

\LA
\LessonHead{Lectio II}
\switchcolumn
\EN
\LessonHead{Lesson II}
\switchcolumn*

\LA
\Cite{Ezek 21:6-11}
\switchcolumn
\EN
\Cite{Ezek 21:6-11}
\switchcolumn*

\LA
\noindent Et tu, fili hóminis, ingemísce in contritióne lumbórum, et in amaritudínibus ingemísce coram eis; cumque díxerint ad te: Quare tu gemis? dices: Pro audítu, quia venit, et tabescet omne cor, et dissolvéntur univérsæ manus, et infirmábitur omnis spíritus, et per cuncta genua fluent aquæ; ecce venit, et fiet, ait Dóminus Deus. Et factus est sermo Dómini ad me, dicens: Fili hóminis, propheta, et dices: Hæc dicit Dóminus Deus: Lóquere: Gládius, gládius exacútus est, et limátus: ut cædat víctimas, exacútus est, ut splendeat limátus est; qui moves sceptrum fílii mei, succidísti omne lignum. Et dedi eum ad levigándum, ut teneátur manu; iste exacútus est gládius, et iste limátus est, ut sit in manu interficiéntis.\par
\switchcolumn
\EN
\noindent And thou, son of man, mourn with the breaking of thy loins, and with bitterness sigh before them. And when they shall say to thee: Why mournest thou? thou shalt say: For that which I hear: because it cometh, and every heart shall melt, and all hands shall be made feeble, and every spirit shall faint, and water shall run down every knee: behold it cometh, and it shall be done, saith the Lord God. And the word of the Lord came to me, saying: Son of man, prophesy, and say: Thus saith the Lord God: Say: The sword, the sword is sharpened, and furbished. It is sharpened to kill victims: it is furbished that it may glitter: thou removest the sceptre of my son, thou hast cut down every tree. And I have given it to be furbished, that it may be handled: this sword is sharpened, and it is furbished, that it may be in the hand of the slayer.\par
\switchcolumn*

\LA
\begin{Resp}\Rsign Aspice, Dómine, de sede sancta tua, et cógita de nobis: inclína, Deus meus, aurem tuam et audi: \Star{} Aperi óculos tuos et vide tribulatiónem nostram.\end{Resp}
\begin{Resp}\Vsign Qui regis Israël, inténde, qui dedúcis velut ovem Joseph.\end{Resp}
\begin{Resp}\Rsign Aperi óculos tuos et vide tribulatiónem nostram.\end{Resp}
\switchcolumn
\EN
\begin{Resp}\Rsign Look down, O Lord, from the dwelling-place of thine holiness, and take thought for us. O my God, incline thine ear, and hear. \Star{} Open thine eyes, and behold our desolation.\end{Resp}
\begin{Resp}\Vsign Give ear, O Shepherd of Israel, Thou That leadest Joseph like a flock.\end{Resp}
\begin{Resp}\Rsign Open thine eyes, and behold our desolation.\end{Resp}
\switchcolumn*

\LA
\LessonHead{Lectio III}
\switchcolumn
\EN
\LessonHead{Lesson III}
\switchcolumn*

\LA
\Cite{Ezek 21:12-15}
\switchcolumn
\EN
\Cite{Ezek 21:12-15}
\switchcolumn*

\LA
\noindent Clama et úlula, fili hóminis, quia hic factus est in pópulo meo, hic in cunctis ducibus Israël qui fúgerant; gládio traditi sunt cum pópulo suo, idcirco plaude super femur, quia probátus est, et hoc cum sceptrum subverterit, et non erit, dicit Dóminus Deus. Tu ergo, fili hóminis, propheta, et pércute manu ad manum: et duplicétur gládius, ac triplicétur gládius interfectórum. Hic est gládius occisiónis magnæ, qui obstupéscere eos facit, et corde tabéscere, et multiplicat ruínas.\par
\switchcolumn
\EN
\noindent Cry, and howl, O son of man, for this sword is upon my people, it is upon all the princes of Israel, that are fled: they are delivered up to the sword with my people, strike therefore upon thy thigh, Because it is tried: and that when it shall overthrow the sceptre, and it shall not be, saith the Lord God. Thou therefore, O son of man, prophesy, and strike thy hands together, and let the sword be doubled, and let the sword of the slain be tripled: this is the sword of a great slaughter, that maketh them stand amazed, And languish in heart, and that multiplieth ruins. In all their gates I have set the dread of the sharp sword, the sword that is furbished to glitter, that is made ready for slaughter.\par
\switchcolumn*

\LA
\begin{Resp}\Rsign Aspice, Dómine, quia facta est desoláta cívitas plena divítiis, sedet in tristítia dómina géntium: \Star{} Non est qui consolétur eam, nisi tu, Deus noster.\end{Resp}
\begin{Resp}\Vsign Plorans plorávit in nocte, et lácrimæ ejus in maxíllis ejus.\end{Resp}
\begin{Resp}\Rsign Non est qui consolétur eam, nisi tu, Deus noster.\end{Resp}
\begin{Resp}\Vsign Glória Patri, et Fílio, \Star{} et Spirítui Sancto.\end{Resp}
\begin{Resp}\Rsign Non est qui consolétur eam, nisi tu, Deus noster.\end{Resp}
\switchcolumn
\EN
\begin{Resp}\Rsign Consider, O Lord, how that the city sitteth solitary that was full of riches; how is she become as a widow, she that was great among the nations; \Star{} She hath none to comfort her, save thee, O our God.\end{Resp}
\begin{Resp}\Vsign She weepeth sore in the night, and her tears are on her cheeks.\end{Resp}
\begin{Resp}\Rsign She hath none to comfort her, save thee, O our God.\end{Resp}
\begin{Resp}\Vsign Glory be to the Father, and to the Son, \Star{} and to the Holy Ghost.\end{Resp}
\begin{Resp}\Rsign She hath none to comfort her, save thee, O our God.\end{Resp}
\switchcolumn*

\LA
\LessonHead{Lectio IV}
\switchcolumn
\EN
\LessonHead{Lesson IV}
\switchcolumn*

\LA
\LessonTitle{De Expositióne sancti Hierónymi Presbýteri in Ezechielem Prophetam.}
\switchcolumn
\EN
\LessonTitle{From the Exposition of the Prophet Ezekiel, written by St. Jerome, Priest (at Bethlehem.)}
\switchcolumn*

\LA
\Source{Liber 7 in Ezechiel. cap. 21.}
\switchcolumn
\EN
\Source{Bk. vii. on Ez. xxi.}
\switchcolumn*

\LA
\noindent Quia supra dixerat: Ipsi dicunt ad me: Numquid non per parábolas lóquitur iste? et apertam pópulus flagitábat senténtiam: idcirco id quod Dóminus per metáphoram, sive parábolam, et, ut alii vertére, proverbium, est locútus, nunc manifestius lóquitur: saltus Nageb et Darom et Theman esse Jerúsalem, et templus illíus, sancta sanctórum et omnem terram Judææ; flammámque quæ combustura sit saltum, intelligi gládium devorántem, qui edúctus sit de vagina sua, ut interfíciat justum et ímpium. Hoc est enim lignum víride, et lignum aridum. Unde et Dóminus: si in ligno, ait, víridi tanta faciunt, in sicco quid facient?\par
\switchcolumn
\EN
\noindent In the foregoing chapter it is written Son of man, set thy face toward the way of Theman, and drop thy word toward Darom, and prophesy against the forest of the field of Nageb. And say to the forest of Nageb Hear the word of the Lord Thus saith the Lord God Behold, I will kindle a fire in thee, and it shall devour every green tree in thee, and every dry tree the flaming flame shall not be quenched, and all faces from Negab to Tsaphon shall be burned therein. And all flesh shall see that I the Lord have kindled it it shall not be quenched. Then said I Ah, Lord God they say of me, Doth he not speak parables? (xx. 46-49.) But now, as the people asked for something clearer, the Lord doth speak more openly that which He had uttered in what is diversely called metaphor, parable, or proverb He showeth how that the forest of the field of Nageb, and Darom, and Theman, are figures of Jerusalem, and the Temple, and the Holy-of-Holies, and of all the land of Judah, and that by the flaming fire which should devour the forest, was to be understood that sword, which should be drawn out of the sheath, and should cut off from the land of Israel the righteous and the wicked. The righteous and the wicked are figured by the green tree and the dry tree. Whence also the Lord saith If they do these things in a green tree, what shall be done in a dry? (Luke xxiii. 31.)\par
\switchcolumn*

\LA
\begin{Resp}\Rsign Super muros tuos, Jerúsalem, constítui custódes; \Star{} Tota die et nocte non tacébunt laudáre nomen Dómini.\end{Resp}
\begin{Resp}\Vsign Prædicábunt pópulis fortitúdinem meam, et annuntiábunt géntibus glóriam meam.\end{Resp}
\begin{Resp}\Rsign Tota die et nocte non tacébunt laudáre nomen Dómini.\end{Resp}
\switchcolumn
\EN
\begin{Resp}\Rsign I have set watchmen upon thy walls, O Jerusalem; \Star{} Which shall never hold their peace day nor night, to praise the name of the Lord.\end{Resp}
\begin{Resp}\Vsign They shall proclaim My might unto the nations, and declare My glory unto the Gentiles.\end{Resp}
\begin{Resp}\Rsign Which shall never hold their peace day nor night, to praise the name of the Lord.\end{Resp}
\switchcolumn*

\LA
\LessonHead{Lectio V}
\switchcolumn
\EN
\LessonHead{Lesson V}
\switchcolumn*

\LA
\noindent Primum dixerat: Vaticinare, vel stilla ad Austrum, Africum et Meridiem et ad saltum Meridianum. Quod quia videbátur obscurum, et dicta Prophetæ pópulus nesciebat, secúndo pónitur manifestius, saltum Meridianum esse Jerúsalem; et omnes infructuosas árbores, ad quarum radíces securis posita sit, intelligi habitatóres ejus, gladiúmque interpretari pro incendio. Tertio jubétur Prophetæ, ut tacéntibus illis nec interrogántibus cur ista vaticinátus sit, fáciat per quæ interrogétur, et respondeat quæ Dóminus locútus est.\par
\switchcolumn
\EN
\noindent The first time He had said: Set thy face toward the South, and drop thy word toward the South wind, and prophesy against the forest of the South. But forasmuch as this seemed dark, and the people knew not what the Prophet said, it is a second time stated more clearly that the forest of the South is Jerusalem and all its unfruitful trees, unto whose roots the axe is being laid, are to be understood as figures of her inhabitants and the fire to be kindled in it, to be interpreted the sword. A third time l is the Prophet commanded that when they should hold their peace, nor ask wherefore he prophesied thus, he should do that by which he should be questioned, and should answer that which the Lord had spoken.\par
\switchcolumn*

\LA
\begin{Resp}\Rsign Muro tuo inexpugnábili circumcínge nos, Dómine, et armis tuæ poténtiæ prótege nos semper: \Star{} Líbera, Dómine, Deus Israël, clamántes ad te.\end{Resp}
\begin{Resp}\Vsign Erue nos in mirabílibus tuis, et da glóriam nómini tuo.\end{Resp}
\begin{Resp}\Rsign Líbera, Dómine, Deus Israël, clamántes ad te.\end{Resp}
\switchcolumn
\EN
\begin{Resp}\Rsign Hedge us about with thy wall that cannot be broken down, O Lord, and shield us continually with the arms of thy might. \Star{} O Lord God of Israel, deliver them that cry unto thee.\end{Resp}
\begin{Resp}\Vsign Deliver us also according to thy marvellous works, and give glory to thy Name.\end{Resp}
\begin{Resp}\Rsign O Lord God of Israel, deliver them that cry unto thee.\end{Resp}
\switchcolumn*

\LA
\LessonHead{Lectio VI}
\switchcolumn
\EN
\LessonHead{Lesson VI}
\switchcolumn*

\LA
\noindent Ingemísce, inquit, ejulare non levi voce nec dolóre moderáto, sed in contritióne lumbórum, ut gémitus tuus ex imis viscéribus et amaritúdine animi proferátur. Et hoc fácies coram eis, ut cum te interrogáverint, cur tanto gémitu conteraris, et quid tibi mali accíderit, ut sic ingemiscas; tu eis meo sermóne respondeas: Idcirco plango et dolórem cordis mei dissimuláre non valeo, quia audítus, qui semper meis áuribus insonúerat, opere complétur et venit, ímminens videlicet Babylonii furentis exercitus; qui cum venerit et vallaverit Jerúsalem, tunc tabescet omne cor, et dissolvéntur univérsæ manus: ut occupante pavóre mentes hóminum, nullus audeat repugnare.\par
\switchcolumn
\EN
\noindent Sigh thou, He saith, cry aloud, not softly nor only half sorrowfully, but with the breaking of thy loins, that thy groaning may come from the depth of thy bowels and from the bitterness of thy soul. And this shalt thou do before them. And when they shall ask thee wherefore thou art afflicted with such lamentation, and what evil hath befallen thee that thou groanest thus, thou shalt answer them with My word, saying I lament, and am not able to hide the grief of mine heart, because that that which hath ever sounded in mine ears will indeed be fulfilled, and cometh, even the host of the wrathful Babylonians which threateneth you and when it shall have come, and shall have made trenches all round about Jerusalem, then every heart shall melt, and all hands shall be feeble, and horror shall take hold of the minds of men, and none shall dare to withstand.\par
\switchcolumn*

\LA
\begin{Resp}\Rsign Sustinúimus pacem, et non venit: quæsívimus bona, et ecce turbátio: cognóvimus, Dómine, peccáta nostra; \Star{} Non in perpétuum obliviscáris nos.\end{Resp}
\begin{Resp}\Vsign Peccávimus, ímpie géssimus, iniquitátem fécimus, Dómine, in omnem justítiam tuam.\end{Resp}
\begin{Resp}\Rsign Non in perpétuum obliviscáris nos.\end{Resp}
\begin{Resp}\Vsign Glória Patri, et Fílio, \Star{} et Spirítui Sancto.\end{Resp}
\begin{Resp}\Rsign Non in perpétuum obliviscáris nos.\end{Resp}
\switchcolumn
\EN
\begin{Resp}\Rsign We looked for peace, and it came not we asked for good, and behold trouble. We acknowledge, O Lord, our wickedness. \Star{} Forget us not for ever.\end{Resp}
\begin{Resp}\Vsign O Lord (our God) we have sinned, we have done ungodly, we have dealt unrighteously in all thine ordinances.\end{Resp}
\begin{Resp}\Rsign Forget us not for ever.\end{Resp}
\begin{Resp}\Vsign Glory be to the Father, and to the Son, \Star{} and to the Holy Ghost.\end{Resp}
\begin{Resp}\Rsign Forget us not for ever.\end{Resp}
\switchcolumn*

\end{paracol}

\Office{Feria II infra Hebdomadam II Novembris}{Monday of the second week of November}{}
\addcontentsline{toc}{subsection}{Feria II infra Hebdomadam II Novembris \textit{— Monday of the second week of November}}
\begin{paracol}{2}
\LA
\LessonHead{Lectio I}
\switchcolumn
\EN
\LessonHead{Lesson I}
\switchcolumn*

\LA
\LessonTitle{De Ezechiéle Prophéta}
\switchcolumn
\EN
\LessonTitle{Lesson from the book of Ezekiel}
\switchcolumn*

\LA
\Cite{Ezek 33:1-5}
\switchcolumn
\EN
\Cite{Ezek 33:1-5}
\switchcolumn*

\LA
\noindent Ecce factum est verbum Dómini ad me, dicens: Fili hóminis, loquere ad fílios pópuli tui, et dices ad eos: Terra, cum induxero super eam gládium, et tulerit pópulus terræ virum unum de novíssimis suis, et constitúerit eum eum super se speculatórem: et ille víderit gládium veniéntem super terram, et cecínerit búccina, et annuntiáverit pópulo; áudiens autem quisquis ille est sónitum búccinæ, et non se observáverit, venerítque gládius, et tulerit eum: sanguis ipsíus super caput ejus erit. Sonum búccinæ audívit, et non se observávit, sanguis ejus in ipso erit; si autem se custodíerit, ánimam suam salvábit.\par
\switchcolumn
\EN
\noindent And the word of the Lord came to me, saying: Son of man, speak to the children of thy people, and say to them: When I bring the sword upon a land, if the people of the land take a man, one of their meanest, and make him a watchman over them: And he see the sword coming upon the land, and sound the trumpet, and tell the people: Then he that heareth the sound of the trumpet, whosoever he be, and doth not look to himself, if the sword come, and cut him off: his blood shall be upon his own head. He heard the sound of the trumpet and did not look to himself, his blood shall be upon him: but if he look to himself, he shall save his life.\par
\switchcolumn*

\LA
\begin{Resp}\Rsign Redémit pópulum suum et liberávit eum, et vénient et exsultábunt in monte Sion et gaudébunt de bonis Dómini super fruménto, vino et óleo, \Star{} Et ultra non esúrient.\end{Resp}
\begin{Resp}\Vsign Erítque ánima eórum quasi hortus irríguus.\end{Resp}
\begin{Resp}\Rsign Et ultra non esúrient.\end{Resp}
\switchcolumn
\EN
\begin{Resp}\Rsign He hath redeemed His people, and ransomed them therefore they shall come and sing in the height of Zion, and shall rejoice in the goodness of the Lord, for wheat, and for wine, and for oil \Star{} And they shall hunger no more.\end{Resp}
\begin{Resp}\Vsign And their soul shall be as a watered garden.\end{Resp}
\begin{Resp}\Rsign And they shall hunger no more.\end{Resp}
\switchcolumn*

\LA
\LessonHead{Lectio II}
\switchcolumn
\EN
\LessonHead{Lesson II}
\switchcolumn*

\LA
\Cite{Ezek 33:6-8}
\switchcolumn
\EN
\Cite{Ezek 33:6-8}
\switchcolumn*

\LA
\noindent Quod si speculátor víderit gládium veniéntem, et non insonúerit búccina, et pópulus se non custodíerit, veneritque gládius, et tulerit de eis ánimam: ille quidem in iniquitáte sua captus est; sánguinem autem ejus de manu speculátoris requíram. Et tu, fili hóminis, speculatórem dedi te dómui Israël: áudiens ergo ex ore meo sermónem, annuntiábis eis ex me. Si, me dicénte ad ímpium: Impie, morte moriéris: non fueris locútus, ut se custódiat ímpius a via sua, ipse ímpius in iniquitáte sua moriétur, sánguinem autem ejus de manu tua requíram.\par
\switchcolumn
\EN
\noindent And if the watchman see the sword coming, and sound not the trumpet: and the people look not to themselves, and the sword come, and cut off a soul from among them: he indeed is taken away in his iniquity, but I will require his blood at the hand of the watchman. So thou, O son of man, I have made thee a watchman to the house of Israel: therefore thou shalt hear the word from my mouth, and shalt tell it them from me. When I say to the wicked: O wicked man, thou shalt surely die: if thou dost not speak to warn the wicked man from his way: that wicked man shall die in his iniquity, but I will require his blood at thy hand.\par
\switchcolumn*

\LA
\begin{Resp}\Rsign Angústiæ mihi sunt úndique, et quid éligam ignóro; \Star{} Mélius est mihi incídere in manus hóminum, quam derelínquere legem Dei mei.\end{Resp}
\begin{Resp}\Vsign Si enim hoc égero, mors mihi est; si autem non égero, non effúgiam manus vestras.\end{Resp}
\begin{Resp}\Rsign Mélius est mihi incídere in manus hóminum, quam derelínquere legem Dei mei.\end{Resp}
\switchcolumn
\EN
\begin{Resp}\Rsign I am straitened on every side, and know not what to choose. \Star{} It is better for me to fall into the hands of men, than to sin against the law of my God.\end{Resp}
\begin{Resp}\Vsign For if I do this thing, it is death unto me and if I do it not, I cannot escape your hands.\end{Resp}
\begin{Resp}\Rsign It is better for me to fall into the hands of men, than to sin against the law of my God.\end{Resp}
\switchcolumn*

\LA
\LessonHead{Lectio III}
\switchcolumn
\EN
\LessonHead{Lesson III}
\switchcolumn*

\LA
\Cite{Ezek 33:9-11}
\switchcolumn
\EN
\Cite{Ezek 33:9-11}
\switchcolumn*

\LA
\noindent Si autem, annuntiánte te ad ímpium ut a viis suis convertátur, non fúerit convérsus a via sua, ipse in iniquitáte sua moriétur, porro tu ánimam tuam liberásti. Tu ergo, fili hóminis, dic ad domum Israël: Sic locúti estis, dicéntes: Iniquitátes nostræ et peccáta nostra super nos sunt, et in ipsis nos tabéscimus: quómodo ergo vívere potérimus? Dic ad eos: Vivo ego, dicit Dóminus Deus, nolo mortem ímpii, sed ut convertátur ímpius a via sua et vivat. Convertímini, convertímini a viis vestris pessimis, et quare moriemini, domus Israël?\par
\switchcolumn
\EN
\noindent But if thou tell the wicked man, that he may be converted from his ways, and he be not converted from his way: he shall die in his iniquity: but thou hast delivered thy soul. Thou therefore, O son of man, say to the house of Israel: Thus you have spoken, saying: Our iniquities, and our sins are upon us, and we pine away in them: how then can we live? Say to them: As I live, saith the Lord God, I desire not the death of the wicked, but that the wicked turn from his way, and live. Turn ye, turn ye from your evil ways: and why will you die, O house of Israel?\par
\switchcolumn*

\LA
\begin{Resp}\Rsign Misit Dóminus Angelum suum et conclúsit ora leónum, \Star{} Et non contaminavérunt, quia coram eo injustítia invénta non est in me.\end{Resp}
\begin{Resp}\Vsign Misit Deus misericórdiam suam et veritátem suam: ánimam meam erípuit de médio catulórum leónum.\end{Resp}
\begin{Resp}\Rsign Et non contaminavérunt, quia coram eo injustítia invénta non est in me.\end{Resp}
\begin{Resp}\Vsign Glória Patri, et Fílio, \Star{} et Spirítui Sancto.\end{Resp}
\begin{Resp}\Rsign Et non contaminavérunt: quia coram eo injustítia invénta non est in me.\end{Resp}
\switchcolumn
\EN
\begin{Resp}\Rsign The Lord hath sent His angel, and hath shut the lions' mouths, that \Star{} They have not hurt me forasmuch as before Him innocency was found in me.\end{Resp}
\begin{Resp}\Vsign God hath sent forth His mercy and His truth, (and delivered) my soul from among the lions' whelps.\end{Resp}
\begin{Resp}\Rsign They have not hurt me forasmuch as before Him innocency was found in me.\end{Resp}
\begin{Resp}\Vsign Glory be to the Father, and to the Son, \Star{} and to the Holy Ghost.\end{Resp}
\begin{Resp}\Rsign They have not hurt me forasmuch as before Him innocency was found in me.\end{Resp}
\switchcolumn*

\end{paracol}

\Office{Feria III infra Hebdomadam II Novembris}{Tuesday of the second week of November}{}
\addcontentsline{toc}{subsection}{Feria III infra Hebdomadam II Novembris \textit{— Tuesday of the second week of November}}
\begin{paracol}{2}
\LA
\LessonHead{Lectio I}
\switchcolumn
\EN
\LessonHead{Lesson I}
\switchcolumn*

\LA
\LessonTitle{De Ezechiéle Prophéta}
\switchcolumn
\EN
\LessonTitle{Lesson from the book of Ezekiel}
\switchcolumn*

\LA
\Cite{Ezek 34:1-4}
\switchcolumn
\EN
\Cite{Ezek 34:1-4}
\switchcolumn*

\LA
\noindent Et factum est verbum Dómini ad me, dicens: Fili hóminis, prophéta de pastóribus Israël: prophéta, et dices pastóribus: Hæc dicit Dóminus Deus: Væ pastóribus Israël, qui pascébant semetípsos! Nonne greges a pastóribus pascúntur? Lac comedebátis, et lanis operiebámini, et quod crassum erat occidebátis: gregem autem meum non pascebátis; quod infirmum fuit non consolidastis, et quod ægrotum non sanastis: quod confractum est non alligastis, et quod abjectum est non reduxístis, et quod períerat non quæsistis, sed cum austeritáte imperabátis eis, et cum poténtia.\par
\switchcolumn
\EN
\noindent And the word of the Lord came to me, saying: Son of man, prophesy concerning the shepherds of Israel: prophesy, and say to the shepherds: Thus saith the Lord God: Woe to the shepherds of Israel, that fed themselves: should not the hocks be fed by the shepherds? You ate the milk, end you clothed yourselves with the wool, and you killed that which was fat: but my flock you did not feed. The weak you have not strengthened, and that which was sick you have not healed, that which was broken you have not bound up, and that which was driven away you have not brought again, neither have you sought that which was lost: but you ruled over them with rigour, and with a high hand.\par
\switchcolumn*

\LA
\begin{Resp}\Rsign A fácie furóris tui, Deus, conturbáta est omnis terra: \Star{} Sed miserére, Dómine, et ne fácias consummatiónem.\end{Resp}
\begin{Resp}\Vsign Dómine, Dóminus noster, quam admirábile est nomen tuum!\end{Resp}
\begin{Resp}\Rsign Sed miserére, Dómine, et ne fácias consummatiónem.\end{Resp}
\switchcolumn
\EN
\begin{Resp}\Rsign Before the face of thine anger, O God, the whole earth is troubled; \Star{} But Thou, O Lord, have mercy, and make not an end utterly.\end{Resp}
\begin{Resp}\Vsign Lord our Ruler, how excellent is thy Name in all the earth!\end{Resp}
\begin{Resp}\Rsign But Thou, O Lord, have mercy, and make not an end utterly.\end{Resp}
\switchcolumn*

\LA
\LessonHead{Lectio II}
\switchcolumn
\EN
\LessonHead{Lesson II}
\switchcolumn*

\LA
\Cite{Ezek 34:5-9}
\switchcolumn
\EN
\Cite{Ezek 34:5-9}
\switchcolumn*

\LA
\noindent Et dispersæ sunt oves meæ, eo quod non esset pastor: et factæ sunt in devoratiónem ómnium bestiárum agri, et dispersæ sunt. Erravérunt greges mei in cunctis móntibus, et in universo colle excélso, et super omnem fáciem terræ dispersi sunt greges mei, et non erat, qui requíreret: non erat, inquam, qui requíreret. Proptérea, pastores, audíte verbum Dómini: Vivo ego, dicit Dóminus Deus, quia, pro eo quod facti sunt greges mei in rapinam, et oves meæ in devoratiónem ómnium bestiárum agri, eo quod non esset pastor; neque enim quæsiérunt pastores mei gregem meum, sed pascébant pastores semetípsos et greges meos non pascébant; proptérea, pastores, audíte verbum Domini.\par
\switchcolumn
\EN
\noindent And my sheep were scattered, because there was no shepherd: and they became the prey of all the beasts of the field, and were scattered. My sheep have wandered in every mountain, and in every high hill: and my flocks were scattered upon the face of the earth, and there was none that sought them, there was none, I say, that sought them. Therefore, ye shepherds, hear the word of the Lord: As I live, saith the Lord God, forasmuch as my flocks have been made a spoil, and my sheep are become a prey to all the beasts of the field, because there was no shepherd: for my shepherds did not seek after my flock, but the shepherds fed themselves, and fed not my flocks: Therefore, ye shepherds, hear the word of the Lord:\par
\switchcolumn*

\LA
\begin{Resp}\Rsign Civitátem istam tu circúmda, Dómine: et Angeli tui custódiant muros ejus. \Star{} Exáudi, Dómine, pópulum tuum cum misericórdia.\end{Resp}
\begin{Resp}\Vsign Avertátur furor tuus, Dómine, a pópulo tuo et a civitáte sancta tua.\end{Resp}
\begin{Resp}\Rsign Exáudi, Dómine, pópulum tuum cum misericórdia.\end{Resp}
\switchcolumn
\EN
\begin{Resp}\Rsign Fence Thou this city; O Lord, and let thine angels keep the walls thereof. \Star{} O Lord, hearken unto thy people with mercy.\end{Resp}
\begin{Resp}\Vsign O Lord, let thine anger be turned away from thy people, and from thine holy city.\end{Resp}
\begin{Resp}\Rsign O Lord, hearken unto thy people with mercy.\end{Resp}
\switchcolumn*

\LA
\LessonHead{Lectio III}
\switchcolumn
\EN
\LessonHead{Lesson III}
\switchcolumn*

\LA
\Cite{Ezek 34:10-12}
\switchcolumn
\EN
\Cite{Ezek 34:10-12}
\switchcolumn*

\LA
\noindent Hæc dicit Dóminus Deus: Ecce ego ipse super pastores, requíram gregem meum de manu eórum et cessare fáciam eos, ut ultra non pascant gregem, nec pascant ámplius pastores semetípsos: et liberábo gregem meum de ore eórum, et non erit ultra eis in escam. Quia hæc dicit Dóminus Deus: Ecce ego ipse requíram oves meas et visitábo eas: sicut visitat pastor gregem suum in die quando fúerit in médio óvium suárum dissipatárum, sic visitábo oves meas, et liberábo eas de ómnibus locis, in quibus dispersæ fuerant in die nubis et caliginis.\par
\switchcolumn
\EN
\noindent Thus saith the Lord God: Behold I myself come upon the shepherds, I will require my flock at their hand, and I will cause them to cease from feeding the flock any more, neither shall the shepherds feed themselves any more: and I will deliver my flock from their mouth, and it shall no more be meat for them. For thus saith the Lord God: Behold I myself will seek my sheep, and will visit them. As the shepherd visiteth his flock in the day when he shall be in the midst of his sheep that were scattered, so will I visit my sheep, and will deliver them out of all the places where they have been scattered in the cloudy and dark day.\par
\switchcolumn*

\LA
\begin{Resp}\Rsign Genti peccatríci, pópulo pleno peccáto miserére, \Star{} Dómine Deus.\end{Resp}
\begin{Resp}\Vsign Esto placábilis super nequítiam pópuli tui.\end{Resp}
\begin{Resp}\Rsign Dómine Deus.\end{Resp}
\begin{Resp}\Vsign Glória Patri, et Fílio, \Star{} et Spirítui Sancto.\end{Resp}
\begin{Resp}\Rsign Dómine Deus.\end{Resp}
\switchcolumn
\EN
\begin{Resp}\Rsign O Lord God! have mercy upon \Star{} The sinful nation, upon the people laden with iniquity.\end{Resp}
\begin{Resp}\Vsign Let it repent thee concerning the transgression of thy people.\end{Resp}
\begin{Resp}\Rsign The people laden with iniquity.\end{Resp}
\begin{Resp}\Vsign Glory be to the Father, and to the Son, \Star{} and to the Holy Ghost.\end{Resp}
\begin{Resp}\Rsign The people laden with iniquity.\end{Resp}
\switchcolumn*

\end{paracol}

\Office{Feria IV infra Hebdomadam II Novembris}{Wednesday of the second week of November}{}
\addcontentsline{toc}{subsection}{Feria IV infra Hebdomadam II Novembris \textit{— Wednesday of the second week of November}}
\begin{paracol}{2}
\LA
\LessonHead{Lectio I}
\switchcolumn
\EN
\LessonHead{Lesson I}
\switchcolumn*

\LA
\LessonTitle{De Ezechiéle Prophéta}
\switchcolumn
\EN
\LessonTitle{Lesson from the book of Ezekiel}
\switchcolumn*

\LA
\Cite{Ezek 40:1-2}
\switchcolumn
\EN
\Cite{Ezek 40:1-2}
\switchcolumn*

\LA
\noindent In vigésimo quinto anno transmigratiónis nostræ, in exórdio anni, décima mensis, quartodécimo anno postquam percússa est cívitas, in ipsa hac die, facta est super me manus Dómini, et adduxit me illuc. In visiónibus Dei addúxit me in terram Israël, et dimísit me super montem excélsum nimis, super quem erat quasi ædifícium civitátis vergéntis ad Austrum, et introdúxit me illuc.\par
\switchcolumn
\EN
\noindent In the five and twentieth year of our captivity, in the beginning of the year, the tenth day of the month, the fourteenth year after the city was destroyed: in the selfsame day the hand of the Lord was upon me, and he brought me thither. In the visions of God he brought me into the land of Israel, and set me upon a very high mountain: upon which there was as the building of a city, bending towards the south.\par
\switchcolumn*

\LA
\begin{Resp}\Rsign Indicábo tibi, homo, quid sit bonum aut quid Dóminus requírat a te: \Star{} Fácere judícium et justítiam et sollícitum ambuláre cum Deo tuo.\end{Resp}
\begin{Resp}\Vsign Spera in Dómino, et fac bonitátem, et inhábita terram.\end{Resp}
\begin{Resp}\Rsign Fácere judícium et justítiam et sollícitum ambuláre cum Deo tuo.\end{Resp}
\switchcolumn
\EN
\begin{Resp}\Rsign I will show thee, O man, what is good, and what doth the Lord require of thee, but to \Star{} Do justice and judgment, and to walk humbly with thy God.\end{Resp}
\begin{Resp}\Vsign Trust in the Lord, and do good, and dwell in the land.\end{Resp}
\begin{Resp}\Rsign Do justice and judgment, and walk humbly with thy God.\end{Resp}
\switchcolumn*

\LA
\LessonHead{Lectio II}
\switchcolumn
\EN
\LessonHead{Lesson II}
\switchcolumn*

\LA
\Cite{Ezek 40:3-4}
\switchcolumn
\EN
\Cite{Ezek 40:3-4}
\switchcolumn*

\LA
\noindent Et ecce vir, cujus erat spécies quasi spécies æris, et funículus líneus in manu ejus, et cálamus mensúræ in manu ejus, stabat autem in porta, et locútus est ad me idem vir: Fili hóminis, vide óculis tuis, et áuribus tuis audi, et pone cor tuum in ómnia quæ ego osténdam tibi, quia ut ostendántur tibi addúctus es huc, annúntia ómnia quæ tu vides dómui Israël.\par
\switchcolumn
\EN
\noindent And he brought me in thither, and behold a man, whose appearance was like the appearance of brass, with a line of flax in his hand, and a measuring reed in his hand, and he stood in the gate. And this man said to me: Son of man, see with thy eyes, and hear with thy ears, and set thy heart upon all that I shall shew thee: for thou art brought hither that they may be shewn to thee: declare all that thou seest, to the house of Israel.\par
\switchcolumn*

\LA
\begin{Resp}\Rsign Angústiæ mihi sunt úndique, et quid éligam ignóro; \Star{} Mélius est mihi incídere in manus hóminum, quam derelínquere legem Dei mei.\end{Resp}
\begin{Resp}\Vsign Si enim hoc égero, mors mihi est; si autem non égero, non effúgiam manus vestras.\end{Resp}
\begin{Resp}\Rsign Mélius est mihi incídere in manus hóminum, quam derelínquere legem Dei mei.\end{Resp}
\switchcolumn
\EN
\begin{Resp}\Rsign I am straitened on every side, and know not what to choose. \Star{} It is better for me to fall into the hands of men, than to sin against the law of my God.\end{Resp}
\begin{Resp}\Vsign For if I do this thing, it is death unto me and if I do it not, I cannot escape your hands.\end{Resp}
\begin{Resp}\Rsign It is better for me to fall into the hands of men, than to sin against the law of my God.\end{Resp}
\switchcolumn*

\LA
\LessonHead{Lectio III}
\switchcolumn
\EN
\LessonHead{Lesson III}
\switchcolumn*

\LA
\Cite{Ezek 40:5-6}
\switchcolumn
\EN
\Cite{Ezek 40:5-6}
\switchcolumn*

\LA
\noindent Et ecce murus forínsecus in circúitu domus úndique: et in manu viri cálamus mensúræ sex cubitórum et palmo; et mensus est latitúdinem ædifícii cálamo uno, altitúdinem quoque cálamo uno. Et venit ad portam quæ respiciébat viam Orientálem, et ascéndit per gradus ejus: et mensus est limen portæ cálamo uno latitúdinem, id est, limen unum cálamo uno in latitúdine.\par
\switchcolumn
\EN
\noindent And behold there was a wall on the outside of the house round about, and in the man's hand a measuring reed of six cubits and a handbreadth: and he measured the breadth of the building one reed, and the height one reed. And he came to the gate that looked toward the east, and he went up the steps thereof: and he measured the breadth of the threshold of the gate one reed, that is, one threshold was one reed broad:\par
\switchcolumn*

\LA
\begin{Resp}\Rsign Misit Dóminus Angelum suum et conclúsit ora leónum, \Star{} Et non contaminavérunt, quia coram eo injustítia invénta non est in me.\end{Resp}
\begin{Resp}\Vsign Misit Deus misericórdiam suam et veritátem suam: ánimam meam erípuit de médio catulórum leónum.\end{Resp}
\begin{Resp}\Rsign Et non contaminavérunt, quia coram eo injustítia invénta non est in me.\end{Resp}
\begin{Resp}\Vsign Glória Patri, et Fílio, \Star{} et Spirítui Sancto.\end{Resp}
\begin{Resp}\Rsign Et non contaminavérunt: quia coram eo injustítia invénta non est in me.\end{Resp}
\switchcolumn
\EN
\begin{Resp}\Rsign The Lord hath sent His angel, and hath shut the lions' mouths, that \Star{} They have not hurt me forasmuch as before Him innocency was found in me.\end{Resp}
\begin{Resp}\Vsign God hath sent forth His mercy and His truth, (and delivered) my soul from among the lions' whelps.\end{Resp}
\begin{Resp}\Rsign They have not hurt me forasmuch as before Him innocency was found in me.\end{Resp}
\begin{Resp}\Vsign Glory be to the Father, and to the Son, \Star{} and to the Holy Ghost.\end{Resp}
\begin{Resp}\Rsign They have not hurt me forasmuch as before Him innocency was found in me.\end{Resp}
\switchcolumn*

\end{paracol}

\Office{Feria VI infra Hebdomadam II Novembris}{Friday of the second week of November}{}
\addcontentsline{toc}{subsection}{Feria VI infra Hebdomadam II Novembris \textit{— Friday of the second week of November}}
\begin{paracol}{2}
\LA
\LessonHead{Lectio I}
\switchcolumn
\EN
\LessonHead{Lesson I}
\switchcolumn*

\LA
\LessonTitle{De Ezechiéle Prophéta}
\switchcolumn
\EN
\LessonTitle{Lesson from the book of Ezekiel}
\switchcolumn*

\LA
\Cite{Ezek 43:1-5}
\switchcolumn
\EN
\Cite{Ezek 43:1-5}
\switchcolumn*

\LA
\noindent Et duxit me ad portam, quæ respiciébat ad viam Orientalem, et ecce glória Dei Israël ingrediebátur per viam Orientalem, et vox erat ei quasi vox aquárum multárum, et terra splendébat a majestate ejus. Et vidi visiónem secúndum spéciem quam víderam, quando venit ut dispérderet civitátem, et spécies secúndum aspéctum quem videram juxta fluvium Chobar, et cécidi super fáciem meam. Et májestas Dómini ingréssa est templum per viam portæ, quæ respiciébat ad Oriéntem; et elevávit me spíritus et introduxit me in atrium interius, et ecce repléta erat glória Dómini domus.\par
\switchcolumn
\EN
\noindent And he brought me to the gate that looked towards the east. And behold the glory of the God of Israel came in by the way of the east: and his voice was like the noise of many waters, and the earth shone with his majesty. And I saw the vision according to the appearance which I had seen when he came to destroy the city: and the appearance was according to the vision which I had seen by the river Chobar: and I fell upon my face. And the majesty of the Lord went into the temple by the way of the gate that looked to the east. And the spirit lifted me up and brought me into the inner court: and behold the house was filled with the glory of the Lord.\par
\switchcolumn*

\LA
\begin{Resp}\Rsign Super muros tuos, Jerúsalem, constítui custódes; \Star{} Tota die et nocte non tacébunt laudáre nomen Dómini.\end{Resp}
\begin{Resp}\Vsign Prædicábunt pópulis fortitúdinem meam, et annuntiábunt géntibus glóriam meam.\end{Resp}
\begin{Resp}\Rsign Tota die et nocte non tacébunt laudáre nomen Dómini.\end{Resp}
\switchcolumn
\EN
\begin{Resp}\Rsign I have set watchmen upon thy walls, O Jerusalem; \Star{} Which shall never hold their peace day nor night, to praise the name of the Lord.\end{Resp}
\begin{Resp}\Vsign They shall proclaim My might unto the nations, and declare My glory unto the Gentiles.\end{Resp}
\begin{Resp}\Rsign Which shall never hold their peace day nor night, to praise the name of the Lord.\end{Resp}
\switchcolumn*

\LA
\LessonHead{Lectio II}
\switchcolumn
\EN
\LessonHead{Lesson II}
\switchcolumn*

\LA
\Cite{Ezek 43:6-8}
\switchcolumn
\EN
\Cite{Ezek 43:6-8}
\switchcolumn*

\LA
\noindent Et audívi loquentem ad me de domo, et vir qui stabat juxta me, dixit ad me: Fili hóminis, locus sólii mei, et locus vestigiórum pedum meórum, ubi hábito in médio filiórum Israël in ætérnum: et non polluent ultra domus Israël nomen sanctum meum, ipsi et reges eórum, in fornicatiónibus suis, et in ruinis regum suórum, et in excélsis. Qui fabricáti sunt limen suum juxta limen meum, et postes suos juxta postes meos, et murus erat inter me et eos: et polluérunt nomen sanctum meum in abominatiónibus, quas fecérunt; propter quod consumpsi eos in ira mea.\par
\switchcolumn
\EN
\noindent And I heard one speaking to me out of the house, and the man that stood by me, Said to me: Son of man, the place of my throne, and the place of the soles of my feet, where I dwell in the midst of the children of Israel for ever: and the house of Israel shall no more profane my holy name, they and their kings by their fornications, and by the carcasses of their kings, and by the high places. They who have set their threshold by my threshold, and their posts by my posts: and there was but a wall between me and them: and they profaned my holy name by the abominations which they committed: for which reason I consumed them in my wrath.\par
\switchcolumn*

\LA
\begin{Resp}\Rsign Muro tuo inexpugnábili circumcínge nos, Dómine, et armis tuæ poténtiæ prótege nos semper: \Star{} Líbera, Dómine, Deus Israël, clamántes ad te.\end{Resp}
\begin{Resp}\Vsign Erue nos in mirabílibus tuis, et da glóriam nómini tuo.\end{Resp}
\begin{Resp}\Rsign Líbera, Dómine, Deus Israël, clamántes ad te.\end{Resp}
\switchcolumn
\EN
\begin{Resp}\Rsign Hedge us about with thy wall that cannot be broken down, O Lord, and shield us continually with the arms of thy might. \Star{} O Lord God of Israel, deliver them that cry unto thee.\end{Resp}
\begin{Resp}\Vsign Deliver us also according to thy marvellous works, and give glory to thy Name.\end{Resp}
\begin{Resp}\Rsign O Lord God of Israel, deliver them that cry unto thee.\end{Resp}
\switchcolumn*

\LA
\LessonHead{Lectio III}
\switchcolumn
\EN
\LessonHead{Lesson III}
\switchcolumn*

\LA
\Cite{Ezek 43:9-11}
\switchcolumn
\EN
\Cite{Ezek 43:9-11}
\switchcolumn*

\LA
\noindent Nunc ergo repellant procul fornicatiónem suam et ruínas regum suórum a me, et habitábo in médio eórum semper. Tu autem, fili hóminis, ostende dómui Israël templum, et confundántur ab iniquitátibus suis, et metiántur fábricam, et erubéscant ex ómnibus quæ fecérunt. Figuram domus, et fábricæ ejus, éxitus et intróitus, et omnem descriptiónem ejus, et univérsa præcepta ejus, cunctumque ordinem ejus, et omnes leges ejus ostende eis, et scribes in óculis eórum, ut custódiant omnes descriptiónes ejus, et præcepta illíus, et faciunt ea.\par
\switchcolumn
\EN
\noindent Now therefore let them put away their fornications, and the carcasses of their kings far from me: and I will dwell in the midst of them for ever. But thou, son of man, shew to the house of Israel the temple, and let them be ashamed of their iniquities, and let them measure the building: And be ashamed of all that they have done. Shew them the form of the house, and of the fashion thereof, the goings out and the comings in, and the whole plan thereof, and all its ordinances, and all its order, and all its laws, and thou shalt write it in their sight: that they may keep the whole form thereof, and its ordinances, and do them.\par
\switchcolumn*

\LA
\begin{Resp}\Rsign Sustinúimus pacem, et non venit: quæsívimus bona, et ecce turbátio: cognóvimus, Dómine, peccáta nostra; \Star{} Non in perpétuum obliviscáris nos.\end{Resp}
\begin{Resp}\Vsign Peccávimus, ímpie géssimus, iniquitátem fécimus, Dómine, in omnem justítiam tuam.\end{Resp}
\begin{Resp}\Rsign Non in perpétuum obliviscáris nos.\end{Resp}
\begin{Resp}\Vsign Glória Patri, et Fílio, \Star{} et Spirítui Sancto.\end{Resp}
\begin{Resp}\Rsign Non in perpétuum obliviscáris nos.\end{Resp}
\switchcolumn
\EN
\begin{Resp}\Rsign We looked for peace, and it came not we asked for good, and behold trouble. We acknowledge, O Lord, our wickedness. \Star{} Forget us not for ever.\end{Resp}
\begin{Resp}\Vsign O Lord (our God) we have sinned, we have done ungodly, we have dealt unrighteously in all thine ordinances.\end{Resp}
\begin{Resp}\Rsign Forget us not for ever.\end{Resp}
\begin{Resp}\Vsign Glory be to the Father, and to the Son, \Star{} and to the Holy Ghost.\end{Resp}
\begin{Resp}\Rsign Forget us not for ever.\end{Resp}
\switchcolumn*

\end{paracol}

\Office{Dominica III Novembris}{Third Sunday of November}{}
\addcontentsline{toc}{subsection}{Dominica III Novembris \textit{— Third Sunday of November}}
\begin{paracol}{2}
\LA
\LessonHead{Lectio I}
\switchcolumn
\EN
\LessonHead{Lesson I}
\switchcolumn*

\LA
\LessonTitle{Incipit liber Daniélis Prophétæ}
\switchcolumn
\EN
\LessonTitle{Lesson from the book of Daniel}
\switchcolumn*

\LA
\Cite{Dan 1:1-4}
\switchcolumn
\EN
\Cite{Dan 1:1-4}
\switchcolumn*

\LA
\noindent Anno tértio regni Jóakim regis Juda, venit Nabuchodónosor, rex Babylónis, in Jerúsalem et obsédit eam: et trádidit Dóminus in manu ejus Jóakim, regem Juda, et partem vasórum domus Dei: et asportávit ea in terra Sénnaar in domum dei sui, et vasa íntulit in domum thesáuri dei sui. Et ait rex Asphenez præpósito eunuchórum, ut introdúceret de fíliis Israël, et de sémine régio et tyrannórum, púeros, in quibus nulla esset mácula, decóros forma, et erudítos omni sapiéntia, cautos sciéntia, et doctos disciplína, et qui possent stare in palátio regis, ut docéret eos lítteras et linguam Chaldæórum.\par
\switchcolumn
\EN
\noindent In the third year of the reign of Joakim king of Juda, Nabuchodonosor king of Babylon came to Jerusalem, and besieged it. And the Lord delivered into his hands Joakim the king of Juda, and part of the vessels of the house of God: and he carried them away into the land of Sennaar, to the house of his god, and the vessels he brought into the treasure house of his god. And the king spoke to Asphenez the master of the eunuchs, that he should bring in some of the children of Israel, and of the king's seed and of the princes, Children in whom there was no blemish, well favoured, and skilful in all wisdom, acute in knowledge, and instructed in science, and such as might stand in the king's palace, that he might teach them the learning, and the tongue of the Chaldeans.\par
\switchcolumn*

\LA
\begin{Resp}\Rsign Vidi Dóminum sedéntem super sólium excélsum et elevátum, et plena erat omnis terra majestáte ejus: \Star{} Et ea, quæ sub ipso erant, replébant templum.\end{Resp}
\begin{Resp}\Vsign Séraphim stabant super illud: sex alæ uni, et sex alæ álteri.\end{Resp}
\begin{Resp}\Rsign Et ea, quæ sub ipso erant, replébant templum.\end{Resp}
\switchcolumn
\EN
\begin{Resp}\Rsign I saw the Lord sitting upon a throne, high and lifted up, and the whole earth was full of His glory; \Star{} And His train filled the temple.\end{Resp}
\begin{Resp}\Vsign Above it stood the Seraphim each one had six wings.\end{Resp}
\begin{Resp}\Rsign And His train filled the temple.\end{Resp}
\switchcolumn*

\LA
\LessonHead{Lectio II}
\switchcolumn
\EN
\LessonHead{Lesson II}
\switchcolumn*

\LA
\Cite{Dan 1:5-9}
\switchcolumn
\EN
\Cite{Dan 1:5-9}
\switchcolumn*

\LA
\noindent Et constítuit eis rex annónam per síngulos dies de cibis suis, et de vino unde bibébat ipse, ut enutríti tribus annis, póstea starent in conspéctu regis. Fuérunt ergo inter eos de fíliis Juda, Dániel, Ananías, Mísaël, et Azarías. Et impósuit eis præpósitus eunuchórum nómina: Daniéli, Baltássar; Ananíæ, Sidrach; Misaéli, Misach; et Azaríæ, Abdénago. Propósuit autem Dániel in corde suo, ne polluerétur de mensa regis, neque de vino potus ejus: et rogávit eunuchórum præpósitum ne contaminarétur. Dedit autem Deus Daniéli grátiam et misericórdiam in conspéctu príncipis eunuchórum.\par
\switchcolumn
\EN
\noindent And the king appointed them a daily provision, of his own meat, and of the wine of which he drank himself, that being nourished three years, afterwards they might stand before the king. Now there were among them of the children of Juda, Daniel, Ananias, Misael, and Azarias. And the master of the eunuchs gave them names: to Daniel, Baltassar: to Ananias, Sidrach: to Misael, Misach: and to Azarias, Abdenago. But Daniel purposed in his heart that he would not be defiled with the king's table, nor with the wine which he drank: and he requested the master of the eunuchs that he might not be defiled. And God gave to Daniel grace and mercy in the sight of the prince of the eunuchs.\par
\switchcolumn*

\LA
\begin{Resp}\Rsign Aspice, Dómine, de sede sancta tua, et cógita de nobis: inclína, Deus meus, aurem tuam et audi: \Star{} Aperi óculos tuos et vide tribulatiónem nostram.\end{Resp}
\begin{Resp}\Vsign Qui regis Israël, inténde, qui dedúcis velut ovem Joseph.\end{Resp}
\begin{Resp}\Rsign Aperi óculos tuos et vide tribulatiónem nostram.\end{Resp}
\switchcolumn
\EN
\begin{Resp}\Rsign Look down, O Lord, from the dwelling-place of thine holiness, and take thought for us. O my God, incline thine ear, and hear. \Star{} Open thine eyes, and behold our desolation.\end{Resp}
\begin{Resp}\Vsign Give ear, O Shepherd of Israel, Thou That leadest Joseph like a flock.\end{Resp}
\begin{Resp}\Rsign Open thine eyes, and behold our desolation.\end{Resp}
\switchcolumn*

\LA
\LessonHead{Lectio III}
\switchcolumn
\EN
\LessonHead{Lesson III}
\switchcolumn*

\LA
\Cite{Dan 1:10-15}
\switchcolumn
\EN
\Cite{Dan 1:10-15}
\switchcolumn*

\LA
\noindent Et ait princeps eunuchórum ad Daniélem: Tímeo ego dóminum meum regem, qui constítuit vobis cibum et potum: qui, si víderit vultus vestros macilentióres præ céteris adolescéntibus coǽvis vestris, condemnábitis caput meum regi. Et dixit Dániel ad Málasar, quem constitúerat princeps eunuchórum super Daniélem, Ananíam, Misaélem et Azaríam: Tenta nos, óbsecro, servos tuos, diébus decem, et dentur nobis legúmina ad vescéndum, et aqua ad bibéndum: et contempláre vultus nostros, et vultus puerórum, qui vescúntur cibo régio: et, sicut víderis, fácies cum servis tuis. Qui, audíto sermóne hujuscémodi, tentávit eos diébus decem. Post dies autem decem, apparuérunt vultus eórum melióres, et corpulentióres præ ómnibus púeris, qui vescebántur cibo régio.\par
\switchcolumn
\EN
\noindent And the prince of the eunuchs said to Daniel: I fear my lord the king, who hath appointed you meat and drink: who if he should see your faces leaner than those of the other youths your equals, you shall endanger my head to the king. And Daniel said to Malasar, whom the prince of the eunuchs had appointed over Daniel, Ananias, Misael, and Azarias: Try, I beseech thee, thy servants for ten days, and let pulse be given us to eat, and water to drink: And look upon our faces, and the faces of the children that eat of the king's meat: and as thou shalt see, deal with thy servants. And when he had heard these words, he tried them for ten days. And after ten days their faces appeared fairer and fatter than all the children that ate of the king's meat.\par
\switchcolumn*

\LA
\begin{Resp}\Rsign Aspice, Dómine, quia facta est desoláta cívitas plena divítiis, sedet in tristítia dómina géntium: \Star{} Non est qui consolétur eam, nisi tu, Deus noster.\end{Resp}
\begin{Resp}\Vsign Plorans plorávit in nocte, et lácrimæ ejus in maxíllis ejus.\end{Resp}
\begin{Resp}\Rsign Non est qui consolétur eam, nisi tu, Deus noster.\end{Resp}
\begin{Resp}\Vsign Glória Patri, et Fílio, \Star{} et Spirítui Sancto.\end{Resp}
\begin{Resp}\Rsign Non est qui consolétur eam, nisi tu, Deus noster.\end{Resp}
\switchcolumn
\EN
\begin{Resp}\Rsign Consider, O Lord, how that the city sitteth solitary that was full of riches; how is she become as a widow, she that was great among the nations; \Star{} She hath none to comfort her, save thee, O our God.\end{Resp}
\begin{Resp}\Vsign She weepeth sore in the night, and her tears are on her cheeks.\end{Resp}
\begin{Resp}\Rsign She hath none to comfort her, save thee, O our God.\end{Resp}
\begin{Resp}\Vsign Glory be to the Father, and to the Son, \Star{} and to the Holy Ghost.\end{Resp}
\begin{Resp}\Rsign She hath none to comfort her, save thee, O our God.\end{Resp}
\switchcolumn*

\LA
\LessonHead{Lectio IV}
\switchcolumn
\EN
\LessonHead{Lesson IV}
\switchcolumn*

\LA
\LessonTitle{Ex libro sancti Athanásii Epíscopi ad Vírgines}
\switchcolumn
\EN
\LessonTitle{From the Book addressed To Virgins by St. Athanasius, Pope (of Alexandria.)}
\switchcolumn*

\LA
\Source{Liber de Virginitate, post initium.}
\switchcolumn
\EN
\Source{Bk. ii.}
\switchcolumn*

\LA
\noindent Si accédant áliqui, e dicant tibi: Ne frequénter jejúnes, ne imbecíllior fias; ne credas illis, neque auscúltes: per istos enim inimícus hæc súggerit. Reminíscere ejus quod scriptum est, quod, cum tres púeri, et Dániel, et álii adolescéntuli, captívi ducti essent a Nabuchodónosor rege Babylónis, jussúmque esset ut de ipsíus mensa régia coméderent et de vino bíberent; Dániel et tres púeri illi noluérunt póllui ex mensa regis, sed dixérunt eunúcho qui eos curándos suscéperat: Da nobis de semínibus terræ, et vescémur. Quibus ait eunúchus: Tímeo ego regem, qui constítuit vobis cibum et potum, ne forte fácies vestræ appáreant regi squalidióres præ céteris púeris qui régia mensa alúntur, et púniat me.\par
\switchcolumn
\EN
\noindent If any should come and say unto thee, Fast not so often, lest thou injure thine health, believe them not, neither listen to them. They are but the tools of the great enemy to suggest such a thing unto thee. Remember how it is written that when the three children, and Daniel, and the other lads, were led captives by Nebuchadnezzar King of Babylon, and it was commanded them to eat of his Royal table, and to drink of his wine, Daniel and those three children would not defile themselves with the King's table, but said unto the eunuch into whose keeping they had been given, Give us of the fruits of the earth, and we will eat. And the eunuch answered them, I fear my lord the King, who hath appointed your meat and your drink, lest perchance your faces should appear unto the King worse - liking than the other children, who are fed from his Royal table, and he should punish me.\par
\switchcolumn*

\LA
\begin{Resp}\Rsign Super muros tuos, Jerúsalem, constítui custódes; \Star{} Tota die et nocte non tacébunt laudáre nomen Dómini.\end{Resp}
\begin{Resp}\Vsign Prædicábunt pópulis fortitúdinem meam, et annuntiábunt géntibus glóriam meam.\end{Resp}
\begin{Resp}\Rsign Tota die et nocte non tacébunt laudáre nomen Dómini.\end{Resp}
\switchcolumn
\EN
\begin{Resp}\Rsign I have set watchmen upon thy walls, O Jerusalem; \Star{} Which shall never hold their peace day nor night, to praise the name of the Lord.\end{Resp}
\begin{Resp}\Vsign They shall proclaim My might unto the nations, and declare My glory unto the Gentiles.\end{Resp}
\begin{Resp}\Rsign Which shall never hold their peace day nor night, to praise the name of the Lord.\end{Resp}
\switchcolumn*

\LA
\LessonHead{Lectio V}
\switchcolumn
\EN
\LessonHead{Lesson V}
\switchcolumn*

\LA
\noindent Cui dixérunt illi: Tenta servos tuos dies decem, et da nobis de semínibus. Et dedit eis legúmina ad vescéndum, et aquam ad bibéndum; et introdúxit eos in conspéctu regis, et visæ sunt fácies ipsórum speciosióres præter céteros púeros, qui régiæ mensæ cibis nutriebántur. Vidésne quid fáciat jejúnium? Morbos sanat, distillatiónes córporis exsíccat, dǽmones fugat, pravas cogitatiónes expéllit, mentem clariórem reddit, cor mundum éfficit, corpus sanctíficat, dénique ad thronum Dei hóminem sistit. Et ne putes hæc témere dici; habes hujus rei testimónium in Evangéliis a Salvatóre prolátum. Cum enim quæsivíssent discípuli quonam modo immúndi spíritus ejiceréntur, respóndit Dóminus: Hoc genus non ejícitur, nisi in oratióne et jejúnio.\par
\switchcolumn
\EN
\noindent Then they said unto him Prove thy servants ten days, and give us herbs. And he gave them pulse to eat and water to drink and, when he brought them in before the King, their countenances appeared fairer than all the children which did eat the portion of the King's meat. Seest thou what fasting doth It healeth diseases, it drieth up the humours of the body, it scareth away devils, it purgeth forth unclean thoughts, it maketh the intellect clearer, it purifieth the heart, it sanctifieth the body, and in the end it leadeth a man unto the throne of God. Think not that this is rash talking. Thou hast the testimony of this in the Gospels under the sanction of the Saviour Himself. His disciples asked Him why they could not cast out an evil spirit, and He said unto them This kind can come forth by nothing but by prayer and fasting. (Mark ix. 28.)\par
\switchcolumn*

\LA
\begin{Resp}\Rsign Muro tuo inexpugnábili circumcínge nos, Dómine, et armis tuæ poténtiæ prótege nos semper: \Star{} Líbera, Dómine, Deus Israël, clamántes ad te.\end{Resp}
\begin{Resp}\Vsign Erue nos in mirabílibus tuis, et da glóriam nómini tuo.\end{Resp}
\begin{Resp}\Rsign Líbera, Dómine, Deus Israël, clamántes ad te.\end{Resp}
\switchcolumn
\EN
\begin{Resp}\Rsign Hedge us about with thy wall that cannot be broken down, O Lord, and shield us continually with the arms of thy might. \Star{} O Lord God of Israel, deliver them that cry unto thee.\end{Resp}
\begin{Resp}\Vsign Deliver us also according to thy marvellous works, and give glory to thy Name.\end{Resp}
\begin{Resp}\Rsign O Lord God of Israel, deliver them that cry unto thee.\end{Resp}
\switchcolumn*

\LA
\LessonHead{Lectio VI}
\switchcolumn
\EN
\LessonHead{Lesson VI}
\switchcolumn*

\LA
\noindent Quisquis ígitur ab immúndo spíritu vexátur, si hoc animadvértat, et hoc phármaco utátur, jejúnio inquam, statim spíritus malus oppréssus abscédet, vim jejúnii métuens. Valde enim dǽmones oblectántur crápula et ebrietáte et córporis cómmodis. Magna vis in jejúnio, et magna ac præclára fiunt per illud. Alióquin unde hómines tam mirífica præstárent, et signa per eos fíerent, et sanitátem infírmis per ipsos largirétur Deus, nisi plane ob exercitatiónes spirituáles, et humilitátem ánimi, et conversatiónem bonam? Jejúnium enim Angelórum cibus est: et qui eo útitur, órdinis angélici censéndus est.\par
\switchcolumn
\EN
\noindent If any man therefore be troubled with an unclean spirit, if he bethink him of this, and have recourse to this remedy, namely, fasting, the evil spirit will be forthwith compelled to leave him from dread of the power of fasting. Devils take great delight in fulness, and drunkenness, and bodily comfort. There is great power in fasting, and great and glorious things are wrought thereby. How cometh it that men work such wonders, and that signs are done by them, and that God through them giveth health to the sick, unless it be from their ghostly exercises, and the meekness of their souls, and their godly conversation To fast is to banquet with Angels, and he that fasteth is to be reckoned, so far, among the Angelic host.\par
\switchcolumn*

\LA
\begin{Resp}\Rsign Sustinúimus pacem, et non venit: quæsívimus bona, et ecce turbátio: cognóvimus, Dómine, peccáta nostra; \Star{} Non in perpétuum obliviscáris nos.\end{Resp}
\begin{Resp}\Vsign Peccávimus, ímpie géssimus, iniquitátem fécimus, Dómine, in omnem justítiam tuam.\end{Resp}
\begin{Resp}\Rsign Non in perpétuum obliviscáris nos.\end{Resp}
\begin{Resp}\Vsign Glória Patri, et Fílio, \Star{} et Spirítui Sancto.\end{Resp}
\begin{Resp}\Rsign Non in perpétuum obliviscáris nos.\end{Resp}
\switchcolumn
\EN
\begin{Resp}\Rsign We looked for peace, and it came not we asked for good, and behold trouble. We acknowledge, O Lord, our wickedness. \Star{} Forget us not for ever.\end{Resp}
\begin{Resp}\Vsign O Lord (our God) we have sinned, we have done ungodly, we have dealt unrighteously in all thine ordinances.\end{Resp}
\begin{Resp}\Rsign Forget us not for ever.\end{Resp}
\begin{Resp}\Vsign Glory be to the Father, and to the Son, \Star{} and to the Holy Ghost.\end{Resp}
\begin{Resp}\Rsign Forget us not for ever.\end{Resp}
\switchcolumn*

\end{paracol}

\Office{Feria II infra Hebdomadam III Novembris}{Monday of the third week of November}{}
\addcontentsline{toc}{subsection}{Feria II infra Hebdomadam III Novembris \textit{— Monday of the third week of November}}
\begin{paracol}{2}
\LA
\LessonHead{Lectio I}
\switchcolumn
\EN
\LessonHead{Lesson I}
\switchcolumn*

\LA
\LessonTitle{De Daniéle Prophéta}
\switchcolumn
\EN
\LessonTitle{Lesson from the book of Daniel}
\switchcolumn*

\LA
\Cite{Dan 2:31-35}
\switchcolumn
\EN
\Cite{Dan 2:31-35}
\switchcolumn*

\LA
\noindent Tu, rex, vidébas, et ecce quasi státua una grandis: státua illa magna, et statúra sublímis stabat contra te, et intúitus ejus erat terríbilis. Hujus státuæ caput ex auro óptimo erat, pectus autem et brácchia de argénto, porro venter et fémora ex ære, tíbiæ autem férreæ: pedum quædam pars erat férrea, quædam autem fíctilis. Vidébas ita, donec abscíssus est lapis de monte sine mánibus: et percússit státuam in pédibus ejus férreis et fictílibus, et commínuit eos. Tunc contríta sunt páriter ferrum, testa, æs, argéntum et aurum et redácta quasi in favíllam æstívæ áreæ, quæ rapta sunt vento, nullúsque locus invéntus est eis; lapis autem, qui percússerat státuam, factus est mons magnus, et implévit univérsam terram.\par
\switchcolumn
\EN
\noindent Thou, O king, sawest, and behold there was as it were a great statue: this statue, which was great and high, tall of stature, stood before thee, and the look thereof was terrible. The head of this statue was of fine gold, but the breast and the arms of silver, and the belly and the thighs of brass: And the legs of iron, the feet part of iron and part of clay. Thus thou sawest, till a stone was cut out of a mountain without hands: and it struck the statue upon the feet thereof that were of iron and of clay, and broke them in pieces. Then was the iron, the clay, the brass, the silver, and the gold broken to pieces together, and became like the chaff of a summer's thrashingfloor, and they were carried away by the wind: and there was no place found for them: but the stone that struck the statue, became a great mountain, and filled the whole earth.\par
\switchcolumn*

\LA
\begin{Resp}\Rsign Redémit pópulum suum et liberávit eum, et vénient et exsultábunt in monte Sion et gaudébunt de bonis Dómini super fruménto, vino et óleo, \Star{} Et ultra non esúrient.\end{Resp}
\begin{Resp}\Vsign Erítque ánima eórum quasi hortus irríguus.\end{Resp}
\begin{Resp}\Rsign Et ultra non esúrient.\end{Resp}
\switchcolumn
\EN
\begin{Resp}\Rsign He hath redeemed His people, and ransomed them therefore they shall come and sing in the height of Zion, and shall rejoice in the goodness of the Lord, for wheat, and for wine, and for oil \Star{} And they shall hunger no more.\end{Resp}
\begin{Resp}\Vsign And their soul shall be as a watered garden.\end{Resp}
\begin{Resp}\Rsign And they shall hunger no more.\end{Resp}
\switchcolumn*

\LA
\LessonHead{Lectio II}
\switchcolumn
\EN
\LessonHead{Lesson II}
\switchcolumn*

\LA
\Cite{Dan 2:36-40}
\switchcolumn
\EN
\Cite{Dan 2:36-40}
\switchcolumn*

\LA
\noindent Hoc est sómnium. Interpretatiónem quoque ejus dicémus coram te, rex. Tu rex regum es: et Deus cæli regnum, et fortitúdinem, et impérium, et glóriam dedit tibi: et ómnia, in quibus hábitant fílii hóminum, et béstiæ agri: vólucres quoque cæli dedit in manu tua, et sub dicióne tua univérsa constítuit; tu es ergo caput áureum. Et post te consúrget regnum áliud minus te argénteum: et regnum tértium áliud ǽreum, quod imperábit univérsæ terræ, et regnum quartum erit velut ferrum. Quómodo ferrum commínuit, et domat ómnia, sic commínuet, et cónteret ómnia hæc.\par
\switchcolumn
\EN
\noindent This is the dream: we will also tell the interpretation thereof before thee, O king. Thou art a king of kings: and the God of heaven hath given thee a kingdom, and strength, and power, and glory: And all places wherein the children of men, and the beasts of the field do dwell: he hath also given the birds of the air into thy hand, and hath put all things under thy power: thou therefore art the head of gold. And after thee shall rise up another kingdom, inferior to thee, of silver: and another third kingdom of brass, which shall rule over all the world. And the fourth kingdom shall be as iron. As iron breaketh into pieces, and subdueth all things, so shall that break and destroy all these.\par
\switchcolumn*

\LA
\begin{Resp}\Rsign Angústiæ mihi sunt úndique, et quid éligam ignóro; \Star{} Mélius est mihi incídere in manus hóminum, quam derelínquere legem Dei mei.\end{Resp}
\begin{Resp}\Vsign Si enim hoc égero, mors mihi est; si autem non égero, non effúgiam manus vestras.\end{Resp}
\begin{Resp}\Rsign Mélius est mihi incídere in manus hóminum, quam derelínquere legem Dei mei.\end{Resp}
\switchcolumn
\EN
\begin{Resp}\Rsign I am straitened on every side, and know not what to choose. \Star{} It is better for me to fall into the hands of men, than to sin against the law of my God.\end{Resp}
\begin{Resp}\Vsign For if I do this thing, it is death unto me and if I do it not, I cannot escape your hands.\end{Resp}
\begin{Resp}\Rsign It is better for me to fall into the hands of men, than to sin against the law of my God.\end{Resp}
\switchcolumn*

\LA
\LessonHead{Lectio III}
\switchcolumn
\EN
\LessonHead{Lesson III}
\switchcolumn*

\LA
\Cite{Dan 2:41-44}
\switchcolumn
\EN
\Cite{Dan 2:41-44}
\switchcolumn*

\LA
\noindent Porro, quia vidísti pedum, et digitórum partem testæ fíguli, et partem férream, regnum divísum erit; quod tamen de plantário ferri oriétur, secúndum quod vidísti ferrum mistum testæ ex luto. Et dígitos pedum ex parte férreos, et ex parte fíctiles, ex parte regnum erit sólidum, et ex parte contrítum. Quod autem vidísti ferrum mistum testæ ex luto, commiscebúntur quidem humáno sémine; sed non adhærébunt sibi, sícuti ferrum miscéri non potest testæ. In diébus autem regnórum illórum suscitábit Deus cæli regnum, quod in ætérnum non dissipábitur, et regnum ejus álteri pópulo non tradétur: commínuet autem, et consúmet univérsa regna hæc, et ipsum stabit in ætérnum.\par
\switchcolumn
\EN
\noindent And whereas thou sawest the feet, and the toes, part of potter's clay, and part of iron: the kingdom shall be divided, but yet it shall take its origin from the iron, according as thou sawest the iron mixed with the miry clay. And as the toes of the feet were part of iron, and part of clay, the kingdom shall be partly strong, and partly broken. And whereas thou sawest the iron mixed with miry clay, they shall be mingled indeed together with the seed of man, but they shall not stick fast one to another, as iron cannot be mixed with clay. But in the days of those kingdoms the God of heaven will set up a kingdom that shall never be destroyed, and his kingdom shall not be delivered up to another people, and it shall break in pieces, and shall consume all these kingdoms, and itself shall stand for ever.\par
\switchcolumn*

\LA
\begin{Resp}\Rsign Misit Dóminus Angelum suum et conclúsit ora leónum, \Star{} Et non contaminavérunt, quia coram eo injustítia invénta non est in me.\end{Resp}
\begin{Resp}\Vsign Misit Deus misericórdiam suam et veritátem suam: ánimam meam erípuit de médio catulórum leónum.\end{Resp}
\begin{Resp}\Rsign Et non contaminavérunt, quia coram eo injustítia invénta non est in me.\end{Resp}
\begin{Resp}\Vsign Glória Patri, et Fílio, \Star{} et Spirítui Sancto.\end{Resp}
\begin{Resp}\Rsign Et non contaminavérunt: quia coram eo injustítia invénta non est in me.\end{Resp}
\switchcolumn
\EN
\begin{Resp}\Rsign The Lord hath sent His angel, and hath shut the lions' mouths, that \Star{} They have not hurt me forasmuch as before Him innocency was found in me.\end{Resp}
\begin{Resp}\Vsign God hath sent forth His mercy and His truth, (and delivered) my soul from among the lions' whelps.\end{Resp}
\begin{Resp}\Rsign They have not hurt me forasmuch as before Him innocency was found in me.\end{Resp}
\begin{Resp}\Vsign Glory be to the Father, and to the Son, \Star{} and to the Holy Ghost.\end{Resp}
\begin{Resp}\Rsign They have not hurt me forasmuch as before Him innocency was found in me.\end{Resp}
\switchcolumn*

\end{paracol}

\Office{Feria III infra Hebdomadam III Novembris}{Tuesday of the third week of November}{}
\addcontentsline{toc}{subsection}{Feria III infra Hebdomadam III Novembris \textit{— Tuesday of the third week of November}}
\begin{paracol}{2}
\LA
\LessonHead{Lectio I}
\switchcolumn
\EN
\LessonHead{Lesson I}
\switchcolumn*

\LA
\LessonTitle{De Daniéle Prophéta}
\switchcolumn
\EN
\LessonTitle{Lesson from the book of Daniel}
\switchcolumn*

\LA
\Cite{Dan 3:14-15}
\switchcolumn
\EN
\Cite{Dan 3:14-15}
\switchcolumn*

\LA
\noindent Pronuntiánsque Nabuchodónosor rex, ait eis: Veréne, Sidrach, Misach et Abdénago, deos meos non cólitis, et státuam áuream, quam constítui, non adorátis? Nunc ergo si estis paráti, quacúmque hora audiéritis sónitum tubæ, fístulæ, cítharæ, sambúcæ, et psaltérii, et symphóniæ, omnísque géneris musicórum, prostérnite vos et adoráte státuam, quam feci. Quod si non adoravéritis, eádem hora mittémini in fornácem ignis ardéntis, et quis est Deus, qui erípiet vos de manu mea?\par
\switchcolumn
\EN
\noindent And Nabuchodonosor the king spoke to them, and said: Is it true, O Sidrach, Misach, and Abdenago, that you do not worship my gods, nor adore the golden statue that I have set up? Now therefore if you be ready at what hour soever you shall hear the sound of the trumpet, flute, harp, sackbut, and psaltery, and symphony, and of all kind of music, prostrate yourselves, and adore the statue which I have made: but if you do not adore, you shall be cast the same hour into the furnace of burning fire: and who is the God that shall deliver you out of my hand?\par
\switchcolumn*

\LA
\begin{Resp}\Rsign A fácie furóris tui, Deus, conturbáta est omnis terra: \Star{} Sed miserére, Dómine, et ne fácias consummatiónem.\end{Resp}
\begin{Resp}\Vsign Dómine, Dóminus noster, quam admirábile est nomen tuum!\end{Resp}
\begin{Resp}\Rsign Sed miserére, Dómine, et ne fácias consummatiónem.\end{Resp}
\switchcolumn
\EN
\begin{Resp}\Rsign Before the face of thine anger, O God, the whole earth is troubled; \Star{} But Thou, O Lord, have mercy, and make not an end utterly.\end{Resp}
\begin{Resp}\Vsign Lord our Ruler, how excellent is thy Name in all the earth!\end{Resp}
\begin{Resp}\Rsign But Thou, O Lord, have mercy, and make not an end utterly.\end{Resp}
\switchcolumn*

\LA
\LessonHead{Lectio II}
\switchcolumn
\EN
\LessonHead{Lesson II}
\switchcolumn*

\LA
\Cite{Dan 3:16-19}
\switchcolumn
\EN
\Cite{Dan 3:16-19}
\switchcolumn*

\LA
\noindent Respondéntes Sidrach, Misach et Abdénago, dixérunt regi Nabuchodónosor: Non opórtet nos de hac re respondére tibi: ecce enim Deus noster, quem cólimus, potest erípere nos de camíno ignis ardéntis, et de mánibus tuis, o rex, liberáre. Quod si nolúerit, notum sit tibi, rex, quia deos tuos non cólimus, et státuam áuream, quam erexísti, non adorámus. Tunc Nabuchodónosor replétus est furóre, et aspéctus faciéi illíus immutátus est super Sidrach, Misach et Abdénago, et præcépit ut succenderétur fornax séptuplum quam succéndi consuéverat.\par
\switchcolumn
\EN
\noindent Sidrach, Misach, and Abdenago answered and said to king Nabuchodonosor: We have no occasion to answer thee concerning this matter. For behold our God, whom we worship, is able to save us from the furnace of burning fire, and to deliver us out of thy hands, O king. But if he will not, be it known to thee, O king, that we will not worship thy gods, nor adore the golden statue which thou hast set up. Then was Nabuchodonosor filled with fury: and the countenance of his face was changed against Sidrach, Misach, and Abdenago, and he commanded that the furnace should be heated seven times more than it had been accustomed to be heated.\par
\switchcolumn*

\LA
\begin{Resp}\Rsign Civitátem istam tu circúmda, Dómine: et Angeli tui custódiant muros ejus. \Star{} Exáudi, Dómine, pópulum tuum cum misericórdia.\end{Resp}
\begin{Resp}\Vsign Avertátur furor tuus, Dómine, a pópulo tuo et a civitáte sancta tua.\end{Resp}
\begin{Resp}\Rsign Exáudi, Dómine, pópulum tuum cum misericórdia.\end{Resp}
\switchcolumn
\EN
\begin{Resp}\Rsign Fence Thou this city; O Lord, and let thine angels keep the walls thereof. \Star{} O Lord, hearken unto thy people with mercy.\end{Resp}
\begin{Resp}\Vsign O Lord, let thine anger be turned away from thy people, and from thine holy city.\end{Resp}
\begin{Resp}\Rsign O Lord, hearken unto thy people with mercy.\end{Resp}
\switchcolumn*

\LA
\LessonHead{Lectio III}
\switchcolumn
\EN
\LessonHead{Lesson III}
\switchcolumn*

\LA
\Cite{Dan 3:21-24}
\switchcolumn
\EN
\Cite{Dan 3:21-24}
\switchcolumn*

\LA
\noindent Et conféstim viri illi vincti, cum braccis suis, et tiáris, et calceaméntis, et véstibus, missi sunt in médium fornácis ignis ardéntis; nam jússio regis urgébat. Fornax autem succénsa erat nimis; porro viros illos, qui míserant Sidrach, Misach et Abdénago, interfécit flamma ignis. Viri autem hi tres, id est Sidrach, Misach et Abdénago, cecidérunt in médio camíno ignis ardéntis, colligáti. Et ambulábant in médio flammæ, laudántes Deum, et benedicéntes Dómino.\par
\switchcolumn
\EN
\noindent And immediately these men were bound and were cast into the furnace of burning fire, with their coats, and their caps, and their shoes, and their garments. For the king's commandment was urgent, and the furnace was heated exceedingly. And the flame of the fire slew those men that had cast in Sidrach, Misach, and Abdenago. But these three men, that is, Sidrach, Misach, and Abdenago, fell down bound in the midst of the furnace of burning fire. And they walked in the midst of the flame, praising God and blessing the Lord.\par
\switchcolumn*

\LA
\begin{Resp}\Rsign Genti peccatríci, pópulo pleno peccáto miserére, \Star{} Dómine Deus.\end{Resp}
\begin{Resp}\Vsign Esto placábilis super nequítiam pópuli tui.\end{Resp}
\begin{Resp}\Rsign Dómine Deus.\end{Resp}
\begin{Resp}\Vsign Glória Patri, et Fílio, \Star{} et Spirítui Sancto.\end{Resp}
\begin{Resp}\Rsign Dómine Deus.\end{Resp}
\switchcolumn
\EN
\begin{Resp}\Rsign O Lord God! have mercy upon \Star{} The sinful nation, upon the people laden with iniquity.\end{Resp}
\begin{Resp}\Vsign Let it repent thee concerning the transgression of thy people.\end{Resp}
\begin{Resp}\Rsign The people laden with iniquity.\end{Resp}
\begin{Resp}\Vsign Glory be to the Father, and to the Son, \Star{} and to the Holy Ghost.\end{Resp}
\begin{Resp}\Rsign The people laden with iniquity.\end{Resp}
\switchcolumn*

\end{paracol}

\Office{Feria IV infra Hebdomadam III Novembris}{Wednesday of the third week of November}{}
\addcontentsline{toc}{subsection}{Feria IV infra Hebdomadam III Novembris \textit{— Wednesday of the third week of November}}
\begin{paracol}{2}
\LA
\LessonHead{Lectio I}
\switchcolumn
\EN
\LessonHead{Lesson I}
\switchcolumn*

\LA
\LessonTitle{De Daniéle Prophéta}
\switchcolumn
\EN
\LessonTitle{Lesson from the book of Daniel}
\switchcolumn*

\LA
\Cite{Dan 4:16-19}
\switchcolumn
\EN
\Cite{Dan 4:16-19}
\switchcolumn*

\LA
\noindent Respóndit Baltássar, et dixit: Dómine mi, sómnium his, qui te odérunt, et interpretátio ejus hóstibus tuis sit. Arborem, quam vidísti sublímem atque robústam, cujus altitúdo pertíngit ad cælum, et aspéctus illíus in omnem terram, et rami ejus pulchérrimi, et fructus ejus nímius, et esca ómnium in ea, subter eam habitántes béstiæ agri, et in ramis ejus commorántes aves cæli, tu es, rex, qui magnificátus es, et invaluísti, et magnitúdo tua crevit, et pervénit usque ad cælum, et potéstas tua in términos univérsæ terræ.\par
\switchcolumn
\EN
\noindent Baltassar answered, and said: My lord, the dream be to them that hate thee, and the interpretation thereof to thy enemies. The tree which thou sawest which was high and strong, whose height reached to the skies, and the sight thereof into all the earth: And the branches thereof were most beautiful, and its fruit exceeding much, and in it was food for all, under which the beasts of the field dwelt, and the birds of the air had their abode in its branches. It is thou, O king, who art grown great and become mighty: for thy greatness hath grown, and hath reached to heaven, and thy power unto the ends of the earth.\par
\switchcolumn*

\LA
\begin{Resp}\Rsign Indicábo tibi, homo, quid sit bonum aut quid Dóminus requírat a te: \Star{} Fácere judícium et justítiam et sollícitum ambuláre cum Deo tuo.\end{Resp}
\begin{Resp}\Vsign Spera in Dómino, et fac bonitátem, et inhábita terram.\end{Resp}
\begin{Resp}\Rsign Fácere judícium et justítiam et sollícitum ambuláre cum Deo tuo.\end{Resp}
\switchcolumn
\EN
\begin{Resp}\Rsign I will show thee, O man, what is good, and what doth the Lord require of thee, but to \Star{} Do justice and judgment, and to walk humbly with thy God.\end{Resp}
\begin{Resp}\Vsign Trust in the Lord, and do good, and dwell in the land.\end{Resp}
\begin{Resp}\Rsign Do justice and judgment, and walk humbly with thy God.\end{Resp}
\switchcolumn*

\LA
\LessonHead{Lectio II}
\switchcolumn
\EN
\LessonHead{Lesson II}
\switchcolumn*

\LA
\Cite{Dan 4:20-22}
\switchcolumn
\EN
\Cite{Dan 4:20-22}
\switchcolumn*

\LA
\noindent Quod autem vidit rex vígilem, et sanctum descéndere de cælo, et dícere: Succídite árborem, et dissipáte illam, áttamen germen radícum ejus in terra dimíttite, et vinciátur ferro et ære in herbis foris, et rore cæli conspergátur, et cum feris sit pábulum ejus, donec septem témpora muténtur super eum, hæc est interpretátio senténtiæ Altíssimi, quæ pervénit super dóminum meum regem: Eícient te ab homínibus, et cum béstiis ferísque erit habitátio tua, et fœnum ut bos cómedes, et rore cæli infundéris.\par
\switchcolumn
\EN
\noindent And whereas the king saw a watcher, and a holy one come down from heaven, and say: Cut down the tree and destroy it, but leave the stump of the roots thereof in the earth, and let it be bound with iron and brass among the grass without, and let it be sprinkled with the dew of heaven, and let his feeding be with the wild beasts, till seven times pass over him. This is the interpretation of the sentence of the most High, which is come upon my lord the king. They shall cast thee out from among men, and thy dwelling shall be with cattle and with wild beasts, and thou shalt eat grass as an ox, and shalt be wet with the dew of heaven:\par
\switchcolumn*

\LA
\begin{Resp}\Rsign Angústiæ mihi sunt úndique, et quid éligam ignóro; \Star{} Mélius est mihi incídere in manus hóminum, quam derelínquere legem Dei mei.\end{Resp}
\begin{Resp}\Vsign Si enim hoc égero, mors mihi est; si autem non égero, non effúgiam manus vestras.\end{Resp}
\begin{Resp}\Rsign Mélius est mihi incídere in manus hóminum, quam derelínquere legem Dei mei.\end{Resp}
\switchcolumn
\EN
\begin{Resp}\Rsign I am straitened on every side, and know not what to choose. \Star{} It is better for me to fall into the hands of men, than to sin against the law of my God.\end{Resp}
\begin{Resp}\Vsign For if I do this thing, it is death unto me and if I do it not, I cannot escape your hands.\end{Resp}
\begin{Resp}\Rsign It is better for me to fall into the hands of men, than to sin against the law of my God.\end{Resp}
\switchcolumn*

\LA
\LessonHead{Lectio III}
\switchcolumn
\EN
\LessonHead{Lesson III}
\switchcolumn*

\LA
\Cite{Dan 4:22-25}
\switchcolumn
\EN
\Cite{Dan 4:22-25}
\switchcolumn*

\LA
\noindent Septem quoque témpora mutabúntur super te, donec scias quod dominétur Excélsus super regnum hóminum, et cuicúmque volúerit, det illud. Quod autem præcépit ut relinquerétur germen radícum ejus, id est árboris: regnum tuum tibi manébit, postquam cognóveris potestátem esse cæléstem. Quam ob rem, rex, consílium meum pláceat tibi, et peccáta tua eleemósynis rédime, et iniquitátes tuas misericórdiis páuperum; fórsitan ignóscet delíctis tuis. Omnia hæc venérunt super Nabuchodónosor regem.\par
\switchcolumn
\EN
\noindent And seven times shall pass over thee, till thou know that the most High ruleth over the kingdom of men, and giveth it to whomsoever he will. But whereas he commanded, that the stump of the roots thereof, that is, of the tree, should be left: thy kingdom shall remain to thee after thou shalt have known that power is from heaven. Wherefore, O king, let my counsel be acceptable to thee, and redeem thou thy sins with alms, and thy iniquities with works of mercy to the poor: perhaps he will forgive thy offences. All these things came upon king Nabuchodonosor.\par
\switchcolumn*

\LA
\begin{Resp}\Rsign Misit Dóminus Angelum suum et conclúsit ora leónum, \Star{} Et non contaminavérunt, quia coram eo injustítia invénta non est in me.\end{Resp}
\begin{Resp}\Vsign Misit Deus misericórdiam suam et veritátem suam: ánimam meam erípuit de médio catulórum leónum.\end{Resp}
\begin{Resp}\Rsign Et non contaminavérunt, quia coram eo injustítia invénta non est in me.\end{Resp}
\begin{Resp}\Vsign Glória Patri, et Fílio, \Star{} et Spirítui Sancto.\end{Resp}
\begin{Resp}\Rsign Et non contaminavérunt: quia coram eo injustítia invénta non est in me.\end{Resp}
\switchcolumn
\EN
\begin{Resp}\Rsign The Lord hath sent His angel, and hath shut the lions' mouths, that \Star{} They have not hurt me forasmuch as before Him innocency was found in me.\end{Resp}
\begin{Resp}\Vsign God hath sent forth His mercy and His truth, (and delivered) my soul from among the lions' whelps.\end{Resp}
\begin{Resp}\Rsign They have not hurt me forasmuch as before Him innocency was found in me.\end{Resp}
\begin{Resp}\Vsign Glory be to the Father, and to the Son, \Star{} and to the Holy Ghost.\end{Resp}
\begin{Resp}\Rsign They have not hurt me forasmuch as before Him innocency was found in me.\end{Resp}
\switchcolumn*

\end{paracol}

\Office{Feria V infra Hebdomadam III Novembris}{Thursday of the third week of November}{}
\addcontentsline{toc}{subsection}{Feria V infra Hebdomadam III Novembris \textit{— Thursday of the third week of November}}
\begin{paracol}{2}
\LA
\LessonHead{Lectio I}
\switchcolumn
\EN
\LessonHead{Lesson I}
\switchcolumn*

\LA
\LessonTitle{De Daniéle Prophéta}
\switchcolumn
\EN
\LessonTitle{Lesson from the book of Daniel}
\switchcolumn*

\LA
\Cite{Dan 5:1-6}
\switchcolumn
\EN
\Cite{Dan 5:1-6}
\switchcolumn*

\LA
\noindent Baltássar rex fecit grande convívium optimátibus suis mille: et unusquísque secúndum suam bibébat ætátem. Præcépit ergo jam temuléntus ut afferréntur vasa áurea et argéntea, quæ asportáverat Nabuchodónosor pater ejus de templo, quod fuit in Jerúsalem, ut bíberent in eis rex, et optimátes ejus, uxorésque ejus, et concubínæ. Tunc alláta sunt vasa áurea, et argéntea, quæ asportáverat de templo, quod fúerat in Jerúsalem: et bibérunt in eis rex, et optimátes ejus, uxóres et concubínæ illíus. Bibébant vinum, et laudábant deos suos áureos et argénteos, ǽreos, férreos, ligneósque et lapídeos. In eádem hora apparuérunt dígiti, quasi manus hóminis scribéntis contra candelábrum in superfície paríetis aulæ régiæ: et rex aspiciébat artículos manus scribéntis. Tunc fácies regis commutáta est, et cogitatiónes ejus conturbábant eum: et compáges renum ejus solvebántur, et génua ejus ad se ínvicem collidebántur.\par
\switchcolumn
\EN
\noindent Baltasar the king made a great feast for a thousand of his nobles: and every one drank according to his age. And being now drunk he commanded that they should bring the vessels of gold and silver which Nabuchodonosor his father had brought away out of the temple, that was in Jerusalem, that the king and his nobles, and his wives and his concubines, might drink in them. Then were the golden and silver vessels brought, which he had brought away out of the temple that was in Jerusalem: and the king and his nobles, his wives and his concubines, drank in them. They drank wine, and praised their gods of gold, and of silver, of brass, of iron, and of wood, and of stone. In the same hour there appeared fingers, as it were of the hand of a man, writing over against the candlestick upon the surface of the wall of the king's palace: and the king beheld the joints of the hand that wrote. Then was the king's countenance changed, and his thoughts troubled him.\par
\switchcolumn*

\LA
\begin{Resp}\Rsign Vidi Dóminum sedéntem super sólium excélsum et elevátum, et plena erat omnis terra majestáte ejus: \Star{} Et ea, quæ sub ipso erant, replébant templum.\end{Resp}
\begin{Resp}\Vsign Séraphim stabant super illud: sex alæ uni, et sex alæ álteri.\end{Resp}
\begin{Resp}\Rsign Et ea, quæ sub ipso erant, replébant templum.\end{Resp}
\switchcolumn
\EN
\begin{Resp}\Rsign I saw the Lord sitting upon a throne, high and lifted up, and the whole earth was full of His glory; \Star{} And His train filled the temple.\end{Resp}
\begin{Resp}\Vsign Above it stood the Seraphim each one had six wings.\end{Resp}
\begin{Resp}\Rsign And His train filled the temple.\end{Resp}
\switchcolumn*

\LA
\LessonHead{Lectio II}
\switchcolumn
\EN
\LessonHead{Lesson II}
\switchcolumn*

\LA
\Cite{Dan 5:13-17}
\switchcolumn
\EN
\Cite{Dan 5:13-17}
\switchcolumn*

\LA
\noindent Igitur introdúctus est Dániel coram rege: ad quem præfátus rex ait: Tu es Dániel de fíliis captivitátis Judæ, quem addúxit pater meus rex de Judǽa? audívi de te, quóniam spíritum deórum hábeas, et sciéntia, intellegentiáque ac sapiéntia amplióres invéntæ sunt in te. Et nunc introgréssi sunt in conspéctu meo sapiéntes magi, ut scriptúram hanc légerent, et interpretatiónem ejus indicárent mihi: et nequivérunt sensum hujus sermónis edícere. Porro ego audívi de te, quod possis obscúra interpretári, et ligáta dissólvere: si ergo vales scriptúram légere, et interpretatiónem ejus indicáre mihi, púrpura vestiéris, et torquem áuream circa collum tuum habébis, et tértius in regno meo princeps eris. Ad quæ respóndens Dániel, ait coram rege: Múnera tua sint tibi, et dona domus tuæ álteri da: scriptúram autem legam tibi, rex, et interpretatiónem ejus osténdam tibi.\par
\switchcolumn
\EN
\noindent Then Daniel was brought in before the king. And the king spoke, and said to him: Art thou Daniel of the children of the captivity of Juda, whom my father the king brought out of Judea? I have heard of thee, that thou hast the spirit of the gods, and excellent knowledge, and understanding, and wisdom are found in thee. And now the wise men the magicians have come in before me, to read this writing, and shew me the interpretation thereof: and they could not declare to me the meaning of this writing. But I have heard of thee, that thou canst interpret obscure things, and resolve difficult things: now if thou art able to read the writing, and to shew me the interpretation thereof, thou shalt be clothed with purple, and shalt have a chain of gold about thy neck, and shalt be the third prince in my kingdom. To which Daniel made answer, and said before the king: thy rewards be to thyself, and the gifts of thy house give to another: but the writing I will read to thee, O king, and shew thee the interpretation thereof.\par
\switchcolumn*

\LA
\begin{Resp}\Rsign Aspice, Dómine, de sede sancta tua, et cógita de nobis: inclína, Deus meus, aurem tuam et audi: \Star{} Aperi óculos tuos et vide tribulatiónem nostram.\end{Resp}
\begin{Resp}\Vsign Qui regis Israël, inténde, qui dedúcis velut ovem Joseph.\end{Resp}
\begin{Resp}\Rsign Aperi óculos tuos et vide tribulatiónem nostram.\end{Resp}
\switchcolumn
\EN
\begin{Resp}\Rsign Look down, O Lord, from the dwelling-place of thine holiness, and take thought for us. O my God, incline thine ear, and hear. \Star{} Open thine eyes, and behold our desolation.\end{Resp}
\begin{Resp}\Vsign Give ear, O Shepherd of Israel, Thou That leadest Joseph like a flock.\end{Resp}
\begin{Resp}\Rsign Open thine eyes, and behold our desolation.\end{Resp}
\switchcolumn*

\LA
\LessonHead{Lectio III}
\switchcolumn
\EN
\LessonHead{Lesson III}
\switchcolumn*

\LA
\Cite{Dan 5:25-31}
\switchcolumn
\EN
\Cite{Dan 5:25-31}
\switchcolumn*

\LA
\noindent Hæc est autem scriptúra, quæ digésta est: Mane, Thecel, Phares. Et hæc est interpretátio sermónis. Mane: numerávit Deus regnum tuum, et complévit illud. Thecel: appénsus es in statéra, et invéntus es minus habens. Phares: divísum est regnum tuum, et datum est Medis, et Persis. Tunc, jubénte rege, indútus est Dániel púrpura, et circúmdata est torques áurea collo ejus: et prædicátum est de eo quod habéret potestátem tértius in regno suo. Eádem nocte interféctus est Baltássar rex Chaldǽus. Et Daríus Medus succéssit in regnum, annos natus sexagínta duos.\par
\switchcolumn
\EN
\noindent And this is the writing that is written: MANE, THECEL, PHARES. And this is the interpretation of the word. MANE: God hath numbered thy kingdom, and hath finished it. THECEL: thou art weighed in the balance, and art found wanting. PHARES: thy kingdom is divided, and is given to the Medes and Persians. Then by the king's command Daniel was clothed with purple, and a chain of gold was put about his neck: and it was proclaimed of him that he had power as the third man in the kingdom. The same night Baltasar the Chaldean king was slain. And Darius the Mede succeeded to the kingdom, being threescore and two years old.\par
\switchcolumn*

\LA
\begin{Resp}\Rsign Aspice, Dómine, quia facta est desoláta cívitas plena divítiis, sedet in tristítia dómina géntium: \Star{} Non est qui consolétur eam, nisi tu, Deus noster.\end{Resp}
\begin{Resp}\Vsign Plorans plorávit in nocte, et lácrimæ ejus in maxíllis ejus.\end{Resp}
\begin{Resp}\Rsign Non est qui consolétur eam, nisi tu, Deus noster.\end{Resp}
\begin{Resp}\Vsign Glória Patri, et Fílio, \Star{} et Spirítui Sancto.\end{Resp}
\begin{Resp}\Rsign Non est qui consolétur eam, nisi tu, Deus noster.\end{Resp}
\switchcolumn
\EN
\begin{Resp}\Rsign Consider, O Lord, how that the city sitteth solitary that was full of riches; how is she become as a widow, she that was great among the nations; \Star{} She hath none to comfort her, save thee, O our God.\end{Resp}
\begin{Resp}\Vsign She weepeth sore in the night, and her tears are on her cheeks.\end{Resp}
\begin{Resp}\Rsign She hath none to comfort her, save thee, O our God.\end{Resp}
\begin{Resp}\Vsign Glory be to the Father, and to the Son, \Star{} and to the Holy Ghost.\end{Resp}
\begin{Resp}\Rsign She hath none to comfort her, save thee, O our God.\end{Resp}
\switchcolumn*

\end{paracol}

\Office{Feria VI infra Hebdomadam III Novembris}{Friday of the third week of November}{}
\addcontentsline{toc}{subsection}{Feria VI infra Hebdomadam III Novembris \textit{— Friday of the third week of November}}
\begin{paracol}{2}
\LA
\LessonHead{Lectio I}
\switchcolumn
\EN
\LessonHead{Lesson I}
\switchcolumn*

\LA
\LessonTitle{De Daniéle Prophéta}
\switchcolumn
\EN
\LessonTitle{Lesson from the book of Daniel}
\switchcolumn*

\LA
\Cite{Dan 6:11-15}
\switchcolumn
\EN
\Cite{Dan 6:11-15}
\switchcolumn*

\LA
\noindent Viri ergo illi curiósius inquiréntes invenérunt Daniélem orántem, et obsecrántem Deum suum. Et accedéntes locúti sunt regi super edícto: Rex, numquid non constituísti ut omnis homo qui rogáret quemquam de diis et homínibus usque ad dies trigínta, nisi te, rex, mitterétur in lacum leónum? Ad quos respóndens rex, ait: Verus est sermo juxta decrétum Medórum atque Persárum, quod prævaricári non licet. Tunc respondéntes dixérunt coram rege: Dániel de fíliis captivitátis Juda, non curávit de lege tua, et de edícto quod constituísti: sed tribus tempóribus per diem orat obsecratióne sua. Quod verbum cum audísset rex, satis contristátus est: et pro Daniéle pósuit cor ut liberáret eum, et usque ad occásum solis laborábat ut erúeret illum. Viri autem illi, intellegéntes regem, dixérunt ei: Scito, rex, quia lex Medórum atque Persárum est ut omne decrétum, quod constitúerit rex, non líceat immutári.\par
\switchcolumn
\EN
\noindent Wherefore those men carefully watching him, found Daniel praying and making supplication to his God. And they came and spoke to the king concerning the edict: O king, hast thou not decreed, that every man that should make a request to any of the gods, or men, for thirty days, but to thyself, O king, should be cast into the den of the lions? And the king answered them, saying: The word is true according to the decree of the Medes and Persians, which it is not lawful to violate. Then they answered, and said before the king: Daniel, who is of the children of the captivity of Juda, hath not regarded thy law, nor the decree that thou hast made: but three times a day he maketh his prayer. Now when the king had heard these words, he was very much grieved, and in behalf of Daniel he set his heart to deliver him and even till sunset he laboured to save him. But those men perceiving the king's design, said to him: Know thou, O king, that the law of the Medes and Persians is, that no decree which the king hath made, may be altered.\par
\switchcolumn*

\LA
\begin{Resp}\Rsign Super muros tuos, Jerúsalem, constítui custódes; \Star{} Tota die et nocte non tacébunt laudáre nomen Dómini.\end{Resp}
\begin{Resp}\Vsign Prædicábunt pópulis fortitúdinem meam, et annuntiábunt géntibus glóriam meam.\end{Resp}
\begin{Resp}\Rsign Tota die et nocte non tacébunt laudáre nomen Dómini.\end{Resp}
\switchcolumn
\EN
\begin{Resp}\Rsign I have set watchmen upon thy walls, O Jerusalem; \Star{} Which shall never hold their peace day nor night, to praise the name of the Lord.\end{Resp}
\begin{Resp}\Vsign They shall proclaim My might unto the nations, and declare My glory unto the Gentiles.\end{Resp}
\begin{Resp}\Rsign Which shall never hold their peace day nor night, to praise the name of the Lord.\end{Resp}
\switchcolumn*

\LA
\LessonHead{Lectio II}
\switchcolumn
\EN
\LessonHead{Lesson II}
\switchcolumn*

\LA
\Cite{Dan 6:16-20}
\switchcolumn
\EN
\Cite{Dan 6:16-20}
\switchcolumn*

\LA
\noindent Tunc rex præcépit, et adduxérunt Daniélem, et misérunt eum in lacum leónum. Dixítque rex Daniéli: Deus tuus, quem colis semper, ipse liberábit te. Allatúsque est lapis unus, et pósitus est super os laci: quem obsignávit rex ánnulo suo, et ánnulo optimátum suórum, ne quid fíeret contra Daniélem. Et ábiit rex in domum suam, et dormívit incœnátus, cibíque non sunt alláti coram eo, ínsuper et somnus recéssit ab eo. Tunc rex primo dilúculo consúrgens, festínus ad lacum leónum perréxit: appropinquánsque lácui, Daniélem voce lacrimábili inclamávit, et affátus est eum: Dániel serve Dei vivéntis, Deus tuus, cui tu servis semper, putásne váluit te liberáre a leónibus?\par
\switchcolumn
\EN
\noindent Then the king commanded, and they brought Daniel, and cast him into the den of the lions. And the king said to Daniel: thy God, whom thou always servest, he will deliver thee. And a stone was brought, and laid upon the mouth of the den: which the king sealed with his own ring, and with the ring of his nobles, that nothing should be done against Daniel. And the king went away to his house and laid himself down without taking supper, and meat was not set before him, and even sleep departed from him. Then the king rising very early in the morning, went in haste to the lions' den: And coming near to the den, cried with a lamentable voice to Daniel, and said to him: Daniel, servant of the living God, hath thy God, whom thou servest always, been able, thinkest thou, to deliver thee from the lions?\par
\switchcolumn*

\LA
\begin{Resp}\Rsign Muro tuo inexpugnábili circumcínge nos, Dómine, et armis tuæ poténtiæ prótege nos semper: \Star{} Líbera, Dómine, Deus Israël, clamántes ad te.\end{Resp}
\begin{Resp}\Vsign Erue nos in mirabílibus tuis, et da glóriam nómini tuo.\end{Resp}
\begin{Resp}\Rsign Líbera, Dómine, Deus Israël, clamántes ad te.\end{Resp}
\switchcolumn
\EN
\begin{Resp}\Rsign Hedge us about with thy wall that cannot be broken down, O Lord, and shield us continually with the arms of thy might. \Star{} O Lord God of Israel, deliver them that cry unto thee.\end{Resp}
\begin{Resp}\Vsign Deliver us also according to thy marvellous works, and give glory to thy Name.\end{Resp}
\begin{Resp}\Rsign O Lord God of Israel, deliver them that cry unto thee.\end{Resp}
\switchcolumn*

\LA
\LessonHead{Lectio III}
\switchcolumn
\EN
\LessonHead{Lesson III}
\switchcolumn*

\LA
\Cite{Dan 6:21-24}
\switchcolumn
\EN
\Cite{Dan 6:21-24}
\switchcolumn*

\LA
\noindent Et Dániel regi respóndens ait: Rex, in ætérnum vive! Deus meus misit ángelum suum, et conclúsit ora leónum, et non nocuérunt mihi: quia coram eo justítia invénta est in me: sed et coram te, rex, delíctum non feci. Tunc veheménter rex gavísus est super eo, et Daniélem præcépit edúci de lacu: eductúsque est Dániel de lacu, et nulla lǽsio invénta est in eo, quia crédidit Deo suo. Jubénte autem rege, addúcti sunt viri illi, qui accusáverant Daniélem: et in lacum leónum missi sunt, ipsi, et fílii, et uxóres eórum: et non pervenérunt usque ad paviméntum laci, donec arríperent eos leónes, et ómnia ossa eórum comminuérunt.\par
\switchcolumn
\EN
\noindent And Daniel answering the king, said: O king, live for ever: My God hath sent his angel, and hath shut up the mouths of the lions, and they have not hurt me: forasmuch as before him justice hath been found in me: yea and before thee, O king, I have done no offence. Then was the king exceeding glad for him, and he commanded that Daniel should be taken out of the den: and Daniel was taken out of the den, and no hurt was found in him, because he believed in his God. And by the king's commandment, those men were brought that had accused Daniel: and they were cast into the lions' den, they and their children, and their wives: and they did not reach the bottom of the den, before the lions caught them, and broke all their bones in pieces.\par
\switchcolumn*

\LA
\begin{Resp}\Rsign Sustinúimus pacem, et non venit: quæsívimus bona, et ecce turbátio: cognóvimus, Dómine, peccáta nostra; \Star{} Non in perpétuum obliviscáris nos.\end{Resp}
\begin{Resp}\Vsign Peccávimus, ímpie géssimus, iniquitátem fécimus, Dómine, in omnem justítiam tuam.\end{Resp}
\begin{Resp}\Rsign Non in perpétuum obliviscáris nos.\end{Resp}
\begin{Resp}\Vsign Glória Patri, et Fílio, \Star{} et Spirítui Sancto.\end{Resp}
\begin{Resp}\Rsign Non in perpétuum obliviscáris nos.\end{Resp}
\switchcolumn
\EN
\begin{Resp}\Rsign We looked for peace, and it came not we asked for good, and behold trouble. We acknowledge, O Lord, our wickedness. \Star{} Forget us not for ever.\end{Resp}
\begin{Resp}\Vsign O Lord (our God) we have sinned, we have done ungodly, we have dealt unrighteously in all thine ordinances.\end{Resp}
\begin{Resp}\Rsign Forget us not for ever.\end{Resp}
\begin{Resp}\Vsign Glory be to the Father, and to the Son, \Star{} and to the Holy Ghost.\end{Resp}
\begin{Resp}\Rsign Forget us not for ever.\end{Resp}
\switchcolumn*

\end{paracol}

\Office{Sabbato infra Hebdomadam III Novembris}{Saturday of the third week of November}{}
\addcontentsline{toc}{subsection}{Sabbato infra Hebdomadam III Novembris \textit{— Saturday of the third week of November}}
\begin{paracol}{2}
\LA
\LessonHead{Lectio I}
\switchcolumn
\EN
\LessonHead{Lesson I}
\switchcolumn*

\LA
\LessonTitle{De Daniéle Prophéta}
\switchcolumn
\EN
\LessonTitle{Lesson from the book of Daniel}
\switchcolumn*

\LA
\Cite{Dan 9:1-5}
\switchcolumn
\EN
\Cite{Dan 9:1-5}
\switchcolumn*

\LA
\noindent In anno primo Daríi fílii Assuéri de sémine Medórum, qui imperávit super regnum Chaldæórum, anno uno regni ejus, ego Dániel intelléxi in libris númerum annórum, de quo factus est sermo Dómini ad Jeremíam prophétam, ut compleréntur desolatiónis Jerúsalem septuagínta anni. Et pósui fáciem meam ad Dóminum Deum meum rogáre et deprecári in jejúniis, sacco, et cínere. Et orávi Dóminum Deum meum, et conféssus sum, et dixi: Obsecro, Dómine Deus magne et terríbilis, custódiens pactum, et misericórdiam diligéntibus te, et custodiéntibus mandáta tua: peccávimus, iniquitátem fécimus, ímpie égimus, et recéssimus: et declinávimus a mandátis tuis ac judíciis.\par
\switchcolumn
\EN
\noindent In the first year of Darius the son of Assuerus of the seed of the Medes, who reigned over the kingdom of the Chaldeans: The first year of his reign, I Daniel understood by books the number of the years, concerning which the word of the Lord came to Jeremias the prophet, that seventy years should be accomplished of the desolation of Jerusalem. And I set my face to the Lord my God, to pray and make supplication with fasting, and sackcloth, and ashes. And I prayed to the Lord my God, and I made my confession, and said: I beseech thee, O Lord God, great and terrible, who keepest the covenant, and mercy to them that love thee, and keep thy commandments. We have sinned, we have committed iniquity, we have done wickedly, and have revolted: and we have gone aside from thy commandments, and thy judgments.\par
\switchcolumn*

\LA
\begin{Resp}\Rsign Laudábilis pópulus, \Star{} Quem Dóminus exercítuum benedíxit dicens: Opus mánuum meárum tu es, heréditas mea Israël.\end{Resp}
\begin{Resp}\Vsign Beáta gens, cujus est Dóminus Deus, pópulus eléctus in hereditátem.\end{Resp}
\begin{Resp}\Rsign Quem Dóminus exercítuum benedíxit dicens: Opus mánuum meárum tu es, heréditas mea Israël.\end{Resp}
\switchcolumn
\EN
\begin{Resp}\Rsign Blessed is the people \Star{} Whom the Lord of hosts hath blessed, saying O Israel thou art the work of Mine own hands, thou art Mine own inheritance.\end{Resp}
\begin{Resp}\Vsign Blessed is the nation whose God is the Lord, and the people whom He hath chosen for His own inheritance.\end{Resp}
\begin{Resp}\Rsign Whom the Lord of hosts hath blessed, saying O Israel thou art the work of Mine own hands, thou art Mine own inheritance.\end{Resp}
\switchcolumn*

\LA
\LessonHead{Lectio II}
\switchcolumn
\EN
\LessonHead{Lesson II}
\switchcolumn*

\LA
\Cite{Dan 9:21-24}
\switchcolumn
\EN
\Cite{Dan 9:21-24}
\switchcolumn*

\LA
\noindent Adhuc me loquénte in oratióne, ecce vir Gábriel, quem víderam in visióne a princípio, cito volans tétigit me in témpore sacrifícii vespertíni. Et dócuit me, et locútus est mihi, dixítque: Dániel, nunc egréssus sum ut docérem te, et intellégeres. Ab exórdio precum tuárum egréssus est sermo: ego autem veni ut indicárem tibi, quia vir desideriórum es: tu ergo animadvérte sermónem, et intéllege visiónem. Septuagínta hebdómades abbreviátæ sunt super pópulum tuum et super urbem sanctam tuam, ut consummétur prævaricátio, et finem accípiat peccátum, et deleátur iníquitas, et adducátur justítia sempitérna, et impleátur vísio et prophetía, et ungátur Sanctus sanctórum.\par
\switchcolumn
\EN
\noindent As I was yet speaking in prayer, behold the man Gabriel, whom I had seen in the vision at the beginning, flying swiftly touched me at the time of the evening sacrifice. And he instructed me, and spoke to me, and said: O Daniel, I am now come forth to teach thee, and that thou mightest understand. From the beginning of thy prayers the word came forth: and I am come to shew it to thee, because thou art a man of desires: therefore do thou mark the word, and understand the vision. Seventy weeks are shortened upon thy people, and upon thy holy city, that transgression may be finished, and sin may have an end, and iniquity may be abolished; and everlasting justice may be brought; and vision and prophecy may be fulfilled; and the saint of saints may be anointed.\par
\switchcolumn*

\LA
\begin{Resp}\Rsign Angústiæ mihi sunt úndique, et quid éligam ignóro; \Star{} Mélius est mihi incídere in manus hóminum, quam derelínquere legem Dei mei.\end{Resp}
\begin{Resp}\Vsign Si enim hoc égero, mors mihi est; si autem non égero, non effúgiam manus vestras.\end{Resp}
\begin{Resp}\Rsign Mélius est mihi incídere in manus hóminum, quam derelínquere legem Dei mei.\end{Resp}
\switchcolumn
\EN
\begin{Resp}\Rsign I am straitened on every side, and know not what to choose. \Star{} It is better for me to fall into the hands of men, than to sin against the law of my God.\end{Resp}
\begin{Resp}\Vsign For if I do this thing, it is death unto me and if I do it not, I cannot escape your hands.\end{Resp}
\begin{Resp}\Rsign It is better for me to fall into the hands of men, than to sin against the law of my God.\end{Resp}
\switchcolumn*

\LA
\LessonHead{Lectio III}
\switchcolumn
\EN
\LessonHead{Lesson III}
\switchcolumn*

\LA
\Cite{Dan 9:25-27}
\switchcolumn
\EN
\Cite{Dan 9:25-27}
\switchcolumn*

\LA
\noindent Scito ergo, et animadvérte: ab éxitu sermónis, ut íterum ædificétur Jerúsalem, usque ad christum ducem, hebdómades septem, et hebdómades sexagínta duæ erunt: et rursum ædificábitur platéa, et muri in angústia témporum. Et post hebdómades sexagínta duas occidétur Christus: et non erit ejus pópulus qui eum negatúrus est. Et civitátem et sanctuárium dissipábit pópulus cum duce ventúro: et finis ejus vástitas, et post finem belli statúta desolátio. Confirmábit autem pactum multis hebdómada una: et in dimídio hebdómadis defíciet hóstia et sacrifícium: et erit in templo abominátio desolatiónis: et usque ad consummatiónem et finem perseverábit desolátio.\par
\switchcolumn
\EN
\noindent Know thou therefore, and take notice: that from the going forth of the word, to build up Jerusalem again, unto Christ the prince, there shall be seven weeks, and sixty-two weeks: and the street shall be built again, and the walls in straitness of times. And after sixty-two weeks Christ shall be slain: and the people that shall deny him shall not be his. And a people with their leader that shall come, shall destroy the city and the sanctuary: and the end thereof shall be waste, and after the end of the war the appointed desolation. And he shall confirm the covenant with many, in one week: and in the half of the week the victim and the sacrifice shall fall: and there shall be in the temple the abomination of desolation: and the desolation shall continue even to the consummation, and to the end.\par
\switchcolumn*

\LA
\begin{Resp}\Rsign Misit Dóminus Angelum suum et conclúsit ora leónum, \Star{} Et non contaminavérunt, quia coram eo injustítia invénta non est in me.\end{Resp}
\begin{Resp}\Vsign Misit Deus misericórdiam suam et veritátem suam: ánimam meam erípuit de médio catulórum leónum.\end{Resp}
\begin{Resp}\Rsign Et non contaminavérunt, quia coram eo injustítia invénta non est in me.\end{Resp}
\begin{Resp}\Vsign Glória Patri, et Fílio, \Star{} et Spirítui Sancto.\end{Resp}
\begin{Resp}\Rsign Et non contaminavérunt: quia coram eo injustítia invénta non est in me.\end{Resp}
\switchcolumn
\EN
\begin{Resp}\Rsign The Lord hath sent His angel, and hath shut the lions' mouths, that \Star{} They have not hurt me forasmuch as before Him innocency was found in me.\end{Resp}
\begin{Resp}\Vsign God hath sent forth His mercy and His truth, (and delivered) my soul from among the lions' whelps.\end{Resp}
\begin{Resp}\Rsign They have not hurt me forasmuch as before Him innocency was found in me.\end{Resp}
\begin{Resp}\Vsign Glory be to the Father, and to the Son, \Star{} and to the Holy Ghost.\end{Resp}
\begin{Resp}\Rsign They have not hurt me forasmuch as before Him innocency was found in me.\end{Resp}
\switchcolumn*

\end{paracol}

\Office{Dominica IV Novembris}{Fourth Sunday of November}{}
\addcontentsline{toc}{subsection}{Dominica IV Novembris \textit{— Fourth Sunday of November}}
\begin{paracol}{2}
\LA
\LessonHead{Lectio I}
\switchcolumn
\EN
\LessonHead{Lesson I}
\switchcolumn*

\LA
\LessonTitle{Incipit liber Osée Prophétæ}
\switchcolumn
\EN
\LessonTitle{Lesson from the book of Osee}
\switchcolumn*

\LA
\Cite{Osee 1:1-3}
\switchcolumn
\EN
\Cite{Osee 1:1-3}
\switchcolumn*

\LA
\noindent Verbum Dómini, quod factum est ad Osée, fílium Béeri, in diébus Ozíæ, Jóathan, Achaz, Ezechíæ, regum Juda; et in diébus Jeróboam, fílii Joas, regis Israël. Princípium loquéndi Dómino in Osée. Et dixit Dóminus ad Osée: Vade, sume tibi uxórem fornicatiónum, quia fórnicans fornicábitur terra a Dómino. Et ábiit, et accépit Gomer, fíliam Debélaim: et concépit, et péperit ei fílium.\par
\switchcolumn
\EN
\noindent The word of the Lord, that came to Osee the son of Beeri, in the days of Ozias, Joathan, Achaz, and Ezechias kings of Juda, and in the days of Jeroboam the son of Joas king of Israel. The beginning of the Lord's speaking by Osee: and the Lord said to Osee: Go, take thee a wife of fornications, and have of her children of fornications: for the land by fornication shall depart from the Lord. So he went, and took Gomer the daughter of Debelaim: and she conceived and bore him a son.\par
\switchcolumn*

\LA
\begin{Resp}\Rsign Vidi Dóminum sedéntem super sólium excélsum et elevátum, et plena erat omnis terra majestáte ejus: \Star{} Et ea, quæ sub ipso erant, replébant templum.\end{Resp}
\begin{Resp}\Vsign Séraphim stabant super illud: sex alæ uni, et sex alæ álteri.\end{Resp}
\begin{Resp}\Rsign Et ea, quæ sub ipso erant, replébant templum.\end{Resp}
\switchcolumn
\EN
\begin{Resp}\Rsign I saw the Lord sitting upon a throne, high and lifted up, and the whole earth was full of His glory; \Star{} And His train filled the temple.\end{Resp}
\begin{Resp}\Vsign Above it stood the Seraphim each one had six wings.\end{Resp}
\begin{Resp}\Rsign And His train filled the temple.\end{Resp}
\switchcolumn*

\LA
\LessonHead{Lectio II}
\switchcolumn
\EN
\LessonHead{Lesson II}
\switchcolumn*

\LA
\Cite{Osee 1:4-7}
\switchcolumn
\EN
\Cite{Osee 1:4-7}
\switchcolumn*

\LA
\noindent Et dixit Dóminus ad eum: Voca nomen ejus Jézrahel, quóniam adhuc módicum, et visitábo sánguinem Jézrahel super domum Jehu, et quiéscere fáciam regnum domus Israël, et in illa die cónteram arcum Israël in valle Jézrahel. Et concépit adhuc, et péperit fíliam. Et dixit ei: Voca nomen ejus Absque misericórdia, quia non addam ultra miseréri dómui Israël, sed oblivióne oblivíscar eórum, et dómui Juda miserébor et salvábo eos in Dómino Deo suo; et non salvábo eos in arcu et gládio, et in bello, et in equis, et in equítibus.\par
\switchcolumn
\EN
\noindent And the Lord said to him: Call his name Jezrahel: for yet a little while, and I will visit the blood of Jezrahel upon the house of Jehu, and I will cause to cease the kingdom of the house of Israel. And in that day I will break in pieces the bow of Israel in the valley of Jezrahel. And she conceived again, and bore a daughter, and he said to him: Call her name, Without mercy: for I will not add any more to have mercy on the house of Israel, but I will utterly forget them. And I will have mercy on the house of Juda, and I will save them by the Lord their God: and I will not save them by bow, nor by sword, nor by battle, nor by horses, nor by horsemen.\par
\switchcolumn*

\LA
\begin{Resp}\Rsign Aspice, Dómine, de sede sancta tua, et cógita de nobis: inclína, Deus meus, aurem tuam et audi: \Star{} Aperi óculos tuos et vide tribulatiónem nostram.\end{Resp}
\begin{Resp}\Vsign Qui regis Israël, inténde, qui dedúcis velut ovem Joseph.\end{Resp}
\begin{Resp}\Rsign Aperi óculos tuos et vide tribulatiónem nostram.\end{Resp}
\switchcolumn
\EN
\begin{Resp}\Rsign Look down, O Lord, from the dwelling-place of thine holiness, and take thought for us. O my God, incline thine ear, and hear. \Star{} Open thine eyes, and behold our desolation.\end{Resp}
\begin{Resp}\Vsign Give ear, O Shepherd of Israel, Thou That leadest Joseph like a flock.\end{Resp}
\begin{Resp}\Rsign Open thine eyes, and behold our desolation.\end{Resp}
\switchcolumn*

\LA
\LessonHead{Lectio III}
\switchcolumn
\EN
\LessonHead{Lesson III}
\switchcolumn*

\LA
\Cite{Osee 1:8-11}
\switchcolumn
\EN
\Cite{Osee 1:8-11}
\switchcolumn*

\LA
\noindent Et ablactávit eam quæ erat Absque misericórdia, et concépit, et péperit fílium. Et dixit: Voca nomen ejus, Non pópulus meus, quia vos non pópulus meus, et ego non ero vester. Et erit númerus filiórum Israël quasi aréna maris, quæ sine mensúra est, et non numerábitur. Et erit, in loco ubi dicétur eis: Non pópulus meus vos; dicétur eis: Fílii Dei vivéntis. Et congregabúntur fílii Juda et fílii Israël páriter: et ponent síbimet caput unum, et ascéndent de terra, quia magnus dies Jézrahel.\par
\switchcolumn
\EN
\noindent And she weaned her that was called Without mercy. And she conceived, and bore a son. And he said: Call his name, Not my people: for you are not my people, and I will not be yours. And the number of the children of Israel shall be as the sand of the sea, that is without measure, and shall not be numbered. And it shall be in the place where it shall be said to them: You are not my people: it shall be said to them: Ye are the sons of the living God. And the children of Juda, and the children of Israel shall be gathered together: and they shall appoint themselves one head, and shall come up out of the land: for great is the day of Jezrahel.\par
\switchcolumn*

\LA
\begin{Resp}\Rsign Aspice, Dómine, quia facta est desoláta cívitas plena divítiis, sedet in tristítia dómina géntium: \Star{} Non est qui consolétur eam, nisi tu, Deus noster.\end{Resp}
\begin{Resp}\Vsign Plorans plorávit in nocte, et lácrimæ ejus in maxíllis ejus.\end{Resp}
\begin{Resp}\Rsign Non est qui consolétur eam, nisi tu, Deus noster.\end{Resp}
\begin{Resp}\Vsign Glória Patri, et Fílio, \Star{} et Spirítui Sancto.\end{Resp}
\begin{Resp}\Rsign Non est qui consolétur eam, nisi tu, Deus noster.\end{Resp}
\switchcolumn
\EN
\begin{Resp}\Rsign Consider, O Lord, how that the city sitteth solitary that was full of riches; how is she become as a widow, she that was great among the nations; \Star{} She hath none to comfort her, save thee, O our God.\end{Resp}
\begin{Resp}\Vsign She weepeth sore in the night, and her tears are on her cheeks.\end{Resp}
\begin{Resp}\Rsign She hath none to comfort her, save thee, O our God.\end{Resp}
\begin{Resp}\Vsign Glory be to the Father, and to the Son, \Star{} and to the Holy Ghost.\end{Resp}
\begin{Resp}\Rsign She hath none to comfort her, save thee, O our God.\end{Resp}
\switchcolumn*

\LA
\LessonHead{Lectio IV}
\switchcolumn
\EN
\LessonHead{Lesson IV}
\switchcolumn*

\LA
\LessonTitle{Ex libro sancti Augustíni Epíscopi de Civitáte Dei}
\switchcolumn
\EN
\LessonTitle{From the Book Upon the City of God, written by St. Augustine, Bishop (of Hippo.)}
\switchcolumn*

\LA
\Source{Liber 18, cap. 28.}
\switchcolumn
\EN
\Source{Bk. xviii. ch. 28}
\switchcolumn*

\LA
\noindent Osée prophéta, quanto profúndius quidem lóquitur, tanto operósius penetrátur. Sed áliquid inde suméndum est, et hic ex nostra promissióne ponéndum. Et erit, inquit, in loco quo dictum est eis, Non pópulus meus vos; vocabúntur et ipsi fílii Dei vivi. Hoc testimónium prophéticum de vocatióne pópuli Géntium, qui prius non pertinébat ad Deum, étiam Apóstoli intellexérunt.\par
\switchcolumn
\EN
\noindent AS to the Prophet Hosea, the deeper his meaning, the harder to pierce. But somewhat may be gotten out of him, and, as I promised, I will give it here. He saith And it shall come to pass that, in the place where it shall be said unto them, Ye are not My people, there it shall be said unto them Ye are the sons of the living God. This was understood even by the Apostles as a Prophetic witness to the call of the Gentiles, who erst had not been God's people. (Rom. ix. 24-26.)\par
\switchcolumn*

\LA
\begin{Resp}\Rsign Super muros tuos, Jerúsalem, constítui custódes; \Star{} Tota die et nocte non tacébunt laudáre nomen Dómini.\end{Resp}
\begin{Resp}\Vsign Prædicábunt pópulis fortitúdinem meam, et annuntiábunt géntibus glóriam meam.\end{Resp}
\begin{Resp}\Rsign Tota die et nocte non tacébunt laudáre nomen Dómini.\end{Resp}
\switchcolumn
\EN
\begin{Resp}\Rsign I have set watchmen upon thy walls, O Jerusalem; \Star{} Which shall never hold their peace day nor night, to praise the name of the Lord.\end{Resp}
\begin{Resp}\Vsign They shall proclaim My might unto the nations, and declare My glory unto the Gentiles.\end{Resp}
\begin{Resp}\Rsign Which shall never hold their peace day nor night, to praise the name of the Lord.\end{Resp}
\switchcolumn*

\LA
\LessonHead{Lectio V}
\switchcolumn
\EN
\LessonHead{Lesson V}
\switchcolumn*

\LA
\noindent Et quia ipse quoque pópulus Géntium spiritáliter est in fíliis Abrahæ, ac per hoc recte dícitur Israël; proptérea séquitur, et dicit: Et congregabúntur fílii Juda et fílii Israël in idípsum, et ponent síbimet principátum unum, et ascéndent a terra. Hoc si adhuc velímus expónere, elóquii prophétici obtundétur sapor. Recolátur tamen lapis ille anguláris, et duo illi paríetes, unus ex Judǽis, alter ex Géntibus: ille nómine filiórum Juda, iste nómine filiórum Israël, eídem uni principátui suo in idípsum inniténtes, et ascendéntes agnoscántur in terra.\par
\switchcolumn
\EN
\noindent And since the converted Gentiles are the spiritual children of Abraham, and are therefore rightly called Israelites, therefore he goeth on, and saith: Then shall the children of Judah and the children of Israel be gathered together, and appoint themselves one head, and they shall come up out of the land. If we went on expounding this, we should water down the flavour of the prophetic draught. Let there be remembered, however, that Corner Stone, and let there be acknowledged those twain walls, (which It bindeth in one,) the Jews and the Gentiles, one called the children of Judah and the other the children of Israel, bound together under One Head, and coming up out of the land.\par
\switchcolumn*

\LA
\begin{Resp}\Rsign Muro tuo inexpugnábili circumcínge nos, Dómine, et armis tuæ poténtiæ prótege nos semper: \Star{} Líbera, Dómine, Deus Israël, clamántes ad te.\end{Resp}
\begin{Resp}\Vsign Erue nos in mirabílibus tuis, et da glóriam nómini tuo.\end{Resp}
\begin{Resp}\Rsign Líbera, Dómine, Deus Israël, clamántes ad te.\end{Resp}
\switchcolumn
\EN
\begin{Resp}\Rsign Hedge us about with thy wall that cannot be broken down, O Lord, and shield us continually with the arms of thy might. \Star{} O Lord God of Israel, deliver them that cry unto thee.\end{Resp}
\begin{Resp}\Vsign Deliver us also according to thy marvellous works, and give glory to thy Name.\end{Resp}
\begin{Resp}\Rsign O Lord God of Israel, deliver them that cry unto thee.\end{Resp}
\switchcolumn*

\LA
\LessonHead{Lectio VI}
\switchcolumn
\EN
\LessonHead{Lesson VI}
\switchcolumn*

\LA
\noindent Istos autem carnáles Israëlítas, qui nunc nolunt crédere in Christum, póstea creditúros, id est fílios eórum, (nam útique isti in suum locum moriéndo transíbunt) idem prophéta testátur, dicens: Quóniam diébus multis sedébunt fílii Israël sine rege, sine príncipe, sine sacrifício, sine altári, sine sacerdótio, sine manifestatiónibus. Quis non vídeat, nunc sic esse Judǽos?\par
\switchcolumn
\EN
\noindent Concerning them that are now Israelites according to the flesh, that will not now believe in Christ, but shall believe hereafter, (that is, their children shall believe, for these shall die, and go to their own place,) this same Prophet giveth witness, where he saith The children of Israel shall abide many days without a King, and without a Prince, and without a sacrifice, and without an Altar, and without a Priest, and without oracles. (iii. 4.) To whom is it not manifest that such is the state of the Jews now\par
\switchcolumn*

\LA
\begin{Resp}\Rsign Sustinúimus pacem, et non venit: quæsívimus bona, et ecce turbátio: cognóvimus, Dómine, peccáta nostra; \Star{} Non in perpétuum obliviscáris nos.\end{Resp}
\begin{Resp}\Vsign Peccávimus, ímpie géssimus, iniquitátem fécimus, Dómine, in omnem justítiam tuam.\end{Resp}
\begin{Resp}\Rsign Non in perpétuum obliviscáris nos.\end{Resp}
\begin{Resp}\Vsign Glória Patri, et Fílio, \Star{} et Spirítui Sancto.\end{Resp}
\begin{Resp}\Rsign Non in perpétuum obliviscáris nos.\end{Resp}
\switchcolumn
\EN
\begin{Resp}\Rsign We looked for peace, and it came not we asked for good, and behold trouble. We acknowledge, O Lord, our wickedness. \Star{} Forget us not for ever.\end{Resp}
\begin{Resp}\Vsign O Lord (our God) we have sinned, we have done ungodly, we have dealt unrighteously in all thine ordinances.\end{Resp}
\begin{Resp}\Rsign Forget us not for ever.\end{Resp}
\begin{Resp}\Vsign Glory be to the Father, and to the Son, \Star{} and to the Holy Ghost.\end{Resp}
\begin{Resp}\Rsign Forget us not for ever.\end{Resp}
\switchcolumn*

\end{paracol}

\Office{Feria II infra Hebdomadam IV Novembris}{Monday of the fourth week of November}{}
\addcontentsline{toc}{subsection}{Feria II infra Hebdomadam IV Novembris \textit{— Monday of the fourth week of November}}
\begin{paracol}{2}
\LA
\LessonHead{Lectio I}
\switchcolumn
\EN
\LessonHead{Lesson I}
\switchcolumn*

\LA
\LessonTitle{De Osée Prophéta}
\switchcolumn
\EN
\LessonTitle{Lesson from the book of Osee}
\switchcolumn*

\LA
\Cite{Osee 4:1-3}
\switchcolumn
\EN
\Cite{Osee 4:1-3}
\switchcolumn*

\LA
\noindent Audíte verbum Dómini, fílii Israël, quia judícium Dómino cum habitatóribus terræ: non est enim véritas, et non est misericórdia, et non est sciéntia Dei in terra. Maledíctum, et mendácium, et homicídium, et furtum, et adultérium inundavérunt, et sanguis sánguinem tétigit. Propter hoc lugébit terra, et infirmábitur omnis qui hábitat in ea, in béstia agri, et in vólucre cæli; sed et pisces maris congregabúntur.\par
\switchcolumn
\EN
\noindent Hear the word of the Lord, ye children of Israel, for the Lord shall enter into judgment with the inhabitants of the land: for there is no truth, and there is no mercy, and there is no knowledge of God in the land. Cursing, and lying, and killing, and theft, and adultery have overflowed, and blood hath touched blood. Therefore shall the land mourn, and every one that dwelleth in it shall languish with the beasts of the field, and with the fowls of the air: yea, the fishes of the sea also shall be gathered together.\par
\switchcolumn*

\LA
\begin{Resp}\Rsign Redémit pópulum suum et liberávit eum, et vénient et exsultábunt in monte Sion et gaudébunt de bonis Dómini super fruménto, vino et óleo, \Star{} Et ultra non esúrient.\end{Resp}
\begin{Resp}\Vsign Erítque ánima eórum quasi hortus irríguus.\end{Resp}
\begin{Resp}\Rsign Et ultra non esúrient.\end{Resp}
\switchcolumn
\EN
\begin{Resp}\Rsign He hath redeemed His people, and ransomed them therefore they shall come and sing in the height of Zion, and shall rejoice in the goodness of the Lord, for wheat, and for wine, and for oil \Star{} And they shall hunger no more.\end{Resp}
\begin{Resp}\Vsign And their soul shall be as a watered garden.\end{Resp}
\begin{Resp}\Rsign And they shall hunger no more.\end{Resp}
\switchcolumn*

\LA
\LessonHead{Lectio II}
\switchcolumn
\EN
\LessonHead{Lesson II}
\switchcolumn*

\LA
\Cite{Osee 4:4-6}
\switchcolumn
\EN
\Cite{Osee 4:4-6}
\switchcolumn*

\LA
\noindent Verúmtamen unusquísque non júdicet, et non arguátur vir: pópulus enim tuus sicut hi qui contradícunt sacerdóti. Et córrues hódie, et córruet étiam prophéta tecum. Nocte tacére feci matrem tuam. Contícuit pópulus meus, eo quod non habúerit sciéntiam: quia tu sciéntiam repulísti, repéllam te, ne sacerdótio fungáris mihi; et oblíta es legis Dei tui, oblivíscar filiórum tuórum et ego.\par
\switchcolumn
\EN
\noindent But yet let not any man judge: and let not a man be rebuked: for thy people are as they that contradict the priest. And thou shalt fall to day, and the prophet also shall fall with thee: in the night I have made thy mother to be silent. My people have been silent, because they had no knowledge: because thou hast rejected knowledge, I will reject thee, that thou shalt not do the office of priesthood to me: and thou hast forgotten the law of thy God, I also will forget thy children.\par
\switchcolumn*

\LA
\begin{Resp}\Rsign Angústiæ mihi sunt úndique, et quid éligam ignóro; \Star{} Mélius est mihi incídere in manus hóminum, quam derelínquere legem Dei mei.\end{Resp}
\begin{Resp}\Vsign Si enim hoc égero, mors mihi est; si autem non égero, non effúgiam manus vestras.\end{Resp}
\begin{Resp}\Rsign Mélius est mihi incídere in manus hóminum, quam derelínquere legem Dei mei.\end{Resp}
\switchcolumn
\EN
\begin{Resp}\Rsign I am straitened on every side, and know not what to choose. \Star{} It is better for me to fall into the hands of men, than to sin against the law of my God.\end{Resp}
\begin{Resp}\Vsign For if I do this thing, it is death unto me and if I do it not, I cannot escape your hands.\end{Resp}
\begin{Resp}\Rsign It is better for me to fall into the hands of men, than to sin against the law of my God.\end{Resp}
\switchcolumn*

\LA
\LessonHead{Lectio III}
\switchcolumn
\EN
\LessonHead{Lesson III}
\switchcolumn*

\LA
\Cite{Osee 4:7-10}
\switchcolumn
\EN
\Cite{Osee 4:7-10}
\switchcolumn*

\LA
\noindent Secúndum multitúdinem eórum sic peccavérunt mihi: glóriam eórum in ignomíniam commutábo. Peccáta pópuli mei cómedent, et ad iniquitátem eórum sublevábunt ánimas eórum. Et erit sicut pópulus, sic sacérdos; et visitábo super eum vias ejus, et cogitatiónes ejus reddam ei. Et cómedent, et non saturabúntur; fornicáti sunt, et non cessavérunt: quóniam Dóminum dereliquérunt in non custodiéndo.\par
\switchcolumn
\EN
\noindent According to the multitude of them so have they sinned against me: I will change their glory into shame. They shall eat the sins of my people, and shall lift up their souls to their iniquity. And there shall be like people like priest: and I will visit their ways upon them, and I will repay them their devices. And they shall eat and shall not be filled: they have committed fornication, and have not ceased: because they have forsaken the Lord in not observing his law.\par
\switchcolumn*

\LA
\begin{Resp}\Rsign Misit Dóminus Angelum suum et conclúsit ora leónum, \Star{} Et non contaminavérunt, quia coram eo injustítia invénta non est in me.\end{Resp}
\begin{Resp}\Vsign Misit Deus misericórdiam suam et veritátem suam: ánimam meam erípuit de médio catulórum leónum.\end{Resp}
\begin{Resp}\Rsign Et non contaminavérunt, quia coram eo injustítia invénta non est in me.\end{Resp}
\begin{Resp}\Vsign Glória Patri, et Fílio, \Star{} et Spirítui Sancto.\end{Resp}
\begin{Resp}\Rsign Et non contaminavérunt: quia coram eo injustítia invénta non est in me.\end{Resp}
\switchcolumn
\EN
\begin{Resp}\Rsign The Lord hath sent His angel, and hath shut the lions' mouths, that \Star{} They have not hurt me forasmuch as before Him innocency was found in me.\end{Resp}
\begin{Resp}\Vsign God hath sent forth His mercy and His truth, (and delivered) my soul from among the lions' whelps.\end{Resp}
\begin{Resp}\Rsign They have not hurt me forasmuch as before Him innocency was found in me.\end{Resp}
\begin{Resp}\Vsign Glory be to the Father, and to the Son, \Star{} and to the Holy Ghost.\end{Resp}
\begin{Resp}\Rsign They have not hurt me forasmuch as before Him innocency was found in me.\end{Resp}
\switchcolumn*

\end{paracol}

\Office{Feria III infra Hebdomadam IV Novembris}{Tuesday of the fourth week of November}{}
\addcontentsline{toc}{subsection}{Feria III infra Hebdomadam IV Novembris \textit{— Tuesday of the fourth week of November}}
\begin{paracol}{2}
\LA
\LessonHead{Lectio I}
\switchcolumn
\EN
\LessonHead{Lesson I}
\switchcolumn*

\LA
\LessonTitle{Incipit Joël Prophéta}
\switchcolumn
\EN
\LessonTitle{Lesson from the book of Joel}
\switchcolumn*

\LA
\Cite{Joël 1:1-4}
\switchcolumn
\EN
\Cite{Joel 1:1-4}
\switchcolumn*

\LA
\noindent Verbum Dómini, quod factum est ad Joël, fílium Phátuel. Audíte hoc, senes, et áuribus percípite, omnes habitatóres terræ: si factum est istud in diébus vestris, aut in diébus patrum vestrórum? Super hoc fíliis vestris narráte, et fílii vestri fíliis suis, et fílii eórum generatióni álteræ. Resíduum erúcæ comédit locústa, et resíduum locústæ comédit bruchus, et resíduum bruchi comédit rubígo.\par
\switchcolumn
\EN
\noindent The word of the Lord that came to Joel the son of Phatuel. Hear this, ye old men, and give ear, all ye inhabitants of the land: did this ever happen in your days, or in the days of your fathers? Tell ye of this to your children, and let your children tell their children, and their children to another generation. That which the palmerworm hath left, the locust hath eaten: and that which the locust hath left, the bruchus hath eaten: and that which the bruchus hath left, the mildew hath destroyed.\par
\switchcolumn*

\LA
\begin{Resp}\Rsign A fácie furóris tui, Deus, conturbáta est omnis terra: \Star{} Sed miserére, Dómine, et ne fácias consummatiónem.\end{Resp}
\begin{Resp}\Vsign Dómine, Dóminus noster, quam admirábile est nomen tuum!\end{Resp}
\begin{Resp}\Rsign Sed miserére, Dómine, et ne fácias consummatiónem.\end{Resp}
\switchcolumn
\EN
\begin{Resp}\Rsign Before the face of thine anger, O God, the whole earth is troubled; \Star{} But Thou, O Lord, have mercy, and make not an end utterly.\end{Resp}
\begin{Resp}\Vsign Lord our Ruler, how excellent is thy Name in all the earth!\end{Resp}
\begin{Resp}\Rsign But Thou, O Lord, have mercy, and make not an end utterly.\end{Resp}
\switchcolumn*

\LA
\LessonHead{Lectio II}
\switchcolumn
\EN
\LessonHead{Lesson II}
\switchcolumn*

\LA
\Cite{Joël 1:5-7}
\switchcolumn
\EN
\Cite{Joel 1:5-7}
\switchcolumn*

\LA
\noindent Expergiscímini, ébrii, et flete et ululáte, omnes qui bíbitis vinum in dulcédine, quóniam périit ab ore vestro. Gens enim ascéndit super terram meam, fortis et innumerábilis: dentes ejus ut dentes leónis, et moláres ejus ut cátuli leónis. Pósuit víneam meam in desértum, et ficum meam decorticávit; nudans spoliávit eam, et projécit; albi facti sunt rami ejus.\par
\switchcolumn
\EN
\noindent Awake, ye that are drunk, and weep, and mourn all ye that take delight in drinking sweet wine: for it is cut off from your mouth. For a nation is come up upon my land, strong and without number: his teeth are like the teeth of a lion: and his cheek teeth as of a lion's whelp. He hath laid my vineyard waste, and hath pilled off the bark of my fig tree: he hath stripped it bare, and cast it away; the branches thereof are made white.\par
\switchcolumn*

\LA
\begin{Resp}\Rsign Civitátem istam tu circúmda, Dómine: et Angeli tui custódiant muros ejus. \Star{} Exáudi, Dómine, pópulum tuum cum misericórdia.\end{Resp}
\begin{Resp}\Vsign Avertátur furor tuus, Dómine, a pópulo tuo et a civitáte sancta tua.\end{Resp}
\begin{Resp}\Rsign Exáudi, Dómine, pópulum tuum cum misericórdia.\end{Resp}
\switchcolumn
\EN
\begin{Resp}\Rsign Fence Thou this city; O Lord, and let thine angels keep the walls thereof. \Star{} O Lord, hearken unto thy people with mercy.\end{Resp}
\begin{Resp}\Vsign O Lord, let thine anger be turned away from thy people, and from thine holy city.\end{Resp}
\begin{Resp}\Rsign O Lord, hearken unto thy people with mercy.\end{Resp}
\switchcolumn*

\LA
\LessonHead{Lectio III}
\switchcolumn
\EN
\LessonHead{Lesson III}
\switchcolumn*

\LA
\Cite{Joël 1:8-11}
\switchcolumn
\EN
\Cite{Joel 1:8-11}
\switchcolumn*

\LA
\noindent Plange, quasi virgo accíncta sacco super virum pubertátis suæ. Périit sacrifícium et libátio de domo Dómini; luxérunt sacerdótes, minístri Dómini. Depopuláta est régio, luxit humus, quóniam devastátum est tríticum, confúsum est vinum, elánguit óleum, confúsi sunt agrícolæ, ululavérunt vinitóres super fruménto et hórdeo, quia périit messis agri.\par
\switchcolumn
\EN
\noindent Lament like a virgin girded with sackcloth for the husband of her youth. Sacrifice and libation is cut off from the house of the Lord: the priests, the Lord's ministers, have mourned: The country is destroyed, the ground hath mourned: for the corn is wasted, the wine is confounded, the oil hath languished. The husbandmen are ashamed, the vinedressers have howled for the wheat, and for the barley, because the harvest of the field is perished.\par
\switchcolumn*

\LA
\begin{Resp}\Rsign Genti peccatríci, pópulo pleno peccáto miserére, \Star{} Dómine Deus.\end{Resp}
\begin{Resp}\Vsign Esto placábilis super nequítiam pópuli tui.\end{Resp}
\begin{Resp}\Rsign Dómine Deus.\end{Resp}
\begin{Resp}\Vsign Glória Patri, et Fílio, \Star{} et Spirítui Sancto.\end{Resp}
\begin{Resp}\Rsign Dómine Deus.\end{Resp}
\switchcolumn
\EN
\begin{Resp}\Rsign O Lord God! have mercy upon \Star{} The sinful nation, upon the people laden with iniquity.\end{Resp}
\begin{Resp}\Vsign Let it repent thee concerning the transgression of thy people.\end{Resp}
\begin{Resp}\Rsign The people laden with iniquity.\end{Resp}
\begin{Resp}\Vsign Glory be to the Father, and to the Son, \Star{} and to the Holy Ghost.\end{Resp}
\begin{Resp}\Rsign The people laden with iniquity.\end{Resp}
\switchcolumn*

\end{paracol}

\Office{Feria IV infra Hebdomadam IV Novembris}{Wednesday of the fourth week of November}{}
\addcontentsline{toc}{subsection}{Feria IV infra Hebdomadam IV Novembris \textit{— Wednesday of the fourth week of November}}
\begin{paracol}{2}
\LA
\LessonHead{Lectio I}
\switchcolumn
\EN
\LessonHead{Lesson I}
\switchcolumn*

\LA
\LessonTitle{De Joéle Prophéta}
\switchcolumn
\EN
\LessonTitle{Lesson from the book of Joel}
\switchcolumn*

\LA
\Cite{Joël 3:1-3}
\switchcolumn
\EN
\Cite{Joel 3:1-3}
\switchcolumn*

\LA
\noindent Quia ecce in diébus illis, et in témpore illo, cum convértero captivitátem Juda et Jerúsalem, congregábo omnes gentes, et dedúcam eas in vallem Jósaphat; et disceptábo cum eis ibi super pópulo meo, et hæreditáte mea Israël, quos dispersérunt in natiónibus, et terram meam divisérunt. Et super pópulum meum misérunt sortem; et posuérunt púerum in prostíbulo, et puéllam vendidérunt pro vino ut bíberent.\par
\switchcolumn
\EN
\noindent For behold in those days, and in that time when I shall bring back the captivity of Juda and Jerusalem: I will gather together all nations, and will bring them down into the valley of Josaphat: and I will plead with them there for my people, and for my inheritance Israel, whom they have scattered among the nations, and have parted my land. And they have cast lots upon my people: and the boy they have put in the stews, and the girl they have sold for wine, that they might drink.\par
\switchcolumn*

\LA
\begin{Resp}\Rsign Indicábo tibi, homo, quid sit bonum aut quid Dóminus requírat a te: \Star{} Fácere judícium et justítiam et sollícitum ambuláre cum Deo tuo.\end{Resp}
\begin{Resp}\Vsign Spera in Dómino, et fac bonitátem, et inhábita terram.\end{Resp}
\begin{Resp}\Rsign Fácere judícium et justítiam et sollícitum ambuláre cum Deo tuo.\end{Resp}
\switchcolumn
\EN
\begin{Resp}\Rsign I will show thee, O man, what is good, and what doth the Lord require of thee, but to \Star{} Do justice and judgment, and to walk humbly with thy God.\end{Resp}
\begin{Resp}\Vsign Trust in the Lord, and do good, and dwell in the land.\end{Resp}
\begin{Resp}\Rsign Do justice and judgment, and walk humbly with thy God.\end{Resp}
\switchcolumn*

\LA
\LessonHead{Lectio II}
\switchcolumn
\EN
\LessonHead{Lesson II}
\switchcolumn*

\LA
\Cite{Joël 3:4-7}
\switchcolumn
\EN
\Cite{Joel 3:4-7}
\switchcolumn*

\LA
\noindent Verum quid mihi et vobis, Tyrus et Sidon, et omnis términus Palæstinórum? numquid ultiónem vos reddétis mihi? et si ulciscímini vos contra me, cito velóciter reddam vicissitúdinem vobis super caput vestrum. Argéntum enim meum et aurum tulístis, et desiderabília mea et pulchérrima intulístis in delúbra vestra. Et fílios Juda et fílios Jerúsalem vendidístis fíliis Græcórum, ut longe facerétis eos de fínibus suis. Ecce ego suscitábo eos de loco in quo vendidístis eos, et convértam retributiónem vestram in caput vestrum.\par
\switchcolumn
\EN
\noindent But what have you to do with me, O Tyre, and Sidon, and all the coast of the Philistines? will you revenge yourselves on me? and if you revenge yourselves on me, I will very soon return you a recompense upon your own head. For you have taken away my silver and my gold: and my desirable and most beautiful things you have carried into your temples. And the children of Juda, and the children of Jerusalem you have sold to the children of the Greeks, that you might remove them far off from their own country. Behold, I will raise them up out of the place wherein you have sold them: and I will return your recompense upon your own heads.\par
\switchcolumn*

\LA
\begin{Resp}\Rsign Angústiæ mihi sunt úndique, et quid éligam ignóro; \Star{} Mélius est mihi incídere in manus hóminum, quam derelínquere legem Dei mei.\end{Resp}
\begin{Resp}\Vsign Si enim hoc égero, mors mihi est; si autem non égero, non effúgiam manus vestras.\end{Resp}
\begin{Resp}\Rsign Mélius est mihi incídere in manus hóminum, quam derelínquere legem Dei mei.\end{Resp}
\switchcolumn
\EN
\begin{Resp}\Rsign I am straitened on every side, and know not what to choose. \Star{} It is better for me to fall into the hands of men, than to sin against the law of my God.\end{Resp}
\begin{Resp}\Vsign For if I do this thing, it is death unto me and if I do it not, I cannot escape your hands.\end{Resp}
\begin{Resp}\Rsign It is better for me to fall into the hands of men, than to sin against the law of my God.\end{Resp}
\switchcolumn*

\LA
\LessonHead{Lectio III}
\switchcolumn
\EN
\LessonHead{Lesson III}
\switchcolumn*

\LA
\Cite{Joël 3:8-12}
\switchcolumn
\EN
\Cite{Joel 3:8-12}
\switchcolumn*

\LA
\noindent Et vendam fílios vestros et fílias vestras in mánibus filiórum Juda, et venúndabunt eos Sabǽis, genti longínquæ, quia Dóminus locútus est. Clamáte hoc in géntibus, sanctificáte bellum, suscitáte robústos: accédant, ascéndant omnes viri bellatóres. Concídite arátra vestra in gládios, et ligónes vestros in lánceas. Infírmus dicat: Quia fortis ego sum. Erúmpite, et veníte, omnes gentes de circúitu, et congregámini; ibi occúmbere fáciet Dóminus robústos tuos. Consúrgant, et ascéndant gentes in vallem Jósaphat, quia ibi sedébo ut júdicem omnes gentes in circúitu.\par
\switchcolumn
\EN
\noindent And I will sell your sons, and your daughters by the hands of the children of Juda, and they shall sell them to the Sabeans, a nation far off, for the Lord hath spoken it. Proclaim ye this among the nations: prepare war, rouse up the strong: let them come, let all the men of war come up. Cut your ploughshares into swords, and your spades into spears. Let the weak say: I am strong. Break forth, and come, all ye nations, from round about, and gather yourselves together: there will the Lord cause all thy strong ones to fall down. Let them arise, and let the nations come up into the valley of Josaphat: for there I will sit to judge all nations round about.\par
\switchcolumn*

\LA
\begin{Resp}\Rsign Misit Dóminus Angelum suum et conclúsit ora leónum, \Star{} Et non contaminavérunt, quia coram eo injustítia invénta non est in me.\end{Resp}
\begin{Resp}\Vsign Misit Deus misericórdiam suam et veritátem suam: ánimam meam erípuit de médio catulórum leónum.\end{Resp}
\begin{Resp}\Rsign Et non contaminavérunt, quia coram eo injustítia invénta non est in me.\end{Resp}
\begin{Resp}\Vsign Glória Patri, et Fílio, \Star{} et Spirítui Sancto.\end{Resp}
\begin{Resp}\Rsign Et non contaminavérunt: quia coram eo injustítia invénta non est in me.\end{Resp}
\switchcolumn
\EN
\begin{Resp}\Rsign The Lord hath sent His angel, and hath shut the lions' mouths, that \Star{} They have not hurt me forasmuch as before Him innocency was found in me.\end{Resp}
\begin{Resp}\Vsign God hath sent forth His mercy and His truth, (and delivered) my soul from among the lions' whelps.\end{Resp}
\begin{Resp}\Rsign They have not hurt me forasmuch as before Him innocency was found in me.\end{Resp}
\begin{Resp}\Vsign Glory be to the Father, and to the Son, \Star{} and to the Holy Ghost.\end{Resp}
\begin{Resp}\Rsign They have not hurt me forasmuch as before Him innocency was found in me.\end{Resp}
\switchcolumn*

\end{paracol}

\Office{Feria V infra Hebdomadam IV Novembris}{Thursday of the fourth week of November}{}
\addcontentsline{toc}{subsection}{Feria V infra Hebdomadam IV Novembris \textit{— Thursday of the fourth week of November}}
\begin{paracol}{2}
\LA
\LessonHead{Lectio I}
\switchcolumn
\EN
\LessonHead{Lesson I}
\switchcolumn*

\LA
\LessonTitle{Incipit Amos Prophéta}
\switchcolumn
\EN
\LessonTitle{Lesson from the book of Amos}
\switchcolumn*

\LA
\Cite{Amos 1:1-2}
\switchcolumn
\EN
\Cite{Amos 1:1-2}
\switchcolumn*

\LA
\noindent Verba Amos, qui fuit in pastóribus de Thécue, quæ vidit super Israël in diébus Ozíæ, regis Juda, et in diébus Jeróboam, fílii Joas, regis Israël, ante duos annos terræmótus. Et dixit: Dóminus de Sion rúgiet, et de Jerúsalem dabit vocem suam; et luxérunt speciósa pastórum, et exsiccátus est vertex Carméli.\par
\switchcolumn
\EN
\noindent The words of Amos, who was among herdsmen of Thecua: which he saw concerning Israel in the days of Ozias king of Juda, and in the days of Jeroboam the son of Joas king of Israel two years before the earthquake. And he said: The Lord will roar from Sion, and utter his voice from Jerusalem: and the beautiful places of the shepherds have mourned, and the top of Carmel is withered.\par
\switchcolumn*

\LA
\begin{Resp}\Rsign Vidi Dóminum sedéntem super sólium excélsum et elevátum, et plena erat omnis terra majestáte ejus: \Star{} Et ea, quæ sub ipso erant, replébant templum.\end{Resp}
\begin{Resp}\Vsign Séraphim stabant super illud: sex alæ uni, et sex alæ álteri.\end{Resp}
\begin{Resp}\Rsign Et ea, quæ sub ipso erant, replébant templum.\end{Resp}
\switchcolumn
\EN
\begin{Resp}\Rsign I saw the Lord sitting upon a throne, high and lifted up, and the whole earth was full of His glory; \Star{} And His train filled the temple.\end{Resp}
\begin{Resp}\Vsign Above it stood the Seraphim each one had six wings.\end{Resp}
\begin{Resp}\Rsign And His train filled the temple.\end{Resp}
\switchcolumn*

\LA
\LessonHead{Lectio II}
\switchcolumn
\EN
\LessonHead{Lesson II}
\switchcolumn*

\LA
\Cite{Amos 1:3-5}
\switchcolumn
\EN
\Cite{Amos 1:3-5}
\switchcolumn*

\LA
\noindent Hæc dicit Dóminus: Super tribus sceléribus Damásci, et super quátuor non convértam eum, eo quod trituráverint in plaustris férreis Gálaad. Et mittam ignem in domum Azaël, et devorábit domos Bénadad. Et cónteram vectem Damásci: et dispérdam habitatórem de campo idóli, et tenéntem sceptrum de domo voluptátis: et transferétur pópulus Sýriæ Cyrénen, dicit Dóminus.\par
\switchcolumn
\EN
\noindent Thus saith the Lord: For three crimes of Damascus, and for four I will not convert it: because they have thrashed Galaad with iron wains. And I will send a fire into the house of Azael, and it shall devour the houses of Benadad. And I will break the bar of Damascus: and I will cut off the inhabitants from the plain of the idol, and him that holdeth the sceptre from the house of pleasure: and the people of Syria shall be carried away to Cyrene, saith the Lord.\par
\switchcolumn*

\LA
\begin{Resp}\Rsign Aspice, Dómine, de sede sancta tua, et cógita de nobis: inclína, Deus meus, aurem tuam et audi: \Star{} Aperi óculos tuos et vide tribulatiónem nostram.\end{Resp}
\begin{Resp}\Vsign Qui regis Israël, inténde, qui dedúcis velut ovem Joseph.\end{Resp}
\begin{Resp}\Rsign Aperi óculos tuos et vide tribulatiónem nostram.\end{Resp}
\switchcolumn
\EN
\begin{Resp}\Rsign Look down, O Lord, from the dwelling-place of thine holiness, and take thought for us. O my God, incline thine ear, and hear. \Star{} Open thine eyes, and behold our desolation.\end{Resp}
\begin{Resp}\Vsign Give ear, O Shepherd of Israel, Thou That leadest Joseph like a flock.\end{Resp}
\begin{Resp}\Rsign Open thine eyes, and behold our desolation.\end{Resp}
\switchcolumn*

\LA
\LessonHead{Lectio III}
\switchcolumn
\EN
\LessonHead{Lesson III}
\switchcolumn*

\LA
\Cite{Amos 1:6-8}
\switchcolumn
\EN
\Cite{Amos 1:6-8}
\switchcolumn*

\LA
\noindent Hæc dicit Dóminus: Super tribus sceléribus Gazæ, et super quátuor non convértam eum, eo quod transtúlerint captivitátem perféctam, ut conclúderent eam in Idumǽa. Et mittam ignem in murum Gazæ, et devorábit ædes ejus. Et dispérdam habitatórem de Azóto, et tenéntem sceptrum de Ascalóne: et convértam manum meam super Accaron, et períbunt réliqui Philisthinórum, dicit Dóminus Deus.\par
\switchcolumn
\EN
\noindent Thus saith the Lord: For three crimes of Gaza, and for four I will not convert it: because they have carried away a perfect captivity to shut them up in Edom. And I will send a fire on the wall of Gaza, and it shall devour the houses thereof. And I will cut off the inhabitant from Azotus, and him that holdeth the sceptre from Ascalon: and I will turn my hand against Accaron, and the rest of the Philistines shall perish, saith the Lord God.\par
\switchcolumn*

\LA
\begin{Resp}\Rsign Aspice, Dómine, quia facta est desoláta cívitas plena divítiis, sedet in tristítia dómina géntium: \Star{} Non est qui consolétur eam, nisi tu, Deus noster.\end{Resp}
\begin{Resp}\Vsign Plorans plorávit in nocte, et lácrimæ ejus in maxíllis ejus.\end{Resp}
\begin{Resp}\Rsign Non est qui consolétur eam, nisi tu, Deus noster.\end{Resp}
\begin{Resp}\Vsign Glória Patri, et Fílio, \Star{} et Spirítui Sancto.\end{Resp}
\begin{Resp}\Rsign Non est qui consolétur eam, nisi tu, Deus noster.\end{Resp}
\switchcolumn
\EN
\begin{Resp}\Rsign Consider, O Lord, how that the city sitteth solitary that was full of riches; how is she become as a widow, she that was great among the nations; \Star{} She hath none to comfort her, save thee, O our God.\end{Resp}
\begin{Resp}\Vsign She weepeth sore in the night, and her tears are on her cheeks.\end{Resp}
\begin{Resp}\Rsign She hath none to comfort her, save thee, O our God.\end{Resp}
\begin{Resp}\Vsign Glory be to the Father, and to the Son, \Star{} and to the Holy Ghost.\end{Resp}
\begin{Resp}\Rsign She hath none to comfort her, save thee, O our God.\end{Resp}
\switchcolumn*

\end{paracol}

\Office{Feria VI infra Hebdomadam IV Novembris}{Friday of the fourth week of November}{}
\addcontentsline{toc}{subsection}{Feria VI infra Hebdomadam IV Novembris \textit{— Friday of the fourth week of November}}
\begin{paracol}{2}
\LA
\LessonHead{Lectio I}
\switchcolumn
\EN
\LessonHead{Lesson I}
\switchcolumn*

\LA
\LessonTitle{Incipit Abdías Prophéta}
\switchcolumn
\EN
\LessonTitle{Lesson from the book of Abdias}
\switchcolumn*

\LA
\Cite{Abd 1:1-4}
\switchcolumn
\EN
\Cite{Abd 1:1-4}
\switchcolumn*

\LA
\noindent Vísio Abdíæ. Hæc dicit Dóminus Deus ad Edom: Audítum audívimus a Dómino, et legátum ad gentes misit: súrgite, et consurgámus advérsus eum in prǽlium. Ecce párvulum dedi te in géntibus: contemptíbilis tu es valde. Supérbia cordis tui éxtulit te, habitántem in scissúris petrárum, exaltántem sólium tuum; qui dicis in corde tuo: Quis détrahet me in terram? Si exaltátus fúeris ut áquila, et si inter sídera posúeris nidum tuum, inde détraham te, dicit Dóminus.\par
\switchcolumn
\EN
\noindent The vision of Abdias. Thus saith the Lord God to Edom: We have heard a rumour from the Lord, and he hath sent an ambassador to the nations: Arise, and let us rise up to battle against him. Behold I have made thee small among the nations: thou art exceeding contemptible. The pride of thy heart hath lifted thee up, who dwellest in the clefts of the rocks, and settest up thy throne on high: who sayest in thy heart: Who shall bring me down to the ground Though thou be exalted as an eagle, and though thou set thy nest among the stars: thence will I bring thee down, saith the Lord.\par
\switchcolumn*

\LA
\begin{Resp}\Rsign Super muros tuos, Jerúsalem, constítui custódes; \Star{} Tota die et nocte non tacébunt laudáre nomen Dómini.\end{Resp}
\begin{Resp}\Vsign Prædicábunt pópulis fortitúdinem meam, et annuntiábunt géntibus glóriam meam.\end{Resp}
\begin{Resp}\Rsign Tota die et nocte non tacébunt laudáre nomen Dómini.\end{Resp}
\switchcolumn
\EN
\begin{Resp}\Rsign I have set watchmen upon thy walls, O Jerusalem; \Star{} Which shall never hold their peace day nor night, to praise the name of the Lord.\end{Resp}
\begin{Resp}\Vsign They shall proclaim My might unto the nations, and declare My glory unto the Gentiles.\end{Resp}
\begin{Resp}\Rsign Which shall never hold their peace day nor night, to praise the name of the Lord.\end{Resp}
\switchcolumn*

\LA
\LessonHead{Lectio II}
\switchcolumn
\EN
\LessonHead{Lesson II}
\switchcolumn*

\LA
\Cite{Abd 1:5-7}
\switchcolumn
\EN
\Cite{Abd 1:5-7}
\switchcolumn*

\LA
\noindent Si fures introíssent ad te, si latrónes per noctem, quómodo conticuísses? nonne furáti essent sufficiéntia sibi? Si vindemiatóres introíssent ad te, numquid saltem racémum reliquíssent tibi? Quómodo scrutáti sunt Esau; investigavérunt abscóndita ejus? Usque ad términum emisérunt te: omnes viri fœ́deris tui illusérunt tibi: invaluérunt advérsum te viri pacis tuæ, qui cómedunt tecum, ponent insídias subter te; non est prudéntia in eo.\par
\switchcolumn
\EN
\noindent If thieves had gone in to thee, if robbers by night, how wouldst thou have held thy peace? would they not have stolen till they had enough if the grapegatherers had come in to thee, would they not have left thee at the least a cluster? How have they searched Esau, how have they sought out his hidden things? They have sent thee out even to the border: all the men of thy confederacy have deceived thee: the men of thy peace have prevailed against thee: they that eat with thee shall lay snares under thee: there is no wisdom in him.\par
\switchcolumn*

\LA
\begin{Resp}\Rsign Muro tuo inexpugnábili circumcínge nos, Dómine, et armis tuæ poténtiæ prótege nos semper: \Star{} Líbera, Dómine, Deus Israël, clamántes ad te.\end{Resp}
\begin{Resp}\Vsign Erue nos in mirabílibus tuis, et da glóriam nómini tuo.\end{Resp}
\begin{Resp}\Rsign Líbera, Dómine, Deus Israël, clamántes ad te.\end{Resp}
\begin{Resp}\Vsign Glória Patri, et Fílio, \Star{} et Spirítui Sancto.\end{Resp}
\begin{Resp}\Rsign Líbera, Dómine, Deus Israël, clamántes ad te.\end{Resp}
\switchcolumn
\EN
\begin{Resp}\Rsign Hedge us about with thy wall that cannot be broken down, O Lord, and shield us continually with the arms of thy might. \Star{} O Lord God of Israel, deliver them that cry unto thee.\end{Resp}
\begin{Resp}\Vsign Deliver us also according to thy marvellous works, and give glory to thy Name.\end{Resp}
\begin{Resp}\Rsign O Lord God of Israel, deliver them that cry unto thee.\end{Resp}
\begin{Resp}\Vsign Glory be to the Father, and to the Son, \Star{} and to the Holy Ghost.\end{Resp}
\begin{Resp}\Rsign O Lord God of Israel, deliver them that cry unto thee.\end{Resp}
\switchcolumn*

\end{paracol}

\Office{Sabbato infra Hebdomadam IV Novembris}{Saturday of the fourth week of November}{}
\addcontentsline{toc}{subsection}{Sabbato infra Hebdomadam IV Novembris \textit{— Saturday of the fourth week of November}}
\begin{paracol}{2}
\LA
\LessonHead{Lectio I}
\switchcolumn
\EN
\LessonHead{Lesson I}
\switchcolumn*

\LA
\LessonTitle{Incipit Jonas Prophéta}
\switchcolumn
\EN
\LessonTitle{Lesson from the book of Jonas}
\switchcolumn*

\LA
\Cite{Jonas 1:1-4}
\switchcolumn
\EN
\Cite{Jonah 1:1-4}
\switchcolumn*

\LA
\noindent Et factum est verbum Dómini ad Jonam, fílium Amáthi, dicens: Surge, et vade in Níniven, civitátem grandem, et prǽdica in ea, quia ascéndit malítia ejus coram me. Et surréxit Jonas, ut fúgeret in Tharsis a fácie Dómini, et descéndit in Joppen: et invénit navem eúntem in Tharsis, et dedit naulum ejus, et descéndit in eam ut iret cum eis in Tharsis a fácie Dómini. Dóminus autem misit ventum magnum in mare: et facta est tempéstas magna in mari, et navis periclitabátur cónteri.\par
\switchcolumn
\EN
\noindent Now the word of the Lord came to Jonas the son of Amathi, saying: Arise, and go to Ninive the great city, and preach in it: for the wickedness thereof is come up before me. And Jonas rose up to flee into Tharsis from the face of the Lord, and he went down to Joppe, and found a ship going to Tharsis: and he paid the fare thereof, and went down into it, to go with them to Tharsis from the face of the Lord. But the Lord sent a great wind into the sea: and a great tempest was raised in the sea, and the ship was in danger to be broken.\par
\switchcolumn*

\LA
\begin{Resp}\Rsign Laudábilis pópulus, \Star{} Quem Dóminus exercítuum benedíxit dicens: Opus mánuum meárum tu es, heréditas mea Israël.\end{Resp}
\begin{Resp}\Vsign Beáta gens, cujus est Dóminus Deus, pópulus eléctus in hereditátem.\end{Resp}
\begin{Resp}\Rsign Quem Dóminus exercítuum benedíxit dicens: Opus mánuum meárum tu es, heréditas mea Israël.\end{Resp}
\switchcolumn
\EN
\begin{Resp}\Rsign Blessed is the people \Star{} Whom the Lord of hosts hath blessed, saying O Israel thou art the work of Mine own hands, thou art Mine own inheritance.\end{Resp}
\begin{Resp}\Vsign Blessed is the nation whose God is the Lord, and the people whom He hath chosen for His own inheritance.\end{Resp}
\begin{Resp}\Rsign Whom the Lord of hosts hath blessed, saying O Israel thou art the work of Mine own hands, thou art Mine own inheritance.\end{Resp}
\switchcolumn*

\LA
\LessonHead{Lectio II}
\switchcolumn
\EN
\LessonHead{Lesson II}
\switchcolumn*

\LA
\Cite{Jonas 1:5-7}
\switchcolumn
\EN
\Cite{Jonah 1:5-7}
\switchcolumn*

\LA
\noindent Et timuérunt nautæ, et clamavérunt viri ad deum suum, et misérunt vasa quæ erant in navi, in mare, ut alleviarétur ab eis; et Jonas descéndit ad interióra navis, et dormiébat sopóre gravi. Et accéssit ad eum gubernátor, et dixit ei: Quid tu sopóre deprímeris? surge, ínvoca Deum tuum, si forte recógitet Deus de nobis, et non pereámus. Et dixit vir ad collégam suum: Veníte et mittámus sortes, et sciámus quare hoc malum sit nobis. Et misérunt sortes, et cécidit sors super Jonam.\par
\switchcolumn
\EN
\noindent And the mariners were afraid, and the men cried to their god: and they cast forth the wares that were in the ship, into the sea, to lighten it of them: and Jonas went down into the inner part of the ship, and fell into a deep sleep. And the shipmaster came to him, and said to him: Why art thou fast asleep? rise up, call upon thy God, if so be that God will think of us, that we may not perish. And they said every one to his fellow: Come, and let us cast lots, that we may know why this evil is upon us. And they cast lots, and the lot fell upon Jonas.\par
\switchcolumn*

\LA
\begin{Resp}\Rsign Angústiæ mihi sunt úndique, et quid éligam ignóro; \Star{} Mélius est mihi incídere in manus hóminum, quam derelínquere legem Dei mei.\end{Resp}
\begin{Resp}\Vsign Si enim hoc égero, mors mihi est; si autem non égero, non effúgiam manus vestras.\end{Resp}
\begin{Resp}\Rsign Mélius est mihi incídere in manus hóminum, quam derelínquere legem Dei mei.\end{Resp}
\switchcolumn
\EN
\begin{Resp}\Rsign I am straitened on every side, and know not what to choose. \Star{} It is better for me to fall into the hands of men, than to sin against the law of my God.\end{Resp}
\begin{Resp}\Vsign For if I do this thing, it is death unto me and if I do it not, I cannot escape your hands.\end{Resp}
\begin{Resp}\Rsign It is better for me to fall into the hands of men, than to sin against the law of my God.\end{Resp}
\switchcolumn*

\LA
\LessonHead{Lectio III}
\switchcolumn
\EN
\LessonHead{Lesson III}
\switchcolumn*

\LA
\Cite{Jonas 1:8-12}
\switchcolumn
\EN
\Cite{Jonah 1:8-12}
\switchcolumn*

\LA
\noindent Et dixérunt ad eum: Indica nobis cujus causa malum istud sit nobis: quod est opus tuum? quæ terra tua, et quo vadis? vel ex quo pópulo es tu? Et dixit ad eos: Hebrǽus ego sum, et Dóminum Deum cæli ego tímeo, qui fecit mare et áridam. Et timuérunt viri timóre magno, et dixérunt ad eum: Quid hoc fecísti? cognovérunt enim viri quod a fácie Dómini fúgeret, quia indicáverat eis. Et dixérunt ad eum: Quid faciémus tibi, et cessábit mare a nobis? quia mare ibat, et intumescébat. Et dixit ad eos: Tóllite me, et míttite in mare, et cessábit mare a vobis: scio enim ego quóniam propter me tempéstas hæc grandis venit super vos.\par
\switchcolumn
\EN
\noindent And they said to him: Tell us for what cause this evil is upon us, what is thy business? of what country art thou? and whither goest thou? or of what people art thou? And he said to them: I am a Hebrew, and I fear the Lord the God of heaven, who made both the sea and the dry land. And the men were greatly afraid, and they said to him: Why hast thou done this? (for the men knew that he fled from the face of the Lord: because he had told them.) And they said to him: What shall we do to thee, that the sea may be calm to us? for the sea flowed and swelled. And he said to them: Take me up, and cast me into the sea, and the sea shall be calm to you: for I know that for my sake this great tempest is upon you.\par
\switchcolumn*

\LA
\begin{Resp}\Rsign Misit Dóminus Angelum suum et conclúsit ora leónum, \Star{} Et non contaminavérunt, quia coram eo injustítia invénta non est in me.\end{Resp}
\begin{Resp}\Vsign Misit Deus misericórdiam suam et veritátem suam: ánimam meam erípuit de médio catulórum leónum.\end{Resp}
\begin{Resp}\Rsign Et non contaminavérunt, quia coram eo injustítia invénta non est in me.\end{Resp}
\begin{Resp}\Vsign Glória Patri, et Fílio, \Star{} et Spirítui Sancto.\end{Resp}
\begin{Resp}\Rsign Et non contaminavérunt: quia coram eo injustítia invénta non est in me.\end{Resp}
\switchcolumn
\EN
\begin{Resp}\Rsign The Lord hath sent His angel, and hath shut the lions' mouths, that \Star{} They have not hurt me forasmuch as before Him innocency was found in me.\end{Resp}
\begin{Resp}\Vsign God hath sent forth His mercy and His truth, (and delivered) my soul from among the lions' whelps.\end{Resp}
\begin{Resp}\Rsign They have not hurt me forasmuch as before Him innocency was found in me.\end{Resp}
\begin{Resp}\Vsign Glory be to the Father, and to the Son, \Star{} and to the Holy Ghost.\end{Resp}
\begin{Resp}\Rsign They have not hurt me forasmuch as before Him innocency was found in me.\end{Resp}
\switchcolumn*

\end{paracol}

\Office{Dominica V Novembris}{Fifth Sunday of November}{}
\addcontentsline{toc}{subsection}{Dominica V Novembris \textit{— Fifth Sunday of November}}
\begin{paracol}{2}
\LA
\LessonHead{Lectio I}
\switchcolumn
\EN
\LessonHead{Lesson I}
\switchcolumn*

\LA
\LessonTitle{Incipit Michǽas Prophéta}
\switchcolumn
\EN
\LessonTitle{Lesson from the book of Micheas}
\switchcolumn*

\LA
\Cite{Mic 1:1-3}
\switchcolumn
\EN
\Cite{Mic 1:1-3}
\switchcolumn*

\LA
\noindent Verbum Dómini, quod factum est ad Michǽam Morasthíten, in diébus Jóathan, Achaz, et Ezechíæ, regum Juda, quod vidit super Samaríam et Jerúsalem. Audíte, pópuli omnes, et atténdat terra, et plenitúdo ejus, et sit Dóminus Deus vobis in testem, Dóminus de templo sancto suo. Quia ecce Dóminus egrediétur de loco suo, et descéndet, et calcábit super excélsa terræ.\par
\switchcolumn
\EN
\noindent The word of the Lord that came to Micheas the Morasthite, in the days of Joathan, Achaz, and Ezechias, kings of Juda: which he saw concerning Samaria and Jerusalem. Hear, all ye people: and let the earth give ear, and all that is therein: and let the Lord God be a witness to you, the Lord from his holy temple. For behold the Lord will come forth out of his place: and he will come down, and will tread upon the high places of the earth.\par
\switchcolumn*

\LA
\begin{Resp}\Rsign Vidi Dóminum sedéntem super sólium excélsum et elevátum, et plena erat omnis terra majestáte ejus: \Star{} Et ea, quæ sub ipso erant, replébant templum.\end{Resp}
\begin{Resp}\Vsign Séraphim stabant super illud: sex alæ uni, et sex alæ álteri.\end{Resp}
\begin{Resp}\Rsign Et ea, quæ sub ipso erant, replébant templum.\end{Resp}
\switchcolumn
\EN
\begin{Resp}\Rsign I saw the Lord sitting upon a throne, high and lifted up, and the whole earth was full of His glory; \Star{} And His train filled the temple.\end{Resp}
\begin{Resp}\Vsign Above it stood the Seraphim each one had six wings.\end{Resp}
\begin{Resp}\Rsign And His train filled the temple.\end{Resp}
\switchcolumn*

\LA
\LessonHead{Lectio II}
\switchcolumn
\EN
\LessonHead{Lesson II}
\switchcolumn*

\LA
\Cite{Mic 1:4-6}
\switchcolumn
\EN
\Cite{Mic 1:4-6}
\switchcolumn*

\LA
\noindent Et consuméntur montes subtus eum, et valles scindéntur, sicut cera a fácie ignis, et sicut aquæ, quæ decúrrunt in præceps. In scélere Jacob omne istud et in peccátis domus Israël. Quod scelus Jacob? nonne Samaría? Et quæ excélsa Judæ? nonne Jerúsalem? Et ponam Samaríam quasi acérvum lápidum in agro, cum plantátur vínea, et détraham in vallem lápides ejus, et fundaménta ejus revelábo.\par
\switchcolumn
\EN
\noindent And the mountains shall be melted under him: and the valleys shall be cleft, as wax before the fire, and as waters that run down a steep place. For the wickedness of Jacob is all this, and for the sins of the house of Israel. What is the wickedness of Jacob? is it not Samaria? and what are the high places of Juda? are they not Jerusalem? And I will make Samaria as a heap of stones in the field when a vineyard is planted: and I will bring down the stones thereof into the valley, and will lay her foundations bare.\par
\switchcolumn*

\LA
\begin{Resp}\Rsign Aspice, Dómine, de sede sancta tua, et cógita de nobis: inclína, Deus meus, aurem tuam et audi: \Star{} Aperi óculos tuos et vide tribulatiónem nostram.\end{Resp}
\begin{Resp}\Vsign Qui regis Israël, inténde, qui dedúcis velut ovem Joseph.\end{Resp}
\begin{Resp}\Rsign Aperi óculos tuos et vide tribulatiónem nostram.\end{Resp}
\switchcolumn
\EN
\begin{Resp}\Rsign Look down, O Lord, from the dwelling-place of thine holiness, and take thought for us. O my God, incline thine ear, and hear. \Star{} Open thine eyes, and behold our desolation.\end{Resp}
\begin{Resp}\Vsign Give ear, O Shepherd of Israel, Thou That leadest Joseph like a flock.\end{Resp}
\begin{Resp}\Rsign Open thine eyes, and behold our desolation.\end{Resp}
\switchcolumn*

\LA
\LessonHead{Lectio III}
\switchcolumn
\EN
\LessonHead{Lesson III}
\switchcolumn*

\LA
\Cite{Mic 1:7-9}
\switchcolumn
\EN
\Cite{Mic 1:7-9}
\switchcolumn*

\LA
\noindent Et ómnia sculptília ejus concidéntur, et omnes mercédes ejus comburéntur igne, et ómnia idóla ejus ponam in perditiónem, quia de mercédibus meretrícis congregáta sunt, et usque ad mercédem meretrícis reverténtur. Super hoc plangam et ululábo, vadam spoliátus, et nudus; fáciam planctum velut dracónum, et luctum quasi struthiónum: quia desperáta est plaga ejus, quia venit usque ad Judam; tétigit portam pópuli mei usque ad Jerúsalem.\par
\switchcolumn
\EN
\noindent And all her graven things shall be cut in pieces, and all her wages shall be burnt with fire, and I will bring to destruction all her idols: for they were gathered together of the hire of a harlot, and unto the hire of a harlot they shall return. Therefore will I lament and howl: I will go stripped and naked: I will make a wailing like the dragons, and a mourning like the ostriches. Because her wound is desperate, because it is come even to Juda, it hath touched the gate of my people even to Jerusalem.\par
\switchcolumn*

\LA
\begin{Resp}\Rsign Aspice, Dómine, quia facta est desoláta cívitas plena divítiis, sedet in tristítia dómina géntium: \Star{} Non est qui consolétur eam, nisi tu, Deus noster.\end{Resp}
\begin{Resp}\Vsign Plorans plorávit in nocte, et lácrimæ ejus in maxíllis ejus.\end{Resp}
\begin{Resp}\Rsign Non est qui consolétur eam, nisi tu, Deus noster.\end{Resp}
\begin{Resp}\Vsign Glória Patri, et Fílio, \Star{} et Spirítui Sancto.\end{Resp}
\begin{Resp}\Rsign Non est qui consolétur eam, nisi tu, Deus noster.\end{Resp}
\switchcolumn
\EN
\begin{Resp}\Rsign Consider, O Lord, how that the city sitteth solitary that was full of riches; how is she become as a widow, she that was great among the nations; \Star{} She hath none to comfort her, save thee, O our God.\end{Resp}
\begin{Resp}\Vsign She weepeth sore in the night, and her tears are on her cheeks.\end{Resp}
\begin{Resp}\Rsign She hath none to comfort her, save thee, O our God.\end{Resp}
\begin{Resp}\Vsign Glory be to the Father, and to the Son, \Star{} and to the Holy Ghost.\end{Resp}
\begin{Resp}\Rsign She hath none to comfort her, save thee, O our God.\end{Resp}
\switchcolumn*

\LA
\LessonHead{Lectio IV}
\switchcolumn
\EN
\LessonHead{Lesson IV}
\switchcolumn*

\LA
\noindent Cum te appetítus inváserit peccándi, velim cógites horríbile illud et intolerábile Christi tribúnal, in quo præsidébit judex in alto et excélso throno; astábit autem omnis creatúra, ad gloriósum illíus conspéctum contremíscens. Adducéndi étiam nos sumus sínguli, eórum quæ in vita gessérimus ratiónem redditúri. Mox illis qui multa mala in vita perpetrárint, terríbiles quidam et deformes assistent angeli, igneos vultus præ se feréntes atque ignem spirántes, ea re propositi et voluntátis acerbitátem ostendéntes, nocti vultu similes, propter mærórem et ódium in humánum genus.\par
\switchcolumn
\EN
\noindent From the Sermon of St. Basil the Great, (Archbishop of Caesarea-in-Pontus,) upon the Thirty-third Psalm.\par
\switchcolumn*

\LA
\begin{Resp}\Rsign Super muros tuos, Jerúsalem, constítui custódes; \Star{} Tota die et nocte non tacébunt laudáre nomen Dómini.\end{Resp}
\begin{Resp}\Vsign Prædicábunt pópulis fortitúdinem meam, et annuntiábunt géntibus glóriam meam.\end{Resp}
\begin{Resp}\Rsign Tota die et nocte non tacébunt laudáre nomen Dómini.\end{Resp}
\switchcolumn
\EN
\begin{Resp}\Rsign I have set watchmen upon thy walls, O Jerusalem; \Star{} Which shall never hold their peace day nor night, to praise the name of the Lord.\end{Resp}
\begin{Resp}\Vsign They shall proclaim My might unto the nations, and declare My glory unto the Gentiles.\end{Resp}
\begin{Resp}\Rsign Which shall never hold their peace day nor night, to praise the name of the Lord.\end{Resp}
\switchcolumn*

\LA
\LessonHead{Lectio V}
\switchcolumn
\EN
\LessonHead{Lesson V}
\switchcolumn*

\LA
\noindent Ad hæc cógites profúndum bárathrum, inextricábiles ténebras, ignem caréntem splendore, uréndi quidem vim habéntem, sed privátum lúmine: deínde vermium genus venénum immittens, et carnem vorans, inexplebíliter edens neque umquam saturitátem sentiens, intolerábiles dolóres corrosióne ipsa infígens: postremo, quod suppliciórum ómnium gravíssimum est, oppróbrium illud et confusiónem sempiternam. Hæc time, et hoc timóre correptus ánimam a peccatórum concupiscéntia tamquam freno quodam réprime.\par
\switchcolumn
\EN
\noindent Think again of the bottomless pit, the impenetrable darkness, the lightless fire, burning, but not glowing the poisonous mass of worms, preying upon the flesh, ever feeding, and never filled, causing by their gnawing unbearable agony lastly, the greatest punishment of all, shame and confusion for ever. Have a dread of these things, and let that dread correct thee, and be as a curb to thy mind to hold it in from the hankering after sin.\par
\switchcolumn*

\LA
\begin{Resp}\Rsign Muro tuo inexpugnábili circumcínge nos, Dómine, et armis tuæ poténtiæ prótege nos semper: \Star{} Líbera, Dómine, Deus Israël, clamántes ad te.\end{Resp}
\begin{Resp}\Vsign Erue nos in mirabílibus tuis, et da glóriam nómini tuo.\end{Resp}
\begin{Resp}\Rsign Líbera, Dómine, Deus Israël, clamántes ad te.\end{Resp}
\switchcolumn
\EN
\begin{Resp}\Rsign Hedge us about with thy wall that cannot be broken down, O Lord, and shield us continually with the arms of thy might. \Star{} O Lord God of Israel, deliver them that cry unto thee.\end{Resp}
\begin{Resp}\Vsign Deliver us also according to thy marvellous works, and give glory to thy Name.\end{Resp}
\begin{Resp}\Rsign O Lord God of Israel, deliver them that cry unto thee.\end{Resp}
\switchcolumn*

\LA
\LessonHead{Lectio VI}
\switchcolumn
\EN
\LessonHead{Lesson VI}
\switchcolumn*

\LA
\noindent Hunc timórem Dómini se docturum prophéta promisit. Docére autem non simpliciter promisit, sed eos qui eum audíre volúerint: non eos qui longius prolapsi sunt, sed qui salútem appeténtes accurrunt: non alienos a promissiónibus, sed ex baptismate filiórum adoptiónis verbo ipsi consiliatos atque conjunctos. Proptérea, Veníte, inquit, hoc est, per bona ópera accédite ad me, fílii, quippe qui per regeneratiónem fílii lucis effici digni facti estis: audíte, qui aures cordis habetis apértas; timórem Dómini docébo vos: illum scilicet, quem paulo ante oratióne nostra descrípsimus.\par
\switchcolumn
\EN
\noindent This fear of the Lord the Prophet hath promised to teach. But he hath not promised to teach it to all, but only to such as will hear him not to such as have fallen far away, but to such as run to him, hungry for salvation, not to such as have no part in the promises, but to such as by baptism are born children of adoption, set at peace and oneness with the Word. Come, ye children, saith he, that is to say, Draw nigh unto me by good works, all ye who by the new birth have become the worthy children of light, hearken unto me, all ye who have the ears of your heart opened, I will teach you the fear of the Lord, even the fear of that Being of Whom we have just been speaking.\par
\switchcolumn*

\LA
\begin{Resp}\Rsign Sustinúimus pacem, et non venit: quæsívimus bona, et ecce turbátio: cognóvimus, Dómine, peccáta nostra; \Star{} Non in perpétuum obliviscáris nos.\end{Resp}
\begin{Resp}\Vsign Peccávimus, ímpie géssimus, iniquitátem fécimus, Dómine, in omnem justítiam tuam.\end{Resp}
\begin{Resp}\Rsign Non in perpétuum obliviscáris nos.\end{Resp}
\begin{Resp}\Vsign Glória Patri, et Fílio, \Star{} et Spirítui Sancto.\end{Resp}
\begin{Resp}\Rsign Non in perpétuum obliviscáris nos.\end{Resp}
\switchcolumn
\EN
\begin{Resp}\Rsign We looked for peace, and it came not we asked for good, and behold trouble. We acknowledge, O Lord, our wickedness. \Star{} Forget us not for ever.\end{Resp}
\begin{Resp}\Vsign O Lord (our God) we have sinned, we have done ungodly, we have dealt unrighteously in all thine ordinances.\end{Resp}
\begin{Resp}\Rsign Forget us not for ever.\end{Resp}
\begin{Resp}\Vsign Glory be to the Father, and to the Son, \Star{} and to the Holy Ghost.\end{Resp}
\begin{Resp}\Rsign Forget us not for ever.\end{Resp}
\switchcolumn*

\end{paracol}

\Office{Feria II infra Hebdomadam V Novembris}{Monday of the fifth week of November}{}
\addcontentsline{toc}{subsection}{Feria II infra Hebdomadam V Novembris \textit{— Monday of the fifth week of November}}
\begin{paracol}{2}
\LA
\LessonHead{Lectio I}
\switchcolumn
\EN
\LessonHead{Lesson I}
\switchcolumn*

\LA
\LessonTitle{Incipit Nahum Prophéta}
\switchcolumn
\EN
\LessonTitle{Lesson from the book of Nahum}
\switchcolumn*

\LA
\Cite{Nah 1:1-4}
\switchcolumn
\EN
\Cite{Nah 1:1-4}
\switchcolumn*

\LA
\noindent Onus Nínive. Liber visiónis Nahum Elcesǽi. Deus æmulátor et ulcíscens Dóminus, ulcíscens Dóminus, et habens furórem; ulcíscens Dóminus in hostes suos, et iráscens ipse inimícis suis. Dóminus pátiens, et magnus fortitúdine, et mundans non fáciet innocéntem; Dóminus in tempestáte et túrbine viæ ejus, et nébulæ pulvis pedum ejus; íncrepans mare, et exsíccans illud, et ómnia flúmina ad desértum dedúcens.\par
\switchcolumn
\EN
\noindent The burden of Ninive. The book of the vision of Nahum the Elcesite. The Lord is a jealous God, and a revenger: the Lord is a revenger, and hath wrath: the Lord taketh vengeance on his adversaries, and he is angry with his enemies. The Lord is patient, and great in power, and will not cleanse and acquit the guilty. The Lord's ways are in a tempest, and a whirlwind, and clouds are the dust of his feet. He rebuketh the sea, and drieth it up: and bringeth all the rivers to be a desert.\par
\switchcolumn*

\LA
\begin{Resp}\Rsign Redémit pópulum suum et liberávit eum, et vénient et exsultábunt in monte Sion et gaudébunt de bonis Dómini super fruménto, vino et óleo, \Star{} Et ultra non esúrient.\end{Resp}
\begin{Resp}\Vsign Erítque ánima eórum quasi hortus irríguus.\end{Resp}
\begin{Resp}\Rsign Et ultra non esúrient.\end{Resp}
\switchcolumn
\EN
\begin{Resp}\Rsign He hath redeemed His people, and ransomed them therefore they shall come and sing in the height of Zion, and shall rejoice in the goodness of the Lord, for wheat, and for wine, and for oil \Star{} And they shall hunger no more.\end{Resp}
\begin{Resp}\Vsign And their soul shall be as a watered garden.\end{Resp}
\begin{Resp}\Rsign And they shall hunger no more.\end{Resp}
\switchcolumn*

\LA
\LessonHead{Lectio II}
\switchcolumn
\EN
\LessonHead{Lesson II}
\switchcolumn*

\LA
\Cite{Nah 1:4-6}
\switchcolumn
\EN
\Cite{Nah 1:4-6}
\switchcolumn*

\LA
\noindent Infirmátus est Basan et Carmélus, et flos Líbani elánguit, montes commóti sunt ab eo, et colles desoláti sunt: et contrémuit terra a fácie ejus, et orbis et omnes habitántes in eo. Ante fáciem indignatiónis ejus quis stabit? et quis resístet in ira furóris ejus? Indignátio ejus effúsa est ut ignis, et petræ dissolútæ sunt ab eo.\par
\switchcolumn
\EN
\noindent Basan languisheth and Carmel: and the dower of Libanus fadeth away. The mountains tremble at him, and the hills are made desolate: and the earth hath quaked at his presence, and the world, and all that dwell therein. Who can stand before the face of his indignation? and who shall resist in the fierceness of his anger? his indignation is poured out like fire: and the rocks are melted by him.\par
\switchcolumn*

\LA
\begin{Resp}\Rsign Angústiæ mihi sunt úndique, et quid éligam ignóro; \Star{} Mélius est mihi incídere in manus hóminum, quam derelínquere legem Dei mei.\end{Resp}
\begin{Resp}\Vsign Si enim hoc égero, mors mihi est; si autem non égero, non effúgiam manus vestras.\end{Resp}
\begin{Resp}\Rsign Mélius est mihi incídere in manus hóminum, quam derelínquere legem Dei mei.\end{Resp}
\switchcolumn
\EN
\begin{Resp}\Rsign I am straitened on every side, and know not what to choose. \Star{} It is better for me to fall into the hands of men, than to sin against the law of my God.\end{Resp}
\begin{Resp}\Vsign For if I do this thing, it is death unto me and if I do it not, I cannot escape your hands.\end{Resp}
\begin{Resp}\Rsign It is better for me to fall into the hands of men, than to sin against the law of my God.\end{Resp}
\switchcolumn*

\LA
\LessonHead{Lectio III}
\switchcolumn
\EN
\LessonHead{Lesson III}
\switchcolumn*

\LA
\Cite{Nah 1:7-10}
\switchcolumn
\EN
\Cite{Nah 1:7-10}
\switchcolumn*

\LA
\noindent Bonus Dóminus, et confórtans in die tribulatiónis, et sciens sperántes in se. Et in dilúvio prætereúnte consummatiónem fáciet loci ejus, et inimícos ejus persequéntur ténebræ. Quid cogitátis contra Dóminum? Consummatiónem ipse fáciet; non consúrget duplex tribulátio. Quia, sicut spinæ se ínvicem complectúntur, sic convívium eórum páriter potántium: consuméntur quasi stípula ariditáte plena.\par
\switchcolumn
\EN
\noindent The Lord is good and giveth strength in the day of trouble: and knoweth them that hope in him. But with a flood that passeth by, he will make an utter end of the place thereof: and darkness shall pursue his enemies. What do ye devise against the Lord? he will make an utter end: there shall not rise a double affliction. For as thorns embrace one another: so while they are feasting and drinking together, they shall be consumed as stubble that is fully dry.\par
\switchcolumn*

\LA
\begin{Resp}\Rsign Misit Dóminus Angelum suum et conclúsit ora leónum, \Star{} Et non contaminavérunt, quia coram eo injustítia invénta non est in me.\end{Resp}
\begin{Resp}\Vsign Misit Deus misericórdiam suam et veritátem suam: ánimam meam erípuit de médio catulórum leónum.\end{Resp}
\begin{Resp}\Rsign Et non contaminavérunt, quia coram eo injustítia invénta non est in me.\end{Resp}
\begin{Resp}\Vsign Glória Patri, et Fílio, \Star{} et Spirítui Sancto.\end{Resp}
\begin{Resp}\Rsign Et non contaminavérunt: quia coram eo injustítia invénta non est in me.\end{Resp}
\switchcolumn
\EN
\begin{Resp}\Rsign The Lord hath sent His angel, and hath shut the lions' mouths, that \Star{} They have not hurt me forasmuch as before Him innocency was found in me.\end{Resp}
\begin{Resp}\Vsign God hath sent forth His mercy and His truth, (and delivered) my soul from among the lions' whelps.\end{Resp}
\begin{Resp}\Rsign They have not hurt me forasmuch as before Him innocency was found in me.\end{Resp}
\begin{Resp}\Vsign Glory be to the Father, and to the Son, \Star{} and to the Holy Ghost.\end{Resp}
\begin{Resp}\Rsign They have not hurt me forasmuch as before Him innocency was found in me.\end{Resp}
\switchcolumn*

\end{paracol}

\Office{Feria III infra Hebdomadam V Novembris}{Tuesday of the fifth week of November}{}
\addcontentsline{toc}{subsection}{Feria III infra Hebdomadam V Novembris \textit{— Tuesday of the fifth week of November}}
\begin{paracol}{2}
\LA
\LessonHead{Lectio I}
\switchcolumn
\EN
\LessonHead{Lesson I}
\switchcolumn*

\LA
\LessonTitle{Incipit Hábacuc Prophéta}
\switchcolumn
\EN
\LessonTitle{Lesson from the book of Habacuc}
\switchcolumn*

\LA
\Cite{Hab 1:1-4}
\switchcolumn
\EN
\Cite{Hab 1:1-4}
\switchcolumn*

\LA
\noindent Onus quod vidit Hábacuc prophéta. Usquequo, Dómine, clamábo, et non exáudies? vociferábor ad te, vim pátiens, et non salvábis? Quare ostendísti mihi iniquitátem et labórem, vidére prædam et injustítiam contra me? et factum est judícium, et contradíctio poténtior? Propter hoc laceráta est lex, et non pervénit usque ad finem judícium; quia ímpius prǽvalet advérsus justum, proptérea egréditur judícium pervérsum.\par
\switchcolumn
\EN
\noindent The burden that Habacuc the prophet saw. How long, O Lord, shall I cry, and thou wilt not hear? shall I cry out to thee suffering violence, and thou wilt not save? Why hast thou shewn me iniquity and grievance, to see rapine and injustice before me? and there is a judgment, but opposition is more powerful. Therefore the law is torn in pieces, and judgment cometh not to the end: because the wicked prevaileth against the just, therefore wrong judgment goeth forth.\par
\switchcolumn*

\LA
\begin{Resp}\Rsign A fácie furóris tui, Deus, conturbáta est omnis terra: \Star{} Sed miserére, Dómine, et ne fácias consummatiónem.\end{Resp}
\begin{Resp}\Vsign Dómine, Dóminus noster, quam admirábile est nomen tuum!\end{Resp}
\begin{Resp}\Rsign Sed miserére, Dómine, et ne fácias consummatiónem.\end{Resp}
\switchcolumn
\EN
\begin{Resp}\Rsign Before the face of thine anger, O God, the whole earth is troubled; \Star{} But Thou, O Lord, have mercy, and make not an end utterly.\end{Resp}
\begin{Resp}\Vsign Lord our Ruler, how excellent is thy Name in all the earth!\end{Resp}
\begin{Resp}\Rsign But Thou, O Lord, have mercy, and make not an end utterly.\end{Resp}
\switchcolumn*

\LA
\LessonHead{Lectio II}
\switchcolumn
\EN
\LessonHead{Lesson II}
\switchcolumn*

\LA
\Cite{Hab 1:5-7}
\switchcolumn
\EN
\Cite{Hab 1:5-7}
\switchcolumn*

\LA
\noindent Aspícite in géntibus, et vidéte, admirámini et obstupéscite; quia opus factum est in diébus vestris, quod nemo credet cum narrábitur; quia ecce ego suscitábo Chaldǽos, gentem amáram et velócem, ambulántem super latitúdinem terræ, ut possídeat tabernácula non sua. Horríbilis et terríbilis est: ex semetípsa judícium et onus ejus egrediétur.\par
\switchcolumn
\EN
\noindent Behold ye among the nations, and see: wonder, and be astonished: for a work is done in your days, which no man will believe when it shall be told. For behold, I will raise up the Chaldeans, a bitter and swift nation, marching upon the breadth of the earth, to possess the dwelling places that are not their own. They are dreadful, and terrible: from themselves shall their judgment, and their burden proceed.\par
\switchcolumn*

\LA
\begin{Resp}\Rsign Civitátem istam tu circúmda, Dómine: et Angeli tui custódiant muros ejus. \Star{} Exáudi, Dómine, pópulum tuum cum misericórdia.\end{Resp}
\begin{Resp}\Vsign Avertátur furor tuus, Dómine, a pópulo tuo et a civitáte sancta tua.\end{Resp}
\begin{Resp}\Rsign Exáudi, Dómine, pópulum tuum cum misericórdia.\end{Resp}
\switchcolumn
\EN
\begin{Resp}\Rsign Fence Thou this city; O Lord, and let thine angels keep the walls thereof. \Star{} O Lord, hearken unto thy people with mercy.\end{Resp}
\begin{Resp}\Vsign O Lord, let thine anger be turned away from thy people, and from thine holy city.\end{Resp}
\begin{Resp}\Rsign O Lord, hearken unto thy people with mercy.\end{Resp}
\switchcolumn*

\LA
\LessonHead{Lectio III}
\switchcolumn
\EN
\LessonHead{Lesson III}
\switchcolumn*

\LA
\Cite{Hab 1:8-10}
\switchcolumn
\EN
\Cite{Hab 1:8-10}
\switchcolumn*

\LA
\noindent Levióres pardis equi ejus, et velocióres lupis vespertínis: et diffundéntur équites ejus; équites namque ejus de longe vénient, volábunt quasi áquila festínans ad comedéndum. Omnes ad prædam vénient, fácies eórum ventus urens; et congregábit quasi arénam captivitátem. Et ipse de régibus triumphábit, et tyránni ridículi ejus erunt; ipse super omnem munitiónem ridébit, et comportábit ággerem, et cápiet eam.\par
\switchcolumn
\EN
\noindent Their horses are lighter than leopards, and swifter than evening wolves; and their horsemen shall be spread abroad: for their horsemen shall come from afar, they shall fly as an eagle that maketh haste to eat. They shall all come to the prey, their face is like a burning wind: and they shall gather together captives as the sand. And their prince shall triumph over kings, and princes shall be his laughingstock: and he shall laugh at every stronghold, and shall cast up a mount, and shall take it.\par
\switchcolumn*

\LA
\begin{Resp}\Rsign Genti peccatríci, pópulo pleno peccáto miserére, \Star{} Dómine Deus.\end{Resp}
\begin{Resp}\Vsign Esto placábilis super nequítiam pópuli tui.\end{Resp}
\begin{Resp}\Rsign Dómine Deus.\end{Resp}
\begin{Resp}\Vsign Glória Patri, et Fílio, \Star{} et Spirítui Sancto.\end{Resp}
\begin{Resp}\Rsign Dómine Deus.\end{Resp}
\switchcolumn
\EN
\begin{Resp}\Rsign O Lord God! have mercy upon \Star{} The sinful nation, upon the people laden with iniquity.\end{Resp}
\begin{Resp}\Vsign Let it repent thee concerning the transgression of thy people.\end{Resp}
\begin{Resp}\Rsign The people laden with iniquity.\end{Resp}
\begin{Resp}\Vsign Glory be to the Father, and to the Son, \Star{} and to the Holy Ghost.\end{Resp}
\begin{Resp}\Rsign The people laden with iniquity.\end{Resp}
\switchcolumn*

\end{paracol}

\Office{Feria IV infra Hebdomadam V Novembris}{Wednesday of the fifth week of November}{}
\addcontentsline{toc}{subsection}{Feria IV infra Hebdomadam V Novembris \textit{— Wednesday of the fifth week of November}}
\begin{paracol}{2}
\LA
\LessonHead{Lectio I}
\switchcolumn
\EN
\LessonHead{Lesson I}
\switchcolumn*

\LA
\LessonTitle{Incipit Sophonías Prophéta}
\switchcolumn
\EN
\LessonTitle{Lesson from the book of Sophonias}
\switchcolumn*

\LA
\Cite{Soph 1:1-3}
\switchcolumn
\EN
\Cite{Zeph 1:1-3}
\switchcolumn*

\LA
\noindent Verbum Dómini, quod factum est ad Sophoníam, fílium Chusi, fílii Godolíæ, fílii Amaríæ, fílii Ezecíæ, in diébus Josíæ, fílii Amon, regis Judæ. Cóngregans congregábo ómnia a fácie terræ, dicit Dóminus, cóngregans hóminem et pecus, cóngregans volatília cæli et pisces maris; et ruínæ impiórum erunt, et dispérdam hómines a fácie terræ, dicit Dóminus.\par
\switchcolumn
\EN
\noindent The word of the Lord that came to Sophonias the son of Chusi, the son of Godolias, the son of Amarias, the son of Ezechias, in the days of Josias the son of Amon king of Juda. Gathering, I will gather together all things from off the face of the land, saith the Lord: I will gather man, and beast, I will gather the birds of the air, and the fishes of the sea: and the ungodly shall meet with ruin: and I will destroy men from off the face of the land, saith the Lord.\par
\switchcolumn*

\LA
\begin{Resp}\Rsign Indicábo tibi, homo, quid sit bonum aut quid Dóminus requírat a te: \Star{} Fácere judícium et justítiam et sollícitum ambuláre cum Deo tuo.\end{Resp}
\begin{Resp}\Vsign Spera in Dómino, et fac bonitátem, et inhábita terram.\end{Resp}
\begin{Resp}\Rsign Fácere judícium et justítiam et sollícitum ambuláre cum Deo tuo.\end{Resp}
\switchcolumn
\EN
\begin{Resp}\Rsign I will show thee, O man, what is good, and what doth the Lord require of thee, but to \Star{} Do justice and judgment, and to walk humbly with thy God.\end{Resp}
\begin{Resp}\Vsign Trust in the Lord, and do good, and dwell in the land.\end{Resp}
\begin{Resp}\Rsign Do justice and judgment, and walk humbly with thy God.\end{Resp}
\switchcolumn*

\LA
\LessonHead{Lectio II}
\switchcolumn
\EN
\LessonHead{Lesson II}
\switchcolumn*

\LA
\Cite{Soph 1:4-6}
\switchcolumn
\EN
\Cite{Zeph 1:4-6}
\switchcolumn*

\LA
\noindent Et exténdam manum meam super Judam et super omnes habitántes Jerúsalem: et dispérdam de loco hoc relíquias Baal, et nómina ædituórum cum sacerdótibus, et eos qui adórant super tecta milítiam cæli, et adórant et jurant in Dómino, et jurant in Melchom, et qui avertúntur de post tergum Dómini, et qui non quæsiérunt Dóminum, nec investigavérunt eum.\par
\switchcolumn
\EN
\noindent And I will stretch out my hand upon Juda, and upon all the inhabitants of Jerusalem: and I will destroy out of this place the remnant of Baal, and the names of the wardens of the temples with the priests: And them that worship the host of heaven upon the tops of houses, and them that adore, and swear by the Lord, and swear by Melchom. And them that turn away from following after the Lord, and that have not sought the Lord, nor searched after him.\par
\switchcolumn*

\LA
\begin{Resp}\Rsign Angústiæ mihi sunt úndique, et quid éligam ignóro; \Star{} Mélius est mihi incídere in manus hóminum, quam derelínquere legem Dei mei.\end{Resp}
\begin{Resp}\Vsign Si enim hoc égero, mors mihi est; si autem non égero, non effúgiam manus vestras.\end{Resp}
\begin{Resp}\Rsign Mélius est mihi incídere in manus hóminum, quam derelínquere legem Dei mei.\end{Resp}
\switchcolumn
\EN
\begin{Resp}\Rsign I am straitened on every side, and know not what to choose. \Star{} It is better for me to fall into the hands of men, than to sin against the law of my God.\end{Resp}
\begin{Resp}\Vsign For if I do this thing, it is death unto me and if I do it not, I cannot escape your hands.\end{Resp}
\begin{Resp}\Rsign It is better for me to fall into the hands of men, than to sin against the law of my God.\end{Resp}
\switchcolumn*

\LA
\LessonHead{Lectio III}
\switchcolumn
\EN
\LessonHead{Lesson III}
\switchcolumn*

\LA
\Cite{Soph 1:7-9}
\switchcolumn
\EN
\Cite{Zeph 1:7-9}
\switchcolumn*

\LA
\noindent Siléte a fácie Dómini Dei, quia juxta est dies Dómini: quia præparávit Dóminus hóstiam, sanctificávit vocátos suos. Et erit, in die hóstiæ Dómini, visitábo super príncipes, et super fílios regis, et super omnes qui indúti sunt veste peregrína; et visitábo super omnem qui arrogánter ingréditur super limen in die illa, qui complent domum Dómini Dei sui iniquitáte et dolo.\par
\switchcolumn
\EN
\noindent Be silent before the face of the Lord God: for the day of the Lord is near, for the Lord hath prepared a victim, he hath sanctified his guests. And it shall come to pass in the day of the victim of the Lord, that I will visit upon the princes, and upon the king's sons, and upon all such as are clothed with strange apparel. And I will visit in that day upon every one that entereth arrogantly over the threshold: them that fill the house of the Lord their God with iniquity and deceit.\par
\switchcolumn*

\LA
\begin{Resp}\Rsign Misit Dóminus Angelum suum et conclúsit ora leónum, \Star{} Et non contaminavérunt, quia coram eo injustítia invénta non est in me.\end{Resp}
\begin{Resp}\Vsign Misit Deus misericórdiam suam et veritátem suam: ánimam meam erípuit de médio catulórum leónum.\end{Resp}
\begin{Resp}\Rsign Et non contaminavérunt, quia coram eo injustítia invénta non est in me.\end{Resp}
\begin{Resp}\Vsign Glória Patri, et Fílio, \Star{} et Spirítui Sancto.\end{Resp}
\begin{Resp}\Rsign Et non contaminavérunt: quia coram eo injustítia invénta non est in me.\end{Resp}
\switchcolumn
\EN
\begin{Resp}\Rsign The Lord hath sent His angel, and hath shut the lions' mouths, that \Star{} They have not hurt me forasmuch as before Him innocency was found in me.\end{Resp}
\begin{Resp}\Vsign God hath sent forth His mercy and His truth, (and delivered) my soul from among the lions' whelps.\end{Resp}
\begin{Resp}\Rsign They have not hurt me forasmuch as before Him innocency was found in me.\end{Resp}
\begin{Resp}\Vsign Glory be to the Father, and to the Son, \Star{} and to the Holy Ghost.\end{Resp}
\begin{Resp}\Rsign They have not hurt me forasmuch as before Him innocency was found in me.\end{Resp}
\switchcolumn*

\end{paracol}

\Office{Feria V infra Hebdomadam V Novembris}{Thursday of the fifth week of November}{}
\addcontentsline{toc}{subsection}{Feria V infra Hebdomadam V Novembris \textit{— Thursday of the fifth week of November}}
\begin{paracol}{2}
\LA
\LessonHead{Lectio I}
\switchcolumn
\EN
\LessonHead{Lesson I}
\switchcolumn*

\LA
\LessonTitle{Incipit Aggǽus Prophéta}
\switchcolumn
\EN
\LessonTitle{Lesson from the book of Aggeus}
\switchcolumn*

\LA
\Cite{Agg 1:1-2}
\switchcolumn
\EN
\Cite{Agg 1:1-2}
\switchcolumn*

\LA
\noindent In anno secúndo Daríi regis, in mense sexto, in die una mensis, factum est verbum Dómini in manu Aggǽi prophétæ, ad Zoróbabel, fílium Saláthiel, ducem Juda, et ad Jesum, fílium Jósedec, sacerdótem magnum, dicens: Hæc ait Dóminus exercítuum, dicens: Pópulus iste dicit: Nondum venit tempus domus Dómini ædificándæ.\par
\switchcolumn
\EN
\noindent In the second year of Darius the king, in the sixth month, in the first day of the month, the word of the Lord came by the hand of Aggeus the prophet, to Zorobabel the son of Salathiel, governor of Juda, and to Jesus the son of Josedec the high priest, saying: Thus saith the Lord of hosts, saying: This people saith: The time is not yet come for building the house of the Lord.\par
\switchcolumn*

\LA
\begin{Resp}\Rsign Vidi Dóminum sedéntem super sólium excélsum et elevátum, et plena erat omnis terra majestáte ejus: \Star{} Et ea, quæ sub ipso erant, replébant templum.\end{Resp}
\begin{Resp}\Vsign Séraphim stabant super illud: sex alæ uni, et sex alæ álteri.\end{Resp}
\begin{Resp}\Rsign Et ea, quæ sub ipso erant, replébant templum.\end{Resp}
\switchcolumn
\EN
\begin{Resp}\Rsign I saw the Lord sitting upon a throne, high and lifted up, and the whole earth was full of His glory; \Star{} And His train filled the temple.\end{Resp}
\begin{Resp}\Vsign Above it stood the Seraphim each one had six wings.\end{Resp}
\begin{Resp}\Rsign And His train filled the temple.\end{Resp}
\switchcolumn*

\LA
\LessonHead{Lectio II}
\switchcolumn
\EN
\LessonHead{Lesson II}
\switchcolumn*

\LA
\Cite{Agg 1:3-6}
\switchcolumn
\EN
\Cite{Agg 1:3-6}
\switchcolumn*

\LA
\noindent Et factum est verbum Dómini in manu Aggǽi prophétæ, dicens: Numquid tempus vobis est ut habitétis in dómibus laqueátis, et domus ista desérta? Et nunc hæc dicit Dóminus exercítuum: Pónite corda vestra super vias vestras: seminástis multum, et intulístis parum; comedístis, et non estis satiáti; bibístis, et non estis inebriáti; operuístis vos, et non estis calefácti; et qui mercédes congregávit, misit eas in sácculum pertúsum.\par
\switchcolumn
\EN
\noindent And the word of the Lord came by the hand of Aggeus the prophet, saying: Is it time for you to dwell in ceiled houses, and this house lie desolate? And now thus saith the Lord of hosts: Set your hearts to consider your ways. You have sowed much, and brought in little: you have eaten, but have not had enough: you have drunk, but have not been filled with drink: you have clothed yourselves, but have not been warmed: and he that hath earned wages, put them into a bag with holes.\par
\switchcolumn*

\LA
\begin{Resp}\Rsign Aspice, Dómine, de sede sancta tua, et cógita de nobis: inclína, Deus meus, aurem tuam et audi: \Star{} Aperi óculos tuos et vide tribulatiónem nostram.\end{Resp}
\begin{Resp}\Vsign Qui regis Israël, inténde, qui dedúcis velut ovem Joseph.\end{Resp}
\begin{Resp}\Rsign Aperi óculos tuos et vide tribulatiónem nostram.\end{Resp}
\switchcolumn
\EN
\begin{Resp}\Rsign Look down, O Lord, from the dwelling-place of thine holiness, and take thought for us. O my God, incline thine ear, and hear. \Star{} Open thine eyes, and behold our desolation.\end{Resp}
\begin{Resp}\Vsign Give ear, O Shepherd of Israel, Thou That leadest Joseph like a flock.\end{Resp}
\begin{Resp}\Rsign Open thine eyes, and behold our desolation.\end{Resp}
\switchcolumn*

\LA
\LessonHead{Lectio III}
\switchcolumn
\EN
\LessonHead{Lesson III}
\switchcolumn*

\LA
\Cite{Agg 1:7-10}
\switchcolumn
\EN
\Cite{Agg 1:7-10}
\switchcolumn*

\LA
\noindent Hæc dicit Dóminus exercítuum: Pónite corda vestra super vias vestras, ascéndite in montem, portáte ligna, et ædificáte domum: et acceptábilis mihi erit, et glorificábor, dicit Dóminus. Respexístis ad ámplius, et ecce factum est minus; et intulístis in domum, et exsufflávi illud. Quam ob causam? dicit Dóminus exercítuum, quia domus mea desérta est, et vos festinátis unusquísque in domum suam. Propter hoc super vos prohíbiti sunt cæli ne darent rorem, et terra prohíbita est ne daret germen suum.\par
\switchcolumn
\EN
\noindent Thus saith the Lord of hosts: Set your hearts upon your ways: Go up to the mountain, bring timber, and build the house: and it shall be acceptable to me, and I shall be glorified, saith the Lord. You have looked for more, and behold it became less, and you brought it home, and I blowed it away: why, saith the Lord of hosts? because my house is desolate, and you make haste every man to his own house. Therefore the heavens over you were stayed from giving dew, and the earth was hindered from yielding her fruits:\par
\switchcolumn*

\LA
\begin{Resp}\Rsign Aspice, Dómine, quia facta est desoláta cívitas plena divítiis, sedet in tristítia dómina géntium: \Star{} Non est qui consolétur eam, nisi tu, Deus noster.\end{Resp}
\begin{Resp}\Vsign Plorans plorávit in nocte, et lácrimæ ejus in maxíllis ejus.\end{Resp}
\begin{Resp}\Rsign Non est qui consolétur eam, nisi tu, Deus noster.\end{Resp}
\begin{Resp}\Vsign Glória Patri, et Fílio, \Star{} et Spirítui Sancto.\end{Resp}
\begin{Resp}\Rsign Non est qui consolétur eam, nisi tu, Deus noster.\end{Resp}
\switchcolumn
\EN
\begin{Resp}\Rsign Consider, O Lord, how that the city sitteth solitary that was full of riches; how is she become as a widow, she that was great among the nations; \Star{} She hath none to comfort her, save thee, O our God.\end{Resp}
\begin{Resp}\Vsign She weepeth sore in the night, and her tears are on her cheeks.\end{Resp}
\begin{Resp}\Rsign She hath none to comfort her, save thee, O our God.\end{Resp}
\begin{Resp}\Vsign Glory be to the Father, and to the Son, \Star{} and to the Holy Ghost.\end{Resp}
\begin{Resp}\Rsign She hath none to comfort her, save thee, O our God.\end{Resp}
\switchcolumn*

\end{paracol}

\Office{Feria VI infra Hebdomadam V Novembris}{Friday of the fifth week of November}{}
\addcontentsline{toc}{subsection}{Feria VI infra Hebdomadam V Novembris \textit{— Friday of the fifth week of November}}
\begin{paracol}{2}
\LA
\LessonHead{Lectio I}
\switchcolumn
\EN
\LessonHead{Lesson I}
\switchcolumn*

\LA
\LessonTitle{Incipit Zacharías Prophéta}
\switchcolumn
\EN
\LessonTitle{Lesson from the book of Zacharias}
\switchcolumn*

\LA
\Cite{Zach 1:1-3}
\switchcolumn
\EN
\Cite{Zech 1:1-3}
\switchcolumn*

\LA
\noindent In mense octávo, in anno secúndo Daríi regis, factum est verbum Dómini ad Zacharíam fílium Barachíæ fílii Addo prophétam, dicens: Irátus est Dóminus super patres vestros iracúndia; et dices ad eos: Hæc dicit Dóminus exercítuum: Convertímini ad me, ait Dóminus exercítuum, et convértar ad vos, dicit Dóminus exercítuum.\par
\switchcolumn
\EN
\noindent In the eighth month, in the second year of king Darius, the word of the Lord came to Zacharias the son of Barachias, the son of Addo, the prophet, saying: The Lord hath been exceeding angry with your fathers. And thou shalt say to them: Thus saith the Lord of hosts: Turn ye to me, saith the Lord of hosts: and I will turn to you, saith the Lord of hosts.\par
\switchcolumn*

\LA
\begin{Resp}\Rsign Super muros tuos, Jerúsalem, constítui custódes; \Star{} Tota die et nocte non tacébunt laudáre nomen Dómini.\end{Resp}
\begin{Resp}\Vsign Prædicábunt pópulis fortitúdinem meam, et annuntiábunt géntibus glóriam meam.\end{Resp}
\begin{Resp}\Rsign Tota die et nocte non tacébunt laudáre nomen Dómini.\end{Resp}
\switchcolumn
\EN
\begin{Resp}\Rsign I have set watchmen upon thy walls, O Jerusalem; \Star{} Which shall never hold their peace day nor night, to praise the name of the Lord.\end{Resp}
\begin{Resp}\Vsign They shall proclaim My might unto the nations, and declare My glory unto the Gentiles.\end{Resp}
\begin{Resp}\Rsign Which shall never hold their peace day nor night, to praise the name of the Lord.\end{Resp}
\switchcolumn*

\LA
\LessonHead{Lectio II}
\switchcolumn
\EN
\LessonHead{Lesson II}
\switchcolumn*

\LA
\Cite{Zach 1:4-5}
\switchcolumn
\EN
\Cite{Zech 1:4-5}
\switchcolumn*

\LA
\noindent Ne sitis sicut patres vestri, ad quos clamábant prophétæ prióres, dicéntes: Hæc dicit Dóminus exercítuum: Convertímini de viis vestris malis, et de cogitatiónibus vestris péssimis; et non audiérunt, neque attendérunt ad me, dicit Dóminus. Patres vestri, ubi sunt? et prophétæ numquid in sempitérnum vivent?\par
\switchcolumn
\EN
\noindent Be not as your fathers, to whom the former prophets have cried, saying: Thus saith the Lord of hosts: Turn ye from your evil ways, and from your wicked thoughts: but they did not give ear, neither did they hearken to me, saith the Lord. Your fathers, where are they? and the prophets, shall they live always?\par
\switchcolumn*

\LA
\begin{Resp}\Rsign Muro tuo inexpugnábili circumcínge nos, Dómine, et armis tuæ poténtiæ prótege nos semper: \Star{} Líbera, Dómine, Deus Israël, clamántes ad te.\end{Resp}
\begin{Resp}\Vsign Erue nos in mirabílibus tuis, et da glóriam nómini tuo.\end{Resp}
\begin{Resp}\Rsign Líbera, Dómine, Deus Israël, clamántes ad te.\end{Resp}
\switchcolumn
\EN
\begin{Resp}\Rsign Hedge us about with thy wall that cannot be broken down, O Lord, and shield us continually with the arms of thy might. \Star{} O Lord God of Israel, deliver them that cry unto thee.\end{Resp}
\begin{Resp}\Vsign Deliver us also according to thy marvellous works, and give glory to thy Name.\end{Resp}
\begin{Resp}\Rsign O Lord God of Israel, deliver them that cry unto thee.\end{Resp}
\switchcolumn*

\LA
\LessonHead{Lectio III}
\switchcolumn
\EN
\LessonHead{Lesson III}
\switchcolumn*

\LA
\Cite{Zach 1:6}
\switchcolumn
\EN
\Cite{Zech 1:6}
\switchcolumn*

\LA
\noindent Verúmtamen verba mea, et legítima mea, quæ mandávi servis meis prophétis, numquid non comprehendérunt patres vestros, et convérsi sunt, et dixérunt: Sicut cogitávit Dóminus exercítuum fácere nobis secúndum vias nostras, et secúndum adinventiónes nostras, fecit nobis?\par
\switchcolumn
\EN
\noindent But yet my words, and my ordinances, which I gave in charge to my servants the prophets, did they not take hold of your fathers, and they returned, and said: As the Lord of hosts thought to do to us according to our ways, and according to our devices, so he hath done to us.\par
\switchcolumn*

\LA
\begin{Resp}\Rsign Sustinúimus pacem, et non venit: quæsívimus bona, et ecce turbátio: cognóvimus, Dómine, peccáta nostra; \Star{} Non in perpétuum obliviscáris nos.\end{Resp}
\begin{Resp}\Vsign Peccávimus, ímpie géssimus, iniquitátem fécimus, Dómine, in omnem justítiam tuam.\end{Resp}
\begin{Resp}\Rsign Non in perpétuum obliviscáris nos.\end{Resp}
\begin{Resp}\Vsign Glória Patri, et Fílio, \Star{} et Spirítui Sancto.\end{Resp}
\begin{Resp}\Rsign Non in perpétuum obliviscáris nos.\end{Resp}
\switchcolumn
\EN
\begin{Resp}\Rsign We looked for peace, and it came not we asked for good, and behold trouble. We acknowledge, O Lord, our wickedness. \Star{} Forget us not for ever.\end{Resp}
\begin{Resp}\Vsign O Lord (our God) we have sinned, we have done ungodly, we have dealt unrighteously in all thine ordinances.\end{Resp}
\begin{Resp}\Rsign Forget us not for ever.\end{Resp}
\begin{Resp}\Vsign Glory be to the Father, and to the Son, \Star{} and to the Holy Ghost.\end{Resp}
\begin{Resp}\Rsign Forget us not for ever.\end{Resp}
\switchcolumn*

\end{paracol}

\Office{Sabbato infra Hebdomadam V Novembris}{Saturday of the fifth week of November}{}
\addcontentsline{toc}{subsection}{Sabbato infra Hebdomadam V Novembris \textit{— Saturday of the fifth week of November}}
\begin{paracol}{2}
\LA
\LessonHead{Lectio I}
\switchcolumn
\EN
\LessonHead{Lesson I}
\switchcolumn*

\LA
\LessonTitle{Incipit Malachías Prophéta}
\switchcolumn
\EN
\LessonTitle{Lesson from the book of Malachias}
\switchcolumn*

\LA
\Cite{Mal 1:1-4}
\switchcolumn
\EN
\Cite{Mal 1:1-4}
\switchcolumn*

\LA
\noindent Onus verbi Dómini ad Israël in manu Malachíæ. Diléxi vos, dicit Dóminus, et dixístis: In quo dilexísti nos? Nonne frater erat Esau Jacob? dicit Dóminus, et diléxi Jacob, Esau autem ódio hábui, et pósui montes ejus in solitúdinem, et hereditátem ejus in dracónes desérti? Quod si díxerit Idumǽa: Destrúcti sumus, sed reverténtes ædificábimus quæ destrúcta sunt: hæc dicit Dóminus exercítuum: Isti ædificábunt, et ego déstruam; et vocabúntur términi impietátis, et pópulus cui irátus est Dóminus usque in ætérnum.\par
\switchcolumn
\EN
\noindent The burden of the word of the Lord to Israel by the hand of Malachias. I have loved you, saith the Lord: and you have said: Wherein hast thou loved us? Was not Esau brother to Jacob, saith the Lord, and I have loved Jacob, But have hated Esau? and I have made his mountains a wilderness, and given his inheritance to the dragons of the desert. But if Edom shall say: We are destroyed, but we will return and build up what hath been destroyed: thus saith the Lord of hosts: They shall build up, and I will throw down: and they shall be called the borders of wickedness, and the people with whom the Lord is angry for ever.\par
\switchcolumn*

\LA
\begin{Resp}\Rsign Laudábilis pópulus, \Star{} Quem Dóminus exercítuum benedíxit dicens: Opus mánuum meárum tu es, heréditas mea Israël.\end{Resp}
\begin{Resp}\Vsign Beáta gens, cujus est Dóminus Deus, pópulus eléctus in hereditátem.\end{Resp}
\begin{Resp}\Rsign Quem Dóminus exercítuum benedíxit dicens: Opus mánuum meárum tu es, heréditas mea Israël.\end{Resp}
\switchcolumn
\EN
\begin{Resp}\Rsign Blessed is the people \Star{} Whom the Lord of hosts hath blessed, saying O Israel thou art the work of Mine own hands, thou art Mine own inheritance.\end{Resp}
\begin{Resp}\Vsign Blessed is the nation whose God is the Lord, and the people whom He hath chosen for His own inheritance.\end{Resp}
\begin{Resp}\Rsign Whom the Lord of hosts hath blessed, saying O Israel thou art the work of Mine own hands, thou art Mine own inheritance.\end{Resp}
\switchcolumn*

\LA
\LessonHead{Lectio II}
\switchcolumn
\EN
\LessonHead{Lesson II}
\switchcolumn*

\LA
\Cite{Mal 1:5-7}
\switchcolumn
\EN
\Cite{Mal 1:5-7}
\switchcolumn*

\LA
\noindent Et óculi vestri vidébunt, et vos dicétis: Magnificétur Dóminus super términum Israël. Fílius honórat patrem, et servus dóminum suum; si ergo Pater ego sum, ubi est honor meus? et, si Dóminus ego sum, ubi est timor meus? dicit Dóminus exercítuum, ad vos, o sacerdótes, qui despícitis nomen meum, et dixístis: In quo despéximus nomen tuum? Offértis super altáre meum panem pollútum, et dícitis: In quo pollúimus te? In eo quod dícitis: Mensa Dómini despécta est.\par
\switchcolumn
\EN
\noindent And your eyes shall see, and you shall say: The Lord be magnified upon the border of Israel. The son honoureth the father, and the servant his master: if then I be a father, where is my honour? and if I be a master, where is my fear? saith the Lord of hosts. To you, O priests, that despise my name, and have said: Wherein have we despised thy name? You offer polluted bread upon my altar, and you say: Wherein have we polluted thee? In that you say: The table of the Lord is contemptible.\par
\switchcolumn*

\LA
\begin{Resp}\Rsign Angústiæ mihi sunt úndique, et quid éligam ignóro; \Star{} Mélius est mihi incídere in manus hóminum, quam derelínquere legem Dei mei.\end{Resp}
\begin{Resp}\Vsign Si enim hoc égero, mors mihi est; si autem non égero, non effúgiam manus vestras.\end{Resp}
\begin{Resp}\Rsign Mélius est mihi incídere in manus hóminum, quam derelínquere legem Dei mei.\end{Resp}
\switchcolumn
\EN
\begin{Resp}\Rsign I am straitened on every side, and know not what to choose. \Star{} It is better for me to fall into the hands of men, than to sin against the law of my God.\end{Resp}
\begin{Resp}\Vsign For if I do this thing, it is death unto me and if I do it not, I cannot escape your hands.\end{Resp}
\begin{Resp}\Rsign It is better for me to fall into the hands of men, than to sin against the law of my God.\end{Resp}
\switchcolumn*

\LA
\LessonHead{Lectio III}
\switchcolumn
\EN
\LessonHead{Lesson III}
\switchcolumn*

\LA
\Cite{Mal 1:8-11}
\switchcolumn
\EN
\Cite{Mal 1:8-11}
\switchcolumn*

\LA
\noindent Si offerátis cæcum ad immolándum, nonne malum est? et, si offerátis claudum et lánguidum, nonne malum est? Offer illud duci tuo, si placúerit ei, aut si suscéperit fáciem tuam, dicit Dóminus exercítuum. Et nunc deprecámini vultum Dei, ut misereátur vestri (de manu enim vestra factum est hoc), si quómodo suscípiat fácies vestras, dicit Dóminus exercítuum. Quis est in vobis qui claudat óstia, et incéndat altáre meum gratuíto? Non est mihi volúntas in vobis, dicit Dóminus exercítuum, et munus non suscípiam de manu vestra; ab ortu enim solis usque ad occásum, magnum est nomen meum in géntibus, et in omni loco sacrificátur: et offértur nómini meo oblátio munda, quia magnum est nomen meum in géntibus, dicit Dóminus exercítuum.\par
\switchcolumn
\EN
\noindent If you offer the blind for sacrifice, is it not evil? and if you offer the lame and the sick, is it not evil? offer it to thy prince, if he will be pleased with it, or if he will regard thy face, saith the Lord of hosts. And now beseech ye the face of God, that he may have mercy on you, (for by your hand hath this been done,) if by any means he will receive your faces, saith the Lord of hosts. Who is there among you, that will shut the doors, and will kindle the fire on my altar gratis? I have no pleasure in you, saith the Lord of hosts: and I will not receive a gift of your hand. For from the rising of the sun even to the going down, my name is great among the Gentiles, and in every place there is sacrifice, and there is offered to my name a clean oblation: for my name is great among the Gentiles, saith the Lord of hosts.\par
\switchcolumn*

\LA
\begin{Resp}\Rsign Misit Dóminus Angelum suum et conclúsit ora leónum, \Star{} Et non contaminavérunt, quia coram eo injustítia invénta non est in me.\end{Resp}
\begin{Resp}\Vsign Misit Deus misericórdiam suam et veritátem suam: ánimam meam erípuit de médio catulórum leónum.\end{Resp}
\begin{Resp}\Rsign Et non contaminavérunt, quia coram eo injustítia invénta non est in me.\end{Resp}
\begin{Resp}\Vsign Glória Patri, et Fílio, \Star{} et Spirítui Sancto.\end{Resp}
\begin{Resp}\Rsign Et non contaminavérunt: quia coram eo injustítia invénta non est in me.\end{Resp}
\switchcolumn
\EN
\begin{Resp}\Rsign The Lord hath sent His angel, and hath shut the lions' mouths, that \Star{} They have not hurt me forasmuch as before Him innocency was found in me.\end{Resp}
\begin{Resp}\Vsign God hath sent forth His mercy and His truth, (and delivered) my soul from among the lions' whelps.\end{Resp}
\begin{Resp}\Rsign They have not hurt me forasmuch as before Him innocency was found in me.\end{Resp}
\begin{Resp}\Vsign Glory be to the Father, and to the Son, \Star{} and to the Holy Ghost.\end{Resp}
\begin{Resp}\Rsign They have not hurt me forasmuch as before Him innocency was found in me.\end{Resp}
\switchcolumn*

\end{paracol}

\SeasonTitle{Dominicæ quæ superfuerunt post Epiphaniam}{Sundays after Epiphany resumed after Pentecost}

\addcontentsline{toc}{section}{Dominicæ quæ superfuerunt post Epiphaniam \textit{— Sundays after Epiphany resumed after Pentecost}}

\Office{Dominica III Post Epiphaniam}{Third Sunday after Epiphany}{Semiduplex}
\addcontentsline{toc}{subsection}{Dominica III Post Epiphaniam \textit{— Third Sunday after Epiphany}}
\begin{paracol}{2}
\LA
\LessonHead{Lectio I}
\switchcolumn
\EN
\LessonHead{Lesson I}
\switchcolumn*

\LA
\LessonTitle{Incipit Epístola beáti Pauli Apóstoli ad Gálatas}
\switchcolumn
\EN
\LessonTitle{Lesson from the letter of St. Paul the Apostle to the Galatians}
\switchcolumn*

\LA
\Cite{Gal 1:1-5}
\switchcolumn
\EN
\Cite{Gal 1:1-5}
\switchcolumn*

\LA
\noindent Paulus Apóstolus non ab homínibus, neque per hóminem, sed per Jesum Christum, et Deum Patrem, qui suscitávit eum a mórtuis: et qui mecum sunt omnes fratres, ecclésiis Galátiæ. Grátia vobis, et pax a Deo Patre, et Dómino nostro Jesu Christo, qui dedit semetípsum pro peccátis nostris, ut eríperet nos de præsénti sǽculo nequam, secúndum voluntátem Dei et Patris nostri, cui est glória in sǽcula sæculórum. Amen.\par
\switchcolumn
\EN
\noindent Paul, an apostle, not of men, neither by man, but by Jesus Christ, and God the Father, who raised him from the dead, And all the brethren who are with me, to the churches of Galatia. Grace be to you, and peace from God the Father, and from our Lord Jesus Christ, Who gave himself for our sins, that he might deliver us from this present wicked world, according to the will of God and our Father: To whom is glory for ever and ever. Amen.\par
\switchcolumn*

\LA
\begin{Resp}\Rsign Dómine, ne in ira tua árguas me, neque in furóre tuo corrípias me: \Star{} Miserére mei, Dómine, quóniam infírmus sum.\end{Resp}
\begin{Resp}\Vsign Timor et tremor venérunt super me, et contexérunt me ténebræ.\end{Resp}
\begin{Resp}\Rsign Miserére mei, Dómine, quóniam infírmus sum.\end{Resp}
\switchcolumn
\EN
\begin{Resp}\Rsign O Lord, rebuke me not in thine anger, neither chasten me in thine hot displeasure. \Star{} Have mercy upon me, O Lord, for I am weak.\end{Resp}
\begin{Resp}\Vsign Fearfulness and trembling are come upon me, and darkness hath overwhelmed me.\end{Resp}
\begin{Resp}\Rsign Have mercy upon me, O Lord, for I am weak.\end{Resp}
\switchcolumn*

\LA
\LessonHead{Lectio II}
\switchcolumn
\EN
\LessonHead{Lesson II}
\switchcolumn*

\LA
\Cite{Gal 1:6-10}
\switchcolumn
\EN
\Cite{Gal 1:6-10}
\switchcolumn*

\LA
\noindent Miror quod sic tam cito transferímini ab eo, qui vos vocávit in grátiam Christi in áliud Evangélium: quod non est áliud, nisi sunt áliqui, qui vos contúrbant, et volunt convértere Evangélium Christi. Sed licet nos, aut Angelus de cælo evangelízet vobis prætérquam quod evangelizávimus vobis, anáthema sit. Sicut prædíximus, et nunc íterum dico: Si quis vobis evangelizáverit præter id, quod accepístis, anáthema sit. Modo enim homínibus suádeo, an Deo? An quæro homínibus placére? Si adhuc homínibus placérem, Christi servus non essem.\par
\switchcolumn
\EN
\noindent I wonder that you are so soon removed from him that called you into the grace of Christ, unto another gospel. Which is not another, only there are some that trouble you, and would pervert the gospel of Christ. But though we, or an angel from heaven, preach a gospel to you besides that which we have preached to you, let him be anathema. As we said before, so now I say again: If any one preach to you a gospel, besides that which you have received, let him be anathema. For do I now persuade men, or God? Or do I seek to please men? If I yet pleased men, I should not be the servant of Christ.\par
\switchcolumn*

\LA
\begin{Resp}\Rsign Deus, qui sedes super thronum, et júdicas æquitátem, esto refúgium páuperum in tribulatióne: \Star{} Quia tu solus labórem et dolórem consíderas.\end{Resp}
\begin{Resp}\Vsign Tibi enim derelíctus est pauper, pupíllo tu eris adjútor.\end{Resp}
\begin{Resp}\Rsign Quia tu solus labórem et dolórem consíderas.\end{Resp}
\switchcolumn
\EN
\begin{Resp}\Rsign O God, Which satest in the throne judging right, be Thou a refuge for the poor, a refuge in times of trouble. \Star{} For Thou alone beholdest mischief and spite.\end{Resp}
\begin{Resp}\Vsign The poor leaveth himself unto thee; Thou wilt be the helper of the fatherless.\end{Resp}
\begin{Resp}\Rsign For Thou alone beholdest mischief and spite.\end{Resp}
\switchcolumn*

\LA
\LessonHead{Lectio III}
\switchcolumn
\EN
\LessonHead{Lesson III}
\switchcolumn*

\LA
\Cite{Gal 1:11-14}
\switchcolumn
\EN
\Cite{Gal 1:11-14}
\switchcolumn*

\LA
\noindent Notum enim vobis fácio, fratres, Evangélium, quod evangelizátum est a me, quia non est secúndum hóminem: neque enim ego ab hómine accépi illud, neque dídici, sed per revelatiónem Jesu Christi. Audístis enim conversatiónem meam aliquándo in Judaísmo: quóniam supra modum persequébar Ecclésiam Dei, et expugnábam illam, et proficiébam in Judaísmo supra multos coætáneos meos in génere meo, abundántius æmulátor exístens paternárum meárum traditiónum.\par
\switchcolumn
\EN
\noindent For I give you to understand, brethren, that the gospel which was preached by me is not according to man. For neither did I receive it of man, nor did I learn it; but by the revelation of Jesus Christ. For you have heard of my conversation in time past in the Jews' religion: how that, beyond measure, I persecuted the church of God, and wasted it. And I made progress in the Jews' religion above many of my equals in my own nation, being more abundantly zealous for the traditions of my fathers.\par
\switchcolumn*

\LA
\begin{Resp}\Rsign A dextris est mihi Dóminus, ne commóvear: \Star{} Propter hoc dilatátum est cor meum, et exsultávit lingua mea.\end{Resp}
\begin{Resp}\Vsign Dóminus pars hereditátis meæ, et cálicis mei.\end{Resp}
\begin{Resp}\Rsign Propter hoc dilatátum est cor meum, et exsultávit lingua mea.\end{Resp}
\begin{Resp}\Vsign Glória Patri, et Fílio, \Star{} et Spirítui Sancto.\end{Resp}
\begin{Resp}\Rsign Propter hoc dilatátum est cor meum, et exsultávit lingua mea.\end{Resp}
\switchcolumn
\EN
\begin{Resp}\Rsign The Lord is at my right hand, I shall never be moved. \Star{} Therefore my heart is glad, and my tongue rejoiceth.\end{Resp}
\begin{Resp}\Vsign The Lord is the portion of mine inheritance, and of my cup.\end{Resp}
\begin{Resp}\Rsign Therefore my heart is glad, and my tongue rejoiceth.\end{Resp}
\begin{Resp}\Vsign Glory be to the Father, and to the Son, \Star{} and to the Holy Ghost.\end{Resp}
\begin{Resp}\Rsign Therefore my heart is glad, and my tongue rejoiceth.\end{Resp}
\switchcolumn*

\LA
\LessonHead{Lectio IV}
\switchcolumn
\EN
\LessonHead{Lesson IV}
\switchcolumn*

\LA
\LessonTitle{De Expositióne sancti Augustíni Epíscopi in Epístolam ad Gálatas}
\switchcolumn
\EN
\LessonTitle{From the Exposition of the Epistle to the Galatians by St. Augustine, Bishop (of Hippo.)}
\switchcolumn*

\LA
\Source{Præfátio tom. 4}
\switchcolumn
\EN
\Source{Preface, Bk. iv.}
\switchcolumn*

\LA
\noindent Causa, propter quam scribit Apóstolus ad Gálatas, hæc est: ut intélligant grátiam Dei id secum ágere, ut sub lege jam non sint. Cum enim prædicáta eis esset Evangélii grátia, non defuérunt quidam ex circumcisióne, quamvis Christiáni nómine, nondum tamen tenéntes ipsum gratiæ benefícium, et adhuc voléntes esse sub onéribus legis, quæ Dóminus Deus imposúerat non justítiæ serviéntibus, sed peccáto, justam scílicet legem injústis homínibus dando ad demonstranda peccáta eórum, non auferenda. Non enim aufert peccáta, nisi grátia fidei, quæ per dilectiónem operátur.\par
\switchcolumn
\EN
\noindent The reason of the Apostle's writing to the Galatians was this that they might understand that the grace of God had worked in them that they were no longer under the law. For when the grace of the Gospel was preached to them, there had not been wanting to them some of them of the circumcision, Christians indeed in name, but who had not yet apprehended that great benefit of grace, and desiring still to be bound with burdens of the law burdens which the Lord God had laid, not upon such as serve righteousness, but upon such as serve sin, laying, that is to say, upon the unrighteous a righteous law, whereby their unrighteousness was made manifest, not taken away. For there is not anything which taketh away sin, save only the grace of faith which worketh by love.\par
\switchcolumn*

\LA
\begin{Resp}\Rsign Notas mihi fecísti, Dómine, vias vitæ: \Star{} Adimplébis me lætítia cum vultu tuo: delectatiónes in déxtera tua usque in finem.\end{Resp}
\begin{Resp}\Vsign Tu es qui restítues hereditátem meam mihi.\end{Resp}
\begin{Resp}\Rsign Adimplébis me lætítia cum vultu tuo: delectatiónes in déxtera tua usque in finem.\end{Resp}
\switchcolumn
\EN
\begin{Resp}\Rsign O Lord, Thou hast shown me the path of life. \Star{} Thou shalt fill me with joy in thy presence, at thy right hand there are pleasures for evermore.\end{Resp}
\begin{Resp}\Vsign Thou art He That shalt restore mine inheritance unto me.\end{Resp}
\begin{Resp}\Rsign Thou shalt fill me with joy in thy presence, at thy right hand there are pleasures for evermore.\end{Resp}
\switchcolumn*

\LA
\LessonHead{Lectio V}
\switchcolumn
\EN
\LessonHead{Lesson V}
\switchcolumn*

\LA
\noindent Sub hac ergo grátia jam Gálatas constitútos illi volébant constitúere sub onéribus legis, asseverántes nihil eis prodésse Evangélium, nisi circumcideréntur, et céteras carnáles Judáici ritus observatiónes subírent. Et ídeo Paulum Apóstolum suspéctum habére cœ́perant, a quo illis Evangélium prædicátum erat, tamquam non tenéntem disciplínam ceterórum Apostolórum, qui gentes cogébant judaizáre.\par
\switchcolumn
\EN
\noindent The men of the circumcision would have the Galatians, who were under grace, to be under the burdens of the law, persuading them that the Gospel profited them nothing, unless they should be circumcised, and take on them the other outward observances of the Jews' religion. Whence the Galatians began to have doubts of the Apostle Paul, by whom the Gospel had been preached to them, as one that held not the doctrine of the other Apostles, who compelled the Gentiles to come under the law.\par
\switchcolumn*

\LA
\begin{Resp}\Rsign Díligam te, Dómine, virtus mea: Dóminus firmaméntum meum, \Star{} Et refúgium meum.\end{Resp}
\begin{Resp}\Vsign Liberátor meus, Deus meus, adjútor meus.\end{Resp}
\begin{Resp}\Rsign Et refúgium meum.\end{Resp}
\switchcolumn
\EN
\begin{Resp}\Rsign I will love thee, O Lord, my strength; the Lord is my rock \Star{} And my fortress.\end{Resp}
\begin{Resp}\Vsign My Deliverer, my God, mine Helper.\end{Resp}
\begin{Resp}\Rsign And my fortress.\end{Resp}
\switchcolumn*

\LA
\LessonHead{Lectio VI}
\switchcolumn
\EN
\LessonHead{Lesson VI}
\switchcolumn*

\LA
\noindent Talis quidem quǽstio est et in Epístola ad Romános: verúmtamen vidétur áliquid interésse, quod ibi contentiónem ipsam dírimit, litémque compónit, quæ inter eos, qui ex Judǽis, et eos, qui ex Géntibus credíderant, orta erat: cum illi tamquam ex méritis óperum legis, sibi rédditum Evangélii prǽmium arbitraréntur, quod prǽmium incircumcísis tamquam imméritis nolébant dari: illi contra Judǽis se præférre gestírent, tamquam interfectóribus Dómini. In hac vero epístola ad eos scribit, qui jam commóti erant auctoritáte illórum, qui ex Judǽis erant, et ad observatiónes legis cogébant.\par
\switchcolumn
\EN
\noindent The same question is discussed in the Epistle to the Romans, but with this difference in that case the Apostle putteth an end to the discussion, and stilleth the strife which had arisen between the Jewish and the Gentile converts, in consequence of the Jews holding that they had earned the knowledge of the Gospel as a reward for their observance of the law, and grudging the same knowledge to the uncircumcised, as to men who had done nothing to deserve it; and the Gentiles, on the contrary, maintaining that they were superior to the Jews, in that they were not the murderers of the Lord. Now, in this Epistle to the Galatians, the Apostle addresseth himself to those who were troubled by the authority claimed by them who were of the circumcision, and sought to bring into subjection to the law them who were of the uncircumcision.\par
\switchcolumn*

\LA
\begin{Resp}\Rsign Dómini est terra, et plenitúdo ejus: \Star{} Orbis terrárum, et univérsi qui hábitant in eo.\end{Resp}
\begin{Resp}\Vsign Ipse super mária fundávit eam, et super flúmina præparávit illam.\end{Resp}
\begin{Resp}\Rsign Orbis terrárum, et univérsi qui hábitant in eo.\end{Resp}
\begin{Resp}\Vsign Glória Patri, et Fílio, \Star{} et Spirítui Sancto.\end{Resp}
\begin{Resp}\Rsign Orbis terrárum, et univérsi qui hábitant in eo.\end{Resp}
\switchcolumn
\EN
\begin{Resp}\Rsign The earth is the Lord's, and the fulness thereof \Star{} The world, and they that dwell therein.\end{Resp}
\begin{Resp}\Vsign For He hath founded it upon the seas, and established it upon the floods.\end{Resp}
\begin{Resp}\Rsign The world, and they that dwell therein.\end{Resp}
\begin{Resp}\Vsign Glory be to the Father, and to the Son, \Star{} and to the Holy Ghost.\end{Resp}
\begin{Resp}\Rsign The world, and they that dwell therein.\end{Resp}
\switchcolumn*

\LA
\LessonHead{Lectio VII}
\switchcolumn
\EN
\LessonHead{Lesson VII}
\switchcolumn*

\LA
\LessonTitle{Léctio sancti Evangélii secúndum Matthǽum}
\switchcolumn
\EN
\LessonTitle{From the Holy Gospel according to Matthew}
\switchcolumn*

\LA
\Cite{Matt 8:1-13}
\switchcolumn
\EN
\Cite{Matt 8:1-13}
\switchcolumn*

\LA
\noindent In illo témpore: Cum descendísset Jesus de monte, secútæ sunt eum turbæ multæ: et ecce leprósus véniens, adorábat eum. Et réliqua.\par
\switchcolumn
\EN
\noindent At that time, when Jesus was come down from the mountain, great multitudes followed him: And behold a leper came and adored him. And so on.\par
\switchcolumn*

\LA
\LessonTitle{Homilía sancti Hierónymi Presbýteri}
\switchcolumn
\EN
\LessonTitle{Homily by St. Jerome, Priest (at Bethlehem.)}
\switchcolumn*

\LA
\Source{Liber 1 Comment. in cap. 8 Matth.}
\switchcolumn
\EN
\Source{Bk. i, Comm. on Matth. viii}
\switchcolumn*

\LA
\noindent De monte Dómino descendénte, occúrrunt turbæ, quia ad altióra ascéndere non valuérunt. Et primus ei occúrrit leprósus: necdum enim póterat cum lepra tam multíplicem in monte Salvatóris audíre sermónem. Et notándum, quod hic primus speciáliter curátus sit: secúndo, puer centuriónis: tértio, socrus Petri fébriens in Caphárnaum: quarto loco, qui obláti sunt ei a dæmónio vexáti: quorum spíritus verbo eiciébat, quando et omnes male habéntes curávit.\par
\switchcolumn
\EN
\noindent When the Lord was come down from the mountain, great multitudes followed Him. They were not able to follow Him when He went up. And first there came a leper. This poor creature's disease had prevented him from hearing the Saviour's long sermon on the Mount. Let it be noted that he is the first person specially named as being healed. The second was the Centurion's servant; the third was Peter's wife's mother, who was sick of a fever at Capernaum; the fourth were they who were brought unto Christ as being troubled with evil spirits, from whom He by His word cast out the evil spirits, at the same time that He healed all that were sick.\par
\switchcolumn*

\LA
\begin{Resp}\Rsign Ad te, Dómine, levávi ánimam meam: \Star{} Deus meus, in te confído, non erubéscam.\end{Resp}
\begin{Resp}\Vsign Custódi ánimam meam, et éripe me.\end{Resp}
\begin{Resp}\Rsign Deus meus, in te confído, non erubéscam.\end{Resp}
\switchcolumn
\EN
\begin{Resp}\Rsign Unto thee, O Lord, do I lift up my soul. \Star{} O my God, I trust in thee, let me not be ashamed.\end{Resp}
\begin{Resp}\Vsign O keep my soul and deliver me.\end{Resp}
\begin{Resp}\Rsign O my God, I trust in thee, let me not be ashamed.\end{Resp}
\switchcolumn*

\LA
\LessonHead{Lectio VIII}
\switchcolumn
\EN
\LessonHead{Lesson VIII}
\switchcolumn*

\LA
\noindent Et ecce leprósus véniens adorábat eum, dicens. Recte post prædicatiónem atque doctrínam, signi offértur occásio, ut per virtútem miráculi, prætéritus apud audiéntes sermo firmétur. Dómine, si vis, potes me mundáre. Qui voluntátem rogat, de virtúte non dúbitat. Et exténdens Jesus manum tétigit eum, dicens: Volo, mundáre. Extendénte manum Dómino, statim lepra fugit. Simúlque consídera, quam húmilis, et sine jactántia respónsio. Ille díxerat, Si vis: Dóminus respóndit, Volo. Ille præmíserat, Potes me mundáre: Dóminus jungit, et dicit, Mundáre. Non ergo, ut pleríque Latinórum putant, jungéndum est, et legéndum, Volo mundáre: sed separátim, ut primum dicat, Volo; deínde ímperet, Mundáre.\par
\switchcolumn
\EN
\noindent And, behold, there came a leper, and worshipped Him, saying: Properly after preaching and doctrine cometh occasion for a sign, that the power of the miracle might confirm in the hearers the truth of the teaching that had gone before. Lord, if Thou wilt, Thou canst make me clean. He that prayeth the Lord to have the will, doubteth not but that He hath the power. And Jesus put forth His hand, and touched him, saying: I will; be thou clean. As soon as the Lord put forth His Hand, the leprosy departed. Let us remark how lowly and unbragging is the Lord's language. The leper had said, If Thou wilt; the Lord answereth, I will. The leper, Thou canst make me clean; the Lord, Be thou clean. Most Latin readers, misled by the identity of form in that language between the Present Infinitive Active and the Second Person Singular Present Imperative Passive of the Verb, read Christ's answer as if it were, I will to make thee clean. This is wrong. The sentences are separate. First cometh the expression of volition, I will, then the command, Be thou clean.\par
\switchcolumn*

\LA
\begin{Resp}\Rsign Duo Séraphim clamábant alter ad álterum: \Star{} Sanctus, sanctus, sanctus Dóminus Deus Sábaoth: * Plena est omnis terra glória ejus.\end{Resp}
\begin{Resp}\Vsign Tres sunt qui testimónium dant in cælo: Pater, Verbum, et Spíritus Sanctus: et hi tres unum sunt.\end{Resp}
\begin{Resp}\Rsign Sanctus, sanctus, sanctus Dóminus Deus Sábaoth.\end{Resp}
\begin{Resp}\Vsign Glória Patri, et Fílio, \Star{} et Spirítui Sancto.\end{Resp}
\begin{Resp}\Rsign Plena est omnis terra glória ejus.\end{Resp}
\switchcolumn
\EN
\begin{Resp}\Rsign One Seraph cried unto another: \Star{} Holy, Holy, Holy is the Lord God of hosts; the whole earth is full of His glory.\end{Resp}
\begin{Resp}\Vsign There are Three That bear record in heaven: the Father, the Word, and the Holy Ghost, and these Three are One.\end{Resp}
\begin{Resp}\Rsign Holy, Holy, Holy is the Lord God of hosts.\end{Resp}
\begin{Resp}\Vsign Glory be to the Father, and to the Son, \Star{} and to the Holy Ghost.\end{Resp}
\begin{Resp}\Rsign The whole earth is full of His glory.\end{Resp}
\switchcolumn*

\LA
\LessonHead{Lectio IX}
\switchcolumn
\EN
\LessonHead{Lesson IX}
\switchcolumn*

\LA
\noindent Et ait illi Jesus: Vide, némini díxeris. Et revéra quid erat necésse ut sermóne jactáret, quod córpore præferébat? Sed vade, osténde te sacerdóti. Várias ob causas mittit eum ad sacerdótem: primum propter humilitátem, ut sacerdótibus deférre honórem videátur. Erat enim lege præcéptum, ut, qui mundáti fúerant a lepra, offérrent múnera sacerdótibus. Deínde, ut mundátum vidéntes leprósum, aut créderent Salvatóri, aut non créderent: si créderent, salvaréntur; si non créderent, inexcusábiles forent. Et simul, ne, quod in eo sæpíssime criminabántur, legem viderétur infríngere.\par
\switchcolumn
\EN
\noindent And Jesus saith unto him: See thou tell no man. What need was there to tell what his body showed? But go thy way, show thyself to the Priest. There were diverse reasons why Christ should send him to the Priest. First, for humility's sake, that He might show reverence to God's Priest. Then there was a command in the law that they that were cleansed of leprosy should make an offering to the Priests. Moreover, that, when the Priests saw the leper cleansed, they might either believe in the Saviour, or refuse to believe; if they believed, that they might be saved, and, if they believed not, that they might have no excuse. Lastly, that He might give no ground for the accusation that was so often brought against Him, that He was unobservant of the law.\par
\switchcolumn*

\LA
\TeDeum{Te Deum}
\switchcolumn
\EN
\TeDeum{Te Deum}
\switchcolumn*

\end{paracol}

\Office{Feria II infra Hebdomadam III post Epiphaniam}{Monday of the Third Week after Epiphany}{Feria}
\addcontentsline{toc}{subsection}{Feria II infra Hebdomadam III post Epiphaniam \textit{— Monday of the Third Week after Epiphany}}
\begin{paracol}{2}
\LA
\RefLine{Lectiones I–III: ex Commúni Evangelistarum (C1a).}
\switchcolumn
\EN
\RefLine{Lessons I–III: from the Commune Evangelistarum (C1a).}
\switchcolumn*

\LA
\LessonHead{Lectio I}
\switchcolumn
\EN
\LessonHead{Lesson I}
\switchcolumn*

\LA
\LessonTitle{De Epístola ad Gálatas}
\switchcolumn
\EN
\LessonTitle{Lesson from the letter of St. Paul the Apostle to the Galatians}
\switchcolumn*

\LA
\Cite{Gal 3:1-6}
\switchcolumn
\EN
\Cite{Gal 3:1-6}
\switchcolumn*

\LA
\noindent O insensáti Gálatæ, quis vos fascinávit non obedíre veritáti, ante quorum óculos Jesus Christus præscríptus est, in vobis crucifíxus? Hoc solum a vobis volo díscere: ex opéribus legis Spíritum accepístis, an ex audítu fídei? sic stulti estis, ut cum Spíritu cœpéritis, nunc carne consummémini? tanta passi estis sine causa? si tamen sine causa. Qui ergo tríbuit vobis Spíritum, et operátur virtútes in vobis: ex opéribus legis, an ex audítu fídei? Sicut scriptum est: Abraham crédidit Deo, et reputátum est illi ad justítiam:\par
\switchcolumn
\EN
\noindent O senseless Galatians, who hath bewitched you that you should not obey the truth, before whose eyes Jesus Christ hath been set forth, crucified among you? This only would I learn of you: Did you receive the Spirit by the works of the law, or by the hearing of faith? Are you so foolish, that, whereas you began in the Spirit, you would now be made perfect by the flesh? Have you suffered so great things in vain? If it be yet in vain. He therefore who giveth to you the Spirit, and worketh miracles among you; doth he do it by the works of the law, or by the hearing of the faith? As it is written: Abraham believed God, and it was reputed to him unto justice.\par
\switchcolumn*

\LA
\begin{Resp}\Rsign Quam magna multitúdo dulcédinis tuæ, Dómine, \Star{} Quam abscondísti timéntibus te!\end{Resp}
\begin{Resp}\Vsign Et perfecísti eis qui sperant in te, Dómine, in conspéctu filiórum hóminum.\end{Resp}
\begin{Resp}\Rsign Quam abscondísti timéntibus te!\end{Resp}
\switchcolumn
\EN
\begin{Resp}\Rsign O how great is thy goodness, O Lord \Star{} Which Thou hast laid up for them that fear thee!\end{Resp}
\begin{Resp}\Vsign Which Thou, O Lord, hast wrought for them that trust in thee before the sons of men!\end{Resp}
\begin{Resp}\Rsign Which Thou hast laid up for them that fear thee!\end{Resp}
\switchcolumn*

\LA
\LessonHead{Lectio II}
\switchcolumn
\EN
\LessonHead{Lesson II}
\switchcolumn*

\LA
\Cite{Gal 3:7-10}
\switchcolumn
\EN
\Cite{Gal 3:7-10}
\switchcolumn*

\LA
\noindent Cognóscite ergo quia qui ex fide sunt, ii sunt fílii Abrahæ. Próvidens autem Scriptúra quia ex fide justíficat gentes Deus, prænuntiávit Abrahæ: Quia benedicéntur in te omnes gentes. Igitur qui ex fide sunt, benedicéntur cum fidéli Abraham. Quicúmque enim ex opéribus legis sunt, sub maledícto sunt. Scriptum est enim: Maledíctus omnis qui non permánserit in ómnibus quæ scripta sunt in libro legis ut fáciat ea.\par
\switchcolumn
\EN
\noindent Know ye therefore, that they who are of faith, the same are the children of Abraham. And the scripture, foreseeing, that God justifieth the Gentiles by faith, told unto Abraham before: In thee shall all nations be blessed. Therefore they that are of faith, shall be blessed with faithful Abraham. For as many as are of the works of the law, are under a curse. For it is written: Cursed is every one, that abideth not in all things, which are written in the book of the law to do them.\par
\switchcolumn*

\LA
\begin{Resp}\Rsign Adjútor meus esto, Deus: \Star{} Ne derelínquas me.\end{Resp}
\begin{Resp}\Vsign Neque despícias me, Deus, salutáris meus.\end{Resp}
\begin{Resp}\Rsign Ne derelínquas me.\end{Resp}
\switchcolumn
\EN
\begin{Resp}\Rsign O God, be Thou my helper. \Star{} Neither leave me.\end{Resp}
\begin{Resp}\Vsign Nor forsake me O God of my salvation.\end{Resp}
\begin{Resp}\Rsign Neither leave me.\end{Resp}
\switchcolumn*

\LA
\LessonHead{Lectio III}
\switchcolumn
\EN
\LessonHead{Lesson III}
\switchcolumn*

\LA
\Cite{Gal 3:11-14}
\switchcolumn
\EN
\Cite{Gal 3:11-14}
\switchcolumn*

\LA
\noindent Quóniam autem in lege nemo justificátur apud Deum, maniféstum est: quia justus ex fide vivit. Lex autem non est ex fide, sed: Qui fécerit ea, vivet in illis. Christus nos redémit de maledícto legis, factus pro nobis maledíctum: quia scriptum est: Maledíctus omnis qui pendet in ligno: ut in géntibus benedíctio Abrahæ fíeret in Christo Jesu, ut pollicitatiónem Spíritus accipiámus per fidem.\par
\switchcolumn
\EN
\noindent But that in the law no man is justified with God, it is manifest: because the just man liveth by faith. But the law is not of faith: but, He that doth those things, shall live in them. Christ hath redeemed us from the curse of the law, being made a curse for us: for it is written: Cursed is every one that hangeth on a tree: That the blessing of Abraham might come on the Gentiles through Christ Jesus: that we may receive the promise of the Spirit by faith.\par
\switchcolumn*

\LA
\begin{Resp}\Rsign Benedícam Dóminum in omni témpore: \Star{} Semper laus ejus in ore meo.\end{Resp}
\begin{Resp}\Vsign In Dómino laudábitur ánima mea, áudiant mansuéti, et læténtur.\end{Resp}
\begin{Resp}\Rsign Semper laus ejus in ore meo.\end{Resp}
\begin{Resp}\Vsign Glória Patri, et Fílio, \Star{} et Spirítui Sancto.\end{Resp}
\begin{Resp}\Rsign Semper laus ejus in ore meo.\end{Resp}
\switchcolumn
\EN
\begin{Resp}\Rsign I will bless the Lord at all times. \Star{} His praise shall continually be in my mouth.\end{Resp}
\begin{Resp}\Vsign My soul shall make her boast in the Lord; the humble shall hear thereof and be glad.\end{Resp}
\begin{Resp}\Rsign His praise shall continually be in my mouth.\end{Resp}
\begin{Resp}\Vsign Glory be to the Father, and to the Son, \Star{} and to the Holy Ghost.\end{Resp}
\begin{Resp}\Rsign His praise shall continually be in my mouth.\end{Resp}
\switchcolumn*

\end{paracol}

\Office{Feria III infra Hebdomadam III post Epiphaniam}{Tuesday of the Third Week after Epiphany}{Feria}
\addcontentsline{toc}{subsection}{Feria III infra Hebdomadam III post Epiphaniam \textit{— Tuesday of the Third Week after Epiphany}}
\begin{paracol}{2}
\LA
\RefLine{Lectiones I–III: de Scriptura occurrente.}
\switchcolumn
\EN
\RefLine{Lessons I–III: from the Scripture of the day.}
\switchcolumn*

\LA
\LessonHead{Lectio I}
\switchcolumn
\EN
\LessonHead{Lesson I}
\switchcolumn*

\LA
\LessonTitle{De Epístola ad Gálatas}
\switchcolumn
\EN
\LessonTitle{Lesson from the letter of St. Paul the Apostle to the Galatians}
\switchcolumn*

\LA
\Cite{Gal 5:1-5}
\switchcolumn
\EN
\Cite{Gal 5:1-5}
\switchcolumn*

\LA
\noindent State, et nolíte íterum jugo servitútis continéri. Ecce ego Paulus dico vobis: quóniam si circumcidámini, Christus vobis nihil próderit. Testíficor autem rursus omni hómini circumcidénti se, quóniam débitor est univérsæ legis faciéndæ. Evacuáti estis a Christo, qui in lege justificámini: a grátia excidístis. Nos enim spíritu ex fide, spem justítiæ exspectámus.\par
\switchcolumn
\EN
\noindent Stand fast, and be not held again under the yoke of bondage. Behold, I Paul tell you, that if you be circumcised, Christ shall profit you nothing. And I testify again to every man circumcising himself, that he is a debtor to the whole law. You are made void of Christ, you who are justified in the law: you are fallen from grace. For we in spirit, by faith, wait for the hope of justice.\par
\switchcolumn*

\LA
\begin{Resp}\Rsign Auribus pércipe, Deus, lácrimas meas: ne síleas a me, remítte mihi: \Star{} Quóniam íncola ego sum apud te, et peregrínus.\end{Resp}
\begin{Resp}\Vsign Compláceat tibi, ut erípias me: Dómine, ad adjuvándum me festína.\end{Resp}
\begin{Resp}\Rsign Quóniam íncola ego sum apud te, et peregrínus.\end{Resp}
\switchcolumn
\EN
\begin{Resp}\Rsign O God, give ear unto my tears; hold not thy peace, but, O, spare me! \Star{} For I am a stranger with thee, and a sojourner.\end{Resp}
\begin{Resp}\Vsign Be pleased, O Lord, to deliver me; O Lord, look upon me to help me.\end{Resp}
\begin{Resp}\Rsign For I am a stranger with thee, and a sojourner.\end{Resp}
\switchcolumn*

\LA
\LessonHead{Lectio II}
\switchcolumn
\EN
\LessonHead{Lesson II}
\switchcolumn*

\LA
\Cite{Gal 5:6-10}
\switchcolumn
\EN
\Cite{Gal 5:6-10}
\switchcolumn*

\LA
\noindent Nam in Christo Jesu neque circumcísio áliquid valet, neque præpútium: sed fides, quæ per caritátem operátur. Currebátis bene: quis vos impedívit veritáti non obedíre? persuásio hæc non est ex eo, qui vocat vos. Módicum ferméntum totam massam corrúmpit. Ego confído in vobis in Dómino, quod nihil áliud sapiétis: qui autem contúrbat vos, portábit judícium, quicúmque est ille.\par
\switchcolumn
\EN
\noindent For in Christ Jesus neither circumcision availeth any thing, nor uncircumcision: but faith that worketh by charity. You did run well, who hath hindered you, that you should not obey the truth? This persuasion is not from him that calleth you. A little leaven corrupteth the whole lump. I have confidence in you in the Lord: that you will not be of another mind: but he that troubleth you, shall bear the judgment, whosoever he be.\par
\switchcolumn*

\LA
\begin{Resp}\Rsign Státuit Dóminus supra petram pedes meos, et diréxit gressus meos Deus: \Star{} Et misit in os meum cánticum novum.\end{Resp}
\begin{Resp}\Vsign Exaudívit preces meas: et edúxit me de lacu misériæ.\end{Resp}
\begin{Resp}\Rsign Et misit in os meum cánticum novum.\end{Resp}
\switchcolumn
\EN
\begin{Resp}\Rsign The Lord hath set my feet upon a rock, and ordered my goings. \Star{} And He hath put a new song in my mouth.\end{Resp}
\begin{Resp}\Vsign He heard my cry He brought me up also out of an horrible pit.\end{Resp}
\begin{Resp}\Rsign And He hath put a new song in my mouth.\end{Resp}
\switchcolumn*

\LA
\LessonHead{Lectio III}
\switchcolumn
\EN
\LessonHead{Lesson III}
\switchcolumn*

\LA
\Cite{Gal 5:11-17}
\switchcolumn
\EN
\Cite{Gal 5:11-17}
\switchcolumn*

\LA
\noindent Ego autem, fratres, si circumcisiónem adhuc prǽdico: quid adhuc persecutiónem pátior? ergo evacuátum est scándalum crucis. Utinam et abscindántur qui vos contúrbant. Vos enim in libertátem vocáti estis, fratres: tantum ne libertátem in occasiónem detis carnis, sed per caritátem Spíritus servíte ínvicem. Omnis enim lex in uno sermóne implétur: Díliges próximum tuum sicut teípsum. Quod si ínvicem mordétis, et coméditis: vidéte ne ab ínvicem consumámini. Dico autem: Spíritu ambuláte, et desidéria carnis non perficiétis. Caro enim concupíscit advérsus spíritum, spíritus autem advérsus carnem: hæc enim sibi ínvicem adversántur, ut non quæcúmque vultis, illa faciátis.\par
\switchcolumn
\EN
\noindent And I, brethren, if I yet preach circumcision, why do I yet suffer persecution? Then is the scandal of the cross made void. I would they were even cut off, who trouble you. For you, brethren, have been called unto liberty: only make not liberty an occasion to the flesh, but by charity of the spirit serve one another. For all the law is fulfilled in one word: Thou shalt love thy neighbour as thyself. But if you bite and devour one another; take heed you be not consumed one of another. I say then, walk in the spirit, and you shall not fulfil the lusts of the flesh. For the flesh lusteth against the spirit: and the spirit against the flesh; for these are contrary one to another: so that you do not the things that you would.\par
\switchcolumn*

\LA
\begin{Resp}\Rsign Ego dixi, Dómine, miserére mei: \Star{} Sana ánimam meam, quia peccávi tibi.\end{Resp}
\begin{Resp}\Vsign Ab ómnibus iniquitátibus meis éripe me, Dómine.\end{Resp}
\begin{Resp}\Rsign Sana ánimam meam, quia peccávi tibi.\end{Resp}
\begin{Resp}\Vsign Glória Patri, et Fílio, \Star{} et Spirítui Sancto.\end{Resp}
\begin{Resp}\Rsign Sana ánimam meam, quia peccávi tibi.\end{Resp}
\switchcolumn
\EN
\begin{Resp}\Rsign I said: Lord, be merciful unto me. \Star{} Heal my soul, for I have sinned against thee.\end{Resp}
\begin{Resp}\Vsign Deliver me from all mine iniquities, O Lord.\end{Resp}
\begin{Resp}\Rsign Heal my soul, for I have sinned against thee.\end{Resp}
\begin{Resp}\Vsign Glory be to the Father, and to the Son, \Star{} and to the Holy Ghost.\end{Resp}
\begin{Resp}\Rsign Heal my soul, for I have sinned against thee.\end{Resp}
\switchcolumn*

\end{paracol}

\Office{Feria IV infra Hebdomadam III post Epiphaniam}{Wednesday of the Third Week after Epiphany}{Feria}
\addcontentsline{toc}{subsection}{Feria IV infra Hebdomadam III post Epiphaniam \textit{— Wednesday of the Third Week after Epiphany}}
\begin{paracol}{2}
\LA
\RefLine{Lectiones I–III: de Scriptura occurrente.}
\switchcolumn
\EN
\RefLine{Lessons I–III: from the Scripture of the day.}
\switchcolumn*

\LA
\LessonHead{Lectio I}
\switchcolumn
\EN
\LessonHead{Lesson I}
\switchcolumn*

\LA
\LessonTitle{Incipit Epístola beáti Pauli Apóstoli ad Ephésios}
\switchcolumn
\EN
\LessonTitle{Lesson from the letter of St. Paul the Apostle to the Ephesians}
\switchcolumn*

\LA
\Cite{Eph 1:1-4}
\switchcolumn
\EN
\Cite{Eph 1:1-4}
\switchcolumn*

\LA
\noindent Paulus Apóstolus Jesu Christi per voluntátem Dei, ómnibus sanctis qui sunt Ephesi, et fidélibus in Christo Jesu. Grátia vobis, et pax a Deo Patre nostro, et Dómino Jesu Christo. Benedíctus Deus et Pater Dómini nostri Jesu Christi, qui benedíxit nos in omni benedictióne spirituáli in cæléstibus in Christo, sicut elégit nos in ipso ante mundi constitutiónem, ut essémus sancti et immaculáti in conspéctu ejus in caritáte.\par
\switchcolumn
\EN
\noindent Paul, an apostle of Jesus Christ, by the will of God, to all the saints who are at Ephesus, and to the faithful in Christ Jesus. Grace be to you, and peace from God the Father, and from the Lord Jesus Christ. Blessed by the God and Father of our Lord Jesus Christ, who hath blessed us with spiritual blessings in heavenly places, in Christ: As he chose us in him before the foundation of the world, that we should be holy and unspotted in his sight in charity.\par
\switchcolumn*

\LA
\begin{Resp}\Rsign Ne perdíderis me cum iniquitátibus meis: \Star{} Neque in finem irátus resérves mala mea.\end{Resp}
\begin{Resp}\Vsign Non intres in judícium cum servo tuo, Dómine.\end{Resp}
\begin{Resp}\Rsign Neque in finem irátus resérves mala mea.\end{Resp}
\switchcolumn
\EN
\begin{Resp}\Rsign Cut me not off in the midst of my sins. \Star{} Nor keep thy wrath against me for my latter end.\end{Resp}
\begin{Resp}\Vsign Enter not into judgment with thy servant, O Lord.\end{Resp}
\begin{Resp}\Rsign Nor keep thy wrath against me for my latter end.\end{Resp}
\switchcolumn*

\LA
\LessonHead{Lectio II}
\switchcolumn
\EN
\LessonHead{Lesson II}
\switchcolumn*

\LA
\Cite{Eph 1:5-10}
\switchcolumn
\EN
\Cite{Eph 1:5-10}
\switchcolumn*

\LA
\noindent Qui prædestinávit nos in adoptiónem filiórum per Jesum Christum in ipsum: secúndum propósitum voluntátis suæ, in laudem glóriæ grátiæ suæ, in qua gratificávit nos in dilécto Fílio suo. In quo habémus redemptiónem per sánguinem ejus, remissiónem peccatórum secúndum divítias grátiæ ejus, quæ superabundávit in nobis in omni sapiéntia et prudéntia: ut notum fáceret nobis sacraméntum voluntátis suæ, secúndum beneplácitum ejus, quod propósuit in eo, in dispensatióne plenitúdinis témporum, instauráre ómnia in Christo, quæ in cælis et quæ in terra sunt, in ipso;\par
\switchcolumn
\EN
\noindent Who hath predestinated us unto the adoption of children through Jesus Christ unto himself: according to the purpose of his will: Unto the praise of the glory of his grace, in which he hath graced us in his beloved son. In whom we have redemption through his blood, the remission of sins, according to the riches of his grace, Which hath superabounded in us in all wisdom and prudence, That he might make known unto us the mystery of his will, according to his good pleasure, which he hath purposed in him, In the dispensation of the fulness of times, to re-establish all things in Christ, that are in heaven and on earth, in him.\par
\switchcolumn*

\LA
\begin{Resp}\Rsign Parátum cor meum, Deus, parátum cor meum: \Star{} Cantábo et psalmum dicam Dómino.\end{Resp}
\begin{Resp}\Vsign Exsúrge, glória mea, exsúrge, psaltérium et cíthara, exsúrgam dilúculo.\end{Resp}
\begin{Resp}\Rsign Cantábo et psalmum dicam Dómino.\end{Resp}
\switchcolumn
\EN
\begin{Resp}\Rsign My heart is ready, O God, my heart is ready. \Star{} I will sing and give praise to the Lord.\end{Resp}
\begin{Resp}\Vsign Awake up, my glory, awake, psaltery and harp! I will awake early.\end{Resp}
\begin{Resp}\Rsign I will sing and give praise to the Lord.\end{Resp}
\switchcolumn*

\LA
\LessonHead{Lectio III}
\switchcolumn
\EN
\LessonHead{Lesson III}
\switchcolumn*

\LA
\Cite{Eph 1:11-14}
\switchcolumn
\EN
\Cite{Eph 1:11-14}
\switchcolumn*

\LA
\noindent In quo étiam et nos sorte vocáti sumus prædestináti secúndum propósitum ejus qui operátur ómnia secúndum consílium voluntátis suæ: ut simus in laudem glóriæ ejus nos, qui ante sperávimus in Christo; in quo et vos, cum audissétis verbum veritátis, Evangélium salútis vestræ, in quo et credéntes signáti estis Spíritu promissiónis Sancto, qui est pignus hereditátis nostræ, in redemptiónem acquisitiónis, in laudem glóriæ ipsíus.\par
\switchcolumn
\EN
\noindent In whom we also are called by lot, being predestinated according to the purpose of him who worketh all things according to the counsel of his will. That we may be unto the praise of his glory, we who before hoped Christ: In whom you also, after you had heard the word of truth, (the gospel of your salvation;) in whom also believing, you were signed with the holy Spirit of promise, Who is the pledge of our inheritance, unto the redemption of acquisition, unto the praise of his glory.\par
\switchcolumn*

\LA
\begin{Resp}\Rsign Adjútor meus, tibi psallam, quia, Deus, suscéptor meus es: \Star{} Deus meus, misericórdia mea.\end{Resp}
\begin{Resp}\Vsign Lætábor, et exsultábo in te, psallam nómini tuo, Altíssime.\end{Resp}
\begin{Resp}\Rsign Deus meus, misericórdia mea.\end{Resp}
\begin{Resp}\Vsign Glória Patri, et Fílio, \Star{} et Spirítui Sancto.\end{Resp}
\begin{Resp}\Rsign Deus meus, misericórdia mea.\end{Resp}
\switchcolumn
\EN
\begin{Resp}\Rsign Unto thee, O my Strength, will I sing, for God is my defence, \Star{} The God of my mercy.\end{Resp}
\begin{Resp}\Vsign I will be glad and rejoice in thee, I will sing praise to thy Name, O Thou Most High.\end{Resp}
\begin{Resp}\Rsign The God of my mercy.\end{Resp}
\begin{Resp}\Vsign Glory be to the Father, and to the Son, \Star{} and to the Holy Ghost.\end{Resp}
\begin{Resp}\Rsign The God of my mercy.\end{Resp}
\switchcolumn*

\end{paracol}

\Office{Feria V infra Hebdomadam III post Epiphaniam}{Thursday of the Third Week after Epiphany}{Feria}
\addcontentsline{toc}{subsection}{Feria V infra Hebdomadam III post Epiphaniam \textit{— Thursday of the Third Week after Epiphany}}
\begin{paracol}{2}
\LA
\RefLine{Lectiones I–III: de Scriptura occurrente.}
\switchcolumn
\EN
\RefLine{Lessons I–III: from the Scripture of the day.}
\switchcolumn*

\LA
\LessonHead{Lectio I}
\switchcolumn
\EN
\LessonHead{Lesson I}
\switchcolumn*

\LA
\LessonTitle{De Epístola ad Ephésios}
\switchcolumn
\EN
\LessonTitle{Lesson from the letter of St. Paul the Apostle to the Ephesians}
\switchcolumn*

\LA
\Cite{Eph 4:1-6}
\switchcolumn
\EN
\Cite{Eph 4:1-6}
\switchcolumn*

\LA
\noindent Obsecro ítaque vos ego vinctus in Dómino, ut digne ambulétis vocatióne, qua vocáti estis, cum omni humilitáte, et mansuetúdine, cum patiéntia, supportántes ínvicem in caritáte, sollíciti serváre unitátem Spíritus in vínculo pacis. Unum corpus, et unus Spíritus, sicut vocáti estis in una spe vocatiónis vestræ. Unus Dóminus, una fides, unum baptísma. Unus Deus et Pater ómnium, qui est super omnes, et per ómnia, et in ómnibus nobis.\par
\switchcolumn
\EN
\noindent I therefore, a prisoner in the Lord, beseech you that you walk worthy of the vocation in which you are called, With all humility and mildness, with patience, supporting one another in charity. Careful to keep the unity of the Spirit in the bond of peace. One body and one Spirit; as you are called in one hope of your calling. One Lord, one faith, one baptism. One God and Father of all, who is above all, and through all, and in us all.\par
\switchcolumn*

\LA
\begin{Resp}\Rsign Deus, in te sperávi, Dómine, non confúndar in ætérnum: in justítia tua líbera me, \Star{} Et éripe me.\end{Resp}
\begin{Resp}\Vsign Inclína ad me aurem tuam, et salva me.\end{Resp}
\begin{Resp}\Rsign Et éripe me.\end{Resp}
\switchcolumn
\EN
\begin{Resp}\Rsign In thee, O God, do I put my trust; let me never be put to confusion, O Lord deliver me in thy righteousness, \Star{} And cause me to escape.\end{Resp}
\begin{Resp}\Vsign Incline thine ear unto me, deliver me speedily.\end{Resp}
\begin{Resp}\Rsign And cause me to escape.\end{Resp}
\switchcolumn*

\LA
\LessonHead{Lectio II}
\switchcolumn
\EN
\LessonHead{Lesson II}
\switchcolumn*

\LA
\Cite{Eph 4:7-10}
\switchcolumn
\EN
\Cite{Eph 4:7-10}
\switchcolumn*

\LA
\noindent Unicuíque autem nostrum data est grátia secúndum mensúram donatiónis Christi. Propter quod dicit: Ascéndens in altum, captívam duxit captivitátem: dedit dona homínibus. Quod autem ascéndit, quid est, nisi quia et descéndit primum in inferióres partes terræ? Qui descéndit, ipse est et qui ascéndit super omnes cælos, ut impléret ómnia.\par
\switchcolumn
\EN
\noindent But to every one of us is given grace, according to the measure of the giving of Christ. Wherefore he saith: Ascending on high, he led captivity captive; he gave gifts to men. Now that he ascended, what is it, but because he also descended first into the lower parts of the earth? He that descended is the same also that ascended above all the heavens, that he might fill all things.\par
\switchcolumn*

\LA
\begin{Resp}\Rsign Repleátur os meum laude tua, ut hymnum dicam glóriæ tuæ, tota die magnitúdinem tuam: noli me proícere in témpore senectútis: \Star{} Dum defécerit in me virtus mea, ne derelínquas me.\end{Resp}
\begin{Resp}\Vsign Gaudébunt lábia mea cum cantávero tibi.\end{Resp}
\begin{Resp}\Rsign Dum defécerit in me virtus mea, ne derelínquas me.\end{Resp}
\switchcolumn
\EN
\begin{Resp}\Rsign Let my mouth be filled with thy praise, that I may sing of thy glory, all the day long of thy greatness. Cast me not off in the time of old age; \Star{} Forsake me not when my strength faileth.\end{Resp}
\begin{Resp}\Vsign My lips shall be fain when I sing unto thee.\end{Resp}
\begin{Resp}\Rsign Forsake me not when my strength faileth.\end{Resp}
\switchcolumn*

\LA
\LessonHead{Lectio III}
\switchcolumn
\EN
\LessonHead{Lesson III}
\switchcolumn*

\LA
\Cite{Eph 4:11-15}
\switchcolumn
\EN
\Cite{Eph 4:11-15}
\switchcolumn*

\LA
\noindent Et ipse dedit quosdam quidem apóstolos, quosdam autem prophétas, álios vero evangelístas, álios autem pastóres et doctóres, ad consummatiónem sanctórum in opus ministérii, in ædificatiónem córporis Christi: donec occurrámus omnes in unitátem fídei, et agnitiónis Fílii Dei, in virum perféctum, in mensúram ætátis plenitúdinis Christi: ut jam non simus párvuli fluctuántes, et circumferámur omni vento doctrínæ in nequítia hóminum, in astútia ad circumventiónem erróris. Veritátem autem faciéntes in caritáte, crescámus in illo per ómnia, qui est caput Christus:\par
\switchcolumn
\EN
\noindent And he gave some apostles, and some prophets, and other some evangelists, and other some pastors and doctors, For the perfecting of the saints, for the work of the ministry, for the edifying of the body of Christ: Until we all meet into the unity of faith, and of the knowledge of the Son of God, unto a perfect man, unto the measure of the age of the fulness of Christ; That henceforth we be no more children tossed to and fro, and carried about with every wind of doctrine by the wickedness of men, by cunning craftiness, by which they lie in wait to deceive But doing the truth in charity, we may in all things grow up in him who is the head, even Christ:\par
\switchcolumn*

\LA
\begin{Resp}\Rsign Gaudébunt lábia mea cum cantávero tibi: \Star{} Et ánima mea, quam redemísti, Dómine.\end{Resp}
\begin{Resp}\Vsign Sed et lingua mea meditábitur justítiam tuam, tota die laudem tuam.\end{Resp}
\begin{Resp}\Rsign Et ánima mea, quam redemísti, Dómine.\end{Resp}
\begin{Resp}\Vsign Glória Patri, et Fílio, \Star{} et Spirítui Sancto.\end{Resp}
\begin{Resp}\Rsign Et ánima mea, quam redemísti, Dómine.\end{Resp}
\switchcolumn
\EN
\begin{Resp}\Rsign My lips shall be fain when I sing unto thee; \Star{} And my soul, which Thou, O Lord, hast redeemed.\end{Resp}
\begin{Resp}\Vsign My tongue shall also talk of thy righteousness, all the day long of thy praise.\end{Resp}
\begin{Resp}\Rsign And my soul, which Thou, O Lord, hast redeemed.\end{Resp}
\begin{Resp}\Vsign Glory be to the Father, and to the Son, \Star{} and to the Holy Ghost.\end{Resp}
\begin{Resp}\Rsign And my soul, which Thou, O Lord, hast redeemed.\end{Resp}
\switchcolumn*

\end{paracol}

\Office{Dominica IV Post Epiphaniam}{Fourth Sunday after Epiphany}{Semiduplex}
\addcontentsline{toc}{subsection}{Dominica IV Post Epiphaniam \textit{— Fourth Sunday after Epiphany}}
\begin{paracol}{2}
\LA
\LessonHead{Lectio I}
\switchcolumn
\EN
\LessonHead{Lesson I}
\switchcolumn*

\LA
\LessonTitle{Incipit Epístola beáti Pauli Apóstoli ad Philippénses}
\switchcolumn
\EN
\LessonTitle{Lesson from the letter of St. Paul the Apostle to the Philippians}
\switchcolumn*

\LA
\Cite{Phil 1:1-7}
\switchcolumn
\EN
\Cite{Phil 1:1-7}
\switchcolumn*

\LA
\noindent Paulus et Timótheus, servi Jesu Christi, ómnibus sanctis in Christo Jesu, qui sunt Philíppis, cum epíscopis et diacónibus. Grátia vobis, et pax a Deo Patre nostro, et Dómino Jesu Christo. Grátias ago Deo meo in omni memória vestri, semper in cunctis oratiónibus meis pro ómnibus vobis, cum gáudio deprecatiónem fáciens, super communicatióne vestra in Evangélio Christi a prima die usque nunc. Confídens hoc ipsum, quia qui cœpit in vobis opus bonum, perfíciet usque in diem Christi Jesu: sicut est mihi justum hoc sentíre pro ómnibus vobis: eo quod hábeam vos in corde, et in vínculis meis, et in defensióne, et confirmatióne Evangélii, sócios gáudii mei omnes vos esse.\par
\switchcolumn
\EN
\noindent Paul and Timothy, the servants of Jesus Christ; to all the saints in Christ Jesus, who are at Philippi, with the bishops and deacons. Grace be unto you, and peace from God our Father, and from the Lord Jesus Christ. I give thanks to my God in every remembrance of you, Always in all my prayers making supplication for you all, with joy; For your communication in the gospel of Christ from the first day until now. Being confident of this very thing, that he, who hath begun a good work in you, will perfect it unto the day of Christ Jesus. As it is meet for me to think this for you all, for that I have you in my heart; and that in my bands, and in the defence and confirmation of the gospel, you all are partakers of my joy.\par
\switchcolumn*

\LA
\begin{Resp}\Rsign Dómine, ne in ira tua árguas me, neque in furóre tuo corrípias me: \Star{} Miserére mei, Dómine, quóniam infírmus sum.\end{Resp}
\begin{Resp}\Vsign Timor et tremor venérunt super me, et contexérunt me ténebræ.\end{Resp}
\begin{Resp}\Rsign Miserére mei, Dómine, quóniam infírmus sum.\end{Resp}
\switchcolumn
\EN
\begin{Resp}\Rsign O Lord, rebuke me not in thine anger, neither chasten me in thine hot displeasure. \Star{} Have mercy upon me, O Lord, for I am weak.\end{Resp}
\begin{Resp}\Vsign Fearfulness and trembling are come upon me, and darkness hath overwhelmed me.\end{Resp}
\begin{Resp}\Rsign Have mercy upon me, O Lord, for I am weak.\end{Resp}
\switchcolumn*

\LA
\LessonHead{Lectio II}
\switchcolumn
\EN
\LessonHead{Lesson II}
\switchcolumn*

\LA
\Cite{Phil 1:8-14}
\switchcolumn
\EN
\Cite{Phil 1:8-14}
\switchcolumn*

\LA
\noindent Testis enim mihi est Deus, quómodo cúpiam omnes vos in viscéribus Jesu Christi. Et hoc oro, ut cáritas vestra magis ac magis abúndet in sciéntia, et in omni sensu: ut probétis potióra, ut sitis sincéri, et sine offénsa in diem Christi, repléti fructu justítiæ per Jesum Christum, in glóriam et laudem Dei. Scire autem vos volo fratres, quia quæ circa me sunt, magis ad proféctum venérunt Evangélii: ita ut víncula mea manifésta fíerent in Christo in omni prætório, et in céteris ómnibus, et plures e frátribus in Dómino confidéntes vínculis meis, abundántius audérent sine timóre verbum Dei loqui.\par
\switchcolumn
\EN
\noindent For God is my witness, how I long after you all in the bowels of Jesus Christ. And this I pray, that your charity may more and more abound in knowledge, and in all understanding: That you may approve the better things, that you may be sincere and without offence unto the day of Christ, Filled with the fruit of justice, through Jesus Christ, unto the glory and praise of God. Now, brethren, I desire you should know, that the things which have happened to me, have fallen out rather to the furtherance of the gospel: So that my bands are made manifest in Christ, in all the court, and in all other places; And many of the brethren in the Lord, growing confident by my bands, are much more bold to speak the word of God without fear.\par
\switchcolumn*

\LA
\begin{Resp}\Rsign Deus, qui sedes super thronum, et júdicas æquitátem, esto refúgium páuperum in tribulatióne: \Star{} Quia tu solus labórem et dolórem consíderas.\end{Resp}
\begin{Resp}\Vsign Tibi enim derelíctus est pauper, pupíllo tu eris adjútor.\end{Resp}
\begin{Resp}\Rsign Quia tu solus labórem et dolórem consíderas.\end{Resp}
\switchcolumn
\EN
\begin{Resp}\Rsign O God, Which satest in the throne judging right, be Thou a refuge for the poor, a refuge in times of trouble. \Star{} For Thou alone beholdest mischief and spite.\end{Resp}
\begin{Resp}\Vsign The poor leaveth himself unto thee; Thou wilt be the helper of the fatherless.\end{Resp}
\begin{Resp}\Rsign For Thou alone beholdest mischief and spite.\end{Resp}
\switchcolumn*

\LA
\LessonHead{Lectio III}
\switchcolumn
\EN
\LessonHead{Lesson III}
\switchcolumn*

\LA
\Cite{Phil 1:15-18}
\switchcolumn
\EN
\Cite{Phil 1:15-18}
\switchcolumn*

\LA
\noindent Quidam quidem et propter invídiam et contentiónem: quidam autem et propter bonam voluntátem Christum prǽdicant: quidam ex caritáte, sciéntes quóniam in defensiónem Evangélii pósitus sum. Quidam autem ex contentióne Christum annúntiant non sincére, existimántes pressúram se suscitáre vínculis meis. Quid enim? Dum omni modo sive per occasiónem, sive per veritátem, Christus annuntiétur: et in hoc gáudeo, sed et gaudébo.\par
\switchcolumn
\EN
\noindent Some indeed, even out of envy and contention; but some also for good will preach Christ. Some out of charity, knowing that I am set for the defence of the gospel. And some out of contention preach Christ not sincerely: supposing that they raise affliction to my bands. But what then? So that by all means, whether by occasion, or by truth, Christ be preached: in this also I rejoice, yea, and will rejoice.\par
\switchcolumn*

\LA
\begin{Resp}\Rsign A dextris est mihi Dóminus, ne commóvear: \Star{} Propter hoc dilatátum est cor meum, et exsultávit lingua mea.\end{Resp}
\begin{Resp}\Vsign Dóminus pars hereditátis meæ, et cálicis mei.\end{Resp}
\begin{Resp}\Rsign Propter hoc dilatátum est cor meum, et exsultávit lingua mea.\end{Resp}
\begin{Resp}\Vsign Glória Patri, et Fílio, \Star{} et Spirítui Sancto.\end{Resp}
\begin{Resp}\Rsign Propter hoc dilatátum est cor meum, et exsultávit lingua mea.\end{Resp}
\switchcolumn
\EN
\begin{Resp}\Rsign The Lord is at my right hand, I shall never be moved. \Star{} Therefore my heart is glad, and my tongue rejoiceth.\end{Resp}
\begin{Resp}\Vsign The Lord is the portion of mine inheritance, and of my cup.\end{Resp}
\begin{Resp}\Rsign Therefore my heart is glad, and my tongue rejoiceth.\end{Resp}
\begin{Resp}\Vsign Glory be to the Father, and to the Son, \Star{} and to the Holy Ghost.\end{Resp}
\begin{Resp}\Rsign Therefore my heart is glad, and my tongue rejoiceth.\end{Resp}
\switchcolumn*

\LA
\LessonHead{Lectio IV}
\switchcolumn
\EN
\LessonHead{Lesson IV}
\switchcolumn*

\LA
\LessonTitle{Ex libro Morálium sancti Gregórii Papæ}
\switchcolumn
\EN
\LessonTitle{From the Book of Moral (Reflections on Job) written by Pope St. Gregory (the Great.)}
\switchcolumn*

\LA
\Source{Lib. 4. Cap. 30.}
\switchcolumn
\EN
\Source{iv, 30}
\switchcolumn*

\LA
\noindent Replémus refectiónibus corpus, ne extenuátum defíciat; extenuámus abstinéntia, ne nos replétum premat: vegetámus hoc mótibus, ne situ immobilitátis intéreat; sed cítius hoc collocándo sístimus, ne ipsa sua vegetatióne succúmbat: adjuméntis hoc véstium tégimus, ne frigus intérimat; et quæsíta adjuménta projícimus, ne calor exúrat. Tot ígitur diversitátibus occurréntes, quid ágimus, nisi corruptibilitáti sérvimus, ut saltem multiplícitas impénsi obséquii corpus sustíneat, quod anxíetas infírmæ mutabilitátis gravat?\par
\switchcolumn
\EN
\noindent We refresh the body lest it should grow too weak and fail us; we chasten it by abstinence, lest it should wax gross, and become lord over us; we strengthen it with exercise, lest it perish by the not using; and straightway we give it rest, lest it faint through weariness; we succour it with raiment, lest the cold should blight it; and we strip it of the raiment wherewith we have clothed it, lest the heat should afflict it. In all these so many offices what do we but serve the corruptible? Upon what is all this care spent but upon that wherover hangeth the doom of weakness and change?\par
\switchcolumn*

\LA
\begin{Resp}\Rsign Notas mihi fecísti, Dómine, vias vitæ: \Star{} Adimplébis me lætítia cum vultu tuo: delectatiónes in déxtera tua usque in finem.\end{Resp}
\begin{Resp}\Vsign Tu es qui restítues hereditátem meam mihi.\end{Resp}
\begin{Resp}\Rsign Adimplébis me lætítia cum vultu tuo: delectatiónes in déxtera tua usque in finem.\end{Resp}
\switchcolumn
\EN
\begin{Resp}\Rsign O Lord, Thou hast shown me the path of life. \Star{} Thou shalt fill me with joy in thy presence, at thy right hand there are pleasures for evermore.\end{Resp}
\begin{Resp}\Vsign Thou art He That shalt restore mine inheritance unto me.\end{Resp}
\begin{Resp}\Rsign Thou shalt fill me with joy in thy presence, at thy right hand there are pleasures for evermore.\end{Resp}
\switchcolumn*

\LA
\LessonHead{Lectio V}
\switchcolumn
\EN
\LessonHead{Lesson V}
\switchcolumn*

\LA
\noindent Unde bene per Paulum dícitur: Vanitáti enim subjécta est creatúra non volens, sed propter eum qui subjécit eam in spe: quia et ipsa creatúra liberábitur a servitúte corruptiónis, in libertátem glóriæ filiórum Dei. Vanitáti quippe creatúra non volens súbditur: quia homo, qui ingénitæ constántiæ statum volens deséruit, pressus justæ mortalitátis póndere, nolens mutabilitátis suæ corruptióni servit. Sed creatúra hæc tunc a servitúte corruptiónis erípitur, cum ad filiórum Dei glóriam incorrúpta resurgéndo sublevátur.\par
\switchcolumn
\EN
\noindent Therefore saith Paul tells: For the creature was made subject to vanity, not willingly, but by reason of Him Who hath subjected the same in hope because the creature itself also shall be delivered from the bondage of corruption into the glorious liberty of the children of God. (Rom. viii. 20.) The creature was made subject to vanity, not willingly for when man had of his own free will abdicated his state of unchangeable blessedness, the just sentence of death was passed upon him, and whether he willed or not, he became subject to the state of change and corruption. But the creature itself also shall be delivered from the bondage of corruption when it shall rise again incorruptible and be made partaker of the glory of the children of God.\par
\switchcolumn*

\LA
\begin{Resp}\Rsign Díligam te, Dómine, virtus mea: Dóminus firmaméntum meum, \Star{} Et refúgium meum.\end{Resp}
\begin{Resp}\Vsign Liberátor meus, Deus meus, adjútor meus.\end{Resp}
\begin{Resp}\Rsign Et refúgium meum.\end{Resp}
\switchcolumn
\EN
\begin{Resp}\Rsign I will love thee, O Lord, my strength; the Lord is my rock \Star{} And my fortress.\end{Resp}
\begin{Resp}\Vsign My Deliverer, my God, mine Helper.\end{Resp}
\begin{Resp}\Rsign And my fortress.\end{Resp}
\switchcolumn*

\LA
\LessonHead{Lectio VI}
\switchcolumn
\EN
\LessonHead{Lesson VI}
\switchcolumn*

\LA
\noindent Hic ítaque elécti moléstia vincti sunt, quia adhuc corruptiónis suæ pœna deprimúntur: sed cum corruptíbili carne exúimur, quasi ab his, quibus nunc astríngimur, moléstiæ vínculis relaxámur. Præsentári namque jam Deo cúpimus, sed adhuc mortális córporis obligatióne præpedímur. Jure ergo vincti dícimur, quia adhuc incéssum nostri desidérii ad Deum líberum non habémus. Unde bene Paulus, ætérna desíderans, sed tamen adhuc corruptiónis suæ sárcinam portans, vinctus clamat: Cúpio dissólvi, et esse cum Christo. Dissólvi enim non quǽreret, nisi se proculdúbio vinctum vidéret.\par
\switchcolumn
\EN
\noindent Where, then, the elect are still subject to sorrow, being yet bound by the sentence of corruption; but when we shall have put off this corruptible we shall be loosed from that sentence, and shall sorrow no more. For though we earnestly desire to appear before God, we are still hindered by the burden of this dying body. Rightly then are we called prisoners, since we are not free to go whither we will, that is to say, to God; and rightly did the prisoner Paul, yearning after the things which are eternal, and still weighed down with the burden of this corruptible, rightly did he cry out I have a desire to depart and to be with Christ. (Phil. i. 23.) He would not have felt this keenness if he had not felt himself bound down.\par
\switchcolumn*

\LA
\begin{Resp}\Rsign Dómini est terra, et plenitúdo ejus: \Star{} Orbis terrárum, et univérsi qui hábitant in eo.\end{Resp}
\begin{Resp}\Vsign Ipse super mária fundávit eam, et super flúmina præparávit illam.\end{Resp}
\begin{Resp}\Rsign Orbis terrárum, et univérsi qui hábitant in eo.\end{Resp}
\begin{Resp}\Vsign Glória Patri, et Fílio, \Star{} et Spirítui Sancto.\end{Resp}
\begin{Resp}\Rsign Orbis terrárum, et univérsi qui hábitant in eo.\end{Resp}
\switchcolumn
\EN
\begin{Resp}\Rsign The earth is the Lord's, and the fulness thereof \Star{} The world, and they that dwell therein.\end{Resp}
\begin{Resp}\Vsign For He hath founded it upon the seas, and established it upon the floods.\end{Resp}
\begin{Resp}\Rsign The world, and they that dwell therein.\end{Resp}
\begin{Resp}\Vsign Glory be to the Father, and to the Son, \Star{} and to the Holy Ghost.\end{Resp}
\begin{Resp}\Rsign The world, and they that dwell therein.\end{Resp}
\switchcolumn*

\LA
\LessonHead{Lectio VII}
\switchcolumn
\EN
\LessonHead{Lesson VII}
\switchcolumn*

\LA
\LessonTitle{Léctio sancti Evangélii secúndum Matthǽum}
\switchcolumn
\EN
\LessonTitle{From the Holy Gospel according to Matthew}
\switchcolumn*

\LA
\Cite{Matt 8:23-27}
\switchcolumn
\EN
\Cite{Matt 8:23-27}
\switchcolumn*

\LA
\noindent In illo témpore: Ascendénte Jesu in navículam, secúti sunt eum discípuli ejus: et ecce motus magnus factus est in mari, ita ut navícula operirétur flúctibus: ipse vero dormiébat. Et réliqua.\par
\switchcolumn
\EN
\noindent At that time, when Jesus entered into the boat, his disciples followed him: And behold a great tempest arose in the sea, so that the boat was covered with waves, but he was asleep. And so on.\par
\switchcolumn*

\LA
\LessonTitle{Homilía sancti Hierónymi Presbýteri}
\switchcolumn
\EN
\LessonTitle{Homily by St. Jerome, Priest (at Bethlehem.)}
\switchcolumn*

\LA
\Source{Liber 1 Comment. in cap. 8 Matth.}
\switchcolumn
\EN
\Source{Bk. i Comm. on Matth. viii}
\switchcolumn*

\LA
\noindent Quintum signum fecit, quando ascéndens navem de Caphárnaum, ventis imperávit et mari. Sextum, quando in regióne Gerasenórum dedit potestátem dæmónibus in porcos. Séptimum, quando ingrédiens civitátem suam, paralýticum secúndum curávit in léctulo. Primus enim paralýticus est puer centuriónis.\par
\switchcolumn
\EN
\noindent The fifth sign that He did was when He took ship at Capernaum, and commanded the winds and the sea the sixth, when, in the country of the Gergesenes, He suffered the devils to enter into the swine the seventh, when, as He came into His own city, He cured the man sick of the palsy lying on a bed. The first man sick of the palsy that He cured was the centurion's servant.\par
\switchcolumn*

\LA
\begin{Resp}\Rsign Ad te, Dómine, levávi ánimam meam: \Star{} Deus meus, in te confído, non erubéscam.\end{Resp}
\begin{Resp}\Vsign Custódi ánimam meam, et éripe me.\end{Resp}
\begin{Resp}\Rsign Deus meus, in te confído, non erubéscam.\end{Resp}
\switchcolumn
\EN
\begin{Resp}\Rsign Unto thee, O Lord, do I lift up my soul. \Star{} O my God, I trust in thee, let me not be ashamed.\end{Resp}
\begin{Resp}\Vsign O keep my soul and deliver me.\end{Resp}
\begin{Resp}\Rsign O my God, I trust in thee, let me not be ashamed.\end{Resp}
\switchcolumn*

\LA
\LessonHead{Lectio VIII}
\switchcolumn
\EN
\LessonHead{Lesson VIII}
\switchcolumn*

\LA
\noindent Ipse vero dormiébat: et accessérunt ad eum, et suscitavérunt eum, dicéntes: Dómine, salva nos. Hujus signi typum in Jona légimus, quando céteris periclitántibus, ipse secúrus est, et dormit, et suscitátur; et império ac sacraménto passiónis suæ líberat suscitántes. Tunc surgens imperávit ventis et mari. Ex hoc loco intellígimus, quod omnes creatúræ séntiant Creatórem. Quas enim increpávit, et quibus imperávit, sentiunt imperántem: non erróre hæreticórum qui ómnia putant animántia, sed majestáte Conditóris, quæ apud nos insensibília, illi sensibília sunt.\par
\switchcolumn
\EN
\noindent But He was asleep; and His disciples came to Him, and awoke Him, saying Lord, save us. There is a type of this in the history of Jonah, who, when the storm arose, was lying fast asleep, and whom the sailors woke to help them; who also saved the sailors by commanding them to throw him into the sea, the said casting of him into the sea, being, as we know, a figure of Christ's Passion. Then He arose and rebuked the winds and the sea. From these words we understand that all things, which have been made, are sentient to their Maker. All things which He rebuketh or commandeth, hear His voice. This is not the error of the heretics who will have it that everything is quick, but part of the majesty of the Creator, Who maketh to feel Him things which we cannot make to feel us.\par
\switchcolumn*

\LA
\begin{Resp}\Rsign Duo Séraphim clamábant alter ad álterum: \Star{} Sanctus, sanctus, sanctus Dóminus Deus Sábaoth: * Plena est omnis terra glória ejus.\end{Resp}
\begin{Resp}\Vsign Tres sunt qui testimónium dant in cælo: Pater, Verbum, et Spíritus Sanctus: et hi tres unum sunt.\end{Resp}
\begin{Resp}\Rsign Sanctus, sanctus, sanctus Dóminus Deus Sábaoth.\end{Resp}
\begin{Resp}\Vsign Glória Patri, et Fílio, \Star{} et Spirítui Sancto.\end{Resp}
\begin{Resp}\Rsign Plena est omnis terra glória ejus.\end{Resp}
\switchcolumn
\EN
\begin{Resp}\Rsign One Seraph cried unto another: \Star{} Holy, Holy, Holy is the Lord God of hosts; the whole earth is full of His glory.\end{Resp}
\begin{Resp}\Vsign There are Three That bear record in heaven: the Father, the Word, and the Holy Ghost, and these Three are One.\end{Resp}
\begin{Resp}\Rsign Holy, Holy, Holy is the Lord God of hosts.\end{Resp}
\begin{Resp}\Vsign Glory be to the Father, and to the Son, \Star{} and to the Holy Ghost.\end{Resp}
\begin{Resp}\Rsign The whole earth is full of His glory.\end{Resp}
\switchcolumn*

\LA
\LessonHead{Lectio IX}
\switchcolumn
\EN
\LessonHead{Lesson IX}
\switchcolumn*

\LA
\noindent Porro hómines miráti sunt, dicéntes: Qualis est hic, quia venti et mare obédiunt ei? Non discípuli, sed nautæ, et céteri, qui in navi erant, mirabántur. Sin autem quis contentióse volúerit, eos, qui mirabántur, fuísse discípulos: respondébimus, recte hómines appellátos, qui necdum nóverant poténtiam Salvatóris.\par
\switchcolumn
\EN
\noindent But the men marvelled, saying What manner of man is this, that even the winds and the sea obey Him? It was not His disciples that marvelled, but the sailors, and the others that were in the ship. If however, any one willeth to withstand this our interpretation and to maintain that it was the disciples who marvelled, we are ready to answer them that they who knew not before the power of the Saviour deserve to be stripped of the title of disciples, and to be called simply the men.\par
\switchcolumn*

\LA
\TeDeum{Te Deum}
\switchcolumn
\EN
\TeDeum{Te Deum}
\switchcolumn*

\end{paracol}

\Office{Dominica V Post Epiphaniam}{Fifth Sunday after Epiphany}{Semiduplex}
\addcontentsline{toc}{subsection}{Dominica V Post Epiphaniam \textit{— Fifth Sunday after Epiphany}}
\begin{paracol}{2}
\LA
\RefLine{Lectiones I–VI: de Scriptura occurrente.}
\switchcolumn
\EN
\RefLine{Lessons I–VI: from the Scripture of the day.}
\switchcolumn*

\LA
\LessonHead{Lectio I}
\switchcolumn
\EN
\LessonHead{Lesson I}
\switchcolumn*

\LA
\LessonTitle{Incipit Epístola prima beáti Pauli Apóstoli ad Timótheum}
\switchcolumn
\EN
\LessonTitle{Lesson from the first letter of St. Paul the Apostle to Timothy}
\switchcolumn*

\LA
\Cite{1 Tim 1:1-4}
\switchcolumn
\EN
\Cite{1 Tim 1:1-4}
\switchcolumn*

\LA
\noindent Paulus Apóstolus Jesu Christi secúndum impérium Dei Salvatóris nostri, et Christi Jesu spei nostræ: Timótheo dilécto fílio in fide. Grátia, misericórdia, et pax a Deo Patre, et Christo Jesu Dómino nostro. Sicut rogávi te ut remanéres Ephesi, cum irem in Macedóniam, ut denuntiáres quibúsdam ne áliter docérent, neque inténderent fábulis, et genealogíis interminátis: quæ quæstiónes præstant magis quam ædificatiónem Dei, quæ est in fide.\par
\switchcolumn
\EN
\noindent Paul, an apostle of Jesus Christ, according to the commandment of God our Saviour, and of Christ Jesus our hope: To Timothy, his beloved son in faith. Grace, mercy, and peace from God the Father, and from Christ Jesus our Lord. As I desired thee to remain at Ephesus when I went into Macedonia, that thou mightest charge some not to teach otherwise, Not to give heed to fables and endless genealogies: which furnish questions rather than the edification of God, which is in faith.\par
\switchcolumn*

\LA
\begin{Resp}\Rsign Dómine, ne in ira tua árguas me, neque in furóre tuo corrípias me: \Star{} Miserére mei, Dómine, quóniam infírmus sum.\end{Resp}
\begin{Resp}\Vsign Timor et tremor venérunt super me, et contexérunt me ténebræ.\end{Resp}
\begin{Resp}\Rsign Miserére mei, Dómine, quóniam infírmus sum.\end{Resp}
\switchcolumn
\EN
\begin{Resp}\Rsign O Lord, rebuke me not in thine anger, neither chasten me in thine hot displeasure. \Star{} Have mercy upon me, O Lord, for I am weak.\end{Resp}
\begin{Resp}\Vsign Fearfulness and trembling are come upon me, and darkness hath overwhelmed me.\end{Resp}
\begin{Resp}\Rsign Have mercy upon me, O Lord, for I am weak.\end{Resp}
\switchcolumn*

\LA
\LessonHead{Lectio II}
\switchcolumn
\EN
\LessonHead{Lesson II}
\switchcolumn*

\LA
\Cite{1 Tim 1:5-11}
\switchcolumn
\EN
\Cite{1 Tim 1:5-11}
\switchcolumn*

\LA
\noindent Finis autem præcépti est cáritas de corde puro, et consciéntia bona, et fide non ficta. A quibus quidam aberrántes, convérsi sunt in vanilóquium, voléntes esse legis doctóres, non intellegéntes neque quæ loquúntur, neque de quibus affírmant. Scimus autem quia bona est lex, si quis ea legítime utátur: sciens hoc quia lex justo non est pósita, sed injústis, et non súbditis, ímpiis et peccatóribus, scelerátis, et contaminátis, parricídis et matricídis, homicídis, fornicáriis, masculórum concubitóribus, plagiáriis, mendácibus et perjúris, et si quid áliud sanæ doctrínæ adversátur, quæ est secúndum Evangélium glóriæ beáti Dei, quod créditum est mihi.\par
\switchcolumn
\EN
\noindent Now the end of the commandment is charity, from a pure heart, and a good conscience, and an unfeigned faith. From which things some going astray, are turned aside unto vain babbling: Desiring to be teachers of the law, understanding neither the things they say, nor whereof they affirm. But we know that the law is good, if a man use it lawfully: Knowing this, that the law is not made for the just man, but for the unjust and disobedient, for the ungodly, and for sinners, for the wicked and defiled, for murderers of fathers, and murderers of mothers, for manslayers, For fornicators, for them who defile themselves with mankind, for menstealers, for liars, for perjured persons, and whatever other thing is contrary to sound doctrine, Which is according to the gospel of the glory of the blessed God, which hath been committed to my trust.\par
\switchcolumn*

\LA
\begin{Resp}\Rsign Deus, qui sedes super thronum, et júdicas æquitátem, esto refúgium páuperum in tribulatióne: \Star{} Quia tu solus labórem et dolórem consíderas.\end{Resp}
\begin{Resp}\Vsign Tibi enim derelíctus est pauper, pupíllo tu eris adjútor.\end{Resp}
\begin{Resp}\Rsign Quia tu solus labórem et dolórem consíderas.\end{Resp}
\switchcolumn
\EN
\begin{Resp}\Rsign O God, Which satest in the throne judging right, be Thou a refuge for the poor, a refuge in times of trouble. \Star{} For Thou alone beholdest mischief and spite.\end{Resp}
\begin{Resp}\Vsign The poor leaveth himself unto thee; Thou wilt be the helper of the fatherless.\end{Resp}
\begin{Resp}\Rsign For Thou alone beholdest mischief and spite.\end{Resp}
\switchcolumn*

\LA
\LessonHead{Lectio III}
\switchcolumn
\EN
\LessonHead{Lesson III}
\switchcolumn*

\LA
\Cite{1 Tim 1:12-16}
\switchcolumn
\EN
\Cite{1 Tim 1:12-16}
\switchcolumn*

\LA
\noindent Grátias ago ei, qui me confortávit Christo Jesu Dómino nostro, quia fidélem me existimávit, ponens in ministério: qui prius blasphémus fui, et persecútor, et contumeliósus: sed misericórdiam Dei consecútus sum, quia ignórans feci in incredulitáte. Superabundávit autem grátia Dómini nostri cum fide et dilectióne, quæ est in Christo Jesu. Fidélis sermo, et omni acceptióne dignus: quod Christus Jesus venit in hunc mundum peccatóres salvos fácere, quorum primus ego sum. Sed ídeo misericórdiam consecútus sum: ut in me primo osténderet Christus Jesus omnem patiéntiam ad informatiónem eórum, qui creditúri sunt illi, in vitam ætérnam.\par
\switchcolumn
\EN
\noindent I give thanks to him who hath strengthened me, even to Christ Jesus our Lord, for that he hath counted me faithful, putting me in the ministry; Who before was a blasphemer, and a persecutor, and contumelious. But I obtained the mercy of God, because I did it ignorantly in unbelief. Now the grace of our Lord hath abounded exceedingly with faith and love, which is in Christ Jesus. A faithful saying, and worthy of all acceptation, that Christ Jesus came into this world to save sinners, of whom I am the chief. But for this cause have I obtained mercy: that in me first Christ Jesus might shew forth all patience, for the information of them that shall believe in him unto life everlasting.\par
\switchcolumn*

\LA
\begin{Resp}\Rsign A dextris est mihi Dóminus, ne commóvear: \Star{} Propter hoc dilatátum est cor meum, et exsultávit lingua mea.\end{Resp}
\begin{Resp}\Vsign Dóminus pars hereditátis meæ, et cálicis mei.\end{Resp}
\begin{Resp}\Rsign Propter hoc dilatátum est cor meum, et exsultávit lingua mea.\end{Resp}
\begin{Resp}\Vsign Glória Patri, et Fílio, \Star{} et Spirítui Sancto.\end{Resp}
\begin{Resp}\Rsign Propter hoc dilatátum est cor meum, et exsultávit lingua mea.\end{Resp}
\switchcolumn
\EN
\begin{Resp}\Rsign The Lord is at my right hand, I shall never be moved. \Star{} Therefore my heart is glad, and my tongue rejoiceth.\end{Resp}
\begin{Resp}\Vsign The Lord is the portion of mine inheritance, and of my cup.\end{Resp}
\begin{Resp}\Rsign Therefore my heart is glad, and my tongue rejoiceth.\end{Resp}
\begin{Resp}\Vsign Glory be to the Father, and to the Son, \Star{} and to the Holy Ghost.\end{Resp}
\begin{Resp}\Rsign Therefore my heart is glad, and my tongue rejoiceth.\end{Resp}
\switchcolumn*

\LA
\LessonHead{Lectio IV}
\switchcolumn
\EN
\LessonHead{Lesson IV}
\switchcolumn*

\LA
\LessonTitle{Sermo sancti Augustíni Epíscopi}
\switchcolumn
\EN
\LessonTitle{From the Sermons of St. Augustine, Bishop (of Hippo.)}
\switchcolumn*

\LA
\Source{De Verbis Apóstoli, Sermo 8 sub initio}
\switchcolumn
\EN
\Source{On the words of the Apostles, 8}
\switchcolumn*

\LA
\noindent Humánus sermo, et omni acceptióne dignus, quia Christus Jesus venit in hunc mundum peccatóres salvos fácere. Atténde Evangélium: Venit enim Fílius hóminis quærére, et salváre, quod períerat. Si homo non periísset, Fílius hóminis non venísset. Ergo períerat homo: venit Deus homo, et invéntus est homo. Períerat homo per líberam voluntátem: venit Deus homo per grátiam liberatrícem.\par
\switchcolumn
\EN
\noindent This is a saying made for man, and worthy of all acceptation, that Christ Jesus came into this world to save sinners. Listen to the words of the Gospel The Son of man is come to seek, and to save that which was lost. If man had not been lost, the Son of man would not have come. Wherefore, man had been lost; God came made Man, and man was found; man had perished by his own free will God made Man came by grace which setteth free.\par
\switchcolumn*

\LA
\begin{Resp}\Rsign Notas mihi fecísti, Dómine, vias vitæ: \Star{} Adimplébis me lætítia cum vultu tuo: delectatiónes in déxtera tua usque in finem.\end{Resp}
\begin{Resp}\Vsign Tu es qui restítues hereditátem meam mihi.\end{Resp}
\begin{Resp}\Rsign Adimplébis me lætítia cum vultu tuo: delectatiónes in déxtera tua usque in finem.\end{Resp}
\switchcolumn
\EN
\begin{Resp}\Rsign O Lord, Thou hast shown me the path of life. \Star{} Thou shalt fill me with joy in thy presence, at thy right hand there are pleasures for evermore.\end{Resp}
\begin{Resp}\Vsign Thou art He That shalt restore mine inheritance unto me.\end{Resp}
\begin{Resp}\Rsign Thou shalt fill me with joy in thy presence, at thy right hand there are pleasures for evermore.\end{Resp}
\switchcolumn*

\LA
\LessonHead{Lectio V}
\switchcolumn
\EN
\LessonHead{Lesson V}
\switchcolumn*

\LA
\noindent Quæris, quid váleat ad malum líberum arbítrium? Recóle hóminem peccántem. Quæris, quid váleat ad auxílium Deus et homo? Atténde in eo grátiam liberántem. Nusquam pótuit sic osténdi, quantum váleat volúntas hóminis usurpáta per supérbiam, ad uténdum sine adjutório Dei: malum non pótuit plus, et maniféstius éxprimi, quam in hómine primo. Ecce perit primus homo, et ubi esset, nisi venísset secúndus homo? quia et ille homo, ídeo et iste homo; et ídeo humánus sermo.\par
\switchcolumn
\EN
\noindent Dost thou ask how free-will availeth to evil? Call to mind a sinner: Dost thou ask what God made Man availeth to help? Consider in Him the grace which setteth free. There is no example which so showeth what availeth the free will of man, when it is taken possession of by pride, to use it without God's help, of evil is there no greater and plainer example, than the first man. The first man fell and where had he been if the second Man had not come? As the first was man, so was the second Man, and therefore is this saying a saying. made for man.\par
\switchcolumn*

\LA
\begin{Resp}\Rsign Díligam te, Dómine, virtus mea: Dóminus firmaméntum meum, \Star{} Et refúgium meum.\end{Resp}
\begin{Resp}\Vsign Liberátor meus, Deus meus, adjútor meus.\end{Resp}
\begin{Resp}\Rsign Et refúgium meum.\end{Resp}
\switchcolumn
\EN
\begin{Resp}\Rsign I will love thee, O Lord, my strength; the Lord is my rock \Star{} And my fortress.\end{Resp}
\begin{Resp}\Vsign My Deliverer, my God, mine Helper.\end{Resp}
\begin{Resp}\Rsign And my fortress.\end{Resp}
\switchcolumn*

\LA
\LessonHead{Lectio VI}
\switchcolumn
\EN
\LessonHead{Lesson VI}
\switchcolumn*

\LA
\noindent Prorsus nusquam sic appáret benígnitas grátiæ, et liberálitas omnipoténtiæ Dei, quam in hómine mediatóre Dei et hóminum, hómine Christo Jesu. Quid enim dícimus, fratres mei? In fide cathólica nutrítis loquor, vel in pacem cathólicam lucrátis. Nóvimus et tenémus, mediatórem Dei et hóminum, hóminem Christum Jesum, in quantum homo erat, ejus esse natúræ, cujus et nos sumus. Non enim altérius natúræ caro nostra, et caro illíus: nec altérius natúræ ánima nostra, et ánima illíus. Hanc suscépit natúram, quam salvándam esse judicávit.\par
\switchcolumn
\EN
\noindent Neither is there any example which so showeth what availeth the tenderness of the grace and the abundance of the All-might of God, as the Man That is the Mediator between God and men, the Man Christ Jesus. For what do we say, my brethren? I speak to them that have been bred up in the Catholic Church, or who have been reconciled to that Church. We know and hold that the Mediator between God and men, the Man Christ Jesus, as touching upon His Manhood, is of the same nature as we. For our flesh is not of one nature, and His Flesh of another nature, neither our soul of one nature and His Soul of another nature. He took upon Himself the same nature which He had freely ordained to save.\par
\switchcolumn*

\LA
\begin{Resp}\Rsign Dómini est terra, et plenitúdo ejus: \Star{} Orbis terrárum, et univérsi qui hábitant in eo.\end{Resp}
\begin{Resp}\Vsign Ipse super mária fundávit eam, et super flúmina præparávit illam.\end{Resp}
\begin{Resp}\Rsign Orbis terrárum, et univérsi qui hábitant in eo.\end{Resp}
\begin{Resp}\Vsign Glória Patri, et Fílio, \Star{} et Spirítui Sancto.\end{Resp}
\begin{Resp}\Rsign Orbis terrárum, et univérsi qui hábitant in eo.\end{Resp}
\switchcolumn
\EN
\begin{Resp}\Rsign The earth is the Lord's, and the fulness thereof \Star{} The world, and they that dwell therein.\end{Resp}
\begin{Resp}\Vsign For He hath founded it upon the seas, and established it upon the floods.\end{Resp}
\begin{Resp}\Rsign The world, and they that dwell therein.\end{Resp}
\begin{Resp}\Vsign Glory be to the Father, and to the Son, \Star{} and to the Holy Ghost.\end{Resp}
\begin{Resp}\Rsign The world, and they that dwell therein.\end{Resp}
\switchcolumn*

\LA
\LessonHead{Lectio VII}
\switchcolumn
\EN
\LessonHead{Lesson VII}
\switchcolumn*

\LA
\LessonTitle{Léctio sancti Evangélii secúndum Matthǽum}
\switchcolumn
\EN
\LessonTitle{From the Holy Gospel according to Matthew}
\switchcolumn*

\LA
\Cite{Matt 13:24-30}
\switchcolumn
\EN
\Cite{Matt 13:24-30}
\switchcolumn*

\LA
\noindent In illo témpore: Dixit Jesus turbis parábolam hanc: Símile factum est regnum cælórum hómini, qui seminávit bonum semen in agro suo. Et réliqua.\par
\switchcolumn
\EN
\noindent At that time, Jesus said to the people a parable: The kingdom of heaven is likened to a man that sowed good seeds in his field. And so on.\par
\switchcolumn*

\LA
\LessonTitle{Homilía S. Augustíni Epíscopi}
\switchcolumn
\EN
\LessonTitle{Homily by St. Augustine, Bishop (of Hippo.)}
\switchcolumn*

\LA
\Source{Liber Quæst. Evang. in Matth. cap. 11, tom. 4}
\switchcolumn
\EN
\Source{Quest. Evang. Matth. xi, Bk. 4}
\switchcolumn*

\LA
\noindent Cum negligéntius ágerent præpósiti Ecclésiæ, aut cum dormitiónem mortis accíperent Apóstoli, venit diábolus, et superseminávit eos, quos malos fílios Dóminus interpretátur. Sed quǽritur: utrum hærétici sint, an male vivéntes cathólici? Possunt enim dici fílii mali étiam hærétici, quia ex eódem Evangélii sémine, et Christi nómine procreáti, pravis opiniónibus ad falsa dógmata convertúntur.\par
\switchcolumn
\EN
\noindent When the Shepherds of the Church wax careless, and since the Apostles sleep the sleep of death, cometh the devil, and soweth them whom the Lord calleth a seed of evil-doers. Now, are these seed of evil-doers the heretics, or Catholics of bad lives? It is possible to call even the heretics a seed of evil-doers because they have sprung up from the seed of the Gospel, and been begotten in the Name of Christ, though afterwards they have turned after crooked ways and lying doctrines.\par
\switchcolumn*

\LA
\begin{Resp}\Rsign Ad te, Dómine, levávi ánimam meam: \Star{} Deus meus, in te confído, non erubéscam.\end{Resp}
\begin{Resp}\Vsign Custódi ánimam meam, et éripe me.\end{Resp}
\begin{Resp}\Rsign Deus meus, in te confído, non erubéscam.\end{Resp}
\switchcolumn
\EN
\begin{Resp}\Rsign Unto thee, O Lord, do I lift up my soul. \Star{} O my God, I trust in thee, let me not be ashamed.\end{Resp}
\begin{Resp}\Vsign O keep my soul and deliver me.\end{Resp}
\begin{Resp}\Rsign O my God, I trust in thee, let me not be ashamed.\end{Resp}
\switchcolumn*

\LA
\LessonHead{Lectio VIII}
\switchcolumn
\EN
\LessonHead{Lesson VIII}
\switchcolumn*

\LA
\noindent Sed quod dicit eos in médio trítici seminátos, quasi vidéntur illi significári, qui uníus communiónis sunt. Verúmtamen quóniam Dóminus agrum ipsum, non Ecclésiam, sed hunc mundum interpretátus est: bene intelligúntur hærétici, quia non societáte uníus Ecclésiæ, vel uníus fídei, sed societáte solíus nóminis christiáni in hoc mundo permiscéntur bonis. At illi, qui in eádem fide mali sunt, pálea pótius quam zizánia reputántur: quia pálea étiam fundaméntum ipsum habet cum fruménto, radicémque commúnem.\par
\switchcolumn
\EN
\noindent But whereas it is written that they were sown in the midst of the wheat, we ought haply to understand that they are of one communion with the righteous. Nevertheless, forasmuch as the Lord saith, The field is the world, (and not, the Church,) we may well understand that the seed of evil doers are the heretics, since in this world they are mingled together with the good, not in one common Communion, but only under one common name of Christian. But they which are of one faith with the good seed, and yet are themselves worthless, may more fitly be likened to straw than to tares, since the straw springeth from one soil and one root with the good ear.\par
\switchcolumn*

\LA
\begin{Resp}\Rsign Duo Séraphim clamábant alter ad álterum: \Star{} Sanctus, sanctus, sanctus Dóminus Deus Sábaoth: * Plena est omnis terra glória ejus.\end{Resp}
\begin{Resp}\Vsign Tres sunt qui testimónium dant in cælo: Pater, Verbum, et Spíritus Sanctus: et hi tres unum sunt.\end{Resp}
\begin{Resp}\Rsign Sanctus, sanctus, sanctus Dóminus Deus Sábaoth.\end{Resp}
\begin{Resp}\Vsign Glória Patri, et Fílio, \Star{} et Spirítui Sancto.\end{Resp}
\begin{Resp}\Rsign Plena est omnis terra glória ejus.\end{Resp}
\switchcolumn
\EN
\begin{Resp}\Rsign One Seraph cried unto another: \Star{} Holy, Holy, Holy is the Lord God of hosts; the whole earth is full of His glory.\end{Resp}
\begin{Resp}\Vsign There are Three That bear record in heaven: the Father, the Word, and the Holy Ghost, and these Three are One.\end{Resp}
\begin{Resp}\Rsign Holy, Holy, Holy is the Lord God of hosts.\end{Resp}
\begin{Resp}\Vsign Glory be to the Father, and to the Son, \Star{} and to the Holy Ghost.\end{Resp}
\begin{Resp}\Rsign The whole earth is full of His glory.\end{Resp}
\switchcolumn*

\LA
\LessonHead{Lectio IX}
\switchcolumn
\EN
\LessonHead{Lesson IX}
\switchcolumn*

\LA
\noindent In illa plane sagéna, qua concludúntur et mali et boni pisces, non absúrde mali cathólici intelligúntur. Aliud est enim mare, quod magis mundum istum signíficat: áliud sagéna, quæ uníus fídei, vel uníus Ecclésiæ communiónem vidétur osténdere. Inter hæréticos et malos cathólicos hoc ínterest, quod hærétici falsa credunt: illi autem, vera credéntes, non vivunt ita ut credunt.\par
\switchcolumn
\EN
\noindent However, as touching the net cast into the sea, and enclosing a great multitude of fishes, both bad and good, we may well understand that by the bad are meant Catholics of bad lives. For the sea is one thing whereby we may understand to be signified the world; and the net another, which seemeth to signify our faith, or the Communion of one Church. Between heretics and sinful Catholics there is this difference, that heretics believe a lie, and sinful Catholics believe the truth, but live not as they believe.\par
\switchcolumn*

\LA
\TeDeum{Te Deum}
\switchcolumn
\EN
\TeDeum{Te Deum}
\switchcolumn*

\end{paracol}

\Office{Dominica VI Post Epiphaniam}{Sixth Sunday after Epiphany}{Semiduplex}
\addcontentsline{toc}{subsection}{Dominica VI Post Epiphaniam \textit{— Sixth Sunday after Epiphany}}
\begin{paracol}{2}
\LA
\RefLine{Lectiones I–VI: de Scriptura occurrente.}
\switchcolumn
\EN
\RefLine{Lessons I–VI: from the Scripture of the day.}
\switchcolumn*

\LA
\LessonHead{Lectio I}
\switchcolumn
\EN
\LessonHead{Lesson I}
\switchcolumn*

\LA
\LessonTitle{Incipit Epístola beáti Pauli Apóstoli ad Hebrǽos}
\switchcolumn
\EN
\LessonTitle{Lesson from the letter of St. Paul the Apostle to the Hebrews}
\switchcolumn*

\LA
\Cite{Heb 1:1-4}
\switchcolumn
\EN
\Cite{Heb 1:1-4}
\switchcolumn*

\LA
\noindent Multifáriam, multísque modis olim Deus loquens pátribus in prophétis: novíssime, diébus istis locútus est nobis in Fílio, quem constítuit herédem universórum, per quem fecit et sǽcula: qui cum sit splendor glóriæ, et figúra substántiæ ejus, portánsque ómnia verbo virtútis suæ, purgatiónem peccatórum fáciens, sedet ad déxteram majestátis in excélsis: tanto mélior ángelis efféctus, quanto differéntius præ illis nomen hereditávit.\par
\switchcolumn
\EN
\noindent God, who, at sundry times and in diverse manners, spoke in times past to the fathers by the prophets, last of all, In these days hath spoken to us by his Son, whom he hath appointed heir of all things, by whom also he made the world. Who being the brightness of his glory, and the figure of his substance, and upholding all things by the word of his power, making purgation of sins, sitteth on the right hand of the majesty on high. Being made so much better than the angels, as he hath inherited a more excellent name than they.\par
\switchcolumn*

\LA
\begin{Resp}\Rsign Dómine, ne in ira tua árguas me, neque in furóre tuo corrípias me: \Star{} Miserére mei, Dómine, quóniam infírmus sum.\end{Resp}
\begin{Resp}\Vsign Timor et tremor venérunt super me, et contexérunt me ténebræ.\end{Resp}
\begin{Resp}\Rsign Miserére mei, Dómine, quóniam infírmus sum.\end{Resp}
\switchcolumn
\EN
\begin{Resp}\Rsign O Lord, rebuke me not in thine anger, neither chasten me in thine hot displeasure. \Star{} Have mercy upon me, O Lord, for I am weak.\end{Resp}
\begin{Resp}\Vsign Fearfulness and trembling are come upon me, and darkness hath overwhelmed me.\end{Resp}
\begin{Resp}\Rsign Have mercy upon me, O Lord, for I am weak.\end{Resp}
\switchcolumn*

\LA
\LessonHead{Lectio II}
\switchcolumn
\EN
\LessonHead{Lesson II}
\switchcolumn*

\LA
\Cite{Heb 1:5-9}
\switchcolumn
\EN
\Cite{Heb 1:5-9}
\switchcolumn*

\LA
\noindent Cui enim dixit aliquándo angelórum: Fílius meus es tu, ego hódie génui te? Et rursum: Ego ero illi in patrem, et ipse erit mihi in fílium? Et cum íterum introdúcit primogénitum in orbem terræ, dicit: Et adórent eum omnes ángeli Dei. Et ad ángelos quidem dicit: Qui facit ángelos suos spíritus, et minístros suos flammam ignis. Ad Fílium autem: Thronus tuus Deus in sǽculum sǽculi: virga æquitátis, virga regni tui. Dilexísti justítiam, et odísti iniquitátem: proptérea unxit te Deus, Deus tuus, óleo exultatiónis præ particípibus tuis.\par
\switchcolumn
\EN
\noindent For to which of the angels hath he said at any time, Thou art my Son, to day have I begotten thee? And again, I will be to him a Father, and he shall be to me a Son? And again, when he bringeth in the first begotten into the world, he saith: And let all the angels of God adore him. And to the angels indeed he saith: He that maketh his angels spirits, and his ministers a flame of fire. But to the Son: thy throne, O God, is for ever and ever: a sceptre of justice is the sceptre of thy kingdom. Thou hast loved justice, and hated iniquity: therefore God, thy God, hath anointed thee with the oil of gladness above thy fellows.\par
\switchcolumn*

\LA
\begin{Resp}\Rsign Deus, qui sedes super thronum, et júdicas æquitátem, esto refúgium páuperum in tribulatióne: \Star{} Quia tu solus labórem et dolórem consíderas.\end{Resp}
\begin{Resp}\Vsign Tibi enim derelíctus est pauper, pupíllo tu eris adjútor.\end{Resp}
\begin{Resp}\Rsign Quia tu solus labórem et dolórem consíderas.\end{Resp}
\switchcolumn
\EN
\begin{Resp}\Rsign O God, Which satest in the throne judging right, be Thou a refuge for the poor, a refuge in times of trouble. \Star{} For Thou alone beholdest mischief and spite.\end{Resp}
\begin{Resp}\Vsign The poor leaveth himself unto thee; Thou wilt be the helper of the fatherless.\end{Resp}
\begin{Resp}\Rsign For Thou alone beholdest mischief and spite.\end{Resp}
\switchcolumn*

\LA
\LessonHead{Lectio III}
\switchcolumn
\EN
\LessonHead{Lesson III}
\switchcolumn*

\LA
\Cite{Heb 1:10-14}
\switchcolumn
\EN
\Cite{Heb 1:10-14}
\switchcolumn*

\LA
\noindent Et: Tu in princípio, Dómine, terram fundásti: et ópera mánuum tuárum sunt cæli. Ipsi períbunt, tu autem permanébis, et omnes ut vestiméntum veteráscent: et velut amíctum mutábis eos, et mutabúntur: tu autem idem ipse es, et anni tui non defícient. Ad quem autem angelórum dixit aliquándo: Sede a dextris meis, quoadúsque ponam inimícos tuos scabéllum pedum tuórum? Nonne omnes sunt administratórii spíritus, in ministérium missi propter eos, qui hæreditátem cápient salútis?\par
\switchcolumn
\EN
\noindent And: Thou in the beginning, O Lord, didst found the earth: and the works of thy hands are the heavens. They shall perish, but thou shalt continue: and they shall all grow old as a garment. And as a vesture shalt thou change them, and they shall be changed: but thou art the selfsame, and thy years shall not fail. But to which of the angels said he at any time: Sit on my right hand, until I make thy enemies thy footstool? Are they not all ministering spirits, sent to minister for them, who shall receive the inheritance of salvation?\par
\switchcolumn*

\LA
\begin{Resp}\Rsign A dextris est mihi Dóminus, ne commóvear: \Star{} Propter hoc dilatátum est cor meum, et exsultávit lingua mea.\end{Resp}
\begin{Resp}\Vsign Dóminus pars hereditátis meæ, et cálicis mei.\end{Resp}
\begin{Resp}\Rsign Propter hoc dilatátum est cor meum, et exsultávit lingua mea.\end{Resp}
\begin{Resp}\Vsign Glória Patri, et Fílio, \Star{} et Spirítui Sancto.\end{Resp}
\begin{Resp}\Rsign Propter hoc dilatátum est cor meum, et exsultávit lingua mea.\end{Resp}
\switchcolumn
\EN
\begin{Resp}\Rsign The Lord is at my right hand, I shall never be moved. \Star{} Therefore my heart is glad, and my tongue rejoiceth.\end{Resp}
\begin{Resp}\Vsign The Lord is the portion of mine inheritance, and of my cup.\end{Resp}
\begin{Resp}\Rsign Therefore my heart is glad, and my tongue rejoiceth.\end{Resp}
\begin{Resp}\Vsign Glory be to the Father, and to the Son, \Star{} and to the Holy Ghost.\end{Resp}
\begin{Resp}\Rsign Therefore my heart is glad, and my tongue rejoiceth.\end{Resp}
\switchcolumn*

\LA
\LessonHead{Lectio IV}
\switchcolumn
\EN
\LessonHead{Lesson IV}
\switchcolumn*

\LA
\LessonTitle{Sermo sancti Athanásii Epíscopi}
\switchcolumn
\EN
\LessonTitle{From the Sermons of St. Athanasius, Pope (of Alexandria.)}
\switchcolumn*

\LA
\Source{Orat. 2 contra Ariános, post medium.}
\switchcolumn
\EN
\Source{2nd against the Arians.}
\switchcolumn*

\LA
\noindent Si persónam, rem, tempus apostólici dicti cognóscerent hærétici, numquam humána in Deitátem transferéntes, tam ímpie et stulte advérsus Christum sese habuíssent. Id intuéri licébit, si inítium lectiónis dénuo repetítum probe excípias. Dicit enim Apóstolus: Multifáriam, multísque modis olim Deus locútus est pátribus nostris per Prophétas: últimis autem diébus locútus est nobis in Fílio. Atque ita paulo post dicit: Perfécta ab eo nostrórum peccatórum purificatióne, ipsum sedére ad déxteram majestátis in excélsis, tanto meliórem Angelis factum, quanto præstántius nomen præ illis sortítus est. De eo ígitur témpore, quo nobis per Fílium locútus est, cum peccatórum purgátio fíeret, apostólicum dictum mentiónem facit. Quando autem nobis locútus est in Fílio, aut quando purgátio peccatórum facta, aut quando natus est homo, nisi post Prophétas, idque in últimis diébus?\par
\switchcolumn
\EN
\noindent If the heretics had but known the person, the matter, and the times of the Apostle who spoke, they would never have spoken of Godhead as if It were human, nor borne themselves so wickedly, and withal so foolishly against Christ. It will be permitted to us to return, and to take again the first words of the Lesson. The Apostle then saith God, Who at sundry times and in diverse manners, spake in time past unto the fathers by the Prophets, hath in these last days spoken unto us by His Son and again, a little farther on: When the Son had purged our sins, He sat down on the right hand of the Majesty on high being made so much better than the angels as He hath by inheritance obtained a more excellent name than they. The Apostle here expressly nameth the times wherein God hath spoken unto us by His Son, and wherein the Same His Son hath purged our sins; for when hath He spoken unto us by His Son, when did the Son purge our sins, or when was He born a Man, but since God spake unto the Fathers by the Prophets, namely, in these last days?\par
\switchcolumn*

\LA
\begin{Resp}\Rsign Notas mihi fecísti, Dómine, vias vitæ: \Star{} Adimplébis me lætítia cum vultu tuo: delectatiónes in déxtera tua usque in finem.\end{Resp}
\begin{Resp}\Vsign Tu es qui restítues hereditátem meam mihi.\end{Resp}
\begin{Resp}\Rsign Adimplébis me lætítia cum vultu tuo: delectatiónes in déxtera tua usque in finem.\end{Resp}
\switchcolumn
\EN
\begin{Resp}\Rsign O Lord, Thou hast shown me the path of life. \Star{} Thou shalt fill me with joy in thy presence, at thy right hand there are pleasures for evermore.\end{Resp}
\begin{Resp}\Vsign Thou art He That shalt restore mine inheritance unto me.\end{Resp}
\begin{Resp}\Rsign Thou shalt fill me with joy in thy presence, at thy right hand there are pleasures for evermore.\end{Resp}
\switchcolumn*

\LA
\LessonHead{Lectio V}
\switchcolumn
\EN
\LessonHead{Lesson V}
\switchcolumn*

\LA
\noindent Deínde cum narrátio institúta esset de humána Verbi dispensatióne, deque últimis tempóribus: consequénter commemorávit, Deum neque superióribus ætátibus tacuísse, sed locútum esse per Prophétas: et postquam Prophétæ suo offício perfúncti sunt, et lex per Angelos pronuntiáta est, et Fílius étiam ad nos descéndit, et ad ministrándum accéssit; tunc demum necessário subíntulit: Tanto mélior Angelis factus: osténdere volens, quanto Fílius præ servo excéllit, tanto functióne officióque servórum, Fílii administratiónem meliórem fuísse.\par
\switchcolumn
\EN
\noindent The Apostle, about to enter on the subject of the Word's human dispensation and the last days, naturally mentioneth first that God had not up to those days been silent, but had spoken unto the fathers by the Prophets and, after the Prophets had discharged their office, and the law had been given by the ministry of angels, that the Son also came down unto us to minister and then he addeth, being made so much better than the angels, to show that as the Son differeth from a servant, so is the ministry of the Son better than the duty and office of servants.\par
\switchcolumn*

\LA
\begin{Resp}\Rsign Díligam te, Dómine, virtus mea: Dóminus firmaméntum meum, \Star{} Et refúgium meum.\end{Resp}
\begin{Resp}\Vsign Liberátor meus, Deus meus, adjútor meus.\end{Resp}
\begin{Resp}\Rsign Et refúgium meum.\end{Resp}
\switchcolumn
\EN
\begin{Resp}\Rsign I will love thee, O Lord, my strength; the Lord is my rock \Star{} And my fortress.\end{Resp}
\begin{Resp}\Vsign My Deliverer, my God, mine Helper.\end{Resp}
\begin{Resp}\Rsign And my fortress.\end{Resp}
\switchcolumn*

\LA
\LessonHead{Lectio VI}
\switchcolumn
\EN
\LessonHead{Lesson VI}
\switchcolumn*

\LA
\noindent Functiónem ígitur discérnens Apóstolus, tum véterem, tum novam, magna dicéndi libertáte útitur, ad Judǽos scribens et loquens. Propter hoc ígitur non in univérsum ex própria comparatiónis ratióne dixit, quod major aut honorátior esset: ne quis quasi de ejúsdem géneris, et cum eo commúnibus rebus hæc verba intellígeret: sed ídeo meliórem illum dixit, ut discrímen natúræ Fílii ad res creátas indicáret.\par
\switchcolumn
\EN
\noindent The Apostle, therefore, seeing the difference between the new ministry and the old, maketh very bold in writing and speaking to the Jews. For this cause, therefore, he doth not compare the details of the two ministries, and then come to the general conclusion that the new was greater or more honourable than the old, (lest any should understand that the two ministries were of the same kind, and that the conclusion that the new is better is arrived at by comparing the degrees in each of things which they had in common,) but he saith that the Son was made better, to distinguish at once and completely the nature of the Son from the nature of things created.\par
\switchcolumn*

\LA
\begin{Resp}\Rsign Dómini est terra, et plenitúdo ejus: \Star{} Orbis terrárum, et univérsi qui hábitant in eo.\end{Resp}
\begin{Resp}\Vsign Ipse super mária fundávit eam, et super flúmina præparávit illam.\end{Resp}
\begin{Resp}\Rsign Orbis terrárum, et univérsi qui hábitant in eo.\end{Resp}
\begin{Resp}\Vsign Glória Patri, et Fílio, \Star{} et Spirítui Sancto.\end{Resp}
\begin{Resp}\Rsign Orbis terrárum, et univérsi qui hábitant in eo.\end{Resp}
\switchcolumn
\EN
\begin{Resp}\Rsign The earth is the Lord's, and the fulness thereof \Star{} The world, and they that dwell therein.\end{Resp}
\begin{Resp}\Vsign For He hath founded it upon the seas, and established it upon the floods.\end{Resp}
\begin{Resp}\Rsign The world, and they that dwell therein.\end{Resp}
\begin{Resp}\Vsign Glory be to the Father, and to the Son, \Star{} and to the Holy Ghost.\end{Resp}
\begin{Resp}\Rsign The world, and they that dwell therein.\end{Resp}
\switchcolumn*

\LA
\LessonHead{Lectio VII}
\switchcolumn
\EN
\LessonHead{Lesson VII}
\switchcolumn*

\LA
\LessonTitle{Léctio sancti Evangélii secúndum Matthǽum}
\switchcolumn
\EN
\LessonTitle{From the Holy Gospel according to Matthew}
\switchcolumn*

\LA
\Cite{Matt 13:31-35}
\switchcolumn
\EN
\Cite{Matt 13:31-32}
\switchcolumn*

\LA
\noindent In illo témpore: Dixit Jesus turbis parábolam hanc: Símile est regnum cælórum grano sinápis, quod accípiens homo seminávit in agro suo. Et réliqua.\par
\switchcolumn
\EN
\noindent At that time, Jesus said to the people a parable: The kingdom of heaven is like to a grain of mustard seed, which a man took and sowed in his field. And so on.\par
\switchcolumn*

\LA
\LessonTitle{Homilía S. Hierónymi Presbýteri}
\switchcolumn
\EN
\LessonTitle{Homily by St. Jerome, Priest (at Bethlehem.)}
\switchcolumn*

\LA
\Source{Lib. 2. Comment. in cap. 13. Matth.}
\switchcolumn
\EN
\Source{Book II Comment. on Matth. xiii.}
\switchcolumn*

\LA
\noindent Regnum cælórum prædicátio Evangélii est, et notítia Scripturárum, quæ ducit ad vitam, et de qua dícitur ad Judǽos: Auferétur a vobis regnum Dei, et dábitur genti faciénti fructus ejus. Símile est ergo hujuscémodi regnum grano sinápis, quod accípiens homo seminávit in agro suo. Homo qui séminat in agro suo, a plerísque Salvátor intellégitur, quod in ánimis credéntium séminet: ab áliis ipse homo séminans in agro suo, hoc est in semetípso, et in corde suo.\par
\switchcolumn
\EN
\noindent The kingdom of heaven is the proclamation of the Gospel, and that knowledge of the Scriptures, which leadeth unto life, and whereof it is said to the Jews: The kingdom of God shall be taken from you, and given to a nation bringing forth the fruits thereof. (Matth. xxi. 43) Therefore is this kingdom like to a grain of mustard-seed, which a man took and sowed in his field. By the man that sowed it in his field, many understand to be meant the Saviour, because He is the Sower That soweth in the souls of believers; others understand every man that soweth good seed in his own field, that is, in himself and in his own heart.\par
\switchcolumn*

\LA
\begin{Resp}\Rsign Laudábilis pópulus, \Star{} Quem Dóminus exercítuum benedíxit dicens: Opus mánuum meárum tu es, heréditas mea Israël.\end{Resp}
\begin{Resp}\Vsign Beáta gens, cujus est Dóminus Deus, pópulus eléctus in hereditátem.\end{Resp}
\begin{Resp}\Rsign Quem Dóminus exercítuum benedíxit dicens: Opus mánuum meárum tu es, heréditas mea Israël.\end{Resp}
\switchcolumn
\EN
\begin{Resp}\Rsign Blessed is the people \Star{} Whom the Lord of hosts hath blessed, saying O Israel thou art the work of Mine own hands, thou art Mine own inheritance.\end{Resp}
\begin{Resp}\Vsign Blessed is the nation whose God is the Lord, and the people whom He hath chosen for His own inheritance.\end{Resp}
\begin{Resp}\Rsign Whom the Lord of hosts hath blessed, saying O Israel thou art the work of Mine own hands, thou art Mine own inheritance.\end{Resp}
\switchcolumn*

\LA
\LessonHead{Lectio VIII}
\switchcolumn
\EN
\LessonHead{Lesson VIII}
\switchcolumn*

\LA
\noindent Quis est iste, qui séminat, nisi sensus noster et ánimus; qui suscípiens granum prædicatiónis, et fovens seméntem, humóre fídei facit in agro sui péctoris pulluláre? Prædicátio Evangélii mínima est ómnibus disciplínis. Ad primam quippe doctrínam, fidem non habet veritátis, hóminem Deum, Christum mórtuum, et scándalum crucis prǽdicans. Confer hujuscémodi doctrínam dogmátibus philosophórum, et libris eórum, et splendóri eloquéntiæ, et compositióni sermónum: et vidébis quanto minor sit céteris semínibus seméntis Evangélii.\par
\switchcolumn
\EN
\noindent Who is he that soweth, but our own mind and soul, which take the grain from preaching, and by nourishing it in the soil, cause it to sprout in the field of our own breast? The preaching of the Gospel is the least of all doctrines. He that preacheth, for his first lesson, God made man, Christ dead, and the stumbling-block of the Cross, receiveth at first but little credit. Compare such teaching as this with the doctrines of the Philosophers, with their books, their magnificent eloquence, and their rounded sentences, and thou shalt see how the grain of the Gospel, when it is sown, is the humblest of all seeds.\par
\switchcolumn*

\LA
\begin{Resp}\Rsign Duo Séraphim clamábant alter ad álterum: \Star{} Sanctus, sanctus, sanctus Dóminus Deus Sábaoth: * Plena est omnis terra glória ejus.\end{Resp}
\begin{Resp}\Vsign Tres sunt qui testimónium dant in cælo: Pater, Verbum, et Spíritus Sanctus: et hi tres unum sunt.\end{Resp}
\begin{Resp}\Rsign Sanctus, sanctus, sanctus Dóminus Deus Sábaoth.\end{Resp}
\begin{Resp}\Vsign Glória Patri, et Fílio, \Star{} et Spirítui Sancto.\end{Resp}
\begin{Resp}\Rsign Plena est omnis terra glória ejus.\end{Resp}
\switchcolumn
\EN
\begin{Resp}\Rsign One Seraph cried unto another: \Star{} Holy, Holy, Holy is the Lord God of hosts; the whole earth is full of His glory.\end{Resp}
\begin{Resp}\Vsign There are Three That bear record in heaven: the Father, the Word, and the Holy Ghost, and these Three are One.\end{Resp}
\begin{Resp}\Rsign Holy, Holy, Holy is the Lord God of hosts.\end{Resp}
\begin{Resp}\Vsign Glory be to the Father, and to the Son, \Star{} and to the Holy Ghost.\end{Resp}
\begin{Resp}\Rsign The whole earth is full of His glory.\end{Resp}
\switchcolumn*

\LA
\LessonHead{Lectio IX}
\switchcolumn
\EN
\LessonHead{Lesson IX}
\switchcolumn*

\LA
\noindent Sed illa cum créverint, nihil mordax, nihil vívidum, nihil vitále demónstrant: sed totum fláccidum marcidúmque et mollítum ebúllit in ólera et in herbas, quæ cito aréscunt et córruunt. Hæc autem prædicátio, quæ parva videbátur in princípio, cum vel in ánima credéntis, vel in tot mundo sata fúerit, non exsúrgit in ólera, sed crescit in árborem: ita ut vólucres cæli (quas vel ánimas credéntium, vel fortitúdines, Dei servítio mancipátas, sentíre debémus) véniant et hábitent in ramis ejus. Ramos puto evangélicæ árboris, quæ de grano sinápis créverit, dógmatum esse diversitátes, in quibus supradictárum vólucrum unaquǽque requiéscit.\par
\switchcolumn
\EN
\noindent But when the doctrines of men grow up, there is therein nothing piercing, nothing healthy, nothing life-giving. The plant is drooping, and delicate, and soft. There are herbs and grass whereof it may truly be said that the grass withereth and the flower fadeth. (Isa. xl. 8.) But the grain of Gospel seed, though, when it was sown, it seemed to be the least of all seeds, when once it is rooted in the soul of man, or in the whole world, groweth not into an herb, but becometh a tree so that the birds of the air (whereby we may understand, either the souls of believers, or the (angelic) powers bound to the service of God,) come and lodge in the branches thereof. I consider that the branches of the Gospel tree, which groweth from the grain of mustard-seed, are the diverse developments of doctrine, on which the birds above mentioned find resting-places.\par
\switchcolumn*

\LA
\TeDeum{Te Deum}
\switchcolumn
\EN
\TeDeum{Te Deum}
\switchcolumn*

\end{paracol}

\PartTitle{Proprium de Sanctis}{Proper of Saints}

\addcontentsline{toc}{chapter}{Proprium de Sanctis \textit{— Proper of Saints}}

\SeasonTitle{Mensis Augustus}{August}

\addcontentsline{toc}{section}{Mensis Augustus \textit{— August}}

\par\needspace{10\baselineskip}\addvspace{1.2em}{\centering\small\textsc{Die 28 Augusti} · \textsc{28 August}\par}
\Office{S. Augustini Episcopi et Confessoris et Ecclesiæ Doctoris}{St. Augustine of Hippo, Bishop, Confessor, and Doctor of the Church}{Duplex}
\addcontentsline{toc}{subsection}{Die 28 Augusti · S. Augustini Episcopi et Confessoris et Ecclesiæ Doctoris \textit{— St. Augustine of Hippo, Bishop, Confessor, and Doctor of the Church}}
\begin{paracol}{2}
\LA
\RefLine{Lectiones I–III: de Scriptura occurrente.}
\switchcolumn
\EN
\RefLine{Lessons I–III: from the Scripture of the day.}
\switchcolumn*

\LA
\RefLine{Lectiones VII–IX: ex Commúni Doctoris Pontificis (C4a).}
\switchcolumn
\EN
\RefLine{Lessons VII–IX: from the Common of a Doctor Bishop (C4a).}
\switchcolumn*

\LA
\LessonHead{Lectio IV}
\switchcolumn
\EN
\LessonHead{Lesson IV}
\switchcolumn*

\LA
\noindent Augustinus, Tagaste in Africa honestis parentibus natus, ac puer docilitate ingenii æquales longe superans, brevi omnibus doctrina antecelluit. Adolescens, dum esset Carthagine, in Manichæorum hæresim incidit. Postea Romam profectus, inde Mediolanum missus ut rhetoricam doceret, cum ibi frequens Ambrósii episcopi esset auditor, ejus opera incensus studio catholicæ fidei, annos natus triginta tres ab ipso baptizatur. Reversus in Africam, cum religione vitæ sanctimoniam conjungens, a Valerio notæ sanctitatis episcopo Hipponensi presbyter factus est. Quo tempore familiam instituit religiosorum, quibuscum victu communi eodemque cultu utens, eos ad apostolicæ vitæ doctrinæque disciplinam diligentissime erudiebat. Sed cum vigeret Manichæorum hæresis, vehementius in illam invehi cœpit, Fortunatumque hæresiarcham confutavit.\par
\switchcolumn
\EN
\noindent Augustine was born of honourable parents at Tagaste in Africa, (upon the 13th day of November, in the year of our Lord 354.) As a boy his great intellectual sharpness caused him to distance all his companions in learning. When he was living at Carthage as a young man, he fell into the heresy of the Manichaeans. He afterwards went to Rome, and was thence sent to Milan to teach Rhetoric. At Milan he often went to hear the sermons of Bishop Ambrose, by whose labours he was drawn to the Catholic Church, and by whom he was baptized (on Holy Saturday, 387,) at the age of thirty-three. After his return to Africa, (in 388,) Valerius, the illustrious and saintly Bishop of Hippo, finding him to unite holiness of life, with Catholic profession, made him a Priest, (about the end of 390.) At this time he founded a sort of family of godly men, who lived and worshipped in common with him, and whom he earnestly formed upon the model of the Apostolic life and teaching. The Manichaean heresy flaming forth with violence, he began strongly to attack it, and confounded the arch-heretic Fortunatus.\par
\switchcolumn*

\LA
\begin{Resp}\Rsign Invéni David servum meum, óleo sancto meo unxi eum: \Star{} Manus enim mea auxiliábitur ei.\end{Resp}
\begin{Resp}\Vsign Nihil profíciet inimícus in eo, et fílius iniquitátis non nocébit ei.\end{Resp}
\begin{Resp}\Rsign Manus enim mea auxiliábitur ei.\end{Resp}
\switchcolumn
\EN
\begin{Resp}\Rsign I have found David My servant, with My holy oil have I anointed him \Star{} For My hand shall help him.\end{Resp}
\begin{Resp}\Vsign The enemy shall prevail nothing against him, nor the son of wickedness afflict him.\end{Resp}
\begin{Resp}\Rsign For My hand shall help him.\end{Resp}
\switchcolumn*

\LA
\LessonHead{Lectio V}
\switchcolumn
\EN
\LessonHead{Lesson V}
\switchcolumn*

\LA
\noindent Hac Augustíni pietate commotus Valerius, eum adjutorem adhibuit episcopalis officii. Nihil illo fuit humilius, nihil continentius. Lectus ac vestitus moderatus, vulgaris mensa, quam semper sacra vel lectione vel disputatione condiebat. Tanta benignitate fuit in pauperes, ut, cum non esset alia facultas, sacra vasa frangeret ad eorum inopiam sustentandam. Feminarum, et in eis sororis, et fratris filiæ, contubernium familiaritatemque vitavit: quippe qui diceret, etsi propinquæ mulieres suspectæ non essent, tamen quæ ad eas ventitarent, posse suspicionem efficere. Nullum finem fecit prædicandi Dei verbum, nisi gravi morbo oppressus. Hæreticos perpetuo insectatus et coram et scriptis, ac nullo loco passus consistere, Africam a Manichæorum, Donatistarum, Pelagianorum, aliorumque præterea hæreticorum errore magna ex parte liberavit.\par
\switchcolumn
\EN
\noindent Valerius, moved by the godly zeal of Augustine, (in December 395,) joined him with himself as an assistant in his duties of Bishop (and dying in the year following, was succeeded by him.) He was lowly and pure in the highest degree. His furniture and dress were plain, and his food of the commonest sort, which he always seasoned when at table by either reading some religious book, or arguing upon some religious subject. His tenderness to the poor was such that, failing all other resources, he broke up the hallowed vessels to relieve their wants. It was his rule not to dwell or be very close friends with any woman, a rule which he did not relax even in the case of his sister and niece, for he was accustomed to say, that although no scandal could arise in the case of such near kinswomen, yet it might arise concerning the women friends who sought their company. He never ceased to preach the Word of God, until he was disabled by heavy sickness. He was always an hard follower after heretics, and by his words and his writings never suffered them to rest anywhere. In great measure, he purged Africa of the Manicheans, Donatists, Pelagians, and other heretics.\par
\switchcolumn*

\LA
\begin{Resp}\Rsign Pósui adjutórium super poténtem, et exaltávi eléctum de plebe mea: \Star{} Manus enim mea auxiliábitur ei.\end{Resp}
\begin{Resp}\Vsign Invéni David servum meum, óleo sancto meo unxi eum.\end{Resp}
\begin{Resp}\Rsign Manus enim mea auxiliábitur ei.\end{Resp}
\switchcolumn
\EN
\begin{Resp}\Rsign I have laid help upon one that is mighty, and have exalted one chosen out of My people \Star{} For My hand shall help him.\end{Resp}
\begin{Resp}\Vsign I have found David My servant, with My holy oil have I anointed him.\end{Resp}
\begin{Resp}\Rsign For My hand shall help him.\end{Resp}
\switchcolumn*

\LA
\LessonHead{Lectio VI}
\switchcolumn
\EN
\LessonHead{Lesson VI}
\switchcolumn*

\LA
\noindent Tam multa pie, subtiliter et copiose scripsit, ut christianam doctrinam maxime illustrarit. Quem in primis secuti sunt, qui postea theologicam disciplinam via et ratione tradiderunt. Wandalis Africam bello vastantibus, et Hipponem tertium jam mensem obsidentibus, in febrim incidit. Itaque cum discessum e vita sibi instare intelligeret, Psalmos David qui ad pœnitentiam pertinent, in conspectu positos profusis lacrimis legebat. Solebat autem dicere neminem, etsi nullius sceleris sibi conscius esset, committere debere, ut sine pænitentia migraret e vita. Ergo sensibus integris, in oratione defixus, astantibus fratribus, quos ad caritatem, pietatem, virtutesque omnes erat adhortatus, migravit in cælum. Vixit annos septuaginta sex, in episcopatu ad triginta sex. Cujus corpus primum in Sardiniam delatum, deínde a Luitprando, Longobardorum rege, magno pretio redemptum, Ticinum translatum est, ibique honorifice conditum.\par
\switchcolumn
\EN
\noindent He wrote so much, and that with such godliness and understanding, that he is to be held among the very chiefest of them by whom the teachings of Christianity have been shown forth. He is one of the first of those whom later theologians have followed, in method, and in argument. He fell sick of a fever what time the Vandals were laying Africa waste, and when they were busy in the third month of besieging Hippo. When he understood that his departure from this present life was at hand, he caused the Psalms of David which most speak the language of repentance to be placed before him, and read them with tears, for he was wont to say that even if a man's conscience were to accuse him of no sin, he should not dare to leave this world except as a penitent. His senses remained vigorous to the last, and it was while rapt in prayer, in the presence of the brethren whom he had exhorted to love, godliness, and all goodness, that he departed for heaven, (upon the 28th of August, 430.) He lived 76 years, whereof he had been a Bishop nearly thirty six. His body was first carried to Sardinia, but Luitprand, King of the Lombards, afterwards bought it for a great price, and took it to Ticino, where it is honourably buried.\par
\switchcolumn*

\LA
\begin{Resp}\Rsign Iste est, qui ante Deum magnas virtútes operátus est, et omnis terra doctrína ejus repléta est: \Star{} Ipse intercédat pro peccátis ómnium populórum.\end{Resp}
\begin{Resp}\Vsign Iste est, qui contémpsit vitam mundi, et pervénit ad cæléstia regna.\end{Resp}
\begin{Resp}\Rsign Ipse intercédat pro peccátis ómnium populórum.\end{Resp}
\begin{Resp}\Vsign Glória Patri, et Fílio, \Star{} et Spirítui Sancto.\end{Resp}
\begin{Resp}\Rsign Ipse intercédat pro peccátis ómnium populórum.\end{Resp}
\switchcolumn
\EN
\begin{Resp}\Rsign This is he which wrought great wonders before God, and the whole earth is full of his teaching \Star{} May he pray for all people, that their sins may be forgiven unto them\end{Resp}
\begin{Resp}\Vsign This is he which loved not his life in this world, and hath attained unto the kingdom of heaven.\end{Resp}
\begin{Resp}\Rsign May he pray for all people, that their sins may be forgiven unto them\end{Resp}
\begin{Resp}\Vsign Glory be to the Father, and to the Son, \Star{} and to the Holy Ghost.\end{Resp}
\begin{Resp}\Rsign May he pray for all people, that their sins may be forgiven unto them\end{Resp}
\switchcolumn*

\end{paracol}

\par\needspace{10\baselineskip}\addvspace{1.2em}{\centering\small\textsc{Die 29 Augusti} · \textsc{29 August}\par}
\Office{In Decollatione S. Joannis Baptistæ}{Beheading of St. John the Baptist}{Duplex majus}
\addcontentsline{toc}{subsection}{Die 29 Augusti · In Decollatione S. Joannis Baptistæ \textit{— Beheading of St. John the Baptist}}
\begin{paracol}{2}
\LA
\RefLine{Lectio IX: de S. Sabinæ Martyris (08-29cc).}
\switchcolumn
\EN
\RefLine{Lesson IX: of St. Sabina, Martyr (08-29cc).}
\switchcolumn*

\LA
\LessonHead{Lectio I}
\switchcolumn
\EN
\LessonHead{Lesson I}
\switchcolumn*

\LA
\LessonTitle{Incipit liber Jeremíæ Prophétæ}
\switchcolumn
\EN
\LessonTitle{Lesson from the book of Jeremias}
\switchcolumn*

\LA
\Cite{Jer 1:1-5}
\switchcolumn
\EN
\Cite{Jer 1:1-5}
\switchcolumn*

\LA
\noindent Verba Jeremíæ fílii Helcíæ, de sacerdótibus, qui fuérunt in Anathoth, in terra Bénjamin. Quod factum est verbum Dómini ad eum in diébus Josíæ fílii Amon regis Juda, in tertiodécimo anno regni ejus. Et factum est in diébus Jóakim fílii Josíæ regis Juda, usque ad consummatiónem undécimi anni Sedecíæ fílii Josíæ regis Juda, usque ad transmigratiónem Jerúsalem, in mense quinto. Et factum est verbum Dómini ad me, dicens: Priúsquam te formárem in útero, novi te: et ántequam exíres de vulva, sanctificávi te, et prophétam in géntibus dedi te.\par
\switchcolumn
\EN
\noindent The words of Jeremias the son of Helcias, of the priests that were in Anathoth, in the land of Benjamin. The word of the Lord which came to him in the days of Josias the son of Amon king of Juda, in the thirteenth year of his reign. And which came to him in the days of Joakim the son of Josias king of Juda, unto the end of the eleventh year of Sedecias the son of Josias king of Juda, even unto the carrying away of Jerusalem captive, in the fifth month. And the word of the Lord came to me, saying: Before I formed thee in the bowels of thy mother, I knew thee: and before thou camest forth out of the womb, I sanctified thee, and made thee a prophet unto the nations.\par
\switchcolumn*

\LA
\begin{Resp}\Rsign Misit Heródes rex manus, ac ténuit Joánnem et vinxit eum in cárcere, quia metuébat eum propter Herodíadem, \Star{} Quam túlerat fratri suo Philíppo uxórem.\end{Resp}
\begin{Resp}\Vsign Arguébat Heródem Joánnes propter Herodíadem.\end{Resp}
\begin{Resp}\Rsign Quam túlerat fratri suo Philíppo uxórem.\end{Resp}
\switchcolumn
\EN
\begin{Resp}\Rsign Herod the King sent forth, and laid hold upon John, and bound him in prison, for he feared him, for Herodias' sake, \Star{} His brother Philip's wife, for he had married her.\end{Resp}
\begin{Resp}\Vsign For John had rebuked Herod, for Herodias' sake.\end{Resp}
\begin{Resp}\Rsign His brother Philip's wife, for he had married her.\end{Resp}
\switchcolumn*

\LA
\LessonHead{Lectio II}
\switchcolumn
\EN
\LessonHead{Lesson II}
\switchcolumn*

\LA
\Cite{Jer 1:6-10}
\switchcolumn
\EN
\Cite{Jer 1:6-10}
\switchcolumn*

\LA
\noindent Et dixi, A a a, Dómine Deus: ecce néscio loqui, quia puer ego sum. Et dixit Dóminus ad me: Noli dícere, Puer sum: quóniam ad ómnia, quæ mittam te, ibis: et univérsa, quæcúmque mandávero tibi, loquéris. Ne tímeas a fácie eórum: quia tecum ego sum, ut éruam te, dicit Dóminus. Et misit Dóminus manum suam, et tétigit os meum: et dixit Dóminus ad me: Ecce dedi verba mea in ore tuo: ecce constítui te hódie super gentes, et super regna, ut evéllas, et déstruas, et dispérdas, et díssipes, et ædífices, et plantes.\par
\switchcolumn
\EN
\noindent And I said: Ah, ah, ah, Lord God: behold, I cannot speak, for I am a child. And the Lord said to me: Say not: I am a child: for thou shalt go to all that I shall send thee: and whatsoever I shall command thee, thou shalt speak. Be not afraid at their presence: for I am with thee to deliver thee, saith the Lord. And the Lord put forth his hand, and touched my mouth: and the Lord said to me: Behold I have given my words in thy mouth: Lo, I have set thee this day over the nations, and over the kingdoms, to root up, and pull down, and to waste, and to destroy, and to build, and to plant.\par
\switchcolumn*

\LA
\begin{Resp}\Rsign Joánnes Baptísta arguébat Heródem \Star{} Propter Herodíadem, quam túlerat fratri suo vivénti uxórem.\end{Resp}
\begin{Resp}\Vsign Misso Heródes spiculatóre, præcépit amputári caput Joánnis in cárcere.\end{Resp}
\begin{Resp}\Rsign Propter Herodíadem, quam tulerat fratri suo vivénti uxórem.\end{Resp}
\switchcolumn
\EN
\begin{Resp}\Rsign John the Baptist had rebuked Herod, \Star{} For Herodias' sake, his brother's wife, whom he had married while his brother was yet alive.\end{Resp}
\begin{Resp}\Vsign Herod sent an executioner, and commanded to behead John in the prison.\end{Resp}
\begin{Resp}\Rsign For Herodias' sake, his brother's wife, whom he had married while his brother was yet alive.\end{Resp}
\switchcolumn*

\LA
\LessonHead{Lectio III}
\switchcolumn
\EN
\LessonHead{Lesson III}
\switchcolumn*

\LA
\Cite{Jer 1:17-19}
\switchcolumn
\EN
\Cite{Jer 1:17-19}
\switchcolumn*

\LA
\noindent Tu ergo accínge lumbos tuos, et surge, et lóquere ad eos ómnia quæ ego præcípio tibi. Ne formídes a fácie eórum: nec enim timére te fáciam vultum eórum. Ego quippe dedi te hódie in civitátem munítam, et in colúmnam férream, et in murum ǽreum, super omnem terram, régibus Juda, princípibus ejus, et sacerdótibus, et pópulo terræ. Et bellábunt advérsum te, et non prævalébunt: quia ego tecum sum, ait Dóminus, ut líberem te.\par
\switchcolumn
\EN
\noindent Thou therefore gird up thy loins, and arise, and speak to them all that I command thee. Be not afraid at their presence: for I will make thee not to fear their countenance. For behold I have made thee this day a fortified city, and a pillar of iron, and a wall of brass, over all the land, to the kings of Juda, to the princes thereof, and to the priests, and to the people of the land. And they shall fight against thee, and shall not prevail: for I am with thee, saith the Lord, to deliver thee.\par
\switchcolumn*

\LA
\begin{Resp}\Rsign Puéllæ saltánti imperávit mater: \Star{} Nihil áliud petas, nisi caput Joánnis. * Et contristátus est rex propter jusjurándum et propter simul discumbéntes.\end{Resp}
\begin{Resp}\Vsign Ait puélla matri suæ: Quid petam? At illa ait.\end{Resp}
\begin{Resp}\Rsign Nihil áliud petas, nisi caput Joánnis.\end{Resp}
\begin{Resp}\Vsign Glória Patri, et Fílio, \Star{} et Spirítui Sancto.\end{Resp}
\begin{Resp}\Rsign Et contristátus est rex propter jusjurándum et propter simul discumbéntes.\end{Resp}
\switchcolumn
\EN
\begin{Resp}\Rsign The damsel danced, and her mother charged her, saying: \Star{} See thou ask nothing but only the head of John. And the king was sorry, for his oath's sake, and for their sakes which sat with him.\end{Resp}
\begin{Resp}\Vsign The damsel said unto her mother What shall I ask? And she said:\end{Resp}
\begin{Resp}\Rsign See thou ask nothing, but only the head of John. And the king was sorry, for his oath's sake, and for their sakes which sat with him.\end{Resp}
\begin{Resp}\Vsign Glory be to the Father, and to the Son, \Star{} and to the Holy Ghost.\end{Resp}
\begin{Resp}\Rsign And the king was sorry, for his oath's sake, and for their sakes which sat with him.\end{Resp}
\switchcolumn*

\LA
\LessonHead{Lectio IV}
\switchcolumn
\EN
\LessonHead{Lesson IV}
\switchcolumn*

\LA
\LessonTitle{Ex libro sancti Ambrósii Epíscopi de Virgínibus}
\switchcolumn
\EN
\LessonTitle{From the Book upon Virgins, written by St. Ambrose, Bishop (of Milan.)}
\switchcolumn*

\LA
\Source{Lib. 3 post initium}
\switchcolumn
\EN
\Source{ii. 6}
\switchcolumn*

\LA
\noindent Quóniam beáti Joánnis Baptístæ non strictim prætereúnda est recordátio, ínterest ut quis et a quibus et quam ob causam, quo modo et quo témpore sit occísus, advértere debeámus. Ab adúlteris justus occíditur, et a reis in júdicem capitális scéleris pœna convértitur. Deínde prǽmium saltatrícis, mors est Prophétæ. Postrémo (quod étiam omnes bárbari horrére consuevérunt) inter épulas atque convívia consummándæ crudelitátis profértur edíctum; et a convívio ad cárcerem, de cárcere ad convívium ferális flagítii circumfértur obséquium. Quanta in uno facínore sunt crímina!\par
\switchcolumn
\EN
\noindent We must not hurry by the record of the Blessed Baptist John. We must ask what he was, and by whom, and why, and how, and when he was slain. He was a righteous man murdered by adulterers. The guilty passed upon their judge the sentence of death. Moreover, the death of the Prophet was the fee of a dancing-girl. And lastly, there was a feature about it from which even savages shrink; the order for completing the atrocity was given amid the merriment of a dinner-party. From banquet to prison, from prison to banquet, that was the course run by the servants of the murderer. How many horrors does this simple crime embrace within its details?\par
\switchcolumn*

\LA
\begin{Resp}\Rsign Justus germinábit sicut lílium: \Star{} Et florébit in ætérnum ante Dóminum.\end{Resp}
\begin{Resp}\Vsign Plantátus in domo Dómini, in átriis domus Dei nostri.\end{Resp}
\begin{Resp}\Rsign Et florébit in ætérnum ante Dóminum.\end{Resp}
\switchcolumn
\EN
\begin{Resp}\Rsign The righteous shall grow as the lily \Star{} Yea, he shall flourish in the presence of the Lord for ever.\end{Resp}
\begin{Resp}\Vsign Those that be planted in the house of the Lord, shall flourish in the courts of the house of our God.\end{Resp}
\begin{Resp}\Rsign Yea, he shall flourish in the presence of the Lord for ever.\end{Resp}
\switchcolumn*

\LA
\LessonHead{Lectio V}
\switchcolumn
\EN
\LessonHead{Lesson V}
\switchcolumn*

\LA
\noindent Quis non, cum e convívio ad carcerem cursari vidéret, putaret Prophétam jussum esse dimitti? Quis, inquam, cum audísset natalem esse Herodis, solemne convivium, puellæ optiónem eligéndi quod vellet datam; missum ad Joánnem ob solutiónem non arbitrarétur? Quid crudelitáti cum deliciis? quid cum funéribus voluptati? Rápitur ad pœnam Prophéta convivali témpore, convivali præcépto, quo non cuperet vel absolvi: perímitur gládio, caput ejus affertur in disco. Hoc crudelitáti férculum debebátur, quo insatiáta épulis féritas vescerétur.\par
\switchcolumn
\EN
\noindent Who is there, that, on seeing the messenger hasten from the dinner-table to the prison, would not have forthwith concluded that he carried an order for the Prophet's release. If any one had heard that it was Herod's birth-day, and that he was giving a great feast, and that he had offered a damsel the choice of whatever she listed, and that thereupon a messenger had been sent to John's dungeon. If any one, I say, had heard this, what would he have supposed? He would have concluded that the damsel had asked and obtained John's freedom. What have executions in common with dinners, or death with gaiety? While the banquet was going on, the Prophet was hurried to death, by an order from the reveller whom he had not troubled even by a prayer for release. He was slain with the sword, and his head was served up in a plate. This was the new dish demanded by a cruelty which the Feast had been powerless to glut.\par
\switchcolumn*

\LA
\begin{Resp}\Rsign Iste cognóvit justítiam, et vidit mirabília magna, et exorávit Altíssimum: \Star{} Et invéntus est in número Sanctórum.\end{Resp}
\begin{Resp}\Vsign Iste est, qui contempsit vitam mundi, et pervenit ad cæléstia regna.\end{Resp}
\begin{Resp}\Rsign Et invéntus est in número Sanctórum.\end{Resp}
\switchcolumn
\EN
\begin{Resp}\Rsign This is he which knew righteousness, and saw great wonders, and made his prayer unto the Most High \Star{} And he is numbered among the Saints.\end{Resp}
\begin{Resp}\Vsign This is he which loved not his life in this world, and is come unto an everlasting kingdom.\end{Resp}
\begin{Resp}\Rsign And he is numbered among the Saints.\end{Resp}
\switchcolumn*

\LA
\LessonHead{Lectio VI}
\switchcolumn
\EN
\LessonHead{Lesson VI}
\switchcolumn*

\LA
\noindent Intuére, rex acerbíssime, tuo spectacula digna convivio. Pórrige déxteram, ne quid sævítiæ tuæ desit; ut inter digitos tuos rivi défluant sacri cruoris. Et, quóniam non exsaturari épulis fames, non restingui póculis pótuit inauditæ sævítiæ sitis; bibe sánguinem scaturiéntibus adhuc venis exsecti cápitis profluentem. Cerne óculos in ipsa morte sceleris tui testes, aversántes conspéctum deliciárum. Claudúntur lúmina non tam mortis necessitate quam horrore luxuriæ. Os aureum illud exsangue, cujus senténtiam ferre non poteras, conticescit, et adhuc timétur.\par
\switchcolumn
\EN
\noindent Look, savage King, look at a decoration which suiteth well with thy banquet. Put out thine hand, so as to lose no part of the luxury of cruelty, and let the streams of the sacred blood run between thy fingers. Thine hunger the dinner hath been unable to satisfy, thy cups have not been able to quench thine inhuman thirst. Suck, suck the blood which the still palpitating veins are discharging from the place where the neck has been severed. Look at the eyes. Even in death they remain the eyes of a witness of thine uncleanness, but they are closing themselves upon the spectacle of thy pleasures. Those eyes indeed are shutting but it seems not so much from the laws of natural death, as from horror at the scene of thine enjoyment. The golden mouth, whose bloodless lips are silent now, can repeat no more the denunciation which thou couldest not bear to hear, and still thou art afraid of it.\par
\switchcolumn*

\LA
\begin{Resp}\Rsign Honestum fecit illum Dóminus, et custodívit eum ab inimícis, et a seductóribus tutávit illum: \Star{} Et dedit illi claritátem ætérnam.\end{Resp}
\begin{Resp}\Vsign Descendítque cum illo in fóveam, et in vinculis non derelíquit eum.\end{Resp}
\begin{Resp}\Rsign Et dedit illi claritátem ætérnam.\end{Resp}
\begin{Resp}\Vsign Glória Patri, et Fílio, \Star{} et Spirítui Sancto.\end{Resp}
\begin{Resp}\Rsign Et dedit illi claritátem ætérnam.\end{Resp}
\switchcolumn
\EN
\begin{Resp}\Rsign The Lord made him honourable, and defended him from his enemies, and kept him safe from those that lay in wait for him, \Star{} And gave him perpetual glory.\end{Resp}
\begin{Resp}\Vsign He went down with him into the pit, and left him not in bonds.\end{Resp}
\begin{Resp}\Rsign And gave him perpetual glory.\end{Resp}
\begin{Resp}\Vsign Glory be to the Father, and to the Son, \Star{} and to the Holy Ghost.\end{Resp}
\begin{Resp}\Rsign And gave him perpetual glory.\end{Resp}
\switchcolumn*

\LA
\LessonHead{Lectio VII}
\switchcolumn
\EN
\LessonHead{Lesson VII}
\switchcolumn*

\LA
\LessonTitle{Léctio sancti Evangélii secúndum Marcum}
\switchcolumn
\EN
\LessonTitle{From the Holy Gospel according to Mark}
\switchcolumn*

\LA
\Cite{Marc 6:17-29}
\switchcolumn
\EN
\Cite{Mark 6:17-29}
\switchcolumn*

\LA
\noindent In illo témpore: Misit Heródes ac ténuit Joánnem, et vinxit eum in cárcere propter Herodíadem, uxórem Philíppi fratris sui, quia duxerat eam. Et réliqua.\par
\switchcolumn
\EN
\noindent At that time Herod had sent forth, and laid hold upon John, and bound him in prison, for Herodias' sake, his brother Philip's wife, for he had married her. And so on.\par
\switchcolumn*

\LA
\LessonTitle{Homilía sancti Augustíni Epíscopi}
\switchcolumn
\EN
\LessonTitle{Homily by St. Augustine, Bishop (of Hippo.)}
\switchcolumn*

\LA
\Source{Sermo 10 in novis. Serm.}
\switchcolumn
\EN
\Source{New Sermons.}
\switchcolumn*

\LA
\noindent Cum sanctum Evangélium legerétur, crudele spectaculum ante óculos nostros constitutum est: caput sancti Joánnis in disco, feralis missus crudelitátis propter ódium veritátis. Puella saltat, et sævit mater; et inter lascivias et delicias convivántium témere jurátur, et ímpie, quod jurátur, implétur. Factum est Joánni, quod ipse prædixerat; de Dómino enim Jesu Christo dixerat: Illum oportet crescere, me autem minui. Iste minutus est in cápite, ille crevit in cruce. Odium péperit véritas. Non pótuit æquo animo tolerári quod homo Dei sanctus monebat; qui útique salútem eórum quærebat, quos sic monebat. Respondérunt illi mala pro bonis.\par
\switchcolumn
\EN
\noindent The reading of the Holy Gospel hath set a scene of cruelty before our eyes even the head of St. John in a charger; a message of death sent forth to discharge the bloody commands of one that hateth the truth; a damsel dancing, and a mother rabid; a rash oath sworn in the midst of uncleanness and the revels of a supper, and a wicked fulfillment of the oath so sworn. It befell unto John according to his own saying. For he had said concerning the Lord Jesus Christ: He must increase, but I must decrease (John iii. 30,) so John decreased by an head, and Christ's height was made higher upon the Cross. The truth drew hatred. It could not be borne in patience that the holy man of God should utter a rebuke, albeit he sought by his rebuke nothing but the soul's health of them to whom he addressed it. They repaid him evil for good.\par
\switchcolumn*

\LA
\begin{Resp}\Rsign Desidérium ánimæ ejus tribuísti ei Dómine, \Star{} Et voluntáte labiórum ejus non fraudásti eum.\end{Resp}
\begin{Resp}\Vsign Quóniam prævenísti eum in benedictiónibus dulcédinis: posuísti in cápite ejus corónam de lápide pretióso.\end{Resp}
\begin{Resp}\Rsign Et voluntáte labiórum ejus non fraudásti eum.\end{Resp}
\switchcolumn
\EN
\begin{Resp}\Rsign Lord, Thou hast given him his heart's desire. \Star{} And hast not withholden the request of his lips.\end{Resp}
\begin{Resp}\Vsign For Thou hast prevented him with the blessings of sweetness: Thou hast set a crown of precious stones upon his head.\end{Resp}
\begin{Resp}\Rsign And hast not withholden the request of his lips.\end{Resp}
\switchcolumn*

\LA
\LessonHead{Lectio VIII}
\switchcolumn
\EN
\LessonHead{Lesson VIII}
\switchcolumn*

\LA
\noindent Quid enim ille diceret, nisi quo plenus erat? et quid illi respondérent, nisi quo pleni erant? Ille tríticum seminávit, sed spinas invénit. Dicebat regi: Non licet tibi habere uxórem fratris tui. Vincebat enim regem libído: tenebat apud se prohíbitam uxórem fratris sui. Sed eum tamen sic libebat, ut non sæviret; honorábat eum, a quo verum audiebat. Sed mulier detestábilis ódium concipiebat, quod aliquándo dato témpore páreret. Quando autem parturiebat, péperit filiam, filiam saltántem.\par
\switchcolumn
\EN
\noindent For what could he say but that whereof he was full? And what could they answer him but that whereof they were full? He sowed wheat, and found thorns. He had said unto the King: It is not lawful for thee to have thy brother's wife. Lust had got the better of the King, and he kept a woman whom it was not lawful for him to have, even his brother's wife. But she pleased him, so that his cruelty was lulled. He respected the Saint who had spoken the truth to him. But the horrible woman conceived hatred, and by-and-by brought it forth. When she brought forth, she brought forth a girl, a dancing-girl.\par
\switchcolumn*

\LA
\begin{Resp}\Rsign Stola jucunditátis índuit eum Dóminus: \Star{} Et corónam pulchritúdinis pósuit super caput ejus.\end{Resp}
\begin{Resp}\Vsign Cibávit illum Dóminus pane vitæ et intelléctus: et aqua sapiéntiæ salutáris potávit illum.\end{Resp}
\begin{Resp}\Rsign Et corónam pulchritúdinis pósuit super caput ejus.\end{Resp}
\begin{Resp}\Vsign Glória Patri, et Fílio, \Star{} et Spirítui Sancto.\end{Resp}
\begin{Resp}\Rsign Et corónam pulchritúdinis pósuit super caput ejus.\end{Resp}
\switchcolumn
\EN
\begin{Resp}\Rsign The Lord hath put on him a robe of honour, \Star{} And put about his head a crown of joy.\end{Resp}
\begin{Resp}\Vsign With the bread of life and understanding hath the Lord fed him, and given him the water of health and wisdom to drink.\end{Resp}
\begin{Resp}\Rsign And put about his head a crown of joy.\end{Resp}
\begin{Resp}\Vsign Glory be to the Father, and to the Son, \Star{} and to the Holy Ghost.\end{Resp}
\begin{Resp}\Rsign And put about his head a crown of joy.\end{Resp}
\switchcolumn*

\end{paracol}

\par\needspace{10\baselineskip}\addvspace{1.2em}{\centering\small\textsc{Die 29 Augusti} · \textsc{29 August}\par}
\Office{S. Sabinæ Martyris}{St. Sabina, Martyr}{Simplex}
\addcontentsline{toc}{subsection}{Die 29 Augusti · S. Sabinæ Martyris \textit{— St. Sabina, Martyr}}
\begin{paracol}{2}
\LA
\LessonHead{Lectio IX (pro commemoratione)}
\switchcolumn
\EN
\LessonHead{Lesson IX (when commemorated)}
\switchcolumn*

\LA
\LessonTitle{Commemoratio: S. Sabinæ Martyris}
\switchcolumn
\EN
\LessonTitle{Commemoration of: St. Sabina, Martyr}
\switchcolumn*

\LA
\noindent Sabina, mulier Romana, Valentini viri claríssimi uxor, a Seráphia vírgine christianæ fidei præceptis instituta, post sanctæ Vírginis martyrium, collectas ejus relíquias piis exsequiis sepelívit. Quæ propter eam causam paulo post, Hadriáno imperatóre, comprehensa, Elpidio júdici sistitur. Cui is: Tu ne illa Sabina et génere et matrimonio nobilíssima? At illa, Sum, inquit: sed Dómino meo Jesu Christo grátias ago, qui me, intercessióne Seráphiæ famulæ suæ, e dæmonum potestate liberávit. Quam varie tentatam ut propositum mutaret, cum a fidei constántia movére non posset, præfectus, pronuntiáta senténtia quod deos contemneret, cápitis damnávit. Ejus corpus a Christiánis in eodem sepúlcro cónditum est, in quo ipsa magistram fidei suæ Seráphiam posúerat.\par
\switchcolumn
\EN
\noindent Sabina was a Roman lady, the wife of a distinguished nobleman named Valentine. The Christian faith was taught to her by a maiden named Seraphia. After the martyrdom of this holy virgin, Sabina gathered together her reliques, and buried them with godly service. For this cause she was in a little while arrested, under the Emperor Hadrian, and brought before the Judge Elpidius. Art thou, said he, the same Sabina who is so distinguished for her blood and for her marriage She answered “I am, but I give thanks to my Lord Jesus Christ for having delivered me through the prayers of His hand-maiden Seraphia from the troubling of devils.” Diverse attempts were made to make her change her mind, but when they proved in vain the Praefect passed sentence of death upon her for despising the gods. The Christians laid her body in the same grave in which she had herself laid that of Seraphia, her teacher in the faith.\par
\switchcolumn*

\LA
\TeDeum{Te Deum}
\switchcolumn
\EN
\TeDeum{Te Deum}
\switchcolumn*

\end{paracol}

\par\needspace{10\baselineskip}\addvspace{1.2em}{\centering\small\textsc{Die 30 Augusti} · \textsc{30 August}\par}
\Office{S. Rosæ a Sancta Maria Limanæ Virginis}{St. Rose of Lima, Virgin}{Duplex}
\addcontentsline{toc}{subsection}{Die 30 Augusti · S. Rosæ a Sancta Maria Limanæ Virginis \textit{— St. Rose of Lima, Virgin}}
\begin{paracol}{2}
\LA
\RefLine{Lectiones I–III: de Scriptura occurrente.}
\switchcolumn
\EN
\RefLine{Lessons I–III: from the Scripture of the day.}
\switchcolumn*

\LA
\RefLine{Lectiones VII–VIII: ex Commúni Virginum tantum (C6a).}
\switchcolumn
\EN
\RefLine{Lessons VII–VIII: from the Common of a Virgin not a Martyr (C6a).}
\switchcolumn*

\LA
\RefLine{Lectio IX: de Ss. Felicis et Adaucti Martyrum (08-30o).}
\switchcolumn
\EN
\RefLine{Lesson IX: of St. Felix and Adauctus, Martyrs (08-30o).}
\switchcolumn*

\LA
\LessonHead{Lectio IV}
\switchcolumn
\EN
\LessonHead{Lesson IV}
\switchcolumn*

\LA
\noindent Primus Americæ Meridionalis flos sanctitatis, virgo Rosa, christianis parentibus Limæ progenita, mox ab incunabulis claruit futuræ sanctimoniæ indiciis. Nam vultus infantis mirabiliter in rosæ effugiem transfiguratus, huic nomini occasionem dedit: cui postea Virgo Deipara cognomen adjecit, jubens vocari deinceps Rosam a sancta Maria. Quinquennis votum perpetuæ virginitatis emisit. Adultior, ne a parentibus ad nuptias cogeretur, clam sibimet venustissimam capitis cæsariem præscidit. Jejuniis supra humánum modum addicta, integras Quadragesimas transegit, pane abstinens, ac dietim solis quinque granulis mali citrini victitans.\par
\switchcolumn
\EN
\noindent The first flower of holiness which came to full blossom in South America, was the maiden Rose. She was born at Lima, of a Christian father and mother, (upon the 20th of April, in the year 1586,) and was remarkable from her childhood for marks of saintliness. The occasion of her name was a strange likeness to a rose, which her face assumed when she was a babe. To this name she afterwards added that of the Virgin Mother of God, desiring to be called St. Mary's Rose. At the age of fifteen years she uttered a vow of perpetual virginity. As she grew older, lest her parents should force her to marry, she polled her head of all her hair, which was very beautiful. She fasted to a degree almost superhuman, passing whole Lents without taking bread, and eating day by day only five pips of a lime.\par
\switchcolumn*

\LA
\begin{Resp}\Rsign Propter veritátem, et mansuetúdinem, et justítiam: \Star{} Et dedúcet te mirabíliter déxtera tua.\end{Resp}
\begin{Resp}\Vsign Spécie tua et pulchritúdine tua inténde, próspere procéde, et regna.\end{Resp}
\begin{Resp}\Rsign Et dedúcet te mirabíliter déxtera tua.\end{Resp}
\switchcolumn
\EN
\begin{Resp}\Rsign Because of truth, and meekness, and righteousness; \Star{} And thy right hand shall lead thee wonderfully.\end{Resp}
\begin{Resp}\Vsign In thy comeliness, and thy beauty, go forward, fare prosperously, and reign.\end{Resp}
\begin{Resp}\Rsign And thy right hand shall lead thee wonderfully.\end{Resp}
\switchcolumn*

\LA
\LessonHead{Lectio V}
\switchcolumn
\EN
\LessonHead{Lesson V}
\switchcolumn*

\LA
\noindent Habitu tertii ordinis sancti Dominici assumpto, pristinas vitæ austeritates duplicavit: oblongo asperrimoque cilicio sparsim mimisculas acus innexuit: sub velo coronam densis aculeis introrsus obarmatam interdiu noctuque gestavit. Sanctæ Catharinæ Senensis ardua premens vestigia, catena ferrea, triplici nexu circumducta, lumbos cinxit. Lectulum sibi e truncis nodosis composuit, horumque vacuas commissuras fragminibus testarum implevit. Cellulam sibi angustissimam struxit in extremo horti angulo, ubi cælestium contemplationi dedita, crebris disciplinis, inedia, vigiliis corpusculum extenuans, at spiritu vegetata, larvas dæmonum, frequenti certamine victrix, impavide protrivit ac superavit.\par
\switchcolumn
\EN
\noindent She took the habit of the Third Order of St. Dominic, and then doubled her former severities. She wore a long and very rough hair-cloth, into which she inserted small pins. She wore day and night under her veil a crown, the inner side of which was armed with pricks. In imitation of the hard steps of St Katharine of Sienna, she girded her loins with a threefold iron chain. She made to herself a bed of knotty sticks, and filled the gaps with broken bits of potsherd. She built herself a very small hut in the farthest corner of the garden, where she gave herself up to thoughts of heavenly things, and to punishing her body with often scourging, starvation, and sleeplessness. But she waxed strong in spirit, and though she often had to fight with evil ghosts, she conquered them, fearlessly prostrated them, and triumphed over them.\par
\switchcolumn*

\LA
\begin{Resp}\Rsign Dilexísti justítiam, et odísti iniquitátem: \Star{} Proptérea unxit te Deus, Deus tuus, óleo lætítiæ.\end{Resp}
\begin{Resp}\Vsign Propter veritátem, et mansuetúdinem, et justítiam.\end{Resp}
\begin{Resp}\Rsign Proptérea unxit te Deus, Deus tuus, óleo lætítiæ.\end{Resp}
\switchcolumn
\EN
\begin{Resp}\Rsign Thou hast loved righteousness, and hated iniquity; \Star{} Therefore God, thy God, hath anointed thee with the oil of gladness.\end{Resp}
\begin{Resp}\Vsign Because of truth, and meekness, and righteousness.\end{Resp}
\begin{Resp}\Rsign Therefore God, thy God, hath anointed thee with the oil of gladness.\end{Resp}
\switchcolumn*

\LA
\LessonHead{Lectio VI}
\switchcolumn
\EN
\LessonHead{Lesson VI}
\switchcolumn*

\LA
\noindent Aegritudinum tormentis, domesticorum insultibus, linguarum morsibus dire agitata, nondum satis pro merito se affligi quærebatur. Per quindecim annos ad plusculas horas desolatione spiritus et ariditate miserrime contabescens, forti animo tulit agones omni morte amariores. Exinde cœpit supernis abundare deliciis, illustrari visionibus, colliquescere seraphicis ardoribus. Angelo tutelori, sanctæ Catharinæ Senensi, Virgini Deiparæ inter assiduas apparitiones mire familiaris, a Christo has voces audire meruit: Rosa cordis mei, tu mihi sponsa esto. Denique Sponsi hujus paradiso feliciter invectam, plurimisque ante et post obitum miraculis coruscam, Clemens decimus Pontifex Maximus sanctarum Virginum catalogo ritu solemni adscripsit.\par
\switchcolumn
\EN
\noindent She suffered greatly from painful illnesses, from the maltreatment of the servants, and from slanderous accusations, but still complained that she did not suffer as much as she deserved. For fifteen years she pined in misery from desolation and dryness of spirit, bravely enduring torments worse than any form of death. After this period she began to overflow with consolation, to be enlightened by visions, and to melt with love like a Seraph's. She attained, by the frequency of visions, to a strange personal familiarity with her Guardian Angel, with St. Katharine of Sienna, and with the Virgin Mother of God, and she earned from Christ the words, Rose of My Heart, be thou My bride. She was famous for many miracles, both before and after she departed hence, and was happily transplanted into the Bridegroom's garden, (upon the 24th of August 1617, being aged 31 years.) Pope Clement X. with solemn pomp inscribed her name in the list of holy maidens.\par
\switchcolumn*

\LA
\begin{Resp}\Rsign Afferéntur Regi vírgines post eam, próximæ ejus; \Star{} Afferéntur tibi in lætítia et exsultatióne.\end{Resp}
\begin{Resp}\Vsign Spécie tua et pulchritúdine tua; inténde, próspere procéde, et regna.\end{Resp}
\begin{Resp}\Rsign Afferéntur tibi in lætítia et exsultatióne.\end{Resp}
\begin{Resp}\Vsign Glória Patri, et Fílio, \Star{} et Spirítui Sancto.\end{Resp}
\begin{Resp}\Rsign Afferéntur tibi in lætítia et exsultatióne.\end{Resp}
\switchcolumn
\EN
\begin{Resp}\Rsign After her shall virgins be brought unto the King, her fellows \Star{} They shall be brought unto thee with gladness and rejoicing.\end{Resp}
\begin{Resp}\Vsign In thy comeliness and thy beauty, go forward, fare prosperously, and reign.\end{Resp}
\begin{Resp}\Rsign They shall be brought unto thee with gladness and rejoicing.\end{Resp}
\begin{Resp}\Vsign Glory be to the Father, and to the Son, \Star{} and to the Holy Ghost.\end{Resp}
\begin{Resp}\Rsign They shall be brought unto thee with gladness and rejoicing.\end{Resp}
\switchcolumn*

\end{paracol}

\par\needspace{10\baselineskip}\addvspace{1.2em}{\centering\small\textsc{Die 30 Augusti} · \textsc{30 August}\par}
\Office{Ss. Felicis et Adaucti Martyrum}{St. Felix and Adauctus, Martyrs}{Simplex}
\addcontentsline{toc}{subsection}{Die 30 Augusti · Ss. Felicis et Adaucti Martyrum \textit{— St. Felix and Adauctus, Martyrs}}
\begin{paracol}{2}
\LA
\LessonHead{Lectio IX (pro commemoratione)}
\switchcolumn
\EN
\LessonHead{Lesson IX (when commemorated)}
\switchcolumn*

\LA
\LessonTitle{Commemoratio: Ss. Felicis et Adaucti Martyrum}
\switchcolumn
\EN
\LessonTitle{Commemoration of: St. Felix and Adauctus, Martyrs}
\switchcolumn*

\LA
\noindent Felix, Diocletiano et Maximiano imperatoribus, propter susceptam Christi religionem comprehensus, in Serapidis templum adductus est. Cui sacrificare cum juberetur, os simulacri conspuit. Quo facto, statim serea statua corruit. Quod cum iterum ac tertio in sede Mercurii Dianæque factum esset, impietatis et magicæ artis accusatus, equuleo torquetur. Mox ad secundum ab Urbe lapidem via Ostiensi ducitur, ut securi feriretur. Cui inter viam oblatus quidam Christianus, cum Felicem agnoscens ad martyrium duci videret, Ego quoque, clara voce inquit, eadem, qua iste, lege vivo ego eumdem Jesum Christum colo. Itaque Felicem osculatus, cum eo securi percutitur, tertio Kalendas Septembris. Cujus nomen cum ignotum esset Christianis, is Adaucti nomine nobilitatus est, quod sancto Martyri Felici adauctus sit ad coronam\par
\switchcolumn
\EN
\noindent Felix was arrested in the reign of the Emperors Diocletian and Maximian, on the charge of having embraced the Christian Faith, and was brought to the temple of Serapis. When he was ordered to offer sacrifice, he spat in the face of the brazen idol, which thereupon fell down. When this happened a second and third time in the temples of Mercury and Diana, he was accused of impiety and magic, and tortured upon the rack. It was not long, however, before he was led out to the second mile-stone upon the road to Ostia, to be smitten with the axe. As they were on the way thither, they chanced to meet a certain Christian, who, when he knew that Felix was going to finish his testimony, said aloud, I live by the same law as he doth I worship the same Christ Jesus. And therewith he kissed Felix, and they were beheaded together, upon the 30th day of August. What the name of the second person was the Christians never knew, and he is therefore honoured under the title of Him-who-was-added x that is, added to the company of the Holy Martyr Felix in winning of the crown.\par
\switchcolumn*

\LA
\TeDeum{Te Deum}
\switchcolumn
\EN
\TeDeum{Te Deum}
\switchcolumn*

\end{paracol}

\par\needspace{10\baselineskip}\addvspace{1.2em}{\centering\small\textsc{Die 31 Augusti} · \textsc{31 August}\par}
\Office{S. Raymundi Nonnati Confessoris}{St. Raymond Nonnatus, Confessor}{Duplex}
\addcontentsline{toc}{subsection}{Die 31 Augusti · S. Raymundi Nonnati Confessoris \textit{— St. Raymond Nonnatus, Confessor}}
\begin{paracol}{2}
\LA
\RefLine{Lectiones I–III: de Scriptura occurrente.}
\switchcolumn
\EN
\RefLine{Lessons I–III: from the Scripture of the day.}
\switchcolumn*

\LA
\RefLine{Lectiones VII–IX: ex Commúni Confessoris non pontificis (C5).}
\switchcolumn
\EN
\RefLine{Lessons VII–IX: from the Common of a Confessor not a Bishop (C5).}
\switchcolumn*

\LA
\LessonHead{Lectio IV}
\switchcolumn
\EN
\LessonHead{Lesson IV}
\switchcolumn*

\LA
\noindent Raymundus, Nonnatus cognomento dictus, quia præter communem naturæ legem e mortuæ matris dissecto latere in lucem eductus fuit, Portelli in Catalaunia piis et nobilibus parentibus ortus, ab ipsa infantia futuræ sanctitatis indicia dedit. Nam puerilia oblectamenta, mundique illecebras respuens, ita pietati operam dabat, ut omnes in puero adultam virtutem admirarentur. Crescente vero ætate, litterarum studiis incubuit: sed mox jubente patre vitam ruri agens, sacellum sancti Nicolai in Portelli finibus situm crebro adibat, ut sacram Deiparæ imaginem, quæ in eo summa fidelium veneratione étiam nunc colitur, visitaret. Ibi effusus in preces, ipsam Dei parentem, ut se in filium adoptare viamque salutis ac scientiam sanctorum edocere dignaretur, enixe deprecabatur.\par
\switchcolumn
\EN
\noindent This Raymond is commonly called the Unborn, because his was one of the rare cases in which the child is not brought into the world in the course of nature, but by a surgical operation after the death of the mother. He was the son of godly and noble parents, at Portel, (in the dioecese of Urgel,) in Catalonia. The tokens of his holy after-life appeared even in his childhood. The things that delight children, and the attractions of the world, had no charm for him. He was so earnest in godliness that all men marvelled at his habits of premature old age. As he grew older, he gave himself to the study of letters, but, at the command of his father, turned to farming. He went often to the Chapel of St Nicolas, in the suburbs of Portel, to visit the sacred image of the Mother of God, which is still sought with great tenderness by the faithful. There he poured forth his soul in prayer, and earnestly entreated the Mother of God herself to be pleased to take him for her son, to show him the way wherein it should be safe for him to walk, and to teach him the science of the Saints.\par
\switchcolumn*

\LA
\begin{Resp}\Rsign Honéstum fecit illum Dóminus, et custodívit eum ab inimícis, et a seductóribus tutávit illum: \Star{} Et dedit illi claritátem ætérnam.\end{Resp}
\begin{Resp}\Vsign Justum dedúxit Dóminus per vias rectas, et osténdit illi regnum Dei.\end{Resp}
\begin{Resp}\Rsign Et dedit illi claritátem ætérnam.\end{Resp}
\switchcolumn
\EN
\begin{Resp}\Rsign The Lord made him honourable, and defended him from his enemies, and kept him safe from those that lay in wait for him; \Star{} And gave him perpetual glory.\end{Resp}
\begin{Resp}\Vsign He went down with him into the pit, and left him not in bonds.\end{Resp}
\begin{Resp}\Rsign And gave him perpetual glory.\end{Resp}
\switchcolumn*

\LA
\LessonHead{Lectio V}
\switchcolumn
\EN
\LessonHead{Lesson V}
\switchcolumn*

\LA
\noindent Nec defuit votis ejus benignissima Virgo. Ab ipsa enim intellexit gratissimum sibi fore, si religionem sub titulo de Mercede, seu de Misericordia redemptionis captivorum, ea suggerante nuper fundatam, ingrederetur. Qua monitione percepta, Barcinonem statim profectus, illud tam præcellentis erga proximum caritatis institutum amplexus est. Regulari igitur militiæ adscriptus, virginitatem, quam pridem beatæ Virgini consecraverat, perpetuo coluit, ceterisque virtutibus enituit, caritate præsertim erga Christianos, qui sub potestate paganorum miseram in captivitate vitam degebant. Hos ut redimeret, in Africam missus, cum jam multos a servitute liberasset, ne, consumpta pecunia, aliis item in proximo abnegandæ fidei discrimine constitutis deesset, se ipsum pignori dedit; sed cum ardentissimo salutis animarum desiderio succensus, plures Mahometanos suis concionibus ad Christum converteret, in arctam custodiam a barbaris conjectus, variisque suppliciis cruciatus, mox labiis perforatis et sera ferrea clausis, crudele martyrium diu sustinuit.\par
\switchcolumn
\EN
\noindent And the most gracious Maiden was not deaf to his prayers. From her he understood that it would please her right well, if he would join the Religious Order which had just been founded at her own inspiration, styled of Ransom or of Mercy, for buying up and freeing slaves. As soon as he had received this intimation from her, he went to Barcelona, and entered the Institute so nobly dedicated to love for our neighbour. Once enlisted in the Regular Army, he guarded unspotted for ever the virginity which he had already consecrated to the Blessed Virgin for ever. But he was a bright and shining light of all other good words and works, especially of tender compassion for Christians who were passing a life of grievous bondage in the possession of unbelieving masters. To free such he was sent into Africa, and delivered many. But his money ran short, and as there were still many in imminent danger of denying the faith, he pawned himself. He was enkindled with a most vehement longing for the salvation of souls, and by his exhortations brought diverse Mohammedans to Christ. The Moors therefore threw him into close prison, and put him to diverse tortures, at last making holes through his lips and locking them together with an iron padlock, which horrid cruelty he long bore.\par
\switchcolumn*

\LA
\begin{Resp}\Rsign Amávit eum Dóminus, et ornávit eum: stolam glóriæ índuit eum, \Star{} Et ad portas paradísi coronávit eum.\end{Resp}
\begin{Resp}\Vsign Índuit eum Dóminus lorícam fídei, et ornávit eum.\end{Resp}
\begin{Resp}\Rsign Et ad portas paradísi coronávit eum.\end{Resp}
\switchcolumn
\EN
\begin{Resp}\Rsign The Lord loved him and beautified him; He clothed him with a robe of glory, \Star{} And crowned him at the gates of Paradise.\end{Resp}
\begin{Resp}\Vsign The Lord hath put on him the breast-plate of faith, and hath adorned him.\end{Resp}
\begin{Resp}\Rsign And crowned him at the gates of Paradise.\end{Resp}
\switchcolumn*

\LA
\LessonHead{Lectio VI}
\switchcolumn
\EN
\LessonHead{Lesson VI}
\switchcolumn*

\LA
\noindent Ob hæc et alia fortiter gesta sanctitatis ejus fama longe lateque diffúsa est. Qua permotus Gregorius nonus, in amplissimum sanctæ Romanæ Ecclesiæ Cardinalium collegium Raymundum adscripsit: sed vir Dei in ea dignitate ab omni pompa abhorrens, religiosæ humilitatis tenacissimus semper fuit. Romam vero pergens, statim ac Cardonam pervenit, extremo morbo confectus, ecclesiasticis sacramentis muniri summis precibus postulavit. Cumque morbus ingravesceret, et sacerdos diutius tardaret, Angelorum ministerio, sub specie religiosorum sui ordinis apparentium, salutari viatico refectus fuit. Quo sumpto, et gratiis Deo peractis, migravit ad Dominum Dominica ultima Augusti, anno millesimo ducentesimo quadragesimo. Mortui corpus, cum circa locum sepulturæ contentio orta esset, arcæ inclusum, et mulæ cæcæ impositum, ad sacellum sancti Nicolai Dei nutu delatum fuit, ut ibi tumularetur, ubi prima jecerat sanctioris vitæ fundamenta. Illic constructo sui ordinis cænobio, a confluentibus voti causa ex universa Catalaunia fidelibus populis honoratur, variis miraculis et signis gloriosus.\par
\switchcolumn
\EN
\noindent The account of these, and other brave things that he did, he got the name of a Saint far and wide. Gregory IX. was moved thereby to make Raymond a Cardinal of the Holy Roman Church, but in this place of honour the man of God shrank from all outward show, and clung ever tightly to the lowliness that beseemeth a Religious man. He had started for Rome, (in obedience to the command of the Pope,) but had only got as far as Cardona, (six miles from Barcelona,) when he was seized with his last illness, and earnestly called for the strengthening Sacraments of the Church. But his position became critical, and the Priest had not arrived. Then Angels came unto him, clad in the habit of his own Order, and ministered unto him the wholesome Provision for the last journey. When he had taken It, he gave God thanks, and departed hence to be ever with the Lord. It was the last Lord's Day in August, 1240. After his death there was some dispute arose as to where his body should be buried so they shut it up in a box, and laid it upon a blind mule, and the beast was guided by God to carry it to the chapel of St. Nicolas, that he might be buried where he had laid the foundations of his nobler life. There was built there a Convent of his Order, and the faithful come together thither from all parts of Catalonia to honour him, and he is famous for diverse signs and wonders.\par
\switchcolumn*

\LA
\begin{Resp}\Rsign Iste homo perfécit ómnia quæ locútus est ei Deus, et dixit ad eum: Ingrédere in réquiem meam: \Star{} Quia te vidi justum coram me ex ómnibus géntibus.\end{Resp}
\begin{Resp}\Vsign Iste est, qui contémpsit vitam mundi, et pervénit ad cæléstia regna.\end{Resp}
\begin{Resp}\Rsign Quia te vidi justum coram me ex ómnibus géntibus.\end{Resp}
\begin{Resp}\Vsign Glória Patri, et Fílio, \Star{} et Spirítui Sancto.\end{Resp}
\begin{Resp}\Rsign Quia te vidi justum coram me ex ómnibus géntibus.\end{Resp}
\switchcolumn
\EN
\begin{Resp}\Rsign This is he which did according unto all that God commanded him; and God said unto him: Enter thou into My rest, \Star{} For thee have I seen righteous before Me among all people.\end{Resp}
\begin{Resp}\Vsign This is he which loved not his life in this world, and is come unto an everlasting kingdom.\end{Resp}
\begin{Resp}\Rsign For thee have I seen righteous before Me among all people.\end{Resp}
\begin{Resp}\Vsign Glory be to the Father, and to the Son, \Star{} and to the Holy Ghost.\end{Resp}
\begin{Resp}\Rsign For thee have I seen righteous before Me among all people.\end{Resp}
\switchcolumn*

\end{paracol}

\SeasonTitle{Mensis September}{September}

\addcontentsline{toc}{section}{Mensis September \textit{— September}}

\par\needspace{10\baselineskip}\addvspace{1.2em}{\centering\small\textsc{Die 1 Septembris} · \textsc{1 September}\par}
\Office{S. Ægidii Abbatis}{St. Giles, Abbot}{Simplex}
\addcontentsline{toc}{subsection}{Die 1 Septembris · S. Ægidii Abbatis \textit{— St. Giles, Abbot}}
\begin{paracol}{2}
\LA
\RefLine{Lectiones I–II: de Scriptura occurrente.}
\switchcolumn
\EN
\RefLine{Lessons I–II: from the Scripture of the day.}
\switchcolumn*

\LA
\LessonHead{Lectio III (pro commemoratione)}
\switchcolumn
\EN
\LessonHead{Lesson III (when commemorated)}
\switchcolumn*

\LA
\noindent Ægidius Atheniénsis, regiæ stirpis, a prima ætate divinis litteris et caritátis officiis ita déditus fuit, ut nihil præterea curare viderétur. Itaque, paréntibus mórtuis, totum patrimónium in páuperes erogávit; quin étiam túnicam exuit, ut ægrotum egentem tégeret, qua ille indútus, statim conváluit. Sed multis deinceps clarior miraculis, timens sui nóminis celebritátem, Arelátem ad beátum Cæsárium contendit. A quo post biennium discédens, secessit in eremum; ubi diutius herbárum radicibus et cervæ lacte, quæ statis ad eum horis veniebat, admirábili sanctitáte vixit. Quæ cerva, insequéntibus quodam die cánibus regiis, cum in antrum Ægidii refugísset, Galliæ regem ímpulit, ut ab eo summis precibus péteret, ut in loco speluncæ monastérium éxstrui paterétur. Cujus administratiónem, flagitante rege, invitus suscépit; eoque munere aliquot annis prudenter piéque gesto, migrávit in cælum.\par
\switchcolumn
\EN
\noindent The holy Abbot Giles was by birth an Athenian, and of Royal lineage. From his youth he showed ever such a love for sacred learning and for works of charity, that he seemed to care for nothing else. When his father and mother were dead, he bestowed his whole inheritance upon the poor. He took off even his own coat, to clothe a poor sick man withal, and the sick man was healed forthwith as soon as he put it on him. As Giles became famous for working miracles, he fled from glory among men, and betook him to Arles, (in France,) to the company of blessed Caesarius. After the space of two years he departed thence, and went into the desert, for he lived in wonderful holiness for a long while upon the roots of herbs and the milk of an hind, which came to him at regular hours. This hind was chased one day by the King's hounds, and took refuge in Giles's cave. Thereby the King of France was moved earnestly to entreat of him that he would suffer a monastery to be built in the place where this cave was. At the instant desire of the King, he took the rule of this monastery, albeit himself unwilling, and discharged this duty wisely and godly for some years, until he passed away to heaven.\par
\switchcolumn*

\LA
\TeDeum{Te Deum}
\switchcolumn
\EN
\TeDeum{Te Deum}
\switchcolumn*

\end{paracol}

\par\needspace{10\baselineskip}\addvspace{1.2em}{\centering\small\textsc{Die 2 Septembris} · \textsc{2 September}\par}
\Office{S. Stephani Hungariæ Regis Confessoris}{St. Stephen, King of Hungary, Confessor}{Semiduplex}
\addcontentsline{toc}{subsection}{Die 2 Septembris · S. Stephani Hungariæ Regis Confessoris \textit{— St. Stephen, King of Hungary, Confessor}}
\begin{paracol}{2}
\LA
\RefLine{Lectiones I–III: de Scriptura occurrente.}
\switchcolumn
\EN
\RefLine{Lessons I–III: from the Scripture of the day.}
\switchcolumn*

\LA
\LessonHead{Lectio IV}
\switchcolumn
\EN
\LessonHead{Lesson IV}
\switchcolumn*

\LA
\noindent Stephanus in Hungariam Christi fidem et regium nomen invéxit. Regia coróna a Romano Pontifice impetrata, ejusque jussu in regem inunctus, regnum Sedi apostolicæ obtulit. Varia pietátis domicília Romæ, Jerosolymis, Constantinopoli: in Hungaria archiepiscopátum Strigoniensem, episcopátus decem admirábili religióne et munificéntia fundavit. Par in páuperes amor et liberalitas, quos véluti Christum ipsum complectens, neminem a se mærentem ac vacuum umquam dismisit; quin ad eórum inópiam sublevandam, amplíssimis facultátibus erogatis, domesticam quoque supellectilem exímia benignitate frequenter distríbuit. Suis ínsuper mánibus lavare páuperum pedes, noctu solus et ignotus nosocomía frequentare, decumbéntibus inservire ac cetera caritátis offícia exhibere consuevit. Quarum virtútum merito, illíus déxtera, resoluto cetero corpore, incorrupta permansit.\par
\switchcolumn
\EN
\noindent This Stephen (was the son of Geysa, fourth Duke of the Hungarians, and was born at Gran in the year 977.) He it was who first gave to Hungary the faith of Christ and the name of a kingdom. He obtained the Kingly crown from the Bishop of Rome, and being by command of the same anointed King, he made an offering of his kingdom to the Apostolic See. With wonderful devotion and bounty he founded diverse godly houses at Rome, Jerusalem, and Constantinople, and in Hungary the Archbishopric of Gran and ten other Sees. Toward the poor he had the same love and bounty. He greeted them as though they were Christ Himself, and never sent any one away sorrowing and empty. He spent vast sums in relieving their poverty, and also often parted among them with exceeding tenderness even the furniture of his house. Moreover it was his use to wash the feet of the poor with his own hands, and to go in the night, alone and unknown, to the hospitals, and to wait on them that lay there, and show them other deeds of kindness. It was the reward of these good works that, when the rest of his body decayed, his right hand remained uncorrupt.\par
\switchcolumn*

\LA
\begin{Resp}\Rsign Honéstum fecit illum Dóminus, et custodívit eum ab inimícis, et a seductóribus tutávit illum: \Star{} Et dedit illi claritátem ætérnam.\end{Resp}
\begin{Resp}\Vsign Justum dedúxit Dóminus per vias rectas, et osténdit illi regnum Dei.\end{Resp}
\begin{Resp}\Rsign Et dedit illi claritátem ætérnam.\end{Resp}
\switchcolumn
\EN
\begin{Resp}\Rsign The Lord made him honourable, and defended him from his enemies, and kept him safe from those that lay in wait for him; \Star{} And gave him perpetual glory.\end{Resp}
\begin{Resp}\Vsign He went down with him into the pit, and left him not in bonds.\end{Resp}
\begin{Resp}\Rsign And gave him perpetual glory.\end{Resp}
\switchcolumn*

\LA
\LessonHead{Lectio V}
\switchcolumn
\EN
\LessonHead{Lesson V}
\switchcolumn*

\LA
\noindent Orándi studio noctes pene totas ducebat insomnes; atque in cæléstium rerum contemplatióne defixus, interdum extra sensus raptus, sublimis in áëra ferri visus fuit. Perduellium conspiratiónes ac validiórum hostium ímpetus, miro prorsus modo, non semel oratiónis præsidio evitavit. Susceptum ex Ghisella Bavárica, sancti Henrici imperatóris sorore, quam sibi matrimonio junxerat, Emerícum fílium tanta morum disciplína talíque pietáte enutrívit, quantum ejus póstea sanctitas declaravit. Regni vero negotia ita dispósuit, ut, accitis undique prudentíssimis et sanctíssimis viris, nihil umquam sine illórum consílio molirétur; humillimis interim precibus in cínere et cilício Deum déprecans, ut univérsum Hungariæ regnum, antequam e vita migraret, catholicum vidére mererétur; vere, propter ingens dilatandæ fidei studium, illíus gentis Apóstolus nuncupatus, facta a Romano Pontifice ipsi posterisque régibus præferendæ crucis potestate.\par
\switchcolumn
\EN
\noindent He passed almost whole nights in earnest prayer, and when totally rapt in the thought of heavenly things, he sometimes became beside himself, and was seen to rise off the ground into the air. In more than one instance he strangely escaped through the power of prayer from rebellion, treason, and the onslaughts of mighty foes. He married Gisela of Bavaria, sister to the holy Emperor Henry, and begat on her Emeric, whom he trained up in such manners and godliness, as are shown by his also becoming a Saint. To carry on the business of his kingdom, he gathered together from all quarters the most learned and godly men, and took nothing in hand without their advice. Meanwhile he entreated of God by the most lowly supplications, offered up in sack-cloth and ashes, that, before he departed this life, he might see all Hungary Catholic. On account of his excellent zeal for the spread of the Faith he is called the Apostle of that nation, and the Bishop of Rome gave to him and to his successors the right to have a Cross carried before them.\par
\switchcolumn*

\LA
\begin{Resp}\Rsign Amávit eum Dóminus, et ornávit eum: stolam glóriæ índuit eum, \Star{} Et ad portas paradísi coronávit eum.\end{Resp}
\begin{Resp}\Vsign Índuit eum Dóminus lorícam fídei, et ornávit eum.\end{Resp}
\begin{Resp}\Rsign Et ad portas paradísi coronávit eum.\end{Resp}
\switchcolumn
\EN
\begin{Resp}\Rsign The Lord loved him and beautified him; He clothed him with a robe of glory, \Star{} And crowned him at the gates of Paradise.\end{Resp}
\begin{Resp}\Vsign The Lord hath put on him the breast-plate of faith, and hath adorned him.\end{Resp}
\begin{Resp}\Rsign And crowned him at the gates of Paradise.\end{Resp}
\switchcolumn*

\LA
\LessonHead{Lectio VI}
\switchcolumn
\EN
\LessonHead{Lesson VI}
\switchcolumn*

\LA
\noindent Dei Genitricem, quam ardentíssime venerabátur, amplíssimo in ejus honórem constructo templo, Hungariæ patrónam instítuit; ab eádem vicíssim Vírgine recéptus in cælum ipso suæ Assumptiónis die, quem Húngari e sancti regis instituto Magnæ Dominæ diem appellant. Sacrum ejus corpus, suavíssimo fragrans odore, liquore cælésti scatens, inter multa et varia miracula, Romani Pontificis jussu, nobiliórem in locum translátum est atque honorificentius cónditum. Ejus autem festum Innocentius undecimus Pontifex maximus, quarto Nonas Septembris, ob insignem victoriam ab exercitu Leopoldi primi, Romanórum elécti imperatóris et Hungariæ regis, eádem die in Budæ expugnatióne, ope divina, e Turcis reportatam, celebrándum instituit.\par
\switchcolumn
\EN
\noindent He had a burning zeal to honour the Mother of God. He built a very great Church in her honour, and made her Patroness of Hungary. In return, the same Virgin received him into heaven, (in the year 1038,) upon the day of her own Assumption, which the Hungarians, by the example of the holy King, call “the Great Lady's Day.” His hallowed body yielded the sweetest savour, and reeked with an heavenly liquid, and amid many and diverse wonders it was removed by command of the Bishop of Rome into a more noble place, and more honourably buried. Pope Innocent XI. ordered his Feast to be held upon the 2nd day of September, on account of the famous victory over the Turks which was gained upon this day, (in the year 1686,) when the army of Leopold I., Emperor (elect) of the Romans, and King of Hungary, wrested from them, by the help of God, the city of Buda.\par
\switchcolumn*

\LA
\begin{Resp}\Rsign Iste homo perfécit ómnia quæ locútus est ei Deus, et dixit ad eum: Ingrédere in réquiem meam: \Star{} Quia te vidi justum coram me ex ómnibus géntibus.\end{Resp}
\begin{Resp}\Vsign Iste est, qui contémpsit vitam mundi, et pervénit ad cæléstia regna.\end{Resp}
\begin{Resp}\Rsign Quia te vidi justum coram me ex ómnibus géntibus.\end{Resp}
\begin{Resp}\Vsign Glória Patri, et Fílio, \Star{} et Spirítui Sancto.\end{Resp}
\begin{Resp}\Rsign Quia te vidi justum coram me ex ómnibus géntibus.\end{Resp}
\switchcolumn
\EN
\begin{Resp}\Rsign This is he which did according unto all that God commanded him; and God said unto him: Enter thou into My rest, \Star{} For thee have I seen righteous before Me among all people.\end{Resp}
\begin{Resp}\Vsign This is he which loved not his life in this world, and is come unto an everlasting kingdom.\end{Resp}
\begin{Resp}\Rsign For thee have I seen righteous before Me among all people.\end{Resp}
\begin{Resp}\Vsign Glory be to the Father, and to the Son, \Star{} and to the Holy Ghost.\end{Resp}
\begin{Resp}\Rsign For thee have I seen righteous before Me among all people.\end{Resp}
\switchcolumn*

\LA
\LessonHead{Lectio VII}
\switchcolumn
\EN
\LessonHead{Lesson VII}
\switchcolumn*

\LA
\LessonTitle{Léctio sancti Evangélii secúndum Lucam}
\switchcolumn
\EN
\LessonTitle{From the Holy Gospel according to Luke}
\switchcolumn*

\LA
\Cite{Luc 19:12-26}
\switchcolumn
\EN
\Cite{Luke 19:12-25}
\switchcolumn*

\LA
\noindent In illo témpore: Dixit Jesus discípulis suis parábolam hanc: Homo quidam nóbilis ábiit in regiónem longinquam, accípere sibi regnum, et revérti. Et réliqua.\par
\switchcolumn
\EN
\noindent At that time, Jesus spoke this parable unto His disciples: A certain nobleman went into a far country to receive for himself a kingdom and to return. And so on.\par
\switchcolumn*

\LA
\LessonTitle{Homilía sancti Ambrósii Epíscopi}
\switchcolumn
\EN
\LessonTitle{Homily by St. Ambrose, Bishop (of Milan.)}
\switchcolumn*

\LA
\Source{Lib. 8. in Lucam.}
\switchcolumn
\EN
\Source{Book viii. on Luke.}
\switchcolumn*

\LA
\noindent Bonus ordo, ut vocatúrus gentes, et Judǽos jussúrus intérfici, qui noluérunt regnáre supra se Christum, hanc præmítteret comparatiónem, ne dicerétur: Nihil déderat pópulo Judæórum, unde póterat mélior fíeri? ut quid ab eo qui nihil recépit, exígitur? Non medíocris ista est mna, quam supra múlier evangélica quia non invénit, lucérnam accéndit, lúmine quærit admóto, gratulátur invéntam.\par
\switchcolumn
\EN
\noindent It is well ordered that, being about to call the Gentiles, and to command the destruction of those Jews, who would not have Christ to reign over them, He should put forth first this parable; lest it should be said He had given the Jews no means of becoming better. How can they be asked to repay who have received nothing? That is not a piece of silver of little worth, which, when the woman before mentioned in this Gospel (xv. 8) hath lost, she lighteth a candle, and sweepeth the house, and searcheth diligently until she findeth it.\par
\switchcolumn*

\LA
\begin{Resp}\Rsign Iste est qui ante Deum magnas virtútes operátus est, et de omni corde suo laudávit Dóminum: \Star{} Ipse intercédat pro peccátis ómnium populórum.\end{Resp}
\begin{Resp}\Vsign Ecce homo sine queréla, verus Dei cultor, ábstinens se ab omni ópere malo, et pérmanens in innocéntia sua.\end{Resp}
\begin{Resp}\Rsign Ipse intercédat pro peccátis ómnium populórum.\end{Resp}
\switchcolumn
\EN
\begin{Resp}\Rsign This is he which wrought great wonders before God, and praised the Lord with all his heart. \Star{} May he pray for all people, that their sins may be forgiven unto them.\end{Resp}
\begin{Resp}\Vsign Behold a man without blame, a worshipper of God in truth, keeping himself clean from every evil work, and abiding still in his innocency.\end{Resp}
\begin{Resp}\Rsign May he pray for all people, that their sins may be forgiven unto them.\end{Resp}
\switchcolumn*

\LA
\LessonHead{Lectio VIII}
\switchcolumn
\EN
\LessonHead{Lesson VIII}
\switchcolumn*

\LA
\noindent Dénique ex una decem mnas álius fecit, álius quinque. Fortásse iste morália habet, quia quinque sunt córporis sensus: ille duplícia, id est, mýstica legis, et morália probitátis. Unde et Matthǽus quinque talénta, et duo talénta pósuit: in quinque taléntis ut sint morália, in duóbus utrúmque, mýsticum atque morále. Ita quod número inférius, re ubérius.\par
\switchcolumn
\EN
\noindent A single pound one gained ten and another five pounds. Perchance by him which had the five pounds is signified he which practiseth well, since the body hath five senses, and by him which had the ten, that is, double the other, he which is learned and orthodox in the deep things of doctrine, as well as upright in his practical life. Hence also in Matthew we have five talents and two talents the five talents signifying good practice, and the two talents precept and practice together. So that that which counteth as the greater number is but a fraction of the lesser number.\par
\switchcolumn*

\LA
\begin{Resp}\Rsign Sint lumbi vestri præcíncti, et lucérnæ ardéntes in mánibus vestris: \Star{} Et vos símiles homínibus exspectántibus dóminum suum, quando revertátur a núptiis.\end{Resp}
\begin{Resp}\Vsign Vigiláte ergo, quia nescítis qua hora Dóminus vester ventúrus sit.\end{Resp}
\begin{Resp}\Rsign Et vos símiles homínibus exspectántibus dóminum suum, quando revertátur a núptiis.\end{Resp}
\begin{Resp}\Vsign Glória Patri, et Fílio, \Star{} et Spirítui Sancto.\end{Resp}
\begin{Resp}\Rsign Et vos símiles homínibus exspectántibus dóminum suum, quando revertátur a núptiis.\end{Resp}
\switchcolumn
\EN
\begin{Resp}\Rsign Let your loins be girded about, and your lights burning; \Star{} And ye yourselves like unto men that wait for their lord, when he will return from the wedding.\end{Resp}
\begin{Resp}\Vsign Watch therefore, for ye know not what hour your Lord doth come.\end{Resp}
\begin{Resp}\Rsign And ye yourselves like unto men that wait for their lord, when he will return from the wedding.\end{Resp}
\begin{Resp}\Vsign Glory be to the Father, and to the Son, \Star{} and to the Holy Ghost.\end{Resp}
\begin{Resp}\Rsign And ye yourselves like unto men that wait for their lord, when he will return from the wedding.\end{Resp}
\switchcolumn*

\LA
\LessonHead{Lectio IX}
\switchcolumn
\EN
\LessonHead{Lesson IX}
\switchcolumn*

\LA
\noindent Et hic póssumus decem mnas decem verba intellígere, id est, legis doctrínam; quinque autem mnas, magistéria disciplínæ. Sed legisperítum in ómnibus volo esse perféctum; non enim in sermóne, sed in virtúte est regnum Dei. Bene autem, quia de Judǽis dicit, duo soli multiplicátam pecúniam déferunt, non útique æris, sed dispensatiónis usúris. Ália est enim pecúniæ fœnébris, ália doctrínæ cæléstis usúra.\par
\switchcolumn
\EN
\noindent And here we may also understand by the ten pounds the ten words, that is, the Commandments, and by the five pounds, the enforcement of their teaching. But I would that a lawyer should be in all things perfect. For the kingdom of God is not in word but in power. (1 Cor. iv. 20.) Meet also is it, that, in speaking of Jews, Christ should represent only two as bringing in increased capital, for these talents are talents not of money but of grace, and to increase money by usury is a very different thing from improving heavenly revelation by the like means.\par
\switchcolumn*

\LA
\TeDeum{Te Deum}
\switchcolumn
\EN
\TeDeum{Te Deum}
\switchcolumn*

\end{paracol}

\par\needspace{10\baselineskip}\addvspace{1.2em}{\centering\small\textsc{Die 3 Septembris} · \textsc{3 September}\par}
\Office{S. Pii X Papæ Confessoris}{St. Pius X, Pope and Confessor}{Duplex}
\addcontentsline{toc}{subsection}{Die 3 Septembris · S. Pii X Papæ Confessoris \textit{— St. Pius X, Pope and Confessor}}
\begin{paracol}{2}
\LA
\RefLine{Lectiones I–III: de Scriptura occurrente.}
\switchcolumn
\EN
\RefLine{Lessons I–III: from the Scripture of the day.}
\switchcolumn*

\LA
\LessonHead{Lectio IV}
\switchcolumn
\EN
\LessonHead{Lesson IV}
\switchcolumn*

\LA
\noindent Pius Papa décimus, cui nomen ántea Joséphus Sarto, in Venetórum pago natus est, quem Riése vocant, paréntibus quidem humílibus, sed probitáte ac pietáte conspícuus. Inter Seminárii Patavíni alúmnos adscríptus, ita pietáte ac doctrína profécit, ut condiscípulis exémplo, moderatóribus admiratióni esset. Sacerdótio initiátus, in óppido Tómbolo primum, qua vicárius cooperátor, dein Salatiáni qua párochus, per plures annos adlaborávit; quibus in obeúndis munéribus, tanta caritátis effusióne, tanto sacerdotáli zelo et sanctitáte vitæ excélluit, ut Epíscopus Tarvisínus inter canónicos cathedrális ecclésiæ eum cooptáret, eúmque Cúriæ episcopális cancellárium simúlque Seminárii diœcesáni spirituálem moderatórem renuntiáret. Hæc offícia tam egrégie persecútus, a Leóne tértio décimo, cui erat probatíssimus, Mantuánæ ecclésiæ Antístes fuit renuntiátus.\par
\switchcolumn
\EN
\noindent Pope Pius X, whose name previously was Joseph Sarto, was born in the village of Riese in the Venetian province, to humble parents remarkable for their godliness and piety. He enrolled among the students in the seminary of Padua, where he exhibited such piety and learning that he was both an example to his fellow students and the admiration of his teachers. Upon his ordination to the priesthood, he labored for several years first as curate in the town of Tombolo, then as pastor at Salzano. He applied himself to his duties with such a constant flow of charity and such priestly zeal, and was so distinguished by the holiness of his life, that the Bishop of Treviso appointed him as a canon of the cathedral church and made him the chancellor of the bishop's curia, as well as spiritual director of the diocesan seminary. His performance in these duties was so outstanding and so highly impressed Leo XIII, that he made him bishop of the Church of Mantua.\par
\switchcolumn*

\LA
\begin{Resp}\Rsign Invéni David servum meum, óleo sancto meo unxi eum: \Star{} Manus enim mea auxiliábitur ei.\end{Resp}
\begin{Resp}\Vsign Nihil profíciet inimícus in eo, et fílius iniquitátis non nocébit ei.\end{Resp}
\begin{Resp}\Rsign Manus enim mea auxiliábitur ei.\end{Resp}
\switchcolumn
\EN
\begin{Resp}\Rsign I have found David My servant, with My holy oil have I anointed him \Star{} For My hand shall help him.\end{Resp}
\begin{Resp}\Vsign The enemy shall prevail nothing against him, nor the son of wickedness afflict him.\end{Resp}
\begin{Resp}\Rsign For My hand shall help him.\end{Resp}
\switchcolumn*

\LA
\LessonHead{Lectio V}
\switchcolumn
\EN
\LessonHead{Lesson V}
\switchcolumn*

\LA
\noindent Boni pastóris nullam partem déserens, eo máxime conténdit, ut juvéntus in sortem Dómini vocáta rite ad sacra instituerétur, piæ consociatiónes novis augéscerent increméntis, rítibus divíni cultus plus decóris ac pietátis accéderet. Præcépta quibus cívitas christiána nítitur, áltius proclamáre non désiit, et qui vitam ínopem ipse ducébat, paupéribus numquam omísit afférre levámen. Tot ígitur suffragántibus méritis, inter purpurátos Patres adléctus et Venetiárum Patriárcha creátus est. Dénique post Leónis décimi tértii óbitum, cum Patrum Cardinálium suffrágia in eum coaléscerent, cumque ipse supplicatiónibus et lácrimis tantum munus a se avértere frustra conátus esset, suasiónibus tandem cedens, « accépto in crucem », inquit, et Summi Pontificátus ápicem ut crucem a Deo sibi oblátam, demísso sed forti ánimo suscépit.\par
\switchcolumn
\EN
\noindent Lacking in nothing that makes a good pastor, he labored particularly to teach young men called to the priesthood, as well as fostering the growth of devout associations and the beauty and dignity of divine worship. He would ever affirm and promote the laws upon which Christian civilization depend, and while leading himself a life of poverty, never missed the opportunity to alleviate the burden of poverty in others. Because of his great merits, he was made a cardinal and created Patriarch of Venice. After the death of Pope Leo XIII, when the votes of the College of Cardinals began to increase in his favor, he tried in vain with supplications and tears to be relieved of so heavy a burden. Finally he ceded to their persuasions, saying I accept the cross. Thus he accepted the crown of the supreme pontificate as a cross, offering himself to God, with a resigned but stedfast spirit.\par
\switchcolumn*

\LA
\begin{Resp}\Rsign Pósui adjutórium super poténtem, et exaltávi eléctum de plebe mea: \Star{} Manus enim mea auxiliábitur ei.\end{Resp}
\begin{Resp}\Vsign Invéni David servum meum, óleo sancto meo unxi eum.\end{Resp}
\begin{Resp}\Rsign Manus enim mea auxiliábitur ei.\end{Resp}
\switchcolumn
\EN
\begin{Resp}\Rsign I have laid help upon one that is mighty, and have exalted one chosen out of My people \Star{} For My hand shall help him.\end{Resp}
\begin{Resp}\Vsign I have found David My servant, with My holy oil have I anointed him.\end{Resp}
\begin{Resp}\Rsign For My hand shall help him.\end{Resp}
\switchcolumn*

\LA
\LessonHead{Lectio VI}
\switchcolumn
\EN
\LessonHead{Lesson VI}
\switchcolumn*

\LA
\noindent In Petri cáthedra constitútus, nihil de prístina vitæ ratióne remísit. Humilitáte præsértim, simplicitáte ac paupertáte refúlsit, ita ut in suo testaménto scríbere potúerit: « Pauper natus sum, pauper vixi, pauper mori cúpio ». Humílitas vero ánimi fortitúdinem in eo alébat, cum de Dei glória, Ecclésiæ libertáte, animarúmque salúte agerétur. Vir acérrimi ingénii et propósiti tenax, inter vicésimi ineúntis sæculi procéllas, Ecclésiam fírmiter rexit, et præclaríssimis ornávit institútis. Músicam sacram ad prístinum splendórem ac dignitátem revocávit; sacrórum Bibliórum stúdiis príncipem sedem Romæ constítuit; Románam Cúriam sapiénter reformávit; leges de fidélibus per catechísmum instituéndis restítuit; Eucharísticæ mensæ crebriórem, imo et cotidiánam consuetúdinem indúxit, ejúsque accéssum púeris quoque a primo ratiónis usu apéruit; actiónis cathólicæ increménta sédulo promóvit; sólidæ cleri institutióni provídit, ádditis quoque semináriis per regiónes dispósitis; sacerdótes omnes ad interiórem vitam coléndam alléxit; leges Ecclésiæ in unum corpus redégit; erróres perniciosíssimos, modernísmi appellatióne comprehénsos, damnávit atque evéllit; civíle vétitum, quod dicunt, in Pontíficis Máximi electióne rejécit. Tandem labóribus fractus ac mæróre conféctus ob bellum Europæum tunc exórtum, die vicésima mensis Augústi anni millésimi nongentésimi décimi quarti, ad cæléste præmium evolávit. Eum ubíque terrárum sanctitátis fama clarum miraculísque fulgéntem, Pius Papa duodécimus, cuncto plaudénte orbe, in Sanctórum númerum rétulit.\par
\switchcolumn
\EN
\noindent Placed upon the chair of Peter, he gave up nothing of his former way of life. He shone especially in humility, simplicity and poverty, so that he was able to write in his last testament: I was born in poverty, I lived in poverty, and I wish to die in poverty. His humility, however, nourished his soul with strength, when it concerned the glory of God, the liberty of Holy Church, and the salvation of souls. A man of passionate temperament and of firm purpose, he ruled the Church firmly as it entered into the twentieth century, and adorned it with brilliant teachings. He restored the sacred music to its pristine glory and dignity; he established Rome as the principal centre for the study of the Holy Bible; he ordered the reform of the Roman Curia with great wisdom; he restored the laws concerning the faithful for the instruction of the catechism; he introduced the custom of more frequent and even daily reception of the Holy Eucharist, as well as permitting its reception by children as soon as they reach the age of reason; he zealously promoted the growth of Catholic action; he provided for the sound education of clerics and increased the number of seminaries in their diverse regions; he encouraged every priest in the practice of the interior life; he brought the laws of the Church together into one body; he condemned and suppressed those most pernicious errors known collectively as Modernism; he suppressed the custom of civil veto at the election of a Supreme Pontiff. Finally worn out with his labors and overcome with grief at the European war which had just begun, he went to his heavenly reward on the twentieth day of August in the year 1914. Renowned throughout all the world for the fame of his holiness and miracles, Pope Pius XII, with the approbation of the whole world, numbered him among the Saints.\par
\switchcolumn*

\LA
\begin{Resp}\Rsign Iste est, qui ante Deum magnas virtútes operátus est, et omnis terra doctrína ejus repléta est: \Star{} Ipse intercédat pro peccátis ómnium populórum.\end{Resp}
\begin{Resp}\Vsign Iste est, qui contémpsit vitam mundi, et pervénit ad cæléstia regna.\end{Resp}
\begin{Resp}\Rsign Ipse intercédat pro peccátis ómnium populórum.\end{Resp}
\begin{Resp}\Vsign Glória Patri, et Fílio, \Star{} et Spirítui Sancto.\end{Resp}
\begin{Resp}\Rsign Ipse intercédat pro peccátis ómnium populórum.\end{Resp}
\switchcolumn
\EN
\begin{Resp}\Rsign This is he which wrought great wonders before God, and the whole earth is full of his teaching \Star{} May he pray for all people, that their sins may be forgiven unto them\end{Resp}
\begin{Resp}\Vsign This is he which loved not his life in this world, and hath attained unto the kingdom of heaven.\end{Resp}
\begin{Resp}\Rsign May he pray for all people, that their sins may be forgiven unto them\end{Resp}
\begin{Resp}\Vsign Glory be to the Father, and to the Son, \Star{} and to the Holy Ghost.\end{Resp}
\begin{Resp}\Rsign May he pray for all people, that their sins may be forgiven unto them\end{Resp}
\switchcolumn*

\LA
\LessonHead{Lectio VII}
\switchcolumn
\EN
\LessonHead{Lesson VII}
\switchcolumn*

\LA
\LessonTitle{Léctio sancti Evangélii secúndum Joánnem}
\switchcolumn
\EN
\LessonTitle{From the Holy Gospel according to John}
\switchcolumn*

\LA
\Cite{Joannes 21:15-17}
\switchcolumn
\EN
\Cite{John 21:15-17}
\switchcolumn*

\LA
\noindent In illo témpore: Dixit Jesus Simóni Petro: Simon Joánnis, díligis me plus his? Et réliqua.\par
\switchcolumn
\EN
\noindent At that time: Jesus saith to Simon Peter, Simon, son of Jonas, lovest thou me more than these? And so on, and that which followeth.\par
\switchcolumn*

\LA
\LessonTitle{Homilía sancti Augustíni Epíscopi}
\switchcolumn
\EN
\LessonTitle{A Homily by St. Augustine, the Bishop}
\switchcolumn*

\LA
\Source{Tractatus 123 in Joánnem. in medio}
\switchcolumn
\EN
\Source{Tractatus 123 in Joannem. in medio}
\switchcolumn*

\LA
\noindent Rédditur negatióni trinæ trina conféssio, ne minus amóri lingua sérviat, quam timóri: et plus vocis elicuísse videátur mors ímminens, quam vita præsens. Sit amóris offícium páscere Domínicum gregem, si fuit timóris indícium negáre pastórem. Qui hoc ánimo pascunt oves Christi, ut suas velint esse, non Christi, se convincúntur amáre, non Christum: vel gloriándi vel dominándi, vel acquiréndi cupiditáte; non obediéndi, et subveniéndi, et Deo placéndi caritáte.\par
\switchcolumn
\EN
\noindent To the threefold denial there is now appended a threefold confession, that his tongue may not yield a feebler service to love than to fear, and imminent death may not appear to have elicited more from his lips than present life. Let it be the office of love to feed the Lord's flock, if it was the signal of fear to deny the Shepherd. Those who have this purpose in feeding the flock of Christ, that they may have them as their own, and not as Christ's, are convicted of loving themselves, and not Christ, from the desire either of boasting, or wielding power, or acquiring gain, and not from the love of obeying, serving and pleasing God.\par
\switchcolumn*

\LA
\begin{Resp}\Rsign Amávit eum Dóminus, et ornávit eum: stolam glóriæ índuit eum, \Star{} Et ad portas paradísi coronávit eum.\end{Resp}
\begin{Resp}\Vsign Induit eum Dóminus lorícam fídei, et ornávit eum.\end{Resp}
\begin{Resp}\Rsign Et ad portas paradísi coronávit eum.\end{Resp}
\switchcolumn
\EN
\begin{Resp}\Rsign The Lord loved him and beautified him He clothed him with a robe of glory \Star{} And crowned him at the gates of Paradise.\end{Resp}
\begin{Resp}\Vsign The Lord hath put on him the breast-plate of faith, and hath adorned him.\end{Resp}
\begin{Resp}\Rsign And crowned him at the gates of Paradise.\end{Resp}
\switchcolumn*

\LA
\LessonHead{Lectio VIII}
\switchcolumn
\EN
\LessonHead{Lesson VIII}
\switchcolumn*

\LA
\noindent Contra hos ergo vígilat tótidem inculcáta ista vox Christi, quos Apóstolos gemit sua quærere, non quæ Jesu Christi. Nam quid est áliud, si díligis me, pasce oves meas; quam si dicerétur, si me díligis, non te páscere cógita? Sed oves meas, et sicut meas pasce, non sicut tuas; glóriam meam in eis quære, non tuam; domínium meum, non tuum; lucra mea, non tua; ne sis in eórum societáte qui pértinent ad témpora periculósa, seípsos amántes, et cétera quæ huic malórum inítio connectúntur.\par
\switchcolumn
\EN
\noindent Against such, therefore, there standeth as a wakeful sentinel this thrice inculcated utterance of Christ, of whom the Apostle complaineth that they seek their own, not the things that are of Christ Jesus. For what else signifieth the words: Lovest thou me? Feed my sheep: than if it were said: If thou lovest me, think not of feeding thyself, but feed my sheep as mine, and not as thine own; seek my glory in them, and not thine own; my dominion, and not thine; my gain, and not thine; lest thou be found in the fellowship of them that belong to the perilous times, lovers of their own selves, and all else that is joined on to this beginning of evils.\par
\switchcolumn*

\LA
\begin{Resp}\Rsign Sint lumbi vestri præcíncti, et lucérnæ ardéntes in mánibus vestris: \Star{} Et vos símiles homínibus exspectántibus dóminum suum, quando revertátur a núptiis.\end{Resp}
\begin{Resp}\Vsign Vigiláte ergo, quia nescítis qua hora Dóminus vester ventúrus sit.\end{Resp}
\begin{Resp}\Rsign Et vos símiles homínibus exspectántibus dóminum suum, quando revertátur a núptiis.\end{Resp}
\begin{Resp}\Vsign Glória Patri, et Fílio, \Star{} et Spirítui Sancto.\end{Resp}
\begin{Resp}\Rsign Et vos símiles homínibus exspectántibus dóminum suum, quando revertátur a núptiis.\end{Resp}
\switchcolumn
\EN
\begin{Resp}\Rsign Let your loins be girded about, and your lights burning; \Star{} And ye yourselves like unto men that wait for their lord, when he will return from the wedding.\end{Resp}
\begin{Resp}\Vsign Watch therefore, for ye know not what hour your Lord doth come.\end{Resp}
\begin{Resp}\Rsign And ye yourselves like unto men that wait for their lord, when he will return from the wedding.\end{Resp}
\begin{Resp}\Vsign Glory be to the Father, and to the Son, \Star{} and to the Holy Ghost.\end{Resp}
\begin{Resp}\Rsign And ye yourselves like unto men that wait for their lord, when he will return from the wedding.\end{Resp}
\switchcolumn*

\LA
\LessonHead{Lectio IX}
\switchcolumn
\EN
\LessonHead{Lesson IX}
\switchcolumn*

\LA
\noindent Mérito dícitur Petro: díligis me, et respóndet, amo te, eíque refértur, pasce agnos meos, et hoc íterum, hoc tértio. Ubi étiam demonstrátur unum atque idem esse amórem et dilectiónem; nam étiam Dóminus novíssime non ait, díligis me, sed amas me. Non ergo nos, sed ipsum amémus, et in pascéndis óvibus ejus ea quæ sunt ejus, non ea quæ sunt nostra, quærámus. Néscio enim quo inexplicábili modo, quisquis seípsum, non Deum amat, non se amat: et quisquis Deum, non seípsum, amat, ipse se amat. Qui enim non potest vívere de se, móritur útique amándo se: non ergo se amat, qui ne vivat se amat.\par
\switchcolumn
\EN
\noindent With great propriety, therefore, is Peter addressed: Lovest thou me? and found replying: I love thee; and the command applied to him: Feed my lambs, and this a second and a third time. We have it also demonstrated here that love (amor) and liking (dilectio) are one and the same thing; for the Lord also in the last question said not: Dost thou like me? (Diligis me?): but: Dost thou love me? (Amas me?) Let us, then, love not ourselves but him; and in feeding his sheep, let us be seeking the things which are his, not the things which are our own. For in some inexplicable way, I know not what, every one that loveth himself, and not God, loveth not himself; and whoever loveth God, and not himself, he it is that loveth himself. For he that cannot live by himself will certainly die by loving himself; he therefore loveth not himself that loveth himself to his own loss of life.\par
\switchcolumn*

\LA
\TeDeum{Te Deum}
\switchcolumn
\EN
\TeDeum{Te Deum}
\switchcolumn*

\end{paracol}

\par\needspace{10\baselineskip}\addvspace{1.2em}{\centering\small\textsc{Die 5 Septembris} · \textsc{5 September}\par}
\Office{S. Laurentii Justiniani Episcopi et Confessoris}{St. Lawrence Justinian, Bishop and Confessor}{Semiduplex}
\addcontentsline{toc}{subsection}{Die 5 Septembris · S. Laurentii Justiniani Episcopi et Confessoris \textit{— St. Lawrence Justinian, Bishop and Confessor}}
\begin{paracol}{2}
\LA
\RefLine{Lectiones I–III: de Scriptura occurrente.}
\switchcolumn
\EN
\RefLine{Lessons I–III: from the Scripture of the day.}
\switchcolumn*

\LA
\RefLine{Lectiones VII–IX: ex Commúni unius Confessoris Pontificis (C4).}
\switchcolumn
\EN
\RefLine{Lessons VII–IX: from the Common of a Confessor Bishop (C4).}
\switchcolumn*

\LA
\LessonHead{Lectio IV}
\switchcolumn
\EN
\LessonHead{Lesson IV}
\switchcolumn*

\LA
\noindent Laurentius, ex illustri Justinianorum familia Venetiis natus, eximiam vel puer morum gravitatem præ se tulit. Exacta inter pietatis officia adolescentia, ad castum verbi et animæ connubium a divina sapientia invitatus, de religiosæ vitæ instituto capessendo deliberare cœpit. Novæ ítaque militiæ clam proludens, præter alias corporis affuctationes, super nudos cubabat asseres, sedensque velut arbiter hinc inter sæculi blandimenta, paratasque a matre nuptias, illinc claustrales inter austeritates oculis in Christi patientis crucem conversis: Tu, inquit, es Domine spes mea: ibi posuisti certissimum refugium tuum; ad canonicorum sancti Georgii in Alga congregationem convolavit: ubi novis excogitatis cruciatibus acrius in seipsum, veluti in hostem infensissimum, instaurans bellum, nullam adeo sibi oblectationem indulgebat, ut ne in domesticum umquam hortum, nec in paternam quidem domum, nisi cum morienti matri extrema pietatis officia siccis oculis persolvit, exinde intraverit. Par erat obedientiæ, mansuetudinis, ac præcipue humilitatis studium, cum abjectissima quæque cœnobii munia sibi ultro desumeret, celeberrima per urbis loca, non tam victum quam ludibria emendicaret, illatasque contumelias ac calumnias immotus ac silens perferret; assiduæ præsertim orationis subsidio, qua sæpe per mentis excessum rapiebatur in Deum; tantoque cor ejus æstuabat ardore, ut nutantes étiam sodales ad perseverantiam ac Jesu Christi amórem inflammaret.\par
\switchcolumn
\EN
\noindent This Lawrence was born at Venice, in the year 1380, of the noble family of the Giustiniani, and was an exceedingly sober lad even from his childhood. He passed a godly boyhood, and feeling the Divine Wisdom calling him to a pure marriage between his own soul and the Word of God, he began to think of becoming a monk. He therefore tried in private some of the exercises of this new warfare, and, among other afflictions of his body, used to sleep upon the bare boards. As he thus sat weighing on the one hand the pleasures of the world, and a marriage which his mother wished to bring about for him, and, on the other, the hardness of the cloister, he turned his eyes upon the Cross of the suffering Christ, and said: “Thou, O Lord, art my trust; there hast Thou made my surest refuge.” He entered among the Canons of St-George's-in-the-sea-weed, where he devised new tortures and declared war against himself as his own worst enemy. He allowed himself no enjoyment, so that he would not even go into the private garden of the house, neither did he ever go thenceforth into the house of his own father, except when his mother was dying, and he went there with dry eyes to pay her the last offices of a son's duty and affection. His obedience, gentleness, and especially his lowliness were very great. He went out of his way to take the meanest pieces of work about the house. He used to go to the most public places of the city, seeking, not so much for food as for mockery, and bore unmoved and in silence the insults and slanders which were cast upon him. He found his ever-present help in prayer, wherein he became often beside himself and rapt in God, and such was the warmth that burned in his heart, that he stirred up failing comrades to hold bravely on and to love Jesus Christ.\par
\switchcolumn*

\LA
\begin{Resp}\Rsign Invéni David servum meum, óleo sancto meo unxi eum: \Star{} Manus enim mea auxiliábitur ei.\end{Resp}
\begin{Resp}\Vsign Nihil profíciet inimícus in eo, et fílius iniquitátis non nocébit ei.\end{Resp}
\begin{Resp}\Rsign Manus enim mea auxiliábitur ei.\end{Resp}
\switchcolumn
\EN
\begin{Resp}\Rsign I have found David My servant, with My holy oil have I anointed him \Star{} For My hand shall help him.\end{Resp}
\begin{Resp}\Vsign The enemy shall prevail nothing against him, nor the son of wickedness afflict him.\end{Resp}
\begin{Resp}\Rsign For My hand shall help him.\end{Resp}
\switchcolumn*

\LA
\LessonHead{Lectio V}
\switchcolumn
\EN
\LessonHead{Lesson V}
\switchcolumn*

\LA
\noindent Ab Eugenio quarto patriæ episcopus designatus, quem magna contentione honorem detrectaverat, majori gessit cum laude. Nam consueta vivendi ratione nihil admodum immutata, paupertatem, quam semper coluerat, in mensa, supellectili ac lecto perpetuo retinuit. Modicam domi alebat familiam, quod grandem alteram sibi esse diceret, pauperes Christi significans. Quacumque adiretur hora, præsto omnibus erat, paterna omnes caritate allevabat, non renuens vel ære se alieno gravare, illorum ne inopiæ deesset. Rogatus qua spe id faceret? Domini mei, qui pro me dissolvere facile poterit, respondebat. Spem autem non confundere divina providentia submissis inopinato subsidiis jugiter declarabat. Plura virginum monasteria construxit, quas étiam ad perfectioris vitæ rationem sua vigilantia composuit. Matronis a sæculi pompis et ornatus vanitate revocandis, ecclesiasticæ disciplinæ ac moribus reformandis maximopere studuit; dignus sane qui ab eodem Eugenio gloria et decus præsulum coram cardinalibus vocaretur, et qui a Nicolao quinto ejus successore, translato e Gradensi civitate titulo, primus Venetiarum patriarcha renuntiaretur.\par
\switchcolumn
\EN
\noindent In the year 1433, Eugene IV named him Bishop of Venice, an office which he very earnestly struggled to avoid, and which he discharged with great honour. He changed in no wise way of living, but kept always to his beloved poverty in his table, his furniture, and his bed. He kept but a small household, saying that he had another very large one, in Christ's poor. At what hour soever any one came to see him, he was always ready to receive them, he helped all with the tenderness of a father, not refusing to charge himself with debts, that he might have wherewith to relieve misery. When he was asked with what hope he incurred these liabilities, he answered: “With hope in my Master, Who can easily meet them for me.” And the Providence of God put not his hope to shame, but helped him amply with unexpected funds. He built several Convents of nuns, for whom his watchful care ordered a more perfect way of living. He laboured much to wean married women from worldly folly and display, and to reform the discipline of the Church and the lives of all. He was indeed worthy that Eugene should call him in the presence of the Cardinals “the glory and ornament of the Episcopate,” and that his successor Nicholas V should transfer the title of Patriarch from Grado, and create him, in 1451, the first Patriarch of Venice.\par
\switchcolumn*

\LA
\begin{Resp}\Rsign Pósui adjutórium super poténtem, et exaltávi eléctum de plebe mea: \Star{} Manus enim mea auxiliábitur ei.\end{Resp}
\begin{Resp}\Vsign Invéni David servum meum, óleo sancto meo unxi eum.\end{Resp}
\begin{Resp}\Rsign Manus enim mea auxiliábitur ei.\end{Resp}
\switchcolumn
\EN
\begin{Resp}\Rsign I have laid help upon one that is mighty, and have exalted one chosen out of My people \Star{} For My hand shall help him.\end{Resp}
\begin{Resp}\Vsign I have found David My servant, with My holy oil have I anointed him.\end{Resp}
\begin{Resp}\Rsign For My hand shall help him.\end{Resp}
\switchcolumn*

\LA
\LessonHead{Lectio VI}
\switchcolumn
\EN
\LessonHead{Lesson VI}
\switchcolumn*

\LA
\noindent Lacrimarum dono insignitus, omnipotenti Deo placationis hostiam quotidie offerebat. Quod cum aliquando nocte Dominicæ Nativitatis perageret, Christum Jesum sub pulcherrimi infantis specie videre promeruit; tantumque in eo erat commissi gregis præsidium, ut cælitus aliquando acceptum fuerit, pontificis sui intercessione ac meritis stetisse rempublicam. Prophetiæ spiritu afflatus, plura humanæ cognitioni prorsus impervia prædixit: morbos ac dæmones suis precibus sæpe fugavit: libros étiam cælestem doctrinam ac pietatem spirantes, grammaticæ pene rudis, conscripsit. Denique cum lethalem incidisset in morbum, et commodiorem domestici lectum seni atque ægro pararent, aversatus ejusmodi delicias, tamquam a durissima morientis Domini sui cruce plus nimio abhorrentes, consueto in stramine se jussit deponi, et finem vitæ suæ adventare prænoscens, sublatis in cælum oculis: Venio, inquit, ad te, o bone Jesu; ac die octava Januarii obdormivit in Domino. Pretiosam ejus mortem testati sunt angelici concentus, a Carthusianis quibusdam monachis auditi, et sacrum cadaver per duos ultra menses inhumatum, suavi fragrans odore, et rubescente facie, integrum atque incorruptum, ac nova post mortem patrata miracula: quibus permotus Alexander octavus Pontifex Maximus eum sanctorum numero adscripsit. Innocentius vero duodecimus quintam Septembris diem qua vir sanctus ad pontificiam primo cathedram fuerat evectus, celebrando illius festo assignavit.\par
\switchcolumn
\EN
\noindent He was eminent for the gift of tears, in which he offered up to God every day the Sacrifice of atonement. When he was so doing one Christmas Midnight, he won to see Christ Jesus in the form of a little Child exceeding fair to look upon. Such was his care of the flock committed to his charge, that it was sometime revealed from heaven that the Commonwealth had been saved by the prayers of her Bishop. He was inspired with the spirit of prophecy, and fore-told many things which no wit of man could have perceived. By his prayers he often put diseases and devils to flight. Though very ignorant of letters, he wrote books which breathe heavenly teaching and godliness. When he fell into his last deadly sickness, his servants got ready a more comfortable bed for the suffering old man, but he turned away from such ease as so different from the hardness of the Cross upon which his Master had died. He ordered himself to be laid upon the planks to which he was accustomed, and when he knew that the end of his life was come, he looked up to heaven and said: “O good Jesus, I am coming to thee,” and so fell asleep in the Lord on the 8th day of January, in the year 1455. How precious was his death was attested by this, that some Charterhouse monks heard Angels singing and that the hallowed corpse, remaining unburied for two months, was whole and uncorrupted, always yielding a sweet smell, and rosy in the face. New miracles took place after his death, whereby Pope Alexander VIII was moved to enroll his name among those of the Saints. Innocent XII appointed for his Feast the 5th day of September, being that upon which he had first been enthroned in his Cathedral Church.\par
\switchcolumn*

\LA
\begin{Resp}\Rsign Iste est, qui ante Deum magnas virtútes operátus est, et omnis terra doctrína ejus repléta est: \Star{} Ipse intercédat pro peccátis ómnium populórum.\end{Resp}
\begin{Resp}\Vsign Iste est, qui contémpsit vitam mundi, et pervénit ad cæléstia regna.\end{Resp}
\begin{Resp}\Rsign Ipse intercédat pro peccátis ómnium populórum.\end{Resp}
\begin{Resp}\Vsign Glória Patri, et Fílio, \Star{} et Spirítui Sancto.\end{Resp}
\begin{Resp}\Rsign Ipse intercédat pro peccátis ómnium populórum.\end{Resp}
\switchcolumn
\EN
\begin{Resp}\Rsign This is he which wrought great wonders before God, and the whole earth is full of his teaching \Star{} May he pray for all people, that their sins may be forgiven unto them\end{Resp}
\begin{Resp}\Vsign This is he which loved not his life in this world, and hath attained unto the kingdom of heaven.\end{Resp}
\begin{Resp}\Rsign May he pray for all people, that their sins may be forgiven unto them\end{Resp}
\begin{Resp}\Vsign Glory be to the Father, and to the Son, \Star{} and to the Holy Ghost.\end{Resp}
\begin{Resp}\Rsign May he pray for all people, that their sins may be forgiven unto them\end{Resp}
\switchcolumn*

\end{paracol}

\par\needspace{10\baselineskip}\addvspace{1.2em}{\centering\small\textsc{Die 8 Septembris} · \textsc{8 September}\par}
\Office{In Nativitate Beatæ Mariæ Virginis}{In Nativitate Beatæ Mariæ Virginis}{Duplex II classis cum Octava simplici}
\addcontentsline{toc}{subsection}{Die 8 Septembris · In Nativitate Beatæ Mariæ Virginis \textit{— In Nativitate Beatæ Mariæ Virginis}}
\begin{paracol}{2}
\LA
\RefLine{Lectio IX: de S. Hadriani Martyris (09-08cc).}
\switchcolumn
\EN
\RefLine{Lesson IX: of St. Adrian, Martyr (09-08cc).}
\switchcolumn*

\LA
\LessonHead{Lectio I}
\switchcolumn
\EN
\LessonHead{Lesson I}
\switchcolumn*

\LA
\LessonTitle{Incípiunt Cántica canticórum}
\switchcolumn
\EN
\LessonTitle{Lesson from the book of Canticles}
\switchcolumn*

\LA
\Cite{Cant 1:1-4}
\switchcolumn
\EN
\Cite{Song 1:1-4}
\switchcolumn*

\LA
\noindent Osculétur me ósculo oris sui, quia melióra sunt úbera tua vino, fragrántia unguéntis óptimis. Oleum effúsum nomen tuum; ídeo adolescéntulæ dilexérunt te. Trahe me: post te currémus in odórem unguentórum tuórum. Introdúxit me rex in cellária sua; exsultábimus et lætábimur in te mémores úberum tuórum super vinum. Recti díligunt te. Nigra sum, sed formósa, fíliæ Jerúsalem, sicut tabernácula Cedar, sicut pelles Salomónis.\par
\switchcolumn
\EN
\noindent Let him kiss me with the kiss of his mouth: for thy breasts are better than wine, Smelling sweet of the best ointments. thy name is as oil poured out: therefore young maidens have loved thee. Draw me: we will run after thee to the odour of thy ointments. The king hath brought me into his storerooms: we will be glad and rejoice in thee, remembering thy breasts more than wine: the righteous love thee. I am black but beautiful, O ye daughters of Jerusalem, as the tents of Cedar, as the curtains of Solomon.\par
\switchcolumn*

\LA
\begin{Resp}\Rsign Hódie nata est beáta Virgo María ex progénie David; \Star{} Per quam salus mundi credéntibus appáruit, cujus vita gloriósa lucem dedit sǽculo.\end{Resp}
\begin{Resp}\Vsign Nativitátem beátæ Maríæ Vírginis cum gáudio celebrémus.\end{Resp}
\begin{Resp}\Rsign Per quam salus mundi credéntibus appáruit, cujus vita gloriósa lucem dedit sǽculo.\end{Resp}
\switchcolumn
\EN
\begin{Resp}\Rsign This day was the Blessed Virgin Mary born of the lineage of David. \Star{} The same is she through whom the salvation of the world hath been manifested before the eyes of all believers. This is she whose glorious life hath given light to the world.\end{Resp}
\begin{Resp}\Vsign Let us keep with rejoicing the Birthday of the Blessed Virgin Mary.\end{Resp}
\begin{Resp}\Rsign The same is she through whom the salvation of the world hath been manifested before the eyes of all believers. This is she whose glorious life hath given light to the world.\end{Resp}
\switchcolumn*

\LA
\LessonHead{Lectio II}
\switchcolumn
\EN
\LessonHead{Lesson II}
\switchcolumn*

\LA
\Cite{Cant 1:5-9}
\switchcolumn
\EN
\Cite{Song 1:5-9}
\switchcolumn*

\LA
\noindent Nolíte me consideráre quod fusca sim, quia decolorávit me sol. Fílii matris meæ pugnavérunt contra me, posuérunt me custódem in víneis, víneam meam non custodívi. Indica mihi, quem díligit ánima mea, ubi pascas, ubi cubes in merídie, ne vagári incípiam post greges sodálium tuórum. Si ignóras te, o pulchérrima inter mulíeres, egrédere et abi post vestígia gregum, et pasce hædos tuos juxta tabernácula pastórum. Equitátui meo in cúrribus pharaónis assimilávi te, amíca mea. Pulchræ sunt genæ tuæ sicut túrturis, collum tuum sicut monília.\par
\switchcolumn
\EN
\noindent Do not consider me that I am brown, because the sun hath altered my colour: the sons of my mother have fought against me, they have made me the keeper in the vineyards: my vineyard I have not kept. show me, O thou whom my soul loveth, where thou feedest, where thou liest in the midday, lest I begin to wander after the flocks of thy companions. If thou know not thyself, O fairest among women, go forth, and follow after the steps of the flocks, and feed thy kids beside the tents of the shepherds. To my company of horsemen, in Pharao's chariots, have I likened thee, O my love. thy cheeks are beautiful as the turtledove's, thy neck as jewels.\par
\switchcolumn*

\LA
\begin{Resp}\Rsign Beatíssimæ Vírginis Maríæ Nativitátem devotíssime celebrémus, \Star{} Ut ipsa pro nobis intercédat ad Dóminum Jesum Christum.\end{Resp}
\begin{Resp}\Vsign Cum jucunditáte Nativitátem beátæ Maríæ Vírginis devotíssime celebrémus.\end{Resp}
\begin{Resp}\Rsign Ut ipsa pro nobis intercédat ad Dóminum Jesum Christum.\end{Resp}
\switchcolumn
\EN
\begin{Resp}\Rsign Let us keep right heartily the Birthday of the Most Blessed Virgin Mary, \Star{} That she may pray for us to our Lord Jesus Christ.\end{Resp}
\begin{Resp}\Vsign Let us keep with right hearty rejoicing the Birthday of the Blessed Virgin Mary.\end{Resp}
\begin{Resp}\Rsign That she may pray for us to our Lord Jesus Christ.\end{Resp}
\switchcolumn*

\LA
\LessonHead{Lectio III}
\switchcolumn
\EN
\LessonHead{Lesson III}
\switchcolumn*

\LA
\Cite{Cant 1:10-16}
\switchcolumn
\EN
\Cite{Song 1:10-16}
\switchcolumn*

\LA
\noindent Murénulas áureas faciémus tibi vermiculátas argénto. Dum esset rex in accúbitu suo, nardus mea dedit odórem suum. Fascículus myrrhæ diléctus meus mihi, inter úbera mea commorábitur. Botrus Cypri diléctus meus mihi in víneis Engáddi. Ecce tu pulchra es, amíca mea, ecce tu pulchra es; óculi tui columbárum. Ecce tu pulcher es, dilécte mi, et decórus. Léctulus noster flóridus, tigna domórum nostrárum cédrina, laqueária nostra cypréssina.\par
\switchcolumn
\EN
\noindent We will make thee chains of gold, inlaid with silver. While the king was at his repose, my spikenard sent forth the odour thereof. A bundle of myrrh is my beloved to me, he shall abide between my breasts. A cluster of cypress my love is to me, in the vineyards of Engaddi. Behold thou art fair, O my love, behold thou art fair, thy eyes are as those of doves. Behold thou art fair, my beloved, and comely. Our bed is flourishing. The beams of our houses are of cedar, our rafters of cypress trees.\par
\switchcolumn*

\LA
\begin{Resp}\Rsign Gloriósæ Vírginis Maríæ ortum digníssimum recolámus, \Star{} Cujus Dóminus humilitátem respéxit, quæ, Angelo nuntiánte, concépit Salvatórem mundi.\end{Resp}
\begin{Resp}\Vsign Beatíssimæ Vírginis Maríæ Nativitátem devotíssime celebrémus.\end{Resp}
\begin{Resp}\Rsign Cujus Dóminus humilitátem respéxit, quæ, Angelo nuntiánte, concépit Salvatórem mundi.\end{Resp}
\begin{Resp}\Vsign Glória Patri, et Fílio, \Star{} et Spirítui Sancto.\end{Resp}
\begin{Resp}\Rsign Cujus Dóminus humilitátem respéxit, quæ, Angelo nuntiánte, concépit Salvatórem mundi.\end{Resp}
\switchcolumn
\EN
\begin{Resp}\Rsign Let us tell again of the right worthy Birth of the Blessed Virgin Mary. \Star{} The same is she whose lowliness the Lord regarded, she who by the message of an Angel conceived the Saviour of the world.\end{Resp}
\begin{Resp}\Vsign Let us keep right earnestly the Birthday of the most Blessed Virgin Mary.\end{Resp}
\begin{Resp}\Rsign The same is she whose lowliness the Lord regarded, she who by the message of an Angel conceived the Saviour of the world.\end{Resp}
\begin{Resp}\Vsign Glory be to the Father, and to the Son, \Star{} and to the Holy Ghost.\end{Resp}
\begin{Resp}\Rsign The same is she whose lowliness the Lord regarded, she who by the message of an Angel conceived the Saviour of the world.\end{Resp}
\switchcolumn*

\LA
\LessonHead{Lectio IV}
\switchcolumn
\EN
\LessonHead{Lesson IV}
\switchcolumn*

\LA
\LessonTitle{Sermo sancti Augustíni Epíscopi}
\switchcolumn
\EN
\LessonTitle{From the Sermons of St. Augustine, Bishop (of Hippo.)}
\switchcolumn*

\LA
\Source{Sermo 18 de Sanctis, qui est 2 de Annunt. Dominica}
\switchcolumn
\EN
\Source{18th on the Saints.}
\switchcolumn*

\LA
\noindent Adest nobis, dilectíssimi, optátus dies beátæ ac venerábilis semper Vírginis Maríæ; ídeo cum summa exsultatióne gáudeat terra nostra, tantæ Vírginis illustráta natáli. Hæc est enim flos campi, de qua ortum est pretiósum lílium convállium, per cujus partum mutátur natúra protoplastórum, delétur et culpa. Præcísum est in ea illud Hevæ infelicitátis elógium, quo dícitur: In dolóre páries fílios tuos; quia ista in lætítia Dóminum péperit.\par
\switchcolumn
\EN
\noindent Dearly beloved brethren, the day for which we have longed, the Feast-day of the Blessed and Worshipful and Ever-Virgin Mary, that day is come. Let our land laugh and sing with merriment, bathed in the glory of this great Virgin's rising. She is the flower of the fields on which the priceless lily of the valleys hath blossomed. This is she whose delivery changed the nature that we draw from our first parents, and cleansed away their offence. At her that dolorous sentence which was pronounced over Eve ended its course to her it was never said: “In sorrow thou shalt bring forth children.” (Gen. iii. 16.) She brought forth a Child, even the Lord, but she brought Him forth, not in sorrow, but in joy.\par
\switchcolumn*

\LA
\begin{Resp}\Rsign Natívitas gloriósæ Vírginis Maríæ ex sémine Abrahæ, ortæ de tribu Juda, clara ex stirpe David; \Star{} Cujus vita ínclita cunctas illústrat ecclésias.\end{Resp}
\begin{Resp}\Vsign Hódie nata est beáta Virgo María ex progénie David.\end{Resp}
\begin{Resp}\Rsign Cujus vita ínclita cunctas illústrat ecclésias.\end{Resp}
\switchcolumn
\EN
\begin{Resp}\Rsign This day was born the glorious Virgin Mary, a child of the seed of Abraham, a daughter of the tribe of Judah, a Princess of the lineage of David. \Star{} This is she whose famous life still sheddeth lustre upon all the Churches.\end{Resp}
\begin{Resp}\Vsign This day was the Blessed Virgin Mary born of the lineage of David.\end{Resp}
\begin{Resp}\Rsign This is she whose famous life still sheddeth lustre upon all the Churches.\end{Resp}
\switchcolumn*

\LA
\LessonHead{Lectio V}
\switchcolumn
\EN
\LessonHead{Lesson V}
\switchcolumn*

\LA
\noindent Heva enim luxit, ista exsultávit: Heva lácrimas, María gáudium in ventre portávit; quia illa peccatórem, ista édidit innocéntem. Mater géneris nostri pœnam íntulit mundo, Génetrix Dómini nostri salútem íntulit mundo. Auctrix peccáti Heva, auctrix mériti María. Heva occidéndo óbfuit, María vivificándo prófuit. Illa percússit, ista sanávit. Pro inobediéntia enim obediéntia commutátur, fides pro perfídia compensátur.\par
\switchcolumn
\EN
\noindent Eve wept, but Mary laughed. Eve's womb was big with tears, but Mary's womb was big with gladness. Eve gave birth to a sinner, but Mary gave birth to the sinless One. The mother of our race brought punishment into the world, but the Mother of our Lord brought salvation into the world. Eve was the foundress of sin, but Mary was the foundress of righteousness. Eve welcomed death, but Mary helped in life. Eve smote, but Mary healed. For Eve's disobedience, Mary offered obedience and for Eve's unbelief, Mary offered faith.\par
\switchcolumn*

\LA
\begin{Resp}\Rsign Cum jucunditáte Nativitátem beátæ Maríæ celebrémus, \Star{} Ut ipsa pro nobis intercédat ad Dóminum Jesum Christum.\end{Resp}
\begin{Resp}\Vsign Corde et ánimo Christo canámus glóriam in hac sacra solemnitáte præcélsæ Genetrícis Dei Maríæ.\end{Resp}
\begin{Resp}\Rsign Ut ipsa pro nobis intercédat ad Dóminum Jesum Christum.\end{Resp}
\switchcolumn
\EN
\begin{Resp}\Rsign Let us keep with rejoicing the Birthday of the Blessed Mary, \Star{} That she may pray for us to our Lord Jesus Christ.\end{Resp}
\begin{Resp}\Vsign With all our heart and with all our soul let us sing praise to Christ on this the solemn Feastday of Mary the mighty Mother of God.\end{Resp}
\begin{Resp}\Rsign That she may pray for us to our Lord Jesus Christ.\end{Resp}
\switchcolumn*

\LA
\LessonHead{Lectio VI}
\switchcolumn
\EN
\LessonHead{Lesson VI}
\switchcolumn*

\LA
\noindent Plaudat nunc órganis María, et inter velóces artículos týmpana puérperæ cóncrepent. Cóncinant lætántes chori, et alternántibus módulis dulcísona cármina misceántur. Audíte ígitur, quemádmodum tympanístria nostra cantáverit; ait enim: Magníficat ánima mea Dóminum: et exsultávit spíritus meus in Deo, salutári meo. Quia respéxit humilitátem ancíllæ suæ: ecce enim ex hoc beátam me dicent omnes generatiónes. Quia fecit mihi magna qui potens est. Causam ígitur invalescéntis erráti, miráculum novi partus evícit; et Hevæ planctum Maríæ cantus exclúsit.\par
\switchcolumn
\EN
\noindent Let Mary now make a loud noise upon the organ, and between its quick notes let the rattling of the Mother's timbrel be heard. Let the gladsome choirs sing with her, and their sweet hymns mingle with the changing music. Hearken to what a song her timbrel will make accompaniment. She saith: “My soul doth magnify the Lord, and my spirit hath rejoiced in God my Saviour. For He hath regarded the lowliness of His hand-maiden, for, behold, from henceforth all generations shall call me blessed for He That is Mighty hath done to me great things.” The new miracle of Mary's delivery hath effaced the curse of the frail backslider, and the singing of Mary hath silenced the wailing of Eve.\par
\switchcolumn*

\LA
\begin{Resp}\Rsign Natívitas tua, Dei Génetrix Virgo, gáudium annuntiávit univérso mundo; \Star{} Ex te enim ortus est sol justítiæ, Christus Deus noster: * Qui, solvens maledictiónem, dedit benedictiónem; et confúndens mortem, donávit nobis vitam sempitérnam.\end{Resp}
\begin{Resp}\Vsign Benedícta tu in muliéribus, et benedíctus fructus ventris tui.\end{Resp}
\begin{Resp}\Rsign Ex te enim ortus est sol justítiæ, Christus Deus noster.\end{Resp}
\begin{Resp}\Vsign Glória Patri, et Fílio, \Star{} et Spirítui Sancto.\end{Resp}
\begin{Resp}\Rsign Qui, solvens maledictiónem, dedit benedictiónem; et confúndens mortem, donávit nobis vitam sempitérnam.\end{Resp}
\switchcolumn
\EN
\begin{Resp}\Rsign Thy Birth, O Virgin Mother of God, was a message of joy to the whole world; \Star{} For out of thee rose the Sun of righteousness, even Christ our God, Who hath taken away the curse and brought a blessing, confounded death, and given unto us everlasting life.\end{Resp}
\begin{Resp}\Vsign Blessed art thou among women, and blessed is the Fruit of thy womb.\end{Resp}
\begin{Resp}\Rsign For out of thee rose the Sun of righteousness, even Christ our God.\end{Resp}
\begin{Resp}\Vsign Glory be to the Father, and to the Son, \Star{} and to the Holy Ghost.\end{Resp}
\begin{Resp}\Rsign Who hath taken away the curse and brought a blessing, confounded death, and given unto us everlasting life.\end{Resp}
\switchcolumn*

\LA
\LessonHead{Lectio VII}
\switchcolumn
\EN
\LessonHead{Lesson VII}
\switchcolumn*

\LA
\LessonTitle{Léctio sancti Evangélii secúndum Matthǽum}
\switchcolumn
\EN
\LessonTitle{From the Holy Gospel according to Matthew}
\switchcolumn*

\LA
\Cite{Matt 1:1-16}
\switchcolumn
\EN
\Cite{Matt 1:1-17}
\switchcolumn*

\LA
\noindent Liber generatiónis Jesu Christi, fílii David, fílii Abraham. Abraham génuit Isaac, Isaac autem génuit Jacob. Et réliqua.\par
\switchcolumn
\EN
\noindent The Book of the generation of Jesus Christ, the Son of David, the son of Abraham. Abraham begat Isaac, and Isaac begat Jacob. And so on.\par
\switchcolumn*

\LA
\LessonTitle{Homilía sancti Hierónymi Presbýteri}
\switchcolumn
\EN
\LessonTitle{Homily by St. Jerome, Priest (at Bethlehem)}
\switchcolumn*

\LA
\Source{Liber 1 Comment. in Matth., in initium}
\switchcolumn
\EN
\Source{Bk. i, Comm. on Matth.}
\switchcolumn*

\LA
\noindent In Isaía légimus: Generatiónem ejus quis enarrábit? Non ergo putémus Evangelístam Prophétæ esse contrárium, ut quod ille impossíbile dixit effátu, hic narráre incípiat: quia ibi de generatióne Divinitátis, hic de incarnatióne est dictum. A carnálibus autem cœpit, ut per hóminem, Deum díscere incipiámus. Fílii David, fílii Abraham. Ordo præpósterus, sed necessário commutátus. Si enim primum posuísset Abraham et póstea David, rursus ei repeténdus fúerat Abraham, ut generatiónis séries texerétur.\par
\switchcolumn
\EN
\noindent In Isaiah (liii. 8) we read: “Who shall declare His generation?” Let us not think that there is any contradiction between the Prophet and the Evangelist, because the Prophet saith that this thing cannot be done, and the Evangelist beginneth by doing it. The one speaketh of the generation of the Divine (Word by the Eternal Father,) the other of the (family in which the) Incarnation (took place.) Matthew beginneth with carnal things, that by learning of men we may go on to learn of God. “The Son of David, the son of Abraham.” The reversal of the order in these clauses is a needful change. If Abraham had been put first and David afterwards, Abraham would have had to be taken again, in order to marshal the pedigree properly.\par
\switchcolumn*

\LA
\begin{Resp}\Rsign Beátam me dicent omnes generatiónes, \Star{} Quia fecit mihi Dóminus magna qui potens est, et sanctum nomen ejus.\end{Resp}
\begin{Resp}\Vsign Et misericórdia ejus a progénie in progénies timéntibus eum.\end{Resp}
\begin{Resp}\Rsign Quia fecit mihi Dóminus magna qui potens est, et sanctum nomen ejus.\end{Resp}
\switchcolumn
\EN
\begin{Resp}\Rsign All generations shall call me blessed, \Star{} For the Lord that is mighty hath magnified me, and holy is his Name.\end{Resp}
\begin{Resp}\Vsign And his mercy is on them that fear him throughout all generations.\end{Resp}
\begin{Resp}\Rsign For the Lord that is mighty hath magnified me, and holy is his Name.\end{Resp}
\switchcolumn*

\LA
\LessonHead{Lectio VIII}
\switchcolumn
\EN
\LessonHead{Lesson VIII}
\switchcolumn*

\LA
\noindent Ideo autem, céteris prætermíssis, horum fílium nuncupávit, quia ad hos tantum facta est de Christo repromíssio. Ad Abraham: In sémine, inquit, tuo benedicéntur omnes gentes, quod est Christus. Ad David: De fructu ventris tui ponam super sedem tuam. Judas autem génuit Phares et Zaram de Thamar. Notándum, in genealogía Salvatóris nullam sanctárum assúmi mulíerum, sed eas quas Scriptúra reprehéndit; ut, qui propter peccatóres vénerat, de peccatóribus nascens, ómnium peccáta deléret. Unde et in consequéntibus Ruth Moabítis pónitur, et Bethsabée uxor Uríæ.\par
\switchcolumn
\EN
\noindent Matthew first calleth Christ the Son of these twain Abraham and David without making mention of the others, because unto these twain only was promise of Christ made unto Abraham, where it is said: “In thy seed” (that is, in Christ) “shall all the nations of the earth be blessed,” (Gen. xxii. 18) and unto David, in the words “Of the fruit of thy body will I set upon thy throne.” (Ps. cxxxi, 11) “And Judas begat Phares and Zara of Thamar.” It is to be remarked that in the genealogy of the Saviour none of the holy women are named, but those women only are named against whom the Scripture hath to say something amiss. He Who came to save sinners was born of sinners, that He might wash away all sin. Afterwards are named Ruth, who was a Moabitess, and Bathsheba, who had been the wife of Uriah.\par
\switchcolumn*

\LA
\begin{Resp}\Rsign Felix namque es, sacra Virgo María, et omni laude digníssima: \Star{} Quia ex te ortus est sol justítiæ, * Christus, Deus noster.\end{Resp}
\begin{Resp}\Vsign Ora pro pópulo, intérveni pro clero, intercéde pro devóto femíneo sexu: séntiant omnes tuum juvámen, quicúmque célebrant tuam sanctam Nativitátem.\end{Resp}
\begin{Resp}\Rsign Quia ex te ortus est sol justítiæ, Christus, Deus noster.\end{Resp}
\begin{Resp}\Vsign Glória Patri, et Fílio, \Star{} et Spirítui Sancto.\end{Resp}
\begin{Resp}\Rsign Christus, Deus noster.\end{Resp}
\switchcolumn
\EN
\begin{Resp}\Rsign O Holy Virgin Mary, happy indeed art thou, and right worthy of all praise \Star{} For out of thee rose the Sun of righteousness, * Even Christ our God.\end{Resp}
\begin{Resp}\Vsign Pray for the people, plead for the clergy, make intercession for all women vowed to God. May all that are keeping this day in honour of thy Nativity feel the might of thine assistance.\end{Resp}
\begin{Resp}\Rsign For out of thee rose the Sun of righteousness, even Christ our God.\end{Resp}
\begin{Resp}\Vsign Glory be to the Father, and to the Son, \Star{} and to the Holy Ghost.\end{Resp}
\begin{Resp}\Rsign Even Christ our God.\end{Resp}
\switchcolumn*

\end{paracol}

\par\needspace{10\baselineskip}\addvspace{1.2em}{\centering\small\textsc{Die 8 Septembris} · \textsc{8 September}\par}
\Office{S. Hadriani Martyris}{St. Adrian, Martyr}{Simplex}
\addcontentsline{toc}{subsection}{Die 8 Septembris · S. Hadriani Martyris \textit{— St. Adrian, Martyr}}
\begin{paracol}{2}
\LA
\LessonHead{Lectio IX (pro commemoratione)}
\switchcolumn
\EN
\LessonHead{Lesson IX (when commemorated)}
\switchcolumn*

\LA
\LessonTitle{Commemoratio: S. Hadriani Martyris}
\switchcolumn
\EN
\LessonTitle{Commemoration of: St. Adrian, Martyr}
\switchcolumn*

\LA
\noindent Hadrianus, jussu Maximiáni imperatóris apud Nicomedíam pérsequens Christiános, cum sæpius eórum in fidei confessióne et tormentórum perpessióne constantiam demirátus esset, veheménter ea re commótus ad Christum sese convértit. Quam ob rem cum aliis viginti tribus Christiánis conjéctus est in carcerem; ubi eum visitans Natália uxor, quæ et ipsa antea in Christum credíderat, ad martyrium incendit. Itaque, e custódia edúctus, tamdiu flagellis cæsus est, donec intestína diffluerent. Postremo, fractis crúribus, mánibus pedibúsque præcisis, una cum multis aliis martyrii certamen feliciter absolvit.\par
\switchcolumn
\EN
\noindent Adrian was a man who was employed by the Emperor Maximian to persecute the Christians of Nicomedia. The firmness with which they owned their faith and endured their torments oftentimes excited his wonder, and at last so powerfully moved him that he himself turned to Christ. For this he was thrown into prison along with three-and-twenty other Christians. There he was visited by Natalia his wife, who also herself already had believed in Christ, and by her urged on to lift up his testimony. When he was brought out of prison he was lashed until his bowels fell out. His shins were then broken, and his hands and feet cut off, whereafter, in company with many others, he brought to an happy end the conflict of martyrdom.\par
\switchcolumn*

\LA
\TeDeum{Te Deum}
\switchcolumn
\EN
\TeDeum{Te Deum}
\switchcolumn*

\end{paracol}

\par\needspace{10\baselineskip}\addvspace{1.2em}{\centering\small\textsc{Die 9 Septembris} · \textsc{9 September}\par}
\Office{S. Gorgonii Martyris}{S. Gorgonius, Martyr}{Simplex}
\addcontentsline{toc}{subsection}{Die 9 Septembris · S. Gorgonii Martyris \textit{— S. Gorgonius, Martyr}}
\begin{paracol}{2}
\LA
\RefLine{Lectiones I–II: de Scriptura occurrente.}
\switchcolumn
\EN
\RefLine{Lessons I–II: from the Scripture of the day.}
\switchcolumn*

\LA
\LessonHead{Lectio III (pro commemoratione)}
\switchcolumn
\EN
\LessonHead{Lesson III (when commemorated)}
\switchcolumn*

\LA
\LessonTitle{Commemoratio S. Gorgonii Martyris}
\switchcolumn
\EN
\LessonTitle{Commemoratio for the Holy Martyr Gorgonitis.}
\switchcolumn*

\LA
\noindent Gorgonius, Nicomedíæ natus, Diocletiáni imperatóris cubicularius, Dorothéo colléga suo adjutore, réliquos omnes cubículi minístros ad Christi fidem perdúxit. Uterque autem, cum vidísset quodam die Mártyrem coram Diocletiáno acerbíssime cruciari, ejus exemplo martyrii amóre incénsus est. Itaque ambórum hæc vox erupit: Quid est, imperátor, quod hujus condemnáta senténtia, quæ nobis cum eo communis est, unum illum punis? istius nostra étiam est fides, idem propositum. Eos ígitur vinctos imperátor flagellis cóncidi jubet, ita ut toto corpore cutis dirumperétur, et in plagas acetum infundi sale permixtum; mox revinctis in cratícula súbjici ímperat vim candéntium carbónum. Denique, varie torti, suspendio necáti sunt. Ac sancti Gorgonii corpus aliquándo Romam portátum, inter duas Lauros via Latina sepúltum, póstea a Gregório quarto summo Pontifice in basilicam Principis Apostolórum translátum est.\par
\switchcolumn
\EN
\noindent Gorgonius was a native of Nicomedia, and one of the chamberlains of the Emperor Diocletian. He, with the help of a fellow-chamberlain named Dorotheus, brought all the other chamber-servants to believe in Christ. Both of them one day saw a martyr hideously tortured in the presence of Diocletian, and the example of his testimony roused them both up to desire the same, and they both said: Why, O Emperor, dost thou punish this man only, by condemning an opinion which we share with him? His belief is our belief. Our will is the same. The Emperor thereupon ordered them to be bound and scourged till their bodies were perfectly flayed, and a mixture of vinegar and salt poured into the wounds. Soon after he commanded them to be bound again and grilled on bars over hot coals. Finally, after a variety of tortures, they were hanged. The body of the holy Gorgonius was some time brought to Rome, and buried between the two laurel-trees upon the Latin Way, but, afterwards, during the Pontificate of Gregory IV., it was brought into the Church of the Prince of the Apostles.\par
\switchcolumn*

\LA
\TeDeum{Te Deum}
\switchcolumn
\EN
\TeDeum{Te Deum}
\switchcolumn*

\end{paracol}

\par\needspace{10\baselineskip}\addvspace{1.2em}{\centering\small\textsc{Die 10 Septembris} · \textsc{10 September}\par}
\Office{S. Nicolai de Tolentino Confessoris}{St. Nicholas of Tolentino Confessor}{Duplex}
\addcontentsline{toc}{subsection}{Die 10 Septembris · S. Nicolai de Tolentino Confessoris \textit{— St. Nicholas of Tolentino Confessor}}
\begin{paracol}{2}
\LA
\RefLine{Lectiones I–III: de Scriptura occurrente.}
\switchcolumn
\EN
\RefLine{Lessons I–III: from the Scripture of the day.}
\switchcolumn*

\LA
\LessonHead{Lectio IV}
\switchcolumn
\EN
\LessonHead{Lesson IV}
\switchcolumn*

\LA
\noindent Nicolaus, Tolentinas a diuturno illíus civitátis domicílio appellatus, in oppido sancti Angeli in Piceno est natus piis paréntibus; qui, liberórum desidério Bárium voti causa profecti, ibique a sancto Nicoláo de futura prole confirmáti, quem suscepérunt fílium de illíus nómine appellarunt. Is ab infántia multárum virtútum, sed abstinéntiæ in primis, specimen dedit. Nam anno vix septimo, beátum ipsum Nicolaum imitatus, complures hebdomadæ dies jejunare cœpit; eamque póstea consuetúdinem retinuit, solo pane et aqua contentus.\par
\switchcolumn
\EN
\noindent This Nicholas is called Nicholas of Tolentino, because he lived in that town for most part of his life. He was born at St. Angelo, (a place near Fermo,) in the March of Ancona, (about the year 1245.) His parents were godly people, and in their desire to have children, vowed and made a pilgrimage to the shrine of St Nicholas at Bari, where they were assured of their wish, and therefore gave the name of Nicholas to the son whom they received. From his childhood the lad gave many good signs, but especially as regarded abstinence. In his seventh year, in imitation of his blessed name-sake, he began to fast upon several days in the week, which custom he always kept, and was content with only bread and water.\par
\switchcolumn*

\LA
\begin{Resp}\Rsign Honéstum fecit illum Dóminus, et custodívit eum ab inimícis, et a seductóribus tutávit illum: \Star{} Et dedit illi claritátem ætérnam.\end{Resp}
\begin{Resp}\Vsign Justum dedúxit Dóminus per vias rectas, et osténdit illi regnum Dei.\end{Resp}
\begin{Resp}\Rsign Et dedit illi claritátem ætérnam.\end{Resp}
\switchcolumn
\EN
\begin{Resp}\Rsign The Lord made him honourable, and defended him from his enemies, and kept him safe from those that lay in wait for him; \Star{} And gave him perpetual glory.\end{Resp}
\begin{Resp}\Vsign He went down with him into the pit, and left him not in bonds.\end{Resp}
\begin{Resp}\Rsign And gave him perpetual glory.\end{Resp}
\switchcolumn*

\LA
\LessonHead{Lectio V}
\switchcolumn
\EN
\LessonHead{Lesson V}
\switchcolumn*

\LA
\noindent Adulta ætate, jam clericali milítiæ adscriptus et canonicus factus, cum quodam die concionatórem ordinis Eremitárum sancti Augustíni de mundi contemptu dicéntem audísset, eo sermóne inflammatus, statim eumdem ordinem est ingréssus. In quo tam exactam religiosæ vitæ ratiónem coluit, ut áspero vestítu, verbéribus et férrea catena corpus domans, atque a carne et omni fere obsónio ábstinens, caritate, humilitate, patiéntia ceterisque virtútibus, aliis prælucéret.\par
\switchcolumn
\EN
\noindent After he reached man's estate, he enlisted himself in the army of the clergy, and was preferred to a Canonry. One day he chanced to hear a sermon upon contempt of the world delivered by a preacher of the Order of Hermits of St. Augustine, and was so moved by it that he forthwith entered that Order. As a Friar he was most strictly observant of that way of life. He subdued his body with rough clothing, stripes, and an iron chain. He never ate meat, and seldom any relish to his meals. And he was a burning and shining light of love, lowliness, long-suffering, and all other graces.\par
\switchcolumn*

\LA
\begin{Resp}\Rsign Amávit eum Dóminus, et ornávit eum: stolam glóriæ índuit eum, \Star{} Et ad portas paradísi coronávit eum.\end{Resp}
\begin{Resp}\Vsign Índuit eum Dóminus lorícam fídei, et ornávit eum.\end{Resp}
\begin{Resp}\Rsign Et ad portas paradísi coronávit eum.\end{Resp}
\switchcolumn
\EN
\begin{Resp}\Rsign The Lord loved him and beautified him; He clothed him with a robe of glory, \Star{} And crowned him at the gates of Paradise.\end{Resp}
\begin{Resp}\Vsign The Lord hath put on him the breast-plate of faith, and hath adorned him.\end{Resp}
\begin{Resp}\Rsign And crowned him at the gates of Paradise.\end{Resp}
\switchcolumn*

\LA
\LessonHead{Lectio VI}
\switchcolumn
\EN
\LessonHead{Lesson VI}
\switchcolumn*

\LA
\noindent Orándi assiduum studium, quamvis sátanæ insídiis varie vexátus et flagellis intérdum cæsus, non intermittebat. Demum, sex ante óbitum mensibus, singulis noctibus angelicum concentum audívit; cujus suavitate cum jam paradisi gaudia prægustaret, crebro illud Apóstoli repetebat: Cupio dissolvi, et esse cum Christo. Denique óbitus sui diem frátribus prædixit, qui fuit quarto Idus Septembris. Miraculis multis étiam post mortem cláruit; quibus rite et ordine cógnitis, ab Eugenio Papa quarto in Sanctórum númerum est relatus.\par
\switchcolumn
\EN
\noindent He persisted in constant and earnest prayer, notwithstanding many troubles from the assaults of Satan, who sometimes even flogged him. Every night for six months before his death he heard Angels singing with such sweetness, that it was a fore-taste of the happiness of heaven, and he would often repeat the words of the Apostle I have a desire to depart and to be with Christ (Phil. i, 23) Lastly, he foretold to his brethren the day of his death, which was the 10th day of September (1306.) After his death also he was famous for miracles, and when due investigation had been made thereof, Pope Eugene IV enrolled his name among those of the Saints.\par
\switchcolumn*

\LA
\begin{Resp}\Rsign Iste homo perfécit ómnia quæ locútus est ei Deus, et dixit ad eum: Ingrédere in réquiem meam: \Star{} Quia te vidi justum coram me ex ómnibus géntibus.\end{Resp}
\begin{Resp}\Vsign Iste est, qui contémpsit vitam mundi, et pervénit ad cæléstia regna.\end{Resp}
\begin{Resp}\Rsign Quia te vidi justum coram me ex ómnibus géntibus.\end{Resp}
\begin{Resp}\Vsign Glória Patri, et Fílio, \Star{} et Spirítui Sancto.\end{Resp}
\begin{Resp}\Rsign Quia te vidi justum coram me ex ómnibus géntibus.\end{Resp}
\switchcolumn
\EN
\begin{Resp}\Rsign This is he which did according unto all that God commanded him; and God said unto him: Enter thou into My rest, \Star{} For thee have I seen righteous before Me among all people.\end{Resp}
\begin{Resp}\Vsign This is he which loved not his life in this world, and is come unto an everlasting kingdom.\end{Resp}
\begin{Resp}\Rsign For thee have I seen righteous before Me among all people.\end{Resp}
\begin{Resp}\Vsign Glory be to the Father, and to the Son, \Star{} and to the Holy Ghost.\end{Resp}
\begin{Resp}\Rsign For thee have I seen righteous before Me among all people.\end{Resp}
\switchcolumn*

\LA
\LessonHead{Lectio VII}
\switchcolumn
\EN
\LessonHead{Lesson VII}
\switchcolumn*

\LA
\LessonTitle{Léctio sancti Evangélii secúndum Lucam}
\switchcolumn
\EN
\LessonTitle{From the Holy Gospel according to Luke}
\switchcolumn*

\LA
\Cite{Luc 12:32-34}
\switchcolumn
\EN
\Cite{Luke 12:32-35}
\switchcolumn*

\LA
\noindent In illo témpore: Dixit Jesus discípulis suis: Nolíte timére pusíllus grex, quia complácuit Patri vestro dare vobis regnum. Et réliqua.\par
\switchcolumn
\EN
\noindent At that time, Jesus said unto His disciples: Fear not, little flock, for it is your Father's good pleasure to give you the kingdom. And so on.\par
\switchcolumn*

\LA
\LessonTitle{Homilía sancti Bedæ Venerábilis Presbýteri}
\switchcolumn
\EN
\LessonTitle{Homily by the Venerable Bede, Priest at Jarrow, and Doctor of the Church.}
\switchcolumn*

\LA
\Source{Lib. 4, Cap. 54, in Luc. 12}
\switchcolumn
\EN
\Source{Bk. iv. Ch. 54 on Luke xii.}
\switchcolumn*

\LA
\noindent Pusíllum gregem electórum, vel ob comparatiónem majóris númeri reprobórum, vel pótius ob humilitátis devotiónem nóminat: quia vidélicet Ecclésiam suam quantálibet numerositáte jam dilatátam, tamen usque ad finem mundi humilitáte vult créscere, et ad promíssum regnum humilitáte perveníre. Ideóque ejus labóres blande consolátus, quam regnum Dei tantum quærére prǽcipit, eídem regnum a Patre dandum complácita benignitáte promíttit.\par
\switchcolumn
\EN
\noindent The elect are called a little flock, perchance because the reprobate are far more in number than they, but, more probably, because they love to be lowly, since it is God's will that however much His Church should grow in numbers, she should grow with lowliness even unto the end of the world, and should enter lowly into that kingdom which is hers by His promise. That kingdom He promiseth to her here, when He biddeth her to seek only the kingdom of God, and, to comfort her in her travail, He doth so sweetly and so graciously say that her Father will give it to her.\par
\switchcolumn*

\LA
\begin{Resp}\Rsign Iste est qui ante Deum magnas virtútes operátus est, et de omni corde suo laudávit Dóminum: \Star{} Ipse intercédat pro peccátis ómnium populórum.\end{Resp}
\begin{Resp}\Vsign Ecce homo sine queréla, verus Dei cultor, ábstinens se ab omni ópere malo, et pérmanens in innocéntia sua.\end{Resp}
\begin{Resp}\Rsign Ipse intercédat pro peccátis ómnium populórum.\end{Resp}
\switchcolumn
\EN
\begin{Resp}\Rsign This is he which wrought great wonders before God, and praised the Lord with all his heart. \Star{} May he pray for all people, that their sins may be forgiven unto them.\end{Resp}
\begin{Resp}\Vsign Behold a man without blame, a worshipper of God in truth, keeping himself clean from every evil work, and abiding still in his innocency.\end{Resp}
\begin{Resp}\Rsign May he pray for all people, that their sins may be forgiven unto them.\end{Resp}
\switchcolumn*

\LA
\LessonHead{Lectio VIII}
\switchcolumn
\EN
\LessonHead{Lesson VIII}
\switchcolumn*

\LA
\noindent Véndite quæ possidétis, et date eleemósynam. Nolíte, inquit, timére, ne propter regnum Dei militántibus, hujus vitæ necessária desint: quin étiam posséssa propter eleemósynam véndite. Quod tunc digne fit, quando quis semel pro Dómino suis ómnibus spretis, nihilóminus post hæc labóre mánuum, unde et victum transígere, et eleemósynam dare queat, operátur. Unde gloriátur Apóstolus, dicens: Argéntum, et aurum, aut vestem nullíus concupívi: ipsi scitis, quóniam ad ea quæ mihi opus erant, et his qui mecum sunt, ministravérunt manus istæ. Ómnia osténdi vobis, quóniam sic laborántes opórtet suscípere infírmos.\par
\switchcolumn
\EN
\noindent Sell what ye have and give alms. Fear not, He saith, lest, while ye fight for the kingdom of God, ye should lack such things as are needful for this life, nay rather, sell even that which ye have, and give alms. This doth, whosoever for the Lord's sake leaveth all that he hath, and then worketh with his hands, that so he may have to eat, and withal to give alms. In this doth the Apostle boast himself, saying I have coveted no man's silver, or gold, or apparel, as ye yourselves know for these hands have ministered unto my necessities, and to them that were with me. I have showed you all things, how that so labouring ye ought to support the weak. (Acts xx. 33, 34, 35.)\par
\switchcolumn*

\LA
\begin{Resp}\Rsign Sint lumbi vestri præcíncti, et lucérnæ ardéntes in mánibus vestris: \Star{} Et vos símiles homínibus exspectántibus dóminum suum, quando revertátur a núptiis.\end{Resp}
\begin{Resp}\Vsign Vigiláte ergo, quia nescítis qua hora Dóminus vester ventúrus sit.\end{Resp}
\begin{Resp}\Rsign Et vos símiles homínibus exspectántibus dóminum suum, quando revertátur a núptiis.\end{Resp}
\begin{Resp}\Vsign Glória Patri, et Fílio, \Star{} et Spirítui Sancto.\end{Resp}
\begin{Resp}\Rsign Et vos símiles homínibus exspectántibus dóminum suum, quando revertátur a núptiis.\end{Resp}
\switchcolumn
\EN
\begin{Resp}\Rsign Let your loins be girded about, and your lights burning; \Star{} And ye yourselves like unto men that wait for their lord, when he will return from the wedding.\end{Resp}
\begin{Resp}\Vsign Watch therefore, for ye know not what hour your Lord doth come.\end{Resp}
\begin{Resp}\Rsign And ye yourselves like unto men that wait for their lord, when he will return from the wedding.\end{Resp}
\begin{Resp}\Vsign Glory be to the Father, and to the Son, \Star{} and to the Holy Ghost.\end{Resp}
\begin{Resp}\Rsign And ye yourselves like unto men that wait for their lord, when he will return from the wedding.\end{Resp}
\switchcolumn*

\LA
\LessonHead{Lectio IX}
\switchcolumn
\EN
\LessonHead{Lesson IX}
\switchcolumn*

\LA
\noindent Fácite vobis sácculos qui non veteráscunt: eleemósynas vidélicet operándo quarum merces in ætérnum máneat. Ubi non hoc præcéptum esse putándum est, ut nil pecúniæ reservétur a sanctis, vel suis scílicet, vel páuperum úsibus suggeréndæ: cum et ipse Dóminus, cui ministrábant Ángeli, tamen ad informándam Ecclésiam suam lóculos habuísse legátur, et a fidélibus obláta consérvans, et suórum necessitátibus aliísque indigéntibus tríbuens; sed ne Deo propter ista serviátur, et ob inópiæ timórem justítia deserátur.\par
\switchcolumn
\EN
\noindent Provide yourselves bags which wax not old that is to say, by almsgiving, the reward thereof remaineth for ever. Nevertheless, we must not think here that this commandment forbiddeth the Saints to keep money for their own use, and for helping of the poor. The Lord Himself, to Whom Angels ministered, had a bag, and kept therein that which the faithful people gave unto Him (John xii. 6,) to relieve therewith the need of His disciples, and other poor folk. But we are commanded not to serve God for gain, nor to work unrighteousness for fear of poverty.\par
\switchcolumn*

\LA
\TeDeum{Te Deum}
\switchcolumn
\EN
\TeDeum{Te Deum}
\switchcolumn*

\end{paracol}

\par\needspace{10\baselineskip}\addvspace{1.2em}{\centering\small\textsc{Die 11 Septembris} · \textsc{11 September}\par}
\Office{Ss. Proti et Hyacinthi Martyrum}{Sts. Protus and Hyacinth Martyrs}{Simplex}
\addcontentsline{toc}{subsection}{Die 11 Septembris · Ss. Proti et Hyacinthi Martyrum \textit{— Sts. Protus and Hyacinth Martyrs}}
\begin{paracol}{2}
\LA
\RefLine{Lectiones I–II: de Scriptura occurrente.}
\switchcolumn
\EN
\RefLine{Lessons I–II: from the Scripture of the day.}
\switchcolumn*

\LA
\LessonHead{Lectio III (pro commemoratione)}
\switchcolumn
\EN
\LessonHead{Lesson III (when commemorated)}
\switchcolumn*

\LA
\LessonTitle{Commemoratio pro Ss. Mart. Proto et Hyacintho.}
\switchcolumn
\EN
\LessonTitle{Commemoration for the holy Martyrs Protus and Hyacinth}
\switchcolumn*

\LA
\noindent Protus et Hyacinthus fratres, beátæ Eugeniæ Vírginis eunuchi, una cum illa ab Héleno episcopo baptizati, ac stúdiis déditi divinárum litterárum, aliquámdiu in Ægypto inter ascetas mira humilitate et vitæ sanctitáte vixérunt. Sed póstea, sanctam Vírginem Eugeniam Romam prosecuti, Gallieno imperatóre, in Urbe propter christianæ fidei professiónem comprehénsi sunt. A quibus cum nullo modo impetrari posset ut, christianam religiónem deseréntes, deos cólerent, acerbis verbéribus cæsi, secúri feriúntur tertio Idus Septembris.\par
\switchcolumn
\EN
\noindent Protus and Hyacinth were brethren, eunuchs of the blessed Virgin Eugenia, and were baptized along with her by Bishop Helenus. They gave themselves to the study of God's Word, and dwelt for a while in wonderful lowliness and holiness of life in a monastery in Egypt. However, they afterwards followed the holy Virgin Eugenia to Rome, in the reign of the Emperor Gallienus, and were arrested in that city for professing the Christian faith. By no means could they be brought to leave the Christian religion and to worship the gods, and they were therefore severely scourged and beheaded, upon the 11th day of September.\par
\switchcolumn*

\LA
\TeDeum{Te Deum}
\switchcolumn
\EN
\TeDeum{Te Deum}
\switchcolumn*

\end{paracol}

\par\needspace{10\baselineskip}\addvspace{1.2em}{\centering\small\textsc{Die 12 Septembris} · \textsc{12 September}\par}
\Office{S. Nominis Beatæ Mariæ Virginis}{S. Nominis Beatæ Mariæ Virginis}{Duplex majus}
\addcontentsline{toc}{subsection}{Die 12 Septembris · S. Nominis Beatæ Mariæ Virginis \textit{— S. Nominis Beatæ Mariæ Virginis}}
\begin{paracol}{2}
\LA
\RefLine{Lectiones I–III: ex In Festis Beatae Mariae Virginis (C11).}
\switchcolumn
\EN
\RefLine{Lessons I–III: from the In Festis Beatae Mariae Virginis (C11).}
\switchcolumn*

\LA
\LessonHead{Lectio IV}
\switchcolumn
\EN
\LessonHead{Lesson IV}
\switchcolumn*

\LA
\LessonTitle{Sermo sancti Bernárdi Abbátis}
\switchcolumn
\EN
\LessonTitle{From the Sermons of St. Bernard, Abbot (of Clairvaux)}
\switchcolumn*

\LA
\Source{Ex Homil. 2 super Missus est, circa finem}
\switchcolumn
\EN
\Source{Second Homily on Luke i, 26}
\switchcolumn*

\LA
\noindent Et nomen, inquit, Vírginis Maria. Loquámur pauca, et super hoc nómine, quod interpretátum maris stella dícitur, et Matri Vírgini valde conveniénter aptátur. Ipsa namque aptíssime síderi comparátur, quia sicut sine sui corruptióne sidus suum emíttit rádium, sic absque sui læsióne Virgo parturívit Fílium. Nec síderi rádius suam mínuit claritátem, nec Vírgini Fílius suam integritátem. Ipsa est ígitur nóbilis illa stella ex Jacob orta, cujus rádius univérsum orbem illúminat, cujus splendor et præfúlget in supérnis, et ínferos pénetrat; terras étiam perlústrans, et calefáciens magis mentes quam córpora, fovet virtútes, éxcoquit vítia. Ipsa, inquam, est præclára et exímia stella super hoc mare magnum et spatiósum necessário subleváta, micans méritis, illústrans exémplis.\par
\switchcolumn
\EN
\noindent It is said: And the virgin's name was Mary. Let us speak a few words upon this name, which signifieth, being interpreted, Star of the Sea, and suiteth very well the Maiden Mother, who may very meetly be likened unto a star. A star giveth forth her rays without any harm to herself, and the Virgin brought forth her Son without any hurt to her virginity. The light of a star taketh nothing away from the star itself, and the birth of her offspring took nothing away from the Virginity of Mary. She is that noble star which was to come out of Jacob, \textasciitilde{}(Num. xxiv. 17,) whose brightness still sheddeth lustre upon all the earth, whose rays are most brilliant in heaven, and shine even unto hell, lighting up earth midway, and warming souls rather than bodies, fostering good and scaring away evil. She, I say, is a clear and shining star, twinkling with excellencies, and resplendent with example, needfully set to look down upon the surface of this great and wide sea.\par
\switchcolumn*

\LA
\begin{Resp}\Rsign Sicut cedrus exaltáta sum in Líbano, et sicut cypréssus in monte Sion: quasi myrrha elécta, \Star{} Dedi suavitátem odóris.\end{Resp}
\begin{Resp}\Vsign Et sicut cinnamómum et bálsamum aromatízans.\end{Resp}
\begin{Resp}\Rsign Dedi suavitátem odóris.\end{Resp}
\switchcolumn
\EN
\begin{Resp}\Rsign I was exalted like a cedar in Lebanon, and as a cypress-tree upon Mount Zion. Like the best myrrh \Star{} I yielded a pleasant odour.\end{Resp}
\begin{Resp}\Vsign Like cinnamon and sweet balsam.\end{Resp}
\begin{Resp}\Rsign I yielded a pleasant odour.\end{Resp}
\switchcolumn*

\LA
\LessonHead{Lectio V}
\switchcolumn
\EN
\LessonHead{Lesson V}
\switchcolumn*

\LA
\noindent O quisquis te intélligis in hujus sǽculi proflúvio magis inter procéllas et tempestátes fluctuáre, quam per terram ambuláre; ne avértas óculos a fulgóre hujus síderis, si non vis óbrui procéllis. Si insúrgant venti tentatiónum, si incúrras scópulos tribulatiónum, réspice stellam, voca Maríam. Si jactáris supérbiæ undis, si ambitiónis, si detractiónis, si æmulatiónis, réspice stellam, voca Maríam. Si iracúndia aut avarítia aut carnis illécebra navículam concússerit mentis, réspice ad Maríam. Si críminum immanitáte turbátus, consciéntiæ fœditáte confúsus, judícii horróre pertérritus, bárathro incípias absorbéri tristítiæ, desperatiónis abýsso, cógita Maríam.\par
\switchcolumn
\EN
\noindent Thou, whosoever thou art, that knowest thyself to be here not so much walking upon firm ground, as battered to and fro by the gales and storms of this life's ocean, if thou wouldest not be overwhelmed by the tempest, keep thine eyes fixed upon this star's clear shining. If the hurricanes of temptation rise against thee, or thou art running upon the rocks of trouble, look to the star, call on Mary. If the waves of pride, or ambition, or slander, or envy toss thee, look to the star, call on Mary. If the billows of anger or avarice, or the enticements of the flesh beat against thy soul's bark, look to Mary. If the enormity of thy sins trouble thee, if the foulness of thy conscience confound thee, if the dread of judgment appall thee, if thou begin to slip into the deep of despondency, into the pit of despair, think of Mary.\par
\switchcolumn*

\LA
\begin{Resp}\Rsign Quæ est ista quæ procéssit sicut sol, et formósa tamquam Jerúsalem? \Star{} Vidérunt eam fíliæ Sion, et beátam dixérunt, et regínæ laudavérunt eam.\end{Resp}
\begin{Resp}\Vsign Et sicut dies verni circúmdabant eam flores rosárum et lília convállium.\end{Resp}
\begin{Resp}\Rsign Vidérunt eam fíliæ Sion, et beátam dixérunt, et regínæ laudavérunt eam.\end{Resp}
\switchcolumn
\EN
\begin{Resp}\Rsign Who is this that cometh up like the sun? This, comely as Jerusalem? \Star{} The daughters of Zion saw her, and called her blessed: the queens also, and they praised her.\end{Resp}
\begin{Resp}\Vsign And about her it was as the flower of roses in the spring of the year, and lilies of the valleys.\end{Resp}
\begin{Resp}\Rsign The daughters of Zion saw her, and called her blessed the queens also, and they praised her.\end{Resp}
\switchcolumn*

\LA
\LessonHead{Lectio VI}
\switchcolumn
\EN
\LessonHead{Lesson VI}
\switchcolumn*

\LA
\noindent In perículis, in angústiis, in rebus dúbiis Maríam cógita, Maríam ínvoca. Non recédat ab ore, non recédat a corde; et, ut ímpetres ejus oratiónis suffrágium, non déseras conversatiónis exémplum. Ipsam sequens, non dévias; ipsam rogans, non despéras; ipsam cógitans, non erras; ipsa tenénte, non córruis; ipsa protegénte, non métuis; ipsa duce, non fatigáris; ipsa propítia, pérvenis: et sic in temetípso experíris quam mérito dictum sit: Et nomen Vírginis María. - Quod quidem venerábile nomen, jamprídem in quibúsdam christiáni orbis pártibus speciáli ritu cultum, Innocéntius undécimus Románus Póntifex, ob insígnem victóriam sub ejusdem Vírginis Maríæ præsídio de immaníssimo Turcárum tyránno, cervícibus pópuli christiáni insultánti, Viénnæ in Austria partam, et in perenne tanti benefícii monuméntum, in Ecclésia universáli síngulis annis celebrári præcépit.\par
\switchcolumn
\EN
\noindent In danger, in difficulty, or in doubt, think on Mary, call on Mary. Let her not be away from thy mouth or from thine heart, and that thou mayest not lack the succour of her prayers, turn not aside from the example of her conversation. If thou follow her, thou wilt never go astray. If thou pray to her, thou wilt never have need to despair. If thou keep her in mind, thou wilt never wander. If she hold thee, thou wilt never fall. If she lead thee, thou wilt never be weary. If she help thee, thou wilt reach home safe at the last and so thou wilt prove in thyself how meetly it is said And the virgin's name was Mary. \{Here ends St Bernard.) Particular honours were already paid to this worshipful name in diverse parts of the Christian world, but the Bishop of Rome, Innocent XI, ordered this Feast in honour of it to be held every year throughout the whole Church, upon the Lord's Day within the Octave of the Birthday of the Blessed Virgin Mary, as an everlasting thanksgiving for the great blessing that, under her protection, the brutal Sultan of the Turks, who was trampling upon the necks of the Christian population, was thoroughly beaten before the walls of Vienna, (upon the 12th day of September, in the year 1683.)\par
\switchcolumn*

\LA
\begin{Resp}\Rsign Ornátam monílibus fíliam Jerúsalem Dóminus concupívit: \Star{} Et vidéntes eam fíliæ Sion, beatíssimam prædicavérunt, dicéntes: * Unguéntum effúsum nomen tuum.\end{Resp}
\begin{Resp}\Vsign Astitit regína a dextris tuis in vestítu deauráto, circúmdata varietáte.\end{Resp}
\begin{Resp}\Rsign Et vidéntes eam fíliæ Sion, beatíssimam prædicavérunt, dicéntes:\end{Resp}
\begin{Resp}\Vsign Glória Patri, et Fílio, \Star{} et Spirítui Sancto.\end{Resp}
\begin{Resp}\Rsign Unguéntum effúsum nomen tuum.\end{Resp}
\switchcolumn
\EN
\begin{Resp}\Rsign When the Lord beheld the daughter of Jerusalem adorned with her jewels, He greatly desired her beauty \Star{} And when the daughters of Zion saw her, they cried out that she was most blessed, saying thy name is as ointment poured forth.\end{Resp}
\begin{Resp}\Vsign Upon thy right hand did stand the Queen in a vesture of gold wrought about with diverse colours.\end{Resp}
\begin{Resp}\Rsign And when the daughters of Zion saw her, they cried out that she was most blessed, saying:\end{Resp}
\begin{Resp}\Vsign Glory be to the Father, and to the Son, \Star{} and to the Holy Ghost.\end{Resp}
\begin{Resp}\Rsign Thy name is as ointment poured forth.\end{Resp}
\switchcolumn*

\LA
\LessonHead{Lectio VII}
\switchcolumn
\EN
\LessonHead{Lesson VII}
\switchcolumn*

\LA
\LessonTitle{Léctio sancti Evangélii secúndum Lucam}
\switchcolumn
\EN
\LessonTitle{From the Holy Gospel according to Luke}
\switchcolumn*

\LA
\Cite{Luc 1:26-38}
\switchcolumn
\EN
\Cite{Luke 1:26-38}
\switchcolumn*

\LA
\noindent In illo témpore: Missus est Angelus Gabriel a Deo in civitátem Galilææ, cui nomen Nazareth, ad Vírginem desponsatam viro, cui nomen erat Joseph, de domo David, et nomen Vírginis Maria. Et réliqua.\par
\switchcolumn
\EN
\noindent At that time: The Angel Gabriel was sent from God unto a city of Galilee, named Nazareth, to a Virgin espoused to a man whose name was Joseph, of the house of David; and the Virgin's name was Mary. And so on.\par
\switchcolumn*

\LA
\LessonTitle{Homilía sancti Petri Chrysólogi}
\switchcolumn
\EN
\LessonTitle{Homily by St. Peter Chrysologus, Archbishop (of Ravenna)}
\switchcolumn*

\LA
\Source{Sermo 142 de Annuntiatione}
\switchcolumn
\EN
\Source{142nd on the Annunciation}
\switchcolumn*

\LA
\noindent Audístis hódie, fratres caríssimi, Angelum cum muliere de hóminis reparatióne tractántem. Audístis agi, ut homo cursibus eisdem, quibus delapsus fuerat ad mortem, rediret ad vitam. Agit, agit cum Maria Angelus de salúte, quia cum Heva angelus egerat de ruína. Audístis Angelum de carnis nostræ limo templum divinæ majestátis arte ineffábili construentem. Audístis in terris Deum in cælis hóminem sacraménto incomprehensibili collocari. Audístis inaudíta ratióne in uno corpore Deum hominémque misceri. Audístis frágilem nostræ carnis natúram ad portandam totam Deitátis glóriam angelica exhortatióne roborari.\par
\switchcolumn
\EN
\noindent Dearly beloved brethren, ye have this day heard how an Angel treated with a woman touching the regeneration of mankind. Ye have heard how it was arranged that man should return to life by the same mean whereby he had fallen into death. The Angel treateth, treateth with Mary concerning salvation, because an angel had treated with Eve concerning destruction. Ye have heard how an Angel set about to raise with unspeakable building a temple of the Divine Majesty out of the dust of the earth. Ye have heard how by a mystery which cannot be understood, God got a place on earth and man a place in heaven. Ye have heard how by a working hitherto unheard of, God and man are joined together in one Body. Ye have heard how at the message of an angel, the weak nature whereof our flesh is sharer, became strong to bear the whole glory of the Godhead.\par
\switchcolumn*

\LA
\begin{Resp}\Rsign Felix namque es, sacra Virgo María, et omni laude digníssima: \Star{} Quia ex te ortus est sol justítiæ, Christus, Deus noster.\end{Resp}
\begin{Resp}\Vsign Ora pro pópulo, intérveni pro clero, intercéde pro devóto femíneo sexu: séntiant omnes tuum juvámen, quicúmque célebrant tui sancti nóminis commemoratiónem.\end{Resp}
\begin{Resp}\Rsign Quia ex te ortus est sol justítiæ, Christus, Deus noster.\end{Resp}
\switchcolumn
\EN
\begin{Resp}\Rsign O Holy Virgin Mary, happy indeed art thou, and right worthy of all praise \Star{} For out of thee rose the Sun of righteousness, even Christ our God.\end{Resp}
\begin{Resp}\Vsign Pray for the people, plead for the clergy, make intercession for all women vowed to God. May all that are keeping this Commemoration of thy Name feel the might of thine assistance.\end{Resp}
\begin{Resp}\Rsign For out of thee rose the Sun of righteousness, even Christ our God.\end{Resp}
\switchcolumn*

\LA
\LessonHead{Lectio VIII}
\switchcolumn
\EN
\LessonHead{Lesson VIII}
\switchcolumn*

\LA
\noindent Denique, ne tanto pónderi cæléstis fabricæ in Maria subtilis nostri corporis arena succumberet, et in Vírgine totius generis humani portatúra fructum, virga ténuis frangerétur; fugatúra metum vox Angeli mox præcessit, dicens: Ne timeas, Maria. Ante causam dignitas Vírginis annuntiátur ex nómine; nam Maria Hebræo sermóne, Latine Domina nuncupátur. Vocat ergo Angelus Dominam, ut Dominatoris Genitricem trepidátio déserat servitutis, quam nasci et vocari Dominam ipsa sui Gérminis fecit et impetrávit auctoritas. Ne timeas, Maria, invenísti enim grátiam. Verum est, quia, qui invénit grátiam, nescit timére: Invenísti grátiam.\par
\switchcolumn
\EN
\noindent Then, lest the frail clay of humanity should break down under the weight of God's work, and in Mary the tender stem should snap, which was about to bear the fruit of all mankind, the Angel's first words were a preventive against fear. And the Angel said unto her Fear not, Mary. Even before the matter is revealed, the exalted station of this Virgin is made clear by her very name, for the name Mary is an Hebrew word, and signifieth Lady. The Angel therefore greeteth her as Lady, that the Mother of the Lord may lay aside the fearfulness of His handmaiden, whom the will of her own Offspring had made to be born and to be called a Lady. Fear not, Mary, for thou hast found grace. He that hath found grace, need fear no more. Thou hast found grace.\par
\switchcolumn*

\LA
\begin{Resp}\Rsign Beátam me dicent omnes generatiónes, \Star{} Quia fecit mihi Dóminus magna qui potens est, et sanctum nomen ejus.\end{Resp}
\begin{Resp}\Vsign Et misericórdia ejus a progénie in progénies timéntibus eum.\end{Resp}
\begin{Resp}\Rsign Quia fecit mihi Dóminus magna qui potens est, et sanctum nomen ejus.\end{Resp}
\begin{Resp}\Vsign Glória Patri, et Fílio, \Star{} et Spirítui Sancto.\end{Resp}
\begin{Resp}\Rsign Quia fecit mihi Dóminus magna qui potens est, et sanctum nomen ejus.\end{Resp}
\switchcolumn
\EN
\begin{Resp}\Rsign All generations shall call me blessed. \Star{} For He that is mighty, even the Lord, hath done to me great things; and Holy is His Name.\end{Resp}
\begin{Resp}\Vsign And his mercy is on them that fear Him, from generation to generation.\end{Resp}
\begin{Resp}\Rsign He that is mighty, even the Lord, hath done to me great things; and Holy is His Name.\end{Resp}
\begin{Resp}\Vsign Glory be to the Father, and to the Son, \Star{} and to the Holy Ghost.\end{Resp}
\begin{Resp}\Rsign He that is mighty, even the Lord, hath done to me great things; and Holy is His Name.\end{Resp}
\switchcolumn*

\LA
\LessonHead{Lectio IX}
\switchcolumn
\EN
\LessonHead{Lesson IX}
\switchcolumn*

\LA
\noindent Beata, quæ inter hómines audíre sola méruit præ ómnibus: Invenísti grátiam. Quantam? Quantam supérius díxerat: plenam. Et vere plenam, quæ largo imbre totam fúnderet et infúnderet creatúram: Invenísti enim grátiam apud Deum. Hæc cum dicit, et ipse Angelus mirátur, aut féminam tantum, aut omnes hómines vitam meruísse per féminam: stupet Angelus totum Deum veníre intra virginális úteri angústias, cui tota simul angústa est creatúra. Hinc est quod remorátur Angelus, hinc est quod vírginem vocat de mérito, de grátia compéllat, vix causam prodit audiénti, sane ut sensum promóveat, vix longa trepidatióne compónit.\par
\switchcolumn
\EN
\noindent Blessed is she who alone among mankind, and first among mankind deserved to hear: Thou hast found grace. And how much grace? Even as the Angel had said: Full! Full of grace. Full indeed! Grace like a bountiful shower drenched and soaked her whole being. For thou hast found grace with God. As he saith this, even the Angel doth marvel. He marvelleth that a woman should merit eternal life, that all men should merit it through her. The Angel marvelleth that the whole Godhead, he to whom the entire universe is small, should enter the narrow womb of a virgin. So he delayeth. He calleth her virgin. That was her right. He haileth her as full of grace. Then, with great trepidation he delivereth his message, scarcely able to phrase it so that it could be understood.\par
\switchcolumn*

\LA
\TeDeum{Te Deum}
\switchcolumn
\EN
\TeDeum{Te Deum}
\switchcolumn*

\end{paracol}

\par\needspace{10\baselineskip}\addvspace{1.2em}{\centering\small\textsc{Die 14 Septembris} · \textsc{14 September}\par}
\Office{In Exaltatione Sanctæ Crucis}{In Exaltatione Sanctæ Crucis}{Duplex majus}
\addcontentsline{toc}{subsection}{Die 14 Septembris · In Exaltatione Sanctæ Crucis \textit{— In Exaltatione Sanctæ Crucis}}
\begin{paracol}{2}
\LA
\LessonHead{Lectio I}
\switchcolumn
\EN
\LessonHead{Lesson I}
\switchcolumn*

\LA
\LessonTitle{De libro Númeri}
\switchcolumn
\EN
\LessonTitle{From the Book of Numbers}
\switchcolumn*

\LA
\Cite{Num 21:1-3}
\switchcolumn
\EN
\Cite{Num 21:1-3}
\switchcolumn*

\LA
\noindent Cum audísset Chananǽus rex Arad, qui habitábat ad merídiem, venísse scílicet Israël per exploratórum viam, pugnávit contra illum et victor exsístens duxit ex eo prædam. At Israël, voto se Dómino óbligans, ait: Si tradíderis pópulum istum in manu mea, delébo urbes ejus. Exaudivítque Dóminus preces Israël, et trádidit Chananǽum, quem ille interfécit, subvérsis úrbibus ejus, et vocávit nomen loci illíus Horma, id est, anáthema.\par
\switchcolumn
\EN
\noindent And when king Arad the Chanaanite, who dwelt towards the south, had heard this, to wit, that Israel was come by the way of the spies, he fought against them, and overcoming them carried off their spoils. But Israel binding himself by vow to the Lord, said: If thou wilt deliver this people into my hand, I will utterly destroy their cities. And the Lord heard the prayers of Israel, and delivered up the Chanaanite, and they cut them off and destroyed their cities: and they called the name of that place Horma, that is to say, Anathema.\par
\switchcolumn*

\LA
\begin{Resp}\Rsign Gloriósum diem sacra venerátur Ecclésia, dum triumphále reserátur lignum: \Star{} In quo Redémptor noster, mortis víncula rumpens, cállidum áspidem superávit. (Allelúja, allelúja, allelúja.)\end{Resp}
\begin{Resp}\Vsign In ligno pendens nostræ salútis sémitam Verbum Patris invénit.\end{Resp}
\begin{Resp}\Rsign In quo Redémptor noster, mortis víncula rumpens, cállidum áspidem superávit. (Allelúja, allelúja, allelúja.)\end{Resp}
\switchcolumn
\EN
\begin{Resp}\Rsign Lo the Church, with solemn gladness, hails the day for ever glorious, when the opening earth revealeth that dread tree of mystic triumph. \Star{} On whose boughs her dying Saviour shattered death and crushed the serpent.\end{Resp}
\begin{Resp}\Vsign He the Word of God eternal, on those stately branches hanging, hath for us a new way opened.\end{Resp}
\begin{Resp}\Rsign On whose boughs her dying Saviour shattered death and crushed the serpent.\end{Resp}
\switchcolumn*

\LA
\LessonHead{Lectio II}
\switchcolumn
\EN
\LessonHead{Lesson II}
\switchcolumn*

\LA
\Cite{Num 21:4-6}
\switchcolumn
\EN
\Cite{Num 21:4-6}
\switchcolumn*

\LA
\noindent Profécti sunt autem et de monte Hor per viam quæ ducit ad Mare Rubrum, ut circumírent terram Edom. Et tædére cœpit pópulum itíneris ac labóris. Locutúsque contra Deum et Móysen ait: Cur eduxísti nos de Ægýpto ut morerémur in solitúdine? Deest panis, non sunt aquæ, ánima nostra jam náuseat super cibo isto levíssimo. Quam ob rem misit Dóminus in pópulum ignítos serpéntes.\par
\switchcolumn
\EN
\noindent And they marched from mount Hor, by the way that leadeth to the Red Sea, to compass the land of Edom. And the people began to be weary of their journey and labour: And speaking against God and Moses, they said: Why didst thou bring us out of Egypt, to die in the wilderness? There is no bread, nor have we any waters: our soul now loatheth this very light food. Wherefore the Lord sent among the people fiery serpents.\par
\switchcolumn*

\LA
\begin{Resp}\Rsign Crux fidélis, inter omnes arbor una nóbilis: nulla silva talem profert, fronde, flore, gérmine: \Star{} Dulce lignum, dulces clavos, dulce pondus sustínuit. (Allelúja.)\end{Resp}
\begin{Resp}\Vsign Super ómnia ligna cedrórum tu sola excélsior.\end{Resp}
\begin{Resp}\Rsign Dulce lignum, dulces clavos, dulce pondus sustínuit. (Allelúja.)\end{Resp}
\switchcolumn
\EN
\begin{Resp}\Rsign Faithful Cross, above all other, one and only noble tree None in foliage, none in blossom, none in fruit thy peers may be \Star{} Sweetest wood and sweetest iron, Sweetest weight is hung on thee.\end{Resp}
\begin{Resp}\Vsign Thou art higher than all cedars.\end{Resp}
\begin{Resp}\Rsign Sweetest wood, and sweetest iron, Sweetest weight is hung on thee.\end{Resp}
\switchcolumn*

\LA
\LessonHead{Lectio III}
\switchcolumn
\EN
\LessonHead{Lesson III}
\switchcolumn*

\LA
\Cite{Num 21:6-9}
\switchcolumn
\EN
\Cite{Num 21:6-9}
\switchcolumn*

\LA
\noindent Ad quorum plagas et mortes plurimórum venérunt ad Móysen atque dixérunt: Peccávimus, quia locúti sumus contra Dóminum et te: ora ut tollat a nobis serpéntes. Oravítque Móyses pro pópulo. Et locútus est Dóminus ad eum: Fac serpéntem ǽneum et pone eum pro signo: qui percússus aspéxerit eum, vivet. Fecit ergo Móyses serpéntem ǽneum et pósuit eum pro signo; quem cum percússi aspícerent, sanabántur.\par
\switchcolumn
\EN
\noindent The serpents bit them and killed many of them. Upon which they came to Moses, and said: We have sinned, because we have spoken against the Lord and thee: pray that he may take away these serpents from us. And Moses prayed for the people. And the Lord said to him: Make a brazen serpent, and set it up for a sign: whosoever being struck shall look on it, shall live. Moses therefore made a brazen serpent, and set it up for a sign: which when they that were bitten looked upon, they were healed.\par
\switchcolumn*

\LA
\begin{Resp}\Rsign Hæc est arbor digníssima, in paradísi médio situáta, \Star{} In qua salútis auctor própria morte mortem ómnium superávit.\end{Resp}
\begin{Resp}\Vsign Crux præcellénti decóre fúlgida, quam Heraclíus imperátor concupiscénti ánimo recuperávit.\end{Resp}
\begin{Resp}\Rsign In qua salútis auctor própria morte mortem ómnium superávit.\end{Resp}
\begin{Resp}\Vsign Glória Patri, et Fílio, \Star{} et Spirítui Sancto.\end{Resp}
\begin{Resp}\Rsign In qua salútis auctor própria morte mortem ómnium superávit.\end{Resp}
\switchcolumn
\EN
\begin{Resp}\Rsign This is that noble tree, planted in the midst of the garden \Star{} Whereon the Author of our salvation did by His Own death openly triumph over the death of all men.\end{Resp}
\begin{Resp}\Vsign Even the Cross, whereof the glory is so excellent, and which the Emperor Heraclius did so eagerly rescue.\end{Resp}
\begin{Resp}\Rsign Whereon the Author of our salvation did by His Own death openly triumph over the death of all men.\end{Resp}
\begin{Resp}\Vsign Glory be to the Father, and to the Son, \Star{} and to the Holy Ghost.\end{Resp}
\begin{Resp}\Rsign Whereon the Author of our salvation did by His Own death openly triumph over the death of all men.\end{Resp}
\switchcolumn*

\LA
\LessonHead{Lectio IV}
\switchcolumn
\EN
\LessonHead{Lesson IV}
\switchcolumn*

\LA
\noindent Chósroas, Persárum rex, extrémis Phocæ impérii tempóribus, Ægýpto et Africa occupáta ac Jerosólyma capta multísque ibi cæsis Christianórum míllibus, Christi Dómini Crucem, quam Hélena in monte Calváriæ collocárat, in Pérsidem ábstulit. Itaque Heraclíus, qui Phocæ succésserat, multis belli incómmodis et calamitátibus afféctus, pacem petébat; quam a Chósroa, victóriis insolénte, ne iníquis quidem conditiónibus impetráre póterat. Quare in summo discrímine se assíduis jejúniis et oratiónibus exércens, opem a Deo veheménter implorábat; cujus mónitu exércitu comparáto, signa cum hoste cóntulit, ac tres duces Chósroæ cum tribus exercítibus superávit.\par
\switchcolumn
\EN
\noindent Chosroes of Persia, having, in the last days of the reign of the Emperor Phocas, overrun Egypt and Africa, (in 614,) took Jerusalem, where he slaughtered thousands of Christians and carried off to Persia the Cross of the Lord, which Helen had put upon Mount Calvary. Heraclius, the successor of Phocas, moved by the thought of the hardships and horrid outrages of war, sought for peace, but Chosroes, drunken with conquest, would not allow of it even upon unfair terms. Heraclius therefore, being set in this uttermost strait, earnestly sought help from God by constant fasting and prayer, and through His good inspiration gathered an army, joined battle with the enemy, and prevailed against three of Chosroes' chief captains, and three armies.\par
\switchcolumn*

\LA
\begin{Resp}\Rsign Nos autem gloriári opórtet in Cruce Dómini nostri Jesu Christi, in quo est salus, vita, et resurréctio nostra: \Star{} Per quem salváti et liberáti sumus. (Allelúja.)\end{Resp}
\begin{Resp}\Vsign Tuam Crucem adorámus, Dómine, et recólimus tuam gloriósam passiónem.\end{Resp}
\begin{Resp}\Rsign Per quem salváti et liberáti sumus. (Allelúja.)\end{Resp}
\switchcolumn
\EN
\begin{Resp}\Rsign But us it behoveth to glory in the Cross of our Lord Jesus Christ, in Whom is our salvation, life, and resurrection \Star{} Who hath saved us and redeemed us.\end{Resp}
\begin{Resp}\Vsign O Lord, we worship thy Cross, and make memorial of thy glorious passion.\end{Resp}
\begin{Resp}\Rsign Who hath saved us and redeemed us.\end{Resp}
\switchcolumn*

\LA
\LessonHead{Lectio V}
\switchcolumn
\EN
\LessonHead{Lesson V}
\switchcolumn*

\LA
\noindent Quibus cládibus fractus Chósroas, in fuga, qua traícere Tigrim parábat, Medársen fílium sócium regni desígnat. Sed eam contuméliam cum Síroës, Chósroæ major natu fílius, ferret atróciter, patri simul et fratri necem machinátur; quam paulo post utríque ex fuga retrácto áttulit, regnúmque ab Heraclío impetrávit quibúsdam accéptis conditiónibus, quarum ea prima fuit, ut Crucem Christi Dómini restitúeret. Ergo Crux, quatuórdecim annis postquam vénerat in potestátem Persárum, recépta est. Quam rédiens Jerosólymam Heraclíus solémni celebritáte suis húmeris rétulit in eum montem, quo eam Salvátor túlerat.\par
\switchcolumn
\EN
\noindent Chosroes was broken by these defeats, and when in his flight, (in 628,) he was about crossing the Tigris, he proclaimed his son Medarses partner in his kingdom. Chosroes' eldest son Siroes took this slight to heart, and formed a plot to murder his father and brother, which plot he brought to effect soon after they had come home. Then he got the kingdom from Heraclius upon certain terms, whereof the first was that he should give back the Cross of the Lord Christ. The Cross therefore was received back after that it had been fourteen years in the power of the Persians, and (in 629) Heraclius came to Jerusalem and bore it with solemn pomp unto the Mount whereunto the Saviour had borne it.\par
\switchcolumn*

\LA
\begin{Resp}\Rsign Dum sacrum pignus cǽlitus exaltátur, Christi fides roborátur: \Star{} Adsunt prodígia divína in virga Móysi prímitus figuráta. (Allelúja, allelúja.)\end{Resp}
\begin{Resp}\Vsign Ad Crucis contáctum resúrgunt mórtui, et Dei magnália reserántur.\end{Resp}
\begin{Resp}\Rsign Adsunt prodígia divína in virga Móysi prímitus figuráta. (Allelúja, allelúja.)\end{Resp}
\switchcolumn
\EN
\begin{Resp}\Rsign The Relique true from heaven revealed, hath now the Gospel's figure sealed \Star{} As by the serpent Moses reared, so by the Cross the sick are healed.\end{Resp}
\begin{Resp}\Vsign When the dead touch the Cross they arise, and the wonderful works of God are made manifest.\end{Resp}
\begin{Resp}\Rsign As by the serpent Moses reared, so by the Cross the sick are healed.\end{Resp}
\switchcolumn*

\LA
\LessonHead{Lectio VI}
\switchcolumn
\EN
\LessonHead{Lesson VI}
\switchcolumn*

\LA
\noindent Quod factum illústri miráculo commendátum est. Nam Heraclíus, ut erat auro et gemmis ornátus, insístere coáctus est in porta, quæ ad Calváriæ montem ducébat. Quo enim magis prógredi conabátur, eo magis retinéri videbátur. Cumque ea re et ipse Heraclíus et réliqui omnes obstupéscerent; Zacharías, Jerosolymórum antístes, Vide, inquit, imperátor, ne isto triumpháli ornátu, in Cruce ferénda parum Jesu Christi paupertátem et humilitátem imitére. Tum Heraclíus, abjécto amplíssimo vestítu detractísque cálceis ac plebéjo amíctu indútus, réliquum viæ fácile confécit, et in eódem Calváriæ loco Crucem státuit, unde fúerat a Persis asportáta. Itaque Exaltatiónis sanctæ Crucis solémnitas, quæ hac die quotánnis celebrabátur, illústrior habéri cœpit ob ejus rei memóriam, quod ibídem fúerit repósita ab Heraclío, ubi Salvatóri primum fúerat constitúta.\par
\switchcolumn
\EN
\noindent This event was marked by a famous miracle. Heraclius, who was adorned with gold and jewels, stayed perforce at the gateway which leadeth unto Mount Calvary, and the harder he strove to go forward, the harder he seemed to be held back, whereat both himself and all they that stood by were sore amazed. Then spake Zacharias, Patriarch of Jerusalem, saying: See, O Emperor, that it be not that in carrying the Cross attired in the guise of a Conqueror thou showest too little of the poverty and lowliness of Jesus Christ. Then Heraclius cast away his princely raiment and took off his shoes from his feet, and in the garb of a countryman easily finished his journey, and set up the Cross once more in the same place upon Calvary whence the Persians had carried it away. That the Cross had been put by Heraclius in the same place wherein it had first been planted by the Saviour caused the yearly Feast of the Exaltation of the Holy Cross to become the more famous thenceforward.\par
\switchcolumn*

\LA
\begin{Resp}\Rsign Hoc signum Crucis erit in cælo, cum Dóminus ad judicándum vénerit: \Star{} Tunc manifésta erunt abscóndita cordis nostri. (Allelúja, allelúja.)\end{Resp}
\begin{Resp}\Vsign Cum séderit Fílius hóminis in sede majestátis suæ, et cœ́perit judicáre sǽculum per ignem.\end{Resp}
\begin{Resp}\Rsign Tunc manifésta erunt abscóndita cordis nostri. (Allelúja, allelúja.)\end{Resp}
\begin{Resp}\Vsign Glória Patri, et Fílio, \Star{} et Spirítui Sancto.\end{Resp}
\begin{Resp}\Rsign Tunc manifésta erunt abscóndita cordis nostri. (Allelúja, allelúja.)\end{Resp}
\switchcolumn
\EN
\begin{Resp}\Rsign This Sign of the Cross shall be in heaven, when the Lord cometh to judgment. \Star{} Then shall the secrets of our hearts be made manifest.\end{Resp}
\begin{Resp}\Vsign When the Son of Man shall sit in the throne of His glory, and shall begin to judge the world by fire.\end{Resp}
\begin{Resp}\Rsign Then shall the secrets of our hearts be made manifest.\end{Resp}
\begin{Resp}\Vsign Glory be to the Father, and to the Son, \Star{} and to the Holy Ghost.\end{Resp}
\begin{Resp}\Rsign Then shall the secrets of our hearts be made manifest.\end{Resp}
\switchcolumn*

\LA
\LessonHead{Lectio VII}
\switchcolumn
\EN
\LessonHead{Lesson VII}
\switchcolumn*

\LA
\LessonTitle{Léctio sancti Evangélii secúndum Joánnem}
\switchcolumn
\EN
\LessonTitle{From the Holy Gospel according to John}
\switchcolumn*

\LA
\Cite{Joannes 12:31-36}
\switchcolumn
\EN
\Cite{John 12:31-36}
\switchcolumn*

\LA
\noindent In illo témpore: Dixit Jesus turbis Judæórum: Nunc judícium est mundi, nunc princeps hujus mundi eiciétur foras. Et réliqua.\par
\switchcolumn
\EN
\noindent At that time, Jesus said unto the multitudes of the Jews: Now is the judgment of this world, now shall the prince of this world be cast out. And so on.\par
\switchcolumn*

\LA
\LessonTitle{Homilía sancti Leónis Papæ}
\switchcolumn
\EN
\LessonTitle{Homily by Pope St. Leo the Great.}
\switchcolumn*

\LA
\Source{Sermo 8 de Passione Domini, post medium}
\switchcolumn
\EN
\Source{8th on the Lord's Passion.}
\switchcolumn*

\LA
\noindent Exaltáto, dilectíssimi, per Crucem Christo, non illa tantum spécies aspéctui mentis occúrrat, quæ fuit in óculis impiórum, quibus per Móysen dictum est: Et erit pendens vita tua ante óculos tuos, et timébis die ac nocte, et non credes vitæ tuæ. Isti enim nihil in crucifíxo Dómino præter fácinus suum cogitáre potuérunt, habéntes timórem, non quo fides vera justificátur, sed quo consciéntia iníqua torquétur. Noster vero intelléctus, quem spíritus veritátis illúminat, glóriam Crucis, cælo terráque radiántem, puro ac líbero corde suscípiat; et interióre ácie vídeat, quale sit quod Dóminus, cum de passiónis suæ loquerétur instántia, dixit: Nunc judícium mundi est, nunc princeps hujus mundi eiciétur foras. Et ego, si exaltátus fúero a terra, ómnia traham ad meípsum.\par
\switchcolumn
\EN
\noindent Dearly beloved brethren, when we gaze upon Christ lifted up upon the Cross, the eyes of our mind see more than that which appeared before the wicked, unto whom it was said through Moses: And thy life shall hang in doubt before thee, and thou shalt fear day and night, and shalt have none assurance of thy life. (Deut. xxviii 66.) They saw in the crucified Lord nothing but the work of their own wickedness, and they feared greatly, (Matth. xxvii. 54,) not with that faith which giveth earnest of life by justification, but with that whereby the evil conscience is tortured. But our understanding is enlightened by the Spirit of truth, and with pure and open hearts we see the glory of the Cross shining over heaven and earth, and discern by inward glance what the Lord meant when His Passion was nigh at hand, and He said: Now is the judgment of this world, now shall the prince of this world be cast out. And I, if I be lifted up from the earth, will draw all things unto Me.\par
\switchcolumn*

\LA
\begin{Resp}\Rsign Dulce lignum, dulces clavos, dulce pondus sustínuit: \Star{} Quæ sola digna fuit portáre prétium hujus sǽculi. (Allelúja.)\end{Resp}
\begin{Resp}\Vsign Hoc signum Crucis erit in cælo cum Dóminus ad judicándum vénerit.\end{Resp}
\begin{Resp}\Rsign Quæ sola digna fuit portáre prétium hujus sǽculi. (Allelúja.)\end{Resp}
\switchcolumn
\EN
\begin{Resp}\Rsign Sweetest wood and sweetest iron, Sweetest weight is hung on thee \Star{} Thou alone wast counted worthy this world's ransom to uphold.\end{Resp}
\begin{Resp}\Vsign The sign of the Cross shall be in heaven when the Lord cometh to judgment.\end{Resp}
\begin{Resp}\Rsign Thou alone wast counted worthy this world's ransom to uphold.\end{Resp}
\switchcolumn*

\LA
\LessonHead{Lectio VIII}
\switchcolumn
\EN
\LessonHead{Lesson VIII}
\switchcolumn*

\LA
\noindent O admirábilis poténtia Crucis! o ineffábilis glória Passiónis, in qua et tribúnal Dómini, et judícium mundi, et potéstas est Crucifíxi! Traxísti enim, Dómine, ómnia ad te, et cum expandísses tota die manus tuas ad pópulum non credéntem et contradicéntem, tibi, confiténdæ majestátis tuæ sensum totus mundus accépit. Traxísti, Dómine, ómnia ad te, cum in exsecratiónem Judáici scéleris, unam protulérunt ómnia eleménta senténtiam; cum, obscurátis lumináribus cæli et convérso in noctem die, terra quoque mótibus quaterétur insólitis, univérsaque creatúra impiórum úsui se negáret. Traxísti, Dómine, ómnia ad te, quóniam, scisso templi velo, Sancta sanctórum ab indígnis pontifícibus recessérunt; ut figúra in veritátem, prophetía in manifestatiónem, et lex in Evangélium verterétur.\par
\switchcolumn
\EN
\noindent How wonderful is the power of the Cross! O how unutterable is the glory of the Passion, wherein standeth the Lord's judgment-seat, and the judgment of this world, and the might of the Crucified! Lord! Thou hast drawn all things unto thee! Thou didst spread out thine Hands all the day unto an unbelieving and gainsaying people, (Isa. lxv. 2,) but the world hath felt and owned thy Majesty! Lord! Thou hast drawn all things unto thee! All the elements gave one wild cry of horror at the iniquity of the Jews the lights of the firmament were darkened, day turned into night, earth quaked with strange tremblings, and all God's work refused to serve the guilty. Lord! Thou hast drawn all things unto thee! The veil of the Temple was rent in twain from the top to the bottom, the Holy of Holies denied itself as a Sanctuary for the ministration of unworthy Priests, that the shadow might be changed for the substance, prophecy for realization, and the Law for the Gospel.\par
\switchcolumn*

\LA
\begin{Resp}\Rsign Sicut Móyses exaltávit serpéntem in desérto, ita exaltári opórtet Fílium hóminis: \Star{} Ut omnis qui credit in ipsum, non péreat, sed hábeat vitam ætérnam. (Allelúja.)\end{Resp}
\begin{Resp}\Vsign Non misit Deus Fílium suum in mundum ut júdicet mundum, sed ut salvétur mundus per ipsum.\end{Resp}
\begin{Resp}\Rsign Ut omnis qui credit in ipsum, non péreat, sed hábeat vitam ætérnam. (Allelúja.)\end{Resp}
\begin{Resp}\Vsign Glória Patri, et Fílio, \Star{} et Spirítui Sancto.\end{Resp}
\begin{Resp}\Rsign Ut omnis qui credit in ipsum, non péreat, sed hábeat vitam ætérnam. (Allelúja.)\end{Resp}
\switchcolumn
\EN
\begin{Resp}\Rsign As Moses lifted up the serpent in the wilderness, even so must the Son of Man be lifted up \Star{} That whosoever believeth in Him should not perish, but have everlasting life.\end{Resp}
\begin{Resp}\Vsign God sent not His Son into the world to condemn the world, but that the world through Him might be saved.\end{Resp}
\begin{Resp}\Rsign That whosoever believeth in Him should not perish, but have everlasting life.\end{Resp}
\begin{Resp}\Vsign Glory be to the Father, and to the Son, \Star{} and to the Holy Ghost.\end{Resp}
\begin{Resp}\Rsign That whosoever believeth in Him should not perish, but have everlasting life.\end{Resp}
\switchcolumn*

\LA
\LessonHead{Lectio IX}
\switchcolumn
\EN
\LessonHead{Lesson IX}
\switchcolumn*

\LA
\noindent Traxísti, Dómine, ómnia ad te, ut, quod in uno Judǽæ templo obumbrátis significatiónibus tegebátur, pleno apertóque sacraménto universárum ubíque natiónum devótio celebráret. Nunc étenim et ordo clárior levitárum, et dígnitas ámplior seniórum, et sacrátior est únctio sacerdótum: quia Crux tua ómnium fons benedictiónum, ómnium est causa gratiárum; per quam credéntibus datur virtus de infirmitáte, glória de oppróbrio, vita de morte. Nunc étiam, carnálium sacrificiórum varietáte cessánte, omnes differéntias hostiárum una córporis et sánguinis tui implet oblátio: quóniam tu es verus Agnus Dei, qui tollis peccáta mundi; et ita in te univérsa pérficis mystéria, ut sicut unum est pro omni víctima sacrifícium, ita unum de omni gente sit regnum.\par
\switchcolumn
\EN
\noindent Lord! Thou hast drawn all things unto thee! That which was veiled under types and shadows in the one Jewish Temple, is hailed by the love of all peoples in full and open worship. There is now a higher order of Levites, a more honourable rank of elders, a Priesthood with an holier anointing. thy Cross is a well of blessings for all, and a cause of thanksgiving for all. Thereby for them that believe in thee, weakness is turned into strength, shame into glory, and death into life. The changing ordinance of diverse carnal sacrifices is gone; the one oblation of thy Body and Blood fulfilleth them all. For Thou art the Very Paschal Lamb, Which takest away the sins of the world, and art in thyself all offerings finished. And even as Thou art the One Sacrifice Which taketh the place of all sacrifices, so may thy kingdom be one kingdom established over all peoples.\par
\switchcolumn*

\LA
\TeDeum{Te Deum}
\switchcolumn
\EN
\TeDeum{Te Deum}
\switchcolumn*

\end{paracol}

\par\needspace{10\baselineskip}\addvspace{1.2em}{\centering\small\textsc{Die 15 Septembris} · \textsc{15 September}\par}
\Office{Septem Dolorum Beatæ Mariæ Virginis}{Septem Dolorum Beatæ Mariæ Virginis}{Duplex II classis}
\addcontentsline{toc}{subsection}{Die 15 Septembris · Septem Dolorum Beatæ Mariæ Virginis \textit{— Septem Dolorum Beatæ Mariæ Virginis}}
\begin{paracol}{2}
\LA
\RefLine{Lectio IX: de S. Nicomedis Martyris (09-15cc).}
\switchcolumn
\EN
\RefLine{Lesson IX: of St. Nicomedes, Martyr (09-15cc).}
\switchcolumn*

\LA
\LessonHead{Lectio I}
\switchcolumn
\EN
\LessonHead{Lesson I}
\switchcolumn*

\LA
\LessonTitle{De Jeremía Prophéta}
\switchcolumn
\EN
\LessonTitle{Lesson from the book of Lamentations}
\switchcolumn*

\LA
\Cite{Lam. 1:2; 1:20-21}
\switchcolumn
\EN
\Cite{Lam 1:2; 1:20-21}
\switchcolumn*

\LA
\noindent (Beth) Plorans plorávit in nocte, et lácrimæ ejus in maxíllis ejus: non est qui consolétur eam ex ómnibus caris ejus: omnes amíci ejus sprevérunt eam, et facti sunt ei inimíci. (Res) Vide, Dómine, quóniam tríbulor, conturbátus est venter meus, subvérsum est cor meum in memetípsa, quóniam amaritúdine plena sum. Foris intérfecit gládius, et domi mors símilis est. (Sin) Audiérunt quia ingemísco ego, et non est qui consolétur me.\par
\switchcolumn
\EN
\noindent Weeping she hath wept in the night, and her tears are on her cheeks: there is none to comfort her among all them that were dear to her: all her friends have despised her, and are become her enemies. Behold, O Lord, for I am in distress, my bowels are troubled: my heart is turned within me, for I am full of bitterness: abroad the sword destroyeth, and at home there is death alike. They have heard that I sigh, and there is none to comfort me.\par
\switchcolumn*

\LA
\begin{Resp}\Rsign Símeon, vir justus et timorátus, dixit ad Maríam: \Star{} Tuam ipsíus ánimam pertransíbit gládius.\end{Resp}
\begin{Resp}\Vsign Ne vocétis me pulchram, sed amáram, quia amaritúdine valde replévit me Omnípotens.\end{Resp}
\begin{Resp}\Rsign Tuam ipsíus ánimam pertransíbit gládius.\end{Resp}
\switchcolumn
\EN
\begin{Resp}\Rsign There was a man whose name was Simeon; and the same man was just and devout; and he said unto Mary \Star{} Yea, a sword shall pierce through thine own soul also.\end{Resp}
\begin{Resp}\Vsign Call me not My-pleasantness, but call me Embittered, for the Almighty hath dealt very bitterly with me.\end{Resp}
\begin{Resp}\Rsign Yea, a sword shall pierce through thine own soul also.\end{Resp}
\switchcolumn*

\LA
\LessonHead{Lectio II}
\switchcolumn
\EN
\LessonHead{Lesson II}
\switchcolumn*

\LA
\Cite{Lam. 2:13, 15-16}
\switchcolumn
\EN
\Cite{Lam 2:13, 15-16}
\switchcolumn*

\LA
\noindent (Mem) Cui comparábo te? vel cui assimilábo te, fília Jerúsalem? cui exæquábo te, et consolábor te, virgo fília Sion? Magna est enim velut mare contrítio tua: quis medébitur tui? (Samech) Plausérunt super te mánibus omnes transeúntes per viam: sibilavérunt, et movérunt caput suum super fíliam Jerúsalem: Hǽccine est urbs, dicéntes, perfécti decóris, gáudium univérsæ terræ? (Phe) Aperuérunt super te os suum omnes inimíci tui, sibilavérunt et fremuérunt déntibus, et dixérunt: Devorábimus.\par
\switchcolumn
\EN
\noindent To what shall I compare thee? or to what shall I liken thee, O daughter of Jerusalem? to what shall I equal thee, that I may comfort thee, O virgin daughter of Sion? for great as the sea is thy destruction: who shall heal thee? All they that passed by the way have clapped their hands at thee: they have hissed, and wagged their heads at the daughter of Jerusalem, saying: Is this the city of perfect beauty, the joy of all the earth? All thy enemies have opened their mouth against thee: they have hissed, and gnashed with the teeth, and have said: We will swallow her up.\par
\switchcolumn*

\LA
\begin{Resp}\Rsign Surge, et áccipe Púerum et Matrem ejus, et fuge in Ægýptum; \Star{} Et esto ibi usque dum dicam tibi.\end{Resp}
\begin{Resp}\Vsign Vocávi fílium meum ex Ægýpto, ut véniat salus in Israël.\end{Resp}
\begin{Resp}\Rsign Et esto ibi usque dum dicam tibi.\end{Resp}
\switchcolumn
\EN
\begin{Resp}\Rsign Arise, and take the young Child and His Mother, and flee into Egypt; \Star{} And be thou there until I bring thee word.\end{Resp}
\begin{Resp}\Vsign Out of Egypt have I called My Son, that salvation may come unto Israel.\end{Resp}
\begin{Resp}\Rsign And be thou there until I bring thee word.\end{Resp}
\switchcolumn*

\LA
\LessonHead{Lectio III}
\switchcolumn
\EN
\LessonHead{Lesson III}
\switchcolumn*

\LA
\Cite{Lam. 2:17-18}
\switchcolumn
\EN
\Cite{Lam 2:17-18}
\switchcolumn*

\LA
\noindent (Ain) Fecit Dóminus quæ cogitávit, complévit sermónem suum quem præcéperat a diébus antíquis: destrúxit et non pepércit, et lætificávit super te inimícum et exaltávit cornu hóstium tuórum. (Sade) Clamávit cor eórum ad Dóminum super muros fíliæ Sion: Deduc quasi torréntem lácrimas per diem et noctem; non des réquiem tibi, neque táceat pupílla óculi tui.\par
\switchcolumn
\EN
\noindent The Lord hath done that which he purposed, he hath fulfilled his word, which he commanded in the days of old: he hath destroyed, and hath not spared, and he hath caused the enemy to rejoice over thee, and hath set up the horn of thy adversaries. Their heart cried to the Lord upon the walls of the daughter of Sion: Let tears run down like a torrent day and night: give thyself no rest, and let not the apple of thy eye cease.\par
\switchcolumn*

\LA
\begin{Resp}\Rsign Fili, quid fecísti nobis sic? \Star{} Ego et pater tuus * Doléntes quærebámus te.\end{Resp}
\begin{Resp}\Vsign Quid est quod me quærebátis? In his quæ Patris mei sunt, opórtet me esse.\end{Resp}
\begin{Resp}\Rsign Ego et pater tuus.\end{Resp}
\begin{Resp}\Vsign Glória Patri, et Fílio, \Star{} et Spirítui Sancto.\end{Resp}
\begin{Resp}\Rsign Doléntes quærebámus te.\end{Resp}
\switchcolumn
\EN
\begin{Resp}\Rsign Son, why hast Thou thus dealt with us? \Star{} I and thy father have sought thee sorrowing.\end{Resp}
\begin{Resp}\Vsign How is it that ye sought Me? Wist ye not that I must be about My Father's business?\end{Resp}
\begin{Resp}\Rsign I and thy Father\end{Resp}
\begin{Resp}\Vsign Glory be to the Father, and to the Son, \Star{} and to the Holy Ghost.\end{Resp}
\begin{Resp}\Rsign Have sought thee sorrowing.\end{Resp}
\switchcolumn*

\LA
\LessonHead{Lectio IV}
\switchcolumn
\EN
\LessonHead{Lesson IV}
\switchcolumn*

\LA
\LessonTitle{Sermo sancti Bernárdi Abbátis}
\switchcolumn
\EN
\LessonTitle{From the Sermons of St. Bernard, Abbot (of Clairvaux.)}
\switchcolumn*

\LA
\Source{Sermo de duodecim stellis}
\switchcolumn
\EN
\Source{On the twelve stars}
\switchcolumn*

\LA
\noindent Martýrium Vírginis tam in Simeónis prophetía, quam in ipsa Domínicæ passiónis história commendátur. Pósitus est hic (ait sanctus senex de púero Jesu) in signum cui contradicétur; et tuam ipsíus ánimam (ad Maríam autem dicébat) pertransíbit gládius. Vere tuam, o beáta Mater, ánimam pertransívit. Alióquin non nisi eam pertránsiens, carnem Fílii tui penetráret. Et quidem posteáquam emísit spíritum tuus ille Jesus, ipsíus plane non áttigit ánimam crudélis láncea, quæ ipsíus apéruit latus, sed tuam útique ánimam pertransívit. Ipsíus nimírum ánima jam ibi non erat, sed tua plane inde nequíbat avélli.\par
\switchcolumn
\EN
\noindent The Martyrdom of the Virgin is set before us, not only in the prophecy of Simeon, but also in the story itself of the Lord's Passion. The holy old man said of the Child Jesus (Luke ii. 34,) Behold, this Child is set for the fall and the rising again of many in Israel; and for a sign which shall be spoken against; yea, said he unto Mary, a sword shall pierce through thine own soul also Even so, O Blessed Mother! The sword did indeed pierce through thy soul! for nought could pierce the Body of thy Son, nor pierce thy soul likewise. Yea, and when this Jesus of thine had given up the ghost, and the bloody spear could torture Him no more, thy soul winced as it pierced His dead Side His Own Soul might leave Him, but thine could not.\par
\switchcolumn*

\LA
\begin{Resp}\Rsign Jesum bajulántem sibi crucem \Star{} Sequebátur turba mulíerum, quæ plangébant et lamentabántur eum.\end{Resp}
\begin{Resp}\Vsign Fíliæ Jerúsalem, super vos ipsas flete et super fílios vestros.\end{Resp}
\begin{Resp}\Rsign Sequebátur turba mulíerum, quæ plangébant et lamentabántur eum.\end{Resp}
\switchcolumn
\EN
\begin{Resp}\Rsign Jesus bearing His Cross, went forth. \Star{} And there followed Him a company of women, which bewailed and lamented Him.\end{Resp}
\begin{Resp}\Vsign Daughters of Jerusalem, weep for yourselves and for your children.\end{Resp}
\begin{Resp}\Rsign There followed Him a company of women, which bewailed and lamented Him.\end{Resp}
\switchcolumn*

\LA
\LessonHead{Lectio V}
\switchcolumn
\EN
\LessonHead{Lesson V}
\switchcolumn*

\LA
\noindent Tuam ergo pertransívit ánimam vis dolóris, ut plusquam Mártyrem non immérito prædicémus, in qua nimírum corpóreæ sensum passiónis excésserit compassiónis afféctus. An non tibi plusquam gládius fuit sermo ille, revéra pertránsiens ánimam, et pertíngens usque ad divisiónem ánimæ et spíritus: Múlier, ecce fílius tuus? O commutatiónem! Joánnes tibi pro Jesu tráditur, servus pro Dómino, discípulus pro Magístro, fílius Zebedǽi pro Fílio Dei, homo purus pro Deo vero. Quómodo non tuam affectuosíssimam ánimam pertransíret hæc audítio, quando et nostra, licet sáxea, licet férrea péctora, sola recordátio scindit?\par
\switchcolumn
\EN
\noindent The sword of sorrow pierced through thy soul, so that we may truly call thee more than martyr, in whom the love, that made thee suffer along with thy Son, wrung thy heart more bitterly than any pang of bodily pain could do. Did not that word of His indeed pierce through thy soul, sharper than any two-edged sword, even to the dividing asunder of soul and spirit, (Heb. iv. 12,) Woman, behold thy son! (John xix. 26.) O what a change to thee! Thou art given John for Jesus, the servant for his Lord, the disciple for his master, the son of Zebedee for the Son of God, a mere man for Very God. O how keenly must the hearing of those words have pierced through thy most loving soul, when even our hearts, stony, iron, as they are, are wrung at the memory thereof only!\par
\switchcolumn*

\LA
\begin{Resp}\Rsign Postquam venérunt in locum qui dícitur Calváriæ, ibi crucifixérunt eum: \Star{} Stabat autem juxta crucem Jesu Mater ejus.\end{Resp}
\begin{Resp}\Vsign Tunc beátam illíus ánimam dolóris gládius pertransívit.\end{Resp}
\begin{Resp}\Rsign Stabat autem juxta crucem Jesu Mater ejus.\end{Resp}
\switchcolumn
\EN
\begin{Resp}\Rsign And when they were come to the place which is called Calvary, there they crucified Him. \Star{} Now there stood by the Cross of Jesus His Mother.\end{Resp}
\begin{Resp}\Vsign Then was it that a sword of sorrow pierced through her blessed soul.\end{Resp}
\begin{Resp}\Rsign There stood by the Cross of Jesus His Mother.\end{Resp}
\switchcolumn*

\LA
\LessonHead{Lectio VI}
\switchcolumn
\EN
\LessonHead{Lesson VI}
\switchcolumn*

\LA
\noindent Non mirémini, fratres, quod María Martyr in ánima fuísse dicátur. Mirétur qui non memínerit se audivísse Paulum inter máxima géntium crímina memorántem, quod sine affectióne fuíssent. Longe id fuit a Maríæ viscéribus, longe sit a sérvulis ejus. Sed forte quis dicat: Numquid non eum præscíerat moritúrum? Et indubitánter. Numquid non sperábat contínuo resurrectúrum? Et fidéliter. Super hæc dóluit crucifíxum? Et veheménter. Alióquin quisnam tu, frater, aut unde tibi hæc sapiéntia, ut miréris plus Maríæ Fílium patiéntem? Ille étiam mori córpore pótuit; ista cómmori corde non pótuit? Fecit illud cáritas, qua majórem nemo hábuit; fecit et hoc cáritas, cui post illam símilis áltera non fuit.\par
\switchcolumn
\EN
\noindent Marvel not, my brethren, that Mary should be called a Martyr in spirit. He indeed may marvel who remembereth not what Paul saith, naming the greater sins of the Gentiles, that they were without natural affection, (Rom. i. 31.) Far other were the bowels of Mary, and far other may those of her servants be! But some man perchance will say Did she not know that He was to die? Yea, without doubt, she knew it. Did she not hope that He was soon to rise again? Yea, she most faithfully hoped it. And did she still mourn because He was crucified? Yea, bitterly. But who art thou, my brother, or whence hast thou such wisdom, to marvel less that the Son of Mary suffered than that Mary suffered with Him? He could die in the Body, and could not she die with Him in her heart? His was the deed of that Love, greater than which hath no man (John xv. 13;) hers, of a love, like to which hath no man, save He.\par
\switchcolumn*

\LA
\begin{Resp}\Rsign Joseph ab Arimathǽa \Star{} Pétiit corpus Jesu, quod, de cruce depósitum, * Suo compléxu Mater excépit.\end{Resp}
\begin{Resp}\Vsign Dolens Sunamítis sinu et génibus suis sustínuit mórtuum fílium.\end{Resp}
\begin{Resp}\Rsign Pétiit corpus Jesu, quod, de cruce depósitum.\end{Resp}
\begin{Resp}\Vsign Glória Patri, et Fílio, \Star{} et Spirítui Sancto.\end{Resp}
\begin{Resp}\Rsign Suo compléxu Mater excépit.\end{Resp}
\switchcolumn
\EN
\begin{Resp}\Rsign Joseph of Arimathaea \Star{} And His Mother received It into her arms.\end{Resp}
\begin{Resp}\Vsign The Shunamite took her dead son, and laid him on her knees; and her soul was vexed within her.\end{Resp}
\begin{Resp}\Rsign He begged the Body of Jesus, and he took It down from the Cross.\end{Resp}
\begin{Resp}\Vsign Glory be to the Father, and to the Son, \Star{} and to the Holy Ghost.\end{Resp}
\begin{Resp}\Rsign And His Mother received It into her arms.\end{Resp}
\switchcolumn*

\LA
\LessonHead{Lectio VII}
\switchcolumn
\EN
\LessonHead{Lesson VII}
\switchcolumn*

\LA
\LessonTitle{Léctio sancti Evangélii secúndum Joánnem}
\switchcolumn
\EN
\LessonTitle{From the Holy Gospel according to John}
\switchcolumn*

\LA
\Cite{Joannes 19:25-27}
\switchcolumn
\EN
\Cite{John 19:25-27}
\switchcolumn*

\LA
\noindent In illo témpore: Stabant juxta crucem Jesu Mater ejus, et soror Matris ejus María Cléophæ, et María Magdaléne. Et réliqua.\par
\switchcolumn
\EN
\noindent At that time: There stood by the Cross of Jesus His Mother, and His Mother's sister, Mary (the wife) of Cleophas, and Mary Magdalene. And so on.\par
\switchcolumn*

\LA
\LessonTitle{Homilía sancti Ambrósii Epíscopi}
\switchcolumn
\EN
\LessonTitle{Homily by St. Ambrose, (Bishop of Milan.)}
\switchcolumn*

\LA
\Source{De Instit. Virg. cap. 7}
\switchcolumn
\EN
\Source{On Virgins, 7}
\switchcolumn*

\LA
\noindent Stabat juxta crucem Mater, et, fugiéntibus viris, stabat intrépida. Vidéte utrum pudórem mutáre potúerit Mater Jesu, quæ ánimum non mutávit. Spectábat piis óculis Fílii vúlnera, per quæ sciébat ómnibus futúram redemptiónem. Spectábat non degéneri mater spectáculo, quæ non metúerat peremptórem. Pendébat in cruce Fílius, Mater se persecutóribus offerébat.\par
\switchcolumn
\EN
\noindent There stood by the Cross His Mother. Men had forsaken Him, but she stood there fearless. Behold how the Mother of Jesus could break through her shrinking modesty, but could not belie her heart. With the eyes of a mother's love she gazed upon the Wounds of her Son, those Wounds through Which she knew that redemption for all mankind was flowing. The Mother, who feared not the executioners, was able to endure the sight of their work. Her Son was hanging upon the Cross, and she braved His tormentors.\par
\switchcolumn*

\LA
\begin{Resp}\Rsign Quis tibi sensus fuit, o Mater dolórum, \Star{} Dum Joseph síndone Fílium tuum invólvit, et pósuit eum in monuménto?\end{Resp}
\begin{Resp}\Vsign Consideráte et vidéte, si est dolor, sicut dolor meus.\end{Resp}
\begin{Resp}\Rsign Dum Joseph síndone Fílium tuum invólvit, et pósuit eum in monuménto?\end{Resp}
\switchcolumn
\EN
\begin{Resp}\Rsign O what a sickening at heart was thine, thou Mother of sorrows \Star{} When Joseph wrapped thy Son in linen and laid Him in the sepulchre?\end{Resp}
\begin{Resp}\Vsign Behold, and see if there be any sorrow like unto my sorrow.\end{Resp}
\begin{Resp}\Rsign When Joseph wrapped thy Son in linen and laid Him in the sepulchre.\end{Resp}
\switchcolumn*

\LA
\LessonHead{Lectio VIII}
\switchcolumn
\EN
\LessonHead{Lesson VIII}
\switchcolumn*

\LA
\Source{Epist. 25 ad Ecclesiam Vercellensem, prope finem}
\switchcolumn
\EN
\Source{From the same 25th Epistle to the Church of Vercelli}
\switchcolumn*

\LA
\noindent María Mater Dómini ante crucem Fílii stabat. Nullus me hoc dócuit, nisi sanctus Joánnes Evangelísta. Mundum álii concússum in passióne Dómini conscripsérunt, cælum ténebris obdúctum, refugísse solem, in paradísum latrónem, sed post piam confessiónem, recéptum. Joánnes dócuit, quod álii non docuérunt, quemádmodum in cruce pósitus Matrem appelláverit. Pluris putátur quod victor suppliciórum pietátis offícia Matri exhibébat, quam quod regnum cæléste donábat. Nam si religiósum est quod latróni donátur vénia, multo uberióris pietátis est quod a Fílio Mater tanto afféctu honorátur.\par
\switchcolumn
\EN
\noindent Mary, the Mother of the Lord, stood by the Cross of her Son. My only informant of this fact is the holy Evangelist John. Others have written that when the Lord suffered, the earth quaked, the heavens were veiled in darkness, the sun was hidden, and the thief received, after a good confession, the promise of Paradise. John hath taught us what the others have not taught us. Upon the Cross He called her Mother. It is reckoned (by John) a greater thing that in the moment of triumph over agony, He should have discharged the watchful duty of a Son to His Mother, than that He should have made gift of the kingdom of heaven. For if it be a sacred thing to have forgiven the thief, this so great kindness of the Son to the Mother is to be worshipped as the outcome of a tenderer and more touching love.\par
\switchcolumn*

\LA
\begin{Resp}\Rsign In toto corde tuo gémitus Matris tuæ ne obliviscáris, \Star{} Ut perficiátur propitiátio et benedíctio.\end{Resp}
\begin{Resp}\Vsign Ave, princeps generósa, Martyrúmque prima rosa, Virginúmque lílium.\end{Resp}
\begin{Resp}\Rsign Ut perficiátur propitiátio et benedíctio.\end{Resp}
\begin{Resp}\Vsign Glória Patri, et Fílio, \Star{} et Spirítui Sancto.\end{Resp}
\begin{Resp}\Rsign Ut perficiátur propitiátio et benedíctio.\end{Resp}
\switchcolumn
\EN
\begin{Resp}\Rsign Forget not the sorrows of thy mother with thine whole heart, \Star{} That thine offering and thy blessing may be perfected.\end{Resp}
\begin{Resp}\Vsign Hail, O maid of Royal birth, Noblest martyr-rose of earth, lily of virginity.\end{Resp}
\begin{Resp}\Rsign That thine offering and thy blessing may be perfected.\end{Resp}
\begin{Resp}\Vsign Glory be to the Father, and to the Son, \Star{} and to the Holy Ghost.\end{Resp}
\begin{Resp}\Rsign That thine offering and thy blessing may be perfected.\end{Resp}
\switchcolumn*

\end{paracol}

\par\needspace{10\baselineskip}\addvspace{1.2em}{\centering\small\textsc{Die 15 Septembris} · \textsc{15 September}\par}
\Office{S. Nicomedis Martyris}{St. Nicomedes, Martyr}{Simplex}
\addcontentsline{toc}{subsection}{Die 15 Septembris · S. Nicomedis Martyris \textit{— St. Nicomedes, Martyr}}
\begin{paracol}{2}
\LA
\LessonHead{Lectio IX (pro commemoratione)}
\switchcolumn
\EN
\LessonHead{Lesson IX (when commemorated)}
\switchcolumn*

\LA
\LessonTitle{Commemoratio: S. Nicomedis Martyris}
\switchcolumn
\EN
\LessonTitle{Commemoration of: St. Nicomedes, Martyr}
\switchcolumn*

\LA
\noindent Nicomedes presbyter, persequénte Christiános Domitiáno imperatóre, quod corpus Felículæ Vírginis, propter confessiónem christianæ fidei a Flacco cómite interfectæ, sepelísset, comprehéndi jussus est; ductúsque ad statuas deórum, cum eis sacrificáre velle (quod fácere jubebátur) constanter negavísset, proptérea quod sacrifícium uni Deo vero, qui regnat in cælis, deberétur, plumbátis cæsus, in eo martyrio ánimam Deo réddidit. Cujus corpus cum idem comes in profluentem Tiberim projíci imperásset, Justus, Nicomedis clericus diligénter conquisítum, ad muros Urbis via Nomentana honorífice in sepúlcro cóndidit.\par
\switchcolumn
\EN
\noindent This Nicomede was a Priest who was ordered to be seized during the persecution of the Christians by the Emperor Domitian, because he had buried the body of the Virgin Felicula, who had been slain by the Count Flaccus for confessing the Christian Faith. He was led to the statues of the gods, and forasmuch as he stoutly disobeyed the command to sacrifice to them, since sacrifice is due only to the one true God Who reigneth in heaven, he was flogged with scourges loaded with lead until he sealed his testimony by giving up his spirit to God. The said Count Flaccus ordered his body to be thrown into the floods of the Tiber, but Justus, clerk to Nicomede, sought diligently for it until he found it, and buried it honourably upon the road to Mentana, hard by the walls of the city.\par
\switchcolumn*

\LA
\TeDeum{Te Deum}
\switchcolumn
\EN
\TeDeum{Te Deum}
\switchcolumn*

\end{paracol}

\par\needspace{10\baselineskip}\addvspace{1.2em}{\centering\small\textsc{Die 16 Septembris} · \textsc{16 September}\par}
\Office{Ss. Cornelii Papæ et Cypriani Episcopi, Martyrum}{Sts. Cornelius, Pope and Cyprian, Bishop, Martyrs}{Semiduplex}
\addcontentsline{toc}{subsection}{Die 16 Septembris · Ss. Cornelii Papæ et Cypriani Episcopi, Martyrum \textit{— Sts. Cornelius, Pope and Cyprian, Bishop, Martyrs}}
\begin{paracol}{2}
\LA
\RefLine{Lectiones I–III: de Scriptura occurrente.}
\switchcolumn
\EN
\RefLine{Lessons I–III: from the Scripture of the day.}
\switchcolumn*

\LA
\RefLine{Lectiones VII–VIII: ex Commúni Plurimorum Martyrum Pontificum (C3).}
\switchcolumn
\EN
\RefLine{Lessons VII–VIII: from the Commune Plurimorum Martyrum Pontificum (C3).}
\switchcolumn*

\LA
\RefLine{Lectio IX: de Ss. Euphemiæ, Luciæ et Geminiani Martyrum (09-16cc).}
\switchcolumn
\EN
\RefLine{Lesson IX: of Ss. Euphemia, Lucy and Geminian, Martyrs (09-16cc).}
\switchcolumn*

\LA
\LessonHead{Lectio IV}
\switchcolumn
\EN
\LessonHead{Lesson IV}
\switchcolumn*

\LA
\noindent Cornélius Romanus, Gallo et Volusiáno imperatóribus pontificátum gerens, cum Lucina, femina sanctíssima, córpora Apostolórum Petri et Pauli e catacumbis in opportuniórem locum tránstulit; ac Pauli corpus Lucina in suo prædio via Ostiénsi, prope eum locum, ubi fuerat gládio percússus, collocávit; Cornélius Principis Apostolórum corpus non longe inde, ubi crucifixus fuerat, repósuit. Quod cum ad imperatóres delátum esset, et Pontifice auctore multos fíeri Christiános, mittitur is in exsílium ad Centumcellas; ubi eum sanctus Cyprianus, epíscopus Carthaginiénsis, per litteras est consolatus.\par
\switchcolumn
\EN
\noindent Cornelius was a Roman who held the Papacy during the reign of the Emperors Gallus and Volusian. He, and that most holy Lady Lucina, took the bodies of the Apostles Peter and Paul out of the Catacombs and put them in more convenient places. Lucina laid the body of Paul in a farm of her own upon the road to Ostia, hard by the place where he had received the sword-stroke. Cornelius placed that of the Prince of the Apostles hard by where he had been crucified. When this was told to the Emperors, and likewise that Cornelius was the means of making many Christians, he was banished to Civitavecchia, where Cyprian, the holy Bishop of Carthage, comforted him by letters.\par
\switchcolumn*

\LA
\begin{Resp}\Rsign Sancti tui, Dómine, mirábile consecúti sunt iter, serviéntes præcéptis tuis, ut inveniréntur illǽsi in aquis válidis: \Star{} Terra appáruit árida: et in Mari Rubro via sine impediménto.\end{Resp}
\begin{Resp}\Vsign Quóniam percússit petram, et fluxérunt aquæ, et torréntes inundavérunt.\end{Resp}
\begin{Resp}\Rsign Terra appáruit árida: et in Mari Rubro via sine impediménto.\end{Resp}
\switchcolumn
\EN
\begin{Resp}\Rsign thy Saints, O Lord, have passed a wonderful way, serving thy commandments, that they might be found without hurt in the midst of the mighty waters. \Star{} Dry land appeared, and, out of the Red Sea, a way without impediment.\end{Resp}
\begin{Resp}\Vsign He smote the rock, and the waters gushed out, and the streams overflowed.\end{Resp}
\begin{Resp}\Rsign Dry land appeared, and, out of the Red Sea, a way without impediment.\end{Resp}
\switchcolumn*

\LA
\LessonHead{Lectio V}
\switchcolumn
\EN
\LessonHead{Lesson V}
\switchcolumn*

\LA
\noindent Hoc autem christianæ caritátis offícium cum frequens alter alteri persólveret, deteriórem in partem id accipiéntes imperatóres, accersítum Romam Cornelium tamquam de majestáte reum plumbátis cædi, raptumque ad Martis simulacrum, ei sacrificáre jubent. Quam impietátem cum ille detestarétur, ei caput abscissum est décimo octavo Kalendas Octobris. Cujus corpus beáta Lucina, clericis adjutóribus, humávit in arenaria prædii sui prope cœmetérium Callísti. Vixit in pontificatu annos circiter duos.\par
\switchcolumn
\EN
\noindent They continued thus to write often one to the other, till the Emperors took in bad part these exchanges of Christian love, and sent for Cornelius to Rome. There they commanded him to be lashed with scourges loaded with lead as though he were a traitor, and then to be carried to offer sacrifice before the image of Mars. He firmly refused to commit this great wickedness, and was forthwith beheaded, upon the 14th day of September, (in the year of our Lord 252.) The blessed Lucina, with the help of the clergy, buried his body in the sand-pit on her own farm, near the Cemetery of Callistus. He lived as Pope about two years.\par
\switchcolumn*

\LA
\begin{Resp}\Rsign Vérbera carníficum non timuérunt Sancti Dei, moriéntes pro Christi nómine: \Star{} Ut herédes fíerent in domo Dómini.\end{Resp}
\begin{Resp}\Vsign Tradidérunt córpora sua propter Deum ad supplícia.\end{Resp}
\begin{Resp}\Rsign Ut herédes fíerent in domo Dómini.\end{Resp}
\switchcolumn
\EN
\begin{Resp}\Rsign The Saints of God shrank not from the stripes of the executioners, but died for Christ's Name's sake \Star{} That they might be made joint-heirs in the house of the Lord.\end{Resp}
\begin{Resp}\Vsign They gave their bodies for God's sake to death.\end{Resp}
\begin{Resp}\Rsign That they might be made joint-heirs in the house of the Lord.\end{Resp}
\switchcolumn*

\LA
\LessonHead{Lectio VI}
\switchcolumn
\EN
\LessonHead{Lesson VI}
\switchcolumn*

\LA
\LessonTitle{Ex libro sancti Hierónymi Presbýteri de Scriptóribus ecclesiásticis}
\switchcolumn
\EN

\switchcolumn*

\LA
\Source{Cap. 67}
\switchcolumn
\EN

\switchcolumn*

\LA
\noindent Cyprianus, Afer, primum gloriose rhetoricam docuit. Exinde, suadénte presbytero Cæcílio, a quo et cognoméntum sortitus est, christianus factus, omnem substantiam suam paupéribus erogávit. Ac post non multum témporis, eléctus in presbyterum, étiam epíscopus Carthaginiénsis constitútus est. Hujus ingenii superfluum est indicem téxere, cum sole clarióra sint ejus opera. Passus est sub Valeriáno et Gallieno princípibus, persecutióne octava, eodem die quo Romæ Cornélius, sed non eodem anno.\par
\switchcolumn
\EN
\noindent Cypriam was an African. He was first distinguished as a teacher of Rhetoric. He afterwards became a Christian at the persuasion of the Priest Caecilius, whose surname he took, and parted all his goods among the poor. It was not long before he was chosen a Priest, and then made Bishop of Carthage. It would be idle to enlarge upon his, wit, seeing that his works are as well known as the sun. He suffered under the Emperors Valerian and Gallienus, in the eighth persecution, and upon the same day, though not in the same year, that Cornelius testified at Rome.\par
\switchcolumn*

\LA
\begin{Resp}\Rsign Tamquam aurum in fornáce probávit eléctos Dóminus, et quasi holocáusti hóstiam accépit illos: et in témpore erit respéctus illórum: \Star{} Quóniam donum et pax est eléctis Dei.\end{Resp}
\begin{Resp}\Vsign Qui confídunt in illum, intélligent veritátem, et fidéles in dilectióne acquiéscent illi.\end{Resp}
\begin{Resp}\Rsign Quóniam donum et pax est eléctis Dei.\end{Resp}
\begin{Resp}\Vsign Glória Patri, et Fílio, \Star{} et Spirítui Sancto.\end{Resp}
\begin{Resp}\Rsign Quóniam donum et pax est eléctis Dei.\end{Resp}
\switchcolumn
\EN
\begin{Resp}\Rsign As gold in the furnace hath the Lord tried His chosen ones, and received them as a burnt-offering, and yet a while, and they shall be regarded; \Star{} For the grace of God, and His peace, are with His chosen.\end{Resp}
\begin{Resp}\Vsign They that put their trust in Him shall understand the truth and such as be faithful in love shall abide with Him.\end{Resp}
\begin{Resp}\Rsign For the grace of God, and His peace, are with His chosen.\end{Resp}
\begin{Resp}\Vsign Glory be to the Father, and to the Son, \Star{} and to the Holy Ghost.\end{Resp}
\begin{Resp}\Rsign For the grace of God, and His peace, are with His chosen.\end{Resp}
\switchcolumn*

\end{paracol}

\par\needspace{10\baselineskip}\addvspace{1.2em}{\centering\small\textsc{Die 16 Septembris} · \textsc{16 September}\par}
\Office{Ss. Euphemiæ, Luciæ et Geminiani Martyrum}{Ss. Euphemia, Lucy and Geminian, Martyrs}{Simplex}
\addcontentsline{toc}{subsection}{Die 16 Septembris · Ss. Euphemiæ, Luciæ et Geminiani Martyrum \textit{— Ss. Euphemia, Lucy and Geminian, Martyrs}}
\begin{paracol}{2}
\LA
\LessonHead{Lectio IX (pro commemoratione)}
\switchcolumn
\EN
\LessonHead{Lesson IX (when commemorated)}
\switchcolumn*

\LA
\LessonTitle{Commemoratio: Ss. Euphemiæ, Luciæ et Geminiani Martyrum}
\switchcolumn
\EN
\LessonTitle{Commemoration of: Ss. Euphemia, Lucy and Geminian, Martyrs}
\switchcolumn*

\LA
\noindent Euphemia, Lúcia et Geminianus, in persecutióne Diocletiáni, non eodem loco, sed eodem die, martyrio coronáti sunt. Euphemia Virgo apud Chalcédonem, Prisco proconsule, varia tormentórum genera, virgárum, equulei, rotárum, ignis constanter passa; demum béstiis objecta, una ex iis morsum sancto corpori infigente, ceteris pedes ejus lambéntibus, immaculátum spíritum Deo réddidit. Lúcia, vidua Romana, a fílio Eutropio quod Christum multos annos coluísset accusata, in ollam pice ac plumbo ferventuem dimittitur; unde incólumis evadens, cum ferro plumboque oneráta per Urbem ducerétur, Geminianum, nóbilem virum, constántia fidei et martyrii ad Christum convértit. Quem étiam cum multis aliis ad fidem perductis, varie tortum, gloriosi martyrii, abscisso cápite, socium hábuit. Quorum córpora Maxima, mulier christiana, honorifice sepelívit.\par
\switchcolumn
\EN
\noindent Euphemia, Lucy, and Geminian were all crowned with Martyrdom in the persecution under Diocletian, upon the same day, though not in the same place. Euphemia was a maiden of Chalcedon, who suffered diverse tortures under the Proconsul Priscus. She endured unflinchingly the rods, the rack, the wheels, and the fire, and in the end was thrown to wild beasts. These all licked her feet, save one, which gave her holy body such a bite, that she forthwith resigned her guileless spirit to God. Lucy was a widow at Rome, who was accused by her own son Eutropius, for that she had for many years worshipped Christ. She was put into a vessel of hot pitch and lead, but came forth unhurt. As she was being haled through the city loaded with iron and lead, the sight of her faith and unwavering testification turned to Christ the nobleman Geminian. He was one of many whom she had brought to the faith, and she had him for a comrade in her glorious martyrdom, for he was diverse ways tormented, and then beheaded. Their bodies were given honourable burial by the Christian lady Maxima.\par
\switchcolumn*

\LA
\TeDeum{Te Deum}
\switchcolumn
\EN
\TeDeum{Te Deum}
\switchcolumn*

\end{paracol}

\par\needspace{10\baselineskip}\addvspace{1.2em}{\centering\small\textsc{Die 17 Septembris} · \textsc{17 September}\par}
\Office{Impressionis Stigmatum S. Francisci}{Impression of the Stigmata of St. Francis}{Duplex}
\addcontentsline{toc}{subsection}{Die 17 Septembris · Impressionis Stigmatum S. Francisci \textit{— Impression of the Stigmata of St. Francis}}
\begin{paracol}{2}
\LA
\LessonHead{Lectio I}
\switchcolumn
\EN
\LessonHead{Lesson I}
\switchcolumn*

\LA
\LessonTitle{De Epistola beáti Pauli Apóstoli ad Gálatas}
\switchcolumn
\EN
\LessonTitle{From the letter of St. Paul to the Galatians.}
\switchcolumn*

\LA
\Cite{Gal 5:25-26; 6:1-6}
\switchcolumn
\EN
\Cite{Gal 5:25-26; 6:1-6}
\switchcolumn*

\LA
\noindent Si Spíritu vívimus, spíritu et ambulemus. Non efficiámur inánis glóriæ cúpidi, invicem provocántes, invicem invidéntes. Fratres, etsi præoccupátus fúerit homo in aliquo delícto, vos, qui spirituales estis, hujusmodi instrúite in spíritu lenitátis consíderans te ipsum, ne et tu tenteris. Alter alterius onera portate, et sic adimplébitis legem Christi. Nam, si quis existimat se aliquid esse, cum nihil sit, ipse se sedúcit. Opus autem suum probet unusquísque, et sic in semetípso tantum glóriam habébit et non in altero. Unusquisque enim onus suum portabit. Communicet autem is qui catechizátur verbo, ei qui se catechízat in ómnibus bonis.\par
\switchcolumn
\EN
\noindent If we live in the Spirit, let us also walk in the Spirit. Let us not be made desirous of vain glory, provoking one another, envying on another. Brethren, and if a man be overtaken in any fault, you, who are spiritual, instruct such a one in the spirit of meekness, considering thyself, lest thou also be tempted. Bear ye one another's burdens; and so you shall fulfil the law of Christ. For if any man think himself to be some thing, whereas he is nothing, he deceiveth himself. But let every one prove his own work, and so he shall have glory in himself only, and not in another. For every one shall bear his own burden. And let him that is instructed in the word, communicate to him that instructeth him, in all good things.\par
\switchcolumn*

\LA
\begin{Resp}\Rsign Euge, serve bone et fidélis, quia in pauca fuísti fidélis, supra multa te constítuam: \Star{} Intra in gáudium Dómini tui.\end{Resp}
\begin{Resp}\Vsign Dómine, quinque talénta tradidísti mihi, ecce ália quinque superlucrátus sum.\end{Resp}
\begin{Resp}\Rsign Intra in gáudium Dómini tui.\end{Resp}
\switchcolumn
\EN
\begin{Resp}\Rsign Well done, thou good and faithful servant, thou hast been faithful over a few things. I will make thee ruler over many things: \Star{} Enter thou into the joy of thy Lord.\end{Resp}
\begin{Resp}\Vsign Lord, thou deliveredst unto me five talents; behold, I have gained beside them five talents more.\end{Resp}
\begin{Resp}\Rsign Enter thou into the joy of thy Lord.\end{Resp}
\switchcolumn*

\LA
\LessonHead{Lectio II}
\switchcolumn
\EN
\LessonHead{Lesson II}
\switchcolumn*

\LA
\Cite{Gal 6:7-13}
\switchcolumn
\EN
\Cite{Gal 6:7-13}
\switchcolumn*

\LA
\noindent Nolíte errare: Deus non irridétur. Quæ enim seminaverit homo hæc et metet. Quóniam qui seminat in carne sua, de carne et metet corruptiónem; qui autem seminat in spiritu, de spíritu metet vitam ætérnam. Bonum autem faciéntes non deficiámus; témpore enim suo metémus non deficiéntes. Ergo, dum tempus habemus, operemur bonum ad omnes, maxime autem ad domesticos fidei. Vidéte quálibus litteris scripsi vobis mea manu. Quicúmque enim volunt placére in carne, hi cogunt vos circumcidi, tantum ut crucis Christi persecutiónem non patiántur. Neque enim qui circumcidúntur, legem custódiunt; sed volunt vos circumcidi, ut in carne vestra glorientur.\par
\switchcolumn
\EN
\noindent Be not deceived, God is not mocked. For what things a man shall sow, those also shall he reap. For he that soweth in his flesh, of the flesh also shall reap corruption. But he that soweth in the spirit, of the spirit shall reap life everlasting. And in doing good, let us not fail. For in due time we shall reap, not failing. Therefore, whilst we have time, let us work good to all men, but especially to those who are of the household of the faith. See what a letter I have written to you with my own hand. For as many as desire to please in the flesh, they constrain you to be circumcised, only that they may not suffer the persecution of the cross of Christ. For neither they themselves who are circumcised, keep the law; but they will have you to be circumcised, that they may glory in your flesh.\par
\switchcolumn*

\LA
\begin{Resp}\Rsign Justus germinábit sicut lílium: \Star{} Et florébit in ætérnum ante Dóminum.\end{Resp}
\begin{Resp}\Vsign Plantátus in domo Dómini, in átriis domus Dei nostri.\end{Resp}
\begin{Resp}\Rsign Et florébit in ætérnum ante Dóminum.\end{Resp}
\switchcolumn
\EN
\begin{Resp}\Rsign The righteous shall grow as the lily; \Star{} Yea, he shall flourish in the presence of the Lord for ever.\end{Resp}
\begin{Resp}\Vsign Those that be planted in the house of the Lord, shall flourish in the courts of the house of our God.\end{Resp}
\begin{Resp}\Rsign Yea, he shall flourish in the presence of the Lord for ever.\end{Resp}
\switchcolumn*

\LA
\LessonHead{Lectio III}
\switchcolumn
\EN
\LessonHead{Lesson III}
\switchcolumn*

\LA
\Cite{Gal 6:14-18}
\switchcolumn
\EN
\Cite{Gal 6:14-18}
\switchcolumn*

\LA
\noindent Mihi autem absit gloriari, nisi in cruce Dómini nostri Jesu Christi; per quem mihi mundus crucifixus est, et ego mundo. In Christo enim Jesu neque circumcisio aliquid valet neque præputium, sed nova creatura. Et quicúmque hanc regulam secúti fúerint, pax super illos et misercordia et super Israël Dei. De cetero nemo mihi molestus sit; ego enim Stígmata Dómini Jesu in corpore meo porto. Gratia Dómini nostri Jesu Christi cum spíritu vestro, fratres. Amen.\par
\switchcolumn
\EN
\noindent But God forbid that I should glory, save in the cross of our Lord Jesus Christ; by whom the world is crucified to me, and I to the world. For in Christ Jesus neither circumcision availeth any thing, nor uncircumcision, but a new creature. And whosoever shall follow this rule, peace on them, and mercy, and upon the Israel of God. From henceforth let no man be troublesome to me; for I bear the marks of the Lord Jesus in my body. The grace of our Lord Jesus Christ be with your spirit, brethren. Amen.\par
\switchcolumn*

\LA
\begin{Resp}\Rsign Iste cognóvit justítiam, et vidit mirabília magna, et exorávit Altíssimum: \Star{} Et invéntus est in número Sanctórum.\end{Resp}
\begin{Resp}\Vsign Iste est qui contémpsit vitam mundi, et pervénit ad cæléstia regna.\end{Resp}
\begin{Resp}\Rsign Et invéntus est in número Sanctórum.\end{Resp}
\begin{Resp}\Vsign Glória Patri, et Fílio, \Star{} et Spirítui Sancto.\end{Resp}
\begin{Resp}\Rsign Et invéntus est in número Sanctórum.\end{Resp}
\switchcolumn
\EN
\begin{Resp}\Rsign This is he which knew righteousness, and saw great wonders, and made his prayer unto the Most High; \Star{} And he is numbered among the Saints.\end{Resp}
\begin{Resp}\Vsign This is he which loved not his life in this world, and is come unto an everlasting kingdom.\end{Resp}
\begin{Resp}\Rsign And he is numbered among the Saints.\end{Resp}
\begin{Resp}\Vsign Glory be to the Father, and to the Son, \Star{} and to the Holy Ghost.\end{Resp}
\begin{Resp}\Rsign And he is numbered among the Saints.\end{Resp}
\switchcolumn*

\LA
\LessonHead{Lectio IV}
\switchcolumn
\EN
\LessonHead{Lesson IV}
\switchcolumn*

\LA
\LessonTitle{Ex Commentariis sancti Bonaventuræ Episcopi}
\switchcolumn
\EN
\LessonTitle{From the Readings upon the Life of St. Francis, composed by St. Buona- Ventura, Cardinal Bishop (of Albano.)}
\switchcolumn*

\LA
\Source{Legenda S. Francisci cap. 13}
\switchcolumn
\EN
\Source{Chapter 13}
\switchcolumn*

\LA
\noindent Fidelis revera famulus et minister Christi, Franciscus, biennio antequam spíritum redderet cælo, cum in loco excélso seorsum, qui mons Alverniæ dícitur, quadragenárium ad honórem Archangeli Michaélis jejunium inchoasset, supernæ contemplatiónis dulcédine abundantius sólito superfusus ac cæléstium desideriórum ardentiori flamma succénsus, supernárum cœpit immissiónum cumulatius dona sentire. Dum ígitur seráphicis desideriórum ardóribus sursum agerétur in Deum, et afféctus compassiva teneritúdine in eum transformarétur, cui ex caritate nimia crucifigi complácuit; quodam mane circa festum Exaltatiónis sanctæ Crucis, in látere montis orans, vidit quasi spéciem uníus Séraphim, sex alas tam fúlgidas quam ignítas habentem, de cælórum sublimitate descendere. Qui, volatu celerrimo ad áëris locum viro Dei propinquum perveniens, non solum alatus, sed et crucifixus appáruit; manus quidem et pedes habens extensos et cruci affixos, alas vero sic miro modo hinc inde dispositas, ut duas supra caput erigeret, duas ad volándum extenderet, duabus vero réliquis totum corpus circumplecténdo velaret. Hoc videns veheménter obstupuit, mixtumque dolori gáudium mens ejus incurrit, dum et in gratióso ejus aspectu, sibi tam mirabíliter quam familiáriter apparentis, excessivam quamdam concipiebat lætítiam, et dira conspecta crucis affixio ipsíus ánimam compassívi dolóris gládio pertransívit.\par
\switchcolumn
\EN
\noindent Francis being indeed a faithful servant and minister of Christ, about the space of two years before he gave back his spirit to heaven, withdrew himself into an high mountain apart, even that mountain which is called Mount Alverno, and began to fast for forty days to the honour of the Archangel Michael. To think of the things above gave him sweeter comfort than beforetime he was wont, and the hot longing for heaven was kindled in him, so that he began to feel that the gifts from above were poured forth upon him in such fulness as he had never felt before. The burning of his desire made his heart rise towards God like the heart of a seraph, and his tender answering love yearned to be changed into the likeness of Him Who hath so loved us that He was content to bear the Cross. And it was so that one morning early, about the time of the Feast of the Exaltation of the Holy Cross, he was praying upon the side of the mountain, and there appeared unto him as it had been one of the Seraphim, having six wings, glorious and fiery, flying to him from heaven. It came therefore very swiftly, and stood in the air, hard by the man of God. He beheld then the appearance thereof that it was not winged only, but crucified also. His hands and feet were stretched forth and nailed to a Cross. Twain of his wings were lifted up and joined one to the other over his head, and twain were stretched forth to fly withal, and with twain he wrapped around his body. When Francis saw it, he was sore amazed, and his soul was.filled with sorrow and gladness, for the eyes of him that appeared were full of strange love and tenderness, so that he conceived great rejoicing thereat, but the nailing to the Cross was so exceedingly dreadful, that as he saw it, a sword of sorrow pierced his soul.\par
\switchcolumn*

\LA
\begin{Resp}\Rsign Honéstum fecit illum Dóminus, et custodívit eum ab inimícis, et a seductóribus tutávit illum: \Star{} Et dedit illi claritátem ætérnam.\end{Resp}
\begin{Resp}\Vsign Justum dedúxit Dóminus per vias rectas, et osténdit illi regnum Dei.\end{Resp}
\begin{Resp}\Rsign Et dedit illi claritátem ætérnam.\end{Resp}
\switchcolumn
\EN
\begin{Resp}\Rsign The Lord made him honourable, and defended him from his enemies, and kept him safe from those that lay in wait for him; \Star{} And gave him perpetual glory.\end{Resp}
\begin{Resp}\Vsign He went down with him into the pit, and left him not in bonds.\end{Resp}
\begin{Resp}\Rsign And gave him perpetual glory.\end{Resp}
\switchcolumn*

\LA
\LessonHead{Lectio V}
\switchcolumn
\EN
\LessonHead{Lesson V}
\switchcolumn*

\LA
\noindent Intelléxit quidem, illo docénte interius qui et apparebat exterius, quod, licet passiónis infirmitas cum immortalitate spíritus seráphici nullátenus conveníret, ideo tamen hujusmodi visio suis fuerat præsentáta conspéctibus, ut amícus ipse Christi prænosceret, se, non per martyrium carnis sed per incendium mentis, totum in Christi Jesu crucifixi expressam similitúdinem transformándum. Dispárens ítaque visio, post arcanum ac familiare collóquium, mentem ipsíus seráphico interius inflammávit ardore; carnem vero Crucifixo conformi exterius insignívit effigie, tamquam si ad ignis liquefactivam virtútem præámbulam sigillativa quædam esset impréssio subsecuta. Statim namque in mánibus et pédibus ejus apparere cœpérunt signa clavórum, ipsórum capítibus in inferióri parte mánuum et superiori pedum apparéntibus, et eórum acumínibus exsisténtibus ex adverso. Dextrum quoque latus, quasi lancea transfixum, rubra cicatrice obductum erat, quod sæpe, sánguinem sacrum effúndens, túnicam et femoralia respergebat.\par
\switchcolumn
\EN
\noindent Then He Whom he beheld with his bodily eyes, began to speak silently unto him in his heart, and he understood that albeit the deathless Seraphim cannot suffer or faint, this vision was nevertheless therefore set before him, that he might know that as a friend of Christ he was to be all changed into the likeness of Christ Jesus crucified, not by the martyrdom of the body, but by the fervour of the soul. Then they held together some sweet converse, as of a man with his friend, and the vision passed from him, but his heart was kindled inwardly with the fire of the Seraphim, and his body was outwardly changed into the likeness of Him Who was crucified, even as wax is softened by the fire and taketh the impression of the seal. From thenceforth there were in his hands and feet the marks of the nails. The heads of the nails were seen in the palms of his hands and on the insteps of his feet, and the points came out on the backs of his hands and the soles of his feet. In his right side also was a long raw wound, as though he had been pierced with a spear, from which wound his holy blood oftentimes ran and stained his shirt and breeches.\par
\switchcolumn*

\LA
\begin{Resp}\Rsign Amávit eum Dóminus, et ornávit eum: stolam glóriæ índuit eum, \Star{} Et ad portas paradísi coronávit eum.\end{Resp}
\begin{Resp}\Vsign Índuit eum Dóminus lorícam fídei, et ornávit eum.\end{Resp}
\begin{Resp}\Rsign Et ad portas paradísi coronávit eum.\end{Resp}
\switchcolumn
\EN
\begin{Resp}\Rsign The Lord loved him and beautified him; He clothed him with a robe of glory, \Star{} And crowned him at the gates of Paradise.\end{Resp}
\begin{Resp}\Vsign The Lord hath put on him the breast-plate of faith, and hath adorned him.\end{Resp}
\begin{Resp}\Rsign And crowned him at the gates of Paradise.\end{Resp}
\switchcolumn*

\LA
\LessonHead{Lectio VI}
\switchcolumn
\EN
\LessonHead{Lesson VI}
\switchcolumn*

\LA
\noindent Postquam ígitur novus homo Franciscus novo et stupéndo miraculo cláruit, cum singulari privilegio retroactis sæculis non concesso insignítus appáruit, sacris vidélicet Stigmátibus decoratus, descéndit de monte secum ferens Crucifixi effigiem, non in tábulis lapídeis vel ligneis manu figuratam artíficis, sed in carneis membris descriptam digito Dei vivi. Quóniam sacraméntum regis seraphicus vir abscóndere bonum esse optime norat, secreti regalis cónscius, signácula illa sacra pro viribus occultábat. Verum, quia Dei est ad glóriam suam magna revelare quæ facit, Dóminus ipse, qui signácula illa secrete impresserat, miracula quædam aperte per ipsa monstrávit; ut illórum occúlta et mira vis Stígmatum manifesta patéret claritate signórum. - Porro rem admirábilem ac tantopere testatam atque in pontificiis diplomátibus præcipuis laudibus et favóribus exáltatam, Benedíctus Papa undecimus anniversaria solemnitate celebrari vóluit; quam póstea Paulus quintus Pontifex maximus, ut corda fidelium in Christi crucifixi accenderéntur amórem, ad universam Ecclésiam propagávit.\par
\switchcolumn
\EN
\noindent Thereafter Francis was a new creature, famous for a new and awful sign. The holy marks of the Lord Jesus, whereon living man for twelve centuries had not been allowed to look, were his adornment. He came down from the mount bearing in himself the form of Jesus Crucified, not portrayed upon tables of stone or wood by the hand of any earthly craftsman, but drawn upon his flesh by the finger of the living God. The dying Seraph knew well that it is good to keep close the secret of a king, (Tobit xii. 7,) and knowing the secret of his King, he strove as far as in him lay to keep the sacred marks hidden from men. Nevertheless, forasmuch as it is the will of the Lord God for His Own glory to make manifest the greatness of His Own works, He openly showed forth diverse wonders through these wounds which He had Himself made in secret, so that the hidden and wondrous power of the marks might become known by the fame of the miracles. The foregoing marvellous but thoroughly witnessed facts, which were already spoken of in Papal documents with especial praise and joy, were made, by the pleasure of Pope Benedict XI, the subject of a yearly memorial, which was afterwards extended by Paul V. to the whole Church, in the hope of fanning in the hearts of the faithful the love of Christ Crucified.\par
\switchcolumn*

\LA
\begin{Resp}\Rsign Iste homo perfécit ómnia quæ locútus est ei Deus, et dixit ad eum: Ingrédere in réquiem meam: \Star{} Quia te vidi justum coram me ex ómnibus géntibus.\end{Resp}
\begin{Resp}\Vsign Iste est, qui contémpsit vitam mundi, et pervénit ad cæléstia regna.\end{Resp}
\begin{Resp}\Rsign Quia te vidi justum coram me ex ómnibus géntibus.\end{Resp}
\begin{Resp}\Vsign Glória Patri, et Fílio, \Star{} et Spirítui Sancto.\end{Resp}
\begin{Resp}\Rsign Quia te vidi justum coram me ex ómnibus géntibus.\end{Resp}
\switchcolumn
\EN
\begin{Resp}\Rsign This is he which did according unto all that God commanded him; and God said unto him: Enter thou into My rest, \Star{} For thee have I seen righteous before Me among all people.\end{Resp}
\begin{Resp}\Vsign This is he which loved not his life in this world, and is come unto an everlasting kingdom.\end{Resp}
\begin{Resp}\Rsign For thee have I seen righteous before Me among all people.\end{Resp}
\begin{Resp}\Vsign Glory be to the Father, and to the Son, \Star{} and to the Holy Ghost.\end{Resp}
\begin{Resp}\Rsign For thee have I seen righteous before Me among all people.\end{Resp}
\switchcolumn*

\LA
\LessonHead{Lectio VII}
\switchcolumn
\EN
\LessonHead{Lesson VII}
\switchcolumn*

\LA
\LessonTitle{Léctio sancti Evangélii secúndum Matthǽum}
\switchcolumn
\EN
\LessonTitle{From the Holy Gospel according to Matthew}
\switchcolumn*

\LA
\Cite{Matt 16:24-27}
\switchcolumn
\EN
\Cite{Matt 16:24-27}
\switchcolumn*

\LA
\noindent In illo témpore: Dixit Jesus discípulis suis: Si quis vult post me veníre, ábneget semetípsum, et tollat crucem suam, et sequátur me. Et réliqua.\par
\switchcolumn
\EN
\noindent At that time, Jesus said unto His disciples: If any man will come after Me, let him deny himself, and take up his cross, and follow Me. And so on.\par
\switchcolumn*

\LA
\LessonTitle{Homilía sancti Gregórii Papæ}
\switchcolumn
\EN
\LessonTitle{Homily by Pope St. Gregory (the Great.)}
\switchcolumn*

\LA
\Source{Homilia 32 in Evang.}
\switchcolumn
\EN
\Source{32nd on the Gospels}
\switchcolumn*

\LA
\noindent Quia Dóminus ac Redémptor noster novus homo venit in mundum, nova præcépta dedit mundo. Vitæ étenim nostræ véteri in vítiis enutrítæ contrarietátem oppósuit novitátis suæ. Quid enim vetus, quid carnális homo nóverat, nisi sua retinére, aliéna rápere, si posset; concupíscere, si non posset? Sed cæléstis médicus síngulis quibúsque vítiis obviántia ádhibet medicaménta. Nam sicut arte medicínæ cálida frígidis, frígida cálidis curántur: ita Dóminus noster contrária oppósuit medicaménta peccátis, ut lúbricis continéntiam, tenácibus largitátem, iracúndis mansuetúdinem, elátis præcíperet humilitátem.\par
\switchcolumn
\EN
\noindent Our Lord and Redeemer came into the world a new Man, and gave the world new commandments. For against the ways of our old life, brought and bred up in sin, He set the contrast of His new life. It was the old way, according to the knowledge of the carnal man, for every man to keep his own goods, and, if he were able to do it, to take his neighbour's goods also, and, if he were not able to take them, at least to lust after them. But the Heavenly Physician hath medicines wherewith to meet all the diseases of sin. For, even, as by the art of the physician, things hot are healed by things cold, and things cold by things hot, so doth our Lord set against sin holiness, ordaining for the lecherous purity, for the miserly munificence, for the hot-tempered meekness, and for the proud lowliness.\par
\switchcolumn*

\LA
\begin{Resp}\Rsign Iste est qui ante Deum magnas virtútes operátus est, et de omni corde suo laudávit Dóminum: \Star{} Ipse intercédat pro peccátis ómnium populórum.\end{Resp}
\begin{Resp}\Vsign Ecce homo sine queréla, verus Dei cultor, ábstinens se ab omni ópere malo, et pérmanens in innocéntia sua.\end{Resp}
\begin{Resp}\Rsign Ipse intercédat pro peccátis ómnium populórum.\end{Resp}
\switchcolumn
\EN
\begin{Resp}\Rsign This is he which wrought great wonders before God, and praised the Lord with all his heart. \Star{} May he pray for all people, that their sins may be forgiven unto them.\end{Resp}
\begin{Resp}\Vsign Behold a man without blame, a worshipper of God in truth, keeping himself clean from every evil work, and abiding still in his innocency.\end{Resp}
\begin{Resp}\Rsign May he pray for all people, that their sins may be forgiven unto them.\end{Resp}
\switchcolumn*

\LA
\LessonHead{Lectio VIII}
\switchcolumn
\EN
\LessonHead{Lesson VIII}
\switchcolumn*

\LA
\noindent Certe cum se sequéntibus nova mandáta propóneret, dixit: Nisi quis renuntiáverit ómnibus quæ póssidet, non potest meus esse discípulus. Ac si apérte dicat: Qui per vitam véterem aliéna concupíscitis, per novæ conversatiónis stúdium et vestra largímini. Quid vero in hac lectióne dicat, audiámus: Qui vult post me veníre, ábneget semetípsum. Ibi dícitur, ut abnegémus nostra: hic dícitur, ut abnegémus nos. Et fortásse laboriósum non est hómini relínquere sua; sed valde laboriósum est relínquere semetípsum. Minus quippe est abnegáre quod habet; valde autem multum est abnegáre quod est.\par
\switchcolumn
\EN
\noindent So the Lord, when He would give a new commandment unto them that came to Him, said Whosoever he be of you that forsaketh not all that he hath, he cannot be My disciple, (Luke xiv. 33,) as though He had said openly: All ye that according to the old man lust after your neighbour's goods, must, according to the zeal of the new man, give away even that which is your own. But let us hear again what He saith in this place If any man will come after Me, let him deny himself. First He saith that we must deny to ourselves that which is our own, and now that we must even deny ourselves to ourselves. Perchance it is not hard for a man to give up that which is his own, but it is exceeding hard to give up himself. To deny himself his possessions is little but to deny himself himself is a denial exceeding great.\par
\switchcolumn*

\LA
\begin{Resp}\Rsign Mihi absit gloriari, nisi in cruce Dómini nostri Jesu Christi; \Star{} Per quem mihi mundus crucifíxus est, et ego mundo.\end{Resp}
\begin{Resp}\Vsign Ego enim Stígmata Dómini Jesu in corpore meo porto.\end{Resp}
\begin{Resp}\Rsign Per quem mihi mundus crucifíxus est, et ego mundo.\end{Resp}
\begin{Resp}\Vsign Glória Patri, et Fílio, \Star{} et Spirítui Sancto.\end{Resp}
\begin{Resp}\Rsign Per quem mihi mundus crucifíxus est, et ego mundo.\end{Resp}
\switchcolumn
\EN
\begin{Resp}\Rsign But God forbid that I should glory, save in the cross of our Lord Jesus Christ; \Star{} By whom the world is crucified to me, and I to the world.\end{Resp}
\begin{Resp}\Vsign For I bear the marks of the Lord Jesus in my body.\end{Resp}
\begin{Resp}\Rsign By whom the world is crucified to me, and I to the world.\end{Resp}
\begin{Resp}\Vsign Glory be to the Father, and to the Son, \Star{} and to the Holy Ghost.\end{Resp}
\begin{Resp}\Rsign By whom the world is crucified to me, and I to the world.\end{Resp}
\switchcolumn*

\LA
\LessonHead{Lectio IX}
\switchcolumn
\EN
\LessonHead{Lesson IX}
\switchcolumn*

\LA
\noindent Ad se autem nobis veniéntibus Dóminus præcépit, ut renuntiémus nostris: quia quicúmque ad fídei agónem vénimus, luctámen contra malígnos spíritus súmimus. Nihil autem malígni spíritus in hoc mundo próprium póssident: nudi ergo cum nudis luctári debémus. Nam si vestítus quisque cum nudo luctátur, cítius ad terram dejícitur, quia habet unde teneátur. Quid enim sunt terréna ómnia, nisi quædam córporis induménta? Qui ergo contra diábolum ad certámen próperat, vestiménta abjíciat, ne succúmbat.\par
\switchcolumn
\EN
\noindent Let when we come unto Him the Lord will have us deny to ourselves even ourselves, since as many of us as are entered into the battle of faith, are entered into a contention against evil spirits. But the evil spirits have nothing of their own in this world, and therefore must we wrestle with them, naked with naked. For if he that is clothed, wrestle with him that is naked, he faileth swiftly, because he hath whereon he that is naked taketh hold. And what are all things earthly but things wherewith the soul is clothed upon? whosoever therefore will wrestle with Satan, let him cast away his clothes, lest he be thereby endangered.\par
\switchcolumn*

\LA
\TeDeum{Te Deum}
\switchcolumn
\EN
\TeDeum{Te Deum}
\switchcolumn*

\end{paracol}

\par\needspace{10\baselineskip}\addvspace{1.2em}{\centering\small\textsc{Die 18 Septembris} · \textsc{18 September}\par}
\Office{S. Josephi de Cupertino Confessoris}{St. Joseph of Cupertino Confessor}{Duplex}
\addcontentsline{toc}{subsection}{Die 18 Septembris · S. Josephi de Cupertino Confessoris \textit{— St. Joseph of Cupertino Confessor}}
\begin{paracol}{2}
\LA
\LessonHead{Lectio I}
\switchcolumn
\EN
\LessonHead{Lesson I}
\switchcolumn*

\LA
\LessonTitle{De Epistola secunda beáti Pauli Apóstoli ad Corinthios}
\switchcolumn
\EN
\LessonTitle{Lesson from the second letter of St. Paul the Apostle to the Corinthians}
\switchcolumn*

\LA
\Cite{2 Cor 4:6-11}
\switchcolumn
\EN
\Cite{2 Cor 4:6-11}
\switchcolumn*

\LA
\noindent Deus, qui dixit de ténebris lucem splendéscere, ipse illuxit in córdibus nostris ad illuminatiónem sciéntiæ claritátis Dei in fácie Christi Jesu. Habémus autem thesáurum istum in vasis fictílibus, ut sublimitas sit virtútis Dei et non ex nobis. In ómnibus tribulatiónem patimur, sed non angustiámur; aporiamur, sed non destitúimur; persecutiónem patimur, sed non derelínquimur; dejícimur, sed non perímus; semper mortificatiónem Jesu in corpore nostro circumferéntes, ut et vita Jesu manisfestétur in corpóribus nostris. Semper enim nos, qui vívimus, in mortem trádimur propter Jesum, ut et vita Jesu manifestétur in carne nostra mortali.\par
\switchcolumn
\EN
\noindent God, who commanded the light to shine out of darkness, hath shined in our hearts, to give the light of the knowledge of the glory of God, in the face of Christ Jesus. But we have this treasure in earthen vessels, that the excellency may be of the power of God, and not of us. In all things we suffer tribulation, but are not distressed; we are straitened, but are not destitute; We suffer persecution, but are not forsaken; we are cast down, but we perish not: Always bearing about in our body the mortification of Jesus, that the life also of Jesus may be made manifest in our bodies. For we who live are always delivered unto death for Jesus' sake; that the life also of Jesus may be made manifest in our mortal flesh.\par
\switchcolumn*

\LA
\begin{Resp}\Rsign Euge, serve bone et fidélis, quia in pauca fuísti fidélis, supra multa te constítuam: \Star{} Intra in gáudium Dómini tui.\end{Resp}
\begin{Resp}\Vsign Dómine, quinque talénta tradidísti mihi, ecce ália quinque superlucrátus sum.\end{Resp}
\begin{Resp}\Rsign Intra in gáudium Dómini tui.\end{Resp}
\switchcolumn
\EN
\begin{Resp}\Rsign Well done, thou good and faithful servant, thou hast been faithful over a few things. I will make thee ruler over many things: \Star{} Enter thou into the joy of thy Lord.\end{Resp}
\begin{Resp}\Vsign Lord, thou deliveredst unto me five talents; behold, I have gained beside them five talents more.\end{Resp}
\begin{Resp}\Rsign Enter thou into the joy of thy Lord.\end{Resp}
\switchcolumn*

\LA
\LessonHead{Lectio II}
\switchcolumn
\EN
\LessonHead{Lesson II}
\switchcolumn*

\LA
\Cite{2 Cor 5:1-8}
\switchcolumn
\EN
\Cite{2 Cor 5:1-8}
\switchcolumn*

\LA
\noindent Scimus enim quóniam si terrestris domus nostra hujus habitatiónis dissolvátur, quod ædificatiónem ex Deo habemus, domum non manufactam, ætérnam in cælis. Nam et in hoc ingemíscimus, habitatiónem nostram, quæ de cælo est, superíndui cupiéntes: si tamen vestiti, non nudi inveniámur. Nam et qui sumus in hoc tabernáculo, ingemíscimus graváti: eo quod nólumus exspoliári, sed supervestiri, ut absorbeátur quod mortale est, a vita. Qui autem efficit nos in hoc ipsum, Deus, qui dedit nobis pignus spíritus. Audéntes ígitur semper, sciéntes quóniam dum sumus in corpore, peregrinámur a Dómino: (per fidem enim ambulámus et non per spéciem) audémus autem, et bonam voluntátem habémus magis peregrinari a corpore, et præséntes esse ad Dóminum.\par
\switchcolumn
\EN
\noindent For we know, if our earthly house of this habitation be dissolved, that we have a building of God, a house not made with hands, eternal in heaven. For in this also we groan, desiring to be clothed upon with our habitation that is from heaven. Yet so that we be found clothed, not naked. For we also, who are in this tabernacle, do groan, being burthened; because we would not be unclothed, but clothed upon, that that which is mortal may be swallowed up by life. Now he that maketh us for this very thing, is God, who hath given us the pledge of the Spirit. Therefore having always confidence, knowing that, while we are in the body, we are absent from the Lord. (For we walk by faith, and not by sight.) But we are confident, and have a good will to be absent rather from the body, and to be present with the Lord.\par
\switchcolumn*

\LA
\begin{Resp}\Rsign Justus germinábit sicut lílium: \Star{} Et florébit in ætérnum ante Dóminum.\end{Resp}
\begin{Resp}\Vsign Plantátus in domo Dómini, in átriis domus Dei nostri.\end{Resp}
\begin{Resp}\Rsign Et florébit in ætérnum ante Dóminum.\end{Resp}
\switchcolumn
\EN
\begin{Resp}\Rsign The righteous shall grow as the lily; \Star{} Yea, he shall flourish in the presence of the Lord for ever.\end{Resp}
\begin{Resp}\Vsign Those that be planted in the house of the Lord, shall flourish in the courts of the house of our God.\end{Resp}
\begin{Resp}\Rsign Yea, he shall flourish in the presence of the Lord for ever.\end{Resp}
\switchcolumn*

\LA
\LessonHead{Lectio III}
\switchcolumn
\EN
\LessonHead{Lesson III}
\switchcolumn*

\LA
\Cite{2 Cor 12:1-9}
\switchcolumn
\EN
\Cite{2 Cor 12:1-9}
\switchcolumn*

\LA
\noindent Si gloriari oportet (non éxpedit quidem) véniam autem ad visiónes, et revelatiónes Dómini. Scio hóminem in Christo ante annos quatuordecim (sive in corpore nescio, sive extra corpus nescio, Deus scit:) raptum hujusmodi usque ad tertium cælum. Et scio hujusmodi hóminem, (sive in corpore, sive extra corpus, nescio, Deus scit:) quóniam raptus est in paradisum: et audívit arcana verba, quæ non licet hómini loqui. Pro hujusmodi gloriábor: pro me autem nihil gloriábor nisi in infirmitátibus meis. Nam, et si volúero gloriari, non ero insípiens: veritátem enim dicam: parco autem, ne quis me exsistimet supra id, quod videt in me, aut aliquid audit ex me. Et ne magnitúdo revelatiónum extollat me, datus est mihi stímulus carnis meæ angelus sátanæ, qui me colaphízet. Propter quod ter Dóminum rogávi, ut discederet a me: et dixit mihi: Sufficit tibi grátia mea: nam virtus in infirmitáte perfícitur. Libenter ígitur gloriábor in infirmitátibus meis, ut inhábitet in me virtus Christi.\par
\switchcolumn
\EN
\noindent If I must glory (it is not expedient indeed), but I will come to visions and revelations of the Lord. I know a man in Christ above fourteen years ago (whether in the body, I know not, or out of the body, I know not; God knoweth), such a one caught up to the third heaven. And I know such a man (whether in the body, or out of the body, I know not: God knoweth), That he was caught up into paradise, and heard secret words, which it is not granted to man to utter. For such an one I will glory; but for myself I will glory nothing, but in my infirmities. For though I should have a mind to glory, I shall not be foolish; for I will say the truth. But I forbear, lest any man should think of me above that which he seeth in me, or any thing he heareth from me. And lest the greatness of the revelations should exalt me, there was given me a sting of my flesh, an angel of Satan, to buffet me. For which thing thrice I besought the Lord, that it might depart from me. And he said to me: My grace is sufficient for thee; for power is made perfect in infirmity. Gladly therefore will I glory in my infirmities, that the power of Christ may dwell in me.\par
\switchcolumn*

\LA
\begin{Resp}\Rsign Iste cognóvit justítiam, et vidit mirabília magna, et exorávit Altíssimum: \Star{} Et invéntus est in número Sanctórum.\end{Resp}
\begin{Resp}\Vsign Iste est qui contémpsit vitam mundi, et pervénit ad cæléstia regna.\end{Resp}
\begin{Resp}\Rsign Et invéntus est in número Sanctórum.\end{Resp}
\begin{Resp}\Vsign Glória Patri, et Fílio, \Star{} et Spirítui Sancto.\end{Resp}
\begin{Resp}\Rsign Et invéntus est in número Sanctórum.\end{Resp}
\switchcolumn
\EN
\begin{Resp}\Rsign This is he which knew righteousness, and saw great wonders, and made his prayer unto the Most High; \Star{} And he is numbered among the Saints.\end{Resp}
\begin{Resp}\Vsign This is he which loved not his life in this world, and is come unto an everlasting kingdom.\end{Resp}
\begin{Resp}\Rsign And he is numbered among the Saints.\end{Resp}
\begin{Resp}\Vsign Glory be to the Father, and to the Son, \Star{} and to the Holy Ghost.\end{Resp}
\begin{Resp}\Rsign And he is numbered among the Saints.\end{Resp}
\switchcolumn*

\LA
\LessonHead{Lectio IV}
\switchcolumn
\EN
\LessonHead{Lesson IV}
\switchcolumn*

\LA
\noindent Josephus a Cupertino, oppido in Salentinis diœcesis Neritonénsis, anno reparátæ salútis millesimo sexcentésimo tertio, piis ibidem paréntibus ortus, Deique amóre præventus, puerítiam atque adolescéntiam summa cum simplicitate morúmque innocéntia molestoque morbo, patientíssime tolerato, Deiparæ Vírginis ope liberatus, se totum pietátis opéribus ac excoléndis virtútibus dedit; utque Deo ad majora vocanti se intimius conjungeret, ordini seraphico nomen dare constítuit. Post varios eventus voti tandem compos factus, apud Minores Conventuales in cœnobio Cryptulæ, inter laicos primum ob litterárum imperítiam, deínde inter clericos divina dispositióne connumerátus est. Sacerdotio post solemnia vota initiatus, perfectius sibi vitæ institutum propósuit. Quam ob rem, mundanis quibuscúmque afféctibus terrenisque rebus pene ad vitam necessariis illico a se abdicatis, ciliciis, flagellis, catenis, omni demum asperitátum ac pœnárum genere corpus afflixit; spíritum vero sanctæ oratiónis altissimǽque contemplatiónis assiduitate dulciter enutrívit. Hinc factum est, ut caritas Dei, quæ jam erat in ejus corde a prima ætate diffúsa, miro planeque singulari modo in dies coruscáverit.\par
\switchcolumn
\EN
\noindent This Joseph was born, of godly parents, at Cupertino, a small village of the diocese of Nardo, between Brindisi and Otranto, (six miles from the coast of the Gulf of Tarento, upon the 17th day of June,) in the year of Redemption 1603. The love of God came to him early, and he passed his childhood and youth in great guilelessness and harmlessness. After recovering by the help of the Virgin Mother of God from a long and painful sickness which he bore very quietly, he gave himself altogether to godliness and self-improvement. God called him inwardly to higher things, and to give himself more utterly to His service, he determined in himself to join the “Seraphic” Order. After diverse failures and changes, he obtained his wish among the Friars of the convent of “La Grotella.” He went first as a lay-brother, on account of his ignorance of letters, but God was pleased to allow him afterwards to be taken among the choir brethren. After taking his solemn vows he was ordained Priest, and then set before him to aim at a more perfect life. To this end (as far as in him lay) he thrust from him all earthly affections and all carnal things, even to such as seem almost needful for life. He tormented his body with haircloth, scourging, spiked chains, and every kind of hardship and affliction. He fed his spirit sweetly upon the constant exercise of holy prayer, and gazing upon the highest matters. And so it came to pass that the love of God, which had been enkindled in his heart from his earliest years, burnt forth day by day more strangely and openly.\par
\switchcolumn*

\LA
\begin{Resp}\Rsign Honéstum fecit illum Dóminus, et custodívit eum ab inimícis, et a seductóribus tutávit illum: \Star{} Et dedit illi claritátem ætérnam.\end{Resp}
\begin{Resp}\Vsign Justum dedúxit Dóminus per vias rectas, et osténdit illi regnum Dei.\end{Resp}
\begin{Resp}\Rsign Et dedit illi claritátem ætérnam.\end{Resp}
\switchcolumn
\EN
\begin{Resp}\Rsign The Lord made him honourable, and defended him from his enemies, and kept him safe from those that lay in wait for him; \Star{} And gave him perpetual glory.\end{Resp}
\begin{Resp}\Vsign He went down with him into the pit, and left him not in bonds.\end{Resp}
\begin{Resp}\Rsign And gave him perpetual glory.\end{Resp}
\switchcolumn*

\LA
\LessonHead{Lectio V}
\switchcolumn
\EN
\LessonHead{Lesson V}
\switchcolumn*

\LA
\noindent Eluxit præcipue ardentíssima ejus caritas in extasibus ad Deum suavíssimis stupendisque raptibus, quibus frequenter afficiebátur. Mirum autem, quod, alienato a sensibus animo, statim ab éxtasi eum revocábat sola obediéntia. Hanc quippe virtútem exímio studio prosequebátur, dicere sólitus, se ab ea véluti cæcum circumduci, et mori potius velle, quam non obedire. Paupertátem vero seraphici Patriárchæ ita æmulátus est, ut, morti próximus, prælato suo assérere vere potúerit se nihil habere, quod more religiosórum resignaret. Itaque, mundo sibique mórtuus, vitam Jesu manifestábat in carne sua, quæ dum in aliquibus ex turpitúdine obscœnum flagitium sentiebat, prodigiósum de se efflábat odórem, indícium nitidíssimæ illíus puritátis, quam, immundo spíritu vehementíssimis tentatiónibus frustra obnubilare diu conante, servávit illæsam, tum arcta sensuum custódia, tum jugi corporis maceratióne, tum denique speciali protectióne puríssimæ Vírginis Maríæ, quam matrem suam appellare consuevit, ac véluti matrem dulcíssimam íntimo cordis afféctu venerabátur, eamque ab aliis venerari exoptabat, ut cum ejusdem patrocinio, sicut ipse ajebat, ómnia bona consequerentur.\par
\switchcolumn
\EN
\noindent The chief outcome of this love of God was the strong and marvellous trances whereinto he oftentimes fell. It was, nevertheless, strange to observe that after he had entirely lost his senses he could be called out of the trance by the mere order of his superiors. To be utterly obedient was one of his chief aims, and he was used to say that those who ruled him could lead him about like a blind man, and that it was better to die than not to obey. He so imitated the poverty of the Seraphic Patriarch, that when he was at the point of death, when the Friars use to dispose of anything they have, he was able to tell his Superior that he had absolutely nothing. Thus bearing about in his body the dying of the Lord Jesus, the life also of Jesus was made manifest in his body. When he saw that certain persons had committed a foul sin of uncleanness, there came from him a strong savour, a proof of that snowy and glorious purity which, in spite of the most hideous temptations whereby the unclean spirit wrestled long to darken it, he kept undefiled, partly by an iron bridling of his senses, partly by the stern punishments he inflicted upon his own body, and partly by the extraordinary protection of the pure Virgin Mary, whom he was used to call his own Mother, whom he honoured and worshipped as his most tender Mother in his very heart of hearts, and whom he was eager that all men should honour, because, as he said, if we have her protection, every good thing comes with it.\par
\switchcolumn*

\LA
\begin{Resp}\Rsign Amávit eum Dóminus, et ornávit eum: stolam glóriæ índuit eum, \Star{} Et ad portas paradísi coronávit eum.\end{Resp}
\begin{Resp}\Vsign Índuit eum Dóminus lorícam fídei, et ornávit eum.\end{Resp}
\begin{Resp}\Rsign Et ad portas paradísi coronávit eum.\end{Resp}
\switchcolumn
\EN
\begin{Resp}\Rsign The Lord loved him and beautified him; He clothed him with a robe of glory, \Star{} And crowned him at the gates of Paradise.\end{Resp}
\begin{Resp}\Vsign The Lord hath put on him the breast-plate of faith, and hath adorned him.\end{Resp}
\begin{Resp}\Rsign And crowned him at the gates of Paradise.\end{Resp}
\switchcolumn*

\LA
\LessonHead{Lectio VI}
\switchcolumn
\EN
\LessonHead{Lesson VI}
\switchcolumn*

\LA
\noindent Hæc beáti Josephi sollicitúdo a sua erga próximos caritate prodíbat; tanto enim animárum zelo exardebat, ut ómnium salútem modis ómnibus instantíssime procuraret. Exténdens páriter caritátem suam in próximum, sive páuperem, sive infirmum, sive quacúmque alia tribulatióne vexátum, quantum in ipso erat, illum recreábat. Nec aliéni erant ab ejus caritate, qui objurgatiónibus, probris omnisque generis injuriis ipsum appéterent; nam eadem patiéntia, mansuetúdine, vultusque hilaritate talia excipiebat, qua tot inter ac tantas vicissitúdines replenduit, dum vel moderatórum ordinis vel sacræ Inquisitiónis jussu hac illac errare versaríque coactus est. Quamquam vero pópuli non solum, sed viri príncipes exímiam ejus sanctitátem et superna charísmata admiraréntur, ea nihilóminus erat humilitate, ut, magnum se peccatórem réputans, Deum enixe deprecarétur; ut sua ab eo illustria dona removéret, hómines vero exoraret ut in eum locum mortuum ejus corpus injícerent, ubi memória sui esset prorsus oblitterata. At Deus, qui ponit húmiles in sublíme quique servum suum, dum viveret, cælésti sapiéntia, prophetía, cordium perscrutatióne, curatiónum grátia ceterisque donis cumulatíssime exornaverat, ejus quoque mortem iis, quibus ipse antea prædixerat, loco ac témpore, anno ætátis suæ sexagesimo primo, Auxími in Piceno pretiósam réddidit sepulcrúmque gloriosum. Illum denique, étiam post óbitum miraculis coruscántem, Benedíctus quartus décimus Beatórum, Clemens tertius décimus Sanctórum fastis adscripsit. Ejus autem Offícium et Missam Clemens quartus décimus, ejusdem ordinis, ad universam Ecclésiam extendit.\par
\switchcolumn
\EN
\noindent This eagerness on the part of the blessed Joseph was but one outcome from his love for his neighbours. So great was his zeal for souls, that he vehemently sought in all ways for the salvation of all. When he saw his neighbour in any trouble, whether it were poverty or sickness or any other affliction, his tenderness went out toward him, and he helped him as well as he could. They who reviled, and slandered, and insulted himself were not cut off from his love. He was used to welcome such with great long-suffering, meekness, and cheerfulness of countenance and he preserved the same constantly amid many hardships and changes when he was sent hither and thither by command of the Superiors of his Order, and of the Holy Inquisition. People and princes alike marvelled at the exceeding holiness of his life, and the spiritual gifts poured upon him from above, but he was so lowly, that he sincerely held himself to be chief among sinners, and earnestly besought God to take away from him the more showy of His gifts. Of men he entreated that after his death they would cast his body somewhere where his memory might soonest perish. But God, Who exalteth them of low degree, glorified His servant during life with the gifts of heavenly wisdom, of prophecy, of discerning the hidden thoughts of the heart, of healing, and of other spiritual gifts in marvellous abundance, gave him a precious death, and made the place of his rest glorious. He fell asleep in Jesus upon the very day and at the very place foretold by himself, that is, at Osimo, (between Ancona and Loretto, upon the 18 th day of September,) in the 61st year of his own age, \textasciitilde{}(and in that of salvation 1663.) He was famous for miracles even after his death, and Benedict XIV. enrolled his name among those of the Blessed, and Clement XIII. among those of the Saints. Clement XIV., being himself a member of the same Order, extended the use of the Office and Mass in memory of him to the whole Church.\par
\switchcolumn*

\LA
\begin{Resp}\Rsign Iste homo perfécit ómnia quæ locútus est ei Deus, et dixit ad eum: Ingrédere in réquiem meam: \Star{} Quia te vidi justum coram me ex ómnibus géntibus.\end{Resp}
\begin{Resp}\Vsign Iste est, qui contémpsit vitam mundi, et pervénit ad cæléstia regna.\end{Resp}
\begin{Resp}\Rsign Quia te vidi justum coram me ex ómnibus géntibus.\end{Resp}
\begin{Resp}\Vsign Glória Patri, et Fílio, \Star{} et Spirítui Sancto.\end{Resp}
\begin{Resp}\Rsign Quia te vidi justum coram me ex ómnibus géntibus.\end{Resp}
\switchcolumn
\EN
\begin{Resp}\Rsign This is he which did according unto all that God commanded him; and God said unto him: Enter thou into My rest, \Star{} For thee have I seen righteous before Me among all people.\end{Resp}
\begin{Resp}\Vsign This is he which loved not his life in this world, and is come unto an everlasting kingdom.\end{Resp}
\begin{Resp}\Rsign For thee have I seen righteous before Me among all people.\end{Resp}
\begin{Resp}\Vsign Glory be to the Father, and to the Son, \Star{} and to the Holy Ghost.\end{Resp}
\begin{Resp}\Rsign For thee have I seen righteous before Me among all people.\end{Resp}
\switchcolumn*

\LA
\LessonHead{Lectio VII}
\switchcolumn
\EN
\LessonHead{Lesson VII}
\switchcolumn*

\LA
\LessonTitle{Léctio sancti Evangélii secúndum Matthǽum}
\switchcolumn
\EN
\LessonTitle{From the Holy Gospel according to Matthew}
\switchcolumn*

\LA
\Cite{Matt 22:1-14}
\switchcolumn
\EN
\Cite{Matt 22:1-14}
\switchcolumn*

\LA
\noindent In illo témpore: Loquebátur Jesus princípibus sacerdotum et pharisæis in parábolis dicens: Simile factum est regnum cælórum hómini regi, qui fecit nuptias fílio suo. Et réliqua.\par
\switchcolumn
\EN
\noindent At that time Jesus spake unto the chief priests and Pharisees by parables, and said The kingdom of heaven is like unto a certain king, which made a marriage for his son. And so on.\par
\switchcolumn*

\LA
\LessonTitle{Homilía sancti Gregórii Papæ}
\switchcolumn
\EN
\LessonTitle{Homily by Pope St. Gregory (the Great.)}
\switchcolumn*

\LA
\Source{Liber 2 Homiliar. Hom. 38, circa medium}
\switchcolumn
\EN
\Source{Bk. ii. Horn. 38, 9}
\switchcolumn*

\LA
\noindent Quia jam, largiénte Dómino, nuptiárum domum, id est, sanctam Ecclésiam intrastis, solerter, fratres, aspicite, ne aliquid de mentis vestræ habitu rex ingrédiens reprehéndat. Cum magno enim cordis timóre pensándum est quod prótinus subditur: Intrávit autem rex, ut vidéret discumbéntes, et vidit ibi hóminem non vestítum veste nuptiali. Quid, fratres caríssimi, éxprimi per nuptialem vestem putamus? Si enim vestem nuptialem baptisma vel fidem dicimus, quis sine baptismate et fide has nuptias intrávit? Eo enim ipso foris est, qui necdum credidit. Quid ergo debémus intellígere per nuptialem vestem, nisi caritátem? Intrat enim ad nuptias, sed cum nuptiali veste non intrat, qui, in sancta Ecclésia assistens, fidem habet, sed caritátem non habet. Recte enim caritas, nuptialis vestis vocátur, quia hanc in se conditor noster hábuit, dum ad sociandæ sibi Ecclésiæ nuptias venit.\par
\switchcolumn
\EN
\noindent Dearly beloved brethren, ye have already entered, at the Lord's bidding, into the house where the marriagefeast is being held, that is to say, into the Holy Church, and look ye well to it, that when the King cometh in to see the guests, he see nothing amiss in your soul's wedding-garment. For indeed it is with great searchings of heart that we are behoven to consider that which so soon cometh. “ And when the King came in to see the guests, he saw there a man which had not on a wedding-garment.” Dearly beloved brethren, what are we to think is signified by this wedding garment Is it baptism or is it faith? But without baptism, or without faith, who could be seated at the marriage-feast? He that believeth not would still be without the house. What then, except love, must we understand by the wedding-garment He who hath faith and is in the Holy Church, but hath not charity, cometh in unto the wedding indeed, but hath not a wedding-garment. And charity is well called the wedding -garment, for it is the garment wherewithal our Maker decked Himself when He came to wed the Church unto Himself.\par
\switchcolumn*

\LA
\begin{Resp}\Rsign Iste est qui ante Deum magnas virtútes operátus est, et de omni corde suo laudávit Dóminum: \Star{} Ipse intercédat pro peccátis ómnium populórum.\end{Resp}
\begin{Resp}\Vsign Ecce homo sine queréla, verus Dei cultor, ábstinens se ab omni ópere malo, et pérmanens in innocéntia sua.\end{Resp}
\begin{Resp}\Rsign Ipse intercédat pro peccátis ómnium populórum.\end{Resp}
\switchcolumn
\EN
\begin{Resp}\Rsign This is he which wrought great wonders before God, and praised the Lord with all his heart. \Star{} May he pray for all people, that their sins may be forgiven unto them.\end{Resp}
\begin{Resp}\Vsign Behold a man without blame, a worshipper of God in truth, keeping himself clean from every evil work, and abiding still in his innocency.\end{Resp}
\begin{Resp}\Rsign May he pray for all people, that their sins may be forgiven unto them.\end{Resp}
\switchcolumn*

\LA
\LessonHead{Lectio VIII}
\switchcolumn
\EN
\LessonHead{Lesson VIII}
\switchcolumn*

\LA
\noindent Sola quippe dilectióne Dei actum est, ut ejus Unigénitus mentes sibi electórum hóminum uníret. Unde et Joánnes dicit: Sic enim diléxit Deus mundum, ut Fílium suum unigénitum daret pro nobis. Qui ergo per caritátem venit ad hómines, eamdem caritátem innotuit vestem esse nuptialem. Omnis ergo vestrum, qui in Ecclésia positus Deo credidit, jam ad nuptias intrávit; sed cum nuptiali veste non venit, si caritátis grátiam non custódit. Et certe, fratres, si quis ad carnales nuptias esset invitatus, vestem mutaret, congaudere se sponso et sponsæ ex ipso sui habitus decore osténderet, inter gaudéntes et festa celebrántes despectis vestibus apparére erubésceret. Nos ad Dei nuptias venímus, et cordis vestem mutare dissimulamus. Congaudent Angeli, cum ad cælum assumúntur electi. Qua ergo mente hæc spiritualia festa conspícimus, qui nuptialem vestem, id est, caritátem, quæ sola nos speciósos éxhibet, non habemus?\par
\switchcolumn
\EN
\noindent It was the work of God's love alone that His Only - begotten Son should wed Himself unto the souls of the elect. Whence indeed John saith “God so loved the world, that He gave His Only-begotten Son, that whosoever believeth in Him should not perish, but have everlasting life.” (iii. 16.) He therefore Whom love brought among men, showeth that the same love is His wedding-garment. Each one therefore of you who is in the Church and believeth in God, hath already come in unto the marriage-feast, but if he keep not the grace of charity, he is come in thither not having a weddinggarment. In sooth, my brethren, if one be asked to an earthly marriage, he changeth his attire, to show even by his garments that he rejoiceth in the joy of the Bride and Bridegroom, and he would be ashamed to appear in unseemly raiment among the guests that are feasting and making merry. We are come unto God's marriagefeast, and we make pretence to change the vesture of our hearts. There is joy among the angels when the elect are taken to heaven. With what face shall we look upon this spiritual feast if we come in thither not having charity, the only wedding-garment wherein we can appear comely\par
\switchcolumn*

\LA
\begin{Resp}\Rsign Sint lumbi vestri præcíncti, et lucérnæ ardéntes in mánibus vestris: \Star{} Et vos símiles homínibus exspectántibus dóminum suum, quando revertátur a núptiis.\end{Resp}
\begin{Resp}\Vsign Vigiláte ergo, quia nescítis qua hora Dóminus vester ventúrus sit.\end{Resp}
\begin{Resp}\Rsign Et vos símiles homínibus exspectántibus dóminum suum, quando revertátur a núptiis.\end{Resp}
\begin{Resp}\Vsign Glória Patri, et Fílio, \Star{} et Spirítui Sancto.\end{Resp}
\begin{Resp}\Rsign Et vos símiles homínibus exspectántibus dóminum suum, quando revertátur a núptiis.\end{Resp}
\switchcolumn
\EN
\begin{Resp}\Rsign Let your loins be girded about, and your lights burning; \Star{} And ye yourselves like unto men that wait for their lord, when he will return from the wedding.\end{Resp}
\begin{Resp}\Vsign Watch therefore, for ye know not what hour your Lord doth come.\end{Resp}
\begin{Resp}\Rsign And ye yourselves like unto men that wait for their lord, when he will return from the wedding.\end{Resp}
\begin{Resp}\Vsign Glory be to the Father, and to the Son, \Star{} and to the Holy Ghost.\end{Resp}
\begin{Resp}\Rsign And ye yourselves like unto men that wait for their lord, when he will return from the wedding.\end{Resp}
\switchcolumn*

\LA
\LessonHead{Lectio IX}
\switchcolumn
\EN
\LessonHead{Lesson IX}
\switchcolumn*

\LA
\noindent Sciéndum vero est quia, sicut in duobus lignis, superiore vidélicet et inferióre, vestis téxitur: ita in duobus præceptis caritas habétur, in dilectióne scilicet Dei et próximi. Scriptum quippe est: Diliges Dóminum Deum tuum ex toto corde tuo et ex tota ánima tua, et próximum tuum sicut teípsum. Qua in re notándum est, quia in dilectióne próximi mensúra amoris pónitur, cum dícitur: Diliges próximum tuum sicut teípsum. Dei autem diléctio nulla mensúra constringitur, cum dícitur: Diliges Dóminum Deum tuum ex toto corde tuo, ex tota ánima tua, ex tota virtúte tua. Non enim jubétur quisque quantum diligat, sed ex quanto, cum dícitur, Ex toto; quia ille veraciter Deum díligit, qui sibi de se nihil relinquit. Duo ergo necesse est ut caritátis præcepta custódiat, quisquis habere in nuptiis vestem nuptialem curat.\par
\switchcolumn
\EN
\noindent We must know that as every garment is woven upon two beams, an upper and a lower, so love is bound unto two commandments, the one bidding us to love God, and the other to love our neighbour. For thus is it written “Thou shalt love the Lord thy God with all thy heart, and with all thy soul, and with all thy mind. This is the first and great commandment. And the second is like unto it Thou shalt love thy neighbour as thyself. On these two commandments hang all the law and the Prophets.” (Matth. xxii. 37-39-) In the which we are to see that bounds are set to that love wherewith we are to love our neighbour, for it is said “Thou shalt love thy neighbour as thyself.” But to the love wherewith we are to love God are set no bounds, for it is said “Thou shalt love the Lord thy God with all thy heart, and with all thy soul, and with all thy mind.” A man is not commanded to what point he is to love God, but from what point, even as it is said, “with all” for he only truly loveth God, who leaveth nothing for himself. We are behoven therefore to keep two commandments touching love, if we would be seen at the marriage with a wedding-garment.\par
\switchcolumn*

\LA
\TeDeum{Te Deum}
\switchcolumn
\EN
\TeDeum{Te Deum}
\switchcolumn*

\end{paracol}

\par\needspace{10\baselineskip}\addvspace{1.2em}{\centering\small\textsc{Die 19 Septembris} · \textsc{19 September}\par}
\Office{S. Januarii Episcopi et Sociorum Martyrum}{St. Januarius, Bishop and Companions, Martyrs}{Duplex}
\addcontentsline{toc}{subsection}{Die 19 Septembris · S. Januarii Episcopi et Sociorum Martyrum \textit{— St. Januarius, Bishop and Companions, Martyrs}}
\begin{paracol}{2}
\LA
\RefLine{Lectiones I–III: de Scriptura occurrente.}
\switchcolumn
\EN
\RefLine{Lessons I–III: from the Scripture of the day.}
\switchcolumn*

\LA
\RefLine{Lectio IX: de Scriptura occurrente.}
\switchcolumn
\EN
\RefLine{Lesson IX: from the Scripture of the day.}
\switchcolumn*

\LA
\LessonHead{Lectio IV}
\switchcolumn
\EN
\LessonHead{Lesson IV}
\switchcolumn*

\LA
\noindent Januarius, Benevénti epíscopus, Diocletiáno et Maximiáno in Christiános sæviéntibus, ad Timotheum Campaniæ præsidem ob christianæ fidei professiónem Nolam perducitur. Ibi, ejus constántia varie tentata, in ardentem fornacem conjéctus, ita illæsus evasit, ut ne vestiméntum aut capillum quidem flamma violáverit. Hinc præses, accénsus iracúndia, Mártyris corpus imperat usque eo dístrahi, quoad nervórum compages artuumque solvántur. Festus interea ejus diaconus et Desiderius lector, comprehénsi vinctique, una cum episcopo ante rhedam præsidis Puteleos pertrahúntur, et in eumdem carcerem, in quo Sosius Misenas et Próculus Puteolánus diáconus, Eutyches et Acutius láici, ad bestias damnati, detinebántur, simul conjiciúntur.\par
\switchcolumn
\EN
\noindent At that time time the Emperors Diocletian and Maximian were furiously raging against Christians, Januarius, Bishop of Benevento, was taken to Nola, to Timothy, President of Campania, on the charge of professing the Christian faith. There his firmness was tried diverse ways, and he was cast into a burning fiery furnace, but came forth thence unhurt, for neither upon his raiment nor upon the hairs of his head did the flame take any hold. Thereupon the wrath of the President was enkindled, and he commanded the martyr to be torn limb from limb. But in the meanwhile Januarius' Deacon Festus and his Reader Desiderius were taken, and the whole three were led in bonds to Puzzuoli in front of the President's carriage, and there thrown into the same prison wherein were already held four other Christians condemned to be devoured by wild beasts, that is to say, Sosius, a Deacon of Miseno Proculus, a Deacon of Puzzuoli and two laymen, named respectively Eutyches and Acutius.\par
\switchcolumn*

\LA
\begin{Resp}\Rsign Sancti tui, Dómine, mirábile consecúti sunt iter, serviéntes præcéptis tuis, ut inveniréntur illǽsi in aquis válidis: \Star{} Terra appáruit árida: et in Mari Rubro via sine impediménto.\end{Resp}
\begin{Resp}\Vsign Quóniam percússit petram, et fluxérunt aquæ, et torréntes inundavérunt.\end{Resp}
\begin{Resp}\Rsign Terra appáruit árida: et in Mari Rubro via sine impediménto.\end{Resp}
\switchcolumn
\EN
\begin{Resp}\Rsign thy Saints, O Lord, have passed a wonderful way, serving thy commandments, that they might be found without hurt in the midst of the mighty waters. \Star{} Dry land appeared, and, out of the Red Sea, a way without impediment.\end{Resp}
\begin{Resp}\Vsign He smote the rock, and the waters gushed out, and the streams overflowed.\end{Resp}
\begin{Resp}\Rsign Dry land appeared, and, out of the Red Sea, a way without impediment.\end{Resp}
\switchcolumn*

\LA
\LessonHead{Lectio V}
\switchcolumn
\EN
\LessonHead{Lesson V}
\switchcolumn*

\LA
\noindent Postero die omnes in amphitheatro feris objecti sunt; quæ, naturalis oblitæ feritátis, ad Januarii pedes se prostravére. Id Timótheus mágicis cantiónibus tribuens, cum senténtiam cápitis in Christi Mártyres pronuntiásset, óculis repénte captus, orante mox beato Januario, lumen recepit; quo miraculo hóminum míllia fere quinque Christi fidem suscepérunt. Verum ingrátus judex, níhilo placatior factus beneficio, sed conversióne tantæ multitúdinis actus in rábiem; véritus maxime príncipum decreta, sanctum Episcopum cum sociis gládio pércuti jussit.\par
\switchcolumn
\EN
\noindent The next day all seven were exposed to the wild beasts in the amphitheatre, but these creatures forgot their natural fierceness, and lay down at the feet of Januarius. Timothy would have it that this came from charms, and commanded the witnesses of Christ to be beheaded. Thereupon he became of a sudden blind, until Januarius had prayed for him by the which miracle nearly five thousand persons were turned to Christ. But this good turn roused up no gratitude in the President, yea, rather, the conversion of so many drave him wild, and in his hot fear to obey the decrees of the Emperors he commanded that the holy Bishop and his companions should be smitten with the sword.\par
\switchcolumn*

\LA
\begin{Resp}\Rsign Vérbera carníficum non timuérunt Sancti Dei, moriéntes pro Christi nómine: \Star{} Ut herédes fíerent in domo Dómini.\end{Resp}
\begin{Resp}\Vsign Tradidérunt córpora sua propter Deum ad supplícia.\end{Resp}
\begin{Resp}\Rsign Ut herédes fíerent in domo Dómini.\end{Resp}
\switchcolumn
\EN
\begin{Resp}\Rsign The Saints of God shrank not from the stripes of the executioners, but died for Christ's Name's sake \Star{} That they might be made joint-heirs in the house of the Lord.\end{Resp}
\begin{Resp}\Vsign They gave their bodies for God's sake to death.\end{Resp}
\begin{Resp}\Rsign That they might be made joint-heirs in the house of the Lord.\end{Resp}
\switchcolumn*

\LA
\LessonHead{Lectio VI}
\switchcolumn
\EN
\LessonHead{Lesson VI}
\switchcolumn*

\LA
\noindent Horum córpora finítimæ urbes, pro suo quæque studio certum sibi patronum ex iis apud Deum adoptándi, sepelienda curarunt. Januarii corpus Neapolitani divino admónitu extulére; quod, primo Benevéntum, inde ad monastérium montis Vírginis, postremo Neapolim translátum et in majori ecclésia cónditum, multis miraculis cláruit. Sed illud in primis memorándum, quod erumpéntes olim e monte Vesuvio flammárum globos, nec vicínis modo sed longinquis étiam regiónibus vastitátis metum afferéntes, exstinxit. Præclárum illud quoque, quod ejus sanguis, qui in ampulla vitrea concretus asservátur, cum in conspéctu cápitis ejusdem Mártyris pónitur, admirándum in modum colliquefíeri et ebullire, perinde atque recens effúsus, ad hæc usque témpora cernitur.\par
\switchcolumn
\EN
\noindent The cities of those coasts strove to obtain their bodies for honourable burial, so as to make sure of having in them advocates with God. By God's will the relics of Januarius were taken to Naples at last, after having been carried from Puzzuoli to Benevento, and from Benevento to Monte Vergine; when they were brought thence to Naples, they were laid in the chief Church there, and there have been famous on account of many miracles. Among these is remarkable the stopping of eruptions of Mount Vesuvius, whereby both that neighbourhood and also places afar off have been like to have been brought to desolation. It is also well known, and is the plain fact, seen even unto this day, that when the blood of Januarius, kept dried up in a small glass phial, is put in sight of the head of the same martyr, it is used to melt and bubble in a very strange way, as though it had but freshly been shed.\par
\switchcolumn*

\LA
\begin{Resp}\Rsign Tamquam aurum in fornáce probávit eléctos Dóminus, et quasi holocáusti hóstiam accépit illos: et in témpore erit respéctus illórum: \Star{} Quóniam donum et pax est eléctis Dei.\end{Resp}
\begin{Resp}\Vsign Qui confídunt in illum, intélligent veritátem, et fidéles in dilectióne acquiéscent illi.\end{Resp}
\begin{Resp}\Rsign Quóniam donum et pax est eléctis Dei.\end{Resp}
\begin{Resp}\Vsign Glória Patri, et Fílio, \Star{} et Spirítui Sancto.\end{Resp}
\begin{Resp}\Rsign Quóniam donum et pax est eléctis Dei.\end{Resp}
\switchcolumn
\EN
\begin{Resp}\Rsign As gold in the furnace hath the Lord tried His chosen ones, and received them as a burnt-offering, and yet a while, and they shall be regarded; \Star{} For the grace of God, and His peace, are with His chosen.\end{Resp}
\begin{Resp}\Vsign They that put their trust in Him shall understand the truth and such as be faithful in love shall abide with Him.\end{Resp}
\begin{Resp}\Rsign For the grace of God, and His peace, are with His chosen.\end{Resp}
\begin{Resp}\Vsign Glory be to the Father, and to the Son, \Star{} and to the Holy Ghost.\end{Resp}
\begin{Resp}\Rsign For the grace of God, and His peace, are with His chosen.\end{Resp}
\switchcolumn*

\LA
\LessonHead{Lectio VII}
\switchcolumn
\EN
\LessonHead{Lesson VII}
\switchcolumn*

\LA
\LessonTitle{Léctio sancti Evangélii secúndum Matthǽum}
\switchcolumn
\EN
\LessonTitle{From the Holy Gospel according to Matthew}
\switchcolumn*

\LA
\Cite{Matt 24:3-13}
\switchcolumn
\EN
\Cite{Matt 24:3-13}
\switchcolumn*

\LA
\noindent In illo témpore: Sedénte Jesu super montem Oliveti, accessérunt ad eum discípuli secréto dicéntes: Dic nobis quando hæc erunt? Et réliqua.\par
\switchcolumn
\EN
\noindent At that time As Jesus sat upon the Mount of Olives, His disciples came unto Him privately, saying Tell us, when shall these things be And so on.\par
\switchcolumn*

\LA
\LessonTitle{Homilía sancti Hilarii Epíscopi}
\switchcolumn
\EN
\LessonTitle{Homily by St. Hilary, Bishop (of Poitiers.)}
\switchcolumn*

\LA
\Source{Comment. in Matth. cap. 24}
\switchcolumn
\EN

\switchcolumn*

\LA
\noindent Discipuli Dóminum intérrogant quando hæc fíerent, quodve signum et adventus sui et consummatiónis sæculi nóscerent. Et, quia tria hæc in unum quæsíta sunt, distinctis et témporis et intelligéntiæ significatiónibus separántur. Respondétur ígitur primum de civitátis occasu, et confirmántur veritáte doctrinæ, ne quis fallax ignorántibus posset obrépere; venturi enim erant étiam eórum témpore, qui se Christum essent nuncupatúri. Ut ígitur fides pestífero mendácio détrahi posset, admonítio præcessit.\par
\switchcolumn
\EN

\switchcolumn*

\LA
\begin{Resp}\Rsign Propter testaméntum Dómini, et leges patérnas, Sancti Dei perstitérunt in amóre fraternitátis: \Star{} Quia unus fuit semper spíritus in eis, et una fides.\end{Resp}
\begin{Resp}\Vsign Ecce quam bonum et quam jucúndum habitáre fratres in unum.\end{Resp}
\begin{Resp}\Rsign Quia unus fuit semper spíritus in eis, et una fides.\end{Resp}
\switchcolumn
\EN
\begin{Resp}\Rsign Because of the covenant of the Lord, and the laws of their fathers, the Saints of God abode in brotherly love, \Star{} For one spirit and one faith was ever in them.\end{Resp}
\begin{Resp}\Vsign Behold how good and how pleasant it is for brethren to dwell together in unity.\end{Resp}
\begin{Resp}\Rsign For one spirit and one faith was ever in them.\end{Resp}
\switchcolumn*

\LA
\LessonHead{Lectio VIII}
\switchcolumn
\EN
\LessonHead{Lesson VIII}
\switchcolumn*

\LA

\switchcolumn
\EN
\LessonTitle{“See that ye be not troubled,” saith the Lord, “for all these things must come to pass, but the end is not yet. (For nation shall rise up against nation, and kingdom against kingdom, and there shall be famines, and pestilences, and earthquakes in diverse places. All these are the beginning of sorrows.) Then shall they deliver you up to be afflicted, and shall kill you and ye shall be hated of all nations for My Name's sake. And then shall many be offended, and shall betray one another, and shall hate one another. And many false prophets shall rise and shall deceive many.” (One such false prophet was Nicolas, one of the seven Deacons.) “And because iniquity shall abound, the love of many shall wax cold.”}
\switchcolumn*

\LA
\noindent Confírmat ígitur eos ad tolerántiam passiónum, fugæ, verberatiónis, intéritus et publici in eos odii propter nomen ejus. Atque his quidem vexatiónibus multi turbabúntur, et, tantis insurgéntibus malis scandalizabúntur, et usque in mutuum ódium excitabúntur. Et falsi prophetæ erunt (ut Nicolaus unus ex septem diacónibus fuit), multosque ementíta veritáte pervertent; et, abundante nequitia, caritas refrigéscet.\par
\switchcolumn
\EN

\switchcolumn*

\LA
\begin{Resp}\Rsign Sancti mei, qui in carne pósiti, certámen habuístis: \Star{} Mercédem labóris ego reddam vobis.\end{Resp}
\begin{Resp}\Vsign Veníte benedícti Patris mei, percípite regnum.\end{Resp}
\begin{Resp}\Rsign Mercédem labóris ego reddam vobis.\end{Resp}
\begin{Resp}\Vsign Glória Patri, et Fílio, \Star{} et Spirítui Sancto.\end{Resp}
\begin{Resp}\Rsign Mercédem labóris ego reddam vobis.\end{Resp}
\switchcolumn
\EN
\begin{Resp}\Rsign O ye My Saints, who, being in the flesh, didst have striving \Star{} I will render unto you a reward of your labours.\end{Resp}
\begin{Resp}\Vsign Come, ye blessed of My Father, inherit the kingdom\end{Resp}
\begin{Resp}\Rsign I will render unto you a reward of your labours.\end{Resp}
\begin{Resp}\Vsign Glory be to the Father, and to the Son, \Star{} and to the Holy Ghost.\end{Resp}
\begin{Resp}\Rsign I will render unto you a reward of your labours.\end{Resp}
\switchcolumn*

\LA
\LessonHead{Lectio IX}
\switchcolumn
\EN
\LessonHead{Lesson IX}
\switchcolumn*

\LA
\noindent Sed usque in finem perseverántibus salus reserváta est; ac tum, per omnes orbis partes viris apostolicis dispersis, Evangélii véritas prædicábitur. Et, cum univérsis fúerit cognítio sacraménti cæléstis invecta, tum Jerúsalem occasus et finis incumbet; ut prædicatiónis fidem, et infidelium pœna et metus civitátis érutæ consequátur. Hæc ígitur in eam, ut fuerant prædicta, perfécta sunt; et, lapidatis, fugatis, peremptis Apóstolis, fame, bello, captivitáte consumpta est. Ac tum fuit digna non esse, cum, ejectis prædicatóribus Christi, indignam Dei prædicatióne se præbuit.\par
\switchcolumn
\EN
\noindent But that shall endure unto the end, the same shall be saved. And this Gospel shall be preached in all the world, for a witness unto all nations and then shall the end come.” When the knowledge of the heavenly revelation had been carried everywhere, then should come the fall and end of Jerusalem then should the punishment of them that had not believed, and the awful example of the city that had been destroyed, bear out the truth of the preacher. When she had stoned, and hunted down, and murdered the Apostles, then should she be consumed by famine, and war, and slavery. And indeed she would then have shown herself unworthy to be any longer, having shown by casting out the preachers of Christ that she was unworthy that any should speak to her of God.\par
\switchcolumn*

\LA
\TeDeum{Te Deum}
\switchcolumn
\EN
\TeDeum{Te Deum}
\switchcolumn*

\end{paracol}

\par\needspace{10\baselineskip}\addvspace{1.2em}{\centering\small\textsc{Die 20 Septembris} · \textsc{20 September}\par}
\Office{Ss. Eustachii et Sociorum Martyrum}{S. Eustachius and Companions, Martyrs}{Duplex}
\addcontentsline{toc}{subsection}{Die 20 Septembris · Ss. Eustachii et Sociorum Martyrum \textit{— S. Eustachius and Companions, Martyrs}}
\begin{paracol}{2}
\LA
\RefLine{Lectiones I–III: de Scriptura occurrente.}
\switchcolumn
\EN
\RefLine{Lessons I–III: from the Scripture of the day.}
\switchcolumn*

\LA
\RefLine{Lectio IX: de In Vigilia S. Matthæi Ap. (09-20o).}
\switchcolumn
\EN
\RefLine{Lesson IX: of In Vigilia S. Matthew Apostle (09-20o).}
\switchcolumn*

\LA
\LessonHead{Lectio IV}
\switchcolumn
\EN
\LessonHead{Lesson IV}
\switchcolumn*

\LA
\noindent Eustachius, qui et Placidus, genere, opibus et militari glória inter Romanos insígnis, sub Trajano imperatóre magistri militum titulum meruit. Cum vero sese aliquándo in venatióne exercéret ac fugiéntem miræ magnitúdinis cervum insequerétur, vidit repénte inter consistentis feræ córnua excélsam atque fulgentem Christi Dómini e cruce pendentis imáginem. Cujus voce ad immortalis vitæ prædam invitatus, una cum uxore Theopista ac duobus párvulis fíliis Agapito et Theopisto, christianæ milítiæ nomen dedit.\par
\switchcolumn
\EN
\noindent Eustache, whose name before his Baptism was Placidus was a Roman, alike well-known on account of his noble birth, his great earthly wealth, and his eminent distinction as a soldier. He gained, under the Emperor Trajan, the post of military commander. Once upon a time he was hunting, and following an extraordinarily large stag, when the beast stood still, and Eustace saw between his horns a tall and glorious figure of the Lord Christ hanging upon the Cross, whence came a voice bidding him to follow after life eternal. Thereupon Eustace and his wife Theopista, and their two little sons Agapitus and Theopistus, enlisted themselves as soldiers under the Great Captain, Christ.\par
\switchcolumn*

\LA
\begin{Resp}\Rsign Sancti tui, Dómine, mirábile consecúti sunt iter, serviéntes præcéptis tuis, ut inveniréntur illǽsi in aquis válidis: \Star{} Terra appáruit árida: et in Mari Rubro via sine impediménto.\end{Resp}
\begin{Resp}\Vsign Quóniam percússit petram, et fluxérunt aquæ, et torréntes inundavérunt.\end{Resp}
\begin{Resp}\Rsign Terra appáruit árida: et in Mari Rubro via sine impediménto.\end{Resp}
\switchcolumn
\EN
\begin{Resp}\Rsign thy Saints, O Lord, have passed a wonderful way, serving thy commandments, that they might be found without hurt in the midst of the mighty waters. \Star{} Dry land appeared, and, out of the Red Sea, a way without impediment.\end{Resp}
\begin{Resp}\Vsign He smote the rock, and the waters gushed out, and the streams overflowed.\end{Resp}
\begin{Resp}\Rsign Dry land appeared, and, out of the Red Sea, a way without impediment.\end{Resp}
\switchcolumn*

\LA
\LessonHead{Lectio V}
\switchcolumn
\EN
\LessonHead{Lesson V}
\switchcolumn*

\LA
\noindent Mox ad visiónis prístinæ locum, sicut ei Dóminus præceperat, regréssus, illum prænuntiántem audívit quanta sibi deinceps, pro ejus gloria, perferenda essent. Quocirca incredibiles calamitátes mira patiéntia perpéssus, brevi in summam egestátem redáctus est. Cumque clam se subducere cogerétur, in itinere cónjugem primum, deínde étiam líberos sibi miserabíliter ereptos ingemuit. Tantis obvolutus ærumnis, in regióne longinqua villicum agens longo témpore delituit, donec, cælésti voce recreátus ac nova occasióne a Trajano conquisítus, íterum bello præfícitur.\par
\switchcolumn
\EN
\noindent In a little while he went back, according as the Lord had commanded him, to the place where he had seen the first vision, and there he heard from God how much he was to bear for His glory. It was not long after that he had great losses and became exceedingly poor, but he bore it very patiently. Then he was constrained to fly away privily, and on the journey was grievously afflicted in that, first, his wife and then his children were parted from him and carried he knew not whither. Under the weight of these sorrows he lay hid a long while in a far-off place, working as the steward of a land -owner, until the voice of God called him forth, and Trajan sought for him again to make him a captain in his army.\par
\switchcolumn*

\LA
\begin{Resp}\Rsign Vérbera carníficum non timuérunt Sancti Dei, moriéntes pro Christi nómine: \Star{} Ut herédes fíerent in domo Dómini.\end{Resp}
\begin{Resp}\Vsign Tradidérunt córpora sua propter Deum ad supplícia.\end{Resp}
\begin{Resp}\Rsign Ut herédes fíerent in domo Dómini.\end{Resp}
\switchcolumn
\EN
\begin{Resp}\Rsign The Saints of God shrank not from the stripes of the executioners, but died for Christ's Name's sake \Star{} That they might be made joint-heirs in the house of the Lord.\end{Resp}
\begin{Resp}\Vsign They gave their bodies for God's sake to death.\end{Resp}
\begin{Resp}\Rsign That they might be made joint-heirs in the house of the Lord.\end{Resp}
\switchcolumn*

\LA
\LessonHead{Lectio VI}
\switchcolumn
\EN
\LessonHead{Lesson VI}
\switchcolumn*

\LA
\noindent Illa in expeditióne, liberis simul cum uxore insperato receptis, victor Urbem ingénti ómnium gratulatióne ingréditur. Sed paulo post inánibus diis pro parta victoria sacrificáre jussus, constantíssime rénuit. Cumque variis artibus ad Christi fidem ejurandam frustra tentarétur, una cum uxore et liberis, leónibus objícitur. Horum mansuetúdine concitátus imperátor, æneum in taurum subjectis flammis candéntem eos immítti jubet, ubi divinis in laudibus consummato martyrio, duodecimo Kalendas Octobris ad sempitérnam felicitátem convolárunt. Quorum illæsa córpora, religióse a fidélibus sepúlta, póstmodum ad ecclésiam eórum nómine erectam honorofice transláta sunt.\par
\switchcolumn
\EN
\noindent While he was with the army he found his wife and children once more, by an unexpected happiness, and re-entered the city (of Rome) as a conquering soldier amid the loud applause of all men, but thereupon, when he was commanded to offer sacrifices of thanksgiving for the victory to the gods that are no gods, he stoutly refused. They tried him in vain with diverse cajoleries to make him deny Christ, but could not, and he and his wife and little ones were thrown to the lions. When these beasts would not touch them, the Emperor's fury was kindled, and he commanded them all to be shut up in the brazen image of a bull, which was heated with fire underneath. There they praised God until their testimony was ended, and they departed hence to be perfectly blessed for ever and ever, upon the 20th day of September. Their bodies were buried whole by the faithful, with deep reverence, and were afterwards honourably carried to a Church built in their name.\par
\switchcolumn*

\LA
\begin{Resp}\Rsign Tamquam aurum in fornáce probávit eléctos Dóminus, et quasi holocáusti hóstiam accépit illos: et in témpore erit respéctus illórum: \Star{} Quóniam donum et pax est eléctis Dei.\end{Resp}
\begin{Resp}\Vsign Qui confídunt in illum, intélligent veritátem, et fidéles in dilectióne acquiéscent illi.\end{Resp}
\begin{Resp}\Rsign Quóniam donum et pax est eléctis Dei.\end{Resp}
\begin{Resp}\Vsign Glória Patri, et Fílio, \Star{} et Spirítui Sancto.\end{Resp}
\begin{Resp}\Rsign Quóniam donum et pax est eléctis Dei.\end{Resp}
\switchcolumn
\EN
\begin{Resp}\Rsign As gold in the furnace hath the Lord tried His chosen ones, and received them as a burnt-offering, and yet a while, and they shall be regarded; \Star{} For the grace of God, and His peace, are with His chosen.\end{Resp}
\begin{Resp}\Vsign They that put their trust in Him shall understand the truth and such as be faithful in love shall abide with Him.\end{Resp}
\begin{Resp}\Rsign For the grace of God, and His peace, are with His chosen.\end{Resp}
\begin{Resp}\Vsign Glory be to the Father, and to the Son, \Star{} and to the Holy Ghost.\end{Resp}
\begin{Resp}\Rsign For the grace of God, and His peace, are with His chosen.\end{Resp}
\switchcolumn*

\LA
\LessonHead{Lectio VII}
\switchcolumn
\EN
\LessonHead{Lesson VII}
\switchcolumn*

\LA
\LessonTitle{Léctio sancti Evangélii secúndum Lucam}
\switchcolumn
\EN
\LessonTitle{From the Holy Gospel according to Luke}
\switchcolumn*

\LA
\Cite{Luc 6:17-23}
\switchcolumn
\EN
\Cite{Luke 6:17-23}
\switchcolumn*

\LA
\noindent In illo témpore: Descéndens Jesus de monte, stetit in loco campéstri, et turba discipulórum ejus, et multitúdo copiósa plebis ab omni Judǽa, et Jerúsalem, et marítima, et Tyri, et Sidónis. Et réliqua.\par
\switchcolumn
\EN
\noindent At that time, Jesus came down from the mountain, and stood in the plain, and the company of His disciples, and a great multitude of people out of all Judea, and Jerusalem, and from the sea coast of Tyre and Sidon. And so on.\par
\switchcolumn*

\LA
\LessonTitle{Homilía sancti Ambrósii Epíscopi}
\switchcolumn
\EN
\LessonTitle{Homily by St Ambrose, Bishop (of Milan.)}
\switchcolumn*

\LA
\Source{Lib. 5. in Lucam. Cap. 6., post init.}
\switchcolumn
\EN
\Source{Bk. v. on Luke vi.}
\switchcolumn*

\LA
\noindent Advérte ómnia diligénter, quómodo et cum Apóstolis ascéndat et descéndat ad turbas. Quómodo enim turba nisi in húmili Christum vidéret? Non séquitur ad excélsa, non ascéndit ad sublímia. Dénique ubi descéndit, invénit infírmos: in excélsis enim infírmi esse non possunt. Hinc étiam Matthǽus docet in inferióribus débiles esse sanátos. Prius enim unusquísque sanándus est, ut paulátim virtútibus procedéntibus ascéndere possit ad montem; et ídeo quemque in inferióribus sanat, hoc est, a libídine révocat, injúriam cæcitátis avértit. Ad vúlnera nostra descéndit: ut usu quodam et cópia suæ natúræ, compartícipes nos fáciat esse regni cæléstis.\par
\switchcolumn
\EN
\noindent Mark well how Jesus goeth upward with His disciples, and downward to the multitude. How should the multitude behold Christ, save in a lower place? Such go not up to the things which are above; such attain not to the things which are high. And when Jesus cometh down, He findeth such as are diseased for such like go not up to the heights. Hence also Matthew saith that there were there all sick people, (iv. 23.) Of these every man had need of healing, that, when he had received strength, by and by, he might go up into the mountain. And therefore, being Himself come down, He healeth them in the plain, that is to say, He calleth them away from their lust, and freeth them of their blindness. He cometh down to our wounds, to the end that by a certain use of His nature, and by the abundance thereof, He might make us joint-heirs of the kingdom of heaven.\par
\switchcolumn*

\LA
\begin{Resp}\Rsign Propter testaméntum Dómini, et leges patérnas, Sancti Dei perstitérunt in amóre fraternitátis: \Star{} Quia unus fuit semper spíritus in eis, et una fides.\end{Resp}
\begin{Resp}\Vsign Ecce quam bonum et quam jucúndum habitáre fratres in unum.\end{Resp}
\begin{Resp}\Rsign Quia unus fuit semper spíritus in eis, et una fides.\end{Resp}
\switchcolumn
\EN
\begin{Resp}\Rsign Because of the covenant of the Lord, and the laws of their fathers, the Saints of God abode in brotherly love, \Star{} For one spirit and one faith was ever in them.\end{Resp}
\begin{Resp}\Vsign Behold how good and how pleasant it is for brethren to dwell together in unity.\end{Resp}
\begin{Resp}\Rsign For one spirit and one faith was ever in them.\end{Resp}
\switchcolumn*

\LA
\LessonHead{Lectio VIII}
\switchcolumn
\EN
\LessonHead{Lesson VIII}
\switchcolumn*

\LA
\noindent Beáti páuperes, quia vestrum est regnum Dei. Quátuor tantum beatitúdines sanctus Lucas Domínicas posuit, octo vero sanctus Matthǽus: sed in illis octo istæ quátuor sunt, et in quátuor istis illæ octo. Hic enim quátuor velut virtútes ampléxus est cardináles: ille in illis octo mýsticum númerum reserávit. Pro octáva enim multi inscribúntur Psalmi: et mandátum áccipis octo illis partem dare fortásse benedictiónibus. Sicut enim spei nostræ octáva perféctio est, ita octáva summa virtútum est.\par
\switchcolumn
\EN
\noindent Blessed be ye poor, for your's is the kingdom of God. Saint Luke giveth us but four of the Lord's Beatitudes, and Saint Matthew eight but in those eight are contained these four, and in these four those eight. For in these four are embraced the cardinal virtues and in those eight they are set forth in a number full of mystery. It is written at the head of more than one of the Psalms that they are for the octave, and thou hast received the commandment Give a portion to seven, and also to eight to seven or eight what? Perchance degrees of blessedness. For as this eighth Beatitude doth name the most glorious realization of our hope the kingdom of Heaven so doth it also name the most royal exertion of our strength blessed are they which are persecuted.\par
\switchcolumn*

\LA
\begin{Resp}\Rsign Sancti mei, qui in carne pósiti, certámen habuístis: \Star{} Mercédem labóris ego reddam vobis.\end{Resp}
\begin{Resp}\Vsign Veníte benedícti Patris mei, percípite regnum.\end{Resp}
\begin{Resp}\Rsign Mercédem labóris ego reddam vobis.\end{Resp}
\begin{Resp}\Vsign Glória Patri, et Fílio, \Star{} et Spirítui Sancto.\end{Resp}
\begin{Resp}\Rsign Mercédem labóris ego reddam vobis.\end{Resp}
\switchcolumn
\EN
\begin{Resp}\Rsign O ye My Saints, who, being in the flesh, didst have striving \Star{} I will render unto you a reward of your labours.\end{Resp}
\begin{Resp}\Vsign Come, ye blessed of My Father, inherit the kingdom\end{Resp}
\begin{Resp}\Rsign I will render unto you a reward of your labours.\end{Resp}
\begin{Resp}\Vsign Glory be to the Father, and to the Son, \Star{} and to the Holy Ghost.\end{Resp}
\begin{Resp}\Rsign I will render unto you a reward of your labours.\end{Resp}
\switchcolumn*

\end{paracol}

\par\needspace{10\baselineskip}\addvspace{1.2em}{\centering\small\textsc{Die 20 Septembris} · \textsc{20 September}\par}
\Office{In Vigilia S. Matthæi Ap.}{In Vigilia S. Matthew Apostle}{Vigilia}
\addcontentsline{toc}{subsection}{Die 20 Septembris · In Vigilia S. Matthæi Ap. \textit{— In Vigilia S. Matthew Apostle}}
\begin{paracol}{2}
\LA
\LessonHead{Lectio I}
\switchcolumn
\EN
\LessonHead{Lesson I}
\switchcolumn*

\LA
\LessonTitle{Commemoratio: In Vigilia S. Matthæi Ap. Léctio sancti Evangélii secúndum Lucam}
\switchcolumn
\EN
\LessonTitle{Commemoration of: In Vigilia S. Matthew Apostle Lesson from the Gospel of St. Luke}
\switchcolumn*

\LA
\Cite{Luc 5:27-32}
\switchcolumn
\EN
\Cite{Luke 5:27-32}
\switchcolumn*

\LA
\noindent In illo témpore: Vidit Jesus publicanum, nómine Levi, sedentem ad telónium, et ait illi : Séquere me. Et réliqua.\par
\switchcolumn
\EN
\noindent At that time, Jesus saw a publican, named Levi, sitting at the receipt of custom, and He said unto him: Follow me. And so on.\par
\switchcolumn*

\LA
\LessonTitle{Homilía sancti Ambrósii Epíscopi}
\switchcolumn
\EN
\LessonTitle{Homily by St. Ambrose, Bishop (of Milan)}
\switchcolumn*

\LA
\Source{Liber 5 Comment. in Lucæ Cap. 5 post initium}
\switchcolumn
\EN
\Source{Book v, Comment. on Luke v}
\switchcolumn*

\LA
\noindent Mystica est hæc vocátio publicani, quem sequi jubet, non corporis gressu, sed mentis afféctu. Itaque ille, prius avara de mercédibus, dura de labóribus periculisque nautárum emoluménta convértens, verbo vocatus, propria derelinquit, qui rapiebat aliena ; ac, vile illud sedile destítuens, toto post Dóminum vestigio mentis incedit. Convivii quoque magni éxhibet apparátum ; qui enim domicílio Christum récipit interno, maximis delectatiónibus exuberántium páscitur voluptátum.\par
\switchcolumn
\EN
\noindent There is a mystery in this calling of the publican, whom He biddeth to follow Him, not so much by bodily steps as by change of heart. Hitherto Levi had been making greedy gains from merchandise, cruel riches at the cost of sailors' toils and dangers; but now, at the call of a word, he, who had been plundering other men's goods, leaveth his own. He leaveth that base station, and followeth hard after the Lord with all his heart. And Levi made Him a great feast in his own house. He that welcometh Christ into his home, feasteth upon the excellency of all pleasures.\par
\switchcolumn*

\LA
\TeDeum{Te Deum}
\switchcolumn
\EN
\TeDeum{Te Deum}
\switchcolumn*

\end{paracol}

\par\needspace{10\baselineskip}\addvspace{1.2em}{\centering\small\textsc{Die 21 Septembris} · \textsc{21 September}\par}
\Office{S. Matthæi Apostoli et Evangelistæ}{St. Matthew, Apostle and Evangelist}{Duplex II classis}
\addcontentsline{toc}{subsection}{Die 21 Septembris · S. Matthæi Apostoli et Evangelistæ \textit{— St. Matthew, Apostle and Evangelist}}
\begin{paracol}{2}
\LA
\RefLine{Lectiones I–III: ex Commúni Evangelistarum (C1a).}
\switchcolumn
\EN
\RefLine{Lessons I–III: from the Commune Evangelistarum (C1a).}
\switchcolumn*

\LA
\RefLine{Lectio IX: de Scriptura occurrente.}
\switchcolumn
\EN
\RefLine{Lesson IX: from the Scripture of the day.}
\switchcolumn*

\LA
\LessonHead{Lectio IV}
\switchcolumn
\EN
\LessonHead{Lesson IV}
\switchcolumn*

\LA
\noindent Matthǽus, qui et Levi, Apóstolus et Evangelísta, Caphárnai cum ad telónium sedéret, a Christo vocátus, statim secútus est ipsum; quem étiam cum réliquis discípulis convívio excépit. Post Christi resurrectiónem, ántequam in provínciam proficiscerétur, quæ ei ad prædicándum obtígerat, primus in Judǽa, propter eos qui ex circumcisióne credíderant, Evangélium Jesu Christi Hebráice scripsit. Mox in Æthiópiam proféctus, Evangélium prædicávit, ac prædicatiónem multis miráculis confirmávit.\par
\switchcolumn
\EN
\noindent It came to pass one day at Capernaum, that Christ went forth, and saw a publican, named Levi, sitting at the receipt of custom; and He said unto him: Follow Me. And he left all, rose up, and followed Him. And Levi made Him a great feast in his own house. This Levi is the Apostle and Evangelist Matthew. After that Christ was risen again from the dead, and while he was yet in Judea, before he set forth for that land which had fallen to the lot of his preaching, he wrote the Gospel of Jesus Christ in the Hebrew tongue, for the sake of them of the circumcision who had believed. His was the first written of the four Gospels. Thereafter he went to Ethiopia, and there preached the Gospel, confirming his preaching with many miracles.\par
\switchcolumn*

\LA
\begin{Resp}\Rsign Vidi conjúnctos viros, habéntes spléndidas vestes, et Ángelus Dómini locútus est ad me, dicens: \Star{} Isti sunt viri sancti facti amíci Dei.\end{Resp}
\begin{Resp}\Vsign Vidi Ángelum Dei fortem, volántem per médium cælum, voce magna clamántem et dicéntem.\end{Resp}
\begin{Resp}\Rsign Isti sunt viri sancti facti amíci Dei.\end{Resp}
\switchcolumn
\EN
\begin{Resp}\Rsign I saw men standing together, clad in shining raiment, and the Angel of the Lord spoke unto me, saying: \Star{} These men are holy, for they are the friends of God.\end{Resp}
\begin{Resp}\Vsign I saw a strong Angel of God fly into the midst of heaven, saying with a loud voice:\end{Resp}
\begin{Resp}\Rsign These men are holy, for they are the friends of God.\end{Resp}
\switchcolumn*

\LA
\LessonHead{Lectio V}
\switchcolumn
\EN
\LessonHead{Lesson V}
\switchcolumn*

\LA
\noindent Illo ígitur in primis miráculo, quo regis fíliam a mórtuis excitávit, regem patrem et uxórem ejus cum univérsa província ad Christi fidem convértit. Rege mórtuo, Hírtacus, ejus succéssor, cum Iphigeníam, régiam fíliam, vellet sibi dari in matrimónium; Matthǽum, cujus ópera illa virginitátem Deo vóverat et in sancto propósito perseverábat, ad altáre mystérium celebrántem jussit occídi. Qui undécimo caléndas octóbris munus apostólicum martýrii glória cumulávit. Cujus corpus Salérnum translátum, ac póstmodum in ecclésia ejus nómine dedicáta, Gregório séptimo summo Pontífice, cónditum, ibídem magno hóminum concúrsu ac pietáte cólitur.\par
\switchcolumn
\EN
\noindent Of his miracles, the most notable was that he raised the King's daughter from the dead, and thereby brought to believe in Christ the King her father, his wife, and all that region. After that the King was dead, Hirtacus, who came after him, was fain to take his daughter Iphigenia to wife, but by the exhortation of Matthew she had made vow of her maidenhood to God, and stood firm to that holy resolution, for which cause Hirtacus commanded to slay the Apostle at the Altar while he was performing the mystery. He crowned the dignity of the Apostleship with the glory of martyrdom upon the 21 st day of September. His body had been brought to Salerno, where it was afterwards buried in a Church dedicated in his name during the papacy of Gregory VII, and there it is held in great worship and sought to by great gatherings of people.\par
\switchcolumn*

\LA
\begin{Resp}\Rsign Beáti estis, cum maledíxerint vobis hómines, et persecúti vos fúerint, et díxerint omne malum advérsum vos, mentiéntes, propter me: \Star{} Gaudéte et exsultáte, quóniam merces vestra copiósa est in cælis.\end{Resp}
\begin{Resp}\Vsign Cum vos óderint hómines, et cum separáverint vos, et exprobráverint, et ejécerint nomen vestrum tamquam malum propter Fílium hóminis.\end{Resp}
\begin{Resp}\Rsign Gaudéte et exsultáte, quóniam merces vestra copiósa est in cælis.\end{Resp}
\switchcolumn
\EN
\begin{Resp}\Rsign Blessed are ye when men shall revile you, and persecute you, and shall say all manner of evil against you falsely, for My sake; \Star{} Rejoice, and be exceeding glad, for great is your reward in heaven.\end{Resp}
\begin{Resp}\Vsign When men shall hate you, and when they shall separate you from their company, and shall reproach you, and cast out your name as evil, for the Son of Man's sake.\end{Resp}
\begin{Resp}\Rsign Rejoice, and be exceeding glad, for great is your reward in heaven.\end{Resp}
\switchcolumn*

\LA
\LessonHead{Lectio VI}
\switchcolumn
\EN
\LessonHead{Lesson VI}
\switchcolumn*

\LA
\LessonTitle{De Expositióne sancti Gregórii Papæ super Ezechiélem Prophétam}
\switchcolumn
\EN
\LessonTitle{From the Exposition of the Book of the Prophet Ezekiel by Pope St. Gregory (the Great.)}
\switchcolumn*

\LA
\Source{Homilia 3, Lib. 1}
\switchcolumn
\EN
\Source{Hom. 3, Bk. I}
\switchcolumn*

\LA
\noindent Sancta quátuor animália, quæ prophetíæ spíritu futúra prævidéntur, subtíli narratióne describúntur, cum dícitur: Quátuor fácies uni, et quátuor pennæ uni. Quid per fáciem, nisi notítia; et quid per pennas, nisi volátus exprímitur? Per fáciem quippe unusquísque cognóscitur: per pennas vero in altum ávium córpora sublevántur. Fácies ítaque ad fidem pértinet, penna ad contemplatiónem. Per fidem namque ab omnipoténti Deo cognóscimur, sicut ipse de suis óvibus dicit: Ego sum pastor bonus, et cognósco oves meas, et cognóscunt me meæ. Qui rursus ait: Ego scio quos elégerim. Per contemplatiónem vero, qua super nosmetípsos tóllimur, quasi in áëra levámur.\par
\switchcolumn
\EN
\noindent The Prophet writeth very minutely touching the four holy living creatures, whom he saw in the spirit as being to come. He saith: “Every one had four faces, and every one had four wings.” What signifieth the face save likeness whereby we are known? or wings, save the power to fly since it is by the face that man is known from man, and by their wings that the birds' bodies are carried up into the air. So the face pertaineth to certitude, and the wings to contemplation. With certitude we are known of God Almighty, Who saith: “I am the Good Shepherd, and know My sheep, and am known of Mine.” (John x. 14.) And again: “I know whom I have chosen.” (xiii. 18.) And by contemplation, whereby we rise above ourselves, we as it were fly heavenwards.\par
\switchcolumn*

\LA
\begin{Resp}\Rsign Isti sunt triumphatóres et amíci Dei, qui contemnéntes jussa príncipum, meruérunt prǽmia ætérna: \Star{} Modo coronántur, et accípiunt palmam.\end{Resp}
\begin{Resp}\Vsign Isti sunt qui venérunt ex magna tribulatióne, et lavérunt stolas suas in sánguine Agni.\end{Resp}
\begin{Resp}\Rsign Modo coronántur, et accípiunt palmam.\end{Resp}
\begin{Resp}\Vsign Glória Patri, et Fílio, \Star{} et Spirítui Sancto.\end{Resp}
\begin{Resp}\Rsign Modo coronántur, et accípiunt palmam.\end{Resp}
\switchcolumn
\EN
\begin{Resp}\Rsign These are they which have conquered, and are become the friends of God, who recked not of the commandments of princes, and earned the everlasting reward. \Star{} And now have they crowns on their heads, and palms in their hands.\end{Resp}
\begin{Resp}\Vsign These are they which came out of great tribulation, and have washed their robes in the blood of the Lamb.\end{Resp}
\begin{Resp}\Rsign And now have they crowns on their heads, and palms in their hands.\end{Resp}
\begin{Resp}\Vsign Glory be to the Father, and to the Son, \Star{} and to the Holy Ghost.\end{Resp}
\begin{Resp}\Rsign And now have they crowns on their heads, and palms in their hands.\end{Resp}
\switchcolumn*

\LA
\LessonHead{Lectio VII}
\switchcolumn
\EN
\LessonHead{Lesson VII}
\switchcolumn*

\LA
\LessonTitle{Léctio sancti Evangélii secúndum Matthǽum}
\switchcolumn
\EN
\LessonTitle{From the Holy Gospel according to Matthew}
\switchcolumn*

\LA
\Cite{Matt 9:9-13}
\switchcolumn
\EN
\Cite{Matt 9:1-13}
\switchcolumn*

\LA
\noindent In illo témpore: Vidit Jesus hóminem sedéntem in telónio, Matthǽum nómine, et ait illi: Séquere me. Et réliqua.\par
\switchcolumn
\EN
\noindent At that time, Jesus saw a man, named Matthew, sitting at the receipt of custom; and He saith unto him Follow Me. And so on.\par
\switchcolumn*

\LA
\LessonTitle{Homilía sancti Hierónymi Presbýteri}
\switchcolumn
\EN
\LessonTitle{Homily by St. Jerome, Priest (at Bethlehem.)}
\switchcolumn*

\LA
\Source{Liber 1 Comment. in Matth. cap. 9}
\switchcolumn
\EN
\Source{Bk. i. Comment, on Matth. ix.}
\switchcolumn*

\LA
\noindent Céteri Evangelístæ, propter verecúndiam et honórem Matthǽi, noluérunt eum nómine appelláre vulgáto, sed dixérunt, Levi; dúplici quippe vocábulo fuit. Ipse autem Matthǽus (secúndum illud quod dícitur a Salomóne: Justus accusátor est sui in princípio sermónis; et in álio loco: Dic tu peccáta tua, ut justificéris) Matthǽum se et publicanum nóminat, ut osténdat legéntibus nullum debére salútem desperáre, si ad melióra convérsus sit, cum ipse de publicáno in Apóstolum sit repénte mutátus.\par
\switchcolumn
\EN
\noindent The other Evangelists, out of tenderness towards the reputation and honour of Matthew, have abstained from speaking of him as a publican by his ordinary name, and have called him Levi. Both names were his. But Matthew himself, according to what Solomon saith: The just man is the first to accuse himself, (Prov. xviii. 17), and again, in another place: Declare thou thy sins that thou mayest be justified, doth plainly call himself Matthew the publican, to show unto his readers that none need be hopeless of salvation if he will but strive to do better, since he himself had been all of a sudden changed from a publican into an Apostle.\par
\switchcolumn*

\LA
\begin{Resp}\Rsign Isti sunt qui vivéntes in carne, plantavérunt Ecclésiam sánguine suo: \Star{} Cálicem Dómini bibérunt, et amíci Dei facti sunt.\end{Resp}
\begin{Resp}\Vsign In omnem terram exívit sonus eórum, et in fines orbis terræ verba eórum.\end{Resp}
\begin{Resp}\Rsign Cálicem Dómini bibérunt, et amíci Dei facti sunt.\end{Resp}
\switchcolumn
\EN
\begin{Resp}\Rsign These are they who while yet they lived in the flesh, planted the Church in their own blood; \Star{} They drank of the Lord's cup, and became the friends of God.\end{Resp}
\begin{Resp}\Vsign Their sound is gone out through all the earth, and their words to the ends of the world.\end{Resp}
\begin{Resp}\Rsign They drank of the Lord's cup, and became the friends of God.\end{Resp}
\switchcolumn*

\LA
\LessonHead{Lectio VIII}
\switchcolumn
\EN
\LessonHead{Lesson VIII}
\switchcolumn*

\LA
\noindent Arguit in hoc loco Porphýrius, et Juliánus Augústus, vel imperítiam histórici mentiéntis, vel stultítiam eórum qui statim secúti sint Salvatórem, quasi irrationabíliter quémlibet vocántem hóminem sint secúti; cum tantæ virtútes tántaque signa præcésserint, quæ Apóstolos, ántequam créderent, vidísse non dubium est. Certe fulgor ipse et majéstas Divinitátis occúltæ, quæ étiam in humána fácie relucébat, ex primo ad se vidéntes tráhere póterat aspéctu. Si enim in magnéte lápide et súccinis hæc esse vis dícitur, ut ánulos et stípulam et festúcas sibi cópulent; quanto magis Dóminus ómnium creaturárum ad se tráhere póterat, quos vocábat?\par
\switchcolumn
\EN
\noindent Porphyry and the Emperor Julian (the Apostate) will have it that the account of this call of Matthew is either a stupid blunder on the part of a lying writer, or else that it showeth what fools they were who followed the Saviour, to go senselessly after any one who called them. But there can be no doubt that before the Apostles believed they had considered the great signs and works of power which had gone before. Moreover, the glory and majesty of the hidden God, which shone somewhat through the Face of the Man Christ Jesus, were enough to draw them which gazed thereon, even at first sight. For if there be in a stone a magnetic power which can make rings and straws and rods come and cleave thereunto, how much more must not the Lord of all creatures have been able to draw unto Himself them whom He called?\par
\switchcolumn*

\LA
\begin{Resp}\Rsign Isti sunt viri sancti, quos elégit Dóminus in caritáte non ficta, et dedit illis glóriam sempitérnam: \Star{} Quorum doctrína fulget Ecclésia, ut sole luna.\end{Resp}
\begin{Resp}\Vsign Sancti per fidem vicérunt regna: operáti sunt justítiam.\end{Resp}
\begin{Resp}\Rsign Quorum doctrína fulget Ecclésia, ut sole luna.\end{Resp}
\begin{Resp}\Vsign Glória Patri, et Fílio, \Star{} et Spirítui Sancto.\end{Resp}
\begin{Resp}\Rsign Quorum doctrína fulget Ecclésia, ut sole luna.\end{Resp}
\switchcolumn
\EN
\begin{Resp}\Rsign These men are saints, whom the Lord hath chosen in love unfeigned, and hath given them glory everlasting. These are they \Star{} By the light of whose teaching the Church is glorified, even as the moon is glorified by the light of the sun.\end{Resp}
\begin{Resp}\Vsign The saints through faith subdued kingdoms, wrought righteousness.\end{Resp}
\begin{Resp}\Rsign By the light of whose teaching the Church is glorified, even as the moon is glorified by the light of the sun.\end{Resp}
\begin{Resp}\Vsign Glory be to the Father, and to the Son, \Star{} and to the Holy Ghost.\end{Resp}
\begin{Resp}\Rsign By the light of whose teaching the Church is glorified, even as the moon is glorified by the light of the sun.\end{Resp}
\switchcolumn*

\LA
\LessonHead{Lectio IX}
\switchcolumn
\EN
\LessonHead{Lesson IX}
\switchcolumn*

\LA
\noindent Et factum est, discumbénte eo in domo, ecce multi publicáni et peccatóres veniéntes discumbébant cum Jesu. Vidébant publicanum, a peccátis ad melióra convérsum, locum invenísse pœniténtiæ; et ob id étiam ipsi non despérant salútem. Neque vero in pristínis vítiis permanéntes véniunt ad Jesum, ut pharisǽi et scribæ múrmurant, sed pœniténtiam agéntes, ut sequens Dómini sermo signíficat, dicens: Misericórdiam volo, et non sacrifícium; non enim veni vocáre justos, sed peccatóres. Ibat autem Dóminus ad convívia peccatórum, ut occasiónem habéret docéndi, et spirituáles invitatóribus suis præbéret cibos.\par
\switchcolumn
\EN
\noindent And it came to pass, as Jesus sat at meat in the house, behold, many publicans and sinners came and sat down with Him. They saw how that a publican who had turned to better things had found a place of repentance, and therefore they also hoped for salvation. It was not, as the Scribes and Pharisees complained, sinners clinging to their sinfulness who came to Jesus, but sinners repenting, as indeed appeareth from the next words of the Lord, where He saith: I will have mercy and not sacrifice; for I am not come to call the righteous, but sinners to repentance. The Lord went to eat with sinners to the end that He might have occasion to teach, and to break spiritual bread unto them which bade Him.\par
\switchcolumn*

\LA
\TeDeum{Te Deum}
\switchcolumn
\EN
\TeDeum{Te Deum}
\switchcolumn*

\end{paracol}

\par\needspace{10\baselineskip}\addvspace{1.2em}{\centering\small\textsc{Die 22 Septembris} · \textsc{22 September}\par}
\Office{S. Thomæ de Villanova Episcopi et Confessoris}{St. Thomas of Villanova, Bishop and Confessor}{Duplex}
\addcontentsline{toc}{subsection}{Die 22 Septembris · S. Thomæ de Villanova Episcopi et Confessoris \textit{— St. Thomas of Villanova, Bishop and Confessor}}
\begin{paracol}{2}
\LA
\RefLine{Lectiones I–III: de Scriptura occurrente.}
\switchcolumn
\EN
\RefLine{Lessons I–III: from the Scripture of the day.}
\switchcolumn*

\LA
\RefLine{Lectiones VII–VIII: ex Commúni unius Confessoris Pontificis (C4).}
\switchcolumn
\EN
\RefLine{Lessons VII–VIII: from the Common of a Confessor Bishop (C4).}
\switchcolumn*

\LA
\RefLine{Lectio IX: de Ss. Mauritii et Sociorum Martyrum (09-22o).}
\switchcolumn
\EN
\RefLine{Lesson IX: of Ss. Maurice and companions, Martyrs (09-22o).}
\switchcolumn*

\LA
\LessonHead{Lectio IV}
\switchcolumn
\EN
\LessonHead{Lesson IV}
\switchcolumn*

\LA
\noindent Thomas, in oppido Fontisplani Toletanæ diœceseos in Hispánia natus anno Dómini millesimo quadringentésimo octogesimo octavo, ab optimis paréntibus, ineunte vita, pietátem et singularem in páuperes misericórdiam accépit. Cujus adhuc puer complura dedit exempla; sed illud in primis nóbile, quod, ut nudos operíret, propriis vestibus non semel seípsum exuit. Exacta puerítia, Complúto, quo missus fuerat ut alumnus in collegio majori sancti Ildefonsi litteris operam daret, patris óbitu revocatus, universam hereditátem egenis virgínibus aléndis dicávit; eodemque statim reversus est, et, sacræ theologíæ cursu confecto, adeo doctrina excelluit, ut, in eadem universitáte cathedram ascendere jussus, philosophicas theologicasque quæstiónes mirabíliter explanaverit, interim assiduis precibus sciéntiam Sanctórum et rectam vitæ morúmque normam a Dómino vehementíssime postulans. Quare, divino instinctu, Eremitárum sancti Augustíni amplexus est institutum.\par
\switchcolumn
\EN
\noindent Thomas of Villanova was born of excellent parentage, in the town of Fuenlana, in the Diocese of Toledo in Spain, in the year of our Lord 1488, and was early taught godliness, and a special pitifulness towards the needy. Of this grace he gave many examples while he was still a lad, whereof it is an eminent one that he more than once stripped himself of raiment of his own, in order to clothe the naked. He was become a man when the death of his father called him from Alcala, whither he had been sent to work as a student in the great College of St. Alonzo. He gave all the inheritance which fell to him to feed poor unmarried women, and forthwith returned to Alcala, and finished his course in Theology. He was so eminent in learning that he was commanded to take a Professorship in that University, and delivered remarkable Lectures upon Philosophy and Theology. Meanwhile he ceased not earnestly to entreat of the Lord in prayer the knowledge of the Saints, and to know what was the path of life whereunto he was called. In course of time, by the inspiration of God, he entered the Institute of Hermits of St. Augustine.\par
\switchcolumn*

\LA
\begin{Resp}\Rsign Invéni David servum meum, óleo sancto meo unxi eum: \Star{} Manus enim mea auxiliábitur ei.\end{Resp}
\begin{Resp}\Vsign Nihil profíciet inimícus in eo, et fílius iniquitátis non nocébit ei.\end{Resp}
\begin{Resp}\Rsign Manus enim mea auxiliábitur ei.\end{Resp}
\switchcolumn
\EN
\begin{Resp}\Rsign I have found David My servant, with My holy oil have I anointed him \Star{} For My hand shall help him.\end{Resp}
\begin{Resp}\Vsign The enemy shall prevail nothing against him, nor the son of wickedness afflict him.\end{Resp}
\begin{Resp}\Rsign For My hand shall help him.\end{Resp}
\switchcolumn*

\LA
\LessonHead{Lectio V}
\switchcolumn
\EN
\LessonHead{Lesson V}
\switchcolumn*

\LA
\noindent Religiónem proféssus, ómnibus religiosi hóminis virtútibus et ornamentis excelluit, humilitate, patiéntia, continéntia, sed ardentíssima caritate summe conspicuus, inter varios et assiduos labóres oratióni rerumque divinárum meditatióni invicto spíritu semper intentus. Prædicándi onus, útpote sanctimónia et doctrina præstans, subire jussus, cælésti aspirante grátia, innumerábiles e vitiórum cœno in viam salútis edúxit. Regéndis deínde frátribus admotus, prudéntiam, æquitátem et mansuetúdinem pari sedulitate ac severitáte conjunxit; adeo, ut priscam sui ordinis disciplínam multis in locis vel firmaverit vel restitúerit.\par
\switchcolumn
\EN
\noindent In the Order wherein he had professed, he was marked for all that maketh a good and edifying Friar, for lowliness, for long-suffering, for cleanness-of-heart, but, above all, for the warmth of his charity. Amid diverse and hard works, he let his spirit never faint from prayer and study of the things of God. On account of his holiness and learning he was bidden to undertake the work of preaching, and, by the help of God's grace, was the means of drawing countless souls out of the slough of sin into the way that leadeth unto life. Being raised to rule over his brethren, he so joined wisdom, justice, and gentleness with watchfulness and firmness, that he either established or restored in many places the original discipline of his Order.\par
\switchcolumn*

\LA
\begin{Resp}\Rsign Pósui adjutórium super poténtem, et exaltávi eléctum de plebe mea: \Star{} Manus enim mea auxiliábitur ei.\end{Resp}
\begin{Resp}\Vsign Invéni David servum meum, óleo sancto meo unxi eum.\end{Resp}
\begin{Resp}\Rsign Manus enim mea auxiliábitur ei.\end{Resp}
\switchcolumn
\EN
\begin{Resp}\Rsign I have laid help upon one that is mighty, and have exalted one chosen out of My people \Star{} For My hand shall help him.\end{Resp}
\begin{Resp}\Vsign I have found David My servant, with My holy oil have I anointed him.\end{Resp}
\begin{Resp}\Rsign For My hand shall help him.\end{Resp}
\switchcolumn*

\LA
\LessonHead{Lectio VI}
\switchcolumn
\EN
\LessonHead{Lesson VI}
\switchcolumn*

\LA
\noindent Granaténsis archiepíscopus designátus, mira humilitate et constántia insigne munus rejécit. Verum, non multo post, Valentinam ecclésiam, superiórum auctoritate coactus, gubernandam suscépit; quam annis ferme undecim ita rexit, ut sanctíssimi et vigilantíssimi pastoris partes expleverit. Ceterum, consueta vivéndi ratióne nihil admodum immutata, inexplébili caritáti multo magis indulsit, cum amplos ecclésiæ réditus in egénos dispérsit, ne lectulo quidem sibi relícto: nam eum, in quo decumbebat cum in cælum evocarétur, ab eodem commodátum hábuit, cui paulo ante eleemósynæ loco donaverat. Obdormívit in Dómino, sexto Idus Septembris, annos natus octo et sexagínta. Servi sui sanctitátem adhuc viventis, et exínde post mortem, miraculis Deus testatam vóluit; præsertim, cum horreum, fruménto paupéribus distributo, pénitus vacuum, repénte plenum invéntum est, et cum ad ejus sepúlcrum puer mórtuus revixit. Quibus aliisque non paucis fulgentem signis, Alexander septimus, Pontifex maximus, Sanctórum número adscripsit.\par
\switchcolumn
\EN
\noindent He was named to the Archbishopric of Granada, but, with excellent lowliness and firmness, he refused to take so high a place. However, not long after, he was forced by the commands of his superiors to accept the government of the Church of Valencia, which he discharged for nearly eleven years with the reputation of a most holy and watchful shepherd of souls. His elevation changed nowise his way of life, except to give greater scope to his wonderful charity by placing the revenues of a wealthy Church at his disposal to distribute to the poor. He did not leave himself even a bed that on which he was lying when he was called to heaven, he had only on loan from a person to whom he had shortly before given it as an alms. He fell asleep in the Lord upon the 8th day of September, in the 69th year of his own age, (and of our Lord 1555.) God was pleased to approve the holiness of His servant by miracles, both during his life and after his death, whereof are specially remarked that when he had utterly emptied his barn by giving away all his corn to the poor, it was suddenly found full again, and that a dead boy was raised to life at his grave. Finding him famous for these signs, and not a few others, Pope Alexander VII. enrolled him in the list of the Saints, and ordered that his memory should be held in remembrance upon the 18th day of September.\par
\switchcolumn*

\LA
\begin{Resp}\Rsign Iste est, qui ante Deum magnas virtútes operátus est, et omnis terra doctrína ejus repléta est: \Star{} Ipse intercédat pro peccátis ómnium populórum.\end{Resp}
\begin{Resp}\Vsign Iste est, qui contémpsit vitam mundi, et pervénit ad cæléstia regna.\end{Resp}
\begin{Resp}\Rsign Ipse intercédat pro peccátis ómnium populórum.\end{Resp}
\begin{Resp}\Vsign Glória Patri, et Fílio, \Star{} et Spirítui Sancto.\end{Resp}
\begin{Resp}\Rsign Ipse intercédat pro peccátis ómnium populórum.\end{Resp}
\switchcolumn
\EN
\begin{Resp}\Rsign This is he which wrought great wonders before God, and the whole earth is full of his teaching \Star{} May he pray for all people, that their sins may be forgiven unto them\end{Resp}
\begin{Resp}\Vsign This is he which loved not his life in this world, and hath attained unto the kingdom of heaven.\end{Resp}
\begin{Resp}\Rsign May he pray for all people, that their sins may be forgiven unto them\end{Resp}
\begin{Resp}\Vsign Glory be to the Father, and to the Son, \Star{} and to the Holy Ghost.\end{Resp}
\begin{Resp}\Rsign May he pray for all people, that their sins may be forgiven unto them\end{Resp}
\switchcolumn*

\end{paracol}

\par\needspace{10\baselineskip}\addvspace{1.2em}{\centering\small\textsc{Die 22 Septembris} · \textsc{22 September}\par}
\Office{Ss. Mauritii et Sociorum Martyrum}{Ss. Maurice and companions, Martyrs}{Simplex}
\addcontentsline{toc}{subsection}{Die 22 Septembris · Ss. Mauritii et Sociorum Martyrum \textit{— Ss. Maurice and companions, Martyrs}}
\begin{paracol}{2}
\LA
\LessonHead{Lectio IX (pro commemoratione)}
\switchcolumn
\EN
\LessonHead{Lesson IX (when commemorated)}
\switchcolumn*

\LA
\LessonTitle{Commemoratio: Ss. Mauritii et Sociorum Martyrum}
\switchcolumn
\EN
\LessonTitle{Commemoration of: Ss. Maurice and companions, Martyrs}
\switchcolumn*

\LA
\noindent Cum Maximiánus imperátor, ducto in Gálliam exercitu, in fínibus Sedunórum sacrifícii causa constitísset, Thebæórum legio, ne se impiórum sacrórum societáte contamináret, a réliquis cópiis discéssit. Quare imperátor misit ad eos mílites, qui suo nómine nuntiárent ut, si salvi esse vellent, in castra ad sacrifícia redírent: qui se christiána religióne impedíri respondérunt. Quorum respónsum ille indigníssime ferens, majori iracúndia quam antea exársit. Itaque, immíssa parte exércitus in Thebǽos, décimum quemque eórum primum occídi jussit; quod martýrium sua sponte, maximéque hortatóre Maurítio, ferre maluérunt quam imperáta fácere nefárii imperatóris. Ac deínde réliquos omnes, constantíssime Christum prædicántes, décimo Kaléndas Octóbris ab univérso exércitu trucidári ímperat.\par
\switchcolumn
\EN
\noindent When the Emperor Maximian led his army into Gaul, he stopped at the frontiers of the Seduni to offer a sacrifice. The Theban legion, that they might not be defiled by any share in the unhallowed rites, withdrew themselves from the rest of his army. Therefore the Emperor sent soldiers unto them to bid them in his name, if they valued their lives, come back into the camp to the sacrifice. They answered that the Christian religion did not allow them. This answer enkindled in him greater wrath than before. He therefore despatched a part of his army to the Thebans. with orders to begin by killing one man in every ten of them. By their own will, and at the urgent exhortation of Maurice, they chose rather to endure this martyrdom than to do the commandment of the unrighteous Emperor. At the last, the Emperor, upon the 22nd day of September, bade his whole army fall upon them and slay them all. They confessed Christ bravely even to the end.\par
\switchcolumn*

\LA
\TeDeum{Te Deum}
\switchcolumn
\EN
\TeDeum{Te Deum}
\switchcolumn*

\end{paracol}

\par\needspace{10\baselineskip}\addvspace{1.2em}{\centering\small\textsc{Die 23 Septembris} · \textsc{23 September}\par}
\Office{S. Lini Papæ et Martyris}{St. Linus, Pope and Martyr}{Semiduplex}
\addcontentsline{toc}{subsection}{Die 23 Septembris · S. Lini Papæ et Martyris \textit{— St. Linus, Pope and Martyr}}
\begin{paracol}{2}
\LA
\RefLine{Lectiones I–III: de Scriptura occurrente.}
\switchcolumn
\EN
\RefLine{Lessons I–III: from the Scripture of the day.}
\switchcolumn*

\LA
\RefLine{Lectiones VII–VIII: ex Commúni unius Summi Pontificis et Martyris (C2b).}
\switchcolumn
\EN
\RefLine{Lessons VII–VIII: from the Common of a Single Martyr Pope (C2b).}
\switchcolumn*

\LA
\RefLine{Lectio IX: de S. Theclæ Virginis et Martyris (09-23cc).}
\switchcolumn
\EN
\RefLine{Lesson IX: of St. Thecla, Virgin \& Martyr (09-23cc).}
\switchcolumn*

\LA
\LessonHead{Lectio IV}
\switchcolumn
\EN
\LessonHead{Lesson IV}
\switchcolumn*

\LA
\noindent Linus Pontifex, Volaterris in Etruria natus, primus post Petrum gubernávit Ecclésiam. Cujus tanta fides et sanctitas fuit, ut non solum dæmones ejiceret, sed étiam mórtuos revocaret ad vitam. Scripsit res gestas beáti Petri, et ea maxime quæ ab illo acta sunt contra Simónem magum. Sancívit ne qua mulier, nisi velato cápite, in ecclésiam introíret. Huic Pontifici caput amputátum est ob constantiam fidei, jussu Saturnini ímpii et ingratíssimi consularis, cujus filiam a dæmonum vexatióne liberaverat. Sepultus est in Vaticano prope sepúlcrum Principis Apostolórum, nono Kalendas Octobris. Sedit annos undecim, menses duos, dies viginti tres, creatis, bis mense Decembri, epíscopis quindecim, presbyteris decem et octo.\par
\switchcolumn
\EN
\noindent Pope Linus was by birth a native of Velletri in Tuscany, and was the first after Peter who governed the Church. His faith and holiness were such that he not only cast out devils, but also raised the dead. He wrote the acts of Blessed Peter, and especially the history of his strife with Simon Magus. He forbade women to enter the Church without having a veil upon their heads. His own head was cut off, on account of his firmness in confessing Christ, by command of the godless Consular Saturninus, an unthankful wretch whose own daughter he had delivered from being tormented by a devil. He was buried upon the Vatican Mount, hard by the grave of the Prince of the Apostles, upon the 23rd day of September. He sat as Pope eleven years, two months, and twenty - three days. He held two December ordinations, wherein he made fifteen Bishops, and eighteen Priests.\par
\switchcolumn*

\LA
\begin{Resp}\Rsign Honéstum fecit illum Dóminus, et custodívit eum ab inimícis, et a seductóribus tutávit illum: \Star{} Et dedit illi claritátem ætérnam.\end{Resp}
\begin{Resp}\Vsign Descendítque cum illo in fóveam, et in vínculis non derelíquit eum.\end{Resp}
\begin{Resp}\Rsign Et dedit illi claritátem ætérnam.\end{Resp}
\switchcolumn
\EN
\begin{Resp}\Rsign The Lord made him honourable, and defended him from his enemies, and kept him safe from those that lay in wait for him, \Star{} And gave him perpetual glory.\end{Resp}
\begin{Resp}\Vsign He went down with him into the pit, and left him not in bonds.\end{Resp}
\begin{Resp}\Rsign And gave him perpetual glory.\end{Resp}
\switchcolumn*

\LA
\LessonHead{Lectio V}
\switchcolumn
\EN
\LessonHead{Lesson V}
\switchcolumn*

\LA
\LessonTitle{De Expositióne sancti Ambrósii Epíscopi in Psalmum centésimum décimum octavum.}
\switchcolumn
\EN
\LessonTitle{From the Exposition of the hundred -and - eighteenth Psalm by St. Ambrose, Bishop \textasciitilde{}(of Milan)}
\switchcolumn*

\LA
\Source{Sermone 21}
\switchcolumn
\EN
\Source{1st Sermon}
\switchcolumn*

\LA
\noindent Príncipes persecúti sunt me gratis: et a verbis tuis trepidávit cor meum. Bene hoc Martyr dicit, quod injúste persecutiónum torménta sustineat; qui nihil rapúerit, nullum violentus oppresserit, nullius sánguinem fuderit, nullius torum putaverit esse violándum; qui nihil légibus debeat, et gravióra latronum sustinere cogátur supplícia; qui loquátur juste, et non audiátur; qui loquátur plena salútis, et impugnétur, ut possit dicere: Cum loquébar illis, impugnábant me gratis. Gratis ígitur persecutiónem pátitur, qui impugnátur sine crimine; impugnátur ut noxius cum sit in tali confessióne laudábilis; impugnátur quasi veneficus, qui in nómine Dómini gloriátur, cum pietas virtútum ómnium fundaméntum sit.\par
\switchcolumn
\EN
\noindent Princes have persecuted me without a cause; but my heart standeth in awe of thy word. These are rightly the words of a martyr, who beareth unjustly the torments of the persecutors, who hath robbed no man, who hath violently oppressed no man, who hath shed the blood of no man, who hath imagined to defile the bed of no man, who is debtor to the laws in nothing, and who is punished more grievously than if he were a robber who speaketh righteousness, and there is none that will hear who speaketh salvation, and all men fight against him who is able to say When I spoke unto them, they fought against me without a cause. (Ps. cxix. 7.) They fight against him without a cause, who can lay no sin to his charge; they fight against him as an evildoer, who is by their own acknowledgement righteous they fight against him as a warlock, who glorieth in the name of the Lord, and who doeth all things well because he doeth all things for God's sake.\par
\switchcolumn*

\LA
\begin{Resp}\Rsign Desidérium ánimæ ejus tribuísti ei Dómine, \Star{} Et voluntáte labiórum ejus non fraudásti eum.\end{Resp}
\begin{Resp}\Vsign Quóniam prævenísti eum in benedictiónibus dulcédinis: posuísti in cápite ejus corónam de lápide pretióso.\end{Resp}
\begin{Resp}\Rsign Et voluntáte labiórum ejus non fraudásti eum.\end{Resp}
\switchcolumn
\EN
\begin{Resp}\Rsign Lord, Thou hast given him his heart's desire. \Star{} And hast not withholden the request of his lips.\end{Resp}
\begin{Resp}\Vsign For Thou hast prevented him with the blessings of sweetness: Thou hast set a crown of precious stones upon his head.\end{Resp}
\begin{Resp}\Rsign And hast not withholden the request of his lips.\end{Resp}
\switchcolumn*

\LA
\LessonHead{Lectio VI}
\switchcolumn
\EN
\LessonHead{Lesson VI}
\switchcolumn*

\LA
\noindent Vere frustra impugnátur, qui apud ímpios et infidos impietátis arcessitur, cum fidei sit magister. Verum, qui gratis impugnátur, fortis debet esse et constans; quómodo ergo subtexuit: Et a verbis tuis trepidávit cor meum? Trepidare infirmitátis est, timoris atque formidinis. Sed est étiam infirmitas ad salútem, est étiam timor sanctórum: Timéte Dóminum, omnes sancti ejus; et, Beátus vir, qui timet Dóminum. Qua ratióne beatus? Quia in mandatis ejus cupit nimis.\par
\switchcolumn
\EN
\noindent They fight against him in vain who is accused of ungodliness among the ungodly and the unfaithful, because he teacheth Faith. Verily, him that is fought against without a cause it behoveth to be strong and patient. Wherefore then saith he My heart standeth in awe of thy word? Awe is the mark of the weak, the timid, and the fearful. But there is also a weakness unto salvation, there is a fear which is an holy fear. O fear the Lord, all ye His Saints. (Ps. xxxiii. 10.) And again: Blessed is the man that feareth the Lord. (Ps. cxi. i.) And wherefore is he blessed? because he delighteth greatly in His commandments.\par
\switchcolumn*

\LA
\begin{Resp}\Rsign Stola jucunditátis índuit eum Dóminus: \Star{} Et corónam pulchritúdinis pósuit super caput ejus.\end{Resp}
\begin{Resp}\Vsign Cibávit illum Dóminus pane vitæ et intelléctus: et aqua sapiéntiæ salutáris potávit illum.\end{Resp}
\begin{Resp}\Rsign Et corónam pulchritúdinis pósuit super caput ejus.\end{Resp}
\begin{Resp}\Vsign Glória Patri, et Fílio, \Star{} et Spirítui Sancto.\end{Resp}
\begin{Resp}\Rsign Et corónam pulchritúdinis pósuit super caput ejus.\end{Resp}
\switchcolumn
\EN
\begin{Resp}\Rsign The Lord hath put on him a robe of honour, \Star{} And put about his head a crown of joy.\end{Resp}
\begin{Resp}\Vsign With the bread of life and understanding hath the Lord fed him, and given him the water of health and wisdom to drink.\end{Resp}
\begin{Resp}\Rsign And put about his head a crown of joy.\end{Resp}
\begin{Resp}\Vsign Glory be to the Father, and to the Son, \Star{} and to the Holy Ghost.\end{Resp}
\begin{Resp}\Rsign And put about his head a crown of joy.\end{Resp}
\switchcolumn*

\end{paracol}

\par\needspace{10\baselineskip}\addvspace{1.2em}{\centering\small\textsc{Die 23 Septembris} · \textsc{23 September}\par}
\Office{S. Theclæ Virginis et Martyris}{St. Thecla, Virgin \& Martyr}{Simplex}
\addcontentsline{toc}{subsection}{Die 23 Septembris · S. Theclæ Virginis et Martyris \textit{— St. Thecla, Virgin \& Martyr}}
\begin{paracol}{2}
\LA
\LessonHead{Lectio IX (pro commemoratione)}
\switchcolumn
\EN
\LessonHead{Lesson IX (when commemorated)}
\switchcolumn*

\LA
\LessonTitle{Commemoratio: S. Theclæ Virginis et Martyris}
\switchcolumn
\EN
\LessonTitle{Commemoration of: St. Thecla, Virgin \& Martyr}
\switchcolumn*

\LA
\noindent Thecla virgo, ex illustribus paréntibus Iconii nata, a Paulo Apóstolo fidei præceptis instituta, miris sanctórum Patrum laudibus celebrátur. Quæ décimum octavum annum agens, Thamíride sponso relícto, cum eam paréntes, quod Christiana esset, accusassent, in ardentem rogum, qui, nisi Christo renuntiáret, ei parátus erat, prius signo crucis armata, seipsam injecit. Sed igne, plúvia quæ repénte exorta est, exstincto, Antiochíam venit; ubi, feris objecta et tauris in diversa incitátis alligata, mox conjecta in fossam plenam serpéntibus, ex ómnibus Jesu Christi grátia liberátur. Cujus ardore fidei et vitæ sanctitáte multi ad Christum convérsi sunt. Iterum in pátriam rédiens, in montem sola secessit. Deínde, multis virtútibus et miraculis insígnis, nonagenaria migrávit ad Dóminum, ac Seleucíæ sepúlta est.\par
\switchcolumn
\EN
\noindent This virgin Thecla was the daughter of noble parents at Iconium, and a disciple of the Apostle Paul. She is the subject of extraordinary praises by the holy Fathers. In the eighteenth year of her age, she parted from one Thamiris, to whom she had been betrothed, and her kindred accused her of being a Christian. A pile was set a-fire for her, unless she should deny Christ, but she made the sign of the Cross, and willingly entered it, and rain came, and put out the fire. She came to Antioch, where they threw her to wild beasts; and strove to tear her asunder, by tying her to oxen driven different ways; and cast her into a pit with many snakes; but by the mercy of Jesus Christ she was delivered from all. The warmth of her faith and the holiness of her life brought many to Christ. She returned into her own country, and withdrew to be an hermit, alone on a certain mountain, and passed away to be with the Lord, aged ninety years, and famous for many good works and miracles. She was buried at Seleucia.\par
\switchcolumn*

\LA
\TeDeum{Te Deum}
\switchcolumn
\EN
\TeDeum{Te Deum}
\switchcolumn*

\end{paracol}

\par\needspace{10\baselineskip}\addvspace{1.2em}{\centering\small\textsc{Die 24 Septembris} · \textsc{24 September}\par}
\Office{Beatæ Mariæ Virginis de Mercede}{Beatæ Mariæ Virginis de Mercede}{Duplex majus}
\addcontentsline{toc}{subsection}{Die 24 Septembris · Beatæ Mariæ Virginis de Mercede \textit{— Beatæ Mariæ Virginis de Mercede}}
\begin{paracol}{2}
\LA
\RefLine{Lectiones I–III: ex In Festis Beatae Mariae Virginis (C11).}
\switchcolumn
\EN
\RefLine{Lessons I–III: from the In Festis Beatae Mariae Virginis (C11).}
\switchcolumn*

\LA
\RefLine{Lectiones VII–IX: ex In Festis Beatae Mariae Virginis (C11).}
\switchcolumn
\EN
\RefLine{Lessons VII–IX: from the In Festis Beatae Mariae Virginis (C11).}
\switchcolumn*

\LA
\LessonHead{Lectio IV}
\switchcolumn
\EN
\LessonHead{Lesson IV}
\switchcolumn*

\LA
\noindent Quo témpore major feliciorque Hispaniárum pars diro Saracenórum opprimebátur jugo, innumerique fidéles sub immani servitute, maximo cum periculo christianæ fidei abjurandæ, amittendæque salútis ætérnæ, infeliciter detinebántur, beatíssima cælórum Regína, tot tantisque benigniter occurrens malis, nimiam caritátem suam in iis rediméndis osténdit. Nam sancto Petro Nolasco, pietáte et opibus florenti, qui, sanctis vacans meditatiónibus, jugiter animo recogitábat qua ratióne tot Christianórum ærumnis sub Maurórum captivitáte degéntium succurri posset, ípsamet beatíssima Virgo, serena fronte, se conspiciéndam dedit; et acceptíssimum sibi ac unigenito suo Fílio fore dixit, si suum in honórem institueréntur ordo religiosórum, quibus cura incumberet captivos e Turcárum tyrannide liberandi. Qua cælésti visióne vir Dei recreatus, mirum est quo caritátis ardore flagrare cœperit, hoc unum servans in corde suo, ut ipse ac instituénda ab eo religio maximam illam caritátem sedulo exercerent, ut quisque ánimam suam póneret pro amicis et próximis suis.\par
\switchcolumn
\EN
\noindent In the early part of the thirteenth century of the era of our Lord, the greatest and fairest part of Spain lay crushed under the yoke of the Saracens, and countless numbers of the faithful were held in brutal slavery, with the most lively danger of being made to deny the Christian faith and of losing everlasting salvation. Amid such sorrows the most Blessed Queen of heaven came mercifully to the rescue, and showed how the greatness of her motherly love was fain for their redemption. Holy Peter Nolasco, in the full bloom of the treasures of godliness as well as rich in earthly wealth, was earnestly pondering with himself how he could succour so many suffering Christians dwelling in bondage to the Moors. To him appeared with gracious visage the Most Blessed Virgin, and bade him know that it would be well pleasing in her own sight, and in the sight of her Only-begotten Son, that an Order of Religious men should be founded in her honour, whose work it should be to redeem prisoners from Mohammedan slavery. Strengthened by this heavenly vision, the man of God began to burn with wonderful charity, nursing in his heart the one desire that he himself and the Order which he should found might exercise that love, greater than which hath no man, that a man lay down his life for his friends. (John xv. 13.)\par
\switchcolumn*

\LA
\begin{Resp}\Rsign Sicut cedrus exaltáta sum in Líbano, et sicut cypréssus in monte Sion: quasi myrrha elécta, \Star{} Dedi suavitátem odóris.\end{Resp}
\begin{Resp}\Vsign Et sicut cinnamómum et bálsamum aromatízans.\end{Resp}
\begin{Resp}\Rsign Dedi suavitátem odóris.\end{Resp}
\switchcolumn
\EN
\begin{Resp}\Rsign I was exalted like a cedar in Lebanon, and as a cypress-tree upon Mount Zion. Like the best myrrh \Star{} I yielded a pleasant odour.\end{Resp}
\begin{Resp}\Vsign Like cinnamon and sweet balsam.\end{Resp}
\begin{Resp}\Rsign I yielded a pleasant odour.\end{Resp}
\switchcolumn*

\LA
\LessonHead{Lectio V}
\switchcolumn
\EN
\LessonHead{Lesson V}
\switchcolumn*

\LA
\noindent Ea ipsa nocte éadem Virgo sanctíssima beato Raymundo de Péñafort et Jacobo Aragoniæ regi appáruit, idípsum de religiosis instituéndis ádmonens, suadensque ut opem pro constructióne tanti operis ferrent. Petrus autem statim ad Raymundi pedes, qui ipsi erat a sacris confessiónibus, ádvolans, ei rem omnem apéruit; quem étiam cælitus instructum réperit, ejusque directióni se humillime subjécit. At, superveniens Jacobus rex, quam et ipse acceperat a beatíssima Vírgine, revelatiónem éxsequi státuit. Unde, collátis inter se consíliis, et consentiéntibus ánimis, in honórem ejusdem Vírginis Matris ordinem institúere aggressi sunt, sub invocatióne sanctæ Maríæ de Mercede redemptiónis captivórum.\par
\switchcolumn
\EN
\noindent Upon the same night the same most holy Virgin appeared to the Blessed Raymund de Penafuerte, and to James, King of Aragon, charging them concerning the founding of the Order, and desiring them to help in raising up so great a work. Peter betook himself forthwith to the feet of Raymund, who was his confessor, and laid the matter before him, whom also he found taught from heaven, and to whose governance he right humbly submitted himself. Then came King James, who appointed to carry out this revelation, which himself also had received from the Most Blessed Virgin. The three took counsel together, and all with one consent entered upon the institution of an Order in honour of the said Virgin Mother, to be placed under the invocation of St. Mary of Ransom, for the redemption of captives.\par
\switchcolumn*

\LA
\begin{Resp}\Rsign Quæ est ista quæ procéssit sicut sol, et formósa tamquam Jerúsalem? \Star{} Vidérunt eam fíliæ Sion, et beátam dixérunt, et regínæ laudavérunt eam.\end{Resp}
\begin{Resp}\Vsign Et sicut dies verni circúmdabant eam flores rosárum et lília convállium.\end{Resp}
\begin{Resp}\Rsign Vidérunt eam fíliæ Sion, et beátam dixérunt, et regínæ laudavérunt eam.\end{Resp}
\switchcolumn
\EN
\begin{Resp}\Rsign Who is this that cometh up like the sun? This, comely as Jerusalem? \Star{} The daughters of Zion saw her, and called her blessed: the queens also, and they praised her.\end{Resp}
\begin{Resp}\Vsign And about her it was as the flower of roses in the spring of the year, and lilies of the valleys.\end{Resp}
\begin{Resp}\Rsign The daughters of Zion saw her, and called her blessed the queens also, and they praised her.\end{Resp}
\switchcolumn*

\LA
\LessonHead{Lectio VI}
\switchcolumn
\EN
\LessonHead{Lesson VI}
\switchcolumn*

\LA
\noindent Die ígitur décima Augusti anno Dómini millesimo ducentésimo décimo octavo, rex idem Jacobus eam institutiónem, jamprídem ab iisdem sanctis viris conceptam, éxsequi státuit, sodálibus quarto voto adstrictis manéndi in pignus sub paganórum potestate, si pro Christianórum liberatióne opus fúerit. Quibus rex ipsa arma sua regia in péctore deferre concessit, et a Gregório nono illud tam præcellentis erga próximum caritátis institutum et religiónem confírmari curávit. Sed et ipse Deus per Vírginem Matrem increméntum dedit, ut talis institútio celerius ac felicius totum per orbem divulgarétur, sanctisque viris florúerit caritate ac pietáte insígnibus, qui eleemósynas a Christi fidelibus collectas in prétium redemptiónis suórum proximórum expénderent, seque ipsos intérdum darent in redemptiónem multórum. Ut autem tanti beneficii et institutiónis, débitæ Deo et Virgini Matri referántur gratiæ, Sedes apostolica hanc peculiarem festivitátem celebrari et Offícium recitari indulsit, cum alia fere innumera eídem ordini privilegia páriter contulísset.\par
\switchcolumn
\EN
\noindent Upon the 10th of August, in the year of our Lord 1218, the above-named King James decreed the establishment of this Order, thus already conceived by these holy men. The brethren take, (in addition to the vows of Poverty, Chastity, and Obedience,) a fourth vow, whereby they bind themselves to remain in pawn with the unbelievers, if need so require, for the liberation of Christians. The King granted them the right to bear on their breasts his own Royal blazon, and obtained from Gregory IX. the confirmation of this Institute and Order so nobly marked by brotherly charity. God Himself, through the Virgin Mother, gave the increase, causing this Institute speedily and prosperously to spread through all the world, and to blossom with holy men, great in love and godliness, to spend in the redemption of their neighbours the alms which are committed to them by Christ's faithful people, to that end, and some whiles to give themselves up for the ransom of many. That due thanks might be rendered to God and to the Virgin Mother for the great blessing of this Institute, the See Apostolic among other well-nigh countless favours bestowed upon it, permitted that this special Feast-day should be kept and this Office said.\par
\switchcolumn*

\LA
\begin{Resp}\Rsign Ornátam monílibus fíliam Jerúsalem Dóminus concupívit: \Star{} Et vidéntes eam fíliæ Sion, beatíssimam prædicavérunt, dicéntes: * Unguéntum effúsum nomen tuum.\end{Resp}
\begin{Resp}\Vsign Astitit regína a dextris tuis in vestítu deauráto, circúmdata varietáte.\end{Resp}
\begin{Resp}\Rsign Et vidéntes eam fíliæ Sion, beatíssimam prædicavérunt, dicéntes:\end{Resp}
\begin{Resp}\Vsign Glória Patri, et Fílio, \Star{} et Spirítui Sancto.\end{Resp}
\begin{Resp}\Rsign Unguéntum effúsum nomen tuum.\end{Resp}
\switchcolumn
\EN
\begin{Resp}\Rsign When the Lord beheld the daughter of Jerusalem adorned with her jewels, He greatly desired her beauty \Star{} And when the daughters of Zion saw her, they cried out that she was most blessed, saying thy name is as ointment poured forth.\end{Resp}
\begin{Resp}\Vsign Upon thy right hand did stand the Queen in a vesture of gold wrought about with diverse colours.\end{Resp}
\begin{Resp}\Rsign And when the daughters of Zion saw her, they cried out that she was most blessed, saying:\end{Resp}
\begin{Resp}\Vsign Glory be to the Father, and to the Son, \Star{} and to the Holy Ghost.\end{Resp}
\begin{Resp}\Rsign Thy name is as ointment poured forth.\end{Resp}
\switchcolumn*

\end{paracol}

\par\needspace{10\baselineskip}\addvspace{1.2em}{\centering\small\textsc{Die 26 Septembris} · \textsc{26 September}\par}
\Office{Ss. Cypriani et Justinæ Virginis, Martyrum}{Sts. Cyprian and Justina, Martyrs}{Simplex}
\addcontentsline{toc}{subsection}{Die 26 Septembris · Ss. Cypriani et Justinæ Virginis, Martyrum \textit{— Sts. Cyprian and Justina, Martyrs}}
\begin{paracol}{2}
\LA
\RefLine{Lectio I: de Scriptura occurrente.}
\switchcolumn
\EN
\RefLine{Lesson I: from the Scripture of the day.}
\switchcolumn*

\LA
\LessonHead{Lectio II}
\switchcolumn
\EN
\LessonHead{Lesson II}
\switchcolumn*

\LA
\noindent Cypriánus, primum magus, póstea Martyr, cum Justínam, christiánam Vírginem, quam júvenis quidam ardénter amábat, cantiónibus ac venefíciis ad ejus libídinis assénsum allícere conarétur, dæmonem consúluit quanam id re cónsequi posset. Cui dæmon respóndit nullam illi artem processúram advérsus eos, qui vere Christum cólerent. Quo respónso commótus Cypriánus, veheménter dolére cœpit vitæ superióris institútum. Itaque, relíctis mágicis ártibus, se totum ad Christi Dómini fidem convértit.\par
\switchcolumn
\EN
\noindent Cyprian was firstly a warlock and lastly a Martyr. A certain young man having a violent lust after a Christian maiden named Justina, employed him to excite her to join in this lewdness, by dint of incantations and philters. Cyprian thereupon asked counsel of the devil, how he might best gain that end. But the devil answered him that these arts are only thrown away upon true worshippers of Christ. This answer troubled Cyprian, and he began to repent heartily of the course of life he had hitherto led. And then he forsook his arts magic, and gave himself wholly up to the faith of the Lord Christ.\par
\switchcolumn*

\LA
\begin{Resp}\Rsign Viri sancti gloriósum sánguinem fudérunt pro Dómino, amavérunt Christum in vita sua, imitáti sunt eum in morte sua: \Star{} Et ídeo corónas triumpháles meruérunt.\end{Resp}
\begin{Resp}\Vsign Unus spíritus, et una fides erat in eis.\end{Resp}
\begin{Resp}\Rsign Et ídeo corónas triumpháles meruérunt.\end{Resp}
\begin{Resp}\Vsign Glória Patri, et Fílio, \Star{} et Spirítui Sancto.\end{Resp}
\begin{Resp}\Rsign Et ídeo corónas triumpháles meruérunt.\end{Resp}
\switchcolumn
\EN
\begin{Resp}\Rsign These men are holy, who have gloriously shed their blood for the Lord's sake, yea, who loved Christ in their lives, and were made like unto Him in their flesh, \Star{} And therefore they have earned crowns of victory.\end{Resp}
\begin{Resp}\Vsign One spirit, and one faith was in them.\end{Resp}
\begin{Resp}\Rsign And therefore they have earned crowns of victory.\end{Resp}
\begin{Resp}\Vsign Glory be to the Father, and to the Son, \Star{} and to the Holy Ghost.\end{Resp}
\begin{Resp}\Rsign And therefore they have earned crowns of victory.\end{Resp}
\switchcolumn*

\LA
\LessonHead{Lectio III (pro commemoratione)}
\switchcolumn
\EN
\LessonHead{Lesson III (when commemorated)}
\switchcolumn*

\LA
\noindent Cypriánus, primum magus, póstea Martyr, cum Justínam, christiánam Vírginem, quam júvenis quidam ardénter amábat, cantiónibus ac venefíciis ad ejus libídinis assénsum allícere conarétur, dæmonem consúluit quanam id re cónsequi posset. Cui dæmon respóndit nullam illi artem processúram advérsus eos, qui vere Christum cólerent. Quo respónso commótus Cypriánus, veheménter dolére cœpit vitæ superióris institútum. Itaque, relíctis mágicis ártibus, se totum ad Christi Dómini fidem convértit. Quam ob causam una cum Vírgine Justína comprehénsus est, et ambo cólaphis flagellísque cæsi sunt; mox in cárcerem conjécti, si forte senténtiam commutárent. Verum, inde póstea emíssi, cum in christiána religióne constantíssimi reperiréntur, in sartáginem plenam fervéntis picis, ádipis et ceræ injécti sunt. Demum Nicomedíæ secúri feriúntur. Quorum projécta córpora cum sex dies inhumáta jacuíssent noctu quidam nautæ clam ea in navem impósita Romam portavérunt. Ac primum in prædio Rufínæ nóbilis féminæ sepúlta sunt; póstea, transláta in Urbem, in basílica Constantiniána cóndita sunt prope baptistérium.\par
\switchcolumn
\EN
\noindent Cyprian was firstly a warlock and lastly a Martyr. A certain young man having a violent lust after a Christian maiden named Justina, employed him to excite her to join in this lewdness, by dint of incantations and philters. Cyprian thereupon asked counsel of the devil, how he might best gain that end. But the devil answered him that these arts are only thrown away upon true worshippers of Christ. This answer troubled Cyprian, and he began to repent heartily of the course of life he had hitherto led. And then he forsook his arts magic, and gave himself wholly up to the faith of the Lord Christ. For this cause, he and the Virgin Justina were arrested together, beaten with blows and scourging, and cast into prison, if haply they might change their mind. Being brought out of the prison, but still standing fast in their Christian religion, they were dipped in a vessel full of hot pitch, fat, and wax, and in the end beheaded, at Nicomedia, (on the 26th day of September, in the year of our Lord 304.) Their bodies were thrown out, and lay unburied for the space of six days, at the end of which time some sailors took them secretly by night on board a ship and carried them to Rome. They were first buried on the farm of the noble Lady Rufina, but afterwards brought into the city, where they lie hard by the Baptistery in the Church of (the Saviour, built by) Constantine.\par
\switchcolumn*

\LA
\TeDeum{Te Deum}
\switchcolumn
\EN
\TeDeum{Te Deum}
\switchcolumn*

\end{paracol}

\par\needspace{10\baselineskip}\addvspace{1.2em}{\centering\small\textsc{Die 27 Septembris} · \textsc{27 September}\par}
\Office{Ss. Cosmæ et Damiani Martyrum}{Sts. Cosmas and Damian, Martyrs}{Semiduplex}
\addcontentsline{toc}{subsection}{Die 27 Septembris · Ss. Cosmæ et Damiani Martyrum \textit{— Sts. Cosmas and Damian, Martyrs}}
\begin{paracol}{2}
\LA
\RefLine{Lectiones I–III: de Scriptura occurrente.}
\switchcolumn
\EN
\RefLine{Lessons I–III: from the Scripture of the day.}
\switchcolumn*

\LA
\LessonHead{Lectio IV}
\switchcolumn
\EN
\LessonHead{Lesson IV}
\switchcolumn*

\LA
\noindent Cosmas et Damiánus, fratres Arabes, in Ægéa urbe nati, nóbiles médici, imperatóribus Diocletiáno et Maximiáno, non magis medicínæ sciéntia quam Christi virtúte, morbis étiam insanabílibus medebántur. Quorum religiónem cum Lysias præféctus cognovísset, addúci eos ad se jubet, ac de vivéndi institúto et de fídei professióne interrogátos, cum se et Christiános esse, et christiánam fidem esse ad salútem necessáriam, líbere prædicárent, deos venerári ímperat; et, si id recúsent, minátur cruciátus et necem acerbíssimam.\par
\switchcolumn
\EN
\noindent Cosmas and Damian, who were eminent physicians in the time of the Emperors Diocletian and Maximian, were brothers, and Arabs by race, but born in the city of Aegae (in Cilicia.) Not more by their knowledge of medicine than by the power of Christ they healed diseases which had been hopeless for others. When the Praefect Lysias learnt to what faith they belonged, he commanded them to be brought before him, and questioned them as to their way of life, and the confession of their religion; and then, forasmuch as they freely owned themselves Christians and the Christian faith needful to salvation, he commanded them to worship the gods, under threats of torments and a most cruel death.\par
\switchcolumn*

\LA
\begin{Resp}\Rsign Sancti tui, Dómine, mirábile consecúti sunt iter, serviéntes præcéptis tuis, ut inveniréntur illǽsi in aquis válidis: \Star{} Terra appáruit árida: et in Mari Rubro via sine impediménto.\end{Resp}
\begin{Resp}\Vsign Quóniam percússit petram, et fluxérunt aquæ, et torréntes inundavérunt.\end{Resp}
\begin{Resp}\Rsign Terra appáruit árida: et in Mari Rubro via sine impediménto.\end{Resp}
\switchcolumn
\EN
\begin{Resp}\Rsign thy Saints, O Lord, have passed a wonderful way, serving thy commandments, that they might be found without hurt in the midst of the mighty waters. \Star{} Dry land appeared, and, out of the Red Sea, a way without impediment.\end{Resp}
\begin{Resp}\Vsign He smote the rock, and the waters gushed out, and the streams overflowed.\end{Resp}
\begin{Resp}\Rsign Dry land appeared, and, out of the Red Sea, a way without impediment.\end{Resp}
\switchcolumn*

\LA
\LessonHead{Lectio V}
\switchcolumn
\EN
\LessonHead{Lesson V}
\switchcolumn*

\LA
\noindent Verum, ut se frustra hæc illis propónere intélligit: Colligáte, inquit, manus et pedes istórum, eósque exquisítis torquéte supplíciis. Quibus jussa exsequéntibus, nihilóminus Cosmas et Damiánus in senténtia persistébant. Quare, ut erant vincti, in profúndum mare jaciúntur. Unde cum salvi ac solúti essent egréssi, mágicis ártibus præféctus factum assígnans, in cárcerem tradit, ac postrídie edúctos in ardéntem rogum ínjici jubet; ubi, cum ab ipsis flamma refúgeret, várie et crudéliter tortos secúri pércuti vóluit. Itaque, in Jesu Christi confessióne, martýrii palmam accepérunt.\par
\switchcolumn
\EN
\noindent He bond their hands and feet together, and put them to the sharpest of the question. And he was obeyed, but nevertheless Cosmas and Damian abode still of the same mind. Therefore they were cast into the depth of the sea, bound as they were, but they came forth again, whole and unbound. The Praefect, therefore, who would have it that it came to pass so by force of art magic, cast them into prison. On the morrow he haled them forth again, and bade cast them upon a great fire, but the flame turned away from them. He was pleased then to have them tormented in diverse and cruel sorts, and lastly, smitten with the axe. Thus did they bear witness for Christ Jesus even until they grasped the palm of their testimony.\par
\switchcolumn*

\LA
\begin{Resp}\Rsign Vérbera carníficum non timuérunt Sancti Dei, moriéntes pro Christi nómine: \Star{} Ut herédes fíerent in domo Dómini.\end{Resp}
\begin{Resp}\Vsign Tradidérunt córpora sua propter Deum ad supplícia.\end{Resp}
\begin{Resp}\Rsign Ut herédes fíerent in domo Dómini.\end{Resp}
\switchcolumn
\EN
\begin{Resp}\Rsign The Saints of God shrank not from the stripes of the executioners, but died for Christ's Name's sake \Star{} That they might be made joint-heirs in the house of the Lord.\end{Resp}
\begin{Resp}\Vsign They gave their bodies for God's sake to death.\end{Resp}
\begin{Resp}\Rsign That they might be made joint-heirs in the house of the Lord.\end{Resp}
\switchcolumn*

\LA
\LessonHead{Lectio VI}
\switchcolumn
\EN
\LessonHead{Lesson VI}
\switchcolumn*

\LA
\LessonTitle{Sermo sancti Augustíni Epíscopi}
\switchcolumn
\EN
\LessonTitle{From the Sermons of St. Augustine, Bishop (of Hippo.)}
\switchcolumn*

\LA
\Source{Serm. 47. de Sanctis.}
\switchcolumn
\EN
\Source{Serm. 47. of the Saints.}
\switchcolumn*

\LA
\noindent Quotiescúmque fratres caríssimi, sanctórum Mártyrum solémnia celebrámus, ita, ipsis intercedéntibus, exspectémus a Dómino cónsequi temporália benefícia, ut ipsos Mártyres imitándo accípere mereámur ætérna. Ab ipsis enim sanctórum Mártyrum in veritáte festivitátum gáudia celebrántur, qui ipsórum Mártyrum exémpla sequúntur. Solemnitátes enim Mártyrum exhortatiónes sunt martyriórum: ut imitári non pígeat, quod celebráre deléctat.\par
\switchcolumn
\EN
\noindent Dearly beloved brethren, as often as we keep the Feasts of the holy Martyrs, we look to obtain of the Lord, by their intercession, such good things in this life that thereby we, following them, may gain better in that which is to come. For they only do truly keep Holiday on the Feasts of the Martyrs, who follow after the Martyrs' example. These Feasts of the Martyrs are the Martyrs' preaching, whereby to stir us up to imitate what we are not loath to honour.\par
\switchcolumn*

\LA
\begin{Resp}\Rsign Tamquam aurum in fornáce probávit eléctos Dóminus, et quasi holocáusti hóstiam accépit illos: et in témpore erit respéctus illórum: \Star{} Quóniam donum et pax est eléctis Dei.\end{Resp}
\begin{Resp}\Vsign Qui confídunt in illum, intélligent veritátem, et fidéles in dilectióne acquiéscent illi.\end{Resp}
\begin{Resp}\Rsign Quóniam donum et pax est eléctis Dei.\end{Resp}
\begin{Resp}\Vsign Glória Patri, et Fílio, \Star{} et Spirítui Sancto.\end{Resp}
\begin{Resp}\Rsign Quóniam donum et pax est eléctis Dei.\end{Resp}
\switchcolumn
\EN
\begin{Resp}\Rsign As gold in the furnace hath the Lord tried His chosen ones, and received them as a burnt-offering, and yet a while, and they shall be regarded; \Star{} For the grace of God, and His peace, are with His chosen.\end{Resp}
\begin{Resp}\Vsign They that put their trust in Him shall understand the truth and such as be faithful in love shall abide with Him.\end{Resp}
\begin{Resp}\Rsign For the grace of God, and His peace, are with His chosen.\end{Resp}
\begin{Resp}\Vsign Glory be to the Father, and to the Son, \Star{} and to the Holy Ghost.\end{Resp}
\begin{Resp}\Rsign For the grace of God, and His peace, are with His chosen.\end{Resp}
\switchcolumn*

\LA
\LessonHead{Lectio VII}
\switchcolumn
\EN
\LessonHead{Lesson VII}
\switchcolumn*

\LA
\LessonTitle{Léctio sancti Evangélii secúndum Lucam}
\switchcolumn
\EN
\LessonTitle{From the Holy Gospel according to Luke}
\switchcolumn*

\LA
\Cite{Luc 6:17-23}
\switchcolumn
\EN
\Cite{Luke 6:17-23}
\switchcolumn*

\LA
\noindent In illo témpore: Descéndens Jesus de monte, stetit in loco campéstri, et turba discipulórum ejus, et multitúdo copiósa plebis ab omni Judǽa, et Jerúsalem, et marítima, et Tyri, et Sidónis. Et réliqua.\par
\switchcolumn
\EN
\noindent At that time, Jesus came down from the mountain, and stood in the plain, and the company of His disciples, and a great multitude of people out of all Judea, and Jerusalem, and from the sea coast of Tyre and Sidon. And so on.\par
\switchcolumn*

\LA
\LessonTitle{Homilía sancti Ambrósii Epíscopi}
\switchcolumn
\EN
\LessonTitle{Homily by St Ambrose, Bishop (of Milan.)}
\switchcolumn*

\LA
\Source{Lib. 5. in Lucam. Cap. 6., post init.}
\switchcolumn
\EN
\Source{Bk. v. on Luke vi.}
\switchcolumn*

\LA
\noindent Advérte ómnia diligénter, quómodo et cum Apóstolis ascéndat et descéndat ad turbas. Quómodo enim turba nisi in húmili Christum vidéret? Non séquitur ad excélsa, non ascéndit ad sublímia. Dénique ubi descéndit, invénit infírmos: in excélsis enim infírmi esse non possunt. Hinc étiam Matthǽus docet in inferióribus débiles esse sanátos. Prius enim unusquísque sanándus est, ut paulátim virtútibus procedéntibus ascéndere possit ad montem; et ídeo quemque in inferióribus sanat, hoc est, a libídine révocat, injúriam cæcitátis avértit. Ad vúlnera nostra descéndit: ut usu quodam et cópia suæ natúræ, compartícipes nos fáciat esse regni cæléstis.\par
\switchcolumn
\EN
\noindent Mark well how Jesus goeth upward with His disciples, and downward to the multitude. How should the multitude behold Christ, save in a lower place? Such go not up to the things which are above; such attain not to the things which are high. And when Jesus cometh down, He findeth such as are diseased for such like go not up to the heights. Hence also Matthew saith that there were there all sick people, (iv. 23.) Of these every man had need of healing, that, when he had received strength, by and by, he might go up into the mountain. And therefore, being Himself come down, He healeth them in the plain, that is to say, He calleth them away from their lust, and freeth them of their blindness. He cometh down to our wounds, to the end that by a certain use of His nature, and by the abundance thereof, He might make us joint-heirs of the kingdom of heaven.\par
\switchcolumn*

\LA
\begin{Resp}\Rsign Propter testaméntum Dómini, et leges patérnas, Sancti Dei perstitérunt in amóre fraternitátis: \Star{} Quia unus fuit semper spíritus in eis, et una fides.\end{Resp}
\begin{Resp}\Vsign Ecce quam bonum et quam jucúndum habitáre fratres in unum.\end{Resp}
\begin{Resp}\Rsign Quia unus fuit semper spíritus in eis, et una fides.\end{Resp}
\switchcolumn
\EN
\begin{Resp}\Rsign Because of the covenant of the Lord, and the laws of their fathers, the Saints of God abode in brotherly love, \Star{} For one spirit and one faith was ever in them.\end{Resp}
\begin{Resp}\Vsign Behold how good and how pleasant it is for brethren to dwell together in unity.\end{Resp}
\begin{Resp}\Rsign For one spirit and one faith was ever in them.\end{Resp}
\switchcolumn*

\LA
\LessonHead{Lectio VIII}
\switchcolumn
\EN
\LessonHead{Lesson VIII}
\switchcolumn*

\LA
\noindent Beáti páuperes, quia vestrum est regnum Dei. Quátuor tantum beatitúdines sanctus Lucas Domínicas posuit, octo vero sanctus Matthǽus: sed in illis octo istæ quátuor sunt, et in quátuor istis illæ octo. Hic enim quátuor velut virtútes ampléxus est cardináles: ille in illis octo mýsticum númerum reserávit. Pro octáva enim multi inscribúntur Psalmi: et mandátum áccipis octo illis partem dare fortásse benedictiónibus. Sicut enim spei nostræ octáva perféctio est, ita octáva summa virtútum est.\par
\switchcolumn
\EN
\noindent Blessed be ye poor, for your's is the kingdom of God. Saint Luke giveth us but four of the Lord's Beatitudes, and Saint Matthew eight but in those eight are contained these four, and in these four those eight. For in these four are embraced the cardinal virtues and in those eight they are set forth in a number full of mystery. It is written at the head of more than one of the Psalms that they are for the octave, and thou hast received the commandment Give a portion to seven, and also to eight to seven or eight what? Perchance degrees of blessedness. For as this eighth Beatitude doth name the most glorious realization of our hope the kingdom of Heaven so doth it also name the most royal exertion of our strength blessed are they which are persecuted.\par
\switchcolumn*

\LA
\begin{Resp}\Rsign Hæc est vera fratérnitas, quæ numquam pótuit violári certámine: qui effúso sánguine secúti sunt Dóminum: \Star{} Contemnéntes aulam régiam, pervenérunt ad regna cæléstia.\end{Resp}
\begin{Resp}\Vsign Ecce quam bonum et quam jucúndum habitáre fratres in unum.\end{Resp}
\begin{Resp}\Rsign Contemnéntes aulam régiam, pervenérunt ad regna cæléstia.\end{Resp}
\begin{Resp}\Vsign Glória Patri, et Fílio, \Star{} et Spirítui Sancto.\end{Resp}
\begin{Resp}\Rsign Contemnéntes aulam régiam, pervenérunt ad regna cæléstia.\end{Resp}
\switchcolumn
\EN
\begin{Resp}\Rsign Theirs is a brotherhood indeed, whose tie no storms availed to sever together they followed the Lord in the shedding of their blood. \Star{} Together they set at nought the Royal Palace; together they attained unto the kingdom of heaven.\end{Resp}
\begin{Resp}\Vsign Behold how good and how pleasant it is for brethren to dwell together in unity.\end{Resp}
\begin{Resp}\Rsign Together they set at nought the Royal Palace; together they attained unto the kingdom of heaven.\end{Resp}
\begin{Resp}\Vsign Glory be to the Father, and to the Son, \Star{} and to the Holy Ghost.\end{Resp}
\begin{Resp}\Rsign Together they set at nought the Royal Palace; together they attained unto the kingdom of heaven.\end{Resp}
\switchcolumn*

\LA
\LessonHead{Lectio IX}
\switchcolumn
\EN
\LessonHead{Lesson IX}
\switchcolumn*

\LA
\noindent Sed prius quæ sunt amplióra videámus. Beáti, inquit, páuperes, quóniam vestrum est regnum Dei. Primam benedictiónem hanc utérque Evangelísta pósuit. Órdine enim prima est, et parens quædam generatióque virtútum: quia qui contempsérit sæculária, ipse merébitur sempitérna: nec potest quisquam méritum regni cæléstis adipísci, qui mundi cupiditáte pressus, emergéndi non habet facultátem.\par
\switchcolumn
\EN
\noindent But let us first consider the fuller of the forms of these Beatitudes. Blessed be ye poor, for your's is the kingdom of God. Both of the Evangelists give to this Beatitude the first place. Yea, surely, for poorness, at least in spirit, is the first in order, the mother, and procreatrix of virtues; since he that setteth no store by temporal things, winneth toward eternal things; neither is any man able to gain the kingdom of heaven, on whom the love of this present world doth so press, that he cannot rid himself thereof.\par
\switchcolumn*

\LA
\TeDeum{Te Deum}
\switchcolumn
\EN
\TeDeum{Te Deum}
\switchcolumn*

\end{paracol}

\par\needspace{10\baselineskip}\addvspace{1.2em}{\centering\small\textsc{Die 28 Septembris} · \textsc{28 September}\par}
\Office{S. Wenceslai Ducis et Martyris}{St. Wenceslaus, Duke and Martyr}{Semiduplex}
\addcontentsline{toc}{subsection}{Die 28 Septembris · S. Wenceslai Ducis et Martyris \textit{— St. Wenceslaus, Duke and Martyr}}
\begin{paracol}{2}
\LA
\RefLine{Lectiones I–III: de Scriptura occurrente.}
\switchcolumn
\EN
\RefLine{Lessons I–III: from the Scripture of the day.}
\switchcolumn*

\LA
\LessonHead{Lectio IV}
\switchcolumn
\EN
\LessonHead{Lesson IV}
\switchcolumn*

\LA
\noindent Wencesláus Bohémiæ dux, Wratisláo patre christiáno, Drahomíra matre gentíli natus, ab ávia Ludmílla fémina sanctíssima pie educátus, omni virtútum génere insígnis, summo stúdio virginitátem per omnem vitam servávit illibátam. Mater, per nefáriam Ludmíllæ necem regni administratiónem assecúta, ímpie cum junióre fílio Bolesláo vivens, concitávit in se prócerum indignatiónem; quare, tyránnici et ímpii regíminis pertæsi, utriúsque excússo jugo, Wencesláum in urbe Pragénsi regem salutárunt.\par
\switchcolumn
\EN
\noindent Wenceslaus, Duke of Bohemia, was the son of a Christian father, Duke Wratislaus I., and an heathen mother named Drahomira. He had for his grandmother a most holy woman, named Ludmilla, who trained him up in godliness. He was a man eminent in all graces, and one who carefully held his virginity unsullied throughout the whole course of his life. His mother seized the supreme power by the foul murder of Ludmilla, and lived foully with her younger son Boleslaus, and the nobles roused thereby to indignation, and wearied with her tyranny and wicked government, cast off the yoke of both of them, and hailed Wenceslaus in the city of Prague as their King.\par
\switchcolumn*

\LA
\begin{Resp}\Rsign Honéstum fecit illum Dóminus, et custodívit eum ab inimícis, et a seductóribus tutávit illum: \Star{} Et dedit illi claritátem ætérnam.\end{Resp}
\begin{Resp}\Vsign Descendítque cum illo in fóveam, et in vínculis non derelíquit eum.\end{Resp}
\begin{Resp}\Rsign Et dedit illi claritátem ætérnam.\end{Resp}
\switchcolumn
\EN
\begin{Resp}\Rsign The Lord made him honourable, and defended him from his enemies, and kept him safe from those that lay in wait for him, \Star{} And gave him perpetual glory.\end{Resp}
\begin{Resp}\Vsign He went down with him into the pit, and left him not in bonds.\end{Resp}
\begin{Resp}\Rsign And gave him perpetual glory.\end{Resp}
\switchcolumn*

\LA
\LessonHead{Lectio V}
\switchcolumn
\EN
\LessonHead{Lesson V}
\switchcolumn*

\LA
\noindent Ille, regnum pietáte magis quam império regens, órphanis, víduis, egénis tanta caritáte subvénit, ut própriis húmeris aliquándo ligna indigéntibus noctu comportárit, paupéribus humándis frequénter affúerit, captívos liberárit, carcéribus deténtos nocte intempésta visitárit, pecúniis et consílio sæpíssime consolátus. Miti ánimo princeps veheménter dolébat quémpiam, etsi reum, morti adjudicári. Summa religióne sacerdótes venerátus, suis mánibus tríticum serébat et vinum exprimébat, quibus in Missæ sacrifício uteréntur. Nocte nudis pédibus super nivem et glaciem circuíbat ecclésias, sanguínea et terram calefaciéntia post se relínquens vestígia.\par
\switchcolumn
\EN
\noindent He ruled his kingdom by his virtues rather than by force. To the orphaned, the widowed, and the destitute he was very charitable, so that some whiles in the winter he carried firewood to the needy on his own shoulders. He helped oftentimes to bury the poor, he set captives free, and went many times to the prisons at the dead of night to comfort with money and advice them that were detained therein. To a Prince of so tender an heart it was a great grief to be behoven to condemn any to death, however guilty. For Priests he had a most earnest respect, and with his own hands sowed the corn and pressed the grapes for the bread and wine which they were to use for the Sacrifice. He would walk round the Church by night with bare feet upon the snow and ice, leaving behind him bloody footprints that warmed the ground.\par
\switchcolumn*

\LA
\begin{Resp}\Rsign Desidérium ánimæ ejus tribuísti ei Dómine, \Star{} Et voluntáte labiórum ejus non fraudásti eum.\end{Resp}
\begin{Resp}\Vsign Quóniam prævenísti eum in benedictiónibus dulcédinis: posuísti in cápite ejus corónam de lápide pretióso.\end{Resp}
\begin{Resp}\Rsign Et voluntáte labiórum ejus non fraudásti eum.\end{Resp}
\switchcolumn
\EN
\begin{Resp}\Rsign Lord, Thou hast given him his heart's desire. \Star{} And hast not withholden the request of his lips.\end{Resp}
\begin{Resp}\Vsign For Thou hast prevented him with the blessings of sweetness: Thou hast set a crown of precious stones upon his head.\end{Resp}
\begin{Resp}\Rsign And hast not withholden the request of his lips.\end{Resp}
\switchcolumn*

\LA
\LessonHead{Lectio VI}
\switchcolumn
\EN
\LessonHead{Lesson VI}
\switchcolumn*

\LA
\noindent Angelos hábuit sui córporis custódes. Cum enim ad singuláre certámen advérsus Radisláum, ducem Curiménsem, eo fine accéderet, ut suórum salúti prospíceret, visi sunt Angeli arma ministrásse et dixísse adversário: Ne férias. Pertérritus hostis, venerabúndus prócidens véniam exorávit. Cum in Germániam proféctus esset, imperátor, conspéctis duóbus Angelis áurea cruce ad se accedéntem ornántibus, e sólio prosíliens brácchiis excépit, régiis insígnibus decorávit, eíque sancti Viti brácchium donávit. Nihilóminus ímpius frater, matre hortánte, convívio excéptum et póstea in ecclésia orántem, parátæ sibi mortis præscium, adhíbitis scéleris comítibus, interfécit. Sanguis per paríetes aspérsus adhuc conspícitur, et, Deo víndice, matrem inhumánam terra absórbuit, interfectóres váriis modis mísere periérunt.\par
\switchcolumn
\EN
\noindent As his Body-guard he had angels. For when Radislaus, Prince of Gurinna, invaded Bohemia, and Wenceslaus, to save the effusion of his people's blood, went out to meet him in single combat, (two) angels were seen serving him with arms, and heard to say to the adversary Strike not. Therefore, his enemy was stricken with terror, fell down in reverence before him, and begged his forgiveness. When he went to Germany, the Emperor saw two angels carrying a golden Cross before him as he drew nigh him, and arose from his throne, embraced him in his arms, created him a King, and gifted him with the arm of the holy (Martyr) Vitus. Nevertheless, his godless brother, at the exhortation of their mother, bade him to a feast, (given on account of the birth of his son,) and when Wenceslaus, with a foreboding of the death prepared for him, went afterwards into the Church, and was praying there, (Boleslaus followed him thither,) together with some accomplices of his crime, and (when they had wounded him,) despatched him (with his own hand, running him through the body with a lance. He suffered a little after midnight, upon the 28th day of September, in the year of our Lord 938.) The stains of his blood may still be seen upon the walls. By the judgment of God, his unnatural mother was swallowed up by the earth, and his murderers, in diverse ways, perished miserably.\par
\switchcolumn*

\LA
\begin{Resp}\Rsign Stola jucunditátis índuit eum Dóminus: \Star{} Et corónam pulchritúdinis pósuit super caput ejus.\end{Resp}
\begin{Resp}\Vsign Cibávit illum Dóminus pane vitæ et intelléctus: et aqua sapiéntiæ salutáris potávit illum.\end{Resp}
\begin{Resp}\Rsign Et corónam pulchritúdinis pósuit super caput ejus.\end{Resp}
\begin{Resp}\Vsign Glória Patri, et Fílio, \Star{} et Spirítui Sancto.\end{Resp}
\begin{Resp}\Rsign Et corónam pulchritúdinis pósuit super caput ejus.\end{Resp}
\switchcolumn
\EN
\begin{Resp}\Rsign The Lord hath put on him a robe of honour, \Star{} And put about his head a crown of joy.\end{Resp}
\begin{Resp}\Vsign With the bread of life and understanding hath the Lord fed him, and given him the water of health and wisdom to drink.\end{Resp}
\begin{Resp}\Rsign And put about his head a crown of joy.\end{Resp}
\begin{Resp}\Vsign Glory be to the Father, and to the Son, \Star{} and to the Holy Ghost.\end{Resp}
\begin{Resp}\Rsign And put about his head a crown of joy.\end{Resp}
\switchcolumn*

\LA
\LessonHead{Lectio VII}
\switchcolumn
\EN
\LessonHead{Lesson VII}
\switchcolumn*

\LA
\LessonTitle{Léctio sancti Evangélii secúndum Matthǽum}
\switchcolumn
\EN
\LessonTitle{From the Holy Gospel according to Matthew}
\switchcolumn*

\LA
\Cite{Matt 10:34-42}
\switchcolumn
\EN
\Cite{Matt 10:34-42}
\switchcolumn*

\LA
\noindent In illo témpore: Dixit Jesus discípulis suis: Nolíte arbitrári quia pacem vénerim míttere in terram: non veni pacem míttere, sed gládium. Et réliqua.\par
\switchcolumn
\EN
\noindent At that time: Jesus said unto his disciples: Think not that I am come to send peace on earth: I came not to send peace, but a sword. And so on, and that which followeth.\par
\switchcolumn*

\LA
\LessonTitle{Homilía sancti Hilárii Epíscopi}
\switchcolumn
\EN
\LessonTitle{A Homily by St. Hilary the Bishop}
\switchcolumn*

\LA
\Source{Comment. in Matth., can. 10}
\switchcolumn
\EN
\Source{Comment. in Matth., can. 10}
\switchcolumn*

\LA
\noindent Quæ ista divísio est? inter prima enim legis præcépta accépimus: Honóra patrem tuum et matrem tuam; et ipse Dóminus ait: Pacem meam do vobis, pacem meam relínquo vobis. Quid sibi vult missus pótius gládius in terram, et separátus a patre fílius, et fília a matre, et nurus advérsus socrum, et hóminis doméstici ejus inimíci? Igitur exínde pública auctóritas impietáti proferétur. Ubíque ódia, ubíque bella, et gládius Dómini inter patrem et fílium, et inter filiam matrémque desæviens.\par
\switchcolumn
\EN
\noindent What meaneth such division in the earth, by reason of the Lord's coming? Amongst the foremost commandments of the Law we have received this: Honour thy father and thy mother. And the Lord himself said: Peace I leave with you, my peace I give unto you. What meaneth this sword sent upon the earth? Why is father set at variance against son, and daughter against mother, and daughter-in-law against mother-in-law? Why shall a man's foes be they of his own household? From all this it might seem that open sanction is given to the neglect of natural duties―hatred everywhere, warfare everywhere, the sword of the Lord furiously raging between father and son, mother and daughter.\par
\switchcolumn*

\LA
\begin{Resp}\Rsign Coróna áurea super caput ejus, \Star{} Expréssa signo sanctitátis, glória honóris, et opus fortitúdinis.\end{Resp}
\begin{Resp}\Vsign Quóniam prævenísti eum in benedictiónibus dulcédinis, posuísti in cápite ejus corónam de lápide pretióso.\end{Resp}
\begin{Resp}\Rsign Expréssa signo sanctitátis, glória honóris, et opus fortitúdinis.\end{Resp}
\switchcolumn
\EN
\begin{Resp}\Rsign The Lord set a crown of gold upon his head, \Star{} Wherein was engraved Holiness, an ornament of honour, a costly work, goodly and beautiful.\end{Resp}
\begin{Resp}\Vsign For thou hast prevented him with the blessings of goodness, and hast set a crown of pure gold upon his head.\end{Resp}
\begin{Resp}\Rsign Wherein was engraved Holiness, an ornament of honour, a costly work, goodly and beautiful.\end{Resp}
\switchcolumn*

\LA
\LessonHead{Lectio VIII}
\switchcolumn
\EN
\LessonHead{Lesson VIII}
\switchcolumn*

\LA
\noindent Gládius telórum ómnium telum acutíssimum est, in quo sit jus potestátis, et judícii sevéritas, et animadvérsio peccatórum. Et, hujus quidem teli nómine, novi Evangélii prædicatiónem appellátam frequens in prophétis auctóritas est. Dei ígitur verbum nuncupátum meminérimus in gládio; qui gládius missus in terram est, idest, prædicátio ejus hóminum córdibus infúsa. Fitque gravis in domo una dissénsio, et doméstica novo hómini erunt inimíca; quia ille, per verbum Dei divísus ab illis, manére et intérior et extérior, id est, et corpus et anima, in spíritus novitáte gaudébit.\par
\switchcolumn
\EN
\noindent The sharpest weapon of all weapons is the sword, whereby is enforced the law of might, the severity of judgment, the punishment of sinners. Furthermore, the word Sword is in common use amongst the Prophets as a figure of the preaching of the new Gospel. Ye are to understand, then, that the word of God is often called a sword; and that this sword is sent upon the earth when the preaching of God's word entereth into the hearts of men. For then at once cometh grave dissénsion in the new man's old home, and thus a man's foes are they of his own household. For he is set at variance against them by the word of God, whereby he joyously abideth in newness of spirit both within and without, that is, in soul and body.\par
\switchcolumn*

\LA
\begin{Resp}\Rsign Hic est vere Martyr, qui pro Christi nómine sánguinem suum fudit: \Star{} Qui minas júdicum non tímuit, nec terrénæ dignitátis glóriam quæsívit, sed ad cæléstia regna pervénit.\end{Resp}
\begin{Resp}\Vsign Justum dedúxit Dóminus per vias rectas, et osténdit illi regnum Dei.\end{Resp}
\begin{Resp}\Rsign Qui minas júdicum non tímuit, nec terrénæ dignitátis glóriam quæsívit, sed ad cæléstia regna pervénit.\end{Resp}
\begin{Resp}\Vsign Glória Patri, et Fílio, \Star{} et Spirítui Sancto.\end{Resp}
\begin{Resp}\Rsign Qui minas júdicum non tímuit, nec terrénæ dignitátis glóriam quæsívit, sed ad cæléstia regna pervénit.\end{Resp}
\switchcolumn
\EN
\begin{Resp}\Rsign This is indeed a Martyr who for the Name of Christ poured forth his life-blood; \Star{} Who feared not the judgment of the world, nor clung to any earthly honours; but set his heart on a heavenly reward.\end{Resp}
\begin{Resp}\Vsign The Lord guided the righteous in right paths, and shewed him the kingdom of God.\end{Resp}
\begin{Resp}\Rsign Who feared not the judgment of the world, nor clung to any earthly honours; but set his heart on a heavenly reward.\end{Resp}
\begin{Resp}\Vsign Glory be to the Father, and to the Son, \Star{} and to the Holy Ghost.\end{Resp}
\begin{Resp}\Rsign Who feared not the judgment of the world, nor clung to any earthly honours; but set his heart on a heavenly reward.\end{Resp}
\switchcolumn*

\LA
\LessonHead{Lectio IX}
\switchcolumn
\EN
\LessonHead{Lesson IX}
\switchcolumn*

\LA
\noindent Pergit deínde eódem præceptórum et intelligéntiæ decúrsu. Nam, posteáquam relinquénda ómnia, quæ in sæculo caríssima sunt, imperáverat, adjécit: Qui non áccipit crucem suam, et séquitur me, non est me dignus; quia Qui Christi sunt, crucifixérunt corpus cum vítiis et concupiscéntia. Et indígnus est Christo, qui non crucem suam, in qua compátimur, commórimur, consepelímur, conresúrgimus, accípiens, Dóminum sit secúrus, in hoc sacraménto fídei, spíritus novitáte victúrus.\par
\switchcolumn
\EN
\noindent The Lord continueth in the same strain of exhortation and discernment. For after bidding us leave whatever we hold dearest in the world, he addeth: He that taketh not up his cross, and followeth after me, is not worthy of me. They that are Christ's have crucified the flesh with the affections and lusts. And hence he is unworthy of Christ who taketh not up his cross, because thereon we suffer, and so thereon we die with the Lord; yea, thereby are we buried with him, and do rise again with him; thereby we follow the Lord to victory, in newness of spirit and in the mystery of faith.\par
\switchcolumn*

\LA
\TeDeum{Te Deum}
\switchcolumn
\EN
\TeDeum{Te Deum}
\switchcolumn*

\end{paracol}

\par\needspace{10\baselineskip}\addvspace{1.2em}{\centering\small\textsc{Die 29 Septembris} · \textsc{29 September}\par}
\Office{In Dedicatione S. Michaëlis Archangelis}{Dedication of the Archbasilica of St. Michael the Archangel}{Duplex I classis}
\addcontentsline{toc}{subsection}{Die 29 Septembris · In Dedicatione S. Michaëlis Archangelis \textit{— Dedication of the Archbasilica of St. Michael the Archangel}}
\begin{paracol}{2}
\LA
\LessonHead{Lectio I}
\switchcolumn
\EN
\LessonHead{Lesson I}
\switchcolumn*

\LA
\LessonTitle{De Daniéle Prophéta}
\switchcolumn
\EN
\LessonTitle{Lesson from the book of Daniel}
\switchcolumn*

\LA
\Cite{Dan 7:9-11}
\switchcolumn
\EN
\Cite{Dan 7:9-11}
\switchcolumn*

\LA
\noindent Aspiciébam donec throni pósiti sunt, et antíquus diérum sedit. Vestiméntum ejus cándidum quasi nix, et capílli cápitis ejus quasi lana munda, thronus ejus flammæ ignis, rotæ ejus ignis accénsus. Flúvius ígneus rapidúsque egrediebátur a fácie ejus; míllia míllium ministrábant ei, et décies míllies centéna míllia assistébant ei. Judícium sedit, et libri apérti sunt. Aspiciébam propter vocem sermónum grándium, quos cornu illud loquebátur; et vidi quóniam interfécta esset béstia, et perísset corpus ejus, et tráditum esset ad comburéndum igni.\par
\switchcolumn
\EN
\noindent I beheld till thrones were placed, and the Ancient of days sat: his garment was white as snow, and the hair of his head like clean wool: his throne like flames of fire: the wheels of it like a burning fire. A swift stream of fire issued forth from before him: thousands of thousands ministered to him, and ten thousand times a hundred thousand stood before him: the judgment sat, and the books were opened. I beheld because of the voice of the great words which that horn spoke: and I saw that the beast was slain, and the body thereof was destroyed, and given to the fire to be burnt:\par
\switchcolumn*

\LA
\begin{Resp}\Rsign Factum est siléntium in cælo, dum commítteret bellum draco cum Michaéle Archángelo: \Star{} Audíta est vox míllia míllium dicéntium: Salus, honor et virtus omnipoténti Deo. (Allelúja.)\end{Resp}
\begin{Resp}\Vsign Míllia míllium ministrábant ei, et décies centéna míllia assistébant ei.\end{Resp}
\begin{Resp}\Rsign Audíta est vox míllia míllium dicéntium: Salus, honor et virtus omnipoténti Deo. (Allelúja.)\end{Resp}
\switchcolumn
\EN
\begin{Resp}\Rsign There was silence in heaven while the dragon fought against Michael the Archangel. \Star{} I heard the voice of thousands of thousands, saying: Salvation, and honour, and power unto God the Almighty.\end{Resp}
\begin{Resp}\Vsign Thousands of thousands ministered unto Him, and ten thousand times hundreds of thousands stood before Him.\end{Resp}
\begin{Resp}\Rsign I heard the voice of thousands of thousands, saying: Salvation, and honour, and power unto God, the Almighty.\end{Resp}
\switchcolumn*

\LA
\LessonHead{Lectio II}
\switchcolumn
\EN
\LessonHead{Lesson II}
\switchcolumn*

\LA
\Cite{Dan 10:4-8}
\switchcolumn
\EN
\Cite{Dan 10:4-8}
\switchcolumn*

\LA
\noindent Die autem vigésima et quarta mensis primi, eram juxta flúvium magnum, qui est Tigris. Et levávi óculos meos, et vidi: et ecce vir unus vestítus líneis, et renes ejus accíncti auro obrýzo; et corpus ejus quasi chrysólithus, et fácies ejus velut spécies fúlguris, et óculi ejus ut lampas ardens, et brácchia ejus, et quæ deórsum sunt usque ad pedes, quasi spécies æris candéntis; et vox sermónum ejus ut vox multitúdinis. Vidi autem ego Dániel solus visiónem; porro viri qui erant mecum non vidérunt; sed terror nímius írruit super eos, et fugérunt in abscónditum. Ego autem, relíctus solus, vidi visiónem grandem hanc, et non remánsit in me fortitúdo, sed et spécies mea immutáta est in me, et emárcui nec hábui quidquam vírium.\par
\switchcolumn
\EN
\noindent And in the four and twentieth day of the first month I was by the great river which is the Tigris. And I lifted up my eyes, and I saw: and behold a man clothed in linen, and his loins were girded with the finest gold: And his body was like the chrysolite, and his face as the appearance of lightning, and his eyes as a burning lamp: and his arms, and all downward even to the feet, like in appearance to glittering brass: and the voice of his word like the voice of a multitude. And I Daniel alone saw the vision: for the men that were with me saw it not: but an exceeding great terror fell upon them, and they fled away, and hid themselves. And I being left alone saw this great vision: and there remained no strength in me, and the appearance of my countenance was changed in me, and I fainted away, and retained no strength.\par
\switchcolumn*

\LA
\begin{Resp}\Rsign Stetit Angelus juxta aram templi, habens thuríbulum áureum in manu sua, et data sunt ei incénsa multa: \Star{} Et ascéndit fumus arómatum de manu Angeli in conspéctu Dómini. (Allelúja.)\end{Resp}
\begin{Resp}\Vsign In conspéctu Angelórum psallam tibi: adorábo ad templum sanctum tuum, et confitébor nómini tuo, Dómine.\end{Resp}
\begin{Resp}\Rsign Et ascéndit fumus arómatum de manu Angeli in conspéctu Dómini. (Allelúja.)\end{Resp}
\switchcolumn
\EN
\begin{Resp}\Rsign An Angel stood at the Altar of the temple, having a golden censer in his hand and there was given unto him much incense; \Star{} And the smoke of the incense ascended up before the Lord, out of the Angel's hand.\end{Resp}
\begin{Resp}\Vsign Before the Angels will I sing praise unto thee; I will worship toward thy holy temple, and praise thy Name, O Lord.\end{Resp}
\begin{Resp}\Rsign And the smoke of the incense ascended up before the Lord, out of the Angel's hand.\end{Resp}
\switchcolumn*

\LA
\LessonHead{Lectio III}
\switchcolumn
\EN
\LessonHead{Lesson III}
\switchcolumn*

\LA
\Cite{Dan 10:9-14}
\switchcolumn
\EN
\Cite{Dan 10:9-14}
\switchcolumn*

\LA
\noindent Et audívi vocem sermónum ejus: et áudiens jacébam consternátus super fáciem meam, et vultus meus hærébat terræ. Et ecce manus tétigit me, et eréxit me super génua mea et super artículos mánuum meárum. Et dixit ad me: Dániel, vir desideriórum, intéllege verba quæ ego loquor ad te, et sta in gradu tuo; nunc enim sum missus ad te. Cumque dixísset mihi sermónem istum, steti tremens. Et ait ad me: Noli metúere, Dániel; quia, ex die primo quo posuísti cor tuum ad intellegéndum, ut te afflígeres in conspéctu Dei tui, exaudíta sunt verba tua, et ego veni propter sermónes tuos. Princeps autem regni Persárum réstitit mihi vigínti et uno diébus; et ecce Míchaël, unus de princípibus primis, venit in adjutórium meum, et ego remánsi ibi juxta regem Persárum. Veni autem ut docérem te quæ ventúra sunt pópulo tuo in novíssimis diébus, quóniam adhuc vísio in dies.\par
\switchcolumn
\EN
\noindent And I heard the voice of his words: and when I heard, I lay in a consternation, upon my face, and my face was close to the ground. And behold a hand touched me, and lifted me up upon my knees, and upon the joints of my hands. And he said to me: Daniel, thou man of desires, understand the words that I speak to thee, and stand upright: for I am sent now to thee. And when he had said this word to me, I stood trembling. And he said to me: Fear not, Daniel: for from the first day that thou didst set thy heart to understand, to afflict thyself in the sight of thy God, thy words have been heard: and I am come for thy words. But the prince of the kingdom of the Persians resisted me one and twenty days: and behold Michael, one of the chief princes, came to help me, and I remained there by the king of the Persians. But I am come to teach thee what things shall befall thy people in the latter days, for as yet the vision is for days.\par
\switchcolumn*

\LA
\begin{Resp}\Rsign In conspéctu Angelórum psallam tibi, et adorábo ad templum sanctum tuum: \Star{} Et confitébor nómini tuo, Dómine. (Allelúja.)\end{Resp}
\begin{Resp}\Vsign Super misericórdia tua et veritáte tua: quóniam magnificásti super nos nomen sanctum tuum.\end{Resp}
\begin{Resp}\Rsign Et confitébor nómini tuo, Dómine. (Allelúja.)\end{Resp}
\begin{Resp}\Vsign Glória Patri, et Fílio, \Star{} et Spirítui Sancto.\end{Resp}
\begin{Resp}\Rsign Et confitébor nómini tuo, Dómine. (Allelúja.)\end{Resp}
\switchcolumn
\EN
\begin{Resp}\Rsign Before the Angels will I sing praise unto thee, and will worship toward thy holy temple \Star{} And I will praise thy Name, O Lord.\end{Resp}
\begin{Resp}\Vsign For thy loving kindness, and for thy truth; for Thou hast glorified thine holy Name on us.\end{Resp}
\begin{Resp}\Rsign And I will praise thy Name, O Lord.\end{Resp}
\begin{Resp}\Vsign Glory be to the Father, and to the Son, \Star{} and to the Holy Ghost.\end{Resp}
\begin{Resp}\Rsign And I will praise thy Name, O Lord.\end{Resp}
\switchcolumn*

\LA
\LessonHead{Lectio IV}
\switchcolumn
\EN
\LessonHead{Lesson IV}
\switchcolumn*

\LA
\LessonTitle{Sermo sancti Gregórii Papæ}
\switchcolumn
\EN
\LessonTitle{From the Sermons of Pope St. Gregory the Great}
\switchcolumn*

\LA
\Source{Homil. 34 in Evang. ante medium}
\switchcolumn
\EN
\Source{34th on the Gospels}
\switchcolumn*

\LA
\noindent Novem Angelórum órdines dícimus, quia vidélicet esse, testánte sacro elóquio, scimus: Angelos, Archángelos, Virtútes, Potestátes, Principátus, Dominatiónes, Thronos, Chérubim atque Séraphim. Esse namque Angelos et Archángelos pene omnes sacri elóquii páginæ testántur. Chérubim vero atque Séraphim sæpe, ut notum est, libri prophetárum loquúntur. Quátuor quoque órdinum nómina Paulus Apóstolus ad Ephésios enúmerat, dicens: Supra omnem Principátum, et Potestátem, et Virtútem, et Dominatiónem. Qui rursus ad Colossénses scribens, ait: Sive Throni, sive Potestátes, sive Principátus, sive Dominatiónes. Dum ergo illis quátuor, quæ ad Ephésios dixit, conjungúntur Throni, quinque sunt órdines; quibus dum Angeli et Archángeli, Chérubim atque Séraphim adjúncta sunt, procul dúbio novem esse Angelórum órdines inveniúntur.\par
\switchcolumn
\EN
\noindent We say that there are nine Orders of Angels, for, by the witness of the holy Word, we know that there be Angels, Archangels, Mights, Powers, Principalities, Dominions, Thrones, Cherubim, and Seraphim. Nearly every page of the holy Word witnesseth that there be Angels and Archangels. The books of the Prophets, as is well known, do oftentimes make mention of Cherubim and Seraphim. Paul, writing to the Ephesians, (i. 21,) counteth up the names of four Orders, where he saith: “The Father of glory raised (Christ) from the dead, and set Him at His own right hand in the heavenly places, far above all Principality, and Power, and Might, and Dominion, (and every name that is named, not only in this world, but also in that which is to come.)” And the same, again, writing to the Colossians, (i. 16,) saith: “By (the Son) were all things created, that are in heaven, and that are in earth, visible and invisible, whether they be Thrones, or Dominions, or Principalities, or Powers; (all things were created by Him and for Him.)” If, then, we add the Thrones to the four Orders of which he spake unto the Ephesians, we have five Orders; and when we add unto them the Angels and the Archangels, the Cherubim and the Seraphim, we find that the Orders of Angels are beyond all doubt nine.\par
\switchcolumn*

\LA
\begin{Resp}\Rsign Hic est Míchaël Archángelus, princeps milítiæ Angelórum, \Star{} Cujus honor præstat benefícia populórum, et orátio perdúcit ad regna cælórum. (Allelúja.)\end{Resp}
\begin{Resp}\Vsign Archángelus Míchaël præpósitus paradísi, quem honoríficant Angelórum cives.\end{Resp}
\begin{Resp}\Rsign Cujus honor præstat benefícia populórum, et orátio perdúcit ad regna cælórum. (Allelúja.)\end{Resp}
\switchcolumn
\EN
\begin{Resp}\Rsign This is Michael, who to battle leads the armies of the skies \Star{} Whosoever on him calleth blessed within his wardship lies. His a prayer whose voice availing aids from earth toward heaven to rise.\end{Resp}
\begin{Resp}\Vsign The Archangel Michael is the Vice-Roy of Paradise, and the Angels that are the dwellers therein, do hold him in worship.\end{Resp}
\begin{Resp}\Rsign Whosoever on him calleth blessed within his wardship lies. His a prayer whose voice availing aids from earth toward heaven to rise.\end{Resp}
\switchcolumn*

\LA
\LessonHead{Lectio V}
\switchcolumn
\EN
\LessonHead{Lesson V}
\switchcolumn*

\LA
\noindent Sciéndum vero quod Angelórum vocábulum nomen est offícii, non natúræ. Nam sancti illi cæléstis pátriæ Spíritus, semper quidem sunt Spíritus, sed semper vocári Angeli nequáquam possunt; quia tunc solum sunt Angeli, cum per eos áliqua nuntiántur. Unde et per Psalmístam dícitur: Qui facit Angelos suos spíritus; ac si paténter dicat: Qui eos, quos semper habet Spíritus, étiam, cum volúerit, Angelos facit. Hi autem qui mínima núntiant, Angeli; qui vero summa annúntiant, Archángeli vocántur. Hinc est enim quod ad Maríam Vírginem non quílibet Angelus, sed Gábriel Archángelus míttitur. Ad hoc quippe ministérium, summum Angelum veníre dignum fúerat, qui summum ómnium nuntiábat. Qui idcírco étiam privátis nomínibus censéntur, ut signétur per vocábula, étiam in operatióne quid váleant. Míchaël namque, Quis ut Deus? Gábriel autem, Fortitúdo Dei; Ráphaël vero dícitur Medicína Dei.\par
\switchcolumn
\EN
\noindent But we must know that the word Angel is the designation, not of a nature, but of an office. Those holy spirits in the heavenly fatherland are always spirits, but they may no wise be always called Angels, (which is, being interpreted, messengers,) for they are Angels only when they are sent as Messengers. Hence also it is said by the Psalmist: (ciii. 5,) Who makest spirits thine Angels! As if it were: Of them who are always with Him as spirits, He doth somewhiles make use as Messengers. They who go on the lesser messages are called Angels they who go on the greater Archangels. Hence it is that unto the Virgin Mary was sent no common Angel, but the Archangel Gabriel. For the delivery of this, the highest message, it was meet that there should be sent the highest Angel. Their individual names also are so given as to signify the kind of ministry wherein each is powerful. Michael signifieth: Who-is-like-unto-God? Gabriel, the Strength-of-God, and Raphael, the Medicine-of-God.\par
\switchcolumn*

\LA
\begin{Resp}\Rsign Venit Míchaël Archángelus cum multitúdine Angelórum, cui trádidit Deus ánimas Sanctórum, \Star{} Ut perdúcat eas in paradísum exsultatiónis. (Allelúja.)\end{Resp}
\begin{Resp}\Vsign Emítte, Dómine, Spíritum sanctum tuum de cælis, spíritum sapiéntiæ et intelléctus.\end{Resp}
\begin{Resp}\Rsign Ut perdúcat eas in paradísum exsultatiónis. (Allelúja.)\end{Resp}
\switchcolumn
\EN
\begin{Resp}\Rsign Where Angels lead the spirits of the blessed dead the glad procession moves with Michael at its head \Star{} To lead them into the garden of Eden.\end{Resp}
\begin{Resp}\Vsign Lord, send thy Holy Spirit from heaven the Spirit of wisdom and understanding.\end{Resp}
\begin{Resp}\Rsign To lead them into the garden of Eden.\end{Resp}
\switchcolumn*

\LA
\LessonHead{Lectio VI}
\switchcolumn
\EN
\LessonHead{Lesson VI}
\switchcolumn*

\LA
\noindent Et quóties miræ virtútis áliquid ágitur, Míchaël mitti perhibétur; ut ex ipso actu et nómine detur intéllegi quia nullus potest fácere, quod fácere prǽvalet Deus. Unde et ille antíquus hostis, qui Deo esse per supérbiam símilis concupívit, dicens: In cælum conscéndam, super astra cæli exaltábo sólium meum, símilis ero Altíssimo; dum in fine mundi in sua virtúte relinquétur extrémo supplício periméndus, cum Michaéle Archángelo præliatúrus esse perhibétur, sicut per Joánnem dícitur: Factum est prǽlium cum Michaéle Archángelo. Ad Maríam quoque Gábriel míttitur, qui Dei Fortitúdo nominátur; illum quippe nuntiáre veniébat, qui ad debellándas aéreas potestátes húmilis apparére dignátus est. Ráphaël quoque interpretátur, ut díximus, Medicína Dei; quia vidélicet, dum Tobíæ óculos quasi per offícium curatiónis tétigit, cæcitátis ejus ténebras tersit.\par
\switchcolumn
\EN
\noindent As often as anything very mighty is to be done, we see that Michael is sent, that by that very thing, and by his name, we may remember that none is able to do as God doeth. Hence that old enemy whose pride hath puffed him up to be fain to be like unto God, even he who said, I will ascend unto heaven, I will exalt my throne above the stars of God. I will be like the Most High, (Isa. xiv. 13, 14,) this old enemy, when at the end of the world he is about to perish in the last death, having no strength but his own, is shown unto us a-fighting with Michael the Archangel, even as saith John, (Apoc. xii. 7): There was war in heaven Michael and his Angels fought against the dragon; and the dragon fought and his angels. Unto Mary is sent Gabriel, whose name is interpreted the Strength of God, for he came to herald the appearing of Him Who was content to appear lowly that He might fight down the powers of the air. Raphael, also, as we have said, signifieth the Medicine-of-God, and it is the name of him who touched as a physician the eyes of Tobias, and cleared away his blindness.\par
\switchcolumn*

\LA
\begin{Resp}\Rsign In témpore illo consúrget Míchaël, qui stat pro fíliis vestris: \Star{} Et véniet tempus, quale non fuit, ex quo gentes esse cœpérunt, usque ad illud. (Allelúja.)\end{Resp}
\begin{Resp}\Vsign In témpore illo salvábitur pópulus tuus omnis, qui invéntus fúerit scriptus in libro vitæ.\end{Resp}
\begin{Resp}\Rsign Et véniet tempus, quale non fuit, ex quo gentes esse cœpérunt, usque ad illud. (Allelúja.)\end{Resp}
\begin{Resp}\Vsign Glória Patri, et Fílio, \Star{} et Spirítui Sancto.\end{Resp}
\begin{Resp}\Rsign Et véniet tempus, quale non fuit, ex quo gentes esse cœpérunt, usque ad illud. (Allelúja.)\end{Resp}
\switchcolumn
\EN
\begin{Resp}\Rsign At that time shall Michael stand up, which standeth for your children. \Star{} And there shall be a time, such as never was since there was a nation even to that same time.\end{Resp}
\begin{Resp}\Vsign At that time thy people shall be delivered, every one that shall be found written in the Book of Life.\end{Resp}
\begin{Resp}\Rsign And there shall be a time, such as never was since there was a nation even to that same time.\end{Resp}
\begin{Resp}\Vsign Glory be to the Father, and to the Son, \Star{} and to the Holy Ghost.\end{Resp}
\begin{Resp}\Rsign And there shall be a time, such as never was since there was a nation even to that same time.\end{Resp}
\switchcolumn*

\LA
\LessonHead{Lectio VII}
\switchcolumn
\EN
\LessonHead{Lesson VII}
\switchcolumn*

\LA
\LessonTitle{Léctio sancti Evangélii secúndum Matthǽum}
\switchcolumn
\EN
\LessonTitle{From the Holy Gospel according to Matthew}
\switchcolumn*

\LA
\Cite{Matt 18:1-10}
\switchcolumn
\EN
\Cite{Matt 18:1-10}
\switchcolumn*

\LA
\noindent In illo témpore: Accessérunt discípuli ad Jesum, dicéntes: Quis, putas, major est in regno cælórum? Et réliqua.\par
\switchcolumn
\EN
\noindent At that time came the disciples unto Jesus, saying: Who is the greatest in the kingdom of heaven? And so on.\par
\switchcolumn*

\LA
\LessonTitle{Homilía sancti Hierónymi Presbýteri}
\switchcolumn
\EN
\LessonTitle{Homily by St. Jerome, Priest (at Bethlehem)}
\switchcolumn*

\LA
\Source{Lib. 3 Comm. in cap. 18 Matth.}
\switchcolumn
\EN
\Source{Bk. iii Comm. on Matth. xviii}
\switchcolumn*

\LA
\noindent Post invéntum statérem, post tribúta réddita, quid sibi vult Apostolórum repentína interrogátio, Quis, putas, major est in regno cælórum? Quia víderant pro Petro et Dómino idem tribútum rédditum, ex æqualitáte prétii arbitráti sunt Petrum ómnibus Apóstolis esse prælátum, qui in redditióne tribúti Dómino fúerat comparátus; ídeo intérrogant quis major sit in regno cælórum. Vidénsque Jesus cogitatiónes eórum, et causas erróris intéllegens, vult desidérium glóriæ, humilitátis contentióne, sanáre.\par
\switchcolumn
\EN
\noindent After the finding of the piece of money in the fish's mouth, after the payment of the tribute, what meaneth this sudden question of the Apostles? Who is the greatest in the kingdom of heaven? They had seen that the same tribute-money was paid for Peter as for the Lord, and from this equality of reckoning they gathered that Peter was Prince of all the Apostles, seeing that he had been appraised at the same price as his Master. Therefore they ask, Who is the greatest in the kingdom of heaven? And Jesus, seeing their thoughts, and understanding wherefore they erred, is fain to take away the desire of glory by the love of lowliness.\par
\switchcolumn*

\LA
\begin{Resp}\Rsign In conspéctu géntium nolíte timére; vos enim in córdibus vestris adoráte et timéte Dóminum; \Star{} Angelus enim ejus vobíscum est. (Allelúja.)\end{Resp}
\begin{Resp}\Vsign Stetit Angelus juxta aram templi, habens thuríbulum áureum in manu sua.\end{Resp}
\begin{Resp}\Rsign Angelus enim ejus vobíscum est. (Allelúja.)\end{Resp}
\switchcolumn
\EN
\begin{Resp}\Rsign Be not ye afraid before the Gentiles, but in your hearts, worship ye the Lord, and fear Him \Star{} For His Angel is with you.\end{Resp}
\begin{Resp}\Vsign An Angel stood at the Altar of the Temple, having a golden censer in his hand.\end{Resp}
\begin{Resp}\Rsign For His Angel is with you.\end{Resp}
\switchcolumn*

\LA
\LessonHead{Lectio VIII}
\switchcolumn
\EN
\LessonHead{Lesson VIII}
\switchcolumn*

\LA
\noindent Si autem manus tua vel pes tuus scandalízat te, abscínde eum et próice abs te. Necésse est quidem veníre scándala; væ tamen ei est hómini, qui, quod necésse est ut fiat in mundo, vítio suo facit ut per se fiat. Igitur omnis truncátur afféctus et univérsa propínquitas amputátur, ne per occasiónem pietátis unusquísque credéntium scándalis páteat. Si, inquit, ita est quis tibi conjúnctus, ut manus, pes, óculus; et est útilis atque sollícitus, et acútus ad perspiciéndum, scándalum autem tibi facit, et propter dissonántiam morum te pértrahit in gehénnam: mélius est ut et propinquitáte ejus et emoluméntis carnálibus cáreas, ne, dum vis lucri fácere cognátos et necessários, causam hábeas ruinárum.\par
\switchcolumn
\EN
\noindent Therefore, if thy hand or thy foot offend thee, cut them off, and cast them from thee for it must needs be that offences come, but woe to that man by whom the offence cometh! because by his sin he maketh, and maketh to be his own work, that which must needs be in the world. Away, then, with every affection and every kinship, lest thy love should throw a stumbling-block before a single believer. Be there any, saith He, who is as near to thee as thine hand, thy foot, or thine eye, useful, careful, far-seeing, but who layeth a stumbling-block before thee, and whose diverse way of life may draw thee to hell it is better for thee to lose such an one and thy worldly happiness with him, than to live surrounded by them that are near and needful to thee, and to pile up unto thyself damnation.\par
\switchcolumn*

\LA
\begin{Resp}\Rsign Míchaël Archángelus venit in adjutórium pópulo Dei, \Star{} Stetit in auxílium pro animábus justis. (Allelúja.)\end{Resp}
\begin{Resp}\Vsign Stetit Angelus juxta aram templi, habens thuríbulum áureum in manu sua.\end{Resp}
\begin{Resp}\Rsign Stetit in auxílium pro animábus justis. (Allelúja.)\end{Resp}
\begin{Resp}\Vsign Glória Patri, et Fílio, \Star{} et Spirítui Sancto.\end{Resp}
\begin{Resp}\Rsign Stetit in auxílium pro animábus justis. (Allelúja.)\end{Resp}
\switchcolumn
\EN
\begin{Resp}\Rsign The Archangel Michael came to help God's people. \Star{} He arose to succour the spirits of the righteous.\end{Resp}
\begin{Resp}\Vsign An Angel stood at the Altar of the Temple, having a golden censer in his hand.\end{Resp}
\begin{Resp}\Rsign He arose to succour the spirits of the righteous.\end{Resp}
\begin{Resp}\Vsign Glory be to the Father, and to the Son, \Star{} and to the Holy Ghost.\end{Resp}
\begin{Resp}\Rsign He arose to succour the spirits of the righteous.\end{Resp}
\switchcolumn*

\LA
\LessonHead{Lectio IX}
\switchcolumn
\EN
\LessonHead{Lesson IX}
\switchcolumn*

\LA
\noindent Dico vobis quia Angeli eórum in cælis semper vident fáciem Patris mei. Supra díxerat, per manum et pedem et óculum, omnes propinquitátes et necessitúdines, quæ scándalum fácere póterant, amputándas; austeritátem, ítaque senténtiæ subjécto præcépto témperat, dicens: Vidéte ne contemnátis unum ex pusíllis istis. Sic, inquit, præcípio severitátem, ut commiscéri cleméntiam dóceam. Quia Angeli eórum in cælis vident semper fáciem Patris. Magna dígnitas animárum, ut unaquǽque hábeat ab ortu nativitátis, in custódiam sui, Angelum delegátum. Unde légimus in Apocalýpsi Joánnis: Angelo Ephesi, et reliquárum ecclesiárum, scribe hæc. Apóstolus quoque prǽcipit velári cápita, in ecclésiis, feminárum propter Angelos.\par
\switchcolumn
\EN
\noindent I say unto you that in heaven their angels do always behold the face of My Father. Above, He had said that every tie of kinship or of convenience which might become a stumbling-block, albeit close and needful as hand, or foot, or eye, was to be cut off, but here He softeneth the hardness of that precept: Take heed that ye despise not one of these little ones. Hardness, saith He, I command not save as teaching tenderness withal; in heaven their Angels do always behold the face of My Father. Oh, how great is the dignity of souls, whereof every one hath from its birth an Angel appointed to guard it! Hence, we read in the Revelation of John unto the Angel of the Church of Ephesus, (and so of the others,) write, (ii. I) And the Apostle (Paul) also saith: The woman (that is, in the Church) ought to have a covering on her head, because of the Angels. (i Cor. xi. 10.)\par
\switchcolumn*

\LA
\TeDeum{Te Deum}
\switchcolumn
\EN
\TeDeum{Te Deum}
\switchcolumn*

\end{paracol}

\par\needspace{10\baselineskip}\addvspace{1.2em}{\centering\small\textsc{Die 30 Septembris} · \textsc{30 September}\par}
\Office{S. Hieronymi Presbyteris Confessoris et Ecclesiæ Doctoris}{St. Jerome, Priest, Confessor and Doctor of the Church}{Duplex}
\addcontentsline{toc}{subsection}{Die 30 Septembris · S. Hieronymi Presbyteris Confessoris et Ecclesiæ Doctoris \textit{— St. Jerome, Priest, Confessor and Doctor of the Church}}
\begin{paracol}{2}
\LA
\RefLine{Lectiones I–III: de Scriptura occurrente.}
\switchcolumn
\EN
\RefLine{Lessons I–III: from the Scripture of the day.}
\switchcolumn*

\LA
\LessonHead{Lectio IV}
\switchcolumn
\EN
\LessonHead{Lesson IV}
\switchcolumn*

\LA
\noindent Hieronymus, Eusebii fílius, Stridone in Dalmatia, Constantio imperatóre, natus, Romæ adoléscens est baptizatus, et in liberálibus disciplinis a Donato et aliis viris doctíssimis eruditus. Tum discéndi studio Gálliam peragrávit; ubi pios aliquot et in divinis litteris eruditos viros coluit, multosque sacros libros sua manu descripsit. Mox se in Græciam cónferens. philosophía et eloquéntia instructus, summórum theologórum consuetúdine flóruit. In primis vero Gregório Nazianzeno Constantinopoli operam dedit; quo doctore se sacras litteras didicisse profitétur. Tum religiónis causa visit Christi Dómini incunábula, totamque lustrávit Palæstinam; quam peregrinatiónem, adhíbitis Hebræórum eruditíssimis, ad sacræ Scripturæ intelligéntiam sibi multum profuísse testátur.\par
\switchcolumn
\EN
\noindent Jerome was the son of one Eusebius, and was born at Sorigni, (a small town upon the confines) of Dalmatia, in the reign of the Emperor Constantius. He was baptised at Rome when a lad, and studied there, under the instruction of Donatus and other very learned personages. He travelled in Gaul for the sake of improving his mind, and there sought the friendship of diverse godly men learned in the Scriptures, and made with his own hand many copies of the holy books. He afterwards betook himself to Greece, where he attained eminence as a philosopher and orator, in the following of the most famous theologians. At Constantinople, in especial, he sat at the feet of Gregory of Nazianzus, from whom he professeth himself to have learnt his theology. Then, for godliness' sake, he went to see the home of the Lord Christ, and so throughout all Palestine. He witnesseth that this pilgrimage, wherein he got the help of the most learned of the Jews for the understanding of the Holy Scriptures, did him much good.\par
\switchcolumn*

\LA
\begin{Resp}\Rsign Honéstum fecit illum Dóminus, et custodívit eum ab inimícis, et a seductóribus tutávit illum: \Star{} Et dedit illi claritátem ætérnam.\end{Resp}
\begin{Resp}\Vsign Justum dedúxit Dóminus per vias rectas, et osténdit illi regnum Dei.\end{Resp}
\begin{Resp}\Rsign Et dedit illi claritátem ætérnam.\end{Resp}
\switchcolumn
\EN
\begin{Resp}\Rsign The Lord made him honourable, and defended him from his enemies, and kept him safe from those that lay in wait for him; \Star{} And gave him perpetual glory.\end{Resp}
\begin{Resp}\Vsign He went down with him into the pit, and left him not in bonds.\end{Resp}
\begin{Resp}\Rsign And gave him perpetual glory.\end{Resp}
\switchcolumn*

\LA
\LessonHead{Lectio V}
\switchcolumn
\EN
\LessonHead{Lesson V}
\switchcolumn*

\LA
\noindent Deínde secessit in vastam Syriæ solitúdinem; ubi quadriennium in lectióne divinórum librórum cælestísque beatitúdinis contemplatióne consumpsit, assidua se abstinéntia, vi lacrimárum et corporis afflictatióne discrucians. Présbyter a Paulino episcopo Antiochíæ factus, Romam de controversiis quorúmdam episcopórum cum Paulino et Epiphanio ad Dámasum Pontificem profectus, ejus ecclesiásticis epistolis scribéndis adjútor fuit. Verum, cum prístinæ solitúdinis desidério tenerétur, in Palæstinam reversus, Bethlehem ad Christi Dómini præsepe, in monasterio quod a Paula Romana exstructum erat, cælestem quamdam vitæ ratiónem instituit; et, quamquam varie morbis doloribúsque tentarétur, tamen corporis incommoda piis labóribus et perpetua lectióne ac scriptióne superabat.\par
\switchcolumn
\EN
\noindent He withdrew himself into the wild deserts of Syria, where he passed four years in studying the Holy Scriptures and in considering the blessedness of heaven, afflicting his body by alway denying himself, by bitter tears, and by chastisement of the flesh. He was ordained Priest by Paulinus, Patriarch of Antioch. He went to Rome on account of the quarrelling of certain Bishops with Paulinus and Epiphanius, and there helped Pope Damasus in the writing of his letters upon Church affairs. But the longing for his old solitude came upon him, and he went back to Palestine, where, in the monastery at Bethlehem, built beside the cradle of the Lord Christ by the Lady Paula of Rome, he set himself to enter on earth upon the life of heaven, serving God in reading and writing without ceasing, regardless of the sufferings of a body tormented by diverse diseases and pains.\par
\switchcolumn*

\LA
\begin{Resp}\Rsign Amávit eum Dóminus, et ornávit eum: stolam glóriæ índuit eum, \Star{} Et ad portas paradísi coronávit eum.\end{Resp}
\begin{Resp}\Vsign Índuit eum Dóminus lorícam fídei, et ornávit eum.\end{Resp}
\begin{Resp}\Rsign Et ad portas paradísi coronávit eum.\end{Resp}
\switchcolumn
\EN
\begin{Resp}\Rsign The Lord loved him and beautified him; He clothed him with a robe of glory, \Star{} And crowned him at the gates of Paradise.\end{Resp}
\begin{Resp}\Vsign The Lord hath put on him the breast-plate of faith, and hath adorned him.\end{Resp}
\begin{Resp}\Rsign And crowned him at the gates of Paradise.\end{Resp}
\switchcolumn*

\LA
\LessonHead{Lectio VI}
\switchcolumn
\EN
\LessonHead{Lesson VI}
\switchcolumn*

\LA
\noindent Tamquam ad oraculum, ex ómnibus orbis terræ partibus, ad ipsum divinæ Scripturæ quæstiónes explicandæ referebántur. Illum Dámasus Pontifex, illum sanctus Augustinus de locis Scripturæ difficillimis sæpe consuluit, propter ejus singularem doctrinam, et linguæ non solum Latinæ et Græcæ, sed Hebráicæ étiam et Chaldáicæ, intelligéntiam; et quod omnes pene scriptores, ejusdem Augustíni testimonio, légerat. Hæreticos acerrimis scriptis exagitávit; piórum et catholicórum patrocinium semper suscépit. Vetus Testaméntum ex Hebræo convértit; novum, jussu Dámasi, Græcæ fidei réddidit, magna étiam ex parte explicávit. Multa præterea Latine réddidit scripta doctórum virórum, et ipse aliis proprii ingenii monuméntis christianam disciplínam illustrávit. Qui, ad summam senectútem perveniens, sanctitáte et doctrina illustris, Honorio imperatóre, migrávit in cælum. Cujus corpus, ad Bethlehem sepúltum, póstea Romam in basilicam sanctæ Maríæ ad Præsepe translátum est.\par
\switchcolumn
\EN
\noindent Hard questions upon the interpretation of the Holy Scripture were sent to him from all parts of the earth, as to an oracle. He was oftentimes consulted by Pope Damasus and by the holy Augustine upon the meaning of the most obscure passages of the Scripture, because of his extraordinary learning, and that he knew not the Latin and Greek tongues only, but also the Hebrew and Chaldee, and, as the same Augustine testifieth, had read nearly all writers. He attacked heretics with keen publications, and ever undertook the defence of the godly and Catholic. He translated the Old Testament from Hebrew into Latin, and, at the command of Damasus, reformed, according to the original Greek, the existing version of the New. Upon great part of the Scriptures he wrote commentaries. He translated likewise into Latin the works of many learned men, and himself contributed to the Christian life many monuments of his own wit. He lived to an extreme old age, and passed away to heaven, famous for learning and holiness, in the reign of the Emperor Honorius, (upon the 30th day of September, in the year of our Lord 420.) His body was buried at Bethlehem, but hath since been brought to Rome, where it lieth in the Church of St. Mary-at-the-Manger.\par
\switchcolumn*

\LA
\begin{Resp}\Rsign Iste homo perfécit ómnia quæ locútus est ei Deus, et dixit ad eum: Ingrédere in réquiem meam: \Star{} Quia te vidi justum coram me ex ómnibus géntibus.\end{Resp}
\begin{Resp}\Vsign Iste est, qui contémpsit vitam mundi, et pervénit ad cæléstia regna.\end{Resp}
\begin{Resp}\Rsign Quia te vidi justum coram me ex ómnibus géntibus.\end{Resp}
\begin{Resp}\Vsign Glória Patri, et Fílio, \Star{} et Spirítui Sancto.\end{Resp}
\begin{Resp}\Rsign Quia te vidi justum coram me ex ómnibus géntibus.\end{Resp}
\switchcolumn
\EN
\begin{Resp}\Rsign This is he which did according unto all that God commanded him; and God said unto him: Enter thou into My rest, \Star{} For thee have I seen righteous before Me among all people.\end{Resp}
\begin{Resp}\Vsign This is he which loved not his life in this world, and is come unto an everlasting kingdom.\end{Resp}
\begin{Resp}\Rsign For thee have I seen righteous before Me among all people.\end{Resp}
\begin{Resp}\Vsign Glory be to the Father, and to the Son, \Star{} and to the Holy Ghost.\end{Resp}
\begin{Resp}\Rsign For thee have I seen righteous before Me among all people.\end{Resp}
\switchcolumn*

\LA
\LessonHead{Lectio VII}
\switchcolumn
\EN
\LessonHead{Lesson VII}
\switchcolumn*

\LA
\LessonTitle{Léctio sancti Evangélii secúndum Matthǽum}
\switchcolumn
\EN
\LessonTitle{From the Holy Gospel according to Matthew}
\switchcolumn*

\LA
\Cite{Matt 5:13-19}
\switchcolumn
\EN
\Cite{Matt 5:13-19}
\switchcolumn*

\LA
\noindent In illo témpore: Dixit Jesus discípulis suis: Vos estis sal terræ. Quod si sal evanúerit, in quo saliétur? Et réliqua.\par
\switchcolumn
\EN
\noindent At that time Jesus said unto His disciples: Ye are the salt of the earth. But if the salt have lost his savour, wherewith shall it be salted? And so on.\par
\switchcolumn*

\LA
\LessonTitle{Homilía S. Hierónymi Presbýteri}
\switchcolumn
\EN
\LessonTitle{Homily by St. Jerome, Priest (at Bethlehem.)}
\switchcolumn*

\LA
\Source{Liber 1 Comment. in cap. 5 Matth.}
\switchcolumn
\EN
\Source{Book i. Comm. on Matth. v.}
\switchcolumn*

\LA
\noindent Sal appellántur Apóstoli et Doctores; quia per illos univérsum hóminum condítur genus. Quod si sal evanúerit, in quo saliétur? Si doctor erraverit, a quo alio doctore emendábitur? Ad níhilum valet ultra, nisi ut mittátur foras, et conculcétur ab homínibus. Exemplum de agricultura sumptum est. Sal étenim, sicut in cibórum condiméntum et ad siccandas carnes necessárium est, ita alium usum non habet. Certe légimus in Scripturis urbes quasdam, ira victórum, sale seminátas, ut nullum in ipsis germen orirétur.\par
\switchcolumn
\EN
\noindent Apostles and teachers are called salt, for it is by them that the whole mass of mankind is seasoned. But if the salt have lost his savour, wherewith shall it be salted? If the teacher have gone astray, by what other teacher is he to be corrected? It is thenceforth good for nothing, but to be cast out, and to be trodden under foot of men. This is a figure taken from farming. Salt is used to savour food withal, and to preserve meat, but it hath no other use. In sooth, we read in the Scriptures of some cities which were sown with salt in the fury of their conquerors, that no bud of life might ever spring there again.\par
\switchcolumn*

\LA
\begin{Resp}\Rsign Iste est qui ante Deum magnas virtútes operátus est, et de omni corde suo laudávit Dóminum: \Star{} Ipse intercédat pro peccátis ómnium populórum.\end{Resp}
\begin{Resp}\Vsign Ecce homo sine queréla, verus Dei cultor, ábstinens se ab omni ópere malo, et pérmanens in innocéntia sua.\end{Resp}
\begin{Resp}\Rsign Ipse intercédat pro peccátis ómnium populórum.\end{Resp}
\switchcolumn
\EN
\begin{Resp}\Rsign This is he which wrought great wonders before God, and praised the Lord with all his heart. \Star{} May he pray for all people, that their sins may be forgiven unto them.\end{Resp}
\begin{Resp}\Vsign Behold a man without blame, a worshipper of God in truth, keeping himself clean from every evil work, and abiding still in his innocency.\end{Resp}
\begin{Resp}\Rsign May he pray for all people, that their sins may be forgiven unto them.\end{Resp}
\switchcolumn*

\LA
\LessonHead{Lectio VIII}
\switchcolumn
\EN
\LessonHead{Lesson VIII}
\switchcolumn*

\LA
\noindent Caveant ergo doctores et epíscopi, et vídeant Poténtes poténter torménta sustinére; nihilque esse remedii, sed majórum ruínam ad tártarum ducere. Vos estis lux mundi. Non potest cívitas abscóndi supra montem pósita; neque accendunt lucérnam, et ponunt eam sub módio, sed super candelabrum, ut luceat ómnibus qui in domo sunt. Docet fiduciam prædicándi, ne Apóstoli abscondántur ob metum, et sint similes lucernæ sub módio; sed tota libertáte se prodant, ut, quod audiérunt in cubiculis, prædicent in tectis.\par
\switchcolumn
\EN
\noindent The teachers and Bishops, then, look well to it, seeing that mighty men shall be mightily tormented, (Wisd. vi. 7.) And there is no help for them, but they fall into hell with a greater crash. Ye are the light of the world. A city that is set on a hill cannot be hid. Neither do men light a candle, and put it under a bushel, but on a candlestick, that it may give light unto all that are in the house. Here He teacheth boldness in preaching, lest the Apostles should shrink away from fear, and be like unto candles under a bushel; but contrariwise should come forward with all freedom, and should proclaim upon the house-tops that which had been spoken in the ear in closets. (Luke xii. 3.)\par
\switchcolumn*

\LA
\begin{Resp}\Rsign In médio Ecclésiæ apéruit os ejus, \Star{} Et implévit eum Dóminus spíritu sapiéntiæ et intelléctus.\end{Resp}
\begin{Resp}\Vsign Jucunditátem et exsultatiónem thesaurizávit super eum.\end{Resp}
\begin{Resp}\Rsign Et implévit eum Dóminus spíritu sapiéntiæ et intelléctus.\end{Resp}
\begin{Resp}\Vsign Glória Patri, et Fílio, \Star{} et Spirítui Sancto.\end{Resp}
\begin{Resp}\Rsign Et implévit eum Dóminus spíritu sapiéntiæ et intelléctus.\end{Resp}
\switchcolumn
\EN
\begin{Resp}\Rsign In the midst of the congregation did the Lord open his mouth. \Star{} And filled him with the spirit of wisdom and understanding.\end{Resp}
\begin{Resp}\Vsign He made him rich with joy and gladness.\end{Resp}
\begin{Resp}\Rsign And filled him with the spirit of wisdom and understanding.\end{Resp}
\begin{Resp}\Vsign Glory be to the Father, and to the Son, \Star{} and to the Holy Ghost.\end{Resp}
\begin{Resp}\Rsign And filled him with the spirit of wisdom and understanding.\end{Resp}
\switchcolumn*

\LA
\LessonHead{Lectio IX}
\switchcolumn
\EN
\LessonHead{Lesson IX}
\switchcolumn*

\LA
\noindent Nolíte putare quóniam veni sólvere legem aut prophétas; non veni sólvere legem aut prophétas; non veni sólvere, sed adimplére. Sive quod de se per alios prophetáta compléverit, sive quod ea, quæ ante propter infirmitátem audiéntium rudia et imperfécta fuerant, sua prædicatióne compléverit, iram tollens et vicem talliónis excludens et occúltam in mente concupiscéntiam damnans. Donec tránseat cælum et terra. Promittúntur nobis cæli novi et terra nova, quæ facturus est Dóminus Deus. Si ergo nova creanda sunt, consequenter vetera transitura.\par
\switchcolumn
\EN
\noindent Think not that I am come to destroy the Law or the Prophets; I am not come to destroy, but to fulfill. The meaning is, either that He was come to fulfill those things which others had prophesied concerning Him, or that He was come to give the full measure of those things which had been spoken darkly and imperfectly on account of the weakness of their hearers, making away with anger, forbidding to take eye for eye and tooth for tooth, and condemning the secret lusting of the heart. Till heaven and earth pass, (one jot or one tittle shall in no wise pass from the Law, till all be fulfilled.) But there are promised unto us new heavens and a new earth, which the Lord God shall make. And if new things are to be created, old things must pass away.\par
\switchcolumn*

\LA
\TeDeum{Te Deum}
\switchcolumn
\EN
\TeDeum{Te Deum}
\switchcolumn*

\end{paracol}

\SeasonTitle{Mensis October}{October}

\addcontentsline{toc}{section}{Mensis October \textit{— October}}

\par\needspace{10\baselineskip}\addvspace{1.2em}{\centering\small\textsc{Die 1 Octobris} · \textsc{1 October}\par}
\Office{S. Remigii Episcopi et Confessoris}{St. Remigius, Bishop and Confessor}{Simplex}
\addcontentsline{toc}{subsection}{Die 1 Octobris · S. Remigii Episcopi et Confessoris \textit{— St. Remigius, Bishop and Confessor}}
\begin{paracol}{2}
\LA
\RefLine{Lectiones I–II: de Scriptura occurrente.}
\switchcolumn
\EN
\RefLine{Lessons I–II: from the Scripture of the day.}
\switchcolumn*

\LA
\LessonHead{Lectio III (pro commemoratione)}
\switchcolumn
\EN
\LessonHead{Lesson III (when commemorated)}
\switchcolumn*

\LA
\noindent Remigius, epíscopus Rheménsis, flóruit Clodovéo rege Francórum; quem étiam baptizávit, et, primus ómnium, doctrina et miraculis Francos ad Christi Dómini fidem perdúxit. Ejus oratióne, mortua puella revixit. Multos interpretátus est divinæ Scripturæ libros. Amplius septuagínta annos summa cum laude Rhemensem ecclésiam administrávit. Cujus vita et mortis sanctitátem multa, quæ consecuta sunt, miracula comprobarunt.\par
\switchcolumn
\EN
\noindent Remigius, Archbishop of Reims, flourished in the time of Chlodwig, King of the Franks, whom he baptised, and was the first who, by his preaching and miracles, brought the Franks to believe in the Lord Christ. At his prayers, a dead maiden was raised to life. He expounded many books of the Holy Scriptures. He ministered to the Church of Rheims with the utmost acceptance for above three score and ten years, and the holiness of his life and death were witnessed by many signs and wonders which befell afterwards.\par
\switchcolumn*

\LA
\TeDeum{Te Deum}
\switchcolumn
\EN
\TeDeum{Te Deum}
\switchcolumn*

\end{paracol}

\par\needspace{10\baselineskip}\addvspace{1.2em}{\centering\small\textsc{Die 2 Octobris} · \textsc{2 October}\par}
\Office{Ss. Angelorum Custodum}{Guardian Angels}{Duplex majus}
\addcontentsline{toc}{subsection}{Die 2 Octobris · Ss. Angelorum Custodum \textit{— Guardian Angels}}
\begin{paracol}{2}
\LA
\LessonHead{Lectio I}
\switchcolumn
\EN
\LessonHead{Lesson I}
\switchcolumn*

\LA
\LessonTitle{De libro Exodi}
\switchcolumn
\EN
\LessonTitle{From the Book of Exodus}
\switchcolumn*

\LA
\Cite{Exod 23:20-23}
\switchcolumn
\EN
\Cite{Exod 23:20-23}
\switchcolumn*

\LA
\noindent Ecce ego mittam Angelum meum, qui præcédat te et custódiat in via et introdúcat in locum quem parávi. Obsérva eum et audi vocem ejus, nec contemnéndum putes, quia non dimíttet, cum peccáveris, et est nomen meum in illo. Quod, si audíeris vocem ejus et féceris ómnia quæ loquor, inimícus ero inimícis tuis et afflígam affligéntes te. Præcedétque te Angelus meus.\par
\switchcolumn
\EN
\noindent Behold I will send my angel, who shall go before thee, and keep thee in thy journey, and bring thee into the place that I have prepared. Take notice of him, and hear his voice, and do not think him one to be contemned: for he will not forgive when thou hast sinned, and my name is in him. But if thou wilt hear his voice, and do all that I speak, I will be an enemy to thy enemies, and will afflict them that afflict thee. And my angel shall go before thee, and shall bring thee in unto the Amorrhite, and the Hethite, and the Pherezite, and the Chanaanite, and the Hevite, and the Jebusite, whom I will destroy.\par
\switchcolumn*

\LA
\begin{Resp}\Rsign Angelis suis Deus mandávit de te, ut custódiant te in ómnibus viis tuis: \Star{} In mánibus portábunt te, ne umquam offéndas ad lápidem pedem tuum.\end{Resp}
\begin{Resp}\Vsign Míllia míllium ministrábant ei, et décies míllies centéna míllia assistébant ei.\end{Resp}
\begin{Resp}\Rsign In mánibus portábunt te, ne umquam offéndas ad lápidem pedem tuum.\end{Resp}
\switchcolumn
\EN
\begin{Resp}\Rsign God hath given His Angels charge over thee, to keep thee in all thy ways. \Star{} They shall bear thee up in their hands, lest haply thou dash thy foot against a stone.\end{Resp}
\begin{Resp}\Vsign Thousands of thousands ministered unto Him, and ten thousand times hundreds of thousands stood before Him.\end{Resp}
\begin{Resp}\Rsign They shall bear thee up in their hands, lest haply thou dash thy foot against a stone.\end{Resp}
\switchcolumn*

\LA
\LessonHead{Lectio II}
\switchcolumn
\EN
\LessonHead{Lesson II}
\switchcolumn*

\LA
\LessonTitle{De Zacharía Prophéta}
\switchcolumn
\EN
\LessonTitle{From Zacharias the Prophet}
\switchcolumn*

\LA
\Cite{Zach 1:7-11}
\switchcolumn
\EN
\Cite{Zech 1:7-11}
\switchcolumn*

\LA
\noindent Factum est verbum Dómini ad Zacharíam, fílium Barachíæ fílii Addo, prophétam, dicens: Vidi per noctem, et ecce Vir ascéndens super equum rufum, et ipse stabat inter myrtéta, quæ erant in profúndo, et post eum equi rufi, várii et albi, et dixi: Quid sunt isti, Dómine mi? Et dixit ad me Angelus qui loquebátur in me; Ego osténdam tibi quid sint hæc. Et respóndit vir qui stabat inter myrtéta et dixit: Isti sunt quos misit Dóminus ut perámbulent terram. Et respondérunt Angelo Dómini, qui stabat inter myrtéta, et dixérunt: Perambulávimus terram, et ecce omnis terra habitátur et quiéscit.\par
\switchcolumn
\EN
\noindent In the four and twentieth day of the eleventh month which is called Sabath, in the second year of Darius, the word of the Lord came to Zacharias the son of Barachias, the son of Addo, the prophet, saying: I saw by night, and behold a man riding upon a red horse, and he stood among the myrtle trees, that were in the bottom: and behind him were horses, red, speckled, and white. And I said: What are these, my Lord? and the angel that spoke in me, said to me: I will shew thee what these are: And the man that stood among the myrtle trees answered, and said: These are they, whom the Lord hath sent to walk through the earth. And they answered the angel of the Lord, that stood among the myrtle trees, and said: We have walked through the earth, and behold all the earth is inhabited, and is at rest.\par
\switchcolumn*

\LA
\begin{Resp}\Rsign Respóndit Angelus Dómini et dixit: Dómine exercítuum, \Star{} Usquequo tu non miseréberis Jerúsalem et úrbium Juda, quibus irátus es?\end{Resp}
\begin{Resp}\Vsign Iste enim septuagésimus annus est.\end{Resp}
\begin{Resp}\Rsign Usquequo tu non miseréberis Jerúsalem et úrbium Juda, quibus irátus es?\end{Resp}
\switchcolumn
\EN
\begin{Resp}\Rsign Then the angel of the Lord answered, and said: O Lord of hosts, \Star{} How long wilt Thou not have mercy on Jerusalem, and on the cities of Judah, against which Thou hast had indignation?\end{Resp}
\begin{Resp}\Vsign These three score and ten years?\end{Resp}
\begin{Resp}\Rsign How long wilt Thou not have mercy on Jerusalem, and on the cities of Judah, against which Thou hast had indignation?\end{Resp}
\switchcolumn*

\LA
\LessonHead{Lectio III}
\switchcolumn
\EN
\LessonHead{Lesson III}
\switchcolumn*

\LA
\Cite{Zach 2:1-5}
\switchcolumn
\EN
\Cite{Zech 2:1-5}
\switchcolumn*

\LA
\noindent Et levávi óculos meos et vidi, et ecce vir, et in manu ejus funículus mensórum; et dixi: Quo tu vadis? Et dixit ad me: Ut métiar Jerúsalem et vídeam quanta sit latitúdo ejus, et quanta longitúdo ejus. Et ecce Angelus qui loquebátur in me, egrediebátur, et Angelus álius egrediebátur in occúrsum ejus et dixit ad eum: Curre, lóquere ad púerum istum dicens: Absque muro habitábitur Jerúsalem, præ multitúdine hóminum et jumentórum in médio ejus. Et ego ero ei, ait Dóminus, murus ignis in circúitu, et in glória ero in médio ejus.\par
\switchcolumn
\EN
\noindent And I lifted up my eyes, and saw, and behold a man, with a measuring line in his hand. And I said: Whither goest thou? and he said to me: To measure Jerusalem, and to see how great is the breadth thereof, and how great the length thereof. And behold the angel that spoke in me went forth, and another angel went out to meet him. And he said to him: Run, speak to this young man, saying: Jerusalem shall be inhabited without walls, by reason of the multitude of men, and of the beasts in the midst thereof. And I will be to it, saith the Lord, a wall of fire round about: and I will be in glory in the midst thereof.\par
\switchcolumn*

\LA
\begin{Resp}\Rsign In conspéctu géntium nolíte timére; vos enim in córdibus vestris adoráte et timéte Dóminum; \Star{} Angelus enim ejus vobíscum est.\end{Resp}
\begin{Resp}\Vsign Stetit Angelus juxta aram templi, habens thuríbulum áureum in manu sua.\end{Resp}
\begin{Resp}\Rsign Angelus enim ejus vobíscum est.\end{Resp}
\begin{Resp}\Vsign Glória Patri, et Fílio, \Star{} et Spirítui Sancto.\end{Resp}
\begin{Resp}\Rsign Angelus enim ejus vobíscum est.\end{Resp}
\switchcolumn
\EN
\begin{Resp}\Rsign When ye see the Gentiles, be not afraid of them, but in your hearts worship and fear the Lord; \Star{} For His Angel is with you.\end{Resp}
\begin{Resp}\Vsign An Angel stood at the Altar of the Temple, having a golden censer in his hand.\end{Resp}
\begin{Resp}\Rsign For His Angel is with you.\end{Resp}
\begin{Resp}\Vsign Glory be to the Father, and to the Son, \Star{} and to the Holy Ghost.\end{Resp}
\begin{Resp}\Rsign For His Angel is with you.\end{Resp}
\switchcolumn*

\LA
\LessonHead{Lectio IV}
\switchcolumn
\EN
\LessonHead{Lesson IV}
\switchcolumn*

\LA
\LessonTitle{Sermo sancti Bernárdi Abbátis}
\switchcolumn
\EN
\LessonTitle{From the Sermons of St. Bernard, Abbot (of Clairvaux.)}
\switchcolumn*

\LA
\Source{In Psalmum Qui habitat}
\switchcolumn
\EN
\Source{On Ps. xc}
\switchcolumn*

\LA
\noindent Angelis suis mandávit de te. Mira dignátio, et vere magna diléctio caritátis. Quis enim? quibus? de quo? quid mandávit? Studióse considerémus, fratres, diligénter commendémus memóriæ hoc tam grande mandátum. Quis enim mandávit? cujus sunt Angeli? cujus mandátis obtémperant? cujus obédiunt voluntáti? Nempe Angelis suis mandávit de te, ut custódiant te in ómnibus viis tuis. Nec cunctántur quin étiam in mánibus tollant te. Summa ergo Majéstas mandávit Angelis, et Angelis suis mandávit. Illis útique sublímibus, tam beátis quam próximis sibi cohæréntibus et vere domésticis mandávit de te. Tu quis es? Quid est homo, quod memor es ejus? aut fílius hóminis, quóniam réputas eum? Quasi vero non sit homo putrédo, et fílius hóminis vermis! Sed quid putas mandávit de te? Ut custódiant te.\par
\switchcolumn
\EN
\noindent He hath given His Angels charge over thee. A wonderful graciousness, and a wonderful outpouring of love. For who hath given charge? And what charge? Unto whom? And over whom? Let us carefully consider, my brethren, let us carefully hold in mind this great charge. For who hath given this charge? To Whom belong the Angels? Whose commandments do they obey, and Whose will do they do? He hath given His Angels charge over thee, to keep thee in all thy ways, and that not carelessly, for they shall bear thee up in their hands. The Highest Majesty, therefore, hath given charge unto Angels, even His Angels. Unto these beings so excellently exalted, so blessed, so near to Himself, even as His own household, unto these hath He given charge over thee. Who art thou? What is man, that Thou art mindful of him? or the son of man, that Thou visitest him? (Ps. viii. 5.) Even as though man were not rottenness, and the son of man, a worm. (Job. xxv. 6.) But what charge hath He given them over thee? To keep thee in all thy ways.\par
\switchcolumn*

\LA
\begin{Resp}\Rsign Vivit ipse Dóminus, \Star{} Quóniam custodívit me Angelus ejus et hinc eúntem et ibi commorántem et inde huc reverténtem.\end{Resp}
\begin{Resp}\Vsign Et non permísit me Dóminus ancíllam suam coinquinári.\end{Resp}
\begin{Resp}\Rsign Quóniam custodívit me Angelus ejus et hinc eúntem et ibi commorántem et inde huc reverténtem.\end{Resp}
\switchcolumn
\EN
\begin{Resp}\Rsign As the Lord liveth \Star{} His Angel hath kept me in my way that I went thither, and in my sojourning there, and in mine homecoming again hither.\end{Resp}
\begin{Resp}\Vsign And the Lord hath not suffered me that am His handmaid to be defiled.\end{Resp}
\begin{Resp}\Rsign His Angel hath kept me in my way that I went thither, and in my sojourning there, and in mine homecoming again hither.\end{Resp}
\switchcolumn*

\LA
\LessonHead{Lectio V}
\switchcolumn
\EN
\LessonHead{Lesson V}
\switchcolumn*

\LA
\noindent Quantam tibi debet hoc verbum inférre reveréntiam, afférre devotiónem, conférre fidúciam! Reveréntiam pro præséntia, devotiónem pro benevoléntia, fidúciam pro custódia. Caute ámbula, ut vidélicet cui adsunt Angeli, sicut eis mandátum est, in ómnibus viis tuis. In quovis diversório, in quovis ángulo, Angelo tuo reveréntiam habe. Tu ne áudeas illo præsénte, quod vidénte me non audéres. An præséntem esse dúbitas, quem non vides? Quid si audíres? quid si tángeres? quid si olfáceres? Vide, quia non solo visu, rerum præséntia comprobétur.\par
\switchcolumn
\EN
\noindent That respect, what thankfulness, what trust, ought this word to work in thee! Respect for their presence, thankfulness for their kindness, trust in their safe keeping. Walk carefully, as one with whom are Angels, as hath been laid in charge upon them, in all thy ways. In every lodging, in every nook, have reverence for thine Angel. Dare not to do in his presence what thou wouldst not dare to do in mine. Or dost thou doubt whether he be indeed present, because thou seest him not? What if thou heardest him? What if thou touchedst him? What if thou smelledst him? Behold, not by sight alone is the presence of things made manifest.\par
\switchcolumn*

\LA
\begin{Resp}\Rsign Angelus Dómini descéndit cum Azaría et sóciis ejus in fornácem, et excússit flammam ignis de fornáce; \Star{} Et non tétigit eos omníno ignis neque contristávit.\end{Resp}
\begin{Resp}\Vsign Benedíctus Deus eórum, qui misit Angelum suum, et éruit servos suos qui credidérunt in eum.\end{Resp}
\begin{Resp}\Rsign Et non tétigit eos omníno ignis neque contristávit.\end{Resp}
\switchcolumn
\EN
\begin{Resp}\Rsign The Angel of the Lord came down into the furnace together with Azariah and his fellows, and smote the flame of the fire out of the furnace; \Star{} So that the fire touched them not at all, neither hurt them.\end{Resp}
\begin{Resp}\Vsign Blessed be their God, Who sent His Angel, and delivered His servants that trusted in Him!\end{Resp}
\begin{Resp}\Rsign So that the fire touched them not at all, neither hurt them.\end{Resp}
\switchcolumn*

\LA
\LessonHead{Lectio VI}
\switchcolumn
\EN
\LessonHead{Lesson VI}
\switchcolumn*

\LA
\noindent In ipso ítaque, fratres, affectuóse diligámus Angelos ejus, tamquam futúros aliquándo coherédes nostros, ínterim vero actóres et tutóres, a Patre pósitos et præpósitos nobis. Quid sub tantis custódibus timeámus? Nec superári nec sedúci, minus autem sedúcere possunt, qui custódiunt nos in ómnibus viis nostris. Fidéles sunt, prudéntes sunt, poténtes sunt: quid trepidámus? Tantum sequámur eos, adhæreámus eis, et in protectióne Dei cæli commorémur. Quóties ergo gravíssima cérnitur urgére tentátio et tribulátio véhemens imminére, ínvoca custódem tuum, ductórem tuum, adjutórem tuum in opportunitátibus, in tribulatióne; incláma eum et dic: Dómine, salva nos, perímus.\par
\switchcolumn
\EN
\noindent Let us also, brethren, dearly love His Angels, as them with whom we are one day to be co-heirs, and who in the meanwhile are leaders and guardians set over us by the Father. With such guardians, whereof shall we be afraid? They that keep us in all our ways, can neither be conquered nor corrupted, far less can they corrupt. They are trusty, they are wary, they are mighty. Whereof shall we be afraid? Only let us follow them, only let us cleave unto them, and we shall abide under the shadow of the God of heaven. As often then as the gloom of temptation threateneth thee, or the sharpness of tribulation hangeth over thee, call upon Him That keepeth thee, thy Shepherd, thy Refuge in times of trouble, call upon Him, and say: “Lord, save us, we perish.”\par
\switchcolumn*

\LA
\begin{Resp}\Rsign In omni tribulatióne eórum non est tribulátus, \Star{} Et Angelus faciéi ejus salvávit eos.\end{Resp}
\begin{Resp}\Vsign In dilectióne sua et indulgéntia sua ipse redémit eos, et portávit eos, et elevávit eos cunctis diébus sæculi.\end{Resp}
\begin{Resp}\Rsign Et Angelus faciéi ejus salvávit eos.\end{Resp}
\begin{Resp}\Vsign Glória Patri, et Fílio, \Star{} et Spirítui Sancto.\end{Resp}
\begin{Resp}\Rsign Et Angelus faciéi ejus salvávit eos.\end{Resp}
\switchcolumn
\EN
\begin{Resp}\Rsign In all their affliction He was not afflicted; \Star{} And the Angel of His presence saved them.\end{Resp}
\begin{Resp}\Vsign In His love and in His pity He redeemed them; and He bore them and carried them all the days of old.\end{Resp}
\begin{Resp}\Rsign And the Angel of His presence saved them.\end{Resp}
\begin{Resp}\Vsign Glory be to the Father, and to the Son, \Star{} and to the Holy Ghost.\end{Resp}
\begin{Resp}\Rsign And the Angel of His presence saved them.\end{Resp}
\switchcolumn*

\LA
\LessonHead{Lectio VII}
\switchcolumn
\EN
\LessonHead{Lesson VII}
\switchcolumn*

\LA
\LessonTitle{Léctio sancti Evangélii secúndum Matthǽum}
\switchcolumn
\EN
\LessonTitle{From the Holy Gospel according to Matthew}
\switchcolumn*

\LA
\Cite{Matt 18:1-10}
\switchcolumn
\EN
\Cite{Matt 18:1-10}
\switchcolumn*

\LA
\noindent In illo témpore: Accessérunt discípuli ad Jesum, dicéntes: Quis, putas, major est in regno cælórum? Et réliqua.\par
\switchcolumn
\EN
\noindent At that time, the disciples came to Jesus, saying: Who thinkest thou is the greater in the kingdom of heaven? And so on.\par
\switchcolumn*

\LA
\LessonTitle{Homilía sancti Hilarii Epíscopi}
\switchcolumn
\EN
\LessonTitle{Homily by St. Hilary, Bishop (of Poitiers.)}
\switchcolumn*

\LA
\Source{Comment. in Matth. can. 18, post initium}
\switchcolumn
\EN
\Source{Comm. on Matth. xviii.}
\switchcolumn*

\LA
\noindent Nónnisi revérsos in natúram puerórum introíre regnum cælórum Dóminus docet: id est, per simplicitátem puerílem vítia córporum nostrórum animǽque revocánda. Púeros autem, credéntes omnes per audiéntiæ fidem nuncupávit. Hi enim patrem sequúntur, matrem amant, próximo velle malum nesciunt, curam opum négligunt; non insoléscunt, non odérunt, non mentiúntur, dictis credunt, et quod áudiunt, verum habent. Reverténdum ígitur est ad simplicitátem infántium; quia, in ea collocáti, spéciem humilitátis Domínicæ circumferémus.\par
\switchcolumn
\EN
\noindent Unless ye become as little children, saith the Lord, ye shall not enter into the kingdom of heaven, that is, unless by the uprooting of bodily and mental depravity, we bring our souls to the innocency of childhood. But He giveth the name of children to all such as believe by the hearing of faith. Children follow their father, love their mother, know not how to wish evil to their neighbours, are not careful for earthly riches; they insult not, they hate not, they lie not, they believe what they are told, and take for truth what they hear. Us then it behoveth to return to the simpleness of little children, for when we are well rooted therein, we shall so far bear about in ourselves an image of the sublime simpleness of the Lord Jesus.\par
\switchcolumn*

\LA
\begin{Resp}\Rsign Machabǽus et qui cum eo erant cognovérunt expugnári præsídia; \Star{} Cum fletu et lácrimis rogábant Dóminum et omnis turba simul, ut bonum Angelum mítteret ad salútem Israël.\end{Resp}
\begin{Resp}\Vsign Cum páriter prompto ánimo procéderent, Jerosólymis appáruit præcédens eos eques in veste cándida.\end{Resp}
\begin{Resp}\Rsign Cum fletu et lácrimis rogábant Dóminum et omnis turba simul, ut bonum Angelum mítteret ad salútem Israël.\end{Resp}
\switchcolumn
\EN
\begin{Resp}\Rsign When Maccabeus and they that were with him heard that (Lysias) besieged the holds, they and all the people with lamentation and tears besought the Lord that He would send; \Star{} A good Angel to deliver Israel.\end{Resp}
\begin{Resp}\Vsign So they went forth together with a willing mind, and, as they were at Jerusalem, there appeared before them on horseback one in white clothing.\end{Resp}
\begin{Resp}\Rsign A good Angel to deliver Israel.\end{Resp}
\switchcolumn*

\LA
\LessonHead{Lectio VIII}
\switchcolumn
\EN
\LessonHead{Lesson VIII}
\switchcolumn*

\LA
\noindent Væ huic mundo ab scándalis. Humílitas passiónis scándalum mundo est. In hoc enim máxime ignorántia detinétur humána, quod sub deformitáte crucis, ætérnæ glóriæ Dóminum nóluit accípere. Et quid mundo tam periculósum, quam non recepísse Christum? Ideo vere necésse esse ait veníre scándala; quia, ad sacraméntum reddéndæ nobis æternitátis, omnis in eo passiónis humílitas esset explénda.\par
\switchcolumn
\EN
\noindent Wee unto the world because of offences! The lowliness of the Passion is an offence unto the world. Such is the state of stupidity to which man's ignorance hath reduced itself, that it turneth away from the Lord of Eternal Glory, because of the unsightliness of the Cross! And what is so certain to bring woe unto the world as to turn away from Christ? And therefore He saith It must needs be that offences come, because his fulfilling the lowliness of the Passion was the predestined mean whereby He was to give us eternal life.\par
\switchcolumn*

\LA
\begin{Resp}\Rsign Tu, Dómine, qui misísti Angelum tuum sub Ezechía rege Juda et interfecísti de castris Sennácherib centum octogínta quinque míllia, \Star{} Et nunc, Dominátor cælórum, mitte Angelum tuum bonum ante nos, in timóre et tremóre magnitúdinis brácchii tui.\end{Resp}
\begin{Resp}\Vsign Ut métuant qui cum blasphémia veniunt advérsus sanctum pópulum tuum.\end{Resp}
\begin{Resp}\Rsign Et nunc, Dominátor cælórum, mitte Angelum tuum bonum ante nos, in timóre et tremóre magnitúdinis brácchii tui.\end{Resp}
\begin{Resp}\Vsign Glória Patri, et Fílio, \Star{} et Spirítui Sancto.\end{Resp}
\begin{Resp}\Rsign Et nunc, Dominátor cælórum, mitte Angelum tuum bonum ante nos, in timóre et tremóre magnitúdinis brácchii tui.\end{Resp}
\switchcolumn
\EN
\begin{Resp}\Rsign O Lord, Thou didst send thine Angel in the time of Hezekiah, King of Judah, and didst slay in the host of Sennacherib an hundred, fourscore, and five thousand. \Star{} Wherefore now also, O Lord of heaven, send thy good Angel before us, for a fear and dread of the might of thine arm.\end{Resp}
\begin{Resp}\Vsign That those be stricken with terror that come with blasphemy against thy holy people.\end{Resp}
\begin{Resp}\Rsign Now also, O Lord of heaven, send thy good Angel before us, for a fear and dread of the might of thine arm.\end{Resp}
\begin{Resp}\Vsign Glory be to the Father, and to the Son, \Star{} and to the Holy Ghost.\end{Resp}
\begin{Resp}\Rsign Now also, O Lord of heaven, send thy good Angel before us, for a fear and dread of the might of thine arm.\end{Resp}
\switchcolumn*

\LA
\LessonHead{Lectio IX}
\switchcolumn
\EN
\LessonHead{Lesson IX}
\switchcolumn*

\LA
\noindent Vidéte ne contemnátis unum de pusíllis istis, qui credunt in me. Aptíssimus vínculum mútui amóris impósuit, ad eos præcípue qui vere in Dómino credidíssent. Pusillórum enim Angeli quotídie Deum vident: quia Fílius hóminis venit salváre quæ pérdita sunt. Ergo et Fílius hóminis salvat, et Deum Angeli vident, et Angeli pusillórum præsunt fidélium oratiónibus. Præésse Angelos absolúta auctóritas est. Salvatórum ígitur per Christum oratiónes Angeli quotídie Deo ófferunt. Ergo periculóse ille contémnitur, cujus desidéria ac postulatiónes ad ætérnum et invisíbilem Deum, ambitióso Angelórum famulátu ac ministério, pervehúntur.\par
\switchcolumn
\EN
\noindent Take heed that ye despise not one of these little ones that believe in Me. He hath laid on us a most meet tie to constrain us to love one another, especially such as indeed believe in the Lord. For I say unto you that in heaven their Angels do always behold the face of My Father Which is in heaven. For the Son of Man is come to save that which was lost. From these words we see, first, that the Son of Man saveth; secondly, that the Angels do see God; and thirdly, that the Angels of these little ones have the wardship over the prayers of the faithful. That the Angels have this wardship is taught us absolutely. The Angels therefore do every day offer to God the prayers, which they which are saved, do make to Him in the Name of Christ. Therefore it is dangerous for a man to despise them, seeing that these are they by whose watchful service and ministry his wishes and requests are presented before the throne of the eternal and unseen God.\par
\switchcolumn*

\LA
\TeDeum{Te Deum}
\switchcolumn
\EN
\TeDeum{Te Deum}
\switchcolumn*

\end{paracol}

\par\needspace{10\baselineskip}\addvspace{1.2em}{\centering\small\textsc{Die 3 Octobris} · \textsc{3 October}\par}
\Office{S. Theresiæ a Jesu Infante Virginis}{St. Thérèse of the Child Jesus, Virgin}{Duplex}
\addcontentsline{toc}{subsection}{Die 3 Octobris · S. Theresiæ a Jesu Infante Virginis \textit{— St. Thérèse of the Child Jesus, Virgin}}
\begin{paracol}{2}
\LA
\RefLine{Lectiones I–III: de Scriptura occurrente.}
\switchcolumn
\EN
\RefLine{Lessons I–III: from the Scripture of the day.}
\switchcolumn*

\LA
\LessonHead{Lectio IV}
\switchcolumn
\EN
\LessonHead{Lesson IV}
\switchcolumn*

\LA
\noindent Terésia a Jesu Infante, Alensonii in Gállia, honestis paréntibus, singulari et assidua erga Deum pietáte conspicuis, orta est. Inde a prima ætate, divino Spíritu prævénta, religiosam vitam ágere cupiebat. Serio autem promisit, se nihil Deo denegatúram, quod ipse ab ea petere viderétur: quam promissiónem fideliter usque ad mortem servare satégit. Quinto ætátis anno, matre amissa, Dei providéntiæ se totam commisit sub vigilanti custódia amantíssimi patris, sororúmque natu majórum: quibus magistris, Teresia ad curréndam perfectiónis viam ut gigas exsultávit. Novennis virgínibus ex Ordine sancti Benedicti Lexoviis excolenda traditur, ibique in rerum divinárum cognitióne excellere visa est. Décimo ætátis anno, arcanus et gravis morbus eam diu cruciávit, a quo prout ipsa enárrat, ope Beatíssimæ Vírginis, quæ eídem subrídens appáruit, et quam, sub titulo Dominæ nostræ a Victoria, per novendialia invocare studuit, divinitus fuit liberata. Tunc, angelico fervore repléta, ad sacrum convivium, in quo Christus sumitur, se diligentíssime præparáre curávit.\par
\switchcolumn
\EN
\noindent Theresa of the Child Jesus was born in Alencon in France. Her parents were estimable people, well known for their piety and their love of God. From her earliest childhood, endowed by a special grace of the Holy Ghost, she yearned to enter the religious life. She promised God with the utmost sincerity that she would deny him nothing he might ask of her. She kept this promise faithfully to the end of her life, although she had to suffer a great deal to keep it. Her mother died when Theresa was but five years old. From then on the child committed herself to the providence of God, under the vigilant care of a most tender father and her elder sisters. Under their teaching Theresa raced as gayly strong as a young giant along the way of perfection. At the age of nine she was sent to school at Lisieux to the Benedictine nuns, where she made remarkable progress in her knowledge of divine things. In her tenth year she was ill for a long time of a serious and mysterious malady. From this, as she herself relateth, she was delivered only by the power of God himself, through the intercession of the Blessed Virgin Mary, who appeared to her with a smiling countenance, and to whom under the title of our Lady of Victories, she was constantly making novenas. Filled with angelic fervour she prepared herself at this time with the utmost care to receive Christ in the sacred banquet of her first Holy Communion.\par
\switchcolumn*

\LA
\begin{Resp}\Rsign Propter veritátem, et mansuetúdinem, et justítiam: \Star{} Et dedúcet te mirabíliter déxtera tua.\end{Resp}
\begin{Resp}\Vsign Spécie tua et pulchritúdine tua inténde, próspere procéde, et regna.\end{Resp}
\begin{Resp}\Rsign Et dedúcet te mirabíliter déxtera tua.\end{Resp}
\switchcolumn
\EN
\begin{Resp}\Rsign Because of truth, and meekness, and righteousness; \Star{} And thy right hand shall lead thee wonderfully.\end{Resp}
\begin{Resp}\Vsign In thy comeliness, and thy beauty, go forward, fare prosperously, and reign.\end{Resp}
\begin{Resp}\Rsign And thy right hand shall lead thee wonderfully.\end{Resp}
\switchcolumn*

\LA
\LessonHead{Lectio V}
\switchcolumn
\EN
\LessonHead{Lesson V}
\switchcolumn*

\LA
\noindent Ut prímitus eucharistico pane fuit refecta, insatiábilem cæléstis hujus cibi famem haurire visa est: unde velut inspirata, Jesum rogabat, ut omnem mundanam consolatiónem in amaritúdinem sibi verteret. Inde tenerrimo in Christum Dóminum et in Ecclésiam amóre exæstuans, nihil antiquius hábuit, quam Carmelitárum Excalceatórum Ordinem íngredi, ut sui abnegatióne, suisque sacrifíciis, sacerdótibus, missionariis, totique Ecclésiæ opem afferret, et innumeras ánimas Jesu Christo lucrifáceret: quod jam morti próxima, apud Deum se factúram pollícita est. Propter ætátis defectum, multas ad religiosam vitam amplecténdam nacta est difficultates, quibus tamen incredibili animi fortitúdine superatis, quindecim annos nata, Lexoviensem Carmelum feliciter ingressa est. Ibi mirábiles Deus in Teresiæ corde ascensiónes dispósuit, quæ, Maríæ Vírginis vitam abscónditam imitata, quasi hortus irriguus, flores ómnium virtútum germinávit, præcipue vero exímiæ in Deum et in próximum caritátis.\par
\switchcolumn
\EN
\noindent After being refreshed for the first time with the Eucharistic Bread Theresa seemed to develop an insatiable hunger for the celestial food. Then, as if by inspiration, she asked Jesus to turn all her earthly consolation into bitterness. After that she burned with a most tender love for Christ the Lord and for his Church. More than anything in the world she wanted to enter the Order of the Discalced Carmelites, where by self-sacrifice she might assist priests, missionaries and the whole Church, and so gain innumerable souls for Christ Jesus. All this, she promised God would do for her, even when apparently she lay at the point of death. Her extreme youth was an obstacle which hindered her entrance upon the religious life. Even this she overcame by her incredible courage of soul. She entered Carmel at Lisieux happily at the age of fifteen. There God fashioned the heart of Theresa in a marvellous way, teaching her to ascend to him step by step. Imitating the hidden life of the Virgin Mary like a well-watered garden she bore flowers of every virtue, especially an abiding love of God and neighbour.\par
\switchcolumn*

\LA
\begin{Resp}\Rsign Dilexísti justítiam, et odísti iniquitátem: \Star{} Proptérea unxit te Deus, Deus tuus, óleo lætítiæ.\end{Resp}
\begin{Resp}\Vsign Propter veritátem, et mansuetúdinem, et justítiam.\end{Resp}
\begin{Resp}\Rsign Proptérea unxit te Deus, Deus tuus, óleo lætítiæ.\end{Resp}
\switchcolumn
\EN
\begin{Resp}\Rsign Thou hast loved righteousness, and hated iniquity; \Star{} Therefore God, thy God, hath anointed thee with the oil of gladness.\end{Resp}
\begin{Resp}\Vsign Because of truth, and meekness, and righteousness.\end{Resp}
\begin{Resp}\Rsign Therefore God, thy God, hath anointed thee with the oil of gladness.\end{Resp}
\switchcolumn*

\LA
\LessonHead{Lectio VI}
\switchcolumn
\EN
\LessonHead{Lesson VI}
\switchcolumn*

\LA
\noindent Quo magis Altíssimo placéret, quum in Sacris Scripturis mónitum illud legísset: Si quis est párvulus véniat ad me; párvula in spíritu esse vóluit, et inde filiali fiducia Deo, tamquam Patri amantíssimo, se perpetuo trádidit. Hanc, spiritualis infantiæ viam, secúndum Evangélii doctrinam, alios docuit, speciatim novítias, quas ex obediéntia ad religiosárum virtútum studium informandas suscépit, atque ita apostolico zelo repléta, mundo, supérbia inflato et vanitátes diligenti, evangelicæ simplicitátis iter patefécit. Sponsus autem Jesus eam patiéndi desidério, tam in anima, quam in corpore, pénitus inflammávit. Insuper Dei caritátem undequáque neglectam animadvertens, summo dolóre affecta, duobus ante óbitum annis, Dei miserentis amori se víctimam óbtulit. Tunc, ut ipsa refert, cæléstis ignis flamma vulneráta est: unde caritate consumpta, in écstasim rapta, ferventíssime ingéminans: Deus meus, te diligo; viginti quatuor annos nata, die trigesima Septembris, anno millesimo octingentésimo nonagesimo septimo, ad Sponsum evolávit. Quod autem móriens promiserat, se perennem rosárum plúviam in terram demissuram, hoc in cælum recepta, innumeris miraculis reapse adimplévit et in dies adimplet. Quam Pius undecimus, Pontifex maximus, Beátis Virgínibus adscriptam, et biennio post, jubilæo máximo recurrente, inter Sanctas relatam, peculiarem ómnium Missiónum Patronam constítuit ac declarávit.\par
\switchcolumn
\EN
\noindent That she might please the most high God to greater degree, when she read in Sacred Scriptures the warning, Whoever is a little one, let him come unto me, she determined to be a little one in spirit. As such she consecrated herself forever with childlike confidence to God, her most loving Father. The way of spiritual childhood, following the teachings of the Gospel, she taught to others, especially to the novices who training in the pursuit of religious virtues she undertook in obedience to her superiors. Overflowing with apostolic zeal she pointed out to a world filled with pride and a love of vanities, the simple way of the Gospels. Meanwhile Jesus, her spouse, inflamed her with a desire to suffer both in soul and in body. Moreover, perceiving that the love of God was everywhere rejected, she became filled with grief and two years before her death, offered herself as a victim of love to the merciful God. She writeth that she was then wounded by a flame of fire from heaven, whereupon she became consumed by love, rapt as it were in ecstasy. Repeating over and over again the fervent words, My God, I love thee, she passed on to her Spouse on the 30th day of September, in the year 1897, at the age of twenty-four years. As she was dying she promised that she would let fall upon earth a ceaseless shower of roses. This promise she hath indeed fulfilled in heaven, and her shower of roses hath continued to this very day. The Sovereign Pontiff Pius XI added her name to the Virgins declared Blessed and two years later, at the time of the great Jubilee, listed her among the Saints. He also appointed and declared her Patroness of all the missions.\par
\switchcolumn*

\LA
\begin{Resp}\Rsign Afferéntur Regi vírgines post eam, próximæ ejus; \Star{} Afferéntur tibi in lætítia et exsultatióne.\end{Resp}
\begin{Resp}\Vsign Spécie tua et pulchritúdine tua; inténde, próspere procéde, et regna.\end{Resp}
\begin{Resp}\Rsign Afferéntur tibi in lætítia et exsultatióne.\end{Resp}
\begin{Resp}\Vsign Glória Patri, et Fílio, \Star{} et Spirítui Sancto.\end{Resp}
\begin{Resp}\Rsign Afferéntur tibi in lætítia et exsultatióne.\end{Resp}
\switchcolumn
\EN
\begin{Resp}\Rsign After her shall virgins be brought unto the King, her fellows \Star{} They shall be brought unto thee with gladness and rejoicing.\end{Resp}
\begin{Resp}\Vsign In thy comeliness and thy beauty, go forward, fare prosperously, and reign.\end{Resp}
\begin{Resp}\Rsign They shall be brought unto thee with gladness and rejoicing.\end{Resp}
\begin{Resp}\Vsign Glory be to the Father, and to the Son, \Star{} and to the Holy Ghost.\end{Resp}
\begin{Resp}\Rsign They shall be brought unto thee with gladness and rejoicing.\end{Resp}
\switchcolumn*

\LA
\LessonHead{Lectio VII}
\switchcolumn
\EN
\LessonHead{Lesson VII}
\switchcolumn*

\LA
\LessonTitle{Léctio sancti Evangélii secúndum Matthǽum}
\switchcolumn
\EN
\LessonTitle{From the Holy Gospel according to Matthew}
\switchcolumn*

\LA
\Cite{Matt 18:1-4}
\switchcolumn
\EN
\Cite{Matt 18:1-4}
\switchcolumn*

\LA
\noindent In illo témpore accessérunt discípuli ad Jesum, dicéntes: Quis, putas, major est in regno cælórum? Et réliqua.\par
\switchcolumn
\EN
\noindent In that time: The disciples came to Jesus, saying: Who thinkest thou is the greater in the kingdom of heaven? And so on.\par
\switchcolumn*

\LA
\LessonTitle{Homilía sancti Leónis Papæ}
\switchcolumn
\EN
\LessonTitle{Homily of St. Leo, Pope}
\switchcolumn*

\LA
\Source{Sermo 37, in Epiph. solemn. 7, cap. 3-4}
\switchcolumn
\EN
\Source{Sermo 37, in Epiphaniae solemn 7, cap 3-4}
\switchcolumn*

\LA
\noindent Tota, dilectíssimi, christianæ sapiéntiæ disciplína, non in abundántia verbi, non in astutia disputándi, neque in appetitu laudis et glóriæ, sed in vera et voluntária humilitate consistit, quam Dóminus Jesus Christus ab útero Matris usque ad supplícium Crucis pro omni fortitúdine et elégit et docuit. Nam, cum discípuli ejus inter se, ut ait Evangelista, disquírerent quis eórum major esset in regno cælórum, vocávit párvulum, et státuit eum in médio eórum, et dixit: Amen, dico vobis, nisi convérsi fuéritis et efficiámini sicut párvuli, non intrabitis in regnum cælórum. Quicúmque ergo humiliaverit se sicut puer iste, hic major erit in regno cælórum. Amat Christus infantiam, quam primum et animo suscépit et corpore. Amat Christus infantiam, humilitátis magistram, innocéntiæ regulam, mansuetúdinis formam. Amat Christus infantiam, ad quam majórum dirigit mores, ad quam senum redúcit ætates; et eos ad suum inclínat exemplum, quos ad regnum sublimat ætérnum.\par
\switchcolumn
\EN
\noindent The whole teaching of Christian wisdom consisteth, dearly beloved, not in an abundance of words, not in skill in disputation, not surely in seeking after praise or glory, but rather in seeking after true and voluntary humility. This was the way which the Lord Jesus Christ chose and taught with all his strength from his Mother's womb to his death on the Cross. When the Lord's disciples, as the Evangelist saith, discussed among themselves which should be the greatest in the kingdom of heaven, he called a little child and set him in the midst of them, and said: Verily, I say unto you, Except ye be converted, and become as little children, ye shall not enter into the kingdom of heaven. Whosoever therefore shall humble himself as this little child, the same is the greatest in the kingdom of heaven. Christ loveth childhood, which state at first he took upon himself, both in his soul and in his body. Christ loveth childhood, for it is the teacher of humility, the rule of innocence, and the type of meekness. Christ loveth childhood, for he formeth the character of the grown man on this model, and bringeth back the latter years of the old to this very state; and he shapeth on this wise those whom he would raise to the kingdom of heaven.\par
\switchcolumn*

\LA
\begin{Resp}\Rsign Hæc est Virgo sápiens, quam Dóminus vigilántem invénit, quæ accéptis lampádibus sumpsit secum óleum: \Star{} Et veniénte Dómino, introívit cum eo ad núptias.\end{Resp}
\begin{Resp}\Vsign Média nocte clamor factus est: Ecce sponsus venit, exíte óbviam ei.\end{Resp}
\begin{Resp}\Rsign Et veniénte Dómino, introívit cum eo ad núptias.\end{Resp}
\switchcolumn
\EN
\begin{Resp}\Rsign This is one of those wise virgins, whom the Lord found watching, for when she took her lamp, she took oil with her. \Star{} And when the Lord came, she went in with him to the marriage.\end{Resp}
\begin{Resp}\Vsign At midnight there was a cry made: Behold, the Bridegroom cometh, go ye out to meet him.\end{Resp}
\begin{Resp}\Rsign And when the Lord came, she went in with him to the marriage.\end{Resp}
\switchcolumn*

\LA
\LessonHead{Lectio VIII}
\switchcolumn
\EN
\LessonHead{Lesson VIII}
\switchcolumn*

\LA
\noindent Ut autem plene valeámus agnóscere quómodo apprehéndi possit tam mira conversio et in puerílem gradum qua nobis mutatióne redeundum sit, doceat nos beátus Paulus et dicat: Nolíte púeri effici sensibus, sed malítia párvuli estote. Non ergo ad lúdicra infantiæ et imperfécta nobis primordia reverténdum est, sed aliquid, quod étiam graves annos déceat, inde suméndum, ut velox sit commótium tránsitus, citus ad pacem recursus: nulla sit memória offensiónis, nulla cupiditas dignitátis; amor sociæ communiónis, æqualitas naturalis. Magnum enim bonum est nocére non nosse et maligna non sápere; quia inférre ac referre injúriam, mundi hujus prudéntia est; némini autem malum pro malo reddere, christianæ est æquanimitátis infantia.\par
\switchcolumn
\EN
\noindent But, that we may be fully able to understand how this marvellous transformation can be accomplished, and by what change we are to return to the state of childhood, let us follow the teaching of the blessed Paul, who saith: Be not children in understanding; howbeit, in malice be ye children, but in understanding be men. Hence we are not to return to the pastimes and imperfect beginnings of childhood, but thence are rather to take whatever is suitable to the full-grown: such as the swift passing of excitement; the speedy restoration of peace; the forgetfulness of injuries; the indifference to dignity; the love of the companionship of their comrades; and the natural evenness of temper. It is indeed a great good not to know and not to have a taste for harm; for to do and to return injuries is the wisdom of this world, but to render no man evil for evil is to possess the childhood of Christian goodwill.\par
\switchcolumn*

\LA
\begin{Resp}\Rsign Média nocte clamor factus est: \Star{} Ecce sponsus venit, exíte óbviam ei.\end{Resp}
\begin{Resp}\Vsign Prudéntes vírgines, aptáte vestras lámpades.\end{Resp}
\begin{Resp}\Rsign Ecce sponsus venit, exíte óbviam ei.\end{Resp}
\begin{Resp}\Vsign Glória Patri, et Fílio, \Star{} et Spirítui Sancto.\end{Resp}
\begin{Resp}\Rsign Ecce sponsus venit, exíte óbviam ei.\end{Resp}
\switchcolumn
\EN
\begin{Resp}\Rsign At midnight there was a cry made: \Star{} Behold, the Bridegroom cometh, go ye out to meet him.\end{Resp}
\begin{Resp}\Vsign Trim your lamps, O ye wise virgins.\end{Resp}
\begin{Resp}\Rsign Behold, the Bridegroom cometh, go ye out to meet him.\end{Resp}
\begin{Resp}\Vsign Glory be to the Father, and to the Son, \Star{} and to the Holy Ghost.\end{Resp}
\begin{Resp}\Rsign Behold, the Bridegroom cometh, go ye out to meet him.\end{Resp}
\switchcolumn*

\LA
\LessonHead{Lectio IX}
\switchcolumn
\EN
\LessonHead{Lesson IX}
\switchcolumn*

\LA
\noindent Ad hanc vos, dilectíssimi, similtúdinem parvulórum mystérium hodiérnæ festivitátis invitat; et hanc vobis humilitátis formam adorátus a Magis puer Salvátor insinuat: qui, ut imitatóribus suis quid glóriæ paráret, osténderet, ortus sui témpore éditos martyrio consecrávit; ut in Bethlehem, ubi Christus natus est, geniti, per communiónem ætátis consortes fíerent passiónis. Amétur ígitur humilitas, et omnis a fidelibus vitétur elátio. Alter álterum sibi præferat, et nemo quod suum est quærat, sed quod alterius; ut, cum in ómnibus abundaverit afféctus benevoléntiæ, in nullo virus inveniátur invidiæ: quóniam qui se exáltat humiliábitur, et qui se humíliat exaltábitur, eodem ipso testante Jesu Christo Dómino nostro, qui cum Patre et Spíritu Sancto vivit et regnat Deus in sǽcula sæculórum. Amen.\par
\switchcolumn
\EN
\noindent The mystery of this day's festival, dearly beloved, calleth you to this imitation of little children. And the Saviour, who was adored by the Magi as a child, recommendeth to you this pattern of humility. To shew what glory he prepareth for them that would imitate him, he consecrated by martyrdom those born at the same time as himself. And thus the children born in Bethlehem, where Christ was born, became sharers of his passion by virtue of sharing the age of his infancy. Let the faithful then love lowliness, and shun all arrogance. Let each one prefer his brother to himself. Let him seek not even what is his own, but only that which will profit his brother. The feeling of charity will abound in all; the poison of envy will be found in none. For he that exalteth himself shall be humbled, and he that humbleth himself shall be exalted. This is the teaching of our Lord Jesus Christ, who with the Father and the Holy Ghost liveth and reigneth, God, for ever and ever. Amen.\par
\switchcolumn*

\LA
\TeDeum{Te Deum}
\switchcolumn
\EN
\TeDeum{Te Deum}
\switchcolumn*

\end{paracol}

\par\needspace{10\baselineskip}\addvspace{1.2em}{\centering\small\textsc{Die 4 Octobris} · \textsc{4 October}\par}
\Office{S. Francisci Confessoris}{St. Francis of Assisi, Confessor}{Duplex majus}
\addcontentsline{toc}{subsection}{Die 4 Octobris · S. Francisci Confessoris \textit{— St. Francis of Assisi, Confessor}}
\begin{paracol}{2}
\LA
\RefLine{Lectiones I–III: de Scriptura occurrente.}
\switchcolumn
\EN
\RefLine{Lessons I–III: from the Scripture of the day.}
\switchcolumn*

\LA
\LessonHead{Lectio IV}
\switchcolumn
\EN
\LessonHead{Lesson IV}
\switchcolumn*

\LA
\noindent Francíscus, Assísii in Umbria natus, patris exémplum secútus, a prima ætáte mercatúram fecit. Qui, quodam die, páuperem, pro Christi amóre flagitántem pecúniam, cum præter consuetúdinem repulísset, repénte eo facto commótus, large ei misericórdiam impertívit; et ex eo die Deo promísit se némini umquam poscénti eleemósynam negatúrum. Cum vero post in gravem morbum incidísset, ex eo aliquándo confirmátus, cœpit ardéntius cólere offícia caritátis; qua in exercitatióne tantum profécit, ut, evangélicæ perfectiónis cúpidus, quidquid habéret, paupéribus largirétur. Quod ferens iníquius pater, eum ad Assisinátem epíscopum duxit, ut coram illo bonis céderet patérnis; qui, rejéctis étiam véstibus, patri concéssit ómnia, illud subjúngens, sibi in pósterum majórem facultátem fore dicéndi: Pater noster, qui es in cælis.\par
\switchcolumn
\EN
\noindent Francis was born at Assisi in Umbria, (in the year of our Lord 11 82.) From his early youth he followed the example of his father, (Peter Bernardone,) and busied himself with merchandise. It befell one day that, contrary to his usage, he had thrust from him a beggar, who cried for money for Christ's sake, when, being cut to the heart with regret, he gave him large alms, and promised to God from that day forth never to deny to any that asked of him. He fell after this into a grievous sickness, and from the time that he was healed thereof, he gave himself more earnestly to works of love for his neighbour. At length he became fain in this sort to be perfect, even as the Lord hath said in the Gospel, \textasciitilde{}(Matth. xix. 21,) and gave to the poor whatsoever he had. His father would not have it so, and brought him before the Bishop of Assisi, that he might renounce all right to any inheritance. He cheerfully gave up all to his father, even to his clothes, telling them that now he should be able with more utter dependence to say Our Father, Who art in heaven.\par
\switchcolumn*

\LA
\begin{Resp}\Rsign Honéstum fecit illum Dóminus, et custodívit eum ab inimícis, et a seductóribus tutávit illum: \Star{} Et dedit illi claritátem ætérnam.\end{Resp}
\begin{Resp}\Vsign Justum dedúxit Dóminus per vias rectas, et osténdit illi regnum Dei.\end{Resp}
\begin{Resp}\Rsign Et dedit illi claritátem ætérnam.\end{Resp}
\switchcolumn
\EN
\begin{Resp}\Rsign The Lord made him honourable, and defended him from his enemies, and kept him safe from those that lay in wait for him; \Star{} And gave him perpetual glory.\end{Resp}
\begin{Resp}\Vsign He went down with him into the pit, and left him not in bonds.\end{Resp}
\begin{Resp}\Rsign And gave him perpetual glory.\end{Resp}
\switchcolumn*

\LA
\LessonHead{Lectio V}
\switchcolumn
\EN
\LessonHead{Lesson V}
\switchcolumn*

\LA
\noindent Cum autem illud ex Evangélio audísset: Nolíte possidére aurum, neque argéntum, neque pecúniam in zonis vestris, non peram in via, neque duas túnicas, neque calceaménta; sibi eam régulam servándam propósuit. Itaque, detráctis cálceis et una conténtus túnica, cum duódecim sócios adhibuísset, órdinem Minórum instítuit. Quare Romam venit anno salútis millésimo ducentésimo nono, ut sui órdinis régula ab apostolica Sede confirmarétur. Quem cum accedéntem ad se summus Póntifex Innocéntius tértius rejecísset; quod in somnis póstea sibi ille, quem repúlerat, collabéntem Lateranénsem basílicam suis húmeris sustinére visus esset, conquisítum accérsi jussit, benignéque accípiens, omnem ejus institutórum ratiónem confirmávit. Francíscus ígitur, dimíssis in omnes orbis terræ partes frátribus ad prædicándum Christi Evangélium, ipse cúpiens sibi áliquam dari martyrii occasiónem, navigávit in Syriam; ubi, a rege soldáno liberalíssime tractátus, cum nihil profíceret, rédiit in Itáliam.\par
\switchcolumn
\EN
\noindent Upon the 24th day of February, in the year 1209,) he heard read the words of the Gospel Provide neither gold, nor silver, nor brass in your purses, nor scrip for your journey, neither two coats, neither shoes. (Matth. x. 9, 10.) Thereupon he determined that that should be his rule of living. He took off his shoes, and contented himself with one coat. When he had gathered twelve comrades, he founded the Order of Friars Minor. He went to Rome in the (same) year, to get from the Apostolic See a confirmation of his Order. When he came Pope Innocent III. thrust him away. Thereafter he dreamt that he saw the Church of the most Holy Saviour falling, and whom he had cast forth bearing it up with his shoulders. He bade therefore that he should be sought for and brought again before him, welcomed him kindly, and approved all the Rule which he had established. Francis therefore sent his Friars into all quarters of the world to preach the Gospel of Christ. He himself was fain to find some occasion of martyrdom, and therefore made a voyage into Syria, (in the year 1219,) but the Sultan treated him with the greatest kindness, offering him many gifts, and, since he could do no good, he returned again to Italy.\par
\switchcolumn*

\LA
\begin{Resp}\Rsign Amávit eum Dóminus, et ornávit eum: stolam glóriæ índuit eum, \Star{} Et ad portas paradísi coronávit eum.\end{Resp}
\begin{Resp}\Vsign Índuit eum Dóminus lorícam fídei, et ornávit eum.\end{Resp}
\begin{Resp}\Rsign Et ad portas paradísi coronávit eum.\end{Resp}
\switchcolumn
\EN
\begin{Resp}\Rsign The Lord loved him and beautified him; He clothed him with a robe of glory, \Star{} And crowned him at the gates of Paradise.\end{Resp}
\begin{Resp}\Vsign The Lord hath put on him the breast-plate of faith, and hath adorned him.\end{Resp}
\begin{Resp}\Rsign And crowned him at the gates of Paradise.\end{Resp}
\switchcolumn*

\LA
\LessonHead{Lectio VI}
\switchcolumn
\EN
\LessonHead{Lesson VI}
\switchcolumn*

\LA
\noindent Multis ígitur exstrúctis suæ famíliæ domicíliis, se in solitúdinem montis Alvérni cóntulit; ubi quadragínta diérum, propter honórem sancti Michaélis Archángeli, jejúnio inchoáto, festo die Exaltatiónis sanctæ Crucis, ei Séraphim crucifíxi effígiem inter alas cóntinens appáruit; qui ejus et mánibus, et pédibus, et láteri vestígia clavórum impréssit. Quæ sanctus Bonaventúra, cum Alexándri quarti summi Pontíficis prædicatióni interésset, narrásse Pontíficem a se visa esse, lítteris commendávit. His insígnibus summi in eum Christi amóris, máximam habébat ómnium admiratiónem. Ac biénnio post, gráviter ægrótans, deférri vóluit in ecclésiam sanctæ Maríæ Angelórum; ut, ubi grátiæ spíritum a Deo accéperat, ibi spíritum vitæ rédderet. Eo in loco fratres ad paupertátem ac patiéntiam, et sanctæ Románæ Ecclésiæ fidem servándam cohortátus, Psalmum illum pronúntians, Voce mea ad Dóminum clamávi; in eo versículo, Me exspéctant justi, donec retríbuas mihi, efflávit ánimam quarto Nonas Octóbris. Quem, miráculis clarum, Gregórius nonus Póntifex máximus in Sanctórum númerum adscrípsit.\par
\switchcolumn
\EN
\noindent Towards the Feast of the Assumption of the Blessed Virgin Mary, (in the year 1224), when there had already been built many houses of Friars of his Order, he withdrew himself into a most secret place upon Mount Alverno, and began to fast for forty days in honour of the holy Archangel Michael. Upon the Feast-day of the Uplifting of the Holy Cross, (as he was praying upon the side of the mountain,) he saw a vision of a crucified Seraph, which left in his hands and feet holes with nails therein, and in his side a great wound. Holy BuonaVentura hath left it in writing that he once heard Pope Alexander IV., when preaching, testify that he had himself seen these marks. It was a sign of such love of Christ toward him as stirred up the great wonder of all men. Two years thereafter he fell sick unto death, and was fain to be carried into the Church of St. Mary-of-the-Angels, that he might give up the breath of life in the same place where God had breathed into him the breath of the life of grace. Being there (laid on the earth, sprinkled with ashes, and covered with an old habit,) he exhorted the Friars to be poor and lowly, and to cleave to the faith of the Holy Church of Rome. (He then caused the Gospel of St. John to be read from the words Now before the feast of the Passover to the end,) after which he began to recite the 141st Psalm: I cried unto the Lord with my voice, and in uttering the words, the righteous wait for me, till Thou deal bountifully with me, he gave up the ghost. It was the 4th day of October, (in the year 1226.) He was famous for miracles, and Pope Gregory IX. added his name to the list of the Saints.\par
\switchcolumn*

\LA
\begin{Resp}\Rsign Iste homo perfécit ómnia quæ locútus est ei Deus, et dixit ad eum: Ingrédere in réquiem meam: \Star{} Quia te vidi justum coram me ex ómnibus géntibus.\end{Resp}
\begin{Resp}\Vsign Iste est, qui contémpsit vitam mundi, et pervénit ad cæléstia regna.\end{Resp}
\begin{Resp}\Rsign Quia te vidi justum coram me ex ómnibus géntibus.\end{Resp}
\begin{Resp}\Vsign Glória Patri, et Fílio, \Star{} et Spirítui Sancto.\end{Resp}
\begin{Resp}\Rsign Quia te vidi justum coram me ex ómnibus géntibus.\end{Resp}
\switchcolumn
\EN
\begin{Resp}\Rsign This is he which did according unto all that God commanded him; and God said unto him: Enter thou into My rest, \Star{} For thee have I seen righteous before Me among all people.\end{Resp}
\begin{Resp}\Vsign This is he which loved not his life in this world, and is come unto an everlasting kingdom.\end{Resp}
\begin{Resp}\Rsign For thee have I seen righteous before Me among all people.\end{Resp}
\begin{Resp}\Vsign Glory be to the Father, and to the Son, \Star{} and to the Holy Ghost.\end{Resp}
\begin{Resp}\Rsign For thee have I seen righteous before Me among all people.\end{Resp}
\switchcolumn*

\LA
\LessonHead{Lectio VII}
\switchcolumn
\EN
\LessonHead{Lesson VII}
\switchcolumn*

\LA
\LessonTitle{Léctio sancti Evangélii secúndum Matthǽum}
\switchcolumn
\EN
\LessonTitle{From the holy Gospel according to Matthew}
\switchcolumn*

\LA
\Cite{Matt 11:25-30}
\switchcolumn
\EN
\Cite{Matt 11:25-30}
\switchcolumn*

\LA
\noindent In illo témpore: Respóndens Jesus dixit: Confíteor tibi, Pater, Dómine cæli et terræ, quia abscondísti hæc a sapiéntibus et prudéntibus, et revelásti ea párvulis. Et réliqua.\par
\switchcolumn
\EN
\noindent At that time, Jesus answered and said: I confess to thee, O Father, Lord of heaven and earth, because thou hast hid these things from the wise and prudent, and hast revealed them to the little ones. And so on.\par
\switchcolumn*

\LA
\LessonTitle{Homilía sancti Augustíni Epíscopi}
\switchcolumn
\EN
\LessonTitle{Homily by St. Augustine, Bishop (of Hippo.)}
\switchcolumn*

\LA
\Source{Sermo 10 de verbis Domini}
\switchcolumn
\EN
\Source{10th Sermon on the Words of the Lord}
\switchcolumn*

\LA
\noindent Veníte ad me, omnes qui laborátis. Quare enim omnes laborámus, nisi quia sumus hómines mortáles, frágiles, infírmi, lútea vasa portántes, quæ fáciunt ínvicem angústias? Sed, si angustiántur vasa carnis, dilaténtur spátia caritátis. Quid ergo dicit, Veníte ad me, omnes qui laborátis, nisi ut non laborétis? Dénique promíssio ejus in promptu est; quóniam laborántes vocávit, quærent forte, qua mercéde vocáti sunt. Et ego vos, inquit, refíciam. Tóllite jugum meum super vos, et díscite a me, non mundum fabricáre, non cuncta visibília et invisibília creáre, non in ipso mundo mirabília fácere et mórtuos suscitáre; sed, Quóniam mitis sum et húmilis corde.\par
\switchcolumn
\EN
\noindent Come unto Me, all ye that labour! And wherefore labour we all, but because we are frail, sickly, dying creatures, burdened with earthen vessels which distress us? But if these fleshly vessels be distressful, let the open expanse of love be free and wide. Come unto Me, all ye that labour! and why? That we may labour no more. His promise is an instant promise, for He calleth such as are labouring. Perchance they will ask Him what shall be their reward? And I, saith He, will give you rest. Take My yoke upon you, and learn of me not how to make the world, not how to create all things visible and invisible, not to work wonders in the earth, nor to raise the dead but for I am meek and lowly in heart.\par
\switchcolumn*

\LA
\begin{Resp}\Rsign Iste est qui ante Deum magnas virtútes operátus est, et de omni corde suo laudávit Dóminum: \Star{} Ipse intercédat pro peccátis ómnium populórum.\end{Resp}
\begin{Resp}\Vsign Ecce homo sine queréla, verus Dei cultor, ábstinens se ab omni ópere malo, et pérmanens in innocéntia sua.\end{Resp}
\begin{Resp}\Rsign Ipse intercédat pro peccátis ómnium populórum.\end{Resp}
\switchcolumn
\EN
\begin{Resp}\Rsign This is he which wrought great wonders before God, and praised the Lord with all his heart. \Star{} May he pray for all people, that their sins may be forgiven unto them.\end{Resp}
\begin{Resp}\Vsign Behold a man without blame, a worshipper of God in truth, keeping himself clean from every evil work, and abiding still in his innocency.\end{Resp}
\begin{Resp}\Rsign May he pray for all people, that their sins may be forgiven unto them.\end{Resp}
\switchcolumn*

\LA
\LessonHead{Lectio VIII}
\switchcolumn
\EN
\LessonHead{Lesson VIII}
\switchcolumn*

\LA
\noindent Magnus esse vis? a mínimo íncipe. Cógitas magnam fábricam constrúere celsitúdinis? de fundaménto prius cógita humilitátis. Et quantam quisque vult et dispónit superimpónere molem ædifícii, quanto erit majus ædifícium, tanto áltius fodit fundaméntum. Et fábrica quidem cum constrúitur, in supérna consúrgit; qui autem fodit fundaméntum, ad ima deprímitur. Ergo et fábrica ante celsitúdinem humiliátur, et fastígium post humiliatiónem erígitur.\par
\switchcolumn
\EN
\noindent Wilt thou be great? Begin by being little. Dost thou think to raise up a lofty building? Then lay the foundations thereof in lowliness. The greater soever, and the more massy, be that which any man thinketh to build, so much the deeper doth he dig his foundation. And when the house is built, it towereth heavenward; but he which layeth the foundation goeth down into the earth. The building, therefore, is low before it is high, and, after it is low, it riseth high to the roof.\par
\switchcolumn*

\LA
\begin{Resp}\Rsign Sint lumbi vestri præcíncti, et lucérnæ ardéntes in mánibus vestris: \Star{} Et vos símiles homínibus exspectántibus dóminum suum, quando revertátur a núptiis.\end{Resp}
\begin{Resp}\Vsign Vigiláte ergo, quia nescítis qua hora Dóminus vester ventúrus sit.\end{Resp}
\begin{Resp}\Rsign Et vos símiles homínibus exspectántibus dóminum suum, quando revertátur a núptiis.\end{Resp}
\begin{Resp}\Vsign Glória Patri, et Fílio, \Star{} et Spirítui Sancto.\end{Resp}
\begin{Resp}\Rsign Et vos símiles homínibus exspectántibus dóminum suum, quando revertátur a núptiis.\end{Resp}
\switchcolumn
\EN
\begin{Resp}\Rsign Let your loins be girded about, and your lights burning; \Star{} And ye yourselves like unto men that wait for their lord, when he will return from the wedding.\end{Resp}
\begin{Resp}\Vsign Watch therefore, for ye know not what hour your Lord doth come.\end{Resp}
\begin{Resp}\Rsign And ye yourselves like unto men that wait for their lord, when he will return from the wedding.\end{Resp}
\begin{Resp}\Vsign Glory be to the Father, and to the Son, \Star{} and to the Holy Ghost.\end{Resp}
\begin{Resp}\Rsign And ye yourselves like unto men that wait for their lord, when he will return from the wedding.\end{Resp}
\switchcolumn*

\LA
\LessonHead{Lectio IX}
\switchcolumn
\EN
\LessonHead{Lesson IX}
\switchcolumn*

\LA
\noindent Quod est fastígium construéndæ fábricæ, quam molímur? quo perventúrum est cacúmen ædifícii? Cito dico, usque ad conspéctum Dei. Vidétis quam excélsum est, quanta res est conspícere Deum. Qui desíderat, et quod dico et quod audit intéllegit. Promíttitur nobis conspéctus Dei, veri Dei, summi Dei. Hoc enim bonum est, vidéntem vidére. Nam, qui colunt falsos deos, fácile illos vident; sed eos vident, qui óculos habent et non vident. Nobis autem promíttitur vísio Dei vivéntis et vidéntis.\par
\switchcolumn
\EN
\noindent What is the roof of the house on which we labour? Whither do its spires rise? I answer you at once; to the presence of God. You see how high it is, yea, what it is to see God. He that will, understandeth what I say, and he heareth. What is promised you is to see God, God, the True, God, the Supreme. Blessed is he who seeth Him by Whom he is seen. Such as worship false gods see them easily, but they see them who have eyes and see not. But unto us it is promised that we shall see that God Who liveth and seeth. (Gen. xvi. 14.)\par
\switchcolumn*

\LA
\TeDeum{Te Deum}
\switchcolumn
\EN
\TeDeum{Te Deum}
\switchcolumn*

\end{paracol}

\par\needspace{10\baselineskip}\addvspace{1.2em}{\centering\small\textsc{Die 5 Octobris} · \textsc{5 October}\par}
\Office{Ss. Placidi et Sociorum Martyrum}{Sts. Placidus and Companions Martyrs}{Simplex}
\addcontentsline{toc}{subsection}{Die 5 Octobris · Ss. Placidi et Sociorum Martyrum \textit{— Sts. Placidus and Companions Martyrs}}
\begin{paracol}{2}
\LA
\RefLine{Lectiones I–II: de Scriptura occurrente.}
\switchcolumn
\EN
\RefLine{Lessons I–II: from the Scripture of the day.}
\switchcolumn*

\LA
\LessonHead{Lectio III (pro commemoratione)}
\switchcolumn
\EN
\LessonHead{Lesson III (when commemorated)}
\switchcolumn*

\LA
\noindent Plácidus Romæ, Tertúllo patre in primis nóbili, natus, puer Deo oblátus et sancto Benedícto tráditus, tantum ejus disciplína et monásticæ vitæ institútis profécit, ut inter præcípuos illíus discípulos numerétur. Ab eo in Sicíliam missus, monastérium et ecclésiam in honórem sancti Joánnis Baptístæ prope Messánæ portum constrúxit, ubi cum mónachis admirábili sanctitáte vixit. Ejus viséndi causa cum eo veníssent Eutychius et Victorínus, illíus fratres, et Flávia, virgo soror, eódem témpore illuc áppulit immánis quidam piráta, Manúcha nómine; qui, capto monastério, cum Plácidum et céteros nullo modo addúcere potuísset ut Christum negárent, ipsum fratrésque illíus ac sorórem crudéliter necári jussit. Cum quibus étiam Donátus, Firmátus diáconus, Faustus aliíque trigínta mónachi, martyrii agónem felíciter consummárunt tértio Nonas Octóbris, anno salútis quingentésimo trigésimo nono.\par
\switchcolumn
\EN
\noindent Placidus was the son of Tertullus, one of the noblest persons of Rome. He was offered to God (by his father) when a child (only seven years of age) and given over to holy Benedict, in whose teaching and Rule of monks he so profited that he was reckoned among the chiefest of his disciples. By him he was sent into Sicily, where he founded near the Port of Messina a Church and monastery in honour of St. John the Baptist, and lived therein with his monks in wonderful holiness. Thither there came to see him his brothers Eutychius and Victorinus and his virgin sister Flavia, and while they were together, there landed there a certain brutal pirate, named Manucha, who took the monastery, and when he could in no wise prevail upon Placidus and the others to deny Christ, he commanded him, his brothers, and his sister to be cruelly murdered. With them Donatus, Firmatus a Deacon, Faustus, and thirty other monks brought the conflict of testimony to the blessed end of martyrdom, upon the fifth day of October, in the year of salvation 539\par
\switchcolumn*

\LA
\TeDeum{Te Deum}
\switchcolumn
\EN
\TeDeum{Te Deum}
\switchcolumn*

\end{paracol}

\par\needspace{10\baselineskip}\addvspace{1.2em}{\centering\small\textsc{Die 6 Octobris} · \textsc{6 October}\par}
\Office{S. Brunonis Confessoris}{St. Bruno, Confessor}{Duplex}
\addcontentsline{toc}{subsection}{Die 6 Octobris · S. Brunonis Confessoris \textit{— St. Bruno, Confessor}}
\begin{paracol}{2}
\LA
\RefLine{Lectiones I–III: de Scriptura occurrente.}
\switchcolumn
\EN
\RefLine{Lessons I–III: from the Scripture of the day.}
\switchcolumn*

\LA
\RefLine{Lectiones VII–IX: ex Commúni Confessoris non pontificis (C5).}
\switchcolumn
\EN
\RefLine{Lessons VII–IX: from the Common of a Confessor not a Bishop (C5).}
\switchcolumn*

\LA
\LessonHead{Lectio IV}
\switchcolumn
\EN
\LessonHead{Lesson IV}
\switchcolumn*

\LA
\noindent Bruno, Carthusianæ religiónis institutor, Coloniæ Agrippinæ natus est. Ab ipsis incunabulis specimen futuræ sanctitátis præferens, morum gravitáte, puerília illíus ætátis, divina favénte grátia, declinans, adeo excelluit, ut jam inde monachórum pater vitæque anachoreticæ futurus instaurátor agnoscerétur. A paréntibus, genere ac virtúte claris, Lutétiam Parisiórum missus, tantum ibi in philosophíæ ac theologíæ stúdiis profecit, ut doctoris ac magistri munus in utraque facultate sit adeptus; nec multo post, ob egregias ipsíus virtútes, ecclésiæ Rheménsis canonicatu potítus.\par
\switchcolumn
\EN
\noindent Bruno, the Founder of the Charterhouse Monks, was born at Cologne, (about the year of our Lord 1030.) From his earliest years he was a very grave child, turning away from childish things, and that so manifestly, that by the grace of God the tokens of holiness already pointed him out as a Father of monks, and a restorer of the life of hermits. His parents, who were eminent for rank and goodness, sent him to Paris, where he studied so well in Philosophy and Theology, that he took the degree of Doctor in both faculties; and a short while after, for his famous graces, he was made a Canon of Rheims.\par
\switchcolumn*

\LA
\begin{Resp}\Rsign Honéstum fecit illum Dóminus, et custodívit eum ab inimícis, et a seductóribus tutávit illum: \Star{} Et dedit illi claritátem ætérnam.\end{Resp}
\begin{Resp}\Vsign Justum dedúxit Dóminus per vias rectas, et osténdit illi regnum Dei.\end{Resp}
\begin{Resp}\Rsign Et dedit illi claritátem ætérnam.\end{Resp}
\switchcolumn
\EN
\begin{Resp}\Rsign The Lord made him honourable, and defended him from his enemies, and kept him safe from those that lay in wait for him; \Star{} And gave him perpetual glory.\end{Resp}
\begin{Resp}\Vsign He went down with him into the pit, and left him not in bonds.\end{Resp}
\begin{Resp}\Rsign And gave him perpetual glory.\end{Resp}
\switchcolumn*

\LA
\LessonHead{Lectio V}
\switchcolumn
\EN
\LessonHead{Lesson V}
\switchcolumn*

\LA
\noindent Elapsis aliquot annis, cum sex aliis familiaribus mundo renuntians, sanctum Hugónem episcopum Gratianopolitánum adiit. Qui, causa eórum adventus cógnita, eosdemque intélligens esse quos eadem nocte véluti septem stellas ad suos pedes corruéntes in somnis viderat, montes suæ diœcesis asperrimos, quos Carthusiános appellant, illis concessit. Illuc Bruno cum sociis, ipso Hugone comitante, secédens, cum per aliquot annos eremiticam vitam egísset, ab Urbano secundo, qui ejusdem Brunónis discipulus fuerat, Romam accersitur. Ejus consílio ac doctrina Pontifex in tot illis Ecclésiæ calamitátibus, per aliquot annos usus est; donec Bruno, recusato Rhegiénsi archiepiscopatu, discedéndi facultátem obtinuit.\par
\switchcolumn
\EN
\noindent After some years, he, and six comrades, forsook the world and betook themselves to Hew, the holy Bishop of Grenoble, who, when he learned the reason of their coming, and believing them to have been figured by seven stars which he had seen that night in a dream falling at his feet, gave them a grant of land in some very wild mountains in his Dioecese, which are called the Chartreuses. Thither Bruno and his companions, together with Hew, withdrew themselves, (in the year 1084,) and led for some years the life of hermits. Pope Urban II., who had formerly been his disciple (at Rheims,) commanded him to come to Rome, (in 1089,) and amid the afflictions which then scourged the Church, held him for some time as his counsellor. But at last Bruno, who had refused the Archbishopric of Reggio, got his leave to go away.\par
\switchcolumn*

\LA
\begin{Resp}\Rsign Amávit eum Dóminus, et ornávit eum: stolam glóriæ índuit eum, \Star{} Et ad portas paradísi coronávit eum.\end{Resp}
\begin{Resp}\Vsign Índuit eum Dóminus lorícam fídei, et ornávit eum.\end{Resp}
\begin{Resp}\Rsign Et ad portas paradísi coronávit eum.\end{Resp}
\switchcolumn
\EN
\begin{Resp}\Rsign The Lord loved him and beautified him; He clothed him with a robe of glory, \Star{} And crowned him at the gates of Paradise.\end{Resp}
\begin{Resp}\Vsign The Lord hath put on him the breast-plate of faith, and hath adorned him.\end{Resp}
\begin{Resp}\Rsign And crowned him at the gates of Paradise.\end{Resp}
\switchcolumn*

\LA
\LessonHead{Lectio VI}
\switchcolumn
\EN
\LessonHead{Lesson VI}
\switchcolumn*

\LA
\noindent Igitur, solitúdinis amóre, eremum quamdam apud Squillácum in Calabriæ fínibus pétiit. Quo in loco, cum ipsum orántem Rogerius comes Calabriæ inter venándum, latrántibus ad illíus spelúncam cánibus, reperísset, sanctitáte viri permotus, illum ac socios fovére et cólere impense cœpit. Nec liberalitas sine præmio fuit; cum enim idem Rogerius Capuam obsidéret, eumque Sergius quidam excubiárum magister pródere statuísset, Bruno, adhuc in dicta erémo vivens, in somnis illi ómnia apériens, ab imminénti periculo cómitem liberávit. Tandem virtútibus ac meritis plenus, nec sanctitáte minus quam doctrinæ fama clarus, obdormívit in Dómino; sepultúsque est in monasterio sancti Stephani, ab ipso Rogerio constructo, ubi háctenus honorifice cólitur.\par
\switchcolumn
\EN
\noindent In his love of the wilderness, he betook himself to a certain desert place in the Diocese of Squillaci, in the uttermost coasts of Calabria, (whither he went in 1090.) He was praying there one day in a cave, when the hounds of Roger, Sovereign Earl of Sicily and Calabria, who was out a-hunting, came and bayed at the door of it. Thus was he found by this Prince, who was moved by his holiness, and began to cherish him and his comrades, and treat them very kindly. The Earl's goodness was rewarded, for when he was one time laying siege to Capua, and one Sergius, who was first groom of his bedchamber, had made a plot to betray him, Bruno, who was still living in the desert above mentioned, appeared to him in a dream, and delivered him from the danger which was hanging over him. At length Bruno, full of graces and good works, and famous for godliness not less than for learning, fell asleep in the Lord, (upon the 6th day of October, in the year 1101,) and was buried in the monastery of St. Stephen, founded by the same Earl Roger, where he is still held in great honour.\par
\switchcolumn*

\LA
\begin{Resp}\Rsign Iste homo perfécit ómnia quæ locútus est ei Deus, et dixit ad eum: Ingrédere in réquiem meam: \Star{} Quia te vidi justum coram me ex ómnibus géntibus.\end{Resp}
\begin{Resp}\Vsign Iste est, qui contémpsit vitam mundi, et pervénit ad cæléstia regna.\end{Resp}
\begin{Resp}\Rsign Quia te vidi justum coram me ex ómnibus géntibus.\end{Resp}
\begin{Resp}\Vsign Glória Patri, et Fílio, \Star{} et Spirítui Sancto.\end{Resp}
\begin{Resp}\Rsign Quia te vidi justum coram me ex ómnibus géntibus.\end{Resp}
\switchcolumn
\EN
\begin{Resp}\Rsign This is he which did according unto all that God commanded him; and God said unto him: Enter thou into My rest, \Star{} For thee have I seen righteous before Me among all people.\end{Resp}
\begin{Resp}\Vsign This is he which loved not his life in this world, and is come unto an everlasting kingdom.\end{Resp}
\begin{Resp}\Rsign For thee have I seen righteous before Me among all people.\end{Resp}
\begin{Resp}\Vsign Glory be to the Father, and to the Son, \Star{} and to the Holy Ghost.\end{Resp}
\begin{Resp}\Rsign For thee have I seen righteous before Me among all people.\end{Resp}
\switchcolumn*

\end{paracol}

\par\needspace{10\baselineskip}\addvspace{1.2em}{\centering\small\textsc{Die 7 Octobris} · \textsc{7 October}\par}
\Office{Sacratissimi Rosarii Beatæ Mariæ Virginis}{Sacratissimi Rosarii Beatæ Mariæ Virginis}{Duplex II classis}
\addcontentsline{toc}{subsection}{Die 7 Octobris · Sacratissimi Rosarii Beatæ Mariæ Virginis \textit{— Sacratissimi Rosarii Beatæ Mariæ Virginis}}
\begin{paracol}{2}
\LA
\RefLine{Lectio IX: de S. Marci Papæ Confessoris (10-07t).}
\switchcolumn
\EN
\RefLine{Lesson IX: of St. Mark, Pope and Confessor (10-07t).}
\switchcolumn*

\LA
\LessonHead{Lectio I}
\switchcolumn
\EN
\LessonHead{Lesson I}
\switchcolumn*

\LA
\LessonTitle{De Libro Ecclesiástici}
\switchcolumn
\EN
\LessonTitle{Lesson from the book of Ecclesiasticus}
\switchcolumn*

\LA
\Cite{Sir 24:11-16}
\switchcolumn
\EN
\Cite{Sir 24:11-16}
\switchcolumn*

\LA
\noindent In ómnibus réquiem quæsívi, et in hereditáte Dómini morábor. Tunc præcépit, et dixit mihi Creátor ómnium: et, qui creávit me, requiévit in tabernáculo meo, et dixit mihi: In Jacob inhábita, et in Israël hereditáre, et in eléctis meis mitte radíces. Ab inítio et ante sǽcula creáta sum, et usque ad futúrum sǽculum non désinam, et in habitatióne sancta coram ipso ministrávi. Et sic in Sion firmáta sum, et in civitáte sanctificáta simíliter requiévi, et in Jerúsalem potéstas mea. Et radicávi in pópulo honorificáto, et in parte Dei mei heréditas illíus, et in plenitúdine sanctórum deténtio mea.\par
\switchcolumn
\EN
\noindent And by my power I have trodden under my feet the hearts of all the high and low: and in all these I sought rest, and I shall abide in the inheritance of the Lord. Then the creator of all things commanded, and said to me: and he that made me, rested in my tabernacle, And he said to me: Let thy dwelling be in Jacob, and thy inheritance in Israel, and take root in my elect. From the beginning, and before the world, was I created, and unto the world to come I shall not cease to be, and in the holy dwelling place I have ministered before him. And so was I established in Sion, and in the holy city likewise I rested, and my power was in Jerusalem. And I took root in an honourable people, and in the portion of my God his inheritance, and my abode is in the full assembly of saints.\par
\switchcolumn*

\LA
\begin{Resp}\Rsign Súmite psaltérium jucúndum in insígni die solemnitátis vestræ: \Star{} Et exsultáte Vírgini adjutríci nostræ.\end{Resp}
\begin{Resp}\Vsign Cantáte ei cánticum novum: annuntiáte inter gentes glóriam ejus.\end{Resp}
\begin{Resp}\Rsign Et exsultáte Vírgini adjutríci nostræ.\end{Resp}
\switchcolumn
\EN
\begin{Resp}\Rsign Take the pleasant harp, in the time appointed, on your solemn feastday; \Star{} And sing aloud unto the Virgin our strength.\end{Resp}
\begin{Resp}\Vsign O sing unto her a new song, declare her glory among the heathen.\end{Resp}
\begin{Resp}\Rsign And sing aloud unto the Virgin our strength.\end{Resp}
\switchcolumn*

\LA
\LessonHead{Lectio II}
\switchcolumn
\EN
\LessonHead{Lesson II}
\switchcolumn*

\LA
\Cite{Sir 24:17-22}
\switchcolumn
\EN
\Cite{Sir 24:17-22}
\switchcolumn*

\LA
\noindent Quasi cedrus exaltáta sum in Líbano, et quasi cypréssus in monte Sion: quasi palma exaltáta sum in Cades, et quasi plantátio rosæ in Jericho: quasi olíva speciósa in campis, et quasi plátanus exaltáta sum juxta aquam in platéis. Sicut cinnamómum et bálsamum aromatízans odórem dedi: quasi myrrha elécta, dedi suavitátem odóris; et quasi storax, et gálbanus, et úngula, et gutta, et quasi Líbanus non incísus vaporávi habitatiónem meam, et quasi bálsamum non mixtum odor meus. Ego quasi terebínthus exténdi ramos meos, et rami mei honóris et grátiæ.\par
\switchcolumn
\EN
\noindent I was exalted like a cedar in Libanus, and as a cypress tree on mount Sion. I was exalted like a palm tree in Cades, and as a rose plant in Jericho: As a fair olive tree in the plains, and as a plane tree by the water in the streets, was I exalted. I gave a sweet smell like cinnamon. and aromatical balm: I yielded a sweet odour like the best myrrh: And I perfumed my dwelling as storax, and galbanum, and onyx, and aloes, and as the frankincense not cut, and my odour is as the purest balm. I have stretched out my branches as the turpentine tree, and my branches are of honour and grace.\par
\switchcolumn*

\LA
\begin{Resp}\Rsign Vidi speciósam ascendéntem désuper rivos aquárum; cujus inæstimábilis odor erat nimis; \Star{} Et sicut dies verni circúmdabant eam flores rosárum et lília convállium.\end{Resp}
\begin{Resp}\Vsign Astitit Regína a dextris tuis in vestítu deauráto, circúmdata varietáte.\end{Resp}
\begin{Resp}\Rsign Et sicut dies verni circúmdabant eam flores rosárum et lília convállium.\end{Resp}
\switchcolumn
\EN
\begin{Resp}\Rsign I saw the fair one going up above the rivers of waters, a priceless savour hung heavy; \Star{} And about her it was as the flower of roses in the spring of the year, and lilies of the valleys.\end{Resp}
\begin{Resp}\Vsign Upon thy right hand did stand the Queen in a vesture of gold, wrought about with diverse colours.\end{Resp}
\begin{Resp}\Rsign And about her it was as the flower of roses in the spring of the year, and lilies of the valleys.\end{Resp}
\switchcolumn*

\LA
\LessonHead{Lectio III}
\switchcolumn
\EN
\LessonHead{Lesson III}
\switchcolumn*

\LA
\Cite{Sir 24:24-31}
\switchcolumn
\EN
\Cite{Sir 24:24-31}
\switchcolumn*

\LA
\noindent Ego mater pulchræ dilectiónis, et timóris, et agnitiónis, et sanctæ spei. In me grátia omnis viæ et veritátis, in me omnis spes vitæ et virtútis. Transíte ad me, omnes qui concupíscitis me, et a generatiónibus meis implémini; spíritus enim meus super mel dulcis, et heréditas mea super mel et favum. Memória mea in generatiónes sæculórum. Qui edunt me, adhuc esúrient, et qui bíbunt me, adhuc sítient. Qui audit me non confundétur, et qui operántur in me non peccábunt. Qui elúcidant me, vitam ætérnam habébunt.\par
\switchcolumn
\EN
\noindent I am the mother of fair love, and of fear, and of knowledge, and of holy hope. In me is all grace of the way and of the truth, in me is all hope of life and of virtue. Come over to me, all ye that desire me, and be filled with my fruits. For my spirit is sweet above honey, and my inheritance above honey and the honeycomb. My memory is unto everlasting generations. They that eat me, shall yet hunger: and they that drink me, shall yet thirst. He that hearkeneth to me, shall not be confounded: and they that work by me, shall not sin. They that explain me shall have life everlasting.\par
\switchcolumn*

\LA
\begin{Resp}\Rsign Quæ est ista quæ procéssit sicut sol, et formósa tamquam Jerúsalem? Vidérunt eam fíliæ Sion, et beátam dixérunt, \Star{} Et regínæ laudavérunt eam.\end{Resp}
\begin{Resp}\Vsign Et sicut dies verni circúmdabant eam flores rosárum et lília convállium.\end{Resp}
\begin{Resp}\Rsign Et regínæ laudavérunt eam.\end{Resp}
\begin{Resp}\Vsign Glória Patri, et Fílio, \Star{} et Spirítui Sancto.\end{Resp}
\begin{Resp}\Rsign Et regínæ laudavérunt eam.\end{Resp}
\switchcolumn
\EN
\begin{Resp}\Rsign Who is this that cometh up like the sun? this, comely as Jerusalem? \Star{} The daughters of Zion saw her, and called her blessed; the queens also, and they praised her.\end{Resp}
\begin{Resp}\Vsign And about her it was as the flower of roses in the spring of the year, and lilies of the valleys.\end{Resp}
\begin{Resp}\Rsign The daughters of Zion saw her, and they called her blessed; the queens also, and they praised her.\end{Resp}
\begin{Resp}\Vsign Glory be to the Father, and to the Son, \Star{} and to the Holy Ghost.\end{Resp}
\begin{Resp}\Rsign The daughters of Zion saw her, and they called her blessed; the queens also, and they praised her.\end{Resp}
\switchcolumn*

\LA
\LessonHead{Lectio IV}
\switchcolumn
\EN
\LessonHead{Lesson IV}
\switchcolumn*

\LA
\noindent Cum Albigénsium hǽresis per Tolosátium regiónem ímpie grassarétur, atque áltius in dies radíces ágeret, sanctus Domínicus, qui nuper Prædicatórum órdinis fundaménta jécerat, ad eam convelléndam totus incúbuit. Id ut præstáret valídius, auxílium beátæ Vírginis, cujus dígnitas illis erróribus impudentíssime petebátur, cuíque datum est cunctas hǽreses interímere in univérso mundo, eníxis précibus implorávit. A qua (ut memóriæ próditum est) cum mónitus esset ut Rosárium pópulis prædicáret, velut singuláre advérsus hǽreses ac vítia præsídium; mirum est quanto mentis fervóre et quam felíci succéssu injúnctum sibi munus sit exsecútus. Est autem Rosárium certa precándi fórmula, qua quíndecim angelicárum salutatiónum décades, oratióne Domínica interjécta, distínguimus, et ad eárum síngulas tótidem nostræ reparatiónis mystéria, pia meditatióne recólimus. Ex eo ergo témpore pius hic orándi modus mirabíliter per sanctum Domínicum promulgári augeríque cœpit. Quem ejúsdem institutórem auctorémque fuísse, summi Pontífices apostólicis lítteris passim affirmárunt.\par
\switchcolumn
\EN
\noindent When the heresy of the Albigenses was making head against God in the County of Toulouse, and striking deeper roots every day, the holy Dominic, who had but just laid the foundations of the Order of Friars Preachers, threw his whole strength into the travail of plucking these blasphemies up. That he might be fitter for the work, he cried for help with his whole soul to that Blessed Maiden, whose glory the falsehoods of the heretics so insolently assailed, and to whom it hath been granted to trample down every heresy throughout the whole earth. It is said that he had from her a word, bidding him preach up the saying of the Rosary among the people, as a strong help against heresy and sin, and it is wonderful with how stout an heart and how good a success he did the work laid upon him. This Rose-garden (or Rosary) is a certain form of prayer, wherein we say one-hundred-and-fifty times the salutation of the Angel, and the Lord's Prayer between every ten times, and, each of the fifteen times that we say the Lord's Prayer, and repeat tenfold the salutation, think of one of fifteen great events in the history of our Redemption. From that time forth this form of godly prayer was extraordinarily spread about by holy Dominic, and waxed common. That this same Dominic was the founder and prime mover thereof hath been said by Popes in diverse letters of the Apostolic See.\par
\switchcolumn*

\LA
\begin{Resp}\Rsign Tu glória Jerúsalem, tu lætítia Israël, tu honorificéntia pópuli nostri, fecísti viríliter: \Star{} Quia cunctas hǽreses sola interemísti.\end{Resp}
\begin{Resp}\Vsign Pulchra es et decóra, terríbilis ut castrórum ácies ordináta.\end{Resp}
\begin{Resp}\Rsign Quia cunctas hǽreses sola interemísti.\end{Resp}
\switchcolumn
\EN
\begin{Resp}\Rsign Thou art the glory of Jerusalem thou art the joy of Israel thou art the honour of our nation thou hast done manfully \Star{} For thou alone hast slain all heresies.\end{Resp}
\begin{Resp}\Vsign Fair and comely art thou, terrible as a fenced camp set in battle array.\end{Resp}
\begin{Resp}\Rsign For thou alone hast slain all heresies.\end{Resp}
\switchcolumn*

\LA
\LessonHead{Lectio V}
\switchcolumn
\EN
\LessonHead{Lesson V}
\switchcolumn*

\LA
\noindent Innumerábiles porro fructus ex hac tam salutári institutióne in christiánam rempúblicam dimanárunt. Inter quos victória illa mérito numerátur, quam sanctíssimus Póntifex Pius quintus et ab eo inflammáti christiáni príncipes apud Echínadas ínsulas de Turcárum tyránno potentíssimo reportárunt. Nam, cum illa ipsa die ea victória reláta sit, qua die sacratíssimi Rosárii sodalitátes per univérsum orbem consuétas supplicatiónes perágerent statutásque preces de more fúnderent, iis précibus haud immérito refértur accépta. Quod quidem cum étiam Gregórius tértius décimus testátus esset; ut pro tam singulári benefício beátæ Vírgini sub appellatióne Rosárii perénnes grátiæ ubíque terrárum haberéntur, in ecclésiis ómnibus, in quibus altáre Rosárii foret, Offícium, ritu dúplici majóri, perpétuo de eo celebrándum indíxit; aliíque Pontífices Rosárium recitántibus ejusdémque Rosárii sodalitátibus indulgéntias pene innúmeras concessére.\par
\switchcolumn
\EN
\noindent From this healthy exercise have grown up numberless good fruits in the Christian Commonwealth. Among these deserveth well to be named that great victory over the Sultan of Turkey, which the most holy Pope Pius V, and the Christian Princes whom he had roused, won at Lepanto, (on the 7th day of October, the first Lord's Day in the month, in the year of our Lord 1571) The day whereon this victory was gained was the very one whereon the Guildbrethren of the most holy Rosary, throughout the whole world, were used to offer their accustomed prayers and appointed supplications, and the event therefore was not unnaturally connected therewith. This being the avowed opinion of Gregory XIII, he ordered that in all Churches where there was, or should be, an Altar of the Rosary, a Feast, in the form of a Greater Double, should be kept for ever upon the first Lord's Day of the month of October, to give unceasing thanks to the Blessed Virgin, under her style of (Queen of) the (Most Holy) Rosary, for that extraordinary mercy of God. Other Popes also have granted almost numberless Indulgences to those who say the Rosary, and to those who join its Guilds.\par
\switchcolumn*

\LA
\begin{Resp}\Rsign Déxtera tua magnificáta est in fortitúdine, déxtera tua confrégit inimícos: \Star{} Submérsi sunt in aquis veheméntibus, et opéruit eos mare.\end{Resp}
\begin{Resp}\Vsign Benedíxit te Dóminus in virtúte sua, quia per te ad níhilum redégit inimícos nostros.\end{Resp}
\begin{Resp}\Rsign Submérsi sunt in aquis veheméntibus, et opéruit eos mare.\end{Resp}
\switchcolumn
\EN
\begin{Resp}\Rsign Thy right hand is become glorious in power; thy right hand hath dashed in pieces the enemy \Star{} They sank as lead in the mighty waters, and the sea covered them.\end{Resp}
\begin{Resp}\Vsign The Lord hath blessed thee by His power, because through thee He hath brought our enemies to nought.\end{Resp}
\begin{Resp}\Rsign They sank as lead in the mighty waters, and the sea covered them.\end{Resp}
\switchcolumn*

\LA
\LessonHead{Lectio VI}
\switchcolumn
\EN
\LessonHead{Lesson VI}
\switchcolumn*

\LA
\noindent Clemens vero undécimus, ánimo réputans insígnem páriter victóriam, anno millésimo septingentésimo décimo sexto in Hungáriæ regno a Cárolo sexto in imperatórem Romanórum elécto de innúmeris Turcárum cópiis relátam, eo die contigísse quo festum Dedicatiónis sanctæ Maríæ ad Nives celebrarétur, atque eódem ferme témpore quo sacratíssimi Rosárii confrátres públicam solemnémque supplicatiónem in alma Urbe, ingénti pópuli concúrsu magnáque religióne, peragéntes, férvidas ad Deum preces pro Turcárum depressióne fúnderent ac poténtem opem Deíparæ Vírginis in auxílium Christianórum humíliter implorárent; eam ob rem, victóriam illam, nec non liberátam paulo post eorúndem Turcárum obsidióne Corcyrénsem ínsulam, ejúsdem beátæ Vírginis patrocínio pie cénsuit adscribéndam. Quam ob rem ut hujus quoque tam insígnis benefícii perénnis semper et memória exstáret et grátia, sacratíssimi Rosárii festum eódem ritu celebrándum ad Ecclésiam univérsam exténdit. Hæc ómnia Benedíctus décimus tértius in Breviário Románo appóni jussit. Leo autem tértius décimus in turbulentíssimis Ecclésiæ tempóribus, diúque preméntium malórum sæva tempestáte, cunctos in orbe fidéles, iterátis apostólicis lítteris ad Mariális Rosárii, præsértim per octóbrem mensem, frequéntiam veheménter incéndit, aucto quoque ánnui festi solemnitátis ritu, additáque litaníis Lauretánis Regínæ sacratíssimi Rosárii invocatióne, et Offício de eádem solemnitáte próprio Ecclésiæ univérsæ concésso. Sanctíssimam ergo Dei Genetrícem cultu hoc eídem gratíssimo júgiter venerémur; ut, quæ tóties Christi fidélibus, Rosárii précibus exoráta, terrénos hostes profligáre dedit ac pérdere, inférnos páriter superáre concédat.\par
\switchcolumn
\EN
\noindent In the year 1716, Charles VI, Elect-Emperor of the Romans, won a famous victory over countless hordes of Turks, (near Timisoara,) in the kingdom of Hungary, upon the day when the Feast of the Dedication of the Church of St. Mary of the Snows was being kept, and almost at the very moment when the Guild-brethren of the most holy Rosary were moving through the streets of Rome in public and solemn procession, amid vast multitudes, all filled with the deepest enthusiasm, calling vehemently upon God for the defeat of the Turks, and entreating the Virgin Mother of God to bring the might of her succour to the help of the Christians. A few days later, (upon the Octave of the Feast of the Assumption,) the Turks raised the siege of Corfu. These mercies Clement XI devoutly ascribed to the helpful prayers of the Blessed Virgin, and that the memory and the sweetness of such a blessing might for all time coming endure gloriously, he extended to the whole Church the observance of the Feast of the most holy Rosary, for the same day and of the same rank, (as it had already been in the places before mentioned.) Benedict XIII commanded the record of all these things to be given a place in the Service-book of the Church of Rome; and Leo XIII, in the most troublous times of the Church and the cruel storm of long pressing evils, by fresh Apostolic letters vehemently urged upon all the faithful throughout the earth the often saying of the Rosary of (the Blessed Virgin) Mary, raised the dignity of the yearly festival, added to the Litany of Loretto the Invocation Queen of the Most Holy Rosary, and granted to the whole Church a special Office for this solemn occasion. Let us all then be earnest in honouring the most holy Mother of God in this form which she liketh so well, that even as the entreaties of Christ's faithful people, approaching her in her Garden of Roses, have so often won her to scatter and destroy their earthly foes, so she may gain for them the victory over their hellish foes likewise.\par
\switchcolumn*

\LA
\begin{Resp}\Rsign Signum magnum appáruit in cælo: Múlier amícta sole, et luna sub pédibus ejus, \Star{} Et in cápite ejus coróna stellárum duódecim.\end{Resp}
\begin{Resp}\Vsign Dábitur cápiti tuo augméntum gratiárum et coróna ínclita próteget te.\end{Resp}
\begin{Resp}\Rsign Et in cápite ejus coróna stellárum duódecim.\end{Resp}
\begin{Resp}\Vsign Glória Patri, et Fílio, \Star{} et Spirítui Sancto.\end{Resp}
\begin{Resp}\Rsign Et in cápite ejus coróna stellárum duódecim.\end{Resp}
\switchcolumn
\EN
\begin{Resp}\Rsign There appeared a great wonder in heaven a woman clothed with the sun, and the moon under her feet, \Star{} And upon her head a crown of twelve stars.\end{Resp}
\begin{Resp}\Vsign Unto thine head shall be given an ornament of grace, and a crown of glory shall cover thee.\end{Resp}
\begin{Resp}\Rsign And upon her head a crown of twelve stars.\end{Resp}
\begin{Resp}\Vsign Glory be to the Father, and to the Son, \Star{} and to the Holy Ghost.\end{Resp}
\begin{Resp}\Rsign And upon her head a crown of twelve stars.\end{Resp}
\switchcolumn*

\LA
\LessonHead{Lectio VII}
\switchcolumn
\EN
\LessonHead{Lesson VII}
\switchcolumn*

\LA
\LessonTitle{Léctio sancti Evangélii secúndum Lucam}
\switchcolumn
\EN
\LessonTitle{From the Holy Gospel according to Luke}
\switchcolumn*

\LA
\Cite{Luc 1:26-38}
\switchcolumn
\EN
\Cite{Luke 1:26-38}
\switchcolumn*

\LA
\noindent In illo témpore: Missus est Angelus Gábriel a Deo in civitátem Galilǽæ, cui nomen Názareth, ad Vírginem desponsátam viro, cui nomen erat Joseph, de domo David, et nomen Vírginis María. Et réliqua.\par
\switchcolumn
\EN
\noindent At that time the angel Gabriel was sent from God into a city of Galilee, called Nazareth, to a virgin espoused to a man whose name was Joseph, of the house of David; and the virgin's name was Mary. And so on.\par
\switchcolumn*

\LA
\LessonTitle{Homilía sancti Bernárdi Abbátis}
\switchcolumn
\EN
\LessonTitle{Homily by St. Bernard, Abbot (of Clairvaux)}
\switchcolumn*

\LA
\Source{Sermo de S. Maria}
\switchcolumn
\EN
\Source{Homily on Holy Mary}
\switchcolumn*

\LA
\noindent Ad commendatiónem grátiæ suæ et ad destructiónem humánæ sapiéntiæ, Deus de fémina sed vírgine dignátus est carnem assúmere, ut símilem símili rédderet, contrárium contrário curáret, pestíferam spinam evélleret, peccáti chirógraphum potentíssime deléret. Heva spina fuit; María rosa éxstitit. Heva spina, vulnerándo; María rosa, ómnium afféctus mulcéndo. Heva spina, infígens ómnibus mortem; María rosa, reddens salutíferam ómnibus sortem. María rosa fuit cándida per virginitátem, rubicúnda per caritátem; cándida carne, rubicúnda mente; cándida virtútem sectándo, rubicúnda vítia calcándo; cándida afféctum purificándo, rubicúnda carnem mortificándo; cándida Deum diligéndo, rubicúnda próximo compatiéndo.\par
\switchcolumn
\EN
\noindent To commend His Own love towards us, and to bring to nought the wisdom of men, God was pleased to take flesh of a woman, albeit a virgin, that He might bring like against like, heal by opposites, pluck out the poisonous thorn, and blot out mightily the handwriting of our sin that was against us. Eve was a thorn, Mary is a rose. Eve is a thorn that pierceth, Mary is a rose that charmeth all the senses. Eve was a thorn that fixed death into all, Mary is a rose that bringeth health to all. Mary was a white rose through her virginity, and a red rose through her love. She was white in her flesh, red in her mind; white in that she followed the path of grace, red in that she trod down sin; white by the purity of her affections, red by the mortification of her body; white by her love for God, red by her compassion for her neighbour.\par
\switchcolumn*

\LA
\begin{Resp}\Rsign Ego quasi vitis fructificávi suavitátem odóris, \Star{} Et flores mei fructus honóris et honestátis.\end{Resp}
\begin{Resp}\Vsign Ego mater pulchræ dilectiónis, et timóris, et agnitiónis, et sanctæ spei.\end{Resp}
\begin{Resp}\Rsign Et flores mei fructus honóris et honestátis.\end{Resp}
\switchcolumn
\EN
\begin{Resp}\Rsign As the vine brought I forth pleasant savour; \Star{} And my flowers are the fruit of honour and seemliness.\end{Resp}
\begin{Resp}\Vsign I am the mother of fair love, and fear, and knowledge, and holy hope.\end{Resp}
\begin{Resp}\Rsign And my flowers are the fruit of honour and seemliness.\end{Resp}
\switchcolumn*

\LA
\LessonHead{Lectio VIII}
\switchcolumn
\EN
\LessonHead{Lesson VIII}
\switchcolumn*

\LA
\Source{Sermo de Aquæductu}
\switchcolumn
\EN
\Source{Homily on the water course}
\switchcolumn*

\LA
\noindent Verbum caro factum est, et hábitat jam in nobis. Hábitat in memória nostra, hábitat in cogitatióne, quia usque ad ipsam descéndit imaginatiónem. Quonam modo, inquis? Nimírum jacens in præsépio, in virgináli grémio cubans, in monte prǽdicans, in oratióne pernóctans, in cruce pendens, in morte pallens, liber inter mórtuos et in inférno ímperans, seu étiam tértia die resúrgens, et Apóstolis loca clavórum, victóriæ signa, demónstrans, novíssime coram eis cæli secréta conscéndens. Quid horum non vere, non pie, non sancte cogitátur? Quidquid horum cógito, Deum cógito; et per ómnia est Deus meus. Hæc ego meditári dixi sapiéntiam, et prudéntiam judicávi eructáre memóriam suavitátis; quóniam in hujuscémodi núcleis virga sacerdotális copiósa prodúxit, quam, in supérnis háuriens, ubérius nobis María refúdit. In supérnis plane et ultra Angelos, quæ Verbum ex ipso Patris corde suscépit.\par
\switchcolumn
\EN
\noindent The Word was made flesh, and dwelleth even now among us. He dwelleth in our memory. He dwelleth in our thought. He hath come down even unto our imagination; and how sayest thou doth he so? By lying in the manger, by nestling in His mother's breast, by preaching upon the mountain, by remaining all night in prayer to God, by hanging upon the Cross, by turning pale in death, by going down free among the dead and triumphing in hell, by rising again the third day, by showing to the Apostles the places of the nails the marks of his victory, by ascending up into heaven while they all beheld Him, of which of these things think we not with truth, with godliness, with holiness? If I think of any of these, I think of God, and He is my God through them all. To think of these things I have decreed to be wisdom, and to set forth the memory of their sweetness I have judged to be prudence. The rod of Aaron the Priest brought forth buds, and bloomed blossoms, and yielded almonds; but these things are the almonds of that Rod which came forth out of the stem of Jesse, the Rod whereof sprang the flower, a Rod which was raised in Mary into places higher than the earthly tabernacle, higher indeed, even into places higher than angels, since she received the Word into herself out of the very heart of the Eternal Father.\par
\switchcolumn*

\LA
\begin{Resp}\Rsign Surge, própera, amíca mea; jam enim hiems tránsiit, imber ábiit et recéssit, \Star{} Flores apparuérunt in terra nostra.\end{Resp}
\begin{Resp}\Vsign Dóminus dabit benignitátem, et terra nostra dabit fructum suum.\end{Resp}
\begin{Resp}\Rsign Flores apparuérunt in terra nostra.\end{Resp}
\begin{Resp}\Vsign Glória Patri, et Fílio, \Star{} et Spirítui Sancto.\end{Resp}
\begin{Resp}\Rsign Flores apparuérunt in terra nostra.\end{Resp}
\switchcolumn
\EN
\begin{Resp}\Rsign Rise up, my love, and make haste. For the winter is past; the rain is over and gone \Star{} The flowers appear on our earth.\end{Resp}
\begin{Resp}\Vsign The Lord shall give that which is good, and our land shall yield her increase.\end{Resp}
\begin{Resp}\Rsign The flowers appear on our earth.\end{Resp}
\begin{Resp}\Vsign Glory be to the Father, and to the Son, \Star{} and to the Holy Ghost.\end{Resp}
\begin{Resp}\Rsign The flowers appear on our earth.\end{Resp}
\switchcolumn*

\end{paracol}

\par\needspace{10\baselineskip}\addvspace{1.2em}{\centering\small\textsc{Die 7 Octobris} · \textsc{7 October}\par}
\Office{S. Marci Papæ Confessoris}{St. Mark, Pope and Confessor}{Simplex}
\addcontentsline{toc}{subsection}{Die 7 Octobris · S. Marci Papæ Confessoris \textit{— St. Mark, Pope and Confessor}}
\begin{paracol}{2}
\LA
\LessonHead{Lectio IX (pro commemoratione)}
\switchcolumn
\EN
\LessonHead{Lesson IX (when commemorated)}
\switchcolumn*

\LA
\LessonTitle{Commemoratio: S. Marci Papæ Confessoris}
\switchcolumn
\EN
\LessonTitle{Commemoration of: St. Mark, Pope and Confessor}
\switchcolumn*

\LA
\noindent Marcus, Romanus, Constantino Magno imperatóre Pontifex, instituit ut epíscopus Ostiénsis, a quo Romanus Pontifex consecrátur, pallio uterétur. Duas Romæ basilicas ædificávit, álteram in Urbe, álteram via Ardeatina; quas Constantinus, auctas, magnis munéribus exornávit. Vixit in pontificatu menses octo, sepultúsque est in cœmeterio Balbinæ.\par
\switchcolumn
\EN
\noindent This Mark was a Roman, who sat as Pope in the reign of the Emperor Constantine the Great. He ordained that the Bishop of Ostia, by whom the Bishop of Rome is consecrated, should use the Pallium. He built two Churches, one in the city and the other on the Ardeatine Way, which Constantine enlarged and richly gifted. Mark lived as Pope eight months, and was buried in the Cemetery of Balbina.\par
\switchcolumn*

\LA
\TeDeum{Te Deum}
\switchcolumn
\EN
\TeDeum{Te Deum}
\switchcolumn*

\end{paracol}

\par\needspace{10\baselineskip}\addvspace{1.2em}{\centering\small\textsc{Die 8 Octobris} · \textsc{8 October}\par}
\Office{S. Birgittæ Viduæ}{St. Birgitta, Widow}{Duplex}
\addcontentsline{toc}{subsection}{Die 8 Octobris · S. Birgittæ Viduæ \textit{— St. Birgitta, Widow}}
\begin{paracol}{2}
\LA
\RefLine{Lectiones I–III: de Scriptura occurrente.}
\switchcolumn
\EN
\RefLine{Lessons I–III: from the Scripture of the day.}
\switchcolumn*

\LA
\RefLine{Lectiones VII–IX: ex Commúni non Virginum non Martyrum (C7a).}
\switchcolumn
\EN
\RefLine{Lessons VII–IX: from the Common of Widows (C7a).}
\switchcolumn*

\LA
\LessonHead{Lectio IV}
\switchcolumn
\EN
\LessonHead{Lesson IV}
\switchcolumn*

\LA
\noindent Birgitta, in Suecia illustribus et piis paréntibus orta, sanctíssime vixit. Cum adhuc in útero gestarétur, a naufragio propter eam mater erepta est. Decennis, post audítum de passióne Dómini sermónem, sequénti nocte Jesum in cruce, recénti sánguine perfusum, vidit, et de eadem passióne secum loquentem. Quo ex témpore in ejusdem meditatióne ita afficiebátur, ut de ea sine lácrimis cogitare deinceps numquam posset.\par
\switchcolumn
\EN
\noindent Bridget was the daughter of princely and godly parents, and was born in Sweden, \textasciitilde{}(in the year of our Lord 1304.) Her life was a very holy one. When she was still in the womb, her mother was for her sake saved from shipwreck. When she was ten years old, she heard a sermon upon the sufferings of the Lord, and the following night she saw Jesus on the Cross, covered with fresh Blood, and heard Him speaking to her of His same sufferings. From that time forth the thought of them touched her so keenly, that she could never again call them to mind without weeping.\par
\switchcolumn*

\LA
\begin{Resp}\Rsign Propter veritátem, et mansuetúdinem, et justítiam: \Star{} Et dedúcet te mirabíliter déxtera tua.\end{Resp}
\begin{Resp}\Vsign Spécie tua et pulchritúdine tua inténde, próspere procéde, et regna.\end{Resp}
\begin{Resp}\Rsign Et dedúcet te mirabíliter déxtera tua.\end{Resp}
\switchcolumn
\EN
\begin{Resp}\Rsign Because of truth, and meekness, and righteousness; \Star{} And thy right hand shall lead thee wonderfully.\end{Resp}
\begin{Resp}\Vsign In thy comeliness, and thy beauty, go forward, fare prosperously, and reign.\end{Resp}
\begin{Resp}\Rsign And thy right hand shall lead thee wonderfully.\end{Resp}
\switchcolumn*

\LA
\LessonHead{Lectio V}
\switchcolumn
\EN
\LessonHead{Lesson V}
\switchcolumn*

\LA
\noindent Ulfoni, Nericiæ principi, in matrimónium tradita, virum ipsum ad pietátis officia, tum optimis exemplis, tum efficacibus verbis, adhortáta est. In filiórum educatióne piíssima; paupéribus, et maxime infirmis, domo ad id muneris dicata, inserviébat quam diligentíssime, illórum pedes sólita lavare et osculari. Cum autem una cum viro suo rediret Compostella, ubi sancti Jacobi Apóstoli sepúlcrum visitaverant, et Atrébati Ulfo gráviter ægrotaret, sanctus Dionysius Birgittæ noctu appáruit, et de mariti salúte aliisque de rebus, quæ futuræ erant, præmonuit.\par
\switchcolumn
\EN
\noindent When she was sixteen years of age she was given in marriage to Ulpho, Prince of Nericia. She moved her husband to godly works, as well by her noble example as by her earnest words. She expended the most motherly care upon the up-bringing of her children. She opened an hospital, in which she carefully tended the poor, especially the sick, and would wash and kiss their feet. She made a pilgrimage with her husband to Compostella, to visit the grave of the holy Apostle James. On their way back Ulpho fell grievously sick at Arras, and St. Denys appeared in the night to Bridget, to tell her as well that her husband would be healed, as diverse other things to come.\par
\switchcolumn*

\LA
\begin{Resp}\Rsign Dilexísti justítiam, et odísti iniquitátem: \Star{} Proptérea unxit te Deus, Deus tuus, óleo lætítiæ.\end{Resp}
\begin{Resp}\Vsign Propter veritátem, et mansuetúdinem, et justítiam.\end{Resp}
\begin{Resp}\Rsign Proptérea unxit te Deus, Deus tuus, óleo lætítiæ.\end{Resp}
\switchcolumn
\EN
\begin{Resp}\Rsign Thou hast loved righteousness, and hated iniquity; \Star{} Therefore God, thy God, hath anointed thee with the oil of gladness.\end{Resp}
\begin{Resp}\Vsign Because of truth, and meekness, and righteousness.\end{Resp}
\begin{Resp}\Rsign Therefore God, thy God, hath anointed thee with the oil of gladness.\end{Resp}
\switchcolumn*

\LA
\LessonHead{Lectio VI}
\switchcolumn
\EN
\LessonHead{Lesson VI}
\switchcolumn*

\LA

\switchcolumn
\EN
\LessonTitle{(In the year 1344) her husband died, after having become a Cistercian monk. Bridget, having heard the voice of Christ in a dream, took upon herself an harder way of life. During her life God made known to her many hidden things. She founded the monastery of Wastein, under the Rule of the Holy Saviour, a Rule which she had received from the Lord Himself. By the command of God she went to Rome, where she stirred up many by her example to seek the love of God. Thence she went to Jerusalem, and then returned again to Rome. From this pilgrimage she caught a fever, of which she lay sick an whole year in sharp sufferings, and then, laden with good works, and after foretelling the day of her own death, she departed from earth to heaven, (upon the 23rd day of July, in the year 1373.) Her body was taken to the monastery of Wastein. She was famous for miracles, and Boniface IX enrolled her name among those of the Saints.}
\switchcolumn*

\LA
\noindent Viro Cisterciénsi monacho facto et paulo post defuncto, Birgitta, audíta Christi voce in somnis, arctiórem vitæ formam est aggressa. Cui deínde arcana multa fuérunt divínitus revelata. Monastérium Vastanense sub regula sancti Salvatoris, ab ipso Dómino accepta, instituit. Romam Dei jussu venit, ubi plurimos ad amórem divinum veheménter accendit. Inde Jerosolymam pétiit, et íterum Romam. Qua ex peregrinatióne cum in febrim incidísset, gravibus per annum integrum afflictáta morbis, cumuláta meritis, prænuntiáto mortis die, migrávit in cælum. Corpus ejus ad Vastanense monastérium translátum est; et miraculis illustrem Bonifatius nonus in Sanctórum númerum retulit.\par
\switchcolumn
\EN

\switchcolumn*

\LA
\begin{Resp}\Rsign Fallax grátia, et vana est pulchritúdo: \Star{} Múlier timens Dóminum, ipsa laudábitur.\end{Resp}
\begin{Resp}\Vsign Date ei de fructu mánuum suárum, et laudent eam in portis ópera ejus.\end{Resp}
\begin{Resp}\Rsign Múlier timens Dóminum, ipsa laudábitur.\end{Resp}
\begin{Resp}\Vsign Glória Patri, et Fílio, \Star{} et Spirítui Sancto.\end{Resp}
\begin{Resp}\Rsign Múlier timens Dóminum, ipsa laudábitur.\end{Resp}
\switchcolumn
\EN
\begin{Resp}\Rsign Favour is deceitful, and beauty is vain \Star{} A woman that feareth God she shall be praised.\end{Resp}
\begin{Resp}\Vsign Give her of the fruit of her hands, and let her own works praise her in the gates.\end{Resp}
\begin{Resp}\Rsign A woman that feareth God she shall be praised.\end{Resp}
\begin{Resp}\Vsign Glory be to the Father, and to the Son, \Star{} and to the Holy Ghost.\end{Resp}
\begin{Resp}\Rsign A woman that feareth God she shall be praised.\end{Resp}
\switchcolumn*

\end{paracol}

\par\needspace{10\baselineskip}\addvspace{1.2em}{\centering\small\textsc{Die 9 Octobris} · \textsc{9 October}\par}
\Office{S. Joannis Leonardi Confessoris}{St. John Leonard, Confessor}{Duplex}
\addcontentsline{toc}{subsection}{Die 9 Octobris · S. Joannis Leonardi Confessoris \textit{— St. John Leonard, Confessor}}
\begin{paracol}{2}
\LA
\RefLine{Lectiones I–III: de Scriptura occurrente.}
\switchcolumn
\EN
\RefLine{Lessons I–III: from the Scripture of the day.}
\switchcolumn*

\LA
\RefLine{Lectio IX: de S. Dionysii Epíscopi, Rustici et Eleutherii Martyrum (10-09t).}
\switchcolumn
\EN
\RefLine{Lesson IX: of Sts. Denis, Bishop, Rusticus et Eleutherius, Martyrs (10-09t).}
\switchcolumn*

\LA
\LessonHead{Lectio IV}
\switchcolumn
\EN
\LessonHead{Lesson IV}
\switchcolumn*

\LA
\noindent Joánnes Leonardi, in oppido Décimi, non longe a Lucénsi urbe, piis et honestis ortus paréntibus, jam inde a prima ætate, solitúdinis et precatiónis amóre, grave quiddam ac matúrum præ se tulit. Annos natus viginti sex a Deo vocátus ad ecclesiásticæ milítiæ nomen dandum, sæcularibus curis íllico nuntium remisit. Ac primo inter púeros latinæ linguæ rudimentis instructus, deínde in litteris et philosophicis ac theologicis disciplinis adeo profecit, ut vix acto quadriennio ad sacerdotium ex obediéntia promotus fúerit. Mox aliquot nactus bonæ índolis nóbiles júvenes, cum eos ad virtútis perfectiónem sedulo exercuísset, insequénti anno congregatiónem instituit clericórum regulárium, quam a Matre Dei, ob incensum erga ipsam suæ devotiónis afféctum, nuncupávit. Horum cura et zelo tanta peracta est animórum commutátio, ut cum in Lucénsi republica, hæreticórum præsertim perfidiosis artibus, ardérent civium odia, profligatique essent mores, brevi témpore primæva Christianórum pietas ibidem revixisse viderétur.\par
\switchcolumn
\EN
\noindent John Leonard was born of pious and respectable parents in the town of Diecimo, not far from the city of Lucca. From very early boyhood he shewed himself mature and serious, with an inclination to solitude and prayer. When he was twenty-six years old God called him to enlist among the soldiers of the Church. John renounced immediately all his worldly interests. At first he had to study elementary Latin with little boys, but he soon advanced in a knowledge of literature, philosophy and theology. After a scant four years, at the command of his superior, he was ordained to the priesthood. Soon afterward he and a group of noble youths, alike inflamed with high ideals, earnestly set about attaining perfection in virtue. The following year they formed the Congregation of the Clerks Regular of the Mother of God, a name chosen because of their intense devotion to her. John and his companions laboured with such diligence in their care of souls, that before long a change of attitude was brought about. In the city state of Lucca, where through the perfidious wiles of the heretics, hateful passion turned fiercely among the citizens, where morals were corrupted, in a very short space of time the primitive piety of the Christians seemed to revive.\par
\switchcolumn*

\LA
\begin{Resp}\Rsign Honéstum fecit illum Dóminus, et custodívit eum ab inimícis, et a seductóribus tutávit illum: \Star{} Et dedit illi claritátem ætérnam.\end{Resp}
\begin{Resp}\Vsign Justum dedúxit Dóminus per vias rectas, et osténdit illi regnum Dei.\end{Resp}
\begin{Resp}\Rsign Et dedit illi claritátem ætérnam.\end{Resp}
\switchcolumn
\EN
\begin{Resp}\Rsign The Lord made him honourable, and defended him from his enemies, and kept him safe from those that lay in wait for him; \Star{} And gave him perpetual glory.\end{Resp}
\begin{Resp}\Vsign He went down with him into the pit, and left him not in bonds.\end{Resp}
\begin{Resp}\Rsign And gave him perpetual glory.\end{Resp}
\switchcolumn*

\LA
\LessonHead{Lectio V}
\switchcolumn
\EN
\LessonHead{Lesson V}
\switchcolumn*

\LA
\noindent Tam salutárium óperum causa incídit Joánnes in acerrimas insectatiónes hóminum nequam, qui recens coactam familiam pérdere omni ope conáti sunt. Sed vir Dei, æquo animo libenter ómnia ferens, impetráta a Summo Pontifice Gregório décimo tertio suæ congregatiónis confirmatióne, apostolici sui labóris fructus constanter servávit. In arduis negotiis componéndis multi epíscopi eo consiliario et adjutore usi sunt, et vel ipse Romanus Pontifex eum delegávit ad intricáta litigia dirimenda, et ad religiosas familias reformandas. Sancto Josepho Calasanctio, ejusque pene collapsæ societati, præsto fuit. Haud levem quoque impéndit operam negotiis nosocomíi Sancti Spíritus in Saxia, et moniálibus oblátis sanctæ Franciscæ Romanæ excolendis.\par
\switchcolumn
\EN
\noindent In his work for the salvation of souls John met most bitter insults from wicked men who tried in every way to destroy the newly gathered family. But the man of God, bearing all things cheerfully and serenely, defended pertinaciously the fruit of his apostolic labours by securing from the Supreme Pontiff, Gregory XIII, papal approbation of his Congregation. Many bishops about to undertake difficult enterprises sought his advice and aid. Even the Holy Father delegated to him the solution of intricate litigation and the reform of religious societies. He stood in support of Saint Joseph Calasanctius when his society was on the verge of collapse. Scarcely less arduous were the hours John devoted to the affairs of the Hospital of the Holy Spirit in the English section of Rome, and to those of the convent of Saint Frances of Rome.\par
\switchcolumn*

\LA
\begin{Resp}\Rsign Amávit eum Dóminus, et ornávit eum: stolam glóriæ índuit eum, \Star{} Et ad portas paradísi coronávit eum.\end{Resp}
\begin{Resp}\Vsign Índuit eum Dóminus lorícam fídei, et ornávit eum.\end{Resp}
\begin{Resp}\Rsign Et ad portas paradísi coronávit eum.\end{Resp}
\switchcolumn
\EN
\begin{Resp}\Rsign The Lord loved him and beautified him; He clothed him with a robe of glory, \Star{} And crowned him at the gates of Paradise.\end{Resp}
\begin{Resp}\Vsign The Lord hath put on him the breast-plate of faith, and hath adorned him.\end{Resp}
\begin{Resp}\Rsign And crowned him at the gates of Paradise.\end{Resp}
\switchcolumn*

\LA
\LessonHead{Lectio VI}
\switchcolumn
\EN
\LessonHead{Lesson VI}
\switchcolumn*

\LA
\noindent Gráviter dolens, gentes adeo plurimas remotis in regiónibus luce Evangélii carére, inflammabátur desidério migrándi in illas oras ad lumen veræ religiónis effundéndum. At cum intellixísset a sancto Philippo Nerio, a quo verus reformator dicebátur, se suamque congregatiónem ad instituéndos Italiæ pópulos destinari, divinæ acquiévit voluntáti ; minime tamen abstinuit quin, si aliquam infidelibus opem afferre posset, experirétur. Hinc ínitis consíliis cum piíssimo præsule Vives, cœtum instituit presbyterórum, quibus propositum esset idoneos informare adolescéntes, in díssitas regiónes subinde mittendos ad fidem propagandam. Quare mérito véluti auctor censétur præclaríssimi illíus instituti, quod Summórum Pontificum ópera amplificátum, proferendæ per univérsum orbem catholicæ fidei mirabíliter inservit. Plura ópera de re sacra et morali conscripsit, cuívis hóminum conditióni accommodatíssima. Denique a sacro ministerio numquam deficiens, in cínere et cilício ad Dóminum migrávit Romæ, die nona Octobris, anno millesimo sexcentésimo nono, ætátis sexagesimo sexto. Quem sanctitátis et miraculis illustrem Pius nonus Pontifex Maximus Beatórum fastis accensuit. Pius vero undecimus, anno millesimo nongentésimo trigesimo octavo, die solemni Paschæ, inter Sanctos adscripsit.\par
\switchcolumn
\EN
\noindent Greatly saddened that so many peoples in far distant places were without the light of the Gospel, John burned with a desire to journey to those countries to spread the light of the true Faith. But when Saint Philip Neri, who called John a true reformer, showed him that he and his Congregation were destined to educate the Italian people, John acquiesced to the will of God. He did not, however, refrain so completely that he did not try to do some work for the infidels. He is therefore, very rightly credited along with the pioneer Vives with being the founder of the movement among the bishops to send well-qualified young men to distant, alien lands to propagate the faith. Wherefore he is very properly regarded as the author of that most illustrious institute which augments the work of the Sovereign Pontiffs and serves to spread the Catholic faith throughout the world. John wrote many works on theology and morality, well adapted to the men of that day. Finally, in sackcloth and ashes, lacking nothing in his sacred ministry, he passed to the Lord in Rome on the 9th day of October, 1609, at the age of sixty-six. He was so famous for sanctity and miracles that Pius IX, the Supreme Pontiff, named him on the calendar of the Blessed. In 1938 on the solemn Feast of Easter, Pius XI enrolled him among the Saints.\par
\switchcolumn*

\LA
\begin{Resp}\Rsign Iste homo perfécit ómnia quæ locútus est ei Deus, et dixit ad eum: Ingrédere in réquiem meam: \Star{} Quia te vidi justum coram me ex ómnibus géntibus.\end{Resp}
\begin{Resp}\Vsign Iste est, qui contémpsit vitam mundi, et pervénit ad cæléstia regna.\end{Resp}
\begin{Resp}\Rsign Quia te vidi justum coram me ex ómnibus géntibus.\end{Resp}
\begin{Resp}\Vsign Glória Patri, et Fílio, \Star{} et Spirítui Sancto.\end{Resp}
\begin{Resp}\Rsign Quia te vidi justum coram me ex ómnibus géntibus.\end{Resp}
\switchcolumn
\EN
\begin{Resp}\Rsign This is he which did according unto all that God commanded him; and God said unto him: Enter thou into My rest, \Star{} For thee have I seen righteous before Me among all people.\end{Resp}
\begin{Resp}\Vsign This is he which loved not his life in this world, and is come unto an everlasting kingdom.\end{Resp}
\begin{Resp}\Rsign For thee have I seen righteous before Me among all people.\end{Resp}
\begin{Resp}\Vsign Glory be to the Father, and to the Son, \Star{} and to the Holy Ghost.\end{Resp}
\begin{Resp}\Rsign For thee have I seen righteous before Me among all people.\end{Resp}
\switchcolumn*

\LA
\LessonHead{Lectio VII}
\switchcolumn
\EN
\LessonHead{Lesson VII}
\switchcolumn*

\LA
\LessonTitle{Léctio sancti Evangélii secúndum Lucam}
\switchcolumn
\EN
\LessonTitle{From the Holy Gospel according to Luke}
\switchcolumn*

\LA
\Cite{Luc 10:1-9}
\switchcolumn
\EN
\Cite{Luke 10:1-9}
\switchcolumn*

\LA
\noindent In illo témpore: Designávit Dóminus et álios septuagínta duos: et misit illos binos ante fáciem suam, in omnem civitátem et locum, quo erat ipse ventúrus. Et réliqua.\par
\switchcolumn
\EN
\noindent At that time, the Lord appointed other seventy-two also, and sent them two and two before His face into every city and place, whither He Himself would come. And so on.\par
\switchcolumn*

\LA
\LessonTitle{Homilía sancti Gregórii Papæ}
\switchcolumn
\EN
\LessonTitle{Homily by Pope St. Gregory (the Great.)}
\switchcolumn*

\LA
\Source{Homilia 17 in Evangelia}
\switchcolumn
\EN
\Source{17th Homily on the Gospels}
\switchcolumn*

\LA
\noindent Dóminus et Salvátor noster, fratres caríssimi, aliquándo nos sermónibus, aliquándo vero opéribus ádmonet. Ipsa étenim facta ejus præcépta sunt: quia dum áliquid tácitus facit, quid ágere debeámus innotéscit. Ecce enim binos in prædicatiónem discípulos mittit: quia duo sunt præcépta caritátis, Dei vidélicet amor, et próximi: et minus quam inter duos cáritas habéri non potest. Nemo enim próprie ad semetípsum habére caritátem dícitur: sed diléctio in álterum tendit, ut cáritas esse possit.\par
\switchcolumn
\EN
\noindent Dearly beloved brethren, our Lord and Saviour doth sometimes admonish us by words, and sometimes by works. Yea, His very works do themselves teach us for that which He doth silently His example still moveth us to copy. Behold how He sendeth forth His disciples to preach by two and two since there are two commandments to love, that is, a commandment to love God, and a commandment to love our neighbour and where there are not two, the one, being alone, hath not whereon to do the Lord's commandment. And no man can properly be said to love himself: for love tendeth outward toward our neighbour, if it be the love whereto the Gospel doth oblige us.\par
\switchcolumn*

\LA
\begin{Resp}\Rsign Iste est qui ante Deum magnas virtútes operátus est, et de omni corde suo laudávit Dóminum: \Star{} Ipse intercédat pro peccátis ómnium populórum.\end{Resp}
\begin{Resp}\Vsign Ecce homo sine queréla, verus Dei cultor, ábstinens se ab omni ópere malo, et pérmanens in innocéntia sua.\end{Resp}
\begin{Resp}\Rsign Ipse intercédat pro peccátis ómnium populórum.\end{Resp}
\switchcolumn
\EN
\begin{Resp}\Rsign This is he which wrought great wonders before God, and praised the Lord with all his heart. \Star{} May he pray for all people, that their sins may be forgiven unto them.\end{Resp}
\begin{Resp}\Vsign Behold a man without blame, a worshipper of God in truth, keeping himself clean from every evil work, and abiding still in his innocency.\end{Resp}
\begin{Resp}\Rsign May he pray for all people, that their sins may be forgiven unto them.\end{Resp}
\switchcolumn*

\LA
\LessonHead{Lectio VIII}
\switchcolumn
\EN
\LessonHead{Lesson VIII}
\switchcolumn*

\LA
\noindent Ecce enim binos ad prædicándum discípulos Dóminus mittit: quátenus hoc nobis tácitus ínnuat, quia qui caritátem erga álterum non habet, prædicatiónis offícium suscípere nullátenus debet. Bene autem dícitur, quia misit eos ante fáciem suam in omnem civitátem et locum, quo erat ipse ventúrus. Prædicatóres enim suos Dóminus séquitur: quia prædicátio prǽvenit, et tunc ad mentis nostræ habitáculum Dóminus venit, quando verba exhortatiónis præcúrrunt: atque per hoc véritas in mente suscípitur.\par
\switchcolumn
\EN
\noindent Behold, the Lord sendeth forth His disciples to preach by two and two and thus doing, He doth silently teach us that whosoever loveth not his neighbour, such a one it behoveth not to take upon him the office of a preacher. Well also is it said that He sent them before His face into every city and place whither He Himself would come. The Lord followeth His preachers first cometh preaching, and then the Lord Himself cometh to the house of our mind, whither the word of exhortation hath come before and so cometh the truth into our mind.\par
\switchcolumn*

\LA
\begin{Resp}\Rsign Sint lumbi vestri præcíncti, et lucérnæ ardéntes in mánibus vestris: \Star{} Et vos símiles homínibus exspectántibus dóminum suum, quando revertátur a núptiis.\end{Resp}
\begin{Resp}\Vsign Vigiláte ergo, quia nescítis qua hora Dóminus vester ventúrus sit.\end{Resp}
\begin{Resp}\Rsign Et vos símiles homínibus exspectántibus dóminum suum, quando revertátur a núptiis.\end{Resp}
\begin{Resp}\Vsign Glória Patri, et Fílio, \Star{} et Spirítui Sancto.\end{Resp}
\begin{Resp}\Rsign Et vos símiles homínibus exspectántibus dóminum suum, quando revertátur a núptiis.\end{Resp}
\switchcolumn
\EN
\begin{Resp}\Rsign Let your loins be girded about, and your lights burning; \Star{} And ye yourselves like unto men that wait for their lord, when he will return from the wedding.\end{Resp}
\begin{Resp}\Vsign Watch therefore, for ye know not what hour your Lord doth come.\end{Resp}
\begin{Resp}\Rsign And ye yourselves like unto men that wait for their lord, when he will return from the wedding.\end{Resp}
\begin{Resp}\Vsign Glory be to the Father, and to the Son, \Star{} and to the Holy Ghost.\end{Resp}
\begin{Resp}\Rsign And ye yourselves like unto men that wait for their lord, when he will return from the wedding.\end{Resp}
\switchcolumn*

\end{paracol}

\par\needspace{10\baselineskip}\addvspace{1.2em}{\centering\small\textsc{Die 9 Octobris} · \textsc{9 October}\par}
\Office{S. Dionysii Epíscopi, Rustici et Eleutherii Martyrum}{Sts. Denis, Bishop, Rusticus et Eleutherius, Martyrs}{Semiduplex}
\addcontentsline{toc}{subsection}{Die 9 Octobris · S. Dionysii Epíscopi, Rustici et Eleutherii Martyrum \textit{— Sts. Denis, Bishop, Rusticus et Eleutherius, Martyrs}}
\begin{paracol}{2}
\LA
\LessonHead{Lectio IX (pro commemoratione)}
\switchcolumn
\EN
\LessonHead{Lesson IX (when commemorated)}
\switchcolumn*

\LA
\LessonTitle{Commemoratio: S. Dionysii Epíscopi, Rustici et Eleutherii Martyrum}
\switchcolumn
\EN
\LessonTitle{Commemoration of: Sts. Denis, Bishop, Rusticus et Eleutherius, Martyrs}
\switchcolumn*

\LA
\noindent Dionysius Atheniénsis, unus ex Areopagítis judícibus, cum adhuc in gentilitátis erróre versarétur, eo die, quo Christus Dóminus Cruci affixus est, solem præter natúram defécisse animadvertens, exclamasse tráditur: Aut Deus natúræ pátitur, aut mundi máchina dissolvitur. Cum autem Paulus Apóstolus in Areopágo Christum annuntiásset, Dionysius fidem christianam amplexus, ab eodem Apóstolo Atheniensium ecclésiæ præfectus est. Postea, ut traditur, Romam veniens, et a Cleménte Pontifice missus in Gálliam, Lutétiam usque Parisiórum, cum Rustico presbytero et Eleutherio diacono, Evangélium prædicávit. Ibi omnes, quod Christum prædicarent, a Fescennio præfecto apprehénsi, variis tormentis cruciántur, et demum secúri feriúntur septimo Idus Octobris.\par
\switchcolumn
\EN
\noindent Denis was an Athenian, one of the Judges of the Court of the Areopagus, and a man of varied and deep learning. There is a story concerning him that on the day when the Lord Christ was nailed to the Cross, and when he saw the unnatural eclipse of the sun, Denis said: Either the God of nature is suffering, or the frame-work of the world is breaking up. When the Apostle Paul came to Athens, and was taken and brought unto the Areopagus, and gave an account of the faith which he preached, affirming that Christ had risen from the dead, and that all the dead likewise are to live again, some mocked, and others said We will hear thee again of this matter. So Paul departed from among them. Howbeit, certain men clave unto him, and believed among the which was Dionysius the Areopagite. \textasciitilde{}(Acts xvii. 3234) Denis was baptized by the Apostle, and set over the Church of the Athenians. He came afterwards to Rome, and was sent by Pope Clement into Gaul, to preach the Gospel. There followed him to Paris one Rusticus a Priest, and Eleutherius a Deacon. He turned many to Christ, and was therefore hided with rods by command of Fescennius the Praefect, and, forasmuch as he still went on bravely preaching Christ, he was tortured with fire upon a grating, and put to diverse other torments, and his comrades likewise. They bore their torments bravely and cheerfully, and then Denis, being over an hundred years of age, and his two comrades with him, suffered by the axe upon the 9th day of October. This is that Denis concerning whom the old story is told that after his head was cut off he took it in his hands and walked two thousand paces, carrying it all the while. He was the author of some marvelous books, clear proofs of a mind fixed in heaven, upon The Names of God, upon The Orders in Heaven and in the Church, upon The Mystic Theology, and diverse others.\par
\switchcolumn*

\LA
\TeDeum{Te Deum}
\switchcolumn
\EN
\TeDeum{Te Deum}
\switchcolumn*

\end{paracol}

\par\needspace{10\baselineskip}\addvspace{1.2em}{\centering\small\textsc{Die 10 Octobris} · \textsc{10 October}\par}
\Office{S. Francisci Borgiæ Confessoris}{St. Francis Borgia, Confessor}{Semiduplex}
\addcontentsline{toc}{subsection}{Die 10 Octobris · S. Francisci Borgiæ Confessoris \textit{— St. Francis Borgia, Confessor}}
\begin{paracol}{2}
\LA
\RefLine{Lectiones I–III: de Scriptura occurrente.}
\switchcolumn
\EN
\RefLine{Lessons I–III: from the Scripture of the day.}
\switchcolumn*

\LA
\LessonHead{Lectio IV}
\switchcolumn
\EN
\LessonHead{Lesson IV}
\switchcolumn*

\LA
\noindent Franciscus, Gandiæ dux quartus, Joánne Borgia et Joánna Aragónia Ferdinándi Catholici nepte génitus, post puerílem ætátem inter domesticos mira innocéntia et pietáte transactam, in aula primum Cároli quinti Cæsaris, mox in Catalauniæ administratióne, admirabílior fuit christianæ virtútis et vitæ austerioris exemplis. Ad Granatense sepúlcrum Isabellam imperatricem cum detulísset, in ejus vultu, fœde commutato, mortalium ómnium caducitátem rélegens, voto se adstrinxit, rebus ómnibus, cum primum licéret, abjectis, regum Regi unice inserviéndi. Inde tantum virtútis increméntum fecit, ut, inter negotiórum turbas, religiosæ perfectiónis simillimam imáginem réferens, miraculum príncipum appellarétur.\par
\switchcolumn
\EN
\noindent Francis, fourth Duke of Gandia, was the son of John Borgia, Duke of Gandia, and of Joan of Aragon, daughter of Alphonso, natural son to Ferdinand V. surnamed the Catholic, King of Aragon. (He was born at Gandia, in the kingdom of Valencia, in the year of our Lord 1510.) He passed his boyhood at home in great innocence and godliness, and was still more remarkable for his Christian graces and the hardness of his living, at the Court of the Emperor Charles V., and as Vice-Roy of Catalonia. (On May the 1st, 1539) died the Empress Isabella, and Francis, (as her master of the horse,) was commanded to attend her body to Granada, where it was to be buried. (At Granada the coffin was opened, in order that he might swear to the magistrates of the city that it was indeed the body of the late Empress,) and the sight of the awful change which death had made in her countenance so thrilled him with the thought of our mortality and corruption, that he bound himself by vow, as soon as he lawfully might, to give up all things, and to serve the King of kings only. From that time he so advanced in Christian graces, that his life might be called the miracle of princes, showing, in the midst of a vast mass of business, an image of perfection attained in a cloister. greatly endeared him to Princes and Popes, and besides founding or enlarging very many houses in diverse places, he sent brethren into the kingdom of Poland, into the islands of the Ocean, and into the provinces of Mexico and Peru, and into other lands also Apostolic men who spread the Roman Catholic faith by their preaching, their sweat, and their blood.\par
\switchcolumn*

\LA
\begin{Resp}\Rsign Honéstum fecit illum Dóminus, et custodívit eum ab inimícis, et a seductóribus tutávit illum: \Star{} Et dedit illi claritátem ætérnam.\end{Resp}
\begin{Resp}\Vsign Justum dedúxit Dóminus per vias rectas, et osténdit illi regnum Dei.\end{Resp}
\begin{Resp}\Rsign Et dedit illi claritátem ætérnam.\end{Resp}
\switchcolumn
\EN
\begin{Resp}\Rsign The Lord made him honourable, and defended him from his enemies, and kept him safe from those that lay in wait for him; \Star{} And gave him perpetual glory.\end{Resp}
\begin{Resp}\Vsign He went down with him into the pit, and left him not in bonds.\end{Resp}
\begin{Resp}\Rsign And gave him perpetual glory.\end{Resp}
\switchcolumn*

\LA
\LessonHead{Lectio V}
\switchcolumn
\EN
\LessonHead{Lesson V}
\switchcolumn*

\LA
\noindent Mortua Eleonora de Castro cónjuge ingréssus est societátem Jesu, ut in ea latéret securius et præcluderet dignitátibus áditum, interposita voti religióne; dignus, quem et viri príncipes complures in amplecténdo severiori instituto fúerint secuti, et Cárolus quintus ipse in abdicando imperio hortatórem sibi aut ducem exstitisse non diffiterétur. In eo arctioris vitæ studio Franciscus jejuniis, catenis férreis, asperrimo cilício, cruentis longisque verberatiónibus, somno brevíssimo, corpus ad extremam usque maciem redégit, nullis præterea parcens labóribus ad sui victoriam et ad salútem animárum. Tot ítaque instructus virtútibus, a sancto Ignátio primum generalis commissarius in Hispaniis, nec multo post præpositus generalis tertius a societáte universa, licet invítus, elígitur. Quo in munere princípibus ac summis Pontificibus prudéntia ac morum sanctitáte apprime carus, præter complura vel cóndita vel aucta ubíque domicília, socios in regnum Poloniæ, in insulas Oceani, in Mexicanam et Peruánam provincias invéxit; missis quoque in alias regiónes apostolicis viris, qui prædicatióne, sudóribus, sánguine fidem catholicam Romanam propagarunt.\par
\switchcolumn
\EN
\noindent His wife, Eleanora de Castro, died (on the 27th of March, 1546,) and (in 1551) he entered the Society of Jesus, that therein he might hide himself more safely, and bar by the obligation of a vow the path to dignities. He was the worthy leader of many princes who have embraced a life of hardship, and Charles V. himself when he resigned the Empire did not deny that he had been moved and shown the way by Francis. In his struggle after austerity Francis, by fasting, by iron chains, by the roughest of hair -cloth, by long and bloody flagellations, and by denying himself any but very little sleep, reduced his body to the last degree, but would still spare no toil to overcome himself and to save souls. Thus full of ghostly strength, he was appointed by holy Ignatius, \textasciitilde{}(in the year 1554,) Commissary-General of the Society in Spain, Portugal, and the Indies, and (on the 2nd of July, 1565,) notwithstanding all the precautions he could take to prevent it, he was chosen by the general Congregation of the Society to be General, being the third who held that office. In this position his wisdom and holiness of life\par
\switchcolumn*

\LA
\begin{Resp}\Rsign Amávit eum Dóminus, et ornávit eum: stolam glóriæ índuit eum, \Star{} Et ad portas paradísi coronávit eum.\end{Resp}
\begin{Resp}\Vsign Índuit eum Dóminus lorícam fídei, et ornávit eum.\end{Resp}
\begin{Resp}\Rsign Et ad portas paradísi coronávit eum.\end{Resp}
\switchcolumn
\EN
\begin{Resp}\Rsign The Lord loved him and beautified him; He clothed him with a robe of glory, \Star{} And crowned him at the gates of Paradise.\end{Resp}
\begin{Resp}\Vsign The Lord hath put on him the breast-plate of faith, and hath adorned him.\end{Resp}
\begin{Resp}\Rsign And crowned him at the gates of Paradise.\end{Resp}
\switchcolumn*

\LA
\LessonHead{Lectio VI}
\switchcolumn
\EN
\LessonHead{Lesson VI}
\switchcolumn*

\LA
\noindent De se ita demisse sentiebat, ut peccatóris nomen sibi proprium fáceret. Romanam purpuram, a summis Pontificibus sæpius oblatam, invicta humilitátis constántia recusávit. Vérrere sordes, emendicáre victum ostiatim, ægris ministrare in nosocomíis, mundi ac sui contemptor, in deliciis hábuit. Singulis diébus multas continenter horas, frequenter octo, quandoque decem, dabat cæléstium contemplatióni. Centies quotídie de genu Deum adorábat. Numquam a sacrificando abstinuit, prodebatque sese divinus quo æstuábat ardor, ejus vultu, sacram Hostiam offeréntis aut concionántis, intérdum radiante. Sanctíssimum Christi corpus in Eucharistia latens ubi asservarétur, instinctu cælésti sentiebat. Cardinali Alexandrino, ad conjungendos contra Turcas, christiános príncipes, legato comes additus a beato Pio quinto, arduum iter, fractis jam pene viribus, suscépit ex obediéntia; in qua et vitæ cursum Romæ, ut optarat, feliciter consummávit, anno ætátis suæ sexagesimo secundo, salútis vero millesimo quingentésimo septuagesimo secundo. A sancta Teresia, quæ ejus utebátur consíliis, vir sanctus, a Gregório décimo tertio fidélis adminíster appellátus; demum a Cleménte décimo, plúribus magnisque clarus miraculis, in Sanctórum númerum est adscriptus.\par
\switchcolumn
\EN
\noindent He thought so little of himself that he gave himself the nickname of Francis the sinner. By the Popes he was oftentimes offered the dignity of Cardinal of the Roman Church, but the lowly firmness with which he refused it could never be overcome. In his cheap esteem of the world and of himself his chief pleasures were to clean the house, to beg for food from door to door, and to wait upon the sick in hospitals. He spent many hours every day, oftentimes eight and sometimes ten, in prayer and meditation. An hundred times every day he worshipped God upon his knees. He never missed the opportunity of offering the Holy Liturgy, and the fire from God which burnt within him sometimes shone forth in his countenance when he was lifting the Sacred Host, or preaching. By an inward power given him from God he could tell where the most Holy Body of Christ, under the Eucharistic veils, was kept. (In 1570, the year before the victory of Lepanto,) the blessed Pius V. sent Francis with the Cardinal Alexandrini on an embassy (into France, Spain, and Portugal,) to unite the Christian Princes against Turkey. His vital strength was then nearly worn out, but, through obedience, he undertook the toil of the journey. He became much worse during the travelling, and on his return brought to a blessed end at Rome, as had been his desire, the pilgrimage of this life, (a little before midnight between the last day of September and the first of October,) in the sixty-second year of his own life, and that of salvation 1572. Holy Teresa, who used his advice, called him an holy man, and Gregory XIII., a faithful servant. He was famous for many and great signs and wonders, and Clement X. at last numbered him among the Saints.\par
\switchcolumn*

\LA
\begin{Resp}\Rsign Iste homo perfécit ómnia quæ locútus est ei Deus, et dixit ad eum: Ingrédere in réquiem meam: \Star{} Quia te vidi justum coram me ex ómnibus géntibus.\end{Resp}
\begin{Resp}\Vsign Iste est, qui contémpsit vitam mundi, et pervénit ad cæléstia regna.\end{Resp}
\begin{Resp}\Rsign Quia te vidi justum coram me ex ómnibus géntibus.\end{Resp}
\begin{Resp}\Vsign Glória Patri, et Fílio, \Star{} et Spirítui Sancto.\end{Resp}
\begin{Resp}\Rsign Quia te vidi justum coram me ex ómnibus géntibus.\end{Resp}
\switchcolumn
\EN
\begin{Resp}\Rsign This is he which did according unto all that God commanded him; and God said unto him: Enter thou into My rest, \Star{} For thee have I seen righteous before Me among all people.\end{Resp}
\begin{Resp}\Vsign This is he which loved not his life in this world, and is come unto an everlasting kingdom.\end{Resp}
\begin{Resp}\Rsign For thee have I seen righteous before Me among all people.\end{Resp}
\begin{Resp}\Vsign Glory be to the Father, and to the Son, \Star{} and to the Holy Ghost.\end{Resp}
\begin{Resp}\Rsign For thee have I seen righteous before Me among all people.\end{Resp}
\switchcolumn*

\LA
\LessonHead{Lectio VII}
\switchcolumn
\EN
\LessonHead{Lesson VII}
\switchcolumn*

\LA
\LessonTitle{Léctio sancti Evangélii secúndum Matthǽum.}
\switchcolumn
\EN
\LessonTitle{From the Holy Gospel according to Matthew}
\switchcolumn*

\LA
\Cite{Matt 19:27-29}
\switchcolumn
\EN
\Cite{Matt 19:27-29}
\switchcolumn*

\LA
\noindent In illo témpore: Dixit Petrus ad Jesum: Ecce nos relíquimus ómnia, et secúti sumus te: quid ergo erit nobis? Et réliqua.\par
\switchcolumn
\EN
\noindent At that time, Peter said unto Jesus: Behold, we have forsaken all, and followed thee what shall we have therefore? And so on.\par
\switchcolumn*

\LA
\LessonTitle{Homilía sancti Hierónymi Presbýteri.}
\switchcolumn
\EN
\LessonTitle{Homily by St. Jerome, Priest (at Bethlehem.)}
\switchcolumn*

\LA
\Source{Lib. 3. in Matth. cap. 19.}
\switchcolumn
\EN
\Source{Bk. iii. on Matth xix.}
\switchcolumn*

\LA
\noindent Grandis fidúcia! Petrus piscátor erat, dives non fúerat, cibos manu et arte quærébat: et tamen lóquitur confidénter: Relíquimus ómnia. Et quia non súfficit tantum relínquere, jungit quod perféctum est: Et secúti sumus te. Fécimus quod jussísti: quid ígitur nobis dabis prǽmii? Jesus autem dixit illis: Amen dico vobis, quod vos qui secúti estis me, in regeneratióne, cum séderit Fílius hóminis in sede majestátis suæ, sedébitis et vos super sedes duódecim, judicántes duódecim tribus Ísraël. Non dixit: Qui reliquístis ómnia; hoc enim et Crates fecit philósophus, et multi álii divítias contempsérunt: sed, Qui secúti estis me: quod próprie Apostolórum est, atque credéntium.\par
\switchcolumn
\EN
\noindent Peter was a fisherman, he was not rich, he earned his bread by his hand and skill, and nevertheless he is thus bold, and saith confidently: We have forsaken all. And because it sufficeth not to forsake only, he addeth that which to do is to be perfect: and followed thee. We have done that which Thou hast commanded us, what reward therefore wilt Thou give us? And Jesus said unto them Amen I say unto you, that ye which have followed Me, in the regeneration, when the Son of Man shall sit in the throne of His glory, ye also shall sit upon twelve thrones, judging the twelve tribes of Israel. He said not, Ye which have forsaken all, for this did even Crates the philosopher, and they which have set nothing by riches are many, but, Ye which have followed Me. This did the Apostles, and this do believers do.\par
\switchcolumn*

\LA
\begin{Resp}\Rsign Iste est qui ante Deum magnas virtútes operátus est, et de omni corde suo laudávit Dóminum: \Star{} Ipse intercédat pro peccátis ómnium populórum.\end{Resp}
\begin{Resp}\Vsign Ecce homo sine queréla, verus Dei cultor, ábstinens se ab omni ópere malo, et pérmanens in innocéntia sua.\end{Resp}
\begin{Resp}\Rsign Ipse intercédat pro peccátis ómnium populórum.\end{Resp}
\switchcolumn
\EN
\begin{Resp}\Rsign This is he which wrought great wonders before God, and praised the Lord with all his heart. \Star{} May he pray for all people, that their sins may be forgiven unto them.\end{Resp}
\begin{Resp}\Vsign Behold a man without blame, a worshipper of God in truth, keeping himself clean from every evil work, and abiding still in his innocency.\end{Resp}
\begin{Resp}\Rsign May he pray for all people, that their sins may be forgiven unto them.\end{Resp}
\switchcolumn*

\LA
\LessonHead{Lectio VIII}
\switchcolumn
\EN
\LessonHead{Lesson VIII}
\switchcolumn*

\LA
\noindent In regeneratióne, cum séderit Fílius hóminis in sede majestátis suæ (quando et mórtui de corruptióne resúrgent incorrúpti) sedébitis et vos in sóliis judicántium, condemnántes duódecim tribus Ísraël: quia vobis credéntibus, illi crédere noluérunt. Et omnis qui relíquerit domum, vel fratres, aut soróres, aut patrem, aut matrem, aut uxórem, aut fílios, aut agros propter nomen meum, céntuplum accípiet, et vitam ætérnam possidébit. Locus iste cum illa senténtia cóngruit, in qua Salvátor lóquitur: Non veni pacem míttere, sed gládium. Veni enim separáre hóminem a patre suo, et matrem a fília, et nurum a socru: et inimíci hóminis doméstici ejus. Qui ergo propter fidem Christi, et prædicatiónem Evangélii, omnes afféctus contémpserint, atque divítias et sǽculi voluptátes: isti céntuplum recípient, et vitam ætérnam possidébunt.\par
\switchcolumn
\EN
\noindent In the regeneration, when the Son of Man shall sit in the throne of His glory, and when the dead shall rise again from corruption incorruptible, (i Cor. xv. 53,) ye also shall sit upon twelve thrones of judgment, condemning the twelve tribes of Israel, because, when ye believed in Me, they would not. (John iii. 18.) And every one that hath forsaken houses, or brethren, or sisters, or father, or mother, or wife, or children, or lands, for My Name's sake, shall receive an hundredfold, and shall inherit everlasting life. This place agreeth well with that other where the Saviour saith I came not to send peace, but a sword. For I am come to set a man at variance against his father, and the daughter against her mother, and the daughter-in-law against her mother-in-law; and a man's foes shall be they of his own household. (Matth. x. 34.) Every one, therefore, that hath set no store by affection, and riches, and the pleasures of the world, for Christ's faith's sake, and the preaching of the Gospel, shall receive an hundred-fold, and shall inherit everlasting life.\par
\switchcolumn*

\LA
\begin{Resp}\Rsign Sint lumbi vestri præcíncti, et lucérnæ ardéntes in mánibus vestris: \Star{} Et vos símiles homínibus exspectántibus dóminum suum, quando revertátur a núptiis.\end{Resp}
\begin{Resp}\Vsign Vigiláte ergo, quia nescítis qua hora Dóminus vester ventúrus sit.\end{Resp}
\begin{Resp}\Rsign Et vos símiles homínibus exspectántibus dóminum suum, quando revertátur a núptiis.\end{Resp}
\begin{Resp}\Vsign Glória Patri, et Fílio, \Star{} et Spirítui Sancto.\end{Resp}
\begin{Resp}\Rsign Et vos símiles homínibus exspectántibus dóminum suum, quando revertátur a núptiis.\end{Resp}
\switchcolumn
\EN
\begin{Resp}\Rsign Let your loins be girded about, and your lights burning; \Star{} And ye yourselves like unto men that wait for their lord, when he will return from the wedding.\end{Resp}
\begin{Resp}\Vsign Watch therefore, for ye know not what hour your Lord doth come.\end{Resp}
\begin{Resp}\Rsign And ye yourselves like unto men that wait for their lord, when he will return from the wedding.\end{Resp}
\begin{Resp}\Vsign Glory be to the Father, and to the Son, \Star{} and to the Holy Ghost.\end{Resp}
\begin{Resp}\Rsign And ye yourselves like unto men that wait for their lord, when he will return from the wedding.\end{Resp}
\switchcolumn*

\LA
\LessonHead{Lectio IX}
\switchcolumn
\EN
\LessonHead{Lesson IX}
\switchcolumn*

\LA
\noindent Ex occasióne hujus senténtiæ quidam introdúcunt mille annos post resurrectiónem, dicéntes, tunc nobis céntuplum ómnium rerum quas dimísimus, et vitam ætérnam esse reddéndam: non intelligéntes, quod, si in céteris digna sit repromíssio, in uxóribus appáreat turpitúdo: ut qui unam pro Dómino dimíserit centum recípiat in futúro. Sensus ergo iste est: Qui carnália pro Salvatóre dimíserit, spirituália recípiet: quæ comparatióne et mérito sui ita erunt, quasi si parvo número centenárius númerus comparétur.\par
\switchcolumn
\EN
\noindent By reason of these words, an hundredfold, some will have that there shall be a thousand years after the resurrection, wherein they that have forsaken all things shall receive an hundredfold of those things which they have forsaken, and shall inherit everlasting life. Such men consider not that though in other things this were worthy, as touching wives it is unseemly for it becometh us not to think that he that hath forsaken one wife in this world, shall receive an hundred wives in that which is to come. But the meaning is this, that every one that for the Saviour's sake hath forsaken earthly things, shall receive spiritual things which things, being rightly weighed against earthly things, are as though an hundredfold were weighed against one.\par
\switchcolumn*

\LA
\TeDeum{Te Deum}
\switchcolumn
\EN
\TeDeum{Te Deum}
\switchcolumn*

\end{paracol}

\par\needspace{10\baselineskip}\addvspace{1.2em}{\centering\small\textsc{Die 11 Octobris} · \textsc{11 October}\par}
\Office{Maternitatis Beatæ Mariæ Virginis}{Maternitatis Beatæ Mariæ Virginis}{Duplex II classis}
\addcontentsline{toc}{subsection}{Die 11 Octobris · Maternitatis Beatæ Mariæ Virginis \textit{— Maternitatis Beatæ Mariæ Virginis}}
\begin{paracol}{2}
\LA
\LessonHead{Lectio I}
\switchcolumn
\EN
\LessonHead{Lesson I}
\switchcolumn*

\LA
\LessonTitle{De libro Ecclesiástici}
\switchcolumn
\EN
\LessonTitle{Lesson from the book of Ecclesiasticus}
\switchcolumn*

\LA
\Cite{Sir 24:5-11}
\switchcolumn
\EN
\Cite{Sir 24:5-11}
\switchcolumn*

\LA
\noindent Ego ex ore Altíssimi prodívi, primogénita ante omnem creatúram: ego feci in cælis ut orirétur lumen indefíciens, et sicut nébula texi omnem terram: ego in altíssimis habitávi et thronus meus in colúmna nubis. Gyrum cæli circuívi sola, et profúndum abýssi penetrávi, in flúctibus maris ambulávi, et in omni terra steti: et in omni pópulo, et in omni gente primátum hábui: et ómnium excelléntium et humílium corda virtúte calcávi: et in his ómnibus réquiem quæsívi, et in hereditáte Dómini morábor.\par
\switchcolumn
\EN
\noindent I came out of the mouth of the most High, the firstborn before all creatures: I made that in the heavens there should rise light that never faileth, and as a cloud I covered all the earth: I dwelt in the highest places, and my throne is in a pillar of a cloud. I alone have compassed the circuit of heaven, and have penetrated into the bottom of the deep, and have walked in the waves of the sea, And have stood in all the earth: and in every people, And in every nation I have had the chief rule: And by my power I have trodden under my feet the hearts of all the high and low: and in all these I sought rest, and I shall abide in the inheritance of the Lord.\par
\switchcolumn*

\LA
\begin{Resp}\Rsign Felix es, sacra Virgo María et omni laude digníssima; \Star{} Ex qua ortus est sol justítiæ, Christus Deus noster, per quem salváti et redémpti sumus.\end{Resp}
\begin{Resp}\Vsign Maternitátem beátæ Maríæ Vírginis cum gáudio celebrémus.\end{Resp}
\begin{Resp}\Rsign Ex qua ortus est sol justítiæ, Christus Deus noster, per quem salváti et redémpti sumus.\end{Resp}
\switchcolumn
\EN
\begin{Resp}\Rsign O Holy Virgin Mary, happy art thou, and right worthy of all praise; \Star{} For out of thee rose the Sun of righteousness, even Christ our God, by whom we are saved and redeemed.\end{Resp}
\begin{Resp}\Vsign Let us celebrate with joy the Motherhood of the Blessed Virgin Mary.\end{Resp}
\begin{Resp}\Rsign For out of thee rose the Sun of righteousness, even Christ our God, by whom we are saved and redeemed.\end{Resp}
\switchcolumn*

\LA
\LessonHead{Lectio II}
\switchcolumn
\EN
\LessonHead{Lesson II}
\switchcolumn*

\LA
\Cite{Sir 24:12-16}
\switchcolumn
\EN
\Cite{Sir 24:12-16}
\switchcolumn*

\LA
\noindent Tunc præcépit, et dixit mihi Creátor ómnium, et qui creávit me, requiévit in tabernáculo meo, et dixit mihi: in Jacob inhábita, et in Israël hereditáre, et in eléctis meis mitte radíces. Ab inítio, et ante sǽcula creáta sum, et usque ad futúrum sǽculum non désinam, et in habitatióne sancta coram ipso ministrávi. Et sic in Sion firmáta sum, et in civitáte sanctificáta simíliter requiévi, et in Jerúsalem potéstas mea. Et radicávi in pópulo honorificáto, et in parte Dei mei heréditas illíus, et in plenitúdine sanctórum deténtio mea.\par
\switchcolumn
\EN
\noindent Then the creator of all things commanded, and said to me: and he that made me, rested in my tabernacle, And he said to me: Let thy dwelling be in Jacob, and thy inheritance in Israel, and take root in my elect. From the beginning, and before the world, was I created, and unto the world to come I shall not cease to be, and in the holy dwelling place I have ministered before him. And so was I established in Sion, and in the holy city likewise I rested, and my power was in Jerusalem. And I took root in an honourable people, and in the portion of my God his inheritance, and my abode is in the full assembly of saints.\par
\switchcolumn*

\LA
\begin{Resp}\Rsign Sine tactu pudóris invénta es Mater Salvatóris: \Star{} Qui cælum terrámque regit, in tua se clausit víscera factus homo.\end{Resp}
\begin{Resp}\Vsign Benedícta tu in muliéribus, et benedíctus fructus ventris tui.\end{Resp}
\begin{Resp}\Rsign Qui cælum terrámque regit, in tua se clausit víscera factus homo.\end{Resp}
\switchcolumn
\EN
\begin{Resp}\Rsign From thee, still maiden undefiled, the Saviour came a little Child: \Star{} He the Lord who ruleth o'er earth and o'er heaven for ever, being made man, was enclosed in the blest sides of thy womb.\end{Resp}
\begin{Resp}\Vsign Blessed art thou among women, and blessed is the Fruit of thy womb.\end{Resp}
\begin{Resp}\Rsign He the Lord who ruleth o'er earth and o'er heaven for ever, being made man, was enclosed in the blest sides of thy womb.\end{Resp}
\switchcolumn*

\LA
\LessonHead{Lectio III}
\switchcolumn
\EN
\LessonHead{Lesson III}
\switchcolumn*

\LA
\Cite{Sir 24:17-23}
\switchcolumn
\EN
\Cite{Sir 24:17-23}
\switchcolumn*

\LA
\noindent Quasi cedrus exaltáta sum in Líbano, et quasi cypréssus in monte Sion: quasi palma exaltáta sum in Cades, et quasi plantátio rosæ in Jéricho: quasi olíva speciósa in campis, et quasi plátanus exaltáta sum juxta aquam in platéis. Sicut cinnamómum et bálsamum aromatízans odórem dedi; quasi myrrha elécta dedi suavitátem odóris, et quasi storax, et gálbanus, et úngula, et gutta, et quasi Líbanus non incísus vaporávi habitatiónem meam, et quasi bálsamum non mixtum odor meus. Ego quasi terebínthus exténdi ramos meos, et rami mei honóris et grátiæ. Ego quasi vitis fructificávi suavitátem odóris.\par
\switchcolumn
\EN
\noindent I was exalted like a cedar in Libanus, and as a cypress tree on mount Sion. I was exalted like a palm tree in Cades, and as a rose plant in Jericho: As a fair olive tree in the plains, and as a plane tree by the water in the streets, was I exalted. I gave a sweet smell like cinnamon. and aromatical balm: I yielded a sweet odour like the best myrrh: And I perfumed my dwelling as storax, and galbanum, and onyx, and aloes, and as the frankincense not cut, and my odour is as the purest balm. I have stretched out my branches as the turpentine tree, and my branches are of honour and grace. As the vine I have brought forth a pleasant odour: and my flowers are the fruit of honour and riches.\par
\switchcolumn*

\LA
\begin{Resp}\Rsign Multæ fíliæ congregavérunt divítias, tu supergréssa es univérsas: \Star{} Speciósa facta es et suávis in delíciis tuis, sancta Dei Génetrix.\end{Resp}
\begin{Resp}\Vsign Séntiant omnes tuum juvámen, quicúmque célebrant tuam sanctam Maternitátem.\end{Resp}
\begin{Resp}\Rsign Speciósa facta es et suávis in delíciis tuis, sancta Dei Génetrix.\end{Resp}
\begin{Resp}\Vsign Glória Patri, et Fílio, \Star{} et Spirítui Sancto.\end{Resp}
\begin{Resp}\Rsign Speciósa facta es et suávis in delíciis tuis, sancta Dei Génetrix.\end{Resp}
\switchcolumn
\EN
\begin{Resp}\Rsign Many daughters have gotten riches, but thou excellest them all: \Star{} O holy Mother of God, thou art become beautiful and gentle in thy gladness.\end{Resp}
\begin{Resp}\Vsign May all that are keeping holy-day in honour of thy blessed Motherhood feel the might of thine assistance.\end{Resp}
\begin{Resp}\Rsign O holy Mother of God, thou art become beautiful and gentle in thy gladness.\end{Resp}
\begin{Resp}\Vsign Glory be to the Father, and to the Son, \Star{} and to the Holy Ghost.\end{Resp}
\begin{Resp}\Rsign O holy Mother of God, thou art become beautiful and gentle in thy gladness.\end{Resp}
\switchcolumn*

\LA
\LessonHead{Lectio IV}
\switchcolumn
\EN
\LessonHead{Lesson IV}
\switchcolumn*

\LA
\LessonTitle{Sermo sancti Leónis Papæ}
\switchcolumn
\EN
\LessonTitle{A Sermon of St. Leo the Pope}
\switchcolumn*

\LA
\Source{Sermo 1 de Nativitate Domini}
\switchcolumn
\EN
\Source{Sermon 1 on the Nativity of the Lord}
\switchcolumn*

\LA
\noindent Virgo régia Davídicæ stirpis elígitur, quæ sacro gravidánda fetu, divínam humanámque prolem prius concíperet mente quam córpore: et ne supérni ignára consílii ad inusitátos pavéret affátus, quod in ea operándum erat a Spíritu Sancto, collóquio dídicit angélico, nec damnum crédidit pudóris Dei Génetrix mox futúra. Cur enim de conceptiónis novitáte despéret, cui efficiéntia de Altíssimi virtúte promíttitur? Confirmátur credéntis fides étiam præeúntis attestatióne miráculi. Donátur Elísabeth inopináta fecúnditas, ut qui concéptum déderat stérili, datúrus non dubitarétur et Vírgini. Verbum ígitur Dei Fílius, qui in princípio erat apud Deum, per quem facta sunt ómnia, et sine quo factum est nihil, propter liberándum hóminem ab ætérna morte, factus est homo.\par
\switchcolumn
\EN
\noindent His Mother was chosen a Virgin of the kingly lineage of David, and when she was to grow heavy with the sacred Child, her soul had already conceived him before her body. She learned the counsel of God announced to her by the Angel, lest the unwonted events should alarm her. The future Mother of God knew what was to be wrought in her by the Holy Ghost, and that her modesty was absolutely safe. For why should she, unto whom was promised all sufficient strength through the power of the Highest, have felt hopeless merely because of the unexampled character of such a conception? She believeth, and her belief is confirmed by the attestation of a miracle which hath already been wrought. The fruitfulness of Elizabeth, before unhoped for, is brought forward that she might not doubt that he who had given conception unto her that was barren, would give the same unto her that was Virgin. And so the Word of God, the Son of God, who was in the beginning with God, by whom all things were made, and without whom was not anything made that was made, to deliver man from eternal death, was made man.\par
\switchcolumn*

\LA
\begin{Resp}\Rsign Gloriósæ Vírginis Maríæ Maternitátem digníssimam recolámus: \Star{} Cujus Dóminus humilitátem respéxit, quæ Angelo nuntiánte concépit Salvatórem mundi.\end{Resp}
\begin{Resp}\Vsign Christo canámus glóriam in hac sacra solemnitáte mirábilis Genetrícis Dei.\end{Resp}
\begin{Resp}\Rsign Cujus Dóminus humilitátem respéxit, quæ Angelo nuntiánte concépit Salvatórem mundi.\end{Resp}
\switchcolumn
\EN
\begin{Resp}\Rsign Let us tell again of the right worthy Motherhood of the glorious Virgin Mary: \Star{} The same is she whose lowliness the Lord regarded, she who by the message of an Angel conceived the Saviour of the world.\end{Resp}
\begin{Resp}\Vsign Let us sing praise to Christ on this the solemn Feastday of the wondrous Mother of God.\end{Resp}
\begin{Resp}\Rsign The same is she whose lowliness the Lord regarded, she who by the message of an Angel conceived the Saviour of the world.\end{Resp}
\switchcolumn*

\LA
\LessonHead{Lectio V}
\switchcolumn
\EN
\LessonHead{Lesson V}
\switchcolumn*

\LA
\Source{Sermo 2 de Nativitate Domini}
\switchcolumn
\EN
\Source{Sermon 2 on the Nativity of the Lord}
\switchcolumn*

\LA
\noindent Ingréditur hæc ínfima Jesus Christus Dóminus noster de cæli sede descéndens, et a patérna glória non recédens, novo órdine, nova nativitáte generátus. Novo órdine, quia invisíbilis in suis, visíbilis factus est in nostris: incomprehensíbilis vóluit comprehéndi: ante témpora manens, esse cœpit in témpore. Nova autem nativitáte génitus est: concéptus a Vírgine, natus ex Vírgine, sine patérnæ carnis concupiscéntia, sine matérnæ integritátis injúria: quia futúrum hóminum Salvatórem talis ortus decébat, qui et in se habéret humánæ substántiæ natúram, et humánæ carnis inquinaménta nescíret. Orígo dissímilis, sed natúra consímilis; humáno usu et consuetúdine, quod crédimus, caret: sed divína potestáte subníxum est, quod Virgo concéperit, Virgo pepérerit, Virgo permánserit.\par
\switchcolumn
\EN
\noindent Our Lord Jesus Christ, descending from his throne in heaven, but leaving not that glory which he hath with the Father, cometh into this lower world by being born after a new order and in a new birth. He cometh after a new order, in that he who is unseen among his own, was seen among us; the Incomprehensible was fain to be comprehended, and he that is from everlasting to everlasting began to be in time. He was the Offspring of a new birth; conceived of a maiden, without the passion of any fleshly father, without any breach of his Mother's virginity, since such a birth beseemed the coming Saviour of mankind, who was to have in him the nature of man's being, and to be free of any defilement of man's flesh. Though he sprung not as we spring, yet is his nature as our nature; we believe that he is free from the use and custom of men; but it was the power of God which wrought that a maiden should conceive, that a maiden should bring forth, and yet abide a maiden still.\par
\switchcolumn*

\LA
\begin{Resp}\Rsign Benedícta fília tu a Dómino, quia per te fructum vitæ communicávimus: \Star{} Sola sine exémplo placuísti Dómino nostro Jesu Christo.\end{Resp}
\begin{Resp}\Vsign Nostras deprecatiónes ne despícias in necessitátibus nostris, sed a perículis cunctis líbera nos, sancta Dei Génetrix.\end{Resp}
\begin{Resp}\Rsign Sola sine exémplo placuísti Dómino nostro Jesu Christo.\end{Resp}
\switchcolumn
\EN
\begin{Resp}\Rsign Blessed art thou of the Lord, O daughter, for through thee have we been given to eat of the tree of life: \Star{} Thou, without example before thee, didst make thyself well-pleasing in the sight of our Lord Jesus Christ.\end{Resp}
\begin{Resp}\Vsign Despise not our petitions in our necessities, but deliver us from all dangers, O holy Mother of God.\end{Resp}
\begin{Resp}\Rsign Thou, without example before thee, didst make thyself well-pleasing in the sight of our Lord Jesus Christ.\end{Resp}
\switchcolumn*

\LA
\LessonHead{Lectio VI}
\switchcolumn
\EN
\LessonHead{Lesson VI}
\switchcolumn*

\LA
\noindent Ex Actis Pii Papæ undécimi\par
\switchcolumn
\EN
\noindent From the Acts of Pope Pius XI\par
\switchcolumn*

\LA
\LessonTitle{Cum anno millésimo nongentésimo trigésimo primo, univérso orbe cathólico plaudénte, solémnia celebraréntur expléti sǽculi décimi quinti, postquam in Ephesína sýnodo beáta María Virgo, de qua natus est Jesus, contra Nestórii hǽresim Mater Dei a Pátribus, Cælestíno Papa præeúnte, conclamáta est, Summus Póntifex Pius undécimus faustíssimi evéntus memóriam, perénni suæ pietátis testimónio perpetuándam vóluit. Itaque quod jam in Urbe exstábat nóbile Ephesínæ proclamatiónis monuméntum, triumphálem arcum in Basílica sanctæ Maríæ Majóris in Exquíliis, a decessóre suo Xysto tértio mirábili ópere musívo ornátum, témporis injúria fatiscéntem felíciter restituéndum una cum ala transvérsa Basílicæ munificéntia sua curávit. Lítteris vero encýclicis, Œcuménici Concílii Ephesíni genuínis lineaméntis descríptis, ineffábile divínæ Maternitátis beátæ Maríæ Vírginis privilégium, pie copioséque illustrávit, ut tam excélsi mystérii doctrína áltius fidélium ánimis insidéret. Insímul autem benedíctam inter omnes mulíeres, Maríam Matrem Dei Nazarenámque Famíliam nobilíssimum præ ómnibus exémplum propósuit imitándum tum dignitátis ac sanctitúdinis casti connúbii tum educatiónis juventúti sancte tradéndæ. Demum, ut neque litúrgicum deésset monuméntum jussit ut festum divínæ Maternitátis beátæ Maríæ Vírginis cum Missa et Offício própriis die undécima octóbris sub ritu dúplici secúndæ classis quotánnis ab univérsa Ecclésia celebrarétur.}
\switchcolumn
\EN
\LessonTitle{In the year 1931, amid the applause of the whole Catholic world, solemn rites were celebrated to mark the completion of the fifteen centuries which had elapsed since the Council of Ephesus, moving against the Nestorian heresy, had acclaimed the blessed Virgin Mary, of whom Jesus was born, as Mother of God. This acclamation had been made by the Fathers of the Church under the leadership of Pope Celestine. Pius XI, as Supreme Pontiff, wished to commemorate the notable event and to give lasting proof of his devotion to Mary. Now there had existed for many years in Rome a grand memorial to the proclamation of Ephesus, the triumphal arch in the basilica of Saint Mary Major on the Esquiline Hill. This monument had already been adorned by a previous pontiff, Sixtus III, with mosaics of marvelous workmanship, now falling to pieces from the decay of the passing ages. Pius XI, therefore, out of his own munificence, caused these to be restored most exquisitely and with them the transept of the basilica. In an Encyclical Letter Pius set forth also the true history of the Council of Ephesus, and expounded fervently and at great length the doctrine of the prerogatives of the Blessed Virgin Mary as Mother of God. He did this that the doctrine of this lofty mystery might sink more deeply into the hearts of the faithful. In it he set forth Mary, the Mother of God, blessed among women, and the most holy Family of Nazareth as the exemplars to be followed above all others, as models of the dignity and holiness of chaste wedlock, as patterns of the holy education to be given youth. Finally that no liturgical detail be lacking, he decreed that the feast of the Divine Motherhood of the Blessed Virgin Mary be celebrated annually on the 11th day of October by the universal Church with a proper Mass and Office under the rite of a double of the second class.}
\switchcolumn*

\LA
\begin{Resp}\Rsign Benedícta tu inter mulíeres, et benedíctus fructus ventris tui: \Star{} Unde hoc mihi, ut véniat Mater Dómini mei ad me?\end{Resp}
\begin{Resp}\Vsign Respéxit humilitátem ancíllæ suæ, et fecit mihi magna, qui potens est.\end{Resp}
\begin{Resp}\Rsign Unde hoc mihi, ut véniat Mater Dómini mei ad me?\end{Resp}
\begin{Resp}\Vsign Glória Patri, et Fílio, \Star{} et Spirítui Sancto.\end{Resp}
\begin{Resp}\Rsign Unde hoc mihi, ut véniat Mater Dómini mei ad me?\end{Resp}
\switchcolumn
\EN
\begin{Resp}\Rsign Blessed art thou among women, and blessed is the Fruit of thy womb: \Star{} Whence is this to me, that the Mother of my Lord should come to me?\end{Resp}
\begin{Resp}\Vsign He hath regarded the lowliness of his handmaiden, and he that is mighty hath magnified me.\end{Resp}
\begin{Resp}\Rsign Whence is this to me, that the Mother of my Lord should come to me?\end{Resp}
\begin{Resp}\Vsign Glory be to the Father, and to the Son, \Star{} and to the Holy Ghost.\end{Resp}
\begin{Resp}\Rsign Whence is this to me, that the Mother of my Lord should come to me?\end{Resp}
\switchcolumn*

\LA
\LessonHead{Lectio VII}
\switchcolumn
\EN
\LessonHead{Lesson VII}
\switchcolumn*

\LA
\LessonTitle{Léctio sancti Evangélii secúndum Lucam}
\switchcolumn
\EN
\LessonTitle{From the Holy Gospel according to Luke}
\switchcolumn*

\LA
\Cite{Luc 2:43-51}
\switchcolumn
\EN
\Cite{Luke 2:43-51}
\switchcolumn*

\LA
\noindent In illo témpore: Cum redírent, remánsit puer Jesus in Jerúsalem, et non cognovérunt paréntes ejus. Et réliqua.\par
\switchcolumn
\EN
\noindent In that time, having fulfilled the days, when they returned, the child Jesus remained in Jerusalem; and his parents knew it not. And so on.\par
\switchcolumn*

\LA
\LessonTitle{Homilía sancti Bernárdi Abbátis}
\switchcolumn
\EN
\LessonTitle{A Homily of St. Bernard the Abbot}
\switchcolumn*

\LA
\Source{Hom. 1 de Laud. Virg. Matris}
\switchcolumn
\EN
\Source{Homilia 1. de Laud. Virg. Matris}
\switchcolumn*

\LA
\noindent Deum et Dóminum Angelórum María suum fílium appéllat, dicens: Fili, quid fecísti nobis sic? Quis hoc áudeat Angelórum? Súfficit eis, et pro magno habent, quod cum sint spíritus ex conditióne, ex grátia facti sint et vocáti Angeli, testánte David: Qui facit Angelos suos spíritus. María vero matrem se agnóscens, majestátem illam, cui illi cum reveréntia sérviunt, cum fidúcia suum núncupat Fílium: nec dedignátur nuncupári Deus, quod esse dignátus est. Nam paulo post subdit Evangelísta: Et erat súbditus illis. Quis? Quibus? Deus homínibus? Deus, inquam, cui Angeli súbditi sunt, cui Principátus et Potestátes obédiunt, súbditus erat Maríæ.\par
\switchcolumn
\EN
\noindent Son, why hast thou thus dealt with us? Mary called God, the Lord of Angels, her son. Which of the angels would have dared to do so? It is enough for them, and they reckon it is a great thing, that, being naturally spirits, they should receive the grace of being made and called angels, as witness David: Who maketh spirits his angels. But Mary, knowing herself to be his Mother, doth boldly apply the word Son to that Majesty whom the angels do serve with awe; neither doth God despise to be called what he hath made himself. For a little after, the Evangelist saith: And he was subject unto them. Who to whom? God to men. I say that God, unto whom the angels are subject, and who is obeyed by the Principalities and Powers, was subject to Mary.\par
\switchcolumn*

\LA
\begin{Resp}\Rsign Beáta es Virgo María, Dei Génetrix, quæ credidísti Dómino: perfécta sunt in te, quæ dicta sunt tibi: \Star{} Proptérea benedíxit te Deus in ætérnum.\end{Resp}
\begin{Resp}\Vsign Diffúsa est grátia in lábiis tuis: intercéde pro nobis ad Dóminum Deum nostrum.\end{Resp}
\begin{Resp}\Rsign Proptérea benedíxit te Deus in ætérnum.\end{Resp}
\switchcolumn
\EN
\begin{Resp}\Rsign Blessed art thou, O Virgin Mary, Mother of God, thou that hast believed the Lord, for there hath been a performance in thee of those things which were told thee: \Star{} Therefore God hath blessed thee for ever.\end{Resp}
\begin{Resp}\Vsign Grace is poured into thy lips; plead for us with the Lord our God.\end{Resp}
\begin{Resp}\Rsign Therefore God hath blessed thee for ever.\end{Resp}
\switchcolumn*

\LA
\LessonHead{Lectio VIII}
\switchcolumn
\EN
\LessonHead{Lesson VIII}
\switchcolumn*

\LA
\noindent Miráre utrúmlibet, et élige quod ámplius miréris, sive Fílii benigníssimam dignatiónem, sive Matris excellentíssimam dignitátem. Utrímque stupor, utrímque miráculum. Et quod Deus féminæ obtémperet, humílitas absque exémplo: et quod Deo fémina principétur, sublímitas sine sócio. In láudibus Vírginum singuláriter cánitur quod sequúntur Agnum quocúmque íerit. Quibus ergo láudibus júdicas dignam, quæ étiam præit? Disce, homo, obedíre: disce, terra, subdi; disce, pulvis, obtemperáre. De Auctóre tuo loquens Evangelísta: Et erat, inquit, súbditus illis. Erubésce, supérbe cinis! Deus se humíliat, et tu te exáltas? Deus se homínibus subdit, et tu, dominári géstiens homínibus, tuo te præpónis Auctóri?\par
\switchcolumn
\EN
\noindent Marvel thou at both these things, and choose whether to marvel most at the sublime condescension of the Son, or at the sublime dignity of Mary. Either is amazing, either marvelous. That God should obey this woman, is a lowliness without parallel; that this woman should rule over God, an exaltation without match. In praise of virgins, and of virgins only, is it sung that These are they which follow the Lamb whithersoever he goeth. Of what praise then thinkest thou that she must be worthy who even leadeth the Lamb? O man, learn to obey. O earth, learn to submit. O dust, learn to keep down. It is of thy Maker that the Evangelist saith: And he was subject unto them. Blush, O proud ashes. God humbleth himself; and dost thou exalt thyself? God is subject unto men; and wilt thou, by striving to rule over men, set thyself before thy Maker?\par
\switchcolumn*

\LA
\begin{Resp}\Rsign Congratulámini mihi, omnes qui dilígitis Dóminum: quia cum essem párvula, plácui Altíssimo: \Star{} Et de meis viscéribus génui Deum et hóminem.\end{Resp}
\begin{Resp}\Vsign Beátam me dicent omnes generatiónes, quia ancíllam húmilem respéxit Deus.\end{Resp}
\begin{Resp}\Rsign Et de meis viscéribus génui Deum et hóminem.\end{Resp}
\begin{Resp}\Vsign Glória Patri, et Fílio, \Star{} et Spirítui Sancto.\end{Resp}
\begin{Resp}\Rsign Et de meis viscéribus génui Deum et hóminem.\end{Resp}
\switchcolumn
\EN
\begin{Resp}\Rsign Rejoice with me, all you that love the Lord, for in my lowliness I pleased the Most High, \Star{} And from my womb I brought forth God-made-Man.\end{Resp}
\begin{Resp}\Vsign All generations shall call me blessed, for God hath regarded the lowliness of his handmaiden.\end{Resp}
\begin{Resp}\Rsign And from my womb I brought forth God-made-Man.\end{Resp}
\begin{Resp}\Vsign Glory be to the Father, and to the Son, \Star{} and to the Holy Ghost.\end{Resp}
\begin{Resp}\Rsign And from my womb I brought forth God-made-Man.\end{Resp}
\switchcolumn*

\LA
\LessonHead{Lectio IX}
\switchcolumn
\EN
\LessonHead{Lesson IX}
\switchcolumn*

\LA
\noindent Felix María, cui nec humílitas défuit, nec virgínitas! Et quidem singuláris virgínitas, quam non temerávit, sed honorávit fecúnditas. Et nihilóminus speciális humílitas, quam non ábstulit, sed éxtulit fecúnda virgínitas: et incomparábilis prorsus fecúnditas, quam virgínitas simul comitátur et humílitas. Quid horum non mirábile? quid non incomparábile? quid non singuláre? Mirum vero, si non hǽsitas, in horum ponderatióne, quid tua júdices dígnius admiratióne, utrum vidélicet pótius stupénda sit fecúnditas in Vírgine, an in Matre intégritas: sublímitas in prole, an cum tanta sublimitáte humílitas: nisi quod indubitánter horum síngulis præferénda sunt simul cuncta, et incomparabíliter excelléntius est atque felícius ómnia percepísse, quam áliqua. Et quid mirum, si Deus, qui mirábilis cérnitur, et légitur in Sanctis suis, mirabiliórem se exhíbuit in Matre sua? Venerámini ergo, cónjuges, in carne corruptíbili carnis integritátem: vos sacræ vírgines in Vírgine fecunditátem. Imitámini, omnes hómines, Dei Matris humilitátem.\par
\switchcolumn
\EN
\noindent O happy Mary, lowly and virgin; and wondrous virginity, which motherhood destroyed not, but exalted; and wondrous lowliness, which the fruitful virginity took not away, but ennobled; and wondrous motherhood, which was both virgin and lowly. Which of them is not wondrous? which of them is not unexampled? and which of them doth not stand alone? The wonder would be if thou wert not puzzled at which to wonder most? motherhood in a virgin, or virginity in a mother; a motherhood so exalted, or lowliness in such exaltation. But indeed more marvelous than any one of these things is the combination of them all, and without all comparison, it is more excellent and more blessed to have received them all, than to have received any one of them alone. What wonder is it that God, of whom we see and read, that He is wonderful in his holy places, should have shown himself wonderful in his Mother? O ye that be married, honor this incorruption in corruptible flesh; O holy maidens, gaze in wonder at motherhood in a maid; O, all mankind, take pattern by the lowliness of the Mother of God.\par
\switchcolumn*

\LA
\TeDeum{Te Deum}
\switchcolumn
\EN
\TeDeum{Te Deum}
\switchcolumn*

\end{paracol}

\par\needspace{10\baselineskip}\addvspace{1.2em}{\centering\small\textsc{Die 13 Octobris} · \textsc{13 October}\par}
\Office{S. Eduardi Regis Confessoris}{St. Edward the Confessor, King}{Semiduplex}
\addcontentsline{toc}{subsection}{Die 13 Octobris · S. Eduardi Regis Confessoris \textit{— St. Edward the Confessor, King}}
\begin{paracol}{2}
\LA
\RefLine{Lectiones I–III: de Scriptura occurrente.}
\switchcolumn
\EN
\RefLine{Lessons I–III: from the Scripture of the day.}
\switchcolumn*

\LA
\RefLine{Lectiones VII–IX: ex Commúni Confessoris non pontificis (C5).}
\switchcolumn
\EN
\RefLine{Lessons VII–IX: from the Common of a Confessor not a Bishop (C5).}
\switchcolumn*

\LA
\LessonHead{Lectio IV}
\switchcolumn
\EN
\LessonHead{Lesson IV}
\switchcolumn*

\LA
\noindent Eduardus, cognomento Confessor, nepos sancti Eduardi Regis et Mártyris, Anglo-Sáxonum regum ultimus, quem futurum regem Brithualdo viro sanctíssimo in mentis excéssu Dóminus demonstrávit, decennis a Danis Angliam vastántibus quæsítus ad necem, exsulare cogitur apud avúnculum, Normanniæ ducem. Ubi in mediis vitiórum illécebris talem se exhibuit integritate vitæ morúmque innocéntia, ut ómnibus admiratióni esset. Eluxit in eo vel tum mira pietas in Deum ac res divinas, fuitque ingenio mitíssimo atque ab omni dominándi cupiditate alieno. Cujus ea vox fertur, Malle se regno carére, quod sine cæde et sánguine obtineri non possit.\par
\switchcolumn
\EN
\noindent Edward, surnamed the Confessor, was the nephew of the holy King Edward the Martyr, and himself the last Anglo-Saxon King. That he should succeed to the Kingdom was shown by the Lord in a trance to a most holy man named Brithwald. When he was ten years old the Danes, who were ravaging England, sought him, to put him to death, and he was driven into exile to dwell with his mother's brother, (Richard II.) Duke of Normandy, at whose Court (and that of his successors, Richard III., Robert surnamed the Devil, and William the Bastard) he lived among all the allurements of vice a life of such uprightness and innocency as made all men to marvel. He was a burning and shining light for love of God and the things of God, very gentle-hearted, and quite free from any lust for power. Of him the saying is preserved, That he would liefer not be a King than win a kingdom through slaughter and blood.\par
\switchcolumn*

\LA
\begin{Resp}\Rsign Honéstum fecit illum Dóminus, et custodívit eum ab inimícis, et a seductóribus tutávit illum: \Star{} Et dedit illi claritátem ætérnam.\end{Resp}
\begin{Resp}\Vsign Justum dedúxit Dóminus per vias rectas, et osténdit illi regnum Dei.\end{Resp}
\begin{Resp}\Rsign Et dedit illi claritátem ætérnam.\end{Resp}
\switchcolumn
\EN
\begin{Resp}\Rsign The Lord made him honourable, and defended him from his enemies, and kept him safe from those that lay in wait for him; \Star{} And gave him perpetual glory.\end{Resp}
\begin{Resp}\Vsign He went down with him into the pit, and left him not in bonds.\end{Resp}
\begin{Resp}\Rsign And gave him perpetual glory.\end{Resp}
\switchcolumn*

\LA
\LessonHead{Lectio V}
\switchcolumn
\EN
\LessonHead{Lesson V}
\switchcolumn*

\LA
\noindent Exstinctis mox tyrannis, qui frátribus suis vitam et regnum eripúerant, revocátur in pátriam. Ubi, summis ómnium votis et gratulatióne, regno potítus, ad hostílium irárum delenda vestígia totum se convértit, a sacris exorsus ac Divórum templis, quorum alia a fundamentis eréxit, alia refécit auxitque redítibus ac privilegiis; in eam curam potíssimum intentus, ut refloresceret collapsa religio. Ab aulæ procéribus compulsum ad nuptias, constans est assertio scriptórum, cum vírgine sponsa virginitátem in matrimonio servasse. Tantus in eo fuit in Christum amor et fides, ut illum aliquándo inter Missárum solemnia vidére merúerit, blando vultu et divina luce fulgentem. Ob profusam caritátem, orphanórum et egenórum pater passim dicebátur, numquam lætior quam cum regios thesáuros exhausísset in páuperes.\par
\switchcolumn
\EN
\noindent When the (Danish) tyrants, who had robbed his brothers Edmund and Alfred of life and kingdom, were passed away, Edward was called back into his own country and with the hearty good-will and rejoicing of all, took the kingdom (in the year 1042, being then about forty years old.) He set himself to repair the breaches which wars had made, and began with the things of God. Of the Churches of the Saints, he built some altogether, and renewed others and gifted them with incomes and privileges, being chiefly fain that religion should rise from the low estate whereinto it had fallen. He was brought by the nobles of his Court to marry, but it is constantly said by all writers that in matrimony he remained a virgin with a virgin bride. So great was his love toward Christ, and so strong his faith, that somewhiles when the Mass was in saying, he won to see Him, with countenance full of grace, and glory of God's light. By reason of the abundance of his charity he was styled everywhere the father of orphans and of the poor, and he was never happier than when he had spent upon the needy the whole of his kingly treasures.\par
\switchcolumn*

\LA
\begin{Resp}\Rsign Amávit eum Dóminus, et ornávit eum: stolam glóriæ índuit eum, \Star{} Et ad portas paradísi coronávit eum.\end{Resp}
\begin{Resp}\Vsign Índuit eum Dóminus lorícam fídei, et ornávit eum.\end{Resp}
\begin{Resp}\Rsign Et ad portas paradísi coronávit eum.\end{Resp}
\switchcolumn
\EN
\begin{Resp}\Rsign The Lord loved him and beautified him; He clothed him with a robe of glory, \Star{} And crowned him at the gates of Paradise.\end{Resp}
\begin{Resp}\Vsign The Lord hath put on him the breast-plate of faith, and hath adorned him.\end{Resp}
\begin{Resp}\Rsign And crowned him at the gates of Paradise.\end{Resp}
\switchcolumn*

\LA
\LessonHead{Lectio VI}
\switchcolumn
\EN
\LessonHead{Lesson VI}
\switchcolumn*

\LA
\noindent Prophetíæ dono illustris, de Angliæ futuro statu multa cælitus prævidit; et illud in primis memorábile, quod Sweyni Danórum regis, in mare demersi, mortem, dum Angliam invadéndi animo classem conscenderet, eodem quo accidit momento, divinitus intelléxit. Joánnem Evangelistam mirifice coluit, nihil cuiquam, quod ejus nómine peterétur, negare sólitus. Cui olim sub lácera veste suo nómine stipem roganti, cum nummi deéssent, detractum ex dígito ánulum porréxit; quem Divus non ita multo post Eduardo remisit, una cum núntio secuturæ mortis. Quare rex, indictis pro se precibus, ipso ab Evangelista prædicto die piíssime obiit, Nonis videlicet Januarii, anno salútis millesimo sexagesimo sexto. quem sequénti sæculo Alexander Papa tertius, miraculis clarum, Sanctórum fastis adscripsit. At ejus memóriam Innocentius undecimus Officio publico per universam Ecclésiam eo die celebrari præcepit, quo annis ab óbitu sex et trigínta translátum ejus corpus, incorruptum et suavem spirans odórem repertum est.\par
\switchcolumn
\EN
\noindent He was famous for the gift of prophecy, and foretold by inspiration from heaven many things that were to befall England. Of this gift the following is a remarkable instance. Sweyn, King of the Danes, was embarking on ship-board with the mind to invade England, when he fell into the sea and was drowned, and God made known his death to Edward at the very same moment that it happened. He had a wonderful love toward John the Evangelist, so that he was used never to refuse anything for the which he was asked in his name. The Evangelist appeared to him once while in tattered raiment, and, in his own name, asked him for an alms. It befell that the King had no money, wherefore he took a ring from off his finger and gifted him therewith. Not long afterward, the Evangelist sent the same ring back to him by a pilgrim, with a message concerning his death, which was then at hand. The King therefore commanded that prayers should be made for him, and then fell blessedly asleep in the Lord, upon the very day which had been foretold to him by the Evangelist, that is to say, upon the 5th day of January, in the year of salvation 1066. He was famous for miracles, and in 1161 Pope Alexander III. numbered him among the Saints. But Innocent XI. commanded that his memory should be celebrated with a public Office throughout the whole Church, upon the 13th day of October, being that day whereon in the year 1102 his body had been lifted, and found uncorrupt and sweetsavoured.\par
\switchcolumn*

\LA
\begin{Resp}\Rsign Iste homo perfécit ómnia quæ locútus est ei Deus, et dixit ad eum: Ingrédere in réquiem meam: \Star{} Quia te vidi justum coram me ex ómnibus géntibus.\end{Resp}
\begin{Resp}\Vsign Iste est, qui contémpsit vitam mundi, et pervénit ad cæléstia regna.\end{Resp}
\begin{Resp}\Rsign Quia te vidi justum coram me ex ómnibus géntibus.\end{Resp}
\begin{Resp}\Vsign Glória Patri, et Fílio, \Star{} et Spirítui Sancto.\end{Resp}
\begin{Resp}\Rsign Quia te vidi justum coram me ex ómnibus géntibus.\end{Resp}
\switchcolumn
\EN
\begin{Resp}\Rsign This is he which did according unto all that God commanded him; and God said unto him: Enter thou into My rest, \Star{} For thee have I seen righteous before Me among all people.\end{Resp}
\begin{Resp}\Vsign This is he which loved not his life in this world, and is come unto an everlasting kingdom.\end{Resp}
\begin{Resp}\Rsign For thee have I seen righteous before Me among all people.\end{Resp}
\begin{Resp}\Vsign Glory be to the Father, and to the Son, \Star{} and to the Holy Ghost.\end{Resp}
\begin{Resp}\Rsign For thee have I seen righteous before Me among all people.\end{Resp}
\switchcolumn*

\end{paracol}

\par\needspace{10\baselineskip}\addvspace{1.2em}{\centering\small\textsc{Die 14 Octobris} · \textsc{14 October}\par}
\Office{S. Callisti Papæ et Martyris}{St. Callistus, Pope and Martyr}{Duplex}
\addcontentsline{toc}{subsection}{Die 14 Octobris · S. Callisti Papæ et Martyris \textit{— St. Callistus, Pope and Martyr}}
\begin{paracol}{2}
\LA
\RefLine{Lectiones I–III: de Scriptura occurrente.}
\switchcolumn
\EN
\RefLine{Lessons I–III: from the Scripture of the day.}
\switchcolumn*

\LA
\RefLine{Lectiones VII–IX: ex Commúni unius Summi Pontificis et Martyris (C2b).}
\switchcolumn
\EN
\RefLine{Lessons VII–IX: from the Common of a Single Martyr Pope (C2b).}
\switchcolumn*

\LA
\LessonHead{Lectio IV}
\switchcolumn
\EN
\LessonHead{Lesson IV}
\switchcolumn*

\LA
\noindent Callistus, Romanus, præfuit Ecclésiæ, Antonino Heliogábalo imperatóre. Constítuit quatuor anni Tempora, quibus jejunium, ex apostolica traditióne acceptum, ab ómnibus servarétur. Ædificávit basilicam sanctæ Maríæ trans Tíberim, et in via Appia vetus cœmetérium ampliávit, in quo multi sancti Sacerdótes et Mártyres sepúlti sunt; unde ab eo Callísti cœmetérium appellátur.\par
\switchcolumn
\EN
\noindent Callistus was a Roman, and ruled the Church in the time of the Emperor Antoninus Heliogabalus. He confirmed the institution of the Ember Fasts, the observance of which hath been received by tradition from the Apostles. He built the Church of St. Mary-beyond-the-Tiber, and enlarged the old burying-place on the Appian Way, wherein are buried so many holy Priests and Martyrs, and which hath since been called, on account of this enlargement, the Cemetery of Kallistus.\par
\switchcolumn*

\LA
\begin{Resp}\Rsign Honéstum fecit illum Dóminus, et custodívit eum ab inimícis, et a seductóribus tutávit illum: \Star{} Et dedit illi claritátem ætérnam.\end{Resp}
\begin{Resp}\Vsign Descendítque cum illo in fóveam, et in vínculis non derelíquit eum.\end{Resp}
\begin{Resp}\Rsign Et dedit illi claritátem ætérnam.\end{Resp}
\switchcolumn
\EN
\begin{Resp}\Rsign The Lord made him honourable, and defended him from his enemies, and kept him safe from those that lay in wait for him, \Star{} And gave him perpetual glory.\end{Resp}
\begin{Resp}\Vsign He went down with him into the pit, and left him not in bonds.\end{Resp}
\begin{Resp}\Rsign And gave him perpetual glory.\end{Resp}
\switchcolumn*

\LA
\LessonHead{Lectio V}
\switchcolumn
\EN
\LessonHead{Lesson V}
\switchcolumn*

\LA
\noindent Ejusdem pietátis fuit, quod beáti Calepodii Presbýteri et Mártyris corpus, jactátum in Tíberim, conquiri diligenter curávit, et, invéntum, honorifice sepelívit. Palmatium consulari, Simplícium senatoria dignitate illustres, Felicem et Blandam, qui deínde omnes martyrium subíere, cum baptismo lustrasset, missus est in carcerem, ubi Privátum militem, ulcéribus plenum, admirabíliter sanitáti restitutum, Christo adjunxit; pro quo idem, recens adhuc a fide suscepta, plumbátis usque ad mortem cæsus occubuit.\par
\switchcolumn
\EN
\noindent It was by his reverence that the body of the blessed Priest and Martyr Calepodius, which had been cast into the Tiber, was carefully looked for, and, when it had been found, honourably buried. He baptized Palmatius, of Consular, and Simplicius, of Senatorial rank, and likewise Felix and Blanda, all of whom in the end underwent martyrdom. On this account he was thrown into prison, where he wonderfully healed a soldier named Privatus, who was full of sores, and so gained him to Christ; and this Privatus had hardly received the faith, before he was lashed to death with scourges loaded with lead.\par
\switchcolumn*

\LA
\begin{Resp}\Rsign Desidérium ánimæ ejus tribuísti ei Dómine, \Star{} Et voluntáte labiórum ejus non fraudásti eum.\end{Resp}
\begin{Resp}\Vsign Quóniam prævenísti eum in benedictiónibus dulcédinis: posuísti in cápite ejus corónam de lápide pretióso.\end{Resp}
\begin{Resp}\Rsign Et voluntáte labiórum ejus non fraudásti eum.\end{Resp}
\switchcolumn
\EN
\begin{Resp}\Rsign Lord, Thou hast given him his heart's desire. \Star{} And hast not withholden the request of his lips.\end{Resp}
\begin{Resp}\Vsign For Thou hast prevented him with the blessings of sweetness: Thou hast set a crown of precious stones upon his head.\end{Resp}
\begin{Resp}\Rsign And hast not withholden the request of his lips.\end{Resp}
\switchcolumn*

\LA
\LessonHead{Lectio VI}
\switchcolumn
\EN
\LessonHead{Lesson VI}
\switchcolumn*

\LA
\noindent Sedit Callistus annos quinque, mensem unum, dies duodecim. Ordinatiónibus quinque, mense Decembri, creávit presbyteros sexdecim, diaconos quatuor, episcopos octo. Post longam famem crebrasque verberatiónes præceps jactus in púteum, atque ita martyrio coronátus sub Alexandro imperatóre, illátus est in cœmetérium Calepodii, via Aurelia, tertio ab Urbe lápide, pridie Idus Octobris. Ejus postmodum corpus in basilicam sanctæ Maríæ trans Tiberim, ab ipso ædificatam, delátum, sub ara majori, maxima veneratióne cólitur.\par
\switchcolumn
\EN
\noindent Callistus sat as Pope five years, one month, and twelve days. He held five Ordinations in the month of December, wherein he ordained sixteen Priests, four Deacons, and eight Bishops. After being long starved, and repeatedly flogged, he was pitched head-foremost down a well, and so crowned with martyrdom, under the Emperor Alexander. His body was carried to the Cemetery of Calepodius on the Aurelian Way, at the third mile-stone from the city, upon the 14th day of October, but was afterwards taken to the Church of St. Mary-beyond-the-Tiber, which had been built by himself. There it lieth beneath the High Altar, and is held in great reverence of all men.\par
\switchcolumn*

\LA
\begin{Resp}\Rsign Stola jucunditátis índuit eum Dóminus: \Star{} Et corónam pulchritúdinis pósuit super caput ejus.\end{Resp}
\begin{Resp}\Vsign Cibávit illum Dóminus pane vitæ et intelléctus: et aqua sapiéntiæ salutáris potávit illum.\end{Resp}
\begin{Resp}\Rsign Et corónam pulchritúdinis pósuit super caput ejus.\end{Resp}
\begin{Resp}\Vsign Glória Patri, et Fílio, \Star{} et Spirítui Sancto.\end{Resp}
\begin{Resp}\Rsign Et corónam pulchritúdinis pósuit super caput ejus.\end{Resp}
\switchcolumn
\EN
\begin{Resp}\Rsign The Lord hath put on him a robe of honour, \Star{} And put about his head a crown of joy.\end{Resp}
\begin{Resp}\Vsign With the bread of life and understanding hath the Lord fed him, and given him the water of health and wisdom to drink.\end{Resp}
\begin{Resp}\Rsign And put about his head a crown of joy.\end{Resp}
\begin{Resp}\Vsign Glory be to the Father, and to the Son, \Star{} and to the Holy Ghost.\end{Resp}
\begin{Resp}\Rsign And put about his head a crown of joy.\end{Resp}
\switchcolumn*

\end{paracol}

\par\needspace{10\baselineskip}\addvspace{1.2em}{\centering\small\textsc{Die 15 Octobris} · \textsc{15 October}\par}
\Office{S. Teresiæ Virginis}{St. Teresa of Avila, Virgin}{Duplex}
\addcontentsline{toc}{subsection}{Die 15 Octobris · S. Teresiæ Virginis \textit{— St. Teresa of Avila, Virgin}}
\begin{paracol}{2}
\LA
\RefLine{Lectiones I–III: de Scriptura occurrente.}
\switchcolumn
\EN
\RefLine{Lessons I–III: from the Scripture of the day.}
\switchcolumn*

\LA
\RefLine{Lectiones VII–IX: ex Commúni Virginum tantum (C6a).}
\switchcolumn
\EN
\RefLine{Lessons VII–IX: from the Common of a Virgin not a Martyr (C6a).}
\switchcolumn*

\LA
\LessonHead{Lectio IV}
\switchcolumn
\EN
\LessonHead{Lesson IV}
\switchcolumn*

\LA
\noindent Teresia Virgo, nata est Abulæ in Hispánia paréntibus tum genere tum pietáte præcláris. Ab iis divini timoris lacte educata, admirándum futuræ sanctitátis in tenerrima adhuc ætate specimen dedit. Nam, cum sanctórum Mártyrum acta perlegeret, adeo in ejus meditatióne Sancti Spíritus ignis exársit, ut, domo aufúgiens, in Africam trajiceret, ubi vitam pro glória Jesu Christi et animárum salúte profúnderet. A patruo revocata, ardens martyrii desidérium eleemosynis aliisque piis opéribus compensávit, jugibus lácrimis deplorans optimam sibi sortem fuísse præreptam. Mortua matre, cum a beatíssima Vírgine peteret ut se matrem esse monstraret, pii voti compos effecta est; semper perinde ac fília patrocinio Deiparæ pérfruens. Vigesimum ætátis annum agens, ad moniáles sanctæ Maríæ de Monte Carmelo se cóntulit. Ibi, per duodeviginti annos gravíssimis morbis et variis tentatiónibus vexata, constantíssime meruit in castris christianæ pœniténtiæ, nullo refecta pabulo cæléstium eárum consolatiónum quibus solet étiam in terris sanctitas abundare.\par
\switchcolumn
\EN
\noindent The virgin Theresa was the daughter of a father and mother, equally honourable on account of their birth and of their godliness, (and was born) at Avila (in the kingdom of Old Castile) in Spain, (on the 28th day of March, in the year of our Lord 1515.) She was brought up from the dawn of her life in the fear of God, and when still only seven years old she gave a startling fore-cast of the holy earnestness of her later years. The reading of the acts of the holy martyrs so inflamed and excited her imagination, that she ran away from her father's house, with the design of going to Morocco and the hope there to lay down her life for the glory of Christ Jesus and the salvation of souls. (Upon the bridge over the Adaja, near the town,) she was met by an uncle and brought back to her mother, and was fain to slake her thirst for martyrdom by giving to the poor all the alms she could, and by other godly exercises, though still ever bewailing with tears that the highest prize had been snatched from her. (In the twelfth year of her age,) her mother died, and she besought the Most Blessed Virgin to be a mother to her in her stead. This she gained; thenceforth she lived always as a daughter under the shelter of the Mother of God. In the twentieth year of her age she withdrew herself among the nuns of St. Mary of Mount Carmel. There she dwelt for two-and-twenty years, tormented by grievous sicknesses and diverse temptations, and so bravely served her time in the hardest ranks of Christ's army, starved even of that comforting knowledge of God's reconciled love, wherein His holy children are so commonly used even upon earth to rejoice.\par
\switchcolumn*

\LA
\begin{Resp}\Rsign Propter veritátem, et mansuetúdinem, et justítiam: \Star{} Et dedúcet te mirabíliter déxtera tua.\end{Resp}
\begin{Resp}\Vsign Spécie tua et pulchritúdine tua inténde, próspere procéde, et regna.\end{Resp}
\begin{Resp}\Rsign Et dedúcet te mirabíliter déxtera tua.\end{Resp}
\switchcolumn
\EN
\begin{Resp}\Rsign Because of truth, and meekness, and righteousness; \Star{} And thy right hand shall lead thee wonderfully.\end{Resp}
\begin{Resp}\Vsign In thy comeliness, and thy beauty, go forward, fare prosperously, and reign.\end{Resp}
\begin{Resp}\Rsign And thy right hand shall lead thee wonderfully.\end{Resp}
\switchcolumn*

\LA
\LessonHead{Lectio V}
\switchcolumn
\EN
\LessonHead{Lesson V}
\switchcolumn*

\LA
\noindent Angelicis ditáta virtútibus, non modo propriam, sed publicam étiam salútem sollicita caritate curávit. Quare severiórem veterum Carmelitárum regulam, Deo afflante et Pio quarto approbante, primum muliéribus, deínde viris observandam propósuit. Efflóruit in eo consílio omnípotens miserentis Dómini benedíctio; nam duo supra trigínta monasteria inops virgo pótuit ædificare, ómnibus humanis destituta auxíliis, quinimmo adversántibus plerumque sæculi princípibus. Infidelium et hæreticórum ténebras perpetuis deflebat lácrimis, atque, ad placandam divinæ ultiónis iram, voluntarios proprii corporis cruciátus Deo pro eórum salúte dicábat. Tanto autem divini amoris incendio cor ejus conflagrávit, ut mérito víderit Angelum igníto jáculo sibi præcordia transverberántem, et audíerit Christum, data déxtera, dicéntem sibi: Deinceps, ut vera sponsa, meum zelábis honórem. Eo consiliante, maxime arduum votum emisit efficiéndi semper quidquid perfectius esse intellígeret. Multa cæléstis sapiéntiæ documénta conscripsit, quibus fidelium mentes ad supernæ pátriæ desidérium maxime excitántur.\par
\switchcolumn
\EN
\noindent Strengthened in the graces of an angel, the wideness of her love embraced in its tender care the salvation of other souls as well as of her own. To this end, under the blessing of God, and the approbation of Pius IV., she set, first before women and then before men, the observance of the stern Rule of the Old Carmelites. The blessing of the Almighty and merciful Lord did indeed rest most evidently upon this design. This penniless virgin, helped by no man, and in the teeth of many that were great in this world, was enabled to build two-and-thirty houses. The darkness of unbelievers and misbelievers drew from her unceasing tears, and she willingly gave up her own body to God to be tortured, to soften the fury of His indignation against them. His own love so blazed in her heart that she attained to see an Angel run her through with a fiery spear, and Christ Himself take her by the hand, and to hear Him say Henceforth thou shalt love Mine honour as a wife indeed. At His inspiration she took the extremely difficult vow to do always that which should seem to her to be most perfect. She wrote much, full of heavenly wisdom, whereby the minds of the faithful are enkindled to long for the Fatherland above.\par
\switchcolumn*

\LA
\begin{Resp}\Rsign Dilexísti justítiam, et odísti iniquitátem: \Star{} Proptérea unxit te Deus, Deus tuus, óleo lætítiæ.\end{Resp}
\begin{Resp}\Vsign Propter veritátem, et mansuetúdinem, et justítiam.\end{Resp}
\begin{Resp}\Rsign Proptérea unxit te Deus, Deus tuus, óleo lætítiæ.\end{Resp}
\switchcolumn
\EN
\begin{Resp}\Rsign Thou hast loved righteousness, and hated iniquity; \Star{} Therefore God, thy God, hath anointed thee with the oil of gladness.\end{Resp}
\begin{Resp}\Vsign Because of truth, and meekness, and righteousness.\end{Resp}
\begin{Resp}\Rsign Therefore God, thy God, hath anointed thee with the oil of gladness.\end{Resp}
\switchcolumn*

\LA
\LessonHead{Lectio VI}
\switchcolumn
\EN
\LessonHead{Lesson VI}
\switchcolumn*

\LA
\noindent Cum autem assidua ederet exempla virtútum, tam anxio castigándi corporis desidério æstuabat, ut, quamvis secus suadérent morbi quibus afflictabátur, corpus ciliciis, catenis, urticárum manípulis aliisque asperrimis flagellis sæpe cruciaret, et aliquándo inter spinas volutaret, sic Deum álloqui sólita: Dómine, aut pati aut mori; se semper miserrima morte pereuntem existimans, quámdiu a cælésti ætérnæ vitæ fonte abesset. Prophetíæ dono excelluit, eamque divinis charismátibus tam liberáliter locupletábat Dóminus, ut sæpius exclámans péterent beneficiis in se divinis modum imponi, nec tam céleri oblivióne culpárum suárum memóriam aboleri. Intolerábili ígitur divini amoris incendio potius quam vi morbi, Albæ cum decúmberet, prænuntiáto suæ mortis die, ecclesiásticis sacramentis muníta, alumnos ad pacem, caritátem et regularem observantiam adhortata, sub colúmbæ specie puríssimam ánimam Deo réddidit, annos nata sexagínta septem annos millesimo quingentésimo octogesimo secundo, Idibus Octobris, juxta Kalendarii Romani emendatiónem. Ei moriénti adesse visus est inter Angelórum ágmina Christus Jesus; et arbor árida, cellæ próxima, statim efflóruit. Ejus corpus, usque ad hanc diem incorruptum, odorato liquore circumfusum, pia veneratióne cólitur. Miraculis cláruit ante et post óbitum, eamque Gregórius décimus quintus in Sanctórum númerum rétulit.\par
\switchcolumn
\EN
\noindent Earnest as were the examples of graces which she had shown, and grievous as was the state of her body, afflicted by disease, she still burnt with the desire of tormenting it. She tortured it with sackcloth, chains of spikes, handfuls of nettles, and heavy scourging. She rolled herself sometimes among thorns, and was used to cry to God “Lord! to suffer or to die.” As long as she remained exiled from the heavenly Fountain of eternal life, her life was to her a lingering death. She was eminent for the gift of prophecy, and God did indeed so pour forth His bounties upon her, that she often cried to Him in entreaty not to bless her so as to make her forget her sins. It was worn out rather by the fever of her love than by the wasting of disease that she sank upon her deathbed at Alba. She foretold the day of her own death, received the Sacraments of the Church, and exhorted her disciples to peace, love, and strictness in observing the Rule, and then her soul, like a pure dove, winged its flight to rest with God, on the 15th day of October in the year 1582, being then 67 years of age. At her death she had a vision of Christ Jesus surrounded by Angels. A dead tree hard by the cell instantly broke into foliage. Her body is untouched by corruption even unto this day, and lieth in a sort of perfumed oil, regarded with godly reverence. She was famous for miracles both before and after her death, and was numbered by Gregory XV. among the Saints.\par
\switchcolumn*

\LA
\begin{Resp}\Rsign Afferéntur Regi vírgines post eam, próximæ ejus; \Star{} Afferéntur tibi in lætítia et exsultatióne.\end{Resp}
\begin{Resp}\Vsign Spécie tua et pulchritúdine tua; inténde, próspere procéde, et regna.\end{Resp}
\begin{Resp}\Rsign Afferéntur tibi in lætítia et exsultatióne.\end{Resp}
\begin{Resp}\Vsign Glória Patri, et Fílio, \Star{} et Spirítui Sancto.\end{Resp}
\begin{Resp}\Rsign Afferéntur tibi in lætítia et exsultatióne.\end{Resp}
\switchcolumn
\EN
\begin{Resp}\Rsign After her shall virgins be brought unto the King, her fellows \Star{} They shall be brought unto thee with gladness and rejoicing.\end{Resp}
\begin{Resp}\Vsign In thy comeliness and thy beauty, go forward, fare prosperously, and reign.\end{Resp}
\begin{Resp}\Rsign They shall be brought unto thee with gladness and rejoicing.\end{Resp}
\begin{Resp}\Vsign Glory be to the Father, and to the Son, \Star{} and to the Holy Ghost.\end{Resp}
\begin{Resp}\Rsign They shall be brought unto thee with gladness and rejoicing.\end{Resp}
\switchcolumn*

\end{paracol}

\par\needspace{10\baselineskip}\addvspace{1.2em}{\centering\small\textsc{Die 16 Octobris} · \textsc{16 October}\par}
\Office{S. Hedwigis Viduæ}{St. Hedwig, Widow}{Semiduplex}
\addcontentsline{toc}{subsection}{Die 16 Octobris · S. Hedwigis Viduæ \textit{— St. Hedwig, Widow}}
\begin{paracol}{2}
\LA
\RefLine{Lectiones I–III: de Scriptura occurrente.}
\switchcolumn
\EN
\RefLine{Lessons I–III: from the Scripture of the day.}
\switchcolumn*

\LA
\RefLine{Lectiones VII–IX: ex Commúni non Virginum non Martyrum (C7a).}
\switchcolumn
\EN
\RefLine{Lessons VII–IX: from the Common of Widows (C7a).}
\switchcolumn*

\LA
\LessonHead{Lectio IV}
\switchcolumn
\EN
\LessonHead{Lesson IV}
\switchcolumn*

\LA
\noindent Hedwigis, regiis clara natálibus, innocéntia tamen vitæ longe clarior, sanctæ Elisabethæ fíliæ regis Hungariæ matértera, Bertholdi et Agnetis Moraviæ marchiónum fília, animi ab ineunte ætate moderatiónem prótulit. Adhuc enim puellula puerílibus abstinuit, et, duodennis Henrico Poloniæ duci a paréntibus nuptui trádita, thálami fide sancte servata, prolem inde susceptam in Dei timóre erudívit. Ut autem commodius Deo vacaret, ex pari voto et consénsu unánimi, ad separatiónem thori virum indúxit. Quo defuncto, ipsa in monasterio Trebnicénsi, Deo, quem assiduis precibus exoraverat, inspirante, Cisterciensem devota sumpsit habitum; in eoque, contemplatióni intenta, divinis Officiis et Missárum solemniis a solis ortu ad merídiem usque assidua assistens, antiquum humani generis hostem fortis contempsit.\par
\switchcolumn
\EN
\noindent Hedwiga, a Princess, in whom the splendour of her family was outshone by the radiant innocency of her life, was the daughter of Bertold and Agnes, Marquess and Marchioness of Moravia, and sister to Gertrude, wife of Andrew, King of Hungary, and mother of the holy Elizabeth of Thuringia. From her earliest childhood she was a very grave child, and had already done with childish things when, at twelve years of age, she was given in marriage by her father and mother to Henry, Grand Prince of Poland. In marriage she kept the bed in all holiness undefined, and brought up in the fear of God the children that were therein begotten of her. (After the birth of her sixth child,) she was fain to give herself more continually to God, and induced her husband to agree to a mutual vow of separation of bed-fellowship. After his death (in 1238,) by the inspiration of God, whom she besought in unceasing prayer, she clad herself for godliness' sake in the habit of a Cistercian nun in the monastery (which had been finished) at Trebnitz (in 1219.) She continued absorbed in God. She remained engaged in the Divine Office and hearing Masses from sunrise till noon, and trod mightily under foot the old enemy of man.\par
\switchcolumn*

\LA
\begin{Resp}\Rsign Propter veritátem, et mansuetúdinem, et justítiam: \Star{} Et dedúcet te mirabíliter déxtera tua.\end{Resp}
\begin{Resp}\Vsign Spécie tua et pulchritúdine tua inténde, próspere procéde, et regna.\end{Resp}
\begin{Resp}\Rsign Et dedúcet te mirabíliter déxtera tua.\end{Resp}
\switchcolumn
\EN
\begin{Resp}\Rsign Because of truth, and meekness, and righteousness; \Star{} And thy right hand shall lead thee wonderfully.\end{Resp}
\begin{Resp}\Vsign In thy comeliness, and thy beauty, go forward, fare prosperously, and reign.\end{Resp}
\begin{Resp}\Rsign And thy right hand shall lead thee wonderfully.\end{Resp}
\switchcolumn*

\LA
\LessonHead{Lectio V}
\switchcolumn
\EN
\LessonHead{Lesson V}
\switchcolumn*

\LA
\noindent Sæculi autem commercia, ni divina vel animárum salútem attingerent, audíre vel loqui non sustinuit. Prudéntia in agéndis sic emicuit, ut neque excéssus esset in modo nec error in ordine, comis alioqui et mansueta in próximum. Grandem autem de se triumphum jejuniis et vigiliis, vestiumque asperitáte, austéra carnem macerans, reportávit; hinc sublimióribus flores virtútibus christiánis, consiliórum gravitáte, animique candore et quiete, in exímium religiosæ pietátis evasit exemplar. Omnibus se ultro subjicere atque vilióra præ ceteris moniálibus alacriter munia subire, paupéribus étiam flexo genu ministrare, leprosórum pedes ablúere et osculari, ipsi familiare erat; neque illórum úlcera sanie manantia, sui victrix, abhorruit.\par
\switchcolumn
\EN
\noindent She could not bear to hear talk of worldly things, unless they had to do with the things of God or the saving of souls. She was very wise in business, not doing too much, nor unseasonably, and withal courteous and gentle toward all men. She got a great victory over herself by maltreating her flesh with fasting, watching, and rough clothing. She was an example of the higher Christian graces and of a godly nun, by the wisdom of her counsels, and the straightforwardness and peacefulness of her mind. It was her use to rank herself after all others, and cheerfully to undertake lower offices than those of the other nuns. She ministered to the poor even upon her knees, and washed and kissed the feet of lepers, having such command over herself as not to recoil from their sores oozing with matter.\par
\switchcolumn*

\LA
\begin{Resp}\Rsign Dilexísti justítiam, et odísti iniquitátem: \Star{} Proptérea unxit te Deus, Deus tuus, óleo lætítiæ.\end{Resp}
\begin{Resp}\Vsign Propter veritátem, et mansuetúdinem, et justítiam.\end{Resp}
\begin{Resp}\Rsign Proptérea unxit te Deus, Deus tuus, óleo lætítiæ.\end{Resp}
\switchcolumn
\EN
\begin{Resp}\Rsign Thou hast loved righteousness, and hated iniquity; \Star{} Therefore God, thy God, hath anointed thee with the oil of gladness.\end{Resp}
\begin{Resp}\Vsign Because of truth, and meekness, and righteousness.\end{Resp}
\begin{Resp}\Rsign Therefore God, thy God, hath anointed thee with the oil of gladness.\end{Resp}
\switchcolumn*

\LA
\LessonHead{Lectio VI}
\switchcolumn
\EN
\LessonHead{Lesson VI}
\switchcolumn*

\LA
\noindent Mira fuit ejus patiéntia animique constántia; præcipue vero in morte Henrici ducis Silesiæ, sui, quem materne diligebat, fílii in bello a Tartaris cæsi, enituit; potius enim grátias Deo, quam fílio lácrimas réddidit. Miraculórum denique glória percrebuit; púerum enim demersum et molendini rotis allisum et prorsus attritum invocata, vitæ restituit; alíaque præstitit, ut, rite iis Clemens quartus probátis, Sanctórum número eam adscripserit, ejusque festum in Polónia, ubi præcipua veneratióne uti patrona cólitur, die décima quinta Octobris celebrari concesserit. Quod deínde ut in tota Ecclésia fieret, Innocentius undecimus ampliávit.\par
\switchcolumn
\EN
\noindent Her long-suffering and endurance were very marvellous, especially when her son Henry, Duke of Silesia, to whom she bore a mother's love, was killed by the Tartars (in 1241.) His death drew from her rather thanksgiving to God than tears for him. (She died upon the 15th day of October, in the year 1243.) She was famous for miracles. One while, being called on, she restored to life a boy who had fallen into the water, been dashed against the wheels of a mill, and wholly crushed. This and the like being duly proved, Clement IV. numbered her name among those of the Saints, and allowed her Feastday to be kept in Poland, in which country, being Patroness, she hath most honour, upon the 15th of October; which permission was given to the whole Church by Innocent XI. (rubrica tridentina) for the 17th day of the same month.\par
\switchcolumn*

\LA
\begin{Resp}\Rsign Fallax grátia, et vana est pulchritúdo: \Star{} Múlier timens Dóminum, ipsa laudábitur.\end{Resp}
\begin{Resp}\Vsign Date ei de fructu mánuum suárum, et laudent eam in portis ópera ejus.\end{Resp}
\begin{Resp}\Rsign Múlier timens Dóminum, ipsa laudábitur.\end{Resp}
\begin{Resp}\Vsign Glória Patri, et Fílio, \Star{} et Spirítui Sancto.\end{Resp}
\begin{Resp}\Rsign Múlier timens Dóminum, ipsa laudábitur.\end{Resp}
\switchcolumn
\EN
\begin{Resp}\Rsign Favour is deceitful, and beauty is vain \Star{} A woman that feareth God she shall be praised.\end{Resp}
\begin{Resp}\Vsign Give her of the fruit of her hands, and let her own works praise her in the gates.\end{Resp}
\begin{Resp}\Rsign A woman that feareth God she shall be praised.\end{Resp}
\begin{Resp}\Vsign Glory be to the Father, and to the Son, \Star{} and to the Holy Ghost.\end{Resp}
\begin{Resp}\Rsign A woman that feareth God she shall be praised.\end{Resp}
\switchcolumn*

\end{paracol}

\par\needspace{10\baselineskip}\addvspace{1.2em}{\centering\small\textsc{Die 17 Octobris} · \textsc{17 October}\par}
\Office{S. Margaritæ Mariæ Alacoque Virginis}{St. Marguerite-Marie Alacoque, Virgin}{Duplex}
\addcontentsline{toc}{subsection}{Die 17 Octobris · S. Margaritæ Mariæ Alacoque Virginis \textit{— St. Marguerite-Marie Alacoque, Virgin}}
\begin{paracol}{2}
\LA
\RefLine{Lectiones I–III: de Scriptura occurrente.}
\switchcolumn
\EN
\RefLine{Lessons I–III: from the Scripture of the day.}
\switchcolumn*

\LA
\LessonHead{Lectio IV}
\switchcolumn
\EN
\LessonHead{Lesson IV}
\switchcolumn*

\LA
\noindent Margarita Maria Alacóque, in pago diœcesis Augustodunénsis, honesto genere nata, jam inde a teneris annis futuræ sanctitátis indícia præbuit. In Deiparam Vírginem et in augustum Eucharistiæ sacraméntum amóre flagrans, adolescentula Deo virginitátem devovit, id exoptans unice ut ad christianas virtútes vitam compóneret. In deliciis habebat prolixas preces rerumque cæléstium contemplatiónem, sui contemptum, patiéntiam in adversis, corporis afflictatiónem, caritátem in próximos, præsertim egénos; summoque studio nitebátur ut sanctíssima divini Redemptoris exempla pro viribus referret.\par
\switchcolumn
\EN
\noindent Margaret Mary Alacoque was born of a respectable family in a village in the diocese of Autun. From her earliest years she gave signs of holiness. Filled with a burning love of the Virgin Mother of God, and of the august mystery of the Eucharist, while still a young girl she dedicated her virginity to God. Above all else she strove to realize in her life the performance of Christian virtues. She delighted to spend continuous hours in prayers and in meditation upon the things of heaven. She was humble, and patient in adversity. She practised bodily penance. She was charitable towards her neighbours, especially the poor. By every means within her power she strove diligently to imitate the most holy example left by our divine Redeemer.\par
\switchcolumn*

\LA
\begin{Resp}\Rsign Propter veritátem, et mansuetúdinem, et justítiam: \Star{} Et dedúcet te mirabíliter déxtera tua.\end{Resp}
\begin{Resp}\Vsign Spécie tua et pulchritúdine tua inténde, próspere procéde, et regna.\end{Resp}
\begin{Resp}\Rsign Et dedúcet te mirabíliter déxtera tua.\end{Resp}
\switchcolumn
\EN
\begin{Resp}\Rsign Because of truth, and meekness, and righteousness; \Star{} And thy right hand shall lead thee wonderfully.\end{Resp}
\begin{Resp}\Vsign In thy comeliness, and thy beauty, go forward, fare prosperously, and reign.\end{Resp}
\begin{Resp}\Rsign And thy right hand shall lead thee wonderfully.\end{Resp}
\switchcolumn*

\LA
\LessonHead{Lectio V}
\switchcolumn
\EN
\LessonHead{Lesson V}
\switchcolumn*

\LA
\noindent Ordinem Visitatiónis ingressa, statim religiosæ vitæ fulgóre nitére cœpit. Altioris dono oratiónis a Deo est decorata, aliisque gratiæ munéribus et crebris visiónibus. Harum celeberrima fuit cum ante Eucharistiam precanti Jesus semetípsum conspiciéndum obtulit, et divinum Cor in aperto péctore flammis incensum ac spinis constrictum osténdit, præcepitque ut, ob talem caritátem et ad ingratórum hóminum injurias expiandas, illa publicum Cordi suo cultum, magnis propositis cæléstis thesáuri præmiis, instituéndum curaret. Cunctanti ex humilitate seque tantæ rei ímparem profiténti amantíssimus Salvátor addit ánimum, simulque exímia sanctitáte virum, Claudium de la Colombière, ducem et adjutórem designat; eamque spe fovet illíus summæ utilitátis, quæ póstea e divini Cordis cultu in Ecclésiam dimanávit.\par
\switchcolumn
\EN
\noindent Margaret entered the Order of the Visitation. There her life became immediately a shining example to others. God endowed her highly with the gift of prayer. He gave her other favours, such as frequent visitations. The most famous of these was that one when Jesus appeared to her as she knelt in prayer before the Blessed Sacrament. Opening his breast, He revealed his divine Heart glowing with flames and encircled with a crown of thorns. He bade her in return for his excessive love and in atonement for the insults of ungrateful men, to seek to have established public adoration of his Heart. This devotion he promised to enrich with treasures of heavenly grace. When, out of humility, she hesitated to undertake so great a task, the loving Saviour encouraged her. At the same time he pointed out Claude de la Colombiere, a man of great holiness, as one who could guide and help her. Our Lord also comforted her with the assurance that very great blessings would accrue afterwards to the Church from the worship of his divine Heart.\par
\switchcolumn*

\LA
\begin{Resp}\Rsign Dilexísti justítiam, et odísti iniquitátem: \Star{} Proptérea unxit te Deus, Deus tuus, óleo lætítiæ.\end{Resp}
\begin{Resp}\Vsign Propter veritátem, et mansuetúdinem, et justítiam.\end{Resp}
\begin{Resp}\Rsign Proptérea unxit te Deus, Deus tuus, óleo lætítiæ.\end{Resp}
\switchcolumn
\EN
\begin{Resp}\Rsign Thou hast loved righteousness, and hated iniquity; \Star{} Therefore God, thy God, hath anointed thee with the oil of gladness.\end{Resp}
\begin{Resp}\Vsign Because of truth, and meekness, and righteousness.\end{Resp}
\begin{Resp}\Rsign Therefore God, thy God, hath anointed thee with the oil of gladness.\end{Resp}
\switchcolumn*

\LA
\LessonHead{Lectio VI}
\switchcolumn
\EN
\LessonHead{Lesson VI}
\switchcolumn*

\LA
\noindent Ut jussa Redemptoris impleret Margarita omni diligéntia studebat. Nec tamen illi defuére molestiæ plures atque acres contumeliæ ab iis qui eam vano mentis errori obnoxiam esse dictitábant. Quæ ómnia æquo animo tulit, immo apponebat lucro, existimans se per oppróbria et dolóres hóstiam Deo gratam fore, et majora ad propositum suum auxília consecuturam. Religiosæ perfectiónis laude florens et per æternárum rerum contemplatiónem in dies síngulos cælésti sponso conjúnctior, ad eum evolávit, anno ætátis suæ quadragesimo tertio, reparátæ salútis millesimo sexcentésimo nonagesimo. Miraculis insignem Benedíctus décimus quintus Sanctis adscripsit; ejusque offícium Pius undecimus Pontifex maximus ad universam Ecclésiam extendit.\par
\switchcolumn
\EN
\noindent Margaret strove ardently to fulfill the Redeemer's command. Vexations, even bitter insults, were her portion from some who maintained that she was subject to mental aberrations. She not only bore these sufferings patiently, she even profited by them, offering herself in anguish and reproach as a victim acceptable to God, bearing all things as a more sure means of accomplishing her purpose. Renowned for her religious perfection, becoming each day more closely united with her divine spouse by contemplation of celestial things, she took flight to him in the forty-third year of her age, and in the year of restored salvation 1690. She became famous for miracles. Benedict XV added her to the list of the saints; Pius XI extended her office to the universal Church.\par
\switchcolumn*

\LA
\begin{Resp}\Rsign Afferéntur Regi vírgines post eam, próximæ ejus; \Star{} Afferéntur tibi in lætítia et exsultatióne.\end{Resp}
\begin{Resp}\Vsign Spécie tua et pulchritúdine tua; inténde, próspere procéde, et regna.\end{Resp}
\begin{Resp}\Rsign Afferéntur tibi in lætítia et exsultatióne.\end{Resp}
\begin{Resp}\Vsign Glória Patri, et Fílio, \Star{} et Spirítui Sancto.\end{Resp}
\begin{Resp}\Rsign Afferéntur tibi in lætítia et exsultatióne.\end{Resp}
\switchcolumn
\EN
\begin{Resp}\Rsign After her shall virgins be brought unto the King, her fellows \Star{} They shall be brought unto thee with gladness and rejoicing.\end{Resp}
\begin{Resp}\Vsign In thy comeliness and thy beauty, go forward, fare prosperously, and reign.\end{Resp}
\begin{Resp}\Rsign They shall be brought unto thee with gladness and rejoicing.\end{Resp}
\begin{Resp}\Vsign Glory be to the Father, and to the Son, \Star{} and to the Holy Ghost.\end{Resp}
\begin{Resp}\Rsign They shall be brought unto thee with gladness and rejoicing.\end{Resp}
\switchcolumn*

\LA
\LessonHead{Lectio VII}
\switchcolumn
\EN
\LessonHead{Lesson VII}
\switchcolumn*

\LA
\LessonTitle{Léctio sancti Evangélii secúndum Matthǽum}
\switchcolumn
\EN
\LessonTitle{From the Holy Gospel according to Matthew}
\switchcolumn*

\LA
\Cite{Matt 11:25-30}
\switchcolumn
\EN
\Cite{Matt 11:25-30}
\switchcolumn*

\LA
\noindent In illo témpore, respondens Jesus dixit: Confiteor tibi, Pater, Dómine cæli et terræ, quia abscondísti hæc a sapiéntibus et prudéntibus et revelásti ea párvulis. Et réliqua.\par
\switchcolumn
\EN
\noindent At that time: Jesus answered and said: I thank thee, O Father, Lord of heaven and earth, because thou hast hid these things from the wise and prudent, and hast revealed them unto babes. And so on, and that which followeth.\par
\switchcolumn*

\LA
\LessonTitle{Homilía sancti Francísci Salésii Epíscopi}
\switchcolumn
\EN
\LessonTitle{Homily of St. Francis de Sales, Bishop}
\switchcolumn*

\LA
\Source{Sermo 23 in die Pentecostes alter, circa medium}
\switchcolumn
\EN
\Source{Sermo 23 in die Pentecostes alter, circa medium}
\switchcolumn*

\LA
\noindent Nulla alia est vera sciéntia, nisi ea quæ a Spíritu Sancto datur, sed hæc humílibus tantummodo tribúitur. Nonne magnos vídimus theólogos qui mira dixérunt de virtútibus, sed ut eas non exercerent? E contra complures vídimus feminas, quæ de virtútibus dissérere nesciébant, sed virtútum ópera digne nóverant adimplére. Eas enim Spíritus Sanctus sapiéntes effécit, quia et timórem Dómini et pietátem et humilitátem habebant.\par
\switchcolumn
\EN
\noindent There is no other true knowledge given but that which is given by the Holy Ghost. And this is granted to the humble only. Have we not known great theologians who spake marvellously of the virtues but did not practice them? Have we not seen, on the contrary, many women who did not know how to discourse on the virtues, practice the works of virtue worthily? The Holy Spirit hath made these wise because they had fear of the Lord, piety and humility.\par
\switchcolumn*

\LA
\begin{Resp}\Rsign Hæc est Virgo sápiens, quam Dóminus vigilántem invénit, quæ accéptis lampádibus sumpsit secum óleum: \Star{} Et veniénte Dómino, introívit cum eo ad núptias.\end{Resp}
\begin{Resp}\Vsign Média nocte clamor factus est: Ecce sponsus venit, exíte óbviam ei.\end{Resp}
\begin{Resp}\Rsign Et veniénte Dómino, introívit cum eo ad núptias.\end{Resp}
\switchcolumn
\EN
\begin{Resp}\Rsign This is one of those wise virgins, whom the Lord found watching, for when she took her lamp, she took oil with her. \Star{} And when the Lord came, she went in with him to the marriage.\end{Resp}
\begin{Resp}\Vsign At midnight there was a cry made: Behold, the Bridegroom cometh, go ye out to meet him.\end{Resp}
\begin{Resp}\Rsign And when the Lord came, she went in with him to the marriage.\end{Resp}
\switchcolumn*

\LA
\LessonHead{Lectio VIII}
\switchcolumn
\EN
\LessonHead{Lesson VIII}
\switchcolumn*

\LA
\Source{Fragmentum sermo 16, in III Dominica post Pentecostem, initio}
\switchcolumn
\EN
\Source{Fragmentum sermo 16, in III Dominica post Pentecostem, initio}
\switchcolumn*

\LA
\noindent Dóminus noster, magnus et præclaríssimus ómnium nostrárum infirmitátum medicus, antequam in hunc mundum veníret, per prophétas suos palam nuntiáverat: Quod confractum fúerit alligábo, et quod infirmum fúerit consolidábo. Et deínde suo ipse ore clamávit dicens: Veníte ad me omnes qui laborátis et oneráti estis, et ego reficiam vos. Quid ígitur mirum si ipsum ab ægrotis, a peccatóribus et publicanis circúmdatum cernimus? Nonne médici glória est ab ægrotis exquiri?\par
\switchcolumn
\EN
\noindent Our Lord is the great and excellent physician of all our infirmities. Before he came into the world he announced openly to his prophets, I will bind up that which was broken, and I will strengthen that which is weak. Finally with his own lips he invited us, saying, Come unto me, all ye that labour and are burdened, and I will give you rest. What wonder is it, then, if we see him surrounded by the sick, by sinners, and by publicans. Is it not the glory of the physician to be sought out by the sick?\par
\switchcolumn*

\LA
\begin{Resp}\Rsign Média nocte clamor factus est: \Star{} Ecce sponsus venit, exíte óbviam ei.\end{Resp}
\begin{Resp}\Vsign Prudéntes vírgines, aptáte vestras lámpades.\end{Resp}
\begin{Resp}\Rsign Ecce sponsus venit, exíte óbviam ei.\end{Resp}
\begin{Resp}\Vsign Glória Patri, et Fílio, \Star{} et Spirítui Sancto.\end{Resp}
\begin{Resp}\Rsign Ecce sponsus venit, exíte óbviam ei.\end{Resp}
\switchcolumn
\EN
\begin{Resp}\Rsign At midnight there was a cry made: \Star{} Behold, the Bridegroom cometh, go ye out to meet him.\end{Resp}
\begin{Resp}\Vsign Trim your lamps, O ye wise virgins.\end{Resp}
\begin{Resp}\Rsign Behold, the Bridegroom cometh, go ye out to meet him.\end{Resp}
\begin{Resp}\Vsign Glory be to the Father, and to the Son, \Star{} and to the Holy Ghost.\end{Resp}
\begin{Resp}\Rsign Behold, the Bridegroom cometh, go ye out to meet him.\end{Resp}
\switchcolumn*

\LA
\LessonHead{Lectio IX}
\switchcolumn
\EN
\LessonHead{Lesson IX}
\switchcolumn*

\LA
\Source{Fragmentum serm. 10 pro Feria II post Pascha, in fine}
\switchcolumn
\EN
\Source{Fragmentum serm. 10 pro Feria II post Pascha, in fine}
\switchcolumn*

\LA
\noindent Fert ille nostras misérias et jas nobilitat, apponit misériam Cordi suo, osténdit latus. Sed eum redamémus oportet, alioquin qui præ amóre osténdit vulnera, semel osténdet præ ira et indignatióne. Fac, o bone Jesu, ut pacem, quam offers, accipiámus, videamúsque vulnera tua, ut quandóquidem manent fides, spes, caritas, fide radicati, spe gaudéntes et caritate fervéntes, exspectémus beátam spem et advéntum tuum, ita ut in illo Te, Agnum ad déxteram non leónem ad sinistram videámus; ac pro fide visiónem, pro spe possessiónem et pro caritate imperfécta perféctam habeámus, in qua gaudébimus in sæcula sæculórum. Amen.\par
\switchcolumn
\EN
\noindent The Lord beareth our miseries and ennobleth them. He layeth our miseries to his Heart; he sheweth his side. It is meet then, that we should make some return to him, lest he who now sheweth us his wounds out of love, may one day shew them with wrath and indignation. Grant, O good Jesu, that we may receive the peace which thou dost offer and let us see thy wounds. Since there remain faith, hope and charity, grant that, rooted in faith, rejoicing in hope, glowing in charity, we may await thy coming in the blessed expectation that we may see thee as the Lamb upon the right, and not as the lion on the left. May clear sight take the place of hope, and, to our imperfect charity may there succeed that perfect charity in which we may rejoice forever and ever. Amen.\par
\switchcolumn*

\LA
\TeDeum{Te Deum}
\switchcolumn
\EN
\TeDeum{Te Deum}
\switchcolumn*

\end{paracol}

\par\needspace{10\baselineskip}\addvspace{1.2em}{\centering\small\textsc{Die 18 Octobris} · \textsc{18 October}\par}
\Office{S. Lucæ Evangelistæ}{St. Luke the Evangelist}{Duplex II classis}
\addcontentsline{toc}{subsection}{Die 18 Octobris · S. Lucæ Evangelistæ \textit{— St. Luke the Evangelist}}
\begin{paracol}{2}
\LA
\RefLine{Lectiones I–III: ex Commúni Evangelistarum (C1a).}
\switchcolumn
\EN
\RefLine{Lessons I–III: from the Commune Evangelistarum (C1a).}
\switchcolumn*

\LA
\RefLine{Lectiones VII–IX: ex Commúni Evangelistarum (C1a).}
\switchcolumn
\EN
\RefLine{Lessons VII–IX: from the Commune Evangelistarum (C1a).}
\switchcolumn*

\LA
\LessonHead{Lectio IV}
\switchcolumn
\EN
\LessonHead{Lesson IV}
\switchcolumn*

\LA
\LessonTitle{Ex libro sancti Hierónymi Presbýteri de Scriptóribus Ecclesiásticis}
\switchcolumn
\EN
\LessonTitle{From the Book on Ecclesiastical Writers, written by St. Jerome, Priest (at Bethlehem)}
\switchcolumn*

\LA
\Source{Cap. 7}
\switchcolumn
\EN
\Source{Cap. 7}
\switchcolumn*

\LA
\noindent Lucas médicus Antiochénsis ut ejus scripta índicant, Græci sermónis non ignárus, fuit sectátor Apóstoli Pauli, et omnis peregrinatiónis ejus comes. Scripsit Evangélium, de quo idem Paulus: Mísimus, inquit, cum illo fratrem, cujus laus est in Evangélio per omnes ecclésias. Et ad Colossénses: Salútat vos Lucas, médicus caríssimus. Et ad Timótheum: Lucas est mecum solus. Aliud quoque édidit volúmen egrégium, quod título, Acta Apostolórum, prænotátur: cujus história usque ad biénnium Romæ commorántis Pauli pervénit, id est, usque ad quartum Nerónis annum. Ex quo intellégimus in eádem urbe librum esse compósitum.\par
\switchcolumn
\EN
\noindent Luke was a physician of Antioch, who, as appeareth from his writings, knew the Greek language. He was a follower of the Apostle Paul, and his fellow traveller in all his wanderings. He wrote a Gospel, whereof the same Paul saith We have sent with him the brother, whose praise is in the Gospel throughout all the Churches (2 Cor. viii, 18) Of him, he writeth unto the Colossians, (iv, 14): Luke, the beloved physician, greeteth you. And again, unto Timothy, (2 Cor. iv, 11): Only Luke is with me. He also published another excellent book entitled The Acts of the Apostles, wherein the history is brought down to Paul's two years sojourn at Rome, that is to say, until the fourth year of Nero, from which we gather that it was at Rome that the said book was composed.\par
\switchcolumn*

\LA
\begin{Resp}\Rsign Vidi conjúnctos viros, habéntes spléndidas vestes, et Ángelus Dómini locútus est ad me, dicens: \Star{} Isti sunt viri sancti facti amíci Dei.\end{Resp}
\begin{Resp}\Vsign Vidi Ángelum Dei fortem, volántem per médium cælum, voce magna clamántem et dicéntem.\end{Resp}
\begin{Resp}\Rsign Isti sunt viri sancti facti amíci Dei.\end{Resp}
\switchcolumn
\EN
\begin{Resp}\Rsign I saw men standing together, clad in shining raiment, and the Angel of the Lord spoke unto me, saying: \Star{} These men are holy, for they are the friends of God.\end{Resp}
\begin{Resp}\Vsign I saw a strong Angel of God fly into the midst of heaven, saying with a loud voice:\end{Resp}
\begin{Resp}\Rsign These men are holy, for they are the friends of God.\end{Resp}
\switchcolumn*

\LA
\LessonHead{Lectio V}
\switchcolumn
\EN
\LessonHead{Lesson V}
\switchcolumn*

\LA
\noindent Igitur períodos Pauli et Theclæ, et totam baptizáti Leónis fábulam, inter apócryphas scriptúras computámus. Quale enim est, ut indivíduus comes Apóstoli, inter céteras ejus res, hoc solum ignoráverit? Sed et Tertulliánus, vicínus eórum témporum, refert presbýterum quemdam in Asia, amatórem Apóstoli Pauli, convíctum a Joánne quod auctor esset libri, et conféssum se hoc Pauli amóre fecísse, et ob id loco excidísse. Quidam suspicántur, quotiescúmque in epístolis suis Paulus dicit, Juxta Evangélium meum, de Lucæ significáre volúmine.\par
\switchcolumn
\EN
\noindent The silence of Luke is one of the reasons why we reckon among Apocryphal books The Acts of Paul and Thecla, and the whole story about the baptism of Leo. For why should the fellow traveller of the Apostle, who knew other things, be ignorant only of this? At the same time there is against these documents the statement of Tertullian, almost a contemporary writer, that the Apostle John convicted a certain Priest in Asia, who was a great admirer of the Apostle Paul, of having written them, and that the said Priest owned that he had been induced to compose them through his admiration for Paul, and that he was deposed in consequence. There are some persons who suspect that when Paul in his Epistles useth the phrase: “According to my Gospel” (Rom. ii, 16; Tim. ii, 8,) he meaneth the Gospel written by Luke.\par
\switchcolumn*

\LA
\begin{Resp}\Rsign Beáti estis, cum maledíxerint vobis hómines, et persecúti vos fúerint, et díxerint omne malum advérsum vos, mentiéntes, propter me: \Star{} Gaudéte et exsultáte, quóniam merces vestra copiósa est in cælis.\end{Resp}
\begin{Resp}\Vsign Cum vos óderint hómines, et cum separáverint vos, et exprobráverint, et ejécerint nomen vestrum tamquam malum propter Fílium hóminis.\end{Resp}
\begin{Resp}\Rsign Gaudéte et exsultáte, quóniam merces vestra copiósa est in cælis.\end{Resp}
\switchcolumn
\EN
\begin{Resp}\Rsign Blessed are ye when men shall revile you, and persecute you, and shall say all manner of evil against you falsely, for My sake; \Star{} Rejoice, and be exceeding glad, for great is your reward in heaven.\end{Resp}
\begin{Resp}\Vsign When men shall hate you, and when they shall separate you from their company, and shall reproach you, and cast out your name as evil, for the Son of Man's sake.\end{Resp}
\begin{Resp}\Rsign Rejoice, and be exceeding glad, for great is your reward in heaven.\end{Resp}
\switchcolumn*

\LA
\LessonHead{Lectio VI}
\switchcolumn
\EN
\LessonHead{Lesson VI}
\switchcolumn*

\LA
\noindent Lucam autem, non solum ab Apóstolo Paulo didicísse Evangélium, qui cum Dómino in carne non fúerat, sed et a céteris Apóstolis; quod ipse quoque in princípio sui volúminis declárat, dicens: Sicut tradidérunt nobis, qui a princípio ipsi vidérunt et minístri fuérunt sermónis. Igitur Evangélium, sicut audíerat, scripsit; Acta vero Apostolórum, sicut víderat ipse, compósuit. Vixit octogínta et quátuor annos, uxórem non habens. Sepultus est Constantinópoli, ad quam urbem, vigésimo Constantíni anno, ossa ejus cum relíquiis Andréæ Apóstoli transláta sunt de Achája.\par
\switchcolumn
\EN
\noindent However, Luke learned his Gospel not from the Apostle Paul only, who had not companied with the Lord in the flesh, but also from other Apostles, as himself declareth at the beginning of his work, where he saith: “They delivered them unto us, which from the beginning were eyewitnesses and ministers of the word.” (i. 2) According to what he had heard, therefore, did he write his Gospel. As to the Acts of the Apostles, he composed them from his own personal knowledge. He was never married. He lived eighty-four years. He is buried at Constantinople, whither his bones were brought from Achaia in the twentieth year of Constantine, together with the relics of the Apostle Andrew.\par
\switchcolumn*

\LA
\begin{Resp}\Rsign Isti sunt triumphatóres et amíci Dei, qui contemnéntes jussa príncipum, meruérunt prǽmia ætérna: \Star{} Modo coronántur, et accípiunt palmam.\end{Resp}
\begin{Resp}\Vsign Isti sunt qui venérunt ex magna tribulatióne, et lavérunt stolas suas in sánguine Agni.\end{Resp}
\begin{Resp}\Rsign Modo coronántur, et accípiunt palmam.\end{Resp}
\begin{Resp}\Vsign Glória Patri, et Fílio, \Star{} et Spirítui Sancto.\end{Resp}
\begin{Resp}\Rsign Modo coronántur, et accípiunt palmam.\end{Resp}
\switchcolumn
\EN
\begin{Resp}\Rsign These are they which have conquered, and are become the friends of God, who recked not of the commandments of princes, and earned the everlasting reward. \Star{} And now have they crowns on their heads, and palms in their hands.\end{Resp}
\begin{Resp}\Vsign These are they which came out of great tribulation, and have washed their robes in the blood of the Lamb.\end{Resp}
\begin{Resp}\Rsign And now have they crowns on their heads, and palms in their hands.\end{Resp}
\begin{Resp}\Vsign Glory be to the Father, and to the Son, \Star{} and to the Holy Ghost.\end{Resp}
\begin{Resp}\Rsign And now have they crowns on their heads, and palms in their hands.\end{Resp}
\switchcolumn*

\end{paracol}

\par\needspace{10\baselineskip}\addvspace{1.2em}{\centering\small\textsc{Die 19 Octobris} · \textsc{19 October}\par}
\Office{S. Petri de Alcantara Confessoris}{St. Peter of Alcantara, Confessor}{Duplex}
\addcontentsline{toc}{subsection}{Die 19 Octobris · S. Petri de Alcantara Confessoris \textit{— St. Peter of Alcantara, Confessor}}
\begin{paracol}{2}
\LA
\RefLine{Lectiones I–III: de Scriptura occurrente.}
\switchcolumn
\EN
\RefLine{Lessons I–III: from the Scripture of the day.}
\switchcolumn*

\LA
\LessonHead{Lectio IV}
\switchcolumn
\EN
\LessonHead{Lesson IV}
\switchcolumn*

\LA
\noindent Petrus, Alcántaræ in Hispánia nobílibus paréntibus natus, a teneris annis futuræ sanctitátis indícia prǽbuit. Décimo sexto ætátis anno ordinem Minórum ingréssus, se ómnium virtutum exemplar exhibuit. Tum munus concionatoris ex obediéntia exercens, innumeros a vitiis ad veram pœniténtiam traduxit. Primævum sancti Francisci institutum exactíssime reparáre cupiens, ope divina fretus et apostolica munítus auctoritate, angustíssimum et pauperrimum cœnobium juxta Pedrósum fundávit; quod vitæ genus asperrimum, ibi feliciter cœptum, per diversas Hispaniæ provincias usque ad Indias mirifice propagátum fuit. Sanctæ Teresiæ, cujus probaverat spíritum, in promovenda Carmelitárum reformatióne adjútor fuit. Ipsa autem, a Deo edocta quod Petri nómine nihil quisquam péteret quin prótinus exaudirétur, ejus precibus se commendare et ipsum, adhuc viventem, sanctum appellare consuevit.\par
\switchcolumn
\EN
\noindent Peter was born at Alcantara, (a small town in the Province of Estramadura,) in Spain, (in the year of our Lord 1499.) His father, (Alphonso Garavito, was a lawyer and Governor of the town,) and his mother (was) of good extraction. The holiness of his life was foreshadowed from his earliest years. In the sixteenth year of his age he entered the Order of Friars Minor, wherein he showed himself a pattern to all. He undertook the work of preaching in obedience to his Superiors, and thereby brought many to turn away from sin to true repentance. He conceived a great desire to bring back the observance of the Rule of St. Francis to the uttermost straitness of old times, and to that end, supported by God's help, and armed with the approval of the Apostolic See, he founded (in the year 1555) a new stern and poor house near Pedraso, from which the harder way of life, therein happily begun, spread marvellously through diverse Provinces of Spain even to the Indies. He was an helper to holy Theresa, with whom he was like-minded, in bringing about the Reformation of the Carmelites. She was taught of God that no one should ask anything in the name of Peter without being heard, and was used to ask him to pray for her, and to call him a Saint while as he was yet alive.\par
\switchcolumn*

\LA
\begin{Resp}\Rsign Honéstum fecit illum Dóminus, et custodívit eum ab inimícis, et a seductóribus tutávit illum: \Star{} Et dedit illi claritátem ætérnam.\end{Resp}
\begin{Resp}\Vsign Justum dedúxit Dóminus per vias rectas, et osténdit illi regnum Dei.\end{Resp}
\begin{Resp}\Rsign Et dedit illi claritátem ætérnam.\end{Resp}
\switchcolumn
\EN
\begin{Resp}\Rsign The Lord made him honourable, and defended him from his enemies, and kept him safe from those that lay in wait for him; \Star{} And gave him perpetual glory.\end{Resp}
\begin{Resp}\Vsign He went down with him into the pit, and left him not in bonds.\end{Resp}
\begin{Resp}\Rsign And gave him perpetual glory.\end{Resp}
\switchcolumn*

\LA
\LessonHead{Lectio V}
\switchcolumn
\EN
\LessonHead{Lesson V}
\switchcolumn*

\LA
\noindent Principium obsequia, qui ipsum velut oraculum consulébant, summa humilitate declinans, Cárolo quinto imperatóri a confessiónibus esse recusávit. Paupertátis rigidíssimus custos, una tunia, qua nulla deterior esset, contentus erat. Puritátem ita coluit, ut a fratre in extremo morbo sibi inserviénte ne leviter quidem tangi passus sit. Corpus suum perpetuis vigiliis, jejuniis, flagellis, frigore, nuditate atque omni genere asperitátum in servitútem redégit; cum quo pactum iníerat ne ullam in hoc sæculo ei requiem præbéret. Caritas Dei et próximi, in ejus corde diffúsa, tantum quandoque excitábat incendium, ut e cellæ angustiis in apertum campum prosilire, aërisque refrigerio conceptum ardórem temperare cogerétur.\par
\switchcolumn
\EN
\noindent He humbly excused himself from accepting the courtesies of princes, by whom his advice was sought as that of an oracle, and declined to become the Confessor of the Emperor Charles V. He was a very careful keeper to poverty, and contented himself with a single tunic than which none was worse. Purity he carried to such a point that when he was lying sick of his last illness, he would not allow the brother who ministered to him to touch him, how lightly soever. He brought his body into bondage by unceasing watching, fasting, scourging, cold, nakedness, and all manner of hardships, having made it a promise never to allow it any rest in this world. The love of God and his neighbour, which was shed abroad in his heart, somewhiles burnt so that he was fain to run from his cell into the open air to cool himself.\par
\switchcolumn*

\LA
\begin{Resp}\Rsign Amávit eum Dóminus, et ornávit eum: stolam glóriæ índuit eum, \Star{} Et ad portas paradísi coronávit eum.\end{Resp}
\begin{Resp}\Vsign Índuit eum Dóminus lorícam fídei, et ornávit eum.\end{Resp}
\begin{Resp}\Rsign Et ad portas paradísi coronávit eum.\end{Resp}
\switchcolumn
\EN
\begin{Resp}\Rsign The Lord loved him and beautified him; He clothed him with a robe of glory, \Star{} And crowned him at the gates of Paradise.\end{Resp}
\begin{Resp}\Vsign The Lord hath put on him the breast-plate of faith, and hath adorned him.\end{Resp}
\begin{Resp}\Rsign And crowned him at the gates of Paradise.\end{Resp}
\switchcolumn*

\LA
\LessonHead{Lectio VI}
\switchcolumn
\EN
\LessonHead{Lesson VI}
\switchcolumn*

\LA
\noindent Gratia contemplatiónis admirábilis in eo fuit: qua cum assidue spíritus reficerétur, intérdum áccidit ut ab omni cibo et potu plúribus diébus abstinúerit. In áëra frequenter sublatus, miro fulgóre coruscare visus est. Rápidos flúvios sicco pede trajecit. Fratres in extrema penuria, cælitus deláta alimónia, cibávit. Báculus, ab ipso terræ defixus, mox in viridem ficulneam excrevit. Cum noctu iter ágeret, densa nive cadente, dírutam domum sine tecto ingréssus est, eique nix in áëre péndula pro tecto fuit, ne illíus copia suffocarétur. Dono prophetíæ ac discretiónis spirituum imbutum fuísse sancta Teresia testátur. Denique, annum agens sexagesimum tertium, hora qua prædixerat, migrávit ad Dóminum, mirábili visióne Sanctorúmque præséntia confortatus. Quem eodem momento in cælum ferri beáta Teresia procul distans vidit; cui póstea appárens dixit: O felix pœniténtia, quæ tantam mihi promeruit glóriam! Post mortem vero plurimis miraculis cláruit, et a Cleménte nono Sanctórum número adscriptus est.\par
\switchcolumn
\EN
\noindent He was marvellous how his thoughts became altogether rapt in God, so that somewhiles it befell that he neither ate nor drank for the space of several days. He was oftentimes seen to rise into the air, shining with an unearthly glory. He passed dry-shod over torrents. When his brethren were in the last state of need, he fed them with food from heaven. A staff which he fixed in the earth grew presently into a green fig-tree. Once while he was travelling by night in the midst of an heavy snow-storm, and took refuge in a ruined and roofless house, then the falling snow made a roof over him lest he should be overwhelmed. Holy Theresa beareth witness that he had the gift of prophecy and of the discerning of spirits. At length, in the 63rd year of his own age, (and of salvation 1562,) at the hour which he had himself foretold, (upon the 18th day of October,) he passed away to be for ever with the Lord, cheered in his last moments by a wonderful vision and by the presence of Saints. At the instant of his death, blessed Theresa, then afar off, saw him carried to heaven. He appeared to her afterwards, and said O what happy penance, to have won for me such glory! After his death he became famous for very many miracles, and Clement IX. inscribed his name among those of the Saints.\par
\switchcolumn*

\LA
\begin{Resp}\Rsign Iste homo perfécit ómnia quæ locútus est ei Deus, et dixit ad eum: Ingrédere in réquiem meam: \Star{} Quia te vidi justum coram me ex ómnibus géntibus.\end{Resp}
\begin{Resp}\Vsign Iste est, qui contémpsit vitam mundi, et pervénit ad cæléstia regna.\end{Resp}
\begin{Resp}\Rsign Quia te vidi justum coram me ex ómnibus géntibus.\end{Resp}
\begin{Resp}\Vsign Glória Patri, et Fílio, \Star{} et Spirítui Sancto.\end{Resp}
\begin{Resp}\Rsign Quia te vidi justum coram me ex ómnibus géntibus.\end{Resp}
\switchcolumn
\EN
\begin{Resp}\Rsign This is he which did according unto all that God commanded him; and God said unto him: Enter thou into My rest, \Star{} For thee have I seen righteous before Me among all people.\end{Resp}
\begin{Resp}\Vsign This is he which loved not his life in this world, and is come unto an everlasting kingdom.\end{Resp}
\begin{Resp}\Rsign For thee have I seen righteous before Me among all people.\end{Resp}
\begin{Resp}\Vsign Glory be to the Father, and to the Son, \Star{} and to the Holy Ghost.\end{Resp}
\begin{Resp}\Rsign For thee have I seen righteous before Me among all people.\end{Resp}
\switchcolumn*

\LA
\LessonHead{Lectio VII}
\switchcolumn
\EN
\LessonHead{Lesson VII}
\switchcolumn*

\LA
\LessonTitle{Léctio sancti Evangélii secúndum Lucam}
\switchcolumn
\EN
\LessonTitle{From the Holy Gospel according to Luke}
\switchcolumn*

\LA
\Cite{Luc 12:32-34}
\switchcolumn
\EN
\Cite{Luke 12:32-35}
\switchcolumn*

\LA
\noindent In illo témpore: Dixit Jesus discípulis suis: Nolíte timére pusíllus grex, quia complácuit Patri vestro dare vobis regnum. Et réliqua.\par
\switchcolumn
\EN
\noindent At that time, Jesus said unto His disciples: Fear not, little flock, for it is your Father's good pleasure to give you the kingdom. And so on.\par
\switchcolumn*

\LA
\LessonTitle{Homilía sancti Bedæ Venerábilis Presbýteri}
\switchcolumn
\EN
\LessonTitle{Homily by the Venerable Bede, Priest at Jarrow, and Doctor of the Church.}
\switchcolumn*

\LA
\Source{Lib. 4, Cap. 54, in Luc. 12}
\switchcolumn
\EN
\Source{Bk. iv. Ch. 54 on Luke xii.}
\switchcolumn*

\LA
\noindent Pusíllum gregem electórum, vel ob comparatiónem majóris númeri reprobórum, vel pótius ob humilitátis devotiónem nóminat: quia vidélicet Ecclésiam suam quantálibet numerositáte jam dilatátam, tamen usque ad finem mundi humilitáte vult créscere, et ad promíssum regnum humilitáte perveníre. Ideóque ejus labóres blande consolátus, quam regnum Dei tantum quærére prǽcipit, eídem regnum a Patre dandum complácita benignitáte promíttit.\par
\switchcolumn
\EN
\noindent The elect are called a little flock, perchance because the reprobate are far more in number than they, but, more probably, because they love to be lowly, since it is God's will that however much His Church should grow in numbers, she should grow with lowliness even unto the end of the world, and should enter lowly into that kingdom which is hers by His promise. That kingdom He promiseth to her here, when He biddeth her to seek only the kingdom of God, and, to comfort her in her travail, He doth so sweetly and so graciously say that her Father will give it to her.\par
\switchcolumn*

\LA
\begin{Resp}\Rsign Iste est qui ante Deum magnas virtútes operátus est, et de omni corde suo laudávit Dóminum: \Star{} Ipse intercédat pro peccátis ómnium populórum.\end{Resp}
\begin{Resp}\Vsign Ecce homo sine queréla, verus Dei cultor, ábstinens se ab omni ópere malo, et pérmanens in innocéntia sua.\end{Resp}
\begin{Resp}\Rsign Ipse intercédat pro peccátis ómnium populórum.\end{Resp}
\switchcolumn
\EN
\begin{Resp}\Rsign This is he which wrought great wonders before God, and praised the Lord with all his heart. \Star{} May he pray for all people, that their sins may be forgiven unto them.\end{Resp}
\begin{Resp}\Vsign Behold a man without blame, a worshipper of God in truth, keeping himself clean from every evil work, and abiding still in his innocency.\end{Resp}
\begin{Resp}\Rsign May he pray for all people, that their sins may be forgiven unto them.\end{Resp}
\switchcolumn*

\LA
\LessonHead{Lectio VIII}
\switchcolumn
\EN
\LessonHead{Lesson VIII}
\switchcolumn*

\LA
\noindent Véndite quæ possidétis, et date eleemósynam. Nolíte, inquit, timére, ne propter regnum Dei militántibus, hujus vitæ necessária desint: quin étiam posséssa propter eleemósynam véndite. Quod tunc digne fit, quando quis semel pro Dómino suis ómnibus spretis, nihilóminus post hæc labóre mánuum, unde et victum transígere, et eleemósynam dare queat, operátur. Unde gloriátur Apóstolus, dicens: Argéntum, et aurum, aut vestem nullíus concupívi: ipsi scitis, quóniam ad ea quæ mihi opus erant, et his qui mecum sunt, ministravérunt manus istæ. Ómnia osténdi vobis, quóniam sic laborántes opórtet suscípere infírmos.\par
\switchcolumn
\EN
\noindent Sell what ye have and give alms. Fear not, He saith, lest, while ye fight for the kingdom of God, ye should lack such things as are needful for this life, nay rather, sell even that which ye have, and give alms. This doth, whosoever for the Lord's sake leaveth all that he hath, and then worketh with his hands, that so he may have to eat, and withal to give alms. In this doth the Apostle boast himself, saying I have coveted no man's silver, or gold, or apparel, as ye yourselves know for these hands have ministered unto my necessities, and to them that were with me. I have showed you all things, how that so labouring ye ought to support the weak. (Acts xx. 33, 34, 35.)\par
\switchcolumn*

\LA
\begin{Resp}\Rsign Sint lumbi vestri præcíncti, et lucérnæ ardéntes in mánibus vestris: \Star{} Et vos símiles homínibus exspectántibus dóminum suum, quando revertátur a núptiis.\end{Resp}
\begin{Resp}\Vsign Vigiláte ergo, quia nescítis qua hora Dóminus vester ventúrus sit.\end{Resp}
\begin{Resp}\Rsign Et vos símiles homínibus exspectántibus dóminum suum, quando revertátur a núptiis.\end{Resp}
\begin{Resp}\Vsign Glória Patri, et Fílio, \Star{} et Spirítui Sancto.\end{Resp}
\begin{Resp}\Rsign Et vos símiles homínibus exspectántibus dóminum suum, quando revertátur a núptiis.\end{Resp}
\switchcolumn
\EN
\begin{Resp}\Rsign Let your loins be girded about, and your lights burning; \Star{} And ye yourselves like unto men that wait for their lord, when he will return from the wedding.\end{Resp}
\begin{Resp}\Vsign Watch therefore, for ye know not what hour your Lord doth come.\end{Resp}
\begin{Resp}\Rsign And ye yourselves like unto men that wait for their lord, when he will return from the wedding.\end{Resp}
\begin{Resp}\Vsign Glory be to the Father, and to the Son, \Star{} and to the Holy Ghost.\end{Resp}
\begin{Resp}\Rsign And ye yourselves like unto men that wait for their lord, when he will return from the wedding.\end{Resp}
\switchcolumn*

\LA
\LessonHead{Lectio IX}
\switchcolumn
\EN
\LessonHead{Lesson IX}
\switchcolumn*

\LA
\noindent Fácite vobis sácculos qui non veteráscunt: eleemósynas vidélicet operándo quarum merces in ætérnum máneat. Ubi non hoc præcéptum esse putándum est, ut nil pecúniæ reservétur a sanctis, vel suis scílicet, vel páuperum úsibus suggeréndæ: cum et ipse Dóminus, cui ministrábant Ángeli, tamen ad informándam Ecclésiam suam lóculos habuísse legátur, et a fidélibus obláta consérvans, et suórum necessitátibus aliísque indigéntibus tríbuens; sed ne Deo propter ista serviátur, et ob inópiæ timórem justítia deserátur.\par
\switchcolumn
\EN
\noindent Provide yourselves bags which wax not old that is to say, by almsgiving, the reward thereof remaineth for ever. Nevertheless, we must not think here that this commandment forbiddeth the Saints to keep money for their own use, and for helping of the poor. The Lord Himself, to Whom Angels ministered, had a bag, and kept therein that which the faithful people gave unto Him (John xii. 6,) to relieve therewith the need of His disciples, and other poor folk. But we are commanded not to serve God for gain, nor to work unrighteousness for fear of poverty.\par
\switchcolumn*

\LA
\TeDeum{Te Deum}
\switchcolumn
\EN
\TeDeum{Te Deum}
\switchcolumn*

\end{paracol}

\par\needspace{10\baselineskip}\addvspace{1.2em}{\centering\small\textsc{Die 20 Octobris} · \textsc{20 October}\par}
\Office{S. Joannis Cantii Confessoris}{St. John Cantius, Confessor}{Duplex}
\addcontentsline{toc}{subsection}{Die 20 Octobris · S. Joannis Cantii Confessoris \textit{— St. John Cantius, Confessor}}
\begin{paracol}{2}
\LA
\RefLine{Lectiones I–III: de Scriptura occurrente.}
\switchcolumn
\EN
\RefLine{Lessons I–III: from the Scripture of the day.}
\switchcolumn*

\LA
\RefLine{Lectiones VII–IX: ex Commúni Confessoris non pontificis (C5).}
\switchcolumn
\EN
\RefLine{Lessons VII–IX: from the Common of a Confessor not a Bishop (C5).}
\switchcolumn*

\LA
\LessonHead{Lectio IV}
\switchcolumn
\EN
\LessonHead{Lesson IV}
\switchcolumn*

\LA
\noindent Joánnes, in oppido Kenty Cracoviénsis diœcesis, a quo Cantii cognomen duxit, Stanislao et Anna piis et honestis paréntibus natus, morum suavitáte, innocéntia, gravitáte, ab ipsa infántia spem fecit maximæ virtútis. In universitáte Cracoviénsi philosophíæ ac theologíæ primum auditor, tum, per omnes academíæ gradus ascendendo, professor ac doctor, sacra, quam annis multis trádidit, doctrina mentes audiéntium non illustrábat modo, sed et ad omnem pietátem inflammabat, simul docens scilicet et fáciens. Sacerdos factus, nihil de litterárum studio remíttens, studium auxit christianæ perfectiónis. Utque passim offéndi Deum maxime dolebat, sic eum sibi et pópulo placare, oblato quotídie non sine multis lácrimis incruento sacrifício, satagebat. Ilkusiensem parochiam annis aliquot egregie administrávit; sed animárum periculo commótus, póstea dimísit, ac, postulante academía, ad prístinum docéndi offícium rediit.\par
\switchcolumn
\EN
\noindent This John was the son of godly and respectable parents named Stanislaus and Anne, and was born (in the year of our Lord 1397,) in the town of Kenty, a place in the diocese of Krakow in Poland, from which he took the Latin name of Cantius. By his gentleness, innocency, and seriousness he gave great hopes even from his childhood. He studied Philosophy and Theology in the University of Krakow, wherein he rose step by step to be a Professor and teacher of those sciences wherein he lectured many years, not only enlightening the minds of his hearers, but stirring up in them all godliness, instructing them by example as well as by word. Having taken Priests' orders, he ceased not to busy himself with letters, but added thereto the striving after Christian perfection. He grieved exceedingly that God should be offended on all hands, and offered up to Him, day by day, not without many tears, the Unbloody Sacrifice for a propitiation for himself and for his people. He was for some years a faithful Parish Priest at Ilkusi, but after a while gave it up for fear of the danger of souls, and accepted the call of the University to take up again his Professorship.\par
\switchcolumn*

\LA
\begin{Resp}\Rsign Honéstum fecit illum Dóminus, et custodívit eum ab inimícis, et a seductóribus tutávit illum: \Star{} Et dedit illi claritátem ætérnam.\end{Resp}
\begin{Resp}\Vsign Justum dedúxit Dóminus per vias rectas, et osténdit illi regnum Dei.\end{Resp}
\begin{Resp}\Rsign Et dedit illi claritátem ætérnam.\end{Resp}
\switchcolumn
\EN
\begin{Resp}\Rsign The Lord made him honourable, and defended him from his enemies, and kept him safe from those that lay in wait for him; \Star{} And gave him perpetual glory.\end{Resp}
\begin{Resp}\Vsign He went down with him into the pit, and left him not in bonds.\end{Resp}
\begin{Resp}\Rsign And gave him perpetual glory.\end{Resp}
\switchcolumn*

\LA
\LessonHead{Lectio V}
\switchcolumn
\EN
\LessonHead{Lesson V}
\switchcolumn*

\LA
\noindent Quotidie témporis ab studio supérerat, partim salúti proximórum sacris præsertim conciónibus curandæ, partim oratióni dabat, in qua cæléstibus quandoque visiónibus et colloquiis dignátus fertur. Christi vero passióne sic afficiebátur, ut in ea contemplanda totas intérdum noctes dúceret insomnes, ejusque causa mélius recolendæ Jerosolymam peregrinátus sit; ubi, et martyrii desidério flagrans, Turcis ipsis Christum crucifixum prædicáre non dubitávit. Quater étiam ad Apostolórum limina, pedes atque viaria onustus sárcina, Romam venit, tum ut Sedem apostolicam, cui maxime addictus fuit, honoraret, tum ut sui (sic enim ajebat) purgatorii pœnas expósita illic quotídie peccatórum venia redímeret. Quo in itinere a latrónibus olim spoliátus et num quid haberet præterea interrogátus, cum negasset, aureos deínde aliquot suo insutos pallio recordatus, fugiéntibus hos étiam clamans óbtulit latrónibus; qui, viri sancti candórem simul et largitátem admirati, étiam ablatos ultro reddidére. Alienæ famæ ne quis detraheret, descriptis, beáti Augustíni exemplo, in pariete versiculis, se atque alios perpetuo vóluit admónitos. Famélicos de suo étiam obsonio satiábat; nudos autem non emptis modo sed detractis quoque sibi vestibus et calceis operiebat, demisso ipse interim usque ad terram pallio, ne domum núdipes redire viderétur.\par
\switchcolumn
\EN
\noindent What time was left him over from his work, he gave up partly to the profit of his neighbour, more especially in preaching, and partly in prayer, wherein he is said sometimes to have had heavenly visions and messages. The sufferings of Christ took such hold upon him, that he sometimes passed whole nights without sleep in thinking thereon, and that he might more keenly realize them, he made a pilgrimage to Jerusalem. There he was seized with such a passionate longing to be a martyr, that he preached Christ crucified even to the Turks. He went four times to Rome to the thresholds of the Apostles, on foot, and laden with a wallet, partly to do honour to the Apostolic See, for which he had a great reverence, and partly (to use his own expression) that he might clear off the pains of his own purgatory by use of the Pardons for sin which are there daily offered. In one of these journeys he was set upon by highway robbers, who plundered him, and having asked him if he had any more, whereto he answered, Nay, left him and fled. Then he remembered that he had some gold pieces sewn up in his clothes. So he ran after the robbers with shouts, and offered them these also, but they were so amazed at the simplicity and charity of the holy man, that they gave him back even that which they had already taken. To hinder scandal-mongering, he wrote up upon the walls, after the example of holy Augustine, certain texts, to be an unceasing warning to himself and others. He gave his own bread to the hungry, and clothed the naked, not with bought raiment only, but by stripping himself of his own garments and shoes, himself meanwhile letting down his own cloak to trail upon the ground, lest any should see that he returned home barefoot.\par
\switchcolumn*

\LA
\begin{Resp}\Rsign Amávit eum Dóminus, et ornávit eum: stolam glóriæ índuit eum, \Star{} Et ad portas paradísi coronávit eum.\end{Resp}
\begin{Resp}\Vsign Índuit eum Dóminus lorícam fídei, et ornávit eum.\end{Resp}
\begin{Resp}\Rsign Et ad portas paradísi coronávit eum.\end{Resp}
\switchcolumn
\EN
\begin{Resp}\Rsign The Lord loved him and beautified him; He clothed him with a robe of glory, \Star{} And crowned him at the gates of Paradise.\end{Resp}
\begin{Resp}\Vsign The Lord hath put on him the breast-plate of faith, and hath adorned him.\end{Resp}
\begin{Resp}\Rsign And crowned him at the gates of Paradise.\end{Resp}
\switchcolumn*

\LA
\LessonHead{Lectio VI}
\switchcolumn
\EN
\LessonHead{Lesson VI}
\switchcolumn*

\LA
\noindent Brevis illi somnus, atque humi; vestis, quæ nuditátem, cibus, qui mortem dumtaxat, arcéret. Virginalem pudicítiam, velut lílium inter spinas, áspero cilício, flagellis atque jejuniis custodívit. Quin et per annos ante óbitum trigínta circiter et quinque, ab esu carnium perpetuo abstinuit. Tandem diérum juxta ac meritórum plenus, cum vicinæ, quam præsensit, morti se diu diligenterque præparasset, ne qua re ámplius tenerétur, si quid domi supérerat, id omnino paupéribus distríbuit. Tum Ecclésiæ sacramentis rite munítus, dissolvi jam cupiens et esse cum Christo, pridie Nativitátis ejus, in cælum evolávit, miraculis ante et post mortem clarus. Mortuus ad próximam academíæ ecclésiam sanctæ Annæ delátus est, ibique honorifice sepúltus. Auctáque in dies pópuli veneratióne ac frequéntia, inter primarios Poloniæ ac Lithuaniæ patronos religiosíssime cólitur. Novisque coruscans miraculis, a Cleménte décimo tertio Pontifice máximo, décimo septimo Kalendas Augusti, anno millesimo septingentésimo sexagesimo septimo, solemni ritu, Sanctórum fastis adscriptus est.\par
\switchcolumn
\EN
\noindent He slept very little, and that upon the ground; his clothing was enough only to clothe his nakedness, and his food to keep him alive. He kept his virgin purity guarded like a lily among thorns by rough hair-cloth, scourging, and fasting. For about thirty-five years before his death he never tasted flesh- meat. At length, when he was full of days and good works, he felt that death was near, and made himself ready to meet it by a long and careful preparation, and to be the freer, he gave to the poor everything that was left in his house. Strengthened by the Sacraments of the Church, and having a desire to depart, and to be with Christ, he took flight to heaven upon the 24th day of December, (in the year of our Lord 1473.) He was famous for miracles both before and after his death. His body was carried into the University Church of St. Anne, hard by his dwelling, and there honourably buried. The popular reverence and the crowds around his sepulchre grew greater day by day, till he hath come to be held in honour as one of the chiefest holy defenders of Poland and Lithuania. At the glory of more wonders, Pope Clement XIII, upon the 16th day of July, in the year 1767, with solemn pomp, enrolled his name among those of the Saints.\par
\switchcolumn*

\LA
\begin{Resp}\Rsign Iste homo perfécit ómnia quæ locútus est ei Deus, et dixit ad eum: Ingrédere in réquiem meam: \Star{} Quia te vidi justum coram me ex ómnibus géntibus.\end{Resp}
\begin{Resp}\Vsign Iste est, qui contémpsit vitam mundi, et pervénit ad cæléstia regna.\end{Resp}
\begin{Resp}\Rsign Quia te vidi justum coram me ex ómnibus géntibus.\end{Resp}
\begin{Resp}\Vsign Glória Patri, et Fílio, \Star{} et Spirítui Sancto.\end{Resp}
\begin{Resp}\Rsign Quia te vidi justum coram me ex ómnibus géntibus.\end{Resp}
\switchcolumn
\EN
\begin{Resp}\Rsign This is he which did according unto all that God commanded him; and God said unto him: Enter thou into My rest, \Star{} For thee have I seen righteous before Me among all people.\end{Resp}
\begin{Resp}\Vsign This is he which loved not his life in this world, and is come unto an everlasting kingdom.\end{Resp}
\begin{Resp}\Rsign For thee have I seen righteous before Me among all people.\end{Resp}
\begin{Resp}\Vsign Glory be to the Father, and to the Son, \Star{} and to the Holy Ghost.\end{Resp}
\begin{Resp}\Rsign For thee have I seen righteous before Me among all people.\end{Resp}
\switchcolumn*

\end{paracol}

\par\needspace{10\baselineskip}\addvspace{1.2em}{\centering\small\textsc{Die 21 Octobris} · \textsc{21 October}\par}
\Office{S. Hilarionis Abbatis}{St. Hilarion, Abbot}{Simplex}
\addcontentsline{toc}{subsection}{Die 21 Octobris · S. Hilarionis Abbatis \textit{— St. Hilarion, Abbot}}
\begin{paracol}{2}
\LA
\RefLine{Lectiones I–II: de Scriptura occurrente.}
\switchcolumn
\EN
\RefLine{Lessons I–II: from the Scripture of the day.}
\switchcolumn*

\LA
\LessonHead{Lectio III (pro commemoratione)}
\switchcolumn
\EN
\LessonHead{Lesson III (when commemorated)}
\switchcolumn*

\LA
\noindent Hilarion, ortus Tabáthæ in Palæstina ex paréntibus infidelibus, Alexandríam missus studiórum causa, ibi morum et ingenii laude flóruit; ac, Jesu Christi suscepta religióne, in fide et caritate mirabíliter profecit. Frequens enim erat in ecclésia, assiduus in jejúnio et oratióne; omnes voluptátum illécebras et terrenárum rerum cupiditates contemnebat. Cum autem Antonii nomen in Ægypto celeberrimum esset, ejus vidéndi studio in solitúdinem conténdit; apud quem duobus mensibus omnem ejus vitæ ratiónem didicit. Domum reversus, mórtuis paréntibus, facultates suas paupéribus dilargitus est; necdum quintum décimum annum egréssus, rediit in solitúdinem, ubi, exstructa exígua casa, quæ vix ipsum cáperet, humi cubábat. Nec vero saccum, quo semel amíctus est, umquam aut lavit aut mutávit, cum supervacáneum esse diceret, mundítias in cilício quærere. In sanctárum Litterárum lectióne et meditatióne multus erat. Paucas ficus et succum herbárum ad victum adhibebat; nec illis ante solis occásum vescebátur. Continéntia et humilitate fuit incredibili. Quibus aliisque virtútibus varias horribilesque tentatiónes diaboli superávit, et innumerábiles dæmones in multis orbis terræ partibus ex hóminum corpóribus ejécit. Qui, octogesimum annum agens, multis ædificátis monastériis, et clarus miraculis, in morbum incidit; cujus vi cum extremo pene spíritu conflictarétur, dicebat: Egredere, quid times? egredere, ánima mea, quid dubitas? septuagínta prope annis servísti Christo, et mortem times? Quibus in verbis spíritum exhalávit.\par
\switchcolumn
\EN
\noindent Hilarion was born of heathens at Tabatha in Palestine, (about the year of our Lord 291.) He was sent to study at Alexandria, where he bore a fair name for life and wit. There he embraced the religion of Jesus Christ, and made wonderful head-way in faith and love. He went oftentimes to Church, was careful in fasting and prayer, and set no price upon the pleasures and lusts of the world. When the name of Antony became famous in Egypt, Hilarion made a journey into the desert on purpose to see him. There he dwelt with him two months, to the end that he might learn all his way of life, and then returned home. After the death of his father and mother, he gave all that he had to the poor. Before he had completed the fifteenth year of his age, he went into the desert, and built there a little house, scarcely big enough to hold him, and wherein he was used to sleep on the ground. The piece of sackcloth wherewith alone he clad himself he never washed and never changed, saying that hair-cloth was a thing not worth the trouble of cleanliness. He took great interest in reading and meditating on the Holy Scriptures. His food was a few figs and some porridge of vegetables, and this he ate not before set of sun. His self-control and lowliness were beyond belief. By these and other arms he overcame diverse and fearful attacks of the devil, and drave out countless evil spirits from the bodies of men in many parts of the world. He had built many monasteries, and was famous for miracles, when, in the eightieth year of his age, he fell sick. When he was gasping for his last breath, he said Go out what art thou afraid of? Go out, my soul! wherefore shrinkest thou? Thou hast served Christ hard on seventy years and art thou afraid of death? And so with these words he gave up the Ghost.\par
\switchcolumn*

\LA
\TeDeum{Te Deum}
\switchcolumn
\EN
\TeDeum{Te Deum}
\switchcolumn*

\end{paracol}

\par\needspace{10\baselineskip}\addvspace{1.2em}{\centering\small\textsc{Die 24 Octobris} · \textsc{24 October}\par}
\Office{S. Raphaëlis Archangeli}{St. Raphael the Archangel;Duplex majus}{Duplex majus}
\addcontentsline{toc}{subsection}{Die 24 Octobris · S. Raphaëlis Archangeli \textit{— St. Raphael the Archangel;Duplex majus}}
\begin{paracol}{2}
\LA
\LessonHead{Lectio I}
\switchcolumn
\EN
\LessonHead{Lesson I}
\switchcolumn*

\LA
\LessonTitle{De libro Tobíæ}
\switchcolumn
\EN
\LessonTitle{From the Book of Tobias}
\switchcolumn*

\LA
\Cite{Tob 12:1-4}
\switchcolumn
\EN
\Cite{Tob 12:1-4}
\switchcolumn*

\LA
\noindent Vocávit ad se Tobías fílium suum, dixítque ei: Quid póssumus dare viro isti sancto, qui venit tecum? Respóndens Tobías, dixit patri suo: Pater, quam mercédem dábimus ei? aut quid dignum póterit esse benefíciis ejus? Me duxit et redúxit sanum, pecúniam a Gabélo ipse recépit, uxórem ipse me habére fecit, et dæmónium ab ea ipse compéscuit: gáudium paréntibus ejus fecit, meípsum a devoratióne piscis erípuit, te quoque vidére fecit lumen cæli, et bonis ómnibus per eum repléti sumus. Quid illi ad hæc potérimus dignum dare? Sed peto te, pater mi, ut roges eum, si forte dignábitur medietátem de ómnibus quæ alláta sunt, sibi assúmere.\par
\switchcolumn
\EN
\noindent Then Tobias called to him his son, and said to him: What can we give to this holy man, that is come with thee? Tobias answering, said to his father: Father, what wages shall we give him? or what can be worthy of his benefits? He conducted me and brought me safe again, he received the money of Gabelus, he caused me to have my wife, and he chased from her the evil spirit, he gave joy to her parents, myself he delivered from being devoured by the fish, thee also he hath made to see the light of heaven, and we are filled with all good things through him. What can we give him sufficient for these things? But I beseech thee, my father, to desire him, that he would vouchsafe to accept one half of all things that have been brought.\par
\switchcolumn*

\LA
\begin{Resp}\Rsign In illo témpore exaudítæ sunt preces ambórum in conspéctu glóriæ summi Dei: \Star{} Et missus est Angelus Dómini sanctus Ráphaël, ut curáret eos ambos, quorum uno témpore sunt oratiónes in conspéctu Dómini recitátæ.\end{Resp}
\begin{Resp}\Vsign Tobías et Sara in tribulatióne pósiti cum lácrimis oráre cœpérunt.\end{Resp}
\begin{Resp}\Rsign Et missus est Angelus Dómini sanctus Ráphaël, ut curáret eos ambos, quorum uno témpore sunt oratiónes in conspéctu Dómini recitátæ.\end{Resp}
\switchcolumn
\EN
\begin{Resp}\Rsign At that time the prayers of them twain were heard in the presence of the glory of the Most High God; \Star{} And holy Raphael, the Angel of the Lord, was sent to heal the twain of them, whose prayers had been uttered at one time in the presence of the Lord.\end{Resp}
\begin{Resp}\Vsign Tobias and Sara were in tribulation, and began to pray with tears.\end{Resp}
\begin{Resp}\Rsign And holy Raphael, the Angel of the Lord, was sent to heal the twain of them, whose prayers had been uttered at one time in the presence of the Lord.\end{Resp}
\switchcolumn*

\LA
\LessonHead{Lectio II}
\switchcolumn
\EN
\LessonHead{Lesson II}
\switchcolumn*

\LA
\Cite{Tob 12:5-13}
\switchcolumn
\EN
\Cite{Tob 12:5-13}
\switchcolumn*

\LA
\noindent Et vocántes eum, pater scílicet et fílius, tulérunt eum in partem; et rogáre cœpérunt ut dignarétur dimídiam partem ómnium quæ attúlerant, accéptam habére. Tunc dixit eis occúlte: Benedícite Deum cæli, et coram ómnibus vivéntibus confitémini ei, quia fecit vobíscum misericórdiam suam. Etenim sacraméntum regis abscóndere bonum est, ópera autem Dei reveláre et confitéri honoríficum est. Bona est orátio cum jejúnio, et eleemósyna magis quam thesáuros auri recóndere; quóniam eleemósyna a morte líberat, et ipsa est quæ purgat peccáta, et facit inveníre misericórdiam et vitam ætérnam. Qui autem fáciunt peccátum et iniquitátem, hostes sunt ánimæ suæ. Manifésto ergo vobis veritátem et non abscóndam a vobis occúltum sermónem. Quando orábas cum lácrimis, et sepeliébas mórtuos, et derelinquébas prándium tuum, et mórtuos abscondébas per diem in domo tua, et nocte sepeliébas eos, ego óbtuli oratiónem tuam Dómino. Et quia accéptus eras Deo, necésse fuit ut tentátio probáret te.\par
\switchcolumn
\EN
\noindent So the father and the son, calling him, took him aside: and began to desire him that he would vouchsafe to accept of half of all things that they had brought. Then he said to them secretly: Bless ye the God of heaven, give glory to him in the sight of all that live, because he hath shewn his mercy to you. For it is good to hide the secret of a king: but honourable to reveal and confess the works of God. Prayer is good with fasting and alms more than to lay up treasures of gold: For alms delivereth from death, and the same is that which purgeth away sins, and maketh to find mercy and life everlasting. But they that commit sin and iniquity, are enemies to their own soul. I discover then the truth unto you, and I will not hide the secret from you. When thou didst pray with tears, and didst bury the dead, and didst leave thy dinner, and hide the dead by day in thy house, and bury them by night, I offered thy prayer to the Lord. And because thou wast acceptable to God, it was necessary that temptation should prove thee.\par
\switchcolumn*

\LA
\begin{Resp}\Rsign Egréssus Tobías invénit júvenem spléndidum stantem præcínctum, et quasi parátum ad ambulándum, et salutávit eum, et dixit: \Star{} Unde te habémus, bone júvenis?\end{Resp}
\begin{Resp}\Vsign Et, ignórans quod Dómini Angelus esset, salutávit eum, et dixit.\end{Resp}
\begin{Resp}\Rsign Unde te habémus, bone júvenis?\end{Resp}
\switchcolumn
\EN
\begin{Resp}\Rsign When Tobias went forth he found a very well-favoured young man standing with his loins girded, and as it were ready to set forth, and he saluted him and said: \Star{} Whence art thou, O good young man?\end{Resp}
\begin{Resp}\Vsign And he knew not that it was an Angel of the Lord, and saluted him, and said:\end{Resp}
\begin{Resp}\Rsign Whence art thou, O good young man?\end{Resp}
\switchcolumn*

\LA
\LessonHead{Lectio III}
\switchcolumn
\EN
\LessonHead{Lesson III}
\switchcolumn*

\LA
\Cite{Tob 12:14-22}
\switchcolumn
\EN
\Cite{Tob 12:14-22}
\switchcolumn*

\LA
\noindent Et nunc misit me Dóminus ut curárem te, et Saram uxórem fílii tui a dæmónio liberárem; ego enim sum Ráphaël Angelus, unus ex septem qui astámus ante Dóminum. Cumque hæc audíssent, turbáti sunt, et treméntes cecidérunt super terram in fáciem suam. Dixítque eis Angelus: Pax vobis, nolíte timére. Etenim cum essem vobíscum, per voluntátem Dei eram: ipsum benedícite, et cantáte illi. Vidébar quidem vobíscum manducáre, et bíbere: sed ego cibo invisíbili, et potu, qui ab homínibus vidéri non potest, utor. Tempus est ergo ut revértar ad eum, qui me misit: vos autem benedícite Deum, et narráte ómnia mirabília ejus. Et cum hæc dixísset, ab aspéctu eórum ablátus est, et ultra eum vidére non potuérunt. Tunc prostráti per horas tres in fáciem, benedixérunt Deum: et exsurgéntes narravérunt ómnia mirabília ejus.\par
\switchcolumn
\EN
\noindent And now the Lord hath sent me to heal thee, and to deliver Sara thy son's wife from the devil. For I am the angel Raphael, one of the seven, who stand before the Lord. And when they had heard these things, they were troubled, and being seized with fear they fell upon the ground on their face. And the angel said to them: Peace be to you, fear not. For when I was with you, I was there by the will of God: bless ye him, and sing praises to him. I seemed indeed to eat and to drink with you: but I use an invisible meat and drink, which cannot be seen by men. It is time therefore that I return to him that sent me: but bless ye God, and publish all his wonderful works. And when he had said these things, he was taken from their sight, and they could see him no more. Then they lying prostrate for three hours upon their face, blessed God: and rising up, they told all his wonderful works.\par
\switchcolumn*

\LA
\begin{Resp}\Rsign Ingréssus Angelus ad Tobíam, salutávit eum, et dixit: Gaudium sit tibi semper: \Star{} Forti animo esto, in próximo enim est, ut a Deo cureris.\end{Resp}
\begin{Resp}\Vsign Et respondens Tobías, ait: Quale gáudium mihi erit, qui in ténebris sedeo, et lumen cæli non video?\end{Resp}
\begin{Resp}\Rsign Forti animo esto, in próximo enim est, ut a Deo cureris.\end{Resp}
\begin{Resp}\Vsign Glória Patri, et Fílio, \Star{} et Spirítui Sancto.\end{Resp}
\begin{Resp}\Rsign Forti animo esto, in próximo enim est, ut a Deo cureris.\end{Resp}
\switchcolumn
\EN
\begin{Resp}\Rsign The Angel went in unto Tobias and saluted him and said: Joy be ever with thee; \Star{} Be of good courage, for it is at hand that God should heal thee.\end{Resp}
\begin{Resp}\Vsign And Tobias answered and said: What joy shall I have while I sit in darkness, and see not the light of heaven?\end{Resp}
\begin{Resp}\Rsign Be of good courage, for it is at hand that God should heal thee.\end{Resp}
\begin{Resp}\Vsign Glory be to the Father, and to the Son, \Star{} and to the Holy Ghost.\end{Resp}
\begin{Resp}\Rsign Be of good courage, for it is at hand that God should heal thee.\end{Resp}
\switchcolumn*

\LA
\LessonHead{Lectio IV}
\switchcolumn
\EN
\LessonHead{Lesson IV}
\switchcolumn*

\LA
\LessonTitle{Sermo Sancti Bonaventuræ Epíscopi}
\switchcolumn
\EN
\LessonTitle{From the Sermons of St. Bonaventure, the Bishop}
\switchcolumn*

\LA
\Source{De Sanctis Angelis Sermo 5, sub fine}
\switchcolumn
\EN
\Source{De Sanctis Angelis Sermo 5, sub fine}
\switchcolumn*

\LA
\noindent Ráphaël interpretátur medicina Dei. Et debémus notare quod eductio a malo est per tria beneficia, a Raphaéle nobis colláta medicante nos. Educit ergo nos Ráphaël médicus ab infirmitáte animi, inducéndo nos ad amaritúdinem contritiónis; unde in Tobía dicit Ráphaël: Ubi introíeris domum tuam, lini super óculos ejus ex felle. Sic fecit, et vidit. Quare Ráphaël non pótuit ipse fácere? Quia Angelus non dat compunctiónem, sed osténdit viam. Per fel intellígitur amaritúdo contritiónis, quæ sanat óculos interioris mentis; Psalmus: Qui sanat contritos corde. Hoc est optimum collyrium. In Júdicum secundo dícitur quod Angelus ascéndit ad locum fléntium, et dixit pópulo: Eduxi vos de terra Ægypti, feci vobis tot et tanta bona; et flevit omnis pópulus, ita ut locus ille appellarétur locus fléntium. Caríssimi, Angeli tota die narrant nobis benefícia Dei, et reducunt ea nobis ad memóriam: Quis est qui te creávit, qui te redémit? Quid fecísti, quem offendísti? Hoc si consideráveris, nullum habes remédium nisi flere.\par
\switchcolumn
\EN
\noindent Raphael means God's healing, and we should observe that the way out of evil lies through three blessing bestowed on us through Raphael's healing. For Raphael the Healer leads us forth from weakness of soul and brings us into the bitterness of contrition; whence in the book of Tobit Raphael says: When thou enterest into thy house, anoint his eyes with gall. He did so, and he saw. Why could not Raphael do this himself? Because an Angel does not bestow compunction, but he shows the way to it. By gall, the bitterness of contrition is to be understood, for it heals the inner eyes of the mind. The Psalmist says: He healeth those that are broken in heart. This is the best eye salve. In the second chapter of the Book of Judges it is written that an angel came up to the place of the weepers, and said to the people: I made you to go up out of Egypt; I have done all of these good things for you: and all the people wept, so that the place was called the Place of the Weepers. Dearly beloved, all day long the Angels are telling us of God's blessings, and recalling them to our minds: Who hath created thee? Who hath redeemed thee? What hast thou done? Whom hast thou offended? When thou considerest this, thine only remedy is to weep.\par
\switchcolumn*

\LA
\begin{Resp}\Rsign Interrogávit Tobías Angelum: De qua domo, aut de qua tribu es tu? Qui respondens, ait: \Star{} Ego sum Azarías, Ananíæ magni fílius.\end{Resp}
\begin{Resp}\Vsign Genus quæris mercenarii, an ipsum mercenárium, qui cum fílio tuo eat? Sed ne forte sollícitus sis.\end{Resp}
\begin{Resp}\Rsign Ego sum Azarías, Ananíæ magni fílius.\end{Resp}
\switchcolumn
\EN
\begin{Resp}\Rsign Tobias asked the Angel: Of what house, or of what tribe, art thou? And he answered and said: \Star{} I am Azarias, the son of Ananias the Great.\end{Resp}
\begin{Resp}\Vsign Dost thou seek for a tribe or family, or an hired man to go with thy son? But be not careful overmuch.\end{Resp}
\begin{Resp}\Rsign I am Azarias, the son of Ananias the Great.\end{Resp}
\switchcolumn*

\LA
\LessonHead{Lectio V}
\switchcolumn
\EN
\LessonHead{Lesson V}
\switchcolumn*

\LA
\noindent Secundo Ráphaël edúcit de servitute diaboli, persuadéndo nobis memóriam passiónis Christi; in cujus figuram dictum est Tobíæ sexto: Cordis ejus particulam si super carbónes ponas, fumus ejus éxtricat omne genus dæmoniórum. Dícitur Tobíæ octavo quod pósuit Tobías particulam cordis super carbónes, et Ráphaël religávit dæmónium in desérto superioris Ægypti. Quid est hoc? Non poterat Ráphaël religare dæmónium nisi ponerétur cor super carbónes? Numquid cor piscis dabat Angelo tantam virtútem? Nequaquam! Nihil posset, nisi ibi mystérium esset. In hoc enim nobis datur intelligi quod nihil est quod ita nos líberet hódie a servitute diaboli sicut passio Christi, quæ processit ex radíce cordis sive caritátis. Cor enim fons est caloris cunctæ vitæ. Si ergo Cor Christi, hoc est passiónem quam sustinuit, procedéntem ex radíce caritátis et fonte caloris, ponas super carbónes, hoc est super inflammatam memóriam; statim dæmon religábitur, ut tibi nocére non possit.\par
\switchcolumn
\EN
\noindent Secondly, Raphael brings us forth from our slavery to the devil, by recalling\par
Christ's Passion to our mind. There is a figure of the Passion in the sixth chapter of Tobit, where it is said that if a piece of the heart is put on the embers, its smoke will drive out all kinds of devils. It is said in the eighth chapter of Tobit that Tobias put a piece of the heart on the embers, and Raphael bound the devil in the utmost parts of Egypt. What does this mean? Was Raphael unable to bind the devil unless the heart was put on the embers? Surely a fish's heart could not give such strength to an Angel? By no means! It would be impossible unless there was a mystery in this. We are to understand that nothing so frees us from the devil's slavery as the Passion of Christ, which issued from the root of his heart, or of his love. For the heart is the fountain, or hot centre, of all life. Therefore if thou puttest the heart of Christ, that is the Passion, that he underwent, issuing from the root of charity and the fountain of heat, on the embers, that is, on the flame of memory: then immediately the devil will be bound, and he will be unable to hurt thee.\par
\switchcolumn*

\LA
\begin{Resp}\Rsign Exívit Tobías ut lavaret pedes suos, et ecce piscis immanis exívit ad devorándum eum: qui expavéscens, clamávit voce magna, dicens: Dómine, invadit me. Et dixit ei Angelus: Apprehénde branchiam ejus, et trahe eum ad te. \Star{} Exéntera hunc piscem, et cor ejus, et fel, et jecur repóne tibi: sunt enim necessaria ad medicaménta utíliter.\end{Resp}
\begin{Resp}\Vsign Attraxit autem Tobías piscem in siccum, et palpitare cœpit ante pedes ejus; et ait Angelus ei.\end{Resp}
\begin{Resp}\Rsign Exéntera hunc piscem, et cor ejus, et fel, et jecur repóne tibi: sunt enim necessaria ad medicaménta utíliter.\end{Resp}
\switchcolumn
\EN
\begin{Resp}\Rsign Tobias went out to wash his feet, and behold a great fish went out, and would have devoured him; and he was sore afraid, and cried out with a loud voice, and said: O sir, he assaileth me. And the Angel said unto him: Take him by his gill, and draw him unto thee. \Star{} Open this fish, and take his heart, and his liver, and his gall, and put them up safely by thee, for they are useful and needful for drugs.\end{Resp}
\begin{Resp}\Vsign And Tobias drew the fish unto the dry ground, and it began to gasp at his feet, and the Angel said unto him.\end{Resp}
\begin{Resp}\Rsign Open this fish, and take his heart, and his liver, and his gall, and put them safely by thee, for they are useful and needful for drugs.\end{Resp}
\switchcolumn*

\LA
\LessonHead{Lectio VI}
\switchcolumn
\EN
\LessonHead{Lesson VI}
\switchcolumn*

\LA
\noindent Tertio liberat nos a contrarietáte Dei, quam incurrimus per offensam Dei, et hoc inducéndo nos ad instantiam oratiónis; et hoc est quod dixit angelus Ráphaël Tobíæ duodecimo: Quando orabas cum lácrimis, ego óbtuli oratiónem tuam Dómino. Ipsi enim Angeli reconciliant nos Deo, quantum possunt. Accusatores nostri coram Deo sunt dæmones. Angeli autem excusant nos, quando ófferunt oratiónes nostras, ad quas devote faciendas nos inducunt; Apocalypsis octavo: Ascéndit fumus arómatum in conspéctu Dómini de manu Angeli. Arómata ista suáviter redoléntia sunt oratiónes Sanctórum. Vis placare Deum, quem offendísti? Ora devote. Offérunt Deo oratiónem tuam, ut te Deo reconcilient. Dícitur in Luca quod Christus, factus in agonía, prolixius orabat, et appáruit Angelus Dómini confortans eum. Et hoc totum factum est propter nos, quia non indiguit confortatióne sua, sed ut ostenderétur quod libenter assistunt devote orántibus, et libenter juvant eos et ipsos confortant, et oratiónes eórum Deo ófferunt. - Festum sancti Raphaëlis Archangeli Benedíctus Papa décimus quintus ad universam Ecclésiam extendit.\par
\switchcolumn
\EN
\noindent Thirdly Raphael delivers us from God's displeasure, which we incur by transgressing against God; he does this by leading us to instant prayer: and this is what Raphael said to Tobit in the twelfth chapter: When thou didst pray with tears, I did bring your prayers before the Lord. For the Angels reconcile us to God, as far as they are able. The devils are accusers before God, but the angels excuse us, when they offer our prayers, and when they urge us to pray devoutly. In the eighth chapter of the Revelation of St. John it is said: The smoke of the incense ascended up before God out of the Angels hand. This sweet smelling incense is the prayer of the saints. Does thou wish to appease God, whom thou hast offended ? Then pray devoutly. They offer thy prayer, so that they may reconcile thee to God. It is said in the Gospel of St. Luke, that Christ being in an agony, prayed more earnestly, and there appeared an Angel of the Lord strengthening him. And all of this was done for our sake: for he had no need of strengthening, but it was to show that the Angels willingly assist those who pray devoutly, and help and strengthen them, and bring their prayers before God.\par
\switchcolumn*

\LA
\begin{Resp}\Rsign Ubi introíeris domum tuam, dixit Angelus Ráphaël ad Tobíam, statim adora Dóminum Deum tuum, et, grátias agens ei, accede ad patrem tuum, et osculare eum: \Star{} Statimque lini super óculos ejus ex felle isto piscis, quem portas tecum; sicas enim quóniam mox aperiéntur óculi ejus, et vidébit pater tuus lumen cæli, et in aspectu tuo gaudébit.\end{Resp}
\begin{Resp}\Vsign Tolle tecum ex felle isto piscis; erit enim necessárium.\end{Resp}
\begin{Resp}\Rsign Statimque lini super óculos ejus ex felle isto piscis, quem portas tecum; sicas enim quóniam mox aperiéntur óculi ejus, et vidébit pater tuus lumen cæli, et in aspectu tuo gaudébit.\end{Resp}
\begin{Resp}\Vsign Glória Patri, et Fílio, \Star{} et Spirítui Sancto.\end{Resp}
\begin{Resp}\Rsign Statimque lini super óculos ejus ex felle isto piscis, quem portas tecum; sicas enim quóniam mox aperiéntur óculi ejus, et vidébit pater tuus lumen cæli, et in aspectu tuo gaudébit.\end{Resp}
\switchcolumn
\EN
\begin{Resp}\Rsign The Angel Raphael said unto Tobias: When thou enterest into thine house, straightway worship the Lord thy God; and when thou hast given thanks unto Him, draw near unto thy father and kiss him. \Star{} And forthwith smear upon his eyes some of the gall of the fish which thou carriest with thee; for know that his eyes will straightway be opened, and thy father shall see the light of heaven, and shall behold thee, and be glad.\end{Resp}
\begin{Resp}\Vsign Take with thee some of the gall of the fish, for it will be needful.\end{Resp}
\begin{Resp}\Rsign And forthwith smear upon his eyes some of the gall of the fish which thou carriest with thee; for know that his eyes will straightway be opened, and thy father shall see the light of heaven, and shall behold thee, and be glad.\end{Resp}
\begin{Resp}\Vsign Glory be to the Father, and to the Son, \Star{} and to the Holy Ghost.\end{Resp}
\begin{Resp}\Rsign And forthwith smear upon his eyes some of the gall of the fish which thou carriest with thee; for know that his eyes will straightway be opened, and thy father shall see the light of heaven, and shall behold thee, and be glad.\end{Resp}
\switchcolumn*

\LA
\LessonHead{Lectio VII}
\switchcolumn
\EN
\LessonHead{Lesson VII}
\switchcolumn*

\LA
\LessonTitle{Léctio sancti Evangélii secúndum Joánnem}
\switchcolumn
\EN
\LessonTitle{From the Holy Gospel according to John}
\switchcolumn*

\LA
\Cite{Joannes 5:1-4}
\switchcolumn
\EN
\Cite{John 5:1-4}
\switchcolumn*

\LA
\noindent In illo témpore: Erat dies festus Judæórum, et ascéndit Jesus Jerosolymam. Et réliqua.\par
\switchcolumn
\EN
\noindent At that time: There was a feast of the Jews, and Jesus went up to Jerusalem. And so on.\par
\switchcolumn*

\LA
\LessonTitle{Homilía sancti Joánnis Chrysóstomi}
\switchcolumn
\EN
\LessonTitle{Homily by St. John Chrystostom}
\switchcolumn*

\LA
\Source{Homilia 36, alias 35, in Joánnem, num. 1}
\switchcolumn
\EN
\Source{Hom. 35 or 36 on John}
\switchcolumn*

\LA
\noindent Quis hic curatiónis modus? quale mystérium subindicátur? Neque enim sine causa hæc scripta sunt; sed futura nobis quasi in figura et imagine describit, ne, si res stupenda accideret inexspéctata, auditórum multórum fidem aliquátenus labefactáret. Quænam ígitur hæc descriptio? Futurum baptisma dandum erat, plenum virtúte et grátia maxima, baptisma, quod peccáta ómnia ablúeret, quod ex mórtuis vivos redderet. Hæc ergo ut in imagine depingúntur in piscina et in aliis multis. Et primo quidem aquam dedit, quæ córporum máculas ablúeret, sordesque non veras, sed tales existimátas, ex funere nempe, ex lepra, et similes; multaque vidére est eádem de causa in veteri lege per aquam mundata.\par
\switchcolumn
\EN
\noindent What manner of mystery is this? What does it reveal? for it is not written without a purpose: but the future is foretold in figure and similitude lest the unexpected occurrence of a wondrous event should in any way disturb the faith of the hearers. What, then is here described? In the future a baptism was to be given, full of power and of the greatest grace, a baptism that would wash away all sin, that would restore life to dead men. These facts then are depicted figuratively, in the pool and in all the other circumstances. And first in this figure the water is set forth, which washes away stains of the body, and those things that are not actually dirty, though they were thought to be so, such as coming in contact with a corpse, or a leper, and such like: it is to be seen that many things under the old law are cleansed by water, in accordance with this idea.\par
\switchcolumn*

\LA
\begin{Resp}\Rsign Benedícite Deum cæli, dixit Angelus Ráphaël et coram ómnibus vivéntibus confitémini ei: \Star{} Quia fecit vobíscum misericórdiam suam.\end{Resp}
\begin{Resp}\Vsign Ipsum benedícite, et cantáte illi, et narráte ómnia mirabília ejus.\end{Resp}
\begin{Resp}\Rsign Quia fecit vobíscum misericórdiam suam.\end{Resp}
\switchcolumn
\EN
\begin{Resp}\Rsign The Angel Raphael said: Bless ye the God of heaven, and confess Him before all living; \Star{} For He hath had mercy upon you.\end{Resp}
\begin{Resp}\Vsign Bless Him, and praise Him, and tell of all His marvellous works.\end{Resp}
\begin{Resp}\Rsign For He hath had mercy upon you.\end{Resp}
\switchcolumn*

\LA
\LessonHead{Lectio VIII}
\switchcolumn
\EN
\LessonHead{Lesson VIII}
\switchcolumn*

\LA
\noindent Sed ad propositum jam redeamus. Primo ítaque, ut diximus, córporum maculas, deínde varias infirmitátes per aquam solvi curat. Ut enim nos Deus ad baptismi grátiam propius reduceret, non jam maculas solum, sed et morbos sanat. Imagines enim quæ propius ad veritátem accedunt, et in baptismate, et in passióne, et in aliis magis conspicuæ sunt quam vetustiores. Quemádmodum enim qui prope regem sunt satéllites, remotióribus sunt honoratiores; ita et in figuris factum est. Et Angelus descéndens turbábat aquam, et sanándi vim indebat ipsi, ut discerent Judæi, Angelórum Dóminum multo magis posse ánimæ morbos omnes curare. Sed, quemádmodum hic aquárum natúra non simpliciter curábat (alioquin enim semper id fáceret), sed Angeli operatióne id fiebat; sic in nobis non aqua simpliciter operátur, sed, postquam Spíritus grátiam accéperit, tunc ómnia solvit peccáta.\par
\switchcolumn
\EN
\noindent But let us now return to the subject. First defilements of the body are washed out, and then God heals various infirmities by means of water. For because God would bring us nearer to the grace of baptism, he not only cleanses defilements but heals diseases. For those figures which come nearest to the reality, in baptism, the Passion, and others are seen more clearly than the older ones. For it is the same with those who form a king's bodyguard: they are more splendidly apparelled than those at the other end of his equipage. And the angel came down and troubled the water, and imbued it with healing power; so that the Jews might learn that the Lord of Angels had far more power to heal all the diseases of the soul. But just as here it was not simply the natural property of the water that healed, (otherwise it would have always have done so) but it happened through the work of the Angel: so it is not simply water that works on us, but after the water has received the grace of the Spirit, then it looses us from all sin.\par
\switchcolumn*

\LA
\begin{Resp}\Rsign Tempus est ut revértar ad eum, qui me misit, dixit Angelus Ráphaël; \Star{} Vos autem benedícite Dóminum, et narráte ómnia mirabília ejus.\end{Resp}
\begin{Resp}\Vsign Confitémini ei coram ómnibus vivéntibus, quia fecit vobíscum misericórdiam suam.\end{Resp}
\begin{Resp}\Rsign Vos autem benedícite Dóminum, et narráte ómnia mirabília ejus.\end{Resp}
\begin{Resp}\Vsign Glória Patri, et Fílio, \Star{} et Spirítui Sancto.\end{Resp}
\begin{Resp}\Rsign Vos autem benedícite Dóminum, et narráte ómnia mirabília ejus.\end{Resp}
\switchcolumn
\EN
\begin{Resp}\Rsign It is time for me to return unto Him That sent me, saith the Angel Raphael; \Star{} But bless ye the Lord, and tell of all His marvellous works.\end{Resp}
\begin{Resp}\Vsign Confess Him before all living, for He hath had mercy upon you.\end{Resp}
\begin{Resp}\Rsign Bless ye the Lord, and tell of all His marvellous works.\end{Resp}
\begin{Resp}\Vsign Glory be to the Father, and to the Son, \Star{} and to the Holy Ghost.\end{Resp}
\begin{Resp}\Rsign Bless ye the Lord, and tell of all His marvellous works.\end{Resp}
\switchcolumn*

\LA
\LessonHead{Lectio IX}
\switchcolumn
\EN
\LessonHead{Lesson IX}
\switchcolumn*

\LA
\noindent Circa hanc piscinam jacebat multitúdo magna infirmórum, cæcórum, claudórum, aridórum, aquæ motum exspectántium. Sed tunc infirmitas impediménto erat quóminus is qui vellet, sanarétur; nunc autem unusquísque potestátem accedéndi habet. Non enim Angelus est qui aquam movet, sed Angelórum Dóminus ómnia éfficit. Nec dicere póssumus: Dum ego accedo, alius ante me descéndit. Sed, si totus orbis venerit, grátia non consúmitur, neque vis vel operátio déficit, sed semper éadem manet. Ac, quemádmodum solares radii quotídie illúminant, nec absumúntur, neque, quod multis subministréntur, lucis quídpiam amittunt; sic, immo multo minus, Spíritus operátio minúitur a multitúdine accipiéntium. Hoc autem factum est ut qui díscerent in aqua curándos esse corporis morbos, et hac in re diu exercitáti essent, facílius crederent étiam morbos animi posse curari.\par
\switchcolumn
\EN
\noindent Around this pool lay a great multitude of impotent folk, of blind, halt, withered, waiting for the moving of the water. But then those who wished for healing were prevented by infirmity from receiving it, while now everyone is able to come forward. For the Angel does not trouble the water, but the Lord of Angels brings all things to come to pass. We may no longer say, While I am coming, another steppeth down before me. But even should the whole world come, the grace would not be used up, neither would its power nor effectual working come to an end. For as the sun's rays shine forth each day without burning out, and as they spread abroad without losing any of their light: far less is the operation of the Spirit diminished by the multitudes who receive it. Now this miracle took place so that as men learned that could heal bodily diseases, and as they become accustomed to this fact over a period of time: so they might they the more readily believe that it could also heal diseases of the soul.\par
\switchcolumn*

\LA
\TeDeum{Te Deum}
\switchcolumn
\EN
\TeDeum{Te Deum}
\switchcolumn*

\end{paracol}

\par\needspace{10\baselineskip}\addvspace{1.2em}{\centering\small\textsc{Die 25 Octobris} · \textsc{25 October}\par}
\Office{Ss. Chrysanthi et Dariæ Martyrum}{Sts. Chrysanthus and Daria, Martyrs}{Simplex}
\addcontentsline{toc}{subsection}{Die 25 Octobris · Ss. Chrysanthi et Dariæ Martyrum \textit{— Sts. Chrysanthus and Daria, Martyrs}}
\begin{paracol}{2}
\LA
\RefLine{Lectiones I–II: de Scriptura occurrente.}
\switchcolumn
\EN
\RefLine{Lessons I–II: from the Scripture of the day.}
\switchcolumn*

\LA
\LessonHead{Lectio III (pro commemoratione)}
\switchcolumn
\EN
\LessonHead{Lesson III (when commemorated)}
\switchcolumn*

\LA
\noindent Chrysanthus et Daría cónjuges, nobili genere nati, fide étiam clariores, quam Daría, maríti opera, cum baptismo susceperat; Romæ innumerábilem hóminum multitúdinem, hæc mulíerum, ille virórum, ad Christum convertérunt. Quare Celerinus præfectus comprehensos trádidit Cláudio tribuno, qui jussit a milítibus Chrysanthum vinctum cruciátibus torqueri; sed víncula ómnia resoluta sunt, mox cómpedes, in quos conjéctus fuerat, confracti. Deínde, bovis corio inclusum, in ardentíssimo sole constituunt. Tum, pédibus ac mánibus catena constrictis, in obscurum carcerem detrudunt; ubi, solutis catenis, claríssima lux locum illustrávit. Daría veri in lupánar compulsa, leónis tutela, dum in oratióne defixa est, a contumelia divinitus defensa est. Denique in arenariam, quæ est via Salaria, uterque ductus, effossa terra, lapidibus óbruti, parem martyrii corónam adepti sunt.\par
\switchcolumn
\EN
\noindent Chrysanthus and Daria were an husband and wife, of noble birth, but glorious rather for their faith, which the wife learnt from the husband. They brought to Christ a great number of persons at Rome, she women, and he men. Therefore the Praefect Celerinus caused them to be taken, and gave them over to Claudius the Tribune, who bade Chrysanthus to be tormented by the soldiers, all bound as he was, but all his bonds brake, and so likewise the shackles wherein his feet were afterwards fastened. Then was Chrysanthus sewn up in an ox-hide and set in the full heat of the sun, and thereafter chained hand and foot and cast into a dark prison, but the chains dropped off from him, and the place was filled with light. Meanwhile Daria was haled to a brothel, but God kept her from insult, a lion guarding her, and herself always rapt in prayer. Lastly they were both of them led to a sand-pit upon the Salarian Way, where they were thrown alive into a hole, and buried in stones, and so were not divided in winning the victory of Martyrdom.\par
\switchcolumn*

\LA
\TeDeum{Te Deum}
\switchcolumn
\EN
\TeDeum{Te Deum}
\switchcolumn*

\end{paracol}

\par\needspace{10\baselineskip}\addvspace{1.2em}{\centering\small\textsc{Die 26 Octobris} · \textsc{26 October}\par}
\Office{S. Evaristi Papæ et Martyris}{St. Evaristus, Pope and Martyr}{Simplex}
\addcontentsline{toc}{subsection}{Die 26 Octobris · S. Evaristi Papæ et Martyris \textit{— St. Evaristus, Pope and Martyr}}
\begin{paracol}{2}
\LA
\RefLine{Lectiones I–II: de Scriptura occurrente.}
\switchcolumn
\EN
\RefLine{Lessons I–II: from the Scripture of the day.}
\switchcolumn*

\LA
\LessonHead{Lectio III (pro commemoratione)}
\switchcolumn
\EN
\LessonHead{Lesson III (when commemorated)}
\switchcolumn*

\LA
\noindent Evarístus, Græcus, ex Judæo patre, Trajano imperatóre, pontificátum gessit. Qui ecclesiárum titulos urbis Romæ presbyteris divísit, et ordinávit ut septem diáconi episcopum custodírent dum evangelicæ prædicatiónis officio fungerétur. Idem constítuit, ex traditióne apostolica, ut matrimónium publice celebrétur et sacerdotis benedíctio adhibeátur. Præfuit Ecclésiæ annos novem, menses tres, presbyteris decem et septem, diáconis duobus, epíscopis quindecim, quater mense Decembri, ordinatis. Martyrio coronátus, prope sepúlcrum Principis Apostolórum in Vaticano sepúltus est septimo Kalendas Novembris.\par
\switchcolumn
\EN
\noindent Evaristus was by birth a Greek Jew, and held the Popedom in the reign of the Emperor Trajan. He it was who divided among the Priests the titles of the Churches in the city of Rome, and commanded that seven Deacons should attend the Bishop when he was executing his office of Gospel preaching. He it was who commanded, in accordance with the tradition of the Apostles, that marriages should be celebrated openly, and that a Priest should be asked to invoke a blessing thereon. He ruled the Church for nine years and three months. He held four Ordinations in the month of December, wherein he ordained seventeen Priests, two Deacons, and fifteen Bishops. Having finished his testimony, he was buried upon the Vatican, hard by the grave of the Prince of the Apostles, upon the 26th day of October, (in the year of our Lord 112.)\par
\switchcolumn*

\LA
\TeDeum{Te Deum}
\switchcolumn
\EN
\TeDeum{Te Deum}
\switchcolumn*

\end{paracol}

\par\needspace{10\baselineskip}\addvspace{1.2em}{\centering\small\textsc{Die 27 Octobris} · \textsc{27 October}\par}
\Office{In Vigilia Ss. Simonis et Judæ Ap.}{Vigil of Sts. Simon and Jude, Apostles}{Vigilia}
\addcontentsline{toc}{subsection}{Die 27 Octobris · In Vigilia Ss. Simonis et Judæ Ap. \textit{— Vigil of Sts. Simon and Jude, Apostles}}
\begin{paracol}{2}
\LA
\LessonHead{Lectio I}
\switchcolumn
\EN
\LessonHead{Lesson I}
\switchcolumn*

\LA
\LessonTitle{Léctio sancti Evangélii secúndum Joánnem}
\switchcolumn
\EN
\LessonTitle{From the holy Gospel according to John}
\switchcolumn*

\LA
\Cite{Joannes 15:1-7}
\switchcolumn
\EN
\Cite{John 15:1-7}
\switchcolumn*

\LA
\noindent In illo témpore: Dixit Jesus discípulis suis: Ego sum vitis vera: et Pater meus agrícola est. Et réliqua.\par
\switchcolumn
\EN
\noindent In that time Jesus said to his disciples: I an the true vine; and my Father is the husbandman. And so on.\par
\switchcolumn*

\LA
\LessonTitle{Homilía sancti Augustíni Epíscopi}
\switchcolumn
\EN
\LessonTitle{Homily by St. Augustine, Bishop (of Hippo.)}
\switchcolumn*

\LA
\Source{Tract. 80 in Joannem}
\switchcolumn
\EN
\Source{Tract 80 on John}
\switchcolumn*

\LA
\noindent Iste locus evangélicus, fratres, ubi se dicit Dóminus vitem, et discípulos suos pálmites, secúndum hoc dicit, quod est caput Ecclésiæ, nosque membra ejus, mediátor Dei et hóminum, homo Christus Jesus. Uníus quippe natúræ sunt vitis et pálmites. Propter quod cum esset Deus, cujus natúræ non sumus, factus est homo, ut in illo esset vitis humána natúra, cujus et nos hómines pálmites esse possémus.\par
\switchcolumn
\EN
\noindent Dearly beloved brethren, this passage of the Gospel, wherein the Lord saith that He is the vine, and that His disciples are the branches, is to be taken in that sense wherein it is also said, that He is the Head of the Church, (Eph. v. 23,) and that we are the members of Him (30) Who is the Mediator between God and men, the man Christ Jesus (1 Tim. ii. 5.) The vine and his branches are of one and the same nature. Therefore, seeing that He was God, of which nature we are not, He was made man, to the end that He might have in Himself this vine, that is, the manhood, whereof we men can be made branches.\par
\switchcolumn*

\LA
\begin{Resp}\Rsign Refúlsit sol in clípeos áureos, et resplenduérunt montes ab eis: \Star{} Et fortitúdo géntium dissipáta est.\end{Resp}
\begin{Resp}\Vsign Erat enim exércitus magnus valde et fortis: et appropiávit Judas et exércitus ejus in prǽlio.\end{Resp}
\begin{Resp}\Rsign Et fortitúdo géntium dissipáta est.\end{Resp}
\switchcolumn
\EN
\begin{Resp}\Rsign The sun shone upon the shields of gold, and the mountains glistered therewith; \Star{} And the army of the heathens was spread abroad.\end{Resp}
\begin{Resp}\Vsign For the army was very great and mighty then Judas and his host drew near and entered into battle.\end{Resp}
\begin{Resp}\Rsign And the army of the heathens was spread abroad.\end{Resp}
\switchcolumn*

\LA
\LessonHead{Lectio II}
\switchcolumn
\EN
\LessonHead{Lesson II}
\switchcolumn*

\LA
\noindent Quid ergo est, Ego sum vitis vera? Numquid ut ádderet, vera, hoc ad eam vitem rétulit, unde ista similitúdo transláta est? Sic enim dícitur vitis per similitúdinem, non per proprietátem: quemádmodum dícitur ovis, agnus, leo, petra, lapis anguláris, et cétera hujúsmodi, quæ magis ipsa sunt vera, ex quibus ducúntur istæ similitúdines, non proprietátes. Sed cum dicit, Ego sum vitis vera: ab illa se útique discérnit, cui dícitur: Quómodo convérsa es in amaritúdinem vitis aliéna? Nam quo pacto est vitis vera, quæ exspectáta est ut fáceret uvam, fecit autem spinas?\par
\switchcolumn
\EN
\noindent Why saith He: I am the true vine? As touching this word true, hath He not here regard to that other parable of a vine, the like figure whereto He doth here apply to Himself? (Jer. ii. 21.) Here is He called a vine, not plainly, but in parable, as also He is called elsewhere a sheep, (Isa. liii. 7, Acts viii. 32,) a lamb, (John i. 36,) a lion, (Apoc. v. 5,) a rock, (1 Cor. x. 4,) a corner-stone, \textasciitilde{}(Eph. ii. 20,) and other things of the like kind. But these things are in themselves that which they seem to be, albeit He is called by their names, not plainly, but in a parable, and herein are they different from that vine, whereof in this place He taketh on Him the name. For when He saith: I am the true vine, doth He not make distinction between Himself, and that which indeed seemed to be a vine, but to which it is said: How art thou turned into the degenerate plant of a strange vine unto Me? (Jer. ii. 21.) For by what title shall that plant be called other than a false vine, whereto they looked that she should bring forth grapes, and she brought forth thorns?\par
\switchcolumn*

\LA
\begin{Resp}\Rsign Ornavérunt fáciem templi corónis áureis, et dedicavérunt altáre Dómino: \Star{} Et facta est lætítia magna in pópulo.\end{Resp}
\begin{Resp}\Vsign In hymnis et confessiónibus benedicébant Dóminum.\end{Resp}
\begin{Resp}\Rsign Et facta est lætítia magna in pópulo.\end{Resp}
\switchcolumn
\EN
\begin{Resp}\Rsign They decked the fore-front of the Temple with crowns of gold, and dedicated the Altar unto the Lord. \Star{} And there was very great gladness among the people.\end{Resp}
\begin{Resp}\Vsign They praised the Lord with Psalms and thanksgiving.\end{Resp}
\begin{Resp}\Rsign And there was very great gladness among the people.\end{Resp}
\switchcolumn*

\LA
\LessonHead{Lectio III}
\switchcolumn
\EN
\LessonHead{Lesson III}
\switchcolumn*

\LA
\noindent Ego sum, inquit, vitis vera: et Pater meus agrícola est. Numquid unum sunt agrícola et vitis? Secúndum hoc ergo vitis Christus, secúndum quod ait: Pater major me est. Secúndum autem id, quod ait: Ego, et Pater unum sumus; et ipse agrícola est: nec talis, quales sunt, qui extrínsecus operándo éxhibent ministérium: sed talis, ut det étiam intrínsecus increméntum. Nam neque qui plantat est áliquid, neque qui rigat: sed, qui increméntum dat, Deus. Sed útique Deus est Christus, quia Deus erat Verbum: unde ipse, et Pater unum sunt. Et, si Verbum caro factum est, quod non erat, manet quod erat.\par
\switchcolumn
\EN
\noindent He saith: I am the true vine, and My Father is the husbandman. Is the vine one with the husbandman? These words then are to be taken in that sense wherein He also saith: My Father is greater than I. (John xiv. 28.) In this sense is He the vine, and the Father is the husbandman. But again, in regard to those words: I and the Father are one, and again and: My Father is the husbandman, we understand that They are not the vine and the husbandman, after the manner of a vine, and the husbandman that from without doth care for and keep it, but after the manner of a vine and Him That from within doth make it to bring forth fruit. For neither is he that planteth anything, neither he that watereth, but God that giveth the increase. (1 Cor. iii. 7.) But Christ is God, for the Word was God. (John i. 1.) Therefore He and the Father are one: and, albeit the Word was made flesh, (John i. 14,) which, before, He was not, He ceased not to be still That Which He was.\par
\switchcolumn*

\LA
\begin{Resp}\Rsign In hymnis et confessiónibus benedicébant Dóminum, \Star{} Qui magna fecit in Israël, et victóriam dedit illis Dóminus omnípotens.\end{Resp}
\begin{Resp}\Vsign Ornavérunt fáciem templi corónis áureis, et dedicavérunt altáre Dómino.\end{Resp}
\begin{Resp}\Rsign Qui magna fecit in Israël, et victóriam dedit illis Dóminus omnípotens.\end{Resp}
\begin{Resp}\Vsign Glória Patri, et Fílio, \Star{} et Spirítui Sancto.\end{Resp}
\begin{Resp}\Rsign Qui magna fecit in Israël, et victóriam dedit illis Dóminus omnípotens.\end{Resp}
\switchcolumn
\EN
\begin{Resp}\Rsign They praised the Lord with psalms and thanksgiving \Star{} Who had done so great things in Israel, and given them the victory the Lord Almighty.\end{Resp}
\begin{Resp}\Vsign They decked the fore-front of the Temple with crowns of gold, and dedicated the Altar unto the Lord.\end{Resp}
\begin{Resp}\Rsign Who had done so great things for Israel, and given them the victory the Lord Almighty.\end{Resp}
\begin{Resp}\Vsign Glory be to the Father, and to the Son, \Star{} and to the Holy Ghost.\end{Resp}
\begin{Resp}\Rsign Who had done so great things for Israel, and given them the victory the Lord Almighty.\end{Resp}
\switchcolumn*

\end{paracol}

\par\needspace{10\baselineskip}\addvspace{1.2em}{\centering\small\textsc{Die 28 Octobris} · \textsc{28 October}\par}
\Office{Ss. Simonis et Judæ Apostolorum}{Sts. Simon and Jude, Apostles}{Duplex II classis}
\addcontentsline{toc}{subsection}{Die 28 Octobris · Ss. Simonis et Judæ Apostolorum \textit{— Sts. Simon and Jude, Apostles}}
\begin{paracol}{2}
\LA
\LessonHead{Lectio I}
\switchcolumn
\EN
\LessonHead{Lesson I}
\switchcolumn*

\LA
\LessonTitle{Incipit Epístola cathólica beáti Judæ Apóstoli}
\switchcolumn
\EN
\LessonTitle{Beginning of the letter of the blessed Apostle Jude}
\switchcolumn*

\LA
\Cite{Judas 1:1-4}
\switchcolumn
\EN
\Cite{Jude 1:1-4}
\switchcolumn*

\LA
\noindent Judas, Jesu Christi servus, frater autem Jacóbi, his qui sunt in Deo Patre diléctis, et Christo Jesu conservátis, et vocátis. Misericórdia vobis, et pax, et cáritas adimpleátur. Caríssimi, omnem sollicitúdinem fáciens scribéndi vobis de commúni vestra salúte, necésse hábui scríbere vobis: déprecans supercertári semel tráditæ sanctis fídei. Subintroiérunt enim quidam hómines (qui olim præscrípti sunt in hoc judícium) ímpii, Dei nostri grátiam transferéntes in luxúriam, et solum Dominatórem, et Dóminum nostrum Jesum Christum negántes.\par
\switchcolumn
\EN
\noindent Jude, the servant of Jesus Christ, and brother of James: to them that are beloved in God the Father, and preserved in Jesus Christ, and called. Mercy unto you, and peace, and charity be fulfilled. Dearly beloved, taking all care to write unto you concerning your common salvation, I was under a necessity to write unto you: to beseech you to contend earnestly for the faith once delivered to the saints. For certain men are secretly entered in, (who were written of long ago unto this judgment,) ungodly men, turning the grace of our Lord God into riotousness, and denying the only sovereign Ruler, and our Lord Jesus Christ.\par
\switchcolumn*

\LA
\begin{Resp}\Rsign Ecce ego mitto vos sicut oves in médio lupórum, dicit Dóminus: \Star{} Estóte ergo prudéntes sicut serpéntes, et símplices sicut colúmbæ.\end{Resp}
\begin{Resp}\Vsign Dum lucem habétis, crédite in lucem, ut fílii lucis sitis.\end{Resp}
\begin{Resp}\Rsign Estóte ergo prudéntes sicut serpéntes, et símplices sicut colúmbæ.\end{Resp}
\switchcolumn
\EN
\begin{Resp}\Rsign Behold, I send you as sheep in the midst of wolves, said the Lord: \Star{} Be ye, therefore, wise as serpents and simple as doves.\end{Resp}
\begin{Resp}\Vsign Whilst you have the light, believe in the light, that you may be the children of light.\end{Resp}
\begin{Resp}\Rsign Be ye therefore wise as serpents and simple as doves.\end{Resp}
\switchcolumn*

\LA
\LessonHead{Lectio II}
\switchcolumn
\EN
\LessonHead{Lesson II}
\switchcolumn*

\LA
\Cite{Judas 1:5-8}
\switchcolumn
\EN
\Cite{Jude 1:5-8}
\switchcolumn*

\LA
\noindent Commonére autem vos volo, sciéntes semel ómnia, quóniam Jesus pópulum de terra Ægýpti salvans, secúndo eos, qui non credidérunt, pérdidit: ángelos vero, qui non servavérunt suum principátum, sed dereliquérunt suum domicílium, in judícium magni diéi, vínculis ætérnis sub calígine reservávit. Sicut Sódoma, et Gomórrha, et finítimæ civitátes símili modo exfornicátæ, et abeúntes post carnem álteram, factæ sunt exémplum, ignis ætérni pœnam sustinéntes. Simíliter et hi carnem quidem máculant, dominatiónem autem spernunt, majestátem autem blasphémant.\par
\switchcolumn
\EN
\noindent I will therefore admonish you, though ye once knew all things, that Jesus, having saved the people out of the land of Egypt, did afterwards destroy them that believed not: And the angels who kept not their principality, but forsook their own habitation, he hath reserved under darkness in everlasting chains, unto the judgment of the great day. As Sodom and Gomorrha, and the neighbouring cities, in like manner, having given themselves to fornication, and going after other flesh, were made an example, suffering the punishment of eternal fire. In like manner these men also defile the flesh, and despise dominion, and blaspheme majesty.\par
\switchcolumn*

\LA
\begin{Resp}\Rsign Tóllite jugum meum super vos, dicit Dóminus, et díscite a me, quia mitis sum et húmilis corde: \Star{} Jugum enim meum suáve est, et onus meum leve.\end{Resp}
\begin{Resp}\Vsign Et inveniétis réquiem animábus vestris.\end{Resp}
\begin{Resp}\Rsign Jugum enim meum suáve est, et onus meum leve.\end{Resp}
\switchcolumn
\EN
\begin{Resp}\Rsign Take up my yoke upon you, said the Lord, and learn of me, because I am meek, and humble of heart. \Star{} For my yoke is sweet and my burden light.\end{Resp}
\begin{Resp}\Vsign And you shall find rest to your souls.\end{Resp}
\begin{Resp}\Rsign For my yoke is sweet and my burden light.\end{Resp}
\switchcolumn*

\LA
\LessonHead{Lectio III}
\switchcolumn
\EN
\LessonHead{Lesson III}
\switchcolumn*

\LA
\Cite{Judas 1:9-13}
\switchcolumn
\EN
\Cite{Jude 1:9-13}
\switchcolumn*

\LA
\noindent Cum Míchael Archángelus cum diábolo dísputans altercarétur de Móysi córpore, non est ausus judícium inférre blasphémiæ: sed dixit: Imperet tibi Dóminus. Hi autem quæcúmque quidem ignórant, blasphémant: quæcúmque autem naturáliter, tamquam muta animália, norunt, in his corrumpúntur. Væ illis, quia in via Cain abiérunt, et erróre Bálaam mercéde effúsi sunt, et in contradictióne Core periérunt! Hi sunt in épulis suis máculæ, convivántes sine timóre, semetípsos pascéntes, nubes sine aqua, quæ a ventis circumferéntur, árbores autumnáles, infructuósæ, bis mórtuæ, eradicátæ, fluctus feri maris, despumántes suas confusiónes, sídera errántia: quibus procélla tenebrárum serváta est in ætérnum.\par
\switchcolumn
\EN
\noindent When Michael the archangel, disputing with the devil, contended about the body of Moses, he durst not bring against him the judgment of railing speech, but said: The Lord command thee. But these men blaspheme whatever things they know not: and what things soever they naturally know, like dumb beasts, in these they are corrupted. Woe unto them, for they have gone in the way of Cain: and after the error of Balaam they have for reward poured out themselves, and have perished in the contradiction of Core. These are spots in their banquets, feasting together without fear, feeding themselves, clouds without water, which are carried about by winds, trees of the autumn, unfruitful, twice dead, plucked up by the roots, Raging waves of the sea, foaming out their own confusion; wandering stars, to whom the storm of darkness is reserved for ever.\par
\switchcolumn*

\LA
\begin{Resp}\Rsign Dum stetéritis ante reges et prǽsides, nolíte cogitáre quómodo aut quid loquámini: \Star{} Dábitur enim vobis in illa hora, quid loquámini.\end{Resp}
\begin{Resp}\Vsign Non enim vos estis qui loquímini: sed Spíritus Patris vestri, qui lóquitur in vobis.\end{Resp}
\begin{Resp}\Rsign Dábitur enim vobis in illa hora, quid loquámini.\end{Resp}
\begin{Resp}\Vsign Glória Patri, et Fílio, \Star{} et Spirítui Sancto.\end{Resp}
\begin{Resp}\Rsign Dábitur enim vobis in illa hora, quid loquámini.\end{Resp}
\switchcolumn
\EN
\begin{Resp}\Rsign But when they shall deliver you up to the judges, take no thought how or what to speak: \Star{} for it shall be given you in that hour what to speak:\end{Resp}
\begin{Resp}\Vsign For it is not you that speak, but the spirit of your Father that speaketh in you.\end{Resp}
\begin{Resp}\Rsign For it shall be given you in that hour what to speak.\end{Resp}
\begin{Resp}\Vsign Glory be to the Father, and to the Son, \Star{} and to the Holy Ghost.\end{Resp}
\begin{Resp}\Rsign For it shall be given you in that hour what to speak.\end{Resp}
\switchcolumn*

\LA
\LessonHead{Lectio IV}
\switchcolumn
\EN
\LessonHead{Lesson IV}
\switchcolumn*

\LA
\noindent Simon Chananǽus, qui et Zelótes, et Thaddǽus, qui et Judas Jacóbi appellátur in Evangélio, uníus ex cathólicis Epístolis scriptor; hic Mesopotámiam, ille Ægýptum evangélica prædicatióne peragrávit. Póstea in Pérsidem conveniéntes, cum innumerábiles fílios Jesu Christo peperíssent fidémque in vastíssimis illis regiónibus et efferátis géntibus disseminássent, doctrína et miráculis, ac dénique glorióso martýrio, simul sanctíssimum Jesu Christi nomen illustrárunt.\par
\switchcolumn
\EN
\noindent Simon the Canaanite, called also Zelotes, went through Egypt preaching the Gospel, whileas the like was done in Mesopotamia by Thaddaeus, called also in the Gospel Judas the brother of James, and the writer of one of the Catholic Epistles. They met together afterwards in Persia, where they begat countless children in Jesus Christ, spread the faith far and wide in those lands, amid raging heathens, and glorified together by their teaching and miracles, and, in the end, by a glorious martyrdom, the most holy name of Jesus Christ.\par
\switchcolumn*

\LA
\begin{Resp}\Rsign Vidi conjúnctos viros, habéntes spléndidas vestes, et Ángelus Dómini locútus est ad me, dicens: \Star{} Isti sunt viri sancti facti amíci Dei.\end{Resp}
\begin{Resp}\Vsign Vidi Ángelum Dei fortem, volántem per médium cælum, voce magna clamántem et dicéntem.\end{Resp}
\begin{Resp}\Rsign Isti sunt viri sancti facti amíci Dei.\end{Resp}
\switchcolumn
\EN
\begin{Resp}\Rsign I saw men standing together, clad in shining raiment, and the Angel of the Lord spoke unto me, saying: \Star{} These men are holy, for they are the friends of God.\end{Resp}
\begin{Resp}\Vsign I saw a strong Angel of God fly into the midst of heaven, saying with a loud voice:\end{Resp}
\begin{Resp}\Rsign These men are holy, for they are the friends of God.\end{Resp}
\switchcolumn*

\LA
\LessonHead{Lectio V}
\switchcolumn
\EN
\LessonHead{Lesson V}
\switchcolumn*

\LA
\LessonTitle{Sermo sancti Gregórii Papæ}
\switchcolumn
\EN
\LessonTitle{Sermon of St. Gregory, Pope}
\switchcolumn*

\LA
\Source{Homilia 30 in Evangelia, post medium}
\switchcolumn
\EN
\Source{Sermon 30 on the Gospels}
\switchcolumn*

\LA
\noindent Scriptum est: Spíritus Dómini ornávit cælos. Ornaménta enim cælórum sunt virtútes prædicántium. Quæ vidélicet ornaménta Paulus enúmerat, dicens: Álii datur per Spíritum sermo sapiéntiæ, álii sermo sciéntiæ secúndum eúndem Spíritum, álteri fides in eódem Spíritu, álii grátia sanitátum in uno Spíritu, álii operátio virtútum, álii prophetía, álii discrétio spirítuum, álii génera linguárum, álii interpretátio sermónum. Hæc autem ómnia operátur unus atque idem Spíritus, dívidens síngulis prout vult.\par
\switchcolumn
\EN
\noindent It is written: By His Spirit the Lord hath adorned the heavens. (Job xxvi. 13.) Now the ornament of the heavens are the godly powers of preachers, and this ornament, what it is, Paul teacheth us thus To one is given by the Spirit the word of wisdom, to another the word of knowledge by the same Spirit; to another faith by the same Spirit; to another the gifts of healing by the same Spirit, to another the working of miracles, to another prophecy, to another discerning of spirits, to another diverse kinds of tongues, to another the interpretation of tongues. But all these worketh that one and the self-same Spirit, dividing to every man severally as He will. (1 Cor. xii. 8.)\par
\switchcolumn*

\LA
\begin{Resp}\Rsign Beáti estis, cum maledíxerint vobis hómines, et persecúti vos fúerint, et díxerint omne malum advérsum vos, mentiéntes, propter me: \Star{} Gaudéte et exsultáte, quóniam merces vestra copiósa est in cælis.\end{Resp}
\begin{Resp}\Vsign Cum vos óderint hómines, et cum separáverint vos, et exprobráverint, et ejécerint nomen vestrum tamquam malum propter Fílium hóminis.\end{Resp}
\begin{Resp}\Rsign Gaudéte et exsultáte, quóniam merces vestra copiósa est in cælis.\end{Resp}
\switchcolumn
\EN
\begin{Resp}\Rsign Blessed are ye when men shall revile you, and persecute you, and shall say all manner of evil against you falsely, for My sake; \Star{} Rejoice, and be exceeding glad, for great is your reward in heaven.\end{Resp}
\begin{Resp}\Vsign When men shall hate you, and when they shall separate you from their company, and shall reproach you, and cast out your name as evil, for the Son of Man's sake.\end{Resp}
\begin{Resp}\Rsign Rejoice, and be exceeding glad, for great is your reward in heaven.\end{Resp}
\switchcolumn*

\LA
\LessonHead{Lectio VI}
\switchcolumn
\EN
\LessonHead{Lesson VI}
\switchcolumn*

\LA
\noindent Quot ergo sunt bona prædicántium, tot sunt ornaménta cælórum. Hinc rursus scriptum est: Verbo Dómini cæli firmáti sunt. Verbum enim Dómini Fílius est Patris. Sed eósdem cælos, vidélicet sanctos Apóstolos, ut tota simul sancta Trínitas ostendátur operáta, repénte de Sancti Spíritus divinitáte adjúngitur: Et Spíritu oris ejus omnis virtus eórum. Cælórum ergo virtus de Spíritu sumpta est: quia mundi hujus potestátibus contraíre non præsúmerent, nisi eos Sancti Spíritus fortitúdo solidásset. Quales namque doctóres sanctæ Ecclésiæ ante advéntum hujus Spíritus fúerint, scimus: et post advéntum illíus, cujus fortitúdinis facti sint, conspícimus.\par
\switchcolumn
\EN
\noindent So much power then as have preachers, so much ornament have the heavens. Wherefore again it is written By the word of the Lord were the heavens made. \textasciitilde{}(Ps. xxxii. 6.) For the Word of the Lord is the Son of the Father. But, to the end that all the Holy Trinity may be made manifest as the Maker of the heavens, that is, of the Apostles, it is straightway added touching God the Holy Ghost: you and all the host of them by the Breath of His mouth. Therefore the might of the same heavens is the might of the Spirit, for they had not braved the powers of this world, unless the strength of the Holy Ghost had comforted them. For we know what manner of men the Teachers of the Holy Church were before the coming of this Spirit and since He came we see in Whose strength they are made strong.\par
\switchcolumn*

\LA
\begin{Resp}\Rsign Isti sunt triumphatóres et amíci Dei, qui contemnéntes jussa príncipum, meruérunt prǽmia ætérna: \Star{} Modo coronántur, et accípiunt palmam.\end{Resp}
\begin{Resp}\Vsign Isti sunt qui venérunt ex magna tribulatióne, et lavérunt stolas suas in sánguine Agni.\end{Resp}
\begin{Resp}\Rsign Modo coronántur, et accípiunt palmam.\end{Resp}
\begin{Resp}\Vsign Glória Patri, et Fílio, \Star{} et Spirítui Sancto.\end{Resp}
\begin{Resp}\Rsign Modo coronántur, et accípiunt palmam.\end{Resp}
\switchcolumn
\EN
\begin{Resp}\Rsign These are they which have conquered, and are become the friends of God, who recked not of the commandments of princes, and earned the everlasting reward. \Star{} And now have they crowns on their heads, and palms in their hands.\end{Resp}
\begin{Resp}\Vsign These are they which came out of great tribulation, and have washed their robes in the blood of the Lamb.\end{Resp}
\begin{Resp}\Rsign And now have they crowns on their heads, and palms in their hands.\end{Resp}
\begin{Resp}\Vsign Glory be to the Father, and to the Son, \Star{} and to the Holy Ghost.\end{Resp}
\begin{Resp}\Rsign And now have they crowns on their heads, and palms in their hands.\end{Resp}
\switchcolumn*

\LA
\LessonHead{Lectio VII}
\switchcolumn
\EN
\LessonHead{Lesson VII}
\switchcolumn*

\LA
\LessonTitle{Léctio sancti Evangélii secúndum Joánnem}
\switchcolumn
\EN
\LessonTitle{From the holy Gospel according to John}
\switchcolumn*

\LA
\Cite{Joannes 15:17-25}
\switchcolumn
\EN
\Cite{John 15:17-25}
\switchcolumn*

\LA
\noindent In illo témpore: Dixit Jesus discípulis suis: Hæc mando vobis, ut diligátis ínvicem. Si mundus vos odit, scitóte quia me priórem vobis ódio hábuit. Et réliqua.\par
\switchcolumn
\EN
\noindent at that time, Jesus said to his disciples: These things I command you, that you love one another. If the world hate you, know ye, that it hath hated me before you. And so on.\par
\switchcolumn*

\LA
\LessonTitle{Homilía sancti Augustíni Epíscopi}
\switchcolumn
\EN
\LessonTitle{Homily by St. Augustine, Bishop (of Hippo.)}
\switchcolumn*

\LA
\Source{Tractatus 87 in Joannem}
\switchcolumn
\EN
\Source{87th Tract on John}
\switchcolumn*

\LA
\noindent In lectióne evangélica quæ hanc antecédit, díxerat Dóminus: Non vos me elegístis; sed ego elégi vos et pósui vos, ut eátis, et fructum afferátis, et fructus vester máneat: ut quodcúmque petiéritis Patrem in nómine meo, det vobis. Hic autem dicit: Hæc mando vobis, ut diligátis ínvicem. Ac per hoc intellégere debémus hunc esse fructum nostrum, de quo ait: Ego vos elégi, ut eátis, et fructum afferátis, et fructus vester máneat. Et quod adjúnxit, Ut quodcúmque petiéritis Patrem in nómine meo, det vobis, tunc útique dabit nobis, si diligámus ínvicem; cum et hoc ipsum ipse déderit nobis, qui nos elégit non habéntes fructum, quia non eum nos elegerámus, et pósuit nos ut fructum afferámus, hoc est, ínvicem diligámus.\par
\switchcolumn
\EN
\noindent In the reading from the Gospel, the last before this, the Lord had said: Ye have not chosen Me, but I have chosen you, and ordained you, that ye should go, and bring forth fruit, and that your fruit should remain; that whatsoever ye shall ask of the Father in My Name, He may give it you. And here He saith: These things I command you, that ye love one another. And by this it is that we must understand what fruit from us it is, whereof He saith: I have chosen, that ye should go, and bring forth fruit, and that your fruit should remain, and so the words added That whatsoever ye shall ask of the Father in My Name, He may give it you. He will give unto us when we love one another, since this (mutual love) is itself the gift of Him Who hath chosen us when as yet we were fruitless, since it hath not been we who have chosen Him, (but He Who hath chosen us,) and ordained us, that we should go, and bring forth fruit, that is to say, should love one another.\par
\switchcolumn*

\LA
\begin{Resp}\Rsign Isti sunt qui vivéntes in carne, plantavérunt Ecclésiam sánguine suo: \Star{} Cálicem Dómini bibérunt, et amíci Dei facti sunt.\end{Resp}
\begin{Resp}\Vsign In omnem terram exívit sonus eórum, et in fines orbis terræ verba eórum.\end{Resp}
\begin{Resp}\Rsign Cálicem Dómini bibérunt, et amíci Dei facti sunt.\end{Resp}
\switchcolumn
\EN
\begin{Resp}\Rsign These are they who while yet they lived in the flesh, planted the Church in their own blood; \Star{} They drank of the Lord's cup, and became the friends of God.\end{Resp}
\begin{Resp}\Vsign Their sound is gone out through all the earth, and their words to the ends of the world.\end{Resp}
\begin{Resp}\Rsign They drank of the Lord's cup, and became the friends of God.\end{Resp}
\switchcolumn*

\LA
\LessonHead{Lectio VIII}
\switchcolumn
\EN
\LessonHead{Lesson VIII}
\switchcolumn*

\LA
\noindent Cáritas ergo est fructus noster, quam defínit Apóstolus, De corde puro, et consciéntia bona, et fide non ficta. Hac dilígimus ínvicem, hac dilígimus Deum; neque enim vera dilectióne diligerémus ínvicem, nisi diligéntes Deum. Díligit enim unusquísque próximum suum tamquam seípsum, si díligit Deum. Nam, si non díligit Deum, non díligit seípsum; in his enim duóbus præcéptis caritátis tota lex pendet et prophétæ. Hic est fructus noster. De fructu ítaque nobis mandans, Hæc mando, inquit, vobis, ut diligátis ínvicem. Unde et Apóstolus Paulus, cum contra ópera carnis commendáre fructum spíritus vellet, a cápite hoc pósuit: Fructus, inquit, spíritus, cáritas est; ac deínde cétera, tamquam ex isto cápite exórta et religáta contéxuit, quæ sunt, gáudium, pax, longanímitas, benígnitas, bónitas, fides, mansuetúdo, continéntia, cástitas.\par
\switchcolumn
\EN
\noindent Love then, is the fruit which we should bring forth, and the Apostle Paul telleth us (1 Tim. i. 5) that this love is love out of a pure heart, and of a good conscience, and of faith unfeigned. This is the love wherewith we love our neighbour, the love wherewith we love God, for we do not really love our neighbour unless we love God. For if any man love God, he loveth his neighbour as himself, since he that loveth not God loveth not himself. For on these two commandments hangeth all the law and the Prophets. Love, then, is the fruit which we should bring forth. And concerning this fruit, the Lord giveth us this commandment: These things (saith He) I command you, that ye love one another. Hence also the Apostle Paul (Gal. v. 22) when he is about praising up the fruits of the Spirit as opposed to the works of the flesh, saith first of all: The fruit of the Spirit is love. And from that as the beginning he draweth out a string of other fruits, as thence begotten and thereto bound, namely, joy, peace, long-suffering, gentleness, goodness, faith, meekness, temperance, chastity.\par
\switchcolumn*

\LA
\begin{Resp}\Rsign Isti sunt viri sancti, quos elégit Dóminus in caritáte non ficta, et dedit illis glóriam sempitérnam: \Star{} Quorum doctrína fulget Ecclésia, ut sole luna.\end{Resp}
\begin{Resp}\Vsign Sancti per fidem vicérunt regna: operáti sunt justítiam.\end{Resp}
\begin{Resp}\Rsign Quorum doctrína fulget Ecclésia, ut sole luna.\end{Resp}
\begin{Resp}\Vsign Glória Patri, et Fílio, \Star{} et Spirítui Sancto.\end{Resp}
\begin{Resp}\Rsign Quorum doctrína fulget Ecclésia, ut sole luna.\end{Resp}
\switchcolumn
\EN
\begin{Resp}\Rsign These men are saints, whom the Lord hath chosen in love unfeigned, and hath given them glory everlasting. These are they \Star{} By the light of whose teaching the Church is glorified, even as the moon is glorified by the light of the sun.\end{Resp}
\begin{Resp}\Vsign The saints through faith subdued kingdoms, wrought righteousness.\end{Resp}
\begin{Resp}\Rsign By the light of whose teaching the Church is glorified, even as the moon is glorified by the light of the sun.\end{Resp}
\begin{Resp}\Vsign Glory be to the Father, and to the Son, \Star{} and to the Holy Ghost.\end{Resp}
\begin{Resp}\Rsign By the light of whose teaching the Church is glorified, even as the moon is glorified by the light of the sun.\end{Resp}
\switchcolumn*

\LA
\LessonHead{Lectio IX}
\switchcolumn
\EN
\LessonHead{Lesson IX}
\switchcolumn*

\LA
\noindent Quis autem bene gaudet, qui bonum non díligit unde gaudet? Quis pacem veram, nisi cum illo potest habére, quem veráciter díligit? Quis est longánimis in bono ópere perseveránter manéndo, nisi férveat diligéndo? Quis est benígnus, nisi díligat cui opitulátur? Quis bonus, nisi diligéndo efficiátur? Quis salúbriter fidélis, nisi ea fide quæ per dilectiónem operátur? Quis utíliter mansuétus, cui non diléctio moderétur? Quis ab eo cóntinet unde turpátur, nisi díligat unde honestátur? Mérito ítaque Magíster bonus dilectiónem sic sæpe comméndat, tamquam sola præcipiénda sit, sine qua non possunt prodésse cétera bona, et quæ non potest habéri sine céteris bonis, quibus homo effícitur bonus.\par
\switchcolumn
\EN
\noindent Who is really joyful that loveth not the cause of his joy? Who can really be at one with another, unless he loveth that other? Who is cheerful under long toil for a good work, unless he loveth the aim? Who is kind, unless he love the object of his tenderness? Who is good, unless by the persuasion of love? Who is truly faithful, unless by the faith which worketh by love? Who is gentle to any use, unless love move him? Who turneth away from baseness unless he love honour? Well, then, doth the Good Master so often command us to love, as though that commandment were all-sufficient, for love is that gift without which all other good things avail nothing, and which cannot be without having every other good gift which maketh a good man good.\par
\switchcolumn*

\LA
\TeDeum{Te Deum}
\switchcolumn
\EN
\TeDeum{Te Deum}
\switchcolumn*

\end{paracol}

\par\needspace{10\baselineskip}\addvspace{1.2em}{\centering\small\textsc{Die 31 Octobris} · \textsc{31 October}\par}
\Office{In Vigilia Omnium Sanctorum}{Vigil of All Saints}{Vigilia}
\addcontentsline{toc}{subsection}{Die 31 Octobris · In Vigilia Omnium Sanctorum \textit{— Vigil of All Saints}}
\begin{paracol}{2}
\LA
\LessonHead{Lectio I}
\switchcolumn
\EN
\LessonHead{Lesson I}
\switchcolumn*

\LA
\LessonTitle{Léctio sancti Evangélii secúndum Lucam}
\switchcolumn
\EN
\LessonTitle{From the Holy Gospel according to Luke}
\switchcolumn*

\LA
\Cite{Luc 6:17-23}
\switchcolumn
\EN
\Cite{Luke 6:17-23}
\switchcolumn*

\LA
\noindent In illo témpore: Descéndens Jesus de monte, stetit in loco campéstri, et turba discipulórum ejus, et multitúdo copiósa plebis ab omni Judǽa, et Jerúsalem, et marítima, et Tyri, et Sidónis. Et réliqua.\par
\switchcolumn
\EN
\noindent At that time, Jesus came down from the mountain, and stood in the plain, and the company of His disciples, and a great multitude of people out of all Judea, and Jerusalem, and from the sea coast of Tyre and Sidon. And so on.\par
\switchcolumn*

\LA
\LessonTitle{Homilía sancti Ambrósii Epíscopi}
\switchcolumn
\EN
\LessonTitle{Homily by St Ambrose, Bishop (of Milan.)}
\switchcolumn*

\LA
\Source{Lib. 5. in Lucam. Cap. 6., post init.}
\switchcolumn
\EN
\Source{Bk. v. on Luke vi.}
\switchcolumn*

\LA
\noindent Advérte ómnia diligénter, quómodo et cum Apóstolis ascéndat et descéndat ad turbas. Quómodo enim turba nisi in húmili Christum vidéret? Non séquitur ad excélsa, non ascéndit ad sublímia. Dénique ubi descéndit, invénit infírmos: in excélsis enim infírmi esse non possunt. Hinc étiam Matthǽus docet in inferióribus débiles esse sanátos. Prius enim unusquísque sanándus est, ut paulátim virtútibus procedéntibus ascéndere possit ad montem; et ídeo quemque in inferióribus sanat, hoc est, a libídine révocat, injúriam cæcitátis avértit. Ad vúlnera nostra descéndit: ut usu quodam et cópia suæ natúræ, compartícipes nos fáciat esse regni cæléstis.\par
\switchcolumn
\EN
\noindent Mark well how Jesus goeth upward with His disciples, and downward to the multitude. How should the multitude behold Christ, save in a lower place? Such go not up to the things which are above; such attain not to the things which are high. And when Jesus cometh down, He findeth such as are diseased for such like go not up to the heights. Hence also Matthew saith that there were there all sick people, (iv. 23.) Of these every man had need of healing, that, when he had received strength, by and by, he might go up into the mountain. And therefore, being Himself come down, He healeth them in the plain, that is to say, He calleth them away from their lust, and freeth them of their blindness. He cometh down to our wounds, to the end that by a certain use of His nature, and by the abundance thereof, He might make us joint-heirs of the kingdom of heaven.\par
\switchcolumn*

\LA
\begin{Resp}\Rsign Refúlsit sol in clípeos áureos, et resplenduérunt montes ab eis: \Star{} Et fortitúdo géntium dissipáta est.\end{Resp}
\begin{Resp}\Vsign Erat enim exércitus magnus valde et fortis: et appropiávit Judas et exércitus ejus in prǽlio.\end{Resp}
\begin{Resp}\Rsign Et fortitúdo géntium dissipáta est.\end{Resp}
\switchcolumn
\EN
\begin{Resp}\Rsign The sun shone upon the shields of gold, and the mountains glistered therewith; \Star{} And the army of the heathens was spread abroad.\end{Resp}
\begin{Resp}\Vsign For the army was very great and mighty then Judas and his host drew near and entered into battle.\end{Resp}
\begin{Resp}\Rsign And the army of the heathens was spread abroad.\end{Resp}
\switchcolumn*

\LA
\LessonHead{Lectio II}
\switchcolumn
\EN
\LessonHead{Lesson II}
\switchcolumn*

\LA
\noindent Beáti páuperes, quia vestrum est regnum Dei. Quátuor tantum beatitúdines sanctus Lucas Domínicas posuit, octo vero sanctus Matthǽus: sed in illis octo istæ quátuor sunt, et in quátuor istis illæ octo. Hic enim quátuor velut virtútes ampléxus est cardináles: ille in illis octo mýsticum númerum reserávit. Pro octáva enim multi inscribúntur Psalmi: et mandátum áccipis octo illis partem dare fortásse benedictiónibus. Sicut enim spei nostræ octáva perféctio est, ita octáva summa virtútum est.\par
\switchcolumn
\EN
\noindent Blessed be ye poor, for your's is the kingdom of God. Saint Luke giveth us but four of the Lord's Beatitudes, and Saint Matthew eight but in those eight are contained these four, and in these four those eight. For in these four are embraced the cardinal virtues and in those eight they are set forth in a number full of mystery. It is written at the head of more than one of the Psalms that they are for the octave, and thou hast received the commandment Give a portion to seven, and also to eight to seven or eight what? Perchance degrees of blessedness. For as this eighth Beatitude doth name the most glorious realization of our hope the kingdom of Heaven so doth it also name the most royal exertion of our strength blessed are they which are persecuted.\par
\switchcolumn*

\LA
\begin{Resp}\Rsign Ornavérunt fáciem templi corónis áureis, et dedicavérunt altáre Dómino: \Star{} Et facta est lætítia magna in pópulo.\end{Resp}
\begin{Resp}\Vsign In hymnis et confessiónibus benedicébant Dóminum.\end{Resp}
\begin{Resp}\Rsign Et facta est lætítia magna in pópulo.\end{Resp}
\switchcolumn
\EN
\begin{Resp}\Rsign They decked the fore-front of the Temple with crowns of gold, and dedicated the Altar unto the Lord. \Star{} And there was very great gladness among the people.\end{Resp}
\begin{Resp}\Vsign They praised the Lord with Psalms and thanksgiving.\end{Resp}
\begin{Resp}\Rsign And there was very great gladness among the people.\end{Resp}
\switchcolumn*

\LA
\LessonHead{Lectio III}
\switchcolumn
\EN
\LessonHead{Lesson III}
\switchcolumn*

\LA
\noindent Sed prius quæ sunt amplióra videámus. Beáti, inquit, páuperes, quóniam vestrum est regnum Dei. Primam benedictiónem hanc utérque Evangelísta pósuit. Órdine enim prima est, et parens quædam generatióque virtútum: quia qui contempsérit sæculária, ipse merébitur sempitérna: nec potest quisquam méritum regni cæléstis adipísci, qui mundi cupiditáte pressus, emergéndi non habet facultátem.\par
\switchcolumn
\EN
\noindent But let us first consider the fuller of the forms of these Beatitudes. Blessed be ye poor, for your's is the kingdom of God. Both of the Evangelists give to this Beatitude the first place. Yea, surely, for poorness, at least in spirit, is the first in order, the mother, and procreatrix of virtues; since he that setteth no store by temporal things, winneth toward eternal things; neither is any man able to gain the kingdom of heaven, on whom the love of this present world doth so press, that he cannot rid himself thereof.\par
\switchcolumn*

\LA
\begin{Resp}\Rsign In hymnis et confessiónibus benedicébant Dóminum, \Star{} Qui magna fecit in Israël, et victóriam dedit illis Dóminus omnípotens.\end{Resp}
\begin{Resp}\Vsign Ornavérunt fáciem templi corónis áureis, et dedicavérunt altáre Dómino.\end{Resp}
\begin{Resp}\Rsign Qui magna fecit in Israël, et victóriam dedit illis Dóminus omnípotens.\end{Resp}
\begin{Resp}\Vsign Glória Patri, et Fílio, \Star{} et Spirítui Sancto.\end{Resp}
\begin{Resp}\Rsign Qui magna fecit in Israël, et victóriam dedit illis Dóminus omnípotens.\end{Resp}
\switchcolumn
\EN
\begin{Resp}\Rsign They praised the Lord with psalms and thanksgiving \Star{} Who had done so great things in Israel, and given them the victory the Lord Almighty.\end{Resp}
\begin{Resp}\Vsign They decked the fore-front of the Temple with crowns of gold, and dedicated the Altar unto the Lord.\end{Resp}
\begin{Resp}\Rsign Who had done so great things for Israel, and given them the victory the Lord Almighty.\end{Resp}
\begin{Resp}\Vsign Glory be to the Father, and to the Son, \Star{} and to the Holy Ghost.\end{Resp}
\begin{Resp}\Rsign Who had done so great things for Israel, and given them the victory the Lord Almighty.\end{Resp}
\switchcolumn*

\end{paracol}

\SeasonTitle{Mensis November}{November}

\addcontentsline{toc}{section}{Mensis November \textit{— November}}

\par\needspace{10\baselineskip}\addvspace{1.2em}{\centering\small\textsc{Die 1 Novembris} · \textsc{1 November}\par}
\Office{Omnium Sanctorum}{All Saints}{Duplex I classis cum Octava communi}
\addcontentsline{toc}{subsection}{Die 1 Novembris · Omnium Sanctorum \textit{— All Saints}}
\begin{paracol}{2}
\LA
\LessonHead{Lectio I}
\switchcolumn
\EN
\LessonHead{Lesson I}
\switchcolumn*

\LA
\LessonTitle{De libro Apocalýpsis beáti Joánnis Apóstoli}
\switchcolumn
\EN
\LessonTitle{Lesson from the Apocalypse of St. John}
\switchcolumn*

\LA
\Cite{Apo 4:2-8}
\switchcolumn
\EN
\Cite{Rev 4:2-8}
\switchcolumn*

\LA
\noindent Et ecce sedes pósita erat in cælo, et supra sedem sedens. Et qui sedébat símilis erat aspéctui lápidis jáspidis, et sárdinis: et iris erat in circúitu sedis símilis visióni smarágdinæ. Et in circúitu sedis sedília vigínti quátuor: et super thronos vigínti quátuor senióres sedéntes, circumamícti vestiméntis albis, et in capítibus eórum corónæ áureæ. Et de throno procedébant fúlgura, et voces, et tonítrua: et septem lámpades ardéntes ante thronum, qui sunt septem spíritus Dei. Et in conspéctu sedis tamquam mare vítreum símile crystállo: et in médio sedis, et in circúitu sedis quátuor animália plena óculis ante et retro. Et ánimal primum símile leóni, et secúndum ánimal símile vítulo, et tértium ánimal habens fáciem quasi hóminis, et quartum ánimal símile áquilæ volánti. Et quátuor animália, síngula eórum habébant alas senas: et in circúitu, et intus plena sunt óculis: et réquiem non habébant die ac nocte, dicéntia: Sanctus, Sanctus, Sanctus Dóminus Deus omnípotens, qui erat, et qui est, et qui ventúrus est.\par
\switchcolumn
\EN
\noindent And behold there was a throne set in heaven, and upon the throne one sitting. And he that sat, was to the sight like the jasper and the sardine stone; and there was a rainbow round about the throne, in sight like unto an emerald. And round about the throne were four and twenty seats; and upon the seats, four and twenty ancients sitting, clothed in white garments, and on their heads were crowns of gold. And from the throne proceeded lightnings, and voices, and thunders; and there were seven lamps burning before the throne, which are the seven spirits of God. And in the sight of the throne was, as it were, a sea of glass like to crystal; and in the midst of the throne, and round about the throne, were four living creatures, full of eyes before and behind. And the first living creature was like a lion: and the second living creature like a calf: and the third living creature, having the face, as it were, of a man: and the fourth living creature was like an eagle flying. And the four living creatures had each of them six wings; and round about and within they are full of eyes. And they rested not day and night, saying: Holy, holy, holy, Lord God Almighty, who was, and who is, and who is to come.\par
\switchcolumn*

\LA
\begin{Resp}\Rsign Vidi Dóminum sedéntem super sólium excélsum et elevátum, et plena erat omnis terra majestáte ejus: \Star{} Et ea, quæ sub ipso erant, replébant templum.\end{Resp}
\begin{Resp}\Vsign Séraphim stabant super illud: sex alæ uni, et sex alæ álteri.\end{Resp}
\begin{Resp}\Rsign Et ea, quæ sub ipso erant, replébant templum.\end{Resp}
\switchcolumn
\EN
\begin{Resp}\Rsign I saw the Lord sitting upon a throne, high and lifted up, and the whole earth was full of His glory; \Star{} And His train filled the temple.\end{Resp}
\begin{Resp}\Vsign Above it stood the Seraphim each one had six wings.\end{Resp}
\begin{Resp}\Rsign And His train filled the temple.\end{Resp}
\switchcolumn*

\LA
\LessonHead{Lectio II}
\switchcolumn
\EN
\LessonHead{Lesson II}
\switchcolumn*

\LA
\Cite{Apo 5:1-8}
\switchcolumn
\EN
\Cite{Rev 5:1-8}
\switchcolumn*

\LA
\noindent Et vidi in déxtera sedéntis supra thronum, librum scriptum intus et foris, signátum sigíllis septem. Et vidi ángelum fortem, prædicántem voce magna: Quis est dignus aperíre librum, et sólvere signácula ejus? Et nemo póterat neque in cælo, neque in terra, neque subtus terram aperíre librum, neque respícere illum. Et ego flebam multum, quóniam nemo dignus invéntus est aperíre librum, nec vidére eum. Et unus de senióribus dixit mihi: Ne fléveris: ecce vicit leo de tribu Juda, radix David, aperíre librum, et sólvere septem signácula ejus. Et vidi: et ecce in médio throni et quátuor animálium, et in médio seniórum, Agnum stantem tamquam occísum, habéntem córnua septem, et óculos septem: qui sunt septem spíritus Dei, missi in omnem terram. Et venit: et accépit de déxtera sedéntis in throno librum. Et cum aperuísset librum, quátuor animália, et vigínti quátuor senióres cecidérunt coram Agno, habéntes sínguli cítharas, et phíalas áureas plenas odoramentórum, quæ sunt oratiónes sanctórum:\par
\switchcolumn
\EN
\noindent And I saw in the right hand of him that sat on the throne, a book written within and without, sealed with seven seals. And I saw a strong angel, proclaiming with a loud voice: Who is worthy to open the book, and to loose the seals thereof? And no man was able, neither in heaven, nor on earth, nor under the earth, to open the book, nor to look on it. And I wept much, because no man was found worthy to open the book, nor to see it. And one of the ancients said to me: Weep not; behold the lion of the tribe of Juda, the root of David, hath prevailed to open the book, and to loose the seven seals thereof. And I saw: and behold in the midst of the throne and of the four living creatures, and in the midst of the ancients, a Lamb standing as it were slain, having seven horns and seven eyes: which are the seven Spirits of God, sent forth into all the earth. And he came and took the book out of the right hand of him that sat on the throne. And when he had opened the book, the four living creatures, and the four and twenty ancients fell down before the Lamb, having every one of them harps, and golden vials full of odours, which are the prayers of saints:\par
\switchcolumn*

\LA
\begin{Resp}\Rsign Beáta es Virgo María, Dei Génetrix, quæ credidísti Dómino: perfécta sunt in te quæ dicta sunt tibi: ecce exaltáta es super choros Angelórum: \Star{} Intercéde pro nobis ad Dóminum, Deum nostrum.\end{Resp}
\begin{Resp}\Vsign Ave, María, grátia plena, Dóminus tecum.\end{Resp}
\begin{Resp}\Rsign Intercéde pro nobis ad Dóminum, Deum nostrum.\end{Resp}
\switchcolumn
\EN
\begin{Resp}\Rsign O Virgin Mary, Mother of God, blessed art thou that didst believe the Lord, for there hath been a performance of those things which were told thee from the Lord. Behold, thou art exalted over choirs of Angels. \Star{} Plead for us with the Lord our God.\end{Resp}
\begin{Resp}\Vsign Hail, Mary, full of grace, the Lord is with thee.\end{Resp}
\begin{Resp}\Rsign Plead for us with the Lord our God.\end{Resp}
\switchcolumn*

\LA
\LessonHead{Lectio III}
\switchcolumn
\EN
\LessonHead{Lesson III}
\switchcolumn*

\LA
\Cite{Apo 5:9-14}
\switchcolumn
\EN
\Cite{Rev 5:9-14}
\switchcolumn*

\LA
\noindent Et cantábant cánticum novum, dicéntes: Dignus es, Dómine, accípere librum, et aperíre signácula ejus: quóniam occísus es, et redemísti nos Deo in sánguine tuo ex omni tribu, et lingua, et pópulo, et natióne: et fecísti nos Deo nostro regnum, et sacerdótes: et regnábimus super terram. Et vidi, et audívi vocem angelórum multórum in circúitu throni, et animálium, et seniórum: et erat númerus eórum míllia míllium, dicéntium voce magna: Dignus est Agnus, qui occísus est, accípere virtútem, et divinitátem, et sapiéntiam, et fortitúdinem, et honórem, et glóriam, et benedictiónem. Et omnem creatúram, quæ in cælo est, et super terram, et sub terra, et quæ sunt in mari, et quæ in eo: omnes audívi dicéntes: Sedénti in throno, et Agno, benedíctio et honor, et glória, et potéstas in sǽcula sæculórum. Et quátuor animália dicébant: Amen. Et vigínti quátuor senióres cecidérunt in fácies suas: et adoravérunt vivéntem in sǽcula sæculórum.\par
\switchcolumn
\EN
\noindent And they sung a new canticle, saying: Thou art worthy, O Lord, to take the book, and to open the seals thereof; because thou wast slain, and hast redeemed us to God, in thy blood, out of every tribe, and tongue, and people, and nation. And hast made us to our God a kingdom and priests, and we shall reign on the earth. And I beheld, and I heard the voice of many angels round about the throne, and the living creatures, and the ancients; and the number of them was thousands of thousands, Saying with a loud voice: The Lamb that was slain is worthy to receive power, and divinity, and wisdom, and strength, and honour, and glory, and benediction. And every creature, which is in heaven, and on the earth, and under the earth, and such as are in the sea, and all that are in them: I heard all saying: To him that sitteth on the throne, and to the Lamb, benediction, and honour, and glory, and power, for ever and ever. And the four living creatures said: Amen. And the four and twenty ancients fell down on their faces, and adored him that liveth for ever and ever.\par
\switchcolumn*

\LA
\begin{Resp}\Rsign In conspéctu Angelórum psallam tibi: \Star{} Et adorábo ad templum sanctum tuum, et confitébor nómini tuo, Dómine.\end{Resp}
\begin{Resp}\Vsign Super misericórdia tua et veritáte tua, quóniam magnificásti super nos nomen sanctum tuum.\end{Resp}
\begin{Resp}\Rsign Et adorábo ad templum sanctum tuum, et confitébor nómini tuo, Dómine.\end{Resp}
\begin{Resp}\Vsign Glória Patri, et Fílio, \Star{} et Spirítui Sancto.\end{Resp}
\begin{Resp}\Rsign Et adorábo ad templum sanctum tuum, et confitébor nómini tuo, Dómine.\end{Resp}
\switchcolumn
\EN
\begin{Resp}\Rsign Before the Angels will I sing praise unto thee, and will worship toward thy holy temple. \Star{} And I will praise thy Name, O Lord.\end{Resp}
\begin{Resp}\Vsign For thy lovingkindness and for thy truth; for Thou hast glorified thine holy Name on us.\end{Resp}
\begin{Resp}\Rsign And I will praise thy Name, O Lord.\end{Resp}
\begin{Resp}\Vsign Glory be to the Father, and to the Son, \Star{} and to the Holy Ghost.\end{Resp}
\begin{Resp}\Rsign And I will praise thy Name, O Lord.\end{Resp}
\switchcolumn*

\LA
\LessonHead{Lectio IV}
\switchcolumn
\EN
\LessonHead{Lesson IV}
\switchcolumn*

\LA
\LessonTitle{Sermo sancti Bedæ Venerábilis Presbýteri}
\switchcolumn
\EN
\LessonTitle{From the Sermons of the Venerable Bede, Priest (at Jarrow.)}
\switchcolumn*

\LA
\Source{Sermo 18 de Sanctis}
\switchcolumn
\EN
\Source{18th upon the Saints}
\switchcolumn*

\LA
\noindent Hódie, dilectíssimi, ómnium Sanctórum sub una solemnitátis lætítia celebrámus festivitátem: quorum societáte cælum exsúltat, quorum patrocíniis terra lætátur, triúmphis Ecclésia sancta coronátur, quorum conféssio quanto in passióne fórtior, tanto est clárior in honóre: quia, dum crevit pugna, crevit et pugnántium glória, et martýrii triúmphus multíplici passiónum génere adornátur: perque gravióra torménta, gravióra fuére et prǽmia: dum cathólica mater Ecclésia, quæ per totum orbem longe latéque diffúsa est, in ipso cápite suo Christo Jesu edúcta est, contumélias, cruces et mortem non timére; magis magísque roboráta, non resisténdo, sed perferéndo, univérsis, quos ágmine ínclito carcer pœnális inclúsit, pari et símili calóre virtútis ad geréndum certámen, glóriam triumphálem inspirávit.\par
\switchcolumn
\EN
\noindent Dearly beloved brethren: This day we keep, with one great cry of joy, a Feast in memory of all God's holy children; His children, whose presence is a gladness to heaven; His children, whose prayers are a blessing to earth; His children, whose victories are the crown of the Holy Church; His chosen, whose testifying is the more glorious in honour, as the agony in which it was given was the sterner in intensity, for as the dreader grew the battle, so the grander grew the fighters, and the triumph of martyrdom waxed the more incisive by the multiplicity of suffering, and the heavier the torment the heavier the prize. And it is our Mother, the Catholic Church, spread far and wide throughout all this planet, it is she that hath learnt, in Christ Jesus her Head, not to fear shame, nor cross, nor death, but hath waxed lealer and lealer, and, not by fighting, but by enduring, hath breathed into all that noble band who have come up to the bitter starting-post the hope of conquest and glory which hath warmed them manfully to accept the race.\par
\switchcolumn*

\LA
\begin{Resp}\Rsign Præcúrsor Dómini venit, de quo ipse testátur: \Star{} Nullus major inter natos mulíerum Joánne Baptísta.\end{Resp}
\begin{Resp}\Vsign Hic est enim prophéta, et plus quam prophéta, de quo Salvátor ait.\end{Resp}
\begin{Resp}\Rsign Nullus major inter natos mulíerum Joánne Baptísta.\end{Resp}
\switchcolumn
\EN
\begin{Resp}\Rsign The Fore-runner of the Lord cometh, to whom He Himself bare witness, saying: \Star{} Among them that are born of women there hath not risen a greater than John the Baptist.\end{Resp}
\begin{Resp}\Vsign A Prophet? Yea, and much more than a Prophet. This is he of whom the Saviour saith:\end{Resp}
\begin{Resp}\Rsign Among them that are born of women there hath not risen a greater than John the Baptist.\end{Resp}
\switchcolumn*

\LA
\LessonHead{Lectio V}
\switchcolumn
\EN
\LessonHead{Lesson V}
\switchcolumn*

\LA
\noindent O vere beáta mater Ecclésia, quam sic honor divínæ dignatiónis illúminat, quam vincéntium gloriósus Mártyrum sanguis exórnat, quam inviolátæ confessiónis cándida índuit virgínitas! Flóribus ejus nec rosæ nec lília desunt. Certent nunc caríssimi, sínguli, ut ad utrósque honóres amplíssimam accípiant dignitátem, corónas, vel de virginitáte cándidas, vel de passióne purpúreas. In cæléstibus castris pax et ácies habent flores suos, quibus mílites Christi coronántur.\par
\switchcolumn
\EN
\noindent If a verity thou art blessed, O my Mother the Church! The blaze of God's mercy beateth full upon thee; thine adornment is the glorious blood of victorious Martyrs, and thy raiment the virgin whiteness of untarnished orthodoxy. thy garlands lack neither roses nor lilies. And now, dearly beloved brethren, let each one of us strive to gain the goodly crown of one sort or the other, either the glistening whiteness of purity, or the red dye of suffering. In the army in heaven peace and war have both chaplets of their own, to crown Christ's soldiers withal.\par
\switchcolumn*

\LA
\begin{Resp}\Rsign Isti sunt qui vivéntes in carne, plantavérunt Ecclésiam sánguine suo: \Star{} Cálicem Dómini bibérunt, et amíci Dei facti sunt.\end{Resp}
\begin{Resp}\Vsign In omnem terram exívit sonus eórum, et in fines orbis terræ verba eórum.\end{Resp}
\begin{Resp}\Rsign Cálicem Dómini bibérunt, et amíci Dei facti sunt.\end{Resp}
\switchcolumn
\EN
\begin{Resp}\Rsign These are they who while yet they lived in the flesh, planted the Church in their own blood; \Star{} They drank of the Lord's cup, and became the friends of God.\end{Resp}
\begin{Resp}\Vsign Their sound is gone out through all the earth, and their words to the ends of the world.\end{Resp}
\begin{Resp}\Rsign They drank of the Lord's cup, and became the friends of God.\end{Resp}
\switchcolumn*

\LA
\LessonHead{Lectio VI}
\switchcolumn
\EN
\LessonHead{Lesson VI}
\switchcolumn*

\LA
\noindent Dei enim ineffábilis et imménsa bónitas étiam hoc provídit, ut labórum quidem tempus et agónis non exténderet, nec longum fáceret, aut ætérnum, sed breve, et ut ita dicam, momentáneum: ut in hac brevi et exígua vita agónes essent et labóres; in illa vero quæ ætérna est, corónæ et prǽmia meritórum: ut labóres quidem cito finiréntur, meritórum vero prǽmia sine fine durárent: ut post hujus mundi ténebras visúri essent candidíssimam lucem, et acceptúri majórem passiónum cunctárum acerbitátibus beatitúdinem, testánte hoc idem Apóstolo, ubi ait: Non sunt condígnæ passiónes hujus témporis ad superventúram glóriam, quæ revelábitur in nobis.\par
\switchcolumn
\EN
\noindent Moreover, to this also hath the unutterable and boundless goodness of God seen, that He spreadeth not the time of working and wrestling, neither maketh it long, nor everlasting, and, as it were, but for a moment, so that in this short and scanty life there is wrestling and working, but the crown and the prize is in a life which is eternal. So the work is soon over, but the wage is paid for ever. And when the night of this world is over, the Saints are to see the clearness of the essential light, and to receive a blessedness outweighing the pangs of any torment, as testifieth the Apostle Paul, where he saith: The sufferings of this present time are not worthy to be compared with the glory which shall be revealed in us. (Rom. viii. 18.)\par
\switchcolumn*

\LA
\begin{Resp}\Rsign Sancti mei, qui in carne pósiti, certámen habuístis: \Star{} Mercédem labóris ego reddam vobis.\end{Resp}
\begin{Resp}\Vsign Veníte benedícti Patris mei, percípite regnum.\end{Resp}
\begin{Resp}\Rsign Mercédem labóris ego reddam vobis.\end{Resp}
\begin{Resp}\Vsign Glória Patri, et Fílio, \Star{} et Spirítui Sancto.\end{Resp}
\begin{Resp}\Rsign Mercédem labóris ego reddam vobis.\end{Resp}
\switchcolumn
\EN
\begin{Resp}\Rsign O ye My Saints, who, being in the flesh, didst have striving \Star{} I will render unto you a reward of your labours.\end{Resp}
\begin{Resp}\Vsign Come, ye blessed of My Father, inherit the kingdom\end{Resp}
\begin{Resp}\Rsign I will render unto you a reward of your labours.\end{Resp}
\begin{Resp}\Vsign Glory be to the Father, and to the Son, \Star{} and to the Holy Ghost.\end{Resp}
\begin{Resp}\Rsign I will render unto you a reward of your labours.\end{Resp}
\switchcolumn*

\LA
\LessonHead{Lectio VII}
\switchcolumn
\EN
\LessonHead{Lesson VII}
\switchcolumn*

\LA
\LessonTitle{Léctio sancti Evangélii secúndum Matthǽum}
\switchcolumn
\EN
\LessonTitle{From the Holy Gospel according to Matthew}
\switchcolumn*

\LA
\Cite{Matt 5:1-12}
\switchcolumn
\EN
\Cite{Matt 5:1-12}
\switchcolumn*

\LA
\noindent In illo témpore: Videns Jesus turbas, ascéndit in montem, et cum sedísset, accessérunt ad eum discípuli ejus. Et réliqua.\par
\switchcolumn
\EN
\noindent At that time: Jesus, seeing the multitudes, he went up into a mountain, and when he was set down, his disciples came unto him. And so on.\par
\switchcolumn*

\LA
\LessonTitle{Homilía sancti Augustíni Epíscopi}
\switchcolumn
\EN
\LessonTitle{Homily by St. Augustine, Bishop (of Hippo.)}
\switchcolumn*

\LA
\Source{Lib. 1 de Sermone Domini in monte, sub initium}
\switchcolumn
\EN
\Source{Bk. I on the Lord's Sermon}
\switchcolumn*

\LA
\noindent Si quǽritur quid signíficet mons, bene intellégitur significáre majóra præcépta justítiæ, quia minóra erant quæ Judǽis data sunt. Unus tamen Deus per sanctos prophétas et fámulos suos, secúndum ordinatíssimam distributiónem témporum, dedit minóra præcépta pópulo, quem adhuc timóre alligári oportébat: et per Fílium suum majóra pópulo, quem caritáte jam liberári convénerat. Cum autem minóra minóribus, majóra majóribus dantur, ab eo dantur, qui solus novit congruéntem suis tempóribus géneri humáno exhibére medicínam.\par
\switchcolumn
\EN
\noindent If it be asked what is signified by the mountain, the said mountain may well be understood to figure the higher and greater commandments of righteousness, since those that have been given to the Jews are the lesser. The one God, in an excellent order of times, gave, by His holy Prophets and servants, His lesser commandments unto the people whom it still behoved to be bound by fear, but by His Son He gave the greater unto the people whom it now beseemed to set free by love. But whether it be the lesser to the lesser, or the greater to the greater, all are alike the gift of Him Who alone knoweth what is in each epoch the seasonable medicine of mankind.\par
\switchcolumn*

\LA
\begin{Resp}\Rsign Sint lumbi vestri præcíncti, et lucérnæ ardéntes in mánibus vestris: \Star{} Et vos símiles homínibus exspectántibus dóminum suum, quando revertátur a núptiis.\end{Resp}
\begin{Resp}\Vsign Vigiláte ergo, quia nescítis qua hora Dóminus vester ventúrus sit.\end{Resp}
\begin{Resp}\Rsign Et vos símiles homínibus exspectántibus dóminum suum, quando revertátur a núptiis.\end{Resp}
\switchcolumn
\EN
\begin{Resp}\Rsign Let your loins be girded about, and your lights burning; \Star{} And ye yourselves like unto men that wait for their lord, when he will return from the wedding.\end{Resp}
\begin{Resp}\Vsign Watch therefore, for ye know not what hour your Lord doth come.\end{Resp}
\begin{Resp}\Rsign And ye yourselves like unto men that wait for their lord, when he will return from the wedding.\end{Resp}
\switchcolumn*

\LA
\LessonHead{Lectio VIII}
\switchcolumn
\EN
\LessonHead{Lesson VIII}
\switchcolumn*

\LA
\noindent Nec mirum est, quod dantur præcépta majóra propter regnum cælórum, et minóra data sunt propter regnum terrénum, ab eódem uno Deo, qui fecit cælum et terram. De hac ergo justítia, quæ major est, per prophétam dícitur: Justítia tua sicut montes Dei: et hoc bene signíficat, quod ab uno magístro, solo docéndis tantis rebus idóneo, docétur in monte. Sedens autem docet, quod pértinet ad dignitátem magistérii. Et accédunt ad eum discípuli ejus, ut audiéndis illíus verbis hi essent étiam córpore vicinióres, qui præcéptis adimpléndis étiam ánimo propinquábant. Et apériens os suum docébat eos, dicens. Ista circumlocútio, qua scríbitur, Et apériens os suum, fortássis ipsa mora comméndat aliquánto longiórem futúrum esse sermónem. Nisi forte non vacet, quod nunc eum dictum est aperuísse os suum, quod ipse in lege véteri aperíre soléret ora prophetárum.\par
\switchcolumn
\EN
\noindent Neither is it marvel that the greater commandments be given touching the kingdom of heaven, and the lesser touching a commonwealth upon earth, since both are alike the gifts of that one God Who is the Maker alike of heaven and of earth. The higher and greater righteousness, then, is that whereof the Prophet saith: Thy righteousness is like the mountains of God. (Ps. xxxv. 7.) Thus is that Teacher, Who alone can give such teaching, mystically represented as teaching upon a mountain. And when He was set. The attitude of sitting while teaching appertaineth to the majesty of His instruction. His disciples came unto Him nearer in the body, to hear those precepts, by the fulfillment of which they should be nearer in spirit. And He opened His Mouth, and taught them, saying these words appear redundant to the sense. It may possibly be that this more pompous introduction is adopted on account of the exceptional length of the discourse to follow. But it may also be that these words are not really redundant, but the pointed declaration that He now opened His Own Mouth, Who, under the Old Law, had been used to open the mouths of the Prophets.\par
\switchcolumn*

\LA
\begin{Resp}\Rsign Média nocte clamor factus est: \Star{} Ecce sponsus venit, exíte óbviam ei.\end{Resp}
\begin{Resp}\Vsign Prudéntes vírgines, aptáte vestras lámpades.\end{Resp}
\begin{Resp}\Rsign Ecce sponsus venit, exíte óbviam ei.\end{Resp}
\begin{Resp}\Vsign Glória Patri, et Fílio, \Star{} et Spirítui Sancto.\end{Resp}
\begin{Resp}\Rsign Ecce sponsus venit, exíte óbviam ei.\end{Resp}
\switchcolumn
\EN
\begin{Resp}\Rsign At midnight there was a cry made: \Star{} Behold, the Bridegroom cometh, go ye out to meet him.\end{Resp}
\begin{Resp}\Vsign Trim your lamps, O ye wise virgins.\end{Resp}
\begin{Resp}\Rsign Behold, the Bridegroom cometh, go ye out to meet him.\end{Resp}
\begin{Resp}\Vsign Glory be to the Father, and to the Son, \Star{} and to the Holy Ghost.\end{Resp}
\begin{Resp}\Rsign Behold, the Bridegroom cometh, go ye out to meet him.\end{Resp}
\switchcolumn*

\LA
\LessonHead{Lectio IX}
\switchcolumn
\EN
\LessonHead{Lesson IX}
\switchcolumn*

\LA
\noindent Quid ergo dicit? Beáti páuperes spíritu: quóniam ipsórum est regnum cælórum. Légimus scriptum de appetitióne rerum temporálium: Omnia vánitas, et præsúmptio spíritus. Præsúmptio autem spíritus, audáciam et supérbiam signíficat. Vulgo étiam magnos spíritus supérbi habére dicúntur: et recte, quandóquidem spíritus étiam ventus vocátur. Unde scriptum est: Ignis, grando, nix, glácies, spíritus procellárum. Quis vero nésciat supérbos inflátos dici, tamquam vento disténtos? Unde est étiam illud Apóstoli: Sciéntia inflat, cáritas vero ædíficat. Quaprópter recte hic intellegúntur páuperes spíritu, húmiles et timéntes Deum, id est, non habéntes inflántem spíritum.\par
\switchcolumn
\EN
\noindent And now, what saith He? Blessed are the poor in spirit, for theirs is the kingdom of heaven. We have read where it is written concerning the lusting after temporal things: The wandering of the desire is vanity and presumption of spirit. (Eccl. vi. 9.) Presumption of spirit signifieth rashness and pride. We are used to say of proud people that they are men of high spirit, and we say well, since spirit is only one of the Latin names for wind. It is so used, for instance, in Ps. cxlviii. 8, Fire, hail, snow, ice, stormy wind. Who hath not heard the proud spoken of as puffed up, as if they were blown out with wind? Hence, alas, the Apostle saith: Knowledge puffeth up, but charity edifieth. (1 Cor. viii. 1.) By the poor in spirit, who are here called blessed, are rightly to be understood such as are lowly and fear God, that is, have not got minds puffed up with windy vanity.\par
\switchcolumn*

\LA
\TeDeum{Te Deum}
\switchcolumn
\EN
\TeDeum{Te Deum}
\switchcolumn*

\end{paracol}

\par\needspace{10\baselineskip}\addvspace{1.2em}{\centering\small\textsc{Die 2 Novembris} · \textsc{2 November}\par}
\Office{In Commemoratione Omnium Fidelium Defunctorum}{In Commemoratione Omnium Fidelium Defunctorum}{Duplex}
\addcontentsline{toc}{subsection}{Die 2 Novembris · In Commemoratione Omnium Fidelium Defunctorum \textit{— In Commemoratione Omnium Fidelium Defunctorum}}
\begin{paracol}{2}
\LA
\LessonHead{Lectio I}
\switchcolumn
\EN
\LessonHead{Lesson I}
\switchcolumn*

\LA
\Cite{Job 7:16-21}
\switchcolumn
\EN
\Cite{Job 7:16-21}
\switchcolumn*

\LA
\noindent Parce mihi, Dómine; nihil enim sunt dies mei. Quid est homo, quia magníficas eum? aut quid appónis erga eum cor tuum? Vísitas eum dilúculo, et súbito probas illum. Usquequo non parcis mihi, nec dimíttis me, ut glútiam salívam meam? Peccávi, quid fáciam tibi, o custos hóminum? quare posuísti me contrárium tibi, et factus sum mihimetípsi gravis? Cur non tollis peccátum meum, et quare non aufers iniquitátem meam? Ecce nunc in púlvere dórmiam: et, si mane me quæsíeris, non subsístam.\par
\switchcolumn
\EN
\noindent Spare me, for my days are nothing. What is a man that thou shouldst magnify him? or why dost thou set thy heart upon him? Thou visitest him early in the morning, and thou provest him suddenly. How long wilt thou not spare me, nor suffer me to swallow down my spittle? I have sinned: what shall I do to thee, O keeper of men? why hast thou set me opposite to thee, and I am become burdensome to myself? Why dost thou not remove my sin, and why dost thou not take away my iniquity? Behold now I shall sleep in the dust: and if thou seek me in the morning, I shall not be.\par
\switchcolumn*

\LA
\begin{Resp}\Rsign Credo quod Redémptor meus vivit, et in novíssimo die de terra surrectúrus sum, \Star{} Et in carne mea vidébo Deum Salvatórem meum.\end{Resp}
\begin{Resp}\Vsign Quem visúrus sum ego ipse, et non álius; et óculi mei conspectúri sunt.\end{Resp}
\begin{Resp}\Rsign Et in carne mea vidébo Deum Salvatórem meum.\end{Resp}
\switchcolumn
\EN
\begin{Resp}\Rsign I believe that my Redeemer liveth, and that I shall stand up from the earth at the latter day, \Star{} And in my flesh shall I see God my Saviour.\end{Resp}
\begin{Resp}\Vsign Whom I shall see for myself, and mine eyes shall behold, and not another.\end{Resp}
\begin{Resp}\Rsign And in my flesh shall I see God my Saviour.\end{Resp}
\switchcolumn*

\LA
\LessonHead{Lectio II}
\switchcolumn
\EN
\LessonHead{Lesson II}
\switchcolumn*

\LA
\Cite{Job 14:1-6}
\switchcolumn
\EN
\Cite{Job 14:1-6}
\switchcolumn*

\LA
\noindent Homo natus de mulíere, brevi vivens témpore, replétur multis misériis. Qui quasi flos egréditur et contéritur, et fugit velut umbra, et nunquam in eódem statu pérmanet. Et dignum ducis super hujuscémodi aperíre óculos tuos, et addúcere eum tecum in judícium? Quis potest fácere mundum de immúndo concéptum sémine? Nonne tu qui solus es? Breves dies hóminis sunt: númerus ménsium ejus apud te est: constituísti términos ejus, qui præteríri non póterunt. Recéde páululum ab eo, ut quiéscat, donec optáta véniat, sicut mercenárii, dies ejus.\par
\switchcolumn
\EN
\noindent Man born of a woman, living for a short time, is filled with many miseries. Who cometh forth like a flower, and is destroyed, and fleeth as a shadow, and never continueth in the same state. And dost thou think it meet to open thy eyes upon such an one, and to bring him into judgment with thee? Who can make him clean that is conceived of unclean seed? is it not thou who only art? The days of man are short, and the number of his months is with thee: thou hast appointed his bounds which cannot be passed. Depart a little from him, that he may rest, until his wished for day come, as that of the hireling\par
\switchcolumn*

\LA
\begin{Resp}\Rsign Qui Lázarum resuscitásti a monuménto fœ́tidum, \Star{} Tu eis, Dómine, dona réquiem, et locum indulgéntiæ.\end{Resp}
\begin{Resp}\Vsign Qui ventúrus es judicáre vivos et mórtuos, et sǽculum per ignem.\end{Resp}
\begin{Resp}\Rsign Tu eis, Dómine, dona réquiem, et locum indulgéntiæ.\end{Resp}
\switchcolumn
\EN
\begin{Resp}\Rsign Thou Who didst call up Lazarus from the grave after that he had begun to stink \Star{} Do Thou, O Lord, grant them rest and a place of forgiveness.\end{Resp}
\begin{Resp}\Vsign Thou Who shalt come to judge the quick and dead, and the world by fire\end{Resp}
\begin{Resp}\Rsign Do Thou, O Lord, grant them rest and a place of forgiveness.\end{Resp}
\switchcolumn*

\LA
\LessonHead{Lectio III}
\switchcolumn
\EN
\LessonHead{Lesson III}
\switchcolumn*

\LA
\Cite{Job 19:20-27}
\switchcolumn
\EN
\Cite{Job 19:20-27}
\switchcolumn*

\LA
\noindent Pelli meæ, consúmptis cárnibus, adhǽsit os meum, et derelícta sunt tantúmmodo lábia circa dentes meos. Miserémini mei, miserémini mei saltem vos, amíci mei, quia manus Dómini tétigit me. Quare persequímini me sicut Deus, et cárnibus meis saturámini? Quis mihi tríbuat ut scribántur sermónes mei? quis mihi det ut exaréntur in libro stylo férreo et plumbi lámina, vel celte sculpántur in sílice? Scio enim quod redémptor meus vivit, et in novíssimo die de terra surrectúrus sum: et rursum circúmdabor pelle mea, et in carne mea vidébo Deum meum: quem visúrus sum ego ipse, et óculi mei conspectúri sunt, et non álius: repósita est hæc spes mea in sinu meo.\par
\switchcolumn
\EN
\noindent The flesh being consumed. My bone hath cleaved to my skin, and nothing but lips are left about my teeth. Have pity on me, have pity on me, at least you my friends, because the hand of the Lord hath touched me. Why do you persecute me as God, and glut yourselves with my flesh? Who will grant me that my words may be written? Who will grant me that they may be marked down in a book? With an iron pen and in a plate of lead, or else be graven with an instrument in flint stone. For I know that my Redeemer liveth, and in the last day I shall rise out of the earth. And I shall be clothed again with my skin, and in my flesh I will see my God. Whom I myself shall see, and my eyes shall behold, and not another: this my hope is laid up in my bosom.\par
\switchcolumn*

\LA
\begin{Resp}\Rsign Dómine, quando véneris judicáre terram, ubi me abscóndam a vultu iræ tuæ? \Star{} Quia peccávi nimis in vita mea.\end{Resp}
\begin{Resp}\Vsign Commíssa mea pavésco, et ante te erubésco: dum véneris judicáre, noli me condemnáre.\end{Resp}
\begin{Resp}\Rsign Quia peccávi nimis in vita mea. \$Requiem\end{Resp}
\begin{Resp}\Rsign Quia peccávi nimis in vita mea.\end{Resp}
\switchcolumn
\EN
\begin{Resp}\Rsign Lord, when Thou comest to judge the earth, where shall I hide myself from the face of thy wrath? \Star{} For I have sinned greatly in my life.\end{Resp}
\begin{Resp}\Vsign I dread my sins, I blush before thee I see the great tribunal set! In fear and terror I implore thee, Forgive when soul and Judge are met!\end{Resp}
\begin{Resp}\Rsign For I have sinned greatly in my life. \$Requiem\end{Resp}
\begin{Resp}\Rsign For I have sinned greatly in my life.\end{Resp}
\switchcolumn*

\LA
\LessonHead{Lectio IV}
\switchcolumn
\EN
\LessonHead{Lesson IV}
\switchcolumn*

\LA
\LessonTitle{Ex Libro sancti Augustíni Epíscopi de cura pro mórtuis gerénda}
\switchcolumn
\EN
\LessonTitle{From the book of St. Augustine, the Bishop, on the Care for the Deceased}
\switchcolumn*

\LA
\Source{Cap 2 et 3}
\switchcolumn
\EN
\Source{Cap. 2 \& 3}
\switchcolumn*

\LA
\noindent Curátio fúneris, condítio sepultúræ, pompa exsequiárum magis sunt vivórum solácia quam subsídia mortuórum. Nec ídeo tamen contemnénda et abiciénda sunt córpora defunctórum, maximéque justórum ac fidélium, quibus tamquam órganis et vasis ad ómnia bona ópera sancte usus est spíritus. Si enim patérna vestis et ánulus, ac si quid hujúsmodi, tanto cárius est pósteris, quanto erga paréntes major afféctus; nullo modo ipsa spernénda sunt córpora, quæ útique multo familiárius atque conjúnctius quam quǽlibet induménta gestámus. Hæc enim non ad ornaméntum vel adjutórium, quod adhibétur extrínsecus, sed ad ipsam natúram hóminis pértinent. Unde et antiquórum justórum fúnera officiósa pietáte curáta sunt, et exséquiæ celebrátæ, et sepultúra provísa; ipsíque, cum víverent, de sepeliéndis vel étiam transferéndis suis corpóribus fíliis mandavérunt.\par
\switchcolumn
\EN
\noindent If this be true, doubtless also the providing for the interment of bodies a place at the Memorials of Saints, is a mark of a good human affection towards the remains of one's friends. Yet it follows not that the bodies of the departed are to be despised and flung aside, and above all of just and faithful men, which bodies as organs and vessels to all good works their spirit has holily used. For if a father's garment and ring, and whatever such like, is the more dear to those whom they leave behind, the greater their affection is towards their parents, in no wise are the bodies themselves to be spurned, which truly we wear in more familiar and close conjunction than any of our putting on. For these pertain not to ornament or aid which is applied from without, but to the very nature of man. Whence also the funerals of the just men of old were with dutiful piety cared for, and their obsequies celebrated, and sepulture provided: and themselves while living did touching burial or even translation of their bodies give charge to their sons.\par
\switchcolumn*

\LA
\begin{Resp}\Rsign Meménto mei, Deus, quia ventus est vita mea, \Star{} Nec aspíciat me visus hóminis.\end{Resp}
\begin{Resp}\Vsign De profúndis clamávi ad te, Dómine: Dómine, exáudi vocem meam.\end{Resp}
\begin{Resp}\Rsign Nec aspíciat me visus hóminis.\end{Resp}
\switchcolumn
\EN
\begin{Resp}\Rsign Remember, O God, that my life is wind. \Star{} The eye of him that hath seen me shall see me no more.\end{Resp}
\begin{Resp}\Vsign Out of the depths have I cried unto thee, O Lord! Lord, hear my voice.\end{Resp}
\begin{Resp}\Rsign The eye of him that hath seen me shall see me no more.\end{Resp}
\switchcolumn*

\LA
\LessonHead{Lectio V}
\switchcolumn
\EN
\LessonHead{Lesson V}
\switchcolumn*

\LA
\Source{Cap. 4}
\switchcolumn
\EN
\Source{Cap. 4}
\switchcolumn*

\LA
\noindent Recordántis et precántis afféctus cum defúnctis a fidélibus caríssimis exhibétur, eum prodésse non dúbium est iis, qui cum in córpore víverent, tália sibi post hanc vitam prodésse meruérunt. Verum, etsi áliqua necéssitas vel humári córpora, vel in sacris locis humári nulla data facultáte permíttat, non sunt prætermitténdæ supplicatiónes pro spirítibus mortuórum: quas faciéndas pro ómnibus in christiána et cathólica societáte defúnctis, étiam tácitis eórum nomínibus, sub generáli commemoratióne suscépit Ecclésia; ut quibus ad ista desunt paréntes, aut fílii, aut quicúmque cognáti vel amíci, ab una eis exhibeántur pia matre commúni. Si autem deéssent istæ supplicatiónes, quæ fiunt recta fide ac pietáte pro mórtuis, puto quod nihil prodésset spirítibus eórum, quámlibet in locis sanctis exánima córpora poneréntur.\par
\switchcolumn
\EN
\noindent And when this affection is exhibited to the departed by faithful men who were most dear to them, there is no doubt that it profits them who while living in the body merited that such things should profit them after this life. But even if some necessity should through absence of all facility not allow bodies to be interred, or in such places interred, yet should there be no pretermitting of supplications for the spirits of the dead: which supplications, that they should be made for all in Christian and catholic fellowship departed, even without mentioning of their names, under a general commemoration, the Church has charged herself withal; to the intent that they which lack, for these offices, parents or sons or whatever kindred or friends, may have the same afforded unto them by the one pious mother which is common to all. But if there were lack of these supplications, which are made with right faith and piety for the dead, I account that it should not a whit profit their spirits, howsoever in holy places the lifeless bodies should be deposited.\par
\switchcolumn*

\LA
\begin{Resp}\Rsign Hei mihi, Dómine, quia peccávi nimis in vita mea: Quid fáciam, miser? ubi fúgiam, nisi ad te, Deus meus? \Star{} Miserére mei, dum véneris in novíssimo die.\end{Resp}
\begin{Resp}\Vsign Anima mea turbáta est valde, sed tu, Dómine, succúrre ei.\end{Resp}
\begin{Resp}\Rsign Miserére mei, dum véneris in novíssimo die.\end{Resp}
\switchcolumn
\EN
\begin{Resp}\Rsign Woe is me, O Lord! for I have sinned greatly in my life. I am smitten what shall I do? Whither shall I flee but unto thee, O my God? \Star{} Have mercy upon me, when Thou comest at the latter day.\end{Resp}
\begin{Resp}\Vsign My soul is sore vexed, but Thou, O Lord, help me.\end{Resp}
\begin{Resp}\Rsign Have mercy upon me, when Thou comest at the latter day.\end{Resp}
\switchcolumn*

\LA
\LessonHead{Lectio VI}
\switchcolumn
\EN
\LessonHead{Lesson VI}
\switchcolumn*

\LA
\Source{Cap. 18}
\switchcolumn
\EN
\Source{Cap. 18}
\switchcolumn*

\LA
\noindent Quæ cum ita sint, non existimémus ad mórtuos, pro quibus curam gérimus, perveníre, nisi quod pro eis sive altáris, sive oratiónum, sive eleemosynárum sacrifíciis solémniter supplicámus: quamvis non pro quibus fiunt, ómnibus prosint; sed iis tantum pro quibus, dum vivunt, comparátur, ut prosint. Sed quia non discérnimus qui sunt, opórtet ea pro regenerátis ómnibus fácere, ut nullus eórum prætermittátur, ad quos hæc benefícia possint et débeant perveníre. Mélius enim supérerunt ista eis, quibus nec obsunt nec prosunt; quam eis déerunt, quibus prosunt. Diligéntius tamen facit hæc quisque pro necessáriis suis, quo pro illo fiat simíliter a suis. Córpori autem humándo quidquid impénditur, non est præsídium salútis, sed humanitátis offícium, secúndum afféctum quo nemo unquam carnem suam ódio habet. Unde opórtet ut quam potest pro carne próximi curam gerat, cum ille inde recésserit, qui gerébat. Et si hæc fáciunt, qui carnis resurrectiónem non credunt, quanto magis debent fácere qui credunt; ut córpori mórtuo, sed tamen resurrectúro et in æternitáte mansúro, impénsum ejúsmodi offícium sit étiam quodámmodo ejúsdem fídei testimónium!\par
\switchcolumn
\EN
\noindent Which things being so, let us not think that to the dead for whom we have a care, any thing reaches save what by sacrifices either of the altar, or of prayers, or of alms, we solemnly supplicate: although not to all for whom they are done be they profitable, but to them only by whom while they live it is obtained that they should be profitable. But forasmuch as we discern not who these be, it is meet to do them for all regenerate persons, that none of them may be passed by to whom these benefits may and ought to reach. For better it is that these things shall be superfluously done to them whom they neither hinder nor help, than lacking to them whom they help. More diligently however does each man these things for his own near and dear friends, in order that they may be likewise done unto him by his. But as for the burying of the body, whatever is bestowed on that, is no aid of salvation, but an office of humanity, according to that affection by which no man ever hates his own flesh. Whence it is fitting that he take what care he is able for the flesh of his neighbor, when he is gone that bare it. And if they do these things who believe not the resurrection of the flesh, how much more are they beholden to do the same who do believe; that so, an office of this kind bestowed upon a body, dead but yet to rise again and to remain to eternity, may also be in some sort a testimony of the same faith?\par
\switchcolumn*

\LA
\begin{Resp}\Rsign Ne recordéris peccáta mea, Dómine, \Star{} Dum véneris judicáre sǽculum per ignem.\end{Resp}
\begin{Resp}\Vsign Dírige, Dómine, Deus meus, in conspéctu tuo viam meam.\end{Resp}
\begin{Resp}\Rsign Dum véneris judicáre sǽculum per ignem. \$Requiem\end{Resp}
\begin{Resp}\Rsign Dum véneris judicáre sǽculum per ignem.\end{Resp}
\switchcolumn
\EN
\begin{Resp}\Rsign Hold not my sins in remembrance, O Lord, \Star{} When thou comest to judge the world by fire.\end{Resp}
\begin{Resp}\Vsign Make my way straight before thy face, O Lord my God.\end{Resp}
\begin{Resp}\Rsign When thou comest to judge the world by fire. \$Requiem\end{Resp}
\begin{Resp}\Rsign When thou comest to judge the world by fire.\end{Resp}
\switchcolumn*

\LA
\LessonHead{Lectio VII}
\switchcolumn
\EN
\LessonHead{Lesson VII}
\switchcolumn*

\LA
\LessonTitle{De Epístola prima beáti Pauli Apóstoli ad Corínthios}
\switchcolumn
\EN
\LessonTitle{From the first letter of St. Paul to the Corinthians}
\switchcolumn*

\LA
\Cite{1 Cor 15:12-22}
\switchcolumn
\EN
\Cite{1 Cor 15:12-22}
\switchcolumn*

\LA
\noindent Si Christus prædicátur quod resurréxit a mórtuis, quómodo quidam dicunt in vobis, quóniam resurréctio mortuórum non est? Si autem resurréctio mortuórum non est: neque Christus resurréxit. Si autem Christus non resurréxit, inánis est ergo prædicátio nostra, inánis est et fides vestra: invenímur autem et falsi testes Dei: quóniam testimónium díximus advérsus Deum quod suscitáverit Christum, quem non suscitávit, si mórtui non resúrgunt. Nam si mórtui non resúrgunt, neque Christus resurréxit. Quod si Christus non resurréxit, vana est fides vestra: adhuc enim estis in peccátis vestris. Ergo et qui dormiérunt in Christo, periérunt. Si in hac vita tantum in Christo sperántes sumus, miserabilióres sumus ómnibus homínibus. Nunc autem Christus resurréxit a mórtuis primítiæ dormiéntium, quóniam quidem per hóminem mors, et per hóminem resurréctio mortuórum. Et sicut in Adam omnes moriúntur, ita et in Christo omnes vivificabúntur.\par
\switchcolumn
\EN
\noindent Now if Christ be preached, that he arose again from the dead, how do some among you say, that there is no resurrection of the dead? But if there be no resurrection of the dead, then Christ is not risen again. And if Christ be not risen again, then is our preaching vain, and your faith is also vain. Yea, and we are found false witnesses of God: because we have given testimony against God, that he hath raised up Christ; whom he hath not raised up, if the dead rise not again. For if the dead rise not again, neither is Christ risen again. And if Christ be not risen again, your faith is vain, for you are yet in your sins. Then they also that are fallen asleep in Christ, are perished. If in this life only we have hope in Christ, we are of all men most miserable. But now Christ is risen from the dead, the firstfruits of them that sleep For by a man came death, and by a man the resurrection of the dead. And as in Adam all die, so also in Christ all shall be made alive.\par
\switchcolumn*

\LA
\begin{Resp}\Rsign Peccántem me cotídie, et non me pœniténtem, timor mortis contúrbat me: \Star{} Quia in inférno nulla est redémptio, miserére mei, Deus, et salva me.\end{Resp}
\begin{Resp}\Vsign Deus, in nómine tuo salvum me fac, et in virtúte tua líbera me.\end{Resp}
\begin{Resp}\Rsign Quia in inférno nulla est redémptio, miserére mei, Deus, et salva me.\end{Resp}
\switchcolumn
\EN
\begin{Resp}\Rsign Forasmuch as I sin daily, and repent not, the fear of death troubleth me. \Star{} O God, have mercy upon me, and save me, for in hell there is no redemption.\end{Resp}
\begin{Resp}\Vsign Save me, O God, by thy name, and judge me in thy strength.\end{Resp}
\begin{Resp}\Rsign O God, have mercy upon me, and save me, for in hell there is no redemption.\end{Resp}
\switchcolumn*

\LA
\LessonHead{Lectio VIII}
\switchcolumn
\EN
\LessonHead{Lesson VIII}
\switchcolumn*

\LA
\Cite{1 Cor 15:35-44}
\switchcolumn
\EN
\Cite{1 Cor 15:35-44}
\switchcolumn*

\LA
\noindent Sed dicet áliquis: Quómodo resúrgunt mórtui? qualíve córpore vénient? Insípiens, tu quod séminas non vivificátur, nisi prius moriátur: et quod séminas, non corpus, quod futúrum est, séminas, sed nudum granum, ut puta trítici, aut alicújus ceterórum. Deus autem dat illi corpus sicut vult, et unicuíque séminum próprium corpus. Non omnis caro, éadem caro: sed ália quidem hóminum, ália vero pécorum, ália vólucrum, ália autem píscium. Et córpora cæléstia, et córpora terréstria: sed ália quidem cæléstium glória, ália autem terréstrium. Alia cláritas solis, ália cláritas lunæ, et ália cláritas stellárum; stella enim a stella differt in claritáte: sic et resurréctio mortuórum. Seminátur in corruptióne, surget in incorruptióne. Seminátur in ignobilitáte, surget in glória: seminátur in infirmitáte, surget in virtúte: seminátur corpus animále, surget corpus spiritále.\par
\switchcolumn
\EN
\noindent But some man will say: How do the dead rise again? or with what manner of body shall they come? Senseless man, that which thou sowest is not quickened, except it die first. And that which thou sowest, thou sowest not the body that shall be; but bare grain, as of wheat, or of some of the rest. But God giveth it a body as he will: and to every seed its proper body. All flesh is not the same flesh: but one is the flesh of men, another of beasts, another of birds, another of fishes. And there are bodies celestial, and bodies terrestrial: but, one is the glory of the celestial, and another of the terrestrial. One is the glory of the sun, another the glory of the moon, and another the glory of the stars. For star differeth from star in glory. So also is the resurrection of the dead. It is sown in corruption, it shall rise in incorruption. It is sown in dishonour, it shall rise in glory. It is sown in weakness, it shall rise in power. It is sown a natural body, it shall rise a spiritual body.\par
\switchcolumn*

\LA
\begin{Resp}\Rsign Dómine, secúndum actum meum noli me judicáre: nihil dignum in conspéctu tuo egi; ídeo déprecor majestátem tuam, \Star{} Ut tu, Deus, déleas iniquitátem meam.\end{Resp}
\begin{Resp}\Vsign Amplius lava me, Dómine, ab injustítia mea, et a delícto meo munda me.\end{Resp}
\begin{Resp}\Rsign Ut tu, Deus, déleas iniquitátem meam.\end{Resp}
\switchcolumn
\EN
\begin{Resp}\Rsign O Lord, judge me not according to my works; for I have done nothing that can be counted in respect of thee. I beseech thy majesty therefore, that thou wouldest \Star{} Blot out my transgressions, O God.\end{Resp}
\begin{Resp}\Vsign Lord, wash me thoroughly from mine iniquity and cleanse me from my sin.\end{Resp}
\begin{Resp}\Rsign Blot out my transgressions, O God.\end{Resp}
\switchcolumn*

\LA
\LessonHead{Lectio IX}
\switchcolumn
\EN
\LessonHead{Lesson IX}
\switchcolumn*

\LA
\Cite{1 Cor 15:51-58}
\switchcolumn
\EN
\Cite{1 Cor 15:51-58}
\switchcolumn*

\LA
\noindent Ecce mystérium vobis dico: Omnes quidem resurgémus, sed non omnes immutábimur. In moménto, in ictu óculi, in novíssima tuba: canet enim tuba, et mórtui resúrgent incorrúpti: et nos immutábimur. Opórtet enim corruptíbile hoc indúere incorruptiónem: et mortále hoc indúere immortalitátem. Cum autem mortále hoc indúerit immortalitátem, tunc fiet sermo, qui scriptus est: Absórpta est mors in victória. Ubi est mors victória tua? ubi est mors stímulus tuus? Stímulus autem mortis peccátum est: virtus vero peccáti lex. Deo autem grátias, qui dedit nobis victóriam per Dóminum nostrum Jesum Christum. Itaque fratres mei dilécti, stábiles estóte, et immóbiles: abundántes in ópere Dómini semper, sciéntes quod labor vester non est inánis in Dómino.\par
\switchcolumn
\EN
\noindent Behold, I tell you a mystery. We shall all indeed rise again: but we shall not all be changed. In a moment, in the twinkling of an eye, at the last trumpet: for the trumpet shall sound, and the dead shall rise again incorruptible: and we shall be changed. For this corruptible must put on incorruption; and this mortal must put on immortality. And when this mortal hath put on immortality, then shall come to pass the saying that is written: Death is swallowed up in victory. O death, where is thy victory? O death, where is thy sting? Now the sting of death is sin: and the power of sin is the law. But thanks be to God, who hath given us the victory through our Lord Jesus Christ. Therefore, my beloved brethren, be ye steadfast and unmoveable; always abounding in the work of the Lord, knowing that your labour is not in vain in the Lord.\par
\switchcolumn*

\LA
\begin{Resp}\Rsign Líbera me, Dómine, de morte ætérna in die illa treménda, \Star{} Quando cæli movéndi sunt et terra, * Dum véneris judicáre sǽculum per ignem.\end{Resp}
\begin{Resp}\Vsign Tremens factus sum ego et tímeo, dum discússio vénerit atque ventúra ira.\end{Resp}
\begin{Resp}\Rsign Quando cæli movéndi sunt et terra.\end{Resp}
\begin{Resp}\Vsign Dies illa, dies iræ, calamitátis et misériæ, dies magna et amára valde.\end{Resp}
\begin{Resp}\Rsign Dum véneris judicáre sǽculum per ignem. \$Requiem\end{Resp}
\begin{Resp}\Rsign Líbera me, Dómine, de morte ætérna in die illa treménda, * Quando cæli movéndi sunt et terra, * Dum véneris judicáre sǽculum per ignem.\end{Resp}
\switchcolumn
\EN
\begin{Resp}\Rsign Deliver me, O Lord, from eternal death in that awful day \Star{} when the heavens and the earth shall be shaken, and Thou shalt come to judge the world by fire.\end{Resp}
\begin{Resp}\Vsign Quaking and dread take hold upon me, when I look for the coming of the trial and the wrath to come.\end{Resp}
\begin{Resp}\Rsign When the heavens and the earth shall be shaken.\end{Resp}
\begin{Resp}\Vsign That day is a day of wrath, of wasteness and desolation, a great day and exceeding bitter.\end{Resp}
\begin{Resp}\Rsign When Thou shalt come to judge the world by fire. \$Requiem\end{Resp}
\begin{Resp}\Rsign Deliver me, O Lord, from eternal death in that awful day * when the heavens and the earth shall be shaken, and Thou shalt come to judge the world by fire.\end{Resp}
\switchcolumn*

\end{paracol}

\par\needspace{10\baselineskip}\addvspace{1.2em}{\centering\small\textsc{Die 3 Novembris} · \textsc{3 November}\par}
\Office{Tertia die infra Octavam Omnium Sanctorum}{Third Day Within the Octave of All Saints}{Semiduplex}
\addcontentsline{toc}{subsection}{Die 3 Novembris · Tertia die infra Octavam Omnium Sanctorum \textit{— Third Day Within the Octave of All Saints}}
\begin{paracol}{2}
\LA
\RefLine{Lectiones I–III: ex Commúni Evangelistarum (C1a).}
\switchcolumn
\EN
\RefLine{Lessons I–III: from the Commune Evangelistarum (C1a).}
\switchcolumn*

\LA
\LessonHead{Lectio IV}
\switchcolumn
\EN
\LessonHead{Lesson IV}
\switchcolumn*

\LA
\LessonTitle{De Sermóne sancti Bedæ Venerábilis Presbýteri}
\switchcolumn
\EN
\LessonTitle{The Lesson is taken from the Sermons of St. Venerable Bede, the Priest}
\switchcolumn*

\LA
\Source{Sermo 18 de Sanctis.}
\switchcolumn
\EN
\Source{Ex Sermone 18. de Sanctis.}
\switchcolumn*

\LA
\noindent Nulla erit tunc usquam discordia, sed cuncta consona, cuncta convenientia, quia omnium erit Sanctorum una concordia, pax cuncta et lætitia continet. Tranquilla sunt omnia et quieta: jugis splendor, non iste qui nunc est, sed tanto clarior, quanto felicior: quia civitas, ut legitur, illa non egebit lumine solis, sed Dominus omnipotens illuminabit eam, et lucerna ejus est Agnus. Ibi Sancti fulgebunt ut stellæ in perpetuas æternitates, et sicut splendor firmamenti, qui erudiunt multos.\par
\switchcolumn
\EN
\noindent Never shall there be discord anywhere there, but all things in harmony. For everywhere there, things are in such concord that all the Saints are at unity with each other in one peace and joy. Everywhere there, all things are tranquil and quiet. Perpetual is the splendour there; not like unto the sunlight which we know here, but a light which is the brighter, as it is the more blessed. For that city, as saith Scripture, needeth not the light of the sun, because the Lord Almighty doth enlighten it by the Lamb which is the Light thereof. There the Saints shall shine like as the brightness of the firmament, and they that have turned many to righteousness, as the stars, for ever and ever.\par
\switchcolumn*

\LA
\begin{Resp}\Rsign Præcúrsor Dómini venit, de quo ipse testátur: \Star{} Nullus major inter natos mulíerum Joánne Baptísta.\end{Resp}
\begin{Resp}\Vsign Hic est enim prophéta, et plus quam prophéta, de quo Salvátor ait.\end{Resp}
\begin{Resp}\Rsign Nullus major inter natos mulíerum Joánne Baptísta.\end{Resp}
\switchcolumn
\EN
\begin{Resp}\Rsign The Fore-runner of the Lord cometh, to whom He Himself bare witness, saying: \Star{} Among them that are born of women there hath not risen a greater than John the Baptist.\end{Resp}
\begin{Resp}\Vsign A Prophet? Yea, and much more than a Prophet. This is he of whom the Saviour saith:\end{Resp}
\begin{Resp}\Rsign Among them that are born of women there hath not risen a greater than John the Baptist.\end{Resp}
\switchcolumn*

\LA
\LessonHead{Lectio V}
\switchcolumn
\EN
\LessonHead{Lesson V}
\switchcolumn*

\LA
\noindent Quapropter nox ibi nulla, nullæ tenebræ, concursus nubium nullus, nec frigoris aut ardoris asperitas ulla: sed talis quædam erit rerum temperies, qualem nec oculus vidit, nec auris audivit, nec in cor hominis ascendit, nisi illorum qui ea perfrui digni inveniuntur, quorum nomina scripta sunt in libro vitæ, qui et laverunt stolas suas in sanguine Agni, et sunt ante sedem Dei, serviuntque ei die ac nocte. Non est senectus ibi, nec senectutis miseria, dum omnes occurrunt in virum perfectum, in mensuram ætatis plenitudinis Christi.\par
\switchcolumn
\EN
\noindent And so there is no night there, no darkness, no gathering of clouds, no asperity of heat or cold. But such is the nature of things there as no eye hath seen, nor ear heard, neither hath it entered into the heart of man, except of those only who have been found worthy to enjoy it, whose names are written in the book of life; and who have washed their robes in the blood of the Lamb, and are before the throne of God, serving him day and night. There is no old age anywhere there, nor misery of old age, for all are come to perfect manhood, to the measure of the stature of the fulness of Christ.\par
\switchcolumn*

\LA
\begin{Resp}\Rsign Isti sunt qui vivéntes in carne, plantavérunt Ecclésiam sánguine suo: \Star{} Cálicem Dómini bibérunt, et amíci Dei facti sunt.\end{Resp}
\begin{Resp}\Vsign In omnem terram exívit sonus eórum, et in fines orbis terræ verba eórum.\end{Resp}
\begin{Resp}\Rsign Cálicem Dómini bibérunt, et amíci Dei facti sunt.\end{Resp}
\switchcolumn
\EN
\begin{Resp}\Rsign These are they who while yet they lived in the flesh, planted the Church in their own blood; \Star{} They drank of the Lord's cup, and became the friends of God.\end{Resp}
\begin{Resp}\Vsign Their sound is gone out through all the earth, and their words to the ends of the world.\end{Resp}
\begin{Resp}\Rsign They drank of the Lord's cup, and became the friends of God.\end{Resp}
\switchcolumn*

\LA
\LessonHead{Lectio VI}
\switchcolumn
\EN
\LessonHead{Lesson VI}
\switchcolumn*

\LA
\noindent Verum super hæc omnia est, consociari Angelorum et Archangelorum cœtibus: Thronis étiam et Dominationibus, Principatibus et Potestatibus, omniumque cælestium supernarum Virtutum contuberniis perfrui, et intueri agmina Sanctorum splendidius sideribus micantia, Patriarcharum fide fulgentia, Prophetarum spe lætantia, Apostolorum in duodecim tribubus Israël orbem judicantia, Mártyrum purpureis victoriæ coronis lucentia, Virginum quoque choros candentia serta gestantes inspicere.\par
\switchcolumn
\EN
\noindent But far above all these things is the fellowship there. That is, to enjoy the companionship of the heavenly citizens: to look upon the choirs of Angels and Archangels, of Thrones and Dominions, Principalities, Powers, and all the heavenly Virtues on high: and to behold the army of the Saints shining more gloriously than the stars; of the Patriarchs glowing with faith; of the Prophets rejoicing in hope; of the Apostles judging the world reformed into twelve tribes of the new Israel; of the Martyrs resplendent in their ruddy crowns of victory; and of the Virgins wearing garlands of the purest white.\par
\switchcolumn*

\LA
\begin{Resp}\Rsign Sancti mei, qui in carne pósiti, certámen habuístis: \Star{} Mercédem labóris ego reddam vobis.\end{Resp}
\begin{Resp}\Vsign Veníte benedícti Patris mei, percípite regnum.\end{Resp}
\begin{Resp}\Rsign Mercédem labóris ego reddam vobis.\end{Resp}
\begin{Resp}\Vsign Glória Patri, et Fílio, \Star{} et Spirítui Sancto.\end{Resp}
\begin{Resp}\Rsign Mercédem labóris ego reddam vobis.\end{Resp}
\switchcolumn
\EN
\begin{Resp}\Rsign O ye My Saints, who, being in the flesh, didst have striving \Star{} I will render unto you a reward of your labours.\end{Resp}
\begin{Resp}\Vsign Come, ye blessed of My Father, inherit the kingdom\end{Resp}
\begin{Resp}\Rsign I will render unto you a reward of your labours.\end{Resp}
\begin{Resp}\Vsign Glory be to the Father, and to the Son, \Star{} and to the Holy Ghost.\end{Resp}
\begin{Resp}\Rsign I will render unto you a reward of your labours.\end{Resp}
\switchcolumn*

\LA
\LessonHead{Lectio VII}
\switchcolumn
\EN
\LessonHead{Lesson VII}
\switchcolumn*

\LA
\LessonTitle{Léctio sancti Evangélii secúndum Matthǽum.}
\switchcolumn
\EN
\LessonTitle{From the Holy Gospel according to Matthew}
\switchcolumn*

\LA
\Cite{Matt 5:1-12}
\switchcolumn
\EN
\Cite{Matt 5:1-12}
\switchcolumn*

\LA
\noindent In illo témpore: Videns Jesus turbas, ascendit in montem, et cum sedisset, accesserunt ad eum discipuli ejus. Et réliqua.\par
\switchcolumn
\EN
\noindent At that time, Jesus, seeing the multitudes, went up into a mountain, and when he was set down, his disciples came unto him. And so on.\par
\switchcolumn*

\LA
\LessonTitle{De Homilia S. Augustíni Epíscopi.}
\switchcolumn
\EN
\LessonTitle{Homily by St. Augustine the Bishop}
\switchcolumn*

\LA
\Source{Ex Lib. 1. de Serm. Domini in monte. Cap. 2.et 3.}
\switchcolumn
\EN
\Source{Ex Lib. 1. de Serm. Domini in monte. Cap. 2.et 3.}
\switchcolumn*

\LA
\noindent Beati mundo corde: quoniam ipsi Deum videbunt. Quam ergo stulti sunt, qui Deum istis exterioribus oculis quærunt, cum corde videatur, sicut alibi scriptum est: Et in simplicitate cordis quærite illum. Hoc est enim mundum cor, quod est simplex cor: et quemadmodum lumen hoc videri non potest, nisi oculis mundis, ita nec Deus videtur, nisi mundum sit illud quo videri potest. Beati pacifici: quoniam ipsi filii Dei vocabuntur. In pace perfectio est, ubi nihil repugnat: et ideo filii Dei pacifici, quoniam nihil in his resistit Deo, et utique filii similitudinem patris habere debent.\par
\switchcolumn
\EN
\noindent Blessed are the pure in heart, for they shall see God. How foolish therefore be they that seek God with their outward eyes, since it is in the heart that he is seen. Thus it is written elsewhere: In simplicity of heart seek him. For a pure heart is one that hath the simplicity of single-mindedness. And just as light can be seen only insofar as the bodily eyes have clear vision, so God can be seen only insofar as the heart is clear by reason of its single-mindedness of purpose. Blessed are the peacemakers for they shall be called the children of God. Peace is perfect where there is no strife. And the children of God are peacemakers because there is nothing in them which withstandeth God. And surely children ought to be like unto their father.\par
\switchcolumn*

\LA
\begin{Resp}\Rsign Sint lumbi vestri præcíncti, et lucérnæ ardéntes in mánibus vestris: \Star{} Et vos símiles homínibus exspectántibus dóminum suum, quando revertátur a núptiis.\end{Resp}
\begin{Resp}\Vsign Vigiláte ergo, quia nescítis qua hora Dóminus vester ventúrus sit.\end{Resp}
\begin{Resp}\Rsign Et vos símiles homínibus exspectántibus dóminum suum, quando revertátur a núptiis.\end{Resp}
\switchcolumn
\EN
\begin{Resp}\Rsign Let your loins be girded about, and your lights burning; \Star{} And ye yourselves like unto men that wait for their lord, when he will return from the wedding.\end{Resp}
\begin{Resp}\Vsign Watch therefore, for ye know not what hour your Lord doth come.\end{Resp}
\begin{Resp}\Rsign And ye yourselves like unto men that wait for their lord, when he will return from the wedding.\end{Resp}
\switchcolumn*

\LA
\LessonHead{Lectio VIII}
\switchcolumn
\EN
\LessonHead{Lesson VIII}
\switchcolumn*

\LA
\noindent Pacifici autem in semetipsis sunt, qui omnes animi sui motus componentes et subjicientes rationi, id est, menti et spiritui, carnalesque concupiscentias habentes edomitas, fiunt regnum Dei: in quo ita sunt ordinata omnia, ut id quod est in homine præcipuum et excellens, hoc imperet ceteris non reluctantibus, quæ sunt nobis bestiisque communia; atque idipsum quod excellit in homine, id est, mens et ratio, subjiciatur potiori, quod est ipsa Veritas, unigenitus Filius Dei. Neque enim imperare inferioribus potest, nisi superiori se ipse subjiciat. Et hæc est pax quæ datur in terra hominibus bonæ voluntatis: hæc vita consummati perfectique sapientis.\par
\switchcolumn
\EN
\noindent Peacemakers are all those that have peace within themselves. That is, those who have set all the passions of their souls in order. This they do by subjecting their passions to reason, namely, to the mind and spirit. Thus, by conquering their fleshly desires, they have made of themselves a little kingdom of God. In this kingdom, everything is ordered into an harmonious whole, by virtue of the fact that the element which chief and pre-eminent in man (namely, mind or reason) ruleth, without resistance, over the elements which we have in common with the beasts. And in this fashion, the element of mind or reason, which is pre-eminent in man, is made subject to something pre-eminent to itself, namely, the Truth, even the Only-Begotten Son of God. For no one is able to govern except he be subject unto the higher powers. Now, herein we have set forth the peace which is given on earth to men of good-will. This is the life of one who is thoroughly and completely wise.\par
\switchcolumn*

\LA
\begin{Resp}\Rsign Média nocte clamor factus est: \Star{} Ecce sponsus venit, exíte óbviam ei.\end{Resp}
\begin{Resp}\Vsign Prudéntes vírgines, aptáte vestras lámpades.\end{Resp}
\begin{Resp}\Rsign Ecce sponsus venit, exíte óbviam ei.\end{Resp}
\begin{Resp}\Vsign Glória Patri, et Fílio, \Star{} et Spirítui Sancto.\end{Resp}
\begin{Resp}\Rsign Ecce sponsus venit, exíte óbviam ei.\end{Resp}
\switchcolumn
\EN
\begin{Resp}\Rsign At midnight there was a cry made: \Star{} Behold, the Bridegroom cometh, go ye out to meet him.\end{Resp}
\begin{Resp}\Vsign Trim your lamps, O ye wise virgins.\end{Resp}
\begin{Resp}\Rsign Behold, the Bridegroom cometh, go ye out to meet him.\end{Resp}
\begin{Resp}\Vsign Glory be to the Father, and to the Son, \Star{} and to the Holy Ghost.\end{Resp}
\begin{Resp}\Rsign Behold, the Bridegroom cometh, go ye out to meet him.\end{Resp}
\switchcolumn*

\LA
\LessonHead{Lectio IX}
\switchcolumn
\EN
\LessonHead{Lesson IX}
\switchcolumn*

\LA
\noindent De hujusmodi regno pacatissimo et ordinatissimo missus est foras princeps hujus sæculi, qui perversis inordinatisque dominatur. Hac pace intrinsecus constituta atque firmata, quascumque persecutiones ille qui foras missus est, forinsecus concitaverit, auget gloriam quæ secundum Deum est; non aliquid in illo ædificio labefactans, sed deficientibus machinis suis innotescere faciens, quanta firmitas intus exstructa sit. Ideo sequitur: Beati, qui persecutionem patiuntur propter justitiam: quoniam ipsorum est regnum cælorum.\par
\switchcolumn
\EN
\noindent From a kingdom such as this, which is a state of complete peace and order, the prince of this world is cast out. For he is one that can rule only through perversity and disorder. When this peace hath been inwardly established and confirmed, then whatever persecutions the prince of this world, who is now cast out, shall stir up from without, he only increaseth the glory which doth redound to God. For he is unable to tear down anything in that which is so upbuilt. And by the failure of his machinations, he doth but make manifest the strength with which it hath been inwardly built. Hence it followeth: Blessed are they that are persecuted for righteousness' sake, for theirs is the kingdom of heaven.\par
\switchcolumn*

\LA
\TeDeum{Te Deum}
\switchcolumn
\EN
\TeDeum{Te Deum}
\switchcolumn*

\end{paracol}

\par\needspace{10\baselineskip}\addvspace{1.2em}{\centering\small\textsc{Die 3 Novembris} · \textsc{3 November}\par}
\Office{In Commemoratione Omnium Fidelium Defunctorum}{In Commemoratione Omnium Fidelium Defunctorum}{Duplex}
\addcontentsline{toc}{subsection}{Die 3 Novembris · In Commemoratione Omnium Fidelium Defunctorum \textit{— In Commemoratione Omnium Fidelium Defunctorum}}
\begin{paracol}{2}
\LA
\LessonHead{Lectio I}
\switchcolumn
\EN
\LessonHead{Lesson I}
\switchcolumn*

\LA
\Cite{Job 7:16-21}
\switchcolumn
\EN
\Cite{Job 7:16-21}
\switchcolumn*

\LA
\noindent Parce mihi, Dómine; nihil enim sunt dies mei. Quid est homo, quia magníficas eum? aut quid appónis erga eum cor tuum? Vísitas eum dilúculo, et súbito probas illum. Usquequo non parcis mihi, nec dimíttis me, ut glútiam salívam meam? Peccávi, quid fáciam tibi, o custos hóminum? quare posuísti me contrárium tibi, et factus sum mihimetípsi gravis? Cur non tollis peccátum meum, et quare non aufers iniquitátem meam? Ecce nunc in púlvere dórmiam: et, si mane me quæsíeris, non subsístam.\par
\switchcolumn
\EN
\noindent Spare me, for my days are nothing. What is a man that thou shouldst magnify him? or why dost thou set thy heart upon him? Thou visitest him early in the morning, and thou provest him suddenly. How long wilt thou not spare me, nor suffer me to swallow down my spittle? I have sinned: what shall I do to thee, O keeper of men? why hast thou set me opposite to thee, and I am become burdensome to myself? Why dost thou not remove my sin, and why dost thou not take away my iniquity? Behold now I shall sleep in the dust: and if thou seek me in the morning, I shall not be.\par
\switchcolumn*

\LA
\begin{Resp}\Rsign Credo quod Redémptor meus vivit, et in novíssimo die de terra surrectúrus sum, \Star{} Et in carne mea vidébo Deum Salvatórem meum.\end{Resp}
\begin{Resp}\Vsign Quem visúrus sum ego ipse, et non álius; et óculi mei conspectúri sunt.\end{Resp}
\begin{Resp}\Rsign Et in carne mea vidébo Deum Salvatórem meum.\end{Resp}
\switchcolumn
\EN
\begin{Resp}\Rsign I believe that my Redeemer liveth, and that I shall stand up from the earth at the latter day, \Star{} And in my flesh shall I see God my Saviour.\end{Resp}
\begin{Resp}\Vsign Whom I shall see for myself, and mine eyes shall behold, and not another.\end{Resp}
\begin{Resp}\Rsign And in my flesh shall I see God my Saviour.\end{Resp}
\switchcolumn*

\LA
\LessonHead{Lectio II}
\switchcolumn
\EN
\LessonHead{Lesson II}
\switchcolumn*

\LA
\Cite{Job 14:1-6}
\switchcolumn
\EN
\Cite{Job 14:1-6}
\switchcolumn*

\LA
\noindent Homo natus de mulíere, brevi vivens témpore, replétur multis misériis. Qui quasi flos egréditur et contéritur, et fugit velut umbra, et nunquam in eódem statu pérmanet. Et dignum ducis super hujuscémodi aperíre óculos tuos, et addúcere eum tecum in judícium? Quis potest fácere mundum de immúndo concéptum sémine? Nonne tu qui solus es? Breves dies hóminis sunt: númerus ménsium ejus apud te est: constituísti términos ejus, qui præteríri non póterunt. Recéde páululum ab eo, ut quiéscat, donec optáta véniat, sicut mercenárii, dies ejus.\par
\switchcolumn
\EN
\noindent Man born of a woman, living for a short time, is filled with many miseries. Who cometh forth like a flower, and is destroyed, and fleeth as a shadow, and never continueth in the same state. And dost thou think it meet to open thy eyes upon such an one, and to bring him into judgment with thee? Who can make him clean that is conceived of unclean seed? is it not thou who only art? The days of man are short, and the number of his months is with thee: thou hast appointed his bounds which cannot be passed. Depart a little from him, that he may rest, until his wished for day come, as that of the hireling\par
\switchcolumn*

\LA
\begin{Resp}\Rsign Qui Lázarum resuscitásti a monuménto fœ́tidum, \Star{} Tu eis, Dómine, dona réquiem, et locum indulgéntiæ.\end{Resp}
\begin{Resp}\Vsign Qui ventúrus es judicáre vivos et mórtuos, et sǽculum per ignem.\end{Resp}
\begin{Resp}\Rsign Tu eis, Dómine, dona réquiem, et locum indulgéntiæ.\end{Resp}
\switchcolumn
\EN
\begin{Resp}\Rsign Thou Who didst call up Lazarus from the grave after that he had begun to stink \Star{} Do Thou, O Lord, grant them rest and a place of forgiveness.\end{Resp}
\begin{Resp}\Vsign Thou Who shalt come to judge the quick and dead, and the world by fire\end{Resp}
\begin{Resp}\Rsign Do Thou, O Lord, grant them rest and a place of forgiveness.\end{Resp}
\switchcolumn*

\LA
\LessonHead{Lectio III}
\switchcolumn
\EN
\LessonHead{Lesson III}
\switchcolumn*

\LA
\Cite{Job 19:20-27}
\switchcolumn
\EN
\Cite{Job 19:20-27}
\switchcolumn*

\LA
\noindent Pelli meæ, consúmptis cárnibus, adhǽsit os meum, et derelícta sunt tantúmmodo lábia circa dentes meos. Miserémini mei, miserémini mei saltem vos, amíci mei, quia manus Dómini tétigit me. Quare persequímini me sicut Deus, et cárnibus meis saturámini? Quis mihi tríbuat ut scribántur sermónes mei? quis mihi det ut exaréntur in libro stylo férreo et plumbi lámina, vel celte sculpántur in sílice? Scio enim quod redémptor meus vivit, et in novíssimo die de terra surrectúrus sum: et rursum circúmdabor pelle mea, et in carne mea vidébo Deum meum: quem visúrus sum ego ipse, et óculi mei conspectúri sunt, et non álius: repósita est hæc spes mea in sinu meo.\par
\switchcolumn
\EN
\noindent The flesh being consumed. My bone hath cleaved to my skin, and nothing but lips are left about my teeth. Have pity on me, have pity on me, at least you my friends, because the hand of the Lord hath touched me. Why do you persecute me as God, and glut yourselves with my flesh? Who will grant me that my words may be written? Who will grant me that they may be marked down in a book? With an iron pen and in a plate of lead, or else be graven with an instrument in flint stone. For I know that my Redeemer liveth, and in the last day I shall rise out of the earth. And I shall be clothed again with my skin, and in my flesh I will see my God. Whom I myself shall see, and my eyes shall behold, and not another: this my hope is laid up in my bosom.\par
\switchcolumn*

\LA
\begin{Resp}\Rsign Dómine, quando véneris judicáre terram, ubi me abscóndam a vultu iræ tuæ? \Star{} Quia peccávi nimis in vita mea.\end{Resp}
\begin{Resp}\Vsign Commíssa mea pavésco, et ante te erubésco: dum véneris judicáre, noli me condemnáre.\end{Resp}
\begin{Resp}\Rsign Quia peccávi nimis in vita mea. \$Requiem\end{Resp}
\begin{Resp}\Rsign Quia peccávi nimis in vita mea.\end{Resp}
\switchcolumn
\EN
\begin{Resp}\Rsign Lord, when Thou comest to judge the earth, where shall I hide myself from the face of thy wrath? \Star{} For I have sinned greatly in my life.\end{Resp}
\begin{Resp}\Vsign I dread my sins, I blush before thee I see the great tribunal set! In fear and terror I implore thee, Forgive when soul and Judge are met!\end{Resp}
\begin{Resp}\Rsign For I have sinned greatly in my life. \$Requiem\end{Resp}
\begin{Resp}\Rsign For I have sinned greatly in my life.\end{Resp}
\switchcolumn*

\LA
\LessonHead{Lectio IV}
\switchcolumn
\EN
\LessonHead{Lesson IV}
\switchcolumn*

\LA
\LessonTitle{Ex Libro sancti Augustíni Epíscopi de cura pro mórtuis gerénda}
\switchcolumn
\EN
\LessonTitle{From the book of St. Augustine, the Bishop, on the Care for the Deceased}
\switchcolumn*

\LA
\Source{Cap 2 et 3}
\switchcolumn
\EN
\Source{Cap. 2 \& 3}
\switchcolumn*

\LA
\noindent Curátio fúneris, condítio sepultúræ, pompa exsequiárum magis sunt vivórum solácia quam subsídia mortuórum. Nec ídeo tamen contemnénda et abiciénda sunt córpora defunctórum, maximéque justórum ac fidélium, quibus tamquam órganis et vasis ad ómnia bona ópera sancte usus est spíritus. Si enim patérna vestis et ánulus, ac si quid hujúsmodi, tanto cárius est pósteris, quanto erga paréntes major afféctus; nullo modo ipsa spernénda sunt córpora, quæ útique multo familiárius atque conjúnctius quam quǽlibet induménta gestámus. Hæc enim non ad ornaméntum vel adjutórium, quod adhibétur extrínsecus, sed ad ipsam natúram hóminis pértinent. Unde et antiquórum justórum fúnera officiósa pietáte curáta sunt, et exséquiæ celebrátæ, et sepultúra provísa; ipsíque, cum víverent, de sepeliéndis vel étiam transferéndis suis corpóribus fíliis mandavérunt.\par
\switchcolumn
\EN
\noindent If this be true, doubtless also the providing for the interment of bodies a place at the Memorials of Saints, is a mark of a good human affection towards the remains of one's friends. Yet it follows not that the bodies of the departed are to be despised and flung aside, and above all of just and faithful men, which bodies as organs and vessels to all good works their spirit has holily used. For if a father's garment and ring, and whatever such like, is the more dear to those whom they leave behind, the greater their affection is towards their parents, in no wise are the bodies themselves to be spurned, which truly we wear in more familiar and close conjunction than any of our putting on. For these pertain not to ornament or aid which is applied from without, but to the very nature of man. Whence also the funerals of the just men of old were with dutiful piety cared for, and their obsequies celebrated, and sepulture provided: and themselves while living did touching burial or even translation of their bodies give charge to their sons.\par
\switchcolumn*

\LA
\begin{Resp}\Rsign Meménto mei, Deus, quia ventus est vita mea, \Star{} Nec aspíciat me visus hóminis.\end{Resp}
\begin{Resp}\Vsign De profúndis clamávi ad te, Dómine: Dómine, exáudi vocem meam.\end{Resp}
\begin{Resp}\Rsign Nec aspíciat me visus hóminis.\end{Resp}
\switchcolumn
\EN
\begin{Resp}\Rsign Remember, O God, that my life is wind. \Star{} The eye of him that hath seen me shall see me no more.\end{Resp}
\begin{Resp}\Vsign Out of the depths have I cried unto thee, O Lord! Lord, hear my voice.\end{Resp}
\begin{Resp}\Rsign The eye of him that hath seen me shall see me no more.\end{Resp}
\switchcolumn*

\LA
\LessonHead{Lectio V}
\switchcolumn
\EN
\LessonHead{Lesson V}
\switchcolumn*

\LA
\Source{Cap. 4}
\switchcolumn
\EN
\Source{Cap. 4}
\switchcolumn*

\LA
\noindent Recordántis et precántis afféctus cum defúnctis a fidélibus caríssimis exhibétur, eum prodésse non dúbium est iis, qui cum in córpore víverent, tália sibi post hanc vitam prodésse meruérunt. Verum, etsi áliqua necéssitas vel humári córpora, vel in sacris locis humári nulla data facultáte permíttat, non sunt prætermitténdæ supplicatiónes pro spirítibus mortuórum: quas faciéndas pro ómnibus in christiána et cathólica societáte defúnctis, étiam tácitis eórum nomínibus, sub generáli commemoratióne suscépit Ecclésia; ut quibus ad ista desunt paréntes, aut fílii, aut quicúmque cognáti vel amíci, ab una eis exhibeántur pia matre commúni. Si autem deéssent istæ supplicatiónes, quæ fiunt recta fide ac pietáte pro mórtuis, puto quod nihil prodésset spirítibus eórum, quámlibet in locis sanctis exánima córpora poneréntur.\par
\switchcolumn
\EN
\noindent And when this affection is exhibited to the departed by faithful men who were most dear to them, there is no doubt that it profits them who while living in the body merited that such things should profit them after this life. But even if some necessity should through absence of all facility not allow bodies to be interred, or in such places interred, yet should there be no pretermitting of supplications for the spirits of the dead: which supplications, that they should be made for all in Christian and catholic fellowship departed, even without mentioning of their names, under a general commemoration, the Church has charged herself withal; to the intent that they which lack, for these offices, parents or sons or whatever kindred or friends, may have the same afforded unto them by the one pious mother which is common to all. But if there were lack of these supplications, which are made with right faith and piety for the dead, I account that it should not a whit profit their spirits, howsoever in holy places the lifeless bodies should be deposited.\par
\switchcolumn*

\LA
\begin{Resp}\Rsign Hei mihi, Dómine, quia peccávi nimis in vita mea: Quid fáciam, miser? ubi fúgiam, nisi ad te, Deus meus? \Star{} Miserére mei, dum véneris in novíssimo die.\end{Resp}
\begin{Resp}\Vsign Anima mea turbáta est valde, sed tu, Dómine, succúrre ei.\end{Resp}
\begin{Resp}\Rsign Miserére mei, dum véneris in novíssimo die.\end{Resp}
\switchcolumn
\EN
\begin{Resp}\Rsign Woe is me, O Lord! for I have sinned greatly in my life. I am smitten what shall I do? Whither shall I flee but unto thee, O my God? \Star{} Have mercy upon me, when Thou comest at the latter day.\end{Resp}
\begin{Resp}\Vsign My soul is sore vexed, but Thou, O Lord, help me.\end{Resp}
\begin{Resp}\Rsign Have mercy upon me, when Thou comest at the latter day.\end{Resp}
\switchcolumn*

\LA
\LessonHead{Lectio VI}
\switchcolumn
\EN
\LessonHead{Lesson VI}
\switchcolumn*

\LA
\Source{Cap. 18}
\switchcolumn
\EN
\Source{Cap. 18}
\switchcolumn*

\LA
\noindent Quæ cum ita sint, non existimémus ad mórtuos, pro quibus curam gérimus, perveníre, nisi quod pro eis sive altáris, sive oratiónum, sive eleemosynárum sacrifíciis solémniter supplicámus: quamvis non pro quibus fiunt, ómnibus prosint; sed iis tantum pro quibus, dum vivunt, comparátur, ut prosint. Sed quia non discérnimus qui sunt, opórtet ea pro regenerátis ómnibus fácere, ut nullus eórum prætermittátur, ad quos hæc benefícia possint et débeant perveníre. Mélius enim supérerunt ista eis, quibus nec obsunt nec prosunt; quam eis déerunt, quibus prosunt. Diligéntius tamen facit hæc quisque pro necessáriis suis, quo pro illo fiat simíliter a suis. Córpori autem humándo quidquid impénditur, non est præsídium salútis, sed humanitátis offícium, secúndum afféctum quo nemo unquam carnem suam ódio habet. Unde opórtet ut quam potest pro carne próximi curam gerat, cum ille inde recésserit, qui gerébat. Et si hæc fáciunt, qui carnis resurrectiónem non credunt, quanto magis debent fácere qui credunt; ut córpori mórtuo, sed tamen resurrectúro et in æternitáte mansúro, impénsum ejúsmodi offícium sit étiam quodámmodo ejúsdem fídei testimónium!\par
\switchcolumn
\EN
\noindent Which things being so, let us not think that to the dead for whom we have a care, any thing reaches save what by sacrifices either of the altar, or of prayers, or of alms, we solemnly supplicate: although not to all for whom they are done be they profitable, but to them only by whom while they live it is obtained that they should be profitable. But forasmuch as we discern not who these be, it is meet to do them for all regenerate persons, that none of them may be passed by to whom these benefits may and ought to reach. For better it is that these things shall be superfluously done to them whom they neither hinder nor help, than lacking to them whom they help. More diligently however does each man these things for his own near and dear friends, in order that they may be likewise done unto him by his. But as for the burying of the body, whatever is bestowed on that, is no aid of salvation, but an office of humanity, according to that affection by which no man ever hates his own flesh. Whence it is fitting that he take what care he is able for the flesh of his neighbor, when he is gone that bare it. And if they do these things who believe not the resurrection of the flesh, how much more are they beholden to do the same who do believe; that so, an office of this kind bestowed upon a body, dead but yet to rise again and to remain to eternity, may also be in some sort a testimony of the same faith?\par
\switchcolumn*

\LA
\begin{Resp}\Rsign Ne recordéris peccáta mea, Dómine, \Star{} Dum véneris judicáre sǽculum per ignem.\end{Resp}
\begin{Resp}\Vsign Dírige, Dómine, Deus meus, in conspéctu tuo viam meam.\end{Resp}
\begin{Resp}\Rsign Dum véneris judicáre sǽculum per ignem. \$Requiem\end{Resp}
\begin{Resp}\Rsign Dum véneris judicáre sǽculum per ignem.\end{Resp}
\switchcolumn
\EN
\begin{Resp}\Rsign Hold not my sins in remembrance, O Lord, \Star{} When thou comest to judge the world by fire.\end{Resp}
\begin{Resp}\Vsign Make my way straight before thy face, O Lord my God.\end{Resp}
\begin{Resp}\Rsign When thou comest to judge the world by fire. \$Requiem\end{Resp}
\begin{Resp}\Rsign When thou comest to judge the world by fire.\end{Resp}
\switchcolumn*

\LA
\LessonHead{Lectio VII}
\switchcolumn
\EN
\LessonHead{Lesson VII}
\switchcolumn*

\LA
\LessonTitle{De Epístola prima beáti Pauli Apóstoli ad Corínthios}
\switchcolumn
\EN
\LessonTitle{From the first letter of St. Paul to the Corinthians}
\switchcolumn*

\LA
\Cite{1 Cor 15:12-22}
\switchcolumn
\EN
\Cite{1 Cor 15:12-22}
\switchcolumn*

\LA
\noindent Si Christus prædicátur quod resurréxit a mórtuis, quómodo quidam dicunt in vobis, quóniam resurréctio mortuórum non est? Si autem resurréctio mortuórum non est: neque Christus resurréxit. Si autem Christus non resurréxit, inánis est ergo prædicátio nostra, inánis est et fides vestra: invenímur autem et falsi testes Dei: quóniam testimónium díximus advérsus Deum quod suscitáverit Christum, quem non suscitávit, si mórtui non resúrgunt. Nam si mórtui non resúrgunt, neque Christus resurréxit. Quod si Christus non resurréxit, vana est fides vestra: adhuc enim estis in peccátis vestris. Ergo et qui dormiérunt in Christo, periérunt. Si in hac vita tantum in Christo sperántes sumus, miserabilióres sumus ómnibus homínibus. Nunc autem Christus resurréxit a mórtuis primítiæ dormiéntium, quóniam quidem per hóminem mors, et per hóminem resurréctio mortuórum. Et sicut in Adam omnes moriúntur, ita et in Christo omnes vivificabúntur.\par
\switchcolumn
\EN
\noindent Now if Christ be preached, that he arose again from the dead, how do some among you say, that there is no resurrection of the dead? But if there be no resurrection of the dead, then Christ is not risen again. And if Christ be not risen again, then is our preaching vain, and your faith is also vain. Yea, and we are found false witnesses of God: because we have given testimony against God, that he hath raised up Christ; whom he hath not raised up, if the dead rise not again. For if the dead rise not again, neither is Christ risen again. And if Christ be not risen again, your faith is vain, for you are yet in your sins. Then they also that are fallen asleep in Christ, are perished. If in this life only we have hope in Christ, we are of all men most miserable. But now Christ is risen from the dead, the firstfruits of them that sleep For by a man came death, and by a man the resurrection of the dead. And as in Adam all die, so also in Christ all shall be made alive.\par
\switchcolumn*

\LA
\begin{Resp}\Rsign Peccántem me cotídie, et non me pœniténtem, timor mortis contúrbat me: \Star{} Quia in inférno nulla est redémptio, miserére mei, Deus, et salva me.\end{Resp}
\begin{Resp}\Vsign Deus, in nómine tuo salvum me fac, et in virtúte tua líbera me.\end{Resp}
\begin{Resp}\Rsign Quia in inférno nulla est redémptio, miserére mei, Deus, et salva me.\end{Resp}
\switchcolumn
\EN
\begin{Resp}\Rsign Forasmuch as I sin daily, and repent not, the fear of death troubleth me. \Star{} O God, have mercy upon me, and save me, for in hell there is no redemption.\end{Resp}
\begin{Resp}\Vsign Save me, O God, by thy name, and judge me in thy strength.\end{Resp}
\begin{Resp}\Rsign O God, have mercy upon me, and save me, for in hell there is no redemption.\end{Resp}
\switchcolumn*

\LA
\LessonHead{Lectio VIII}
\switchcolumn
\EN
\LessonHead{Lesson VIII}
\switchcolumn*

\LA
\Cite{1 Cor 15:35-44}
\switchcolumn
\EN
\Cite{1 Cor 15:35-44}
\switchcolumn*

\LA
\noindent Sed dicet áliquis: Quómodo resúrgunt mórtui? qualíve córpore vénient? Insípiens, tu quod séminas non vivificátur, nisi prius moriátur: et quod séminas, non corpus, quod futúrum est, séminas, sed nudum granum, ut puta trítici, aut alicújus ceterórum. Deus autem dat illi corpus sicut vult, et unicuíque séminum próprium corpus. Non omnis caro, éadem caro: sed ália quidem hóminum, ália vero pécorum, ália vólucrum, ália autem píscium. Et córpora cæléstia, et córpora terréstria: sed ália quidem cæléstium glória, ália autem terréstrium. Alia cláritas solis, ália cláritas lunæ, et ália cláritas stellárum; stella enim a stella differt in claritáte: sic et resurréctio mortuórum. Seminátur in corruptióne, surget in incorruptióne. Seminátur in ignobilitáte, surget in glória: seminátur in infirmitáte, surget in virtúte: seminátur corpus animále, surget corpus spiritále.\par
\switchcolumn
\EN
\noindent But some man will say: How do the dead rise again? or with what manner of body shall they come? Senseless man, that which thou sowest is not quickened, except it die first. And that which thou sowest, thou sowest not the body that shall be; but bare grain, as of wheat, or of some of the rest. But God giveth it a body as he will: and to every seed its proper body. All flesh is not the same flesh: but one is the flesh of men, another of beasts, another of birds, another of fishes. And there are bodies celestial, and bodies terrestrial: but, one is the glory of the celestial, and another of the terrestrial. One is the glory of the sun, another the glory of the moon, and another the glory of the stars. For star differeth from star in glory. So also is the resurrection of the dead. It is sown in corruption, it shall rise in incorruption. It is sown in dishonour, it shall rise in glory. It is sown in weakness, it shall rise in power. It is sown a natural body, it shall rise a spiritual body.\par
\switchcolumn*

\LA
\begin{Resp}\Rsign Dómine, secúndum actum meum noli me judicáre: nihil dignum in conspéctu tuo egi; ídeo déprecor majestátem tuam, \Star{} Ut tu, Deus, déleas iniquitátem meam.\end{Resp}
\begin{Resp}\Vsign Amplius lava me, Dómine, ab injustítia mea, et a delícto meo munda me.\end{Resp}
\begin{Resp}\Rsign Ut tu, Deus, déleas iniquitátem meam.\end{Resp}
\switchcolumn
\EN
\begin{Resp}\Rsign O Lord, judge me not according to my works; for I have done nothing that can be counted in respect of thee. I beseech thy majesty therefore, that thou wouldest \Star{} Blot out my transgressions, O God.\end{Resp}
\begin{Resp}\Vsign Lord, wash me thoroughly from mine iniquity and cleanse me from my sin.\end{Resp}
\begin{Resp}\Rsign Blot out my transgressions, O God.\end{Resp}
\switchcolumn*

\LA
\LessonHead{Lectio IX}
\switchcolumn
\EN
\LessonHead{Lesson IX}
\switchcolumn*

\LA
\Cite{1 Cor 15:51-58}
\switchcolumn
\EN
\Cite{1 Cor 15:51-58}
\switchcolumn*

\LA
\noindent Ecce mystérium vobis dico: Omnes quidem resurgémus, sed non omnes immutábimur. In moménto, in ictu óculi, in novíssima tuba: canet enim tuba, et mórtui resúrgent incorrúpti: et nos immutábimur. Opórtet enim corruptíbile hoc indúere incorruptiónem: et mortále hoc indúere immortalitátem. Cum autem mortále hoc indúerit immortalitátem, tunc fiet sermo, qui scriptus est: Absórpta est mors in victória. Ubi est mors victória tua? ubi est mors stímulus tuus? Stímulus autem mortis peccátum est: virtus vero peccáti lex. Deo autem grátias, qui dedit nobis victóriam per Dóminum nostrum Jesum Christum. Itaque fratres mei dilécti, stábiles estóte, et immóbiles: abundántes in ópere Dómini semper, sciéntes quod labor vester non est inánis in Dómino.\par
\switchcolumn
\EN
\noindent Behold, I tell you a mystery. We shall all indeed rise again: but we shall not all be changed. In a moment, in the twinkling of an eye, at the last trumpet: for the trumpet shall sound, and the dead shall rise again incorruptible: and we shall be changed. For this corruptible must put on incorruption; and this mortal must put on immortality. And when this mortal hath put on immortality, then shall come to pass the saying that is written: Death is swallowed up in victory. O death, where is thy victory? O death, where is thy sting? Now the sting of death is sin: and the power of sin is the law. But thanks be to God, who hath given us the victory through our Lord Jesus Christ. Therefore, my beloved brethren, be ye steadfast and unmoveable; always abounding in the work of the Lord, knowing that your labour is not in vain in the Lord.\par
\switchcolumn*

\LA
\begin{Resp}\Rsign Líbera me, Dómine, de morte ætérna in die illa treménda, \Star{} Quando cæli movéndi sunt et terra, * Dum véneris judicáre sǽculum per ignem.\end{Resp}
\begin{Resp}\Vsign Tremens factus sum ego et tímeo, dum discússio vénerit atque ventúra ira.\end{Resp}
\begin{Resp}\Rsign Quando cæli movéndi sunt et terra.\end{Resp}
\begin{Resp}\Vsign Dies illa, dies iræ, calamitátis et misériæ, dies magna et amára valde.\end{Resp}
\begin{Resp}\Rsign Dum véneris judicáre sǽculum per ignem. \$Requiem\end{Resp}
\begin{Resp}\Rsign Líbera me, Dómine, de morte ætérna in die illa treménda, * Quando cæli movéndi sunt et terra, * Dum véneris judicáre sǽculum per ignem.\end{Resp}
\switchcolumn
\EN
\begin{Resp}\Rsign Deliver me, O Lord, from eternal death in that awful day \Star{} when the heavens and the earth shall be shaken, and Thou shalt come to judge the world by fire.\end{Resp}
\begin{Resp}\Vsign Quaking and dread take hold upon me, when I look for the coming of the trial and the wrath to come.\end{Resp}
\begin{Resp}\Rsign When the heavens and the earth shall be shaken.\end{Resp}
\begin{Resp}\Vsign That day is a day of wrath, of wasteness and desolation, a great day and exceeding bitter.\end{Resp}
\begin{Resp}\Rsign When Thou shalt come to judge the world by fire. \$Requiem\end{Resp}
\begin{Resp}\Rsign Deliver me, O Lord, from eternal death in that awful day * when the heavens and the earth shall be shaken, and Thou shalt come to judge the world by fire.\end{Resp}
\switchcolumn*

\end{paracol}

\par\needspace{10\baselineskip}\addvspace{1.2em}{\centering\small\textsc{Die 4 Novembris} · \textsc{4 November}\par}
\Office{S. Caroli Episcopi et Confessoris}{St. Charles Borromeo, Bishop and Confessor}{Duplex}
\addcontentsline{toc}{subsection}{Die 4 Novembris · S. Caroli Episcopi et Confessoris \textit{— St. Charles Borromeo, Bishop and Confessor}}
\begin{paracol}{2}
\LA
\RefLine{Lectiones I–III: de Scriptura occurrente.}
\switchcolumn
\EN
\RefLine{Lessons I–III: from the Scripture of the day.}
\switchcolumn*

\LA
\RefLine{Lectiones VII–VIII: ex Commúni unius Confessoris Pontificis (C4).}
\switchcolumn
\EN
\RefLine{Lessons VII–VIII: from the Common of a Confessor Bishop (C4).}
\switchcolumn*

\LA
\RefLine{Lectio IX: de Ss. Mart. Vitalis et Agricolæ (11-04cc).}
\switchcolumn
\EN
\RefLine{Lesson IX: of Sts. Vitalis and Agricola, Martyrs (11-04cc).}
\switchcolumn*

\LA
\LessonHead{Lectio IV}
\switchcolumn
\EN
\LessonHead{Lesson IV}
\switchcolumn*

\LA
\noindent Cárolus, Mediolani nobili Borromæórum família natus, quanta futurus esset sanctitáte conspicuus, divina lux super pariéntis matris cubiculum noctu coruscans præsignávit. A puerítia clericali milítiæ adscriptus, abbatía postmodum insignítus, patrem admonuit ne réditus in rem familiarem converteréntur; quorum ipse nactus administratiónem, quidquid supérerat, expendebat in páuperes. Adoléscens liberálibus disciplinis Papíæ operam dedit. Castitátem adeo coluit, ut impudícas étiam mulieres, ad labefactándam ejus pudicítiam pluries immissas, invicta constántia fugaverit. Vigesimum tertium ætátis annum agens, a Pio quarto ejus avunculo in sacrum cardinalium collegium cooptatus, insígni pietátis ac virtútum ómnium splendore præluxit. Mox ab eodem Mediolanénsis archiepíscopus creatus, in eo plurimam operam adhibuit, ut, juxta sacrosanctum Tridentinum concílium, quod ejus potíssimum sollicitúdine jam tum fuerat absolutum, ecclésiam sibi commissam compóneret; atque, ut depravatos plebis suæ mores reformaret, præter iteratam sæpius synodórum celebratiónem, seípsum exímiæ sanctitátis præbuit exemplar. In profligandis hæreticis e partibus Rhætórum et Helvetiórum, quorum plurimos ad christianam fidem convértit, maxime laborávit.\par
\switchcolumn
\EN
\noindent Charles, of the noble family of Borromeo, was born (on the 2nd day of October, in the year of our Lord 1538,) at (the Castle of Arona, fourteen miles from) Milan. In foretoken of his holy life, God caused a bright light to shine by night over the chamber where his mother lay in travail. As soon as his age would allow him, he received the tonsure. When he was twelve years old, he was made Abbot (of the rich Benedictine Abbey of St. Gratinian and St. Felin,) but reminded his father that the revenues thereof were not to be used as mere family property. His father, to whom the administration of these revenues fell during his son's non-age, still gave them forthwith over to him, and whatever was left over, he gave to the poor. While he was young he studied letters at Pavia. He kept his purity thoroughly, so that he scared away the unclean women, of whom many were set upon him, to overthrow his self-control. In the twenty-third year of his age, his uncle Pius IV. made him a Cardinal, in which dignity he was a burning and shining light of godliness and all graces before the whole of the Sacred College. About forty days afterwards the same Pope created him Archbishop of Milan. As such it was his great desire to order the Church committed to his charge in accordance with the requirements of the most holy Council of Trent, which was in great part by his labours brought to a conclusion. To raise up the degraded lives of the people, he oftentimes held Synods, but himself set an example of deep godliness. He worked earnestly to purge the parts about the Alps and borders of Switzerland of heresy, and brought many of the heretics to the Christian faith.\par
\switchcolumn*

\LA
\begin{Resp}\Rsign Invéni David servum meum, óleo sancto meo unxi eum: \Star{} Manus enim mea auxiliábitur ei.\end{Resp}
\begin{Resp}\Vsign Nihil profíciet inimícus in eo, et fílius iniquitátis non nocébit ei.\end{Resp}
\begin{Resp}\Rsign Manus enim mea auxiliábitur ei.\end{Resp}
\switchcolumn
\EN
\begin{Resp}\Rsign I have found David My servant, with My holy oil have I anointed him \Star{} For My hand shall help him.\end{Resp}
\begin{Resp}\Vsign The enemy shall prevail nothing against him, nor the son of wickedness afflict him.\end{Resp}
\begin{Resp}\Rsign For My hand shall help him.\end{Resp}
\switchcolumn*

\LA
\LessonHead{Lectio V}
\switchcolumn
\EN
\LessonHead{Lesson V}
\switchcolumn*

\LA
\noindent Hujus viri caritas præcipue enituit, cum, Uritano principatu véndito, prétium univérsum, ad quadragínta aureórum millia, una die in páuperes erogávit. Nec minore pietáte viginti millia, quæ sibi fúerant legata, distríbuit. Ecclesiásticos proventus, quibus ab avunculo copiose fuerat cumulátus, dimísit, nonnullis retentis, quibus ad proprios usus et egenórum necessitates utebátur. Quo témpore pestis Mediolani grassabátur, domesticam supellectilem, ne relícto sibi lectulo, in eosdem alendos cóntulit, super nuda in pósterum tábula decumbens; eoque morbo laborántes sédulo invisens, paterno reficiebat afféctu, et Ecclésiæ sacraménta propriis ipse mánibus administrans, mirum in modum solabátur. Humillimis interim precibus reconciliator accédens, publica supplicatióne indicta, fune sibi ad collum alligato, nudis pédibus étiam offendiculo cruentátis, crucem bájulans, semetípsum pro peccátis pópuli hóstiam offerens, divinam indignatiónem avértere satagebat. Ecclesiásticæ libertátis fuit acerrimus propugnator. Disciplinæ vero restituéndæ sollícitus, a seditiosis, dum oratióni insísteret, torménti bellici laxáta rota, igneo glóbulo percússus, divina virtúte servátur illæsus.\par
\switchcolumn
\EN
\noindent Charity was the brightest mark of his life. His principality of Oria, (in the kingdom of Naples,) he sold for forty thousand crowns, and gave the whole sum to the poor in one day. Twenty thousand crowns being left him as a legacy (by Virginia, widow of Count Frederick Borromeo,) he gave the whole to the poor. The incomes of the benefices wherewith he had been loaded by his uncle, he spent upon the needs of the poor, except what he used for himself. When the plague grievously raged in Milan, he gave up to the sick poor the furniture of his own house, even to his own bedding, and thenceforward slept upon the boards. He constantly visited the sick, cheered them by his fatherly kindness, and wonderfully comforted them, ministering to them with his own hands the Sacraments of the Church. At the same time he drew near to plead for them with God in lowly entreaty, and ordered a public Procession wherein he walked himself carrying a Cross, with a rope halter round his neck, and his bare feet bleeding from the stones, and fain to turn away the Divine anger by offering himself as a scapegoat for the sins of his people. He was a stout defender of the freedom of the Church. But in the Church he was an earnest reformer of discipline, and once, when he was engaged in prayer, (the paid agent of) some conspirators took a shot at him with a blunderbuss, but, though the ball struck him, the power of God kept him unharmed.\par
\switchcolumn*

\LA
\begin{Resp}\Rsign Pósui adjutórium super poténtem, et exaltávi eléctum de plebe mea: \Star{} Manus enim mea auxiliábitur ei.\end{Resp}
\begin{Resp}\Vsign Invéni David servum meum, óleo sancto meo unxi eum.\end{Resp}
\begin{Resp}\Rsign Manus enim mea auxiliábitur ei.\end{Resp}
\switchcolumn
\EN
\begin{Resp}\Rsign I have laid help upon one that is mighty, and have exalted one chosen out of My people \Star{} For My hand shall help him.\end{Resp}
\begin{Resp}\Vsign I have found David My servant, with My holy oil have I anointed him.\end{Resp}
\begin{Resp}\Rsign For My hand shall help him.\end{Resp}
\switchcolumn*

\LA
\LessonHead{Lectio VI}
\switchcolumn
\EN
\LessonHead{Lesson VI}
\switchcolumn*

\LA
\noindent Abstinéntia fuit admirábili; jejunábat sæpíssime, pane tantum et aqua, solis quandoque lupínis contentus. Nocturnis vigiliis, asperrimo cilício, assiduis flagellis corpus domábat. Humilitátis ac mansuetúdinis studiosíssimus fuit. Oratiónem ac verbi Dei prædicatiónem, gravíssimis licet curis occupatus, numquam intermisit. Multas ecclésias, monasteria, collegia ædificávit. Plura scripsit, ad episcopórum præsertim instructiónem utilíssima; cujus étiam ópera parochórum catechismus pródiit. Demum, in solitúdinem Varalli montis, ubi scultpis imagínibus Dominicæ passiónis mysteria ad vivum repræsentántur, secessit; ibique , diébus aliquot voluntária castigatióne asperam, sed Christi dolórum meditatiónibus suavem vitam ducens, in febrim íncidit. Mediolanum reversus, ingravescénte morbo, cínere ac cilício coopertus, et óculis in Crucifixi imáginem defixis, migrávit in cælum, ætátis anno quadragesimo septimo, Dómini vero millesimo quingentésimo octogesimo quarto, tertio Nonas Novembris. Quem, miraculis clarum, Paulus quintus Pontifex maximus in Sanctórum númerum retulit.\par
\switchcolumn
\EN
\noindent He was remarkable for his abstinence. He very often fasted upon nothing but bread and water, and sometimes nothing but lupines. He tamed his body by depriving himself of sleep, by very rough haircloth, and by constant scourging. He was an earnest practiser of lowliness and meekness. However much he was taken up with business, he never gave himself relaxation from prayer and from preaching the word of God. He built many Churches, convents, and schools. He wrote much matter, useful more especially for the good of Bishops. The publication of the Parish Priests' Catechism was due to his care. In October, 1584, he withdrew himself, for the purpose of making a retreat, to (what is called) the Sacro Monte of Varallo, an hill whereon (many sacred subjects and especially) the incidents of the Lord's sufferings are represented in life-size groups of coloured figures. (On Oct. 24) he was taken ill of a (tertian) ague, \textasciitilde{}(but concealed it,) and lived there for some days a life of torture by voluntary suffering, but of sweetness by thoughts of Christ's woes. After his return to Milan, (which he reached in a litter upon All Souls' Day,) his sickness became hopeless, and early in the night between the (3rd and) 4th days of November, in the 47th year of his own age, and in that of our Lord 1584, covered with ashes and sackcloth, and with his eyes fixed upon the image of Christ crucified, he exchanged earth for heaven. He was famous for miracles, and Pope Paul V. numbered him among the Saints.\par
\switchcolumn*

\LA
\begin{Resp}\Rsign Iste est, qui ante Deum magnas virtútes operátus est, et omnis terra doctrína ejus repléta est: \Star{} Ipse intercédat pro peccátis ómnium populórum.\end{Resp}
\begin{Resp}\Vsign Iste est, qui contémpsit vitam mundi, et pervénit ad cæléstia regna.\end{Resp}
\begin{Resp}\Rsign Ipse intercédat pro peccátis ómnium populórum.\end{Resp}
\begin{Resp}\Vsign Glória Patri, et Fílio, \Star{} et Spirítui Sancto.\end{Resp}
\begin{Resp}\Rsign Ipse intercédat pro peccátis ómnium populórum.\end{Resp}
\switchcolumn
\EN
\begin{Resp}\Rsign This is he which wrought great wonders before God, and the whole earth is full of his teaching \Star{} May he pray for all people, that their sins may be forgiven unto them\end{Resp}
\begin{Resp}\Vsign This is he which loved not his life in this world, and hath attained unto the kingdom of heaven.\end{Resp}
\begin{Resp}\Rsign May he pray for all people, that their sins may be forgiven unto them\end{Resp}
\begin{Resp}\Vsign Glory be to the Father, and to the Son, \Star{} and to the Holy Ghost.\end{Resp}
\begin{Resp}\Rsign May he pray for all people, that their sins may be forgiven unto them\end{Resp}
\switchcolumn*

\end{paracol}

\par\needspace{10\baselineskip}\addvspace{1.2em}{\centering\small\textsc{Die 4 Novembris} · \textsc{4 November}\par}
\Office{Ss. Mart. Vitalis et Agricolæ}{Sts. Vitalis and Agricola, Martyrs}{Simplex}
\addcontentsline{toc}{subsection}{Die 4 Novembris · Ss. Mart. Vitalis et Agricolæ \textit{— Sts. Vitalis and Agricola, Martyrs}}
\begin{paracol}{2}
\LA
\LessonHead{Lectio IX (pro commemoratione)}
\switchcolumn
\EN
\LessonHead{Lesson IX (when commemorated)}
\switchcolumn*

\LA
\LessonTitle{Commemoratio: Ss. Mart. Vitalis et Agricolæ}
\switchcolumn
\EN
\LessonTitle{Commemoration of Sts. Vitalis and Agricola}
\switchcolumn*

\LA
\noindent Vitalis et Agricola ejus dominus, in persecutióne Diocletiáni et Maximiáni, Bononiæ ob Jesu Christi prædicatiónem comprehénsi sunt. Cumque Vitalis, quo magis precibus et minis tentarétur ut senténtiam mutaret, eo magis se Christi cultórem ac servum profiterétur; vario tormentórum genere cruciatus, constanter ómnia pérferens, in oratióne spíritum Deo réddidit. Agricola vero, cum ejus supplícium dilátum esset, si forte, tormentis servi permotus, Christo vellet renuntiáre, ejus exemplo magis confirmátus est. Itaque, cruci affixus, Vitali servo consors et cosius fuit nobilis martyrii. Eórum corpora, ad Judæórum sepúlcra cum essent humata, a sancto Ambrosio invénta, in sacrum celebremque locum transláta sunt.\par
\switchcolumn
\EN
\noindent Vitalis was a slave, and Agricola his owner. They were arrested at Bologna in the persecution under Diocletian and Maximian, for preaching Jesus Christ. Vitalis, the more he was implored and threatened to change his mind so much the more proclaimed himself a worshipper and servant of Christ, and after bravely bearing a course of diverse tortures, gave up his soul in prayer to God. The execution of Agricola had been put off, in the hope that the agonies of his servant might scare him into denying Christ; but the sight only hardened him. He was therefore crucified, and so became sharer and fellow with his slave Vitalis in the glory of testification. Their bodies were laid in the Jews' burying-place, where they were found by St. Ambrose, who removed them to an hallowed and honourable sepulchre.\par
\switchcolumn*

\LA
\TeDeum{Te Deum}
\switchcolumn
\EN
\TeDeum{Te Deum}
\switchcolumn*

\end{paracol}

\par\needspace{10\baselineskip}\addvspace{1.2em}{\centering\small\textsc{Die 5 Novembris} · \textsc{5 November}\par}
\Office{Quinta die infra Octavam Omnium Sanctorum}{Fifth Day Within the Octave of All Saints}{Semiduplex}
\addcontentsline{toc}{subsection}{Die 5 Novembris · Quinta die infra Octavam Omnium Sanctorum \textit{— Fifth Day Within the Octave of All Saints}}
\begin{paracol}{2}
\LA
\RefLine{Lectiones I–III: de Scriptura occurrente.}
\switchcolumn
\EN
\RefLine{Lessons I–III: from the Scripture of the day.}
\switchcolumn*

\LA
\LessonHead{Lectio IV}
\switchcolumn
\EN
\LessonHead{Lesson IV}
\switchcolumn*

\LA
\LessonTitle{De Sermone sancti Bedæ Venerabilis Presbýteri}
\switchcolumn
\EN
\LessonTitle{From the Sermons of the Venerable Bede, Priest (at Jarrow.)}
\switchcolumn*

\LA
\Source{Ex Sermone 18 de Sanctis}
\switchcolumn
\EN
\Source{18th on the Saints.}
\switchcolumn*

\LA
\noindent Ad hanc igitur operum salutarium delectet nos pervenire palmam. Libenter ac prompte certemus, omnes in agone justitiæ Deo et Christo spectante curramus; et qui sæculo et mundo majores esse jam cœpimus, cursum nostrum nulla sæculi cupiditate tardemus. Si expeditos, si celeres in operis agone currentes dies nos ultimus invenerit, nusquam Dominus meritis nostris deerit remunerator.\par
\switchcolumn
\EN
\noindent Therefore, may it be our delight to go on unto this prize of good living. Freely and cheerfully let us strive in the race, running under the eyes of God and of Christ. We have already taken a station above floating and earthly things, and let us allow no love for things fleeting to hamper our running. If the last day shall find us lithe and speedful in the race of good living, we shall never have to complain that our Master is a scanty rewarder of our works.\par
\switchcolumn*

\LA
\begin{Resp}\Rsign Præcúrsor Dómini venit, de quo ipse testátur: \Star{} Nullus major inter natos mulíerum Joánne Baptísta.\end{Resp}
\begin{Resp}\Vsign Hic est enim prophéta, et plus quam prophéta, de quo Salvátor ait.\end{Resp}
\begin{Resp}\Rsign Nullus major inter natos mulíerum Joánne Baptísta.\end{Resp}
\switchcolumn
\EN
\begin{Resp}\Rsign The Fore-runner of the Lord cometh, to whom He Himself bare witness, saying: \Star{} Among them that are born of women there hath not risen a greater than John the Baptist.\end{Resp}
\begin{Resp}\Vsign A Prophet? Yea, and much more than a Prophet. This is he of whom the Saviour saith:\end{Resp}
\begin{Resp}\Rsign Among them that are born of women there hath not risen a greater than John the Baptist.\end{Resp}
\switchcolumn*

\LA
\LessonHead{Lectio V}
\switchcolumn
\EN
\LessonHead{Lesson V}
\switchcolumn*

\LA
\noindent Qui coronam in persecutione purpuream pro passione donabit, ipse in pace vincentibus pro justitiæ meritis dabit et candidam. Nam nec Abraham, nec Isaac, nec Jacob occisi sunt; et tamen fidei et justitiæ meritis honorati inter patriarchas primi esse meruerunt: ad quorum congregatur convivium quisquis fidelis, et justus, et laudabilis invenitur. Memores esse debemus, voluntatem non nostram, sed Dei facere debere: quia qui fecerit ejus voluntatem, manet in æternum, quomodo et ille manet in æternum.\par
\switchcolumn
\EN
\noindent That giveth a red crown for suffering under persecution, the same giveth a white crown to them that under peace, prevail in battles of righteousness. Neither Abraham, nor Isaac, nor Jacob, were slain, and nevertheless in honour for faith and righteousness, they have gained the first place among the Patriarchs, and it is to sit down with them in the kingdom of God that are gathered the faithful, the righteous, and the praiseworthy. We must remember that it is God's will, and not our own will, that we must do, for he that doeth His will abideth for ever, even as He abideth for ever.\par
\switchcolumn*

\LA
\begin{Resp}\Rsign Isti sunt qui vivéntes in carne, plantavérunt Ecclésiam sánguine suo: \Star{} Cálicem Dómini bibérunt, et amíci Dei facti sunt.\end{Resp}
\begin{Resp}\Vsign In omnem terram exívit sonus eórum, et in fines orbis terræ verba eórum.\end{Resp}
\begin{Resp}\Rsign Cálicem Dómini bibérunt, et amíci Dei facti sunt.\end{Resp}
\switchcolumn
\EN
\begin{Resp}\Rsign These are they who while yet they lived in the flesh, planted the Church in their own blood; \Star{} They drank of the Lord's cup, and became the friends of God.\end{Resp}
\begin{Resp}\Vsign Their sound is gone out through all the earth, and their words to the ends of the world.\end{Resp}
\begin{Resp}\Rsign They drank of the Lord's cup, and became the friends of God.\end{Resp}
\switchcolumn*

\LA
\LessonHead{Lectio VI}
\switchcolumn
\EN
\LessonHead{Lesson VI}
\switchcolumn*

\LA
\noindent Quapropter, carissimi, mente integra, fide firma, virtute robusta, caritate perfecta, parati ad omnem voluntatem Dei simus, conservantes fortiter Dominica mandata, in simplicitate innocentiam, in caritate concordiam, modestiam in humilitate, diligentiam in administratione, vigilantiam in adjuvandis laborantibus, misericordiam in fovendis pauperibus, in defendenda veritate constantiam, in disciplinæ severitate censuram, ne aliquid ad exemplum bonorum factorum desit in nobis. Hæc sunt enim vestigia quæ nobis Sancti quique revertentes in patriam reliquerunt, ut illorum semitis inhærentes, sequeremur et gaudia.\par
\switchcolumn
\EN
\noindent Therefore, dearly beloved brethren, with mind clear, faith firm, courage true, love thorough, let us be ready to do whatever God willeth, keeping stoutly all the commandments of the Lord, having innocency in simplicity, peaceableness in love, modesty in lowliness, in ministering diligence, in helping them that toil watchfulness, in succouring the poor mercifulness, in standing up for the truth firmness, in keeping of discipline sternness, lest we be found wanting in any good work. These are the steps which the Saints who have already gone home have left marked for us, that we may be able to keep in their footprints, and so to follow them into their joy.\par
\switchcolumn*

\LA
\begin{Resp}\Rsign Sancti mei, qui in carne pósiti, certámen habuístis: \Star{} Mercédem labóris ego reddam vobis.\end{Resp}
\begin{Resp}\Vsign Veníte benedícti Patris mei, percípite regnum.\end{Resp}
\begin{Resp}\Rsign Mercédem labóris ego reddam vobis.\end{Resp}
\begin{Resp}\Vsign Glória Patri, et Fílio, \Star{} et Spirítui Sancto.\end{Resp}
\begin{Resp}\Rsign Mercédem labóris ego reddam vobis.\end{Resp}
\switchcolumn
\EN
\begin{Resp}\Rsign O ye My Saints, who, being in the flesh, didst have striving \Star{} I will render unto you a reward of your labours.\end{Resp}
\begin{Resp}\Vsign Come, ye blessed of My Father, inherit the kingdom\end{Resp}
\begin{Resp}\Rsign I will render unto you a reward of your labours.\end{Resp}
\begin{Resp}\Vsign Glory be to the Father, and to the Son, \Star{} and to the Holy Ghost.\end{Resp}
\begin{Resp}\Rsign I will render unto you a reward of your labours.\end{Resp}
\switchcolumn*

\LA
\LessonHead{Lectio VII}
\switchcolumn
\EN
\LessonHead{Lesson VII}
\switchcolumn*

\LA
\LessonTitle{Léctio sancti Evangélii secúndum Matthǽum}
\switchcolumn
\EN
\LessonTitle{From the Holy Gospel according to Matthew}
\switchcolumn*

\LA
\Cite{Matt 5:1-12}
\switchcolumn
\EN
\Cite{Matt 5:1-12}
\switchcolumn*

\LA
\noindent In illo témpore: Videns Jesus turbas, ascendit in montem, et cum sedisset, accesserunt ad eum discipuli ejus. Et réliqua.\par
\switchcolumn
\EN
\noindent At that time Jesus seeing the multitudes, he went up into a mountain, and when he was set down, his disciples came unto him. And so on.\par
\switchcolumn*

\LA
\LessonTitle{De Homilia sancti Augustíni Epíscopi}
\switchcolumn
\EN
\LessonTitle{Homily by St. Augustine, Bishop (of Hippo.)}
\switchcolumn*

\LA
\Source{Liber 1 de Sermone Domini in monte, cap. 3 et 4}
\switchcolumn
\EN
\Source{Bk. i. on the Lord's Sermon, Ch. 4}
\switchcolumn*

\LA
\noindent Itaque in hoc tertio gradu, in quo scientia est, lugetur amissio summi boni, quia inhæretur extremis. In quarto autem gradu labor est: ubi vehementer incumbitur, ut sese animus avellat ab eis quibus pestifera dulcedine innexus est. Hic ergo esuritur et sititur justitia, et multum necessaria est fortitudo: quia non relinquitur sine dolore, quod cum delectatione retinetur. Quinto autem gradu perseverantibus in labore datur evadendi consilium: quia, nisi quisque adjuvetur a superiore, nullo modo sibi est idoneus, ut sese tantis miseriarum implicamentis expediat. Est autem justum consilium, ut qui se a potentiori adjuvari vult, adjuvet et infirmiorem, in quo est ipse potentior. Itaque, Beati misericordes: quia ipsorum miserebitur Deus.\par
\switchcolumn
\EN
\noindent First, Blessed are the poor in spirit. Secondly, Blessed are the meek. Thirdly, Blessed are they that mourn. They that are blessed under this third head, having knowledge, do mourn that they possess not yet the Highest Good, which possession belongeth unto the end of their course. But in the fourth place, Blessed are they which do hunger and thirst after righteousness. Here there is that earnest striving, wherewith the mind doth struggle to tear herself away from those things whose deathful sweetness would make her fain to cling unto them. Here is hungering and thirsting after righteousness, and there is sore need of firmness, for what it is a joy to have, it must be a grief to lose. But the fifth head is the declaration that Blessed are the merciful, and in these words a door of comfort and reward is opened unto the toiling. Entangled in such straits a man can be of no use to himself, unless One That is stronger than he help him; and if he be helped of the Stronger, it is but just that he in turn should help such as is weaker than himself. And so, Blessed are the merciful, for, in their turn, they shall obtain mercy from God.\par
\switchcolumn*

\LA
\begin{Resp}\Rsign Sint lumbi vestri præcíncti, et lucérnæ ardéntes in mánibus vestris: \Star{} Et vos símiles homínibus exspectántibus dóminum suum, quando revertátur a núptiis.\end{Resp}
\begin{Resp}\Vsign Vigiláte ergo, quia nescítis qua hora Dóminus vester ventúrus sit.\end{Resp}
\begin{Resp}\Rsign Et vos símiles homínibus exspectántibus dóminum suum, quando revertátur a núptiis.\end{Resp}
\switchcolumn
\EN
\begin{Resp}\Rsign Let your loins be girded about, and your lights burning; \Star{} And ye yourselves like unto men that wait for their lord, when he will return from the wedding.\end{Resp}
\begin{Resp}\Vsign Watch therefore, for ye know not what hour your Lord doth come.\end{Resp}
\begin{Resp}\Rsign And ye yourselves like unto men that wait for their lord, when he will return from the wedding.\end{Resp}
\switchcolumn*

\LA
\LessonHead{Lectio VIII}
\switchcolumn
\EN
\LessonHead{Lesson VIII}
\switchcolumn*

\LA
\noindent Sexto gradu est cordis munditia de bona conscientia bonorum operum, valens ad contemplandum summum illud bonum, quod solo puro et sereno intellectu cerni potest. Postremo est septima ipsa sapientia, id est, contemplatio veritatis, pacificans totum hominem, et suscipiens similitudinem Dei, quæ ita concluditur: Beati pacifici: quoniam ipsi filii Dei vocabuntur. Octava tamquam ad caput redit, quia consummatum perfectumque ostendit et probat. Itaque in prima et in octava nominatum est regnum cælorum: Beati pauperes spiritu: quoniam ipsorum est regnum cælorum; et, Beati qui persecutionem patiuntur propter justitiam: quoniam ipsorum est regnum cælorum.\par
\switchcolumn
\EN
\noindent Blessed are the pure in heart. This sixth benediction is pronounced upon those hearts which by pure, clear consciousness of good works are able to look to that Highest Good, Which only the clear, calm mind can perceive. Lastly cometh in the seventh place that Blessed are the peacemakers, that is to say, blessed are they who cultivate wisdom, which is the contemplation of the True, since it is the fruit of this contemplation of the True to produce profound and utter internal peace in man, and to catch the reflection of the Divine, this being the idea which is expressed in the words: Blessed are the peacemakers, for they shall be called the children of God. The eighth phrase is a return to the first, since it showeth lowliness of spirit in its aspect of completion and crowning; and thence the kingdom of heaven is the reward mentioned in both places. Blessed are the poor in spirit, for their's is the kingdom of heaven. Blessed are they which are persecuted for righteousness' sake, for their's is the kingdom of heaven.\par
\switchcolumn*

\LA
\begin{Resp}\Rsign Média nocte clamor factus est: \Star{} Ecce sponsus venit, exíte óbviam ei.\end{Resp}
\begin{Resp}\Vsign Prudéntes vírgines, aptáte vestras lámpades.\end{Resp}
\begin{Resp}\Rsign Ecce sponsus venit, exíte óbviam ei.\end{Resp}
\begin{Resp}\Vsign Glória Patri, et Fílio, \Star{} et Spirítui Sancto.\end{Resp}
\begin{Resp}\Rsign Ecce sponsus venit, exíte óbviam ei.\end{Resp}
\switchcolumn
\EN
\begin{Resp}\Rsign At midnight there was a cry made: \Star{} Behold, the Bridegroom cometh, go ye out to meet him.\end{Resp}
\begin{Resp}\Vsign Trim your lamps, O ye wise virgins.\end{Resp}
\begin{Resp}\Rsign Behold, the Bridegroom cometh, go ye out to meet him.\end{Resp}
\begin{Resp}\Vsign Glory be to the Father, and to the Son, \Star{} and to the Holy Ghost.\end{Resp}
\begin{Resp}\Rsign Behold, the Bridegroom cometh, go ye out to meet him.\end{Resp}
\switchcolumn*

\LA
\LessonHead{Lectio IX}
\switchcolumn
\EN
\LessonHead{Lesson IX}
\switchcolumn*

\LA
\noindent Cum jam dicitur: Quis nos separabit a caritate Christi? Tribulatio, an angustia, an persecutio, an fames, an nuditas, an periculum, an gladius? Septem sunt ergo, quæ perficiunt: nam octava clarificat, et quod perfectum est, demonstrat, ut per hos gradus perficiantur et ceteri, tamquam a capite rursum exordiens. Videtur ergo mihi étiam septiformis operatio Spiritus sancti, de qua Isaias loquitur, his gradibus sententiisque congruere. Sed interest ordinis: nam ibi enumeratio ab excellentioribus cœpit, hic vero ab inferioribus. Ibi namque incipit a sapientia Dei, et desinit ad timorem Dei: sed initium sapientiæ timor Domini est.\par
\switchcolumn
\EN
\noindent Paul saith Who shall separate us from the love of Christ? Shall tribulation, or distress, or persecution, or famine, or nakedness, or peril, or sword? There are therefore seven things which bring to perfection, for the eighth is the glorification and manifestation of that which is perfected, that from this head others again may begin, and be finished. It seemeth to me also that these heads and sayings have some connection with the seven gifts of the Holy Ghost whereof Isaiah speaketh. But there is a difference of order, for there the highest is taken first, but here the lowest; there the wisdom of God, but here the fear of God, but the beginning of wisdom is the fear of the Lord.\par
\switchcolumn*

\LA
\TeDeum{Te Deum}
\switchcolumn
\EN
\TeDeum{Te Deum}
\switchcolumn*

\end{paracol}

\par\needspace{10\baselineskip}\addvspace{1.2em}{\centering\small\textsc{Die 6 Novembris} · \textsc{6 November}\par}
\Office{Sexta die infra Octavam Omnium Sanctorum}{Sixth Day Within the Octave of All Saints}{Semiduplex}
\addcontentsline{toc}{subsection}{Die 6 Novembris · Sexta die infra Octavam Omnium Sanctorum \textit{— Sixth Day Within the Octave of All Saints}}
\begin{paracol}{2}
\LA
\RefLine{Lectiones I–III: de Scriptura occurrente.}
\switchcolumn
\EN
\RefLine{Lessons I–III: from the Scripture of the day.}
\switchcolumn*

\LA
\LessonHead{Lectio IV}
\switchcolumn
\EN
\LessonHead{Lesson IV}
\switchcolumn*

\LA
\LessonTitle{Sermo sancti Bernardi Abbatis}
\switchcolumn
\EN
\LessonTitle{From the Sermons of St. Bernard, Abbot (of Clairvaux)}
\switchcolumn*

\LA
\Source{Serm. 2 de Festo Omn. Sanct.}
\switchcolumn
\EN
\Source{2nd for All Saints' Day}
\switchcolumn*

\LA
\noindent Quia Sanctorum omnium festivam hodie, dilectissimi, omnique dignissimam devotione memoriam celebramus, operæ pretium puto, de communi eorum felicitate, in qua beata jam requie perfruuntur, et futura quam præstolantur consummatione adjuvante Spiritu sancto, sermonem facere caritati vestræ. Fidelis quippe sermo, et omni acceptione dignus: ut quos solemni veneratione prosequimur, étiam simili conversatione sequamur: quos beatissimos prædicamus, ad eorum beatitudinem tota aviditate curramus: quorum delectamur præconiis, sublevemur eorum patrociniis.\par
\switchcolumn
\EN
\noindent Dearly beloved brethren, since we keep on this day the memory of all the Saints, that memory so joyous and so worthy of all our thoughts, it seemeth to me worth the while, the Holy Ghost helping me, to address to your kind indulgence some remarks upon that happiness which they are all enjoying in blessed restfulness, and that final consummation they are awaiting. It is a faithful saying, and worthy of all acceptation, that if we thus solemnly honour them, we should follow the example of their conversation; if we proclaim them so blessed, we should strive our best to reach the same blessedness; if we are well pleased to hear them praised, we should be bettered by their prayers.\par
\switchcolumn*

\LA
\begin{Resp}\Rsign Præcúrsor Dómini venit, de quo ipse testátur: \Star{} Nullus major inter natos mulíerum Joánne Baptísta.\end{Resp}
\begin{Resp}\Vsign Hic est enim prophéta, et plus quam prophéta, de quo Salvátor ait.\end{Resp}
\begin{Resp}\Rsign Nullus major inter natos mulíerum Joánne Baptísta.\end{Resp}
\switchcolumn
\EN
\begin{Resp}\Rsign The Fore-runner of the Lord cometh, to whom He Himself bare witness, saying: \Star{} Among them that are born of women there hath not risen a greater than John the Baptist.\end{Resp}
\begin{Resp}\Vsign A Prophet? Yea, and much more than a Prophet. This is he of whom the Saviour saith:\end{Resp}
\begin{Resp}\Rsign Among them that are born of women there hath not risen a greater than John the Baptist.\end{Resp}
\switchcolumn*

\LA
\LessonHead{Lectio V}
\switchcolumn
\EN
\LessonHead{Lesson V}
\switchcolumn*

\LA
\Source{Sermo 5. de eodem Festo, circa medium.}
\switchcolumn
\EN

\switchcolumn*

\LA
\noindent Ad quid ergo Sanctis laus nostra? ad quid glorificatio nostra? ad quid nostra hæc ipsa solemnitas? Quo eis terrenos honores, quos juxta veracem Filii promissionem honorificat Pater cælestis? Quo eis præconia nostra? Pleni sunt. Prorsus ita est, dilectissimi: bonorum nostrorum Sancti non egent, nec quidquam eis nostra devotione præstatur. Plane, quod eorum memoriam veneramur, nostra interest, non ipsorum. Vultis scire quantum interest nostra? Ego in me (fateor) ex hac recordatione sentio desiderium vehemens inflammari, et desiderium triplex.\par
\switchcolumn
\EN
\noindent What is it to the Saints that we should praise them? What to them that we should glorify them? What is this our Feast to them? What are honours on earth to them whom, according as the Son hath faithfully promised, His Father is honouring? What are our eulogies to them? They are full. Verily, dearly beloved brethren, of our goods the Saints have no need, and our devotion toward them doth nothing for them. Our honouring their memory hath to do with ourselves and not with them. Would ye know what it hath to do with us? In me I confess that at their remembrance I feel kindled a vehement longing, yea, a three-fold longing.\par
\switchcolumn*

\LA
\begin{Resp}\Rsign Isti sunt qui vivéntes in carne, plantavérunt Ecclésiam sánguine suo: \Star{} Cálicem Dómini bibérunt, et amíci Dei facti sunt.\end{Resp}
\begin{Resp}\Vsign In omnem terram exívit sonus eórum, et in fines orbis terræ verba eórum.\end{Resp}
\begin{Resp}\Rsign Cálicem Dómini bibérunt, et amíci Dei facti sunt.\end{Resp}
\switchcolumn
\EN
\begin{Resp}\Rsign These are they who while yet they lived in the flesh, planted the Church in their own blood; \Star{} They drank of the Lord's cup, and became the friends of God.\end{Resp}
\begin{Resp}\Vsign Their sound is gone out through all the earth, and their words to the ends of the world.\end{Resp}
\begin{Resp}\Rsign They drank of the Lord's cup, and became the friends of God.\end{Resp}
\switchcolumn*

\LA
\LessonHead{Lectio VI}
\switchcolumn
\EN
\LessonHead{Lesson VI}
\switchcolumn*

\LA
\noindent Vulgo dicitur: Quod non videt oculus, cor non dolet. Oculus meus, memoria mea: et cogitare de Sanctis, quodammodo eos videre est. Sic nempe portio nostra in terra viventium; nec modica sane portio, si tamen, ut decet, memoriam affectio comitetur: sic, inquam, conversatio nostra in cælis est. Verumtamen non sic nostra, sicut illorum. Ipsorum enim substantia ibi est, nostra autem desideria: ipsi per præsentiam, nos per memoriam ibi sumus.\par
\switchcolumn
\EN
\noindent It is a common saying that out of sight, out of mind. The memory is a kind of sight, and to think of the Saints, is to call them up before the mind's eye. Such is our portion in the land of the living, but it is not a little portion, if love, (as it ought to do,) be joined with remembrance; it is in such sense that we must say that our conversation is in heaven. (Phil. iii. 20.) Very differently to what is theirs. For they are there actually, where we are only in desire; they in very presence, we only in thought.\par
\switchcolumn*

\LA
\begin{Resp}\Rsign Sancti mei, qui in carne pósiti, certámen habuístis: \Star{} Mercédem labóris ego reddam vobis.\end{Resp}
\begin{Resp}\Vsign Veníte benedícti Patris mei, percípite regnum.\end{Resp}
\begin{Resp}\Rsign Mercédem labóris ego reddam vobis.\end{Resp}
\begin{Resp}\Vsign Glória Patri, et Fílio, \Star{} et Spirítui Sancto.\end{Resp}
\begin{Resp}\Rsign Mercédem labóris ego reddam vobis.\end{Resp}
\switchcolumn
\EN
\begin{Resp}\Rsign O ye My Saints, who, being in the flesh, didst have striving \Star{} I will render unto you a reward of your labours.\end{Resp}
\begin{Resp}\Vsign Come, ye blessed of My Father, inherit the kingdom\end{Resp}
\begin{Resp}\Rsign I will render unto you a reward of your labours.\end{Resp}
\begin{Resp}\Vsign Glory be to the Father, and to the Son, \Star{} and to the Holy Ghost.\end{Resp}
\begin{Resp}\Rsign I will render unto you a reward of your labours.\end{Resp}
\switchcolumn*

\LA
\LessonHead{Lectio VII}
\switchcolumn
\EN
\LessonHead{Lesson VII}
\switchcolumn*

\LA
\LessonTitle{Léctio sancti Evangélii secúndum Matthǽum}
\switchcolumn
\EN
\LessonTitle{From the Holy Gospel according to Matthew}
\switchcolumn*

\LA
\Cite{Matt 5:1-12}
\switchcolumn
\EN
\Cite{Matt 5:1-12}
\switchcolumn*

\LA
\noindent In illo témpore: Videns Jesus turbas, ascendit in montem, et cum sedisset, accesserunt ad eum discipuli ejus. Et réliqua.\par
\switchcolumn
\EN
\noindent At that time, Jesus seeing the multitudes, went up into a mountain, and when he was set down, his disciples came unto him. And so on.\par
\switchcolumn*

\LA
\LessonTitle{De Homilia sancti Augustíni Epíscopi}
\switchcolumn
\EN
\LessonTitle{Homily by St. Augustine, Bishop (of Hippo.)}
\switchcolumn*

\LA
\Source{Liber 1. de Sermone Domini in monte, Cap. 4.}
\switchcolumn
\EN
\Source{Bk. i. on the Lord's Sermon. Ch. 4}
\switchcolumn*

\LA
\noindent Quapropter, si gradatim tamquam ascendentes numeremus, primus ibi est timor Dei, secunda pietas, tertia scientia, quarta fortitudo, quintum consilium, sextus intellectus, septima sapientia. Timor Dei congruit humilibus, de quibus hic dicitur: Beati pauperes spiritu: quoniam ipsorum est regnum cælorum; id est, non inflati, non superbi, de quibus Apostolus dicit: Noli altum sapere, sed time; id est, noli extolli. Pietas congruit mitibus. Qui enim pie quærit, honorat sanctam Scripturam, et non reprehendit quod nondum intelligit, et propterea non resistit; quod est mitem esse. Unde hic dicitur: Beati mites: quoniam ipsi hereditate possidebunt terram.\par
\switchcolumn
\EN
\noindent Wherefore, if we reckon up the Beatitudes as ascending steps, the first is the fear of God; the second, godliness; the third, knowledge; the fourth, firmness; the fifth, counsel; the sixth, understanding; the seventh, wisdom. The fear of God pertaineth unto the lowly, as it is said Blessed are the poor in spirit, for their's is the kingdom of heaven, that is, it is for them that are not puffed up, for them that are not proud, as also saith the Apostle Be not highminded, but fear, (Rom. xi. 20,) that is, Be not puffed up. Godliness pertaineth unto the meek; for he that seeketh after a godly sort, honoureth the Holy Scripture, and when he findeth therein that which he doth not yet understand he blameth not the Scripture, nor gainstandeth. And this is to be meek. Therefore is it said here Blessed are the meek, for they shall inherit the earth.\par
\switchcolumn*

\LA
\begin{Resp}\Rsign Sint lumbi vestri præcíncti, et lucérnæ ardéntes in mánibus vestris: \Star{} Et vos símiles homínibus exspectántibus dóminum suum, quando revertátur a núptiis.\end{Resp}
\begin{Resp}\Vsign Vigiláte ergo, quia nescítis qua hora Dóminus vester ventúrus sit.\end{Resp}
\begin{Resp}\Rsign Et vos símiles homínibus exspectántibus dóminum suum, quando revertátur a núptiis.\end{Resp}
\switchcolumn
\EN
\begin{Resp}\Rsign Let your loins be girded about, and your lights burning; \Star{} And ye yourselves like unto men that wait for their lord, when he will return from the wedding.\end{Resp}
\begin{Resp}\Vsign Watch therefore, for ye know not what hour your Lord doth come.\end{Resp}
\begin{Resp}\Rsign And ye yourselves like unto men that wait for their lord, when he will return from the wedding.\end{Resp}
\switchcolumn*

\LA
\LessonHead{Lectio VIII}
\switchcolumn
\EN
\LessonHead{Lesson VIII}
\switchcolumn*

\LA
\noindent Scientia congruit lugentibus, qui jam cognoverunt in Scripturis, quibus malis vincti teneantur, quæ tamquam bona et utilia ignorantes appetierunt; de quibus hic dicitur: Beati qui lugent nunc. Fortitudo congruit esurientibus et sitientibus: laborant enim desiderantes gaudium de veris bonis, et amórem a terrenis et corporalibus avertere cupientes; de quibus hic dicitur: Beati qui esuriunt et sitiunt justitiam. Consilium congruit misericordibus: hoc enim unum remedium est de tantis malis evadendi, ut dimittamus, sicut nobis dimitti volumus: et adjuvemus, in quo possumus, alios, sicut et nos, in quo non possumus, cupimus adjuvari; de quibus hic dicitur: Beati misericordes: quoniam ipsorum miserebitur Deus.\par
\switchcolumn
\EN
\noindent Is knowledge pertaineth unto them that mourn, who have already learnt from the Scriptures amid what ills they are entangled, even in those things which once in their ignorance they affected as being good and useful. Of such is it said Blessed are ye that weep now. (Luke vi. 21.) Firmness pertaineth unto such as hunger and thirst after righteousness. These are they who toil bravely, animated by the longing for that joy which is caused by real blessedness, and striving therefore to wean their love away from so-called joys whose origin is merely earthly and fleshly. Of them is it said Blessed are they which do hunger and thirst after righteousness. Counsel pertaineth unto the merciful, for our only way of escape from the horrors of our own guilt's punishment is that we should forgive even as we hope to be ourselves forgiven, and should help others as much as we can, even as we would fain be holpen in that wherein we can ourselves do nothing. And of such as so do, it is said Blessed are the merciful, for they shall obtain mercy from God.\par
\switchcolumn*

\LA
\begin{Resp}\Rsign Média nocte clamor factus est: \Star{} Ecce sponsus venit, exíte óbviam ei.\end{Resp}
\begin{Resp}\Vsign Prudéntes vírgines, aptáte vestras lámpades.\end{Resp}
\begin{Resp}\Rsign Ecce sponsus venit, exíte óbviam ei.\end{Resp}
\begin{Resp}\Vsign Glória Patri, et Fílio, \Star{} et Spirítui Sancto.\end{Resp}
\begin{Resp}\Rsign Ecce sponsus venit, exíte óbviam ei.\end{Resp}
\switchcolumn
\EN
\begin{Resp}\Rsign At midnight there was a cry made: \Star{} Behold, the Bridegroom cometh, go ye out to meet him.\end{Resp}
\begin{Resp}\Vsign Trim your lamps, O ye wise virgins.\end{Resp}
\begin{Resp}\Rsign Behold, the Bridegroom cometh, go ye out to meet him.\end{Resp}
\begin{Resp}\Vsign Glory be to the Father, and to the Son, \Star{} and to the Holy Ghost.\end{Resp}
\begin{Resp}\Rsign Behold, the Bridegroom cometh, go ye out to meet him.\end{Resp}
\switchcolumn*

\LA
\LessonHead{Lectio IX}
\switchcolumn
\EN
\LessonHead{Lesson IX}
\switchcolumn*

\LA
\noindent Intellectus congruit mundis corde, tamquam purgato oculo, quo cerni possit, quod corporeus oculus non vidit, nec auris audivit, nec in cor hominis ascendit; de quibus hic dicitur: Beati mundo corde: quoniam ipsi Deum videbunt. Sapientia congruit pacificis, in quibus jam ordinata sunt omnia, nullusque motus adversus rationem rebellis est, sed cuncta obtemperant spiritui hominis cum et ipse obtemperet Deo; de quibus hic dicitur: Beati pacifici. Unum autem præmium, quod est regnum cælorum, pro his gradibus varie nominatum est.\par
\switchcolumn
\EN
\noindent Understanding pertaineth unto the pure in heart, for these are they whose clear eye can see that which the fleshly eye hath not seen, neither the ear heard, neither hath it entered into the heart of man to conceive; and therefore of them it is said: Blessed are the pure in heart, for they shall see God. Wisdom pertaineth unto the peacemakers, even unto them in whom all things are well ordered, and passion no longer maketh insurrection against reason, but all things are subject unto human common sense, even as the same again is made subject unto God. And of such is it said: Blessed are the peacemakers. But for all these forms of blessedness there is one and the same reward, although diversely named, and that reward is the kingdom of heaven.\par
\switchcolumn*

\LA
\TeDeum{Te Deum}
\switchcolumn
\EN
\TeDeum{Te Deum}
\switchcolumn*

\end{paracol}

\par\needspace{10\baselineskip}\addvspace{1.2em}{\centering\small\textsc{Die 7 Novembris} · \textsc{7 November}\par}
\Office{Septima die infra Octavam Omnium Sanctorum}{Seventh Day Within the Octave of All Saints}{Semiduplex}
\addcontentsline{toc}{subsection}{Die 7 Novembris · Septima die infra Octavam Omnium Sanctorum \textit{— Seventh Day Within the Octave of All Saints}}
\begin{paracol}{2}
\LA
\RefLine{Lectiones I–III: de Scriptura occurrente.}
\switchcolumn
\EN
\RefLine{Lessons I–III: from the Scripture of the day.}
\switchcolumn*

\LA
\LessonHead{Lectio IV}
\switchcolumn
\EN
\LessonHead{Lesson IV}
\switchcolumn*

\LA
\LessonTitle{Sermo sancti Joannis Chrysostomi}
\switchcolumn
\EN
\LessonTitle{From the Sermons of St. John Chrysostom, (Patriarch of Constantinople.)}
\switchcolumn*

\LA
\Source{Serm. de Mart, quod aut imitandi sunt, aut non laudandi. Tom. 3}
\switchcolumn
\EN
\Source{On the Martyrs.}
\switchcolumn*

\LA
\noindent Qui Sanctorum merita religiosa caritate miratur, quique justorum glorias frequenti laude colloquitur, eorum mores sanctos atque justitiam imitetur: quoniam quem delectat Sancti alicujus meritum, delectare debet par circa cultum Dei obsequium. Quare aut imitari debet, si laudat: aut laudare non debet, si imitari detrectat: ut qui alium laudat, laudabilem se reddat: et qui Sanctorum merita admiratur, mirabilis ipse vitæ sanctitate reddatur. Nam si propterea justos fidelesque diligimus, quod in ipsis justitiam fidemque suspicimus, possumus nos quoque esse quod sunt, si faciamus ipsi quod faciunt.\par
\switchcolumn
\EN
\noindent He that wondereth with reverential love at the mighty deeds of the holy, he that hath oftentimes on his tongue praises for the glory of the righteous, let such an one copy their holy lives and their righteousness; for if any take pleasure in the work of a Saint, he ought to take pleasure in serving God as that Saint served Him. If he praiseth the Saint, he ought to imitate him, and if he is not ready to imitate him, he ought not to praise him. Let him that praiseth another make himself worthy of a like praise, and if he be in admiration of the Saints, let his own admirable life reflect the holiness of theirs. If we love the good and leal because they are good and leal, let us not forget that we can be what they are, by doing as they did.\par
\switchcolumn*

\LA
\begin{Resp}\Rsign Præcúrsor Dómini venit, de quo ipse testátur: \Star{} Nullus major inter natos mulíerum Joánne Baptísta.\end{Resp}
\begin{Resp}\Vsign Hic est enim prophéta, et plus quam prophéta, de quo Salvátor ait.\end{Resp}
\begin{Resp}\Rsign Nullus major inter natos mulíerum Joánne Baptísta.\end{Resp}
\switchcolumn
\EN
\begin{Resp}\Rsign The Fore-runner of the Lord cometh, to whom He Himself bare witness, saying: \Star{} Among them that are born of women there hath not risen a greater than John the Baptist.\end{Resp}
\begin{Resp}\Vsign A Prophet? Yea, and much more than a Prophet. This is he of whom the Saviour saith:\end{Resp}
\begin{Resp}\Rsign Among them that are born of women there hath not risen a greater than John the Baptist.\end{Resp}
\switchcolumn*

\LA
\LessonHead{Lectio V}
\switchcolumn
\EN
\LessonHead{Lesson V}
\switchcolumn*

\LA
\noindent Neque enim difficile nobis est, quod ab ipsis geritur, imitari, cum sine præcedenti exemplo ab antiquis talia gesta conspicimus: ut non ipsi aliorum æmuli redderentur, sed æmulandæ virtutis seipsos nobis præberent exemplum: ut dum nos ex ipsis, et ex nobis alii proficiunt, sic Christus in suis semper servis in Ecclesia sancta laudetur. Unde ab origine mundi innocens Abel occiditur, Enoch Deo placens transfertur, justus Noë invenitur, Abraham fidelis probatur, Moyses mansuetus dignoscitur, Jesus castus, David lenis, Elias acceptus, Daniel sanctus, tres pueri victores redduntur.\par
\switchcolumn
\EN
\noindent It ought not to be hard for us to copy others, when we see what they of old time did without any examples before them, so that in them who copied not others, but set example for others to copy, and in us who copy them, and in them which take example by us, Christ may be glorified in His holy Church. Thus from the very beginning of the world there have been the harmless Abel who was slain, Enoch who walked with God, and was seen no more, for God took him, Noah who was found righteous, Abraham who was tried and found faithful, Moses who was the meekest of men, Joshua who was chaste, David who was gentle, Elijah who was accepted, Daniel who was holy, and the three Children who were victorious.\par
\switchcolumn*

\LA
\begin{Resp}\Rsign Isti sunt qui vivéntes in carne, plantavérunt Ecclésiam sánguine suo: \Star{} Cálicem Dómini bibérunt, et amíci Dei facti sunt.\end{Resp}
\begin{Resp}\Vsign In omnem terram exívit sonus eórum, et in fines orbis terræ verba eórum.\end{Resp}
\begin{Resp}\Rsign Cálicem Dómini bibérunt, et amíci Dei facti sunt.\end{Resp}
\switchcolumn
\EN
\begin{Resp}\Rsign These are they who while yet they lived in the flesh, planted the Church in their own blood; \Star{} They drank of the Lord's cup, and became the friends of God.\end{Resp}
\begin{Resp}\Vsign Their sound is gone out through all the earth, and their words to the ends of the world.\end{Resp}
\begin{Resp}\Rsign They drank of the Lord's cup, and became the friends of God.\end{Resp}
\switchcolumn*

\LA
\LessonHead{Lectio VI}
\switchcolumn
\EN
\LessonHead{Lesson VI}
\switchcolumn*

\LA
\noindent Apostoli discipuli Christi, credentium magistri habentur: e quibus eruditi Confessores fortissimi pugnant, Martyres perfecti triumphant, et christiani semper exercitus diabolum Deo armati debellant. In istis semper pares virtutes, dissimiles pugnæ, gloriosæ victoriæ. Unde tu Christiane delicatus es miles, si putas te posse sine pugna vincere, sine certamine triumphare. Exsere vires, fortiter dimica, atrociter in prælio isto concerta. Considera pactum, conditionem attende, militiam nosce: pactum quod spopondisti; conditionem qua accessisti: militiam cui nomen dedisti.\par
\switchcolumn
\EN
\noindent The Apostles, the disciples of Christ, are held the teachers of believers. Confessors taught of them fight right manfully, the noble martyrs triumph, and the Christian army armed with the armour of God, ever prevaileth in warfare against the devil. All these have been men of like lealty, diverse warfarings, and glorious victories. And thou, O Christian, art but a carpet-knight, if thou thinkest to conquer without a fight, to triumph without a struggle. Nerve thyself, strive manfully, hit hard in the press. Consider thine engagement, look to thy state, know thine arm, even the engagement which thou hast taken, the state wherein thou art come, and the arm wherewith thou hast enrolled thyself a soldier.\par
\switchcolumn*

\LA
\begin{Resp}\Rsign Sancti mei, qui in carne pósiti, certámen habuístis: \Star{} Mercédem labóris ego reddam vobis.\end{Resp}
\begin{Resp}\Vsign Veníte benedícti Patris mei, percípite regnum.\end{Resp}
\begin{Resp}\Rsign Mercédem labóris ego reddam vobis.\end{Resp}
\begin{Resp}\Vsign Glória Patri, et Fílio, \Star{} et Spirítui Sancto.\end{Resp}
\begin{Resp}\Rsign Mercédem labóris ego reddam vobis.\end{Resp}
\switchcolumn
\EN
\begin{Resp}\Rsign O ye My Saints, who, being in the flesh, didst have striving \Star{} I will render unto you a reward of your labours.\end{Resp}
\begin{Resp}\Vsign Come, ye blessed of My Father, inherit the kingdom\end{Resp}
\begin{Resp}\Rsign I will render unto you a reward of your labours.\end{Resp}
\begin{Resp}\Vsign Glory be to the Father, and to the Son, \Star{} and to the Holy Ghost.\end{Resp}
\begin{Resp}\Rsign I will render unto you a reward of your labours.\end{Resp}
\switchcolumn*

\LA
\LessonHead{Lectio VII}
\switchcolumn
\EN
\LessonHead{Lesson VII}
\switchcolumn*

\LA
\LessonTitle{Léctio sancti Evangélii secúndum Matthǽum}
\switchcolumn
\EN
\LessonTitle{From the Holy Gospel according to Matthew}
\switchcolumn*

\LA
\Cite{Matt 5:1-12}
\switchcolumn
\EN
\Cite{Matt 5:1-12}
\switchcolumn*

\LA
\noindent In illo témpore: Videns Jesus turbas, ascendit in montem, et cum sedisset, accesserunt ad eum discipuli ejus. Et réliqua.\par
\switchcolumn
\EN
\noindent At that time Jesus seeing the multitudes, he went up into a mountain, and when he was set down, his disciples came unto him. And so on.\par
\switchcolumn*

\LA
\LessonTitle{De Homilia sancti Augustíni Epíscopi}
\switchcolumn
\EN
\LessonTitle{Homily by St. Augustine, Bishop (of Hippo.)}
\switchcolumn*

\LA
\Source{Liber 1 de Sermone Domini in monte, cap. 4}
\switchcolumn
\EN
\Source{Bk. i., on the Lord's Sermon.}
\switchcolumn*

\LA
\noindent In primo gradu, sicut oportebat, positum est regnum cælorum, quod est perfectæ summæque sapientiæ animæ rationalis. Sic ítaque dictum est: Beati pauperes spiritu: quoniam ipsorum est regnum cælorum; tamquam diceretur: Initium sapientiæ timor Domini. Mitibus hereditas data est, tamquam testamentum patris cum pietate quærentibus. Beati mites: quoniam ipsi hereditate possidebunt terram. Lugentibus consolatio, tamquam scientibus quid amiserint, et quibus mersi sunt. Beati qui lugent nunc: quoniam ipsi consolabuntur. Esurientibus et sitientibus saturitas, tamquam refectio laborantibus fortiterque certantibus ad salutem. Beati qui esuriunt et sitiunt justitiam, quoniam ipsi saturabuntur.\par
\switchcolumn
\EN
\noindent At the first step in blessedness is set forth, as was behoven, the kingdom of heaven, the realisation of the perfect and highest wisdom of the reasonable soul. Thus is it said Blessed are the poor in spirit, for their's is the kingdom of heaven, as though it were said The fear of the Lord is the beginning of wisdom. Then unto the meek is given an inheritance, as the legacy of a father to dutiful children. Blessed are the meek, for they shall inherit the earth. Thirdly, there is comfort for such as mourn, knowing what they have lost, and what encompasseth them. Blessed are they that mourn (now), for they shall be comforted. Fourthly, the hungry and thirsty are promised that they shall be filled, a refreshment for the strugglers for life, and for the weary. Blessed are they which do hunger and thirst after righteousness, for they shall be filled.\par
\switchcolumn*

\LA
\begin{Resp}\Rsign Sint lumbi vestri præcíncti, et lucérnæ ardéntes in mánibus vestris: \Star{} Et vos símiles homínibus exspectántibus dóminum suum, quando revertátur a núptiis.\end{Resp}
\begin{Resp}\Vsign Vigiláte ergo, quia nescítis qua hora Dóminus vester ventúrus sit.\end{Resp}
\begin{Resp}\Rsign Et vos símiles homínibus exspectántibus dóminum suum, quando revertátur a núptiis.\end{Resp}
\switchcolumn
\EN
\begin{Resp}\Rsign Let your loins be girded about, and your lights burning; \Star{} And ye yourselves like unto men that wait for their lord, when he will return from the wedding.\end{Resp}
\begin{Resp}\Vsign Watch therefore, for ye know not what hour your Lord doth come.\end{Resp}
\begin{Resp}\Rsign And ye yourselves like unto men that wait for their lord, when he will return from the wedding.\end{Resp}
\switchcolumn*

\LA
\LessonHead{Lectio VIII}
\switchcolumn
\EN
\LessonHead{Lesson VIII}
\switchcolumn*

\LA
\noindent Misericordibus misericordia tamquam vero et optimo consilio utentibus, ut hoc eis exhibeatur a potentiore, quod invalidioribus ipsi exhibent. Beati misericordes: quoniam ipsorum miserebitur Deus. Mundis corde facultas est videndi Deum, tamquam purum oculum ad intelligenda æterna gerentibus. Beati mundo corde: quoniam ipsi Deum videbunt. Pacificis Dei similitudo est, tamquam perfecte sapientibus, formatisque ad imaginem Dei per regenerationem renovati hominis. Beati pacifici: quoniam ipsi filii Dei vocabuntur. Et ista quidem in hac vita possunt compleri, sicut completa esse in Apostolis credimus. Nam illa omnimoda et in angelicam formam mutatio, quæ post hanc vitam promittitur, nullis verbis exponi potest.\par
\switchcolumn
\EN
\noindent Mercy is proclaimed unto the merciful, as unto them who have taken the true and best counsel how to obtain from Him That is Mightier than they what they that are weaker than they obtain from them. Blessed are the merciful, for they shall obtain mercy from God. Unto the pure in heart it pertaineth to see God, for their eye is clear to take in the things eternal. Blessed are the pure in heart, for they shall see God. To the peacemakers it is given to be in the likeness and image of God, for these are the perfectly wise, created anew in the image of God, by the regeneration of the new man. Blessed are the peacemakers, for they shall be called the children of God. The foregoing are forms of blessedness which we believe can be thoroughly attained in this life. The Apostles, for instance, did, we believe, attain them. As for that entire change into the likeness of Angels, which is promised us when this life is done, no words can set it forth.\par
\switchcolumn*

\LA
\begin{Resp}\Rsign Média nocte clamor factus est: \Star{} Ecce sponsus venit, exíte óbviam ei.\end{Resp}
\begin{Resp}\Vsign Prudéntes vírgines, aptáte vestras lámpades.\end{Resp}
\begin{Resp}\Rsign Ecce sponsus venit, exíte óbviam ei.\end{Resp}
\begin{Resp}\Vsign Glória Patri, et Fílio, \Star{} et Spirítui Sancto.\end{Resp}
\begin{Resp}\Rsign Ecce sponsus venit, exíte óbviam ei.\end{Resp}
\switchcolumn
\EN
\begin{Resp}\Rsign At midnight there was a cry made: \Star{} Behold, the Bridegroom cometh, go ye out to meet him.\end{Resp}
\begin{Resp}\Vsign Trim your lamps, O ye wise virgins.\end{Resp}
\begin{Resp}\Rsign Behold, the Bridegroom cometh, go ye out to meet him.\end{Resp}
\begin{Resp}\Vsign Glory be to the Father, and to the Son, \Star{} and to the Holy Ghost.\end{Resp}
\begin{Resp}\Rsign Behold, the Bridegroom cometh, go ye out to meet him.\end{Resp}
\switchcolumn*

\LA
\LessonHead{Lectio IX}
\switchcolumn
\EN
\LessonHead{Lesson IX}
\switchcolumn*

\LA
\noindent Beati ergo, qui persecutionem patiuntur propter justitiam: quoniam ipsorum est regnum cælorum. Hæc octava sententia, quæ ad caput redit, perfectumque hominem declarat, significatur fortasse et circumcisione octava die in veteri testamento; et Domini resurrectione post sabbatum, qui est utique octavus, idemque primus dies; et celebratione octavarum feriarum, quas in regeneratione novi hominis celebramus; et numero ipso Pentecostes: nam septenario numero septies multiplicato, quo fiunt quadraginta novem, quasi octavus additur, ut quinquaginta compleantur. Et tamquam redeatur ad caput, quo die missus est Spiritus sanctus, eo in regnum cælorum ducimur, et hereditatem accipimus, et consolamur, et pascimur, et misericordiam consequimur, et mundamur, et pacificamur. Atque ita perfecti omnes extrinsecus illatas molestias pro veritate et justitia sustinemus.\par
\switchcolumn
\EN
\noindent Blessed are they which are persecuted for righteousness' sake, for their's is the kingdom of heaven. In this eighth word, which returneth back again to the fountainhead and setteth forth the perfect crown of human blessedness, is contained perchance a (mystic) connection with the fact that it was upon the eighth day that the old Law commanded that circumcision should be performed, and that it was upon the day next after the Sabbath (being the seventh day) that the Lord rose again, the day whereon He so rose being thus the eighth day, (had a week contained more than seven days,) and the first day (as a matter of fact.) There is perchance also a connection with the fact that we observe eight days \textasciitilde{}(being from the Lord's Day whereon He rose to the Lord's Day in White both inclusive,) in honour of the creation of the new man. And yet again there is perchance a connection with the number contained in (the name of) the Feast of Pentecost, (which is, being interpreted, the Feast of the FiftiethDay.) For this number of fifty days is reckoned (from that of the offering of the Sheaf of the Passover) by counting (seven weeks, that is to say,) seven multiplied by seven, which is forty-nine, and thereto adding one, which joined with seven maketh eight, and so making full fifty. And thus borne backward to our fountainhead, the day whereon the Holy Ghost was sent down, we are borne unto the kingdom of heaven, and inherit the earth, and are comforted, and are filled, and obtain mercy, and are made pure, and are set at peace. And when we have thus been perfected within, we bear for truth's and righteousness' sake any troubles that may come upon us from without.\par
\switchcolumn*

\LA
\TeDeum{Te Deum}
\switchcolumn
\EN
\TeDeum{Te Deum}
\switchcolumn*

\end{paracol}

\par\needspace{10\baselineskip}\addvspace{1.2em}{\centering\small\textsc{Die 8 Novembris} · \textsc{8 November}\par}
\Office{In Octava Omnium Sanctorum}{Octave of All Saints}{Duplex majus}
\addcontentsline{toc}{subsection}{Die 8 Novembris · In Octava Omnium Sanctorum \textit{— Octave of All Saints}}
\begin{paracol}{2}
\LA
\RefLine{Lectiones I–III: de Scriptura occurrente.}
\switchcolumn
\EN
\RefLine{Lessons I–III: from the Scripture of the day.}
\switchcolumn*

\LA
\RefLine{Lectio IX: de Ss. Quatuor Coronatorum Martyrum (11-08r).}
\switchcolumn
\EN
\RefLine{Lesson IX: of Four Crowned Martyrs (11-08r).}
\switchcolumn*

\LA
\LessonHead{Lectio IV}
\switchcolumn
\EN
\LessonHead{Lesson IV}
\switchcolumn*

\LA
\LessonTitle{Ex libro sancti Cypriáni Epíscopi et Martyris de mortalitáte}
\switchcolumn
\EN
\LessonTitle{From the Book upon Death, written by the holy Martyr Cyprian, Bishop (of Carthage)}
\switchcolumn*

\LA
\Source{In fine}
\switchcolumn
\EN
\Source{At end}
\switchcolumn*

\LA
\noindent Considerandum est, fratres dilectissimi, et identidem cogitandum, renuntiasse nos mundo, et tamquam hospites et peregrinos hic interim degere. Amplectamur diem, qui assignat singulos domicilio suo, qui nos istinc ereptos, et laqueis sæcularibus exsolutos, paradiso restituit, et regno cælesti. Quis non peregre constitutus properaret in patriam regredi? Quis non ad suos navigare festinans, ventum prosperum cupidius optaret, ut velociter caros liceret amplecti?\par
\switchcolumn
\EN
\noindent Dearly beloved brethren, we should keep well in our mind and thoughts that we are living here meanwhile as strangers and pilgrims. Let us hail that day which will see us each at home in one of the many mansions, which will see us delivered hence, and disentangled from the nets and snares of things temporal, and put us back into the Garden of Eden, and into the kingdom of heaven. Is there any in a far country but is quick to make his way to his Fatherland? Was ever any in haste to make his voyage homeward, but longed for a fair wind, that he might the sooner embrace his loved ones?\par
\switchcolumn*

\LA
\begin{Resp}\Rsign Præcúrsor Dómini venit, de quo ipse testátur: \Star{} Nullus major inter natos mulíerum Joánne Baptísta.\end{Resp}
\begin{Resp}\Vsign Hic est enim prophéta, et plus quam prophéta, de quo Salvátor ait.\end{Resp}
\begin{Resp}\Rsign Nullus major inter natos mulíerum Joánne Baptísta.\end{Resp}
\switchcolumn
\EN
\begin{Resp}\Rsign The Fore-runner of the Lord cometh, to whom He Himself bare witness, saying: \Star{} Among them that are born of women there hath not risen a greater than John the Baptist.\end{Resp}
\begin{Resp}\Vsign A Prophet? Yea, and much more than a Prophet. This is he of whom the Saviour saith:\end{Resp}
\begin{Resp}\Rsign Among them that are born of women there hath not risen a greater than John the Baptist.\end{Resp}
\switchcolumn*

\LA
\LessonHead{Lectio V}
\switchcolumn
\EN
\LessonHead{Lesson V}
\switchcolumn*

\LA
\noindent Patriam nostram paradisum computamus, parentes Patriarchas habere jam cæpimus: quid non properamus et currimus, ut patriam nostram videre, ut parentes salutare possimus? Magnus illic nos carorum numerus exspectat, parentum, fratrum, filiorum frequens nos et copiosa turba desiderat, jam de sua immortalitate secura, et adhuc de nostra salute sollicita. Ad horum conspectum et complexum venire, quanta et illis et nobis in commune lætitia est! Qualis illic cælestium regnorum voluptas sine timore moriendi, et cum æternitate vivendi! Quam summa et perpetua felicitas!\par
\switchcolumn
\EN
\noindent We reckon Paradise to be our home; already we begin to have the Patriarchs for our kinsmen. Why should we not make haste and run, to see our home, and to greet our kinsfolk? There are a great many of those we love waiting for us there father, and mother, and brothers, and children, there in great company they await us, they who are sure now never to die any more, but not yet sure of us. O, when we come to see them and to embrace them, what gladness will it be both for us and for them! O, what will be the brightness of life in that heavenly kingdom where there is no more fear of death, but the certainty of living everlastingly! O, what consummated, O, what enduring happiness?\par
\switchcolumn*

\LA
\begin{Resp}\Rsign Isti sunt qui vivéntes in carne, plantavérunt Ecclésiam sánguine suo: \Star{} Cálicem Dómini bibérunt, et amíci Dei facti sunt.\end{Resp}
\begin{Resp}\Vsign In omnem terram exívit sonus eórum, et in fines orbis terræ verba eórum.\end{Resp}
\begin{Resp}\Rsign Cálicem Dómini bibérunt, et amíci Dei facti sunt.\end{Resp}
\switchcolumn
\EN
\begin{Resp}\Rsign These are they who while yet they lived in the flesh, planted the Church in their own blood; \Star{} They drank of the Lord's cup, and became the friends of God.\end{Resp}
\begin{Resp}\Vsign Their sound is gone out through all the earth, and their words to the ends of the world.\end{Resp}
\begin{Resp}\Rsign They drank of the Lord's cup, and became the friends of God.\end{Resp}
\switchcolumn*

\LA
\LessonHead{Lectio VI}
\switchcolumn
\EN
\LessonHead{Lesson VI}
\switchcolumn*

\LA
\noindent Illic Apostolorum gloriosus chorus, illic Prophetarum exsultantium numerus, illic Mártyrum innumerabilis populus, ob certaminis et passionis victoriam coronatus. Triumphantes illic Virgines, quæ concupiscentiam carnis et corporis continentiæ robore subegerunt. Remunerati misericordes, qui alimentis et largitionibus pauperum justitiæ opera fecerunt, qui Dominica præcepta servantes ad cælestes thesauros terrena patrimonia transtulerunt. Ad hos, fratres dilectissimi avida cupiditate properemus, et cum his cito esse, ut cito ad Christum venire contingat, optemus.\par
\switchcolumn
\EN
\noindent There is the glorious company of the Apostles, there is the jubilant fellowship of the Prophets, there is the countless army of Martyrs crowned for victory in strife and in suffering. There triumph the virgins who by noble self-control have tamed the desires of the flesh and of the body. There are repaid with mercy the merciful, who by feeding and gifting the needy, have wrought righteousness, have kept the commandments of the Lord, and have exchanged heritages upon earth for treasures in heaven. Thitherward, dearly beloved brethren, let us eagerly run, with such as these soon to be, unto Christ soon to come, let us be fain.\par
\switchcolumn*

\LA
\begin{Resp}\Rsign Sancti mei, qui in carne pósiti, certámen habuístis: \Star{} Mercédem labóris ego reddam vobis.\end{Resp}
\begin{Resp}\Vsign Veníte benedícti Patris mei, percípite regnum.\end{Resp}
\begin{Resp}\Rsign Mercédem labóris ego reddam vobis.\end{Resp}
\begin{Resp}\Vsign Glória Patri, et Fílio, \Star{} et Spirítui Sancto.\end{Resp}
\begin{Resp}\Rsign Mercédem labóris ego reddam vobis.\end{Resp}
\switchcolumn
\EN
\begin{Resp}\Rsign O ye My Saints, who, being in the flesh, didst have striving \Star{} I will render unto you a reward of your labours.\end{Resp}
\begin{Resp}\Vsign Come, ye blessed of My Father, inherit the kingdom\end{Resp}
\begin{Resp}\Rsign I will render unto you a reward of your labours.\end{Resp}
\begin{Resp}\Vsign Glory be to the Father, and to the Son, \Star{} and to the Holy Ghost.\end{Resp}
\begin{Resp}\Rsign I will render unto you a reward of your labours.\end{Resp}
\switchcolumn*

\LA
\LessonHead{Lectio VII}
\switchcolumn
\EN
\LessonHead{Lesson VII}
\switchcolumn*

\LA
\LessonTitle{Léctio sancti Evangélii secúndum Matthǽum}
\switchcolumn
\EN
\LessonTitle{From the Holy Gospel according to Matthew}
\switchcolumn*

\LA
\Cite{Matt 5:1-12}
\switchcolumn
\EN
\Cite{Matt 5:1-12}
\switchcolumn*

\LA
\noindent In illo témpore: Videns Jesus turbas, ascendit in montem, et cum sedisset, accesserunt ad eum discipuli ejus. Et réliqua.\par
\switchcolumn
\EN
\noindent At that time Jesus seeing the multitudes, he went up into a mountain, and when he was set down, his disciples came unto him. And so on.\par
\switchcolumn*

\LA
\LessonTitle{De Homilia sancti Augustíni Epíscopi}
\switchcolumn
\EN
\LessonTitle{Homily by St. Augustine, Bishop (of Hippo.)}
\switchcolumn*

\LA
\Source{Liber 1 de Sermone Domini in monte, cap. 5}
\switchcolumn
\EN
\Source{Bk. i. on the Lord's Sermon.}
\switchcolumn*

\LA
\noindent Beati eritis, inquit, cum vobis maledicent, et persequentur vos, et dicent omne malum adversum vos mentientes, propter me: gaudete et exsultate, quoniam merces vestra multa est in cælis. Animadvertat quisquis delicias hujus sæculi, et facultates rerum temporalium quærit in nomine christiano, intrinsecus esse beatitudinem nostram, sicut de anima ecclesiastica ore prophetico dicitur: Omnis gloria ejus filiæ regis ab intus. Nam extrinsecus maledicta, et persecutiones, et detractiones promittuntur de quibus tamen magna merces est in cælis, quæ sentitur in corde patientium, eorum qui jam possunt dicere: Gloriamur in tribulationibus: scientes quod tribulatio patientiam operatur, patientia autem probationem, probatio vero spem, spes autem non confundit: quia caritas Dei diffúsa est in cordibus nostris per Spiritum sanctum, qui datus est nobis.\par
\switchcolumn
\EN
\noindent Blessed are ye, saith the Lord, when men shall revile you, and persecute you, and shall say all manner of evil against you falsely, for My sake. Rejoice, and be exceeding glad for great is your reward in heaven. If any be seeking under the name of a Christian the pleasures of this world and the possession of temporal goods, let him bethink him that our blessedness is inward, even as the mouth of the Prophet saith concerning the soul of the Church The King's daughter is all glorious within. (Ps. xliv. 15.) Without, she is reviled, and persecution and evil report are her promised portion. And yet for these very things, great is her reward in heaven, as indeed is felt in the hearts of the sufferers, at least of such as are able already to say But we glory in tribulations also; knowing that tribulation worketh patience; and patience, experience; and experience, hope; and hope maketh not ashamed; because the love of God is shed abroad in our hearts by the Holy Ghost, Which is given unto us. \textasciitilde{}(Rom. v. 3-5.)\par
\switchcolumn*

\LA
\begin{Resp}\Rsign Sint lumbi vestri præcíncti, et lucérnæ ardéntes in mánibus vestris: \Star{} Et vos símiles homínibus exspectántibus dóminum suum, quando revertátur a núptiis.\end{Resp}
\begin{Resp}\Vsign Vigiláte ergo, quia nescítis qua hora Dóminus vester ventúrus sit.\end{Resp}
\begin{Resp}\Rsign Et vos símiles homínibus exspectántibus dóminum suum, quando revertátur a núptiis.\end{Resp}
\switchcolumn
\EN
\begin{Resp}\Rsign Let your loins be girded about, and your lights burning; \Star{} And ye yourselves like unto men that wait for their lord, when he will return from the wedding.\end{Resp}
\begin{Resp}\Vsign Watch therefore, for ye know not what hour your Lord doth come.\end{Resp}
\begin{Resp}\Rsign And ye yourselves like unto men that wait for their lord, when he will return from the wedding.\end{Resp}
\switchcolumn*

\LA
\LessonHead{Lectio VIII}
\switchcolumn
\EN
\LessonHead{Lesson VIII}
\switchcolumn*

\LA
\noindent Non enim ista perpeti fructuosum est, sed ista pro Christi nomine non solum æquo animo, sed étiam cum exsultatione tolerare. Nam multi hæretici nomine christiano animas decipientes, multa talia patiuntur: sed ideo excluduntur ab ista mercede, quia non dictum est tantum: Beati qui persecutionem patiuntur, sed additum est: Propter justitiam. Ubi autem sana fides non est, non potest esse justitia: quia justus ex fide vivit. Neque schismatici aliquid sibi ex ista mercede promittant: quia similiter, ubi caritas non est, non potest esse justitia. Dilectio enim proximi malum non operatur: quam si haberent, non dilaniarent corpus Christi, quod est Ecclesia.\par
\switchcolumn
\EN
\noindent To suffer such things is not in itself fruitful; what is fruitful, is to bear them for Christ's Name's sake not calmly only but gladly. There are a great many heretics who mislead souls under the name of Christians, and they suffer such things plentifully, but they are cut out from the reward, for it is not said only Blessed are they which are persecuted, but Blessed are they which are persecuted for righteousness' sake. Where there is not sound faith there cannot be righteousness, for the just shall live by faith. (Heb. x. 38.) Neither let schismatics promise themselves any of that reward, for as righteousness cannot exist where there is no faith, so neither can it exist where there is no love. And schismatics have no love, for love worketh no ill to his neighbour, (Rom. xiii. 10,) and if they had it, they would not tear the Body of Christ, which is the Church.\par
\switchcolumn*

\LA
\begin{Resp}\Rsign Média nocte clamor factus est: \Star{} Ecce sponsus venit, exíte óbviam ei.\end{Resp}
\begin{Resp}\Vsign Prudéntes vírgines, aptáte vestras lámpades.\end{Resp}
\begin{Resp}\Rsign Ecce sponsus venit, exíte óbviam ei.\end{Resp}
\begin{Resp}\Vsign Glória Patri, et Fílio, \Star{} et Spirítui Sancto.\end{Resp}
\begin{Resp}\Rsign Ecce sponsus venit, exíte óbviam ei.\end{Resp}
\switchcolumn
\EN
\begin{Resp}\Rsign At midnight there was a cry made: \Star{} Behold, the Bridegroom cometh, go ye out to meet him.\end{Resp}
\begin{Resp}\Vsign Trim your lamps, O ye wise virgins.\end{Resp}
\begin{Resp}\Rsign Behold, the Bridegroom cometh, go ye out to meet him.\end{Resp}
\begin{Resp}\Vsign Glory be to the Father, and to the Son, \Star{} and to the Holy Ghost.\end{Resp}
\begin{Resp}\Rsign Behold, the Bridegroom cometh, go ye out to meet him.\end{Resp}
\switchcolumn*

\end{paracol}

\par\needspace{10\baselineskip}\addvspace{1.2em}{\centering\small\textsc{Die 8 Novembris} · \textsc{8 November}\par}
\Office{Ss. Quatuor Coronatorum Martyrum}{Four Crowned Martyrs}{Simplex}
\addcontentsline{toc}{subsection}{Die 8 Novembris · Ss. Quatuor Coronatorum Martyrum \textit{— Four Crowned Martyrs}}
\begin{paracol}{2}
\LA
\LessonHead{Lectio IX (pro commemoratione)}
\switchcolumn
\EN
\LessonHead{Lesson IX (when commemorated)}
\switchcolumn*

\LA
\LessonTitle{Commemoratio: Ss. Quatuor Coronatorum Martyrum}
\switchcolumn
\EN
\LessonTitle{Commemoration of: Four Crowned Martyrs}
\switchcolumn*

\LA
\noindent Sevérus, Severiánus, Carpóphorus et Victorínus fratres, in persecutióne Diocletiáni, deórum cultum líbere detestántes, plumbátis cæsi, in verbéribus vitam pro Christi nómine profudérunt. Quorum corpora, cánibus objecta, cum ab illis intacta diu fuíssent, subláta a Christiánis, via Lavicana tertio ab Urbe lápide, in arenaria sepeliúntur, propre sepúlcrum sanctórum Mártyrum Claudii, Nicóstrati, Symphoriáni, Castórii et Simplícii, qui eodem imperatóre passi erant; quod, cum essent summi sculptores, nullo modo adduci potúerant ut idolórum statuas fácerent, et, ad solis simulacrum ducti, ut illud veneraréntur, numquam commissuros se dixérunt ut adorarent ópera mánuum hóminum. Quámobrem in carcerem conjecti, cum ibi multos dies in eodem proposito perstitissent, primum scorpiónibus cæsi, deínde, vivi plumbeis lóculis inclusi, in flumen dejiciúntur. Exstat in Urbe ecclésia sub nómine sanctórum Quatuor Coronatórum, quorum diu ignota nómina, divinitus póstea patefácta sunt; ubi non solum illórum quatuor, sed étiam horum quinque Mártyrum córpora honorifice sepúlta sunt, et eórum festívitas sexto Idus Novembris celebrátur.\par
\switchcolumn
\EN
\noindent In the persecution under Diocletian four brothers named Severus, Severian, Carpophorus, and Victorinus, boldly refused to worship the gods, and were lashed with whips loaded with lead until they gave up their lives for Christ's Name's sake under the strokes. Their bodies were thrown out to be eaten by the dogs, but as they remained untouched after a long while, the Christians took them away, and buried them in a sand-pit upon the Lavican Way at the third milestone from the City, hard by the grave of the holy martyrs Claudius, Nicostratus, Symphorian, Castorius, and Simplicius, who had suffered under the same Emperor, because being excellent sculptors they could nowise be brought to make figures of idols, and when they were brought to the image of the Sun to do reverence to it, they had said they would never worship the works of men's hands. For this reason they were thrown into prison, and when after many days they were still found of the same mind, they were first lashed with scourges armed with hooks, and then soldered up alive in leaden coffins and thrown into the river. There is in the City of Rome a Church called that of the Four Holy Crowned. Their actual names were long unknown, but afterwards made manifest by God. In this Church are honourably buried the bodies of these four, and also those of the other five; and a Festival is held in their honour upon the 8th day of November.\par
\switchcolumn*

\LA
\TeDeum{Te Deum}
\switchcolumn
\EN
\TeDeum{Te Deum}
\switchcolumn*

\end{paracol}

\par\needspace{10\baselineskip}\addvspace{1.2em}{\centering\small\textsc{Die 9 Novembris} · \textsc{9 November}\par}
\Office{In Dedicatione Archibasilicæ Ss. Salvatoris}{In Dedicatione Archibasilicæ Ss. Salvatoris}{Duplex II classis}
\addcontentsline{toc}{subsection}{Die 9 Novembris · In Dedicatione Archibasilicæ Ss. Salvatoris \textit{— In Dedicatione Archibasilicæ Ss. Salvatoris}}
\begin{paracol}{2}
\LA
\RefLine{Lectiones VII–VIII: ex Commúni Dedicationis Ecclesiae (C8).}
\switchcolumn
\EN
\RefLine{Lessons VII–VIII: from the Common of the Dedication of Churches (C8).}
\switchcolumn*

\LA
\RefLine{Lectio IX: de S. Theodori Martyris (11-09cc).}
\switchcolumn
\EN
\RefLine{Lesson IX: of St. Theodore, Martyr (11-09cc).}
\switchcolumn*

\LA
\LessonHead{Lectio I}
\switchcolumn
\EN
\LessonHead{Lesson I}
\switchcolumn*

\LA
\LessonTitle{De libro Apocalýpsis beáti Joánnis Apóstoli}
\switchcolumn
\EN
\LessonTitle{From the Book of Revelation}
\switchcolumn*

\LA
\Cite{Apo 21:9-11}
\switchcolumn
\EN
\Cite{Rev 21:9-11}
\switchcolumn*

\LA
\noindent Et venit unus de septem ángelis habéntibus phíalas plenas septem plagis novíssimis, et locútus est mecum, dicens: Veni, et osténdam tibi sponsam, uxórem Agni. Et sústulit me in spíritu in montem magnum et altum, et osténdit mihi civitátem sanctam Jerúsalem descendéntem de cælo a Deo, habéntem claritátem Dei: et lumen ejus símile lápidi pretióso tamquam lápidi jáspidis, sicut crystállum.\par
\switchcolumn
\EN
\noindent And there came one of the seven angels, who had the vials full of the seven last plagues, and spoke with me, saying: Come, and I will shew thee the bride, the wife of the Lamb. And he took me up in spirit to a great and high mountain: and he shewed me the holy city Jerusalem coming down out of heaven from God, Having the glory of God, and the light thereof was like to a precious stone, as to the jasper stone, even as crystal.\par
\switchcolumn*

\LA
\begin{Resp}\Rsign In dedicatióne templi decantábat pópulus laudem: \Star{} Et in ore eórum dulcis resonábat sonus.\end{Resp}
\begin{Resp}\Vsign Fundáta est domus Dómini supra vérticem móntium, et vénient ad eam omnes gentes.\end{Resp}
\begin{Resp}\Rsign Et in ore eórum dulcis resonábat sonus.\end{Resp}
\switchcolumn
\EN
\begin{Resp}\Rsign When the Temple was dedicated the people sang praise; \Star{} And sweet in their mouths was the sound.\end{Resp}
\begin{Resp}\Vsign The Lord's house is established in the top of the mountains; and all nations shall flow unto it.\end{Resp}
\begin{Resp}\Rsign And sweet in their mouths was the sound.\end{Resp}
\switchcolumn*

\LA
\LessonHead{Lectio II}
\switchcolumn
\EN
\LessonHead{Lesson II}
\switchcolumn*

\LA
\Cite{Apo 21:12-15}
\switchcolumn
\EN
\Cite{Rev 21:12-15}
\switchcolumn*

\LA
\noindent Et habébat murum magnum, et altum, habéntem portas duódecim: et in portis ángelos duódecim, et nómina inscrípta, quæ sunt nómina duódecim tríbuum filiórum Israël: ab oriénte portæ tres, et ab aquilóne portæ tres, et ab austro portæ tres, et ab occásu portæ tres. Et murus civitátis habens fundaménta duódecim, et in ipsis duódecim nómina duódecim apostolórum Agni. Et qui loquebátur mecum, habébat mensúram arundíneam áuream, ut metirétur civitátem, et portas ejus, et murum.\par
\switchcolumn
\EN
\noindent And it had a wall great and high, having twelve gates, and in the gates twelve angels, and names written thereon, which are the names of the twelve tribes of the children of Israel. On the east, three gates: and on the north, three gates: and on the south, three gates: and on the west, three gates. And the wall of the city had twelve foundations, and in them, the twelve names of the twelve apostles of the Lamb. And he that spoke with me, had a measure of a reed of gold, to measure the city and the gates thereof, and the wall.\par
\switchcolumn*

\LA
\begin{Resp}\Rsign Fundáta est domus Dómini supra vérticem móntium, et exaltáta est super omnes colles: \Star{} Et vénient ad eam omnes gentes, et dicent: Glória tibi, Dómine.\end{Resp}
\begin{Resp}\Vsign Veniéntes autem vénient cum exsultatióne, portántes manípulos suos.\end{Resp}
\begin{Resp}\Rsign Et vénient ad eam omnes gentes, et dicent: Glória tibi, Dómine.\end{Resp}
\switchcolumn
\EN
\begin{Resp}\Rsign The Lord's house is established in the top of the mountains, and exalted above the hills; \Star{} And all nations shall flow unto it, and shall say: Glory be to thee, O Lord!\end{Resp}
\begin{Resp}\Vsign They shall doubtless come again with rejoicing, bringing their sheaves with them.\end{Resp}
\begin{Resp}\Rsign And all nations shall flow unto it, and shall say: Glory be to thee, O Lord!\end{Resp}
\switchcolumn*

\LA
\LessonHead{Lectio III}
\switchcolumn
\EN
\LessonHead{Lesson III}
\switchcolumn*

\LA
\Cite{Apo 21:16-18}
\switchcolumn
\EN
\Cite{Rev 21:16-18}
\switchcolumn*

\LA
\noindent Et cívitas in quadro pósita est, et longitúdo ejus tanta est quanta et latitúdo: et mensus est civitátem de arúndine áurea per stádia duódecim míllia: et longitúdo, et altitúdo, et latitúdo ejus æquália sunt. Et mensus est murum ejus centum quadragínta quátuor cubitórum, mensúra hóminis, quæ est ángeli. Et erat structúra muri ejus ex lápide jáspide: ipsa vero cívitas aurum mundum símile vitro mundo.\par
\switchcolumn
\EN
\noindent And the city lieth in a foursquare, and the length thereof is as great as the breadth: and he measured the city with the golden reed for twelve thousand furlongs, and the length and the height and the breadth thereof are equal. And he measured the wall thereof an hundred and forty-four cubits, the measure of a man, which is of an angel. And the building of the wall thereof was of jasper stone: but the city itself pure gold, like to clear glass.\par
\switchcolumn*

\LA
\begin{Resp}\Rsign Bénedic, Dómine, domum istam, quam ædificávi nómini tuo: veniéntium in loco isto \Star{} Exáudi preces in excélso sólio glóriæ tuæ.\end{Resp}
\begin{Resp}\Vsign Dómine, si convérsus fúerit pópulus tuus, et oráverit ad sanctuárium tuum.\end{Resp}
\begin{Resp}\Rsign Exáudi preces in excélso sólio glóriæ tuæ.\end{Resp}
\begin{Resp}\Vsign Glória Patri, et Fílio, \Star{} et Spirítui Sancto.\end{Resp}
\begin{Resp}\Rsign Exáudi preces in excélso sólio glóriæ tuæ.\end{Resp}
\switchcolumn
\EN
\begin{Resp}\Rsign O Lord, bless this house which I have built unto thy name. whosoever shall come unto this place and pray, then \Star{} Hear thou from the excellent throne of thy glory.\end{Resp}
\begin{Resp}\Vsign O Lord, if thy people turn and pray toward thy sanctuary.\end{Resp}
\begin{Resp}\Rsign Hear thou from the excellent throne of thy glory.\end{Resp}
\begin{Resp}\Vsign Glory be to the Father, and to the Son, \Star{} and to the Holy Ghost.\end{Resp}
\begin{Resp}\Rsign Hear Thou from the excellent throne of thy glory.\end{Resp}
\switchcolumn*

\LA
\LessonHead{Lectio IV}
\switchcolumn
\EN
\LessonHead{Lesson IV}
\switchcolumn*

\LA
\noindent Ritus, quos in consecrándis ecclésiis et altáribus Romána servat Ecclésia, beátus Silvéster Papa primus instítuit. Nam, etsi jam ab Apostolórum témpore loca fuérunt Deo dicáta, quæ a quibúsdam oratória, ab áliis ecclésiæ dicebántur, ubi colléctæ fiébant per unam sábbati, et christiánus pópulus oráre, Dei verbum audíre et Eucharístiam súmere sólitus erat; non tamen illa ádeo solémni ritu consecrabántur, nec in eis adhuc in títulum eréctum erat altáre, quod chrísmate delibútum, Dómini nostri Jesu Christi, qui altáre, hóstia et sacérdos noster est, figúram exprímeret.\par
\switchcolumn
\EN
\noindent The Rites whereof the Church of Rome maketh use for the hallowing of Churches and Altars were first instituted by the blessed Pope Sylvester. From the very time of the Apostles there had been places set apart for God, where assemblies took place upon the first day of every week, and where the Christians were used to pray, to hear the word of God, and to receive the Eucharist, which places were by some called Oratories and by others Churches. But these places were not dedicated with so solemn a form, nor did they set up therein an Altar for a pillar, and pour chrism thereon, (Gen. xxviii. 18,) for a figure of our Lord Jesus Christ, Who is Himself our Altar, our Victim, and our Priest.\par
\switchcolumn*

\LA
\begin{Resp}\Rsign Orántibus in loco isto, \Star{} Dimítte peccáta pópuli tui, Deus, et osténde eis viam bonam per quam ámbulent, et da glóriam in loco isto.\end{Resp}
\begin{Resp}\Vsign Qui regis Israël inténde, qui dedúcis velut ovem Joseph, qui sedes super Chérubim.\end{Resp}
\begin{Resp}\Rsign Dimítte peccáta pópuli tui, Deus, et osténde eis viam bonam per quam ámbulent, et da glóriam in loco isto.\end{Resp}
\switchcolumn
\EN
\begin{Resp}\Rsign If they pray toward this place, \Star{} Forgive the sin of thy people, O God, and teach them the good way wherein they should walk, and manifest forth thy glory in this place.\end{Resp}
\begin{Resp}\Vsign Give ear, O Shepherd of Israel, thou that leadest Joseph like a flock, thou that sittest upon the cherubim.\end{Resp}
\begin{Resp}\Rsign Forgive the sin of thy people, O God, and teach them the good way wherein they should walk, and manifest forth thy glory in this place.\end{Resp}
\switchcolumn*

\LA
\LessonHead{Lectio V}
\switchcolumn
\EN
\LessonHead{Lesson V}
\switchcolumn*

\LA
\noindent Sed, ubi Constantínus imperátor per baptísmi sacraméntum sanitátem salutémque consecútus est, tum primum lege ab eo lata concéssum est toto orbe terrárum, Christiáni ut ecclésias ædificárent; quos ille, non solum edícto, sed étiam exémplo ad sacram ædificatiónem est cohortátus. Nam, et in suo Lateranénsi palátio ecclésiam Salvatóri dedicávit, et ei continéntem basílicam nómine sancti Joánnis Baptístæ cóndidit, eo loco quo ipse, baptizátus a sancto Silvéstro, ab infidelitátis lepra mundátus est; quam idem Póntifex consecrávit quinto idus novémbris. Cujus consecratiónis memória celebrátur hodiérno die, quo primum Romæ públice ecclésia consecráta est, et imágo Salvatóris, in paríete depícta, pópulo Románo appáruit.\par
\switchcolumn
\EN
\noindent But when the Emperor Constantine had by the Sacrament of Baptism received health both of body and soul, then first in a law by him published was it allowed to the Christians throughout the whole world to build Churches, to the which holy building he exhorted them by his example as well as by his decree. He dedicated in his own Lateran Palace a Church to the Saviour, and built hard by it a Cathedral in the name of St. John the Baptist, upon the place where he had been baptized by holy Sylvester and cleansed from his leprosy. This Cathedral was hallowed by the said Pope upon the 9th day of November. It is this consecration, the memory whereof is still celebrated upon this day, the first whereon the public consecration of a Church ever took place in Rome, and the image of the Saviour was seen by the Roman people painted upon the wall.\par
\switchcolumn*

\LA
\begin{Resp}\Rsign O quam metuéndus est locus iste: \Star{} Vere non est hic áliud nisi domus Dei et porta cæli.\end{Resp}
\begin{Resp}\Vsign Hæc est domus Dómini fírmiter ædificáta, bene fundáta est supra firmam petram.\end{Resp}
\begin{Resp}\Rsign Vere non est hic áliud nisi domus Dei et porta cæli.\end{Resp}
\switchcolumn
\EN
\begin{Resp}\Rsign How dreadful is this place! \Star{} Surely this is none other but the house of God, and this is the gate of heaven.\end{Resp}
\begin{Resp}\Vsign This is the house of God, stoutly builded, well founded upon a sure rock.\end{Resp}
\begin{Resp}\Rsign Surely this is none other but the house of God, and this is the gate of heaven.\end{Resp}
\switchcolumn*

\LA
\LessonHead{Lectio VI}
\switchcolumn
\EN
\LessonHead{Lesson VI}
\switchcolumn*

\LA
\noindent Quod, si beátus Silvéster póstea in consecratióne altáris Príncipis Apostolórum decrévit ut deínceps nisi ex lápide altária non ædificaréntur; tamen basílicæ Lateranénsis altáre fuit e ligno eréctum. Quod mirum non est; nam, cum a sancto Petro usque ad Silvéstrum, propter persecutiónes, Pontífices certo loco consístere non possent, quocúmque eos necéssitas compulísset, sive in cryptas, sive in cœmetéria, sive in ædes piórum, super illo altári lígneo ad arcæ similitúdinem cóncavo, sacra faciébant. Quo altári sanctus Silvéster, réddita Ecclésiæ pace, honóris causa Príncipis Apostolórum, qui in illo sacrificásse dícitur, et reliquórum Pontíficum, qui usque ad id tempus ad mystéria conficiénda eo usi fúerant, in Lateranénsi prima ecclésia collocáto, sancívit ne quisquam in eo, præter Románum Pontíficem, Missam deínceps celebráret.\par
\switchcolumn
\EN
\noindent The Blessed Sylvester afterwards decreed, when he was consecrating the Altar of the Prince of the Apostles, that Altars were thenceforward to be made of stone only, but notwithstanding this the Lateran Cathedral hath the altar made of wood. This is not surprising. From St. Peter to Sylvester the Popes had not been able, by reason of persecutions, to abide fixedly in one place, and they celebrated the Holy Liturgy in cellars, in burying-places, in the houses of godly persons, or wherever need drove them, upon a wooden altar made like an empty box. When peace was given to the Church, holy Sylvester took this box, and to do honour to the Prince of the Apostles, who is said to have offered sacrifice thereon, and to the other Popes who thereon had been used to execute the mystery even unto that time, set it in the first Church, even the Lateran, and ordained that no one but the Bishop of Rome should celebrate the Liturgy thereon for all time coming. s/\$/\textasciitilde{}/\par
\switchcolumn*

\LA
Eándem ecclésiam incéndiis, vastatiónibus, terræ ínsuper mótibus disjéctam eversámque, ac sédula summórum Pontíficum cura reparátam, nova póstmodum molitióne restitútam, Benedíctus décimus tértius, Póntifex máximus, órdinis Prædicatórum, die vigésima octáva aprílis anni millésimi septingentésimi vigésimi sexti, ritu solémni consecrávit, ejúsque celebritátis memóriam hac die recoléndam státuit. Quod autem Pius nonus perficiéndum censúerat, Leo décimus tértius cellam máximam, vetustáte fatiscéntem, ingénti molitióne producéndam laxandámque curávit; vetus musívum, multis jam ántea pártibus instaurátum, ad antíquum exémplar restítui et in novam ábsidem, ópere cultúque magnífico exornátam, transférri; aulam transvérsam, laqueári et contignatióne reféctis, expolíri jussit, anno millésimo octingentésimo octuagésimo quarto; sacrário, æde canonicórum, perpetuáque ad baptistérium Constantiniánum pórticu adjéctis.\par
\switchcolumn
\EN
The original Lateran Cathedral, cast down and destroyed by fires, pillage, and earthquakes, and renewed by the constant care of the Popes, was at last rebuilt afresh, and solemnly consecrated by Pope Benedict XIII, a Friar Preacher, upon the 28th day of April, in the year 1726, the memory of which Festival he ordained to be kept upon this day. In the year 1884 Leo XIII took in hand a work which had received the sanction of his predecessor Pius IX. The great sanctuary, the walls of which were giving way with age, was lengthened and widened, a task of immense labour. The ancient mosaic had been renewed previously in several places; it was now restored according to the original design, and transferred to the new apse, the embellishment of which was carried out with great magnificence. The transept was redecorated, and its ceiling and woodwork repaired. A sacristy, a residence for the canons, and a portico connecting with the baptistery of Constantine, were added to the existing buildings.\par
\switchcolumn*

\LA
\begin{Resp}\Rsign Mane surgens Jacob, erigébat lápidem in títulum, fundens óleum désuper; votum vovit Dómino: \Star{} Vere locus iste sanctus est, et ego nesciébam.\end{Resp}
\begin{Resp}\Vsign Cumque evigilásset Jacob de somno, ait.\end{Resp}
\begin{Resp}\Rsign Vere locus iste sanctus est, et ego nesciébam.\end{Resp}
\begin{Resp}\Vsign Glória Patri, et Fílio, \Star{} et Spirítui Sancto.\end{Resp}
\begin{Resp}\Rsign Vere locus iste sanctus est, et ego nesciébam.\end{Resp}
\switchcolumn
\EN
\begin{Resp}\Rsign Jacob rose up early in the morning, and set up the stone for a pillar, and poured oil upon the top of it, and vowed a vow unto the Lord. \Star{} Surely this place is holy, and I knew it not.\end{Resp}
\begin{Resp}\Vsign And Jacob awaked out of his sleep, and he said\end{Resp}
\begin{Resp}\Rsign Surely this place is holy, and I knew it not.\end{Resp}
\begin{Resp}\Vsign Glory be to the Father, and to the Son, \Star{} and to the Holy Ghost.\end{Resp}
\begin{Resp}\Rsign Surely this place is holy, and I knew it not.\end{Resp}
\switchcolumn*

\end{paracol}

\par\needspace{10\baselineskip}\addvspace{1.2em}{\centering\small\textsc{Die 9 Novembris} · \textsc{9 November}\par}
\Office{S. Theodori Martyris}{St. Theodore, Martyr}{Simplex}
\addcontentsline{toc}{subsection}{Die 9 Novembris · S. Theodori Martyris \textit{— St. Theodore, Martyr}}
\begin{paracol}{2}
\LA
\LessonHead{Lectio IX (pro commemoratione)}
\switchcolumn
\EN
\LessonHead{Lesson IX (when commemorated)}
\switchcolumn*

\LA
\LessonTitle{Commemoratio: S. Theodori Martyris}
\switchcolumn
\EN
\LessonTitle{Commemoration of: St. Theodore, Martyr}
\switchcolumn*

\LA
\noindent Theodórus, miles christiánus, Maximiáno imperatóre, quod idolórum fanum incendísset, comprehénsus, cum a præfécto legiónis pœna ei remitterétur, si pœ́nitens facti christiánam fidem exsecrarétur, constánter in fídei confessióne persevérans, missus est in cárcerem. Ubi, úngulis excarnificátus, dum costæ nudaréntur, lætus canébat: Benedícam Dóminum in omni témpore. Quare, in ardéntem rogum injéctus, in oratióne et divínis láudibus ánimam Christo réddidit, quinto Idus Novémbris. Cujus corpus Eusébia matróna, síndone involútum, sepelívit in suo prǽdio.\par
\switchcolumn
\EN
\noindent This Theodore was a Christian soldier, who was arrested in the reign of the Emperor Maximian for having set fire to a temple of idols. The Commander of the Legion offered him pardon if he would profess repentance and curse the Christian faith, but, as he refused to swerve as regarding the confession of his belief, he was cast into prison. There he was tormented with iron claws. As they were tearing the flesh off his ribs, he sang joyfully (the 33rd Psalm) I will bless the Lord at all times. Thereafter he was thrown upon an heap of burning wood, and there, still praying and praising God, he gave up his soul to Christ, upon the 9th day of November, (in the year of salvation 304.) The Lady Eusebia wrapped his body in a winding-sheet, and buried it on her own farm.\par
\switchcolumn*

\LA
\TeDeum{Te Deum}
\switchcolumn
\EN
\TeDeum{Te Deum}
\switchcolumn*

\end{paracol}

\par\needspace{10\baselineskip}\addvspace{1.2em}{\centering\small\textsc{Die 10 Novembris} · \textsc{10 November}\par}
\Office{S. Andreæ Avellini Confessoris}{St. Andrew Avellino Confessor}{Duplex}
\addcontentsline{toc}{subsection}{Die 10 Novembris · S. Andreæ Avellini Confessoris \textit{— St. Andrew Avellino Confessor}}
\begin{paracol}{2}
\LA
\RefLine{Lectiones I–III: de Scriptura occurrente.}
\switchcolumn
\EN
\RefLine{Lessons I–III: from the Scripture of the day.}
\switchcolumn*

\LA
\RefLine{Lectiones VII–VIII: ex Commúni Confessoris non pontificis (C5).}
\switchcolumn
\EN
\RefLine{Lessons VII–VIII: from the Common of a Confessor not a Bishop (C5).}
\switchcolumn*

\LA
\RefLine{Lectio IX: de Ss. Tryphonis, Respicii, et Nymphæ Virg., Martyrum (11-10o).}
\switchcolumn
\EN
\RefLine{Lesson IX: of Sts. Tryphon, Respicius, et Nympha, Martyrs (11-10o).}
\switchcolumn*

\LA
\LessonHead{Lectio IV}
\switchcolumn
\EN
\LessonHead{Lesson IV}
\switchcolumn*

\LA
\noindent Andreas Avellinus, dictus antea Lancellottus, apud Castrum Novum Lucaniæ pagum natus, inter ipsa infantiæ primordia, futuræ sanctitátis non obscura præbuit indicia. Adoléscens, ad litteras addiscendas paterna e domo egréssus, lubricam illíus ætátis sémitam inter bonárum artium stúdia ita peregit, ut sapiéntiæ initium, quod est timor Dómini, ob óculos potíssimum habere numquam prætermiserit. Cum egregia proinde forma exímium castitátis studium conjunxit, quo impudícas sæpe mulíerum insidias elusit, intérdum étiam apértam vim propulsávit. Clericali milítiæ jam pridem adscriptus, Neapolim se cóntulit, ut legálibus disciplinis vacaret; ibique jurisprudéntiæ lauream adeptus atque interea ad sacerdotalem dignitátem evectus, causárum patrocinia in foro dumtaxat ecclesiástico proque privátis quibusdam persónis, juxta sacrórum cánonum sanctiónes ágere cœpit. Verum, cum aliquándo inter causam agéndam leve ei mendácium excidísset, mox vero fortúita sacrárum Scripturárum lectióne in illa verba incidísset: Os, quod mentítur, occídit ánimam; tanto ejus culpæ dolóre ac pœniténtiæ correptus est, ut statim ab ejusmodi vitæ instituto sibi recedéndum esse duxerit. Itaque, abdicátis fœnsibus curis, se totum divino cultui sacrisque ministériis mancipávit. Cumque ecclesiásticæ virtútis exemplis emineret, sanctimoniálium regímini a tunc exsisténte archiepiscopo Neapolitano præfectus fuit. Quo in munere cum pravórum hóminum odia subiísset, primo quidem intentátæ sibi necis periculum declinávit; mox vero, per sicárium tribus in fácie acceptis vulnéribus, injuriæ atrocitátem æquo animo pértulit. Tunc, perfectioris vitæ desidério flagrans, ut inter Cléricos regulares adscriberétur, suppliciter postulávit; votique compos factus, ob ingentem quo æstuábat crucis amórem, ut sibi Andreæ nomen imponerétur, precibus impetrávit.\par
\switchcolumn
\EN
\noindent Lancelot Avelino, who afterwards took the name of Andrew, was born at Castro Nuovo, a small town in Lucania, (in the kingdom of Naples, in the year of our Lord 1520.) From his earliest childhood he gave no dark signs of the holiness of his after life. When as a lad he was away from home at school, he so passed the slippery paths of that age, as ever keeping before his eyes, amid the pursuit of earthly knowledge, the true beginning of wisdom, which is the fear of the Lord. \textasciitilde{}(Prov. ix. 10.) He was exceedingly comely, but withal careful in purity, and thereby escaped oftentimes the shameless proposals of women, and somewhiles even resisted open violence. He had already become a clerk when he went to Naples to study law. There he was ordained Priest, and also took his degree in Jurisprudence. He undertook cases only in the Church Courts, and for certain private persons, according to the rules of the Sacred Canons. l Once in pleading a cause, in a matter indeed which was of no weight, a lie escaped him. Almost forthwith thereafter, in reading the Holy Scriptures, he came upon the words: The mouth that lieth killeth the soul (Wisd. i. 11) and so great was the grief and remorse which he felt for his sin that he made up his mind to leave that way of life. He therefore gave up his law business, and set himself altogether to mind the worship of God and the execution of his holy ministry. The eminent pattern which he gave of all the graces proper to a Churchman moved the Archbishop of Naples to commit to him the care of a certain nunnery in that city. The holy man's zeal (for removing all obstacles to the recollection of these spouses of Christ, in which consisteth the very essence of their state and virtue,) stirred up the malice and rage of certain wicked men in the city, (whom he had forbid being admitted to the grate to speak to any of the nuns.) He once narrowly escaped death, with which they threatened him; and another time received three wounds in his face from a bully. These injuries he bore with thorough meekness. Out of an earnest desire of more readily attaining to a perfect disengagement of his heart from all earthly things, he humbly sought and \textasciitilde{}(in 1556) obtained to be admitted into the Order of Regular Clerks, (called Theatins,) and on this occasion, out of the love he bore to the Cross, he entreated that his name might be changed from Lancelot to Andrew.\par
\switchcolumn*

\LA
\begin{Resp}\Rsign Honéstum fecit illum Dóminus, et custodívit eum ab inimícis, et a seductóribus tutávit illum: \Star{} Et dedit illi claritátem ætérnam.\end{Resp}
\begin{Resp}\Vsign Justum dedúxit Dóminus per vias rectas, et osténdit illi regnum Dei.\end{Resp}
\begin{Resp}\Rsign Et dedit illi claritátem ætérnam.\end{Resp}
\switchcolumn
\EN
\begin{Resp}\Rsign The Lord made him honourable, and defended him from his enemies, and kept him safe from those that lay in wait for him; \Star{} And gave him perpetual glory.\end{Resp}
\begin{Resp}\Vsign He went down with him into the pit, and left him not in bonds.\end{Resp}
\begin{Resp}\Rsign And gave him perpetual glory.\end{Resp}
\switchcolumn*

\LA
\LessonHead{Lectio V}
\switchcolumn
\EN
\LessonHead{Lesson V}
\switchcolumn*

\LA
\noindent Arctioris ítaque vitæ curriculum álacri studio ingréssus, in eas maxime virtútis exercitatiónes incúbuit, ad quas sese arduis étiam emissis votis obstrinxit; altero scilicet suæ ipsíus voluntáti jugiter obsisténdi, altero vero in via christianæ perfectiónis semper ulterius progrediendi. Regularis disciplinæ cultor assiduus, et in ea promovenda, cum aliis præesset, studiosíssimus fuit. Quidquid ab institúti sui officii et regulæ præscripto supererat témporis, oratióni et animárum salúti tribuebat. In confessiónibus excipiéndis mira ejus pietas et prudéntia enituit; vicos et oppida Neapoli finítima evangelicis ministériis magno cum animárum lucro frequens lustrábat. Quam ardentem erga próximos sancti viri caritátem signis étiam Dóminus illustrávit. Cum enim, intempesta nocte, ab audíta ægri confessióne domum rediret, ac pluviæ ventorúmque vis prælucéntem facem exstinxísset, non solum ipse cum sociis inter effusíssimos imbres nihil madefactus est; verum etiam, inusitato splendore e suo corpore mirabíliter emicante, sociis inter densíssimas ténebras iter monstrávit. Abstinéntia et patiéntia, nec non abjectióne atque ódio sui summopere præstitit. Necem fratris fílio illatam, imperturbato animo tulit, ac suos ab omni ulciscéndi cupiditate compescuit; immo étiam pro interfectóribus opem et misericórdiam judícium implorávit.\par
\switchcolumn
\EN
\noindent He entered manfully and cheerily upon the harder life, set to work to better himself therein, and to that end made two very grim vows, the first, perpetually to fight against his own will, the second, always to advance to the utmost of his power in Christian perfection. Of the discipline of his Order he was a stern defender, and when he was set over others the observance thereof was his great care. Whatever time the duties of his work and his institute left him, he gave to prayer and the salvation of souls. His godliness and wisdom in hearing of confessions were beautiful. He went many times through the farthest lanes and suburbs of Naples, bringing Gospel ministry with great gain of souls. The greatness of his love toward his neighbour God was pleased to crown even by signs and wonders. One stormy night he was coming home from hearing a sick man's confession, when the rain and wind put out the light which was carried before him, but he and they that were with him not only came dry through the thickest of the rain, but there came also a strange light out of his body and showed them the way in the deepest of the darkness. He was a wonderful instance of selfcontrol, long-suffering, lowliness, and hatred of self. He bore with stillness the murder of his nephew, held in the passion of his kinsfolk to take revenge, and even asked pity for the assassins from the judges.\par
\switchcolumn*

\LA
\begin{Resp}\Rsign Amávit eum Dóminus, et ornávit eum: stolam glóriæ índuit eum, \Star{} Et ad portas paradísi coronávit eum.\end{Resp}
\begin{Resp}\Vsign Índuit eum Dóminus lorícam fídei, et ornávit eum.\end{Resp}
\begin{Resp}\Rsign Et ad portas paradísi coronávit eum.\end{Resp}
\switchcolumn
\EN
\begin{Resp}\Rsign The Lord loved him and beautified him; He clothed him with a robe of glory, \Star{} And crowned him at the gates of Paradise.\end{Resp}
\begin{Resp}\Vsign The Lord hath put on him the breast-plate of faith, and hath adorned him.\end{Resp}
\begin{Resp}\Rsign And crowned him at the gates of Paradise.\end{Resp}
\switchcolumn*

\LA
\LessonHead{Lectio VI}
\switchcolumn
\EN
\LessonHead{Lesson VI}
\switchcolumn*

\LA
\noindent Plúribus in locis Clericórum regulárium ordinem propagávit, eorúmdem domicília Mediolani et Placéntiæ instituit. Illius operam sanctus Cárolus Borromæus et Paulus de Aretio Clericus regularis, cardinales, quibus erat acceptíssimus, in pastoralis muneris curis adhibuérunt. Deiparam Vírginem singulari amóre et cultu prosequebátur. Angelórum colloquio pérfrui meruit, quos, cum divinas laudes persólveret, e regióne concinéntes se audisse testátus est. Denique, post heróica virtútum exempla, prophetíæ quoque dono illustris, quo et secreta cordium et abséntia et futura prospéxit, annis gravis et labóribus fractus, ad aram celebratúrus in verbis illis tertio repetitis: Introíbo ad altáre Dei, repentino apoplexíæ morbo correptus est; mox sacramentis rite munítus, placidíssime inter suos ánimam efflávit. Ejus corpus Neapoli in ecclésia sancti Pauli ad hæc usque témpora eo frequentíssimo pópuli concursu cólitur, quo fuit elátum. Illum denique, insígnibus in vita et post mortem miraculis clarum, Clemens undecimus Pontifex maximus solemni ritu Sanctórum catálogo adscripsit.\par
\switchcolumn
\EN
\noindent He spread in many places the Institute of Regular Clerks, and founded their houses at Milan and Piacenza. The holy Cardinal Charles Borromeo, and the Cardinal Paul of Arezzo, being himself a Regular Clerk, men by both of whom he was well liked, used his help in their care for souls. Toward the Virgin Mother of God he was constant in an extraordinary love and reverence. He won the conversation of Angels, whom he said he used to hear singing when he was praising God. He set an example of the highest graces, even to the gift of prophecy, whereby he saw into men's hearts and knew things afar off or even yet to come. Full of years and worn out with work, he was beginning the Liturgy, when, having repeated thrice the words, I will go unto the Altar of God, he was felled by a stroke of apoplexy, and, duly fortified by the Sacraments, in the arms of his friends, most peacefully gave up his soul to God, (upon the 10th day of November, in the year 1608.) The crowds which flock to his grave in the Church of St. Paul at Naples are still as great as they were when his body was first laid there. He was famous for signs and wonders both during his life and after his death, and Pope Clement XI solemnly enrolled his name among those of the Saints.\par
\switchcolumn*

\LA
\begin{Resp}\Rsign Iste homo perfécit ómnia quæ locútus est ei Deus, et dixit ad eum: Ingrédere in réquiem meam: \Star{} Quia te vidi justum coram me ex ómnibus géntibus.\end{Resp}
\begin{Resp}\Vsign Iste est, qui contémpsit vitam mundi, et pervénit ad cæléstia regna.\end{Resp}
\begin{Resp}\Rsign Quia te vidi justum coram me ex ómnibus géntibus.\end{Resp}
\begin{Resp}\Vsign Glória Patri, et Fílio, \Star{} et Spirítui Sancto.\end{Resp}
\begin{Resp}\Rsign Quia te vidi justum coram me ex ómnibus géntibus.\end{Resp}
\switchcolumn
\EN
\begin{Resp}\Rsign This is he which did according unto all that God commanded him; and God said unto him: Enter thou into My rest, \Star{} For thee have I seen righteous before Me among all people.\end{Resp}
\begin{Resp}\Vsign This is he which loved not his life in this world, and is come unto an everlasting kingdom.\end{Resp}
\begin{Resp}\Rsign For thee have I seen righteous before Me among all people.\end{Resp}
\begin{Resp}\Vsign Glory be to the Father, and to the Son, \Star{} and to the Holy Ghost.\end{Resp}
\begin{Resp}\Rsign For thee have I seen righteous before Me among all people.\end{Resp}
\switchcolumn*

\end{paracol}

\par\needspace{10\baselineskip}\addvspace{1.2em}{\centering\small\textsc{Die 10 Novembris} · \textsc{10 November}\par}
\Office{Ss. Tryphonis, Respicii, et Nymphæ Virg., Martyrum}{Sts. Tryphon, Respicius, et Nympha, Martyrs}{Simplex}
\addcontentsline{toc}{subsection}{Die 10 Novembris · Ss. Tryphonis, Respicii, et Nymphæ Virg., Martyrum \textit{— Sts. Tryphon, Respicius, et Nympha, Martyrs}}
\begin{paracol}{2}
\LA
\LessonHead{Lectio IX (pro commemoratione)}
\switchcolumn
\EN
\LessonHead{Lesson IX (when commemorated)}
\switchcolumn*

\LA
\LessonTitle{Commemoratio: Ss. Tryphonis, Respicii, et Nymphæ Virg., Martyrum}
\switchcolumn
\EN
\LessonTitle{Commemoration of: Sts. Tryphon, Respicius, et Nympha, Martyrs}
\switchcolumn*

\LA
\noindent Tryphon, Decio imperatóre, cum Jesu Christi fidem prædicans, omnes ad ejus cultum perducere conarétur, a Decii satellítibus comprehénsus, primum equuleo torquétur et ungulis férreis excarnificátur; deínde, sublimibus pédibus candéntibus clavis confíxus, fustibus cæditur, et, admotis facibus ardéntibus, aduritur. Quæ cum ómnia fortiter ferentem vidísset Respicius tribunus, ad Christi Dómini fidem convérsus, statim se christianum esse palam proféssus est. Qui, varie cruciatus, una cum Tryphone rapitur ad Jovis simulacrum; quæ statua, Tryphone orante, cóncidit. Quare plumbátis crudelíssime contúsi, nobilíssimum martyrium consecúti sunt quarto Idus Novembris. Eodem die virgo quædam, cui nomen Nympha, cum Jesum Christum verum esse Deum clara voce testarétur, martyrii palmam ad virginitátis corónam adjunxit.\par
\switchcolumn
\EN
\noindent In the reign of the Emperor Decius one Tryphon strove by preaching the faith of Jesus Christ to bring all men to worship Him. For this cause he was taken by the servants of Decius. He was first tormented upon the rack, and flesh stripped from him with iron claws; then red-hot nails were driven into his insteps, he was beaten with cudgels and scarified with lighted torches. The sight of the courage wherewith he bore all, brought the Praefect Respicius to believe in the Lord Christ, and he forthwith declared himself a Christian. He also was diverse ways tormented, and then led along with Tryphon before the statue of Jupiter. When Tryphon prayed, the statue fell down. Then were both Tryphon and Respicius savagely lashed with whips loaded with lead, until they grasped the crown of a most glorious testimony, upon the 10th day of November. Upon the same day a certain maiden named Nympha, having openly confessed that Jesus Christ is very God, added the palm of martyrdom to the crown of virginity.\par
\switchcolumn*

\LA
\TeDeum{Te Deum}
\switchcolumn
\EN
\TeDeum{Te Deum}
\switchcolumn*

\end{paracol}

\par\needspace{10\baselineskip}\addvspace{1.2em}{\centering\small\textsc{Die 11 Novembris} · \textsc{11 November}\par}
\Office{S. Martini Episcopi et Confessoris}{St. Martin, Bishop of Tours, Confessor}{Duplex}
\addcontentsline{toc}{subsection}{Die 11 Novembris · S. Martini Episcopi et Confessoris \textit{— St. Martin, Bishop of Tours, Confessor}}
\begin{paracol}{2}
\LA
\RefLine{Lectio IX: de S. Mennæ Martyris (11-11cc).}
\switchcolumn
\EN
\RefLine{Lesson IX: of S. Menna, Martyr (11-11cc).}
\switchcolumn*

\LA
\LessonHead{Lectio I}
\switchcolumn
\EN
\LessonHead{Lesson I}
\switchcolumn*

\LA
\LessonTitle{De Epístola prima beáti Pauli Apóstoli ad Timótheum}
\switchcolumn
\EN
\LessonTitle{Lesson from the first letter of St. Paul the Apostle to Timothy}
\switchcolumn*

\LA
\Cite{1 Tim 3:1-7}
\switchcolumn
\EN
\Cite{1 Tim 3:1-7}
\switchcolumn*

\LA
\noindent Fidélis sermo: si quis episcopátum desíderat, bonum opus desíderat. Opórtet ergo epíscopum irreprehensíbilem esse, uníus uxóris virum, sóbrium, prudéntem, ornátum, pudícum, hospitálem, doctórem, non vinoléntum, non percussórem, sed modéstum; non litigiósum, non cúpidum, sed suæ dómui bene præpósitum, fílios habéntem súbditos cum omni castitáte. Si quis autem dómui suæ præésse nescit, quómodo Ecclésiæ Dei diligéntiam habébit? Non neóphytum, ne in supérbiam elátus, in judícium íncidat diáboli. Opórtet autem illum et testimónium habére bonum ab iis qui foris sunt, ut non in oppróbrium íncidat, et in láqueum diáboli.\par
\switchcolumn
\EN
\noindent A faithful saying: if a man desire the office of a bishop, he desireth a good work. It behoveth therefore a bishop to be blameless, the husband of one wife, sober, prudent, of good behaviour, chaste, given to hospitality, a teacher, Not given to wine, no striker, but modest, not quarrelsome, not covetous, but One that ruleth well his own house, having his children in subjection with all chastity. But if a man know not how to rule his own house, how shall he take care of the church of God? Not a neophyte: lest being puffed up with pride, he fall into the judgment of the devil. Moreover he must have a good testimony of them who are without: lest he fall into reproach and the snare of the devil.\par
\switchcolumn*

\LA
\begin{Resp}\Rsign Hic est Martínus, eléctus Dei Póntifex, cui Dóminus post Apóstolos tantam grátiam conférre dignátus est, \Star{} Ut in virtúte Trinitátis Deíficæ mererétur fíeri trium mortuórum suscitátor magníficus.\end{Resp}
\begin{Resp}\Vsign Sanctæ Trinitátis fidem Martínus conféssus est.\end{Resp}
\begin{Resp}\Rsign Ut in virtúte Trinitátis Deíficæ mererétur fíeri trium mortuórum suscitátor magníficus.\end{Resp}
\switchcolumn
\EN
\begin{Resp}\Rsign This is that Martin whom God chose to be an High Priest unto Himself, he upon whom the Lord was pleased to bestow favour like as upon His Apostles; \Star{} So that he prevailed gloriously in the power of the Divine Trinity three times to raise the dead to life.\end{Resp}
\begin{Resp}\Vsign Martin confessed the faith of the Holy Trinity.\end{Resp}
\begin{Resp}\Rsign So that he prevailed gloriously in the power of the Divine Trinity three times to raise the dead to life.\end{Resp}
\switchcolumn*

\LA
\LessonHead{Lectio II}
\switchcolumn
\EN
\LessonHead{Lesson II}
\switchcolumn*

\LA
\LessonTitle{De Epístola ad Titum}
\switchcolumn
\EN
\LessonTitle{From the letter of St. Paul the Apostle to Titus}
\switchcolumn*

\LA
\Cite{Titus 1:7-11}
\switchcolumn
\EN
\Cite{Titus 1:7-11}
\switchcolumn*

\LA
\noindent Opórtet enim epíscopum sine crímine esse, sicut Dei dispensatórem: non supérbum, non iracúndum, non vinoléntum, non percussórem, non turpis lucri cúpidum: sed hospitálem, benígnum, sóbrium, justum, sanctum, continéntem, amplecténtem eum, qui secúndum doctrínam est, fidélem sermónem: ut potens sit exhortári in doctrína sana, et eos qui contradícunt, argúere. Sunt enim multi étiam inobediéntes, vaníloqui, et seductóres: máxime qui de circumcisióne sunt: quos opórtet redárgui: qui univérsas domos subvértunt, docéntes quæ non opórtet, turpis lucri grátia.\par
\switchcolumn
\EN
\noindent For a bishop must be without crime, as the steward of God: not proud, not subject to anger, not given to wine, no striker, not greedy of filthy lucre: But given to hospitality, gentle, sober, just, holy, continent: Embracing that faithful word which is according to doctrine, that he may be able to exhort in sound doctrine, and to convince the gainsayers. For there are also many disobedient, vain talkers, and seducers: especially they who are of the circumcision: Who must be reproved, who subvert whole houses, teaching things which they ought not, for filthy lucre’s sake.\par
\switchcolumn*

\LA
\begin{Resp}\Rsign Dómine, si adhuc pópulo tuo sum necessárius, non recúso subíre propter eos labórem: \Star{} Fiat volúntas tua.\end{Resp}
\begin{Resp}\Vsign Oculis ac mánibus in cælum semper inténtus, invíctum ab oratióne spíritum non relaxábat.\end{Resp}
\begin{Resp}\Rsign Fiat volúntas tua.\end{Resp}
\switchcolumn
\EN
\begin{Resp}\Rsign Lord, if I be still needful to thy people, I refuse not to work for them. \Star{} Thy will be done.\end{Resp}
\begin{Resp}\Vsign With eyes and hands lifted up to heaven, he never let his mighty spirit slacken in prayer.\end{Resp}
\begin{Resp}\Rsign Thy will be done.\end{Resp}
\switchcolumn*

\LA
\LessonHead{Lectio III}
\switchcolumn
\EN
\LessonHead{Lesson III}
\switchcolumn*

\LA
\Cite{Titus 2:1-8}
\switchcolumn
\EN
\Cite{Titus 2:1-8}
\switchcolumn*

\LA
\noindent Tu autem lóquere quæ decent sanam doctrínam: senes ut sóbrii sint, pudíci, prudéntes, sani in fide, in dilectióne, in patiéntia: anus simíliter in hábitu sancto, non criminatríces, non multo vino serviéntes, bene docéntes: ut prudéntiam dóceant adolescéntulas, ut viros suos ament, fílios suos díligant, prudéntes, castas, sóbrias, domus curam habéntes, benígnas, súbditas viris suis, ut non blasphemétur verbum Dei. Júvenes simíliter hortáre ut sóbrii sint. In ómnibus teípsum præbe exémplum bonórum óperum, in doctrína, in integritáte, in gravitáte, verbum sanum, irreprehensíbile: ut is qui ex advérso est, vereátur, nihil habens malum dícere de nobis.\par
\switchcolumn
\EN
\noindent But speak thou the things that become sound doctrine: That the aged men be sober, chaste, prudent, sound in faith, in love, in patience. The aged women, in like manner, in holy attire, not false accusers, not given to much wine, teaching well: That they may teach the young women to be wise, to love their husbands, to love their children, To be discreet, chaste, sober, having a care of the house, gentle, obedient to their husbands, that the word of God be not blasphemed. Young men, in like manner, exhort that they be sober. In all things shew thyself an example of good works, in doctrine, in integrity, in gravity, The sound word that can not be blamed: that he, who is on the contrary part, may be afraid, having no evil to say of us.\par
\switchcolumn*

\LA
\begin{Resp}\Rsign O beátum virum Martínum antístitem, \Star{} Qui nec mori tímuit, nec vívere recusávit!\end{Resp}
\begin{Resp}\Vsign Dómine, si adhuc pópulo tuo sum necessárius, non recúso labórem: fiat volúntas tua.\end{Resp}
\begin{Resp}\Rsign Qui nec mori tímuit, nec vívere recusávit!\end{Resp}
\begin{Resp}\Vsign Glória Patri, et Fílio, \Star{} et Spirítui Sancto.\end{Resp}
\begin{Resp}\Rsign Qui nec mori tímuit, nec vívere recusávit!\end{Resp}
\switchcolumn
\EN
\begin{Resp}\Rsign Oh, how blessed a man was Bishop Martin! \Star{} He neither feared to die, nor refused to live.\end{Resp}
\begin{Resp}\Vsign Lord, if I be still needful to thy people, I refuse not to work for them. thy will be done.\end{Resp}
\begin{Resp}\Rsign He neither feared to die, nor refused to live.\end{Resp}
\begin{Resp}\Vsign Glory be to the Father, and to the Son, \Star{} and to the Holy Ghost.\end{Resp}
\begin{Resp}\Rsign He neither feared to die, nor refused to live.\end{Resp}
\switchcolumn*

\LA
\LessonHead{Lectio IV}
\switchcolumn
\EN
\LessonHead{Lesson IV}
\switchcolumn*

\LA
\noindent Martínus, Sabáriæ in Pannónia natus, cum décimum attigísset annum, invítis paréntibus ad ecclésiam confúgiens, in catechumenórum númerum adscríbi vóluit. Quíndecim annos natus in milítiam proféctus, primum in Constántii, deínde Juliáni exércitu militávit. Qui, cum nihil habéret præter arma et vestiméntum quo tegebátur, Ambiáni, páuperi ac nudo, ab eo peténti ut Christi nómine sibi eleemósynam tribúeret, partem chlámydis dedit. Cui sequénti nocte Christus, dimidiáta illa veste indútus, appáruit, hanc mittens vocem: Martínus catechúmenus hac me veste contéxit.\par
\switchcolumn
\EN
\noindent Martin was born at Sabaria in Pannonia. When he was ten years old he went to the Church, in the spite of his (heathen) father and mother, and by his own will was numbered among the Catechumens. At fifteen years of age he joined the army, and served as a soldier first under Constantius and then under Julian. Once at the gate of Amiens a poor man asked him for an alms for Christ's name's sake, and since he had nothing to his hand but his arms and his clothes, he gave him half of his cloak. In the night following Christ appeared to him clad in the half of his cloak, and saying (to the angels who bare Him company) While Martin is yet a Catechumen, he hath clad Me in this garment.\par
\switchcolumn*

\LA
\begin{Resp}\Rsign Oculis ac mánibus in cælum semper inténtus, \Star{} Invíctum ab oratióne spíritum non relaxábat.\end{Resp}
\begin{Resp}\Vsign Dum sacraménta offérret beátus Martínus, globus ígneus appáruit super caput ejus.\end{Resp}
\begin{Resp}\Rsign Invíctum ab oratióne spíritum non relaxábat.\end{Resp}
\switchcolumn
\EN
\begin{Resp}\Rsign With eyes and hands lifted up to heaven \Star{} He never let his mighty spirit slacken in prayer.\end{Resp}
\begin{Resp}\Vsign While as blessed Martin was offering up the mysteries, a ball of fire appeared above his head.\end{Resp}
\begin{Resp}\Rsign He never let his mighty spirit slacken in prayer.\end{Resp}
\switchcolumn*

\LA
\LessonHead{Lectio V}
\switchcolumn
\EN
\LessonHead{Lesson V}
\switchcolumn*

\LA
\noindent Decem et octo annos cum habéret, baptizátus est. Quare, relícta militári vita, ad Hilárium Pictaviénsem epíscopum se cóntulit, a quo in acolythórum númerum redáctus est. Post, factus epíscopus Turonénsis, monastérium ædificávit, ubi cum octogínta mónachis sanctíssime aliquámdiu vixit. Qui, cum póstea ad Candacénsem vicum suæ diœcésis in gravem febrim incidísset, assídua Deum oratióne precabátur, ut se ex illo mortáli cárcere liberáret. Quem audiéntes discípuli, sic rogábant: Cur nos, pater déseris? cui nos míseros derelínquis? Quorum voce commótus Martínus, ita Deum orábat: Dómine, si adhuc pópulo tuo sum necessárius, non recúso labórem.\par
\switchcolumn
\EN
\noindent At eighteen years of age he was baptized. He gave up thereupon the life of a soldier, and betook himself to Hilary, Bishop of Poietiers, by whom he was placed in the order of Acolytes. Being afterwards made Bishop of Tours, he built a monastery wherein he lived in holiness for a while in company of four-score monks. At the last he fell sick of a grievous fever at Cande, a village in his diocese, and besought God in constant prayer to set him free from the prison of this dying body. His disciples heard him and said Father, why wilt thou go away from us? unto whom wilt thou bequeath us in our sorrow? Their words moved Martin, and he said Lord, if I be still needful to thy people, I refuse not to work.\par
\switchcolumn*

\LA
\begin{Resp}\Rsign Beátus Martínus óbitum suum longe ante præscívit, dixítque frátribus, \Star{} Dissolutiónem sui córporis imminére, quia judicábat se jam resólvi.\end{Resp}
\begin{Resp}\Vsign Víribus córporis cœpit repénte destítui, convocatísque discípulis dixit.\end{Resp}
\begin{Resp}\Rsign Dissolutiónem sui córporis imminére, quia judicábat se jam resólvi.\end{Resp}
\switchcolumn
\EN
\begin{Resp}\Rsign Blessed Martin knew of his own death of a long time before it came to pass, and he said unto the brethren: \Star{} That the dissolution of his body was nigh at hand, for he deemed himself to be already breaking up.\end{Resp}
\begin{Resp}\Vsign His bodily strength gave way all of a sudden, and he called his disciples together, and said unto them\end{Resp}
\begin{Resp}\Rsign That the dissolution of his body was nigh at hand, for he deemed himself to be already breaking up.\end{Resp}
\switchcolumn*

\LA
\LessonHead{Lectio VI}
\switchcolumn
\EN
\LessonHead{Lesson VI}
\switchcolumn*

\LA
\noindent Sed, cum eum in illa veheménti febre supínum orántem vidérent discípuli, supplíciter ab eo petiérunt, ut, convérso córpore, tantísper, dum remítteret morbi vis, pronus conquiésceret. Quibus Martínus, Sínite me, inquit, cælum pótius quam terram aspícere, ut, suo jam itínere itúrus ad Dóminum, spíritus dirigátur. Instánte jam morte, viso humáni géneris hoste, Quid, inquit, astas, cruénta béstia? nihil in me funéste repéries. Ea in voce, unum et octogínta annos natus, ánimam Deo réddidit; quam Angelórum chorus excépit, eósque divínas canéntes laudes multi, in primísque sanctus Severínus Coloniénsis epíscopus, audiérunt.\par
\switchcolumn
\EN
\noindent When his disciples saw him, in the height of the fever, lying upon his back and praying, they entreated him to turn over and take a little rest upon his side while the violence of his sickness would allow him. But Martin answered them Suffer me to look heavenward rather than earthward, that my spirit may see the way whereby it is so soon going to the Lord. At the moment of death he saw the enemy of mankind, and cried out: What are you come here for, you bloody brute? You murderer, you'll find nothing in me. With these words on his lips, he gave up his soul to God, being aged eighty years. He was received by a company of Angels, who were heard praising God by many persons, especially by holy Severinus, Bishop of Cologne.\par
\switchcolumn*

\LA
\begin{Resp}\Rsign Dixérunt discípuli ad beátum Martínum: Cur nos, pater, déseris, aut cui nos desolátos relínquis? \Star{} Invádent enim gregem tuum lupi rapáces.\end{Resp}
\begin{Resp}\Vsign Scimus quidem desideráre te Christum, sed salva sunt tibi tua prǽmia: nostri pótius miserére, quos déseris.\end{Resp}
\begin{Resp}\Rsign Invádent enim gregem tuum lupi rapáces.\end{Resp}
\begin{Resp}\Vsign Glória Patri, et Fílio, \Star{} et Spirítui Sancto.\end{Resp}
\begin{Resp}\Rsign Invádent enim gregem tuum lupi rapáces.\end{Resp}
\switchcolumn
\EN
\begin{Resp}\Rsign His disciples said unto blessed Martin: Father, why wilt thou go away from us, and with whom wilt thou leave us orphans? \Star{} For ravening wolves will break in upon thy flock.\end{Resp}
\begin{Resp}\Vsign We know that thou wouldest fain be with Christ, but, sooner or later, thy reward is sure. Rather, then, have pity upon us, whom thou art leaving.\end{Resp}
\begin{Resp}\Rsign For ravening wolves will break in upon thy flock.\end{Resp}
\begin{Resp}\Vsign Glory be to the Father, and to the Son, \Star{} and to the Holy Ghost.\end{Resp}
\begin{Resp}\Rsign For ravening wolves will break in upon thy flock.\end{Resp}
\switchcolumn*

\LA
\LessonHead{Lectio VII}
\switchcolumn
\EN
\LessonHead{Lesson VII}
\switchcolumn*

\LA
\LessonTitle{Léctio sancti Evangélii secúndum Lucam}
\switchcolumn
\EN
\LessonTitle{From the Holy Gospel according to Luke}
\switchcolumn*

\LA
\Cite{Luc 11:33-36}
\switchcolumn
\EN
\Cite{Luke 11:33-36}
\switchcolumn*

\LA
\noindent In illo témpore: Dixit Jesus discípulis suis: Nemo lucérnam accéndit et in abscóndito ponit neque sub módio, sed supra candelábrum, ut qui ingrediúntur, lumen vídeant. Et réliqua.\par
\switchcolumn
\EN
\noindent In that time: Jesus said to his disciples: No man lighteth a candle, and putteth it in a hidden place, nor under a bushel; but upon a candlestick, that they that come in, may see the light. And so on.\par
\switchcolumn*

\LA
\LessonTitle{Homilía sancti Ambrósii Epíscopi}
\switchcolumn
\EN
\LessonTitle{Homily by St. Ambrose, Bishop (of Milan.)}
\switchcolumn*

\LA
\Source{Liber 7 Comment. in Luc. cap. 11, post initium}
\switchcolumn
\EN
\Source{Bk. vii. Comment, on Luke xi.}
\switchcolumn*

\LA
\noindent Quia in superióribus Ecclésiam Synagógæ prǽtulit, hortátur nos ut fidem pótius nostram ad Ecclésiam transferámus. Lucérna enim fides est, juxta quod scriptum est: Lucérna pédibus meis verbum tuum, Dómine. Verbum enim Dei fides nostra est; Verbum Dei lux est; Lucérna est fides: Erat Lux vera, quæ illúminat omnem hóminem veniéntem in hunc mundum. Lucérna autem lucére non potest, nisi aliúnde lumen accéperit.\par
\switchcolumn
\EN
\noindent In that which goeth before, Christ hath set the Church before the synagogue, and He exhorteth us rather to trust in the Church. The candle is faith, even as it is written thy word is a lamp unto my feet, and a light unto my path. (Ps. cxviii. 105.) Our faith is the word of God. The word of God is light. Faith is the candle. It is written concerning the Word of God, that That was the true Light, Which lighteth every man that cometh into this world. (John i. 9.) But a candle cannot shine, unless it be lighted from some other fire.\par
\switchcolumn*

\LA
\begin{Resp}\Rsign O beátum virum, in cujus tránsitu Sanctórum canit númerus, Angelórum exsúltat chorus, \Star{} Omniúmque cæléstium Virtútum occúrrit psalléntium exércitus!\end{Resp}
\begin{Resp}\Vsign Ecclésia virtúte roborátur, sacerdótes Dei revelatióne glorificántur, quem Michaël assúmpsit cum Ángelis.\end{Resp}
\begin{Resp}\Rsign Omniúmque cæléstium Virtútum occúrrit psalléntium exércitus!\end{Resp}
\switchcolumn
\EN
\begin{Resp}\Rsign Blessed indeed was this man, at the time of whose passing the Saints sang in company, a band of Angels shouted aloud for joy \Star{} And an army of all the Powers of heaven came out to meet him, singing praises.\end{Resp}
\begin{Resp}\Vsign His strength is a stay to the Church, his manifestation a glory to the Priests of God; Michael and his Angels took him away.\end{Resp}
\begin{Resp}\Rsign And an army of all the Powers of heaven came out to meet him, singing praises.\end{Resp}
\switchcolumn*

\LA
\LessonHead{Lectio VIII}
\switchcolumn
\EN
\LessonHead{Lesson VIII}
\switchcolumn*

\LA
\noindent Hæc est Lucérna quæ accénditur, virtus scílicet nostræ mentis et sensus, ut drachma illa possit, quæ períerat, reperíri. Nemo ergo fidem sub lege constítuat; lex enim intra mensúram est, ultra mensúram grátia; lex obúmbrat, grátia claríficat. Et ídeo nemo fidem suam intra mensúram legis inclúdat, sed ad Ecclésiam cónferat, in qua septifórmis Spíritus relúcet grátia, quam Princeps ille sacerdótum fulgóre supérnæ divinitátis illúminat, ne eam legis umbra restínguat. Dénique lucérna illa, quam matutínis vespertinísque tempóribus, ritu véteri Judæórum, princeps sacerdótum solébat accéndere, velut sub módio sita legis, evánuit; et cívitas illa Jerúsalem quæ in terris est, quæ occídit prophétas, quasi in conválle fletus pósita delitéscit. Illa autem Jerúsalem quæ in cælo est, in qua mílitat fides nostra, in illo altíssimo ómnium locáta monte, hoc est Christo, non potest ténebris et ruínis hujus mundi abscóndi; sed fulgens candóre Solis ætérni, luce nos grátiæ spiritális illúminat.\par
\switchcolumn
\EN
\noindent Also it is written What woman, having ten pieces of silver, if she lose one piece, doth not light a candle, and sweep the house, and seek diligently till she find it? (Luke xv. 8.) And here the candle lighted to find the lost piece is the strength in our understandings and affections. Let no man therefore seek faith under the law. For the law is by measure, but grace without measure; the law overshadoweth, but grace enlighteneth. And therefore let no man shut up his faith within the measure of the law, but give it unto the Church, the Church, wherein shineth the sevenfold grace of the Spirit, and whereon the Divine glory of the Great High Priest doth strike from heaven, lest the shadow of the law should rest any more at all upon her. Under the old law there was the sevenfold lamp which the Priest of the Jews lighted every morning and every evening, and this was as it were a candle put under a bushel. That Jerusalem which is upon earth, that Jerusalem which killed the Prophets, lieth hid, as it were, in a dark place in the valley of tears. But that Jerusalem which is in heaven, whereof by faith we are soldiers, is a city set upon the highest of all mountains, even upon Christ. Her the darkness and tempests of earth cannot hide, but she blazeth with the glory of the Eternal Sun, and maketh to fall upon us the light of spiritual grace.\par
\switchcolumn*

\LA
\begin{Resp}\Rsign Martínus Ábrahæ sinu lætus excípitur: Martínus, hic pauper et módicus, \Star{} Cælum dives ingréditur, hymnis cæléstibus honorátur.\end{Resp}
\begin{Resp}\Vsign Martínus epíscopus migrávit a sǽculo: vivit in Christo gemma sacerdótum.\end{Resp}
\begin{Resp}\Rsign Cælum dives ingréditur, hymnis cæléstibus honorátur.\end{Resp}
\begin{Resp}\Vsign Glória Patri, et Fílio, \Star{} et Spirítui Sancto.\end{Resp}
\begin{Resp}\Rsign Cælum dives ingréditur, hymnis cæléstibus honorátur.\end{Resp}
\switchcolumn
\EN
\begin{Resp}\Rsign Martin was carried joyfully into Abraham's bosom. Martin, who was poor here and of small estimation, \Star{} Etereth rich into heaven, and the songs of heaven are raised in his honour.\end{Resp}
\begin{Resp}\Vsign Bishop Martin, that jewel of Priests, goeth away from time, liveth in Christ.\end{Resp}
\begin{Resp}\Rsign Entereth rich into heaven, and the songs of heaven are raised in his honour.\end{Resp}
\begin{Resp}\Vsign Glory be to the Father, and to the Son, \Star{} and to the Holy Ghost.\end{Resp}
\begin{Resp}\Rsign Entereth rich into heaven, and the songs of heaven are raised in his honour.\end{Resp}
\switchcolumn*

\end{paracol}

\par\needspace{10\baselineskip}\addvspace{1.2em}{\centering\small\textsc{Die 11 Novembris} · \textsc{11 November}\par}
\Office{S. Mennæ Martyris}{S. Menna, Martyr}{Simplex}
\addcontentsline{toc}{subsection}{Die 11 Novembris · S. Mennæ Martyris \textit{— S. Menna, Martyr}}
\begin{paracol}{2}
\LA
\LessonHead{Lectio IX (pro commemoratione)}
\switchcolumn
\EN
\LessonHead{Lesson IX (when commemorated)}
\switchcolumn*

\LA
\LessonTitle{Commemoratio: S. Mennæ Martyris}
\switchcolumn
\EN
\LessonTitle{Commemoration of: S. Menna, Martyr}
\switchcolumn*

\LA
\noindent Mennas Ægýptius, christiánus miles, in persecutióne Diocletiáni et Maximiáni imperatórum, cum pœniténtiæ causa in solitúdinem secessísset, natáli die imperatórum, quo pópulus célebri spectáculo tenebátur, in theátrum prosíliens, líbera voce Gentílium superstitiónem insectabátur. Quam ob rem comprehénsus, et, Pyrrho prǽside, in metrópoli Cottiénsium Phrýgiæ vinctus, loris crudéliter cœditur. Deínde, equúleo tortus, lampádibus ardéntibus ad corpus admótis plagísque cilício confricátis, tum per tríbulos et virgas férreas mánibus ac pédibus colligátis tractus, plumbátis étiam contúsus, demum gládio interfícitur in ignémque conjícitur. Corpus, inde a Christiánis eréptum, sepúltum est ac póstea Constantinópolim translátum.\par
\switchcolumn
\EN
\noindent Mennas was a Christian Egyptian soldier who had withdrawn himself into a desert place to do penance, but one day, during the persecution under the Emperors Diocletian and Maximian, upon the said Emperors' birthday, when the people were gathered together at a great show, stood forth in the theatre, and reviled with a loud voice the idolatries of the Gentiles. Thereupon he was arrested, and being bound at Cotyasus, the chief city of Phrygia, under the authority of the President Pyrrhus, was first savagely lashed with thongs, then racked, then scarified with fire applied to his naked body, then had his wounds lacerated by rubbing with hair-cloth, then dragged through thorns and iron spikes with hands and feet tied, then lashed again with whips loaded with lead, and lastly slain with the sword, and thrown into a fire. The Christians saved his body thence and buried it, and it hath in after times been carried to Constantinople.\par
\switchcolumn*

\LA
\TeDeum{Te Deum}
\switchcolumn
\EN
\TeDeum{Te Deum}
\switchcolumn*

\end{paracol}

\par\needspace{10\baselineskip}\addvspace{1.2em}{\centering\small\textsc{Die 12 Novembris} · \textsc{12 November}\par}
\Office{S. Martini Papæ et Martyris}{St. Martin I, Pope and Martyr}{Semiduplex}
\addcontentsline{toc}{subsection}{Die 12 Novembris · S. Martini Papæ et Martyris \textit{— St. Martin I, Pope and Martyr}}
\begin{paracol}{2}
\LA
\RefLine{Lectiones I–III: de Scriptura occurrente.}
\switchcolumn
\EN
\RefLine{Lessons I–III: from the Scripture of the day.}
\switchcolumn*

\LA
\LessonHead{Lectio IV}
\switchcolumn
\EN
\LessonHead{Lesson IV}
\switchcolumn*

\LA
\noindent Martinus, Tuderti in Umbria natus, inítio pontificatus, et litteris et legatiónibus missis operam dedit, ut Paulum Constantinopolitanum patriarcham a nefaria hæresi ad catholicæ fidei veritátem revocaret. Qui, Constante imperatóre hæretico fretus, eo améntiæ progréssus fuerat, ut Sedis apostolicæ legatos varie in insulas relegarit. Quo ejus scélere commótus Pontifex, coacto Romæ concílio centum quinque episcopórum, eum condemnávit.\par
\switchcolumn
\EN
\noindent Martin was born at Todi in Tuscany. At the beginning of his Popedom, (in the year 649,) he was careful to send an embassage with letters to Paul, Patriarch of Constantinople, to call upon him to return to the truth of the Catholic faith from the blasphemous heresy (of the Monothelites.) But Paul, being backed up by the heretic Emperor Constans, had become so rabid, that he sent away the messengers of the Apostolic See into diverse places in the islands. This crime moved the Pope to gather together at Rome a council of one - hundred - and five Bishops, by whom Paul was condemned.\par
\switchcolumn*

\LA
\begin{Resp}\Rsign Honéstum fecit illum Dóminus, et custodívit eum ab inimícis, et a seductóribus tutávit illum: \Star{} Et dedit illi claritátem ætérnam.\end{Resp}
\begin{Resp}\Vsign Descendítque cum illo in fóveam, et in vínculis non derelíquit eum.\end{Resp}
\begin{Resp}\Rsign Et dedit illi claritátem ætérnam.\end{Resp}
\switchcolumn
\EN
\begin{Resp}\Rsign The Lord made him honourable, and defended him from his enemies, and kept him safe from those that lay in wait for him, \Star{} And gave him perpetual glory.\end{Resp}
\begin{Resp}\Vsign He went down with him into the pit, and left him not in bonds.\end{Resp}
\begin{Resp}\Rsign And gave him perpetual glory.\end{Resp}
\switchcolumn*

\LA
\LessonHead{Lectio V}
\switchcolumn
\EN
\LessonHead{Lesson V}
\switchcolumn*

\LA
\noindent Quæ causa fuit Constanti mitténdi in Italiam Olympium exarchum, ut Martinum Pontificem interficiéndum aut ad se perducéndum curaret. Igitur Olympius, Romam veniens, lictori mandat ut Pontificem, dum in basilica sanctæ Maríæ ad Præsepe Missárum solemnia celebraret, occíderet. Quod ubi lictor aggréditur, cæcus repénte factus est.\par
\switchcolumn
\EN
\noindent Thereupon Constans sent Olympius into Italy as Exarch, straitly commanding him either to slay Pope Martin, or else to bring him into his Imperial presence. Olympius therefore came to Rome and bade a lictor to kill the Pope while as he was celebrating the Liturgy solemnly in the Cathedral Church of St Mary-at-the-Manger. But when the lictor went thither, he was struck with blindness.\par
\switchcolumn*

\LA
\begin{Resp}\Rsign Desidérium ánimæ ejus tribuísti ei Dómine, \Star{} Et voluntáte labiórum ejus non fraudásti eum.\end{Resp}
\begin{Resp}\Vsign Quóniam prævenísti eum in benedictiónibus dulcédinis: posuísti in cápite ejus corónam de lápide pretióso.\end{Resp}
\begin{Resp}\Rsign Et voluntáte labiórum ejus non fraudásti eum.\end{Resp}
\switchcolumn
\EN
\begin{Resp}\Rsign Lord, Thou hast given him his heart's desire. \Star{} And hast not withholden the request of his lips.\end{Resp}
\begin{Resp}\Vsign For Thou hast prevented him with the blessings of sweetness: Thou hast set a crown of precious stones upon his head.\end{Resp}
\begin{Resp}\Rsign And hast not withholden the request of his lips.\end{Resp}
\switchcolumn*

\LA
\LessonHead{Lectio VI}
\switchcolumn
\EN
\LessonHead{Lesson VI}
\switchcolumn*

\LA
\noindent Constanti autem imperatóri ex eo témpore multæ calamitátes incidérunt; quibus níhilo mélior factus, Theodórum Callíopam ad Urbem mittens, imperat ut Pontifici manus injíciat. A quo per fraudem captus Martinus et Constantinopolim perductus, deínde in Chersonesum relegátus; ibi ob catholicam fidem ærumnis confectus, sextodecimo Kalendas Octobris cessit e vita, clarus miraculis. Cujus corpus, Romam póstea translátum, in ecclésia conditum est, quæ sanctórum Silvestri et Martini nómine dedicáta erat. Præfuit Ecclésiæ annos sex, mensem unum, dies viginti sex. Habuit ordinatiónes duas mense Decembri, quibus creávit presbyteros undecim, diaconos quinque, episcopos per diversa loca trigínta tres.\par
\switchcolumn
\EN
\noindent From that time forth many evils befell the Emperor Constans; but he repented not. He sent the Exarch Theodore Calliopas to Rome, with command to lay hands on the Pope. By him Martin was treacherously taken (on the 17th day of June, 653,) and (forthwith carried to the island of Naxos. On the 17th of September in 654 he was) brought to Constantinople, (where he was kept in prison) till he was sent to the Crimea (on the 15th of May, 655.) There his sufferings for the Catholic faith utterly broke him down, and he left this life for a better, upon the 12th day of November, (in the same year 655.) He was famous for miracles. His body was afterwards brought back to Rome and buried in the Church dedicated under the names of St. Silvester and St. Martin (of Tours.) He ruled the Church for six years, one month, and twenty-six days. He held two ordinations in the month of December, wherein he made eleven Priests, five Deacons, and thirtythree Bishops for diverse places.\par
\switchcolumn*

\LA
\begin{Resp}\Rsign Stola jucunditátis índuit eum Dóminus: \Star{} Et corónam pulchritúdinis pósuit super caput ejus.\end{Resp}
\begin{Resp}\Vsign Cibávit illum Dóminus pane vitæ et intelléctus: et aqua sapiéntiæ salutáris potávit illum.\end{Resp}
\begin{Resp}\Rsign Et corónam pulchritúdinis pósuit super caput ejus.\end{Resp}
\begin{Resp}\Vsign Glória Patri, et Fílio, \Star{} et Spirítui Sancto.\end{Resp}
\begin{Resp}\Rsign Et corónam pulchritúdinis pósuit super caput ejus.\end{Resp}
\switchcolumn
\EN
\begin{Resp}\Rsign The Lord hath put on him a robe of honour, \Star{} And put about his head a crown of joy.\end{Resp}
\begin{Resp}\Vsign With the bread of life and understanding hath the Lord fed him, and given him the water of health and wisdom to drink.\end{Resp}
\begin{Resp}\Rsign And put about his head a crown of joy.\end{Resp}
\begin{Resp}\Vsign Glory be to the Father, and to the Son, \Star{} and to the Holy Ghost.\end{Resp}
\begin{Resp}\Rsign And put about his head a crown of joy.\end{Resp}
\switchcolumn*

\LA
\LessonHead{Lectio VII}
\switchcolumn
\EN
\LessonHead{Lesson VII}
\switchcolumn*

\LA
\LessonTitle{Léctio sancti Evangélii secúndum Matthǽum}
\switchcolumn
\EN
\LessonTitle{From the Holy Gospel according to Matthew}
\switchcolumn*

\LA
\Cite{Matt 16:13-19}
\switchcolumn
\EN
\Cite{Matt 16:13-19}
\switchcolumn*

\LA
\noindent In illo témpore: Venit Jesus in partes Cæsaréæ Philíppi, et interrogábat discipulos suos, dicens: Quem dicunt hómines esse Fílium hóminis? Et réliqua.\par
\switchcolumn
\EN
\noindent At that time: Jesus came into the coasts of Cassarea Philippi, and He asked His disciples, saying: Who do men say that I, the Son of Man, am? And so on.\par
\switchcolumn*

\LA
\LessonTitle{Homilía sancti Leónis Papæ.}
\switchcolumn
\EN
\LessonTitle{Homily by Pope St. Leo (the Great)}
\switchcolumn*

\LA
\Source{Sermo 2 in anniversari assumptionis suæ ante medium}
\switchcolumn
\EN
\Source{2nd on the anniversary of his own election.}
\switchcolumn*

\LA
\noindent Cum, sicut evangelica lectióne reserátum est, interrogasset Dóminus discipulos, quem ipsum (multis diversa opinántibus) crederent; respondissétque beátus Petrus, dicens: Tu es Christus Fílius Dei vivi; Dóminus ait: Beátus es, Simon Bar-Jona, quia caro et sanguis non revelávit tibi, sed Pater meus, qui in cælis est: et ego dico tibi, quia tu es Petrus, et super hanc petram ædificábo Ecclésiam meam, et portæ ínferi non prævalébunt advérsus eam. Et tibi dabo claves regni cælórum: et quodcúmque ligáveris super terram, erit ligátum et in cælis: et quodcúmque solveris super terram, erit solutum et in cælis. Manet ergo disposítio veritátis, et beátus Petrus, in accepta fortitúdine petræ persevérans, suscepta Ecclésiæ gubernácula non relíquit.\par
\switchcolumn
\EN
\noindent When the Lord, as we read in the Evangelist, asked His disciples: Who did men, amid their diverse speculations, believe that He, the Son of Man, was; blessed Peter answered and said: Thou art the Christ, the Son of the living God. And Jesus answered and said unto him: Blessed art thou, Simon Barjona: for flesh and blood hath not revealed it unto thee, but My Father, Which is in: and I say also unto thee, that thou art Peter, and upon this rock I will build My Church, and the gates of hell shall not prevail against it; and I will give unto thee the keys of the kingdom of heaven; and whatsoever thou shalt bind on earth shall be bound in heaven; and whatsoever thou shalt loose on earth shall be loosed in heaven. Thus therefore standeth the ordinance of the Truth, and blessed Peter, abiding still that firm rock which God hath made him, hath never lost that right to rule in the Church which God hath given unto him.\par
\switchcolumn*

\LA
\begin{Resp}\Rsign Coróna áurea super caput ejus, \Star{} Expréssa signo sanctitátis, glória honóris, et opus fortitúdinis.\end{Resp}
\begin{Resp}\Vsign Quóniam prævenísti eum in benedictiónibus dulcédinis, posuísti in cápite ejus corónam de lápide pretióso.\end{Resp}
\begin{Resp}\Rsign Expréssa signo sanctitátis, glória honóris, et opus fortitúdinis.\end{Resp}
\switchcolumn
\EN
\begin{Resp}\Rsign A crown of gold upon his head, \Star{} Wherein is engraved Holiness, an ornament of honour, a costly work.\end{Resp}
\begin{Resp}\Vsign For Thou hast prevented him with the blessings of sweetness, Thou hast set a crown of precious stones upon his head.\end{Resp}
\begin{Resp}\Rsign Wherein is engraved Holiness, an ornament of honour, a costly work.\end{Resp}
\switchcolumn*

\LA
\LessonHead{Lectio VIII}
\switchcolumn
\EN
\LessonHead{Lesson VIII}
\switchcolumn*

\LA
\noindent In univérsa namque Ecclésia, Tu es Christus Fílius Dei vivi, quotídie Petrus dicit; et omnis lingua, quæ confitétur Dóminum, magisterio hujus vocis imbuitur. Hæc fides diabolum vincit et captivórum ejus víncula dissolvit. Hæc érutos mundo, ínserit cælo, et portæ ínferi advérsus eam prævalere non possunt. Tanta enim divinitus soliditate muníta est, ut eam neque hæretica umquam corrumpere pravitas, nec pagána potúerit superare perfidia. His ítaque modis, dilectíssimi, rationábile obsequio celebrétur hodiérna festívitas: ut in persona humilitátis meæ ille intelligátur, ille honorétur, in quo et ómnium pastórum sollicitudo, cum commendatárum sibi óvium custódia persevérat, et cujus étiam dignitas in indigno heréde non déficit.\par
\switchcolumn
\EN
\noindent In the universal Church it is Peter that doth still say every day, Thou art the Christ, the Son of the living God, and every tongue which confesseth that Jesus is Lord is taught that confession by the teaching of Peter. This is the faith that overcometh the devil and looseth the bands of his prisoners. This is the faith which maketh men free of the world and bringeth them to heaven, and the gates of hell are impotent to prevail against it. With such ramparts of salvation hath God fortified this rock, that the contagion of heresy will never be able to infect it, nor idolatry and unbelief to overcome it. This teaching it is, my dearly beloved brethren, which maketh the keeping of this Feast to-day to be our reasonable service, even the teaching which maketh you to know and honour in myself, lowly though I be, that Peter who is still entrusted with the care of all other shepherds and of all the flocks to them committed, and whose authority I have, albeit unworthy to be his heir.\par
\switchcolumn*

\LA
\begin{Resp}\Rsign Dómine, prævenísti eum in benedictiónibus dulcédinis: \Star{} Posuísti in cápite ejus corónam de lápide pretióso.\end{Resp}
\begin{Resp}\Vsign Vitam pétiit a te, et tribuísti ei longitúdinem diérum in sæculum sæculi.\end{Resp}
\begin{Resp}\Rsign Posuísti in cápite ejus corónam de lápide pretióso.\end{Resp}
\begin{Resp}\Vsign Glória Patri, et Fílio, \Star{} et Spirítui Sancto.\end{Resp}
\begin{Resp}\Rsign Posuísti in cápite ejus corónam de lápide pretióso.\end{Resp}
\switchcolumn
\EN
\begin{Resp}\Rsign O Lord, thou hast prevented him with blessings of sweetness: \Star{} Thou hast set on his head a crown of precious stones.\end{Resp}
\begin{Resp}\Vsign He asked life of thee: \Star{} and thou hast given him length of days for ever and ever.\end{Resp}
\begin{Resp}\Rsign Thou hast set on his head a crown of precious stones.\end{Resp}
\begin{Resp}\Vsign Glory be to the Father, and to the Son, \Star{} and to the Holy Ghost.\end{Resp}
\begin{Resp}\Rsign Thou hast set on his head a crown of precious stones.\end{Resp}
\switchcolumn*

\LA
\LessonHead{Lectio IX}
\switchcolumn
\EN
\LessonHead{Lesson IX}
\switchcolumn*

\LA
\noindent Cum ergo cohortatiónes nostras áuribus vestræ sanctitátis adhibemus, ipsum vobis, cujus vice fungimur, loqui crédite: quia et illíus vos afféctu monémus, et non aliud vobis, quam quod docuit, prædicámus; obsecrántes, ut succincti lumbos mentis vestræ, castam et sobriam vitam in Dei timóre ducatis. Corona mea, sicut Apóstolus ait, et gáudium vos estis, si fides vestra, quæ ab inítio Evangélii in universo mundo prædicáta est, in dilectióne et sanctitáte permanserit. Nam licet omnem Ecclésiam, quæ in toto est orbe terrárum, cunctis oporteat florere virtútibus; vos tamen præcipue inter ceteros pópulos decet meritis pietátis excellere, quos in ipsa apostolicæ petræ arce fundatos, et Dóminus noster Jesus Christus cum ómnibus redémit, et beátus Apóstolus Petrus præ ómnibus erudívit.\par
\switchcolumn
\EN
\noindent When therefore, we address our exhortations to your godly ears,\par
1 Believe ye that ye are hearing him speak whose office we are discharging. Yea, it is with his love for you that we warn you, and we preach unto you no other thing than that which he taught, entreating you that ye would gird up the loins of your mind and lead pure and sober lives in the fear of God. 2 My disciples dearly beloved, ye are to me, as the disciples of the Apostle Paul were to him, \textasciitilde{}(Phil. iv. 1,) a crown and a joy, if your faith, which, in the first times of the Gospel, was spoken of throughout the whole world, (Rom. i. 8,) abide still lovely and holy. For, albeit it behoveth the whole Church which is spread throughout all the world, to be strong in righteousness, you it chiefly becometh above all other peoples to excel in worth and godliness, whose house is built upon the very crown of the Rock of the Apostle, and whom not only hath our Lord Jesus Christ, as He hath redeemed all men, but whom also His blessed Apostle Peter hath made the foremost object of his teaching.\par
\switchcolumn*

\LA
\TeDeum{Te Deum}
\switchcolumn
\EN
\TeDeum{Te Deum}
\switchcolumn*

\end{paracol}

\par\needspace{10\baselineskip}\addvspace{1.2em}{\centering\small\textsc{Die 13 Novembris} · \textsc{13 November}\par}
\Office{S. Didaci Confessoris}{St. Didacus, Confessor}{Semiduplex}
\addcontentsline{toc}{subsection}{Die 13 Novembris · S. Didaci Confessoris \textit{— St. Didacus, Confessor}}
\begin{paracol}{2}
\LA
\RefLine{Lectiones I–III: de Scriptura occurrente.}
\switchcolumn
\EN
\RefLine{Lessons I–III: from the Scripture of the day.}
\switchcolumn*

\LA
\LessonHead{Lectio IV}
\switchcolumn
\EN
\LessonHead{Lesson IV}
\switchcolumn*

\LA
\noindent Dídacus, Hispanus, ex oppido sancti Nicolai de Portu diœcesis Hispalénsis, ab ineunte ætate, pii sub sacerdotis disciplína, sanctioris vitæ, solitaria in ecclésia, tirocinium exercuit. Deínde, ut firmius Deo se conjungeret, in convéntu de Arizzáfa fratrum Minórum (quos Observántes vocant) sancti Francisci regulam in statu laicali proféssus est. Magna ibi alacritate húmilis obediéntiæ et regularis observantiæ jugum subiens, contemplatióni in primis déditus, mira Dei luce perfundebátur, adeo ut de rebus cæléstibus, litterárum expers, mirándum in modum et plane divinitus loquerétur.\par
\switchcolumn
\EN
\noindent Diego was a Spaniard, and was born at the little town of San Nicola-del-Porto, in the diocese of Seville. From his childhood he learnt the more holy life under a godly Priest, (who lived hermit) in a lonely Church, and so served his apprenticeship. Afterwards, being fain to be more utterly God's only, he professed himself as a lay brother under the Rule of St. Francis in the convent of the Friars Minor, called Observant, of Arrizafa. There he cheerfully bore the yoke of the lowliest obedience and the strictest observance. He was much given to contemplation, and a wonderful light from God shone in him, so that, though he was untaught, he could speak touching heavenly things strangely and as it were supernaturally.\par
\switchcolumn*

\LA
\begin{Resp}\Rsign Honéstum fecit illum Dóminus, et custodívit eum ab inimícis, et a seductóribus tutávit illum: \Star{} Et dedit illi claritátem ætérnam.\end{Resp}
\begin{Resp}\Vsign Justum dedúxit Dóminus per vias rectas, et osténdit illi regnum Dei.\end{Resp}
\begin{Resp}\Rsign Et dedit illi claritátem ætérnam.\end{Resp}
\switchcolumn
\EN
\begin{Resp}\Rsign The Lord made him honourable, and defended him from his enemies, and kept him safe from those that lay in wait for him; \Star{} And gave him perpetual glory.\end{Resp}
\begin{Resp}\Vsign He went down with him into the pit, and left him not in bonds.\end{Resp}
\begin{Resp}\Rsign And gave him perpetual glory.\end{Resp}
\switchcolumn*

\LA
\LessonHead{Lectio V}
\switchcolumn
\EN
\LessonHead{Lesson V}
\switchcolumn*

\LA
\noindent Canariis in insulis, ubi frátribus sui ordinis præfuit, multa perpéssus, martyrii æstuans desidério, plures infidéles verbo et exemplo ad Christi fidem convértit. Romam veniens anno jubilæi, Nicolao quinto Pontifice, ægrotórum curæ in convéntu Aræ cæli destinátus, eo caritátis afféctu munus hoc exercuit, ut, Urbe annonæ inópia laborante, ægrotis tamen, quorum aliquándo úlcera étiam lambéndo abstergebat, nihil pénitus necessarii defécerit. Eximia quoque fides et grátia curatiónum in eo eluxit, cum lámpadis, quæ collucebat ante imáginem beatíssimæ Dei Genitrícis, quam summa devotióne colebat, oleo ægros inungens, signo crucis impresso, multórum morbos mirabíliter sanaverit.\par
\switchcolumn
\EN
\noindent In the Canary Islands, where he was warden of the brethren of his Order, he underwent much, earnestly willing to be a martyr, and by his word and example brought many unbelievers to Christ. He came to Rome in the year of the Jubilee, \textasciitilde{}(being that of our Lord 1450,) in the reign of Pope Nicolas V., and there was set to tend the. sick in the Convent of Ara Coeli, which work he did with such love, that although the city was plagued with a famine, the sufferers (whose sores he would sometimes cleanse even with his tongue) scarcely lacked anything needful. He was a burning and shining light of faith, and had the gift of healing, taking the oil from the lamp which burned before the image of the most blessed Mother of God, to whom he was earnestly devoted, and anointing the sick therewith, whereupon many were marvellously cured.\par
\switchcolumn*

\LA
\begin{Resp}\Rsign Amávit eum Dóminus, et ornávit eum: stolam glóriæ índuit eum, \Star{} Et ad portas paradísi coronávit eum.\end{Resp}
\begin{Resp}\Vsign Índuit eum Dóminus lorícam fídei, et ornávit eum.\end{Resp}
\begin{Resp}\Rsign Et ad portas paradísi coronávit eum.\end{Resp}
\switchcolumn
\EN
\begin{Resp}\Rsign The Lord loved him and beautified him; He clothed him with a robe of glory, \Star{} And crowned him at the gates of Paradise.\end{Resp}
\begin{Resp}\Vsign The Lord hath put on him the breast-plate of faith, and hath adorned him.\end{Resp}
\begin{Resp}\Rsign And crowned him at the gates of Paradise.\end{Resp}
\switchcolumn*

\LA
\LessonHead{Lectio VI}
\switchcolumn
\EN
\LessonHead{Lesson VI}
\switchcolumn*

\LA
\noindent Demum, Complúti finem sibi vitæ adesse intélligens, lácera et obsoleta indútus túnica, conjectis in crucem óculis, singulari devotióne illis verbis ex sacro hymno pronuntiátis: Dulce lignum, dulces clavos, dulcia ferens póndera, quæ fuísti digna portare Regem cælórum et Dóminum, ánimam Deo réddidit pridie Idus Novembris, anno Dóminis supra millesimum quadringentésimo sexagesimo tertio. Ejus corpus, cum menses non paucos (ut pio confluéntium desidério fieret satis) insepúltum mansísset, quasi jam incorruptiónem indúerit, odórem suavíssimum efflávit. Illum, multis et illustribus miraculis clarum, Xystus quintus, Pontifex maximus, Sanctórum número adscripsit.\par
\switchcolumn
\EN
\noindent He was at Alcala when he understood that the end of his life was at hand. Clothed in a ragged cast-away habit, he fixed his eyes upon the Cross, and said with extraordinary earnestness 'Sweet the nails, and sweet the iron, Sweet the Weight That hung on thee,' thou that wast chosen to up-bear the Lord, the King of heaven, and so he gave up his soul to God, upon the 12th day of November, in the year of our Lord 1463. To satisfy the godly wishes of the multitude, his body was kept unburied for not a few months, and lay in a right sweet savour, as though the corruptible had already put on incorruption. He was famous for many and great miracles, and Pope Sixtus V. enrolled him in the number of the Saints.\par
\switchcolumn*

\LA
\begin{Resp}\Rsign Iste homo perfécit ómnia quæ locútus est ei Deus, et dixit ad eum: Ingrédere in réquiem meam: \Star{} Quia te vidi justum coram me ex ómnibus géntibus.\end{Resp}
\begin{Resp}\Vsign Iste est, qui contémpsit vitam mundi, et pervénit ad cæléstia regna.\end{Resp}
\begin{Resp}\Rsign Quia te vidi justum coram me ex ómnibus géntibus.\end{Resp}
\begin{Resp}\Vsign Glória Patri, et Fílio, \Star{} et Spirítui Sancto.\end{Resp}
\begin{Resp}\Rsign Quia te vidi justum coram me ex ómnibus géntibus.\end{Resp}
\switchcolumn
\EN
\begin{Resp}\Rsign This is he which did according unto all that God commanded him; and God said unto him: Enter thou into My rest, \Star{} For thee have I seen righteous before Me among all people.\end{Resp}
\begin{Resp}\Vsign This is he which loved not his life in this world, and is come unto an everlasting kingdom.\end{Resp}
\begin{Resp}\Rsign For thee have I seen righteous before Me among all people.\end{Resp}
\begin{Resp}\Vsign Glory be to the Father, and to the Son, \Star{} and to the Holy Ghost.\end{Resp}
\begin{Resp}\Rsign For thee have I seen righteous before Me among all people.\end{Resp}
\switchcolumn*

\LA
\LessonHead{Lectio VII}
\switchcolumn
\EN
\LessonHead{Lesson VII}
\switchcolumn*

\LA
\LessonTitle{Léctio sancti Evangélii secúndum Lucam}
\switchcolumn
\EN
\LessonTitle{From the Holy Gospel according to Luke}
\switchcolumn*

\LA
\Cite{Luc 12:32-34}
\switchcolumn
\EN
\Cite{Luke 12:32-35}
\switchcolumn*

\LA
\noindent In illo témpore: Dixit Jesus discípulis suis: Nolíte timére pusíllus grex, quia complácuit Patri vestro dare vobis regnum. Et réliqua.\par
\switchcolumn
\EN
\noindent At that time, Jesus said unto His disciples: Fear not, little flock, for it is your Father's good pleasure to give you the kingdom. And so on.\par
\switchcolumn*

\LA
\LessonTitle{Homilía sancti Bedæ Venerábilis Presbýteri}
\switchcolumn
\EN
\LessonTitle{Homily by the Venerable Bede, Priest at Jarrow, and Doctor of the Church.}
\switchcolumn*

\LA
\Source{Lib. 4, Cap. 54, in Luc. 12}
\switchcolumn
\EN
\Source{Bk. iv. Ch. 54 on Luke xii.}
\switchcolumn*

\LA
\noindent Pusíllum gregem electórum, vel ob comparatiónem majóris númeri reprobórum, vel pótius ob humilitátis devotiónem nóminat: quia vidélicet Ecclésiam suam quantálibet numerositáte jam dilatátam, tamen usque ad finem mundi humilitáte vult créscere, et ad promíssum regnum humilitáte perveníre. Ideóque ejus labóres blande consolátus, quam regnum Dei tantum quærére prǽcipit, eídem regnum a Patre dandum complácita benignitáte promíttit.\par
\switchcolumn
\EN
\noindent The elect are called a little flock, perchance because the reprobate are far more in number than they, but, more probably, because they love to be lowly, since it is God's will that however much His Church should grow in numbers, she should grow with lowliness even unto the end of the world, and should enter lowly into that kingdom which is hers by His promise. That kingdom He promiseth to her here, when He biddeth her to seek only the kingdom of God, and, to comfort her in her travail, He doth so sweetly and so graciously say that her Father will give it to her.\par
\switchcolumn*

\LA
\begin{Resp}\Rsign Iste est qui ante Deum magnas virtútes operátus est, et de omni corde suo laudávit Dóminum: \Star{} Ipse intercédat pro peccátis ómnium populórum.\end{Resp}
\begin{Resp}\Vsign Ecce homo sine queréla, verus Dei cultor, ábstinens se ab omni ópere malo, et pérmanens in innocéntia sua.\end{Resp}
\begin{Resp}\Rsign Ipse intercédat pro peccátis ómnium populórum.\end{Resp}
\switchcolumn
\EN
\begin{Resp}\Rsign This is he which wrought great wonders before God, and praised the Lord with all his heart. \Star{} May he pray for all people, that their sins may be forgiven unto them.\end{Resp}
\begin{Resp}\Vsign Behold a man without blame, a worshipper of God in truth, keeping himself clean from every evil work, and abiding still in his innocency.\end{Resp}
\begin{Resp}\Rsign May he pray for all people, that their sins may be forgiven unto them.\end{Resp}
\switchcolumn*

\LA
\LessonHead{Lectio VIII}
\switchcolumn
\EN
\LessonHead{Lesson VIII}
\switchcolumn*

\LA
\noindent Véndite quæ possidétis, et date eleemósynam. Nolíte, inquit, timére, ne propter regnum Dei militántibus, hujus vitæ necessária desint: quin étiam posséssa propter eleemósynam véndite. Quod tunc digne fit, quando quis semel pro Dómino suis ómnibus spretis, nihilóminus post hæc labóre mánuum, unde et victum transígere, et eleemósynam dare queat, operátur. Unde gloriátur Apóstolus, dicens: Argéntum, et aurum, aut vestem nullíus concupívi: ipsi scitis, quóniam ad ea quæ mihi opus erant, et his qui mecum sunt, ministravérunt manus istæ. Ómnia osténdi vobis, quóniam sic laborántes opórtet suscípere infírmos.\par
\switchcolumn
\EN
\noindent Sell what ye have and give alms. Fear not, He saith, lest, while ye fight for the kingdom of God, ye should lack such things as are needful for this life, nay rather, sell even that which ye have, and give alms. This doth, whosoever for the Lord's sake leaveth all that he hath, and then worketh with his hands, that so he may have to eat, and withal to give alms. In this doth the Apostle boast himself, saying I have coveted no man's silver, or gold, or apparel, as ye yourselves know for these hands have ministered unto my necessities, and to them that were with me. I have showed you all things, how that so labouring ye ought to support the weak. (Acts xx. 33, 34, 35.)\par
\switchcolumn*

\LA
\begin{Resp}\Rsign Sint lumbi vestri præcíncti, et lucérnæ ardéntes in mánibus vestris: \Star{} Et vos símiles homínibus exspectántibus dóminum suum, quando revertátur a núptiis.\end{Resp}
\begin{Resp}\Vsign Vigiláte ergo, quia nescítis qua hora Dóminus vester ventúrus sit.\end{Resp}
\begin{Resp}\Rsign Et vos símiles homínibus exspectántibus dóminum suum, quando revertátur a núptiis.\end{Resp}
\begin{Resp}\Vsign Glória Patri, et Fílio, \Star{} et Spirítui Sancto.\end{Resp}
\begin{Resp}\Rsign Et vos símiles homínibus exspectántibus dóminum suum, quando revertátur a núptiis.\end{Resp}
\switchcolumn
\EN
\begin{Resp}\Rsign Let your loins be girded about, and your lights burning; \Star{} And ye yourselves like unto men that wait for their lord, when he will return from the wedding.\end{Resp}
\begin{Resp}\Vsign Watch therefore, for ye know not what hour your Lord doth come.\end{Resp}
\begin{Resp}\Rsign And ye yourselves like unto men that wait for their lord, when he will return from the wedding.\end{Resp}
\begin{Resp}\Vsign Glory be to the Father, and to the Son, \Star{} and to the Holy Ghost.\end{Resp}
\begin{Resp}\Rsign And ye yourselves like unto men that wait for their lord, when he will return from the wedding.\end{Resp}
\switchcolumn*

\LA
\LessonHead{Lectio IX}
\switchcolumn
\EN
\LessonHead{Lesson IX}
\switchcolumn*

\LA
\noindent Fácite vobis sácculos qui non veteráscunt: eleemósynas vidélicet operándo quarum merces in ætérnum máneat. Ubi non hoc præcéptum esse putándum est, ut nil pecúniæ reservétur a sanctis, vel suis scílicet, vel páuperum úsibus suggeréndæ: cum et ipse Dóminus, cui ministrábant Ángeli, tamen ad informándam Ecclésiam suam lóculos habuísse legátur, et a fidélibus obláta consérvans, et suórum necessitátibus aliísque indigéntibus tríbuens; sed ne Deo propter ista serviátur, et ob inópiæ timórem justítia deserátur.\par
\switchcolumn
\EN
\noindent Provide yourselves bags which wax not old that is to say, by almsgiving, the reward thereof remaineth for ever. Nevertheless, we must not think here that this commandment forbiddeth the Saints to keep money for their own use, and for helping of the poor. The Lord Himself, to Whom Angels ministered, had a bag, and kept therein that which the faithful people gave unto Him (John xii. 6,) to relieve therewith the need of His disciples, and other poor folk. But we are commanded not to serve God for gain, nor to work unrighteousness for fear of poverty.\par
\switchcolumn*

\LA
\TeDeum{Te Deum}
\switchcolumn
\EN
\TeDeum{Te Deum}
\switchcolumn*

\end{paracol}

\par\needspace{10\baselineskip}\addvspace{1.2em}{\centering\small\textsc{Die 14 Novembris} · \textsc{14 November}\par}
\Office{S. Josaphat Episcopi et Martyris}{St. Josaphat, Bishop and Martyr}{Duplex}
\addcontentsline{toc}{subsection}{Die 14 Novembris · S. Josaphat Episcopi et Martyris \textit{— St. Josaphat, Bishop and Martyr}}
\begin{paracol}{2}
\LA
\RefLine{Lectiones I–III: de Scriptura occurrente.}
\switchcolumn
\EN
\RefLine{Lessons I–III: from the Scripture of the day.}
\switchcolumn*

\LA
\LessonHead{Lectio IV}
\switchcolumn
\EN
\LessonHead{Lesson IV}
\switchcolumn*

\LA
\noindent Jósaphat Kuncewítius, nobílibus et catholicis paréntibus Vladimíriæ in Volhinia natus, cum puerulus matrem de Christi passióne loquentem audíret, jaculo e látere imaginis Jesu crucifixi immisso, vulnus in corde suscépit. Dei amóre incénsus, adeo oratióni aliisque piis opéribus instare cœpit, ut provectióribus adolescéntibus exemplo et admiratióni esset. Vicennis inter claustrales sancti Basilíi alumnos monasticam regulam proféssus, mirum quos in evangelica perfectióne progréssus fécerit. Nudis pédibus, frigidíssima licet sæviénte regiónis híeme, incedébat; carnes numquam, vinum nonnisi ex obediéntia adhibuit, asperrimoque cilício ad obitum usque corpus afflixit. Castitátis florem, quem ab adolescéntia Virgini Deiparæ vóverat, inviolátum servávit. Virtútis doctrinæque ejus brevi sic fama percrébuit, ut, quamvis júnior, Bytenii monasterio præfectus sit mox Vilnénsis archimandríta, ac demum archiepíscopus Polocénsis, invitus quidem, sed Catholicis gestiéntibus, fúerit renuntiatus.\par
\switchcolumn
\EN
\noindent Josaphat Kuncewitz was born of noble Catholic parents at Vladimir in Volhynia. Once as a child, as he listed to his mother tell the story of the Passion, a dart came forth from the side of Christ on the crucifix and wounded the boy in the heart. Set on fire with love of God, he devoted himself to prayer and works of charity with such zeal that he became the admiration and the model for youths far older than he. When Josaphat was twenty years old he was professed among the cloistered followers of the monastic rule of Saint Basil. Almost at once he made remarkable progress in evangelical perfection. He went barefoot, even in the severe winters of that country. He never ate meat, and drank wine only when obliged to do so under obedience. He disciplined his body by wearing rough hair-shirts until the day of his death. He kept unspotted the flower of chastity which in his youth he had dedicated to the Virgin Mother of God. He became so celebrated for virtue and learning that despite his youth he was made superior of the monastery at Byten, and the Archimandrite of Vilnius. Finally much against his will, but to the very great joy of the Catholic people, he was made Archbishop of Polotsk.\par
\switchcolumn*

\LA
\begin{Resp}\Rsign Honéstum fecit illum Dóminus, et custodívit eum ab inimícis, et a seductóribus tutávit illum: \Star{} Et dedit illi claritátem ætérnam.\end{Resp}
\begin{Resp}\Vsign Descendítque cum illo in fóveam, et in vínculis non derelíquit eum.\end{Resp}
\begin{Resp}\Rsign Et dedit illi claritátem ætérnam.\end{Resp}
\switchcolumn
\EN
\begin{Resp}\Rsign The Lord made him honourable, and defended him from his enemies, and kept him safe from those that lay in wait for him, \Star{} And gave him perpetual glory.\end{Resp}
\begin{Resp}\Vsign He went down with him into the pit, and left him not in bonds.\end{Resp}
\begin{Resp}\Rsign And gave him perpetual glory.\end{Resp}
\switchcolumn*

\LA
\LessonHead{Lectio V}
\switchcolumn
\EN
\LessonHead{Lesson V}
\switchcolumn*

\LA
\noindent Hac dignitate auctus, nihil de priori vivéndi ratióne remíttens, nonnisi divinum cultum et creditárum sibi óvium salútem cordi hábuit. Catholicæ unitátis ac veritátis strenuus propugnator, totis viribus adlaborávit, ut schismaticos hæreticosque ad communiónem cum beáti Petri Sede reduceret. Summum Pontificem ejusque potestátis plenitúdinem ab impudentíssimis impiórum calumniis et erróribus, qua conciónibus, qua scriptis pietáte ac doctrina refertis, defendere numquam déstitit. Episcopalem jurisdictiónem et Ecclésiæ bona a láicis usurpáta vindicávit. Incredibile dictu est quot hæreticos in sinum matris Ecclésiæ revocáverit. Uniónis vero Græcæ Ecclésiæ cum Latina Josephátum promotórem exstitisse præclaríssimum, étiam pontifícia oracula diserte testántur. Ad hæc, et templi Dei decóri instaurando, et sacrárum vírginum exstruéndis ædibus, aliisque piis opéribus juvandis, mensæ suæ proventus ultro erogávit. In páuperes adeo effúsus, ut, cum olim inópiæ cujusdam víduæ sublevandæ nihil occurreret, episcopale pallium, seu omophórion, oppignorari jusserit.\par
\switchcolumn
\EN
\noindent In the years following the promotion to this dignity, Josaphat did not relax in any way his austere mode of living. Nothing was so close to his heart as service to God and the salvation of the flock entrusted to his care. He was a vigorous champion of Catholic unity and truth. He laboured to the utmost of his ability to win back schismatics and heretics to unity with the See of blessed Peter. Both by preaching and writing he defended the Supreme Pontiff and the doctrine of the Pope's plenitude of power. He directed these works, full of piety and learning against most shameful calumnies and the errors of wicked men. Josaphat vindicated episcopal rights and restored ecclesiastical property seized by laymen. He won back an incredible number of heretics to the bosom of holy Mother Church. How successfully he laboured to re-establish communion between the Greek and Latin Churches is told in Papal commendations. He gladly spent the revenues set aside for his maintenance to rebuild God's house, to erect convents for consecrated virgins, and to carry on other charitable works. So generous was Josaphat towards the poor that in one instance when he did not have money enough to supply the needs of a certain widow, he pawned his omophorion, that is, his episcopal pallium.\par
\switchcolumn*

\LA
\begin{Resp}\Rsign Desidérium ánimæ ejus tribuísti ei Dómine, \Star{} Et voluntáte labiórum ejus non fraudásti eum.\end{Resp}
\begin{Resp}\Vsign Quóniam prævenísti eum in benedictiónibus dulcédinis: posuísti in cápite ejus corónam de lápide pretióso.\end{Resp}
\begin{Resp}\Rsign Et voluntáte labiórum ejus non fraudásti eum.\end{Resp}
\switchcolumn
\EN
\begin{Resp}\Rsign Lord, Thou hast given him his heart's desire. \Star{} And hast not withholden the request of his lips.\end{Resp}
\begin{Resp}\Vsign For Thou hast prevented him with the blessings of sweetness: Thou hast set a crown of precious stones upon his head.\end{Resp}
\begin{Resp}\Rsign And hast not withholden the request of his lips.\end{Resp}
\switchcolumn*

\LA
\LessonHead{Lectio VI}
\switchcolumn
\EN
\LessonHead{Lesson VI}
\switchcolumn*

\LA
\noindent Tot catholicæ fidei increménta perditissimórum hóminum adeo excitavérunt odia, ut, conspiratióne ínita, Christi athletam ad necem quærerent; quam sibi imminére ipse in suo ad pópulum sermóne prænuntiávit. Cum ítaque Vitépscum pastoralis visitatiónis grátia profectus esset, illi archiepiscopales invadunt ædes, obvios quosque feriunt ac cædunt. Tum vir mitíssimus quæréntibus sponte occurrit, eosque amice compellans, Filíoli, inquit, quare familiares meos cæditis? Si quid contra me habetis, ecce adsum. Hinc, ímpetu facto, eum verbéribus contundunt, telis confodiunt, ac demum, immani secúri necátum, in flumen projíciunt, die duodecima Novembris anni sexcentésimi vicesimi tertii supra millesimum, ætátis ejus quadragesimi tertii. Corpus, mirábili luce circumfusum, ex imo flúminis alveo elátum est. Sanguis Mártyris parricidis ipsis in primis prófuit, qui fere omnes, cápitis damnati, ejurato schísmate, suum scelus detestáti sunt. Cum tantus Præsul plurimis post obitum coruscaret miraculis, eum Urbanus octavus, Pontifex maximus, Beatórum honóribus decorávit. Pius nonus, tertio Kalendas Julias anni millesimi octingentesimi sexagesimi septimi, cum sæcularia Apostolórum Principum solemnia celebraréntur, coram patrum cardinalium senatu, simulque astántibus fere quingentis, patriarchis, metropolítis et epíscopis cujuscúmque ritus, qui ex toto terrárum orbe convenerant; hunc ecclesiásticæ unitátis assertórem, primum ex orientálibus, solemni ritu in Vaticana basilica Sanctórum ordini accensuit. Cujus Offícium ac Missam Leo décimus tertius, summus Pontifex, ad universam exténdit Ecclésiam.\par
\switchcolumn
\EN
\noindent The great progress made by the Catholic faith so stirred up the anger of certain of its wicked enemies that they conspired to murder this athlete of Christ. In a sermon he foretold to his people what was about to happen. As he was setting out for Vitebsk on a pastoral visit, these enemies broke into the episcopal palace, attacking and wounding every one they found. Undaunted, this most kindly man hurried out to the assassins of his own free will and addressed them mildly. My little children, he said, why do ye strike my servants? If ye have any complaint against me, I am here. Thereupon they rushed at him, overwhelmed him with blows and pierced him through with spears. Finally they slew him a stroke of a great axe and threw his body into the river. This happened on November 12th, 1632, when he was forty-three years old. Later his body, surrounded by a marvellous light, was raised from the deepest part of the river. The blood of this Martyr benefited first of all those murderers of their spiritual father. Sentenced to die for their crime, almost all abjured their schism and repented of their crime. Because this wonderful high priest became famous after his death for many miracles, the Supreme Pontiff, Urban VIII, honoured him with the title Blessed. On the 29th of June, 1867, during the solemn observance of the centenaries of the Princes of the Apostles, in the presence of the college of cardinals, of about five hundred others, patriarchs, metropolitans, and bishops of every rite from all parts of the world, assembled in the Vatican basilica, with all solemn ceremonies, Pius IX canonized the first eastern Christian to uphold the unity of the Church. The Supreme Pontiff, Leo XIII, extended his Mass and Office to the universal Church.\par
\switchcolumn*

\LA
\begin{Resp}\Rsign Stola jucunditátis índuit eum Dóminus: \Star{} Et corónam pulchritúdinis pósuit super caput ejus.\end{Resp}
\begin{Resp}\Vsign Cibávit illum Dóminus pane vitæ et intelléctus: et aqua sapiéntiæ salutáris potávit illum.\end{Resp}
\begin{Resp}\Rsign Et corónam pulchritúdinis pósuit super caput ejus.\end{Resp}
\begin{Resp}\Vsign Glória Patri, et Fílio, \Star{} et Spirítui Sancto.\end{Resp}
\begin{Resp}\Rsign Et corónam pulchritúdinis pósuit super caput ejus.\end{Resp}
\switchcolumn
\EN
\begin{Resp}\Rsign The Lord hath put on him a robe of honour, \Star{} And put about his head a crown of joy.\end{Resp}
\begin{Resp}\Vsign With the bread of life and understanding hath the Lord fed him, and given him the water of health and wisdom to drink.\end{Resp}
\begin{Resp}\Rsign And put about his head a crown of joy.\end{Resp}
\begin{Resp}\Vsign Glory be to the Father, and to the Son, \Star{} and to the Holy Ghost.\end{Resp}
\begin{Resp}\Rsign And put about his head a crown of joy.\end{Resp}
\switchcolumn*

\LA
\LessonHead{Lectio VII}
\switchcolumn
\EN
\LessonHead{Lesson VII}
\switchcolumn*

\LA
\LessonTitle{Léctio sancti Evangélii secúndum Joánnem}
\switchcolumn
\EN
\LessonTitle{From the Holy Gospel according to John}
\switchcolumn*

\LA
\Cite{Joannes 10:11-16}
\switchcolumn
\EN
\Cite{John 10:11-16}
\switchcolumn*

\LA
\noindent In illo témpore: Dixit Jesus pharisǽis: Ego sum pastor bonus. Bonus pastor animam suam dat pro ovibus suis. Et réliqua.\par
\switchcolumn
\EN
\noindent In that time Jesus said to the pharisees: I am the good shepherd. The good shepherd giveth his life for his sheep. And so on.\par
\switchcolumn*

\LA
\LessonTitle{Homilia sancti Joannis Chrysostomi}
\switchcolumn
\EN
\LessonTitle{Homily by St. John Chrysostom, Patriarch (of Constantinople.)}
\switchcolumn*

\LA
\Source{Hom. 59 in Joánnem}
\switchcolumn
\EN
\Source{59th on John.}
\switchcolumn*

\LA
\noindent Magnum quiddam, dilectissimi, magnum, inquam, est Ecclesiæ prælatio, et quæ multa indiget sapientia et fortitudine, qualem Christus proposuit: ut animam pro ovibus ponamus, et numquam illas deseramus: ut lupo generose resistamus. Hæc enim inter pastorem et mercenarium est differentia. Alter propriæ, contemptis ovibus: alter, sua contempta, ovium semper saluti invigilat. Pastoris ergo exemplo demonstrato, deceptores duos meminit; furem mactantem et rapientem oves; et mercenarium permittentem, neque defendentem comissas.\par
\switchcolumn
\EN
\noindent Dearly beloved brethren, the Bishops of the Church hold a great office, an office that needeth much that wisdom and strength whereof Christ hath given us an example. We must learn of Him to lay down our lives for the sheep and never to leave them; and to fight bravely against the wolf. This is the difference between the true shepherd and the hireling. The one leaveth the sheep and seeketh his own safety, but the other recketh not of his own safety, so as he may watch over the sheep. Christ then having given us the pattern of a good shepherd, warneth us against two enemies; first, the thief that cometh not but to kill and to steal, and, secondly, the hireling that standeth by, and defendeth not them that are committed to his charge.\par
\switchcolumn*

\LA
\begin{Resp}\Rsign Coróna áurea super caput ejus, \Star{} Expréssa signo sanctitátis, glória honóris, et opus fortitúdinis.\end{Resp}
\begin{Resp}\Vsign Quóniam prævenísti eum in benedictiónibus dulcédinis, posuísti in cápite ejus corónam de lápide pretióso.\end{Resp}
\begin{Resp}\Rsign Expréssa signo sanctitátis, glória honóris, et opus fortitúdinis.\end{Resp}
\switchcolumn
\EN
\begin{Resp}\Rsign A crown of gold upon his head, \Star{} Wherein is engraved Holiness, an ornament of honour, a costly work.\end{Resp}
\begin{Resp}\Vsign For Thou hast prevented him with the blessings of sweetness, Thou hast set a crown of precious stones upon his head.\end{Resp}
\begin{Resp}\Rsign Wherein is engraved Holiness, an ornament of honour, a costly work.\end{Resp}
\switchcolumn*

\LA
\LessonHead{Lectio VIII}
\switchcolumn
\EN
\LessonHead{Lesson VIII}
\switchcolumn*

\LA
\noindent Quod superiori tempore Ezechiel his verbis insectatus est: Væ pastoribus Israël: nonne pascebant semetipsos? nonne greges pascuntur a pastoribus? Sed illi contrarium faciebant, quod maximæ malitiæ genus est, et plurimorum causa malorum. Idcirco inquit: Neque quod abjectum erat, reducebant: neque quod perierat quærebant: neque confractum alligabant: neque infirmum consolidabant: quoniam se, non gregem pascebant. Idem et Paulus aliis verbis significat: Omnes quæ sua sunt, quærunt, non quæ Jesu Christi.\par
\switchcolumn
\EN
\noindent Ezechiel hath said of old time, (xxxiv. 2): Woe be to the shepherds of Israel! do they not feed themselves? Should not the shepherds feed the flocks? But they did the contrary, a great wickedness and the root of many evils. Therefore, he saith, they brought not back that which was gone astray neither did they search for that which was lost neither did they bind up that which was broken, nor strengthen that which was sick; for they fed themselves, and not the flock. And Paul hath the same in other words, where he saith, (Phil. ii. 21): All seek their own, not the things which are Jesus Christ's.\par
\switchcolumn*

\LA
\begin{Resp}\Rsign Hic est vere Martyr, qui pro Christi nómine sánguinem suum fudit: \Star{} Qui minas júdicum non tímuit, nec terrénæ dignitátis glóriam quæsívit, sed ad cæléstia regna pervénit.\end{Resp}
\begin{Resp}\Vsign Justum dedúxit Dóminus per vias rectas, et osténdit illi regnum Dei.\end{Resp}
\begin{Resp}\Rsign Qui minas júdicum non tímuit, nec terrénæ dignitátis glóriam quæsívit, sed ad cæléstia regna pervénit.\end{Resp}
\begin{Resp}\Vsign Glória Patri, et Fílio, \Star{} et Spirítui Sancto.\end{Resp}
\begin{Resp}\Rsign Qui minas júdicum non tímuit, nec terrénæ dignitátis glóriam quæsívit, sed ad cæléstia regna pervénit.\end{Resp}
\switchcolumn
\EN
\begin{Resp}\Rsign This is a Martyr indeed, who shed his blood for Christ's Name's sake \Star{} Who feared not for the threats of judges, nor sought to be great with the glory of this world, but pressed on unto the kingdom of heaven.\end{Resp}
\begin{Resp}\Vsign The Lord guided the righteous in right paths, and showed him the kingdom of God.\end{Resp}
\begin{Resp}\Rsign Who feared not for the threats of judges, nor sought to be great with the glory of this world, but pressed on unto the kingdom of heaven.\end{Resp}
\begin{Resp}\Vsign Glory be to the Father, and to the Son, \Star{} and to the Holy Ghost.\end{Resp}
\begin{Resp}\Rsign Who feared not for the threats of judges, nor sought to be great with the glory of this world, but pressed on unto the kingdom of heaven.\end{Resp}
\switchcolumn*

\LA
\LessonHead{Lectio IX}
\switchcolumn
\EN
\LessonHead{Lesson IX}
\switchcolumn*

\LA
\noindent Verum ab utroque se dissidere ostendit Christus: ab illis quidem, qui in aliorum perniciem veniunt, cum dicat, se propterea venisse, ut vitam haberent et abundantius haberent: ab his autem, qui oves a lupis rapi negligebant, dicendo se propter eas animam ponere, ne oves perirent. Nam cum Judæi ipsum interimere vellent non propterea destitit a doctrina, neque credentes tradidit, sed perstitit, et pertulit mortem: ideo frequenter inquit: Ego sum pastor bonus. Quæ cum nullo niti testimonio viderentur (quod enim poneret animam suam, non multo post re demonstratum est; quod autem vitam haberent et abundantius haberent, eventurum erat in futuro sæculo) alterum ab altero confirmat.\par
\switchcolumn
\EN
\noindent Christ showeth Himself very different from either the thief or the hireling; whereas the thief cometh to destroy, He came that they might have life, and that they might have it more abundantly. The hireling fleeth, but He layeth down His life for the sheep, that the sheep perish not. When then the Jews went about to kill Him, He ceased not to teach He gave not up them that believed in Him, but stood steadfast and died. Wherefore He hath good title often to say, I am the Good Shepherd. It was but a little while, and He showed us how He could lay down His life for the sheep. And if it appeareth not as yet how they have life, and have it more abundantly, (but it shall appear, in the world which is to come,) we may well be persuaded of the truth of the second promise, who have seen the fulfilment of the first.\par
\switchcolumn*

\LA
\TeDeum{Te Deum}
\switchcolumn
\EN
\TeDeum{Te Deum}
\switchcolumn*

\end{paracol}

\par\needspace{10\baselineskip}\addvspace{1.2em}{\centering\small\textsc{Die 15 Novembris} · \textsc{15 November}\par}
\Office{S. Alberti Magni Episcopi Confessoris et Ecclesiæ Doctoris}{St. Albert the Great, Bishop, Confessor and Doctor of the Church}{Duplex}
\addcontentsline{toc}{subsection}{Die 15 Novembris · S. Alberti Magni Episcopi Confessoris et Ecclesiæ Doctoris \textit{— St. Albert the Great, Bishop, Confessor and Doctor of the Church}}
\begin{paracol}{2}
\LA
\RefLine{Lectiones I–III: de Scriptura occurrente.}
\switchcolumn
\EN
\RefLine{Lessons I–III: from the Scripture of the day.}
\switchcolumn*

\LA
\RefLine{Lectiones VII–IX: ex Commúni Doctoris Pontificis (C4a).}
\switchcolumn
\EN
\RefLine{Lessons VII–IX: from the Common of a Doctor Bishop (C4a).}
\switchcolumn*

\LA
\LessonHead{Lectio IV}
\switchcolumn
\EN
\LessonHead{Lesson IV}
\switchcolumn*

\LA
\noindent Albertus, ob singularem doctrinam cognomento Magnus, Lauíngiæ ad Danubium, in Suevia, natus, a púero diligenter institútus est. Studiórum causa e pátria discédens, dum Patavii morarétur, hortante beato Jordano, generali Magistro Ordinis Prædicatórum, Dominicianæ famíliæ, frustra obsisténte avunculo, adscribi postulávit. Inter fratres adlectus, in ómnibus Deo deditus, religiosa observántia et pietáte enituit, filiali ac tenerrima in beátam Mariam Vírginem devotióne flagrans. Totam vitæ formam, oratióne studium præveniendo, ita dispósuit, ut apostolicam religiónem proféssus, ad prædicatiónem verbi Dei et animárum salútem procurandam idoneus evaderet administer. Mox ad stúdia explenda Colóniam Agrippinam missus, ita profecit, ut omnes fere sæculares sciéntias diligentíssime præ ómnibus suis coævis investigaverit atque auxerit ; et de divinæ legis fonte, testante Alexandro quarto, adeo salutifera fluéntia doctrinæ potávit, ut ejusdem in suo péctore vigéret plenitudo.\par
\switchcolumn
\EN
\noindent Albert, called the Great, because of his extraordinary learning, was born in Swabia, at Lauingen on the Danube, and very carefully educated from boyhood. To further his higher studies he left his native country and went to Padua. At the urging of blessed Jordan, Master General of the Order of Preachers, he asked admission into the family of the Dominicans, in spite of the futile protests of his uncle. After his election to membership among the brethren, Albert was dedicated in all things to God, and was conspicuous for his piety, his strict observance of the rule, and especially for his tender and filial devotion to the Blessed Virgin Mary. Always before study he spent some time in prayer. After his profession of apostolic religion, he so regulated his schedule of life that he became an accomplished preacher of the word of God and an efficient instrument for the salvation of souls. Soon the Order sent Albert to complete his studies at Cologne, where he made such progress in every branch of secular science that he surpassed all his contemporaries in scholarship and achievement. In the meantime, as Alexander IV testifieth, he drank so deeply of the health-giving waters of doctrine, sprung forth from the fountain of the divine law, that his soul was flooded with their abundance.\par
\switchcolumn*

\LA
\begin{Resp}\Rsign Invéni David servum meum, óleo sancto meo unxi eum: \Star{} Manus enim mea auxiliábitur ei.\end{Resp}
\begin{Resp}\Vsign Nihil profíciet inimícus in eo, et fílius iniquitátis non nocébit ei.\end{Resp}
\begin{Resp}\Rsign Manus enim mea auxiliábitur ei.\end{Resp}
\switchcolumn
\EN
\begin{Resp}\Rsign I have found David My servant, with My holy oil have I anointed him \Star{} For My hand shall help him.\end{Resp}
\begin{Resp}\Vsign The enemy shall prevail nothing against him, nor the son of wickedness afflict him.\end{Resp}
\begin{Resp}\Rsign For My hand shall help him.\end{Resp}
\switchcolumn*

\LA
\LessonHead{Lectio V}
\switchcolumn
\EN
\LessonHead{Lesson V}
\switchcolumn*

\LA
\noindent Ut scientiárum thesáuris alios ditaret, lector Hildeshémii, deinceps Friburgi, Ratisbonæ et Argentinæ constitútus est. Omnium in se admiratiónem convértens, cum in Parisiénsi inclyta universitáte sacræ facultáti docéndo decus adderet, magister theologíæ renunciátus est. Doctrinis gentílium philosophórum ad rectæ ratiónis dictámina revocatis, eárum cum fide cohæréntiam clarius commonstrávit. De divinórum intelléctu mira expósuit. Quantum vero univérsas disciplínas, præsertim sacras, férvido ingenio atque indefesso studio provexerit, plurima ejus scripta in omni fere scientiárum genere aperte declarant. Ut studio generali sui Ordinis præesset, Colóniam reversus est, eo successu ut ejus in scholis auctoritas et doctrinæ fama magis magisque vigúerit. Thomam de Aquino discipulum diléctum hábuit, cujus altitúdinem mentis ipse primus perspéxit ac prædicávit. Erga Sacrosanctum altaris Sacraméntum piíssimo ferebátur afféctu, deque eo præclára conscripsit ; rei quoque mysticæ ánimis instituéndis vias ampliores parávit, adeo ut frúgifer tanti magistri zelus quam late in Ecclésia patúerit.\par
\switchcolumn
\EN
\noindent That others might share the rich treasure of his learning, Albert was appointed professor at Hildesheim, then Freiburg, Ratisbon and Strasbourg successively. He became the marvel of all. During the period when he taught sacred theology in the famous University of Paris, he became world-famous, and received the degree of Master of Theology. Examining the teachings of pagan philosophers in the light of sound reason, he demonstrated clearly that they were in fundamental accord with the tenets of the faith. He expounded most brilliantly the thesis on the extent of the power of human understanding to comprehend divine mysteries. How great was his genius, how brilliant his intellect, how zealously he applied himself to study until he had become learned in every branch of scholarship, especially sacred theology, is best shown by his numerous writings. These encompass every known subject. Albert returned to Cologne to become president of the university conducted by his Order. He was so successful that he became ever more widely acknowledged as an authority by the schools; his reputation for learning increased. Among his pupils was his beloved Thomas Aquinas. Albert was first to recognize and acclaim the greatness of that intellect. He had a deep devotion towards the Blessed Sacrament of the Altar and composed some magnificent works upon it. He also pointed out wider fields for the study of the mystical things of the soul. He succeeded so well that the zeal of this great master spread far and wide in the Church.\par
\switchcolumn*

\LA
\begin{Resp}\Rsign Pósui adjutórium super poténtem, et exaltávi eléctum de plebe mea: \Star{} Manus enim mea auxiliábitur ei.\end{Resp}
\begin{Resp}\Vsign Invéni David servum meum, óleo sancto meo unxi eum.\end{Resp}
\begin{Resp}\Rsign Manus enim mea auxiliábitur ei.\end{Resp}
\switchcolumn
\EN
\begin{Resp}\Rsign I have laid help upon one that is mighty, and have exalted one chosen out of My people \Star{} For My hand shall help him.\end{Resp}
\begin{Resp}\Vsign I have found David My servant, with My holy oil have I anointed him.\end{Resp}
\begin{Resp}\Rsign For My hand shall help him.\end{Resp}
\switchcolumn*

\LA
\LessonHead{Lectio VI}
\switchcolumn
\EN
\LessonHead{Lesson VI}
\switchcolumn*

\LA
\noindent Tot inter gravíssima munia, religiosæ vitæ exemplis præfulgens, a frátribus Prior Teutoníæ provinciæ eléctus est. Anágniam vocatus, Gulielmum, Ordines mendicántes ímpio ausu impetentem, coram summo Pontifice Alexandro quarto rétudit, qui Episcopum Ratisbonensem eum póstea constítuit. Curæ sui gregis Albertus se totum impendit, morum humilitate ac paupertátis amóre studiosíssime reténtis. Dimisso officio, ad episcopalis tamen ordinis labóres promptus atque álacer per Germaniam et finítimas regiónes spiritualia ministrávit. Consília requiréntibus quam recta ac salutifera sollicite præbebat, et in sedandis discordiis tam prudentem se osténdit, ut eum non solum Colónia pacis conciliatórem noverit, verum étiam ad díssitas regiónes Præláti ac viri príncipes árbitrum componéndis dissidiis eum sæpe advocáverint. A sancto Ludovíco, Francórum rege, relíquiis Christi Passiónis, quam devotíssime Albertus colebat, donatus est. In altero Concílio Lugdunénsi negotia gravióra peregit. Tandem, senio consumptus, docére déstitit. Contemplatióni exínde intentus, in gáudium Dómini sui intrávit anno millesimo ducentésimo octogesimo. Sacros honores in diœcesibus plúribus atque in Ordine Prædicatórum ei, Romanórum Pontificum auctoritate, jam antea tributos, Pius Papa undecimus cumulans, sancti Alberti Magni festum, addito Doctoris titulo, Sacrórum Rituum Congregatiónis votum libentíssime excípiens, ad Ecclésiam universam exténdit ; et eum Pius duodecimus cultórum scientiárum naturalium cælestem apud Deum Patronum constítuit.\par
\switchcolumn
\EN
\noindent Amid so many very important duties Albert shone as an exemplar of the religious life. His brethren, therefore, selected him to be prior of the German province. He was summoned to Anagni, in the presence of the Supreme Pontiff, Alexander IV to refute that William who had been impiously and arrogantly attacking the mendicant orders. Soon after this the Pope appointed Albert Bishop of Ratisbon. As Bishop, Albert devoted himself almost entirely to the care of his flock. Yet he retained meticulously his humility and love of poverty. Up to the time he resigned his see, Albert was prompt and energetic in fulfilling the duties of his episcopal office. He ministered to the spiritual needs of souls throughout Germany and the neighbouring provinces. He was careful that the advice he gave to those who sought his counsels was wise and salutary. So prudent was he in settling disputes that at Cologne he was called the peacemaker. From far distant places, prelates and princes invited him to act as an arbiter to resolve differences. Saint Louis, King of France, presented him with some relics of the sacred Passion of Christ, and Albert cherished them devoutly all his days. In the second Council of Lyons he was instruméntal in bringing to a successful conclusion several weighty matters. He taught until he was worn out with age. His last days were spent in holy contemplation. In the year 1280 he entered into the joy of his Lord. By the authority of the Roman Pontiffs, the honours of the altar had long since been conferred upon Albert in many dioceses and in the Order of Preachers, when Pius XI, gladly accepting the recommendations of the Congregation of Sacred Rites, extended his Feast to the universal Church, and conferred upon the title of Doctor. Pope Pius XII appointed him the heavenly patron with God of all those who study the natural sciences.\par
\switchcolumn*

\LA
\begin{Resp}\Rsign Iste est, qui ante Deum magnas virtútes operátus est, et omnis terra doctrína ejus repléta est: \Star{} Ipse intercédat pro peccátis ómnium populórum.\end{Resp}
\begin{Resp}\Vsign Iste est, qui contémpsit vitam mundi, et pervénit ad cæléstia regna.\end{Resp}
\begin{Resp}\Rsign Ipse intercédat pro peccátis ómnium populórum.\end{Resp}
\begin{Resp}\Vsign Glória Patri, et Fílio, \Star{} et Spirítui Sancto.\end{Resp}
\begin{Resp}\Rsign Ipse intercédat pro peccátis ómnium populórum.\end{Resp}
\switchcolumn
\EN
\begin{Resp}\Rsign This is he which wrought great wonders before God, and the whole earth is full of his teaching \Star{} May he pray for all people, that their sins may be forgiven unto them\end{Resp}
\begin{Resp}\Vsign This is he which loved not his life in this world, and hath attained unto the kingdom of heaven.\end{Resp}
\begin{Resp}\Rsign May he pray for all people, that their sins may be forgiven unto them\end{Resp}
\begin{Resp}\Vsign Glory be to the Father, and to the Son, \Star{} and to the Holy Ghost.\end{Resp}
\begin{Resp}\Rsign May he pray for all people, that their sins may be forgiven unto them\end{Resp}
\switchcolumn*

\end{paracol}

\par\needspace{10\baselineskip}\addvspace{1.2em}{\centering\small\textsc{Die 16 Novembris} · \textsc{16 November}\par}
\Office{S. Gertrudis Virginis}{St. Gertrude, Virgin}{Duplex}
\addcontentsline{toc}{subsection}{Die 16 Novembris · S. Gertrudis Virginis \textit{— St. Gertrude, Virgin}}
\begin{paracol}{2}
\LA
\RefLine{Lectiones I–III: de Scriptura occurrente.}
\switchcolumn
\EN
\RefLine{Lessons I–III: from the Scripture of the day.}
\switchcolumn*

\LA
\RefLine{Lectiones VII–IX: ex Commúni Virginum tantum (C6a).}
\switchcolumn
\EN
\RefLine{Lessons VII–IX: from the Common of a Virgin not a Martyr (C6a).}
\switchcolumn*

\LA
\LessonHead{Lectio IV}
\switchcolumn
\EN
\LessonHead{Lesson IV}
\switchcolumn*

\LA
\noindent Gertrúdis, Islébii in Saxónia nobili genere nata, quinquénnis in monastério Rodardénsi, ordinis sancti Benedicti, virginitátem suam ac seipsam Jesu Christo óbtulit. Quo ex témpore a mundanis rebus prorsus aliena, virtútique sedulo intenta, cæléstis vitæ genus instituit. Ad humaniórum litterárum notítiam, rerum divinárum cognitiónem adjunxit; quarum meditatióne vehementius ad virtútem incensa, brevi christianam perfectiónem adepta est. De Christo ejusque vitæ mystériis sæpenúmero pio cum animi sensu loquebátur, unamque Dei glóriam cogitans, ad illam vota sua ómnia et actiónes referebat. Quamvis autem multis exímiis natúræ et gratiæ donis a Deo aucta esset, ita tamen sibi ipsa vilescebat, ut, inter præcipua divinæ bonitátis miracula, hoc item memoráret, quod se indigníssimam misericorditer sustinéret.\par
\switchcolumn
\EN
\noindent Gertrude was born of a noble family at Eisleben, in Saxony, (about the year of our Lord 1264.) At five years of age she offered her virginity and herself to Jesus Christ, in the Benedictine nunnery at Rodalsdorf. From that time forth she was utterly estranged from earthly things, ever striving for things higher, and began to lead a kind of heavenly life. To learning in human letters she added knowledge of the things of God. In the thought thereof she earnestly desired, and soon reached, the perfection of a Christian soul. Of Christ, and of the things in His life, she spoke oftentimes with movings of spirit. The glory of God was the one end of all her thoughts, and to that her every longing and her every act were given. Though God had crowned her with so many and so noble gifts both of nature and of grace, her belief regarding herself was so humble that she was used to number as among the greatest of the wonders of His goodness that He had always in His mercy borne with one who was so utterly unworthy.\par
\switchcolumn*

\LA
\begin{Resp}\Rsign Propter veritátem, et mansuetúdinem, et justítiam: \Star{} Et dedúcet te mirabíliter déxtera tua.\end{Resp}
\begin{Resp}\Vsign Spécie tua et pulchritúdine tua inténde, próspere procéde, et regna.\end{Resp}
\begin{Resp}\Rsign Et dedúcet te mirabíliter déxtera tua.\end{Resp}
\switchcolumn
\EN
\begin{Resp}\Rsign Because of truth, and meekness, and righteousness; \Star{} And thy right hand shall lead thee wonderfully.\end{Resp}
\begin{Resp}\Vsign In thy comeliness, and thy beauty, go forward, fare prosperously, and reign.\end{Resp}
\begin{Resp}\Rsign And thy right hand shall lead thee wonderfully.\end{Resp}
\switchcolumn*

\LA
\LessonHead{Lectio V}
\switchcolumn
\EN
\LessonHead{Lesson V}
\switchcolumn*

\LA
\noindent Trigesimum ætátis annum agens, primum Rodardénsis monasterii, ubi religiosam vitam est professa, deínde Elpediáni præses electa, quadragínta annórum spátio, ea caritate, prudéntia, et regularis disciplinæ studio munus obívit, ut cœnobium religiosæ perfectiónis domicílium viderétur. Utrobíque vero, licet esset ómnium mater et magistra, ómnium tamen minima haberi volebat, ac demissióne pari ministram se exhibebat. Quo liberius Deo vacaret, vigiliis, abstinéntia aliisque cruciátibus corpus afflixit; semperque sui símilis, morum innocéntiam, mansuetúdinem ac patiéntiam prætulit singularem. Proximórum salúti omni ope studuit, piæque curæ copiósum fructum rétulit. Divini amoris vi frequéntes patiebátur éxstases, altissimǽque contemplatiónis et divinæ uniónis donum obtinuit.\par
\switchcolumn
\EN
\noindent In the thirtieth year of her age she was elected Abbess of Rodalsdorf, where she had professed herself in the religious life, and afterwards of Heldelfs. This office she bore for forty years in love, wisdom, and zeal for strict observance, so that the house seemed like an ideal example of a sisterhood of perfect nuns. To each one she was a mother and a teacher, and yet would be as the least of all, being in sooth in all lowliness among them as she that served. That she might be more utterly God's only, she tormented her body with sleeplessness, hunger, and other afflictions, but withal ever true to herself, stood forth a pattern of innocency, gentleness, and long-suffering. The salvation of her neighbours was her constant earnest endeavour, and her godly toil bore abundant fruit. The love of God oftentimes threw her into trances, and she was given the grace of the deepest contemplation, even to union of spirit with God.\par
\switchcolumn*

\LA
\begin{Resp}\Rsign Dilexísti justítiam, et odísti iniquitátem: \Star{} Proptérea unxit te Deus, Deus tuus, óleo lætítiæ.\end{Resp}
\begin{Resp}\Vsign Propter veritátem, et mansuetúdinem, et justítiam.\end{Resp}
\begin{Resp}\Rsign Proptérea unxit te Deus, Deus tuus, óleo lætítiæ.\end{Resp}
\switchcolumn
\EN
\begin{Resp}\Rsign Thou hast loved righteousness, and hated iniquity; \Star{} Therefore God, thy God, hath anointed thee with the oil of gladness.\end{Resp}
\begin{Resp}\Vsign Because of truth, and meekness, and righteousness.\end{Resp}
\begin{Resp}\Rsign Therefore God, thy God, hath anointed thee with the oil of gladness.\end{Resp}
\switchcolumn*

\LA
\LessonHead{Lectio VI}
\switchcolumn
\EN
\LessonHead{Lesson VI}
\switchcolumn*

\LA
\noindent Ut meritum acceptíssimæ sibi sponsæ Christus osténderet, in corde Gertrudis jucundam sibi esse mansiónem testátus est. Deiparam Vírginem, véluti matrem et curatricem a Jesu acceptam, pietáte præcipua prosequebátur, ab eáque multa accépit beneficia. Erga diviníssimum Eucharistiæ sacraméntum et passiónem Dómini tanto amóre cum grati animi sensu afficiebátur, ut intérdum ubéribus lácrimis perfunderétur. Justórum ánimas, piacularibus flammis addictas, quotidiánis subsidiis et precibus juvábat. Multa ad confovéndam pietátem scripsit. Divinárum étiam revelatiónum et prophetíæ dono cláruit. Dénique, flagrantíssimo Dei amóre potius quam morbo languéscens, anno Dómini millesimo ducentésimo nonagesimo secundo vita decessit. Miraculis vivens et post mortem a Deo illustráta est.\par
\switchcolumn
\EN
\noindent Christ Himself, to show what such a bride was to Him, revealed that He had in the heart of Gertrude a pleasant dwelling-place. The Virgin Mother of God she ever sought with deep reverence as a mother and warden whom she had received from Jesus Himself, and from her she had many benefits. Toward the most Divine Sacrament of the Eucharist, and the sufferings of the Lord, her soul was moved with love and gratitude, so that she sometimes wept abundantly. She helped with daily gifts and prayers the souls of the just condemned to the purifying fire. She wrote much for the fostering of godliness. She was glorified also by revelations from God, and by the gift of prophecy. Her last illness was rather the wasting of a home-sickness to be with God than a decay of the flesh, and she left this life (to live the undying life in Him, upon the 17th day of November,) in the year of our Lord 1292. God made her bright with miracles both during her life and after her death.\par
\switchcolumn*

\LA
\begin{Resp}\Rsign Afferéntur Regi vírgines post eam, próximæ ejus; \Star{} Afferéntur tibi in lætítia et exsultatióne.\end{Resp}
\begin{Resp}\Vsign Spécie tua et pulchritúdine tua; inténde, próspere procéde, et regna.\end{Resp}
\begin{Resp}\Rsign Afferéntur tibi in lætítia et exsultatióne.\end{Resp}
\begin{Resp}\Vsign Glória Patri, et Fílio, \Star{} et Spirítui Sancto.\end{Resp}
\begin{Resp}\Rsign Afferéntur tibi in lætítia et exsultatióne.\end{Resp}
\switchcolumn
\EN
\begin{Resp}\Rsign After her shall virgins be brought unto the King, her fellows \Star{} They shall be brought unto thee with gladness and rejoicing.\end{Resp}
\begin{Resp}\Vsign In thy comeliness and thy beauty, go forward, fare prosperously, and reign.\end{Resp}
\begin{Resp}\Rsign They shall be brought unto thee with gladness and rejoicing.\end{Resp}
\begin{Resp}\Vsign Glory be to the Father, and to the Son, \Star{} and to the Holy Ghost.\end{Resp}
\begin{Resp}\Rsign They shall be brought unto thee with gladness and rejoicing.\end{Resp}
\switchcolumn*

\end{paracol}

\par\needspace{10\baselineskip}\addvspace{1.2em}{\centering\small\textsc{Die 17 Novembris} · \textsc{17 November}\par}
\Office{S. Gregorii Thaumaturgi Episcopi et Confessoris}{St. Gregory the Wonderworker, Bishop and Confessor}{Semiduplex}
\addcontentsline{toc}{subsection}{Die 17 Novembris · S. Gregorii Thaumaturgi Episcopi et Confessoris \textit{— St. Gregory the Wonderworker, Bishop and Confessor}}
\begin{paracol}{2}
\LA
\RefLine{Lectiones I–III: de Scriptura occurrente.}
\switchcolumn
\EN
\RefLine{Lessons I–III: from the Scripture of the day.}
\switchcolumn*

\LA
\LessonHead{Lectio IV}
\switchcolumn
\EN
\LessonHead{Lesson IV}
\switchcolumn*

\LA
\noindent Gregórius, Neocæsaréæ Ponti epíscopus, sanctitáte doctrinaque illustris, signis vero ac miraculis multo illustrior, quorum multitúdine atque præstántia Thaumatúrgus appellátus est, et sancti Basilíi testimonio cum Moyse, Prophetis et Apóstolis comparátus; montem, qui ecclésiæ ædificatiónem impediebat, oratióne álio tránstulit. Item palúdem, inter fratres causam discordiárum, exsiccávit. Lycum fluvium, perniciose agros inundántem, defixo ad ripam quo sustentabátur baculo, qui statim virentem crevit in árborem, coércuit, ut póstea ultra eum terminum non efflúxerit.\par
\switchcolumn
\EN
\noindent Gregory, Archbishop of Neo-Caesarea, in Pontus, is famous indeed for his holiness and doctrine, but much more so on account of the signs and wonders which he wrought, the number and character of which were so extraordinary that they have gotten him the name of Thaumatourgos, (which is, being interpreted from the Greek, the Wonderworker.) Holy Basil compareth him with Moses, with the Prophets, and with the Apostles, and testifieth that by his prayers he moved a mountain that stood in the way of the building of a Church. Moreover, he dried up a marsh, which was a cause of strife between brothers. Also, when the River Lycus overflowed and wasted the fields, he set his walking-stick on the bank, \textasciitilde{}(which stick forthwith grew into a green tree,) and confined the stream within its bed, so that it never more passed that place again.\par
\switchcolumn*

\LA
\begin{Resp}\Rsign Invéni David servum meum, óleo sancto meo unxi eum: \Star{} Manus enim mea auxiliábitur ei.\end{Resp}
\begin{Resp}\Vsign Nihil profíciet inimícus in eo, et fílius iniquitátis non nocébit ei.\end{Resp}
\begin{Resp}\Rsign Manus enim mea auxiliábitur ei.\end{Resp}
\switchcolumn
\EN
\begin{Resp}\Rsign I have found David My servant, with My holy oil have I anointed him \Star{} For My hand shall help him.\end{Resp}
\begin{Resp}\Vsign The enemy shall prevail nothing against him, nor the son of wickedness afflict him.\end{Resp}
\begin{Resp}\Rsign For My hand shall help him.\end{Resp}
\switchcolumn*

\LA
\LessonHead{Lectio V}
\switchcolumn
\EN
\LessonHead{Lesson V}
\switchcolumn*

\LA
\noindent Sæpíssime dæmones ex idolórum simulacris atque ex hóminum corpóribus ejécit, múltaque alia mirabíliter effécit, quibus innumerábiles hómines traxit ad Jesu Christi fidem, cum étiam prophetico spíritu futura prædiceret. Qui, migratúrus e vita, cum quæsísset quot in civitáte Neocæsariénsi réliqui essent infidéles, responsumque esset tantum esse septemdecim; Deo grátias agens, Tótidem, inquit, erant fidéles, cum cœpi episcopátum. Plura scripsit, quibus etiam, non solum miraculis, Dei Ecclésiam illustrávit.\par
\switchcolumn
\EN
\noindent He oftentimes cast out devils either from heathen idols or from the bodies of men, and did many other marvelous things, whereby he drew countless numbers to believe in Jesus Christ. Also he had the spirit of prophecy, and foretold things to come. When he was at the point of death, he asked how many unbelievers were left in the city of Neo-Caesarea? and when they answered Seventeen, he gave God thanks, and said: Just so many were the faithful when I took the Bishopric. He wrote a great deal, whereby, as well as by his wonders, he hath enlightened the Church of God.\par
\switchcolumn*

\LA
\begin{Resp}\Rsign Pósui adjutórium super poténtem, et exaltávi eléctum de plebe mea: \Star{} Manus enim mea auxiliábitur ei.\end{Resp}
\begin{Resp}\Vsign Invéni David servum meum, óleo sancto meo unxi eum.\end{Resp}
\begin{Resp}\Rsign Manus enim mea auxiliábitur ei.\end{Resp}
\switchcolumn
\EN
\begin{Resp}\Rsign I have laid help upon one that is mighty, and have exalted one chosen out of My people \Star{} For My hand shall help him.\end{Resp}
\begin{Resp}\Vsign I have found David My servant, with My holy oil have I anointed him.\end{Resp}
\begin{Resp}\Rsign For My hand shall help him.\end{Resp}
\switchcolumn*

\LA
\LessonHead{Lectio VI}
\switchcolumn
\EN
\LessonHead{Lesson VI}
\switchcolumn*

\LA
\LessonTitle{Sermo sancti Máximi Epíscopi}
\switchcolumn
\EN
\LessonTitle{From the Sermons of St. Maximus, Bishop (of Turin.)}
\switchcolumn*

\LA
\Source{Ex Homilía 59, quæ est 2 de S. Eusebio}
\switchcolumn
\EN
\Source{59 Homily, being the 2nd on St. Eusebius of Vercelli.}
\switchcolumn*

\LA
\noindent Beáti Patris Gregórii merita jam in tuto posita secúri magnificémus; qui, gubernaculum fidei viríliter tenens, ánchoram spei tranquilla jam in statióne compósuit, et plenam cæléstibus divítiis et ætérnis mercibus navem optato in littore collocávit. Qui, contra omnes adversarios, scutum timoris Dei tamdiu, donec ad victoriam perveníret. Quid enim fuit totus vitæ illíus cursus, nisi uníus cum vígili hoste conflictus?\par
\switchcolumn
\EN
\noindent Our blessed Father Gregory is safe now, and we may safely praise his great deeds. He that kept such a manful hand upon the tiller of faith, hath now cast the anchor of hope in moorings of great calm, and brought his ship, heavy laden with heavenly riches and everlasting merchandise, safe into the heaven where he would be. Thus fareth it now with him who never fainted, but for so long time held up ever the shield of the fear of God against all that did beset him. What was his whole life but one long fight against an enemy that never slept?\par
\switchcolumn*

\LA
\begin{Resp}\Rsign Iste est, qui ante Deum magnas virtútes operátus est, et omnis terra doctrína ejus repléta est: \Star{} Ipse intercédat pro peccátis ómnium populórum.\end{Resp}
\begin{Resp}\Vsign Iste est, qui contémpsit vitam mundi, et pervénit ad cæléstia regna.\end{Resp}
\begin{Resp}\Rsign Ipse intercédat pro peccátis ómnium populórum.\end{Resp}
\begin{Resp}\Vsign Glória Patri, et Fílio, \Star{} et Spirítui Sancto.\end{Resp}
\begin{Resp}\Rsign Ipse intercédat pro peccátis ómnium populórum.\end{Resp}
\switchcolumn
\EN
\begin{Resp}\Rsign This is he which wrought great wonders before God, and the whole earth is full of his teaching \Star{} May he pray for all people, that their sins may be forgiven unto them\end{Resp}
\begin{Resp}\Vsign This is he which loved not his life in this world, and hath attained unto the kingdom of heaven.\end{Resp}
\begin{Resp}\Rsign May he pray for all people, that their sins may be forgiven unto them\end{Resp}
\begin{Resp}\Vsign Glory be to the Father, and to the Son, \Star{} and to the Holy Ghost.\end{Resp}
\begin{Resp}\Rsign May he pray for all people, that their sins may be forgiven unto them\end{Resp}
\switchcolumn*

\LA
\LessonHead{Lectio VII}
\switchcolumn
\EN
\LessonHead{Lesson VII}
\switchcolumn*

\LA
\LessonTitle{Léctio sancti Evangélii secúndum Marcum}
\switchcolumn
\EN
\LessonTitle{From the Holy Gospel according to Mark}
\switchcolumn*

\LA
\Cite{Marc 11:22-24}
\switchcolumn
\EN
\Cite{Mark 11:22-24}
\switchcolumn*

\LA
\noindent In illo témpore: Respóndens Jesus discípulis suis ait illis: Habéte fidem Dei. Amen dico vobis quia quicúmque díxerit huic monti: Tóllere et míttere in mare, et non hæsitáverit in corde suo, sed credíderit quia quodcúmque díxerit fiat, fiet ei. Et réliqua.\par
\switchcolumn
\EN
\noindent And Jesus answering, saith to them: Have the faith of God. Amen I say to you, that whosoever shall say to this mountain, Be thou removed and be cast into the sea, and shall not stagger in his heart, but believe, that whatsoever he saith shall be done; it shall be done unto him. And so on.\par
\switchcolumn*

\LA
\LessonTitle{Homilía sancti Bedæ Venerábilis Presbýteri}
\switchcolumn
\EN
\LessonTitle{Homily by the Venerable Bede, Priest (at Jarrow.)}
\switchcolumn*

\LA
\Source{Liber 3 Comment. in Marc. cap. 11}
\switchcolumn
\EN
\Source{Bk. iii. Comment, on Mark, xi.}
\switchcolumn*

\LA
\noindent Solent Gentiles qui contra Ecclésiam maledicta scripsére, improperare nostris quod non habúerint plenam fidem Dei, quia numquam montes transferre potúerint. Quibus respondéndum est non scripta esse, quæ in Ecclésia sunt gesta, sicut étiam de factis ipsíus Christi et Dómini nostri Scriptúra testátur. Unde et hoc quoque fíeri potuísset, ut mons ablátus de terra mitterétur in mare, si necessitas id fíeri poposcísset. Quómodo légimus factum precibus beáti Patris Gregórii, Neocæsaréæ Ponti antístitis, viri meritis et virtútibus exímii, ut mons in terra tantum loco cederet quantum íncolæ civitátis opus habebant.\par
\switchcolumn
\EN
\noindent The heathen, who have written blasphemies against the Church, are used to cast in our teeth that we have not full faith in God, since we have never been able to move mountains. Such should be answered that we do not possess records of everything that hath come to pass in the Church, any more than, the Scripture being witness, we possess records of all the doings of our Lord Christ Himself. \textasciitilde{}(John xx. 30; xxi. 25.) Mountains may have been removed and cast into the sea, in case of need; a like case, indeed, as we read, was that which came to pass at the prayers of the Blessed Father Gregory, Archbishop of Neo-Caesarea, in Pontus, that right worthy and mighty man, when a mountain was moved from one place on land to another place on land, as the dwellers in the city had need.\par
\switchcolumn*

\LA
\begin{Resp}\Rsign Amávit eum Dóminus, et ornávit eum: stolam glóriæ índuit eum, \Star{} Et ad portas paradísi coronávit eum.\end{Resp}
\begin{Resp}\Vsign Induit eum Dóminus lorícam fídei, et ornávit eum.\end{Resp}
\begin{Resp}\Rsign Et ad portas paradísi coronávit eum.\end{Resp}
\switchcolumn
\EN
\begin{Resp}\Rsign The Lord loved him and beautified him He clothed him with a robe of glory \Star{} And crowned him at the gates of Paradise.\end{Resp}
\begin{Resp}\Vsign The Lord hath put on him the breast-plate of faith, and hath adorned him.\end{Resp}
\begin{Resp}\Rsign And crowned him at the gates of Paradise.\end{Resp}
\switchcolumn*

\LA
\LessonHead{Lectio VIII}
\switchcolumn
\EN
\LessonHead{Lesson VIII}
\switchcolumn*

\LA
\noindent Cum enim, volens ædificáre ecclésiam in loco apto, vidéret eum angustiórem esse quam res exigébat, eo quod ex una parte, rupe maris, ex alia, monte próximo coarctarétur; venit nocte ad locum, et, genibus flexis, admonuit Dóminum promissiónis suæ, ut montem longius juxta fidem petentis ágeret. Et, mane facto, reversus invénit montem tantum spatii reliquisse structóribus ecclésiæ quantum opus habuerant. Poterat ergo hic, poterat alius quis ejusdem meriti vir, si opportunitas exegísset, impetrare a Dómino, mérito fidei, ut étiam mons tollerétur et mitterétur in mare.\par
\switchcolumn
\EN
\noindent Gregory was wishful to build a Church in a meet place, but the site was too narrow, being wedged in between a mountain on the one side and a precipice going down into the sea on the other. He came therefore by night to the place, kneeling down, and reminding the Lord of His promise, and calling upon Him to remove the mountain. And in the morning, when he came thither again, he found that the mountain had been removed back, and as much room left for the builders of the Church as they needed. This man therefore would have been able, and any other man of like grace would have been able, if need were, to obtain of the Lord, by the force of his faith, that even a mountain should be removed, and be cast into the sea.\par
\switchcolumn*

\LA
\begin{Resp}\Rsign Sint lumbi vestri præcíncti, et lucérnæ ardéntes in mánibus vestris: \Star{} Et vos símiles homínibus exspectántibus dóminum suum, quando revertátur a núptiis.\end{Resp}
\begin{Resp}\Vsign Vigiláte ergo, quia nescítis qua hora Dóminus vester ventúrus sit.\end{Resp}
\begin{Resp}\Rsign Et vos símiles homínibus exspectántibus dóminum suum, quando revertátur a núptiis.\end{Resp}
\begin{Resp}\Vsign Glória Patri, et Fílio, \Star{} et Spirítui Sancto.\end{Resp}
\begin{Resp}\Rsign Et vos símiles homínibus exspectántibus dóminum suum, quando revertátur a núptiis.\end{Resp}
\switchcolumn
\EN
\begin{Resp}\Rsign Let your loins be girded about, and your lights burning; \Star{} And ye yourselves like unto men that wait for their lord, when he will return from the wedding.\end{Resp}
\begin{Resp}\Vsign Watch therefore, for ye know not what hour your Lord doth come.\end{Resp}
\begin{Resp}\Rsign And ye yourselves like unto men that wait for their lord, when he will return from the wedding.\end{Resp}
\begin{Resp}\Vsign Glory be to the Father, and to the Son, \Star{} and to the Holy Ghost.\end{Resp}
\begin{Resp}\Rsign And ye yourselves like unto men that wait for their lord, when he will return from the wedding.\end{Resp}
\switchcolumn*

\LA
\LessonHead{Lectio IX}
\switchcolumn
\EN
\LessonHead{Lesson IX}
\switchcolumn*

\LA
\noindent Verum, quia montis nómine nonnumquam diabolus significátur, videlicet propter supérbiam, qua se contra Deum erigit, et esse vult símilis Altíssimo; mons ad præcéptum eórum qui fortes fide sunt, tollitur de terra et in mare projícitur, cum, prædicántibus verbum doctóribus sanctis, immundus spíritus ab eórum corde repéllitur qui ad vitam sunt præordináti, et in turbulentis amarisque infidelium méntibus vesaniam suæ tyrannidis exercére permittitur.\par
\switchcolumn
\EN
\noindent Mystically, however, by a mountain is sometimes signified the devil, on account of the pride whereby he lifteth himself up against God, and would fain be like unto the Most High. And when holy teachers, strong in faith, do preach the Word, this mountain is removed, and cast into the sea, that is to say, the unclean spirit is removed out of the hearts of such as are foreordained unto eternal life, and sent free to exercise the wild rage of his tyranny in the riotous and embittered minds of the unfaithful.\par
\switchcolumn*

\LA
\TeDeum{Te Deum}
\switchcolumn
\EN
\TeDeum{Te Deum}
\switchcolumn*

\end{paracol}

\par\needspace{10\baselineskip}\addvspace{1.2em}{\centering\small\textsc{Die 18 Novembris} · \textsc{18 November}\par}
\Office{In Dedicatione Basilicarum Ss. Apostolorum Petri et Pauli}{In Dedicatione Basilicarum Ss. Apostolorum Petri et Pauli}{Duplex majus}
\addcontentsline{toc}{subsection}{Die 18 Novembris · In Dedicatione Basilicarum Ss. Apostolorum Petri et Pauli \textit{— In Dedicatione Basilicarum Ss. Apostolorum Petri et Pauli}}
\begin{paracol}{2}
\LA
\LessonHead{Lectio I}
\switchcolumn
\EN
\LessonHead{Lesson I}
\switchcolumn*

\LA
\LessonTitle{De libro Apocalýpsis beáti Joánnis Apóstoli.}
\switchcolumn
\EN
\LessonTitle{Lesson from the book of Revelation}
\switchcolumn*

\LA
\Cite{Apo 21:18-20}
\switchcolumn
\EN
\Cite{Rev 21:18-20}
\switchcolumn*

\LA
\noindent Et erat structúra muri ejus ex lápide jáspide; ipsa vero cívitas aurum mundum símile vitro mundo. Et fundaménta muri civitátis omni lápide pretióso ornáta. Fundaméntum primum jaspis, secúndum sapphírus, tértium chalcedónius, quartum smarágdus, quintum sárdonyx, sextum sárdius, séptimum chrysólithus, octávum berýllus, nonum topázius, décimum chrysóprasus, undécimum hyacínthus, duodécimum amethýstus.\par
\switchcolumn
\EN
\noindent And the building of the wall thereof was of jasper stone: but the city itself pure gold, like to clear glass. And the foundations of the wall of the city were adorned with all manner of precious stones. The first foundation was jasper: the second, sapphire: the third, a chalcedony: the fourth, an emerald: The fifth, sardonyx: the sixth, sardius: the seventh, chrysolite: the eighth, beryl: the ninth, a topaz: the tenth, a chrysoprasus: the eleventh, a jacinth: the twelfth, an amethyst.\par
\switchcolumn*

\LA
\begin{Resp}\Rsign In dedicatióne templi decantábat pópulus laudem: \Star{} Et in ore eórum dulcis resonábat sonus.\end{Resp}
\begin{Resp}\Vsign Fundáta est domus Dómini supra vérticem móntium, et vénient ad eam omnes gentes.\end{Resp}
\begin{Resp}\Rsign Et in ore eórum dulcis resonábat sonus.\end{Resp}
\switchcolumn
\EN
\begin{Resp}\Rsign When the Temple was dedicated the people sang praise; \Star{} And sweet in their mouths was the sound.\end{Resp}
\begin{Resp}\Vsign The Lord's house is established in the top of the mountains; and all nations shall flow unto it.\end{Resp}
\begin{Resp}\Rsign And sweet in their mouths was the sound.\end{Resp}
\switchcolumn*

\LA
\LessonHead{Lectio II}
\switchcolumn
\EN
\LessonHead{Lesson II}
\switchcolumn*

\LA
\Cite{Apo 21:21-23}
\switchcolumn
\EN
\Cite{Rev 21:21-23}
\switchcolumn*

\LA
\noindent Et duódecim portæ duódecim margarítæ sunt, per síngulas: et síngulæ portæ erant ex síngulis margarítis: et platéa civitátis aurum mundum, tamquam vitrum perlúcidum. Et templum non vidi in ea; Dóminus enim Deus omnípotens templum illíus est, et Agnus. Et cívitas non eget sole neque luna, ut lúceant in ea; nam cláritas Dei illuminávit eam, et lucérna ejus est Agnus.\par
\switchcolumn
\EN
\noindent And the twelve gates are twelve pearls, one to each: and every several gate was of one several pearl. And the street of the city was pure gold, as it were transparent glass. And I saw no temple therein. For the Lord God Almighty is the temple thereof, and the Lamb. And the city hath no need of the sun, nor of the moon, to shine in it. For the glory of God hath enlightened it, and the Lamb is the lamp thereof.\par
\switchcolumn*

\LA
\begin{Resp}\Rsign Fundáta est domus Dómini supra vérticem móntium, et exaltáta est super omnes colles: \Star{} Et vénient ad eam omnes gentes, et dicent: Glória tibi, Dómine.\end{Resp}
\begin{Resp}\Vsign Veniéntes autem vénient cum exsultatióne, portántes manípulos suos.\end{Resp}
\begin{Resp}\Rsign Et vénient ad eam omnes gentes, et dicent: Glória tibi, Dómine.\end{Resp}
\switchcolumn
\EN
\begin{Resp}\Rsign The Lord's house is established in the top of the mountains, and exalted above the hills; \Star{} And all nations shall flow unto it, and shall say: Glory be to thee, O Lord!\end{Resp}
\begin{Resp}\Vsign They shall doubtless come again with rejoicing, bringing their sheaves with them.\end{Resp}
\begin{Resp}\Rsign And all nations shall flow unto it, and shall say: Glory be to thee, O Lord!\end{Resp}
\switchcolumn*

\LA
\LessonHead{Lectio III}
\switchcolumn
\EN
\LessonHead{Lesson III}
\switchcolumn*

\LA
\Cite{Apo 21:24-27}
\switchcolumn
\EN
\Cite{Rev 21:24-27}
\switchcolumn*

\LA
\noindent Et ambulábunt gentes in lúmine ejus: et reges terræ áfferent glóriam suam et honórem in illam. Et portæ ejus non claudéntur per diem: nox enim non erit illic. Et áfferent glóriam et honórem géntium in illam. Non intrábit in eam áliquod coinquinátum, aut abominatiónem fáciens et mendácium, nisi qui scripti sunt in libro vitæ Agni.\par
\switchcolumn
\EN
\noindent And the nations shall walk in the light of it: and the kings of the earth shall bring their glory and honour into it. And the gates thereof shall not be shut by day: for there shall be no night there. And they shall bring the glory and honour of the nations into it. There shall not enter into it any thing defiled, or that worketh abomination or maketh a lie, but they that are written in the book of life of the Lamb.\par
\switchcolumn*

\LA
\begin{Resp}\Rsign Bénedic, Dómine, domum istam, quam ædificávi nómini tuo: veniéntium in loco isto \Star{} Exáudi preces in excélso sólio glóriæ tuæ.\end{Resp}
\begin{Resp}\Vsign Dómine, si convérsus fúerit pópulus tuus, et oráverit ad sanctuárium tuum.\end{Resp}
\begin{Resp}\Rsign Exáudi preces in excélso sólio glóriæ tuæ.\end{Resp}
\begin{Resp}\Vsign Glória Patri, et Fílio, \Star{} et Spirítui Sancto.\end{Resp}
\begin{Resp}\Rsign Exáudi preces in excélso sólio glóriæ tuæ.\end{Resp}
\switchcolumn
\EN
\begin{Resp}\Rsign O Lord, bless this house which I have built unto thy name. whosoever shall come unto this place and pray, then \Star{} Hear thou from the excellent throne of thy glory.\end{Resp}
\begin{Resp}\Vsign O Lord, if thy people turn and pray toward thy sanctuary.\end{Resp}
\begin{Resp}\Rsign Hear thou from the excellent throne of thy glory.\end{Resp}
\begin{Resp}\Vsign Glory be to the Father, and to the Son, \Star{} and to the Holy Ghost.\end{Resp}
\begin{Resp}\Rsign Hear Thou from the excellent throne of thy glory.\end{Resp}
\switchcolumn*

\LA
\LessonHead{Lectio IV}
\switchcolumn
\EN
\LessonHead{Lesson IV}
\switchcolumn*

\LA
\noindent Ex locis sacris quæ olim apud Christiános veneratiónem habuérunt, illa celebérrima et frequentíssima fuérunt, in quibus cóndita Sanctórum córpora, vel áliquod Mártyrum vestígium aut monuméntum esset. In quorum número sanctórum sanctórum locórum, in primis semper fuit insígnis ea Vaticáni pars, quam sancti Petri Confessiónem appellábant. Nam eo Christiáni ex ómnibus orbis terræ pártibus, tamquam ad fídei petram et Ecclésiæ fundaméntum conveniéntes, locum, Príncipis Apostolórum sepúlcro consecrátum, summa religióne ac pietáte venerabántur.\par
\switchcolumn
\EN
\noindent Among the hallowed places which have from of old time been held in honour among Christians, the most famous and sought after were those where the bodies of the Saints were buried, or where there was some trace or token of the Martyrs. Among these spots so hallowed hath been ever among the most noteworthy that place on the Vatican Hill which is called the Confession of St. Peter. Thither Christians do come from all parts of the earth as unto the rock of faith and the foundation-stone of the Church, and surround with godly reverence and love the spot hallowed by the grave of the Prince of the Apostles.\par
\switchcolumn*

\LA
\begin{Resp}\Rsign Orántibus in loco isto, \Star{} Dimítte peccáta pópuli tui, Deus, et osténde eis viam bonam per quam ámbulent, et da glóriam in loco isto.\end{Resp}
\begin{Resp}\Vsign Qui regis Israël inténde, qui dedúcis velut ovem Joseph, qui sedes super Chérubim.\end{Resp}
\begin{Resp}\Rsign Dimítte peccáta pópuli tui, Deus, et osténde eis viam bonam per quam ámbulent, et da glóriam in loco isto.\end{Resp}
\switchcolumn
\EN
\begin{Resp}\Rsign If they pray toward this place, \Star{} Forgive the sin of thy people, O God, and teach them the good way wherein they should walk, and manifest forth thy glory in this place.\end{Resp}
\begin{Resp}\Vsign Give ear, O Shepherd of Israel, thou that leadest Joseph like a flock, thou that sittest upon the cherubim.\end{Resp}
\begin{Resp}\Rsign Forgive the sin of thy people, O God, and teach them the good way wherein they should walk, and manifest forth thy glory in this place.\end{Resp}
\switchcolumn*

\LA
\LessonHead{Lectio V}
\switchcolumn
\EN
\LessonHead{Lesson V}
\switchcolumn*

\LA
\noindent Illuc Constantínus Magnus imperátor octávo die post suscéptum baptísmum venit, depositóque diadémate, et humi jacens, vim lacrimárum profúdit. Mox, sumpto ligóne ac bidénte, terram éruit; indéque duódecim terræ cóphinis, honóris causa duódecim Apostolórum, ablátis, ac loco basílicæ Príncipis Apostolórum designáto, ecclésiam ædificávit. Quam sanctus Silvéster Papa décimo quarto Kaléndas Decémbris, eo modo quo Lateranénsem ecclésiam quinto Idus Novémbris consecráverat, dedicávit, et in ea altáre lapídeum, chrísmate delibútum, eréxit; atque ex eo témpore sancívit ne deínceps altária nisi ex lápide fíerent.\par
Idem beátus Silvester basilicam sancti Pauli Apóstoli, in via Ostiénsi ab eodem Constantino imperatóre magnificentíssime ædificatam, dedicávit. Quas basilicas idem imperátor multis prædiis attributis locupletávit, ac munéribus amplíssimis exornávit.\par
\switchcolumn
\EN
\noindent Whither came the Emperor Constantine the Great upon the eighth day after his Baptism, and, taking off his crown, cast himself down upon the ground, and wept abundantly. Then presently he took a spade and pick-axe, and began to break up the earth, whereof he carried away twelve baskets-full in honour of the twelve Apostles, and built a Church upon that spot, appointed for the Cathedral Church of the Prince of the Apostles. This Church was hallowed by holy Pope Sylvester upon the 18th day of November, in like manner as he had hallowed the Church of the Lateran upon the 9th day of the same month. In this Church did the Pope set up an altar of stone, and pour ointment thereon, and ordain that from thenceforth no altars should be set up, save of stone. The same Emperor Constantine likewise built a very stately Church upon the road to Ostia, in honour of the holy Apostle Paul, which Church also was hallowed by the blessed Sylvester. These Churches the Emperor enriched by grants of much land, and adorned with exceedingly rich gifts.\par
\switchcolumn*

\LA
\begin{Resp}\Rsign O quam metuéndus est locus iste: \Star{} Vere non est hic áliud nisi domus Dei et porta cæli.\end{Resp}
\begin{Resp}\Vsign Hæc est domus Dómini fírmiter ædificáta, bene fundáta est supra firmam petram.\end{Resp}
\begin{Resp}\Rsign Vere non est hic áliud nisi domus Dei et porta cæli.\end{Resp}
\switchcolumn
\EN
\begin{Resp}\Rsign How dreadful is this place! \Star{} Surely this is none other but the house of God, and this is the gate of heaven.\end{Resp}
\begin{Resp}\Vsign This is the house of God, stoutly builded, well founded upon a sure rock.\end{Resp}
\begin{Resp}\Rsign Surely this is none other but the house of God, and this is the gate of heaven.\end{Resp}
\switchcolumn*

\LA
\LessonHead{Lectio VI}
\switchcolumn
\EN
\LessonHead{Lesson VI}
\switchcolumn*

\LA
\noindent Porro Vaticanam basilicam, vetustáte jampridem collabentem, ac proptérea multórum Pontificum pietáte latius ac magnificentius a fundamentis erectam, Urbanus octavus, hac eádem recurrénte die anni millesimi sexcentésimi vigesimi sexti, solemni ritu consecrávit. Basilicam vero Ostiensem, quum dira incendii vis anno millesimo octingentésimo vigesimo tertio penitus consumpsísset, indefessa quatuor Pontificum cura splendidius quam antea erectam, et ab intéritu véluti vindicátum, Pius nonus, auspicatíssimam nactus occasiónem qua dogma de immaculáta beátæ Maríæ Vírginis Conceptióne, nuper ab ipso proclamátum, ingentem cardinalium et episcopórum númerum ex dissitis étiam catholici orbis regiónibus Romam attraxerat; die décima Decembris anni millesimi octingentésimi quinquagesimi quarti, tanta circúmdatus purpuratórum patrum et antistitum coróna, solemniter dedicávit, ejusque celebritátis memóriam hac die recoléndam decrevit.\par
\switchcolumn
\EN
\noindent The Church of St. Peter upon the Vatican fell in course of time to ruins, and having been rebuilt from the foundations, enlarged and garnished, by the zeal of many Popes, was solemnly consecrated anew by Urban VIII, upon the same day, in the year 1628. The Church of St. Paul upon the road to Ostia was almost entirely consumed by fire in the year 1823, but was rebuilt in a more splendid form and, as it were, raised from the dead, by the unwearied zeal of four successive Popes. In the year 1854 Pius IX seized the happy occasion when the doctrine concerning the Immaculate Conception of the Virgin Mary, which he had just set forth, had drawn together to Rome a great multitude of Cardinals and Bishops from all quarters of the Catholic world, solemnly to dedicate this new Church in their presence upon the 10th day of December in the year aforesaid; but he decreed that the yearly Feast in honour of that dedication should be kept upon this day, being the same as that of the Dedication of the Church of St. Peter.\par
\switchcolumn*

\LA
\begin{Resp}\Rsign Mane surgens Jacob, erigébat lápidem in títulum, fundens óleum désuper; votum vovit Dómino: \Star{} Vere locus iste sanctus est, et ego nesciébam.\end{Resp}
\begin{Resp}\Vsign Cumque evigilásset Jacob de somno, ait.\end{Resp}
\begin{Resp}\Rsign Vere locus iste sanctus est, et ego nesciébam.\end{Resp}
\begin{Resp}\Vsign Glória Patri, et Fílio, \Star{} et Spirítui Sancto.\end{Resp}
\begin{Resp}\Rsign Vere locus iste sanctus est, et ego nesciébam.\end{Resp}
\switchcolumn
\EN
\begin{Resp}\Rsign Jacob rose up early in the morning, and set up the stone for a pillar, and poured oil upon the top of it, and vowed a vow unto the Lord. \Star{} Surely this place is holy, and I knew it not.\end{Resp}
\begin{Resp}\Vsign And Jacob awaked out of his sleep, and he said\end{Resp}
\begin{Resp}\Rsign Surely this place is holy, and I knew it not.\end{Resp}
\begin{Resp}\Vsign Glory be to the Father, and to the Son, \Star{} and to the Holy Ghost.\end{Resp}
\begin{Resp}\Rsign Surely this place is holy, and I knew it not.\end{Resp}
\switchcolumn*

\LA
\LessonHead{Lectio VII}
\switchcolumn
\EN
\LessonHead{Lesson VII}
\switchcolumn*

\LA
\LessonTitle{Léctio sancti Evangélii secúndum Lucam.}
\switchcolumn
\EN
\LessonTitle{From the Holy Gospel according to Luke}
\switchcolumn*

\LA
\Cite{Luc 19:1-10}
\switchcolumn
\EN
\Cite{Luke 19:1-10}
\switchcolumn*

\LA
\noindent In illo témpore: Ingréssus Jesus perambulábat Jéricho. Et ecce vir nómine Zachǽus: et hic princeps erat publicanórum, et ipse dives. Et réliqua.\par
\switchcolumn
\EN
\noindent At that time: Jesus entering in, he walked through Jericho. And behold, there was a man named Zacchaeus, which was the chief among the publicans, and he was rich. And so on.\par
\switchcolumn*

\LA
\LessonTitle{Homilía sancti Gregórii Papæ.}
\switchcolumn
\EN
\LessonTitle{Homily by Pope St. Gregory (the Great)}
\switchcolumn*

\LA
\Source{Lib. 27. Moralium, cap. 27. post medium.}
\switchcolumn
\EN
\Source{Bk. xxviii of Moral Reflections on Job, ch. 27}
\switchcolumn*

\LA
\noindent Si veráciter sapiéntes esse, atque ipsam sapiéntiam contemplári appétimus, stultos nos humíliter cognoscámus. Relinquámus nóxiam sapiéntiam, discámus laudábilem fatuitátem. Hinc quippe scriptum est: Stulta mundi elégit Deus, ut confúndat sapiéntes. Hinc rursum dícitur: Si quis vidétur inter vos sápiens esse in hoc sǽculo, stultus fiat, ut sit sápiens. Hinc evangélicæ históriæ verba testántur: quia Zachǽus cum vidére præ turba nihil posset, sycómori árborem ascéndit, ut transeúntem Dóminum cérneret. Sycómorus quippe ficus fátua dícitur.\par
\switchcolumn
\EN
\noindent If we would be truly wise, and behold wisdom herself, we must humbly acknowledge ourselves to be fools. Let us cast away harmful wisdom, and learn praiseworthy folly. For this reason indeed is it written God hath chosen the foolish things of the world, to confound the wise. (1 Cor. 1:27) And again it is said If any man among you seemeth to be wise in this world, let him become a fool, that he may be wise. (3:18.) And unto this doth the very Gospel bear witness, wherein it is said that Zacchaeus sought to see Jesus, Who He was; and could not for the press, because he was little of stature. And he ran before, and climbed up into a sycamore tree to see Him; for He was to pass that way. For this name Sycamore, being interpreted, signifieth the Foolish Fig.\par
\switchcolumn*

\LA
\begin{Resp}\Rsign Domus mea, domus oratiónis vocábitur, dicit Dóminus: in ea omnis qui petit, áccipit; et qui quærit, ínvenit; \Star{} Et pulsánti aperiétur.\end{Resp}
\begin{Resp}\Vsign Pétite, et accipiétis; quǽrite, et inveniétis.\end{Resp}
\begin{Resp}\Rsign Et pulsánti aperiétur.\end{Resp}
\switchcolumn
\EN
\begin{Resp}\Rsign My house shall be called the house of prayer, saith the Lord. Therein, he that asketh, receiveth; he that seeketh, findeth; \Star{} And to him that knocketh, it shall be opened.\end{Resp}
\begin{Resp}\Vsign Ask, and ye shall receive; seek, and ye shall find.\end{Resp}
\begin{Resp}\Rsign And to him that knocketh, it shall be opened.\end{Resp}
\switchcolumn*

\LA
\LessonHead{Lectio VIII}
\switchcolumn
\EN
\LessonHead{Lesson VIII}
\switchcolumn*

\LA
\noindent Pusíllus ítaque Zachǽus sycómorum súbiit, et Dóminum vidit: quia qui mundi stultítiam humíliter éligunt, ipsi Dei sapiéntiam subtíliter contemplántur. Pusillitátem namque nostram ad vidéndum Dóminum turba prǽpedit: quia infirmitátem humánæ mentis, ne lucem veritátis inténdat, curárum sæculárium tumúltus premit. Sed prudénter sycómorum ascéndimus, si próvide eam quæ divínitus præcípitur, stultítiam mente tenémus. Quid enim in hoc mundo stúltius quam amíssa non quǽrere, posséssa rapiéntibus relaxáre, nullam pro accéptis injúriis injúriam réddere, immo adjúnctis áliis patiéntiam præbére?\par
\switchcolumn
\EN
\noindent Little Zacchaeus therefore accepted the humiliation of having recourse to the sycamore and saw the Lord. They who humbly choose to be fools in the estimation of the world, have a deep insight into the wisdom of God. The press standeth in our way, on account of our little stature, when we are fain to see the Lord; for the toilsome din of worldly business tormenteth our weak minds, so as to hinder our perceiving the light of the truth. But we climb up wisely into the sycamore tree, if we willingly give up our minds to that folly which God giveth unto us. What can be more utter folly (in this world) than not to seek for that we have lost, to leave that whereof we have been robbed in the hands of our despoilers, to take no revenge for wrongs which have been done us, yea, even to offer to him that taketh away our cloak, our coat also, and be patient?\par
\switchcolumn*

\LA
\begin{Resp}\Rsign Lápides pretiósi omnes muri tui, \Star{} Et turres Jerúsalem gemmis ædificabúntur.\end{Resp}
\begin{Resp}\Vsign Portæ Jerúsalem ex sapphíro et smarágdo ædificabúntur, et ex lápide pretióso omnis circúitus muri ejus.\end{Resp}
\begin{Resp}\Rsign Et turres Jerúsalem gemmis ædificabúntur.\end{Resp}
\begin{Resp}\Vsign Glória Patri, et Fílio, \Star{} et Spirítui Sancto.\end{Resp}
\begin{Resp}\Rsign Et turres Jerúsalem gemmis ædificabúntur.\end{Resp}
\switchcolumn
\EN
\begin{Resp}\Rsign All thy walls are of stones most precious. \Star{} The towers of Jerusalem shall be built up with jewels.\end{Resp}
\begin{Resp}\Vsign The gates of Jerusalem shall be built up with the sapphire stone, and the emerald, and all her walls round about with stones most precious.\end{Resp}
\begin{Resp}\Rsign The towers of Jerusalem shall be built up with jewels.\end{Resp}
\begin{Resp}\Vsign Glory be to the Father, and to the Son, \Star{} and to the Holy Ghost.\end{Resp}
\begin{Resp}\Rsign The towers of Jerusalem shall be built up with jewels.\end{Resp}
\switchcolumn*

\LA
\LessonHead{Lectio IX}
\switchcolumn
\EN
\LessonHead{Lesson IX}
\switchcolumn*

\LA
\noindent Quasi enim sycómorum nos ascéndere Dóminus prǽcipit, cum dicit: Qui aufert quæ tua sunt, ne répetas. Et rursum: Si quis te percússerit in déxteram maxíllam, præbe illi et álteram. Per sycómorum Dóminus tránsiens cérnitur: quia per hanc sapiéntem stultítiam, etsi necdum, ut est, sólide, jam tamen per contemplatiónis lumen Dei sapiéntia quasi in tránsitu vidétur, quam vidére néqueunt qui sibi sapiéntes esse vidéntur: quia ad conspiciéndum Dóminum, in eláta cogitatiónum suárum turba deprehénsi, adhuc sycómori árborem non invenérunt.\par
\switchcolumn
\EN
\noindent The Lord biddeth us, as it were, to climb up into the sycamore, where He saith Of him that taketh away thy goods, ask them not again. (Luke 6:30.) And again Whosoever shall smite thee on thy right cheek, turn to him the other also. \textasciitilde{}(Matth. 5:39.) From the boughs of this sycamore tree, the Lord is seen passing by. He may indeed, as yet, not be seen face to Face, but by this wise folly the inward eye may see the Wisdom of God, as it were, passing by, even that Wisdom Which they that are wise in their own conceit cannot see. They are mixed up in the overbearing press of their own imaginations, and have not yet found the sycamore tree where into to climb up, if they would see the Lord.\par
\switchcolumn*

\LA
\TeDeum{Te Deum}
\switchcolumn
\EN
\TeDeum{Te Deum}
\switchcolumn*

\end{paracol}

\par\needspace{10\baselineskip}\addvspace{1.2em}{\centering\small\textsc{Die 19 Novembris} · \textsc{19 November}\par}
\Office{S. Elisabeth Viduæ}{St. Elizabeth, Widow}{Duplex}
\addcontentsline{toc}{subsection}{Die 19 Novembris · S. Elisabeth Viduæ \textit{— St. Elizabeth, Widow}}
\begin{paracol}{2}
\LA
\RefLine{Lectiones I–III: de Scriptura occurrente.}
\switchcolumn
\EN
\RefLine{Lessons I–III: from the Scripture of the day.}
\switchcolumn*

\LA
\RefLine{Lectiones VII–VIII: ex Commúni non Virginum non Martyrum (C7a).}
\switchcolumn
\EN
\RefLine{Lessons VII–VIII: from the Common of Widows (C7a).}
\switchcolumn*

\LA
\RefLine{Lectio IX: de S. Pontiani Papæ et Martyris (11-19o).}
\switchcolumn
\EN
\RefLine{Lesson IX: of St. Pontianus, Pope and Martyr (11-19o).}
\switchcolumn*

\LA
\LessonHead{Lectio IV}
\switchcolumn
\EN
\LessonHead{Lesson IV}
\switchcolumn*

\LA
\noindent Elisabeth, Andréæ regis Hungariæ fília, ab infántia Deum timere cœpit; et, crescens ætate, crevit étiam pietáte. Ludovíco Lantgrávio Hássiæ et Thuríngiæ in cónjugem copulata, non minori cura quæ Dei quam quæ viri sui erant, exsequebátur. Surgens enim noctúrno témpore, oratióni diu incumbebat; ac, variis misericórdiæ officiis dédita, viduis, pupíllis, ægrotis, egéntibus sedulo inserviébat, gravique fame urgente, domus suæ fruménta liberáliter erogábat. Leprosos hospítio suscípiens, manus eórum et pedes osculabátur. Curandis autem et aléndis paupéribus insigne xenodochíum construxit.\par
\switchcolumn
\EN
\noindent Elisabeth, daughter of Andrew II., King of Hungary, (was born in the year 1207.) She began to fear God even from a little child, and grew in grace as she grew in years. (In her fourteenth year) she was married to Lewis, Landgrave of Hesse and Thuringia, and thenceforth gave herself up to the things of her husband, with as much zeal as to the things of God. She rose in the night to make long prayers. She consecrated herself to works of mercy. She waited continually on widows and orphans, the sick and the needy. When a sore famine came (in the year 1225,) she provided corn bountifully from her own house. She founded an house of refuge for lepers, and would even kiss their hands and feet. She built also a great hospital for the suffering and starving poor.\par
\switchcolumn*

\LA
\begin{Resp}\Rsign Propter veritátem, et mansuetúdinem, et justítiam: \Star{} Et dedúcet te mirabíliter déxtera tua.\end{Resp}
\begin{Resp}\Vsign Spécie tua et pulchritúdine tua inténde, próspere procéde, et regna.\end{Resp}
\begin{Resp}\Rsign Et dedúcet te mirabíliter déxtera tua.\end{Resp}
\switchcolumn
\EN
\begin{Resp}\Rsign Because of truth, and meekness, and righteousness; \Star{} And thy right hand shall lead thee wonderfully.\end{Resp}
\begin{Resp}\Vsign In thy comeliness, and thy beauty, go forward, fare prosperously, and reign.\end{Resp}
\begin{Resp}\Rsign And thy right hand shall lead thee wonderfully.\end{Resp}
\switchcolumn*

\LA
\LessonHead{Lectio V}
\switchcolumn
\EN
\LessonHead{Lesson V}
\switchcolumn*

\LA
\noindent Defuncto cónjuge, ut Deo liberius serviret, depositis ómnibus sæcularis glóriæ induméntis, vili túnica induta est, atque, ordinem Pœniténtium sancti Francisci ingressa, patiéntiæ et humilitátis virtúte maxime enituit. Nam, bonis ómnibus exuta, a propriis ædibus ejecta, ab ómnibus derelicta, contumelias, irrisiónes, obtrectatiónes invicto animo tolerávit, adeo ut summopere gaudéret se talia pro Deo pati. Ad ínfima quæque ministeria erga páuperes et ægrotos se abjíciens, eis necessaria procurabat, solis oléribus et legumínibus pro suo victu contenta.\par
\switchcolumn
\EN
\noindent After husband died (on his way to the Holy War, on the eleventh day of September, 1227.) Then Elizabeth, more utterly to be God's only, laid aside all the garments of earthly state, clad herself in mean raiment, and entered the Third Order of St. Francis, wherein she was a burning and shining light of longsuffering and lowliness. (Her brotherin-law) stripped her (and her three little children) of all their goods, and turned them out of their own house. She was deserted by all, and assailed with insults, gibes, and calumnies, but she bore it all with patience, yea, even rejoicing that she suffered such things for God's sake. She gave herself to the meanest services toward the poor and sick, and sought for them the needfuls of life, while she lived herself only on potherbs and vegetables.\par
\switchcolumn*

\LA
\begin{Resp}\Rsign Dilexísti justítiam, et odísti iniquitátem: \Star{} Proptérea unxit te Deus, Deus tuus, óleo lætítiæ.\end{Resp}
\begin{Resp}\Vsign Propter veritátem, et mansuetúdinem, et justítiam.\end{Resp}
\begin{Resp}\Rsign Proptérea unxit te Deus, Deus tuus, óleo lætítiæ.\end{Resp}
\switchcolumn
\EN
\begin{Resp}\Rsign Thou hast loved righteousness, and hated iniquity; \Star{} Therefore God, thy God, hath anointed thee with the oil of gladness.\end{Resp}
\begin{Resp}\Vsign Because of truth, and meekness, and righteousness.\end{Resp}
\begin{Resp}\Rsign Therefore God, thy God, hath anointed thee with the oil of gladness.\end{Resp}
\switchcolumn*

\LA
\LessonHead{Lectio VI}
\switchcolumn
\EN
\LessonHead{Lesson VI}
\switchcolumn*

\LA
\noindent Cum vero in his aliisque plurimis sanctis opéribus vitam religiosíssime transegísset, finis tandem suæ peregrinatiónis advenit, quem domesticis suis ante prædixit. Cumque, defixis in cælum óculis, divinæ contemplatióni vacaret, a Deo mirabíliter recreáta et sacramentis refecta, obdormívit in Dómino. Statimque plurima ad ejus tumulum miracula patráta sunt. Quibus auditis et rite probátis, Gregórius nonus Sanctórum número eam adscripsit.\par
\switchcolumn
\EN
\noindent In these and many other holy works she prayerfully passed the rest of her life, till (in the twenty-fourth year of her age,) the end of her earthly pilgrimage came, as she had already foretold to her servants. With her eyes fixed on heaven, absorbed in the thought of God, by Him wondrously comforted, and strengthened by the Sacraments, she fell asleep in the Lord, (upon the 19th day of November, in the year of salvation 1231.) Forthwith many miracles were wrought at her grave, which being known and duly proved, Gregory IX. numbered her name among those of the Saints.\par
\switchcolumn*

\LA
\begin{Resp}\Rsign Fallax grátia, et vana est pulchritúdo: \Star{} Múlier timens Dóminum, ipsa laudábitur.\end{Resp}
\begin{Resp}\Vsign Date ei de fructu mánuum suárum, et laudent eam in portis ópera ejus.\end{Resp}
\begin{Resp}\Rsign Múlier timens Dóminum, ipsa laudábitur.\end{Resp}
\begin{Resp}\Vsign Glória Patri, et Fílio, \Star{} et Spirítui Sancto.\end{Resp}
\begin{Resp}\Rsign Múlier timens Dóminum, ipsa laudábitur.\end{Resp}
\switchcolumn
\EN
\begin{Resp}\Rsign Favour is deceitful, and beauty is vain \Star{} A woman that feareth God she shall be praised.\end{Resp}
\begin{Resp}\Vsign Give her of the fruit of her hands, and let her own works praise her in the gates.\end{Resp}
\begin{Resp}\Rsign A woman that feareth God she shall be praised.\end{Resp}
\begin{Resp}\Vsign Glory be to the Father, and to the Son, \Star{} and to the Holy Ghost.\end{Resp}
\begin{Resp}\Rsign A woman that feareth God she shall be praised.\end{Resp}
\switchcolumn*

\end{paracol}

\par\needspace{10\baselineskip}\addvspace{1.2em}{\centering\small\textsc{Die 19 Novembris} · \textsc{19 November}\par}
\Office{S. Pontiani Papæ et Martyris}{St. Pontianus, Pope and Martyr}{Simplex}
\addcontentsline{toc}{subsection}{Die 19 Novembris · S. Pontiani Papæ et Martyris \textit{— St. Pontianus, Pope and Martyr}}
\begin{paracol}{2}
\LA
\LessonHead{Lectio IX (pro commemoratione)}
\switchcolumn
\EN
\LessonHead{Lesson IX (when commemorated)}
\switchcolumn*

\LA
\LessonTitle{Commemoratio: S. Pontiani Papæ et Martyris}
\switchcolumn
\EN
\LessonTitle{Commemoration of: St. Pontianus, Pope and Martyr}
\switchcolumn*

\LA
\noindent Pontianus, Romanus, præfuit Ecclésiæ Alexandro imperatóre; qui sanctum Pontificem, propter christianæ fidei confessiónem, in Sardiniam insulam cum Hippólyto presbytero relegávit. Ubi is, pro Christi fide multis calamitátibus afflíctus, tertio Kalendas Novembris e vita migrávit. Ejus corpus, Fabiáno Pontifice, cum clero Romam delátum, in cœmeterio Callísti via Appia sepelítur. Sedit annos quatuor, menses quatuor, dies viginti quinque. Fecit ordinatiónes duas mense Decembri, quibus creávit presbyteros sex, diaconos quinque, episcopos per diversa loca sex.\par
\switchcolumn
\EN
\noindent Pontian was a Roman, who ruled the Church in the reign of the Emperor Alexander. This Emperor banished him into the Island of Sardinia, along with the Priest Hippolytus, on account of their profession of the Christian faith. There he. endured many hardships because of his belief in Christ, and departed this life upon the 19th day of November, (in the year of our Lord 235.) His body was brought to Rome by Pope Fabian and his clergy, and buried in the cemetery of Kallistus, upon the Appian Way. He sat in the seat of Peter four years, four months, and twentyfive days. He held two Ordinations in the month of December, wherein he made six Priests, five Deacons, and six Bishops for diverse places. At Lauds a Commemoration is made of the holy Martyr.\par
\switchcolumn*

\LA
\TeDeum{Te Deum}
\switchcolumn
\EN
\TeDeum{Te Deum}
\switchcolumn*

\end{paracol}

\par\needspace{10\baselineskip}\addvspace{1.2em}{\centering\small\textsc{Die 20 Novembris} · \textsc{20 November}\par}
\Office{S. Felicis de Valois Confessoris}{St. Felix de Valois, Confessor}{Duplex}
\addcontentsline{toc}{subsection}{Die 20 Novembris · S. Felicis de Valois Confessoris \textit{— St. Felix de Valois, Confessor}}
\begin{paracol}{2}
\LA
\RefLine{Lectiones I–III: de Scriptura occurrente.}
\switchcolumn
\EN
\RefLine{Lessons I–III: from the Scripture of the day.}
\switchcolumn*

\LA
\LessonHead{Lectio IV}
\switchcolumn
\EN
\LessonHead{Lesson IV}
\switchcolumn*

\LA
\noindent Felix, Hugo ántea dictus, ex regali Valesiórum família ortus in Gállia, ab ineunte ætate non lévia dedit futuræ sanctitátis indicia, præsertim misericórdiæ erga páuperes. Nam, adhuc infantulus, manu propria, ac si grandior esset et judícii maturitáte polleret, nummos egenis distríbuit; jam grandiúsculus, solebat ex appositis in mensa dápibus ad ipsos mittere, et ferme eo, quod sapidius erat, obsonio paupérculos púeros recreábat; adoléscens, non semel vestibus se expoliávit, ut inopes cooperíret. Ab avunculo Theobaldo, Xamphánæ et Blesii comite, vitam reo mortis impetrávit, prædicens hunc infamem háctenus sicárium, mox sanctíssimis præditum móribus evasurum: verídicum testimónium monstrávit eventus.\par
\switchcolumn
\EN
\noindent Hugo de Valois, who afterwards took the name of Felix, was born (in the year 1127) of the same family of the de Valois which in after times became Kingly. From his earliest childhood he gave tokens, especially by his pity toward the poor, of the holiness of his coming life. When he was still a little lad he distributed money to the poor with his own hand, with the seriousness of an old man. When he was a little bigger he used to send them dishes from the table, and took especial delight in treating poor children with the most toothsome of the sweetmeats. As a boy he took clothes off his own back more than once, to cover the naked. He begged and obtained from his uncle Theobald, Earl of Champagne and Blois, the life of a felon condemned to death, foretelling to him that this blackguard cut-throat would yet become a man of most holy life which did indeed come to pass as he had said.\par
\switchcolumn*

\LA
\begin{Resp}\Rsign Honéstum fecit illum Dóminus, et custodívit eum ab inimícis, et a seductóribus tutávit illum: \Star{} Et dedit illi claritátem ætérnam.\end{Resp}
\begin{Resp}\Vsign Justum dedúxit Dóminus per vias rectas, et osténdit illi regnum Dei.\end{Resp}
\begin{Resp}\Rsign Et dedit illi claritátem ætérnam.\end{Resp}
\switchcolumn
\EN
\begin{Resp}\Rsign The Lord made him honourable, and defended him from his enemies, and kept him safe from those that lay in wait for him; \Star{} And gave him perpetual glory.\end{Resp}
\begin{Resp}\Vsign He went down with him into the pit, and left him not in bonds.\end{Resp}
\begin{Resp}\Rsign And gave him perpetual glory.\end{Resp}
\switchcolumn*

\LA
\LessonHead{Lectio V}
\switchcolumn
\EN
\LessonHead{Lesson V}
\switchcolumn*

\LA
\noindent Post exactam laudabíliter adolescéntiam cœpit ex cæléstis contemplatiónis studio solitúdinem cogitare; prius tamen vóluit Sacris initiari, ut omnem regni, a cujus successióne jure legis Sálicæ non longe distabat, spem sibi præcideret. Sacerdos factus, et prima Missa devotíssime celebrata, non multo post in eremum secessit, ubi, summa abstinéntia víctitans, cæléstium charísmatum abundántia pascebátur. Ibi cum sancto Joánne de Matha Parisiénsi doctore, a quo ex divina inspiratióne quæsítus et invéntus, per aliquot annos sanctíssime vixit; donec ambo per Angelum a Deo admóniti Romam petiérunt, specialem a summo Pontifice vivéndi regulam impetráturi. Facta ígitur Innocentio Papæ tertio inter Missárum solemnia revelatióne religiónis et institúti de rediméndis captivis, ab ipso Pontifice, simul cum socio, cándidis vestibus bicolori cruce signátis induitur, ad eam formam qua Angelus indútus appáruit. Et ínsuper vóluit Pontifex, ut nova religio, juxta tríplicem colórem quo hábitus constat, sanctíssimæ Trinitátis titulo decorarétur.\par
\switchcolumn
\EN
\noindent After a praiseworthy boyhood, he began to think of withdrawing from the world in order to be alone with heavenly thoughts. But he first wished to take orders, to the end that he might clear himself of all expectation of succeeding to the crown, to which, in consequence of the Salic Law, he was somewhat near. He became a Priest, and said his first Mass with deep devotion. Then, in a little while, he withdrew himself into the wilderness, where he lived in extreme abstinence, fed by heavenly grace. Thither, by the inspiration of God, came the holy Doctor John de la Mata of Paris, and found him, and they led an holy life together for several years, until they were both warned of an Angel to go to Rome and seek a special Rule of life from the Pope. Pope Innocent III. while he was solemnly celebrating the Liturgy (on the 28th day of January, 1198,) received in a vision the revelation of the Order and Institute for the redemption of bondsmen, and he forthwith clad Felix and John in white garments marked with a cross of red and blue, made after the likeness of the raiment wherein the Angel had appeared. This Pope also willed that the new Order should bear, as well as the habit of three colours, the name of the Most Holy Trinity.\par
\switchcolumn*

\LA
\begin{Resp}\Rsign Amávit eum Dóminus, et ornávit eum: stolam glóriæ índuit eum, \Star{} Et ad portas paradísi coronávit eum.\end{Resp}
\begin{Resp}\Vsign Índuit eum Dóminus lorícam fídei, et ornávit eum.\end{Resp}
\begin{Resp}\Rsign Et ad portas paradísi coronávit eum.\end{Resp}
\switchcolumn
\EN
\begin{Resp}\Rsign The Lord loved him and beautified him; He clothed him with a robe of glory, \Star{} And crowned him at the gates of Paradise.\end{Resp}
\begin{Resp}\Vsign The Lord hath put on him the breast-plate of faith, and hath adorned him.\end{Resp}
\begin{Resp}\Rsign And crowned him at the gates of Paradise.\end{Resp}
\switchcolumn*

\LA
\LessonHead{Lectio VI}
\switchcolumn
\EN
\LessonHead{Lesson VI}
\switchcolumn*

\LA
\noindent Regula propria, ex summi Pontificis Innocentii confirmatióne, accepta, in diœcesi Meldénsi apud locum qui Cervus Frígidus dícitur, primum ordinis paulo ante a se et socio exstructum cœnobium ampliávit, ubi religiosam observantiam et redemptiónis institutum mirifice coluit, ac inde per alumnos in alias provincias diligentíssime propagávit. Illustrem hic a beáta Vírgine Matre favórem accépit; dormiéntibus síquidem cunctis frátribus et ad matutínas preces in pervigílio Nativitátis Deiparæ media nocte recitandas, Deo sic disponente, non surgéntibus, Felix, de more vigilans et horas præveniens, chorum ingréssus, réperit beátam Vírginem in médio chori, hábitu cruce ordinis insigníto indutam, ac Cælítibus simíliter indutis sociátam. Quibus permixtus Felix, præcinénte Deipara, laudes divinas concinuit ríteque persolvit. Et, quasi jam a terrestri ad cælestem chorum evocarétur, instántis mortis ab Angelo certior factus, fílios ad caritátem erga páuperes et captivos adhortans, ánimam Deo réddidit, ætate ac meritis consummatus, anno post Christum natum ducentésimo duodecimo supra millesimum, sub eodem Pontifice Innocentio tertio.\par
\switchcolumn
\EN
\noindent When they had received the confirmation of their rule from Pope Innocent, John and Felix enlarged the first house of their Order, which they had built a little while before at Cerfroi, in the diocese of Meaux, in France. There Felix wonderfully devoted himself to the promotion of Regular Observance and of the Institute for the redemption of bondsmen, and thence he busily spread the same by sending forth his disciples into other provinces. Here it was that he received an extraordinary favour from the blessed Maiden-Mother. On the night of the Nativity of the Mother of God, the brethren lay all asleep, and by the Providence of God woke not to say Mattins. But Felix was watching, as his custom was, and came betimes into the Choir. There he found the Blessed Virgin in the midst of the Choir, clad in raiment marked with the Cross of his Order, the Cross of red and blue; and with her a company of the heavenly host in like garments. And Felix was mingled among them. And the Mother of God began to sing, and they all sang with her and praised God; and Felix sang with them; and so they finished the Office. So now that he seemed to have been already called away from glorifying God on earth, to glorify Him in heaven, an Angel told Felix that the hour of his death was at hand. When therefore he had exhorted his children to be tender to the poor and to slaves, he gave up his soul to God (upon the 4th day of November) in the year of Christ 1212, in the time of the same Pope Innocent III., being four-score-and-five years old, and full of good works.\par
\switchcolumn*

\LA
\begin{Resp}\Rsign Iste homo perfécit ómnia quæ locútus est ei Deus, et dixit ad eum: Ingrédere in réquiem meam: \Star{} Quia te vidi justum coram me ex ómnibus géntibus.\end{Resp}
\begin{Resp}\Vsign Iste est, qui contémpsit vitam mundi, et pervénit ad cæléstia regna.\end{Resp}
\begin{Resp}\Rsign Quia te vidi justum coram me ex ómnibus géntibus.\end{Resp}
\begin{Resp}\Vsign Glória Patri, et Fílio, \Star{} et Spirítui Sancto.\end{Resp}
\begin{Resp}\Rsign Quia te vidi justum coram me ex ómnibus géntibus.\end{Resp}
\switchcolumn
\EN
\begin{Resp}\Rsign This is he which did according unto all that God commanded him; and God said unto him: Enter thou into My rest, \Star{} For thee have I seen righteous before Me among all people.\end{Resp}
\begin{Resp}\Vsign This is he which loved not his life in this world, and is come unto an everlasting kingdom.\end{Resp}
\begin{Resp}\Rsign For thee have I seen righteous before Me among all people.\end{Resp}
\begin{Resp}\Vsign Glory be to the Father, and to the Son, \Star{} and to the Holy Ghost.\end{Resp}
\begin{Resp}\Rsign For thee have I seen righteous before Me among all people.\end{Resp}
\switchcolumn*

\LA
\LessonHead{Lectio VII}
\switchcolumn
\EN
\LessonHead{Lesson VII}
\switchcolumn*

\LA
\LessonTitle{Léctio sancti Evangélii secúndum Lucam}
\switchcolumn
\EN
\LessonTitle{From the Holy Gospel according to Luke}
\switchcolumn*

\LA
\Cite{Luc 12:32-34}
\switchcolumn
\EN
\Cite{Luke 12:32-35}
\switchcolumn*

\LA
\noindent In illo témpore: Dixit Jesus discípulis suis: Nolíte timére pusíllus grex, quia complácuit Patri vestro dare vobis regnum. Et réliqua.\par
\switchcolumn
\EN
\noindent At that time, Jesus said unto His disciples: Fear not, little flock, for it is your Father's good pleasure to give you the kingdom. And so on.\par
\switchcolumn*

\LA
\LessonTitle{Homilía sancti Bedæ Venerábilis Presbýteri}
\switchcolumn
\EN
\LessonTitle{Homily by the Venerable Bede, Priest at Jarrow, and Doctor of the Church.}
\switchcolumn*

\LA
\Source{Lib. 4, Cap. 54, in Luc. 12}
\switchcolumn
\EN
\Source{Bk. iv. Ch. 54 on Luke xii.}
\switchcolumn*

\LA
\noindent Pusíllum gregem electórum, vel ob comparatiónem majóris númeri reprobórum, vel pótius ob humilitátis devotiónem nóminat: quia vidélicet Ecclésiam suam quantálibet numerositáte jam dilatátam, tamen usque ad finem mundi humilitáte vult créscere, et ad promíssum regnum humilitáte perveníre. Ideóque ejus labóres blande consolátus, quam regnum Dei tantum quærére prǽcipit, eídem regnum a Patre dandum complácita benignitáte promíttit.\par
\switchcolumn
\EN
\noindent The elect are called a little flock, perchance because the reprobate are far more in number than they, but, more probably, because they love to be lowly, since it is God's will that however much His Church should grow in numbers, she should grow with lowliness even unto the end of the world, and should enter lowly into that kingdom which is hers by His promise. That kingdom He promiseth to her here, when He biddeth her to seek only the kingdom of God, and, to comfort her in her travail, He doth so sweetly and so graciously say that her Father will give it to her.\par
\switchcolumn*

\LA
\begin{Resp}\Rsign Iste est qui ante Deum magnas virtútes operátus est, et de omni corde suo laudávit Dóminum: \Star{} Ipse intercédat pro peccátis ómnium populórum.\end{Resp}
\begin{Resp}\Vsign Ecce homo sine queréla, verus Dei cultor, ábstinens se ab omni ópere malo, et pérmanens in innocéntia sua.\end{Resp}
\begin{Resp}\Rsign Ipse intercédat pro peccátis ómnium populórum.\end{Resp}
\switchcolumn
\EN
\begin{Resp}\Rsign This is he which wrought great wonders before God, and praised the Lord with all his heart. \Star{} May he pray for all people, that their sins may be forgiven unto them.\end{Resp}
\begin{Resp}\Vsign Behold a man without blame, a worshipper of God in truth, keeping himself clean from every evil work, and abiding still in his innocency.\end{Resp}
\begin{Resp}\Rsign May he pray for all people, that their sins may be forgiven unto them.\end{Resp}
\switchcolumn*

\LA
\LessonHead{Lectio VIII}
\switchcolumn
\EN
\LessonHead{Lesson VIII}
\switchcolumn*

\LA
\noindent Véndite quæ possidétis, et date eleemósynam. Nolíte, inquit, timére, ne propter regnum Dei militántibus, hujus vitæ necessária desint: quin étiam posséssa propter eleemósynam véndite. Quod tunc digne fit, quando quis semel pro Dómino suis ómnibus spretis, nihilóminus post hæc labóre mánuum, unde et victum transígere, et eleemósynam dare queat, operátur. Unde gloriátur Apóstolus, dicens: Argéntum, et aurum, aut vestem nullíus concupívi: ipsi scitis, quóniam ad ea quæ mihi opus erant, et his qui mecum sunt, ministravérunt manus istæ. Ómnia osténdi vobis, quóniam sic laborántes opórtet suscípere infírmos.\par
\switchcolumn
\EN
\noindent Sell what ye have and give alms. Fear not, He saith, lest, while ye fight for the kingdom of God, ye should lack such things as are needful for this life, nay rather, sell even that which ye have, and give alms. This doth, whosoever for the Lord's sake leaveth all that he hath, and then worketh with his hands, that so he may have to eat, and withal to give alms. In this doth the Apostle boast himself, saying I have coveted no man's silver, or gold, or apparel, as ye yourselves know for these hands have ministered unto my necessities, and to them that were with me. I have showed you all things, how that so labouring ye ought to support the weak. (Acts xx. 33, 34, 35.)\par
\switchcolumn*

\LA
\begin{Resp}\Rsign Sint lumbi vestri præcíncti, et lucérnæ ardéntes in mánibus vestris: \Star{} Et vos símiles homínibus exspectántibus dóminum suum, quando revertátur a núptiis.\end{Resp}
\begin{Resp}\Vsign Vigiláte ergo, quia nescítis qua hora Dóminus vester ventúrus sit.\end{Resp}
\begin{Resp}\Rsign Et vos símiles homínibus exspectántibus dóminum suum, quando revertátur a núptiis.\end{Resp}
\begin{Resp}\Vsign Glória Patri, et Fílio, \Star{} et Spirítui Sancto.\end{Resp}
\begin{Resp}\Rsign Et vos símiles homínibus exspectántibus dóminum suum, quando revertátur a núptiis.\end{Resp}
\switchcolumn
\EN
\begin{Resp}\Rsign Let your loins be girded about, and your lights burning; \Star{} And ye yourselves like unto men that wait for their lord, when he will return from the wedding.\end{Resp}
\begin{Resp}\Vsign Watch therefore, for ye know not what hour your Lord doth come.\end{Resp}
\begin{Resp}\Rsign And ye yourselves like unto men that wait for their lord, when he will return from the wedding.\end{Resp}
\begin{Resp}\Vsign Glory be to the Father, and to the Son, \Star{} and to the Holy Ghost.\end{Resp}
\begin{Resp}\Rsign And ye yourselves like unto men that wait for their lord, when he will return from the wedding.\end{Resp}
\switchcolumn*

\LA
\LessonHead{Lectio IX}
\switchcolumn
\EN
\LessonHead{Lesson IX}
\switchcolumn*

\LA
\noindent Fácite vobis sácculos qui non veteráscunt: eleemósynas vidélicet operándo quarum merces in ætérnum máneat. Ubi non hoc præcéptum esse putándum est, ut nil pecúniæ reservétur a sanctis, vel suis scílicet, vel páuperum úsibus suggeréndæ: cum et ipse Dóminus, cui ministrábant Ángeli, tamen ad informándam Ecclésiam suam lóculos habuísse legátur, et a fidélibus obláta consérvans, et suórum necessitátibus aliísque indigéntibus tríbuens; sed ne Deo propter ista serviátur, et ob inópiæ timórem justítia deserátur.\par
\switchcolumn
\EN
\noindent Provide yourselves bags which wax not old that is to say, by almsgiving, the reward thereof remaineth for ever. Nevertheless, we must not think here that this commandment forbiddeth the Saints to keep money for their own use, and for helping of the poor. The Lord Himself, to Whom Angels ministered, had a bag, and kept therein that which the faithful people gave unto Him (John xii. 6,) to relieve therewith the need of His disciples, and other poor folk. But we are commanded not to serve God for gain, nor to work unrighteousness for fear of poverty.\par
\switchcolumn*

\LA
\TeDeum{Te Deum}
\switchcolumn
\EN
\TeDeum{Te Deum}
\switchcolumn*

\end{paracol}

\par\needspace{10\baselineskip}\addvspace{1.2em}{\centering\small\textsc{Die 21 Novembris} · \textsc{21 November}\par}
\Office{In Præsentatione Beatæ Mariæ Virginis}{In Præsentatione Beatæ Mariæ Virginis}{Duplex majus}
\addcontentsline{toc}{subsection}{Die 21 Novembris · In Præsentatione Beatæ Mariæ Virginis \textit{— In Præsentatione Beatæ Mariæ Virginis}}
\begin{paracol}{2}
\LA
\RefLine{Lectiones I–III: ex In Festis Beatae Mariae Virginis (C11).}
\switchcolumn
\EN
\RefLine{Lessons I–III: from the In Festis Beatae Mariae Virginis (C11).}
\switchcolumn*

\LA
\RefLine{Lectiones VII–IX: ex In Festis Beatae Mariae Virginis (C11).}
\switchcolumn
\EN
\RefLine{Lessons VII–IX: from the In Festis Beatae Mariae Virginis (C11).}
\switchcolumn*

\LA
\LessonHead{Lectio IV}
\switchcolumn
\EN
\LessonHead{Lesson IV}
\switchcolumn*

\LA
\LessonTitle{Ex libro sancti Joánnis Damascéni de fide orthodóxa}
\switchcolumn
\EN
\LessonTitle{From the Book Upon the Orthodox Faith, written by St. John of Damascus}
\switchcolumn*

\LA
\Source{Liber 4, Cap. 15}
\switchcolumn
\EN
\Source{Bk. iv, chap. 15}
\switchcolumn*

\LA
\noindent Jóachim lectíssimam illam ac summis láudibus dignam mulíerem Annam matrimónio sibi copulávit. Verum, quemádmodum prisca illa Anna, cum sterilitátis morbo laboráret, per oratiónem ac promissiónem, Samuélem procreávit; eódem modo hæc étiam, per obsecratiónem et promissiónem, Dei Genetrícem a Deo accépit, ut ne hic quoque cuíquam ex illústribus matrónis céderet. Itaque grátia (nam hoc sonat Annæ vocábulum) Dóminam parit (id enim Maríæ nómine significátur). Vere étenim rerum ómnium conditárum Dómina facta est, cum Creatóris Mater éxstitit. In lucem autem éditur in domo probáticæ Jóachim, atque ad templum addúcitur. Ac deínde, in domo Dei plantáta atque per Spíritum sagináta, instar olívæ frugíferæ virtútum ómnium domicílium effícitur; ut quæ vidélicet ab omni hujúsce vitæ et carnis concupiscéntia mentem abstraxísset, atque ita vírginem una cum córpore ánimam conservásset, ut eam decébat, quæ Deum sinu suo exceptúra erat.\par
\switchcolumn
\EN
\noindent Joachim took to wife that most eminent and praiseworthy woman, Anne. And even as the ancient Hannah, being stricken with barrenness, by prayer and promise became the mother of Samuel, so likewise this woman also through prayer and promise received from God the Mother of God, that in fruitfulness she might not be behind any of the famous matrons. And thus grace (for such is the signification of the name of Anne) is mother of the Lady (for such is the signification of the name of Mary.) And indeed she became the Lady of every creature, since she hath been mother of the Creator. She first saw the light in Joachim's house, hard by the Pool of Bethesda, at Jerusalem, and was carried to the Temple. There planted in the Lord, the dew of His Spirit made her to flourish in the courts of her God, and like a green olive she became a tree, so that all the doves of grace came and lodged in her branches. And so she raised her mind utterly above the lust of life and the lust of the flesh, and kept her soul virgin in her virgin body, as became her that was to receive God into her womb.\par
\switchcolumn*

\LA
\begin{Resp}\Rsign Sicut cedrus exaltáta sum in Líbano, et sicut cypréssus in monte Sion: quasi myrrha elécta, \Star{} Dedi suavitátem odóris.\end{Resp}
\begin{Resp}\Vsign Et sicut cinnamómum et bálsamum aromatízans.\end{Resp}
\begin{Resp}\Rsign Dedi suavitátem odóris.\end{Resp}
\switchcolumn
\EN
\begin{Resp}\Rsign I was exalted like a cedar in Lebanon, and as a cypress-tree upon Mount Zion. Like the best myrrh \Star{} I yielded a pleasant odour.\end{Resp}
\begin{Resp}\Vsign Like cinnamon and sweet balsam.\end{Resp}
\begin{Resp}\Rsign I yielded a pleasant odour.\end{Resp}
\switchcolumn*

\LA
\LessonHead{Lectio V}
\switchcolumn
\EN
\LessonHead{Lesson V}
\switchcolumn*

\LA
\LessonTitle{Ex libro sancti Ambrósii Epíscopi de Virgínibus}
\switchcolumn
\EN
\LessonTitle{From the Book Upon Virgins, written by St. Ambrose, Bishop (of Milan.)}
\switchcolumn*

\LA
\Source{Liber 2, post initium.}
\switchcolumn
\EN
\Source{ii}
\switchcolumn*

\LA
\noindent Talis fuit Maria, ut ejus uníus vita ómnium sit disciplína. Si ígitur auctor non dísplicet, opus probémus; ut, quæcúmque sibi ejus exoptat præmium, imitétur exemplum. Quantæ in una vírgine spécies virtútum émicant! Secretum verecúndiæ, vexíllum fidei, devotiónis obsequium; virgo intra domum, comes ad ministérium, mater ad templum. O quantis illa virgínibus occurret! quantas complexa, ad Dóminum trahet, dicens: Hæc torum fílii mei, hæc thálamos nuptiales immaculáto servávit pudóre!\par
\switchcolumn
\EN
\noindent Such was Mary that her single life offereth an example to all. If then the doer displease us not, let us applaud the deed; if any other woman seek like reward, let her follow after like works. In the one Virgin how many glorious examples do shine forth. Her's was the hidden treasure of modesty, her's the high standard of faith, her's the self-sacrifice of earnestness, her's to be the pattern of maidenhood at home, of kinswomanhood in ministry, of motherhood in the Temple. O to how many virgins hath she been helpful, how many hath she taken in her arms and presented unto the Lord, saying Here is one who, (like me,) hath kept stainlessly clean the wedding -chamber, the marriage -bed of my Son!\par
\switchcolumn*

\LA
\begin{Resp}\Rsign Quæ est ista quæ procéssit sicut sol, et formósa tamquam Jerúsalem? \Star{} Vidérunt eam fíliæ Sion, et beátam dixérunt, et regínæ laudavérunt eam.\end{Resp}
\begin{Resp}\Vsign Et sicut dies verni circúmdabant eam flores rosárum et lília convállium.\end{Resp}
\begin{Resp}\Rsign Vidérunt eam fíliæ Sion, et beátam dixérunt, et regínæ laudavérunt eam.\end{Resp}
\switchcolumn
\EN
\begin{Resp}\Rsign Who is this that cometh up like the sun? This, comely as Jerusalem? \Star{} The daughters of Zion saw her, and called her blessed: the queens also, and they praised her.\end{Resp}
\begin{Resp}\Vsign And about her it was as the flower of roses in the spring of the year, and lilies of the valleys.\end{Resp}
\begin{Resp}\Rsign The daughters of Zion saw her, and called her blessed the queens also, and they praised her.\end{Resp}
\switchcolumn*

\LA
\LessonHead{Lectio VI}
\switchcolumn
\EN
\LessonHead{Lesson VI}
\switchcolumn*

\LA
\noindent Quid ergo éxsequar cibórum parsimóniam, officiórum redundantiam: álterum ultra natúram superfuísse, álterum pene ipsi natúræ defuísse? Illic nulla intermissa témpora, hic congeminatos jejúnio dies. Et, si quando reficiéndi successísset volúntas, cibus plerumque obvius, qui mortem arcéret, non delicias ministraret. Dormire non prius cupiditas quam necessitas fuit; et tamen, cum quiésceret corpus, vigilaret animus, qui frequenter in somnis aut lecta répetit, aut somno interrupta continuat, aut dispósita gerit, aut gerenda prænuntiat.\par
\switchcolumn
\EN
\noindent Why should I go on to speak of the scantiness of her eating, or of the multiplicity of her work? how her labour seemed above human capacity, and her refreshment insufficient for human strength, her toil never missing a moment, her fasting taking two days together. And when she was fain to eat, she took not dainties, but whatsoever food came first to hand that would keep body and soul together. She would not sleep till need was, and even then, while her body rested, her soul watched, for she often talked in her sleep, either repeating things that she had read, or going on with what she was doing before sleep interrupted her, or rehearsing things executed, or talking of things projected.\par
\switchcolumn*

\LA
\begin{Resp}\Rsign Ornátam monílibus fíliam Jerúsalem Dóminus concupívit: \Star{} Et vidéntes eam fíliæ Sion, beatíssimam prædicavérunt, dicéntes: * Unguéntum effúsum nomen tuum.\end{Resp}
\begin{Resp}\Vsign Astitit regína a dextris tuis in vestítu deauráto, circúmdata varietáte.\end{Resp}
\begin{Resp}\Rsign Et vidéntes eam fíliæ Sion, beatíssimam prædicavérunt, dicéntes:\end{Resp}
\begin{Resp}\Vsign Glória Patri, et Fílio, \Star{} et Spirítui Sancto.\end{Resp}
\begin{Resp}\Rsign Unguéntum effúsum nomen tuum.\end{Resp}
\switchcolumn
\EN
\begin{Resp}\Rsign When the Lord beheld the daughter of Jerusalem adorned with her jewels, He greatly desired her beauty \Star{} And when the daughters of Zion saw her, they cried out that she was most blessed, saying thy name is as ointment poured forth.\end{Resp}
\begin{Resp}\Vsign Upon thy right hand did stand the Queen in a vesture of gold wrought about with diverse colours.\end{Resp}
\begin{Resp}\Rsign And when the daughters of Zion saw her, they cried out that she was most blessed, saying:\end{Resp}
\begin{Resp}\Vsign Glory be to the Father, and to the Son, \Star{} and to the Holy Ghost.\end{Resp}
\begin{Resp}\Rsign Thy name is as ointment poured forth.\end{Resp}
\switchcolumn*

\end{paracol}

\par\needspace{10\baselineskip}\addvspace{1.2em}{\centering\small\textsc{Die 22 Novembris} · \textsc{22 November}\par}
\Office{S. Cæciliæ Virginis et Martyris}{St. Cecilia, Virgin and Martyr}{Duplex}
\addcontentsline{toc}{subsection}{Die 22 Novembris · S. Cæciliæ Virginis et Martyris \textit{— St. Cecilia, Virgin and Martyr}}
\begin{paracol}{2}
\LA
\RefLine{Lectiones I–III: ex Commúni Unius Virginis et Martyris (C6).}
\switchcolumn
\EN
\RefLine{Lessons I–III: from the Common of a Virgin Martyr (C6).}
\switchcolumn*

\LA
\LessonHead{Lectio IV}
\switchcolumn
\EN
\LessonHead{Lesson IV}
\switchcolumn*

\LA
\noindent Cæcília, Virgo Romana, nobili genere nata, a prima ætate christianæ fidei præceptis instituta, virginitátem suam Deo vovit. Sed, cum póstea, contra suam voluntátem, data esset in matrimónium Valeriáno, prima nuptiárum nocte hunc cum eo sermónem hábuit: Ego, Valeriane, in Angeli tutela sum, qui virginitátem meam custódit; quare ne quid in me committas, quo ira Dei in te concitétur. Quibus verbis commótus Valerianus, illam attingere non est ausus; quin étiam addidit se in Christum crediturum, si eum Angelum vidéret. Cui Cæcília cum sine baptismo negaret id fíeri posse, incénsus cupiditate vidéndi Angelum, se baptizari velle respóndit. Quare hortatu Vírginis ad Urbanum Papam, qui propter persecutiónem in Mártyrum sepúlcris via Appia latebat, veniens, ab eo baptizátur.\par
\switchcolumn
\EN
\noindent Cecilia was a Roman maiden of noble birth, trained up from her earliest years in the teaching of the Christian faith, and who by vow consecrated her virginity to God. She was afterwards given in marriage, against her will, to Valerian. On the first night she said to him Valerian! I am under the wardship of an Angel, who keepeth me always a maiden. Therefore do nothing unto me, lest the anger of God should be aroused against thee. Valerian was moved at her words, and dared not to touch her. Also he added even this, that he would believe in Christ, if he could see the Angel. Cecilia answered him that that could not be unless he were first baptized, and for the sake of seeing the Angel he was willing. So she bade him go unto Pope Urban, who was hiding in the sepulchre of the Martyrs on the Appian Way on account of the persecution. And he went unto him and was baptized.\par
\switchcolumn*

\LA
\begin{Resp}\Rsign Cilicio Cæcília membra domabat, Deum gemítibus exorabat, \Star{} Tiburtium et Valerianum ad corónas vocábat.\end{Resp}
\begin{Resp}\Vsign Hæc est Virgo sapiens, et una de número prudéntum.\end{Resp}
\begin{Resp}\Rsign Tiburtium et Valerianum ad corónas vocabat.\end{Resp}
\switchcolumn
\EN
\begin{Resp}\Rsign Cecilia brought her body under with haircloth, and besought God with loud crying; \Star{} And called Tiburtius and Valerian to crowns.\end{Resp}
\begin{Resp}\Vsign This is one of the wise virgins, one chosen out of the number of the prudent.\end{Resp}
\begin{Resp}\Rsign And called Tiburtius and Valerian to crowns.\end{Resp}
\switchcolumn*

\LA
\LessonHead{Lectio V}
\switchcolumn
\EN
\LessonHead{Lesson V}
\switchcolumn*

\LA
\noindent Inde ad Cæciliam reversus, orántem et cum ea Angelum divino splendore fulgentem invénit. Quo aspectu obstupefactus, ut primum ex timóre confirmátus est, Tiburtium fratrem suum accersit; qui, a Cæcília Christi fide imbutus et ab eodem Urbano baptizatus, ipse étiam ejusdem Angeli, quem frater ejus viderat, aspectu dignátus est. Uterque autem paulo post, Almachio præfecto, constanter martyrium subiit. Qui mox Cæciliam comprehéndi imperat, ab eáque primum, ubi Tiburtii et Valeriáni facultates sint, exquirit.\par
\switchcolumn
\EN
\noindent Whence he came back to Cecilia, and found her praying, and the Angel with her, shining from the glory of God. As soon as he had recovered from the shock of wonder and fear, he brought his brother Tiburtius, and Cecilia taught him Christ, and he was baptized by the same Pope Urban, and he also was vouchsafed to see the Angel whom his brother had seen. A little while after, both of them bravely suffered martyrdom under the Praefect Almachius, who then caused Cecilia to be taken, and asked of her, first of all, where was the property of Tiburtius and Valerian?\par
\switchcolumn*

\LA
\begin{Resp}\Rsign Cæcíliam intra cubiculum orántem invénit, et juxta eam stantem Angelum Dómini: \Star{} Quem videns Valerianus, nimio terróre correptus est.\end{Resp}
\begin{Resp}\Vsign Angelus Dómini descéndit de cælo, et lumen refulsit in habitáculo.\end{Resp}
\begin{Resp}\Rsign Quem videns Valerianus, nimio terróre correptus est.\end{Resp}
\switchcolumn
\EN
\begin{Resp}\Rsign He found Cecilia praying in her chamber, and standing by her the Angel of the Lord. \Star{} And when Valerian saw him, he feared with a great fear.\end{Resp}
\begin{Resp}\Vsign The Angel of the Lord descended from heaven, and a light shone in all the house.\end{Resp}
\begin{Resp}\Rsign And when Valerian saw him, he feared with a great fear.\end{Resp}
\switchcolumn*

\LA
\LessonHead{Lectio VI}
\switchcolumn
\EN
\LessonHead{Lesson VI}
\switchcolumn*

\LA
\noindent Cui, cum Virgo ómnia illórum paupéribus distributa esse respondísset, eo furóre concitátus est, ut eam, in ipsíus ædes reductam, in balneo comburi jusserit. Quo in loco cum diem noctemque ita fuísset ut ne flamma quidem illam attingeret, eo immissus est carnifex, qui ter secúri ictam, cum caput abscindere non potuísset, semivivam relíquit. Illa triduo post, sextodecimo Kalendas Octobris, Alexandro imperatóre, duplici virginitátis et martyrii palma decorata, evolávit in cælum. Cujus corpus ab ipso Urbano Papa in Callísti cœmeterio sepúltum est, in ejus ædibus ecclésia ipsíus Cæcíliæ consecrata. Ejus et Urbani ac Lucii Pontificum, Tiburtii, Valeriáni et Maximi corpora, a Paschali primo Pontifice inde transláta in Urbem, in eádem sanctæ Cæcíliæ ecclésia cóndita sunt.\par
\switchcolumn
\EN
\noindent To him the Virgin answered that all their goods had been given to the poor. Thereupon he was filled with fury, and commanded her to be taken home, and burnt in the bath. She was in that place a day and a night, but the fire had not harmed her. Then was sent the executioner, who gave her three strokes of the axe, and, as he could not cut off her head, left her half-dead. Three days thereafter, upon the 22nd day of November, in the reign of the Emperor Alexander Severus, she winged her flight for heaven, glorified with the two palms of virginity and martyrdom. Her body was buried in the cemetery of Callistus by the aforenamed Pope Urban, who also consecrated a Church in her name in her own house. Her relics were brought into the city by Pope Paschal I, along with those of Tiburtius, Valerian, and Maximus, and all laid together in the said Church of St Cecilia.\par
\switchcolumn*

\LA
\begin{Resp}\Rsign Dómine Jesu Christe, pastor bone, seminator casti consílii, suscipe seminum fructus quos in Cæcília seminásti: \Star{} Cæcília famula tua quasi apis tibi argumentosa deservit.\end{Resp}
\begin{Resp}\Vsign Nam sponsum, quem quasi leónem ferocem accepit, ad te quasi agnum mansuetíssimum destinávit.\end{Resp}
\begin{Resp}\Rsign Cæcília famula tua quasi apis tibi argumentosa deservit.\end{Resp}
\begin{Resp}\Vsign Glória Patri, et Fílio, \Star{} et Spirítui Sancto.\end{Resp}
\begin{Resp}\Rsign Cæcília famula tua quasi apis tibi argumentosa deservit.\end{Resp}
\switchcolumn
\EN
\begin{Resp}\Rsign O Lord Jesus Christ, the Good Shepherd, Who hast said of abstaining from marriage: He that is able to receive it, let him receive it, accept now the fruits from this seed, which Thou didst sow in the heart of Cecilia. \Star{} Busy like a bee, thine handmaiden Cecilia served thee.\end{Resp}
\begin{Resp}\Vsign For her husband, who was like a raging lion when she took him, she sent unto thee meek as the meekest of lambs.\end{Resp}
\begin{Resp}\Rsign Busy like a bee, thine handmaiden Cecilia served thee.\end{Resp}
\begin{Resp}\Vsign Glory be to the Father, and to the Son, \Star{} and to the Holy Ghost.\end{Resp}
\begin{Resp}\Rsign Busy like a bee, thine handmaiden Cecilia served thee.\end{Resp}
\switchcolumn*

\LA
\LessonHead{Lectio VII}
\switchcolumn
\EN
\LessonHead{Lesson VII}
\switchcolumn*

\LA
\LessonTitle{Léctio sancti Evangélii secúndum Matthǽum}
\switchcolumn
\EN
\LessonTitle{From the Holy Gospel according to Matthew}
\switchcolumn*

\LA
\Cite{Matt 25:1-13}
\switchcolumn
\EN
\Cite{Matt 25:1-13}
\switchcolumn*

\LA
\noindent In illo témpore: Dixit Jesus discípulis suis parábolam hanc: Simile erit regnum cælórum decem virgínibus, quæ, accipiéntes lampades suas, exiérunt obviam sponso et sponsæ. Et réliqua.\par
\switchcolumn
\EN
\noindent In that time, Jesus said to his disciples: Then shall the kingdom of heaven be like to ten virgins, who taking their lamps went out to meet the bridegroom and the bride. And so on.\par
\switchcolumn*

\LA
\LessonTitle{Homilía sancti Joánnis Chrysóstomi}
\switchcolumn
\EN
\LessonTitle{Homily by St. John Chrysostom, Patriarch (of Constantinople.)}
\switchcolumn*

\LA
\Source{Homilia 79 in Matth., post initium}
\switchcolumn
\EN
\Source{Hom 79 on Matth.}
\switchcolumn*

\LA
\noindent Quam ob causam in persona vírginum hanc parábolam profert, nec qualemcúmque personam sine discrimine súbjicit? Magna quædam de virginitate disserúerat dicens: Sunt eunuchi, qui seipsos castravérunt propter regnum cælórum; et, Qui potest cápere, capiat. Nec ignorabat, de virginitate magnam esse vulgo existimatiónem, quippe cum sit ea res natúra sublimis: quod inde patet, quia neque in veteri Testaménto a priscis illis sanctisque viris culta fuit, et in novo nulla legis necessitate jubétur. Non enim id imperávit, sed fidelium voluntáti permisit. Unde et Paulus ait: De virgínibus autem præcéptum Dómini non habeo; et laúdo quidem eum, qui hoc sectátur institutum, nolentem autem non cogo, neque eam rem præcéptum facio.\par
\switchcolumn
\EN
\noindent Wherefore doth the Lord set forth this parable under the figure of virgins, and not make it of acceptation for all men? He had spoken great things touching virginity, saying There be eunuchs, which have made themselves eunuchs for the kingdom of heaven's sake. He that is able to receive it, let him receive it. \textasciitilde{}(Matth. xix. 12.) He knew also that virginity is a thing which is held in great honour among men, being indeed a thing higher than nature, as is plain from this, that under the Old Testament even the Patriarchs and Saints did not practice it, and that under the New Testament it is not enjoined by any commandment of necessity; for the Lord did not make it binding, but left it open to the free choice of the faithful. Whence also Paul saith Concerning virgins I have no commandment of the Lord; yet I give my judgment, as one that hath obtained mercy of the Lord to be faithful. I suppose therefore that this is good for the present distress, that it is good for a man so to be. But and if thou marry, thou hast not sinned, and if a virgin marry, she hath not sinned. (1 Cor. vii. 25, 26, 28.)\par
\switchcolumn*

\LA
\begin{Resp}\Rsign Beáta Cæcília dixit Tiburtio: Hodie te fateor meum esse cognátum, quia amor Dei te fecit esse \Star{} Contemptórem idolórum.\end{Resp}
\begin{Resp}\Vsign Sicut enim amor Dei mihi tuum fratrem conjugem fecit, ita te mihi cognátum fecit esse.\end{Resp}
\begin{Resp}\Rsign Contemptórem idolórum.\end{Resp}
\switchcolumn
\EN
\begin{Resp}\Rsign Blessed Cecilia said unto Tiburtius: Today I call thee my brother, for the love of God hath made thee \Star{} To cast away idols.\end{Resp}
\begin{Resp}\Vsign For even as the love of God hath made thy brother to be my husband, so the same hath made thee to be my brother, (and)\end{Resp}
\begin{Resp}\Rsign To cast away idols.\end{Resp}
\switchcolumn*

\LA
\LessonHead{Lectio VIII}
\switchcolumn
\EN
\LessonHead{Lesson VIII}
\switchcolumn*

\LA
\noindent Quóniam ígitur et magna res erat, ac de ea apud multos magna erat existimátio, ne quis, ea perfécta, se totum perfécisse putaret ac cetera negligeret, hanc parábolam pósuit; ut osténderet, virginitátem, quamvis cetera habeat, si misericórdiæ bonis carúerit, cum fornicatóribus éjici. Ac mérito quidem inhumanum ac misericórdia carentem cum illis cóllocat; fornicator enim córporum, istæ vero pecuniárum cupiditate vincúntur. Non est autem córporum et pecuniæ par cupiditas; sed acrior multo atque vehementior illa córporum est. Quanto ígitur cum imbecilliore luctántur, tanto minus venia dignæ sunt, si vincántur. Idcirco étiam fatuas appellávit; quóniam, majori certamine superato, in facilióre totum perdidérunt.\par
\switchcolumn
\EN
\noindent Virginity then, being a thing in itself so great and so much esteemed among many, lest any man having attained unto it, and kept it undefiled, should think that he hath done all, and so leave the rest undone, the Lord putteth forth this parable, in order to show that if virginity, though it have all else, lack mercy, its owner will but have his portion without among the fornicators, among whom Christ doth justly place the heartless and pitiless celibate. The fornicator is entangled in lust after bodies, the other in lust after money. The lust for bodies and the lust for money are two very different things, whereof the fleshly is by far the keener and the stubborner appetite. They that strive with the weaker enemy are therefore much less excusable if they fall. Wherefore the Lord hath called such virgins foolish, for having first won the stern battle, and then been destroyed in the light one.\par
\switchcolumn*

\LA
\begin{Resp}\Rsign Cæcília me misit ad vos, ut ostendatis mihi sanctum antístitem; \Star{} Quia ad ipsum habeo secreta quæ pérferam.\end{Resp}
\begin{Resp}\Vsign Tunc Valerianus perréxit, et, signo quod acceperat, invénit sanctum Urbanum.\end{Resp}
\begin{Resp}\Rsign Quia ad ipsum habeo secreta quæ pérferam.\end{Resp}
\begin{Resp}\Vsign Glória Patri, et Fílio, \Star{} et Spirítui Sancto.\end{Resp}
\begin{Resp}\Rsign Quia ad ipsum habeo secreta quæ pérferam.\end{Resp}
\switchcolumn
\EN
\begin{Resp}\Rsign Cecilia hath sent me unto you, that ye may show me the holy Bishop; \Star{} For unto him I have a secret message to deliver.\end{Resp}
\begin{Resp}\Vsign Then Valerian went his way, and found the holy Urban by the sign which had been given him.\end{Resp}
\begin{Resp}\Rsign For unto him I have a secret message to deliver.\end{Resp}
\begin{Resp}\Vsign Glory be to the Father, and to the Son, \Star{} and to the Holy Ghost.\end{Resp}
\begin{Resp}\Rsign For unto him I have a secret message to deliver.\end{Resp}
\switchcolumn*

\LA
\LessonHead{Lectio IX}
\switchcolumn
\EN
\LessonHead{Lesson IX}
\switchcolumn*

\LA
\noindent Lámpades autem hoc loco illud ipsum virginitátis domum appellat et sanctimoniæ puritátem; óleum vero benignitátem, eleemosynam, impensum indigéntibus auxílium. Tardante autem sponso dormitavérunt omnes, et dormiérunt. Non parvum témporis rursus spatium interjectum osténdit, ut discipulos, regnum ipsíus mox futurum exspectántes, ab ea opinióne dedúceret; id enim illi sperábant: quapropter crebro ab hujusmodi eos spe révocat. Ad hoc, illud quoque índicat, somnum quemdam esse mortem. Dormiérunt, inquit; media autem nocte clamor factus est: vel, ut in eádem parábola persistat, vel rursus osténdit in nocte futuram esse resurrectiónem. Clamórem étiam Paulus commemorat, dicens: In jussu, in voce Archangeli, in novíssima tuba descéndet de cælo.\par
\switchcolumn
\EN
\noindent Why the lamps spoken of in this parable, the Lord signifieth the actual gift of virginity and holy continency, and by the oil gentleness, almsgiving, and helpfulness toward the needy. While the Bridegroom tarried, they all slumbered and slept. His disciples hoped that His kingdom was to come forthwith. To call them away from this hope, to lead them away from this thought, He showeth them the time of waiting for the Bridegroom to be no very short one. They all slumbered and slept. He calleth death a sleep. And at midnight there was a cry made, Behold, the Bridegroom cometh, go ye out to meet Him. This at midnight is either a continuation of the parable (and so signifieth the awaking of the dead,) or else meaneth that the again-rising to come will actually take place in the night. Of the cry Paul also maketh mention, where he saith The Lord Himself shall descend from heaven with a shout, with the voice of the Archangel, and with the trump of God. (1 Thess. iv. 16.)\par
\switchcolumn*

\LA
\TeDeum{Te Deum}
\switchcolumn
\EN
\TeDeum{Te Deum}
\switchcolumn*

\end{paracol}

\par\needspace{10\baselineskip}\addvspace{1.2em}{\centering\small\textsc{Die 23 Novembris} · \textsc{23 November}\par}
\Office{S. Clementis Papæ et Martyris}{St. Clement, Pope and Martyr}{Duplex}
\addcontentsline{toc}{subsection}{Die 23 Novembris · S. Clementis Papæ et Martyris \textit{— St. Clement, Pope and Martyr}}
\begin{paracol}{2}
\LA
\RefLine{Lectiones I–III: de Scriptura occurrente.}
\switchcolumn
\EN
\RefLine{Lessons I–III: from the Scripture of the day.}
\switchcolumn*

\LA
\RefLine{Lectiones VII–VIII: ex Commúni unius Summi Pontificis et Martyris (C2b).}
\switchcolumn
\EN
\RefLine{Lessons VII–VIII: from the Common of a Single Martyr Pope (C2b).}
\switchcolumn*

\LA
\RefLine{Lectio IX: de S. Felicitatis Martyris (11-23cc).}
\switchcolumn
\EN
\RefLine{Lesson IX: of St. Felicity, Martyr (11-23cc).}
\switchcolumn*

\LA
\LessonHead{Lectio IV}
\switchcolumn
\EN
\LessonHead{Lesson IV}
\switchcolumn*

\LA
\noindent Clemens, Románus, Faustíni fílius, de regióne Cǽlii montis, discípulus beáti Petri, cujus méminit Paulus scribens ad Philippénses: Etiam rogo et te, germáne compar, ádjuva illas quæ mecum laboravérunt in Evangélio, cum Cleménte et céteris adjutóribus meis, quorum nómina sunt in libro vitæ. Hic septem Urbis regiónes divísit septem notáriis, síngulas síngulis attríbuens, qui passiónes Mártyrum et res ab eis gestas, diligentíssime conquisítas, lítteris mandárent. Multa scripsit et ipse accuráte et salutáriter, quibus christiánam religiónem illustrávit.\par
\switchcolumn
\EN
\noindent Clement, the son of Faustinus, was a Roman, from the quarter of the Coelian Mount. He was a disciple of the blessed Peter, and is the same concerning whom Paul saith, writing to the Philippians And I entreat thee also, true yokefellow, help those women which laboured with me in the Gospel, with Clement also, and with other my fellow-labourers, whose names are written in the book of life. \textasciitilde{}(iv. 3.) (He succeeded Cletus as Bishop of Rome.) He it was who divided the seven quarters of the city among seven scribes, one to each, whose duty it was to search out most carefully, and record in writing the sufferings and acts of the Martyrs. He himself also wrote much, and that most orthodox and healthy, whereby he clearly explained the Christian Religion.\par
\switchcolumn*

\LA
\begin{Resp}\Rsign Oránte sancto Cleménte, appáruit ei Agnus Dei, \Star{} De sub cujus pede fons vivus emánat: flúminis ímpetus lætíficat civitátem Dei.\end{Resp}
\begin{Resp}\Vsign Vidi supra montem Agnum stantem.\end{Resp}
\begin{Resp}\Rsign De sub cujus pede fons vivus emánat: flúminis ímpetus lætíficat civitátem Dei.\end{Resp}
\switchcolumn
\EN
\begin{Resp}\Rsign While Holy Clement was at prayer, there appeared unto him the Lamb of God; \Star{} With the river of the water of life proceeding from under His Feet, even that river, the streams whereof make glad the city of God.\end{Resp}
\begin{Resp}\Vsign And I looked, and, lo, a Lamb stood on the mount.\end{Resp}
\begin{Resp}\Rsign With the river of the water of life proceeding from under His Feet, even that river, the streams whereof make glad the city of God.\end{Resp}
\switchcolumn*

\LA
\LessonHead{Lectio V}
\switchcolumn
\EN
\LessonHead{Lesson V}
\switchcolumn*

\LA
\noindent Cum autem doctrína ac vitæ sanctitáte multos ad Christi fidem convérteret, a Trajáno imperatóre relegátus est trans mare Pónticum in solitúdine urbis Chersónæ; in qua duo míllia Christianórum réperit, qui ab eódem Trajáno condemnáti fúerant. Qui cum in eruéndis et secándis marmóribus aquæ penúria laborárent, Clemens, facta oratióne, in vicínum collem ascéndit, in cujus jugo vidit Agnum déxtero pede fontem aquæ dulcis, qui inde scatébat, attingéntem; ubi omnes sitim explevérunt. Eóque miráculo multi infidéles, ad Christi fidem convérsi, Cleméntis étiam sanctitátem venerári cœpérunt.\par
\switchcolumn
\EN
\noindent His teaching and the holiness of his life brought many to believe in Christ, and he was therefore exiled by the Emperor Trajan to Kherson, in the Crimea, where he found two thousand Christians, who had been condemned by the same Trajan. There they all worked in the marble quarries. During their labour they suffered for want of water, and Clement prayed, and then went up an hill hard by, on the top whereof he saw a Lamb standing, touching with its right foot a flowing spring of sweet waters. Therewith they all quenched their thirst, and by this miracle many unbelievers were brought to believe in Christ, and began to honour the holiness of Clement.\par
\switchcolumn*

\LA
\begin{Resp}\Rsign Omnes una voce dixérunt: Ora pro nobis, sancte Clemens; \Star{} Ut digni efficiámur promissiónibus Christi.\end{Resp}
\begin{Resp}\Vsign Non meis méritis ad vos me misit Dóminus vestris corónis partícipem fíeri.\end{Resp}
\begin{Resp}\Rsign Ut digni efficiámur promissiónibus Christi.\end{Resp}
\switchcolumn
\EN
\begin{Resp}\Rsign They all said with one voice Holy Clement, pray for us; \Star{} That we may be made worthy of the promises of Christ.\end{Resp}
\begin{Resp}\Vsign For no worthiness of mine own hath the Lord sent me unto you, to become a partaker in your crowns.\end{Resp}
\begin{Resp}\Rsign That we may be made worthy of the promises of Christ.\end{Resp}
\switchcolumn*

\LA
\LessonHead{Lectio VI}
\switchcolumn
\EN
\LessonHead{Lesson VI}
\switchcolumn*

\LA
\noindent Quibus concitátus Trajánus, misit illuc qui Cleméntem, alligáta ad ejus collum ánchora, in profúndum dejícerent. Quod cum factum esset, Christiánis ad littus orántibus, mare ad tria milliária recéssit; eóque illi accedéntes, ædículam marmóream in templi formam et intus arcam lapídeam, ubi Mártyris corpus cónditum est. Ecclésia étiam in eo ínsulæ loco unde divínitus fons manárat, ejúsdem nómine dedicáta est. Vixit in pontificátu annos novem, menses sex, dies sex. Fecit ordinatiónes duas mense Decémbri, quibus creávit presbýteros decem, diáconos duos, epíscopos per divérsa loca quíndecim.\par
\switchcolumn
\EN
\noindent These things moved Trajan to send a messenger to the Crimea, who tied an anchor about Clement's neck, and cast him into the deep of the sea. After it had been done, while the Christians were praying on the shore, the sea went back three miles, and when they followed it, they found a grotto of marble, in form like a temple, and therein a stone coffin wherein was laid the body of the Martyr, and, hard by, the anchor wherewith he had been sunk. Then were the country people moved to receive the faith of Christ. The body of Clement was afterwards brought to Rome, in the time of Pope Nicholas I, and buried in his own Church. A Church was also built in the Crimea, in the place where God had made the water to break forth. Clement lived as Pope nine years, six months, and six days. He held two Ordinations in the month of December, wherein he made ten Priests, two Deacons, and fifteen Bishops for diverse places.\par
\switchcolumn*

\LA
\begin{Resp}\Rsign Dedísti, Dómine, habitáculum Mártyri tuo Cleménti in mari, in modum templi marmórei angélicis mánibus præparátum: \Star{} Iter præbens pópulo terræ, ut enárrent mirabília tua.\end{Resp}
\begin{Resp}\Vsign Dedísti, Dómine, Sanctis tuis viam in mari, et in flumínibus sémitam.\end{Resp}
\begin{Resp}\Rsign Iter præbens pópulo terræ, ut enárrent mirabília tua.\end{Resp}
\begin{Resp}\Vsign Glória Patri, et Fílio, \Star{} et Spirítui Sancto.\end{Resp}
\begin{Resp}\Rsign Iter præbens pópulo terræ, ut enárrent mirabília tua.\end{Resp}
\switchcolumn
\EN
\begin{Resp}\Rsign Lord, Thou hast given unto thy Martyr Clement a tabernacle in the sea, after the fashion of a temple of marble, builded by the hands of Angels. \Star{} And Thou givest a way thither unto the people on the land, that they may tell of thy marvellous works.\end{Resp}
\begin{Resp}\Vsign Lord, Thou didst give unto thy Saints a way in the sea, and a path through the mighty waters.\end{Resp}
\begin{Resp}\Rsign And Thou gavest a way thither unto the people on the land, that they may tell of thy marvellous works.\end{Resp}
\begin{Resp}\Vsign Glory be to the Father, and to the Son, \Star{} and to the Holy Ghost.\end{Resp}
\begin{Resp}\Rsign And Thou gavest a way thither unto the people on the land, that they may tell of thy marvellous works.\end{Resp}
\switchcolumn*

\end{paracol}

\par\needspace{10\baselineskip}\addvspace{1.2em}{\centering\small\textsc{Die 23 Novembris} · \textsc{23 November}\par}
\Office{S. Felicitatis Martyris}{St. Felicity, Martyr}{Simplex}
\addcontentsline{toc}{subsection}{Die 23 Novembris · S. Felicitatis Martyris \textit{— St. Felicity, Martyr}}
\begin{paracol}{2}
\LA
\LessonHead{Lectio IX (pro commemoratione)}
\switchcolumn
\EN
\LessonHead{Lesson IX (when commemorated)}
\switchcolumn*

\LA
\LessonTitle{Commemoratio: S. Felicitatis Martyris}
\switchcolumn
\EN
\LessonTitle{Commemoration of: St. Felicity, Martyr}
\switchcolumn*

\LA
\Source{Ex Homilía 3 in Evangelica}
\switchcolumn
\EN
\Source{3rd Horn, on the Gospels.}
\switchcolumn*

\LA
\noindent Sermo sancti Gregórii Papæ\par
\switchcolumn
\EN
\noindent The Lesson is taken from the Sermons of Pope St. Gregory (the Great.)\par
\switchcolumn*

\LA
Beáta Felícitas, cujus hódie natalícia celebrámus, septem fílios sic post se tímuit vivos in carne relínquere, sicut carnáles paréntes solent metúere ne mórtuos præmíttant. In persecutiónis enim labóre deprehénsa, filiórum corda in amóre supérnæ pátriæ prædicándo roborávit; et parturívit spíritu quos carne pepérerat, ut prædicatióne páreret Deo quos carne pepérerat mundo. Numquid ergo hanc féminam Mártyrem díxerim? Sed plus quam Mártyrem; quæ, septem pignóribus ad regnum præmíssis, tot ante se mórtuos transmísit. Ad pœnam prima venit, sed pervénit octáva.\par
\switchcolumn
\EN
That blessed woman Felicity, whose Birth-feast we are keeping today, had as much dread of leaving her seven sons living after her in the flesh, as have carnal-minded mothers of seeing them go dead before them. When she was taken in the strong pains of persecution, she braced up the hearts of her children by bidding them cleave to the Fatherland above, and became their mother for the spiritual, as she had aforetime been for the fleshly life, bringing them forth for God by her exhortation, as she had brought them forth for the world by her body. And shall I not call this woman a Martyr? Nay, more than Martyr. The seven whom she trusted to God were seven children sent before her to death. She suffered first and triumphed last.\par
\switchcolumn*

\LA
\TeDeum{Te Deum}
\switchcolumn
\EN
\TeDeum{Te Deum}
\switchcolumn*

\end{paracol}

\par\needspace{10\baselineskip}\addvspace{1.2em}{\centering\small\textsc{Die 24 Novembris} · \textsc{24 November}\par}
\Office{S. Joannis a Cruce Confessoris et Ecclesiæ Doctoris}{St. John of the Cross, Confessor and Doctor of the Church}{Duplex}
\addcontentsline{toc}{subsection}{Die 24 Novembris · S. Joannis a Cruce Confessoris et Ecclesiæ Doctoris \textit{— St. John of the Cross, Confessor and Doctor of the Church}}
\begin{paracol}{2}
\LA
\RefLine{Lectiones I–III: de Scriptura occurrente.}
\switchcolumn
\EN
\RefLine{Lessons I–III: from the Scripture of the day.}
\switchcolumn*

\LA
\RefLine{Lectiones VII–VIII: ex Commúni Doctoris non Pontificis (C5a).}
\switchcolumn
\EN
\RefLine{Lessons VII–VIII: from the Common of a Doctor not a Bishop (C5a).}
\switchcolumn*

\LA
\RefLine{Lectio IX: de S. Chrysogoni Martyris (11-24o).}
\switchcolumn
\EN
\RefLine{Lesson IX: of St. Chrysogonus, Martyr (11-24o).}
\switchcolumn*

\LA
\LessonHead{Lectio IV}
\switchcolumn
\EN
\LessonHead{Lesson IV}
\switchcolumn*

\LA
\noindent Joánnes a Cruce, Fontíberi in Hispánia piis paréntibus natus, a primis annis certo innotuit quam Deiparæ Virgini futurus esset acceptus; nam quinquennis, in púteum lapsus, ejusdem Deiparæ manu sublatus, incolumis evasit. Tanto autem patiéndi desidério flagrávit, ut novennis, spreto molliori lecto, super sarmentis cubare consueverit. Adoléscens hospítio páuperum ægrotántium Metymnæ Campi famulum sese addixit, quibus magno caritátis ardore, vilíssima quæque complectens officia, præsto aderat. Cujus exemplo excitáti ceteri, éadem caritátis múnera ardentius obibant. Verum, ad altióra vocatus, beátæ Maríæ Vírginis de Monte Carmelo institutum amplexus est; ubi, sacerdos ex obediéntia factus, severioris disciplinæ et arctioris vitæ cupidíssimus, primitivam ordinis regulam ex superioris licéntia ita proféssus est, ut, ob jugem Dominicæ passiónis memóriam, bello in se, tamquam in infensíssimum hostem indicto, vigiliis, jejuniis, férreis flagellis omnique pœnárum genere, brevi carnem cum vitiis et concupiscéntiis suis crucifixerit; dignus plane, qui a sancta Teresia inter puriores sanctioresque ánimas, Ecclésiam Dei id témporis illustrántes, recenserétur.\par
\switchcolumn
\EN
\noindent John of the Cross was born of godly parents at Fontibere, (near Avila,) in Spain, (in the year of our Lord 1542.) It began soon to appear that he was foreordained to be an acceptable servant unto the Virgin Mother of God. At five years of age he fell into a well, but the hand of the Mother of God took him up, and saved him from all hurt. So burning was his desire to suffer that when he was nine years old he gave up any softer bed, and used to lie on potsherds. In his youth he devoted himself as a servant in the hospital for the sick poor at Medina del Campo, and embraced with eager charity the meanest offices there, his readiness likewise exciting others to imitate him. (In 1563) he obeyed the call to higher things, and entered the Order of the Blessed Virgin Mary of Mount Carmel, wherein, by command of his Superiors, he received Priest's Orders. By their leave and his own strong desire for the sternest discipline and the strictest life, he adopted the primitive Rule. Full of the memory of what our Lord suffered, he declared war against himself as his own worst enemy, and carried it on by depriving himself of sleep and food, by iron chains, by whips, and by every kind of self-torture. And in a little while he had crucified the flesh, with the affections and lusts thereof. He was indeed worthy that holy Theresa should say of him that he was one of the purest and holiest souls by whom God was then enlightening His Church.\par
\switchcolumn*

\LA
\begin{Resp}\Rsign Honéstum fecit illum Dóminus, et custodívit eum ab inimícis, et a seductóribus tutávit illum: \Star{} Et dedit illi claritátem ætérnam.\end{Resp}
\begin{Resp}\Vsign Justum dedúxit Dóminus per vias rectas, et osténdit illi regnum Dei.\end{Resp}
\begin{Resp}\Rsign Et dedit illi claritátem ætérnam.\end{Resp}
\switchcolumn
\EN
\begin{Resp}\Rsign The Lord made him honourable, and defended him from his enemies, and kept him safe from those that lay in wait for him; \Star{} And gave him perpetual glory.\end{Resp}
\begin{Resp}\Vsign He went down with him into the pit, and left him not in bonds.\end{Resp}
\begin{Resp}\Rsign And gave him perpetual glory.\end{Resp}
\switchcolumn*

\LA
\LessonHead{Lectio V}
\switchcolumn
\EN
\LessonHead{Lesson V}
\switchcolumn*

\LA
\noindent Singulari vitæ austeritáte et ómnium virtútum præsidio munítus, præ assidua rerum divinárum contemplatióne, diuturnas et mirábiles éxtases frequenter patiebátur; tantoque in Deum æstuábat amóre, ut, cum divinus ignis sese intro diutius continere non posset, foras erumpere ejusque vultum irradiare visus sit. Proximórum salúti summopere intentus, tum in verbi Dei prædicatióne, tum in sacramentórum administratióne fuit assiduus. Hinc tot meritis auctus, strictiorisque disciplinæ promovendæ ardore veheménter accénsus, sanctæ Teresiæ comes divinitus datus est, ut, quam ipsa inter sorores primævam Carmeli ordinis observantiam instauráverat, eamdem et inter fratres, Joánne adjutore, restitúeret. Innumeros ítaque una cum Dei famula in divino opere promovéndo perpéssus labóres, cœnobia, quæ ejusdem sanctæ Vírginis cura per totam Hispániam erecta fuerant, nullis vitæ incommodis et periculis térritus, singula perlustrávit. In quibus aliisque quam plurimis, ejus ópera eréctis, restauratam observantiam propagando, verbo et exemplo firmávit; ut mérito primus, post sanctam Teresiam, Carmelitárum excalceatórum ordinis professor et parens habeátur.\par
\switchcolumn
\EN
\noindent The strange hardness of his life, and the might of his graces, joined to the unceasing concentration of his mind on God, had the effect of oftentimes subjecting him to daily and extraordinary trances. So burning was his love of God that the fire sometimes could not be kept bound within, and brake forth, so that his face shone. The salvation of his neighbours was one of his dearest longings, and he was unwearied in preaching the Word of God, and in administering the Sacraments. As strong in so many good works, and glowing with zeal to make discipline harder, he was given by God to be an helpmeet to holy Theresa, and he aided her to set up again the primitive observance among the brethren of the Order of Mount Carmel, as she had already done among the sisters. In doing God's work, he and God's handmaid together went through toils that cannot be numbered. No discomforts or dangers held him back from going throughout all Spain to visit all and each of the convents which the care of that holy Virgin had founded, and in them, and in very many others erected by her means for spreading the renewed observance, he strengthened it by his word and example. He is indeed worthy to be reckoned second only to the holy Theresa as a professor and founder of the Order of bare-footed Carmelites.\par
\switchcolumn*

\LA
\begin{Resp}\Rsign Amávit eum Dóminus, et ornávit eum: stolam glóriæ índuit eum, \Star{} Et ad portas paradísi coronávit eum.\end{Resp}
\begin{Resp}\Vsign Índuit eum Dóminus lorícam fídei, et ornávit eum.\end{Resp}
\begin{Resp}\Rsign Et ad portas paradísi coronávit eum.\end{Resp}
\switchcolumn
\EN
\begin{Resp}\Rsign The Lord loved him and beautified him; He clothed him with a robe of glory, \Star{} And crowned him at the gates of Paradise.\end{Resp}
\begin{Resp}\Vsign The Lord hath put on him the breast-plate of faith, and hath adorned him.\end{Resp}
\begin{Resp}\Rsign And crowned him at the gates of Paradise.\end{Resp}
\switchcolumn*

\LA
\LessonHead{Lectio VI}
\switchcolumn
\EN
\LessonHead{Lesson VI}
\switchcolumn*

\LA
\noindent Virginitátem perpetuo coluit, impudentésque mulieres ejus pudicítiæ insidiari conántes, non modo répulit, sed étiam Christo lucrifecit. In divinis explicandis arcanis æque ac sancta Teresia, apostolicæ Sedis judício, divínitus instructus, libros de mystica theología, cælésti sapiéntia refertos, conscripsit. Semel interrogátus a Christo quid præmii pro tot labóribus pósceret, respóndit: Dómine, pati et contemni pro te. Imperio in dæmones, quos e corpóribus sæpe fugabat, discretióne spirituum, prophetíæ dono, miraculórum glória celebratíssimus, ea semper fuit humilitate, ut sæpius a Dómino flagitaverit eo loco mori, ubi ómnibus esset ignotus. Voti compos factus. Ubedæ, diro morbo et in crure quinque plagis sanie manántibus, ad impléndum patiéndi desidérium constantíssime tolerátis, Ecclésiæ sacramentis pie sancteque susceptis, in Christi crucifixi amplexu, quem semper in corde atque ore habuerat, post illa verba: In manus tuas comméndo spíritum meum, obdormívit in Dómino, die et hora a se prædictis, anno salútis millesimo quingentésimo nonagesimo primo, ætátis quadragesimo nono. Migrántem ejus ánimam splendidíssimus ignis globus excepit; corpus vero suavíssimum odórem spirávit, quod, etiamnum incorruptum, Segoviæ honorifice cólitur. Eum, plurimis ante et post óbitum fulgentem signis, Benedíctus décimus tertius Pontifex maximus in Sanctórum númerum retulit, et Pius undecimus ex Sacrórum Rituum Congregatiónis consulto, universalis Ecclésiæ Doctórem declarávit.\par
\switchcolumn
\EN
\noindent He remained throughout all his life a clean maid, and when some shameless women tried to beguile his modesty, he not only foiled them, but gained them for Christ. In the judgment of the Apostolic See he was as much taught of God as was holy Theresa, for explaining God's hidden mysteries, and he wrote books of mystical theology filled with heavenly wisdom. Christ once asked him what reward he would have for so much work; whereto he answered Lord, that I may suffer, and be disesteemed for thy sake. He was very famous for his power over devils, whom he oftentimes scared out of men's bodies, for discerning of spirits, for the gift of prophecy, and for eminent miracles. He was extraordinarily lowly, and oftentimes entreated of the Lord that he might die in some place where he was unknown. In accordance with his prayer, (he was sent) to Ubeda, (where for three months the Prior imprisoned and cruelly ill-used him during his last sickness.) To crown his love of suffering, he bore uncomplainingly five open sores in his leg, running with water. (At last, upon the 14th day of December,) in the year 1591, being the day, and at the hour foretold by himself, after having in godly and holy wise received the Sacraments of the Church, hugging (the image of) that crucified Saviour of Whom his heart and his mouth had been used to be full, he uttered the words: “Into thy hands I commend my spirit”, and fell asleep in the Lord. As his soul passed away it was received into a glorious cloud of fire. His body yielded a right sweet savour, and is still uncorrupt where it lieth, held in great honour, at Segovia. He was famous for very many miracles both before and since his death, and Pope Benedict XIII numbered his name among those of the Saints.\par
\switchcolumn*

\LA
\begin{Resp}\Rsign Iste homo perfécit ómnia quæ locútus est ei Deus, et dixit ad eum: Ingrédere in réquiem meam: \Star{} Quia te vidi justum coram me ex ómnibus géntibus.\end{Resp}
\begin{Resp}\Vsign Iste est, qui contémpsit vitam mundi, et pervénit ad cæléstia regna.\end{Resp}
\begin{Resp}\Rsign Quia te vidi justum coram me ex ómnibus géntibus.\end{Resp}
\begin{Resp}\Vsign Glória Patri, et Fílio, \Star{} et Spirítui Sancto.\end{Resp}
\begin{Resp}\Rsign Quia te vidi justum coram me ex ómnibus géntibus.\end{Resp}
\switchcolumn
\EN
\begin{Resp}\Rsign This is he which did according unto all that God commanded him; and God said unto him: Enter thou into My rest, \Star{} For thee have I seen righteous before Me among all people.\end{Resp}
\begin{Resp}\Vsign This is he which loved not his life in this world, and is come unto an everlasting kingdom.\end{Resp}
\begin{Resp}\Rsign For thee have I seen righteous before Me among all people.\end{Resp}
\begin{Resp}\Vsign Glory be to the Father, and to the Son, \Star{} and to the Holy Ghost.\end{Resp}
\begin{Resp}\Rsign For thee have I seen righteous before Me among all people.\end{Resp}
\switchcolumn*

\end{paracol}

\par\needspace{10\baselineskip}\addvspace{1.2em}{\centering\small\textsc{Die 24 Novembris} · \textsc{24 November}\par}
\Office{S. Chrysogoni Martyris}{St. Chrysogonus, Martyr}{Simplex}
\addcontentsline{toc}{subsection}{Die 24 Novembris · S. Chrysogoni Martyris \textit{— St. Chrysogonus, Martyr}}
\begin{paracol}{2}
\LA
\LessonHead{Lectio IX (pro commemoratione)}
\switchcolumn
\EN
\LessonHead{Lesson IX (when commemorated)}
\switchcolumn*

\LA
\LessonTitle{Commemoratio: S. Chrysogoni Martyris}
\switchcolumn
\EN
\LessonTitle{Commemoration of S. Chrysogonus, Martyr}
\switchcolumn*

\LA
\noindent Chrysógonus, Diocletiáno imperatóre, Romæ inclusus in carcere, ibi biennium sanctæ Anastasiæ facultátibus vixit; quam etiam, afflictam propter Christum a viro suo Publio, proptereáque a suis oratiónibus per litteras auxílium postulántem, mútuis epistolis est consolatus. Sed, cum imperátor Romam scripsísset ut, reliquis Christiánis qui in vinculis essent interfectis, Chrysógonus Aquilejam ad se mitterétur, eo perductus est. Cui imperátor: Accersívi, inquit, te, Chrysógone, ut honóribus augeam, si modo induxeris ánimum deos cólere. At ille: Ego eum, qui vere est Deus, mente et oratióne véneror; deos autem, qui nihil sunt nisi dæmonum simulacra, odi et éxsecror. Quo responso excandéscens imperátor, ad Aquas Gradatas eum secúri pércuti jubet octavo Kalendas Decembris. Cujus corpus, projectum in mare, paulo post in littore invéntum, Zóilus presbyter in suis ædibus sepelívit.\par
\switchcolumn
\EN
\noindent Chrysogonus was imprisoned at Rome in the reign of the Emperor Diocletian. There he lived for the space of two years upon the alms of the holy Anastasia. She was suffering much persecution from her husband Publius for Christ's Name's sake, and was used to write to Chrysogonus to ask for the help of his prayers, and he in return comforted her by his epistles. Presently the Emperor wrote to Rome commanding the rest of the Christians who were in prison there to be put to death, and Chrysogonus “to be sent to himself at Aquileia. When he was brought thither, he said unto him: I have sent for thee, O Chrysogonus, that I may increase thine honours, if only thou wilt bring thy mind to worship the gods.” Thereto Chrysogonus answered: “With my mind and with my prayers I worship Him Who is God indeed, but such gods as are nothing but images of devils, them I hate and curse.” Then was the Emperor kindled to fury at this answer, and commanded Chrysogonus to be beheaded at Aquae Gradatae upon the 24th day of November. His body was cast into the sea, but found a little while afterwards washed up upon the shore, and the Priest Zoilus took it and buried it in his own house.\par
\switchcolumn*

\LA
\TeDeum{Te Deum}
\switchcolumn
\EN
\TeDeum{Te Deum}
\switchcolumn*

\end{paracol}

\par\needspace{10\baselineskip}\addvspace{1.2em}{\centering\small\textsc{Die 25 Novembris} · \textsc{25 November}\par}
\Office{S. Catharinæ Virginis et Martyris}{St. Catharine of Alexandria, Virgin and Martyr}{Duplex}
\addcontentsline{toc}{subsection}{Die 25 Novembris · S. Catharinæ Virginis et Martyris \textit{— St. Catharine of Alexandria, Virgin and Martyr}}
\begin{paracol}{2}
\LA
\RefLine{Lectiones I–III: de Scriptura occurrente.}
\switchcolumn
\EN
\RefLine{Lessons I–III: from the Scripture of the day.}
\switchcolumn*

\LA
\RefLine{Lectiones VII–IX: ex Commúni Unius Virginis et Martyris (C6).}
\switchcolumn
\EN
\RefLine{Lessons VII–IX: from the Common of a Virgin Martyr (C6).}
\switchcolumn*

\LA
\LessonHead{Lectio IV}
\switchcolumn
\EN
\LessonHead{Lesson IV}
\switchcolumn*

\LA
\noindent Catharina, nobilis virgo Alexandrina, a prima ætate stúdia liberalium artium cum fidei ardore conjungens, brevi ad eam sanctitátis et doctrinæ perfectiónem pervenit, ut, decem et octo annos nata, eruditíssimum quemque superaret. Quæ cum Maximíni jussu multos, propter christianæ religiónis professiónem varie tormentis cruciatos, ad supplícium rapi vidéret, non dubitanter ipsum adiit Maximinum, eique nefariam immanitátem objíciens, sapientíssimis ratiónibus Christi fidem ad salútem necessariam esse affirmavit.\par
\switchcolumn
\EN
\noindent This Catharine was a noble maiden of Alexandria, who from her earliest years joined the study of the liberal arts with fervent faith, and in a short while came to such an height of holiness and learning, that when she was eighteen years of age she prevailed over the chiefest wits. When she saw many diversely tormented and haled to death by command of Maximin, because they professed the Christian religion, she went boldly unto him and rebuked him for his savage cruelty, bringing forward likewise most sage reasons why the faith of Christ should be needful for salvation.\par
\switchcolumn*

\LA
\begin{Resp}\Rsign Propter veritátem, et mansuetúdinem, et justítiam: \Star{} Et dedúcet te mirabíliter déxtera tua.\end{Resp}
\begin{Resp}\Vsign Spécie tua et pulchritúdine tua inténde, próspere procéde, et regna.\end{Resp}
\begin{Resp}\Rsign Et dedúcet te mirabíliter déxtera tua.\end{Resp}
\switchcolumn
\EN
\begin{Resp}\Rsign Because of truth, and meekness, and righteousness; \Star{} And thy right hand shall lead thee wonderfully.\end{Resp}
\begin{Resp}\Vsign In thy comeliness, and thy beauty, go forward, fare prosperously, and reign.\end{Resp}
\begin{Resp}\Rsign And thy right hand shall lead thee wonderfully.\end{Resp}
\switchcolumn*

\LA
\LessonHead{Lectio V}
\switchcolumn
\EN
\LessonHead{Lesson V}
\switchcolumn*

\LA
\noindent Cujus prudéntiam Maximinus admiratus, retineri eam jubet, accersitis undique doctíssimis homínibus, magnisque propositis præmiis, qui convictam Catharinam a Christi fide ad idolórum cultum perduxissent. Quod contra áccidit. Nam plures philosophi, qui ad eam coarguéndam convenerant, vi ac subtilitate ejus disputatiónis tanto Jesu Christi amóre sunt incénsi, ut pro illo mori non dubitaverint. Quam ob rem Maximinus blanditiis ac promissis Catharinam de senténtia dedúcere aggréditur; verum, id frustra fíeri intélligens, verbéribus affectam plumbátisque contusam, dies undecim sine cibo ac potu inclusam tenet in carcere.\par
\switchcolumn
\EN
\noindent Maximian marvelled at her wisdom, and bade keep her, while he gathered together the most learned men from all quarters and offered them great rewards if they would confute Catharine and bring her from believing in Christ to worship idols. But the event fell contrariwise, for many of the philosophers who had come to dispute with her were overcome by the force and skill of her reasoning, so that the love of Christ Jesus was kindled in them, and they were content even to die for His sake. Then did Maximin strive to beguile Catharine with fair words and promises, and when he found it was lost pains, he caused her to be hided, and bruised with lead-laden whips, and so cast into prison, and neither meat nor drink given to her for the space of eleven days.\par
\switchcolumn*

\LA
\begin{Resp}\Rsign Dilexísti justítiam, et odísti iniquitátem: \Star{} Proptérea unxit te Deus, Deus tuus, óleo lætítiæ.\end{Resp}
\begin{Resp}\Vsign Propter veritátem, et mansuetúdinem, et justítiam.\end{Resp}
\begin{Resp}\Rsign Proptérea unxit te Deus, Deus tuus, óleo lætítiæ.\end{Resp}
\switchcolumn
\EN
\begin{Resp}\Rsign Thou hast loved righteousness, and hated iniquity; \Star{} Therefore God, thy God, hath anointed thee with the oil of gladness.\end{Resp}
\begin{Resp}\Vsign Because of truth, and meekness, and righteousness.\end{Resp}
\begin{Resp}\Rsign Therefore God, thy God, hath anointed thee with the oil of gladness.\end{Resp}
\switchcolumn*

\LA
\LessonHead{Lectio VI}
\switchcolumn
\EN
\LessonHead{Lesson VI}
\switchcolumn*

\LA
\noindent Quo témpore Maximíni uxor et Porphyrius belli dux, visendæ Vírginis causa carcerem ingressi, et ejusdem prædicatióne in Jesum Christum credéntes, póstea martyrio coronáti sunt. Interim Catharina edúcitur e custódia; et rota expedítur, crebris et acutis præfixa gládiis, ut Vírginis corpus crudelíssime dilacerarétur. Quæ máchina, brevi Catharinæ oratióne, confracta est; eoque miraculo multi Christi fidem suscepérunt. Ipse Maximinus, in impietáte et crudelitate obstinatior, Catharinam secúri pércuti imperat. Quæ, fortiter dato cápite, ad duplicátum virginitátis et martyrii præmium evolávit septimo Kalendas Decembris; cujus corpus ab Angelis in Sina, Arábiæ monte, mirabíliter collocátum est.\par
\switchcolumn
\EN
\noindent At that time Maximin's wife and Porphyry the Captain of his host, went to the prison to see the damsel, and at her preaching believed in Jesus Christ, and were afterwards crowned with martyrdom. Then was Catharine brought out of ward, and a wheel was set, wherein were fastened many and sharp blades, so that her virgin body might thereby be most direfully cut and torn in pieces, but in a little while, as Catharine prayed, this machine was broken in pieces, at the which marvel many believed in Christ. But Maximin was hardened in his godlessness and cruelty, and commanded to behead Catharine. She bravely offered her neck to the stroke and passed away hence to receive the twain crowns of maidenhood and martyrdom, upon the 25th day of November. Her body was marvelously laid by angels upon Mount Sinai in Arabia.\par
\switchcolumn*

\LA
\begin{Resp}\Rsign Afferéntur Regi vírgines post eam, próximæ ejus; \Star{} Afferéntur tibi in lætítia et exsultatióne.\end{Resp}
\begin{Resp}\Vsign Spécie tua et pulchritúdine tua; inténde, próspere procéde, et regna.\end{Resp}
\begin{Resp}\Rsign Afferéntur tibi in lætítia et exsultatióne.\end{Resp}
\begin{Resp}\Vsign Glória Patri, et Fílio, \Star{} et Spirítui Sancto.\end{Resp}
\begin{Resp}\Rsign Afferéntur tibi in lætítia et exsultatióne.\end{Resp}
\switchcolumn
\EN
\begin{Resp}\Rsign After her shall virgins be brought unto the King, her fellows \Star{} They shall be brought unto thee with gladness and rejoicing.\end{Resp}
\begin{Resp}\Vsign In thy comeliness and thy beauty, go forward, fare prosperously, and reign.\end{Resp}
\begin{Resp}\Rsign They shall be brought unto thee with gladness and rejoicing.\end{Resp}
\begin{Resp}\Vsign Glory be to the Father, and to the Son, \Star{} and to the Holy Ghost.\end{Resp}
\begin{Resp}\Rsign They shall be brought unto thee with gladness and rejoicing.\end{Resp}
\switchcolumn*

\end{paracol}

\par\needspace{10\baselineskip}\addvspace{1.2em}{\centering\small\textsc{Die 26 Novembris} · \textsc{26 November}\par}
\Office{S. Silvestri Abbatis}{St. Sylvester, Abbot}{Duplex}
\addcontentsline{toc}{subsection}{Die 26 Novembris · S. Silvestri Abbatis \textit{— St. Sylvester, Abbot}}
\begin{paracol}{2}
\LA
\RefLine{Lectiones I–III: de Scriptura occurrente.}
\switchcolumn
\EN
\RefLine{Lessons I–III: from the Scripture of the day.}
\switchcolumn*

\LA
\RefLine{Lectio IX: de S. Petri Alexandrini Martyris (11-26t).}
\switchcolumn
\EN
\RefLine{Lesson IX: of St. Peter of Alexandria, Martyr (11-26t).}
\switchcolumn*

\LA
\LessonHead{Lectio IV}
\switchcolumn
\EN
\LessonHead{Lesson IV}
\switchcolumn*

\LA
\noindent Silvester, Auximi in Piceno nobili genere ortus, statim puerílem ætátem litteris ac bonis móribus mirifice exornavit. Adoléscens, Bonóniam ad stúdia jurisprudéntiæ missus a patre, cum sacris litteris, a Deo mónitus, dedísset operam, parentis incurrit indignatiónem; quam æquo animo toto decennio pértulit. Ob egregiam ejus virtútem a canonicis cathedralis Auximanæ ecclésiæ socius honoris eléctus est; in quo munere pópulo oratiónibus, exemplo et conciónibus opem tulit.\par
\switchcolumn
\EN
\noindent This Silvester was born of a noble family at Osimo, in Picenum, and in his childhood was a wonderful example both in regard to letters and good living. When he grew older his father sent him to Bologna to study the law, but God warned him to give himself to divinity, and he thereby incurred the wrath of his father, which he bore with complacency for ten full years. On account of his eminent graces he was elected an honorary canon of the Cathedral of Osimo, in the which dignity he ministered to the people by his prayers, his example, and his sermons.\par
\switchcolumn*

\LA
\begin{Resp}\Rsign Honéstum fecit illum Dóminus, et custodívit eum ab inimícis, et a seductóribus tutávit illum: \Star{} Et dedit illi claritátem ætérnam.\end{Resp}
\begin{Resp}\Vsign Justum dedúxit Dóminus per vias rectas, et osténdit illi regnum Dei.\end{Resp}
\begin{Resp}\Rsign Et dedit illi claritátem ætérnam.\end{Resp}
\switchcolumn
\EN
\begin{Resp}\Rsign The Lord made him honourable, and defended him from his enemies, and kept him safe from those that lay in wait for him; \Star{} And gave him perpetual glory.\end{Resp}
\begin{Resp}\Vsign He went down with him into the pit, and left him not in bonds.\end{Resp}
\begin{Resp}\Rsign And gave him perpetual glory.\end{Resp}
\switchcolumn*

\LA
\LessonHead{Lectio V}
\switchcolumn
\EN
\LessonHead{Lesson V}
\switchcolumn*

\LA
\noindent Inter funus nobilis cujusdam defuncti, in aperto tumulo formosi viri suique propinqui deforme cadáver conspíciens: Ego, inquit, sum, quod hic fuit; quod hic est, ego ero. Et mox, peracto funere, illa sibi Dómini occurrénte senténtia: Qui vult veníre post me, ábneget semetípsum, et tollat crucem suam, et sequátur me; in solitúdinem, majoris perfectiónis studio, secessit, ibique vigiliis, oratiónibus jejuniisque deditus, crudas tantum herbas in cibum sæpius adhibuit. Ut autem magis latéret hómines, varias mutavit sedes; ac demum pervenit ad montem Fanum, locum, quamvis prope Fabrianum, eo tamen témpore desértum, ibique in honórem sanctíssimi patris Benedicti templum eréxit, congregatiónisque Silvestrinórum fundaménta jecit, sub regula et habitu in visióne sibi ab eodem Sancto osténsis.\par
\switchcolumn
\EN
\noindent At the funeral of a certain nobleman he perceived in an open grave the disfigured corpse of a kinsman of his own who had been very comely in his lifetime, and he said to himself, I am what he was, and what he is I shall be. Straightway after the funeral he read the words of the Lord, If any man will come after Me let him deny himself and take up his cross and follow Me (Matth. xvi. 24.) Thereupon he withdrew into the desert to seek after greater perfection, and then gave himself up to watching, praying, and fasting, very often taking no food but uncooked herbs. In order, however, to cut himself off the more from men, he moved from one place to another, and at length came to Mount Fano, which is hard by Fabriano, but was itself then absolutely uninhabited. Then he built a church in honour of the holy Father Benedict, and founded the congregation of Silvestrians, with a rule and dress which were revealed to him in a vision by the holy Patriarch himself.\par
\switchcolumn*

\LA
\begin{Resp}\Rsign Amávit eum Dóminus, et ornávit eum: stolam glóriæ índuit eum, \Star{} Et ad portas paradísi coronávit eum.\end{Resp}
\begin{Resp}\Vsign Índuit eum Dóminus lorícam fídei, et ornávit eum.\end{Resp}
\begin{Resp}\Rsign Et ad portas paradísi coronávit eum.\end{Resp}
\switchcolumn
\EN
\begin{Resp}\Rsign The Lord loved him and beautified him; He clothed him with a robe of glory, \Star{} And crowned him at the gates of Paradise.\end{Resp}
\begin{Resp}\Vsign The Lord hath put on him the breast-plate of faith, and hath adorned him.\end{Resp}
\begin{Resp}\Rsign And crowned him at the gates of Paradise.\end{Resp}
\switchcolumn*

\LA
\LessonHead{Lectio VI}
\switchcolumn
\EN
\LessonHead{Lesson VI}
\switchcolumn*

\LA
\noindent At ínvidens sátanas variis terróribus illíus monachos turbare nitebátur, noctu monasterii jánuas hostíliter invadens. Sed vir Dei, hostis impetum ita repressit, ut monachi in sancto instituto magis confirmaréntur ac patris sanctitátem agnoscerent. Spíritu prophetíæ aliisque donis enituit. Quæ ut semper profúnda humilitate conservavit, ita contra se dæmonis invidiam concitavit; a quo præceps actus per scalas oratórii, et prope interimendus, præsentíssimo Vírginis beneficio incolumitáti redditus est. Quod benefícium perpetua et singulari in illam pietáte commendavit ad ultimum usque vitæ spíritum, quem, fere nonagenarius, sanctitáte et miraculis clarus, Deo réddidit anno salútis millesimo ducentésimo sexagesimo septimo, sexto Kalendas Decembris. Ejus Offícium ac Missam Leo décimus tertius Pontifex maximus ad universam exténdit Ecclésiam.\par
\switchcolumn
\EN
\noindent Satan envied him, strove to trouble his monks by diverse terrors, and made an hostile attack by night upon the gates of his monastery, but the man of God so overcame the assault of the enemy that his monks were the more confirmed in their Institute and recognised the holiness of their father. He shone with the spirit of prophecy and other gifts. These things he always preserved by the deepest lowliness, whereby he so stirred up against him the ill-will of the devil that that evil spirit cast him headlong down the stairs of his oratory, and went near to slay him, but he was restored to soundness by the helpful gift of the Virgin. This help he remembered with an unceasing and singular love toward her until the last breath of his life, the which breath he resigned to God, famous for holiness and miracles, aged almost ninety years, upon the 26th day of November, in the year of salvation 1267. The Supreme Pontiff Leo XIII extended his Office and Mass to the whole Church.\par
\switchcolumn*

\LA
\begin{Resp}\Rsign Iste homo perfécit ómnia quæ locútus est ei Deus, et dixit ad eum: Ingrédere in réquiem meam: \Star{} Quia te vidi justum coram me ex ómnibus géntibus.\end{Resp}
\begin{Resp}\Vsign Iste est, qui contémpsit vitam mundi, et pervénit ad cæléstia regna.\end{Resp}
\begin{Resp}\Rsign Quia te vidi justum coram me ex ómnibus géntibus.\end{Resp}
\begin{Resp}\Vsign Glória Patri, et Fílio, \Star{} et Spirítui Sancto.\end{Resp}
\begin{Resp}\Rsign Quia te vidi justum coram me ex ómnibus géntibus.\end{Resp}
\switchcolumn
\EN
\begin{Resp}\Rsign This is he which did according unto all that God commanded him; and God said unto him: Enter thou into My rest, \Star{} For thee have I seen righteous before Me among all people.\end{Resp}
\begin{Resp}\Vsign This is he which loved not his life in this world, and is come unto an everlasting kingdom.\end{Resp}
\begin{Resp}\Rsign For thee have I seen righteous before Me among all people.\end{Resp}
\begin{Resp}\Vsign Glory be to the Father, and to the Son, \Star{} and to the Holy Ghost.\end{Resp}
\begin{Resp}\Rsign For thee have I seen righteous before Me among all people.\end{Resp}
\switchcolumn*

\LA
\LessonHead{Lectio VII}
\switchcolumn
\EN
\LessonHead{Lesson VII}
\switchcolumn*

\LA
\LessonTitle{Léctio sancti Evangélii secúndum Matthǽum.}
\switchcolumn
\EN
\LessonTitle{From the Holy Gospel according to Matthew}
\switchcolumn*

\LA
\Cite{Matt 19:27-29}
\switchcolumn
\EN
\Cite{Matt 19:27-29}
\switchcolumn*

\LA
\noindent In illo témpore: Dixit Petrus ad Jesum: Ecce nos relíquimus ómnia, et secúti sumus te: quid ergo erit nobis? Et réliqua.\par
\switchcolumn
\EN
\noindent At that time, Peter said unto Jesus: Behold, we have forsaken all, and followed thee what shall we have therefore? And so on.\par
\switchcolumn*

\LA
\LessonTitle{Homilía sancti Hierónymi Presbýteri.}
\switchcolumn
\EN
\LessonTitle{Homily by St. Jerome, Priest (at Bethlehem.)}
\switchcolumn*

\LA
\Source{Lib. 3. in Matth. cap. 19.}
\switchcolumn
\EN
\Source{Bk. iii. on Matth xix.}
\switchcolumn*

\LA
\noindent Grandis fidúcia! Petrus piscátor erat, dives non fúerat, cibos manu et arte quærébat: et tamen lóquitur confidénter: Relíquimus ómnia. Et quia non súfficit tantum relínquere, jungit quod perféctum est: Et secúti sumus te. Fécimus quod jussísti: quid ígitur nobis dabis prǽmii? Jesus autem dixit illis: Amen dico vobis, quod vos qui secúti estis me, in regeneratióne, cum séderit Fílius hóminis in sede majestátis suæ, sedébitis et vos super sedes duódecim, judicántes duódecim tribus Ísraël. Non dixit: Qui reliquístis ómnia; hoc enim et Crates fecit philósophus, et multi álii divítias contempsérunt: sed, Qui secúti estis me: quod próprie Apostolórum est, atque credéntium.\par
\switchcolumn
\EN
\noindent Peter was a fisherman, he was not rich, he earned his bread by his hand and skill, and nevertheless he is thus bold, and saith confidently: We have forsaken all. And because it sufficeth not to forsake only, he addeth that which to do is to be perfect: and followed thee. We have done that which Thou hast commanded us, what reward therefore wilt Thou give us? And Jesus said unto them Amen I say unto you, that ye which have followed Me, in the regeneration, when the Son of Man shall sit in the throne of His glory, ye also shall sit upon twelve thrones, judging the twelve tribes of Israel. He said not, Ye which have forsaken all, for this did even Crates the philosopher, and they which have set nothing by riches are many, but, Ye which have followed Me. This did the Apostles, and this do believers do.\par
\switchcolumn*

\LA
\begin{Resp}\Rsign Iste est qui ante Deum magnas virtútes operátus est, et de omni corde suo laudávit Dóminum: \Star{} Ipse intercédat pro peccátis ómnium populórum.\end{Resp}
\begin{Resp}\Vsign Ecce homo sine queréla, verus Dei cultor, ábstinens se ab omni ópere malo, et pérmanens in innocéntia sua.\end{Resp}
\begin{Resp}\Rsign Ipse intercédat pro peccátis ómnium populórum.\end{Resp}
\switchcolumn
\EN
\begin{Resp}\Rsign This is he which wrought great wonders before God, and praised the Lord with all his heart. \Star{} May he pray for all people, that their sins may be forgiven unto them.\end{Resp}
\begin{Resp}\Vsign Behold a man without blame, a worshipper of God in truth, keeping himself clean from every evil work, and abiding still in his innocency.\end{Resp}
\begin{Resp}\Rsign May he pray for all people, that their sins may be forgiven unto them.\end{Resp}
\switchcolumn*

\LA
\LessonHead{Lectio VIII}
\switchcolumn
\EN
\LessonHead{Lesson VIII}
\switchcolumn*

\LA
\noindent In regeneratióne, cum séderit Fílius hóminis in sede majestátis suæ (quando et mórtui de corruptióne resúrgent incorrúpti) sedébitis et vos in sóliis judicántium, condemnántes duódecim tribus Ísraël: quia vobis credéntibus, illi crédere noluérunt. Et omnis qui relíquerit domum, vel fratres, aut soróres, aut patrem, aut matrem, aut uxórem, aut fílios, aut agros propter nomen meum, céntuplum accípiet, et vitam ætérnam possidébit. Locus iste cum illa senténtia cóngruit, in qua Salvátor lóquitur: Non veni pacem míttere, sed gládium. Veni enim separáre hóminem a patre suo, et matrem a fília, et nurum a socru: et inimíci hóminis doméstici ejus. Qui ergo propter fidem Christi, et prædicatiónem Evangélii, omnes afféctus contémpserint, atque divítias et sǽculi voluptátes: isti céntuplum recípient, et vitam ætérnam possidébunt.\par
\switchcolumn
\EN
\noindent In the regeneration, when the Son of Man shall sit in the throne of His glory, and when the dead shall rise again from corruption incorruptible, (i Cor. xv. 53,) ye also shall sit upon twelve thrones of judgment, condemning the twelve tribes of Israel, because, when ye believed in Me, they would not. (John iii. 18.) And every one that hath forsaken houses, or brethren, or sisters, or father, or mother, or wife, or children, or lands, for My Name's sake, shall receive an hundredfold, and shall inherit everlasting life. This place agreeth well with that other where the Saviour saith I came not to send peace, but a sword. For I am come to set a man at variance against his father, and the daughter against her mother, and the daughter-in-law against her mother-in-law; and a man's foes shall be they of his own household. (Matth. x. 34.) Every one, therefore, that hath set no store by affection, and riches, and the pleasures of the world, for Christ's faith's sake, and the preaching of the Gospel, shall receive an hundred-fold, and shall inherit everlasting life.\par
\switchcolumn*

\LA
\begin{Resp}\Rsign Sint lumbi vestri præcíncti, et lucérnæ ardéntes in mánibus vestris: \Star{} Et vos símiles homínibus exspectántibus dóminum suum, quando revertátur a núptiis.\end{Resp}
\begin{Resp}\Vsign Vigiláte ergo, quia nescítis qua hora Dóminus vester ventúrus sit.\end{Resp}
\begin{Resp}\Rsign Et vos símiles homínibus exspectántibus dóminum suum, quando revertátur a núptiis.\end{Resp}
\begin{Resp}\Vsign Glória Patri, et Fílio, \Star{} et Spirítui Sancto.\end{Resp}
\begin{Resp}\Rsign Et vos símiles homínibus exspectántibus dóminum suum, quando revertátur a núptiis.\end{Resp}
\switchcolumn
\EN
\begin{Resp}\Rsign Let your loins be girded about, and your lights burning; \Star{} And ye yourselves like unto men that wait for their lord, when he will return from the wedding.\end{Resp}
\begin{Resp}\Vsign Watch therefore, for ye know not what hour your Lord doth come.\end{Resp}
\begin{Resp}\Rsign And ye yourselves like unto men that wait for their lord, when he will return from the wedding.\end{Resp}
\begin{Resp}\Vsign Glory be to the Father, and to the Son, \Star{} and to the Holy Ghost.\end{Resp}
\begin{Resp}\Rsign And ye yourselves like unto men that wait for their lord, when he will return from the wedding.\end{Resp}
\switchcolumn*

\end{paracol}

\par\needspace{10\baselineskip}\addvspace{1.2em}{\centering\small\textsc{Die 26 Novembris} · \textsc{26 November}\par}
\Office{S. Petri Alexandrini Martyris}{St. Peter of Alexandria, Martyr}{Simplex}
\addcontentsline{toc}{subsection}{Die 26 Novembris · S. Petri Alexandrini Martyris \textit{— St. Peter of Alexandria, Martyr}}
\begin{paracol}{2}
\LA
\LessonHead{Lectio IX (pro commemoratione)}
\switchcolumn
\EN
\LessonHead{Lesson IX (when commemorated)}
\switchcolumn*

\LA
\LessonTitle{Commemoratio: S. Petri Alexandrini Martyris}
\switchcolumn
\EN
\LessonTitle{Commemoration of: St. Peter of Alexandria, Martyr}
\switchcolumn*

\LA
\noindent Petrus, epíscopus Alexandríæ, post Theonam virum sanctíssimum, sanctitátis et doctrinæ splendore non solum illustravit Ægyptum, sed toti luxit Ecclésiæ Dei. Qui in persecutióne Maximiáni Galerii illam témporum acerbitátem ita pértulit, ut multi, admirábilem ejus patiéntiam intuéntes, plurimum in christiana virtúte proficerent. Is primus Aríum diaconum Alexandrinum, propter schisma Meletianum cui favebat, a fidelium communióne sejunxit. Ad eum, cápitis ab eodem Maximiáno damnátum, in carcere cum Achillas et Alexander presbyteri deprecatores Aríi venissent, respóndit, noctu apparuísse sibi Jesum veste discissa, causamque rei sciscitanti dixisse: Aríus vestem meam, quæ est Ecclésia, dilaceravit. Quibus étiam prædícens fore, ut sibi in episcopatu succéderent, præcepit ne umquam Aríum in communiónem recíperent, quem Deo mortuum esse sciret. Et hanc divinam prænotiónem veram fuísse, non diu post rei probavit eventus. Denique, duodecimo sui episcopátus anno, sexto Kalendas Decembris, abscisso cápite, ad martyrii corónam evolávit.\par
\switchcolumn
\EN
\noindent This Peter succeeded that eminent Saint, Theonas, as Pope of Alexandria, (in the year of our Lord 300,) and the glory of his holiness and teaching hath enlightened not Egypt only, but the whole Church of God. The wondrous patience wherewith he bore the roughness of the times in the persecution under Maximian Galerius caused many greatly to increase in Christian graces. He was the first who cut off Arius, then a Deacon of Alexandria, from the Communion of the faithful, on account of his leaning to the Meletian schism. He was condemned to death by Maximian, and was in prison when there came to him the two Priests Achilles and Alexander to plead for Arius, but Peter told them that Jesus had appeared to him in the night clad in a rent garment, and when he asked what was thereby signified, had said unto him Arius hath torn My vesture, which is the Church. Also, he foretold to them that they should be Popes of Alexandria after him, and strictly commanded them never to receive Arius into Communion, because he knew him to be dead in the sight of God. That this was a true prophecy the event did shortly prove. At length, in the twelfth year of his papacy, upon the 26th day of November, (in the year of salvation 311,) his head was cut off, and he went hence to receive the crown of his testimony.\par
\switchcolumn*

\LA
\TeDeum{Te Deum}
\switchcolumn
\EN
\TeDeum{Te Deum}
\switchcolumn*

\end{paracol}

\par\needspace{10\baselineskip}\addvspace{1.2em}{\centering\small\textsc{Die 29 Novembris} · \textsc{29 November}\par}
\Office{In Vigilia S. Andreæ Apostoli}{Vigil of St. Andrew the Apostle}{Vigilia}
\addcontentsline{toc}{subsection}{Die 29 Novembris · In Vigilia S. Andreæ Apostoli \textit{— Vigil of St. Andrew the Apostle}}
\begin{paracol}{2}
\LA
\LessonHead{Lectio I}
\switchcolumn
\EN
\LessonHead{Lesson I}
\switchcolumn*

\LA
\LessonTitle{Léctio sancti Evangélii secúndum Joánnem}
\switchcolumn
\EN
\LessonTitle{Continuation of the Holy Gospel according to John}
\switchcolumn*

\LA
\Cite{Joannes 1:35-51}
\switchcolumn
\EN
\Cite{John 1:35-51}
\switchcolumn*

\LA
\noindent In illo témpore: Stabat Joannes, et ex discipulis ejus duo. Et respiciens Jesum ambulantem, dicit: Ecce Agnus Dei. Et réliqua.\par
\switchcolumn
\EN
\noindent At that time again John stood, and two of his disciples. And beholding Jesus walking, he saith: Behold the Lamb of God. And so on.\par
\switchcolumn*

\LA
\LessonTitle{Homilia sancti Augustíni Epíscopi}
\switchcolumn
\EN
\LessonTitle{Homily of St. Augustine, bishop of Hippo}
\switchcolumn*

\LA
\Source{Tract. 7 in Joann., post init.}
\switchcolumn
\EN
\Source{7th Tract on John}
\switchcolumn*

\LA
\noindent Quia talis erat Joannes amicus sponsi, non quærebat gloriam suam, sed testimonium perhibebat veritati: numquid voluit apud se remanere discipulos suos, ut non sequerentur Dominum? Magis ipse ostendit discipulis suis quem sequerentur: habebant enim illum tamquam agnum. Et ille: Quid me attenditis: Ego non sum agnus. Ecce Agnus Dei. De quo et superius dixerat: Ecce Agnus Dei. Et quid nobis prodest Agnus Dei? Ecce, ait, qui tollit peccatum mundi. Secuti sunt illum, hoc audito, duo qui erant cum Joanne.\par
\switchcolumn
\EN
\noindent Since John was the friend of the Bridegroom, (iii. 29,) he sought not his own glory, but bare witness to the truth. Would he that his disciples should remain with him rather than that they should follow the Lord? Nay, he showed his disciples Whom they should follow. They thought that he himself was the Lamb; but he saith Why wait ye on me? I am not the Lamb. Behold the Lamb of God! This was He of Whom he had already said above (29) Behold the Lamb of God! And what use to us is the Lamb of God? Behold the Lamb of God, saith John, Which taketh away the sin of the world. And the two disciples heard him speak, and they followed Jesus.\par
\switchcolumn*

\LA
\begin{Resp}\Rsign Laudábilis pópulus, \Star{} Quem Dóminus exercítuum benedíxit dicens: Opus mánuum meárum tu es, heréditas mea Israël.\end{Resp}
\begin{Resp}\Vsign Beáta gens, cujus est Dóminus Deus, pópulus eléctus in hereditátem.\end{Resp}
\begin{Resp}\Rsign Quem Dóminus exercítuum benedíxit dicens: Opus mánuum meárum tu es, heréditas mea Israël.\end{Resp}
\switchcolumn
\EN
\begin{Resp}\Rsign Blessed is the people \Star{} Whom the Lord of hosts hath blessed, saying O Israel thou art the work of Mine own hands, thou art Mine own inheritance.\end{Resp}
\begin{Resp}\Vsign Blessed is the nation whose God is the Lord, and the people whom He hath chosen for His own inheritance.\end{Resp}
\begin{Resp}\Rsign Whom the Lord of hosts hath blessed, saying O Israel thou art the work of Mine own hands, thou art Mine own inheritance.\end{Resp}
\switchcolumn*

\LA
\LessonHead{Lectio II}
\switchcolumn
\EN
\LessonHead{Lesson II}
\switchcolumn*

\LA
\noindent Videamus sequentia. Ecce Agnus Dei. Hoc Joannes. Et audierunt eum duo discipuli loquentem, et secuti sunt Jesum. Non sic illum sequebantur, quasi jam ut inhærerent illi: nam manifestum est, quando illi inhæserunt, quia de navi eos vocavit. In his enim duobus erat Andreas, sicut modo audistis. Andreas autem frater Petri erat. Et novimus in Evangelio quod Petrum et Andream Dominus de navi vocavit, dicens: Venite post me, et faciam vos piscatores hominum. Et ex illo jam inhæserunt illi, ut non recederent.\par
\switchcolumn
\EN
\noindent Let us see what followed. John stood, and two of his disciples; and looking on Jesus as He walked, he saith Behold the Lamb of God! And the two disciples heard him speak, and they followed Jesus. They followed Him, not yet to cleave unto Him, for it is manifest that they clave unto Him only after that He had called them out of the ship. (Matth. iv. 18.) One of the two which heard John speak, and followed Him, was Andrew, Simon Peter's brother. And we know how it is written in the Gospel of Matthew: Jesus walking by the sea of Galilee, saw two brethren, Simon called Peter, and Andrew his brother, casting a net into the sea, for they were fishers. And He saith unto them Follow Me, and I will make you fishers of men. And they straightway left their nets and followed Him.\par
\switchcolumn*

\LA
\begin{Resp}\Rsign Angústiæ mihi sunt úndique, et quid éligam ignóro; \Star{} Mélius est mihi incídere in manus hóminum, quam derelínquere legem Dei mei.\end{Resp}
\begin{Resp}\Vsign Si enim hoc égero, mors mihi est; si autem non égero, non effúgiam manus vestras.\end{Resp}
\begin{Resp}\Rsign Mélius est mihi incídere in manus hóminum, quam derelínquere legem Dei mei.\end{Resp}
\switchcolumn
\EN
\begin{Resp}\Rsign I am straitened on every side, and know not what to choose. \Star{} It is better for me to fall into the hands of men, than to sin against the law of my God.\end{Resp}
\begin{Resp}\Vsign For if I do this thing, it is death unto me and if I do it not, I cannot escape your hands.\end{Resp}
\begin{Resp}\Rsign It is better for me to fall into the hands of men, than to sin against the law of my God.\end{Resp}
\switchcolumn*

\LA
\LessonHead{Lectio III}
\switchcolumn
\EN
\LessonHead{Lesson III}
\switchcolumn*

\LA
\noindent Modo ergo quod illum sequuntur isti duo, non quasi non recessuri sequuntur, sed videre voluerunt ubi habitaret, et facere quod scriptum est: Limen ostiorum ejus exterat pes tuus: surge ad illum venire assidue, et erudire præceptis ejus. Ostendit eis ille ubi maneret: venerunt, et fuerunt cum illo. Quam beatum diem duxerunt, quam beatam noctem! Quis est, qui nobis dicat quæ audierint illi a Domino? Aedificemus et nosmetipsi in corde nostro, et faciamus domum, quo veniat ille, et doceat nos, et colloquatur nobis.\par
\switchcolumn
\EN
\noindent From that time, therefore, was it that they clave unto Him continuously. Then Jesus turned, and saw them following, and saith unto them What seek ye? They said unto Him Rabbi, (which is to say, being interpreted, Master,) where dwellest Thou? So they follow Him now, not as to cleave unto Him for ever, but as to know where He dwelt, and to obey that which is written If thou see a man of understanding, go to him early in the morning, and let thy foot wear the steps of his doors. (Ecclus. vi. 36.) He saith unto them Come and see. They came and saw where He dwelt, and abode with Him that day. O what a blessed day! O what a blessed night! for it was about the tenth hour. Who shall tell what they heard from the Lord? O let us also build an house in our hearts, where He may come, and teach us, and talk with us!\par
\switchcolumn*

\LA
\begin{Resp}\Rsign Misit Dóminus Angelum suum et conclúsit ora leónum, \Star{} Et non contaminavérunt, quia coram eo injustítia invénta non est in me.\end{Resp}
\begin{Resp}\Vsign Misit Deus misericórdiam suam et veritátem suam: ánimam meam erípuit de médio catulórum leónum.\end{Resp}
\begin{Resp}\Rsign Et non contaminavérunt, quia coram eo injustítia invénta non est in me.\end{Resp}
\begin{Resp}\Vsign Glória Patri, et Fílio, \Star{} et Spirítui Sancto.\end{Resp}
\begin{Resp}\Rsign Et non contaminavérunt: quia coram eo injustítia invénta non est in me.\end{Resp}
\switchcolumn
\EN
\begin{Resp}\Rsign The Lord hath sent His angel, and hath shut the lions' mouths, that \Star{} They have not hurt me forasmuch as before Him innocency was found in me.\end{Resp}
\begin{Resp}\Vsign God hath sent forth His mercy and His truth, (and delivered) my soul from among the lions' whelps.\end{Resp}
\begin{Resp}\Rsign They have not hurt me forasmuch as before Him innocency was found in me.\end{Resp}
\begin{Resp}\Vsign Glory be to the Father, and to the Son, \Star{} and to the Holy Ghost.\end{Resp}
\begin{Resp}\Rsign They have not hurt me forasmuch as before Him innocency was found in me.\end{Resp}
\switchcolumn*

\end{paracol}

\par\needspace{10\baselineskip}\addvspace{1.2em}{\centering\small\textsc{Die 30 Novembris} · \textsc{30 November}\par}
\Office{S. Andreæ Apostoli}{St. Andrew the Apostle}{Duplex II classis}
\addcontentsline{toc}{subsection}{Die 30 Novembris · S. Andreæ Apostoli \textit{— St. Andrew the Apostle}}
\begin{paracol}{2}
\LA
\LessonHead{Lectio I}
\switchcolumn
\EN
\LessonHead{Lesson I}
\switchcolumn*

\LA
\LessonTitle{De Epístola beáti Pauli Apóstoli ad Romános}
\switchcolumn
\EN
\LessonTitle{Lesson from the letter of St. Paul the Apostle to the Romans}
\switchcolumn*

\LA
\Cite{Rom 10:4-9}
\switchcolumn
\EN
\Cite{Rom 10:4-9}
\switchcolumn*

\LA
\noindent Finis legis, Christus, ad justítiam omni credénti. Móyses enim scripsit, quóniam justítiam, quæ ex lege est, qui fécerit homo, vivet in ea. Quæ autem ex fide est justítia, sic dicit: Ne díxeris in corde tuo: Quis ascéndet in cælum? id est, Christum dedúcere: aut, Quis descéndet in abýssum? hoc est, Christum a mórtuis revocáre. Sed quid dicit Scriptúra? Prope est verbum in ore tuo, et in corde tuo: hoc est verbum fídei, quod prædicámus. Quia si confiteáris in ore tuo Dóminum Jesum, et in corde tuo credíderis quod Deus illum suscitávit a mórtuis, salvus eris.\par
\switchcolumn
\EN
\noindent For the end of the law is Christ, unto justice to every one that believeth. For Moses wrote, that the justice which is of the law, the man that shall do it, shall live by it. But the justice which is of faith, speaketh thus: Say not in thy heart, Who shall ascend into heaven? that is, to bring Christ down; Or who shall descend into the deep? that is, to bring up Christ again from the dead. But what saith the scripture? The word is nigh thee, even in thy mouth, and in thy heart. This is the word of faith, which we preach. For if thou confess with thy mouth the Lord Jesus, and believe in thy heart that God hath raised him up from the dead, thou shalt be saved.\par
\switchcolumn*

\LA
\begin{Resp}\Rsign Cum perambuláret Dóminus juxta mare Galilǽæ, vidit Petrum et Andréam rétia mitténtes in mare, et vocávit eos, dicens: \Star{} Veníte post me, fáciam vos fíeri piscatóres hóminum.\end{Resp}
\begin{Resp}\Vsign Erant enim piscatóres, et ait illis.\end{Resp}
\begin{Resp}\Rsign Veníte post me, fáciam vos fíeri piscatóres hóminum.\end{Resp}
\switchcolumn
\EN
\begin{Resp}\Rsign The Lord, walking by the Sea of Galilee, saw Peter and Andrew casting their nets into the sea, and He called them saying: \Star{} Follow Me, and I will make you fishers of men.\end{Resp}
\begin{Resp}\Vsign For they were fishers, and He saith unto them:\end{Resp}
\begin{Resp}\Rsign Follow Me, and I will make you fishers of men.\end{Resp}
\switchcolumn*

\LA
\LessonHead{Lectio II}
\switchcolumn
\EN
\LessonHead{Lesson II}
\switchcolumn*

\LA
\Cite{Rom 10:10-15}
\switchcolumn
\EN
\Cite{Rom 10:10-15}
\switchcolumn*

\LA
\noindent Corde enim créditur ad justítiam: ore autem conféssio fit ad salútem. Dicit enim Scriptúra: Omnis qui credit in illum, non confundétur. Non enim est distínctio Judǽi et Græci: nam idem Dóminus ómnium, dives in omnes qui ínvocant illum. Omnis enim quicúmque invocáverit nomen Dómini, salvus erit. Quómodo ergo invocábunt, in quem non credidérunt? aut quómodo credent ei, quem non audiérunt? quómodo autem áudient sine prædicánte? Quómodo vero prædicábunt nisi mittántur? sicut scriptum est: Quam speciósi pedes evangelizántium pacem, evangelizántium bona!\par
\switchcolumn
\EN
\noindent For, with the heart, we believe unto justice; but, with the mouth, confession is made unto salvation. For the scripture saith: Whosoever believeth in him, shall not be confounded. For there is no distinction of the Jew and the Greek: for the same is Lord over all, rich unto all that call upon him. For whosoever shall call upon the name of the Lord, shall be saved. How then shall they call on him, in whom they have not believed? Or how shall they believe him, of whom they have not heard? And how shall they hear, without a preacher? And how shall they preach unless they be sent, as it is written: How beautiful are the feet of them that preach the gospel of peace, of them that bring glad tidings of good things!\par
\switchcolumn*

\LA
\begin{Resp}\Rsign Mox ut vocem Dómini prædicántis audívit beátus Andréas, relíctis rétibus, quorum usu actúque vivébat, \Star{} Ætérnæ vitæ secútus est prǽmia largiéntem.\end{Resp}
\begin{Resp}\Vsign Hic est qui pro amóre Christi pepéndit in cruce, et pro lege ejus sustínuit passiónem.\end{Resp}
\begin{Resp}\Rsign Ætérnæ vitæ secútus est prǽmia largiéntem.\end{Resp}
\switchcolumn
\EN
\begin{Resp}\Rsign As soon as the blessed Andrew heard the voice of the Lord calling him, he left his nets, by the exercise and use whereof he lived; \Star{} And followed Him Who giveth life everlasting.\end{Resp}
\begin{Resp}\Vsign This is that disciple who for the love of Christ hung upon the cross, and suffered for the law of his God.\end{Resp}
\begin{Resp}\Rsign And followed Him Who giveth life everlasting.\end{Resp}
\switchcolumn*

\LA
\LessonHead{Lectio III}
\switchcolumn
\EN
\LessonHead{Lesson III}
\switchcolumn*

\LA
\Cite{Rom 10:16-21}
\switchcolumn
\EN
\Cite{Rom 10:16-21}
\switchcolumn*

\LA
\noindent Sed non omnes obédiunt Evangélio. Isaías enim dicit: Dómine, quis crédidit audítui nostro? Ergo fides ex audítu, audítus autem per verbum Christi. Sed dico: Numquid non audiérunt? Et quidem in omnem terram exívit sonus eórum, et in fines orbis terræ verba eórum. Sed dico: Numquid Israël non cognóvit? Primus Móyses dicit: Ego ad æmulatiónem vos addúcam in non gentem: in gentem insipiéntem, in iram vos mittam. Isaías autem audet, et dicit: Invéntus sum a non quæréntibus me: palam appárui iis, qui me non interrogábant. Ad Israël autem dicit: Tota die expándi manus meas ad pópulum non credéntem, et contradicéntem.\par
\switchcolumn
\EN
\noindent But all do not obey the gospel. For Isaias saith: Lord, who hath believed our report? Faith then cometh by hearing; and hearing by the word of Christ. But I say: Have they not heard? Yes, verily, their sound hath gone forth into all the earth, and their words unto the ends of the whole world. But I say: Hath not Israel known? First, Moses saith: I will provoke you to jealousy by that which is not a nation; by a foolish nation I will anger you. But Isaias is bold, and saith: I was found by them that did not seek me: I appeared openly to them that asked not after me. But to Israel he saith: All the day long have I spread my hands to a people that believeth not, and contradicteth me.\par
\switchcolumn*

\LA
\begin{Resp}\Rsign Doctor bonus et amícus Dei Andréas dúcitur ad crucem, quam a longe aspíciens dixit: Salve, crux, \Star{} Súscipe discípulum ejus, qui pepéndit in te magíster meus Christus.\end{Resp}
\begin{Resp}\Vsign Salve, crux, quæ in córpore Christi dedicáta es, et ex membris ejus tamquam margarítis ornáta.\end{Resp}
\begin{Resp}\Rsign Súscipe discípulum ejus, qui pepéndit in te magíster meus Christus.\end{Resp}
\begin{Resp}\Vsign Glória Patri, et Fílio, \Star{} et Spirítui Sancto.\end{Resp}
\begin{Resp}\Rsign Súscipe discípulum ejus, qui pepéndit in te magíster meus Christus.\end{Resp}
\switchcolumn
\EN
\begin{Resp}\Rsign Andrew the good teacher, the friend of God, was led to the cross, and when he saw it afar off, he said: God bless thee, O cross; \Star{} Welcome to the follower of Him That hung on thee, even my Master Christ.\end{Resp}
\begin{Resp}\Vsign God bless thee, O cross, thou art hallowed by the Body of Christ; His Members make thee goodly as with pearls.\end{Resp}
\begin{Resp}\Rsign Be welcome to the follower of Him That hung on thee, even my Master Christ.\end{Resp}
\begin{Resp}\Vsign Glory be to the Father, and to the Son, \Star{} and to the Holy Ghost.\end{Resp}
\begin{Resp}\Rsign Be welcome to the follower of Him That hung on thee, even my Master Christ.\end{Resp}
\switchcolumn*

\LA
\LessonHead{Lectio IV}
\switchcolumn
\EN
\LessonHead{Lesson IV}
\switchcolumn*

\LA
\noindent Andréas Apóstolus, Bethsáidæ natus, qui est Galilǽæ vicus, frater Petri, discípulus Joánnis Baptistæ, cum eum de Christo dicéntem audísset: Ecce Agnus Dei; secútus Jesum, fratrem quoque suum ad eúndem perdúxit. Cum póstea una cum fratre piscarétur in mari Galilǽæ, ambo a prætereúnte Christo Dómino ante álios Apóstolos vocáti illis verbis: Veníte post me, fáciam vos fíeri piscatóres hóminum; nullam interponéntes moram, et relíctis rétibus, secúti sunt eum. Post cujus passiónem et resurrectiónem, Andréas, cum in Scýthiam Európæ, quæ ei província ad Christi fidem disseminándam obtígerat, venísset, deínde Epírum ac Thráciam peragrásset; doctrína et miráculis innumerábiles hómines ad Christum convértit. Post, Patras Achájæ proféctus, et in ea urbe plúrimis ad veritátem evangélicam perdúctis, Ægéam procónsulem, prædicatióni evangélicæ resisténtem, libérrime increpávit, quod, qui judex hóminum habéri vellet, Christum Deum ómnium júdicem, a dæmónibus elúsus, non agnósceret.\par
\switchcolumn
\EN
\noindent The Apostle Andrew was born at Bethsaida, a town of Galilee, and was the brother of Peter. He was a disciple of John the Baptist, and heard him say of Christ, Behold the Lamb of God, (John i. 35-37, 40,) whereupon he immediately followed Jesus, bringing his brother also with him. Some while after, they were both fishing in the Sea of Galilee, and the Lord Christ, going by, called them both, before any other of the Apostles, in the words, Follow Me, and I will make you fishers of men. They made no delay, but left their nets, and followed Him. \textasciitilde{}(Matth. iv. 18-20.) After the death and Resurrection of Christ, Andrew was allotted Scythia as the province of his preaching, and, after labouring there, he went through Epirus and Thrace, where he turned vast multitudes to Christ by his teaching and miracles. Finally he went to Patras in Achaia, and there also he brought many to the knowledge of Gospel truth. Aegeas the Pro-consul resisted the preaching of the Gospel, and the Apostle freely rebuked him, bidding him know that while he held himself a judge of his fellow men, he was himself hindered by devils from knowing Christ our God, the Judge of all.\par
\switchcolumn*

\LA
\begin{Resp}\Rsign Homo Dei ducebátur ut crucifígerent eum: pópulus autem clamábat voce magna, dicens: \Star{} Innocens ejus sanguis sine causa damnátur.\end{Resp}
\begin{Resp}\Vsign Cumque dúcerent eum ut crucifigerétur, factus est concúrsus populórum clamántium et dicéntium.\end{Resp}
\begin{Resp}\Rsign Innocens ejus sanguis sine causa damnátur.\end{Resp}
\switchcolumn
\EN
\begin{Resp}\Rsign The man of God was led to be crucified, and the people cried with a loud voice, saying: \Star{} The innocent blood of this just person is condemned without a cause.\end{Resp}
\begin{Resp}\Vsign And when they led him out to crucify him, all the people ran together and cried, saying\end{Resp}
\begin{Resp}\Rsign The innocent blood of this just person is condemned without a cause.\end{Resp}
\switchcolumn*

\LA
\LessonHead{Lectio V}
\switchcolumn
\EN
\LessonHead{Lesson V}
\switchcolumn*

\LA
\noindent Tum Ægéas irátus, Désine, inquit, Christum jactáre, cui simília verba nihil profuérunt, quóminus a Judǽis crucifigerétur. Andréam vero de Christo nihilóminus líbere prædicántem quod pro salúte humáni géneris se crucifigéndum obtulísset, ímpia oratióne interpéllat, ac demum hortátur, ut sibi cónsulens, diis velit immoláre. Cui Andréas: Ego omnipoténti Deo, qui unus et verus est, ímmolo cotídie, non taurórum carnes, nec hircórum sánguinem, sed immaculátum Agnum in altári; cujus carnem posteáquam omnis pópulus credéntium manducáverit, Agnus, qui sacrificátus est, ínteger persevérat et vivus. Quam ob rem ira accénsus Ægéas, jubet eum in cárcerem detrúdi unde pópulus Andréam fácile liberásset, nisi ipse sedásset multitúdinem, veheméntius rogans, ne se ad optatíssimam martýrii corónam properántem impedírent.\par
\switchcolumn
\EN
\noindent Then Egeas, being angry, answered him, Boast no more of this thy Christ. He spake words even such as thine, but they availed Him not, and He was crucified by the Jews. Whereto Andrew boldly answered that Christ had given Himself up to die for man's salvation; but the Pro-consul blasphemously interrupted him, and bade him look to himself, and sacrifice to the gods. Then said Andrew, We have an altar, whereon day by day I offer up to God, the Almighty, the One, and the True, not the flesh of bulls nor the blood of goats, but a Lamb without spot and when all they that believe have eaten of the Flesh Thereof, the Lamb That was slain abideth whole and liveth. Then Aegeas being filled with wrath, bound the Apostle in prison. Now, the people would have delivered him, but he himself calmed the multitude, and earnestly besought them not to take away from him the crown of martyrdom, for which he longed and which was now drawing near.\par
\switchcolumn*

\LA
\begin{Resp}\Rsign O bona crux, quæ decórem et pulchritúdinem de membris Dómini suscepísti; áccipe me ab homínibus, et redde me magístro meo: \Star{} Ut per te me recípiat, qui per te me redémit.\end{Resp}
\begin{Resp}\Vsign Beátus Andréas expánsis mánibus ad cælum orábat, dicens: Salva me, bona crux.\end{Resp}
\begin{Resp}\Rsign Ut per te me recípiat, qui per te me redémit.\end{Resp}
\switchcolumn
\EN
\begin{Resp}\Rsign O precious cross, which the Members of my Lord have made so fair and goodly, welcome me from among men, and join me again to my Master, \Star{} That, as by thee He redeemed me, so by thee also He may take me unto Himself.\end{Resp}
\begin{Resp}\Vsign The blessed Andrew stretched forth his hands to heaven and prayed, saying: Precious cross, be my salvation,\end{Resp}
\begin{Resp}\Rsign That, as by thee He redeemed me, so by thee also He may take me unto Himself.\end{Resp}
\switchcolumn*

\LA
\LessonHead{Lectio VI}
\switchcolumn
\EN
\LessonHead{Lesson VI}
\switchcolumn*

\LA
\noindent Igitur paulo post in tribúnal prodúctum, cum Ægéas crucis extolléntem mystéria sibíque suam impietátem exprobrántem diútius ferre non posset, in crucem tolli, et Christi mortem imitári jussit. Addúctus Andréas ad locum martýrii, cum crucem vidísset, longe exclamáre cœpit: O bona crux, quæ decórem ex membris Dómini suscepísti, diu desideráta, sollícite amáta, sine intermissióne quæsíta, et aliquándo cupiénti ánimo præparáta: áccipe me ab homínibus, et redde me magístro meo; ut per te me recípiat, qui per te me redémit. Itaque cruci affíxus est: in qua bíduum vivus pendens, et Christi fidem prædicáre nunquam intermíttens, ad eum migrávit, cujus mortis similitúdinem concupíerat. Quæ ómnia presbýteri et diáconi Achájæ, qui ejus passiónem scripsérunt, se ita ut commemoráta sunt, audísse et vidísse testántur. Ejus ossa primum Constantíno imperatóre Constantinópolim, deínde Amálphim transláta sunt. Caput, Pio secúndo Pontífice, Romam allátum, in basílica sancti Petri collocátum est.\par
\switchcolumn
\EN
\noindent Come a short while after, he was brought before the judgment-seat, where he extolled the mystery of the cross, and rebuked Aegeas for his ungodliness. Then Aegeas could bear with him no longer, but commanded him to be crucified, in imitation of Christ. Andrew, then, was led to the place of martyrdom, and, as soon as he came in sight of the cross, he cried out, O precious cross, which the Members of my Lord have made so goodly, how long have I desired thee! how warmly have I loved thee! how constantly have I sought thee! And, now that thou art come to me, how is my soul drawn to thee! Welcome me from among men, and join me again to my Master, that as by thee He redeemed me, so by thee also He may take me unto Himself. So he was fastened to the cross, whereon he hung living for two days, during which time he ceased not to preach the faith of Christ, and, finally, passed into the Presence of Him the likeness of Whose death he had loved so well. All the above particulars of his last sufferings were written by the Priests and Deacons of Achaia, who bear witness to them of their own knowledge. Under the Emperor Constantine the bones of the Apostle were first taken to Constantinople, whence they were afterwards brought to Amalfi. In the Pontificate of Pope Pius II his head was carried to Rome, where it is kept in the Basilica of St. Peter.\par
\switchcolumn*

\LA
\begin{Resp}\Rsign Expándi manus meas tota die in cruce ad pópulum non credéntem, sed contradicéntem mihi: \Star{} Qui ámbulant vias non bonas, sed post peccáta sua.\end{Resp}
\begin{Resp}\Vsign Deus ultiónum Dóminus, Deus ultiónum líbere egit: exaltáre, qui júdicas terram; redde retributiónem supérbis.\end{Resp}
\begin{Resp}\Rsign Qui ámbulant vias non bonas, sed post peccáta sua.\end{Resp}
\begin{Resp}\Vsign Glória Patri, et Fílio, \Star{} et Spirítui Sancto.\end{Resp}
\begin{Resp}\Rsign Qui ámbulant vias non bonas, sed post peccáta sua.\end{Resp}
\switchcolumn
\EN
\begin{Resp}\Rsign All day long I have stretched forth my hands upon the cross unto a disobedient and gainsaying people; \Star{} Which walketh in a way that is not good, but after their own sins.\end{Resp}
\begin{Resp}\Vsign The Lord God to Whom vengeance belongeth, the God to Whom vengeance belongeth, hath shown Himself: lift up thyself, Thou Judge of the earth, render a reward to the proud.\end{Resp}
\begin{Resp}\Rsign Which walketh in a way that is not good, but after their own sins.\end{Resp}
\begin{Resp}\Vsign Glory be to the Father, and to the Son, \Star{} and to the Holy Ghost.\end{Resp}
\begin{Resp}\Rsign Which walketh in a way that is not good, but after their own sins.\end{Resp}
\switchcolumn*

\LA
\LessonHead{Lectio VII}
\switchcolumn
\EN
\LessonHead{Lesson VII}
\switchcolumn*

\LA
\LessonTitle{Léctio sancti Evangélii secúndum Matthǽum}
\switchcolumn
\EN
\LessonTitle{From the Holy Gospel according to Matthew}
\switchcolumn*

\LA
\Cite{Matt 4:18-22}
\switchcolumn
\EN
\Cite{Matt 4:18-22}
\switchcolumn*

\LA
\noindent In illo témpore: Ambulans Jesus juxta mare Galilǽæ, vidit duos fratres, Simónem, qui vocátur Petrus, et Andréam fratrem ejus, mitténtes rete in mare. Et réliqua.\par
\switchcolumn
\EN
\noindent At that time, Jesus, walking by the sea of Galilee, saw two brethren, Simon who is called Peter, and Andrew his brother, casting a net into the sea (for they were fishers). And so on.\par
\switchcolumn*

\LA
\LessonTitle{Homilía sancti Gregórii Papæ}
\switchcolumn
\EN
\LessonTitle{Homily by Pope St. Gregory (the Great.)}
\switchcolumn*

\LA
\Source{Homilia 5 in Evangelia}
\switchcolumn
\EN
\Source{15th on the Gospels}
\switchcolumn*

\LA
\noindent Audístis, fratres caríssimi, quia ad uníus jussiónis vocem Petrus et Andréas relíctis rétibus secúti sunt Redemptórem. Nulla vero hunc fácere adhuc mirácula víderant, nihil ab eo de prǽmio ætérnæ retributiónis audíerant: et tamen ad unum Dómini præcéptum hoc quod possidére videbántur, oblíti sunt. Quanta nos ejus mirácula vidémus, quot flagéllis afflígimur, quantis minárum asperitátibus deterrémur, et tamen vocántem sequi contémnimus.\par
\switchcolumn
\EN
\noindent Dearly beloved brethren, ye hear how that Peter and Andrew, having once heard the Lord call them, left their nets, and followed their Saviour. As yet they had seen none of His miracles, as yet they had received no promise of their exceeding and eternal reward; nevertheless, at one word of the Lord they forgot all those things which they seemed to have. We have seen many of His miracles; we have received many of His gracious chastenings; many times hath He warned us of the wrath to come and yet Christ calleth and we do not follow.\par
\switchcolumn*

\LA
\begin{Resp}\Rsign Orávit sanctus Andréas, dum respíceret in cælum, et voce magna clamávit et dixit: Tu es Deus meus, quem vidi: ne me patiáris ab ímpio júdice depóni: \Star{} Quia virtútem sanctæ crucis agnóvi.\end{Resp}
\begin{Resp}\Vsign Tu es magíster meus Christus, quem diléxi, quem cognóvi, quem conféssus sum: tantúmmodo in ista voce exáudi me.\end{Resp}
\begin{Resp}\Rsign Quia virtútem sanctæ crucis agnóvi.\end{Resp}
\switchcolumn
\EN
\begin{Resp}\Rsign The holy Andrew lifted up his eyes to heaven, and prayed, and cried with a loud voice, and said: Thou art my God, Whom I have seen; suffer not the unjust judge to take me down from the cross \Star{} For now I know what the power of thy holy Cross is.\end{Resp}
\begin{Resp}\Vsign Thou art Christ my Master, Whom I have loved, Whom I have known, Whom I have confessed in this thing hear me.\end{Resp}
\begin{Resp}\Rsign For now I know what the power of thy holy Cross is.\end{Resp}
\switchcolumn*

\LA
\LessonHead{Lectio VIII}
\switchcolumn
\EN
\LessonHead{Lesson VIII}
\switchcolumn*

\LA
\noindent In cælo jam sedet, qui de conversióne nos ádmonet; jam jugo fídei colla géntium súbdidit, jam mundi glóriam stravit, jam ruínis ejus crebrescéntibus, distrícti sui judícii diem propinquántem denúntiat: et tamen supérba mens nostra adhuc non vult hoc sponte desérere, quod cotídie perdit invíta. Quid ergo, fratres caríssimi, quid in ejus judício dictúri sumus, qui ab amóre præséntis sǽculi nec præcéptis fléctimur, nec verbéribus emendámur?\par
\switchcolumn
\EN
\noindent He who calleth us to be converted is now enthroned in heaven; He hath broken the necks of the Gentiles to the yoke of the faith, He hath laid low the glory of the world, and the wrecks thereof, falling ever more and more to decay, do preach unto us that the coming of that day when He is to be revealed as our Judge is drawing nigh and yet, so stubborn is our mind, that we will not yet freely abandon that which, will we, nill we, we lose day by day. Dearly beloved brethren, what shall we answer at His judgment-seat, we whom no lessons can persuade, and no stripes can break of the love of this present world?\par
\switchcolumn*

\LA
\begin{Resp}\Rsign Videns crucem Andréas exclamávit, dicens: O crux admirábilis, o crux desiderábilis, o crux quæ per totum mundum rútilas: \Star{} Súscipe discípulum Christi, ac per te me recípiat, qui per te móriens me redémit.\end{Resp}
\begin{Resp}\Vsign O bona crux, quæ decórem et pulchritúdinem de membris Dómini suscepísti.\end{Resp}
\begin{Resp}\Rsign Súscipe discípulum Christi, ac per te me recípiat, qui per te móriens me redémit.\end{Resp}
\begin{Resp}\Vsign Glória Patri, et Fílio, \Star{} et Spirítui Sancto.\end{Resp}
\begin{Resp}\Rsign Súscipe discípulum Christi, ac per te me recípiat, qui per te móriens me redémit.\end{Resp}
\switchcolumn
\EN
\begin{Resp}\Rsign When Andrew saw the cross he cried, saying: How wonderful art thou, O cross! O cross, how loveable art thou! O cross, thy bright beams enlighten the darkness of the whole world! \Star{} Welcome a follower of Jesus, that, as by thee He died to redeem me, so by thee also He may take me unto Himself.\end{Resp}
\begin{Resp}\Vsign O precious cross, which the Members of my Lord have made so fair and goodly,\end{Resp}
\begin{Resp}\Rsign Welcome a follower of Jesus, that, as by thee He died to redeem me, so by thee also He may take me unto Himself.\end{Resp}
\begin{Resp}\Vsign Glory be to the Father, and to the Son, \Star{} and to the Holy Ghost.\end{Resp}
\begin{Resp}\Rsign Welcome a follower of Jesus, that, as by thee He died to redeem me, so by thee also He may take me unto Himself.\end{Resp}
\switchcolumn*

\LA
\LessonHead{Lectio IX}
\switchcolumn
\EN
\LessonHead{Lesson IX}
\switchcolumn*

\LA
\noindent Sed fortásse áliquis tácitis sibi cogitatiónibus dicat: Ad vocem Domínicam utérque iste piscátor quid, aut quantum dimísit, qui pene nihil hábuit? Sed hac in re, fratres caríssimi, afféctum debémus pótius pensáre quam censum. Multum relíquit, qui sibi nihil retínuit: multum relíquit, qui quantúmlibet parum, totum deséruit. Certe nos et hábita cum amóre possidémus, et ea, quæ mínime habémus, ex desidério quǽrimus. Multum ergo Petrus et Andréas dimísit, quando utérque étiam desidéria habéndi derelíquit.\par
\switchcolumn
\EN
\noindent Perhaps one perchance will ask in his heart, what Peter or Andrew had to lose by obeying the call of the Lord? Dearly beloved brethren, we must consider here rather the intention than the loss incurred by this obedience. He that keepeth nothing for himself, giveth up much; he that sacrificeth his all, sacrificed what is to him a great deal. Beyond doubt, we cling to whatever we have, and what we have least, that we desire most. Peter and Andrew therefore gave up much when they gave up even the desire of possessing anything.\par
\switchcolumn*

\LA
\TeDeum{Te Deum}
\switchcolumn
\EN
\TeDeum{Te Deum}
\switchcolumn*

\end{paracol}

\SeasonTitle{Mensis December}{December}

\addcontentsline{toc}{section}{Mensis December \textit{— December}}

\par\needspace{10\baselineskip}\addvspace{1.2em}{\centering\small\textsc{Die 2 Decembris} · \textsc{2 December}\par}
\Office{S. Bibianæ Virginis et Martyris}{St. Bibiana, Virgin and Martyr}{Semiduplex}
\addcontentsline{toc}{subsection}{Die 2 Decembris · S. Bibianæ Virginis et Martyris \textit{— St. Bibiana, Virgin and Martyr}}
\begin{paracol}{2}
\LA
\RefLine{Lectiones I–III: de Scriptura occurrente.}
\switchcolumn
\EN
\RefLine{Lessons I–III: from the Scripture of the day.}
\switchcolumn*

\LA
\RefLine{Lectiones VII–IX: ex Commúni Unius Virginis et Martyris (C6), in 2 loco.}
\switchcolumn
\EN
\RefLine{Lessons VII–IX: from the Common of a Virgin Martyr (C6), set 2.}
\switchcolumn*

\LA
\LessonHead{Lectio IV}
\switchcolumn
\EN
\LessonHead{Lesson IV}
\switchcolumn*

\LA
\noindent Bibiana virgo Romana, nobili genere nata, christiana fide nobilior fuit. Ejus enim pater Flavianus sub Juliano Apostata impiissimo tyranno expræfectus, servilibusque notis compunctus, ad aquas Taurinas deportatus, martyr occubuit. Mater Dafrosa, et filiæ primum conclusæ domi, ut inedia conficerentur; mox relegata mater extra Urbem capite plexa est. Mortuis autem piis parentibus, Bibiana cum sorore sua Demetria bonis omnibus exspoliatur; Apronianus Urbis prætor, pecuniis inhians, sorores persequitur; quas humana prorsus ope destitutas, Deo mirabiliter, qui dat escam esurientibus, enutriente, cum vivaciores vegetioresque conspexisset, vehementer est admiratus.\par
\switchcolumn
\EN
\noindent Bibiana was a Roman maiden, distinguished on account of the nobility of her family, but now far more distinguished for her confession of Christ. In the reign of the foul tyrant, Julian the Apostate, her father Flavian, although he was an ex-Praefect, was branded as a slave and banished to Acquapendente, not far from Rome, where he soon died a martyr for his faith. His wife, Dafrosa, and his two daughters, Bibiana and Demetria, were first imprisoned in their own house, with the idea of starving them to death; but the mother was afterwards taken outside the city and beheaded. Bibiana and her sister Demetria, after the death of their holy parents, were stripped of all they had in the world. Apronianus, Praetor of the city, who hankered after their property, continued to persecute them, but although they were destitute of all human support, God, Who giveth bread to the hungry, fed them, and kept them in health, life, and strength, to the wonder of their enemies.\par
\switchcolumn*

\LA
\begin{Resp}\Rsign Propter veritátem, et mansuetúdinem, et justítiam: \Star{} Et dedúcet te mirabíliter déxtera tua.\end{Resp}
\begin{Resp}\Vsign Spécie tua et pulchritúdine tua inténde, próspere procéde, et regna.\end{Resp}
\begin{Resp}\Rsign Et dedúcet te mirabíliter déxtera tua.\end{Resp}
\switchcolumn
\EN
\begin{Resp}\Rsign Because of truth, and meekness, and righteousness; \Star{} And thy right hand shall lead thee wonderfully.\end{Resp}
\begin{Resp}\Vsign In thy comeliness, and thy beauty, go forward, fare prosperously, and reign.\end{Resp}
\begin{Resp}\Rsign And thy right hand shall lead thee wonderfully.\end{Resp}
\switchcolumn*

\LA
\LessonHead{Lectio V}
\switchcolumn
\EN
\LessonHead{Lesson V}
\switchcolumn*

\LA
\noindent Suadet nihilominus Apronianus, ut venerentur deos Gentium, amissas ideo opes, imperatoris gratiam, præclarissimas nuptias consecuturæ. Si secus fecerint, minatur carceres, virgas, secures. At illæ neque blanditiis, neque minis a recta fide declinantes, paratæ potius mori, quam fœdari moribus ethnicorum, prætoris impietatem constantissime detestantur. Quare Demetria ob oculos Bibianæ repente corruens, obiit in Domino: Bibiana Rufinæ mulieri vaferrimæ seducenda traditur, quæ ab incunabulis edocta, christianas leges, et illibatum servare virginitatis florem, se ipsa fortior feminæ superavit insidias, et prætoris astus delusit.\par
\switchcolumn
\EN
\noindent Apronianus then attacked them, to make them worship the gods of the Gentiles, and promised them the restoration of their property, the favour of the Emperor, and a great marriage for each of them, if they would give way, and, on the other hand, imprisonment, stripes, and death. But neither promises nor threats availed, for they remained firm in the faith, being resolved rather to die than to pollute themselves by doing according to the deeds of the heathen; and, as for the iniquity of the Praetor, they loathed it continually. At length the strength of Demetria gave way, and she fell down suddenly, and died in the Lord, before the eyes of her sister Bibiana. Then Bibiana was put into the hands of an artful woman named Rufina, to seduce her if possible; but she had known the law of Christ from her childhood, and kept the lily of her purity undefiled, triumphing over the efforts of that vile person, and disappointing the lust of the Praetor.\par
\switchcolumn*

\LA
\begin{Resp}\Rsign Dilexísti justítiam, et odísti iniquitátem: \Star{} Proptérea unxit te Deus, Deus tuus, óleo lætítiæ.\end{Resp}
\begin{Resp}\Vsign Propter veritátem, et mansuetúdinem, et justítiam.\end{Resp}
\begin{Resp}\Rsign Proptérea unxit te Deus, Deus tuus, óleo lætítiæ.\end{Resp}
\switchcolumn
\EN
\begin{Resp}\Rsign Thou hast loved righteousness, and hated iniquity; \Star{} Therefore God, thy God, hath anointed thee with the oil of gladness.\end{Resp}
\begin{Resp}\Vsign Because of truth, and meekness, and righteousness.\end{Resp}
\begin{Resp}\Rsign Therefore God, thy God, hath anointed thee with the oil of gladness.\end{Resp}
\switchcolumn*

\LA
\LessonHead{Lectio VI}
\switchcolumn
\EN
\LessonHead{Lesson VI}
\switchcolumn*

\LA
\noindent Nihil autem proficiente Rufina, quæ, præter dolosa verba, illam quotidie verberibus affligebat, ut de sancto proposito dimoveret, spe sua frustratus prætor, accensus ira, quod in Bibiana perdidisset operam, a lictoribus eam denudari, vinctisque manibus columnæ alligari, eamque plumbatis cædi jubet, donec efflaret animam. Cujus sacrum corpus objectum canibus biduo jacuit in foro Tauri, illæsum tamen, et divinitus servatum; quod deínde Joannes presbyter sepelivit noctu juxta sepulcrum sororis et matris ad palatium Licinianum, ubi usque in præsens exstat ecclesia Deo, sanctæ Bibianæ nomine, dicata, quam Urbanus octavus instauravit, sanctarum Bibianæ, Demetriæ, et Dafrosæ corporibus in ea repertis, et sub ara maxima collocatis.\par
\switchcolumn
\EN
\noindent Then, when Rufina saw that her false words availed not, she took to blows, and scourged Bibiana daily, but the saint was not staggered in her holy resolution. At last the Praetor, mad with baffled lust, when he found his labour was thrown away, ordered his lictors to strip her naked, hang her up by the hands to a pillar, and flog her to death with whips weighted with lead. When all was over, her sacred body was thrown out for the dogs to eat. It lay two days in the Forum Tauri, but the animals would not touch it; and, at last, a Priest, named John, took it, and buried it by night beside the graves of her mother and sister, near the Licinian Palace. This is the place where there is still a church, dedicated in the name of St. Bibiana. When this church was being restored by Urban VIII., the bodies of these three holy women, Bibiana, Demetria, and Dafrosa, were found, and were re-buried under the High Altar.\par
\switchcolumn*

\LA
\begin{Resp}\Rsign Afferéntur Regi vírgines post eam, próximæ ejus; \Star{} Afferéntur tibi in lætítia et exsultatióne.\end{Resp}
\begin{Resp}\Vsign Spécie tua et pulchritúdine tua; inténde, próspere procéde, et regna.\end{Resp}
\begin{Resp}\Rsign Afferéntur tibi in lætítia et exsultatióne.\end{Resp}
\begin{Resp}\Vsign Glória Patri, et Fílio, \Star{} et Spirítui Sancto.\end{Resp}
\begin{Resp}\Rsign Afferéntur tibi in lætítia et exsultatióne.\end{Resp}
\switchcolumn
\EN
\begin{Resp}\Rsign After her shall virgins be brought unto the King, her fellows \Star{} They shall be brought unto thee with gladness and rejoicing.\end{Resp}
\begin{Resp}\Vsign In thy comeliness and thy beauty, go forward, fare prosperously, and reign.\end{Resp}
\begin{Resp}\Rsign They shall be brought unto thee with gladness and rejoicing.\end{Resp}
\begin{Resp}\Vsign Glory be to the Father, and to the Son, \Star{} and to the Holy Ghost.\end{Resp}
\begin{Resp}\Rsign They shall be brought unto thee with gladness and rejoicing.\end{Resp}
\switchcolumn*

\end{paracol}

\PartTitle{Commune Sanctorum}{Common of the Saints}

\addcontentsline{toc}{chapter}{Commune Sanctorum \textit{— Common of the Saints}}

\Office{Commune Evangelistarum}{Commune Evangelistarum}{}
\addcontentsline{toc}{subsection}{Commune Evangelistarum \textit{— Commune Evangelistarum}}
\begin{paracol}{2}
\LA
\LessonHead{Lectio I}
\switchcolumn
\EN
\LessonHead{Lesson I}
\switchcolumn*

\LA
\LessonTitle{Incipit liber Ezechiélis Prophétæ}
\switchcolumn
\EN
\LessonTitle{Lesson from the book of Ezekiel}
\switchcolumn*

\LA
\Cite{Ezech 1:1-4}
\switchcolumn
\EN
\Cite{Ezek 1:1-4}
\switchcolumn*

\LA
\noindent Et factum est in trigésimo anno, in quarto, in quinta mensis, cum essem in médio captivórum juxta flúvium Chobar, apérti sunt cæli, et vidi visiónes Dei. In quinta mensis, ipse est annus quintus transmigratiónis regis Jóachim, factum est verbum Dómini ad Ezechiélem fílium Buzi sacerdótem, in terra Chaldæórum, secus flumen Chobar: et facta est super eum ibi manus Dómini. Et vidi, et ecce ventus túrbinis veniébat ab Aquilóne, et nubes magna, et ignis invólvens, et splendor in circúitu ejus: et de médio ejus, quasi spécies eléctri, id est, de médio ignis:\par
\switchcolumn
\EN
\noindent Now it came to pass in the thirtieth year, in the fourth month, on the fifth day of the month, when I was in the midst of the captives by the river Chobar, the heavens were opened, and I saw the visions of God. On the fifth day of the month, the same was the fifth year of the captivity of king Joachin, The word of the Lord came to Ezechiel the priest the son of Buzi in the land of the Chaldeans, by the river Chobar: and the hand of the Lord was there upon him. And I saw, and behold a whirlwind came out of the north: and a great cloud, and a fire infolding it, and brightness was about it: and out of the midst thereof, that is, out of the midst of the fire, as it were the resemblance of amber.\par
\switchcolumn*

\LA
\begin{Resp}\Rsign Ecce ego mitto vos sicut oves in médio lupórum, dicit Dóminus: \Star{} Estóte ergo prudéntes sicut serpéntes, et símplices sicut colúmbæ.\end{Resp}
\begin{Resp}\Vsign Dum lucem habétis, crédite in lucem, ut fílii lucis sitis.\end{Resp}
\begin{Resp}\Rsign Estóte ergo prudéntes sicut serpéntes, et símplices sicut colúmbæ.\end{Resp}
\switchcolumn
\EN
\begin{Resp}\Rsign Behold, I send you as sheep in the midst of wolves, said the Lord: \Star{} Be ye, therefore, wise as serpents and simple as doves.\end{Resp}
\begin{Resp}\Vsign Whilst you have the light, believe in the light, that you may be the children of light.\end{Resp}
\begin{Resp}\Rsign Be ye therefore wise as serpents and simple as doves.\end{Resp}
\switchcolumn*

\LA
\LessonHead{Lectio II}
\switchcolumn
\EN
\LessonHead{Lesson II}
\switchcolumn*

\LA
\Cite{Ezech 1:5-9}
\switchcolumn
\EN
\Cite{Ezek 1:5-9}
\switchcolumn*

\LA
\noindent Et in médio ejus similitúdo quátuor animálium, et hic aspéctus eórum, similitúdo hóminis in eis. Quátuor fácies uni, et quátuor pennæ uni. Pedes eórum, pedes recti, et planta pedis eórum quasi planta pedis vítuli: et scintíllæ quasi aspéctus æris candéntis. Et manus hóminis sub pennis eórum, in quátuor pártibus: et fácies et pennas per quátuor partes habébant, junctǽque erant pennæ eórum altérius ad álterum. Non revertebántur cum incéderent, sed unumquódque ante fáciem suam gradiebátur.\par
\switchcolumn
\EN
\noindent And in the midst thereof the likeness of four living creatures: and this was their appearance: there was the likeness of a man in them. Every one had four faces, and every one four wings. Their feet were straight feet, and the sole of their foot was like the sole of a calf's foot, and they sparkled like the appearance of glowing brass. And they had the hands of a man under their wings on their four sides: and they had faces, and wings on the four sides, And the wings of one were joined to the wings of another. They turned not when they went: but every one went straight forward.\par
\switchcolumn*

\LA
\begin{Resp}\Rsign Tóllite jugum meum super vos, dicit Dóminus, et díscite a me, quia mitis sum et húmilis corde: \Star{} Jugum enim meum suáve est, et onus meum leve.\end{Resp}
\begin{Resp}\Vsign Et inveniétis réquiem animábus vestris.\end{Resp}
\begin{Resp}\Rsign Jugum enim meum suáve est, et onus meum leve.\end{Resp}
\switchcolumn
\EN
\begin{Resp}\Rsign Take up my yoke upon you, said the Lord, and learn of me, because I am meek, and humble of heart. \Star{} For my yoke is sweet and my burden light.\end{Resp}
\begin{Resp}\Vsign And you shall find rest to your souls.\end{Resp}
\begin{Resp}\Rsign For my yoke is sweet and my burden light.\end{Resp}
\switchcolumn*

\LA
\LessonHead{Lectio III}
\switchcolumn
\EN
\LessonHead{Lesson III}
\switchcolumn*

\LA
\Cite{Ezech 1:10-12}
\switchcolumn
\EN
\Cite{Ezek 1:10-12}
\switchcolumn*

\LA
\noindent Similitúdo autem vultus eórum, fácies hóminis et fácies leónis a dextris ipsórum quátuor, fácies autem bovis a sinístris ipsórum quátuor, et fácies áquilæ désuper ipsórum quátuor. Fácies eórum et pennæ eórum exténtæ désuper: duæ pennæ singulórum jungebántur, et duæ tegébant córpora eórum: et unumquódque eórum coram fácie sua ambulábat: ubi erat ímpetus spíritus, illuc gradiebántur, nec revertebántur cum ambulárent.\par
\switchcolumn
\EN
\noindent And as for the likeness of their faces: there was the face of a man, and the face of a lion on the right side of all the four: and the face of an ox, on the left side of all the four: and the face of an eagle over all the four. And their faces, and their wings were stretched upward: two wings of every one were joined, and two covered their bodies: And every one of them went straight forward: whither the impulse of the spirit was to go, thither they went: and they turned not when they went.\par
\switchcolumn*

\LA
\begin{Resp}\Rsign Dum stetéritis ante reges et prǽsides, nolíte cogitáre quómodo aut quid loquámini: \Star{} Dábitur enim vobis in illa hora, quid loquámini.\end{Resp}
\begin{Resp}\Vsign Non enim vos estis qui loquímini: sed Spíritus Patris vestri, qui lóquitur in vobis.\end{Resp}
\begin{Resp}\Rsign Dábitur enim vobis in illa hora, quid loquámini.\end{Resp}
\begin{Resp}\Vsign Glória Patri, et Fílio, \Star{} et Spirítui Sancto.\end{Resp}
\begin{Resp}\Rsign Dábitur enim vobis in illa hora, quid loquámini.\end{Resp}
\switchcolumn
\EN
\begin{Resp}\Rsign But when they shall deliver you up to the judges, take no thought how or what to speak: \Star{} for it shall be given you in that hour what to speak:\end{Resp}
\begin{Resp}\Vsign For it is not you that speak, but the spirit of your Father that speaketh in you.\end{Resp}
\begin{Resp}\Rsign For it shall be given you in that hour what to speak.\end{Resp}
\begin{Resp}\Vsign Glory be to the Father, and to the Son, \Star{} and to the Holy Ghost.\end{Resp}
\begin{Resp}\Rsign For it shall be given you in that hour what to speak.\end{Resp}
\switchcolumn*

\LA
\LessonHead{Lectio VII}
\switchcolumn
\EN
\LessonHead{Lesson VII}
\switchcolumn*

\LA
\LessonTitle{Léctio sancti Evangélii secúndum Lucam}
\switchcolumn
\EN
\LessonTitle{From the Holy Gospel according to Luke}
\switchcolumn*

\LA
\Cite{Luc 10:1-9}
\switchcolumn
\EN
\Cite{Luke 10:1-9}
\switchcolumn*

\LA
\noindent In illo témpore: Designávit Dóminus et álios septuagínta duos: et misit illos binos ante fáciem suam, in omnem civitátem et locum, quo erat ipse ventúrus. Et réliqua.\par
\switchcolumn
\EN
\noindent At that time, the Lord appointed other seventy-two also, and sent them two and two before His face into every city and place, whither He Himself would come. And so on.\par
\switchcolumn*

\LA
\LessonTitle{Homilía sancti Gregórii Papæ}
\switchcolumn
\EN
\LessonTitle{Homily by Pope St. Gregory (the Great.)}
\switchcolumn*

\LA
\Source{Homilia 17 in Evangelia}
\switchcolumn
\EN
\Source{17th Homily on the Gospels}
\switchcolumn*

\LA
\noindent Dóminus et Salvátor noster, fratres caríssimi, aliquándo nos sermónibus, aliquándo vero opéribus ádmonet. Ipsa étenim facta ejus præcépta sunt: quia dum áliquid tácitus facit, quid ágere debeámus innotéscit. Ecce enim binos in prædicatiónem discípulos mittit: quia duo sunt præcépta caritátis, Dei vidélicet amor, et próximi: et minus quam inter duos cáritas habéri non potest. Nemo enim próprie ad semetípsum habére caritátem dícitur: sed diléctio in álterum tendit, ut cáritas esse possit.\par
\switchcolumn
\EN
\noindent Dearly beloved brethren, our Lord and Saviour doth sometimes admonish us by words, and sometimes by works. Yea, His very works do themselves teach us for that which He doth silently His example still moveth us to copy. Behold how He sendeth forth His disciples to preach by two and two since there are two commandments to love, that is, a commandment to love God, and a commandment to love our neighbour and where there are not two, the one, being alone, hath not whereon to do the Lord's commandment. And no man can properly be said to love himself: for love tendeth outward toward our neighbour, if it be the love whereto the Gospel doth oblige us.\par
\switchcolumn*

\LA
\begin{Resp}\Rsign Isti sunt qui vivéntes in carne, plantavérunt Ecclésiam sánguine suo: \Star{} Cálicem Dómini bibérunt, et amíci Dei facti sunt.\end{Resp}
\begin{Resp}\Vsign In omnem terram exívit sonus eórum, et in fines orbis terræ verba eórum.\end{Resp}
\begin{Resp}\Rsign Cálicem Dómini bibérunt, et amíci Dei facti sunt.\end{Resp}
\switchcolumn
\EN
\begin{Resp}\Rsign These are they who while yet they lived in the flesh, planted the Church in their own blood; \Star{} They drank of the Lord's cup, and became the friends of God.\end{Resp}
\begin{Resp}\Vsign Their sound is gone out through all the earth, and their words to the ends of the world.\end{Resp}
\begin{Resp}\Rsign They drank of the Lord's cup, and became the friends of God.\end{Resp}
\switchcolumn*

\LA
\LessonHead{Lectio VIII}
\switchcolumn
\EN
\LessonHead{Lesson VIII}
\switchcolumn*

\LA
\noindent Ecce enim binos ad prædicándum discípulos Dóminus mittit: quátenus hoc nobis tácitus ínnuat, quia qui caritátem erga álterum non habet, prædicatiónis offícium suscípere nullátenus debet. Bene autem dícitur, quia misit eos ante fáciem suam in omnem civitátem et locum, quo erat ipse ventúrus. Prædicatóres enim suos Dóminus séquitur: quia prædicátio prǽvenit, et tunc ad mentis nostræ habitáculum Dóminus venit, quando verba exhortatiónis præcúrrunt: atque per hoc véritas in mente suscípitur.\par
\switchcolumn
\EN
\noindent Behold, the Lord sendeth forth His disciples to preach by two and two and thus doing, He doth silently teach us that whosoever loveth not his neighbour, such a one it behoveth not to take upon him the office of a preacher. Well also is it said that He sent them before His face into every city and place whither He Himself would come. The Lord followeth His preachers first cometh preaching, and then the Lord Himself cometh to the house of our mind, whither the word of exhortation hath come before and so cometh the truth into our mind.\par
\switchcolumn*

\LA
\begin{Resp}\Rsign Isti sunt viri sancti, quos elégit Dóminus in caritáte non ficta, et dedit illis glóriam sempitérnam: \Star{} Quorum doctrína fulget Ecclésia, ut sole luna.\end{Resp}
\begin{Resp}\Vsign Sancti per fidem vicérunt regna: operáti sunt justítiam.\end{Resp}
\begin{Resp}\Rsign Quorum doctrína fulget Ecclésia, ut sole luna.\end{Resp}
\begin{Resp}\Vsign Glória Patri, et Fílio, \Star{} et Spirítui Sancto.\end{Resp}
\begin{Resp}\Rsign Quorum doctrína fulget Ecclésia, ut sole luna.\end{Resp}
\switchcolumn
\EN
\begin{Resp}\Rsign These men are saints, whom the Lord hath chosen in love unfeigned, and hath given them glory everlasting. These are they \Star{} By the light of whose teaching the Church is glorified, even as the moon is glorified by the light of the sun.\end{Resp}
\begin{Resp}\Vsign The saints through faith subdued kingdoms, wrought righteousness.\end{Resp}
\begin{Resp}\Rsign By the light of whose teaching the Church is glorified, even as the moon is glorified by the light of the sun.\end{Resp}
\begin{Resp}\Vsign Glory be to the Father, and to the Son, \Star{} and to the Holy Ghost.\end{Resp}
\begin{Resp}\Rsign By the light of whose teaching the Church is glorified, even as the moon is glorified by the light of the sun.\end{Resp}
\switchcolumn*

\LA
\LessonHead{Lectio IX}
\switchcolumn
\EN
\LessonHead{Lesson IX}
\switchcolumn*

\LA
\noindent Hinc namque eísdem prædicatóribus Isaías dicit: Paráte viam Dómini, rectas fácite sémitas Dei nostri. Hinc fíliis Psalmísta ait: Iter fácite ei, qui ascéndit super occásum. Super occásum namque Dóminus ascéndit: quia unde in passióne occúbuit, inde majórem suam glóriam resurgéndo manifestávit. Super occásum vidélicet ascéndit; quia mortem quam pértulit, resurgéndo calcávit. Ei ergo qui ascéndit super occásum, iter fácimus, cum nos ejus glóriam vestris méntibus prædicámus, ut eas et ipse post véniens, per amóris sui præséntiam illústret.\par
\switchcolumn
\EN
\noindent Therefore, to preachers saith Isaiah: “Prepare ye the way of the Lord, make straight an highway for our God.” (xl. 3.) And again the Psalmist saith: “Spread a path before him that rideth upon the West.” (lxvii. 4.) The Lord rideth upon the West above that from which in death He veiled His glory, hath He royally exalted that glory that excelleth, even the glory of His rising again. He rideth upon the West, Who, being risen again from the dead, is throned high above the death to which He bowed. Before Him, therefore, That rideth upon the West, we spread a path, when we set forth His glory before the eyes of your mind, to the end that He Himself may come after, and Himself enlighten the same your minds by His presence and His love.\par
\switchcolumn*

\LA
\TeDeum{Te Deum}
\switchcolumn
\EN
\TeDeum{Te Deum}
\switchcolumn*

\end{paracol}

\Office{Commune unius Summi Pontificis et Martyris}{Common of a Single Martyr Pope}{}
\addcontentsline{toc}{subsection}{Commune unius Summi Pontificis et Martyris \textit{— Common of a Single Martyr Pope}}
\begin{paracol}{2}
\LA
\LessonHead{Lectio VII}
\switchcolumn
\EN
\LessonHead{Lesson VII}
\switchcolumn*

\LA
\LessonTitle{Léctio sancti Evangélii secúndum Matthǽum}
\switchcolumn
\EN
\LessonTitle{From the Holy Gospel according to Matthew}
\switchcolumn*

\LA
\Cite{Matt 16:13-19}
\switchcolumn
\EN
\Cite{Matt 16:13-19}
\switchcolumn*

\LA
\noindent In illo témpore: Venit Jesus in partes Cæsaréæ Philíppi, et interrogábat discípulos suos, dicens: Quem dicunt hómines esse Fílium hóminis? Et réliqua.\par
\switchcolumn
\EN
\noindent At that time: When Jesus came into the coasts of Caesarea Philippi, he asked his disciples, saying, Whom do men say that I the Son of man am? And so on.\par
\switchcolumn*

\LA
\LessonTitle{Homilía sancti Leónis Papæ}
\switchcolumn
\EN
\LessonTitle{A Homily by St. Leo the Pope}
\switchcolumn*

\LA
\Source{Sermo 2 in anniversario assumpt. suæ, ante medium}
\switchcolumn
\EN
\Source{Sermo 2 in anniversario assumpt. suæ ante medium}
\switchcolumn*

\LA
\noindent Cum, sicut evangélica lectióne reserátum est, interrogásset Dóminus discípulos, quem ipsum (multis divérsa opinántibus) créderent; respondissétque beátus Petrus, dicens: Tu es Christus Fílius Dei vivi; Dóminus ait: Beátus es, Simon Bar-Jona, quia caro et sanguis non revelávit tibi, sed Pater meus, qui in cælis est: et ego dico tibi, quia tu es Petrus, et super hanc petram ædificábo Ecclésiam meam, et portæ ínferi non prævalébunt advérsus eam. Et tibi dabo claves regni cælórum: et quodcúmque ligáveris super terram, erit ligátum et in cælis: et quodcúmque sólveris super terram, erit solútum et in cælis. Manet ergo disposítio veritátis, et beátus Petrus, in accépta fortitúdine petræ persevérans, suscépta Ecclésiæ gubernácula non relíquit.\par
\switchcolumn
\EN
\noindent When the Lord, as we read in the Gospel, asked his disciples who did men, amid their diverse speculations, believe him the Son of Man to be, blessed Peter answered and said: Thou art the Christ, the Son of the living God. And the Lord answered and said unto him: Blessed art thou, Simon Bar-Jona: for flesh and blood hath not revealed it unto thee, but my Father, which is in heaven: and I say also unto thee: That thou art Peter, and upon this rock I will build my Church, and the gates of hell shall not prevail against it; and I will give unto thee the keys of the kingdom of heaven; and whatsoever thou shalt bind on earth shall be bound in heaven; and whatsoever thou shalt loose on earth shall be loosed in heaven. But the dispensation of truth perdures, and blessed Peter, persevering in the strength of the rock which he hath received, hath not relinquished the position he assumed at the helm of the Church.\par
\switchcolumn*

\LA
\begin{Resp}\Rsign Coróna áurea super caput ejus, \Star{} Expréssa signo sanctitátis, glória honóris, et opus fortitúdinis.\end{Resp}
\begin{Resp}\Vsign Quóniam prævenísti eum in benedictiónibus dulcédinis, posuísti in cápite ejus corónam de lápide pretióso.\end{Resp}
\begin{Resp}\Rsign Expréssa signo sanctitátis, glória honóris, et opus fortitúdinis.\end{Resp}
\switchcolumn
\EN
\begin{Resp}\Rsign A crown of gold upon his head, \Star{} Wherein is engraved Holiness, an ornament of honour, a costly work.\end{Resp}
\begin{Resp}\Vsign For Thou hast prevented him with the blessings of sweetness, Thou hast set a crown of precious stones upon his head.\end{Resp}
\begin{Resp}\Rsign Wherein is engraved Holiness, an ornament of honour, a costly work.\end{Resp}
\switchcolumn*

\LA
\LessonHead{Lectio VIII}
\switchcolumn
\EN
\LessonHead{Lesson VIII}
\switchcolumn*

\LA
\noindent In univérsa namque Ecclésia, Tu es Christus Fílius Dei vivi, quotídie Petrus dicit; et omnis lingua, quæ confitétur Dóminum, magistério hujus vocis imbúitur. Hæc fides diábolum vincit et captivórum ejus víncula dissólvit. Hæc érutos mundo, ínserit cælo, et portæ ínferi advérsus eam prævalére non possunt. Tanta enim divínitus soliditáte muníta est, ut eam neque hærética umquam corrúmpere právitas, nec pagána potúerit superáre perfídia. His ítaque modis, dilectíssimi, rationábile obséquio celebrétur hodiérna festívitas: ut in persóna humilitátis meæ ille intelligátur, ille honorétur, in quo et ómnium pastórum sollicitúdo, cum commendatárum sibi óvium custódia persevérat, et cujus étiam dígnitas in indígno heréde non déficit.\par
\switchcolumn
\EN
\noindent In the universal Church it is as if Peter were still saying every day: Thou art the Christ, the Son of the living God. For every tongue which confesseth the Lord is taught that confession by the teaching of Peter. This is the Faith that overcometh the devil and looseth the bonds of his prisoners. This is the Faith which maketh men free of the world and bringeth them to heaven, and the gates of hell are impotent to prevail against it. This is the rock which God hath fortified with such ramparts of salvation, that the contagion of heresy will never be able to infect it, nor idolatry and unbelief to overcome it. And therefore, dearly beloved, we celebrate today’s festival with reasonable obedience, that in my humble person he may be acknowledged and honoured who doth continue to care for all the shepherds as well as sheep entrusted unto him, and who doth lose none of his dignity even in an unworthy successor.\par
\switchcolumn*

\LA
\begin{Resp}\Rsign Hic est vere Martyr, qui pro Christi nómine sánguinem suum fudit: \Star{} Qui minas júdicum non tímuit, nec terrénæ dignitátis glóriam quæsívit, sed ad cæléstia regna pervénit.\end{Resp}
\begin{Resp}\Vsign Justum dedúxit Dóminus per vias rectas, et osténdit illi regnum Dei.\end{Resp}
\begin{Resp}\Rsign Qui minas júdicum non tímuit, nec terrénæ dignitátis glóriam quæsívit, sed ad cæléstia regna pervénit.\end{Resp}
\begin{Resp}\Vsign Glória Patri, et Fílio, \Star{} et Spirítui Sancto.\end{Resp}
\begin{Resp}\Rsign Qui minas júdicum non tímuit, nec terrénæ dignitátis glóriam quæsívit, sed ad cæléstia regna pervénit.\end{Resp}
\switchcolumn
\EN
\begin{Resp}\Rsign This is a Martyr indeed, who shed his blood for Christ's Name's sake \Star{} Who feared not for the threats of judges, nor sought to be great with the glory of this world, but pressed on unto the kingdom of heaven.\end{Resp}
\begin{Resp}\Vsign The Lord guided the righteous in right paths, and showed him the kingdom of God.\end{Resp}
\begin{Resp}\Rsign Who feared not for the threats of judges, nor sought to be great with the glory of this world, but pressed on unto the kingdom of heaven.\end{Resp}
\begin{Resp}\Vsign Glory be to the Father, and to the Son, \Star{} and to the Holy Ghost.\end{Resp}
\begin{Resp}\Rsign Who feared not for the threats of judges, nor sought to be great with the glory of this world, but pressed on unto the kingdom of heaven.\end{Resp}
\switchcolumn*

\LA
\LessonHead{Lectio IX}
\switchcolumn
\EN
\LessonHead{Lesson IX}
\switchcolumn*

\LA
\noindent Cum ergo cohortatiónes nostras áuribus vestræ sanctitátis adhibémus, ipsum vobis, cujus vice fúngimur, loqui crédite: quia et illíus vos afféctu monémus, et non áliud vobis, quam quod dócuit, prædicámus; obsecrántes, ut succíncti lumbos mentis vestræ, castam et sóbriam vitam in Dei timóre ducátis. Coróna mea, sicut Apóstolus ait, et gáudium vos estis, si fides vestra, quæ ab inítio Evangélii in univérso mundo prædicáta est, in dilectióne et sanctitáte permánserit. Nam licet omnem Ecclésiam, quæ in toto est orbe terrárum, cunctis opórteat florére virtútibus; vos tamen præcípue inter céteros pópulos decet méritis pietátis excéllere, quos in ipsa apostólicæ petræ arce fundátos, et Dóminus noster Jesus Christus cum ómnibus redémit, et beátus Apóstolus Petrus præ ómnibus erudívit.\par
\switchcolumn
\EN
\noindent When, therefore, we address our exhortations to your godly ears, believe ye that ye are hearing him speak whose office we are discharging. Yea, it is with his love for you that we warn you. And we preach unto you no other thing than that which he taught, entreating you as did he: Gird up the loins of your mind; be sober; be ye holy in all manner of living; pass the time of your sojourning here in the fear of God. My disciples, dearly beloved, ye are to me as the disciples of the Apostle Paul were to him, namely: My crown and joy; if so be that your faith, abide, still in all lowliness and holiness, like unto the first times of the Gospel. For although the whole Church, which is in all the world, should indeed abound in all the virtues, it becometh especially you among all others to excel in acts of piety, founded as ye be on the very citadel of the Apostolic Rock ye who have not only been redeemed with the rest of men by our Lord Jesus Christ, but who have been instructed by the blessed Apostle Peter far beyond all others.\par
\switchcolumn*

\LA
\TeDeum{Te Deum}
\switchcolumn
\EN
\TeDeum{Te Deum}
\switchcolumn*

\end{paracol}

\Office{Commune Plurimorum Martyrum Pontificum}{Commune Plurimorum Martyrum Pontificum}{}
\addcontentsline{toc}{subsection}{Commune Plurimorum Martyrum Pontificum \textit{— Commune Plurimorum Martyrum Pontificum}}
\begin{paracol}{2}
\LA
\LessonHead{Lectio VII}
\switchcolumn
\EN
\LessonHead{Lesson VII}
\switchcolumn*

\LA
\LessonTitle{Léctio sancti Evangélii secúndum Lucam}
\switchcolumn
\EN
\LessonTitle{From the Holy Gospel according to Luke}
\switchcolumn*

\LA
\Cite{Luc 21:9-19}
\switchcolumn
\EN
\Cite{Luke 21:9-19}
\switchcolumn*

\LA
\noindent In illo témpore: Dixit Jesus discípulis suis: Cum audiéritis prǽlia, et seditiónes, nolíte terréri: opórtet primum hæc fíeri, sed nondum statim finis. Et réliqua.\par
\switchcolumn
\EN
\noindent At that time, Jesus said unto His disciples: When ye shall hear of wars and commotions, be not terrified for these things must first come to pass; but the end is not by and by. And so on.\par
\switchcolumn*

\LA
\LessonTitle{Homilía sancti Gregórii Papæ}
\switchcolumn
\EN
\LessonTitle{Homily by Pope St Gregory (the Great.)}
\switchcolumn*

\LA
\Source{Homilia 35 in Evangel.}
\switchcolumn
\EN
\Source{35th on the Gospels.}
\switchcolumn*

\LA
\noindent Dóminus ac Redémptor noster peritúri mundi præcurréntia mala denúntiat, ut eo minus pertúrbent veniéntia, quo fúerint præscíta. Minus enim jácula fériunt, quæ prævidéntur: et nos tolerabílius mundi mala suscípimus, si contra hæc per præsciéntiæ clípeum munímur. Ecce enim dicit: Cum audiéritis prǽlia et seditiónes, nolíte terréri: opórtet enim primum hæc fíeri, sed nondum statim finis. Pensánda sunt verba Redemptóris nostri, per quæ nos áliud intérius, áliud extérius passúros esse denúntiat; bella quippe ad hostes pértinent, seditiónes ad cives. Ut ergo nos índicet intérius exteriúsque turbári, áliud nos fatétur ab hóstibus, áliud a frátribus pérpeti.\par
\switchcolumn
\EN
\noindent Our Lord and Redeemer willeth us to know what shall be the signs that the end of the world is at hand, to the end that ye may be the less terrified, when that cometh whereof ye have already had warning. Darts strike less which are seen coming: and the plagues of the earth will be to us more bearable, if we are harnessed against them with the shield of foreknowledge. Behold, how He saith When ye shall hear of wars and commotions be not terrified for these things must first come to pass; but the end is not by and by. It behoveth us to ponder these words of our Redeemer, wherein He warneth us of suffering, from without, and from within. Wars are the work of a foreign enemy, commotions of the citizens. Therefore, that He may let us know that we shall be troubled from within and from without, He showeth that our wrestling shall be in part against strangers, and in part against our brethren.\par
\switchcolumn*

\LA
\begin{Resp}\Rsign Propter testaméntum Dómini, et leges patérnas, Sancti Dei perstitérunt in amóre fraternitátis: \Star{} Quia unus fuit semper spíritus in eis, et una fides.\end{Resp}
\begin{Resp}\Vsign Ecce quam bonum et quam jucúndum habitáre fratres in unum.\end{Resp}
\begin{Resp}\Rsign Quia unus fuit semper spíritus in eis, et una fides.\end{Resp}
\switchcolumn
\EN
\begin{Resp}\Rsign Because of the covenant of the Lord, and the laws of their fathers, the Saints of God abode in brotherly love, \Star{} For one spirit and one faith was ever in them.\end{Resp}
\begin{Resp}\Vsign Behold how good and how pleasant it is for brethren to dwell together in unity.\end{Resp}
\begin{Resp}\Rsign For one spirit and one faith was ever in them.\end{Resp}
\switchcolumn*

\LA
\LessonHead{Lectio VIII}
\switchcolumn
\EN
\LessonHead{Lesson VIII}
\switchcolumn*

\LA
\noindent Sed his malis præveniéntibus, quia non statim finis sequátur, adjúngit: Surget gens contra gentem, et regnum advérsus regnum: et terræmótus magni erunt per loca, et pestiléntiæ, et fames, terrorésque de cælo: et signa magna erunt. Última tribulátio multis tribulatiónibus prævenítur: et per crebra mala, quæ prævéniunt, indicántur mala perpétua, quæ subsequéntur. Et ídeo post bella et seditiónes non statim finis: quia multa debent mala præcúrrere, ut malum váleant sine fine nuntiáre.\par
\switchcolumn
\EN
\noindent When these woes come, the end is not by and by. And He saith further Nation shall rise against nation, and kingdom against kingdom; and great earthquakes shall be in diverse places, and pestilences, and famines, and fearful sights and great signs shall there be from heaven. Before the last tribulation cometh, shall come many other tribulations: and, by the many woes which shall come first, shall be foreshadowed the everlasting woe which shall come in the end. And therefore, after wars and commotions, the end is not yet by and by many woes must come first, to give warning of the woe that hath no end.\par
\switchcolumn*

\LA
\begin{Resp}\Rsign Sancti mei, qui in carne pósiti, certámen habuístis: \Star{} Mercédem labóris ego reddam vobis.\end{Resp}
\begin{Resp}\Vsign Veníte benedícti Patris mei, percípite regnum.\end{Resp}
\begin{Resp}\Rsign Mercédem labóris ego reddam vobis.\end{Resp}
\begin{Resp}\Vsign Glória Patri, et Fílio, \Star{} et Spirítui Sancto.\end{Resp}
\begin{Resp}\Rsign Mercédem labóris ego reddam vobis.\end{Resp}
\switchcolumn
\EN
\begin{Resp}\Rsign O ye My Saints, who, being in the flesh, didst have striving \Star{} I will render unto you a reward of your labours.\end{Resp}
\begin{Resp}\Vsign Come, ye blessed of My Father, inherit the kingdom\end{Resp}
\begin{Resp}\Rsign I will render unto you a reward of your labours.\end{Resp}
\begin{Resp}\Vsign Glory be to the Father, and to the Son, \Star{} and to the Holy Ghost.\end{Resp}
\begin{Resp}\Rsign I will render unto you a reward of your labours.\end{Resp}
\switchcolumn*

\end{paracol}

\Office{Commune unius Confessoris Pontificis}{Common of a Confessor Bishop}{}
\addcontentsline{toc}{subsection}{Commune unius Confessoris Pontificis \textit{— Common of a Confessor Bishop}}
\begin{paracol}{2}
\LA
\LessonHead{Lectio VII}
\switchcolumn
\EN
\LessonHead{Lesson VII}
\switchcolumn*

\LA
\LessonTitle{Léctio sancti Evangélii secúndum Matthǽum}
\switchcolumn
\EN
\LessonTitle{From the Holy Gospel according to Matthew}
\switchcolumn*

\LA
\Cite{Matt 25:14-23}
\switchcolumn
\EN
\Cite{Matt 25:14-23}
\switchcolumn*

\LA
\noindent In illo témpore: Dixit Jesus discípulis suis parábolam hanc: Homo péregre proficíscens, vocávit servos suos, et trádidit illis bona sua. Et réliqua.\par
\switchcolumn
\EN
\noindent At that time Jesus spake unto His disciples this parable A man, travelling into a far country, called his own servants, and delivered unto them his goods. And so on.\par
\switchcolumn*

\LA
\LessonTitle{Homilía sancti Gregórii Papæ}
\switchcolumn
\EN
\LessonTitle{Homily by Pope St Gregory (the Great.)}
\switchcolumn*

\LA
\Source{Homilia 9 in Evangelia}
\switchcolumn
\EN
\Source{9th on the Gospels.}
\switchcolumn*

\LA
\noindent Léctio sancti Evangélii, fratres caríssimi, solícite consideráre nos ádmonet, ne nos, qui plus céteris in hoc mundo accepísse áliquid cérnimur, ab Auctóre mundi grávius inde judicémur. Cum enim augéntur dona, ratiónes étiam crescunt donórum. Tanto ergo esse humílior atque ad serviéndum Deo prómptior quisque debet ex múnere, quanto se obligatiórem esse cónspicit in reddénda ratióne. Ecce homo, qui péregre proficíscitur, servos suos vocat, eisque ad negótium talénta partítur. Post multum vero témporis positúrus ratiónem revértitur: bene operántes pro apportáto lucro remúnerat, servum vero a bono ópere torpéntem damnat.\par
\switchcolumn
\EN
\noindent Dearly beloved brethren, this Lesson from the Holy Gospel moveth us to take good heed lest we, who are seen in this world to have received more than others, should thereby bring ourselves into greater condemnation from the Maker of this world. To whom much is given, of the same is much required. Therefore, let him that receiveth much, strive to be all the more lowly, and all the more ready to do God service, for his very gifts’ sake, knowing that he will be obliged to give account thereof. Behold, a man, travelling into a far country, calleth his own servants, and delivereth unto them talents, to the end that they may trade therewith. After a long time, the lord of those servants cometh, and reckoneth with them, and to them that have done well He rendereth a reward of their labours, but that servant which was careless of his master’s work He condemneth.\par
\switchcolumn*

\LA
\begin{Resp}\Rsign Amávit eum Dóminus, et ornávit eum: stolam glóriæ índuit eum, \Star{} Et ad portas paradísi coronávit eum.\end{Resp}
\begin{Resp}\Vsign Induit eum Dóminus lorícam fídei, et ornávit eum.\end{Resp}
\begin{Resp}\Rsign Et ad portas paradísi coronávit eum.\end{Resp}
\switchcolumn
\EN
\begin{Resp}\Rsign The Lord loved him and beautified him He clothed him with a robe of glory \Star{} And crowned him at the gates of Paradise.\end{Resp}
\begin{Resp}\Vsign The Lord hath put on him the breast-plate of faith, and hath adorned him.\end{Resp}
\begin{Resp}\Rsign And crowned him at the gates of Paradise.\end{Resp}
\switchcolumn*

\LA
\LessonHead{Lectio VIII}
\switchcolumn
\EN
\LessonHead{Lesson VIII}
\switchcolumn*

\LA
\noindent Quis ítaque iste homo est qui péregre proficíscitur, nisi Redémptor noster, qui, in ea carne quam assúmpserat, ábiit in cælum? Carnis enim locus próprius terra est; quæ quasi ad peregrína dúcitur, dum per Redemptórem nostrum in cælo collocátur. Sed homo iste, péregre proficíscens, servis suis bona sua trádidit; quia fidélibus suis spirituália dona concéssit. Et uni quidem quinque talénta, álii duo, álii vero commísit unum. Quinque étenim sunt córporis sensus, vidélicet: visus, audítus, gustus, odorátus et tactus. Quinque ergo taléntis, donum quinque sénsuum, id est, exteriórum sciéntia, exprímitur. Duóbus vero, intelléctus et operátio designátur. Uníus autem talénti nómine, intelléctus tantúmmodo designátur.\par
\switchcolumn
\EN
\noindent What other, then, is that man travelling into a far country but our Redeemer, Who is gone up from us into heaven in that Flesh Which He had taken into Himself? For the earth is the home of the Flesh, Which travelleth into a far country when our Redeemer giveth It a place in heaven. But that man travelling into a far country delivered unto his servants his goods and so doth our Redeemer give spiritual gifts unto His faithful people. “And unto one he gave five talents, to another two, and to another one.” There are five bodily senses that is, sight, hearing, taste, smell, and touch. By the five talents therefore are signified the five senses, that is, outward knowledge. By the two, wit and work. And by the figure of the one talent, understanding, which is alone.\par
\switchcolumn*

\LA
\begin{Resp}\Rsign Sint lumbi vestri præcíncti, et lucérnæ ardéntes in mánibus vestris: \Star{} Et vos símiles homínibus exspectántibus dóminum suum, quando revertátur a núptiis.\end{Resp}
\begin{Resp}\Vsign Vigiláte ergo, quia nescítis qua hora Dóminus vester ventúrus sit.\end{Resp}
\begin{Resp}\Rsign Et vos símiles homínibus exspectántibus dóminum suum, quando revertátur a núptiis.\end{Resp}
\begin{Resp}\Vsign Glória Patri, et Fílio, \Star{} et Spirítui Sancto.\end{Resp}
\begin{Resp}\Rsign Et vos símiles homínibus exspectántibus dóminum suum, quando revertátur a núptiis.\end{Resp}
\switchcolumn
\EN
\begin{Resp}\Rsign Let your loins be girded about, and your lights burning; \Star{} And ye yourselves like unto men that wait for their lord, when he will return from the wedding.\end{Resp}
\begin{Resp}\Vsign Watch therefore, for ye know not what hour your Lord doth come.\end{Resp}
\begin{Resp}\Rsign And ye yourselves like unto men that wait for their lord, when he will return from the wedding.\end{Resp}
\begin{Resp}\Vsign Glory be to the Father, and to the Son, \Star{} and to the Holy Ghost.\end{Resp}
\begin{Resp}\Rsign And ye yourselves like unto men that wait for their lord, when he will return from the wedding.\end{Resp}
\switchcolumn*

\LA
\LessonHead{Lectio IX}
\switchcolumn
\EN
\LessonHead{Lesson IX}
\switchcolumn*

\LA
\noindent Sed is, qui quinque talénta accéperat, ália quinque lucrátus est: quia sunt nonnúlli, qui, etsi intérna ac mýstica penetráre nésciunt, pro intentióne tamen supérnæ pátriæ docent recta quos possunt; de ipsis exterióribus, quæ accepérunt, duplum taléntum portant; dumque se a carnis petulántia et a terrenárum rerum ámbitu atque a visibílium voluptáte custódiunt, ab his étiam álios admonéndo compéscunt. Et sunt nonnúlli, qui, quasi duóbus taléntis ditáti, intelléctum atque operatiónem percípiunt, subtília de intérnis intélligunt, mira in exterióribus operántur; cumque et intelligéndo et operándo áliis prǽdicant, quasi duplicátum de negótio lucrum repórtant.\par
\switchcolumn
\EN
\noindent And so he that had received five talents, gained other five talents for some there be who, while yet they are not able to go on unto things inward and mystic, do yet so desire our Fatherland which is above, that they teach well all whom they can, and of those very outward things which they have received make gain double. These are they which keep themselves clean from the unruly motions of the flesh, and from the lust of the world, and from the delight of things which are seen, and, by their preaching, keep other men also clean from all these things. And some there are who receive, as their two talents, the power to think and the power to work. These are they which inwardly understand dark things, and outwardly work wonders. And these, since they preach unto others, both through their understanding and their works, gain, as it were, double, for the talents which they have received.\par
\switchcolumn*

\LA
\TeDeum{Te Deum}
\switchcolumn
\EN
\TeDeum{Te Deum}
\switchcolumn*

\end{paracol}

\Office{Commune Doctoris Pontificis}{Common of a Doctor Bishop}{}
\addcontentsline{toc}{subsection}{Commune Doctoris Pontificis \textit{— Common of a Doctor Bishop}}
\begin{paracol}{2}
\LA
\LessonHead{Lectio VII}
\switchcolumn
\EN
\LessonHead{Lesson VII}
\switchcolumn*

\LA
\LessonTitle{Léctio sancti Evangélii secúndum Matthǽum}
\switchcolumn
\EN
\LessonTitle{From the Holy Gospel according to Matthew}
\switchcolumn*

\LA
\Cite{Matt 5:13-19}
\switchcolumn
\EN
\Cite{Matt 5:13-19}
\switchcolumn*

\LA
\noindent In illo témpore: Dixit Jesus discípulis suis: Vos estis sal terræ. Quod si sal evanúerit, in quo saliétur? Et réliqua.\par
\switchcolumn
\EN
\noindent At that time: Jesus said unto his disciples: Ye are the salt of the earth: But if the salt have lost his savour, wherewith shall it be salted? And so on, and that which followeth.\par
\switchcolumn*

\LA
\LessonTitle{Homilía sancti Augustíni Epíscopi}
\switchcolumn
\EN
\LessonTitle{A Homily by St. Augustine the Bishop}
\switchcolumn*

\LA
\Source{Lib. 1 de Sermone Domini in monte, cap. 6}
\switchcolumn
\EN
\Source{Lib. 1 de Sermóne Dómini in monte, cap. 6}
\switchcolumn*

\LA
\noindent Osténdit Dóminus fátuos esse judicándos, qui temporálium bonórum vel cópiam sectántes vel inópiam metuéntes, amíttunt ætérna, quæ nec dari possunt ab homínibus nec auférri. Itaque, si sal infatuátum fúerit, in quo saliétur? Id est, si vos, per quos condiéndi sunt quodámmodo pópuli, metu persecutiónum temporálium amiséritis regna cælórum; qui erunt hómines, per quos a vobis error auferátur, cum vos elégerit Deus, per quos errórem áuferat ceterórum?\par
\switchcolumn
\EN
\noindent The Lord would have us understand how that men do lose their power of savouring others with righteousness when they are willing to place their eternal welfare in jeopardy for the sake of any témporal advantage, like as attainment of ease or luxury, or escape from suffering or toil. For that which is eternal, unlike things of this world, can neither be bestowed by men, nor by them taken away. Hence, when he asketh: If the salt have lost his savour, wherewith shall it be salted? he would have us understand the question to be: If ye, by whom mankind is preserved from corruption, be willing to lose the kingdom of heaven so as to escape trials or persecutions in this world, who is there to preserve you from corruption, seeing ye are they that God hath chosen to preserve all others from corruption?\par
\switchcolumn*

\LA
\begin{Resp}\Rsign Amávit eum Dóminus, et ornávit eum: stolam glóriæ índuit eum, \Star{} Et ad portas paradísi coronávit eum.\end{Resp}
\begin{Resp}\Vsign Induit eum Dóminus lorícam fídei, et ornávit eum.\end{Resp}
\begin{Resp}\Rsign Et ad portas paradísi coronávit eum.\end{Resp}
\switchcolumn
\EN
\begin{Resp}\Rsign The Lord loved him and beautified him He clothed him with a robe of glory \Star{} And crowned him at the gates of Paradise.\end{Resp}
\begin{Resp}\Vsign The Lord hath put on him the breast-plate of faith, and hath adorned him.\end{Resp}
\begin{Resp}\Rsign And crowned him at the gates of Paradise.\end{Resp}
\switchcolumn*

\LA
\LessonHead{Lectio VIII}
\switchcolumn
\EN
\LessonHead{Lesson VIII}
\switchcolumn*

\LA
\noindent Ergo ad níhilum valet sal infatuátum, nisi ut mittátur foras et calcétur ab homínibus. Non ítaque calcátur ab homínibus qui pátitur persecutiónem; sed qui, persecutiónem timéndo, infatuátur. Calcári enim non potest nisi inférior; sed inférior non est, qui, quamvis córpore multa in terra sustíneat, corde tamen fixus in cælo est.\par
\switchcolumn
\EN
\noindent Those that should be the salt of the earth, but have lost their savour, are thenceforth good for nothing, but to be cast out, and to be trodden under foot of men. But no one that suffereth persecution is truly said to be trodden under foot of men. Rather, that one is truly trodden under foot of men who through fear of persecution hath lost the savour of righteousness. For no one can be trodden upon, unless he be beneath him which treadeth upon him. And certainly no one who hath his heart in heaven, no matter how grievously he doth suffer in his body on earth, is rightly said to be beneath anyone who misuseth him.\par
\switchcolumn*

\LA
\begin{Resp}\Rsign In médio Ecclésiæ apéruit os ejus, \Star{} Et implévit eum Dóminus spíritu sapiéntiæ et intelléctus.\end{Resp}
\begin{Resp}\Vsign Jucunditátem et exsultatiónem thesaurizávit super eum.\end{Resp}
\begin{Resp}\Rsign Et implévit eum Dóminus spíritu sapiéntiæ et intelléctus.\end{Resp}
\begin{Resp}\Vsign Glória Patri, et Fílio, \Star{} et Spirítui Sancto.\end{Resp}
\begin{Resp}\Rsign Et implévit eum Dóminus spíritu sapiéntiæ et intelléctus.\end{Resp}
\switchcolumn
\EN
\begin{Resp}\Rsign In the midst of the congregation did the Lord open his mouth. \Star{} And filled him with the spirit of wisdom and understanding.\end{Resp}
\begin{Resp}\Vsign He made him rich with joy and gladness.\end{Resp}
\begin{Resp}\Rsign And filled him with the spirit of wisdom and understanding.\end{Resp}
\begin{Resp}\Vsign Glory be to the Father, and to the Son, \Star{} and to the Holy Ghost.\end{Resp}
\begin{Resp}\Rsign And filled him with the spirit of wisdom and understanding.\end{Resp}
\switchcolumn*

\LA
\LessonHead{Lectio IX}
\switchcolumn
\EN
\LessonHead{Lesson IX}
\switchcolumn*

\LA
\noindent Vos estis lumen mundi. Quómodo dixit supérius sal terræ, sic nunc dicit lumen mundi. Nam, neque supérius ista terra accipiénda est, quam pédibus corpóreis calcámus; sed hómines, qui in terra hábitant, vel étiam peccatóres, quorum condiéndis et exstinguéndis putóribus apostólicum salem Dóminus misit. Et hic mundum non cælum et terram, sed hómines qui sunt in mundo vel díligunt mundum, opórtet intélligi; quibus illuminándis Apóstoli missi sunt. Non potest cívitas abscóndi super montem pósita; id est, fundáta super insígnem magnámque justítiam, quam signíficat étiam ipse mons, in quo dísputat Dóminus.\par
\switchcolumn
\EN
\noindent Ye are the light of the world. And we are to understand the word World in the same sense as the word Earth when he spoke above of the salt of the earth, that is, not that earth whereupon we walk with our bodily feet, but the men which dwell upon the earth; in other words, sinners, for the sweetening and correction of whose corruption, the Lord hath sent his Apostles, as it were, as so much salt. And so by the world we are to understand, not the heaven and the earth, but the men who are in the world and love the world, for the enlightening of whom the Apostles have been sent. A city that is set on a hill cannot be hid: that is, what is founded upon the heights of righteousness, whereof the mountain upon which the Lord gave this discourse was itself a figure, is magnificent in the eyes of all men.\par
\switchcolumn*

\LA
\TeDeum{Te Deum}
\switchcolumn
\EN
\TeDeum{Te Deum}
\switchcolumn*

\end{paracol}

\Office{Commune Confessoris non pontificis}{Common of a Confessor not a Bishop}{}
\addcontentsline{toc}{subsection}{Commune Confessoris non pontificis \textit{— Common of a Confessor not a Bishop}}
\begin{paracol}{2}
\LA
\LessonHead{Lectio VII}
\switchcolumn
\EN
\LessonHead{Lesson VII}
\switchcolumn*

\LA
\LessonTitle{Léctio sancti Evangélii secúndum Lucam}
\switchcolumn
\EN
\LessonTitle{From the Holy Gospel according to Luke}
\switchcolumn*

\LA
\Cite{Luc 12:35-40}
\switchcolumn
\EN
\Cite{Luke 12:35-40}
\switchcolumn*

\LA
\noindent In illo témpore: Dixit Jesus discípulis suis: Sint lumbi vestri præcincti, et lucernæ ardentes in mánibus vestris. Et réliqua.\par
\switchcolumn
\EN
\noindent At that time, Jesus said unto His disciples: Let your loins be girded about, and your lights burning. And so on.\par
\switchcolumn*

\LA
\LessonTitle{Homilía sancti Gregórii Papæ}
\switchcolumn
\EN
\LessonTitle{Homily by Pope St. Gregory (the Great.)}
\switchcolumn*

\LA
\Source{Homilia 13 in Evang.}
\switchcolumn
\EN
\Source{13th on the Gospels.}
\switchcolumn*

\LA
\noindent Sancti Evangélii, fratres caríssimi, apérta vobis est léctio recitata. Sed ne alíquibus ipsa ejus planíties alta fortásse videátur, eam sub brevitáte transcúrrimus, quátenus ejus exposítio ita nesciéntibus fiat cógnita, ut tamen sciéntibus non sit onerósa. Dóminus dicit: Sint lumbi vestri præcíncti. Lumbos enim præcíngimus, cum carnis luxúriam per continéntiam coarctámus. Sed quia minus est, mala non ágere, nisi étiam quisque stúdeat, et bonis opéribus insudáre, prótinus ádditur: Et lucérnæ ardéntes in mánibus vestris. Lucérnas quippe ardéntes in mánibus tenémus, cum per bona ópera próximis nostris lucis exémpla monstrámus. De quibus profécto opéribus Dóminus dicit: Lúceat lux vestra coram homínibus, ut vídeant ópera vestra bona, et gloríficent Patrem vestrum qui in cælis est.\par
\switchcolumn
\EN
\noindent Dearly beloved brethren, the words of the Holy Gospel, which have just been read, lie open before you, and, lest their very plainness should make them seem to some to be hard, we will go through them with such shortness as that neither may they which understand not remain unenlightened, nor they which understand be wearied. The Lord saith Let your loins be girded about. Now, we gird our loins about, when by continency we master the lustful inclination of the flesh. But, forasmuch as it sufficeth not for a man to abstain from evil deeds, if he strive not to join thereto the earnest doing of good works, it is immediately added And your lights burning. Our lights burn when, by good works, we give bright example to our neighbour; concerning which works the Lord saith Let your light so shine before men, that they may see your good works, and glorify your Father Which is in heaven. (Matth. v. 16.)\par
\switchcolumn*

\LA
\begin{Resp}\Rsign Iste est qui ante Deum magnas virtútes operátus est, et de omni corde suo laudávit Dóminum: \Star{} Ipse intercédat pro peccátis ómnium populórum.\end{Resp}
\begin{Resp}\Vsign Ecce homo sine queréla, verus Dei cultor, ábstinens se ab omni ópere malo, et pérmanens in innocéntia sua.\end{Resp}
\begin{Resp}\Rsign Ipse intercédat pro peccátis ómnium populórum.\end{Resp}
\switchcolumn
\EN
\begin{Resp}\Rsign This is he which wrought great wonders before God, and praised the Lord with all his heart. \Star{} May he pray for all people, that their sins may be forgiven unto them.\end{Resp}
\begin{Resp}\Vsign Behold a man without blame, a worshipper of God in truth, keeping himself clean from every evil work, and abiding still in his innocency.\end{Resp}
\begin{Resp}\Rsign May he pray for all people, that their sins may be forgiven unto them.\end{Resp}
\switchcolumn*

\LA
\LessonHead{Lectio VIII}
\switchcolumn
\EN
\LessonHead{Lesson VIII}
\switchcolumn*

\LA
\noindent Duo autem sunt, quæ jubéntur, et lumbos restríngere, et lucérnas tenére: ut et mundítia sit castitátis in córpore, et lumen veritátis in operatióne. Redemptóri étenim nostro unum sine áltero placére nequáquam potest: si aut is qui bona agit, adhuc luxúriæ inquinaménta non déserit: aut is qui castitáte præéminet, necdum se per bona ópera exércet. Nec cástitas ergo magna est sine bono ópere, nec opus bonum est áliquod sine castitáte. Sed et si utrúmque ágitur, restat, ut quisquis ille est, spe ad supérnam pátriam tendat, et nequáquam se a vítiis pro mundi hujus honestáte contíneat.\par
\switchcolumn
\EN
\noindent Here, then, are two commandments, to gird our loins about, and to keep our lights burning the cleanness of purity in our body, and the light of the truth in our works. Whoso hath the one and not the other, pleaseth not thereby our Redeemer; that is, he pleaseth Him not which doth good works, but bridleth not himself from the pollutions of lust, neither he which is eminent in chastity, but exerciseth not himself in good works. Neither is chastity a great thing without good works, nor good works anything without chastity. And if any man do both, it remaineth that he must look by hope toward our Fatherland above, and not have for his reason wherethrough he turneth himself away from vice, the love of honour in this present world.\par
\switchcolumn*

\LA
\begin{Resp}\Rsign Sint lumbi vestri præcíncti, et lucérnæ ardéntes in mánibus vestris: \Star{} Et vos símiles homínibus exspectántibus dóminum suum, quando revertátur a núptiis.\end{Resp}
\begin{Resp}\Vsign Vigiláte ergo, quia nescítis qua hora Dóminus vester ventúrus sit.\end{Resp}
\begin{Resp}\Rsign Et vos símiles homínibus exspectántibus dóminum suum, quando revertátur a núptiis.\end{Resp}
\begin{Resp}\Vsign Glória Patri, et Fílio, \Star{} et Spirítui Sancto.\end{Resp}
\begin{Resp}\Rsign Et vos símiles homínibus exspectántibus dóminum suum, quando revertátur a núptiis.\end{Resp}
\switchcolumn
\EN
\begin{Resp}\Rsign Let your loins be girded about, and your lights burning; \Star{} And ye yourselves like unto men that wait for their lord, when he will return from the wedding.\end{Resp}
\begin{Resp}\Vsign Watch therefore, for ye know not what hour your Lord doth come.\end{Resp}
\begin{Resp}\Rsign And ye yourselves like unto men that wait for their lord, when he will return from the wedding.\end{Resp}
\begin{Resp}\Vsign Glory be to the Father, and to the Son, \Star{} and to the Holy Ghost.\end{Resp}
\begin{Resp}\Rsign And ye yourselves like unto men that wait for their lord, when he will return from the wedding.\end{Resp}
\switchcolumn*

\LA
\LessonHead{Lectio IX}
\switchcolumn
\EN
\LessonHead{Lesson IX}
\switchcolumn*

\LA
\noindent Et vos símiles homínibus exspectántibus dóminum suum, quando revertátur a núptiis: ut cum vénerit et pulsáverit, conféstim apériant ei. Venit quippe Dóminus, cum ad judícium próperat: pulsat vero, cum jam per ægritúdinis moléstias esse mortem vicínam desígnat. Cui conféstim aperímus, si hunc cum amóre suscípimus. Aperíre enim júdici pulsánti non vult, qui exíre de córpore trépidat: et vidére eum, quem contempsísse se méminit, júdicem formídat. Qui autem de sua spe et operatióne secúrus est, pulsanti conféstim áperit, quia lætus júdicem sústinet; et, cum tempus propínquæ mortis advénerit, de glória retributiónis hilaréscit.\par
\switchcolumn
\EN
\noindent And ye yourselves like unto men that wait for their lord, when he will return from the wedding that, when he cometh and knocketh, they may open unto him immediately. The Lord cometh at the hour of judgment He knocketh when, by the pains of sickness, He biddeth us know that death is nigh. To Him open we immediately, if we receive Him in love. Whoso feareth to leave this body, will not open to the Judge when He knocketh, for he dreadeth to see that Judge, Whom he knoweth that he hath despised. But whosoever knoweth that his hope and works are built upon a good foundation, when he heareth the Judge knock, openeth to Him immediately, for to such an one that coming is blessed, yea, when the hour of death is at hand, such an one haileth with gladness a glorious reward.\par
\switchcolumn*

\LA
\TeDeum{Te Deum}
\switchcolumn
\EN
\TeDeum{Te Deum}
\switchcolumn*

\end{paracol}

\Office{Commune Doctoris non Pontificis}{Common of a Doctor not a Bishop}{}
\addcontentsline{toc}{subsection}{Commune Doctoris non Pontificis \textit{— Common of a Doctor not a Bishop}}
\begin{paracol}{2}
\LA
\LessonHead{Lectio VII}
\switchcolumn
\EN
\LessonHead{Lesson VII}
\switchcolumn*

\LA
\LessonTitle{Léctio sancti Evangélii secúndum Matthǽum}
\switchcolumn
\EN
\LessonTitle{From the Holy Gospel according to Matthew}
\switchcolumn*

\LA
\Cite{Matt 5:13-19}
\switchcolumn
\EN
\Cite{Matt 5:13-19}
\switchcolumn*

\LA
\noindent In illo témpore: Dixit Jesus discípulis suis: Vos estis sal terræ. Quod si sal evanúerit, in quo saliétur? Et réliqua.\par
\switchcolumn
\EN
\noindent At that time: Jesus said unto his disciples: Ye are the salt of the earth: But if the salt have lost his savour, wherewith shall it be salted? And so on, and that which followeth.\par
\switchcolumn*

\LA
\LessonTitle{Homilía sancti Augustíni Epíscopi}
\switchcolumn
\EN
\LessonTitle{A Homily by St. Augustine the Bishop}
\switchcolumn*

\LA
\Source{Lib. 1 de Sermone Domini in monte, cap. 6}
\switchcolumn
\EN
\Source{Lib. 1 de Sermóne Dómini in monte, cap. 6}
\switchcolumn*

\LA
\noindent Osténdit Dóminus fátuos esse judicándos, qui temporálium bonórum vel cópiam sectántes vel inópiam metuéntes, amíttunt ætérna, quæ nec dari possunt ab homínibus nec auférri. Itaque, si sal infatuátum fúerit, in quo saliétur? Id est, si vos, per quos condiéndi sunt quodámmodo pópuli, metu persecutiónum temporálium amiséritis regna cælórum; qui erunt hómines, per quos a vobis error auferátur, cum vos elégerit Deus, per quos errórem áuferat ceterórum?\par
\switchcolumn
\EN
\noindent The Lord would have us understand how that men do lose their power of savouring others with righteousness when they are willing to place their eternal welfare in jeopardy for the sake of any témporal advantage, like as attainment of ease or luxury, or escape from suffering or toil. For that which is eternal, unlike things of this world, can neither be bestowed by men, nor by them taken away. Hence, when he asketh: If the salt have lost his savour, wherewith shall it be salted? he would have us understand the question to be: If ye, by whom mankind is preserved from corruption, be willing to lose the kingdom of heaven so as to escape trials or persecutions in this world, who is there to preserve you from corruption, seeing ye are they that God hath chosen to preserve all others from corruption?\par
\switchcolumn*

\LA
\begin{Resp}\Rsign Iste est qui ante Deum magnas virtútes operátus est, et de omni corde suo laudávit Dóminum: \Star{} Ipse intercédat pro peccátis ómnium populórum.\end{Resp}
\begin{Resp}\Vsign Ecce homo sine queréla, verus Dei cultor, ábstinens se ab omni ópere malo, et pérmanens in innocéntia sua.\end{Resp}
\begin{Resp}\Rsign Ipse intercédat pro peccátis ómnium populórum.\end{Resp}
\switchcolumn
\EN
\begin{Resp}\Rsign This is he which wrought great wonders before God, and praised the Lord with all his heart. \Star{} May he pray for all people, that their sins may be forgiven unto them.\end{Resp}
\begin{Resp}\Vsign Behold a man without blame, a worshipper of God in truth, keeping himself clean from every evil work, and abiding still in his innocency.\end{Resp}
\begin{Resp}\Rsign May he pray for all people, that their sins may be forgiven unto them.\end{Resp}
\switchcolumn*

\LA
\LessonHead{Lectio VIII}
\switchcolumn
\EN
\LessonHead{Lesson VIII}
\switchcolumn*

\LA
\noindent Ergo ad níhilum valet sal infatuátum, nisi ut mittátur foras et calcétur ab homínibus. Non ítaque calcátur ab homínibus qui pátitur persecutiónem; sed qui, persecutiónem timéndo, infatuátur. Calcári enim non potest nisi inférior; sed inférior non est, qui, quamvis córpore multa in terra sustíneat, corde tamen fixus in cælo est.\par
\switchcolumn
\EN
\noindent Those that should be the salt of the earth, but have lost their savour, are thenceforth good for nothing, but to be cast out, and to be trodden under foot of men. But no one that suffereth persecution is truly said to be trodden under foot of men. Rather, that one is truly trodden under foot of men who through fear of persecution hath lost the savour of righteousness. For no one can be trodden upon, unless he be beneath him which treadeth upon him. And certainly no one who hath his heart in heaven, no matter how grievously he doth suffer in his body on earth, is rightly said to be beneath anyone who misuseth him.\par
\switchcolumn*

\LA
\begin{Resp}\Rsign In médio Ecclésiæ apéruit os ejus, \Star{} Et implévit eum Dóminus spíritu sapiéntiæ et intelléctus.\end{Resp}
\begin{Resp}\Vsign Jucunditátem et exsultatiónem thesaurizávit super eum.\end{Resp}
\begin{Resp}\Rsign Et implévit eum Dóminus spíritu sapiéntiæ et intelléctus.\end{Resp}
\begin{Resp}\Vsign Glória Patri, et Fílio, \Star{} et Spirítui Sancto.\end{Resp}
\begin{Resp}\Rsign Et implévit eum Dóminus spíritu sapiéntiæ et intelléctus.\end{Resp}
\switchcolumn
\EN
\begin{Resp}\Rsign In the midst of the congregation did the Lord open his mouth. \Star{} And filled him with the spirit of wisdom and understanding.\end{Resp}
\begin{Resp}\Vsign He made him rich with joy and gladness.\end{Resp}
\begin{Resp}\Rsign And filled him with the spirit of wisdom and understanding.\end{Resp}
\begin{Resp}\Vsign Glory be to the Father, and to the Son, \Star{} and to the Holy Ghost.\end{Resp}
\begin{Resp}\Rsign And filled him with the spirit of wisdom and understanding.\end{Resp}
\switchcolumn*

\end{paracol}

\Office{Commune Unius Virginis et Martyris}{Common of a Virgin Martyr}{}
\addcontentsline{toc}{subsection}{Commune Unius Virginis et Martyris \textit{— Common of a Virgin Martyr}}
\begin{paracol}{2}
\LA
\LessonHead{Lectio I}
\switchcolumn
\EN
\LessonHead{Lesson I}
\switchcolumn*

\LA
\noindent De virgínibus præcéptum Dómini non hábeo: consílium autem do, tamquam misericórdiam consecútus a Dómino, ut sim fidélis. Exístimo ergo hoc bonum esse propter instántem necessitátem, quóniam bonum est hómini sic esse. Alligátus es uxóri? noli quǽrere solutiónem. Solútus es ab uxóre? noli quǽrere uxórem. Si autem accéperis uxórem, non peccásti. Et si núpserit virgo, non peccávit. Tribulatiónem tamen carnis habébunt hujúsmodi. Ego autem vobis parco. Hoc ítaque dico, fratres: Tempus breve est: réliquum est, ut et qui habent uxóres, tamquam non habéntes sint; et qui flent, tamquam non flentes; et qui gaudent, tamquam non gaudéntes; et qui emunt, tamquam non possidéntes; et qui utúntur hoc mundo, tamquam non utántur; prǽterit enim figúra hujus mundi.\par
\switchcolumn
\EN
\noindent Now concerning virgins, I have no commandment of the Lord; but I give counsel, as having obtained mercy of the Lord, to be faithful. I think therefore that this is good for the present necessity, that it is good for a man so to be. Art thou bound to a wife? seek not to be loosed. Art thou loosed from a wife? seek not a wife. But if thou take a wife, thou hast not sinned. And if a virgin marry, she hath not sinned: nevertheless, such shall have tribulation of the flesh. But I spare you. This therefore I say, brethren; the time is short; it remaineth, that they also who have wives, be as if they had none; And they that weep, as though they wept not; and they that rejoice, as if they rejoiced not; and they that buy, as though they possessed not; And they that use this world, as if they used it not: for the fashion of this world passeth away.\par
\switchcolumn*

\LA
\begin{Resp}\Rsign Veni Sponsa Christi, áccipe corónam, quam tibi Dóminus præparávit in ætérnum; pro cujus amóre sánguinem tuum fudísti, \Star{} Et cum Ángelis in paradísum introísti.\end{Resp}
\begin{Resp}\Vsign Veni, elécta mea, et ponam in te thronum meum; quia concupívit Rex spéciem tuam.\end{Resp}
\begin{Resp}\Rsign Et cum Ángelis in paradísum introísti.\end{Resp}
\switchcolumn
\EN

\switchcolumn*

\LA
\LessonHead{Lectio II}
\switchcolumn
\EN
\LessonHead{Lesson II}
\switchcolumn*

\LA
\Cite{1 Cor 7:32-35}
\switchcolumn
\EN
\Cite{1 Cor 7:32-35}
\switchcolumn*

\LA
\noindent Volo autem vos sine sollicitúdine esse. Qui sine uxóre est, sollícitus est quæ Dómini sunt, quómodo pláceat Deo. Qui autem cum uxóre est, sollícitus est quæ sunt mundi, quómodo pláceat uxóri; et divísus est. Et múlier innúpta et virgo cógitat quæ Dómini sunt, ut sit sancta córpore et spíritu. Quæ autem nupta est, cógitat quæ sunt mundi, quómodo pláceat viro. Porro hoc ad utilitátem vestram dico, non ut láqueum vobis iníciam, sed ad id, quod honéstum est, et quod facultátem prǽbeat sine impediménto Dóminum obsecrándi.\par
\switchcolumn
\EN
\noindent But I would have you to be without solicitude. He that is without a wife, is solicitous for the things that belong to the Lord, how he may please God. But he that is with a wife, is solicitous for the things of the world, how he may please his wife: and he is divided. And the unmarried woman and the virgin thinketh on the things of the Lord, that she may be holy both in body and in spirit. But she that is married thinketh on the things of the world, how she may please her husband. And this I speak for your profit: not to cast a snare upon you; but for that which is decent, and which may give you power to attend upon the Lord, without impediment.\par
\switchcolumn*

\LA
\begin{Resp}\Rsign Diffúsa est grátia in lábiis tuis, \Star{} Proptérea benedíxit te Deus in ætérnum.\end{Resp}
\begin{Resp}\Vsign Spécie tua et pulchritúdine tua inténde, próspere procéde, et regna.\end{Resp}
\begin{Resp}\Rsign Proptérea benedíxit te Deus in ætérnum.\end{Resp}
\switchcolumn
\EN
\begin{Resp}\Rsign Grace is poured into thy lips, therefore; \Star{} God hath blessed thee for ever.\end{Resp}
\begin{Resp}\Vsign In thy comeliness and thy beauty, go forward, fare prosperously, and reign.\end{Resp}
\begin{Resp}\Rsign God hath blessed thee for ever.\end{Resp}
\switchcolumn*

\LA
\LessonHead{Lectio III}
\switchcolumn
\EN
\LessonHead{Lesson III}
\switchcolumn*

\LA
\Cite{1 Cor 7:36-40}
\switchcolumn
\EN
\Cite{1 Cor 7:36-40}
\switchcolumn*

\LA
\noindent Si quis autem turpem se vidéri exístimat super vírgine sua, quod sit superadúlta, et ita opórtet fíeri; quod vult fáciat: non peccat, si nubat. Nam, qui státuit in corde suo firmus, non habens necessitátem, potestátem autem habens suæ voluntátis, et hoc judicávit in corde suo, serváre vírginem suam, bene facit. Ígitur et qui matrimónio jungit vírginem suam, bene facit; et qui non jungit, mélius facit. Múlier alligáta est legi quanto témpore, vir ejus vivit. Quod, si dormíerit vir ejus, liberáta est; cui vult nubat, tantum in Dómino. Beátior autem erit, si sic permánserit, secúndum meum consílium; puto autem quod et ego Spíritum Dei hábeam.\par
\switchcolumn
\EN
\noindent But if any man think that he seemeth dishonoured, with regard to his virgin, for that she is above the age, and it must so be: let him do what he will; he sinneth not, if she marry. For he that hath determined being steadfast in his heart, having no necessity, but having power of his own will; and hath judged this in his heart, to keep his virgin, doth well. Therefore, both he that giveth his virgin in marriage, doth well; and he that giveth her not, doth better. A woman is bound by the law as long as her husband liveth; but if her husband die, she is at liberty: let her marry to whom she will; only in the Lord. But more blessed shall she be, if she so remain, according to my counsel; and I think that I also have the spirit of God.\par
\switchcolumn*

\LA
\begin{Resp}\Rsign Spécie tua et pulchritúdine tua \Star{} Inténde, próspere procéde, et regna.\end{Resp}
\begin{Resp}\Vsign Diffúsa est grátia in lábiis tuis, proptérea benedíxit te Deus in ætérnum.\end{Resp}
\begin{Resp}\Rsign Inténde, próspere procéde, et regna.\end{Resp}
\begin{Resp}\Vsign Glória Patri, et Fílio, \Star{} et Spirítui Sancto.\end{Resp}
\begin{Resp}\Rsign Inténde, próspere procéde, et regna.\end{Resp}
\switchcolumn
\EN
\begin{Resp}\Rsign In thy comeliness and thy beauty; \Star{} Go forward, fare prosperously, and reign.\end{Resp}
\begin{Resp}\Vsign Grace is poured into thy lips, therefore God hath blessed thee for ever.\end{Resp}
\begin{Resp}\Rsign Go forward, fare prosperously, and reign.\end{Resp}
\begin{Resp}\Vsign Glory be to the Father, and to the Son, \Star{} and to the Holy Ghost.\end{Resp}
\begin{Resp}\Rsign Go forward, fare prosperously, and reign.\end{Resp}
\switchcolumn*

\LA
\LessonHead{Lectio VII}
\switchcolumn
\EN
\LessonHead{Lesson VII}
\switchcolumn*

\LA
\LessonTitle{Léctio sancti Evangélii secúndum Matthǽum}
\switchcolumn
\EN
\LessonTitle{From the Holy Gospel according to St. Matthew}
\switchcolumn*

\LA
\Cite{Matt 25:1-13}
\switchcolumn
\EN
\Cite{Matt 25:1}
\switchcolumn*

\LA
\noindent In illo témpore: Dixit Jesus discípulis suis parábolam hanc: Símile erit regnum cælórum decem virgínibus, quæ, accipiéntes lámpades suas, exiérunt óbviam sponso et sponsæ. Et réliqua.\par
\switchcolumn
\EN
\noindent At that time, Jesus said to His disciples: The Kingdom of heaven shall be likened unto ten virgins, which took their lamps, and went forth to meet the Bridegroom and the Bride. And so on.\par
\switchcolumn*

\LA
\LessonTitle{Homilía sancti Gregórii Papæ}
\switchcolumn
\EN
\LessonTitle{Homily by Pope St. Gregory the Great}
\switchcolumn*

\LA
\Source{Homilía 12 in Evangelia}
\switchcolumn
\EN
\Source{12th on the Gospels}
\switchcolumn*

\LA
\noindent Sæpe vos, fratres caríssimi, admóneo prava ópera fúgere, mundi hujus inquinaménta devitáre, sed hodiérna sancti Evangélii lectióne compéllor dícere, ut, et bona, quæ ágitis, cum magna cautéla teneátis; ne per hoc, quod a vobis rectum géritur, favor aut grátia humána requirátur; ne appetítus laudis subrépat, et, quod foris osténditur, intus a mercéde vacuétur. Ecce enim Redemptóris voce decem vírgines, et omnes dicúntur vírgines, et tamen intra beatitúdinis jánuam non omnes sunt recéptæ; quia eárum quædam, dum de virginitáte sua glóriam foris éxpetunt, in vasis suis óleum habére noluérunt.\par
\switchcolumn
\EN
\noindent Dearly beloved brethren oftentimes do I warn you to fly corrupt conversation, and to keep yourselves unspotted from the world. But the portion which is this day read from the Holy Gospel doth oblige me to say that even to these good things which ye do, ye must needs take all careful heed. Look ye well to it, that, when ye work righteousness, ye do it not as seeking the praise and admiration of men, for if the lust of praise do once creep in, that which seemeth so fair without, loseth its reward within. Behold how the Redeemer speaketh of these ten virgins. He calleth them all virgins, yet entered not all of them into the door of blessedness, for there were some of them who sought outwardly the honour of virginity, but would take no oil within their vessels with their lamps.\par
\switchcolumn*

\LA
\begin{Resp}\Rsign Hæc est Virgo sápiens, quam Dóminus vigilántem invénit, quæ accéptis lampádibus sumpsit secum óleum: \Star{} Et veniénte Dómino, introívit cum eo ad núptias.\end{Resp}
\begin{Resp}\Vsign Média nocte clamor factus est: Ecce sponsus venit, exíte óbviam ei.\end{Resp}
\begin{Resp}\Rsign Et veniénte Dómino, introívit cum eo ad núptias.\end{Resp}
\switchcolumn
\EN
\begin{Resp}\Rsign This is one of those wise virgins, whom the Lord found watching, for when she took her lamp, she took oil with her. \Star{} And when the Lord came, she went in with him to the marriage.\end{Resp}
\begin{Resp}\Vsign At midnight there was a cry made: Behold, the Bridegroom cometh, go ye out to meet him.\end{Resp}
\begin{Resp}\Rsign And when the Lord came, she went in with him to the marriage.\end{Resp}
\switchcolumn*

\LA
\LessonHead{Lectio VIII}
\switchcolumn
\EN
\LessonHead{Lesson VIII}
\switchcolumn*

\LA
\noindent Sed prius quæréndum nobis est quid sit regnum cælórum, aut cur decem virgínibus comparétur, quæ étiam vírgines prudéntes et fátuæ dicántur. Dum enim cælórum regnum constat quia reprobórum nullus ingréditur, étiam fátuis virgínibus cur símile esse perhibétur? Sed sciéndum nobis est quod sæpe in sacro elóquio regnum cælórum præséntis témporis Ecclésia dícitur. De quo álio in loco Dóminus dicit: Mittet Fílius hóminis Ángelos suos, et cólligent de regno ejus ómnia scándala. Neque enim in illo regno beatitúdinis, in quo pax summa est, inveníri scándala póterunt, quæ colligántur.\par
\switchcolumn
\EN
\noindent First of all, it is for us to ask What is the kingdom of Heaven? And wherefore shall the same be likened unto ten virgins, whereof, albeit five were wise, yet five were foolish For if the kingdom of heaven be such that there shall in no wise enter into it anything that defileth, neither whatsoever worketh abomination, or maketh a lie, (Apoc. xxi. 27,) how can it be like unto five virgins which were foolish? But we must know that, in the word of God, the kingdom of heaven doth oftentimes signify the Church as she now is, touching the which the Lord saith in another place “The Son of Man shall send forth His Angels, and they shall gather out of His kingdom all things that offend.” (Matth. xiii. 41.) In that kingdom of Blessedness, wherein peace shall have her perfect reign, there shall be nothing found that offendeth for the angels to gather out.\par
\switchcolumn*

\LA
\begin{Resp}\Rsign Média nocte clamor factus est: \Star{} Ecce sponsus venit, exíte óbviam ei.\end{Resp}
\begin{Resp}\Vsign Prudéntes vírgines, aptáte vestras lámpades.\end{Resp}
\begin{Resp}\Rsign Ecce sponsus venit, exíte óbviam ei.\end{Resp}
\begin{Resp}\Vsign Glória Patri, et Fílio, \Star{} et Spirítui Sancto.\end{Resp}
\begin{Resp}\Rsign Ecce sponsus venit, exíte óbviam ei.\end{Resp}
\switchcolumn
\EN
\begin{Resp}\Rsign At midnight there was a cry made: \Star{} Behold, the Bridegroom cometh, go ye out to meet him.\end{Resp}
\begin{Resp}\Vsign Trim your lamps, O ye wise virgins.\end{Resp}
\begin{Resp}\Rsign Behold, the Bridegroom cometh, go ye out to meet him.\end{Resp}
\begin{Resp}\Vsign Glory be to the Father, and to the Son, \Star{} and to the Holy Ghost.\end{Resp}
\begin{Resp}\Rsign Behold, the Bridegroom cometh, go ye out to meet him.\end{Resp}
\switchcolumn*

\LA
\LessonHead{Lectio IX}
\switchcolumn
\EN
\LessonHead{Lesson IX}
\switchcolumn*

\LA
\noindent In quinque autem córporis sénsibus unusquísque subsístit; geminátus autem quinárius denárium pérficit. Et, quia ex utróque sexu fidélium multitúdo collígitur, sancta Ecclésia decem virgínibus símilis esse denuntiátur. In qua quia mali cum bonis et réprobi cum eléctis admíxti sunt, recte símilis virgínibus prudéntibus et fátuis esse perhibétur. Sunt namque pleríque continéntes, qui ab appetítu se exterióri custódiunt et spe ad interióra rapiúntur, carnem mácerant, et toto desidério ad supérnam pátriam anhélant, ætérna prǽmia éxpetunt, pro labóribus suis recípere laudes humánas nolunt. Hi nimírum glóriam suam non in ore hóminum ponunt, sed intra consciéntiam cóntegunt. Et sunt pleríque, qui corpus per abstinéntiam afflígunt, sed de ipsa sua abstinéntia humános favóres éxpetunt.\par
\switchcolumn
\EN
\noindent The body of every man doth consist of five senses, and five being doubled, is ten. Forasmuch, therefore, as the whole body of the faithful doth consist of two sexes, the Holy Church is likened unto ten virgins. And forasmuch as in the Church the good are for the present mingled with the bad, and the reprobate with the elect, it is rightly said that, of the ten virgins, five are wise and five are foolish. There are many who have self-control, which do keep themselves from lusting after things outward, whose hope beareth them to things inward, who chastise the flesh, who long with intense home-sickness for their Fatherland which is in heaven, who seek an eternal reward, and who will not to receive for their labours the praise of men. These are they who reckon their glory, not in the mouths of men, but in the testimony of their own conscience. And many there be likewise who afflict the body by self-control, and yet who seek for their self-control applause from men.\par
\switchcolumn*

\LA
\TeDeum{Te Deum}
\switchcolumn
\EN
\TeDeum{Te Deum}
\switchcolumn*

\LA
\LessonHead{Lectio VII (in 2 loco)}
\switchcolumn
\EN
\LessonHead{Lesson VII (set 2)}
\switchcolumn*

\LA
\LessonTitle{Léctio sancti Evangélii secúndum Matthǽum}
\switchcolumn
\EN
\LessonTitle{From the Holy Gospel according to Matthew}
\switchcolumn*

\LA
\Cite{Matt 13:44-52}
\switchcolumn
\EN
\Cite{Matt 13:44}
\switchcolumn*

\LA
\noindent In illo témpore: Dixit Jesus discípulis suis parábolam hanc: Simile est regnum cælórum thesáuro abscóndito in agro. Et réliqua.\par
\switchcolumn
\EN
\noindent At that time Jesus spake unto His disciples this parable The kingdom of heaven is like unto treasure hid in a field. And so on.\par
\switchcolumn*

\LA
\LessonTitle{Homilía sancti Gregórii Papæ}
\switchcolumn
\EN
\LessonTitle{Homily by Pope St Gregory the Great}
\switchcolumn*

\LA
\Source{Homilia 11 in Evangelia}
\switchcolumn
\EN
\Source{11th on the Gospels}
\switchcolumn*

\LA
\noindent Cælórum regnum, fratres caríssimi, idcírco terrénis rebus símile dícitur, ut, ex his quæ ánimus novit, surgat ad incógnita, quæ non novit: quátenus exémplo visibílium se ad invisibília rápiat, et, per ea quæ usu dídicit quasi confricátus incaléscat; ut per hoc, quod scit notum dilígere, discat et incógnita amáre. Ecce enim cælórum regnum thesáuro abscóndito in agro comparátur; quem, qui invénit homo, abscóndit, et præ gáudio illíus vadit, et vendit univérsa quæ habet, et emit agrum illum.\par
\switchcolumn
\EN
\noindent Dearly beloved brethren, the kingdom of heaven is likened unto the things of earth, to the end that by the mean of things which we know, our mind may rise to the contemplation of the things which we know not by the example of things which are seen, may fix her gaze on things which are not seen by the touch of things which she useth, may be warmed towards the things which she useth not; by things which she knoweth and loveth, to love also the things which she knoweth not. For, behold, “the kingdom of heaven is likened unto treasure hid in a field, the which when a man hath found, he hideth, and, for joy thereof, goeth and selleth all that he hath and buyeth that field.”\par
\switchcolumn*

\LA
\begin{Resp}\Rsign Hæc est Virgo sápiens, quam Dóminus vigilántem invénit, quæ accéptis lampádibus sumpsit secum óleum: \Star{} Et veniénte Dómino, introívit cum eo ad núptias.\end{Resp}
\begin{Resp}\Vsign Média nocte clamor factus est: Ecce sponsus venit, exíte óbviam ei.\end{Resp}
\begin{Resp}\Rsign Et veniénte Dómino, introívit cum eo ad núptias.\end{Resp}
\switchcolumn
\EN
\begin{Resp}\Rsign This is one of those wise virgins, whom the Lord found watching, for when she took her lamp, she took oil with her. \Star{} And when the Lord came, she went in with him to the marriage.\end{Resp}
\begin{Resp}\Vsign At midnight there was a cry made: Behold, the Bridegroom cometh, go ye out to meet him.\end{Resp}
\begin{Resp}\Rsign And when the Lord came, she went in with him to the marriage.\end{Resp}
\switchcolumn*

\LA
\LessonHead{Lectio VIII (in 2 loco)}
\switchcolumn
\EN
\LessonHead{Lesson VIII (set 2)}
\switchcolumn*

\LA
\noindent Qua in re hoc quoque notándum est, quod invéntus thesáurus abscónditur, ut servétur: quia stúdium cæléstis desidérii a malígnis spirítibus custodíre non súfficit, qui hoc ab humánis láudibus non abscóndit. In præsénti étenim vita, quasi in via sumus, qua ad pátriam pérgimus. Malígni autem spíritus iter nostrum quasi quidam latrúnculi óbsident. Deprædári ergo desíderat, qui thesáurum públice portat in via. Hoc autem dico, non ut próximi ópera nostra bona non vídeant, cum scriptum sit: Vídeant ópera vestra bona, et gloríficent Patrem vestrum qui in cælis est; sed, ut per hoc quod ágimus, laudes extérius non quærámus. Sic autem sit opus in público, quátenus inténtio máneat in occúlto; ut, et de bono ópere próximis præbeámus exémplum, et tamen per intentiónem, qua Deo soli placére quǽrimus, semper optémus secrétum.\par
\switchcolumn
\EN
\noindent And herein we must remark that the treasure, when once it hath been found, is hidden to keep it safe. He whose intimate yearnings after God are not hidden from the praise of men, is open thereby to the attacks of evil spirits. In this life we are, as it were, journeying homewards on a road beset by evil spirits who are like highwaymen. He therefore inviteth robbery who carrieth his treasure ostentatiously. Doubtless our neighbour should be able to see our good works, as it is written: “Let your light so shine before men that they may see your good works, and glorify your Father which is in heaven.” But this is not to be understood to mean that we are to seek the praise of men by what we do. Rather, let us in such wise work in the open that the inner intention of devotion is not advertised. So we shall give an example to our neighbour, and yet keep hidden, except from the sight of God, our purpose of pleasing him.\par
\switchcolumn*

\LA
\begin{Resp}\Rsign Média nocte clamor factus est: \Star{} Ecce sponsus venit, exíte óbviam ei.\end{Resp}
\begin{Resp}\Vsign Prudéntes vírgines, aptáte vestras lámpades.\end{Resp}
\begin{Resp}\Rsign Ecce sponsus venit, exíte óbviam ei.\end{Resp}
\begin{Resp}\Vsign Glória Patri, et Fílio, \Star{} et Spirítui Sancto.\end{Resp}
\begin{Resp}\Rsign Ecce sponsus venit, exíte óbviam ei.\end{Resp}
\switchcolumn
\EN
\begin{Resp}\Rsign At midnight there was a cry made: \Star{} Behold, the Bridegroom cometh, go ye out to meet him.\end{Resp}
\begin{Resp}\Vsign Trim your lamps, O ye wise virgins.\end{Resp}
\begin{Resp}\Rsign Behold, the Bridegroom cometh, go ye out to meet him.\end{Resp}
\begin{Resp}\Vsign Glory be to the Father, and to the Son, \Star{} and to the Holy Ghost.\end{Resp}
\begin{Resp}\Rsign Behold, the Bridegroom cometh, go ye out to meet him.\end{Resp}
\switchcolumn*

\LA
\LessonHead{Lectio IX (in 2 loco)}
\switchcolumn
\EN
\LessonHead{Lesson IX (set 2)}
\switchcolumn*

\LA
\noindent Thesáurus autem cæléste est desidérium; ager vero, in quo thesáurus abscónditur, disciplína stúdii cæléstis. Quem profécto agrum, vénditis ómnibus, cómparat, qui, voluptátibus carnis renúntians, cuncta sua terréna desidéria per disciplínæ cæléstis custódiam calcat: ut nihil jam quod caro blandítur, líbeat; nihil, quod carnálem vitam trucídat, spíritus perhorréscat.\par
\switchcolumn
\EN
\noindent The treasure is the desire for heaven, the field wherein it is hidden is the earnest observance wherewith this desire is surrounded. Whosoever turneth his back upon the enjoyments of the flesh, and by earnest striving heavenward, putteth all earthly lusts under the feet of discipline, so that he smileth back no more when the flesh smileth at him, and shuddereth no more at anything that can only kill the body whosoever doth thus, hath sold all that he had, and bought that field.\par
\switchcolumn*

\LA
\TeDeum{Te Deum}
\switchcolumn
\EN
\TeDeum{Te Deum}
\switchcolumn*

\end{paracol}

\Office{Commune Virginum tantum}{Common of a Virgin not a Martyr}{}
\addcontentsline{toc}{subsection}{Commune Virginum tantum \textit{— Common of a Virgin not a Martyr}}
\begin{paracol}{2}
\LA
\LessonHead{Lectio VII}
\switchcolumn
\EN
\LessonHead{Lesson VII}
\switchcolumn*

\LA
\LessonTitle{Léctio sancti Evangélii secúndum Matthǽum}
\switchcolumn
\EN
\LessonTitle{From the Holy Gospel according to St. Matthew}
\switchcolumn*

\LA
\Cite{Matt 25:1-13}
\switchcolumn
\EN
\Cite{Matt 25:1}
\switchcolumn*

\LA
\noindent In illo témpore: Dixit Jesus discípulis suis parábolam hanc: Símile erit regnum cælórum decem virgínibus, quæ, accipiéntes lámpades suas, exiérunt óbviam sponso et sponsæ. Et réliqua.\par
\switchcolumn
\EN
\noindent At that time, Jesus said to His disciples: The Kingdom of heaven shall be likened unto ten virgins, which took their lamps, and went forth to meet the Bridegroom and the Bride. And so on.\par
\switchcolumn*

\LA
\LessonTitle{Homilía sancti Gregórii Papæ}
\switchcolumn
\EN
\LessonTitle{Homily by Pope St. Gregory the Great}
\switchcolumn*

\LA
\Source{Homilía 12 in Evangelia}
\switchcolumn
\EN
\Source{12th on the Gospels}
\switchcolumn*

\LA
\noindent Sæpe vos, fratres caríssimi, admóneo prava ópera fúgere, mundi hujus inquinaménta devitáre, sed hodiérna sancti Evangélii lectióne compéllor dícere, ut, et bona, quæ ágitis, cum magna cautéla teneátis; ne per hoc, quod a vobis rectum géritur, favor aut grátia humána requirátur; ne appetítus laudis subrépat, et, quod foris osténditur, intus a mercéde vacuétur. Ecce enim Redemptóris voce decem vírgines, et omnes dicúntur vírgines, et tamen intra beatitúdinis jánuam non omnes sunt recéptæ; quia eárum quædam, dum de virginitáte sua glóriam foris éxpetunt, in vasis suis óleum habére noluérunt.\par
\switchcolumn
\EN
\noindent Dearly beloved brethren oftentimes do I warn you to fly corrupt conversation, and to keep yourselves unspotted from the world. But the portion which is this day read from the Holy Gospel doth oblige me to say that even to these good things which ye do, ye must needs take all careful heed. Look ye well to it, that, when ye work righteousness, ye do it not as seeking the praise and admiration of men, for if the lust of praise do once creep in, that which seemeth so fair without, loseth its reward within. Behold how the Redeemer speaketh of these ten virgins. He calleth them all virgins, yet entered not all of them into the door of blessedness, for there were some of them who sought outwardly the honour of virginity, but would take no oil within their vessels with their lamps.\par
\switchcolumn*

\LA
\begin{Resp}\Rsign Hæc est Virgo sápiens, quam Dóminus vigilántem invénit, quæ accéptis lampádibus sumpsit secum óleum: \Star{} Et veniénte Dómino, introívit cum eo ad núptias.\end{Resp}
\begin{Resp}\Vsign Média nocte clamor factus est: Ecce sponsus venit, exíte óbviam ei.\end{Resp}
\begin{Resp}\Rsign Et veniénte Dómino, introívit cum eo ad núptias.\end{Resp}
\switchcolumn
\EN
\begin{Resp}\Rsign This is one of those wise virgins, whom the Lord found watching, for when she took her lamp, she took oil with her. \Star{} And when the Lord came, she went in with him to the marriage.\end{Resp}
\begin{Resp}\Vsign At midnight there was a cry made: Behold, the Bridegroom cometh, go ye out to meet him.\end{Resp}
\begin{Resp}\Rsign And when the Lord came, she went in with him to the marriage.\end{Resp}
\switchcolumn*

\LA
\LessonHead{Lectio VIII}
\switchcolumn
\EN
\LessonHead{Lesson VIII}
\switchcolumn*

\LA
\noindent Sed prius quæréndum nobis est quid sit regnum cælórum, aut cur decem virgínibus comparétur, quæ étiam vírgines prudéntes et fátuæ dicántur. Dum enim cælórum regnum constat quia reprobórum nullus ingréditur, étiam fátuis virgínibus cur símile esse perhibétur? Sed sciéndum nobis est quod sæpe in sacro elóquio regnum cælórum præséntis témporis Ecclésia dícitur. De quo álio in loco Dóminus dicit: Mittet Fílius hóminis Ángelos suos, et cólligent de regno ejus ómnia scándala. Neque enim in illo regno beatitúdinis, in quo pax summa est, inveníri scándala póterunt, quæ colligántur.\par
\switchcolumn
\EN
\noindent First of all, it is for us to ask What is the kingdom of Heaven? And wherefore shall the same be likened unto ten virgins, whereof, albeit five were wise, yet five were foolish For if the kingdom of heaven be such that there shall in no wise enter into it anything that defileth, neither whatsoever worketh abomination, or maketh a lie, (Apoc. xxi. 27,) how can it be like unto five virgins which were foolish? But we must know that, in the word of God, the kingdom of heaven doth oftentimes signify the Church as she now is, touching the which the Lord saith in another place “The Son of Man shall send forth His Angels, and they shall gather out of His kingdom all things that offend.” (Matth. xiii. 41.) In that kingdom of Blessedness, wherein peace shall have her perfect reign, there shall be nothing found that offendeth for the angels to gather out.\par
\switchcolumn*

\LA
\begin{Resp}\Rsign Média nocte clamor factus est: \Star{} Ecce sponsus venit, exíte óbviam ei.\end{Resp}
\begin{Resp}\Vsign Prudéntes vírgines, aptáte vestras lámpades.\end{Resp}
\begin{Resp}\Rsign Ecce sponsus venit, exíte óbviam ei.\end{Resp}
\begin{Resp}\Vsign Glória Patri, et Fílio, \Star{} et Spirítui Sancto.\end{Resp}
\begin{Resp}\Rsign Ecce sponsus venit, exíte óbviam ei.\end{Resp}
\switchcolumn
\EN
\begin{Resp}\Rsign At midnight there was a cry made: \Star{} Behold, the Bridegroom cometh, go ye out to meet him.\end{Resp}
\begin{Resp}\Vsign Trim your lamps, O ye wise virgins.\end{Resp}
\begin{Resp}\Rsign Behold, the Bridegroom cometh, go ye out to meet him.\end{Resp}
\begin{Resp}\Vsign Glory be to the Father, and to the Son, \Star{} and to the Holy Ghost.\end{Resp}
\begin{Resp}\Rsign Behold, the Bridegroom cometh, go ye out to meet him.\end{Resp}
\switchcolumn*

\LA
\LessonHead{Lectio IX}
\switchcolumn
\EN
\LessonHead{Lesson IX}
\switchcolumn*

\LA
\noindent In quinque autem córporis sénsibus unusquísque subsístit; geminátus autem quinárius denárium pérficit. Et, quia ex utróque sexu fidélium multitúdo collígitur, sancta Ecclésia decem virgínibus símilis esse denuntiátur. In qua quia mali cum bonis et réprobi cum eléctis admíxti sunt, recte símilis virgínibus prudéntibus et fátuis esse perhibétur. Sunt namque pleríque continéntes, qui ab appetítu se exterióri custódiunt et spe ad interióra rapiúntur, carnem mácerant, et toto desidério ad supérnam pátriam anhélant, ætérna prǽmia éxpetunt, pro labóribus suis recípere laudes humánas nolunt. Hi nimírum glóriam suam non in ore hóminum ponunt, sed intra consciéntiam cóntegunt. Et sunt pleríque, qui corpus per abstinéntiam afflígunt, sed de ipsa sua abstinéntia humános favóres éxpetunt.\par
\switchcolumn
\EN
\noindent The body of every man doth consist of five senses, and five being doubled, is ten. Forasmuch, therefore, as the whole body of the faithful doth consist of two sexes, the Holy Church is likened unto ten virgins. And forasmuch as in the Church the good are for the present mingled with the bad, and the reprobate with the elect, it is rightly said that, of the ten virgins, five are wise and five are foolish. There are many who have self-control, which do keep themselves from lusting after things outward, whose hope beareth them to things inward, who chastise the flesh, who long with intense home-sickness for their Fatherland which is in heaven, who seek an eternal reward, and who will not to receive for their labours the praise of men. These are they who reckon their glory, not in the mouths of men, but in the testimony of their own conscience. And many there be likewise who afflict the body by self-control, and yet who seek for their self-control applause from men.\par
\switchcolumn*

\LA
\TeDeum{Te Deum}
\switchcolumn
\EN
\TeDeum{Te Deum}
\switchcolumn*

\end{paracol}

\Office{Commune non Virginum non Martyrum}{Common of Widows}{}
\addcontentsline{toc}{subsection}{Commune non Virginum non Martyrum \textit{— Common of Widows}}
\begin{paracol}{2}
\LA
\LessonHead{Lectio VII}
\switchcolumn
\EN
\LessonHead{Lesson VII}
\switchcolumn*

\LA
\LessonTitle{Léctio sancti Evangélii secúndum Matthǽum}
\switchcolumn
\EN
\LessonTitle{From the Holy Gospel according to Matthew}
\switchcolumn*

\LA
\Cite{Matt 13:44-52}
\switchcolumn
\EN
\Cite{Matt 13:44}
\switchcolumn*

\LA
\noindent In illo témpore: Dixit Jesus discípulis suis parábolam hanc: Simile est regnum cælórum thesáuro abscóndito in agro. Et réliqua.\par
\switchcolumn
\EN
\noindent At that time Jesus spake unto His disciples this parable The kingdom of heaven is like unto treasure hid in a field. And so on.\par
\switchcolumn*

\LA
\LessonTitle{Homilía sancti Gregórii Papæ}
\switchcolumn
\EN
\LessonTitle{Homily by Pope St Gregory the Great}
\switchcolumn*

\LA
\Source{Homilia 11 in Evangelia}
\switchcolumn
\EN
\Source{11th on the Gospels}
\switchcolumn*

\LA
\noindent Cælórum regnum, fratres caríssimi, idcírco terrénis rebus símile dícitur, ut, ex his quæ ánimus novit, surgat ad incógnita, quæ non novit: quátenus exémplo visibílium se ad invisibília rápiat, et, per ea quæ usu dídicit quasi confricátus incaléscat; ut per hoc, quod scit notum dilígere, discat et incógnita amáre. Ecce enim cælórum regnum thesáuro abscóndito in agro comparátur; quem, qui invénit homo, abscóndit, et præ gáudio illíus vadit, et vendit univérsa quæ habet, et emit agrum illum.\par
\switchcolumn
\EN
\noindent Dearly beloved brethren, the kingdom of heaven is likened unto the things of earth, to the end that by the mean of things which we know, our mind may rise to the contemplation of the things which we know not by the example of things which are seen, may fix her gaze on things which are not seen by the touch of things which she useth, may be warmed towards the things which she useth not; by things which she knoweth and loveth, to love also the things which she knoweth not. For, behold, “the kingdom of heaven is likened unto treasure hid in a field, the which when a man hath found, he hideth, and, for joy thereof, goeth and selleth all that he hath and buyeth that field.”\par
\switchcolumn*

\LA
\begin{Resp}\Rsign Os suum apéruit sapiéntiæ, et lex cleméntiæ in lingua ejus: considerávit sémitas domus suæ, \Star{} Et panem otiósa non comédit.\end{Resp}
\begin{Resp}\Vsign Gustávit et vidit quia bona est negotiátio ejus: non exstinguétur in nocte lucérna ejus.\end{Resp}
\begin{Resp}\Rsign Et panem otiósa non comédit.\end{Resp}
\switchcolumn
\EN
\begin{Resp}\Rsign She openeth her mouth with wisdom, and in her tongue is the law of kindness. She looketh well to the ways of her household \Star{} And eateth not the bread of idleness.\end{Resp}
\begin{Resp}\Vsign She tasteth and perceiveth that her merchandise is good. Her candle goeth not out by night.\end{Resp}
\begin{Resp}\Rsign And she eateth not the bread of idleness.\end{Resp}
\switchcolumn*

\LA
\LessonHead{Lectio VIII}
\switchcolumn
\EN
\LessonHead{Lesson VIII}
\switchcolumn*

\LA
\noindent Qua in re hoc quoque notándum est, quod invéntus thesáurus abscónditur, ut servétur: quia stúdium cæléstis desidérii a malígnis spirítibus custodíre non súfficit, qui hoc ab humánis láudibus non abscóndit. In præsénti étenim vita, quasi in via sumus, qua ad pátriam pérgimus. Malígni autem spíritus iter nostrum quasi quidam latrúnculi óbsident. Deprædári ergo desíderat, qui thesáurum públice portat in via. Hoc autem dico, non ut próximi ópera nostra bona non vídeant, cum scriptum sit: Vídeant ópera vestra bona, et gloríficent Patrem vestrum qui in cælis est; sed, ut per hoc quod ágimus, laudes extérius non quærámus. Sic autem sit opus in público, quátenus inténtio máneat in occúlto; ut, et de bono ópere próximis præbeámus exémplum, et tamen per intentiónem, qua Deo soli placére quǽrimus, semper optémus secrétum.\par
\switchcolumn
\EN
\noindent And herein we must remark that the treasure, when once it hath been found, is hidden to keep it safe. He whose intimate yearnings after God are not hidden from the praise of men, is open thereby to the attacks of evil spirits. In this life we are, as it were, journeying homewards on a road beset by evil spirits who are like highwaymen. He therefore inviteth robbery who carrieth his treasure ostentatiously. Doubtless our neighbour should be able to see our good works, as it is written: “Let your light so shine before men that they may see your good works, and glorify your Father which is in heaven.” But this is not to be understood to mean that we are to seek the praise of men by what we do. Rather, let us in such wise work in the open that the inner intention of devotion is not advertised. So we shall give an example to our neighbour, and yet keep hidden, except from the sight of God, our purpose of pleasing him.\par
\switchcolumn*

\LA
\begin{Resp}\Rsign Regnum mundi et omnem ornátum sǽculi contémpsi, propter amórem Dómini mei Jesu Christi: \Star{} Quem vidi, quem amávi, in quem crédidi, quem diléxi.\end{Resp}
\begin{Resp}\Vsign Eructávit cor meum verbum bonum: dico ego ópera mea Regi.\end{Resp}
\begin{Resp}\Rsign Quem vidi, quem amávi, in quem crédidi, quem diléxi.\end{Resp}
\begin{Resp}\Vsign Glória Patri, et Fílio, \Star{} et Spirítui Sancto.\end{Resp}
\begin{Resp}\Rsign Quem vidi, quem amávi, in quem crédidi, quem diléxi.\end{Resp}
\switchcolumn
\EN
\begin{Resp}\Rsign The kingdom of this world and all the beauty of life I have esteemed as nothing, for the excellency of the love of Jesus Christ my Lord \Star{} Whom, having seen, I loved Whom, having believed, I longed after.\end{Resp}
\begin{Resp}\Vsign My heart is overflowing with a good matter I speak of my works unto the King.\end{Resp}
\begin{Resp}\Rsign Whom, having seen, I loved Whom, having believed, I longed after.\end{Resp}
\begin{Resp}\Vsign Glory be to the Father, and to the Son, \Star{} and to the Holy Ghost.\end{Resp}
\begin{Resp}\Rsign Whom, having seen, I loved Whom, having believed, I longed after.\end{Resp}
\switchcolumn*

\LA
\LessonHead{Lectio IX}
\switchcolumn
\EN
\LessonHead{Lesson IX}
\switchcolumn*

\LA
\noindent Thesáurus autem cæléste est desidérium; ager vero, in quo thesáurus abscónditur, disciplína stúdii cæléstis. Quem profécto agrum, vénditis ómnibus, cómparat, qui, voluptátibus carnis renúntians, cuncta sua terréna desidéria per disciplínæ cæléstis custódiam calcat: ut nihil jam quod caro blandítur, líbeat; nihil, quod carnálem vitam trucídat, spíritus perhorréscat.\par
\switchcolumn
\EN
\noindent The treasure is the desire for heaven, the field wherein it is hidden is the earnest observance wherewith this desire is surrounded. Whosoever turneth his back upon the enjoyments of the flesh, and by earnest striving heavenward, putteth all earthly lusts under the feet of discipline, so that he smileth back no more when the flesh smileth at him, and shuddereth no more at anything that can only kill the body whosoever doth thus, hath sold all that he had, and bought that field.\par
\switchcolumn*

\LA
\TeDeum{Te Deum}
\switchcolumn
\EN
\TeDeum{Te Deum}
\switchcolumn*

\end{paracol}

\Office{Commune Dedicationis Ecclesiae}{Common of the Dedication of Churches}{Duplex I classis}
\addcontentsline{toc}{subsection}{Commune Dedicationis Ecclesiae \textit{— Common of the Dedication of Churches}}
\begin{paracol}{2}
\LA
\LessonHead{Lectio VII}
\switchcolumn
\EN
\LessonHead{Lesson VII}
\switchcolumn*

\LA
\LessonTitle{Léctio sancti Evangélii secúndum Lucam}
\switchcolumn
\EN
\LessonTitle{From the Holy Gospel according to Luke}
\switchcolumn*

\LA
\Cite{Luc 19:1-10}
\switchcolumn
\EN
\Cite{Luke 19:1-10}
\switchcolumn*

\LA
\noindent In illo témpore: Ingréssus Jesus perambulábat Jéricho. Et ecce vir nómine Zachǽus: et hic princeps erat publicanórum, et ipse dives. Et réliqua.\par
\switchcolumn
\EN
\noindent At that time: Jesus entered and passed through Jericho. And, behold, there was a man named Zacchaeus, which was the chief among the publicans, and he was rich. And so on.\par
\switchcolumn*

\LA
\LessonTitle{Homilía sancti Ambrósii Epíscopi}
\switchcolumn
\EN
\LessonTitle{Homily by St. Ambrose, Bishop (of Milan.)}
\switchcolumn*

\LA
\Source{Lib. 8 in Luc., prope finem}
\switchcolumn
\EN
\Source{Bk. viii. on Luke.}
\switchcolumn*

\LA
\noindent Zachǽus, statúra pusíllus, hoc est, nulla nobilitátis ingénitæ dignitáte sublímis, exíguus méritis sicut pópulus natiónum, audíto Dómini Salvatóris advéntu, quem sui non recéperant, vidére cupiébat. Sed nemo fácile Jesum videt; nemo potest Jesum vidére constitútus in terra. Et, quia non prophétas, non legem habébat, tamquam formæ grátiam naturális, ascéndit in sycómorum, vanitátem scílicet Judæórum vestígio suo próterens, erráta quoque córrigens superióris ætátis. Et ídeo Jesum in interióris domus recépit hospítio.\par
\switchcolumn
\EN
\noindent Zacchaeus was little of stature, that is, he was not raised aloft among men by nobility of birth, and, like the most of the world, he possessed few merits. When he heard that the Lord and Saviour, Who had come unto His Own, and Whom His Own had not received, (John i. 11,) was coming, he desired to see Him. But the sight of Jesus is not easy; to any on the earth it is impossible. And since Zacchaeus had neither the Prophets, nor yet the Law, as a gracious help to his nature, he climbed up into a sycamore tree, raising his feet above the vanity of the Jews, and straightening the crooked branches of his former life, and therefore he received Jesus to lodge within his house.\par
\switchcolumn*

\LA
\begin{Resp}\Rsign Domus mea, domus oratiónis vocábitur, dicit Dóminus: in ea omnis qui petit, áccipit; et qui quærit, ínvenit; \Star{} Et pulsánti aperiétur.\end{Resp}
\begin{Resp}\Vsign Pétite, et accipiétis; quǽrite, et inveniétis.\end{Resp}
\begin{Resp}\Rsign Et pulsánti aperiétur.\end{Resp}
\switchcolumn
\EN
\begin{Resp}\Rsign My house shall be called the house of prayer, saith the Lord. Therein, he that asketh, receiveth; he that seeketh, findeth; \Star{} And to him that knocketh, it shall be opened.\end{Resp}
\begin{Resp}\Vsign Ask, and ye shall receive; seek, and ye shall find.\end{Resp}
\begin{Resp}\Rsign And to him that knocketh, it shall be opened.\end{Resp}
\switchcolumn*

\LA
\LessonHead{Lectio VIII}
\switchcolumn
\EN
\LessonHead{Lesson VIII}
\switchcolumn*

\LA
\noindent Et bene ascéndit in árborem, ut arbor bona bonos fructus fáceret, ac, naturáli excísus oleástro et contra natúram insértus in bonam olívam, fructum posset legis afférre. Radix enim sancta, etsi rami inútiles. Quorum infructuósam glóriam plebs géntium, fide resurrectiónis, quasi quadam córporis elevatióne, transcéndit. Zachǽus ergo in sycómoro, cæcus in via: quorum álterum Dóminus miseratúrus exspéctat, álterum mansiónis suæ claritáte nobílitat; álterum sanatúrus intérrogat, apud álterum se, non invitátus, invítat. Sciébat enim úberem hospítii sui esse mercédem. Sed tamen, etsi nondum vocem invitántis audíerat, jam víderat afféctum.\par
\switchcolumn
\EN
\noindent He did well to climb up into a tree, that a good tree might bring forth good fruits, (Matth. vii. 17,) and that the slip of the wild olive, grafted, contrary to nature, into the good olive, might bring forth the fruits of the law. (Rom. xi. 17, 24.) For the root is holy, however unprofitable the branches. Their barren beauty hath now been overshadowed by the belief of the Gentiles in the Resurrection, as by a material upgrowth. Zacchaeus, then, was in the sycamore tree, and the blind man by the way-side, (xviii. 35.) For the one, Jesus stood waiting to show mercy, and asked him before He healed him, what he would that He should do for him; being unbidden of the other, He bade Himself to be his Guest, knowing how rich was the reward of receiving Him. Nevertheless, albeit He had heard no words of invitation, yet had He seen how his heart went.\par
\switchcolumn*

\LA
\begin{Resp}\Rsign Lápides pretiósi omnes muri tui, \Star{} Et turres Jerúsalem gemmis ædificabúntur.\end{Resp}
\begin{Resp}\Vsign Portæ Jerúsalem ex sapphíro et smarágdo ædificabúntur, et ex lápide pretióso omnis circúitus muri ejus.\end{Resp}
\begin{Resp}\Rsign Et turres Jerúsalem gemmis ædificabúntur.\end{Resp}
\begin{Resp}\Vsign Glória Patri, et Fílio, \Star{} et Spirítui Sancto.\end{Resp}
\begin{Resp}\Rsign Et turres Jerúsalem gemmis ædificabúntur.\end{Resp}
\switchcolumn
\EN
\begin{Resp}\Rsign All thy walls are of stones most precious. \Star{} The towers of Jerusalem shall be built up with jewels.\end{Resp}
\begin{Resp}\Vsign The gates of Jerusalem shall be built up with the sapphire stone, and the emerald, and all her walls round about with stones most precious.\end{Resp}
\begin{Resp}\Rsign The towers of Jerusalem shall be built up with jewels.\end{Resp}
\begin{Resp}\Vsign Glory be to the Father, and to the Son, \Star{} and to the Holy Ghost.\end{Resp}
\begin{Resp}\Rsign The towers of Jerusalem shall be built up with jewels.\end{Resp}
\switchcolumn*

\end{paracol}

\Office{Sanctæ Mariæ Sabbato (a Trinitate usque ad Adventum)}{Our Lady's Saturday (from Trinity to Advent)}{Simplex}
\addcontentsline{toc}{subsection}{Sanctæ Mariæ Sabbato (a Trinitate usque ad Adventum) \textit{— Our Lady's Saturday (from Trinity to Advent)}}
\begin{paracol}{2}
\LA
\RefLine{Lectiones I–II: de Scriptura occurrente.}
\switchcolumn
\EN
\RefLine{Lessons I–II: from the Scripture of the day.}
\switchcolumn*

\LA
\LessonHead{Lectio III (mense Septembri)}
\switchcolumn
\EN
\LessonHead{Lesson III (in September)}
\switchcolumn*

\LA
\LessonTitle{Ex Epístola sancti Leónis Papæ ad Pulchériam Augústam}
\switchcolumn
\EN
\LessonTitle{From the epistle of Pope St. Leo to Pulcheria Augusta}
\switchcolumn*

\LA
\Source{Epist. 13, ante medium}
\switchcolumn
\EN
\Source{Epistle 13 in the first half}
\switchcolumn*

\LA
\noindent Sacraméntum reconcilatiónis nostræ, ante témpora ætérna dispósitum, nullæ implébant figúræ; quia nondum supervénerat Spíritus Sanctus in Vírginem nec virtus Altíssimi obumbráverat ei, ut, et intra intemeráta víscera, ædificánte sibi Sapiéntia domum, Verbum caro fíeret, et, forma Dei ac forma servi in unam conveniénte persónam, Creátor témporum nascerétur in témpore, et, per quem facta sunt ómnia, ipse inter ómnia gignerétur. Nisi enim novus homo, factus in similitúdinem carnis peccáti, nostram suscíperet vetustátem, et, consubstantiális Patri, consubstantiális esse dignarétur et matri, naturámque sibi nostram solus a peccáto liber uníret; sub jugo diáboli generáliter tenerétur humána captívas.\par
\switchcolumn
\EN
\noindent And the fulfillment of the mystery of our atonement, which was ordained from all eternity, was not assisted by any figures because the Holy Spirit had not yet come upon the Virgin, and the power of the Most High had not over-shadowed her: so that Wisdom building herself a house within her undefiled body, the Word became flesh; and the form of God and the form of a slave coming together into one person, the Creator of times was born in time; and He Himself through whom all things were made, was brought forth in the midst of all things. For if the New Man had not been made in the likeness of sinful flesh, and taken on Him our old nature, and being consubstantial with the Father, had deigned to be consubstantial with His mother also, and being alone free from sin, had united our nature to Him the whole human race would be held in bondage beneath the Devil's yoke.\par
\switchcolumn*

\LA
\TeDeum{Te Deum}
\switchcolumn
\EN
\TeDeum{Te Deum}
\switchcolumn*

\LA
\LessonHead{Lectio III (mense Octobri)}
\switchcolumn
\EN
\LessonHead{Lesson III (in October)}
\switchcolumn*

\LA
\LessonTitle{Sermo sancti Bernárdi Abbátis}
\switchcolumn
\EN
\LessonTitle{A sermon of St. Bernard, Abbot}
\switchcolumn*

\LA
\Source{Serm. in cap. 12 Apoc., ante med.}
\switchcolumn
\EN
\Source{Sermon of ch. 12 of the Apocalypse, in the first half}
\switchcolumn*

\LA
\noindent Amplectámur Maríæ vestígia, fratres mei, et devotíssima supplicatióne beátis illíus pédibus provolvámur. Teneámus eam nec dimittámus, donec benedíxerit nobis; potens est enim. Nempe vellus est médium inter rorem et áream: múlier inter solem et lunam: María inter Christum et Ecclésiam constitúta. Sed forte miráris, non tam vellus opértum rore, quam amíctam sole mulíerem. Magna síquidem familiáritas, sed mira omníno vicínitas solis et mulíeris. Quómodo enim in tam veheménti fervóre tam frágilis natúra subsístit? Mérito quidem admiráris, Móyses sancte, et curiósius desíderas intuéri. Verúmtamen solve calceaménta de pédibus tuis, et involúcra pone carnálium cogitatiónum, si accédere concupíscis.\par
\switchcolumn
\EN
\noindent My brothers, let us cast ourselves at Mary's blessed feet with the most devout supplication. Let us embrace them, let us hold fast to her, and not let go until she blesses us; for she has great power. Like the fleece between the dew and the threshing floor, like the woman between the sun and the moon, so Mary stands between Christ and the Church. But perhaps you wonder, not so much at the fleece covered with dew, but at the woman clothed with the sun. Any relationship between the sun and the woman would be striking enough, but this proximity is indeed something to be marveled at. How can so fragile a nature subsist in that excessive heat? You have a right to wonder, saintly Moses; it is only natural that in your curiosity you should want to look more closely. But “remove the sandals from your feet”, put aside the covering of fleshly thoughts, if you wish to draw near.\par
\switchcolumn*

\LA
\TeDeum{Te Deum}
\switchcolumn
\EN
\TeDeum{Te Deum}
\switchcolumn*

\LA
\LessonHead{Lectio III (mense Novembri)}
\switchcolumn
\EN
\LessonHead{Lesson III (in November)}
\switchcolumn*

\LA
\LessonTitle{De Expositióne S. Basilíi Epíscopi in Isaíam Prophétam}
\switchcolumn
\EN
\LessonTitle{From the Commentary of St. Basil, Bishop, on the Prophet Isaias}
\switchcolumn*

\LA
\Source{In cap. 8 post initium}
\switchcolumn
\EN
\Source{On chapter 8 near the beginning}
\switchcolumn*

\LA
\noindent Accéssi, inquit, ad prophetíssam, et in útero accépit et péperit fílium. Quod María prophetíssa fúerit, ad quam próxime accéssit Isaías per prænotiónem spíritus, nemo contradíxerit, qui sit memor verbórum Maríæ, quæ prophético affláta spíritu elocúta est. Quid enim ait? Magníficat ánima mea Dóminum: et exsultávit spíritus meus in Deo, salutári meo. Quia respéxit humilitátem ancíllæ suæ: ecce enim ex hoc beátam me dicent omnes generatiónes. Quod si ánimum accommodáveris univérsis ejus verbis, non útique per dissídium negáveris eam fuísse prophetíssam, quod Dómini Spíritus in eam supervénerit, et virtus Altíssimi obumbráverit ei.\par
\switchcolumn
\EN
\noindent “I went to the prophetess”, he says, “and she conceived and bore a son.” To see Mary in the prophetess whom Isaias approached by spiritual foreknowledge one need only recall the words Mary uttered when she was inspired by the spirit of prophecy: “My soul magnifies the Lord, and my spirit rejoices in God my Saviour; because He has regarded the lowliness of His handmaid; for behold, henceforth all generation shall call me blessed.” In fact, if you bear in mind all that she said, you can only agree that she was indeed the prophetess, since the spirit of Lord came upon her, and the power of the Most High overshadowed her.\par
\switchcolumn*

\LA
\TeDeum{Te Deum}
\switchcolumn
\EN
\TeDeum{Te Deum}
\switchcolumn*

\LA
\LessonHead{Lectio III (mense Decembri)}
\switchcolumn
\EN
\LessonHead{Lesson III (in December)}
\switchcolumn*

\LA
\LessonTitle{Ex libro Officiórum S. Ambrósii Epíscopi}
\switchcolumn
\EN
\LessonTitle{From the Book of Offices of St. Ambrose, Bishop}
\switchcolumn*

\LA
\Source{Lib. 1, cap. 18}
\switchcolumn
\EN
\Source{Book I. ch. 18}
\switchcolumn*

\LA
\noindent Bonus regéndæ castitátis pudor est comes, qui primus, in ipso cognitiónis ingréssu, Dómini Matrem comméndat legéntibus, et tamquam testis lócuples dignam quæ ad tale munus eligerétur ástruit: quod in cubículo, quod sola, quod salutáta ab Angelo tacet et mota est in intróitu ejus, quod ad virílis sexus spéciem peregrínam turbátur aspéctus Vírginis. Itaque, quamvis esset húmilis præ verecúndia tamen non resalutávit, nec ullum respónsum rétulit, nisi ubi de suscipiénda Dómini generatióne cognóvit; ut qualitátem efféctus dísceret, non ut sermónem refélleret.\par
\switchcolumn
\EN
\noindent The good comrade and guardian of chastity is modesty. And this is the first thing that strikes the reader of the Gospel, as he begins to learn about the Mother of our Lord. Her modesty commend her, and like reliable witness, attests that she is worthy of being chosen for so high an office. Alone in her room, she was silent at the Angel's greeting. His entrance had disturbed her, and her countenance reflected the Virgins agitation at seeing this stranger appear in the form of a man. And so, although she was humble, still out of modesty she did not return his greeting, nor did she gave any answer, until she learned that she must agree to becoming the Mother of God, and then she spoke, not to reject his message but to learn how this marvel was to come about.\par
\switchcolumn*

\LA
\TeDeum{Te Deum}
\switchcolumn
\EN
\TeDeum{Te Deum}
\switchcolumn*

\end{paracol}

\Office{In Festis Beatae Mariae Virginis}{In Festis Beatae Mariae Virginis}{}
\addcontentsline{toc}{subsection}{In Festis Beatae Mariae Virginis \textit{— In Festis Beatae Mariae Virginis}}
\begin{paracol}{2}
\LA
\LessonHead{Lectio I}
\switchcolumn
\EN
\LessonHead{Lesson I}
\switchcolumn*

\LA
\LessonTitle{De Parábolis Salomónis}
\switchcolumn
\EN
\LessonTitle{Lesson from the book of Proverbs}
\switchcolumn*

\LA
\Cite{Prov 8:12-17}
\switchcolumn
\EN
\Cite{Prov 8:12-17}
\switchcolumn*

\LA
\noindent Ego sapiéntia hábito in consílio et erudítis intérsum cogitatiónibus. Timor Dómini odit malum: arrogántiam, et supérbiam, et viam pravam, et os bilíngue detéstor. Meum est consílium et ǽquitas, mea est prudéntia, mea est fortitúdo. Per me reges regnant, et legum conditóres justa decérnunt; per me príncipes ímperant, et poténtes decérnunt justítiam. Ego diligéntes me díligo; et qui mane vígilant ad me, invénient me.\par
\switchcolumn
\EN
\noindent I wisdom dwell in counsel, and am present in learned thoughts. The fear of the Lord hateth evil: I hate arrogance, and pride, and every wicked way, and a mouth with a double tongue. Counsel and equity is mine, prudence is mine, strength is mine. By me kings reign, and lawgivers decree just things, By me princes rule, and the mighty decree justice. I love them that love me: and they that in the morning early watch for me, shall find me.\par
\switchcolumn*

\LA
\begin{Resp}\Rsign Sancta et immaculáta virgínitas, quibus te láudibus éfferam, néscio: \Star{} Quia quem cæli cápere non póterant, tuo grémio contulísti.\end{Resp}
\begin{Resp}\Vsign Benedícta tu in muliéribus, et benedíctus fructus ventris tui.\end{Resp}
\begin{Resp}\Rsign Quia quem cæli cápere non póterant, tuo grémio contulisti.\end{Resp}
\switchcolumn
\EN
\begin{Resp}\Rsign O how holy and how spotless is thy virginity; I am too dull to praise thee: \Star{} For thou hast borne in thy breast Him Whom the heavens cannot contain.\end{Resp}
\begin{Resp}\Vsign Blessed art thou among women, and blessed is the fruit of thy womb.\end{Resp}
\begin{Resp}\Rsign For thou hast borne in thy breast Him Whom the heavens cannot contain.\end{Resp}
\switchcolumn*

\LA
\LessonHead{Lectio II}
\switchcolumn
\EN
\LessonHead{Lesson II}
\switchcolumn*

\LA
\Cite{Prov 8:18-25}
\switchcolumn
\EN
\Cite{Prov 8:18-25}
\switchcolumn*

\LA
\noindent Mecum sunt divítiæ et glória, opes supérbæ et justítia. Mélior est enim fructus meus auro et lápide pretióso, et genímina mea argénto elécto. In viis justítiæ ámbulo, in médio semitárum judícii, ut ditem diligéntes me et thesáuros eórum répleam. Dóminus possidébit me in inítio viárum suárum, ántequam quidquam fáceret a princípio. Ab ætérno ordináta sum et ex antíquis, ántequam terra fíeret. Nondum erant abýssi, et ego jam concépta eram; necdum fontes aquárum erúperant, necdum montes gravi mole constíterant; ante colles ego parturiébar.\par
\switchcolumn
\EN
\noindent With me are riches and glory, glorious riches and justice. For my fruit is better than gold and the precious stone, and my blossoms than choice silver. I walk in the way of justice, in the midst of the paths of judgment, That I may enrich them that love me, and may fill their treasures. The Lord possessed me in the beginning of his ways, before he made any thing from the beginning. I was set up from eternity, and of old before the earth was made. The depths were not as yet, and I was already conceived. neither had the fountains of waters as yet sprung out: The mountains with their huge bulk had not as yet been established: before the hills I was brought forth:\par
\switchcolumn*

\LA
\begin{Resp}\Rsign Congratulámini mihi, omnes qui dilígitis Dóminum: quia cum essem párvula, plácui Altíssimo, \Star{} Et de meis viscéribus génui Deum et hóminem.\end{Resp}
\begin{Resp}\Vsign Beátam me dicent omnes generatiónes, quia ancíllam húmilem respéxit Deus.\end{Resp}
\begin{Resp}\Rsign Et de meis viscéribus génui Deum et hóminem.\end{Resp}
\switchcolumn
\EN
\begin{Resp}\Rsign Rejoice with me, all ye that love the Lord, for while I was yet a little one, I pleased the Most High. \Star{} And I have brought forth from my bowels God and man.\end{Resp}
\begin{Resp}\Vsign All generations shall call me blessed, since the Lord hath regarded the lowliness of his handmaiden.\end{Resp}
\begin{Resp}\Rsign And I have brought forth from my bowels God and man.\end{Resp}
\switchcolumn*

\LA
\LessonHead{Lectio III}
\switchcolumn
\EN
\LessonHead{Lesson III}
\switchcolumn*

\LA
\Cite{Prov 8:34-36; 9:1-5}
\switchcolumn
\EN
\Cite{Prov 8:34-36; 9:1-5}
\switchcolumn*

\LA
\noindent Beátus homo qui audit me, et qui vígilat ad fores meas cotídie, et obsérvat ad postes óstii mei. Qui me invénerit, invéniet vitam, et háuriet salútem a Dómino; qui autem in me peccáverit, lædet ánimam suam. Omnes, qui me odérunt, díligunt mortem. Sapiéntia ædificávit sibi domum, excídit colúmnas septem. Immolávit víctimas suas, míscuit vinum et propósuit mensam suam. Misit ancíllas suas ut vocárent ad arcem et ad mœ́nia civitátis: Si quis est párvulus, véniat ad me. Et insipiéntibus locúta est: Veníte, comédite panem meum, et bíbite vinum quod míscui vobis.\par
\switchcolumn
\EN
\noindent Blessed is the man that heareth me, and that watcheth daily at my gates, and waiteth at the posts of my doors. He that shall find me, shall find life, and shall have salvation from the Lord: But he that shall sin against me, shall hurt his own soul. All that hate me love death. Wisdom hath built herself a house, she hath hewn out her seven pillars. She hath slain her victims, mingled her wine, and set forth her table. She hath sent her maids to invite to the tower, and to the walls of the city: Whosoever is a little one, let him come to me. And to the unwise she said: Come, eat my bread, and drink the wine which I have mingled for you\par
\switchcolumn*

\LA
\begin{Resp}\Rsign Beáta es, Virgo María, quæ Dóminum portásti, Creatórem mundi: \Star{} Genuísti qui te fecit, et in ætérnum pérmanes Virgo.\end{Resp}
\begin{Resp}\Vsign Ave, María, grátia plena; Dóminus tecum.\end{Resp}
\begin{Resp}\Rsign Genuísti qui te fecit, et in ætérnum pérmanes Virgo.\end{Resp}
\begin{Resp}\Vsign Glória Patri, et Fílio, \Star{} et Spirítui Sancto.\end{Resp}
\begin{Resp}\Rsign Genuísti qui te fecit, et in ætérnum pérmanes Virgo.\end{Resp}
\switchcolumn
\EN
\begin{Resp}\Rsign Blessed art thou, O Virgin Mary, who hast carried the Lord, the Maker of the world. \Star{} Thou hast borne Him Who created thee, and thou abidest a virgin for ever.\end{Resp}
\begin{Resp}\Vsign Hail, Mary, full of grace. the Lord is with thee.\end{Resp}
\begin{Resp}\Rsign Thou hast borne Him Who created thee, and thou abidest a virgin for ever.\end{Resp}
\begin{Resp}\Vsign Glory be to the Father, and to the Son, \Star{} and to the Holy Ghost.\end{Resp}
\begin{Resp}\Rsign Thou hast borne Him Who created thee, and thou abidest a virgin for ever.\end{Resp}
\switchcolumn*

\LA
\LessonHead{Lectio VII}
\switchcolumn
\EN
\LessonHead{Lesson VII}
\switchcolumn*

\LA
\LessonTitle{Léctio sancti Evangélii secúndum Lucam}
\switchcolumn
\EN
\LessonTitle{From the Holy Gospel according to Luke}
\switchcolumn*

\LA
\Cite{Luc 11:27-28}
\switchcolumn
\EN
\Cite{Luke 11:27-28}
\switchcolumn*

\LA
\noindent In illo témpore: Loquénte Jesu ad turbas, extóllens vocem quædam múlier de turba dixit illi: Beátus venter qui te portávit. Et réliqua.\par
\switchcolumn
\EN
\noindent At that time, as Jesus spoke unto the multitudes, a certain woman of the company lifted up her voice and said unto Him: Blessed is the womb that bore thee. And so on.\par
\switchcolumn*

\LA
\LessonTitle{Homilía sancti Bedæ Venerábilis Presbýteri}
\switchcolumn
\EN
\LessonTitle{Homily by the Venerable Bede}
\switchcolumn*

\LA
\Source{Lib 4. cap. 49. in Luc. 11.}
\switchcolumn
\EN
\Source{Lib 4 Cap 49 in Luke 11}
\switchcolumn*

\LA
\noindent Magnæ devotiónis et fídei hæc múlier osténditur, quæ, scribis et pharisǽis Dóminum tentántibus simul et blasphemántibus, tanta ejus incarnatiónem præ ómnibus sinceritáte cognóscit, tanta fidúcia confitétur, ut et præséntium prócerum calúmniam, et futurórum confúndat hæreticórum perfídiam. Nam, sicut tunc Judǽi, Sancti Spíritus ópera blasphemándo, verum consubstantialémque Patri Dei Fílium negábant; sic hærétici póstea, negándo Maríam semper Vírginem, Sancti Spíritus operánte virtúte, nascitúro cum humánis membris Unigénito Dei, carnis suæ matériam ministrásse, verum consubstantialémque matri Fílium hóminis fatéri non debére dixérunt.\par
\switchcolumn
\EN
\noindent It is plain that this was a woman of great earnestness and faith. The Scribes and Pharisees were at once tempting and blaspheming the Lord, but this woman so clearly grasped His Incarnation, and so bravely confessed the same, that she confounded both the lies of the great men who were present, and the faithlessness of the heretics who were yet to come. Even as the Jews then, blaspheming the works of the Holy Ghost, denied the very Son of God Who is of one substance with the Father, so afterwards did the heretics, by denying that Mary always a Virgin did, under the operation of the Holy Ghost, supply flesh to the Only begotten One of God, when He was about being born in an human Body, even so, I say, did the heretics deny that the Son of Man should be called a true Son, Who is of one substance with His Mother.\par
\switchcolumn*

\LA
\begin{Resp}\Rsign Felix namque es, sacra Virgo María, et omni laude digníssima: \Star{} Quia ex te ortus est sol justítiæ, Christus, Deus noster.\end{Resp}
\begin{Resp}\Vsign Ora pro pópulo, intérveni pro clero, intercéde pro devóto femíneo sexu: séntiant omnes tuum juvámen, quicúmque célebrant tuam sanctam festivitátem.\end{Resp}
\begin{Resp}\Rsign Quia ex te ortus est sol justítiæ, Christus, Deus noster.\end{Resp}
\switchcolumn
\EN
\begin{Resp}\Rsign O Holy Virgin Mary, happy indeed art thou, and right worthy of all praise \Star{} For out of thee rose the Sun of righteousness, even Christ our God.\end{Resp}
\begin{Resp}\Vsign Pray for the people, plead for the clergy, make intercession for all women vowed to God. May all that are keeping this thine holy Feast day feel the might of thine assistance.\end{Resp}
\begin{Resp}\Rsign For out of thee rose the Sun of righteousness, even Christ our God.\end{Resp}
\switchcolumn*

\LA
\LessonHead{Lectio VIII}
\switchcolumn
\EN
\LessonHead{Lesson VIII}
\switchcolumn*

\LA
\noindent Sed, si caro Verbi Dei secúndum carnem nascéntis a carne Vírginis matris pronuntiátur extránea, sine causa venter qui eam portásset, úbera quæ lactássent, beatificántur. Dicit autem Apóstolus: Quia misit Deus Fílium suum, factum ex mulíere, factum sub lege. Neque audiéndi sunt, qui legéndum putant: Natum ex mulíere, factum sub lege, sed Factum ex mulíere; quia concéptus ex útero virgináli, carnem non de níhilo, non aliúnde, sed matérna traxit ex carne. Alióquin nec vere Fílius hóminis dicerétur, qui oríginem non habéret ex hómine. Et nos ígitur, his contra Eutychen dictis, extollámus vocem cum Ecclésia cathólica, cujus hæc múlier typum gessit, extollámus et mentem de médio turbárum, dicamúsque Salvatóri: Beátus venter qui te portávit, et úbera quæ suxísti. Vere enim beáta parens, quæ, sicut quidam ait, Eníxa est puérpera Regem, qui cælum terrámque tenet per sǽcula.\par
\switchcolumn
\EN
\noindent If we shall say that the Flesh, Wherewith the Son of God was born in the flesh, was something outside of the flesh of the Virgin His Mother, without reason should we bless the womb that bore Him, and the paps which He hath sucked. But the Apostle saith: “God sent forth His Son, made of a woman, made under the law”, (Gal. iv. 4,) and they are not to be listened to who read this passage: “Born of a woman, made under the law.” He was made of a woman, for He was conceived in a virgin's womb, and took His Flesh, not from nothing, not from elsewhere, but from the flesh of His Mother. Otherwise, and if He had not been sprung of a woman, He could not with truth be called the Son of man. Let us therefore, denying the doctrine of Eutyches, lift up our voice, along with the Universal Church, whereof that woman was a figure, let us lift up our heart as well as our voice from the company, and say unto the Saviour: Blessed is the womb that bare thee, and the paps which thou hast sucked. Blessed Mother, of whom one hath said thou art His Mother, Who reigns over earth and over heaven for ever.\par
\switchcolumn*

\LA
\begin{Resp}\Rsign Beátam me dicent omnes generatiónes, \Star{} Quia fecit mihi Dóminus magna qui potens est, et sanctum nomen ejus.\end{Resp}
\begin{Resp}\Vsign Et misericórdia ejus a progénie in progénies timéntibus eum.\end{Resp}
\begin{Resp}\Rsign Quia fecit mihi Dóminus magna qui potens est, et sanctum nomen ejus.\end{Resp}
\begin{Resp}\Vsign Glória Patri, et Fílio, \Star{} et Spirítui Sancto.\end{Resp}
\begin{Resp}\Rsign Quia fecit mihi Dóminus magna qui potens est, et sanctum nomen ejus.\end{Resp}
\switchcolumn
\EN
\begin{Resp}\Rsign All generations shall call me blessed. \Star{} For He that is mighty, even the Lord, hath done to me great things; and Holy is His Name.\end{Resp}
\begin{Resp}\Vsign And his mercy is on them that fear Him, from generation to generation.\end{Resp}
\begin{Resp}\Rsign He that is mighty, even the Lord, hath done to me great things; and Holy is His Name.\end{Resp}
\begin{Resp}\Vsign Glory be to the Father, and to the Son, \Star{} and to the Holy Ghost.\end{Resp}
\begin{Resp}\Rsign He that is mighty, even the Lord, hath done to me great things; and Holy is His Name.\end{Resp}
\switchcolumn*

\LA
\LessonHead{Lectio IX}
\switchcolumn
\EN
\LessonHead{Lesson IX}
\switchcolumn*

\LA
\noindent Quinímmo beáti qui áudiunt verbum Dei et custódiunt. Pulchre Salvátor attestatióni mulíeris ánnuit, non eam tantúmmodo quæ Verbum Dei corporáliter generáre merúerat, sed et omnes qui idem Verbum spiritáliter audítu fídei concípere, et boni óperis custódia vel in suo vel in proximórum corde párere et quasi álere studúerint, asséverans esse beátos; quia, et éadem Dei Génetrix, et inde quidem beáta quia Verbi incarnándi minístra facta est temporális, sed inde multo beátior quia ejúsdem semper amándi custos manébat ætérna.\par
\switchcolumn
\EN
\noindent Yea, rather, blessed are they that hear the Word of God and keep it. How nobly doth the Saviour say “Yea” to the woman's blessing, declaring also that not only is she blessed who was meet to give bodily birth to the Word of God, but that all they who spiritually conceive the same Word by the hearing of faith, and, by keeping it through good works, bring it forth and, as it were, carefully nurse it, in their own hearts, and in the hearts of their neighbours, are also blessed. Yea, and that the very Mother of God herself was blessed in being for a while the handmaid of the Word of God made Flesh, but that she was much more blessed in this, that through her love she keepeth Him for ever.\par
\switchcolumn*

\LA
\TeDeum{Te Deum}
\switchcolumn
\EN
\TeDeum{Te Deum}
\switchcolumn*

\end{paracol}
\backmatter
\IndexOfOffices
\end{document}
